\documentclass[11pt,a4paper,oneside]{book}
\usepackage[T1]{fontenc}
\usepackage[utf8]{inputenc}
\usepackage[indonesian]{babel}
\usepackage{lmodern}
\usepackage{amsmath,amsthm,amssymb,amscd}
\usepackage{graphicx}
\usepackage[a4paper,margin=22mm,headheight=15pt]{geometry}
\usepackage{microtype}
\usepackage[hidelinks,pdfusetitle]{hyperref}
\usepackage{bookmark}
% language package supplied by unit wrapper

% UTF-8 input supplied by unit wrapper

%\usepackage[utf8]{inputenc}

\usepackage{ifthen}

%\usepackage{fancyhdr}

%Aus Grundvorspann

%Packages

\usepackage{amscd}

\usepackage{amssymb}

\usepackage{epsfig} %Graphik

\usepackage{verbatim,moreverb} %

\usepackage{enumitem}

\usepackage{eurosym}

\usepackage{epigraph}

\setcounter{MaxMatrixCols}{20}

%Mathematische Einzelbefehle

%[[Projekt:Semantische Vorlagen/Mathematischer Modus/Latex]]

\newcommand{\mathdisp}[2]{\[ #1 \, #2 \]}

\newcommand{\mathdisplayteile}[3]{\[ #1 \] \[ #2 \, #3 \]}

\newcommand{\mathind}[3]{$ #1$, $#2 $#3}

\newcommand{\mathbed}[8]{$ #1 ${\ifthenelse{\equal{#2}{}}{, }{ #2 }}$ #3 ${\ifthenelse{\equal{#5}{}}{}{{\ifthenelse{\equal{#4}{}}{, }{ #4 }}}$ #5 $}{\ifthenelse{\equal{#7}{}}{}{{\ifthenelse{\equal{#6}{}}{, }{ #6 }}}$ #7 $}#8}

\newcommand{\mathbeddisp}[8]{\[ #1 {\ifthenelse{\equal{#2}{}}{,\, }{ \text{ #2 } }} #3 {\ifthenelse{\equal{#5}{}}{}{{\ifthenelse{\equal{#4}{}}{,\, }{ \text{ #4 } }}} #5 }{\ifthenelse{\equal{#7}{}}{}{{\ifthenelse{\equal{#6}{}}{,\, }{ \text{ #6 } }}} #7 }#8 \]}

\newcommand{\mathbedtermdisp}[8]{\[ #1\] #2 $ #3 $#4 $ #5 $#6$ #7 $#8}

\newcommand{\mathkon}[4]{$ #1 $ #2 $ #3 $#4}

\newcommand{\mathkor}[5]{{\ifthenelse{\equal{#1}{}}{}{#1 }}$ #2 $ #3 $ #4 $#5}

\newcommand{\mathkork}[5]{{\ifthenelse{\equal{#1}{}}{}{#1 }}$ #2 $ #3 $ #4 $#5}

\newcommand{\mathlist}[5]{$ #1 $\ifthenelse{\equal{#2}{}}{,}{ #2} $ #3 $\ifthenelse{\equal{#4}{}}{,}{ #4} $ #5 $}

\newcommand{\mathlistdisp}[5]{\[  #1  \ifthenelse{\equal{#2}{}}{,}{ \text{ #2 }}  #3  \ifthenelse{\equal{#4}{}}{,}{\text{ #4 } }  #5 \] }

\newcommand{\alignier}[2]{\begin{align*} #1 \, #2 \end{align*}}

\newcommand{\mathlistvierdisp}[7]{\[ #1,\, #3, \, #5 \text{ #6 } #7  \]}

%[[Projekt:Semantische Vorlagen/Mathematische Vergleichsvorlagen/Latex]]

\newcommand{\mavergleichskette}[4]{$ #1 #2 #3#4$}

\newcommand{\mavergleichskettedisp}[4]{\[  #1 #2 #3 #4 \]}

\newcommand{\mavergleichskettedisplang}[4]{\[  #1 #2 #3 #4 \]}

\newcommand{\mavergleichskettealign}[4]{ \begin{eqnarray*}  #1 #2 #3 #4 \end{eqnarray*}  }

\newcommand{\vergleichskette}[9]{ #1 \, #2 \, #3 \ifthenelse
{\equal {#4}{ }} {} {\, #4 \, #5 }  \ifthenelse {\equal {#6}{ } }{}
{\, #6 \, #7 } \ifthenelse {\equal {#8}{ }}{} {\, #8 \, #9 }  }

\newcommand{\vergleichskettelang}[9]{ #1 \, #2 \, #3 \ifthenelse
{\equal {#4}{ }} {} {\, #4 \, #5 }  \ifthenelse {\equal {#6}{ } }{}
{\, #6 \, #7 } \ifthenelse {\equal {#8}{ }}{} {\, #8 \, #9 }  }

\newcommand{\vergleichskettefortsetzung}[8]{ \, #1 \, #2 \ifthenelse
{\equal {#3}{ }} {} {\, #3 \, #4 }  \ifthenelse {\equal {#5}{ } }{}
{\, #5 \, #6 } \ifthenelse {\equal {#7}{ }}{} {\, #7 \, #8 } }

\newcommand{\vergleichskettealign}[9]{ #1 & #2 & #3 \ifthenelse
{\equal {#4}{ }} {} {\cr & #4 & #5 }  \ifthenelse {\equal {#6}{ }
}{} {\cr & #6 & #7 } \ifthenelse {\equal {#8}{ }}{} {\cr & #8 & #9 }
}

\newcommand{\vergleichskettefortsetzungalign}[8]{ \cr & #1 & #2 \ifthenelse
{\equal {#3}{ }} {} {\cr & #3 & #4 }  \ifthenelse {\equal {#5}{ }
}{} {\cr & #5 & #6 } \ifthenelse {\equal {#7}{ }}{} {\cr & #7 & #8 }
}

\newcommand{\mavergleichskettealigndrucklinks}[4]{ \begin{eqnarray*}  #1 #2 #3 #4 \end{eqnarray*}  }

\newcommand{\vergleichskettealigndrucklinks}[9]{  &  & #1 \cr & #2 & #3 \ifthenelse
{\equal {#4}{ }} {} {\cr & #4 & #5 }  \ifthenelse {\equal {#6}{ }
}{} {\cr & #6 & #7 } \ifthenelse {\equal {#8}{ }}{} {\cr & #8 & #9 }
}

\newcommand{\mavergleichskettealignhandlinks}[4]{ \begin{eqnarray*}  #1 #2 #3 #4 \end{eqnarray*}  }

\newcommand{\vergleichskettealignhandlinks}[9]{ #1 & #2 & #3 \ifthenelse
{\equal {#4}{ }} {} {\cr & #4 & #5 }  \ifthenelse {\equal {#6}{ }
}{} {\cr & #6 & #7 } \ifthenelse {\equal {#8}{ }}{} {\cr & #8 & #9 }
}

\newcommand{\mavergleichskettedisphandlinks}[4]{ \[ #1 #2 #3 #4 \] }

\newcommand{\vergleichskettedisphandlinks}[9]{  #1 \, #2 \, #3 \ifthenelse
{\equal {#4}{ }} {} {\, #4 \, #5 }  \ifthenelse {\equal {#6}{ } }{}
{\, #6 \, #7 } \ifthenelse {\equal {#8}{ }}{} {\, #8 \, #9 }  }

\newcommand{\zeilemitteil}[2]{ \ifthenelse {\equal {#2} {} } {#1} { #1 \cr & & #2}   }

\newcommand{\mavergleichskettek}[4]{$ #1 #2 #3#4$}

\newcommand{\vergleichskettek}[9]{ #1 \, #2 \, #3 \ifthenelse
{\equal {#4}{ }} {} {\, #4 \, #5 }  \ifthenelse {\equal {#6}{ } }{}
{\, #6 \, #7 } \ifthenelse {\equal {#8}{ }}{} {\, #8 \, #9 }  }

\newcommand{\vergleichskettefortsetzungk}[8]{ \, #1 \, #2 \ifthenelse
{\equal {#3}{ }} {} {\, #3 \, #4 }  \ifthenelse {\equal {#5}{ } }{}
{\, #5 \, #6 } \ifthenelse {\equal {#7}{ }}{} {\, #7 \, #8 } }

%[[Projekt:Semantische Vorlagen/Mathematische Strukturvorlagen/Latex]]

\newcommand{\N}   {{\mathbb N}}
\newcommand{\Z}   {{\mathbb Z}}
\newcommand{\Q}   {{\mathbb Q}}
\newcommand{\R}   {{\mathbb R}}
\newcommand{\C}   {{\mathbb C}}
\newcommand{\Complex}  {{\mathbb C}}

\newcommand{\defeq}{:=}

\newcommand{\defeqr}{=:}

\newcommand{\maabb}[4]{${\ifthenelse {\equal {#1}{} } {} {#1 \colon }}\,  #2 \rightarrow #3 #4 $}

\newcommand{\maabbele}[6]{$ {\ifthenelse {\equal {#1}{} } {} {#1 \colon }}\,   #2 \rightarrow  #3 , \, #4 \mapsto #5 #6 $}

\newcommand{\maabbdisp}[4]{\[ {\ifthenelse {\equal {#1}{} } {} {#1 \colon }}\,  #2 \longrightarrow #3 #4 \]}

\newcommand{\maabbeledisp}[6]{\[ {\ifthenelse {\equal {#1}{} } {} {#1 \colon }}\,   #2 \longrightarrow  #3 , \, #4 \longmapsto #5 #6 \]}

%Variante bei ifequal-Problem

\newcommand{\maabbnamedisp}[4]{\[ #1 \colon \, #2 \longrightarrow #3 #4 \]}

\newcommand{\maabbeledispvar}[6]{\[ {\ifthenelse {\equal {#1}{} } {} {#1 \colon }}\,   #2 \longrightarrow  #3 , \, #4 \longmapsto #5 #6  \]}

\newcommand{\maabbelementzeiledisplay}[6]{\[ {\ifthenelse {\equal {#1}{} } {} {#1 \colon }}\, #2 \longrightarrow #3 , \]  	\[ #4 \longmapsto #5 #6 \]}

\newcommand{\maabbelementdoppelzeiledisplay}[6]{\[ {\ifthenelse {\equal {#1}{} } {} {#1 \colon }}\, #2 \longrightarrow #3 , \] \[ #4 \longmapsto \] \[  #5 #6 \]}

\newcommand{ \arccot  } {\operatorname{arccot } }

\newcommand{\betrag}[1]{\left| #1 \right|}

\newcommand{\strukturgarbe}{ {\mathcal O} }

\newcommand{\schnittring}[1]{\Gamma (#1 , {\mathcal O})}

\newcommand{\tensor}{\otimes}

%[[Projekt:Semantische Vorlagen/Textgestaltung/Latex]]

\newcommand{\einrueckung}[1]{\par \hspace{1cm}#1 \par}

\newcommand{\zusatz}[2]{(#1)#2}

\newcommand{\zusatzklammer}[3]{(#1#2)#3}

\newcommand{\zusatzgs}[3]{\-- #1 \-- #3}

\newcommand{\zusatzfussnote}[3]{#3\footnote{#1#2}}

%Feinabstimmung

\newcommand{\latexdruckeinskurz}{\!}

\newcommand{\latexdruckzweikurz}{\! \!}

\newcommand{\latexdruckdreikurz}{\! \! \!}

%[[Projekt:Semantische Vorlagen/Satzumgebungen/Latex]]

\newtheorem{fakt}{Fakta}[section] %so wird abschnittsweise %durchnummeriert. Bei [] wird % einfach nummeriert.

\newtheorem{Fakt}{fakt}[Fakt]

\newtheorem{Proposition}[fakt]{Proposisi}

\newtheorem{Korollar}[fakt]{Korolari}

\newtheorem{Lemma}[fakt]{Lema}

\newtheorem{Satz}[fakt]{Teorema}

\newtheorem{Theorem}[fakt]{Teorema}

\theoremstyle{definition}

\newtheorem{Axiom}[fakt]{Aksioma}

\newtheorem{aufgabe}[fakt]{Soal}

\newtheorem{Aufgabe}[fakt]{Soal}

\newtheorem{klausuraufgabe}{Soal ujian}

\newtheorem{Klausuraufgabe}{Soal ujian}

\newtheorem{beispiel}[fakt]{Contoh}

\newtheorem{Beispiel}[fakt]{Contoh}

\newtheorem{Konstruktion}[fakt]{Konstruksi}

\newtheorem{konstruktion}[fakt]{Konstruksi}

\newtheorem{Verfahren}[fakt]{Verfahren}

\newtheorem{verfahren}[fakt]{Verfahren}

\newtheorem{bemerkung}[fakt]{Catatan}

\newtheorem{Bemerkung}[fakt]{Catatan}

\newtheorem{definition}[fakt]{Definisi}

\newtheorem{Definition}[fakt]{Definisi}

\newtheorem{notation}[fakt]{Notation}

\newtheorem{Notation}[fakt]{Notation}

\newtheorem{Frage}[fakt]{Frage}

\newtheorem{frage}[fakt]{Frage}

\newtheorem{problem}[fakt]{Problem}

\newtheorem{Problem}[fakt]{Problem}

\newtheorem{situation}[fakt]{Situation}

\newtheorem{Situation}[fakt]{Situation}

%Einlesegestaltung

%[[Projekt:Semantische Vorlagen/Umgebungsdesign/Latex]]

%Aufgabengestaltung

\newcommand{\inputaufgabe}[4]{\aufgabevorskip \begin{aufgabe} \punkte {#1} \aufgabepunktskip #2 #3 #4\end{aufgabe} \aufgabenachskip}

\newcommand{\inputaufgabegibtloesung}[4]{\aufgabevorskip \begin{aufgabe}\hspace{-1.8 mm}* \punkte {#1} \aufgabepunktskip #2 #3 #4\end{aufgabe} \aufgabenachskip}

\newcommand{\inputaufgabepunktenummervonhand}[3]{\aufgabevorskip {\bf Aufgabe #1} (#2 Punkte) \aufgabepunktskip #3  \aufgabenachskip}

\newcommand{\inputaufgabeloesung}[2]{ \begin{aufgabe} #1 \end{aufgabe} \aufgabenachskip

\bigskip

Lösung

\bigskip #2

}

\newcommand{\inputaufgabeklausurloesung}[3]{\par \newpage \begin{aufgabe} {\ifthenelse {\equal {#1}{}}{} {\ifthenelse {\equal {#1}{1}} {(1 Punkt)} {(#1 Punkte)}}} \par \bigskip #2 \end{aufgabe} \underline{L\"osung:} \par \bigskip #3 }

\newcommand{\inputaufgabeloesungvar}[2]{\aufgabevorskip Aufgabe: #1 \aufgabenachskip
Lösung: #2}

\newcommand{\inputaufgabepunkteloesung}[3]{\aufgabevorskip \begin{aufgabe} {\ifthenelse {\equal {#1}{}}{} {\ifthenelse {\equal {#1}{1}} {(1 Punkt)} {(#1 Punkte)}}}  \aufgabepunktskip #2 \end{aufgabe} \aufgabenachskip \loesung #3}

\newcommand{\loesung}[1]{\underline{Lösung:} #1}

\newcommand{\punkte}[1]{\ifthenelse {\equal {#1}{}}{} {\ifthenelse {\equal {#1}{1}} {(1 poin)} {(#1 poin)}}}

\newcommand{\inputexercise}[4]{\aufgabevorskip \begin{exercise} \points {#1} \aufgabepunktskip #2 #3 #4\end{exercise} \aufgabenachskip}

\newcommand{\points}[1]{\ifthenelse {\equal {#1}{}}{} {\ifthenelse {\equal {#1}{1}} {(1 point)} {(#1 points)}}}

%Beispielgestaltung

\newcommand{\inputbeispiel}[2]
{\beispielvorskip
\begin{beispiel}
\beispielbenennung{#1}
\beispielbenennungnachskip #2
\end{beispiel}
\beispielnachskip}

\newcommand{\beispielbenennung}[1]{ \ifthenelse {\equal {#1}{} \OR \equal {#1}{   }  }{}{(#1)}     }

%Bemerkunggestaltung

\newcommand{\inputbemerkung}[2]
{\bemerkungvorskip
\begin{bemerkung}
\bemerkungbenennung{#1}
\bemerkungbenennungnachskip #2
\end{bemerkung}
\bemerkungnachskip}

\newcommand{\bemerkungbenennung}[1]{  \ifthenelse {\equal {#1}{} \OR \equal {#1}{   }  }{}{(#1)}}

\newcommand{\inputverfahren}[2]
{\bemerkungvorskip
\begin{verfahren}
\bemerkungbenennung{#1}
\bemerkungbenennungnachskip #2
\end{verfahren}
\bemerkungnachskip}

\newcommand{\inputkonstruktion}[2]
{\bemerkungvorskip
\begin{konstruktion}
\bemerkungbenennung{#1}
\bemerkungbenennungnachskip #2
\end{konstruktion}
\bemerkungnachskip}

\newcommand{\inputfrage}[2]
{\bemerkungvorskip
\begin{frage}
\problembenennung{#1}
\bemerkungbenennungnachskip #2
\end{frage}
\bemerkungnachskip}

\newcommand{\fragebenennung}[1]{  \ifthenelse {\equal {#1}{} \OR \equal {#1}{   }  }{}{(#1)}}

\newcommand{\inputproblem}[2]
{\bemerkungvorskip
\begin{problem}
\problembenennung{#1}
\bemerkungbenennungnachskip #2
\end{problem}
\bemerkungnachskip}

\newcommand{\problembenennung}[1]{  \ifthenelse {\equal {#1}{} \OR \equal {#1}{   }  }{}{(#1)}}

%Situationsgestaltung

\newcommand{\inputsituation}[2]
{\begin{situation}
\situationbenennung{#1} #2
\end{situation}
}

\newcommand{\situationbenennung}[1]{\ifthenelse {\equal {#1}{} \OR \equal {#1}{  }  }{}{(#1)}}

%Definitiongestaltung

\newcommand{\inputdefinition}[2]
{\definitionvorskip
\begin{definition}
\definitionbenennung{#1}
\definitionbenennungnachskip #2
\end{definition}
\definitionnachskip}

\newcommand{\inputaxiom}[2]
{\definitionvorskip
\begin{Axiom}
\definitionbenennung{#1}
\definitionbenennungnachskip #2
\end{Axiom}
\definitionnachskip}

\newcommand{\definitionbenennung}[1]{\ifthenelse {\equal {#1}{} \OR \equal {#1}{  }  }{}{(#1)}}

\newcommand{\inputnotation}[2]
{\definitionvorskip \begin{notation} \definitionbenennung{#1} \definitionbenennungnachskip #2 \end{notation} \definitionnachskip}

%Fakt- und Beweisgestaltung

\renewcommand{\proofname}{\hspace{0cm}{\it Bukti}}

\newcommand{\beweis}[1]{\begin{proof} #1 \end{proof}}

\newcommand{\inputfakt}[4]
{\faktvorskip
\begin{#2}

%\label{#1}   (Absetzen, sonst  kann das folgende dahinter rutschen)

\faktbenennung{#3}
\faktbenennungnachskip #4
\end{#2}
\faktnachskip
}

\newcommand{\inputfaktbeweis}[5]
{\faktvorskip
\begin{#2}

%\label{#1}   (Absetzen, sonst  kann das folgende dahinter rutschen)

\faktbenennung{#3}
\faktbenennungnachskip #4
\end{#2}
\faktnachskip
\beweis{#5}}

\newcommand{\inputfaktbeweisnichtvorgefuehrt}[5]{\inputfaktbeweis {#1} {#2} {#3} {#4} {Dieser Beweis wurde in der Vorlesung nicht vorgef\"uhrt.} }

\newcommand{\inputfaktbeweistrivial}[4]{\inputfaktbeweis {#1} {#2} {#3} {#4} {Das ist trivial.} }

\newcommand{\inputfaktuebergangbeweis}[6]
{\faktvorskip
\begin{#2}

%\label{#1} (Absetzen, sonst kann das
% folgende dahinter rutschen)

\faktbenennung{#3}
\faktbenennungnachskip #4
\end{#2}
\faktnachskip #5 \faktnachskip
\beweis{#6}}

\newcommand{\inputbeweis}[1]{\beweis{#1} }

\newcommand{\faktbenennung}[1]{\ifthenelse {\equal {#1}{} \OR \equal {#1}{  }  }{}{(#1)}}

%Gestaltung der Faktstruktur

\newcommand{\faktsituation}[1]{\faktsituationskip #1}

\newcommand{\faktvoraussetzung}[1]{\faktvoraussetzungskip #1}

\newcommand{\faktvoraussetzungleer}[1]{ \hspace{-0,15cm} }

\newcommand{\faktvoraussetzungpos}[1]{\faktvoraussetzungskip #1}

\newcommand{\faktuebergang}[1]{\faktuebergangskip #1}

\newcommand{\faktuebergangleer}[1]{\faktuebergangskip \hspace{-0,15cm} }

\newcommand{\faktfolgerung}[1]{\faktfolgerungskip #1}

\newcommand{\faktzusatz}[1]{\faktzusatzskip #1}

%[[Projekt:Semantische Vorlagen/Beweisgestaltung/Latex]]

\newcommand{\teilbeweis}[5]{#1#2#3#4#5}

\newcommand{\fallunterscheidung}[5]{#1 #2\ifthenelse {\equal {#3}{}} {} { #3}\ifthenelse {\equal {#4}{}} {} { #4}\ifthenelse {\equal {#5}{}} {} { #5}}

\newcommand{\fallunterscheidungzwei}[2]{#1 #2}

\newcommand{\fallunterscheidungdrei}[3]{#1 #2 #3 }

\newcommand{\fallunterscheidungvier}[4]{#1 #2 #3 #4 }

\newcommand{\fallunterscheidungfuenf}[5]{#1 #2 #3 #4 #5 }

%[[Projekt:Semantische Vorlagen/Stichwortumsetzung/Latex]]

\newcommand{\betonung}[2]{\emph{#1}#2}

\newcommand{\stichwort}[2]{\emph{#1}#2}

\newcommand{\stichwortpraemath}[3]{$#1$-\emph{#2}#3}

\newcommand{\definitionswort}[2]{\emph{#1}#2}

\newcommand{\definitionswortpraemath}[3]{$#1$-\emph{#2}#3}

\newcommand{\definitionswortteil}[2]{\emph{#1}#2}

\newcommand{\definitionswortenp}[2]{\emph{#1}#2}

\newcommand{\definitionsverweis}[3]{#1#3}

\newcommand{\definitionsverweisanfuehrung}[3]{\anfuehrung {#1} {#3} }

%Anführungszeichen

\newcommand{\anfuehrung}[2]{"`{#1}"'#2}

\newcommand{\anfuehrungenglisch}[2]{``{#1}''#2}

%[[Projekt:Semantische Auflistungsvorlagen/Latex]]

%Aufzählungen

\newcommand{\aufzaehlungeins}[1]{\begin{enumerate} \item  #1 \end{enumerate}}

\newcommand{\aufzaehlungzwei}[2]{\begin{enumerate} \item  #1 \item  #2 \end{enumerate}}

\newcommand{\aufzaehlungdrei}[3]{\begin{enumerate} \item  #1 \item  #2 \item  #3  \end{enumerate}}

\newcommand{\aufzaehlungvier}[4]{\begin{enumerate} \item  #1 \item  #2 \item  #3 \item  #4 \end{enumerate}}

\newcommand{\aufzaehlungfuenf}[5]{\begin{enumerate} \item  #1 \item  #2 \item  #3 \item  #4 \item  #5  \end{enumerate}}

\newcommand{\aufzaehlungsechs}[6]{\begin{enumerate} \item  #1 \item  #2 \item  #3 \item  #4 \item  #5 \item  #6 \end{enumerate}}

\newcommand{\aufzaehlungsieben}[7]{\begin{enumerate} \item  #1 \item  #2 \item  #3 \item  #4 \item  #5 \item #6 \item  #7 \end{enumerate}}

\newcommand{\aufzaehlungacht}[8]{\begin{enumerate} \item  #1 \item  #2 \item  #3 \item  #4 \item  #5 \item #6 \item  #7 \item #8 \end{enumerate}}

\newcommand{\aufzaehlungneun}[9]{\begin{enumerate} \item  #1 \item  #2 \item  #3 \item  #4 \item  #5 \item #6 \item  #7  \item  #8  \item  #9 \end{enumerate}}

\newcommand{\itemfuenf}[5]{\item  #1 \item  #2 \item  #3 \item  #4 \item  #5}

\newcommand{\itemsechs}[6]{\item  #1 \item  #2 \item  #3 \item  #4 \item  #5 \item  #6}

\newcommand{\itemsieben}[7]{\item  #1 \item  #2 \item  #3 \item  #4 \item  #5 \item  #6 \item #7}

\newcommand{\aufzaehlungeinsabc}[1]{\begin{enumerate}[label=(\alph*)]  #1 \end{enumerate}}

\newcommand{\aufzaehlungzweiabc}[2]{\begin{enumerate}[label=(\alph*)] \item  #1 \item  #2 \end{enumerate}}

\newcommand{\aufzaehlungdreiabc}[3]{\begin{enumerate}[label=(\alph*)] \item  #1 \item  #2 \item  #3  \end{enumerate}}

\newcommand{\aufzaehlungvierabc}[4]{\begin{enumerate}[label=(\alph*)] \item  #1 \item  #2 \item  #3 \item  #4 \end{enumerate}}

\newcommand{\aufzaehlungfuenfabc}[5]{\begin{enumerate}[label=(\alph*)] \item  #1 \item  #2 \item  #3 \item  #4 \item  #5  \end{enumerate}}

\newcommand{\aufzaehlungsechsabc}[6]{\begin{enumerate}[label=(\alph*)] \item  #1 \item  #2 \item  #3 \item  #4 \item  #5 \item  #6 \end{enumerate}}

\newcommand{\aufzaehlungsiebenabc}[7]{\begin{enumerate}[label=(\alph*)] \item  #1 \item  #2 \item  #3 \item  #4 \item  #5 \item #6 \item  #7 \end{enumerate}}

\newcommand{\aufzaehlungachtabc}[8]{\begin{enumerate}[label=(\alph*)] \item  #1 \item  #2 \item  #3 \item  #4 \item  #5 \item #6 \item  #7 \item #8 \end{enumerate}}

\newcommand{\aufzaehlungneunabc}[9]{\begin{enumerate}[label=(\alph*)] \item  #1 \item  #2 \item  #3 \item  #4 \item  #5 \item #6 \item  #7  \item  #8  \item  #9 \end{enumerate}}

\newcommand{\aufzaehlungzweireihe}[2]{\begin{enumerate} #1 #2 \end{enumerate}}

\newcommand{\aufzaehlungdreibeweis}[3]{(1) #1 \par (2) #2 \par (3) #3 }

\newcommand{\aufzaehlungvierbeweis}[4]{(1) #1 \par (2) #2 \par (3) #3  \par (4) #4}

\newcommand{\aufzaehlungfuenfbeweis}[5]{(1) #1 \par (2) #2 \par (3) #3  \par (4) #4  \par (5) #5  }

%Auflistungen

\newcommand{\listitem}{\par \listskip \listdot}

\newcommand{\auflistungzwei}[2]{\par \listskip \listdot  #1 \par \listskip \listdot  #2 \par \listskip}

\newcommand{\auflistungdrei}[3]{\par \listskip \listdot  #1 \par \listskip \listdot  #2\par \listskip \listdot  #3\par\listskip }

\newcommand{\auflistungvier}[4]{\par \listskip \listdot  #1 \par \listskip \listdot  #2\par \listskip \listdot  #3\par \listskip \listdot  #4  \par \listskip }

\newcommand{\auflistungfuenf}[5]{\par \listskip \listdot  #1 \par \listskip \listdot  #2\par \listskip \listdot  #3\par \listskip \listdot  #4  \par \listskip  \listdot  #5 \par \listskip   }

\newcommand{\auflistungsechs}[6]{\par \listskip \listdot  #1 \par \listskip \listdot  #2\par \listskip \listdot  #3\par \listskip \listdot  #4  \par \listskip  \listdot  #5 \par \listskip \listdot  #6 \par \listskip}

\newcommand{\auflistungsieben}[7]{\par \listskip \listdot  #1 \par \listskip \listdot  #2\par \listskip \listdot  #3\par \listskip \listdot  #4  \par \listskip  \listdot  #5 \par \listskip \listdot  #6 \par \listskip \listdot  #7 \par \listskip }

\newcommand{\auflistungacht}[8]{\par \listskip \listdot  #1 \par \listskip \listdot  #2\par \listskip \listdot  #3\par \listskip \listdot  #4  \par \listskip  \listdot  #5 \par \listskip \listdot  #6 \par \listskip \listdot  #7 \par \listskip \listdot  #8  \par \listskip}

\newcommand{\auflistungneun}[9]{\par \listskip \listdot  #1 \par \listskip \listdot  #2\par \listskip \listdot  #3\par \listskip \listdot  #4  \par \listskip  \listdot  #5 \par \listskip \listdot  #6 \par \listskip \listdot  #7 \par \listskip \listdot  #8  \par \listskip  \listdot  #9 \par \listskip  }

\newcommand{\listitemfuenf}[5]{\listitem  #1 \listitem  #2 \listitem  #3 \listitem  #4 \listitem  #5}

\newcommand{\auflistungzweireihe}[2]{ #1 #2 }

%Gestaltung der Auflistungen

\newcommand{\listdot}{$\bullet \,$}

%Bildgestaltung

\newcommand{\inputbild}[1]{\includegraphics[scale]{#1}}

%Tabellengestaltung

%[[Projekt:Semantische Vorlagen/Zeilenkonfiguration/Latex]]

\newcommand{\zeileundzwei}[2]{#1 & #2}

\newcommand{\zeileunddrei}[3]{#1 & #2 & #3}

\newcommand{\zeileundvier}[4]{#1 & #2 & #3 & #4}

\newcommand{\zeileundfuenf}[5]{#1 & #2 & #3 & #4 & #5}

\newcommand{\zeileundsechs}[6]{#1 & #2 & #3 & #4 & #5 & #6}

\newcommand{\zeileundsieben}[7]{#1 & #2 & #3 & #4 & #5 & #6 & #7}

\newcommand{\zeileundacht}[8]{#1 & #2 & #3 & #4 & #5 & #6 & #7 & #8}

\newcommand{\zeileundneun}[9]{#1 & #2 & #3 & #4 & #5 & #6 & #7 & #8 & #9}

\newcommand{\mazeileundzwei}[2]{$#1$ & $#2$}

\newcommand{\mazeileunddrei}[3]{$#1$ & $ #2$ & $#3$}

\newcommand{\mazeileundvier}[4]{$#1$ & $#2$ & $#3$ & $#4$}

\newcommand{\mazeileundfuenf}[5]{$#1$ & $#2$ & $#3$ & $#4$ & $#5$}

\newcommand{\mazeileundsechs}[6]{$#1$ & $#2$ & $#3$ & $#4$ & $#5$ & $#6$}

\newcommand{\mazeileundsieben}[7]{$#1$ & $#2$ & $#3$ & $#4$ & $#5$ & $#6$ & $#7$}

\newcommand{\mazeileundacht}[8]{$#1$ & $#2$ & $#3$ & $#4$ & $#5$ & $#6$ & $#7$ & $#8$}

\newcommand{\mazeileundneun}[9]{$#1$ & $#2$ & $#3$ & $#4$ & $#5$ & $#6$ & $#7$ & $#8$ & $#9$}

%besser als Listen?
%Liste mit und (&) bzw. Bruch.

\newcommand {\listesechsmaund}[6]{$#1$ & $#2$& $#3$ & $#4$ & $#5$ & $#6$}

\newcommand {\listesiebenmaund}[7]{$#1$ & $#2$& $#3$ & $#4$ & $#5$ & $#6$ & $#7$}

\newcommand{\listezweiund}[2]{#1}

\newcommand{\listedreiund}[3]{#1}

\newcommand{\listesiebenund}[7]{#1}

\newcommand{\listesechsbruch}[6]{#1 \\ \hline #2 \\ \hline #3 \\ \hline #4 \\ \hline #5 \\ \hline #6}

\newcommand{\leitzeileunddrei}[3]{#1}

\newcommand{\madoppelzeileunddrei}[6]{$#1$&$#2$&$#3$\\ \hline $#4$&$#5$&$#6$}

%[[Projekt:Semantische Vorlagen/Tabellen/Latex]]

\newcommand{\tabelleviervier}[4]{\par \bigskip \begin{tabular}{|c|c|c|c|} \hline #1 \\ \hline #2 \\ \hline #3 \\ \hline #4 \\ \hline \end{tabular}}

\newcommand{\tabellefuenfvier}[5]{\par \bigskip \begin{tabular}{|c|c|c|c|c|} \hline #1 \\ \hline #2 \\ \hline #3 \\ \hline #4 \\ \hline #5  \\ \hline \end{tabular}}

%Mit math. Einträgen - ohne Leitzeile

\newcommand{\matabellezweivier}[5]{\begin{center} \begin{tabular}{|c|c|} \hline #2 \\ \hline #3 \\ \hline #4 \\ \hline #5  \\ \hline \end{tabular} \end{center} }

\newcommand{\matabellezweifuenf}[6]{\begin{center} \begin{tabular}{|c|c|} \hline     #2 \\ \hline #3 \\ \hline #4 \\ \hline #5 \\ \hline #6  \\ \hline \end{tabular} \end{center} }

\newcommand{\matabellezweisechs}[7]{\begin{center} \begin{tabular}{|c|c|} \hline     #2 \\ \hline #3 \\ \hline #4 \\ \hline #5 \\ \hline #6 \\ \hline #7   \\ \hline \end{tabular} \end{center} }

\newcommand{\matabellezweisieben}[8]{\begin{center} \begin{tabular}{|c|c|} \hline     #2 \\ \hline #3 \\ \hline #4 \\ \hline #5 \\ \hline #6 \\ \hline #7 \\ \hline #8  \\ \hline \end{tabular} \end{center} }

%Mit math. Einträgen - Mit Leitzeile

\newcommand{\matabellesiebenzwoelf}[3]{\begin{center} \begin{tabular}{|c||c|c|c|c|c|c|c|}  \hline  #1   \\ \hline \hline #2 \\ \hline #3 \\ \hline \end{tabular} \end{center} }

\newcommand{\matabellezwoelfzwoelf}[3]{\begin{center} \begin{tabular}{|c||c|c|c|c|c|c|c|c|c|c|c|c|}  \hline  #1   \\ \hline \hline #2 \\ \hline #3 \\ \hline \end{tabular} \end{center} }

\newcommand{\matabelledreineun}[5]{\begin{center} \begin{tabular}{|c|c|c|}\hline  #1 \\ \hline  #2 \\ \hline #3 \\ \hline   #4 \\ \hline  #5 \\ \hline \end{tabular} \end{center} }

\newcommand{\matabelledreizehn}[6]{\begin{center} \begin{tabular}{|c|c|c|}\hline  #1 \\ \hline  #2 \\ \hline #3 \\ \hline   #4 \\ \hline  #5 \\ \hline  #6 \\ \hline \end{tabular} \end{center} }

\newcommand{\matabelledreielf}[6]{\begin{center} \begin{tabular}{|c|c|c|}\hline  #1 \\ \hline  #2 \\ \hline #3 \\ \hline   #4 \\ \hline  #5 \\ \hline  #6 \\ \hline \end{tabular} \end{center} }

%[[Projekt:Semantische Vorlagen/Klausur/Punktetabellen/Latex]]

\newcommand{\punktetabellezehn}
{\vspace{3ex} \begin{center} \renewcommand{\arraystretch}{1.5} \tabcolsep=0.78ex \begin{tabular}{@{}|c|c|c|c|c|c|c|c|c|c|c|c|c|}
\hline Aufgabe & 1 & 2 & 3 & 4 & 5 & 6 & 7 & 8 & 9 & 10 & $\sum$
\\ \hline m\"ogliche Pkt. & \aeins & \azwei & \adrei & \avier & \afuenf & \asechs   & \asieben & \aacht  & \aneun & \azehn & \aelf
\\ \hline erhaltene Pkt.  & \phantom{11} & \phantom{11} & \phantom{11} & \phantom{11} & \phantom{11}& \phantom{11} & \phantom{11} & \phantom{11} & \phantom{11} & & \\ \hline \end{tabular} \end{center} \vspace{2ex} }

\newcommand{\punktetabelleelf}
{\vspace{3ex} \begin{center} \renewcommand{\arraystretch}{1.5} \tabcolsep=0.78ex \begin{tabular}{@{}|c|c|c|c|c|c|c|c|c|c|c|c|c|c|}
\hline Aufgabe & 1 & 2 & 3 & 4 & 5 & 6 & 7 & 8 & 9 & 10 & 11 & $\sum$
\\ \hline m\"ogliche Pkt. & \aeins & \azwei & \adrei & \avier & \afuenf & \asechs   & \asieben & \aacht  & \aneun & \azehn & \aelf & \azwoelf
\\ \hline erhaltene Pkt.  & \phantom{11} & \phantom{11} & \phantom{11} & \phantom{11} & \phantom{11}& \phantom{11} & \phantom{11} & \phantom{11} & \phantom{11} & \phantom{11} & & \\ \hline \end{tabular} \end{center} \vspace{2ex} }

\newcommand{\punktetabellezwoelf}
{\vspace{3ex} \begin{center} \renewcommand{\arraystretch}{1.5} \tabcolsep=0.78ex \begin{tabular}{@{}|c|c|c|c|c|c|c|c|c|c|c|c|c|c|c|}
\hline Aufgabe & 1 & 2 & 3 & 4 & 5 & 6 & 7 & 8 & 9 & 10 & 11 & 12    & $\sum$
\\ \hline m\"ogliche Pkt. & \aeins & \azwei & \adrei & \avier & \afuenf & \asechs   & \asieben & \aacht  & \aneun & \azehn & \aelf & \azwoelf & \adreizehn
\\ \hline erhaltene Pkt.  & \phantom{11} & \phantom{11} & \phantom{11} & \phantom{11} & \phantom{11}& \phantom{11} & \phantom{11} & \phantom{11}
& \phantom{11} & \phantom{11} & \phantom{11} & & \\ \hline \end{tabular} \end{center} \vspace{2ex} }

\newcommand{\punktetabelledreizehn}
{\vspace{3ex} \begin{center} \renewcommand{\arraystretch}{1.5} \tabcolsep=0.78ex \begin{tabular}{@{}|c|c|c|c|c|c|c|c|c|c|c|c|c|c|c|c|}
\hline Aufgabe & 1 & 2 & 3 & 4 & 5 & 6 & 7 & 8 & 9 & 10 & 11 & 12 & 13   & $\sum$
\\ \hline m\"ogliche Pkt. & \aeins & \azwei & \adrei & \avier & \afuenf & \asechs   & \asieben & \aacht  & \aneun & \azehn & \aelf & \azwoelf & \adreizehn & \avierzehn
\\ \hline erhaltene Pkt.  & \phantom{11} & \phantom{11} & \phantom{11} & \phantom{11} & \phantom{11}& \phantom{11} & \phantom{11} & \phantom{11} & \phantom{11}
& \phantom{11} & \phantom{11} & \phantom{11} & & \\ \hline \end{tabular} \end{center} \vspace{2ex} }

\newcommand{\punktetabellevierzehn}
{\vspace{3ex} \begin{center} \renewcommand{\arraystretch}{1.5} \tabcolsep=0.78ex \begin{tabular}{@{}|c|c|c|c|c|c|c|c|c|c|c|c|c|c|c|c|c|}
\hline Aufgabe & 1 & 2 & 3 & 4 & 5 & 6 & 7 & 8 & 9 & 10 & 11 & 12 & 13 & 14  & $\sum$
\\ \hline m\"ogliche Pkt. & \aeins & \azwei & \adrei & \avier & \afuenf & \asechs   & \asieben & \aacht  & \aneun & \azehn & \aelf & \azwoelf & \adreizehn & \avierzehn & \afuenfzehn
\\ \hline erhaltene Pkt. & \phantom{11} & \phantom{11} & \phantom{11} & \phantom{11} & \phantom{11} & \phantom{11}& \phantom{11} & \phantom{11} & \phantom{11} & \phantom{11}
& \phantom{11} & \phantom{11} & \phantom{11} & & \\ \hline \end{tabular} \end{center} \vspace{2ex} }

\newcommand{\punktetabellefuenfzehn}
{\vspace{3ex} \begin{center} \renewcommand{\arraystretch}{1.5} \tabcolsep=0.78ex \begin{tabular}{@{}|c|c|c|c|c|c|c|c|c|c|c|c|c|c|c|c|c|c|}
\hline Aufgabe & 1 & 2 & 3 & 4 & 5 & 6 & 7 & 8 & 9 & 10 & 11 & 12 & 13 & 14 & 15 & $\sum$
\\ \hline m\"ogliche Pkt. & \aeins & \azwei & \adrei & \avier & \afuenf & \asechs   & \asieben & \aacht  & \aneun & \azehn & \aelf & \azwoelf & \adreizehn & \avierzehn & \afuenfzehn
& \asechzehn
\\ \hline erhaltene Pkt. & \phantom{11} & \phantom{11} & \phantom{11}
& \phantom{11} & \phantom{11} & \phantom{11} & \phantom{11}& \phantom{11} & \phantom{11} & \phantom{11} & \phantom{11}
& \phantom{11} & \phantom{11} & \phantom{11} & & \\ \hline \end{tabular} \end{center} \vspace{2ex} }

\newcommand{\punktetabellesechzehn}
{\vspace{3ex} \begin{center} \renewcommand{\arraystretch}{1.5} \tabcolsep=0.78ex \begin{tabular}{@{}|c|c|c|c|c|c|c|c|c|c|c|c|c|c|c|c|c|c|}
\hline Aufgabe & 1 & 2 & 3 & 4 & 5 & 6 & 7 & 8 & 9 & 10 & 11 & 12 & 13 & 14 & 15 & 16 & $\sum$
\\ \hline m\"ogliche Pkt. & \aeins & \azwei & \adrei & \avier & \afuenf & \asechs   & \asieben & \aacht  & \aneun & \azehn & \aelf & \azwoelf & \adreizehn & \avierzehn & \afuenfzehn
& \asechzehn   & \asiebzehn
\\ \hline erhaltene Pkt. & \phantom{11} & \phantom{11} & \phantom{11} & \phantom{11}
& \phantom{11} & \phantom{11} & \phantom{11} & \phantom{11}& \phantom{11} & \phantom{11} & \phantom{11} & \phantom{11}
& \phantom{11} & \phantom{11} & \phantom{11} & & \\ \hline \end{tabular} \end{center} \vspace{2ex} }

\newcommand{\punktetabellesiebzehn}
{\vspace{3ex} \begin{center} \renewcommand{\arraystretch}{1.5} \tabcolsep=0.78ex \begin{tabular}{@{}|c|c|c|c|c|c|c|c|c|c|c|c|c|c|c|c|c|c|c|c|}
\hline Aufgabe & 1 & 2 & 3 & 4 & 5 & 6 & 7 & 8 & 9 & 10 & 11 & 12 & 13 & 14 & 15 & 16 & 17 & $\sum$
\\ \hline m\"ogliche Pkt. & \aeins & \azwei & \adrei & \avier & \afuenf & \asechs   & \asieben & \aacht  & \aneun & \azehn & \aelf & \azwoelf & \adreizehn & \avierzehn & \afuenfzehn
& \asechzehn   & \asiebzehn & \aachtzehn
\\ \hline erhaltene Pkt. & \phantom{11} & \phantom{11} & \phantom{11} & \phantom{11}
& \phantom{11} & \phantom{11} & \phantom{11} & \phantom{11} & \phantom{11}& \phantom{11} & \phantom{11} & \phantom{11} & \phantom{11}
& \phantom{11} & \phantom{11} & \phantom{11} & & \\ \hline \end{tabular} \end{center} \vspace{2ex} }

\newcommand{\punktetabelleachtzehn}
{\vspace{3ex} \begin{center} \renewcommand{\arraystretch}{1.5} \tabcolsep=0.78ex \begin{tabular}{@{}|c|c|c|c|c|c|c|c|c|c|c|c|c|c|c|c|c|c|c|c|c|}
\hline Aufgabe & 1 & 2 & 3 & 4 & 5 & 6 & 7 & 8 & 9 & 10 & 11 & 12 & 13 & 14 & 15 & 16 & 17 & 18 & $\sum$
\\ \hline m\"ogliche Pkt. & \aeins & \azwei & \adrei & \avier & \afuenf & \asechs   & \asieben & \aacht  & \aneun & \azehn & \aelf & \azwoelf & \adreizehn & \avierzehn & \afuenfzehn
& \asechzehn   & \asiebzehn & \aachtzehn & \aneunzehn
\\ \hline erhaltene Pkt. & \phantom{11} & \phantom{11} & \phantom{11} & \phantom{11} & \phantom{11}
& \phantom{11} & \phantom{11} & \phantom{11} & \phantom{11} & \phantom{11}& \phantom{11} & \phantom{11} & \phantom{11} & \phantom{11}
& \phantom{11} & \phantom{11} & \phantom{11} & & \\ \hline \end{tabular} \end{center} \vspace{2ex} }

\newcommand{\punktetabelleneunzehn}
{\vspace{3ex} \begin{center} \renewcommand{\arraystretch}{1.5} \tabcolsep=0.78ex \begin{tabular}{@{}|c|c|c|c|c|c|c|c|c|c|c|c|c|c|c|c|c|c|c|c|c|c|}
\hline Aufgabe & 1 & 2 & 3 & 4 & 5 & 6 & 7 & 8 & 9 & 10 & 11 & 12 & 13 & 14 & 15 & 16 & 17 & 18 & 19 & $\sum$
\\ \hline m\"ogliche Pkt. & \aeins & \azwei & \adrei & \avier & \afuenf & \asechs   & \asieben & \aacht  & \aneun & \azehn & \aelf & \azwoelf & \adreizehn & \avierzehn & \afuenfzehn
& \asechzehn   & \asiebzehn & \aachtzehn & \aneunzehn  & \azwanzig
\\ \hline erhaltene Pkt. & \phantom{11} & \phantom{11} & \phantom{11} & \phantom{11} & \phantom{11} & \phantom{11}
& \phantom{11} & \phantom{11} & \phantom{11} & \phantom{11} & \phantom{11}& \phantom{11} & \phantom{11} & \phantom{11} & \phantom{11}
& \phantom{11} & \phantom{11} & \phantom{11} & & \\ \hline \end{tabular} \end{center} \vspace{2ex} }

\newcommand{\punktetabellezwanzig}
{\vspace{3ex} \begin{center} \renewcommand{\arraystretch}{1.5} \tabcolsep=0.78ex \begin{tabular}{@{}|c|c|c|c|c|c|c|c|c|c|c|c|c|c|c|c|c|c|c|c|c|c|c|}
\hline Aufgabe & 1 & 2 & 3 & 4 & 5 & 6 & 7 & 8 & 9 & 10 & 11 & 12 & 13 & 14 & 15 & 16 & 17 & 18 & 19 & 20 & $\sum$
\\ \hline m\"ogliche Pkt. & \aeins & \azwei & \adrei & \avier & \afuenf & \asechs   & \asieben & \aacht  & \aneun & \azehn & \aelf & \azwoelf & \adreizehn & \avierzehn & \afuenfzehn
& \asechzehn   & \asiebzehn & \aachtzehn & \aneunzehn  & \azwanzig & \aeinundzwanzig
\\ \hline erhaltene Pkt. & \phantom{11} & \phantom{11} & \phantom{11} & \phantom{11} & \phantom{11} & \phantom{11} & \phantom{11}
& \phantom{11} & \phantom{11} & \phantom{11} & \phantom{11} & \phantom{11}& \phantom{11} & \phantom{11} & \phantom{11} & \phantom{11}
& \phantom{11} & \phantom{11} & \phantom{11} & & \\ \hline \end{tabular} \end{center} \vspace{2ex} }

\newcommand{\punktetabelleeinundzwanzig}
{\vspace{3ex} \begin{center} \renewcommand{\arraystretch}{1.5} \tabcolsep=0.78ex \begin{tabular}{@{}|c|c|c|c|c|c|c|c|c|c|c|c|c|c|c|c|c|c|c|c|c|c|c|c|}
\hline Aufgabe & 1 & 2 & 3 & 4 & 5 & 6 & 7 & 8 & 9 & 10 & 11 & 12 & 13 & 14 & 15 & 16 & 17 & 18 & 19 & 20 & 21 & $\sum$
\\ \hline m\"ogliche Pkt. & \aeins & \azwei & \adrei & \avier & \afuenf & \asechs   & \asieben & \aacht  & \aneun & \azehn & \aelf & \azwoelf & \adreizehn & \avierzehn & \afuenfzehn
& \asechzehn   & \asiebzehn & \aachtzehn & \aneunzehn  & \azwanzig & \aeinundzwanzig  & \azweiundzwanzig
\\ \hline erhaltene Pkt. & \phantom{11} & \phantom{11} & \phantom{11} & \phantom{11} & \phantom{11} & \phantom{11} & \phantom{11}
& \phantom{11} & \phantom{11} & \phantom{11} & \phantom{11} & \phantom{11}& \phantom{11} & \phantom{11} & \phantom{11} & \phantom{11}
& \phantom{11} & \phantom{11} & \phantom{11} & \phantom{11} & & \\ \hline \end{tabular} \end{center} \vspace{2ex} }

\newcommand{\punktetabellezweiundzwanzig}
{\vspace{3ex} \begin{center} \renewcommand{\arraystretch}{1.5} \tabcolsep=0.78ex \begin{tabular}{@{}|c|c|c|c|c|c|c|c|c|c|c|c|c|c|c|c|c|c|c|c|c|c|c|c|c|}
\hline Aufgabe & 1 & 2 & 3 & 4 & 5 & 6 & 7 & 8 & 9 & 10 & 11 & 12 & 13 & 14 & 15 & 16 & 17 & 18 & 19 & 20 & 21 & 22 & $\sum$
\\ \hline  Punkte & \aeins & \azwei & \adrei & \avier & \afuenf & \asechs   & \asieben & \aacht  & \aneun & \azehn & \aelf & \azwoelf & \adreizehn & \avierzehn & \afuenfzehn
& \asechzehn   & \asiebzehn & \aachtzehn & \aneunzehn  & \azwanzig & \aeinundzwanzig  & \azweiundzwanzig  & \adreiundzwanzig
\\ \hline Erhalten & \phantom{11} & \phantom{11} & \phantom{11} & \phantom{11} & \phantom{11} & \phantom{11} & \phantom{11}
& \phantom{11} & \phantom{11} & \phantom{11} & \phantom{11} & \phantom{11}& \phantom{11} & \phantom{11} & \phantom{11}  & \phantom{11} & \phantom{11}
& \phantom{11} & \phantom{11} & \phantom{11} & \phantom{11} & & \\ \hline \end{tabular} \end{center} \vspace{2ex} }

%[[Projekt:Semantische Vorlagen/Wahrheitstabellen/Latex]]

\newcommand{\wahrheitstabelleeins}[4]{
\begin{center}
\begin{tabular}{|c|c|}
\multicolumn{2}{l}{\textbf{#1}}\\
\hline
$p$ & $#2$\\ \hline
w  &	#3 \\ \hline
f  &	#4\\ \hline
\end{tabular}
\end{center} }

\newcommand{\wahrheitstabelleeinszwei}[7]{
\begin{center}
\begin{tabular}{|c|c|c|}
\multicolumn{3}{l}{\textbf{#1}}\\
\hline
$p$ & $#2$ & $#5$ \\ \hline
w  &	#3 &  #4 \\ \hline
f  &	#6 &  #7 \\ \hline
\end{tabular}
\end{center} }

\newcommand{\wahrheitstabelleeinsdrei}[4]{ \begin{center} \begin{tabular}{|c|c|c|c|} \multicolumn{4}{l}{\textbf{#1}} \\  \hline #2  \\  \hline #3  \\  \hline  #4  \\  \hline
\end{tabular} \end{center} }

%Wahrheitstabellen mit zwei Variablen p und q

\newcommand{\wahrheitstabellezweieins}[6]{ \begin{center} \begin{tabular}{|c|c|c|} \multicolumn{3}{l}{\textbf{#1}} \\  \hline  #2  \\  \hline  #3  \\  \hline  #4  \\  \hline  #5  \\  \hline  #6  \\  \hline \end{tabular} \end{center} }

\newcommand{\wahrheitstabellezweizwei}[6]{ \begin{center} \begin{tabular}{|c|c|c|c|} \multicolumn{4}{l}{\textbf{#1}} \\  \hline #2  \\  \hline #3  \\  \hline #4  \\  \hline #5  \\  \hline #6  \\  \hline\end{tabular} \end{center} }

\newcommand{\wahrheitstabellezweidrei}[6]{ \begin{center} \begin{tabular}{|c|c|c|c|c|} \multicolumn{5}{l}{\textbf{#1}} \\  \hline #2  \\  \hline #3  \\  \hline #4  \\  \hline #5  \\  \hline #6  \\  \hline \end{tabular} \end{center} }

\newcommand{\wahrheitstabellezweivier}[6]{ \begin{center} \begin{tabular}{|c|c|c|c|c|c|} \multicolumn{6}{l}{\textbf{#1}} \\  \hline #2  \\  \hline #3  \\  \hline #4  \\  \hline  #5  \\  \hline  #6  \\  \hline \end{tabular} \end{center} }

\newcommand{\wahrheitstabellezweifuenf}[6]{ \begin{center} \begin{tabular}{|c|c|c|c|c|c|c|} \multicolumn{7}{l}{\textbf{#1}} \\  \hline #2  \\  \hline #3  \\  \hline #4  \\  \hline #5  \\  \hline #6  \\  \hline
\end{tabular} \end{center} }

%Wahrheitstabellen mit drei Variablen p,q,r

\newcommand{\wahrheitstabelledreieins}[6]{ \begin{center} \begin{tabular}{|c|c|c|c||} \multicolumn{4}{l}{\textbf{#1}} \\  \hline #2  \\  \hline #3  \\  \hline #4  \\  \hline #5  \\  \hline #6 \\  \hline
\end{tabular} \end{center} }

%[[Projekt:Semantische Vorlagen/Wertetabellen/Latex]]

%Wertetabellen

\newcommand{\wertetabellevorskip}{\par \bigskip}

\newcommand{\wertetabellenachskip}{\par \bigskip}

\newcommand{\wertetabelleeinsausteilzeilen}[4]{\wertetabellevorskip \begin{tabular}{|c|c|} \hline #1 & #2  \\ \hline  #3 & #4  \\ \hline \end{tabular} \wertetabellenachskip }

\newcommand{\wertetabellezweiausteilzeilen}[4]{\wertetabellevorskip \begin{tabular}{|c|c|c|} \hline #1 & #2  \\ \hline  #3 & #4  \\ \hline \end{tabular} \wertetabellenachskip }

\newcommand{\wertetabelledreiausteilzeilen}[4]{\wertetabellevorskip \begin{tabular}{|c|c|c|c|} \hline #1 & #2  \\ \hline  #3 & #4  \\ \hline \end{tabular} \wertetabellenachskip }

\newcommand{\wertetabellevierausteilzeilen}[4]{\wertetabellevorskip \begin{tabular}{|c|c|c|c|c|} \hline #1 & #2  \\ \hline  #3 & #4  \\ \hline \end{tabular} \wertetabellenachskip }

\newcommand{\wertetabellefuenfausteilzeilen}[4]{\wertetabellevorskip \begin{tabular}{|c|c|c|c|c|c|} \hline #1 & #2  \\ \hline  #3 & #4  \\ \hline \end{tabular} \wertetabellenachskip }

\newcommand{\wertetabellesechsausteilzeilen}[6]{\wertetabellevorskip \begin{tabular}{|c|c|c|c|c|c|c|} \hline #1 & #2 & #3 \\ \hline  #4 & #5 & #6 \\ \hline \end{tabular} \wertetabellenachskip }

\newcommand{\wertetabellesiebenausteilzeilen}[6]{\wertetabellevorskip \begin{tabular}{|c|c|c|c|c|c|c|c|} \hline #1 & #2 & #3 \\ \hline  #4 & #5 & #6 \\ \hline \end{tabular} \wertetabellenachskip }

\newcommand{\wertetabelleachtausteilzeilen}[6]{\wertetabellevorskip \begin{tabular}{|c|c|c|c|c|c|c|c|c|} \hline #1 & #2 & #3 \\ \hline  #4 & #5 & #6 \\ \hline \end{tabular} \wertetabellenachskip }

\newcommand{\wertetabelleneunausteilzeilen}[6]{\wertetabellevorskip \begin{tabular}{|c|c|c|c|c|c|c|c|c|c|} \hline #1 & #2 & #3 \\ \hline  #4 & #5 & #6 \\ \hline \end{tabular} \wertetabellenachskip }

\newcommand{\wertetabellezehnausteilzeilen}[6]{\wertetabellevorskip \begin{tabular}{|c|c|c|c|c|c|c|c|c|c|c|} \hline #1 & #2 & #3 \\ \hline  #4 & #5 & #6 \\ \hline \end{tabular} \wertetabellenachskip }

\newcommand{\wertetabelleelfausteilzeilen}[8]{\wertetabellevorskip \begin{tabular}{|c|c|c|c|c|c|c|c|c|c|c|c|} \hline #1 & #2 & #3 & #4 \\ \hline #5 & #6 & #7 & #8 \\ \hline \end{tabular} \wertetabellenachskip }

%[[Projekt:Semantische Vorlagen/Verknüpfungstabellen/Latex]]

%Verknüpfungstabellen

\newcommand{\verknuepfungstabelleeins}[2]{\begin{center}
\begin{tabular}{|c||c|}
\hline #1 \\  \hline \hline #2 \\
\hline
\end{tabular}
\end{center} }

\newcommand{\verknuepfungstabellezwei}[3]{\begin{center}
\begin{tabular}{|c||c|c|}
\hline #1 \\  \hline \hline #2 \\ \hline #3 \\
\hline
\end{tabular}
\end{center} }

\newcommand{\verknuepfungstabelledrei}[4]{\begin{center}
\begin{tabular}{|c||c|c|c|}
\hline #1 \\  \hline \hline #2 \\ \hline #3 \\ \hline #4 \\
\hline
\end{tabular}
\end{center} }

\newcommand{\verknuepfungstabellevier}[5]{\begin{center}
\begin{tabular}{|c||c|c|c|c|}
\hline #1 \\  \hline \hline #2 \\ \hline #3 \\ \hline #4 \\  \hline  #5 \\
\hline
\end{tabular}
\end{center} }

\newcommand{\verknuepfungstabellefuenf}[6]{\begin{center}
\begin{tabular}{|c||c|c|c|c|c|}
\hline #1 \\  \hline \hline #2 \\ \hline #3 \\ \hline #4 \\  \hline  #5 \\ \hline  #6 \\
\hline
\end{tabular}
\end{center} }

\newcommand{\verknuepfungstabellesechs}[7]{\begin{center}
\begin{tabular}{|c||c|c|c|c|c|c|}
\hline #1 \\  \hline \hline #2 \\ \hline #3 \\ \hline #4 \\  \hline  #5 \\ \hline  #6 \\  \hline #7 \\
\hline
\end{tabular}
\end{center} }

\newcommand{\verknuepfungstabellesieben}[8]{\begin{center}
\begin{tabular}{|c||c|c|c|c|c|c|c|c|}
\hline #1 \\  \hline \hline #2 \\ \hline #3 \\ \hline #4 \\  \hline  #5 \\ \hline  #6 \\  \hline #7 \\  \hline  #8 \\
\hline
\end{tabular}
\end{center} }

\newcommand{\verknuepfungstabelleacht}[9]{\begin{center}
\begin{tabular}{|c||c|c|c|c|c|c|c|c|}
\hline #1 \\  \hline \hline #2 \\ \hline #3 \\ \hline #4 \\  \hline  #5 \\ \hline  #6 \\  \hline #7 \\  \hline  #8 \\ \hline  #9 \\
\hline
\end{tabular}
\end{center} }

%[[Projekt:Semantische Vorlagen/Initialisierung für Daten/Latex]]

\newcommand{\aeins}{    }

\newcommand{\azwei}{     }

\newcommand{\adrei}{    }

\newcommand{\avier}{    }

\newcommand{\afuenf}{    }

\newcommand{\asechs}{    }

\newcommand{\asieben}{    }

\newcommand{\aacht}{    }

\newcommand{\aneun}{    }

\newcommand{\azehn}{    }

\newcommand{\aelf}{    }

\newcommand{\azwoelf}{    }

\newcommand{\adreizehn}{     }

\newcommand{\avierzehn}{    }

\newcommand{\afuenfzehn}{    }

\newcommand{\asechzehn}{    }

\newcommand{\asiebzehn}{    }

\newcommand{\aachtzehn}{    }

\newcommand{\aneunzehn}{    }

\newcommand{\azwanzig}{    }

\newcommand{\aeinundzwanzig}{    }

\newcommand{\azweiundzwanzig}{    }

\newcommand{\adreiundzwanzig}{    }

\newcommand{\avierundzwanzig}{    }

\newcommand{\afuenfundzwanzig}{    }

\newcommand{\asechsundzwanzig}{    }

\newcommand{\asiebenundzwanzig}{    }

\newcommand{\aachtundzwanzig}{    }

\newcommand{\aneunundzwanzig}{    }

\newcommand{\leitzeilenull}{}

\newcommand{\leitzeileeins}{}

\newcommand{\leitzeilezwei}{}

\newcommand{\leitzeiledrei}{}

\newcommand{\leitzeilevier}{}

\newcommand{\leitzeilefuenf}{}

\newcommand{\leitzeilesechs}{}

\newcommand{\leitzeilesieben}{}

\newcommand{\leitzeileacht}{}

\newcommand{\leitzeileneun}{}

\newcommand{\leitzeilezehn}{}

\newcommand{\leitzeileelf}{}

\newcommand{\leitzeilezwoelf}{}

\newcommand{\leitspaltenull}{}

\newcommand{\leitspalteeins}{}

\newcommand{\leitspaltezwei}{}

\newcommand{\leitspaltedrei}{}

\newcommand{\leitspaltevier}{}

\newcommand{\leitspaltefuenf}{}

\newcommand{\leitspaltesechs}{}

\newcommand{\leitspaltesieben}{}

\newcommand{\leitspalteacht}{}

\newcommand{\leitspalteneun}{}

\newcommand{\leitspaltezehn}{}

\newcommand{\leitspalteelf}{}

\newcommand{\leitspaltezwoelf}{}

\newcommand{\leitspaltedreizehn}{}

\newcommand{\leitspaltevierzehn}{}

\newcommand{\leitspaltefuenfzehn}{}

\newcommand{\leitspaltesechzehn}{}

\newcommand{\leitspaltesiebzehn}{}

\newcommand{\leitspalteachtzehn}{}

\newcommand{\leitspalteneunzehn}{}

\newcommand{\leitspaltezwanzig}{}

\newcommand{\aeinsxeins}{}

\newcommand{\aeinsxzwei}{}

\newcommand{\aeinsxdrei}{}

\newcommand{\aeinsxvier}{}

\newcommand{\aeinsxfuenf}{}

\newcommand{\aeinsxsechs}{}

\newcommand{\aeinsxsieben}{}

\newcommand{\aeinsxacht}{}

\newcommand{\aeinsxneun}{}

\newcommand{\aeinsxzehn}{}

\newcommand{\aeinsxelf}{}

\newcommand{\aeinsxzwoelf}{}

\newcommand{\azweixeins}{}

\newcommand{\azweixzwei}{}

\newcommand{\azweixdrei}{}

\newcommand{\azweixvier}{}

\newcommand{\azweixfuenf}{}

\newcommand{\azweixsechs}{}

\newcommand{\azweixsieben}{}

\newcommand{\azweixacht}{}

\newcommand{\azweixneun}{}

\newcommand{\azweixzehn}{}

\newcommand{\azweixelf}{}

\newcommand{\azweixzwoelf}{}

\newcommand{\adreixeins}{}

\newcommand{\adreixzwei}{}

\newcommand{\adreixdrei}{}

\newcommand{\adreixvier}{}

\newcommand{\adreixfuenf}{}

\newcommand{\adreixsechs}{}

\newcommand{\adreixsieben}{}

\newcommand{\adreixacht}{}

\newcommand{\adreixneun}{}

\newcommand{\adreixzehn}{}

\newcommand{\adreixelf}{}

\newcommand{\adreixzwoelf}{}

\newcommand{\avierxeins}{}

\newcommand{\avierxzwei}{}

\newcommand{\avierxdrei}{}

\newcommand{\avierxvier}{}

\newcommand{\avierxfuenf}{}

\newcommand{\avierxsechs}{}

\newcommand{\avierxsieben}{}

\newcommand{\avierxacht}{}

\newcommand{\avierxneun}{}

\newcommand{\avierxzehn}{}

\newcommand{\avierxelf}{}

\newcommand{\avierxzwoelf}{}

\newcommand{\afuenfxeins}{}

\newcommand{\afuenfxzwei}{}

\newcommand{\afuenfxdrei}{}

\newcommand{\afuenfxvier}{}

\newcommand{\afuenfxfuenf}{}

\newcommand{\afuenfxsechs}{}

\newcommand{\afuenfxsieben}{}

\newcommand{\afuenfxacht}{}

\newcommand{\afuenfxneun}{}

\newcommand{\afuenfxzehn}{}

\newcommand{\afuenfxelf}{}

\newcommand{\afuenfxzwoelf}{}

\newcommand{\asechsxeins}{}

\newcommand{\asechsxzwei}{}

\newcommand{\asechsxdrei}{}

\newcommand{\asechsxvier}{}

\newcommand{\asechsxfuenf}{}

\newcommand{\asechsxsechs}{}

\newcommand{\asechsxsieben}{}

\newcommand{\asechsxacht}{}

\newcommand{\asechsxneun}{}

\newcommand{\asechsxzehn}{}

\newcommand{\asechsxelf}{}

\newcommand{\asechsxzwoelf}{}

\newcommand{\asiebenxeins}{}

\newcommand{\asiebenxzwei}{}

\newcommand{\asiebenxdrei}{}

\newcommand{\asiebenxvier}{}

\newcommand{\asiebenxfuenf}{}

\newcommand{\asiebenxsechs}{}

\newcommand{\asiebenxsieben}{}

\newcommand{\asiebenxacht}{}

\newcommand{\asiebenxneun}{}

\newcommand{\asiebenxzehn}{}

\newcommand{\asiebenxelf}{}

\newcommand{\asiebenxzwoelf}{}

\newcommand{\aachtxeins}{}

\newcommand{\aachtxzwei}{}

\newcommand{\aachtxdrei}{}

\newcommand{\aachtxvier}{}

\newcommand{\aachtxfuenf}{}

\newcommand{\aachtxsechs}{}

\newcommand{\aachtxsieben}{}

\newcommand{\aachtxacht}{}

\newcommand{\aachtxneun}{}

\newcommand{\aachtxzehn}{}

\newcommand{\aachtxelf}{}

\newcommand{\aachtxzwoelf}{}

\newcommand{\aneunxeins}{}

\newcommand{\aneunxzwei}{}

\newcommand{\aneunxdrei}{}

\newcommand{\aneunxvier}{}

\newcommand{\aneunxfuenf}{}

\newcommand{\aneunxsechs}{}

\newcommand{\aneunxsieben}{}

\newcommand{\aneunxacht}{}

\newcommand{\aneunxneun}{}

\newcommand{\aneunxzehn}{}

\newcommand{\aneunxelf}{}

\newcommand{\aneunxzwoelf}{}

\newcommand{\azehnxeins}{}

\newcommand{\azehnxzwei}{}

\newcommand{\azehnxdrei}{}

\newcommand{\azehnxvier}{}

\newcommand{\azehnxfuenf}{}

\newcommand{\azehnxsechs}{}

\newcommand{\azehnxsieben}{}

\newcommand{\azehnxacht}{}

\newcommand{\azehnxneun}{}

\newcommand{\azehnxzehn}{}

\newcommand{\azehnxelf}{}

\newcommand{\azehnxzwoelf}{}

\newcommand{\aelfxeins}{}

\newcommand{\aelfxzwei}{}

\newcommand{\aelfxdrei}{}

\newcommand{\aelfxvier}{}

\newcommand{\aelfxfuenf}{}

\newcommand{\aelfxsechs}{}

\newcommand{\aelfxsieben}{}

\newcommand{\aelfxacht}{}

\newcommand{\aelfxneun}{}

\newcommand{\aelfxzehn}{}

\newcommand{\aelfxelf}{}

\newcommand{\aelfxzwoelf}{}

\newcommand{\azwoelfxeins}{}

\newcommand{\azwoelfxzwei}{}

\newcommand{\azwoelfxdrei}{}

\newcommand{\azwoelfxvier}{}

\newcommand{\azwoelfxfuenf}{}

\newcommand{\azwoelfxsechs}{}

\newcommand{\azwoelfxsieben}{}

\newcommand{\azwoelfxacht}{}

\newcommand{\azwoelfxneun}{}

\newcommand{\azwoelfxzehn}{}

\newcommand{\azwoelfxelf}{}

\newcommand{\azwoelfxzwoelf}{}

\newcommand{\adreizehnxeins}{}

\newcommand{\adreizehnxzwei}{}

\newcommand{\adreizehnxdrei}{}

\newcommand{\adreizehnxvier}{}

\newcommand{\adreizehnxfuenf}{}

\newcommand{\adreizehnxsechs}{}

\newcommand{\adreizehnxsieben}{}

\newcommand{\adreizehnxacht}{}

\newcommand{\adreizehnxneun}{}

\newcommand{\adreizehnxzehn}{}

\newcommand{\adreizehnxelf}{}

\newcommand{\adreizehnxzwoelf}{}

\newcommand{\avierzehnxeins}{}

\newcommand{\avierzehnxzwei}{}

\newcommand{\avierzehnxdrei}{}

\newcommand{\avierzehnxvier}{}

\newcommand{\avierzehnxfuenf}{}

\newcommand{\avierzehnxsechs}{}

\newcommand{\avierzehnxsieben}{}

\newcommand{\avierzehnxacht}{}

\newcommand{\avierzehnxneun}{}

\newcommand{\avierzehnxzehn}{}

\newcommand{\avierzehnxelf}{}

\newcommand{\avierzehnxzwoelf}{}

\newcommand{\afuenfzehnxeins}{}

\newcommand{\afuenfzehnxzwei}{}

\newcommand{\afuenfzehnxdrei}{}

\newcommand{\afuenfzehnxvier}{}

\newcommand{\afuenfzehnxfuenf}{}

\newcommand{\afuenfzehnxsechs}{}

\newcommand{\afuenfzehnxsieben}{}

\newcommand{\afuenfzehnxacht}{}

\newcommand{\afuenfzehnxneun}{}

\newcommand{\afuenfzehnxzehn}{}

\newcommand{\afuenfzehnxelf}{}

\newcommand{\afuenfzehnxzwoelf}{}

\newcommand{\asechzehnxeins}{}

\newcommand{\asechzehnxzwei}{}

\newcommand{\asechzehnxdrei}{}

\newcommand{\asechzehnxvier}{}

\newcommand{\asechzehnxfuenf}{}

\newcommand{\asechzehnxsechs}{}

\newcommand{\asechzehnxsieben}{}

\newcommand{\asechzehnxacht}{}

\newcommand{\asechzehnxneun}{}

\newcommand{\asechzehnxzehn}{}

\newcommand{\asechzehnxelf}{}

\newcommand{\asechzehnxzwoelf}{}

\newcommand{\asiebzehnxeins}{}

\newcommand{\asiebzehnxzwei}{}

\newcommand{\asiebzehnxdrei}{}

\newcommand{\asiebzehnxvier}{}

\newcommand{\asiebzehnxfuenf}{}

\newcommand{\asiebzehnxsechs}{}

\newcommand{\asiebzehnxsieben}{}

\newcommand{\asiebzehnxacht}{}

\newcommand{\asiebzehnxneun}{}

\newcommand{\asiebzehnxzehn}{}

\newcommand{\asiebzehnxelf}{}

\newcommand{\asiebzehnxzwoelf}{}

\newcommand{\aachtzehnxeins}{}

\newcommand{\aachtzehnxzwei}{}

\newcommand{\aachtzehnxdrei}{}

\newcommand{\aachtzehnxvier}{}

\newcommand{\aachtzehnxfuenf}{}

\newcommand{\aachtzehnxsechs}{}

\newcommand{\aachtzehnxsieben}{}

\newcommand{\aachtzehnxacht}{}

\newcommand{\aachtzehnxneun}{}

\newcommand{\aachtzehnxzehn}{}

\newcommand{\aachtzehnxelf}{}

\newcommand{\aachtzehnxzwoelf}{}

%[[Projekt:Semantische Vorlagen/Tabellen mit benannten Parametern/Latex]]

%Tabellen mit zwei Zeilen

\newcommand{\tabelleleitzweixzwei}{

\begin{center}

\begin{tabular}{|c||c|c|}

\hline  \leitzeilenull  &  \leitzeileeins  &  \leitzeilezwei   \\

\hline \hline  \leitspalteeins  &  $ \aeinsxeins $  &  $ \aeinsxzwei $ \\

\hline  \leitspaltezwei  &  $ \azweixeins $  &  $ \azweixzwei $   \\

\hline

\end{tabular}

\end{center} }

%Tabellen mit drei Zeilen

\newcommand{\tabelleleitdreixdrei}{

\begin{center}

\begin{tabular}{|c||c|c|c||}

\hline  \leitzeilenull  &  \leitzeileeins  &  \leitzeilezwei  &  \leitzeiledrei        \\

\hline \hline  \leitspalteeins  &  $ \aeinsxeins $  &  $ \aeinsxzwei $  &  $ \aeinsxdrei $      \\

\hline  \leitspaltezwei  &  $ \azweixeins $  &  $ \azweixzwei $  &  $ \azweixdrei $       \\

\hline  \leitspaltedrei  &  $ \adreixeins $  &  $ \adreixzwei $  &  $ \adreixdrei $       \\

\hline

\end{tabular}

\end{center} }

\newcommand{\tabelleleitdreixvier}{

\begin{center}

\begin{tabular}{|c||c|c|c|c||}

\hline  \leitzeilenull  &  \leitzeileeins  &  \leitzeilezwei  &  \leitzeiledrei  &  \leitzeilevier      \\

\hline \hline  \leitspalteeins  &  $ \aeinsxeins $  &  $ \aeinsxzwei $  &  $ \aeinsxdrei $  &  $ \aeinsxvier $   \\

\hline  \leitspaltezwei  &  $ \azweixeins $  &  $ \azweixzwei $  &  $ \azweixdrei $  &  $ \azweixvier $     \\

\hline  \leitspaltedrei  &  $ \adreixeins $  &  $ \adreixzwei $  &  $ \adreixdrei $  &  $ \adreixvier $     \\

\hline

\end{tabular}

\end{center} }

\newcommand{\tabelleleitdreixsechs}{

\begin{center}

\begin{tabular}{|c||c|c|c|c|c|c|}

\hline  \leitzeilenull  &  \leitzeileeins  &  \leitzeilezwei  &  \leitzeiledrei  &  \leitzeilevier  &  \leitzeilefuenf   &  \leitzeilesechs    \\

\hline \hline  \leitspalteeins  &  $ \aeinsxeins $  &  $ \aeinsxzwei $  &  $ \aeinsxdrei $  &  $ \aeinsxvier $  &  $ \aeinsxfuenf $  &  $ \aeinsxsechs $  \\

\hline  \leitspaltezwei  &  $ \azweixeins $  &  $ \azweixzwei $  &  $ \azweixdrei $  &  $ \azweixvier $  &  $ \azweixfuenf $   &  $ \azweixsechs $  \\

\hline  \leitspaltedrei  &  $ \adreixeins $  &  $ \adreixzwei $  &  $ \adreixdrei $  &  $ \adreixvier $  &  $ \adreixfuenf $   &  $ \adreixsechs $   \\

\hline

\end{tabular}

\end{center} }

\newcommand{\tabelleleitdreixelf}{

\begin{center}

\begin{tabular}{|c||c|c|c|c|c|c|c|c|c|c|c|}

\hline  \leitzeilenull  &  \leitzeileeins  &  \leitzeilezwei  &  \leitzeiledrei  &  \leitzeilevier  &  \leitzeilefuenf   &  \leitzeilesechs  &  \leitzeilesieben  &  \leitzeileacht  &  \leitzeileneun  &  \leitzeilezehn &  \leitzeileelf    \\

\hline \hline  \leitspalteeins  &  $ \aeinsxeins $  &  $ \aeinsxzwei $  &  $ \aeinsxdrei $  &  $ \aeinsxvier $  &  $ \aeinsxfuenf $  &  $ \aeinsxsechs $  &  $ \aeinsxsieben $  &  $ \aeinsxacht $  &  $ \aeinsxneun $ &  $ \aeinsxzehn $  &  $ \aeinsxelf $  \\

\hline  \leitspaltezwei  &  $ \azweixeins $  &  $ \azweixzwei $  &  $ \azweixdrei $  &  $ \azweixvier $  &  $ \azweixfuenf $   &  $ \azweixsechs $  &  $ \azweixsieben $  &  $ \azweixacht $  &  $ \azweixneun $  &  $ \azweixzehn $ &  $ \azweixelf $     \\

\hline  \leitspaltedrei  &  $ \adreixeins $  &  $ \adreixzwei $  &  $ \adreixdrei $  &  $ \adreixvier $  &  $ \adreixfuenf $   &  $ \adreixsechs $  &  $ \adreixsieben $  &  $ \adreixacht $  &  $ \adreixneun $  &  $ \adreixzehn $     &  $ \adreixelf $   \\

\hline

\end{tabular}

\end{center} }

%Tabellen mit vier Zeilen

\newcommand{\tabelleleitvierxzwei}{

\begin{center}

\begin{tabular}{|c||c|c|}

\hline  \leitzeilenull  &  \leitzeileeins  &  \leitzeilezwei   \\

\hline \hline  \leitspalteeins  &  $ \aeinsxeins $  &  $ \aeinsxzwei $    \\

\hline  \leitspaltezwei  &  $ \azweixeins $  &  $ \azweixzwei $    \\

\hline  \leitspaltedrei  &  $ \adreixeins $  &  $ \adreixzwei $    \\

\hline  \leitspaltevier  &  $ \avierxeins $  &  $ \avierxzwei $    \\

\hline

\end{tabular}

\end{center} }

\newcommand{\tabelleleitvierxdrei}{

\begin{center}

\begin{tabular}{|c||c|c|c|}

\hline  \leitzeilenull  &  \leitzeileeins  &  \leitzeilezwei  &  \leitzeiledrei     \\

\hline \hline  \leitspalteeins  &  $ \aeinsxeins $  &  $ \aeinsxzwei $  &  $ \aeinsxdrei $     \\

\hline  \leitspaltezwei  &  $ \azweixeins $  &  $ \azweixzwei $  &  $ \azweixdrei $      \\

\hline  \leitspaltedrei  &  $ \adreixeins $  &  $ \adreixzwei $  &  $ \adreixdrei $    \\

\hline  \leitspaltevier  &  $ \avierxeins $  &  $ \avierxzwei $  &  $ \avierxdrei $    \\

\hline

\end{tabular}

\end{center} }

\newcommand{\tabelleleitvierxvier}{

\begin{center}

\begin{tabular}{|c||c|c|c|c|}

\hline  \leitzeilenull  &  \leitzeileeins  &  \leitzeilezwei  &  \leitzeiledrei  &  \leitzeilevier    \\

\hline \hline  \leitspalteeins  &  $ \aeinsxeins $  &  $ \aeinsxzwei $  &  $ \aeinsxdrei $  &  $ \aeinsxvier $   \\

\hline  \leitspaltezwei  &  $ \azweixeins $  &  $ \azweixzwei $  &  $ \azweixdrei $  &  $ \azweixvier $     \\

\hline  \leitspaltedrei  &  $ \adreixeins $  &  $ \adreixzwei $  &  $ \adreixdrei $  &  $ \adreixvier $   \\

\hline  \leitspaltevier  &  $ \avierxeins $  &  $ \avierxzwei $  &  $ \avierxdrei $  &  $ \avierxvier $   \\

\hline

\end{tabular}

\end{center} }

\newcommand{\tabelleleittextvierxzwei}{

	\begin{center}

		\begin{tabular}{|c||c|c|}

			\hline  \leitzeilenull  &  \leitzeileeins  &  \leitzeilezwei   \\

			\hline \hline  \leitspalteeins  &   \aeinsxeins   &   \aeinsxzwei    \\

			\hline  \leitspaltezwei  &  \azweixeins   &   \azweixzwei     \\

			\hline  \leitspaltedrei  &  \adreixeins   &   \adreixzwei     \\

			\hline  \leitspaltevier  &  \avierxeins   &   \avierxzwei     \\

			\hline

		\end{tabular}

	\end{center} }

\newcommand{\tabelleleitvierxacht}{

\begin{center}

\begin{tabular}{|c||c|c|c|c|c|c|c|c|}

\hline  \leitzeilenull  &  \leitzeileeins  &  \leitzeilezwei  &  \leitzeiledrei  &  \leitzeilevier  &  \leitzeilefuenf  &  \leitzeilesechs  &  \leitzeilesieben  &  \leitzeileacht  \\

\hline \hline  \leitspalteeins  &  $ \aeinsxeins $  &  $ \aeinsxzwei $  &  $ \aeinsxdrei $  &  $ \aeinsxvier $  &  $ \aeinsxfuenf $  &  $ \aeinsxsechs $  &  $ \aeinsxsieben $  &  $ \aeinsxacht $  \\

\hline  \leitspaltezwei  &  $ \azweixeins $  &  $ \azweixzwei $  &  $ \azweixdrei $  &  $ \azweixvier $   &  $ \azweixfuenf $  &  $ \azweixsechs $  &  $ \azweixsieben $  &  $ \azweixacht $   \\

\hline  \leitspaltedrei  &  $ \adreixeins $  &  $ \adreixzwei $  &  $ \adreixdrei $  &  $ \adreixvier $ &  $ \adreixfuenf $  &  $ \adreixsechs $  &  $ \adreixsieben $  &  $ \adreixacht $   \\

\hline  \leitspaltevier  &  $ \avierxeins $  &  $ \avierxzwei $  &  $ \avierxdrei $  &  $ \avierxvier $  &  $ \avierxfuenf $  &  $ \avierxsechs $  &  $ \avierxsieben $  &  $ \avierxacht $  \\

\hline

\end{tabular}

\end{center} }

%Tabellen mit fünf Zeilen

\newcommand{\tabelleleitfuenfxfuenf}{

\begin{center}

\begin{tabular}{|c||c|c|c|c|c|}

\hline  \leitzeilenull  &  \leitzeileeins  &  \leitzeilezwei  &  \leitzeiledrei  &  \leitzeilevier  &  \leitzeilefuenf   \\

\hline \hline  \leitspalteeins  &  $ \aeinsxeins $  &  $ \aeinsxzwei $  &  $ \aeinsxdrei $  &  $ \aeinsxvier $  &  $ \aeinsxfuenf $   \\

\hline  \leitspaltezwei  &  $ \azweixeins $  &  $ \azweixzwei $  &  $ \azweixdrei $  &  $ \azweixvier $  &  $ \azweixfuenf $     \\

\hline  \leitspaltedrei  &  $ \adreixeins $  &  $ \adreixzwei $  &  $ \adreixdrei $  &  $ \adreixvier $  &  $ \adreixfuenf $      \\

\hline  \leitspaltevier  &  $ \avierxeins $  &  $ \avierxzwei $  &  $ \avierxdrei $  &  $ \avierxvier $  &  $ \avierxfuenf $     \\

\hline  \leitspaltefuenf  &  $ \afuenfxeins $  &  $ \afuenfxzwei $  &  $ \afuenfxdrei $  &  $ \afuenfxvier $  &  $ \afuenfxfuenf $    \\

\hline

\end{tabular}

\end{center} }

%Tabellen mit sechs Zeilen

\newcommand{\tabelleleitsechsxdrei}{

\begin{center}

\begin{tabular}{|c||c|c|c|}

\hline  \leitzeilenull  &  \leitzeileeins  &  \leitzeilezwei  &  \leitzeiledrei    \\

\hline \hline  \leitspalteeins  &  $ \aeinsxeins $  &  $ \aeinsxzwei $  &  $ \aeinsxdrei $    \\

\hline  \leitspaltezwei  &  $ \azweixeins $  &  $ \azweixzwei $  &  $ \azweixdrei $      \\

\hline  \leitspaltedrei  &  $ \adreixeins $  &  $ \adreixzwei $  &  $ \adreixdrei $     \\

\hline  \leitspaltevier  &  $ \avierxeins $  &  $ \avierxzwei $  &  $ \avierxdrei $      \\

\hline  \leitspaltefuenf  &  $ \afuenfxeins $  &  $ \afuenfxzwei $  &  $ \afuenfxdrei $     \\

\hline  \leitspaltesechs  &  $ \asechsxeins $  &  $ \asechsxzwei $  &  $ \asechsxdrei $      \\

\hline

\end{tabular}

\end{center} }

\newcommand{\tabelleleitsechsxfuenf}{

\begin{center}

\begin{tabular}{|c||c|c|c|c|c|}

\hline  \leitzeilenull  &  \leitzeileeins  &  \leitzeilezwei  &  \leitzeiledrei  &  \leitzeilevier  &  \leitzeilefuenf   \\

\hline \hline  \leitspalteeins  &  $ \aeinsxeins $  &  $ \aeinsxzwei $  &  $ \aeinsxdrei $  &  $ \aeinsxvier $  &  $ \aeinsxfuenf $   \\

\hline  \leitspaltezwei  &  $ \azweixeins $  &  $ \azweixzwei $  &  $ \azweixdrei $  &  $ \azweixvier $  &  $ \azweixfuenf $     \\

\hline  \leitspaltedrei  &  $ \adreixeins $  &  $ \adreixzwei $  &  $ \adreixdrei $  &  $ \adreixvier $  &  $ \adreixfuenf $      \\

\hline  \leitspaltevier  &  $ \avierxeins $  &  $ \avierxzwei $  &  $ \avierxdrei $  &  $ \avierxvier $  &  $ \avierxfuenf $     \\

\hline  \leitspaltefuenf  &  $ \afuenfxeins $  &  $ \afuenfxzwei $  &  $ \afuenfxdrei $  &  $ \afuenfxvier $  &  $ \afuenfxfuenf $    \\

\hline  \leitspaltesechs  &  $ \asechsxeins $  &  $ \asechsxzwei $  &  $ \asechsxdrei $  &  $ \asechsxvier $  &  $ \asechsxfuenf $    \\

\hline

\end{tabular}

\end{center} }

\newcommand{\tabelleleitsechsxsechs}{

\begin{center}

\begin{tabular}{|c||c|c|c|c|c|c|}

\hline  \leitzeilenull  &  \leitzeileeins  &  \leitzeilezwei  &  \leitzeiledrei  &  \leitzeilevier  &  \leitzeilefuenf  &  \leitzeilesechs   \\

\hline \hline  \leitspalteeins  &  $ \aeinsxeins $  &  $ \aeinsxzwei $  &  $ \aeinsxdrei $  &  $ \aeinsxvier $  &  $ \aeinsxfuenf $  &  $ \aeinsxsechs $ \\

\hline  \leitspaltezwei  &  $ \azweixeins $  &  $ \azweixzwei $  &  $ \azweixdrei $  &  $ \azweixvier $  &  $ \azweixfuenf $  &  $ \azweixsechs $   \\

\hline  \leitspaltedrei  &  $ \adreixeins $  &  $ \adreixzwei $  &  $ \adreixdrei $  &  $ \adreixvier $  &  $ \adreixfuenf $   &  $ \adreixsechs $   \\

\hline  \leitspaltevier  &  $ \avierxeins $  &  $ \avierxzwei $  &  $ \avierxdrei $  &  $ \avierxvier $  &  $ \avierxfuenf $   &  $ \avierxsechs $  \\

\hline  \leitspaltefuenf  &  $ \afuenfxeins $  &  $ \afuenfxzwei $  &  $ \afuenfxdrei $  &  $ \afuenfxvier $  &  $ \afuenfxfuenf $ &  $ \afuenfxsechs $   \\

\hline  \leitspaltesechs  &  $ \asechsxeins $  &  $ \asechsxzwei $  &  $ \asechsxdrei $  &  $ \asechsxvier $  &  $ \asechsxfuenf $  &  $ \asechsxsechs $  \\

\hline

\end{tabular}

\end{center} }

%Tabellen mit sieben Zeilen

\newcommand{\tabelleleitsiebenxsieben}{

\begin{center}

\begin{tabular}{|c||c|c|c|c|c|c|c|}

\hline  \leitzeilenull  &  \leitzeileeins  &  \leitzeilezwei  &  \leitzeiledrei  &  \leitzeilevier  &  \leitzeilefuenf  &  \leitzeilesechs &  \leitzeilesieben  \\

\hline \hline  \leitspalteeins  &  $ \aeinsxeins $  &  $ \aeinsxzwei $  &  $ \aeinsxdrei $  &  $ \aeinsxvier $  &  $ \aeinsxfuenf $  &  $ \aeinsxsechs $ &  $ \aeinsxsieben $ \\

\hline  \leitspaltezwei  &  $ \azweixeins $  &  $ \azweixzwei $  &  $ \azweixdrei $  &  $ \azweixvier $  &  $ \azweixfuenf $  &  $ \azweixsechs $  &  $ \azweixsieben $ \\

\hline  \leitspaltedrei  &  $ \adreixeins $  &  $ \adreixzwei $  &  $ \adreixdrei $  &  $ \adreixvier $  &  $ \adreixfuenf $   &  $ \adreixsechs $ &  $ \adreixsieben $  \\

\hline  \leitspaltevier  &  $ \avierxeins $  &  $ \avierxzwei $  &  $ \avierxdrei $  &  $ \avierxvier $  &  $ \avierxfuenf $   &  $ \avierxsechs $ &  $ \avierxsieben $ \\

\hline  \leitspaltefuenf  &  $ \afuenfxeins $  &  $ \afuenfxzwei $  &  $ \afuenfxdrei $  &  $ \afuenfxvier $  &  $ \afuenfxfuenf $ &  $ \afuenfxsechs $ &  $ \afuenfxsieben $  \\

\hline  \leitspaltesechs  &  $ \asechsxeins $  &  $ \asechsxzwei $  &  $ \asechsxdrei $  &  $ \asechsxvier $  &  $ \asechsxfuenf $  &  $ \asechsxsechs $ &  $ \asechsxsieben $ \\

\hline  \leitspaltesieben  &  $ \asiebenxeins $  &  $ \asiebenxzwei $  &  $ \asiebenxdrei $  &  $ \asiebenxvier $  &  $ \asiebenxfuenf $  &  $ \asiebenxsechs $ &  $ \asiebenxsieben $ \\

\hline

\end{tabular}

\end{center} }

%Tabellen mit acht Zeilen

\newcommand{\tabelleleitachtxdrei}{

\begin{center}

\begin{tabular}{|c||c|c|c|}

\hline  \leitzeilenull  &  \leitzeileeins  &  \leitzeilezwei  &  \leitzeiledrei     \\

\hline \hline  \leitspalteeins  &  $ \aeinsxeins $  &  $ \aeinsxzwei $  &  $ \aeinsxdrei $    \\

\hline  \leitspaltezwei  &  $ \azweixeins $  &  $ \azweixzwei $  &  $ \azweixdrei $     \\

\hline  \leitspaltedrei  &  $ \adreixeins $  &  $ \adreixzwei $  &  $ \adreixdrei $    \\

\hline  \leitspaltevier  &  $ \avierxeins $  &  $ \avierxzwei $  &  $ \avierxdrei $      \\

\hline  \leitspaltefuenf  &  $ \afuenfxeins $  &  $ \afuenfxzwei $  &  $ \afuenfxdrei $     \\

\hline  \leitspaltesechs  &  $ \asechsxeins $  &  $ \asechsxzwei $  &  $ \asechsxdrei $      \\

\hline  \leitspaltesieben  &  $ \asiebenxeins $  &  $ \asiebenxzwei $  &  $ \asiebenxdrei $     \\

\hline  \leitspalteacht  &  $ \aachtxeins $  &  $ \aachtxzwei $  &  $ \aachtxdrei $      \\

\hline

\end{tabular}

\end{center} }

%Tabellen mit zehn Zeilen

\newcommand{\tabelleleitzehnxdrei}{

\begin{center}

\begin{tabular}{|c||c|c|c|}

\hline  \leitzeilenull  &  \leitzeileeins  &  \leitzeilezwei  &  \leitzeiledrei     \\

\hline \hline  \leitspalteeins  &  $ \aeinsxeins $  &  $ \aeinsxzwei $  &  $ \aeinsxdrei $    \\

\hline  \leitspaltezwei  &  $ \azweixeins $  &  $ \azweixzwei $  &  $ \azweixdrei $     \\

\hline  \leitspaltedrei  &  $ \adreixeins $  &  $ \adreixzwei $  &  $ \adreixdrei $    \\

\hline  \leitspaltevier  &  $ \avierxeins $  &  $ \avierxzwei $  &  $ \avierxdrei $      \\

\hline  \leitspaltefuenf  &  $ \afuenfxeins $  &  $ \afuenfxzwei $  &  $ \afuenfxdrei $     \\

\hline  \leitspaltesechs  &  $ \asechsxeins $  &  $ \asechsxzwei $  &  $ \asechsxdrei $      \\

\hline  \leitspaltesieben  &  $ \asiebenxeins $  &  $ \asiebenxzwei $  &  $ \asiebenxdrei $     \\

\hline  \leitspalteacht  &  $ \aachtxeins $  &  $ \aachtxzwei $  &  $ \aachtxdrei $      \\

\hline  \leitspalteneun  &  $ \aneunxeins $  &  $ \aneunxzwei $  &  $ \aneunxdrei $      \\

\hline  \leitspaltezehn  &  $ \azehnxeins $  &  $ \azehnxzwei $  &  $ \azehnxdrei $      \\

\hline

\end{tabular}

\end{center} }

%Tabellen mit elf Zeilen

\newcommand{\tabelleleitelfxvier}{

\begin{center}

\begin{tabular}{|c||c|c|c|c|}

\hline  \leitzeilenull  &  \leitzeileeins  &  \leitzeilezwei  &  \leitzeiledrei  &  \leitzeilevier  \\

\hline \hline  \leitspalteeins  &  $ \aeinsxeins $  &  $ \aeinsxzwei $  &  $ \aeinsxdrei $  &  $ \aeinsxvier $  \\

\hline  \leitspaltezwei  &  $ \azweixeins $  &  $ \azweixzwei $  &  $ \azweixdrei $  &  $ \azweixvier $   \\

\hline  \leitspaltedrei  &  $ \adreixeins $  &  $ \adreixzwei $  &  $ \adreixdrei $  &  $ \adreixvier $   \\

\hline  \leitspaltevier  &  $ \avierxeins $  &  $ \avierxzwei $  &  $ \avierxdrei $  &  $ \avierxvier $   \\

\hline  \leitspaltefuenf  &  $ \afuenfxeins $  &  $ \afuenfxzwei $  &  $ \afuenfxdrei $  &  $ \afuenfxvier $   \\

\hline  \leitspaltesechs  &  $ \asechsxeins $  &  $ \asechsxzwei $  &  $ \asechsxdrei $  &  $ \asechsxvier $   \\

\hline  \leitspaltesieben  &  $ \asiebenxeins $  &  $ \asiebenxzwei $  &  $ \asiebenxdrei $  &  $ \asiebenxvier $    \\

\hline  \leitspalteacht  &  $ \aachtxeins $  &  $ \aachtxzwei $  &  $ \aachtxdrei $  &  $ \aachtxvier $   \\

\hline  \leitspalteneun  &  $ \aneunxeins $  &  $ \aneunxzwei $  &  $ \aneunxdrei $  &  $ \aneunxvier $   \\

\hline  \leitspaltezehn  &  $ \azehnxeins $  &  $ \azehnxzwei $  &  $ \azehnxdrei $  &  $ \azehnxvier $   \\

\hline  \leitspalteelf  &  $ \aelfxeins $  &  $ \aelfxzwei $  &  $ \aelfxdrei $  &  $ \aelfxvier $   \\

\hline

\end{tabular}

\end{center} }

%Tabellen mit zwölf Zeilen

\newcommand{\tabelleleitzwoelfxsieben}{

\begin{center}

\begin{tabular}{|c||c|c|c|c|c|c|}

\hline  \leitzeilenull  &  \leitzeileeins  &  \leitzeilezwei  &  \leitzeiledrei  &  \leitzeilevier &  \leitzeilefuenf  &  \leitzeilesechs &  \leitzeilesieben \\

\hline \hline  \leitspalteeins  &  $ \aeinsxeins $  &  $ \aeinsxzwei $  &  $ \aeinsxdrei $  &  $ \aeinsxvier $  &  $ \aeinsxfuenf $  &  $ \aeinsxsechs $  &  $ \aeinsxsieben $  \\

\hline  \leitspaltezwei  &  $ \azweixeins $  &  $ \azweixzwei $  &  $ \azweixdrei $  &  $ \azweixvier $  &  $ \azweixfuenf $  &  $ \azweixsechs $  &  $ \azweixsieben $   \\

\hline  \leitspaltedrei  &  $ \adreixeins $  &  $ \adreixzwei $  &  $ \adreixdrei $  &  $ \adreixvier $  &  $ \adreixfuenf $  &  $ \adreixsechs $  &  $ \adreixsieben $   \\

\hline  \leitspaltevier  &  $ \avierxeins $  &  $ \avierxzwei $  &  $ \avierxdrei $  &  $ \avierxvier $   &  $ \avierxfuenf $  &  $ \avierxsechs $  &  $ \avierxsieben $  \\

\hline  \leitspaltefuenf  &  $ \afuenfxeins $  &  $ \afuenfxzwei $  &  $ \afuenfxdrei $  &  $ \afuenfxvier $ &  $ \afuenfxfuenf $  &  $ \afuenfxsechs $  &  $ \afuenfxsieben $  \\

\hline  \leitspaltesechs  &  $ \asechsxeins $  &  $ \asechsxzwei $  &  $ \asechsxdrei $  &  $ \asechsxvier $  &  $ \asechsxfuenf $  &  $ \asechsxsechs $  &  $ \asechsxsieben $  \\

\hline  \leitspaltesieben  &  $ \asiebenxeins $  &  $ \asiebenxzwei $  &  $ \asiebenxdrei $  &  $ \asiebenxvier $  &  $ \asiebenxfuenf $  &  $ \asiebenxsechs $  &  $ \asiebenxsieben $   \\

\hline  \leitspalteacht  &  $ \aachtxeins $  &  $ \aachtxzwei $  &  $ \aachtxdrei $  &  $ \aachtxvier $  &  $ \aachtxfuenf $  &  $ \aachtxsechs $  &  $ \aachtxsieben $ \\

\hline  \leitspalteneun  &  $ \aneunxeins $  &  $ \aneunxzwei $  &  $ \aneunxdrei $  &  $ \aneunxvier $  &  $ \aneunxfuenf $  &  $ \aneunxsechs $  &  $ \aneunxsieben $  \\

\hline  \leitspaltezehn  &  $ \azehnxeins $  &  $ \azehnxzwei $  &  $ \azehnxdrei $  &  $ \azehnxvier $ &  $ \azehnxfuenf $  &  $ \azehnxsechs $  &  $ \azehnxsieben $   \\

\hline  \leitspalteelf  &  $ \aelfxeins $  &  $ \aelfxzwei $  &  $ \aelfxdrei $  &  $ \aelfxvier $  &  $ \aelfxfuenf $  &  $ \aelfxsechs $  &  $ \aelfxsieben $  \\

\hline  \leitspaltezwoelf  &  $ \azwoelfxeins $  &  $ \azwoelfxzwei $  &  $ \azwoelfxdrei $  &  $ \zwoelfxvier $  &  $ \azwoelfxfuenf $  &  $ \azwoelfxsechs $  &  $ \zwoelfxsieben $ \\

\hline

\end{tabular}

\end{center} }

%Tabellen mit sechzehn Zeilen

\newcommand{\tabelleleitsechzehnxzwei}{

\begin{center}

\begin{tabular}{|c||c|c|}

\hline  \leitzeilenull  &  \leitzeileeins  &  \leitzeilezwei   \\

\hline \hline  \leitspalteeins  &  $ \aeinsxeins $  &  $ \aeinsxzwei $ \\

\hline  \leitspaltezwei  &  $ \azweixeins $  &  $ \azweixzwei $  \\

\hline  \leitspaltedrei  &  $ \adreixeins $  &  $ \adreixzwei $  \\

\hline  \leitspaltevier  &  $ \avierxeins $  &  $ \avierxzwei $  \\

\hline  \leitspaltefuenf  &  $ \afuenfxeins $  &  $ \afuenfxzwei $  \\

\hline  \leitspaltesechs  &  $ \asechsxeins $  &  $ \asechsxzwei $  \\

\hline  \leitspaltesieben  &  $ \asiebenxeins $  &  $ \asiebenxzwei $  \\

\hline  \leitspalteacht  &  $ \aachtxeins $  &  $ \aachtxzwei $  \\

\hline  \leitspalteneun  &  $ \aneunxeins $  &  $ \aneunxzwei $   \\

\hline  \leitspaltezehn  &  $ \azehnxeins $  &  $ \azehnxzwei $  \\

\hline  \leitspalteelf  &  $ \aelfxeins $  &  $ \aelfxzwei $   \\

\hline  \leitspaltezwoelf  &  $ \azwoelfxeins $  &  $ \azwoelfxzwei $   \\

\hline  \leitspaltedreizehn  &  $ \adreizehnxeins $  &  $ \adreizehnxzwei $   \\

\hline  \leitspaltevierzehn  &  $ \avierzehnxeins $  &  $ \avierzehnxzwei $  \\

\hline  \leitspaltefuenfzehn  &  $ \afuenfzehnxeins $  &  $ \afuenfzehnxzwei $   \\

\hline  \leitspaltesechzehn  &  $ \asechzehnxeins $  &  $ \asechzehnxzwei $   \\

\hline

\end{tabular}

\end{center} }

\newcommand{\tabelleleitsechzehnxsechs}{

\begin{center}

\begin{tabular}{|c||c|c|c|c|c|c|}

\hline  \leitzeilenull  &  \leitzeileeins  &  \leitzeilezwei  &  \leitzeiledrei  &  \leitzeilevier  &  \leitzeilefuenf  &  \leitzeilesechs  \\

\hline \hline  \leitspalteeins  &  $ \aeinsxeins $  &  $ \aeinsxzwei $  &  $ \aeinsxdrei $  &  $ \aeinsxvier $  &  $ \aeinsxfuenf $  &  $ \aeinsxsechs $  \\

\hline  \leitspaltezwei  &  $ \azweixeins $  &  $ \azweixzwei $  &  $ \azweixdrei $  &  $ \azweixvier $  &  $ \azweixfuenf $  &  $ \azweixsechs $  \\

\hline  \leitspaltedrei  &  $ \adreixeins $  &  $ \adreixzwei $  &  $ \adreixdrei $  &  $ \adreixvier $  &  $ \adreixfuenf $  &  $ \adreixsechs $  \\

\hline  \leitspaltevier  &  $ \avierxeins $  &  $ \avierxzwei $  &  $ \avierxdrei $  &  $ \avierxvier $  &  $ \avierxfuenf $  &  $ \avierxsechs $  \\

\hline  \leitspaltefuenf  &  $ \afuenfxeins $  &  $ \afuenfxzwei $  &  $ \afuenfxdrei $  &  $ \afuenfxvier $  &  $ \afuenfxfuenf $  &  $ \afuenfxsechs $  \\

\hline  \leitspaltesechs  &  $ \asechsxeins $  &  $ \asechsxzwei $  &  $ \asechsxdrei $  &  $ \asechsxvier $  &  $ \asechsxfuenf $  &  $ \asechsxsechs $  \\

\hline  \leitspaltesieben  &  $ \asiebenxeins $  &  $ \asiebenxzwei $  &  $ \asiebenxdrei $  &  $ \asiebenxvier $  &  $ \asiebenxfuenf $  &  $ \asiebenxsechs $  \\

\hline  \leitspalteacht  &  $ \aachtxeins $  &  $ \aachtxzwei $  &  $ \aachtxdrei $  &  $ \aachtxvier $  &  $ \aachtxfuenf $  &  $ \aachtxsechs $  \\

\hline  \leitspalteneun  &  $ \aneunxeins $  &  $ \aneunxzwei $  &  $ \aneunxdrei $  &  $ \aneunxvier $  &  $ \aneunxfuenf $  &  $ \aneunxsechs $  \\

\hline  \leitspaltezehn  &  $ \azehnxeins $  &  $ \azehnxzwei $  &  $ \azehnxdrei $  &  $ \azehnxvier $  &  $ \azehnxfuenf $  &  $ \azehnxsechs $  \\

\hline  \leitspalteelf  &  $ \aelfxeins $  &  $ \aelfxzwei $  &  $ \aelfxdrei $  &  $ \aelfxvier $  &  $ \aelfxfuenf $  &  $ \aelfxsechs $  \\

\hline  \leitspaltezwoelf  &  $ \azwoelfxeins $  &  $ \azwoelfxzwei $  &  $ \azwoelfxdrei $  &  $ \azwoelfxvier $  &  $ \azwoelfxfuenf $  &  $ \azwoelfxsechs $  \\

\hline  \leitspaltedreizehn  &  $ \adreizehnxeins $  &  $ \adreizehnxzwei $  &  $ \adreizehnxdrei $  &  $ \adreizehnxvier $  &  $ \adreizehnxfuenf $  &  $ \adreizehnxsechs $  \\

\hline  \leitspaltevierzehn  &  $ \avierzehnxeins $  &  $ \avierzehnxzwei $  &  $ \avierzehnxdrei $  &  $ \avierzehnxvier $  &  $ \avierzehnxfuenf $  &  $ \avierzehnxsechs $  \\

\hline  \leitspaltefuenfzehn  &  $ \afuenfzehnxeins $  &  $ \afuenfzehnxzwei $  &  $ \afuenfzehnxdrei $  &  $ \afuenfzehnxvier $  &  $ \afuenfzehnxfuenf $  &  $ \afuenfzehnxsechs $  \\

\hline  \leitspaltesechzehn  &  $ \asechzehnxeins $  &  $ \asechzehnxzwei $  &  $ \asechzehnxdrei $  &  $ \asechzehnxvier $  &  $ \asechzehnxfuenf $  &  $ \asechzehnxsechs $  \\

\hline

\end{tabular}

\end{center} }

%Tabellen mit achtzehn Zeilen

\newcommand{\tabelleleitachtzehnxneun}{

\begin{center}

\begin{tabular}{|c||c|c|c|c|c|c|c|c|c|}

\hline  \leitzeilenull  &  \leitzeileeins  &  \leitzeilezwei  &  \leitzeiledrei  &  \leitzeilevier  &  \leitzeilefuenf  &  \leitzeilesechs &  \leitzeilesieben  &  \leitzeileacht  &  \leitzeileneun \\

\hline \hline  \leitspalteeins  &  $ \aeinsxeins $  &  $ \aeinsxzwei $  &  $ \aeinsxdrei $  &  $ \aeinsxvier $  &  $ \aeinsxfuenf $  &  $ \aeinsxsechs $ &  $ \aeinsxsieben $  &  $ \aeinsxacht $  &  $ \aeinsxneun $  \\

\hline  \leitspaltezwei  &  $ \azweixeins $  &  $ \azweixzwei $  &  $ \azweixdrei $  &  $ \azweixvier $  &  $ \azweixfuenf $  &  $ \azweixsechs $ &  $ \azweixsieben $  &  $ \azweixacht $  &  $ \azweixneun $  \\

\hline  \leitspaltedrei  &  $ \adreixeins $  &  $ \adreixzwei $  &  $ \adreixdrei $  &  $ \adreixvier $  &  $ \adreixfuenf $  &  $ \adreixsechs $ &  $ \adreixsieben $  &  $ \adreixacht $  &  $ \adreixneun $  \\

\hline  \leitspaltevier  &  $ \avierxeins $  &  $ \avierxzwei $  &  $ \avierxdrei $  &  $ \avierxvier $  &  $ \avierxfuenf $  &  $ \avierxsechs $  &  $ \avierxsieben $  &  $ \avierxacht $  &  $ \avierxneun $ \\

\hline  \leitspaltefuenf  &  $ \afuenfxeins $  &  $ \afuenfxzwei $  &  $ \afuenfxdrei $  &  $ \afuenfxvier $  &  $ \afuenfxfuenf $  &  $ \afuenfxsechs $ &  $ \afuenfxsieben $  &  $ \afuenfxacht $  &  $ \afuenfxneun $  \\

\hline  \leitspaltesechs  &  $ \asechsxeins $  &  $ \asechsxzwei $  &  $ \asechsxdrei $  &  $ \asechsxvier $  &  $ \asechsxfuenf $  &  $ \asechsxsechs $ &  $ \asechsxsieben $  &  $ \asechsxacht $  &  $ \asechsxneun $  \\

\hline  \leitspaltesieben  &  $ \asiebenxeins $  &  $ \asiebenxzwei $  &  $ \asiebenxdrei $  &  $ \asiebenxvier $  &  $ \asiebenxfuenf $  &  $ \asiebenxsechs $ &  $ \asiebenxsieben $  &  $ \asiebenxacht $  &  $ \asiebenxneun $  \\

\hline  \leitspalteacht  &  $ \aachtxeins $  &  $ \aachtxzwei $  &  $ \aachtxdrei $  &  $ \aachtxvier $  &  $ \aachtxfuenf $  &  $ \aachtxsechs $ &  $ \aachtxsieben $  &  $ \aachtxacht $  &  $ \aachtxneun $  \\

\hline  \leitspalteneun  &  $ \aneunxeins $  &  $ \aneunxzwei $  &  $ \aneunxdrei $  &  $ \aneunxvier $  &  $ \aneunxfuenf $  &  $ \aneunxsechs $ &  $ \aneunxsieben $  &  $ \aneunxacht $  &  $ \aneunxneun $  \\

\hline  \leitspaltezehn  &  $ \azehnxeins $  &  $ \azehnxzwei $  &  $ \azehnxdrei $  &  $ \azehnxvier $  &  $ \azehnxfuenf $  &  $ \azehnxsechs $ &  $ \azehnxsieben $  &  $ \azehnxacht $  &  $ \azehnxneun $  \\

\hline  \leitspalteelf  &  $ \aelfxeins $  &  $ \aelfxzwei $  &  $ \aelfxdrei $  &  $ \aelfxvier $  &  $ \aelfxfuenf $  &  $ \aelfxsechs $ &  $ \aelfxsieben $  &  $ \aelfxacht $  &  $ \aelfxneun $  \\

\hline  \leitspaltezwoelf  &  $ \azwoelfxeins $  &  $ \azwoelfxzwei $  &  $ \azwoelfxdrei $  &  $ \azwoelfxvier $  &  $ \azwoelfxfuenf $  &  $ \azwoelfxsechs $ &  $ \azwoelfxsieben $  &  $ \azwoelfxacht $  &  $ \azwoelfxneun $  \\

\hline  \leitspaltedreizehn  &  $ \adreizehnxeins $  &  $ \adreizehnxzwei $  &  $ \adreizehnxdrei $  &  $ \adreizehnxvier $  &  $ \adreizehnxfuenf $  &  $ \adreizehnxsechs $ &  $ \adreizehnxsieben $  &  $ \adreizehnxacht $  &  $ \adreizehnxneun $  \\

\hline  \leitspaltevierzehn  &  $ \avierzehnxeins $  &  $ \avierzehnxzwei $  &  $ \avierzehnxdrei $  &  $ \avierzehnxvier $  &  $ \avierzehnxfuenf $  &  $ \avierzehnxsechs $  &  $ \avierzehnxsieben $  &  $ \avierzehnxacht $  &  $ \avierzehnxneun $ \\

\hline  \leitspaltefuenfzehn  &  $ \afuenfzehnxeins $  &  $ \afuenfzehnxzwei $  &  $ \afuenfzehnxdrei $  &  $ \afuenfzehnxvier $  &  $ \afuenfzehnxfuenf $  &  $ \afuenfzehnxsechs $ &  $ \afuenfzehnxsieben $  &  $ \afuenfzehnxacht $  &  $ \afuenfzehnxneun $  \\

\hline  \leitspaltesechzehn  &  $ \asechzehnxeins $  &  $ \asechzehnxzwei $  &  $ \asechzehnxdrei $  &  $ \asechzehnxvier $  &  $ \asechzehnxfuenf $  &  $ \asechzehnxsechs $ &  $ \asechzehnxsieben $  &  $ \asechzehnxacht $  &  $ \asechzehnxneun $  \\

\hline  \leitspaltesiebzehn  &  $ \asiebzehnxeins $  &  $ \asiebzehnxzwei $  &  $ \asiebzehnxdrei $  &  $ \asiebzehnxvier $  &  $ \asiebzehnxfuenf $  &  $ \asiebzehnxsechs $ &  $ \asiebzehnxsieben $  &  $ \asiebzehnxacht $  &  $ \asiebzehnxneun $  \\

\hline  \leitspalteachtzehn  &  $ \aachtzehnxeins $  &  $ \aachtzehnxzwei $  &  $ \aachtzehnxdrei $  &  $ \aachtzehnxvier $  &  $ \aachtzehnxfuenf $  &  $ \aachtzehnxsechs $ &  $ \aachtzehnxsieben $  &  $ \aachtzehnxacht $  &  $ \aachtzehnxneun $  \\

\hline

\end{tabular}

\end{center} }

\newcommand{\bruchhilfealign}{\cr & & }

\newcommand{\bruchhilfedisp}{\]  \[ }

\newcommand{\bruchhilfe}{ \- }

%[[Projekt:Semantische Vorlagen/Notlösungen/Latex]]

%
%Bei verschachtelten ifthen-Befehlen
%kann es Probleme geben, daher
%

\newcommand{\faktuebergangpos}[1]{ \faktuebergangskip  {\hspace{-0,55cm} } {#1} }

%Ende des Grundvorspannes

%Globale Strukturierungen

\newcommand{\seitenueberschrift}[1]{\seitenueberschriftvorskip
\begin{center} \large{ #1} \end{center} \seitenueberschriftnachskip}

\newcommand{\zwischenueberschrift}[1]{\section*{#1} }

\newcommand{\zwischenueberschriftaufgaben}[1]{ \zwischenueberschriftvorskip \begin{center} {\bf #1} \end{center} \zwischenueberschriftnachskip}

\newcommand{\unterueberschrift}[1]{ \unterueberschriftvorskip \hspace{0.5cm} {\em #1}  \unterueberschriftnachskip}

\renewcommand{\definitionbenennung}[1]{}

%Grund:Bei Definitionen im Normaltext
%nicht nötig, da Begriff früh kommt

\renewcommand{\beispielbenennung}[1]{}

\renewcommand{\bemerkungbenennung}[1]{}

%Umbrüche bei großer Schrift (entfällt)

%[[Projekt:Semantische Vorlagen/Mathematischer Modus/Umbrüche/Latex]]

\newcommand{\mathdisplaybruch}{}

\newcommand{\mathdisplaybruchinklammer}{}

\newcommand{\mathl}[2]{$#1$#2}

\newcommand{\mathlk}[2]{$#1$#2}

%Nicht durch semantische Vorlage produzierte Befehle

\newcommand{\mathdisplaybruchbruch}{ \]  \[ }

%Gestaltung der vertikalen und horizontalen Abstände

%[[Projekt:Semantische Vorlagen/Latex/Vorlesungsskript/Skips]]

%Abstände bei Überschriften

\newcommand{\seitenueberschriftvorskip}{\par \bigskip \bigskip \bigskip }

\newcommand{\seitenueberschriftnachskip}{\par \bigskip}

\newcommand{\zwischenueberschriftvorskip}{\par \bigskip}

\newcommand{\zwischenueberschriftnachskip}{\par \smallskip}

\newcommand{\unterueberschriftvorskip}{\par \bigskip}

\newcommand{\unterueberschriftnachskip}{\par \smallskip}

%Aufzählungen und Listen

\newcommand{\listskip}{\smallskip}

%Abstände bei mathematischen Strukturen

\newcommand{\aufgabevorskip}{\bigskip}

\newcommand{\aufgabepunktskip}{\par}

\newcommand{\aufgabenachskip}{}

\newcommand{\aufgabeloesungskip}{\par \bigskip}

\newcommand{\beispielvorskip}{}

\newcommand{\beispielnachskip}{}

\newcommand{\beispielbenennungnachskip}{}

\newcommand{\bemerkungvorskip}{}

\newcommand{\bemerkungnachskip}{}

\newcommand{\bemerkungbenennungnachskip}{}

\newcommand{\definitionvorskip}{}

\newcommand{\definitionnachskip}{}

\newcommand{\definitionbenennungnachskip}{}

\newcommand{\faktvorskip}{}

\newcommand{\faktnachskip}{}

\newcommand{\faktbenennungvorskip}{}

\newcommand{\faktbenennungnachskip}{}

\newcommand{\faktsituationskip}{}

\newcommand{\faktvoraussetzungskip}{}

\newcommand{\faktuebergangskip}{}

\newcommand{\faktfolgerungskip}{}

\newcommand{\faktzusatzskip}{}

%Sonstiges

\newcommand{\deckblattabstand}{\vspace{2cm}}

\newcommand{\gesamtskriptnewpage}{ }

\newcommand{\einzeltextnewpage}{ }

%[[Projekt:Semantische Vorlagen/Latex/Vorlesungsskript/Horizontale Abstände]]

\newcommand{\leerzeichen}{ }

%Seitengestaltung

%[[Projekt:Semantische Vorlagen/Latex/Vorlesungsskript/Seitengestaltung]]

%Einzelne Abstände

% page geometry supplied by reader wrapper
% page geometry supplied by reader wrapper

% page geometry supplied by reader wrapper
% page geometry supplied by reader wrapper

% page geometry supplied by reader wrapper
% page geometry supplied by reader wrapper

\setlength{\parindent}{0cm}

\setlength{\parskip}{1ex}

%Bildgestaltung

%[[Projekt:Semantische Vorlagen/Latex/Vorlesungsskript/Bildgestaltung]]

\newcommand{\bild}[1]{\vspace{4mm} #1}

\newcommand{\bildtext}[1]{\begin{center} {\small #1 } \end{center} }

%Lokale Einstellungen

%\input{Lokalername}

\newcommand{\bildeinlesung}[1]{../authority/media/#1}

\renewcommand{\bildeinlesung}[2]{../authority/media/#1.#2}

\newcommand{\bildeinlesungpng}[2]{../authority/media/#1.png}

\newcommand{\bildeinlesungPNG}[2]{../authority/media/#1.png}

\newcommand{\bildeinlesungjpg}[2]{../authority/media/#1.jpg}

\newcommand{\bildeinlesungJPG}[2]{../authority/media/#1.jpg}

\newcommand{\bildeinlesungjpeg}[2]{../authority/media/#1.jpg}

\newcommand{\bildeinlesungsvg}[2]{generated/media/#1.png}

\newcommand{\bildeinlesunggif}[2]{generated/media/#1.png}

\newcommand{\bildeinlesungGIF}[2]{generated/media/#1.png}

\newcommand{\bildeinlesungxcf}[2]{../authority/media/#1.png}

%Klausurgestaltung

%[[Projekt:Semantische Vorlagen/Klausurgestaltung/Latex]]

\newcommand{\fachbereich}{Fachbereich}

\newcommand{\dozent}{Dozent}

\newcommand{\klausurdatum}{Datum}

\newcommand{\klausurgebiet}{Mathematik}

\newcommand{\klausurtyp}{Klausur}

\newcommand{\klausurmodus}{Dauer: Zehn Minuten Einlesezeit und zwei volle Stunden Bearbeitungszeit. Es sind keine Hilfsmittel erlaubt. Alle Antworten sind zu begr\"unden.

Es gibt insgesamt 64 Punkte. Es gilt die Sockelregelung, d.h. die Bewertung pro Aufgabe(nteil) beginnt in der Regel bei der halben Punktzahl. Zum Bestehen braucht man $16$ Punkte, ab $32$ Punkten gibt es eine Eins.

Tragen Sie auf dem Deckblatt und jedem weiteren Blatt Ihren Namen und Ihre
Matrikelnummer leserlich ein. Viel Erfolg!}

\newcommand{\klausurname}{
\renewcommand{\arraystretch}{1.5}
\begin{tabular}{@{}p{3.2cm}p{10cm}}
Name, Vorname: &
..................................................................................\\
Matrikelnummer: &
..................................................................................\\
\end{tabular}
}

\newcommand{\klausurvorspann}[5]{\textbf{\fachbereich}\hfill \textbf{\klausurdatum}
\\
\textbf{\dozent} \hfill
\vspace{4ex}

\begin{center}
{\bfseries{\large \klausurgebiet \vspace{2ex}

\klausurtyp} }
\end{center}

\vspace{2ex}

\klausurmodus

\vspace{2ex}

\klausurname
}

\newcommand{\klausurnote}{\begin{tabular}{@{}p{3.2cm}p{10cm}} Note: & \end{tabular} }

\newcommand{\klausurtaeuschungwarnung}{ \vspace{2ex}{Mir ist bewu\ss t, dass bei einem schweren T\"auschungsversuch s\"amtliche Pr\"u\-fungsvorleis\-tungen in diesem Kurs verfallen. Ein schwerer T\"auschungsversuch liegt insbesondere beim Einsatz von Literatur und von elektronischen Hilfsmitteln vor. \vspace{1ex} } }

\newcommand{\klausureinverstaendnis}{
\vspace{2ex}
Ich erkl\"are mich durch meine Unterschrift einverstanden, dass mein Klausurergebnis mit meiner Matrikelnummer im Internet bekanntgegeben wird.
\vspace{1ex}

\begin{tabular}{@{}p{3.2cm}p{10cm}}
Unterschrift: &
.....................................................................................
\end{tabular}
}

%Variante für Testklausur

\newcommand{\testklausurvorspann}[5]{\textbf{\fachbereich}\hfill \textbf{\klausurdatum} \\ \textbf{\dozent} \hfill \vspace{4ex}

\begin{center} {\bfseries{\large \klausurgebiet \vspace{2ex}

\klausurtyp} } \end{center}

\vspace{2ex}

\testklausurmodus

\vspace{3ex}

\klausurname }

\newcommand{\testklausurmodus}{Dauer: Zehn Minuten Einlesezeit und zwei volle Stunden Bearbeitungszeit.

Es sind keine Hilfsmittel erlaubt.

Alle Antworten sind zu begr\"unden.

Es gibt insgesamt 64 Punkte. Es gilt die Sockelregelung, d.h. die Bewertung pro Aufgabe(nteil) beginnt in der Regel bei der halben Punktzahl.

Zur Orientierung: Zum Bestehen braucht man $16$ Punkte, ab $32$ Punkten gibt es eine Eins.

Tragen Sie auf dem Deckblatt und jedem weiteren Blatt Ihren Namen und Ihre Matrikelnummer leserlich ein.

Viel Erfolg!}

%\renewcommand{\inputaufgabegibtloesung}[4]{\aufgabevorskip \begin{aufgabe} \punkte {#1} \aufgabepunktskip #2 #3 #4\end{aufgabe} \aufgabenachskip}

\newcommand{\klausurmitloesungenvorspann}[5]{\textbf{\fachbereich}\hfill \textbf{\klausurdatum} \\ \textbf{\dozent} \hfill \vspace{4ex}

\begin{center} {\bfseries{\large \klausurgebiet \vspace{2ex}

\klausurtyp \hspace{0mm} mit L\"osungen} } \end{center}

}

%Lizenzierung und Bildverzeichnis

%[[Projekt:Semantische Vorlagen/Latex/Vorlesungsskript/Lizenzierung]]

\newcommand{\Lizenztext}{Teks sumber dibuat oleh Holger Brenner di Wikiversity berbahasa Jerman dan digunakan di sini berdasarkan CC BY-SA 4.0. Media mengikuti lisensi per berkas.}

\newcommand{\Abbildungslizenzierung}{ Erl\"auterung: Die in diesem Text verwendeten Bilder stammen aus Commons (also von http://commons.wikimedia.org) und haben eine Lizenz, die die Verwendung hier erlaubt. Die Bilder werden mit ihren Dateinamen auf Commons angef\"uhrt zusammen mit ihrem Autor bzw. Hochlader und der Lizenz. }

%Konfiguration des Bildverzeichnisses

\newcommand{\bildlizenz}[6]{
\ifthenelse {\equal{#2}{}}
{\addcontentsline{lof}{figure}{Sumber = #1, Kreator = Pengguna #3 di #4, Lisensi = #5 \bildlizenzskip  }}
{ \ifthenelse {\equal{#3}{}}
{ \addcontentsline{lof}{figure}{ Sumber = #1, Kreator = #2, Lisensi = #5 \bildlizenzskip }}
{ \addcontentsline{lof}{figure}{ Sumber = #1, Kreator = #2 (diunggah oleh Pengguna #3 di #4), Lisensi = #5 \bildlizenzskip }} } }

\newcommand{\bildlizenzskip}{ }

%Englische Variante

\newcommand{\ignoriereeng}[1]{}

\ignoriereeng{

%[[Projekt:Semantische Vorlagen/Latex/Englische Zusätze]]

\newtheorem{corollary}[fakt]{Corollary}

\newtheorem{Corollary}[fakt]{Corollary}

\theoremstyle{definition}

\newtheorem{example}[fakt]{Example}

\newtheorem{remark}[fakt]{Remark}

\newcommand{\inputremark}[2] {\bemerkungvorskip \begin{remark} \bemerkungbenennung{#1} \bemerkungbenennungnachskip #2 \end{remark} \bemerkungnachskip}

\newcommand{\inputexample}[2] {\beispielvorskip \begin{example} \beispielbenennung{#1} \beispielbenennungnachskip #2 \end{example} \beispielnachskip}

\newcommand{\inputfaktproof}[5] {\faktvorskip \begin{#2}

%\label{#1} (Absetzen, sonst kann das folgende dahinter rutschen)

\faktbenennung{#3} \faktbenennungnachskip #4 \end{#2} \faktnachskip \beweis{#5}}

\renewcommand{\proofname}{\hspace{-0.64cm}{\it Proof}}

}

%Variante für Artikel

\newcommand{\ignoriereart}[1]{}

\ignoriereart{
\newcommand{\titelueberschrift}[1]{\title[#1] { #1} \maketitle }

\renewcommand{\seitenueberschrift}[1]{ \section {#1} }

\renewcommand{\zwischenueberschrift}[1]{ \subsection {#1} }
}

% Indonesian aliases used by the independent translated reader.
\newtheorem{Teorema}[fakt]{Teorema}
\renewcommand{\proofname}{Bukti}
\newtheorem{Proposisi}[fakt]{Proposisi}
\newtheorem{Lema}[fakt]{Lema}
% Keep each source image and caption together when it fits as a single block.
\renewcommand{\bild}[1]{%
\par\vspace{4mm}%
\noindent\begin{minipage}{\linewidth}#1\end{minipage}%
\par}
% Keep source filenames with underscores safe in the persisted list of figures.
\renewcommand{\bildlizenz}[6]{%
\ifthenelse{\equal{#2}{}}%
{\addcontentsline{lof}{figure}{Sumber = \protect\nolinkurl{#1}, Kreator = Pengguna #3 di #4, Lisensi = #5 \bildlizenzskip}}%
{\ifthenelse{\equal{#3}{}}%
{\addcontentsline{lof}{figure}{Sumber = \protect\nolinkurl{#1}, Kreator = #2, Lisensi = #5 \bildlizenzskip}}%
{\addcontentsline{lof}{figure}{Sumber = \protect\nolinkurl{#1}, Kreator = #2 (diunggah oleh Pengguna #3 di #4), Lisensi = #5 \bildlizenzskip}}}}
% The frozen MediaWiki TeX uses underscore-normalized PNG names while some
% canonical Commons filenames retain spaces. The bounded cumulative build
% creates exact generated aliases and routes PNG macros only through that
% generated directory; canonical source files remain untouched.
\renewcommand{\bildeinlesungpng}[2]{generated/media/#1.png}
\renewcommand{\bildeinlesungPNG}[2]{generated/media/#1.png}
% The frozen source numbers facts and exercises by lecture/worksheet section.
% Book chapters are reader scaffolding and must not enter source identifiers.
\renewcommand{\thefakt}{\arabic{section}.\arabic{fakt}}

\addto\captionsindonesian{%
  \renewcommand{\contentsname}{Daftar Isi}%
  \renewcommand{\figurename}{Gambar}%
  \renewcommand{\listfigurename}{Daftar Gambar}%
}
\hypersetup{pdflang={id-ID},bookmarksopen=true}
\emergencystretch=1em
\ifdefined\pdfinfoomitdate\pdfinfoomitdate=1\fi
\ifdefined\pdftrailerid\pdftrailerid{}\fi
\ifdefined\pdfsuppressptexinfo\pdfsuppressptexinfo=15\fi

\title{Geometri Diferensial dan Manifold Mulus\\\large Edisi Lengkap Bahasa Indonesia}
\author{Holger Brenner\\Terjemahan Bahasa Indonesia independen}
\date{Sumber: \textit{Differentialgeometrie (Osnabrück 2023)}}

\begin{document}
\hypersetup{pageanchor=false}
\maketitle
\hypersetup{pageanchor=true}
\frontmatter
\chapter*{Tentang edisi ini}
Pembaca lengkap ini menerjemahkan seluruh 29 Kuliah dan 29 Lembar Kerja dari kursus Holger Brenner di Wikiversity berbahasa Jerman. Selain inti terjemahan, edisi ini memuat sepuluh formulir ujian resmi beserta solusi sumber yang tersedia, serta dua jembatan dan enam perbaikan solusi yang ditulis khusus untuk edisi ini dan diberi label terpisah. Teks sumber digunakan berdasarkan CC BY-SA 4.0. Terjemahan ini merupakan karya independen dan bukan edisi resmi atau dukungan dari penulis maupun Wikiversity. Setiap gambar mengikuti status hak atau lisensi berkasnya sendiri, yang dicatat pada daftar gambar dan manifest media.

Sumber permanen dan hash revisi dicatat dalam berkas kontrol edisi ini. Rumus, urutan, latihan, penanda solusi, dan atribusi media dipertahankan; ID mesin hanya merupakan lapisan tambahan. Solusi yang disediakan sumber dipisahkan secara eksplisit dari 38 butir solusi asli edisi ini.

Proses terjemahan dan produksi edisi dibantu oleh OpenAI Codex gpt-5.6-sol, Ultra, di bawah arahan pengguna. Kredit penulis, sumber, dan kontributor manusia tetap dipertahankan sebagaimana dicatat dalam edisi dan manifestnya.
\tableofcontents

\mainmatter
\part{Unit 1}
\chapter[Kuliah 1]{Kuliah 1: Analisis dan Geometri}
\setcounter{section}{1}

  \zwischenueberschrift{Hipermuka}

Misalkan
\maabbdisp {h} {\R^n } { \R
} {}
adalah sebuah
\definitionsverweis {fungsi terdiferensialkan}{}{}
dan

\mavergleichskette
{\vergleichskette
{ Y
}
{ = }{ h^{-1}(c)
}
{  }{
}
{    }{
}
{   }{
}
}
{}{}{}
adalah serat di atas

\mavergleichskette
{\vergleichskette
{ c
}
{ \in }{ \R
}
{  }{
}
{    }{
}
{   }{
}
}
{}{}{.}
Himpunan-himpunan bagian ini dipelajari khususnya dalam konteks teorema fungsi implisit. Jika
\definitionsverweis {diferensial total}{}{}
\maabbdisp {\left(D h \right)_{ P}} {\R^n } { \R
} {}
surjektif untuk setiap titik

\mavergleichskette
{\vergleichskette
{ P
}
{ \in }{ Y
}
{  }{
}
{    }{
}
{   }{
}
}
{}{}{}
---dengan kata lain, jika $P$ selalu merupakan
\definitionsverweis {titik reguler}{}{}
dari pemetaan itu---maka $Y$, dalam suatu lingkungan terbuka dari $P$, homeomorfik dengan suatu himpunan terbuka di $\R^{n-1}$. Selain itu, dalam konteks ini kita telah memperkenalkan
\definitionsverweis {ruang singgung pada serat}{}{}
sebagai kernel diferensial total. Ruang ini merupakan subruang vektor berdimensi $(n-1 )$ dari $\R^n$, tetapi biasanya dibayangkan melekat pada titik $P$, yakni sebagai ruang afin-linear. Di titik $P$, ruang singgung ini menyinggung $Y$. Kita menyebut $Y$ sebuah \stichwort {hipermuka} {,} karena ruang singgungnya berdimensi satu lebih rendah daripada ruang sekitarnya. Untuk

\mavergleichskette
{\vergleichskette
{ n
}
{ = }{ 3
}
{  }{
}
{    }{
}
{   }{
}
}
{}{}{}
memang diperoleh sebuah permukaan. Jika kita menggunakan hasil kali dalam standar pada $\R^n$, ruang singgung pada serat juga dapat dipandang sebagai ruang semua vektor yang tegak lurus terhadap
\definitionsverweis {gradien}{}{}

\mathl{\operatorname{Grad} \,  h  (  P )}{}.

\inputbeispiel{}
{

Misalkan

\mavergleichskettedisp
{\vergleichskette
{ h(x,y,z)
}
{ =} { x^2+y^2+z^2
}
{ } {
}
{ } {
}
{ } {
}
}
{}{}{}
dan kita tinjau serat $Y$ dari $h$ di atas $1$, yaitu permukaan bola berjari-jari $1$ dengan pusat di titik asal. Diferensial totalnya adalah
\mathl{\left(  2x  , \,   2y  , \,   2z                 \right)}{,} sehingga di setiap titik permukaan bola setidaknya satu komponennya tidak sama dengan $0$; oleh karena itu, $h$
\definitionsverweis {reguler}{}{}
di setiap titik pada $Y$. Misalkan

\mavergleichskette
{\vergleichskette
{ P
}
{ = }{ \left(  a   , \,   b  , \,  c                 \right)
}
{ \in }{ Y
}
{    }{
}
{   }{
}
}
{}{}{.}
Ruang singgung di $P$ diberikan oleh syarat linear

\mavergleichskettedisp
{\vergleichskette
{ ax+by+cz
}
{ =} { 0
}
{ } {
}
{ } {
}
{ } {
}
}
{}{}{}
dan merupakan himpunan semua vektor yang ortogonal terhadap gradien
\mathl{\begin{pmatrix}  2a \\ 2b\\  2c   \end{pmatrix}}{}.

}

\inputbeispiel{}
{

$\,$

\bild{ \begin{center}
\includegraphics[width=5.5cm]{ \bildeinlesungsvg {3d-function-6} {svg} }
\end{center}
\bildtext {} }

\bildlizenz { 3d-function-6.svg } {} {MartinThoma} {Commons} {CC BY 3.0} {}

Misalkan

\mavergleichskette
{\vergleichskette
{ V
}
{ \subseteq }{ \R^{n-1}
}
{  }{
}
{    }{
}
{   }{
}
}
{}{}{}
adalah sebuah
\definitionsverweis {himpunan terbuka}{}{}
dan misalkan
\maabbdisp {f} { V } { \R
} {}
adalah sebuah
\definitionsverweis {fungsi yang terdiferensialkan secara kontinu}{}{}.
\definitionsverweis {Grafik}{}{}
dari $f$ adalah

\mavergleichskettedisp
{\vergleichskette
{ Y
}
{ =} { \left\{ (x_1  , \ldots , x_{n-1},f( x_1  , \ldots , x_{n-1}))  \mid   x_1  , \ldots , x_{n-1} \in V         \right\}
}
{ \subseteq} { V \times \R
}
{ \subseteq} { \R^n
}
{ } {
}
}
{}{}{,}
yang juga dapat dipandang sebagai himpunan nol
\zusatzklammer {serat di atas $0$} {} {}
dari

\mavergleichskettedisp
{\vergleichskette
{ h(x_1  , \ldots , x_n)
}
{ \defeq} { x_n- f( x_1 , \ldots , x_{n-1})
}
{ } {
}
{ } {
}
{ } {
}
}
{}{}{.}
Turunan parsial dari $h$ adalah

\mathdisp {- { \frac{   \partial   f   }{  \partial   x_1   } } , \ldots , - { \frac{   \partial   f   }{  \partial   x_{n-1}   } } ,1} { . }

Khususnya, $h$
\definitionsverweis {reguler}{}{}
di setiap titik pada $Y$.
\definitionsverweis {Ruang singgung}{}{}
pada $Y$ di sebuah titik

\mavergleichskette
{\vergleichskette
{ P
}
{ \in }{ Y
}
{  }{
}
{    }{
}
{   }{
}
}
{}{}{}
tegak lurus terhadap
\mathl{\begin{pmatrix}  - { \frac{   \partial   f   }{  \partial   x_1   } } ( P)  \\ \vdots \\  -  { \frac{   \partial   f   }{  \partial   x_{n-1}   } } ( P)\\ 1   \end{pmatrix}}{.} Setiap garis dasar diparameterkan
\mathl{\begin{pmatrix}  x_1  \\ \vdots \\ x_{n-1}   \end{pmatrix} + t  \begin{pmatrix} v_1  \\\vdots\\ v_{n-1}   \end{pmatrix}}{} menjadi kurva yang diparameterkan
\maabbeledisp {} { I } { Y
} { t } { \begin{pmatrix}  x_1+tv_1  \\ \vdots \\  x_{n-1} +tv_{n-1}\\f \begin{pmatrix}  x_1+tv_1  \\ \vdots \\  x_{n-1} +tv_{n-1}   \end{pmatrix}   \end{pmatrix}
} {}
pada $Y$, yang turunannya menghasilkan sebuah vektor singgung. Jika lintasan untuk

\mavergleichskettedisp
{\vergleichskette
{ v
}
{ =} { e_i
}
{ } {
}
{ } {
}
{ } {
}
}
{}{}{}
dilambangkan dengan $\gamma_i$, maka

\mavergleichskettedisp
{\vergleichskette
{ \gamma'_i(t)
}
{ =} { \begin{pmatrix} \vdots \\ 1 \\ \vdots\\ { \frac{   \partial   f   }{  \partial   x_i   } }   \end{pmatrix}
}
{ } {
}
{ } {
}
{ } {
}
}
{}{}{}
dan

\mavergleichskettedisp
{\vergleichskette
{ \gamma^{\prime \prime}_i(t)
}
{ =} { \begin{pmatrix} \vdots \\ { \frac{   \partial    }{  \partial   x_i   } } { \frac{   \partial   f   }{  \partial   x_i   } }    \end{pmatrix}
}
{ } {
}
{ } {
}
{ } {
}
}
{}{}{}
menurut
Teorema 49.1 (Analisis (Osnabrück 2021--2023)).

}

  \zwischenueberschrift{Analisis dan Geometri}

Misalkan

\mavergleichskette
{\vergleichskette
{ W
}
{ \subseteq }{ \R^n
}
{  }{
}
{    }{
}
{   }{
}
}
{}{}{}
\definitionsverweis {terbuka}{}{,}
\maabb {h} { W } { \R
} {}
adalah sebuah
\definitionsverweis {fungsi yang terdiferensialkan secara kontinu}{}{}
dan

\mavergleichskette
{\vergleichskette
{ Y
}
{ = }{ h^{-1}(c)
}
{  }{
}
{    }{
}
{   }{
}
}
{}{}{}
adalah
\definitionsverweis {serat}{}{}
di atas

\mavergleichskette
{\vergleichskette
{ c
}
{ \in }{ \R
}
{  }{
}
{    }{
}
{   }{
}
}
{}{}{,}
dengan $h$
\definitionsverweis {reguler}{}{}
di setiap titik pada $Y$. Kita membahas beberapa objek atau konstruksi yang memang ada di ruang linear sekitar $\R^n$, tetapi berhubungan sedemikian erat dengan hipermuka $Y$ sehingga kita ingin memahami dan mempelajarinya semata-mata sebagai objek pada $Y$. Untuk sementara, objek-objek ini memerlukan ruang sekitar karena definisinya mengacu pada analisis berdimensi lebih tinggi, yang saat ini baru tersedia untuk ruang vektor
\zusatzklammer {linear} {} {}.
Salah satu tujuan penting kita adalah mengembangkan konsep-konsep ini hanya pada objek geometris
\zusatzklammer {nonlinear} {} {}
$Y$. Kita akan menyinggung kurva pada hipermuka, ekstremum dengan kendala, dan persamaan diferensial untuk medan vektor tangensial.

\inputbemerkung
{}
{

Misalkan

\mavergleichskette
{\vergleichskette
{ W
}
{ \subseteq }{ \R^n
}
{  }{
}
{    }{
}
{   }{
}
}
{}{}{}
\definitionsverweis {terbuka}{}{,}
\maabb {h} { W } { \R
} {}
adalah sebuah
\definitionsverweis {fungsi yang terdiferensialkan secara kontinu}{}{}
dan

\mavergleichskette
{\vergleichskette
{ Y
}
{ = }{ h^{-1}(c)
}
{  }{
}
{    }{
}
{   }{
}
}
{}{}{}
adalah
\definitionsverweis {serat}{}{}
di atas

\mavergleichskette
{\vergleichskette
{ c
}
{ \in }{ \R
}
{  }{
}
{    }{
}
{   }{
}
}
{}{}{,}
dengan $h$
\definitionsverweis {reguler}{}{}
di setiap titik pada $Y$. Misalkan
\maabbdisp {\gamma} { I } { \R^n
} {}
adalah sebuah
\definitionsverweis {kurva terdiferensialkan}{}{,}
yang seluruhnya terletak di $Y$, dengan $I$ suatu interval terbuka. Maka, pada setiap saat

\mavergleichskette
{\vergleichskette
{ t
}
{ \in }{ I
}
{  }{
}
{    }{
}
{   }{
}
}
{}{}{}
\definitionsverweis {turunan}{}{}
$\gamma'(t)$ berada di ruang singgung pada $Y$ di titik $\gamma(t)$. Hal ini didasarkan pada sifat konstan dari

\mavergleichskette
{\vergleichskette
{ h \circ \gamma
}
{ = }{ c
}
{  }{
}
{    }{
}
{   }{
}
}
{}{}{,}
yang, melalui aturan rantai, memberikan

\mavergleichskettedisp
{\vergleichskette
{ \left(Dh \circ \gamma\right)_{t}
}
{ =} { \left(Dh  \right)_{ \gamma(t) }  \circ \left(D \gamma  \right)_{ t }
}
{ =} { \left(Dh  \right)_{ \gamma(t) } ( \gamma'( t) )
}
{ =} { 0
}
{ } {
}
}
{}{}{,}
dan karenanya

\mavergleichskette
{\vergleichskette
{ \gamma'( t)
}
{ \in }{ \ker  \left(Dh \right)_{ \gamma(t) }
}
{ = }{ T_{\gamma(t) }Y
}
{    }{
}
{   }{
}
}
{}{}{.}
Selain itu, dari
Soal 1.1
diperoleh bahwa setiap vektor singgung di $P$ pada $Y$ dapat direalisasikan oleh suatu kurva terdiferensialkan pada $Y$. Jadi, ruang singgung dapat dijelaskan secara bermakna hanya dengan kurva-kurva terdiferensialkan pada $Y$, meskipun keterdiferensialan kurva tersebut masih mengandaikan ruang sekitar.

}

\inputbemerkung
{}
{

Misalkan

\mavergleichskette
{\vergleichskette
{ W
}
{ \subseteq }{ \R^n
}
{  }{
}
{    }{
}
{   }{
}
}
{}{}{}
\definitionsverweis {terbuka}{}{,}
\maabb {h} { W } { \R
} {}
adalah sebuah
\definitionsverweis {fungsi yang terdiferensialkan secara kontinu}{}{}
dan

\mavergleichskette
{\vergleichskette
{ Y
}
{ = }{ h^{-1}(c)
}
{  }{
}
{    }{
}
{   }{
}
}
{}{}{}
adalah
\definitionsverweis {serat}{}{}
di atas

\mavergleichskette
{\vergleichskette
{ c
}
{ \in }{ \R
}
{  }{
}
{    }{
}
{   }{
}
}
{}{}{,}
dengan $h$
\definitionsverweis {reguler}{}{}
di setiap titik pada $Y$. Kita merangkum kembali teorema tentang ekstremum lokal dengan kendala dalam situasi ini; bandingkan
Akibat 54.3 (Analisis (Osnabrück 2021--2023))
\zusatzklammer {dengan notasi yang berbeda} {} {.}
Selain fungsi $h$, yang menentukan $Y$ dan dengan demikian menentukan kendala, terdapat pula suatu lingkungan terbuka $U$ dari $Y$ beserta sebuah fungsi bernilai real $f$ yang didefinisikan di sana. Yang dicari adalah ekstremum lokal $f$, tetapi bukan pada seluruh $U$---yang biasanya diselidiki dengan gradien dan matriks Hesse---melainkan hanya pada $Y$. Misalnya, kita mungkin ingin mencari suhu maksimum pada permukaan suatu benda sambil mengabaikan bahwa bagian dalamnya lebih panas. Teorema tersebut menyatakan bahwa jika $f$ terdiferensialkan secara kontinu dan memiliki ekstremum lokal di suatu titik

\mavergleichskette
{\vergleichskette
{ P
}
{ \in }{ Y
}
{  }{
}
{    }{
}
{   }{
}
}
{}{}{}
di bawah kendala $Y$, maka gradien
\mathl{\operatorname{Grad} \,  f  (  P  )}{} merupakan kelipatan dari gradien
\mathl{\operatorname{Grad} \,  h  (  P  )}{}. Kriteria keterdiferensialan ini mengacu pada ruang sekitar: baik karena keterdiferensialan itu sendiri didefinisikan dengan mengacu pada ruang sekitar, maupun karena kedua gradien yang dibandingkan menjulur ke dalam ruang sekitar tersebut.

}



\inputfaktbeweis
{Gewöhnliches Differentialgleichungssystem/Hyperfläche/Lösung/Fakt}
{Teorema}
{}
{

\faktsituation {Misalkan

\mavergleichskette
{\vergleichskette
{ W
}
{ \subseteq }{ \R^n
}
{  }{
}
{    }{
}
{   }{
}
}
{}{}{}
\definitionsverweis {terbuka}{}{,}
\maabb {h} { W } { \R
} {}
adalah sebuah
\definitionsverweis {fungsi yang terdiferensialkan secara kontinu}{}{}
dan

\mavergleichskette
{\vergleichskette
{ Y
}
{ = }{ h^{-1}(c)
}
{  }{
}
{    }{
}
{   }{
}
}
{}{}{}
adalah
\definitionsverweis {serat}{}{}
di atas

\mavergleichskette
{\vergleichskette
{ c
}
{ \in }{ \R
}
{  }{
}
{    }{
}
{   }{
}
}
{}{}{,}
dengan $h$
\definitionsverweis {reguler}{}{}
di setiap titik pada $Y$.
Misalkan

\mavergleichskette
{\vergleichskette
{ Y
}
{ \subseteq }{ U
}
{ \subseteq }{ \R^n
}
{    }{
}
{   }{
}
}
{}{}{}
\definitionsverweis {terbuka}{}{}
dan misalkan $F$ adalah sebuah
\definitionsverweis {medan vektor}{}{}
terdiferensialkan pada $U$}
\faktvoraussetzung {dengan sifat bahwa untuk setiap titik

\mavergleichskette
{\vergleichskette
{ P
}
{ \in }{ Y
}
{  }{
}
{    }{
}
{   }{
}
}
{}{}{}
vektor $F(P)$ berada dalam
\definitionsverweis {ruang singgung pada serat}{}{}
di $P$ pada $Y$\zusatzfussnote {Medan vektor seperti ini disebut medan vektor tangensial} {.} {.}}
\faktfolgerung {Maka solusi $v(t)$ dari
\definitionsverweis {masalah nilai awal}{}{}
untuk $F$ dengan syarat awal

\mavergleichskette
{\vergleichskette
{ v(t_0)
}
{ \in }{ Y
}
{  }{
}
{    }{
}
{   }{
}
}
{}{}{}
seluruhnya terletak di $Y$.}
\faktzusatz {}
\faktzusatz {}

}
{

Kita bekerja dengan struktur Euklides pada $\R^n$ dan dengan
\definitionsverweis {gradien}{}{}
$\operatorname{Grad} \,  h$ dari $h$. Berdasarkan asumsi, pada $Y$ gradien ini tidak pernah sama dengan $0$, dan karena kekontinuan, sifat tersebut juga berlaku pada suatu lingkungan terbuka dari $Y$. Dengan memperkecil $U$ jika perlu, kita dapat menganggap bahwa gradien itu tidak memiliki titik nol di seluruh $U$. Pada $U$, kita tinjau medan vektor baru $G$ yang didefinisikan oleh

\mavergleichskettedisp
{\vergleichskette
{ G(P)
}
{ =} { F(P)-  { \frac{   \left\langle  F(P) ,
\operatorname{Grad} \,  h  (  P )    \right\rangle   }{  \Vert {
\operatorname{Grad} \,  h  (  P ) } \Vert^2   } }
\operatorname{Grad} \,  h  (  P )
}
{ } {
}
{ } {
}
{ } {
}
}
{}{}{.}
Untuk

\mavergleichskette
{\vergleichskette
{ P
}
{ \in }{ Y
}
{  }{
}
{    }{
}
{   }{
}
}
{}{}{}
berlaku

\mavergleichskette
{\vergleichskette
{ G(P)
}
{ = }{ F(P)
}
{  }{
}
{    }{
}
{   }{
}
}
{}{}{,}
karena pada $Y$ gradien $h$ tegak lurus terhadap medan vektor tersebut. Selain itu, $G$
\zusatzklammer {berbeda dengan $F$} {} {}
memiliki sifat bahwa untuk semua titik

\mavergleichskette
{\vergleichskette
{ P
}
{ \in }{ U
}
{  }{
}
{    }{
}
{   }{
}
}
{}{}{}
vektor $G(P)$ tegak lurus terhadap gradien, sebab

\mavergleichskettealignhandlinks
{\vergleichskettealignhandlinks
{ \left\langle  G(P) ,
\operatorname{Grad} \,  h  (  P )   \right\rangle
}
{ =} { \left\langle  F(P)- { \frac{   \left\langle  F(P) ,
\operatorname{Grad} \,  h  (  P )    \right\rangle  }{  \Vert {
\operatorname{Grad} \,  h  (  P ) } \Vert^2  } }
\operatorname{Grad} \,  h  (  P )  ,
\operatorname{Grad} \,  h  (  P )   \right\rangle
}
{ =} { \left\langle  F(P) ,
\operatorname{Grad} \,  h  (  P )   \right\rangle - \left\langle  F(P) ,
\operatorname{Grad} \,  h  (  P )   \right\rangle
}
{ =} { 0
}
{ } {
}
}
{}
{}{.}
Sekarang misalkan
\maabbdisp {v} { I} { \R^n
} {}
adalah, menurut
Teorema 56.2 (Analisis (Osnabrück 2021--2023))
\zusatzfussnote {Catatan edisi: Dengan asumsi yang dinyatakan, fungsi h hanya berkelas C¹ sehingga gradiennya sekadar kontinu; hal ini belum membuktikan bahwa medan G lokal Lipschitz sebagaimana diperlukan untuk kesimpulan keunikan Picard--Lindelöf. Langkah ini memerlukan regularitas tambahan atau pembuktian lokal melalui koordinat manifold.} {} {}
,
solusi tunggal dari masalah nilai awal untuk $G$ dengan syarat awal
\mavergleichskette
{\vergleichskette
{ v(t_0)
}
{ = }{ P
}
{ \in }{ Y
}
{    }{
}
{   }{
}
}
{}{}{.}
Maka

\mavergleichskettealign
{\vergleichskettealign
{ \left(Dh \circ v \right)_{t}
}
{ =} { \left(Dh  \right)_{v(t)} \circ  \left(D v \right)_{t}
}
{ =} { \left(Dh  \right)_{v(t)}  \left(  v'(t)  \right)
}
{ =} { \left(Dh  \right)_{v(t)}  \left(  G(v(t))  \right)
}
{ =} { \left\langle
\operatorname{Grad} \,  h  ( v(t) )  ,  G(v(t))    \right\rangle
}
}
{
\vergleichskettefortsetzungalign
{ =} { 0
}
{ } {}
{ } {}
{ } {}
}
{}{.}
Oleh karena itu, $h$ konstan pada citra $v( I)$, dan karena

\mavergleichskette
{\vergleichskette
{ h(v(t_0))
}
{ = }{ c
}
{  }{
}
{    }{
}
{   }{
}
}
{}{}{}
maka

\mavergleichskette
{\vergleichskette
{ h(v(t))
}
{ = }{ c
}
{  }{
}
{    }{
}
{   }{
}
}
{}{}{}
untuk semua

\mavergleichskette
{\vergleichskette
{ t
}
{ \in }{ I
}
{  }{
}
{    }{
}
{   }{
}
}
{}{}{,}
sehingga

\mavergleichskette
{\vergleichskette
{ v(t)
}
{ \in }{ Y
}
{  }{
}
{    }{
}
{   }{
}
}
{}{}{}
untuk semua

\mavergleichskette
{\vergleichskette
{ t
}
{ \in }{ I
}
{  }{
}
{    }{
}
{   }{
}
}
{}{}{.}

}

Dengan regularitas tambahan, cara lain untuk membuktikan pernyataan ini adalah dengan memindahkan medan vektor $F$ melalui suatu difeomorfisme lokal antara

\mavergleichskette
{\vergleichskette
{ P
}
{ \in }{ V
}
{ \subseteq }{ Y
}
{    }{
}
{   }{
}
}
{}{}{}
dan

\mavergleichskette
{\vergleichskette
{ V'
}
{ \subseteq }{ \R^{n-1}
}
{  }{
}
{    }{
}
{   }{
}
}
{}{}{}
ke $\R^{n-1}$, memastikan bahwa medan dalam koordinat bersifat lokal Lipschitz, membuktikan keberadaan solusi di sana, memindahkan kembali solusi itu ke $Y$, lalu merekatkan solusi-solusi parsial tersebut. Tanpa regularitas tersebut, argumen alternatif ini belum lengkap.

Lingkungan terbuka $U$ dari $Y$ diperlukan agar kita dapat berbicara tentang medan vektor terdiferensialkan, meskipun pada akhirnya hanya medan vektor pada $Y$ yang relevan.

  \zwischenueberschrift{Percepatan}

\inputbemerkung
{}
{

Kita tinjau kurva
\maabbeledisp {v} {\R} { \R^2
} { t } { \begin{pmatrix}  \cos  t  \\ \sin  t    \end{pmatrix}
} {,}
yang seluruhnya bergerak pada lingkaran satuan. Dengan demikian, turunannya

\mavergleichskette
{\vergleichskette
{ v'(t)
}
{ = }{ \begin{pmatrix}  - \sin  t  \\ \cos  t    \end{pmatrix}
}
{  }{
}
{    }{
}
{   }{
}
}
{}{}{}
selalu tangensial terhadap lingkaran satuan
\zusatzklammer {hal ini juga mengikuti dari
Catatan 1.1.3} {} {.}
Normanya selalu sama dengan $1$, sehingga lingkaran satuan ditempuh secara seragam dengan besar kecepatan konstan. Namun, arah kecepatannya berubah secara kontinu, dan turunan keduanya adalah

\mavergleichskettedisp
{\vergleichskette
{ v^{\prime \prime}(t)
}
{ =} { \begin{pmatrix}  - \cos  t  \\ - \sin  t    \end{pmatrix}
}
{ =} { - \begin{pmatrix}  \cos  t  \\ \sin  t    \end{pmatrix}
}
{ =} { - v(t)
}
{ } {
}
}
{}{}{.}
Percepatannya adalah negatif dari vektor posisi dan tegak lurus terhadap vektor kecepatan. Jika ruang singgung dari bidang sekitar di sebuah titik $P$ pada lingkaran---yakni $\R^2$ yang dipandang berpusat di $P$---diuraikan menjadi arah yang tangensial terhadap lingkaran
\zusatzklammer {ruang singgung pada serat} {} {}
dan arah yang tegak lurus terhadapnya
\zusatzklammer {ruang normal} {} {}
\zusatzklammer {dalam arti penguraian suatu ruang vektor sebagai
\definitionsverweis {jumlah langsung}{}{}
dari subruang-subruang vektor} {} {,}
maka seluruh percepatan berlangsung di ruang normal; proyeksi ortogonalnya pada ruang singgung selalu $0$. Dalam pengertian ini, percepatan tidak relevan bagi gerak pada lingkaran satuan.

Sekarang kita tinjau secara lebih umum sebuah kurva
\maabbeledisp {w} {\R} { \R^2
} { t } { \begin{pmatrix}  \cos  f(t)  \\ \sin f(t)    \end{pmatrix}
} {,}
dengan $f$ suatu fungsi real yang terdiferensialkan dua kali. Seperti sebelumnya, kurva ini bergerak pada lingkaran satuan, dan turunannya adalah

\mavergleichskettedisp
{\vergleichskette
{ w'(t)
}
{ =} { f'(t) \begin{pmatrix}  - \sin  f(t)  \\ \cos f(t)    \end{pmatrix}
}
{ } {
}
{ } {
}
{ } {
}
}
{}{}{,}
sehingga, seperti sebelumnya, kurva ini selalu tangensial terhadap lingkaran satuan, sedangkan norma kecepatannya sekarang sama dengan $\betrag { f'(t) }$. Turunan keduanya adalah

\mavergleichskettedisp
{\vergleichskette
{ w^{\prime \prime} (t)
}
{ =} { f^{\prime \prime} (t) \begin{pmatrix}  - \sin  f(t)  \\ \cos f(t)    \end{pmatrix} - \left(  f'(t)  \right)^2 \begin{pmatrix}  \cos  f(t)  \\ \sin f(t)    \end{pmatrix}
}
{ } {
}
{ } {
}
{ } {
}
}
{}{}{,}
yang secara langsung memperlihatkan penguraian ke dalam komponen tangensial dan komponen yang ortogonal terhadapnya. Suku pertama, yang memuat turunan kedua fungsi, adalah komponen tangensial. Komponen normalnya adalah
\mathl{- \left(  f'(t)  \right)^2 \begin{pmatrix}  \cos  f(t)  \\ \sin f(t)    \end{pmatrix}}{,} sedangkan percepatan tangensial muncul ketika turunan kedua $f$ tidak sama dengan nol.

}

Untuk sebuah hipermuka

\mavergleichskette
{\vergleichskette
{ Y
}
{ \subseteq }{ \R^n
}
{  }{
}
{    }{
}
{   }{
}
}
{}{}{}
dan sebuah titik

\mavergleichskette
{\vergleichskette
{ P
}
{ \in }{ Y
}
{  }{
}
{    }{
}
{   }{
}
}
{}{}{}
kita selalu dapat menguraikan ruang sekitar $\R^n$, yang dipandang sebagai ruang singgung dari $\R^n$ di $P$, secara ortogonal sebagai

\mavergleichskettedisp
{\vergleichskette
{ \R^n
}
{ =} { T_PY \oplus N_PY
}
{ } {
}
{ } {
}
{ } {
}
}
{}{}{,}
dengan $N_PY$ sebuah garis yang terdiri atas semua titik yang ortogonal terhadap $T_PY$ dan disebut \stichwort {garis normal} {}. Garis normal terdiri atas semua kelipatan gradien
\mathl{\operatorname{Grad} \,  h  (  P )}{} di $P$, jika $h$ menyatakan suatu fungsi yang mendeskripsikan $Y$. Sebuah \stichwort {vektor normal} {} adalah elemen garis normal bernorma $1$. Ada dua vektor normal demikian, dan yang satu merupakan negatif dari yang lain. Melalui
\definitionsverweis {proyeksi ortogonal}{}{}
sepanjang $N_PY$, yakni pemetaan
\maabbdisp {\pi_P} {\R^n } { T_PY
} {}
kita dapat memetakan setiap vektor di $\R^n$ ke sebuah vektor di ruang singgung.

\inputdefinition
{}
{

Misalkan

\mavergleichskette
{\vergleichskette
{ W
}
{ \subseteq }{ \R^n
}
{  }{
}
{    }{
}
{   }{
}
}
{}{}{}
\definitionsverweis {terbuka}{}{,}
\maabb {h} { W } { \R
} {}
adalah sebuah
\definitionsverweis {fungsi yang terdiferensialkan secara kontinu}{}{}
dan

\mavergleichskette
{\vergleichskette
{ Y
}
{ = }{ h^{-1}(c)
}
{  }{
}
{    }{
}
{   }{
}
}
{}{}{}
adalah
\definitionsverweis {serat}{}{}
di atas

\mavergleichskette
{\vergleichskette
{ c
}
{ \in }{ \R
}
{  }{
}
{    }{
}
{   }{
}
}
{}{}{,}
dengan $h$
\definitionsverweis {reguler}{}{}
di setiap titik pada $Y$. Misalkan
\maabbdisp {\gamma} {I} {Y
} {}
merupakan
\zusatzklammer {sebagai pemetaan ke $\R^n$} {} {}
sebuah \definitionsverweis {kurva yang terdiferensialkan dua kali secara kontinu}{}{.}
Maka pemetaan
\maabbeledisp {} {I} { TY
} {t} { \pi_{\gamma(t)}( \gamma^{\prime \prime} (t))
} {,}
dengan $\pi_{\gamma(t)}$ menyatakan proyeksi ortogonal
\maabbdisp {} {\R^n } { T_{\gamma(t)}Y
} {}
disebut
\definitionswort {turunan tangensial kedua}{}
\zusatzklammer {atau
\definitionswort {percepatan tangensial}{}} {} {}
dari $\gamma$.

}

Gagasan tentang turunan kedua yang tangensial terhadap hipermuka dikembangkan lebih lanjut dalam konteks yang disebut koneksi.

  \zwischenueberschrift{Kurva Geodesik}

\inputdefinition
{}
{

Misalkan

\mavergleichskette
{\vergleichskette
{ W
}
{ \subseteq }{ \R^n
}
{  }{
}
{    }{
}
{   }{
}
}
{}{}{}
\definitionsverweis {terbuka}{}{,}
\maabb {h} { W } { \R
} {}
adalah sebuah
\definitionsverweis {fungsi yang terdiferensialkan secara kontinu}{}{}
dan

\mavergleichskette
{\vergleichskette
{ Y
}
{ = }{ h^{-1}(c)
}
{  }{
}
{    }{
}
{   }{
}
}
{}{}{}
adalah
\definitionsverweis {serat}{}{}
di atas

\mavergleichskette
{\vergleichskette
{ c
}
{ \in }{ \R
}
{  }{
}
{    }{
}
{   }{
}
}
{}{}{,}
dengan $h$
\definitionsverweis {reguler}{}{}
di setiap titik pada $Y$. Sebuah
\zusatzklammer {sebagai pemetaan ke $\R^n$} {} {}
\definitionsverweis {kurva yang terdiferensialkan dua kali secara kontinu}{}{}
\maabbdisp {\gamma} {I} {Y
} {}
disebut
\definitionswort {geodesik}{}
\zusatzklammer {atau
\definitionswort {kurva geodesik}{}
pada $Y$} {} {,}
jika
\definitionsverweis {percepatan tangensial}{}{}
kurva itu sama dengan $0$ di mana-mana.

}

Sebuah kurva yang terdiferensialkan dua kali merupakan geodesik tepat ketika turunan keduanya selalu ortogonal terhadap ruang singgung. Gagasan fisika yang mendasarinya adalah gerak sebuah partikel yang bergerak pada $Y$ tanpa percepatan. Secara umum memang terdapat gaya-gaya kendala, seperti gravitasi atau tarikan magnetis, yang memaksa geraknya tetap pada $Y$ dan pada gilirannya memerlukan percepatan di ruang sekitar; namun pada $Y$ sendiri tidak ada percepatan. Jadi, kurva geodesik menggambarkan gerak alami tanpa percepatan pada sebuah hipermuka terdiferensialkan.

\bild{ \begin{center}
\includegraphics[width=5.5cm]{ \bildeinlesungsvg {Great circle passing through two points} {svg} }
\end{center}
\bildtext {} }

\bildlizenz { Great circle passing through two points.svg } {} {HaEr48} {Commons} {CC BY-SA 4.0} {}

\inputdefinition
{}
{

Irisan sebuah
\definitionsverweis {permukaan bola}{}{}
\mathl{S}{} dengan suatu bidang yang melalui pusat bola disebut sebuah \definitionswort {lingkaran besar}{} pada $S$.

}

Pada bola Bumi, khatulistiwa dan lingkaran-lingkaran bujur merupakan lingkaran besar; lingkaran-lingkaran lintang pada umumnya bukan.

\inputbeispiel{}
{

Kita berada pada bola satuan

\mavergleichskettedisp
{\vergleichskette
{ Y
}
{ =} { \left\{ (x,y,z)  \mid   x^2+y^2+z^2 = 1         \right\}
}
{ \subseteq} { \R^3
}
{ } {
}
{ } {
}
}
{}{}{.}
Di

\mavergleichskette
{\vergleichskette
{ P
}
{ \in }{ Y
}
{  }{
}
{    }{
}
{   }{
}
}
{}{}{}
terdapat sebuah partikel yang menerima impuls gerak dalam arah tangensial tertentu yang dinyatakan oleh

\mavergleichskette
{\vergleichskette
{ v
}
{ \in }{ T_PY
}
{  }{
}
{    }{
}
{   }{
}
}
{}{}{.}
Partikel itu terus bergerak dalam arah ini
\zusatzklammer {tanpa diperlambat dan tanpa dipercepat} {} {}
namun harus selalu tetap berada pada bola
\zusatzklammer {karena gravitasi Bumi atau karena permukaan yang kukuh} {} {.}
Misalkan titik tersebut diberikan oleh
\mathl{\begin{pmatrix}  x_1  \\ y_1 \\ z_1   \end{pmatrix}}{} dan vektor singgung tegak lurus yang menyatakan impuls sesaat diberikan oleh

\mavergleichskettedisp
{\vergleichskette
{ \begin{pmatrix}  x_3  \\ y_3 \\ z_3   \end{pmatrix}
}
{ =} { a \begin{pmatrix}  x_2  \\ y_2 \\ z_2   \end{pmatrix}
}
{ } {
}
{ } {
}
{ } {
}
}
{}{}{,}
dengan $\begin{pmatrix}  x_2  \\ y_2 \\ z_2   \end{pmatrix}$ suatu vektor bernorma $1$ yang tegak lurus terhadap vektor posisi. Gerak tersebut berada pada bidang yang ditentukan oleh kedua vektor ini sekaligus pada bola. Salah satu parametrisasi geraknya adalah
\maabbeledisp {\gamma} {\R} { \R^3
} { t } { \cos \left( at \right) \begin{pmatrix}  x_1  \\ y_1 \\ z_1   \end{pmatrix} + \sin  \left( at \right) \begin{pmatrix}  x_2  \\ y_2 \\ z_2   \end{pmatrix}
} {.}
Berlaku

\mavergleichskette
{\vergleichskette
{ \gamma(0)
}
{ = }{ \begin{pmatrix}  x_1  \\ y_1 \\ z_1   \end{pmatrix}
}
{  }{
}
{    }{
}
{   }{
}
}
{}{}{}
dan

\mavergleichskette
{\vergleichskette
{ \gamma'(t)
}
{ = }{ -a \sin  \left( at \right) \begin{pmatrix}  x_1  \\ y_1 \\ z_1   \end{pmatrix} +a \cos \left( at \right) \begin{pmatrix}  x_2  \\ y_2 \\ z_2   \end{pmatrix}
}
{  }{
}
{    }{
}
{   }{
}
}
{}{}{,}
sehingga

\mavergleichskette
{\vergleichskette
{ \gamma'(0)
}
{ = }{ a \begin{pmatrix}  x_2  \\ y_2 \\ z_2   \end{pmatrix}
}
{  }{
}
{    }{
}
{   }{
}
}
{}{}{.}

}

\inputbeispiel{}
{

Kita tinjau selimut silinder

\mavergleichskettedisp
{\vergleichskette
{ Z
}
{ =} { \left\{  \begin{pmatrix}  x  \\ y \\ z   \end{pmatrix} \in \R^3  \mid   x^2+y^2 = 1         \right\}
}
{ } {
}
{ } {
}
{ } {
}
}
{}{}{.}
Secara geometris, kita memperoleh gerakan-gerakan berikut pada silinder, yang percepatannya selalu ortogonal terhadap ruang singgung. Ruang singgung di titik
\mathl{\begin{pmatrix}  x  \\ y \\ z   \end{pmatrix}}{} diberikan oleh
\mathl{\R  \begin{pmatrix}  y  \\ -x\\  0   \end{pmatrix} + \R \begin{pmatrix}  0  \\ 0\\  1   \end{pmatrix}}{}, dan vektor $- \begin{pmatrix}  x  \\ y \\  0   \end{pmatrix}$ tegak lurus terhadapnya. Untuk gerak-gerak berikut, penskalaan parameter waktunya dapat diubah secara afin.
\aufzaehlungdrei{Gerak pada sebuah garis di selimut silinder yang sejajar dengan sumbu silinder:

\mavergleichskettedisp
{\vergleichskette
{ \gamma(t)
}
{ =} { \begin{pmatrix}  x_0  \\ y_0 \\ t   \end{pmatrix}
}
{ } {
}
{ } {
}
{ } {
}
}
{}{}{}
dengan

\mavergleichskette
{\vergleichskette
{ x_0^2 +y_0^2
}
{ = }{ 1
}
{  }{
}
{    }{
}
{   }{
}
}
{}{}{.}
Turunan pertamanya adalah $\begin{pmatrix}  0 \\ 0\\  1   \end{pmatrix}$ dan turunan keduanya adalah $0$.
}{Gerak melingkar yang diberikan oleh

\mavergleichskettedisp
{\vergleichskette
{ \delta(t)
}
{ =} { \begin{pmatrix}  \cos  t  \\ \sin  t \\ c   \end{pmatrix}
}
{ } {
}
{ } {
}
{ } {
}
}
{}{}{.}
Turunan pertamanya adalah $\begin{pmatrix}  - \sin  t  \\ \cos  t \\  0   \end{pmatrix}$ dan turunan keduanya adalah $- \begin{pmatrix}  \cos  t  \\ \sin  t \\  0   \end{pmatrix}$.
}{

\bild{ \begin{center}
\includegraphics[width=5.5cm]{ \bildeinlesungjpg {2019-07-Helix} {jpg} }
\end{center}
\bildtext {} }

\bildlizenz {  2019-07-Helix.jpg } {} {Ag2gaeh} {Commons} {CC BY-SA 4.0} {}

Sebuah kurva ulir
\zusatzklammer {heliks} {} {}
yang diberikan oleh

\mavergleichskettedisp
{\vergleichskette
{ \epsilon(t)
}
{ =} { \begin{pmatrix}  \cos  t  \\ \sin  t \\  e t   \end{pmatrix}
}
{ } {
}
{ } {
}
{ } {
}
}
{}{}{}
dengan

\mavergleichskette
{\vergleichskette
{ e
}
{ \neq }{ 0
}
{  }{
}
{    }{
}
{   }{
}
}
{}{}{.}
Turunan pertamanya adalah $\begin{pmatrix}  - \sin  t  \\ \cos  t \\ e   \end{pmatrix}$ dan turunan keduanya adalah $- \begin{pmatrix}  \cos  t  \\ \sin  t \\  0   \end{pmatrix}$.
}

}

\chapter{Lembar Kerja 1}
\setcounter{section}{1}

  \zwischenueberschrift{Soal latihan}

\inputaufgabegibtloesung
{}
{

Misalkan

\mavergleichskette
{\vergleichskette
{ G
}
{ \subseteq }{ \R^n
}
{  }{
}
{    }{
}
{   }{
}
}
{}{}{}
\definitionsverweis {terbuka}{}{}
dan
\maabbdisp {\varphi} { G } { \R^m
} {}
sebuah
\definitionsverweis {pemetaan terdiferensialkan secara kontinu}{}{,}
yang di titik

\mavergleichskette
{\vergleichskette
{ P
}
{ \in }{ G
}
{  }{
}
{    }{
}
{   }{
}
}
{}{}{}
mempunyai
\definitionsverweis {diferensial total}{}{}
yang surjektif. Misalkan

\mavergleichskette
{\vergleichskette
{ v
}
{ \in }{ T_PZ
}
{  }{
}
{    }{
}
{   }{
}
}
{}{}{}
sebuah vektor di
\definitionsverweis {ruang singgung pada serat}{}{}
$Z$ dari $\varphi$ yang melalui $P$. Buktikan bahwa terdapat
\definitionsverweis {kurva terdiferensialkan secara kontinu}{}{}
\maabbdisp {\gamma} { ]- \epsilon, \epsilon[ } { Z
} {}
\zusatzklammer {untuk suatu

\mavergleichskettek
{\vergleichskettek
{ \epsilon
}
{ > }{ 0
}
{  }{
}
{    }{
}
{   }{
}
}
{}{}{}
yang sesuai} {} {}
dengan

\mavergleichskette
{\vergleichskette
{ \gamma(0)
}
{ = }{ P
}
{  }{
}
{    }{
}
{   }{
}
}
{}{}{}
dan

\mavergleichskettedisp
{\vergleichskette
{ \gamma'(0)
}
{ =} { v
}
{ } {
}
{ } {
}
{ } {
}
}
{}{}{.}

}
{} {}

\inputaufgabe
{}
{

Misalkan

\mavergleichskette
{\vergleichskette
{ W
}
{ \subseteq }{ \R^n
}
{  }{
}
{    }{
}
{   }{
}
}
{}{}{}
\definitionsverweis {terbuka}{}{,}
\maabb {h} { W } { \R
} {}
sebuah
\definitionsverweis {fungsi terdiferensialkan secara kontinu}{}{}
dan

\mavergleichskette
{\vergleichskette
{ Y
}
{ = }{ h^{-1}(0)
}
{  }{
}
{    }{
}
{   }{
}
}
{}{}{}
merupakan
\definitionsverweis {serat}{}{}
di atas $0$, dengan $h$
\definitionsverweis {reguler}{}{}
di setiap titik pada $Y$. Misalkan
\maabbdisp {g} { W } {\R
} {}
sebuah fungsi terdiferensialkan secara kontinu yang tidak mempunyai titik nol pada $Y$. Buktikan bahwa $Y$ juga dapat diperoleh sebagai serat dari

\mavergleichskettedisp
{\vergleichskette
{ \tilde{h}
}
{ =} { g h
}
{ } {
}
{ } {
}
{ } {
}
}
{}{}{,}
bahwa
\definitionsverweis {gradien}{}{}
dari $h$ dan $\tilde{h}$ merupakan kelipatan skalar satu sama lain, dan bahwa
\definitionsverweis {ruang-ruang singgung}{}{}
hanya bergantung pada $Y$.

}
{} {}

\inputaufgabe
{}
{

Misalkan

\mavergleichskette
{\vergleichskette
{ W
}
{ \subseteq }{ \R^n
}
{  }{
}
{    }{
}
{   }{
}
}
{}{}{}
\definitionsverweis {terbuka}{}{,}
\maabb {h} { W } { \R
} {}
sebuah
\definitionsverweis {fungsi terdiferensialkan secara kontinu}{}{}
dan

\mavergleichskette
{\vergleichskette
{ Y
}
{ = }{ h^{-1}(c)
}
{  }{
}
{    }{
}
{   }{
}
}
{}{}{}
merupakan
\definitionsverweis {serat}{}{}
di atas

\mavergleichskette
{\vergleichskette
{ c
}
{ \in }{ \R
}
{  }{
}
{    }{
}
{   }{
}
}
{}{}{,}
dengan $h$
\definitionsverweis {reguler}{}{}
di setiap titik pada $Y$. Misalkan
\maabbdisp {L} {\R^n } { \R^n
} {}
sebuah
\definitionsverweis {isomorfisme linear}{}{,}

\mavergleichskette
{\vergleichskette
{ W'
}
{ = }{ L^{-1}(W)
}
{  }{
}
{    }{
}
{   }{
}
}
{}{}{}
dan

\mavergleichskettedisp
{\vergleichskette
{ Z
}
{ =} { L^{-1}(Y)
}
{ =} { (h \circ L)^{-1}(c)
}
{ } {
}
{ } {
}
}
{}{}{.}
Buktikan bahwa konsep
\definitionsverweis {ruang singgung}{}{,}
medan vektor tangensial, dan solusi persamaan diferensial tangensial yang terkait pada $Y$ dan $Z$ secara umum saling bersesuaian
\zusatzklammer {yakni dapat dipetakan satu sama lain dengan bantuan $L$} {} {.}
Bagaimana dengan gradiennya?

}
{} {}

\inputaufgabe
{}
{

Misalkan

\mavergleichskette
{\vergleichskette
{ Y
}
{ = }{ h^{-1}(c)
}
{  }{
}
{    }{
}
{   }{
}
}
{}{}{}
sebuah
\definitionsverweis {hipermuka terdiferensialkan}{}{}
dan misalkan
\maabbdisp {f} { U} { \R
} {}
sebuah fungsi terdiferensialkan secara kontinu yang didefinisikan pada suatu lingkungan terbuka

\mavergleichskette
{\vergleichskette
{ Y
}
{ \subseteq }{ U
}
{  }{
}
{    }{
}
{   }{
}
}
{}{}{}
Andaikan $f$ di

\mavergleichskette
{\vergleichskette
{ P
}
{ \in }{ Y
}
{  }{
}
{    }{
}
{   }{
}
}
{}{}{}
mempunyai
\definitionsverweis {ekstremum lokal}{}{}
pada $Y$. Buktikan bahwa
\definitionsverweis {komponen tangensial}{}{}
dari $\operatorname{Grad} \,  f  (  P )$ sama dengan $0$.

}
{} {}

\inputaufgabe
{}
{

Misalkan

\mavergleichskette
{\vergleichskette
{ Y
}
{ \subseteq }{ \R^3
}
{  }{
}
{    }{
}
{   }{
}
}
{}{}{}
adalah permukaan bola satuan. Tentukan
\definitionsverweis {dekomposisi ortogonal}{}{}
vektor

\mavergleichskette
{\vergleichskette
{ \begin{pmatrix}  6  \\ 7 \\  -3   \end{pmatrix}
}
{ \in }{ \R^3
}
{  }{
}
{    }{
}
{   }{
}
}
{}{}{}
menjadi komponen tangensial dan komponen ortogonal di titik

\mavergleichskette
{\vergleichskette
{ \begin{pmatrix}  -  { \frac{   1  }{  2 } }  \\ { \frac{   1  }{  2 } } \\  { \frac{   1  }{  \sqrt{2}  } }   \end{pmatrix}
}
{ \in }{ Y
}
{  }{
}
{    }{
}
{   }{
}
}
{}{}{.}

}
{} {}

\inputaufgabe
{}
{

Misalkan

\mavergleichskette
{\vergleichskette
{ Y
}
{ \subseteq }{ \R^3
}
{  }{
}
{    }{
}
{   }{
}
}
{}{}{}
adalah
\definitionsverweis {permukaan bola satuan}{}{.}
Misalkan
\maabbdisp {\gamma} {[0, 2 \pi]} { Y
} {}
merupakan gerak melingkar beraturan yang mengelilingi permukaan bola pada ketinggian ${ \frac{   1  }{  2 } }$.
\aufzaehlungzwei {Parametrisasikan $\gamma$.
} {Tentukan
\definitionsverweis {percepatan tangensial}{}{}
dari $\gamma$ pada setiap saat

\mavergleichskette
{\vergleichskette
{ t
}
{ \in }{ [0, 2 \pi]
}
{  }{
}
{    }{
}
{   }{
}
}
{}{}{.}
}

}
{} {}

\inputaufgabe
{}
{

Misalkan $Y$ adalah
\definitionsverweis {grafik}{}{}
dari
\maabbeledisp {f} {\R^2} { \R
} {(x,y)} { x^2+y^3
} {,}
dan misalkan
\maabbdisp {\gamma} {\R} { Y
} {}
adalah lintasan pada grafik tersebut yang berada di atas kurva dasar
\mathl{t \mapsto (t,t)}{.} Tentukan, untuk setiap

\mavergleichskette
{\vergleichskette
{ t
}
{ \in }{ \R
}
{  }{
}
{    }{
}
{   }{
}
}
{}{}{,}
\definitionsverweis {percepatan tangensial}{}{}
dari $\gamma$ di $\gamma(t)$.

}
{} {}

\inputaufgabe
{}
{

Misalkan

\mavergleichskette
{\vergleichskette
{ Y
}
{ \subseteq }{ \R^n
}
{  }{
}
{    }{
}
{   }{
}
}
{}{}{}
sebuah hiperbidang, yaitu serat dari fungsi linear

\mavergleichskettedisp
{\vergleichskette
{ h(x_1 , \ldots , x_n)
}
{ =} { a_1x_1 + \cdots + a_nx_n
}
{ } {
}
{ } {
}
{ } {
}
}
{}{}{}
di atas

\mavergleichskette
{\vergleichskette
{ c
}
{ \in }{ \R
}
{  }{
}
{    }{
}
{   }{
}
}
{}{}{}
dengan

\mavergleichskette
{\vergleichskette
{ a_i
}
{ \neq }{ 0
}
{  }{
}
{    }{
}
{   }{
}
}
{}{}{}
untuk setidaknya satu indeks $i$. Tafsirkan pengertian
\definitionsverweis {diferensial total}{}{,}
\definitionsverweis {gradien}{}{,}
\definitionsverweis {ruang singgung}{}{,}
dekomposisi menjadi komponen tangensial dan normal,
\definitionsverweis {percepatan tangensial}{}{,}
serta
\definitionsverweis {kurva geodesik}{}{}
dalam kasus khusus ini.

}
{} {}

\inputaufgabe
{}
{

Misalkan

\mavergleichskette
{\vergleichskette
{ Y
}
{ = }{ h^{-1}(c)
}
{  }{
}
{    }{
}
{   }{
}
}
{}{}{}
sebuah
\definitionsverweis {hipermuka terdiferensialkan}{}{,}
yang memuat garis
\mathl{P+ \R v}{.} Buktikan bahwa garis yang dilalui secara seragam tersebut merupakan
\definitionsverweis {kurva geodesik}{}{}
pada $Y$.

}
{} {}

Dua titik

\mathbed {P,Q \in K} {}
 {P \neq Q} {}
 {} {} {} {,}
disebut \stichwort {antipodal} {,} jika garis yang menghubungkannya melalui pusat bola.

\inputaufgabe
{}
{

Misalkan

\mavergleichskette
{\vergleichskette
{ K
}
{ \subseteq }{ \R^3
}
{  }{
}
{    }{
}
{   }{
}
}
{}{}{}
sebuah permukaan bola. Buktikan bahwa setiap dua
\definitionsverweis {titik yang tidak antipodal}{}{}

\mathbed {P,Q \in K} {}
 {P \neq Q} {}
 {} {} {} {,}
terletak pada tepat satu
\definitionsverweis {lingkaran besar}{}{}
dari $K$.

}
{} {}

\bild{ \begin{center}
\includegraphics[width=5.5cm]{ \bildeinlesungsvg {Planned flight map of the Oiseau Blanc} {svg} }
\end{center}
\bildtext {Mengapa pesawat terbang menempuh lintasan melengkung?} }

\bildlizenz { Planned flight map of the Oiseau Blanc.svg } {} {Pethrus} {Commons} {CC BY-SA 3.0} {}

\inputaufgabe
{}
{

Sebuah pesawat akan terbang dari Osnabrück menuju suatu tempat tujuan di belahan Bumi selatan. Mungkinkah rutenya lebih pendek jika pesawat itu mula-mula terbang ke arah utara?

}
{} {}

\inputaufgabe
{}
{

Lucy Sonnenschein merasa bahwa hari-hari terlalu singkat. Karena itu, setiap hari dan selama masih terang, ia memutuskan terbang ke arah barat melintasi satu zona waktu agar setiap harinya berlangsung $25$, bukan $24$, jam.
\aufzaehlungfuenf{Dapatkah Lucy Sonnenschein meningkatkan proporsi jam terang dalam hidupnya dengan strategi ini?
}{Dapatkah Lucy Sonnenschein memperpanjang hidupnya dengan strategi ini?
}{Sekarang Lucy mempunyai teman di seluruh dunia. Setelah berapa hari ia bertemu dengan mereka lagi?
}{Apa yang menyebabkan Lucy mengalami kehilangan waktu yang mengimbangi tambahan waktu hariannya?
}{Apakah hal ini berkaitan dengan teori relativitas?
}

}
{} {}

\inputaufgabe
{}
{

Berikan contoh kurva dua kali terdiferensialkan yang bergerak pada suatu
\definitionsverweis {lingkaran besar}{}{}
dari permukaan bola satuan, tetapi bukan
\definitionsverweis {kurva geodesik}{}{.}

}
{} {}

\inputaufgabe
{}
{

Sebuah \stichwort {segitiga geodesik} {} pada permukaan bola

\mavergleichskettedisp
{\vergleichskette
{ S^2
}
{ =} { \left\{ (x,y,z) \in \R^3  \mid   x^2+y^2+z^2 = 1        \right\}
}
{ } {
}
{ } {
}
{ } {
}
}
{}{}{}
terdiri atas tiga titik berbeda yang tidak terletak pada satu lingkaran besar,

\mavergleichskette
{\vergleichskette
{ A,B,C
}
{ \in }{ S^2
}
{  }{
}
{    }{
}
{   }{
}
}
{}{}{,}
dengan setiap pasangan titik dihubungkan oleh busur yang lebih pendek dari suatu
\definitionsverweis {lingkaran besar}{}{}
\zusatzklammer {setiap pasangan titik tidak antipodal} {} {.}
Buktikan bahwa jumlah sudut dalam segitiga itu lebih besar daripada $\pi$.

}
{} {}

  \zwischenueberschrift{Soal untuk dikumpulkan}

\inputaufgabe
{3}
{

Misalkan $Y$ adalah hipermuka yang diberikan oleh syarat

\mavergleichskettedisp
{\vergleichskette
{ 2x^2 +3y^2+5z^2
}
{ =} { 10
}
{ } {
}
{ } {
}
{ } {
}
}
{}{}{}
di $\R^3$. Tentukan dekomposisi ortogonal vektor

\mavergleichskette
{\vergleichskette
{ \begin{pmatrix}  4  \\ 2 \\  -5   \end{pmatrix}
}
{ \in }{ \R^3
}
{  }{
}
{    }{
}
{   }{
}
}
{}{}{}
menjadi komponen tangensial dan komponen ortogonal di titik

\mavergleichskette
{\vergleichskette
{ \begin{pmatrix}  -1 \\ -1\\  1   \end{pmatrix}
}
{ \in }{ Y
}
{  }{
}
{    }{
}
{   }{
}
}
{}{}{.}

}
{} {}

\inputaufgabe
{4}
{

Misalkan $Y$ adalah
\definitionsverweis {grafik}{}{}
dari
\maabbeledisp {f} {\R^2} { \R
} {(x,y)} { x^2+y^2
} {,}
dan misalkan
\maabbdisp {\gamma} {\R} { Y
} {}
adalah lintasan pada grafik tersebut yang berada di atas kurva dasar
\mathl{x \mapsto (x,1)}{.} Tentukan, untuk setiap

\mavergleichskette
{\vergleichskette
{ x
}
{ \in }{ \R
}
{  }{
}
{    }{
}
{   }{
}
}
{}{}{,}
\definitionsverweis {percepatan tangensial}{}{}
dari $\gamma$ di $\gamma(x)$.

}
{} {}

\inputaufgabe
{4}
{

Pada 11 April 2023 pukul 17.45, secara tiba-tiba dan mengejutkan semua orang, rotasi Bumi pada porosnya dihentikan. Pada saat yang sama, hambatan udara dan semua kerugian akibat gesekan ditiadakan. Semua manusia, hewan, dan benda yang tidak terpasang tetap atau tidak sempat berpegangan akan meneruskan gerak sesaatnya sesuai hukum inersia. Gravitasi tetap bekerja sehingga, untungnya, mereka tidak terlepas ke angkasa. Kapan penghapus papan tulis dari ruang kuliah di Osnabrück
\zusatzklammer {penghapus itu terbang keluar melalui jendela} {} {}
untuk pertama kalinya kembali ke posisi awalnya?

}
{} {}

\inputaufgabe
{4 (3+1)}
{

Misalkan

\mavergleichskette
{\vergleichskette
{ W
}
{ \subseteq }{ \R^n
}
{  }{
}
{    }{
}
{   }{
}
}
{}{}{}
\definitionsverweis {terbuka}{}{,}
\maabb {h} { W } { \R
} {}
sebuah
\definitionsverweis {fungsi terdiferensialkan secara kontinu}{}{}
dan

\mavergleichskette
{\vergleichskette
{ Y
}
{ = }{ h^{-1}(c)
}
{  }{
}
{    }{
}
{   }{
}
}
{}{}{}
merupakan
\definitionsverweis {serat}{}{}
di atas

\mavergleichskette
{\vergleichskette
{ c
}
{ \in }{ \R
}
{  }{
}
{    }{
}
{   }{
}
}
{}{}{,}
dengan $h$
\definitionsverweis {reguler}{}{}
di setiap titik pada $Y$. Misalkan
\maabbdisp {\gamma} { I } { Y
} {}
sebuah kurva dua kali
\definitionsverweis {terdiferensialkan secara kontinu}{}{.}
\aufzaehlungzwei {Buktikan bahwa jika $\gamma$ merupakan
\definitionsverweis {kurva geodesik}{}{,}
maka
\definitionsverweis {norma}{}{}
dari $\gamma'$ adalah konstan.
} {Berikan contoh kurva dengan norma kecepatan konstan yang bukan kurva geodesik.
}

}
{} {}

\inputaufgabe
{3}
{

Misalkan
\maabb {h,g} {\R^n } { \R
} {}
adalah
\definitionsverweis {fungsi-fungsi terdiferensialkan secara kontinu}{}{}
dan

\mavergleichskette
{\vergleichskette
{ Y
}
{ = }{ h^{-1}(c)
}
{  }{
}
{    }{
}
{   }{
}
}
{}{}{}
serta

\mavergleichskette
{\vergleichskette
{ Z
}
{ = }{ g^{-1}(d)
}
{  }{
}
{    }{
}
{   }{
}
}
{}{}{}
adalah
\definitionsverweis {hipermuka terdiferensialkan}{}{}
yang terkait, dengan $h$ di setiap titik pada $Y$ dan $g$ di setiap titik pada $Z$ bersifat
\definitionsverweis {reguler}{}{.}
Misalkan

\mavergleichskette
{\vergleichskette
{ T
}
{ = }{ Y \cap Z
}
{  }{
}
{    }{
}
{   }{
}
}
{}{}{}
dan misalkan
\maabbdisp {\gamma} { I } { T
} {}
sebuah kurva dua kali terdiferensialkan secara kontinu yang, jika dipandang sebagai kurva pada $Y$, merupakan
\definitionsverweis {kurva geodesik}{}{.}
Apakah $\gamma$ juga merupakan kurva geodesik pada $Z$?

}
{} {}

\section*{Solusi yang disediakan oleh sumber}
\addcontentsline{toc}{section}{Solusi yang disediakan oleh sumber}
% Authority: German Wikiversity pageid 131353, revid 1111802, 2026-08-15T16:41:12Z.
% Exact UTF-8 witness SHA-256: b07aac01b6f803481fab9b5331e19480fb9bf0058597159b6152c65c7bc955ce.

\zwischenueberschrift{Solusi untuk Soal 1}

Berlaku

\mavergleichskettedisp
{\vergleichskette
{ v
}
{ \in }{ T_PZ
}
{ = }{ \ker \left((D\varphi)_P\right)
}
{ } {
}
{ } {
}
}
{}{}{.}
Menurut
\definitionsverweis {teorema fungsi implisit}{}{,}
terdapat himpunan terbuka

\mathbed {P \in W} {}
 {W \subseteq G} {}
 {} {} {} {,}
himpunan terbuka

\mavergleichskette
{\vergleichskette
{ V
}
{ \subseteq }{ \R^{n-m}
}
{ } {
}
{ } {
}
{ } {
}
}
{}{}{,}
dan pemetaan terdiferensialkan secara kontinu
\maabbdisp {\psi} { V } { W } {}
sedemikian sehingga

\mavergleichskette
{\vergleichskette
{ \psi(V)
}
{ \subseteq }{ Z \cap W
}
{ } {
}
{ } {
}
{ } {
}
}
{}{}{}
dan $\psi$ menginduksi
\definitionsverweis {bijeksi}{}{}
\maabbdisp {\psi} { V } { Z \cap W } {}.
Misalkan

\mavergleichskette
{\vergleichskette
{ Q
}
{ \in }{ V
}
{ } {
}
{ } {
}
{ } {
}
}
{}{}{}
adalah titik yang memenuhi

\mavergleichskettedisp
{\vergleichskette
{ \psi(Q)
}
{ =} { P
}
{ } {
}
{ } {
}
{ } {
}
}
{}{}{.}
Pemetaan $\psi$ di setiap titik

\mavergleichskette
{\vergleichskette
{ Q
}
{ \in }{ V
}
{ } {
}
{ } {
}
{ } {
}
}
{}{}{}
bersifat
\definitionsverweis {reguler}{}{}
dan
\definitionsverweis {diferensial total}{}{}
dari $\psi$ memenuhi

\mavergleichskettedisp
{\vergleichskette
{ (D\varphi)_P \circ (D\psi)_Q
}
{ =} { 0
}
{ } {
}
{ } {
}
{ } {
}
}
{}{}{.}
Dengan demikian,

\mavergleichskettedisp
{\vergleichskette
{ \operatorname{im} \left((D\psi)_Q\right)
}
{ =} { \ker \left((D\varphi)_P\right)
}
{ =} { T_PZ
}
{ } {
}
{ } {
}
}
{}{}{.}
Karena $\psi$ reguler di $Q$, pemetaan
\maabbdisp {(D\psi)_Q} { \R^{n-m} } { \R^n } {}
bersifat injektif, sedangkan pemetaan
\maabbdisp {(D\psi)_Q} { \R^{n-m} } { T_PZ } {}
bersifat bijektif. Misalkan

\mavergleichskette
{\vergleichskette
{ u
}
{ \in }{ \R^{n-m}
}
{ } {
}
{ } {
}
{ } {
}
}
{}{}{}
adalah prapeta dari $v$, dan misalkan
\maabbeledisp {\theta} { ]-\epsilon,\epsilon[ } { \R^{n-m} } {t} {Q+tu} {,}
dengan $\epsilon$ dipilih cukup kecil agar seluruh citra terletak di dalam $V$. Maka
\maabbdisp {\gamma=\psi\circ\theta} { ]-\epsilon,\epsilon[ } { \R^n } {}
mempunyai sifat

\mavergleichskettedisp
{\vergleichskette
{ \gamma(0)
}
{ =} { \psi(\theta(0))
}
{ =} { \psi(Q)
}
{ =} { P
}
{ } {
}
}
{}{}{}
dan

\mavergleichskettedisp
{\vergleichskette
{ \gamma'(0)
}
{ =} { (D\psi)_Q(\theta'(0))
}
{ =} { (D\psi)_Q(u)
}
{ =} { v
}
{ } {
}
}
{}{}{.}

\part{Unit 2}
\chapter[Kuliah 2]{Kuliah 2: Permukaan Putar, Medan Normal, dan Pemetaan Gauss}
\setcounter{section}{2}

  \zwischenueberschrift{Permukaan Putar}

Kita membahas satu kelas lain dari permukaan di ruang, yaitu permukaan putar.

\bild{ \begin{center}
\includegraphics[width=5.5cm]{ \bildeinlesungsvg {Integral apl rot obsah1} {svg} }
\end{center}
\bildtext {} }

\bildlizenz { Integral apl rot obsah1.svg } {} {Pajs} {Wikipedia bahasa Ceko} {domain publik} {}

Misalkan diberikan suatu
\definitionsverweis {kurva terdiferensialkan}{}{}
\maabbeledisp {\gamma} {]a,b[} {\R^2
} { t } {(x(t),y(t))
} {,}
dengan

\mavergleichskette
{\vergleichskette
{ y(t)
}
{ \geq }{ 0
}
{  }{
}
{    }{
}
{   }{
}
}
{}{}{}
berlaku. Kita tertarik pada
\definitionsverweis {permukaan putar}{}{,}
yang bersesuaian, yaitu himpunan bagian

\mathdisp {\left\{  (x(t), y(t) \cos \alpha , y(t) \sin \alpha)   \mid  t \in ]a,b[ , \, \alpha \in [0,2\pi]       \right\}} {  }

dari
\mathl{\R^3}{,} yang terbentuk ketika lintasan kurva diputar mengelilingi sumbu $x$. Selain itu, kita mengandaikan bahwa $\gamma$ memiliki citra

\mavergleichskette
{\vergleichskette
{ M
}
{ = }{ \gamma \left(  ]a,b[  \right)
}
{  }{
}
{    }{
}
{   }{
}
}
{}{}{}
dan merupakan homeomorfisme ke citra tersebut. Secara khusus, kita meninjau kasus ketika
\maabb {f} {]a,b[} { \R_+
} {}
adalah fungsi terdiferensialkan dan permukaan putarnya berasal dari grafik fungsi tersebut.



\inputfaktbeweis
{Rotationsfläche/Differenzierbare positive Kurve/Graph/Tangentialraum/Fakt}
{Lemma}
{}
{

\faktsituation {Misalkan $I$ adalah sebuah
\definitionsverweis {interval terbuka}{}{}
dan
\maabb {f} { I } { \R_+
} {}
sebuah
\definitionsverweis {fungsi terdiferensialkan}{}{.}}
\faktuebergang {Maka
\definitionsverweis {permukaan putar}{}{}
$Y$
\zusatzklammer {mengelilingi sumbu $x$} {} {}
yang diperoleh dari
\definitionsverweis {grafik}{}{}
$f$ memiliki sifat-sifat berikut.}
\faktfolgerung {\aufzaehlungdrei{$Y$ adalah himpunan nol dari
\maabbeledisp {h} {I \times \R \times \R} {\R
} { \left(  x   , \,   y  , \,  z                 \right) } { y^2+z^2-f(x)^2
} {.}
}{Fungsi $h$ yang mendeskripsikan permukaan itu
\definitionsverweis {reguler}{}{}
di setiap titik $Y$.
}{
\definitionsverweis {Ruang singgung}{}{}
$Y$ di suatu titik
\mathl{\left(  x   , \,   y  , \,  z                 \right)}{} adalah
\zusatzklammer {untuk

\mavergleichskettek
{\vergleichskettek
{ z
}
{ \neq }{ 0
}
{  }{
}
{    }{
}
{   }{
}
}
{}{}{}} {} {}

\mavergleichskettedisp
{\vergleichskette
{ T_PY
}
{ =} { \R \begin{pmatrix}  0  \\ z \\  -y   \end{pmatrix} + \R\begin{pmatrix}  z  \\ 0 \\  f(x)f'(x)   \end{pmatrix}
}
{ } {
}
{ } {
}
{ } {
}
}
{}{}{.}
}}
\faktzusatz {}
\faktzusatz {}

}
{

(1) jelas.
\definitionsverweis {matriks Jacobi}{}{}
dari $h$ adalah
\mathl{2 \left(  -f(x)f'(x)  , \,   y  , \,  z                 \right)}{.} Karena $f$ positif, kasus

\mavergleichskette
{\vergleichskette
{ y
}
{ = }{ z
}
{ = }{ 0
}
{    }{
}
{   }{
}
}
{}{}{}
tidak mungkin terjadi; dengan demikian, $h$ reguler pada $Y$. Jika

\mavergleichskette
{\vergleichskette
{ z
}
{ \neq }{ 0
}
{  }{
}
{    }{
}
{   }{
}
}
{}{}{}
maka vektor-vektor yang dinyatakan di atas
\definitionsverweis {bebas linear}{}{}
dan termasuk dalam
\definitionsverweis {kernel}{}{}
matriks Jacobi.

}



\inputfaktbeweis
{Rotationsfläche/Differenzierbare positive Kurve/Graph/Geodätische/Fakt}
{Lemma}
{}
{

\faktsituation {Misalkan $I$ adalah sebuah
\definitionsverweis {interval terbuka}{}{}
dan
\maabb {f} { I } { \R_+
} {}
sebuah
\definitionsverweis {fungsi terdiferensialkan}{}{}
serta misalkan $Y$ adalah
\definitionsverweis {permukaan putar}{}{}
\zusatzklammer {mengelilingi sumbu $x$} {} {}
yang diperoleh dari
\definitionsverweis {grafik}{}{}
$f$. Misalkan
\maabbeledisp {\gamma} { J } { I
} { t } { \gamma(t)
} {,}
adalah fungsi bijektif terdiferensialkan sedemikian rupa sehingga
\mathl{\left(  \gamma(t)  , \,  f(\gamma(t))                   \right)}{}
merupakan parametrisasi grafik $f$ dengan norma kecepatan konstan.}
\faktuebergang {Maka berlaku sifat-sifat berikut.}
\faktfolgerung {\aufzaehlungzwei {Untuk setiap titik

\mavergleichskette
{\vergleichskette
{ \left(  x   , \,   y  , \,  z                 \right)
}
{ = }{ \left(  x   , \,   f(x) \cos \alpha  , \,   f(x) \sin \alpha                 \right)
}
{ \in }{ Y
}
{    }{
}
{   }{
}
}
{}{}{}
tersebut, kurva
\maabbeledisp {\theta} { J } { Y
} { t } { \left(  \gamma(t)   , \,  f( \gamma(t) ) \cos \alpha  , \,   f( \gamma(t) ) \sin \alpha                 \right)
} {,}
adalah sebuah
\definitionsverweis {geodesik}{}{.}
} {Sebuah kurva lingkaran
\maabbeledisp {\theta} {\R} {I \times \R^2
} {\alpha} { \left(  x   , \,  f( x ) \cos \alpha  , \,   f(x ) \sin \alpha                 \right)
} {,}
adalah geodesik tepat ketika

\mavergleichskette
{\vergleichskette
{ f'(x)
}
{ = }{ 0
}
{  }{
}
{    }{
}
{   }{
}
}
{}{}{}
berlaku.
}}
\faktzusatz {}
\faktzusatz {}

}
{

\aufzaehlungzwei {Kurva $\theta$ bergerak di bidang
\mathl{\R  \begin{pmatrix}  1  \\ 0 \\  0   \end{pmatrix} + \R  \begin{pmatrix}  0  \\ \cos \alpha \\  \sin \alpha   \end{pmatrix}}{.} Setelah melakukan rotasi, tanpa mengurangi keumuman kita dapat mengandaikan

\mavergleichskette
{\vergleichskette
{ \alpha
}
{ = }{ 0
}
{  }{
}
{    }{
}
{   }{
}
}
{}{}{}
berlaku; dengan demikian, kurva tersebut menelusuri grafik $f$ secara langsung. Karena parametrisasi panjang busur,

\mavergleichskettedisp
{\vergleichskette
{ \theta'(t)
}
{ =} { \begin{pmatrix}  \gamma'(t) \\f'(\gamma(t)) \gamma'(t) \\  0   \end{pmatrix}
}
{ } {
}
{ } {
}
{ } {
}
}
{}{}{}
menurut
Soal 2.3
tegak lurus terhadap

\mavergleichskettedisp
{\vergleichskette
{ \theta^{\prime \prime} (t)
}
{ =} { \begin{pmatrix}  \gamma^{\prime \prime} (t) \\f'(\gamma(t)) \gamma^{\prime \prime} (t) +f^{\prime \prime}(\gamma(t))  \left(  \gamma^{\prime } (t)  \right)^2 \\  0   \end{pmatrix}
}
{ } {
}
{ } {
}
{ } {
}
}
{}{}{.}
Artinya, percepatan tegak lurus terhadap ruang singgung
\zusatzklammer {yang dibangun oleh turunan kurva bersama $\begin{pmatrix}  0  \\ 0\\  1   \end{pmatrix}$} {} {}
dan karena itu percepatan tangensial sama dengan $0$. Jadi, kurva tersebut adalah geodesik.
} {Gerak melingkar berlangsung pada bidang yang ditentukan oleh $x$ yang tetap. Turunan
\zusatzklammer {terhadap $\alpha$} {} {}
dari kurva yang diberikan adalah $f(x) \begin{pmatrix}  0  \\ - \sin \alpha \\  \cos \alpha   \end{pmatrix}$, sedangkan percepatannya sama dengan $- f(x) \begin{pmatrix}  0  \\  \cos \alpha \\  \sin \alpha   \end{pmatrix}$. Dengan demikian, di dalam bidang tersebut percepatan tegak lurus terhadap kecepatan kurva, yang merupakan sebuah vektor singgung permukaan putar. Sebuah vektor singgung yang bebas linear terhadapnya diberikan oleh $\begin{pmatrix}  1  \\f'(x)\\  0   \end{pmatrix}$. Akan tetapi, vektor ini selalu tegak lurus terhadap percepatan hanya ketika

\mavergleichskette
{\vergleichskette
{ f'(x)
}
{ = }{ 0
}
{  }{
}
{    }{
}
{   }{
}
}
{}{}{}
berlaku.
}

}

  \zwischenueberschrift{Medan Normal}

\inputdefinition
{}
{

Misalkan

\mavergleichskette
{\vergleichskette
{ W
}
{ \subseteq }{ \R^n
}
{  }{
}
{    }{
}
{   }{
}
}
{}{}{}

\definitionsverweis {terbuka}{}{,}
\maabb {h} { W } { \R
} {}
sebuah
\definitionsverweis {fungsi yang terdiferensialkan secara kontinu}{}{}
dan

\mavergleichskette
{\vergleichskette
{ Y
}
{ = }{ h^{-1}(c)
}
{  }{
}
{    }{
}
{   }{
}
}
{}{}{}

\definitionsverweis {serat}{}{}
di atas

\mavergleichskette
{\vergleichskette
{ c
}
{ \in }{ \R
}
{  }{
}
{    }{
}
{   }{
}
}
{}{}{,}
dengan $h$
\definitionsverweis {reguler}{}{}
di setiap titik $Y$. Misalkan pada suatu lingkungan terbuka

\mavergleichskette
{\vergleichskette
{ Y
}
{ \subseteq }{ U
}
{ \subseteq }{ \R^n
}
{    }{
}
{   }{
}
}
{}{}{}
didefinisikan sebuah
\definitionsverweis {medan vektor}{}{}
yang
\definitionsverweis {kontinu}{}{,}
\maabb {F} {U} {\R^n
} {}
dan memenuhi

\mavergleichskettedisp
{\vergleichskette
{ F(P)
}
{ \in} { N_PY
}
{ } {
}
{ } {
}
{ } {
}
}
{}{}{}
untuk semua

\mavergleichskette
{\vergleichskette
{ P
}
{ \in }{ Y
}
{  }{
}
{    }{
}
{   }{
}
}
{}{}{}
disebut
\definitionswort {medan normal}{}
pada $Y$.

}

Jika selain itu berlaku

\mavergleichskettedisp
{\vergleichskette
{ \Vert {F(P)} \Vert
}
{ =} { 1
}
{ } {
}
{ } {
}
{ } {
}
}
{}{}{}
maka medan tersebut disebut
\stichwort {medan normal satuan} {.} Karena garis normal dalam kasus hipermuka yang sedang ditinjau berdimensi satu, untuk setiap titik

\mavergleichskette
{\vergleichskette
{ P
}
{ \in }{ Y
}
{  }{
}
{    }{
}
{   }{
}
}
{}{}{}
hanya terdapat dua nilai yang mungkin bagi medan normal satuan di titik itu, dan keduanya saling berlawanan. Dalam hal ini, perhatian utama hanya tertuju pada nilai-nilai medan di $Y$; jadi, dua medan normal pada $Y$ dianggap identik jika keduanya berimpit di $Y$. Lingkungan terbuka hanya diperlukan agar kita dapat berbicara tentang keterdiferensialan.



\inputfaktbeweis
{Hyperfläche/Glatt/Normalenfeld/Fakt}
{Lemma}
{}
{

\faktsituation {Misalkan

\mavergleichskette
{\vergleichskette
{ W
}
{ \subseteq }{ \R^n
}
{  }{
}
{    }{
}
{   }{
}
}
{}{}{}

\definitionsverweis {terbuka}{}{,}
\maabb {h} { W } { \R
} {}
sebuah
\definitionsverweis {fungsi yang terdiferensialkan secara kontinu}{}{}
dan

\mavergleichskette
{\vergleichskette
{ Y
}
{ = }{ h^{-1}(c)
}
{  }{
}
{    }{
}
{   }{
}
}
{}{}{}

\definitionsverweis {serat}{}{}
di atas

\mavergleichskette
{\vergleichskette
{ c
}
{ \in }{ \R
}
{  }{
}
{    }{
}
{   }{
}
}
{}{}{,}
dengan $h$
\definitionsverweis {reguler}{}{}
di setiap titik $Y$.}
\faktfolgerung {Maka
\mathl{{ \frac{
\operatorname{Grad} \,  h  }{  \Vert {
\operatorname{Grad} \,  h } \Vert  } }}{} adalah sebuah
\definitionsverweis {medan normal satuan}{}{}
pada $Y$.}
\faktzusatz {}
\faktzusatz {}

}
{

Menurut asumsi,
\definitionsverweis {medan gradien}{}{}
$\operatorname{Grad} \,  h$ tidak pernah nol pada $Y$ dan, karena kontinuitas yang diasumsikan, juga tidak pernah nol pada suatu lingkungan terbuka

\mavergleichskette
{\vergleichskette
{ Y
}
{ \subseteq }{ U
}
{ \subseteq }{ \R^n
}
{    }{
}
{   }{
}
}
{}{}{.}
Jadi, medan vektor yang dinyatakan di atas terdefinisi pada suatu lingkungan terbuka yang memuat $Y$. Ortogonalitas terhadap ruang-ruang singgung pada serat dan normanya yang sama dengan satu sudah jelas.

}

\inputbeispiel{}
{

\bild{ \begin{center}
\includegraphics[width=5.5cm]{ \bildeinlesungpng {Hyperboloid1} {png} }
\end{center}
\bildtext {Sebuah \stichwort {hiperboloid satu lembar} {.}} }

\bildlizenz { Hyperboloid1.png } {} {RokerHRO} {Commons} {domain publik} {}

Kita meninjau fungsi
\maabbele {h} {\R^3} {\R
} {(x,y,z)} { x^2+y^2-z^2-1
} {}
beserta seratnya di atas $0$, yaitu \stichwort {hiperboloid satu lembar} {} $Y$.
\definitionsverweis {Gradien}{}{}
$h$ di suatu titik
\mathl{(x,y,z)}{} diberikan oleh
\mathl{\begin{pmatrix}  2x \\ 2y\\  -2z   \end{pmatrix}}{}
dan
\definitionsverweis {diferensial total}{}{}
diberikan oleh
\mathl{\left(  2x  , \,   2y  , \,   -2z                 \right)}{}; oleh karena itu, $h$
\definitionsverweis {reguler}{}{}
di setiap titik $Y$.
\definitionsverweis {Ruang singgung}{}{}
di suatu titik adalah kernel diferensial total. Salah satu basisnya ialah
\mathl{\begin{pmatrix}  y  \\ -x\\  0   \end{pmatrix} ,  \begin{pmatrix}  z  \\ 0 \\  x   \end{pmatrix}}{}
\zusatzklammer {jika

\mavergleichskettek
{\vergleichskettek
{ x
}
{ = }{ 0
}
{  }{
}
{    }{
}
{   }{
}
}
{}{}{}
berlaku, vektor kedua harus diganti dengan $\begin{pmatrix}  0  \\ z \\  y   \end{pmatrix}$} {} {.}

\definitionsverweis {Medan normal satuan}{}{}
adalah

\mavergleichskettealign
{\vergleichskettealign
{ N( \begin{pmatrix}  x  \\ y \\ z   \end{pmatrix} )
}
{ =} { { \frac{
\operatorname{Grad} \, h(x,y,z)  }{  \Vert {
\operatorname{Grad} \, h(x,y,z) } \Vert  } }
}
{ =} { \left(  x^2+y^2+z^2  \right)^{ - { \frac{   1  }{  2 } } } \begin{pmatrix}  x  \\ y \\  -z   \end{pmatrix}
}
{ } {
}
{ } {
}
}
{}
{}{.}

}

\inputdefinition
{}
{

Misalkan

\mavergleichskette
{\vergleichskette
{ W
}
{ \subseteq }{ \R^n
}
{  }{
}
{    }{
}
{   }{
}
}
{}{}{}

\definitionsverweis {terbuka}{}{,}
\maabb {h} { W } { \R
} {}
sebuah
\definitionsverweis {fungsi yang terdiferensialkan secara kontinu}{}{}
dan

\mavergleichskette
{\vergleichskette
{ Y
}
{ = }{ h^{-1}(c)
}
{  }{
}
{    }{
}
{   }{
}
}
{}{}{}

\definitionsverweis {serat}{}{}
di atas

\mavergleichskette
{\vergleichskette
{ c
}
{ \in }{ \R
}
{  }{
}
{    }{
}
{   }{
}
}
{}{}{,}
dengan $h$
\definitionsverweis {reguler}{}{}
di setiap titik $Y$. Penetapan sebuah
\definitionsverweis {medan normal satuan}{}{}
pada $Y$ disebut
\definitionswort {orientasi}{}
pada $Y$.

}

Untuk setiap orientasi selalu terdapat orientasi negatif atau orientasi yang berlawanan. Jika sebuah fungsi $h$ yang mendeskripsikan $Y$ telah ditetapkan, maka
Lemma 2.4
langsung memberikan orientasi
\mathl{{ \frac{
\operatorname{Grad} \,  h  }{  \Vert {
\operatorname{Grad} \,  h } \Vert  } }}{.} Akan tetapi, $-h$ menghasilkan orientasi yang berlawanan, dan tidak ada cara kanonis untuk memilih salah satunya. Dari medan normal yang tidak pernah nol, kita selalu dapat beralih ke medan normal satuan yang bersesuaian dan dengan demikian memperoleh suatu orientasi.



\inputfaktbeweis
{Hyperfläche/Glatt/Zusammenhängend/Orientierung/Fakt}
{Lemma}
{}
{

\faktsituation {Misalkan

\mavergleichskette
{\vergleichskette
{ W
}
{ \subseteq }{ \R^n
}
{  }{
}
{    }{
}
{   }{
}
}
{}{}{}

\definitionsverweis {terbuka}{}{,}
\maabb {h} { W } { \R
} {}
sebuah
\definitionsverweis {fungsi yang terdiferensialkan secara kontinu}{}{}
dan

\mavergleichskette
{\vergleichskette
{ Y
}
{ = }{ h^{-1}(c)
}
{  }{
}
{    }{
}
{   }{
}
}
{}{}{}

\definitionsverweis {serat}{}{}
di atas

\mavergleichskette
{\vergleichskette
{ c
}
{ \in }{ \R
}
{  }{
}
{    }{
}
{   }{
}
}
{}{}{,}
dengan $h$
\definitionsverweis {reguler}{}{}
di setiap titik $Y$. Misalkan $Y$
\definitionsverweis {terhubung}{}{.}}
\faktfolgerung {Maka terdapat tepat dua
\definitionsverweis {orientasi}{}{}
pada $Y$.}
\faktzusatz {}
\faktzusatz {}

}
{

Menurut
Lemma 2.4,
terdapat medan normal satuan yang merepresentasikan suatu orientasi, dan negatifnya menghasilkan orientasi lain. Keduanya kita sebut berturut-turut $G$ dan $-G$. Sekarang, misalkan
\maabbdisp {N} { Y } { \R^n
} {}
adalah sebarang medan normal satuan kontinu. Kita meninjau himpunan tempat medan-medan itu berimpit,

\mavergleichskettedisp
{\vergleichskette
{ Y_1
}
{ =} { \left\{ P \in Y  \mid   N(P) = G(P)         \right\}
}
{ } {
}
{ } {
}
{ } {
}
}
{}{}{}
dan

\mavergleichskettedisp
{\vergleichskette
{ Y_2
}
{ =} { \left\{ P \in Y  \mid   N(P) = - G(P)         \right\}
}
{ } {
}
{ } {
}
{ } {
}
}
{}{}{.}
Karena untuk setiap titik

\mavergleichskette
{\vergleichskette
{ P
}
{ \in }{ Y
}
{  }{
}
{    }{
}
{   }{
}
}
{}{}{}
berlaku hubungan

\mavergleichskettedisp
{\vergleichskette
{ N(P)
}
{ =} { \pm G(P)
}
{ } {
}
{ } {
}
{ } {
}
}
{}{}{}
dan hanya ada dua pilihan tanda. Dengan demikian, $Y$ adalah gabungan saling lepas dari
\mathkor {} {Y_1} {dan} {Y_2} {.}
Sebagai himpunan tempat fungsi-fungsi kontinu berimpit,
\mathkor {} {Y_1} {dan} {Y_2} {}
\definitionsverweis {tertutup}{}{}
\zusatzklammer {dan dengan demikian juga
\definitionsverweis {terbuka}{}{}
di $Y$} {} {.}
Karena keterhubungan tersebut, diperoleh

\mavergleichskette
{\vergleichskette
{ Y
}
{ = }{ Y_1
}
{  }{
}
{    }{
}
{   }{
}
}
{}{}{}
atau

\mavergleichskette
{\vergleichskette
{ Y
}
{ = }{ Y_2
}
{  }{
}
{    }{
}
{   }{
}
}
{}{}{.}

}

  \zwischenueberschrift{Pemetaan Gauss}

\inputdefinition
{}
{

Misalkan

\mavergleichskette
{\vergleichskette
{ W
}
{ \subseteq }{ \R^n
}
{  }{
}
{    }{
}
{   }{
}
}
{}{}{}

\definitionsverweis {terbuka}{}{,}
\maabb {h} { W } { \R
} {}
sebuah
\definitionsverweis {fungsi yang terdiferensialkan secara kontinu}{}{}
dan

\mavergleichskette
{\vergleichskette
{ Y
}
{ = }{ h^{-1}(c)
}
{  }{
}
{    }{
}
{   }{
}
}
{}{}{}

\definitionsverweis {serat}{}{}
di atas

\mavergleichskette
{\vergleichskette
{ c
}
{ \in }{ \R
}
{  }{
}
{    }{
}
{   }{
}
}
{}{}{,}
dengan $h$
\definitionsverweis {reguler}{}{}
di setiap titik $Y$. Misalkan suatu
\definitionsverweis {orientasi}{}{}
pada $Y$ telah ditetapkan. Pemetaan
\maabbdisp {} { Y } { S^{n-1}
} {,}
yang memetakan setiap titik

\mavergleichskette
{\vergleichskette
{ P
}
{ \in }{ Y
}
{  }{
}
{    }{
}
{   }{
}
}
{}{}{}
ke vektor normal satuannya yang ditetapkan oleh orientasi disebut
\definitionswort {pemetaan Gauss}{}
dari $Y$.

}

Pada umumnya, kita membayangkan vektor singgung dan vektor normal melekat pada titik

\mavergleichskette
{\vergleichskette
{ P
}
{ \in }{ Y
}
{  }{
}
{    }{
}
{   }{
}
}
{}{}{}
tersebut. Namun, di sini vektor normal perlu dipandang sebagai sebuah titik pada permukaan bola berdimensi $n-1$ yang \anfuehrung{netral}{}. Pemetaan Gauss membuka kemungkinan untuk mengaitkan sebarang hipermuka mulus dengan hipermuka yang sangat sederhana, yaitu permukaan bola.



\inputfaktbeweis
{Hyperfläche/Glatt/Gauß-Abbildung/Eigenschaften/Fakt}
{Lemma}
{}
{

\faktsituation {Misalkan

\mavergleichskette
{\vergleichskette
{ W
}
{ \subseteq }{ \R^n
}
{  }{
}
{    }{
}
{   }{
}
}
{}{}{}

\definitionsverweis {terbuka}{}{,}
\maabb {h} { W } { \R
} {}
sebuah
\definitionsverweis {fungsi yang terdiferensialkan secara kontinu}{}{}
dan

\mavergleichskette
{\vergleichskette
{ Y
}
{ = }{ h^{-1}(c)
}
{  }{
}
{    }{
}
{   }{
}
}
{}{}{}

\definitionsverweis {serat}{}{}
di atas

\mavergleichskette
{\vergleichskette
{ c
}
{ \in }{ \R
}
{  }{
}
{    }{
}
{   }{
}
}
{}{}{,}
dengan $h$
\definitionsverweis {reguler}{}{}
di setiap titik $Y$. Misalkan suatu
\definitionsverweis {orientasi}{}{}
pada $Y$ telah ditetapkan.}
\faktuebergang {Maka, untuk
\definitionsverweis {pemetaan Gauss}{}{,}
berlaku pernyataan-pernyataan berikut.}
\faktfolgerung {\aufzaehlungvier{Pemetaan Gauss
\definitionsverweis {kontinu}{}{.}
}{Dalam kasus

\mavergleichskette
{\vergleichskette
{ Y
}
{ = }{ S^{n-1}
}
{  }{
}
{    }{
}
{   }{
}
}
{}{}{}
ini, pemetaan Gauss adalah identitas
\zusatzklammer {jika orientasinya mengarah ke luar} {} {}
atau pemetaan antipodal untuk orientasi sebaliknya.
}{Jika $Y$
\definitionsverweis {terhubung}{}{,}
maka pemetaan Gauss untuk kedua orientasi saling antipodal.
}{Misalkan $Y$ terhubung. Pemetaan Gauss konstan tepat ketika $Y$ merupakan suatu bagian terbuka dari sebuah
\definitionsverweis {subruang afin-linear}{}{}
berdimensi $n-1$.
}}
\faktzusatz {}
\faktzusatz {}

}
{

\aufzaehlungvier{Hal ini mengikuti dari
Lemma 2.4.
}{Jelas.
}{Jelas.
}{Arah sebaliknya jelas. Untuk membuktikan arah maju, misalkan vektor normal satuannya konstan dan sama dengan

\mavergleichskettedisp
{\vergleichskette
{ v
}
{ =} { \begin{pmatrix} a_1  \\\vdots\\ a_n   \end{pmatrix}
}
{ } {
}
{ } {
}
{ } {
}
}
{}{}{.}

\textit{Catatan edisi: sumber mencoba menormalkan fungsi bernilai real $h$
dengan membaginya oleh normanya pada serat nol. Langkah itu tidak terdefinisi
dan tidak membuktikan bahwa $h$ linear. Bukti berikut mengganti langkah tersebut.}

Definisikan fungsi linear

\[
\ell(x)=\langle v,x\rangle .
\]

Untuk setiap $P\in Y$ dan $w\in T_PY$, berlaku

\[
(d\ell)_P(w)=\langle v,w\rangle=0,
\]

karena $v$ merupakan vektor normal pada $Y$. Dengan demikian, diferensial
dari pembatasan $\ell$ pada $Y$ bernilai nol. Manifold mulus $Y$ terhubung dan
terhubung-lintasan secara lokal, sehingga $\ell$ konstan pada $Y$, katakanlah
bernilai $c\in\R$. Oleh sebab itu,

\[
Y\subseteq H=\{x\in\R^n\mid \langle v,x\rangle=c\}.
\]

Untuk setiap $P\in Y$, kedua ruang singgung adalah
$T_PY=T_PH=v^\perp$. Karena $Y$ dan $H$ berdimensi sama, inklusi
$Y\longrightarrow H$ merupakan difeomorfisme lokal. Jadi $Y$ terbuka dalam
hiperbidang afin $H$.}

}

Kita ingat kembali teorema berikut.

\inputfakt{Hyperfläche/Glatt und kompakt/Jede Hyperebene als Tangentialraum/Fakt}{Teorema}{}
{

\faktsituation {Misalkan

\mavergleichskette
{\vergleichskette
{ U
}
{ \subseteq }{\R^n
}
{  }{
}
{    }{
}
{   }{
}
}
{}{}{}
sebuah
\definitionsverweis {himpunan bagian terbuka}{}{}
dan misalkan
\maabbdisp {f} { U } {\R
} {}
sebuah
\definitionsverweis {fungsi yang terdiferensialkan secara kontinu}{}{.}}
\faktvoraussetzung {
\definitionsverweis {Serat}{}{}

\mavergleichskette
{\vergleichskette
{M
}
{ = }{ f^{-1}(b)
}
{  }{
}
{    }{
}
{   }{
}
}
{}{}{}
dari $f$ di atas suatu titik

\mavergleichskette
{\vergleichskette
{b
}
{ \in }{\R
}
{  }{
}
{    }{
}
{   }{
}
}
{}{}{}
bersifat
\definitionsverweis {kompak}{}{}
dan
\definitionsverweis {reguler}{}{}
di setiap titik.}
\faktfolgerung {Maka, untuk setiap subruang
\mathl{(n-1)}{-}dimensional

\mavergleichskette
{\vergleichskette
{V
}
{ \subset }{ \R^n
}
{  }{
}
{    }{
}
{   }{
}
}
{}{}{}
terdapat sekurang-kurangnya satu titik

\mavergleichskette
{\vergleichskette
{a
}
{ \in }{ M
}
{  }{
}
{    }{
}
{   }{
}
}
{}{}{}
sehingga subruang tersebut berimpit dengan
\definitionsverweis {ruang singgung}{}{}

\mathl{T_aM}{.}}
\faktzusatz {}
\faktzusatz {}

}


\inputfaktbeweis
{Hyperfläche/Glatt und kompakt/Gauß-Abbildung surjektiv/Fakt}
{Teorema}
{}
{

\faktsituation {Misalkan

\mavergleichskette
{\vergleichskette
{ W
}
{ \subseteq }{\R^n
}
{  }{
}
{    }{
}
{   }{
}
}
{}{}{}
sebuah
\definitionsverweis {himpunan bagian terbuka}{}{}
dan misalkan
\maabbdisp {h} { W } {\R
} {}
sebuah
\definitionsverweis {fungsi yang terdiferensialkan secara kontinu}{}{.}}
\faktvoraussetzung {
\definitionsverweis {Serat}{}{}

\mavergleichskette
{\vergleichskette
{ Y
}
{ = }{ h^{-1}(c)
}
{  }{
}
{    }{
}
{   }{
}
}
{}{}{}
dari $h$ di atas

\mavergleichskette
{\vergleichskette
{ c
}
{ \in }{\R
}
{  }{
}
{    }{
}
{   }{
}
}
{}{}{}
bersifat
\definitionsverweis {kompak}{}{}
dan
\definitionsverweis {reguler}{}{}
di setiap titik.}
\faktfolgerung {Maka
\definitionsverweis {pemetaan Gauss}{}{}
\maabbdisp {} { Y } { S^{n-1}
} {}
surjektif.}
\faktzusatz {}
\faktzusatz {}

}
{

Melalui relasi ortogonalitas, sebuah vektor normal satuan mendeskripsikan sebuah hiperbidang, yaitu subruang vektor berdimensi $n-1$; vektor normal satuan negatifnya juga mendeskripsikan hiperbidang yang sama. Teorema 2.10
menunjukkan bahwa setiap hiperbidang muncul sebagai ruang singgung pada $Y$. Selain itu, bukti
Teorema 2.10
menunjukkan
\zusatzklammer {jika selain maksimum kita juga meninjau minimum} {} {,}
bahwa kedua vektor normal satuan itu muncul.

}

\chapter{Lembar Kerja 2}
\setcounter{section}{2}

  \zwischenueberschrift{Soal latihan}

\inputaufgabegibtloesung
{}
{

Misalkan
\maabbdisp {f,g} {\R} { \R^n
} {}
\definitionsverweis {kurva-kurva terdiferensialkan}{}{.}
Hitung
\definitionsverweis {turunan}{}{}
fungsi
\maabbeledisp {} {\R} {\R
} { t } { \left\langle f(t) , g(t)  \right\rangle
} {.}
Nyatakan hasilnya menggunakan hasil kali dalam.

}
{} {}

\inputaufgabegibtloesung
{}
{

Misalkan
\maabb {h} { I } { V
} {}
sebuah kurva yang terdiferensialkan dua kali dalam
\definitionsverweis {ruang vektor Euklides}{}{}
$V$. Buktikan bahwa jika

\mavergleichskette
{\vergleichskette
{h'(t)
}
{ \neq }{ 0
}
{  }{
}
{    }{
}
{   }{
}
}
{}{}{,}
maka berlaku persamaan

\mavergleichskettedisp
{\vergleichskette
{ \Vert { h'(t)} \Vert'
}
{ =} { { \frac{   \left\langle h'(t) , h^{\prime \prime} (t)  \right\rangle  }{  \left\langle h'(t) , h'(t)  \right\rangle  } } \Vert { h'(t)} \Vert
}
{ } {
}
{ } {
}
{ } {
}
}
{}{}{.}

}
{} {}

Sebuah
\definitionsverweis {kurva terdiferensialkan}{}{}
\maabbeledisp {f} { [a,b] } {\R^n
} { t } {f(t) = \left( f_1(t)  , \,   \ldots  , \,  f_n(t)                 \right)
} {,}
disebut
\definitionswort {berparametrisasi panjang busur}{,}
jika

\mavergleichskette
{\vergleichskette
{ f_1'(t)^2 + \cdots +  f_n'(t)^2
}
{ = }{ 1
}
{  }{
}
{    }{
}
{   }{
}
}
{}{}{}
untuk setiap $t$.

\inputaufgabe
{}
{

Misalkan
\maabb {\gamma} { I } {\R^n
} {}
sebuah kurva yang dua kali
\definitionsverweis {terdiferensialkan secara kontinu}{}{}
dan
\definitionsverweis {berparametrisasi panjang busur}{}{.}
Buktikan bahwa
\mathkor {} {\gamma'(t)} {dan} {\gamma^{\prime \prime}(t)} {}
saling tegak lurus.

}
{} {}

\inputaufgabe
{}
{

Misalkan

\mavergleichskette
{\vergleichskette
{ P
}
{ \in }{ \R^n
}
{  }{
}
{    }{
}
{   }{
}
}
{}{}{}
sebuah titik dan misalkan

\mavergleichskette
{\vergleichskette
{ I
}
{ = }{ {]{-1},1[}
}
{  }{
}
{    }{
}
{   }{
}
}
{}{}{.}
Kita tinjau himpunan

\mavergleichskettedisp
{\vergleichskette
{ M
}
{ =} { \left\{ f:I \rightarrow \R^n   \mid  f \text{ terdiferensialkan} , \,  f(0) = P       \right\}
}
{ } {
}
{ } {
}
{ } {
}
}
{}{}{.}
Dua
\definitionsverweis {kurva}{}{}

\mavergleichskette
{\vergleichskette
{ f,g
}
{ \in }{ M
}
{  }{
}
{    }{
}
{   }{
}
}
{}{}{}
disebut \stichwort {ekuivalen secara tangensial} {,} jika

\mavergleichskettedisp
{\vergleichskette
{ f'(0)
}
{ =} { g'(0)
}
{ } {
}
{ } {
}
{ } {
}
}
{}{}{.}

a) Buktikan bahwa ini merupakan suatu
\definitionsverweis {relasi ekuivalensi}{}{.}

b) Temukan wakil paling sederhana bagi setiap
\definitionsverweis {kelas ekuivalensi}{}{.}

c) Berikan satu wakil lain untuk setiap kelas.

d) Deskripsikan himpunan kelas-kelas ekuivalensi tersebut
\zusatzklammer {yakni
\definitionsverweis {himpunan hasil bagi}{}{}} {} {.}

}
{} {}

\inputaufgabe
{}
{

Dengan syarat
\zusatzklammer {reguler di setiap titik} {} {,}
berikan sebuah contoh
\definitionsverweis {hipermuka terdiferensialkan}{}{,}
yang tidak
\definitionsverweis {terhubung}{}{.}

}
{} {}

\inputaufgabe
{}
{

Misalkan

\mavergleichskette
{\vergleichskette
{ V
}
{ \subseteq }{ \R^{n-1}
}
{  }{
}
{    }{
}
{   }{
}
}
{}{}{}
sebuah
\definitionsverweis {himpunan terbuka}{}{}
dan misalkan
\maabbdisp {f} { V } {\R
} {}
sebuah fungsi yang
\definitionsverweis {terdiferensialkan secara kontinu}{}{.}
Kita tinjau
\definitionsverweis {grafik}{}{}
$Y$ dari $f$ sebagai hipermuka terdiferensialkan di $\R^n$ dalam pengertian
Contoh 1.2.
Tentukan
\definitionsverweis {medan-medan normal satuan}{}{}
dalam kasus ini.

}
{} {}

\inputaufgabegibtloesung
{}
{

Misalkan
\mathkor {} {r} {dan} {R} {}
bilangan real positif yang memenuhi

\mavergleichskette
{\vergleichskette
{ 0
}
{ < }{ r
}
{ < }{ R
}
{    }{
}
{   }{
}
}
{}{}{.}
Tentukan suatu
\definitionsverweis {medan normal satuan}{}{}
untuk torus
\zusatzklammer {tertanam} {} {}

\mavergleichskettedisp
{\vergleichskette
{ T
}
{ =} { \left\{ (x,y,z) \in \R^3  \mid   \left(  \sqrt{x^2+y^2} -R  \right)^2 +z^2  = r^2         \right\}
}
{ } {
}
{ } {
}
{ } {
}
}
{}{}{.}

}
{} {}

\inputaufgabe
{}
{

Misalkan
\maabbdisp {f} { I } {\R_+
} {}
sebuah
\definitionsverweis {fungsi terdiferensialkan}{}{}
dan misalkan $Y$ adalah permukaan putar yang terkait dengan grafik $f$, dipandang sebagai hipermuka terdiferensialkan di $\R^3$ dalam pengertian
Lemma 2.1.
Tentukan
\definitionsverweis {medan-medan normal satuan}{}{}
dalam kasus ini.

}
{} {}

\inputaufgabe
{}
{

Buktikan bahwa
\definitionsverweis {pemetaan Gauss}{}{}
pada sebuah hiperbidang bersifat konstan. Apakah pernyataan sebaliknya juga berlaku?

}
{} {}

\inputaufgabe
{}
{

Buktikan bahwa
\definitionsverweis {pemetaan Gauss}{}{}
pada sebuah permukaan bola di $\R^n$ yang berpusat di titik asal, untuk setiap pilihan orientasi, diinduksi oleh homoteti skalar pada $\R^n$ (dengan faktor positif untuk orientasi keluar dan faktor negatif untuk orientasi masuk).

}
{} {}

\inputaufgabe
{}
{

Misalkan

\mavergleichskette
{\vergleichskette
{ Y
}
{ \subseteq }{ \R^n
}
{  }{
}
{    }{
}
{   }{
}
}
{}{}{}
sebuah
\definitionsverweis {hipermuka terdiferensialkan}{}{}
dengan
\definitionsverweis {medan normal satuan}{}{}
$N$ dan medan normal satuan yang berlawanan, $-N$. Buktikan bahwa kedua
\definitionsverweis {pemetaan Gauss}{}{}
yang terkait saling berhubungan melalui
\definitionsverweis {pemetaan antipodal}{}{}
\maabbeledisp {} {S^{n-1}} {S^{n-1}
} { P } { -P
} {.}

}
{} {}

\inputaufgabegibtloesung
{}
{

Misalkan

\mavergleichskette
{\vergleichskette
{ \alpha,\beta
}
{ \in }{ \R_+
}
{  }{
}
{    }{
}
{   }{
}
}
{}{}{}
dan misalkan

\mavergleichskettedisp
{\vergleichskette
{ Y
}
{ =} { \left\{ (x,y)\in\R^2  \mid   \alpha x^2+\beta y^2 = 1        \right\}
}
{ } {
}
{ } {
}
{ } {
}
}
{}{}{.}
Orientasikan $Y$ dengan medan normal satuan yang diperoleh dari gradien keluar fungsi
$ (x,y) \longmapsto \alpha x^2+\beta y^2 $.
\aufzaehlungzwei {Tentukan
\definitionsverweis {pemetaan Gauss}{}{}
untuk $Y$.
} { Berikan secara eksplisit suatu
\definitionsverweis {pemetaan invers}{}{}
untuk pemetaan Gauss pada $Y$.
}

}
{} {}

\inputaufgabegibtloesung
{}
{

Misalkan

\mavergleichskette
{\vergleichskette
{ \alpha,\beta
}
{ \in }{ \R_+
}
{  }{
}
{    }{
}
{   }{
}
}
{}{}{}
dan misalkan

\mavergleichskettedisp
{\vergleichskette
{ Y
}
{ =} { \left\{ (x,y)\in\R^2  \mid   \alpha x^2+\beta y^2 = 1        \right\}
}
{ } {
}
{ } {
}
{ } {
}
}
{}{}{.}
Buktikan bahwa
\definitionsverweis {pemetaan Gauss}{}{}
untuk $Y$ bersifat bijektif.

}
{} {}

\inputaufgabe
{}
{

Berikan sebuah contoh
\definitionsverweis {hipermuka terdiferensialkan}{}{}

\mavergleichskette
{\vergleichskette
{ Y
}
{ \subseteq }{ \R^2
}
{  }{
}
{    }{
}
{   }{
}
}
{}{}{}
sedemikian sehingga
\definitionsverweis {pemetaan Gauss}{}{}
bersifat surjektif dan terdapat titik-titik pada $S^1$ yang dicapai tak berhingga kali.

}
{} {}

\inputaufgabe
{}
{

Misalkan

\mavergleichskette
{\vergleichskette
{ W
}
{ \subseteq }{ \R^n
}
{  }{
}
{    }{
}
{   }{
}
}
{}{}{}
\definitionsverweis {terbuka}{}{,}
\maabb {h} { W } { \R
} {}
sebuah
\definitionsverweis {fungsi yang terdiferensialkan secara kontinu}{}{}
dan

\mavergleichskette
{\vergleichskette
{ Y
}
{ = }{ h^{-1}(c)
}
{  }{
}
{    }{
}
{   }{
}
}
{}{}{}
adalah
\definitionsverweis {serat}{}{}
di atas

\mavergleichskette
{\vergleichskette
{ c
}
{ \in }{ \R
}
{  }{
}
{    }{
}
{   }{
}
}
{}{}{,}
dengan $h$
\definitionsverweis {reguler}{}{}
di setiap titik pada $Y$. Misalkan
\maabbdisp {L} {\R^n } { \R^n
} {}
sebuah
\definitionsverweis {isometri linear}{}{,}

\mavergleichskette
{\vergleichskette
{ W'
}
{ = }{ L^{-1}(W)
}
{  }{
}
{    }{
}
{   }{
}
}
{}{}{}
dan

\mavergleichskettedisp
{\vergleichskette
{ Z
}
{ =} { L^{-1}(Y)
}
{ =} { (h \circ L)^{-1}(c)
}
{ } {
}
{ } {
}
}
{}{}{.}
Buktikan bahwa konsep-konsep kurva geodesik, medan normal, medan normal satuan, dan pemetaan Gauss pada $Y$ dan $Z$ saling bersesuaian
\zusatzklammer {yakni dapat dipetakan satu sama lain dengan bantuan $L$} {} {.}

}
{} {}

  \zwischenueberschrift{Soal untuk dikumpulkan}

\inputaufgabe
{3}
{

Tentukan citra parabola standar berorientasi ke atas di bawah
\definitionsverweis {pemetaan Gauss}{}{.}

}
{} {}

\inputaufgabe
{4}
{

Misalkan

\mavergleichskette
{\vergleichskette
{ W
}
{ = }{ \R^3 \setminus \{ (0,0,0) \}
}
{  }{
}
{    }{
}
{   }{
}
}
{}{}{}
dan misalkan

\mavergleichskette
{\vergleichskette
{ Y
}
{ \subseteq }{ W
}
{  }{
}
{    }{
}
{   }{
}
}
{}{}{}
adalah
\definitionsverweis {himpunan nol}{}{}
dari

\mavergleichskettedisp
{\vergleichskette
{ h(x,y,z)
}
{ =} { x^2+y^2-z^2
}
{ } {
}
{ } {
}
{ } {
}
}
{}{}{.}
Tentukan citra $Y$
\zusatzklammer {dengan orientasi yang ditentukan oleh gradien $h$} {} {}
di bawah
\definitionsverweis {pemetaan Gauss}{}{.}

}
{} {}

\inputaufgabe
{5}
{

Misalkan

\mavergleichskette
{\vergleichskette
{ \alpha,\beta, \gamma
}
{ \in }{ \R_+
}
{  }{
}
{    }{
}
{   }{
}
}
{}{}{}
dan misalkan

\mavergleichskettedisp
{\vergleichskette
{ Y
}
{ =} { \left\{ (x,y,z) \in \R^3  \mid   \alpha x^2+ \beta y^2 + \gamma z^2 = 1        \right\}
}
{ } {
}
{ } {
}
{ } {
}
}
{}{}{.}
Buktikan bahwa
\definitionsverweis {pemetaan Gauss}{}{}
untuk $Y$ bersifat bijektif.

}
{} {}

\inputaufgabe
{3}
{

Misalkan

\mavergleichskette
{\vergleichskette
{ W
}
{ \subseteq }{ \R^n
}
{  }{
}
{    }{
}
{   }{
}
}
{}{}{}
\definitionsverweis {terbuka}{}{,}
\maabb {h} { W } { \R
} {}
sebuah
\definitionsverweis {fungsi yang terdiferensialkan secara kontinu}{}{}
dan

\mavergleichskette
{\vergleichskette
{ Y
}
{ = }{ h^{-1}(c)
}
{  }{
}
{    }{
}
{   }{
}
}
{}{}{}
adalah
\definitionsverweis {serat}{}{}
di atas

\mavergleichskette
{\vergleichskette
{ c
}
{ \in }{ \R
}
{  }{
}
{    }{
}
{   }{
}
}
{}{}{,}
dengan $h$
\definitionsverweis {reguler}{}{}
di setiap titik pada $Y$. Misalkan
\maabbdisp {L} {\R^n } { \R^n
} {}
sebuah
\definitionsverweis {isomorfisme linear}{}{,}

\mavergleichskette
{\vergleichskette
{ W'
}
{ = }{ L^{-1}(W)
}
{  }{
}
{    }{
}
{   }{
}
}
{}{}{}
dan

\mavergleichskettedisp
{\vergleichskette
{ Z
}
{ =} { L^{-1}(Y)
}
{ =} { (h \circ L)^{-1}(c)
}
{ } {
}
{ } {
}
}
{}{}{.}
Buktikan bahwa konsep-konsep kurva geodesik, medan normal, medan normal satuan, dan pemetaan Gauss pada $Y$ dan $Z$ secara umum tidak saling bersesuaian
\zusatzklammer {berbeda dengan
Soal 2.14,
yang melibatkan suatu isometri} {} {.}

}
{} {}

\section*{Solusi yang disediakan oleh sumber}
\addcontentsline{toc}{section}{Solusi yang disediakan oleh sumber}
\subsection*{Solusi Soal 2.1}
% Authority: German Wikiversity pageid 139491, revid 1090802, 2026-05-31T14:13:03Z.
% Exact UTF-8 wikitext witness SHA-256: f6787c5b9547ac189e011cb8b1f50ecad82247095e086138bd45e0a4f06647bb.
% Expanded German TeX source SHA-256: 6dcc38f066a8350fdba67145857c168e1b6ca532c07af0a8f34ee2b954ad9432.

Misalkan $f_i(t)$ dan $g_i(t)$ berturut-turut adalah fungsi-fungsi komponen dari
\mathkor {} {f} {dan} {g} {.}
Fungsi yang dimaksud adalah

\mavergleichskettedisp
{\vergleichskette
{  \left\langle f(t) , g(t)  \right\rangle
}
{ =} { \sum_{i = 1}^n  f_i(t) g_i(t)
}
{ } {
}
{ } {
}
{ } {
}
}
{}{}{}
dengan turunan
\zusatzklammer {aturan hasil kali diterapkan pada setiap suku} {} {}

\mathdisp {\sum_{i = 1}^n  f'_i(t) g_i(t)  + \sum_{i = 1}^n  f_i(t) g'_i(t)} { . }

Dengan hasil kali dalam, hasil ini dapat ditulis sebagai

\mathdisp {\left\langle f'(t) ,  g(t)  \right\rangle + \left\langle f(t) ,  g'(t)  \right\rangle} { . }
\subsection*{Solusi Soal 2.2}
% Authority: German Wikiversity pageid 120435, revid 1089268, 2026-05-31T10:02:12Z.
% Exact UTF-8 wikitext witness SHA-256: 0bdd1cde6d1d0dc74757aa824634d10b6b4bf413efb4431538b7d9bac2b6f06c.
% Expanded German TeX source SHA-256: a92e556f50d2216192fc61703eea4bae3d233ab632c8eedb19bd35dec2ed89b0.

\textit{Catatan edisi: solusi sumber memakai basis ortogonal seolah-olah
ortonormal dan menuliskan turunan koordinat yang keliru. Argumen berikut
mengganti langkah tersebut.}

Karena $h'(t)\neq 0$, norma kecepatannya tak nol. Turunkan kuadrat normanya:

\[
\frac{d}{dt}\Vert h'(t)\Vert^2
=2\Vert h'(t)\Vert\bigl(\Vert h'(t)\Vert\bigr)'
=2\left\langle h'(t),h''(t)\right\rangle .
\]

Membagi dengan $2\Vert h'(t)\Vert$ memberikan

\[
\bigl(\Vert h'(t)\Vert\bigr)'
=\frac{\left\langle h'(t),h''(t)\right\rangle}{\Vert h'(t)\Vert}
=\frac{\left\langle h'(t),h''(t)\right\rangle}
       {\left\langle h'(t),h'(t)\right\rangle}\Vert h'(t)\Vert .
\]
\subsection*{Solusi Soal 2.7}
\begingroup\scriptsize
% Authority: German Wikiversity pageid 152882, revid 1094222, 2026-06-14T15:20:09Z.
% Exact UTF-8 wikitext witness SHA-256: dbe0990deca042c24dda79891b6ccc66b57723be69023687cacd7ffa206f6988.
% Expanded German TeX source SHA-256: 90d9007e6a313dcbc4614045f0831b826f7a9af82897e17a36fbce08088ef92b.

Pada himpunan terbuka $U=\{(x,y,z)\in\R^3\mid x^2+y^2>0\}$, berlaku

\mavergleichskettedisp
{\vergleichskette
{ h(x,y,z)
}
{ =} { x^2+y^2 -2 \sqrt{x^2+y^2}R+z^2 + R^2
}
{ } {
}
{ } {
}
{ } {
}
}
{}{}{}
dan

\mavergleichskette
{\vergleichskette
{ T
}
{ = }{ h^{-1} (r^2)
}
{  }{
}
{    }{
}
{   }{
}
}
{}{}{.}
Gradien $h$ adalah

\mavergleichskettedisp
{\vergleichskette
{
\operatorname{Grad} \,  h
}
{ =} { 2 \begin{pmatrix}   x   -  xR \left(  x^2+y^2  \right)^{- { \frac{   1  }{  2 } } }  \\ y   -  y R \left(  x^2+y^2  \right)^{- { \frac{   1  }{  2 } } } \\ z   \end{pmatrix}
}
{ } {
}
{ } {
}
{ } {
}
}
{}{}{.}
Oleh karena itu,
\definitionsverweis {medan normal satuan}{}{}
tersebut adalah

\mavergleichskettedisp
{\vergleichskette
{ N(x,y,z)
}
{ =} { { \frac{   1  }{  \sqrt{ \left(   x   -  xR \left(  x^2+y^2  \right)^{- { \frac{   1  }{  2 } } }  \right)^2 + \left(  y   -  y R \left(  x^2+y^2  \right)^{- { \frac{   1  }{  2 } } }  \right)^2   +z^2   }  } } \begin{pmatrix}   x   -  xR \left(  x^2+y^2  \right)^{- { \frac{   1  }{  2 } } }  \\ y   -  y R \left(  x^2+y^2  \right)^{- { \frac{   1  }{  2 } } } \\ z   \end{pmatrix}
}
{ } {
}
{ } {
}
{ } {
}
}
{}{}{.}
\endgroup
\subsection*{Solusi Soal 2.12}
\begingroup\small
% Authority: German Wikiversity pageid 152865, revid 1095926, 2026-06-15T07:27:34Z.
% Exact UTF-8 wikitext witness SHA-256: 78284fe9d10f41f0503cba0d3295c761d5d168d8de06505226909c7242d4993a.
% Expanded German TeX source SHA-256: 3191adb6fbfaf1be11c0e7a061468a4fd6fe9235cd0dea40771ec97bd10f970f.

\aufzaehlungzwei {Untuk orientasi keluar yang ditentukan dalam soal,
\definitionsverweis {medan normal satuan}{}{}
diperoleh dengan menormalisasi gradien fungsi yang mendeskripsikan elips tersebut, yaitu

\mavergleichskettedisp
{\vergleichskette
{ G(x,y)
}
{ =} { \begin{pmatrix}  2 \alpha x  \\ 2 \beta y    \end{pmatrix}
}
{ } {
}
{ } {
}
{ } {
}
}
{}{}{}
dan dengan menormalisasikannya kita memperoleh

\mavergleichskettedisp
{\vergleichskette
{ N(x,y)
}
{ =} {  { \frac{   1 }{  \sqrt{ \alpha^2 x^2 + \beta^2 y^2 }  } } \begin{pmatrix}   \alpha x  \\  \beta y    \end{pmatrix}
}
{ } {
}
{ } {
}
{ } {
}
}
{}{}{.}
Dengan demikian, pemetaan Gauss tersebut adalah
\maabbeledisp {\varphi} {Y = \left\{  \begin{pmatrix}  x  \\ y   \end{pmatrix}   \in \R^2  \mid   \alpha x^2+\beta y^2 = 1        \right\}  } { S^1
} { \begin{pmatrix}  x  \\ y   \end{pmatrix} } { { \frac{   1 }{  \sqrt{ \alpha^2 x^2 + \beta^2 y^2 }  } } \begin{pmatrix}   \alpha x  \\  \beta y    \end{pmatrix}
} {.}
} {Kita akan menunjukkan bahwa
\maabbeledisp {\psi} {S^1} {Y
} { \begin{pmatrix}  u  \\v   \end{pmatrix} } { { \frac{   1 }{  \sqrt{ { \frac{  u^2 }{ \alpha } }   + { \frac{  v^2 }{ \beta } } }  } }       \begin{pmatrix}   { \frac{   u  }{ \alpha } }   \\ { \frac{   v  }{ \beta } }    \end{pmatrix}
} {,}
merupakan pemetaan invers. Karena

\mavergleichskettedisp
{\vergleichskette
{ { \frac{   1 }{   { \frac{  u^2 }{ \alpha } }   + { \frac{  v^2 }{ \beta } }  } } \left(  \alpha { \frac{  u^2 }{ \alpha^2 } } + \beta{ \frac{  v^2 }{ \beta^2 } }  \right)
}
{ =} { 1
}
{ } {
}
{ } {
}
{ } {
}
}
{}{}{}
citra $\psi$ terletak di $Y$. Selain itu,

\mavergleichskettealign
{\vergleichskettealign
{  \varphi \left(  \psi \begin{pmatrix}  u  \\v   \end{pmatrix}  \right)
}
{ =} { \varphi \left(  { \frac{   1 }{  \sqrt{ { \frac{  u^2 }{ \alpha } }   + { \frac{  v^2 }{ \beta } } }  } }       \begin{pmatrix}   { \frac{   u  }{ \alpha } }   \\ { \frac{   v  }{ \beta } }    \end{pmatrix}    \right)
}
{ =} { \varphi \begin{pmatrix}   { \frac{   1 }{  \alpha    \sqrt{ { \frac{  u^2 }{ \alpha } }   + { \frac{  v^2 }{ \beta } } }  } }   u    \\ { \frac{   1 }{ \beta  \sqrt{ { \frac{  u^2 }{ \alpha } }  + { \frac{  v^2 }{ \beta } } }  } }  v       \end{pmatrix}
}
{ =} { { \frac{   1 }{  \sqrt{  \alpha^2 { \frac{  u^2 }{  \alpha^2   \left(   { \frac{  u^2 }{ \alpha } }  +  { \frac{  v^2 }{ \beta } }   \right)  } }  + \beta^2 { \frac{  v^2 }{  \beta^2 \left(   { \frac{  u^2 }{ \alpha } }  +  { \frac{  v^2 }{ \beta } }   \right)  } }       }  } }  \begin{pmatrix}  { \frac{   1 }{  \sqrt{ { \frac{  u^2 }{ \alpha } }   + { \frac{  v^2 }{ \beta } } }  } }   u    \\   { \frac{   1 }{   \sqrt{ { \frac{  u^2 }{ \alpha } } + { \frac{  v^2 }{ \beta } } }  } } v    \end{pmatrix}
}
{ =} { { \frac{   1 }{  \sqrt{ { \frac{  u^2 }{  \left(   { \frac{  u^2 }{ \alpha } }  +  { \frac{  v^2 }{ \beta } }   \right)  } }   + { \frac{  v^2 }{  \left(   { \frac{  u^2 }{ \alpha } }  +  { \frac{  v^2 }{ \beta } }   \right)  } }       }  } }  \begin{pmatrix}   { \frac{   1 }{  \sqrt{ { \frac{  u^2 }{ \alpha } }   + { \frac{  v^2 }{ \beta } } }  } }   u    \\   { \frac{   1 }{   \sqrt{ { \frac{  u^2 }{ \alpha } }   + { \frac{  v^2 }{ \beta } } }  } } v    \end{pmatrix}
}
}
{
\vergleichskettefortsetzungalign
{ =} { { \frac{    \sqrt{ { \frac{  u^2 }{ \alpha } }  +  { \frac{  v^2 }{ \beta } } } }{  \sqrt{    u^2+ v^2       }  } }  \begin{pmatrix}   { \frac{   1 }{  \sqrt{ { \frac{  u^2 }{ \alpha } }   + { \frac{  v^2 }{ \beta } } }  } }   u    \\   { \frac{   1 }{   \sqrt{ { \frac{  u^2 }{ \alpha } }   + { \frac{  v^2 }{ \beta } } }  } } v    \end{pmatrix}
}
{ =} { \begin{pmatrix}  u  \\v   \end{pmatrix}
}
{ } {
}
{ } {}
}
{}{}
dan

\mavergleichskettealign
{\vergleichskettealign
{  \psi \left(  \varphi \begin{pmatrix}  x  \\ y   \end{pmatrix}  \right)
}
{ =} { \psi \begin{pmatrix}  { \frac{    \alpha x }{  \sqrt{\alpha^2 x^2+ \beta^2 y^2 }  } }   \\ { \frac{   \beta y }{  \sqrt{\alpha^2 x^2+ \beta^2 y^2  } }}    \end{pmatrix}
}
{ =} { { \frac{   1 }{  \sqrt{ { \frac{   \alpha x^2 }{  \alpha^2x^2 + \beta^2y^2  } }  + { \frac{   \beta y^2 }{  \alpha^2x^2 + \beta^2y^2  } }  }  } }  \begin{pmatrix}  { \frac{    x }{  \sqrt{\alpha^2 x^2+ \beta^2 y^2 }  } }   \\ { \frac{   y }{  \sqrt{\alpha^2 x^2+ \beta^2 y^2  } }}    \end{pmatrix}
}
{ =} { { \frac{   \sqrt{\alpha^2 x^2+ \beta^2 y^2 } }{  \sqrt{    \alpha x^2 +  \beta y^2  }  } }  \begin{pmatrix}  { \frac{    x }{  \sqrt{\alpha^2 x^2+ \beta^2 y^2 }  } }   \\ { \frac{   y }{  \sqrt{\alpha^2 x^2+ \beta^2 y^2  } }}    \end{pmatrix}
}
{ =} { \begin{pmatrix}  x  \\ y   \end{pmatrix}
}
}
{}
{}{}
menunjukkan bahwa kedua pemetaan tersebut saling invers.
}
\endgroup
\subsection*{Solusi Soal 2.13}
% Authority: German Wikiversity pageid 152873, revid 1094182, 2026-06-14T15:10:42Z.
% Exact UTF-8 wikitext witness SHA-256: 1050f074ecb58b7a8211d756b7fcd6fe4a0a4db6a024083c31ca3c5469e9fd50.
% Expanded German TeX source SHA-256: 6896eb6a4a6f4c25bcfaab674847dcd78bc079cfd6c58b3b8542c5c220d85e7b.

Untuk memperoleh
\definitionsverweis {medan normal satuan}{}{,}
kita mulai dari gradien fungsi yang mendeskripsikan elips tersebut, yaitu

\mavergleichskettedisp
{\vergleichskette
{ G(x,y)
}
{ =} { \begin{pmatrix}  2 \alpha x  \\ 2 \beta y    \end{pmatrix}
}
{ } {
}
{ } {
}
{ } {
}
}
{}{}{}
dan dengan menormalisasikannya kita memperoleh

\mavergleichskettedisp
{\vergleichskette
{ N(x,y)
}
{ =} {  { \frac{   1 }{  \sqrt{ \alpha^2 x^2 + \beta^2 y^2 }  } } \begin{pmatrix}   \alpha x  \\  \beta y    \end{pmatrix}
}
{ } {
}
{ } {
}
{ } {
}
}
{}{}{.}
Dengan demikian, pemetaan Gauss tersebut adalah
\maabbeledisp {\varphi} {Y = \left\{  \begin{pmatrix}  x  \\ y   \end{pmatrix}   \in \R^2  \mid   \alpha x^2+\beta y^2 = 1        \right\}  } { S^1
} { \begin{pmatrix}  x  \\ y   \end{pmatrix} } { { \frac{   1 }{  \sqrt{ \alpha^2 x^2 + \beta^2 y^2 }  } } \begin{pmatrix}   \alpha x  \\  \beta y    \end{pmatrix}
} {.} Kesurjektifan pemetaan ini mengikuti
Teorema 2.11.
Untuk membuktikan injektivitasnya, kita tinjau syarat

\mavergleichskettedisp
{\vergleichskette
{ { \frac{   1  }{  \sqrt{ \alpha^2 x^2 + \beta^2 y^2 }  } } \begin{pmatrix}   \alpha x  \\  \beta y    \end{pmatrix}
}
{ =} { { \frac{   1  }{  \sqrt{ \alpha^2 z^2 + \beta^2 w^2 }  } } \begin{pmatrix}   \alpha z  \\  \beta w    \end{pmatrix}
}
{ } {
}
{ } {
}
{ } {
}
}
{}{}{.}
Ini berarti bahwa
\mathkor {} {\begin{pmatrix}   \alpha x  \\  \beta y    \end{pmatrix}} {dan} {\begin{pmatrix}   \alpha z  \\  \beta w    \end{pmatrix}} {}
bergantung secara linear dengan faktor positif $c$, sehingga
$ (x,y)=c(z,w) $. Karena kedua titik terletak pada $Y$,

\[
1=\alpha x^2+\beta y^2
=c^2\left(\alpha z^2+\beta w^2\right)=c^2.
\]

Karena $c>0$, diperoleh $c=1$. Jadi berlaku

\mavergleichskettedisp
{\vergleichskette
{ \begin{pmatrix}   \alpha x  \\  \beta y    \end{pmatrix}
}
{ =} { \begin{pmatrix}   \alpha z  \\  \beta w    \end{pmatrix}
}
{ } {
}
{ } {
}
{ } {
}
}
{}{}{}
dan oleh karena itu

\mavergleichskette
{\vergleichskette
{ x
}
{ = }{ z
}
{  }{
}
{    }{
}
{   }{
}
}
{}{}{}
dan

\mavergleichskette
{\vergleichskette
{ y
}
{ = }{ w
}
{  }{
}
{    }{
}
{   }{
}
}
{}{}{.}

\part{Unit 3}
\chapter[Kuliah 3]{Kuliah 3: Kelengkungan Kurva Berparameter Panjang Busur}
\setcounter{section}{3}

  \zwischenueberschrift{Kelengkungan Kurva Berparametrisasi Panjang Busur}

\inputdefinition
{}
{

Suatu
\definitionsverweis {kurva terdiferensialkan}{}{}
\maabbeledisp {f} { [a,b] } {\R^n
} { t } {f(t) = \left( f_1(t)  , \,   \ldots  , \,  f_n(t)                 \right)
} {,}
disebut
\definitionswort {berparametrisasi panjang busur}{,}
jika

\mavergleichskette
{\vergleichskette
{ f_1'(t)^2 + \cdots +  f_n'(t)^2
}
{ = }{ 1
}
{  }{
}
{    }{
}
{   }{
}
}
{}{}{}
berlaku untuk setiap $t$.

}

Menurut
Soal 2.3,
kurva berparametrisasi panjang busur yang terdiferensialkan dua kali secara kontinu memiliki sifat bahwa turunan keduanya $\gamma^{\prime \prime} (t)$ selalu tegak lurus terhadap turunan pertamanya $\gamma^{\prime } (t)$.

Misalkan
\maabbdisp {\gamma} { I } {\R^2
} {}
adalah suatu pemetaan yang dua kali
\definitionsverweis {terdiferensialkan secara kontinu}{}{}
dan merupakan
\definitionsverweis {kurva}{}{},
dan

\mavergleichskette
{\vergleichskette
{ t
}
{ \in }{ I
}
{  }{
}
{    }{
}
{   }{
}
}
{}{}{.}
Salah satu cara alami untuk menghampiri perilaku kurva pada parameter
\mathkor {} {t} {yakni di titik} {\gamma(t)} {}
melampaui pendekatan oleh garis singgungnya adalah dengan menentukan suatu lingkaran yang menyesuaikan diri sebaik mungkin dengan kurva tersebut, atau menentukan gerak melingkar yang cocok dengan kurva hingga turunan kedua. Andaikan kurva tersebut
\definitionsverweis {berparametrisasi panjang busur}{}{.}

\inputdefinition
{}
{

Misalkan
\maabbdisp {\gamma} { I } {\R^2
} {}
adalah suatu pemetaan yang dua kali
\definitionsverweis {terdiferensialkan secara kontinu}{}{},
\definitionsverweis {berparametrisasi panjang busur}{}{}
dan merupakan
\definitionsverweis {kurva}{}{},
dan

\mavergleichskette
{\vergleichskette
{ t_0
}
{ \in }{ I
}
{  }{
}
{    }{
}
{   }{
}
}
{}{}{,}
dengan turunan kedua $\gamma^{\prime \prime} (t_0)$ tidak sama dengan $0$. Lingkaran berjari-jari

\mavergleichskettedisp
{\vergleichskette
{ r
}
{ =} { { \frac{   1  }{  \betrag {  \gamma_1'(t_0) \gamma_2^{\prime \prime} (t_0) - \gamma_2'(t_0) \gamma_1^{\prime \prime} (t_0)  }  } }
}
{ } {
}
{ } {
}
{ } {
}
}
{}{}{}
dan berpusat di

\mavergleichskettedisp
{\vergleichskette
{ M
}
{ =} { \gamma(t_0) +  { \frac{   1  }{  \gamma_1'(t_0) \gamma_2^{\prime \prime} (t_0) - \gamma_2'(t_0) \gamma_1^{\prime \prime} (t_0)  } }  \begin{pmatrix}  - \gamma_2'(t_0) \\ \gamma_1'(t_0)   \end{pmatrix}
}
{ } {
}
{ } {
}
{ } {
}
}
{}{}{}
disebut
\definitionswort {lingkaran kelengkungan}{}
dari $\gamma$ di $t_0$.

}

\inputdefinition
{}
{

Misalkan
\maabbdisp {\gamma} { I } {\R^2
} {}
adalah suatu pemetaan yang dua kali
\definitionsverweis {terdiferensialkan secara kontinu}{}{},
\definitionsverweis {berparametrisasi panjang busur}{}{}
dan merupakan
\definitionsverweis {kurva}{}{},
dan

\mavergleichskette
{\vergleichskette
{ t_0
}
{ \in }{ I
}
{  }{
}
{    }{
}
{   }{
}
}
{}{}{.}
Besaran

\mavergleichskettedisp
{\vergleichskette
{ \kappa (t_0)
}
{ =} { \gamma_1'(t_0) \gamma_2^{\prime \prime} (t_0) - \gamma_2'(t_0) \gamma_1^{\prime \prime} (t_0)
}
{ } {
}
{ } {
}
{ } {
}
}
{}{}{}
disebut
\definitionswort {kelengkungan}{}
kurva di $t_0$.

}

Lingkaran kelengkungan juga disebut \stichwort {lingkaran oskulasi} {.} Dalam kasus kurva berparametrisasi panjang busur,
\mathkor {} {\begin{pmatrix}  \gamma_1'(t_0) \\ \gamma_2'(t_0)    \end{pmatrix}} {dan} {\begin{pmatrix}  \gamma_1^{\prime \prime} (t_0) \\ \gamma_2^{\prime \prime}(t_0)    \end{pmatrix}} {}
saling tegak lurus, dan turunan kedua tidak sama dengan $0$. Karena itu, kedua vektor ini membentuk suatu basis bagi $\R^2$, sehingga determinannya, yang muncul sebagai penyebut dalam definisi lingkaran kelengkungan, tidak sama dengan $0$. Jika determinan ini
\zusatzklammer {dan dengan demikian kelengkungannya} {} {}
positif, basis tersebut merepresentasikan
\definitionsverweis {orientasi standar}{}{}
\zusatzklammer {yaitu orientasi yang diberikan oleh vektor standar
\mathkor {} {e_1} {dan} {e_2} {}} {} {};
jika tidak, basis tersebut merepresentasikan orientasi yang berlawanan. Jika kelengkungan positif, geraknya membelok
\zusatzklammer {dilihat dari arah singgung} {} {}
ke kiri; jika kelengkungan negatif, geraknya membelok ke kanan. Untuk kelengkungan positif, pusat lingkaran kelengkungan dapat dituliskan sebagai

\mavergleichskettedisp
{\vergleichskette
{ M
}
{ =} { \gamma(t_0) + r  \begin{pmatrix}  - \gamma_2'(t_0) \\ \gamma_1'(t_0)   \end{pmatrix}
}
{ } {
}
{ } {
}
{ } {
}
}
{}{}{}
dan untuk kelengkungan negatif sebagai

\mavergleichskettedisp
{\vergleichskette
{ M
}
{ =} { \gamma(t_0) + r \begin{pmatrix}  \gamma_2'(t_0) \\ - \gamma_1'(t_0)   \end{pmatrix}
}
{ } {
}
{ } {
}
{ } {
}
}
{}{}{.}
Menurut
Soal 3.1,
kurva yang dilalui dengan arah berlawanan memiliki kelengkungan dengan tanda yang berlawanan. Sesekali kita juga akan berbicara tentang kelengkungan kurva di titik

\mavergleichskette
{\vergleichskette
{ \gamma(t)
}
{ \in }{ \gamma(I)
}
{ \subseteq }{ \R^2
}
{    }{
}
{   }{
}
}
{}{}{};
hal ini tidak menimbulkan masalah bagi kurva yang injektif atau, kalaupun tidak injektif, hanya dilalui berulang kali secara periodik.

\inputbeispiel{}
{

Kita meninjau parametrisasi trigonometrik lingkaran satuan, yaitu
\maabbeledisp {\gamma} {\R} { \R^2
} { t } { \begin{pmatrix}  \cos  t  \\ \sin  t    \end{pmatrix}
} {.}
Kita memperoleh

\mavergleichskettedisp
{\vergleichskette
{ \gamma'(t)
}
{ =} { \begin{pmatrix}  - \sin  t  \\ \cos  t    \end{pmatrix}
}
{ } {
}
{ } {
}
{ } {
}
}
{}{}{}
dan

\mavergleichskettedisp
{\vergleichskette
{ \gamma^{\prime \prime} (t)
}
{ =} { \begin{pmatrix}  - \cos  t  \\ - \sin  t    \end{pmatrix}
}
{ } {
}
{ } {
}
{ } {
}
}
{}{}{.}
\definitionsverweis {Kelengkungan}{}{}
pada setiap waktu $t$ adalah

\mavergleichskettedisp
{\vergleichskette
{ \gamma_1'(t) \gamma_2^{\prime \prime} (t) - \gamma_2'(t) \gamma_1^{\prime \prime} (t)
}
{ =} { \left(  -  \sin  t   \right)  \left(  -  \sin  t   \right)  -  \cos  t  \left(  - \cos  t   \right)
}
{ =} { 1
}
{ } {
}
{ } {
}
}
{}{}{,}
jari-jari kelengkungannya adalah $1$, dan pusat lingkaran kelengkungannya adalah

\mavergleichskettedisp
{\vergleichskette
{ \begin{pmatrix}  \cos  t  \\ \sin  t    \end{pmatrix}  + \begin{pmatrix}  -\cos  t  \\ - \sin  t    \end{pmatrix}
}
{ =} { \begin{pmatrix}  0  \\ 0    \end{pmatrix}
}
{ } {
}
{ } {
}
{ } {
}
}
{}{}{,}
sehingga lingkaran kelengkungannya selalu merupakan lingkaran satuan.

Jika lingkaran itu dilalui searah jarum jam, yaitu jika kita meninjau kurva
\maabbeledisp {\delta} {\R} { \R^2
} { s } { \begin{pmatrix}  \cos \left(  -s \right)  \\ \sin  \left(  -s \right)    \end{pmatrix} = \begin{pmatrix}  \cos  s  \\ - \sin  s    \end{pmatrix}
} {,}
maka kelengkungannya adalah $-1$, sedangkan jari-jari kelengkungan dan lingkaran kelengkungannya sama seperti sebelumnya.

}

Definisi lingkaran kelengkungan sudah dapat digunakan apabila $\gamma'(t_0)$ memiliki norma $1$
\zusatzklammer {jadi hanya di $t_0$} {} {},
tanpa mengharuskan kurva berparametrisasi panjang busur, asalkan
\mathkor {} {\begin{pmatrix}  \gamma_1'(t_0) \\ \gamma_2'(t_0)    \end{pmatrix}} {dan} {\begin{pmatrix}  \gamma_1^{\prime \prime} (t_0) \\ \gamma_2^{\prime \prime}(t_0)    \end{pmatrix}} {}
\definitionsverweis {bebas linear}{}{}.
Jika kedua vektor tersebut bergantung linear, jari-jari kelengkungannya juga dikatakan tak hingga.

\inputbeispiel{}
{

\bild{ \begin{center}
\includegraphics[width=5.5cm]{ \bildeinlesungsvg {Parabola circle} {svg} }
\end{center}
\bildtext {} }

\bildlizenz { Parabola circle.svg } {} {IkamusumeFan} {Commons} {CC-by-sa 4,0} {}

Kita meninjau parabola standar sebagai kurva
\maabbeledisp {\gamma} {\R} {\R^2
} { t } { \begin{pmatrix}  t  \\t^2   \end{pmatrix}
} {}
di titik nol

\mavergleichskette
{\vergleichskette
{ t_0
}
{ = }{ 0
}
{  }{
}
{    }{
}
{   }{
}
}
{}{}{.}
Kita memperoleh

\mavergleichskettedisp
{\vergleichskette
{ \gamma'(t)
}
{ =} { \begin{pmatrix}  1  \\ 2t   \end{pmatrix}
}
{ } {
}
{ } {
}
{ } {
}
}
{}{}{.}
Secara umum kurva ini tidak
\definitionsverweis {berparametrisasi panjang busur}{}{,}
tetapi di titik nol kurva ini memang demikian. Turunan keduanya adalah

\mavergleichskettedisp
{\vergleichskette
{ \gamma^{\prime \prime}(t)
}
{ =} { \begin{pmatrix}  0  \\ 2   \end{pmatrix}
}
{ } {
}
{ } {
}
{ } {
}
}
{}{}{}
dan tidak bergantung pada waktu. Karena itu,
\definitionsverweis {jari-jari kelengkungan}{}{}
adalah

\mavergleichskettedisp
{\vergleichskette
{ r
}
{ =} { { \frac{   1  }{  \betrag {  \gamma_1'(0) \gamma_2^{\prime \prime} (0) - \gamma_2'(0) \gamma_1^{\prime \prime} (0)  }  } }
}
{ =} { { \frac{   1  }{  2 } }
}
{ } {
}
{ } {
}
}
{}{}{}
dan pusat
\definitionsverweis {lingkaran kelengkungan}{}{}
adalah
\mathl{\begin{pmatrix}  0  \\ { \frac{   1  }{  2 } }     \end{pmatrix}}{.}

}



\inputfaktbeweis
{Bogenparametrisierte Kurve/Ebene/Krümmungskreis/Durchlauf/Fakt}
{Lemma}
{}
{

\faktsituation {Misalkan
\maabbdisp {\gamma} { I } {\R^2
} {}
adalah suatu pemetaan yang dua kali
\definitionsverweis {terdiferensialkan secara kontinu}{}{},
\definitionsverweis {berparametrisasi panjang busur}{}{}
dan merupakan
\definitionsverweis {kurva}{}{},
dan

\mavergleichskette
{\vergleichskette
{ t_0
}
{ \in }{ I
}
{  }{
}
{    }{
}
{   }{
}
}
{}{}{}
dengan

\mavergleichskettedisp
{\vergleichskette
{ \gamma^{\prime \prime} (t_0)
}
{ \neq} { 0
}
{ } {
}
{ } {
}
{ } {
}
}
{}{}{},
dan misalkan $K$ adalah
\definitionsverweis {lingkaran kelengkungan}{}{}
dengan pusat $M$ dan jari-jari $r$. Jika pasangan $\gamma'(t_0), \gamma^{\prime \prime} (t_0)$ merepresentasikan
\definitionsverweis {orientasi standar}{}{,}
pilih

\mavergleichskette
{\vergleichskette
{ \alpha
}
{ \in }{ [0, 2 \pi[
}
{  }{
}
{    }{
}
{   }{
}
}
{}{}{}
sedemikian sehingga

\mavergleichskettedisp
{\vergleichskette
{ \begin{pmatrix}  \gamma_1'(t_0) \\ \gamma_2'(t_0)   \end{pmatrix}
}
{ =} { \begin{pmatrix}  - \sin  \left(  \alpha + { \frac{  t_0  }{  r  } }   \right)  \\ \cos \left(  \alpha + { \frac{  t_0  }{  r  } }   \right)   \end{pmatrix}
}
{ } {
}
{ } {
}
{ } {
}
}
{}{}{.}
Jika pasangan $\gamma'(t_0), \gamma^{\prime \prime} (t_0)$ tidak merepresentasikan
\definitionsverweis {orientasi standar}{}{,}
pilih

\mavergleichskette
{\vergleichskette
{ \alpha
}
{ \in }{ [0, 2 \pi[
}
{  }{
}
{    }{
}
{   }{
}
}
{}{}{}
sedemikian sehingga

\mavergleichskettedisp
{\vergleichskette
{ \begin{pmatrix}  \gamma_1'(t_0) \\ \gamma_2'(t_0)   \end{pmatrix}
}
{ =} { \begin{pmatrix}  - \sin  \left(  \alpha + { \frac{  t_0  }{  r  } }   \right)  \\ - \cos \left(  \alpha + { \frac{  t_0  }{  r  } }   \right)   \end{pmatrix}
}
{ } {
}
{ } {
}
{ } {
}
}
{}{}{.}}
\faktfolgerung {Maka
\zusatzklammer {dalam kasus berorientasi standar} {} {},
\maabbeledisp {\delta} {\R} {\R^2
} { t } { M+r \begin{pmatrix}  \cos \left( \alpha + { \frac{   t  }{  r  } }  \right)  \\ \sin  \left( \alpha + { \frac{   t  }{  r  } }  \right)    \end{pmatrix}
} {,}
atau
\zusatzklammer {dalam kasus yang tidak berorientasi standar} {} {},
\maabbeledisp {\delta} {\R} {\R^2
} { t } { M+r \begin{pmatrix}  \cos \left( \alpha + { \frac{   t  }{  r  } }  \right)  \\ - \sin  \left( \alpha + { \frac{   t  }{  r  } }  \right)    \end{pmatrix}
} {,}
merupakan gerak berparametrisasi panjang busur pada lingkaran kelengkungan yang di $t_0$ cocok dengan $\gamma$ hingga turunan kedua.}
\faktzusatz {}
\faktzusatz {}

}
{

Kita meninjau kasus berorientasi standar. Kita memperoleh

\mavergleichskettedisp
{\vergleichskette
{ \delta (t_0)
}
{ =} { M+r \begin{pmatrix}  \cos \left( \alpha +  { \frac{  t_0  }{  r  } }  \right)  \\ \sin  \left(  \alpha + { \frac{  t_0  }{  r  } }  \right)    \end{pmatrix}
}
{ =} { \gamma(t_0) + r\begin{pmatrix}  - \gamma_2'(t_0) \\ \gamma_1'(t_0)   \end{pmatrix} +r \begin{pmatrix}  \gamma_2'(t_0) \\ - \gamma_1'(t_0)   \end{pmatrix}
}
{ =} { \gamma(t_0)
}
{ } {
}
}
{}{}{.}
Selanjutnya,

\mavergleichskettedisp
{\vergleichskette
{ \delta' (t_0)
}
{ =} { \begin{pmatrix}  - \sin  \left( \alpha +  { \frac{  t_0  }{  r  } }  \right)  \\ \cos \left(  \alpha + { \frac{  t_0  }{  r  } }   \right)    \end{pmatrix}
}
{ =} { \begin{pmatrix}  \gamma_1'(t_0) \\ \gamma_2'(t_0)   \end{pmatrix}
}
{ =} { \gamma'(t_0)
}
{ } {
}
}
{}{}{,}
sehingga, khususnya, $\delta$ berparametrisasi panjang busur. Vektor $\begin{pmatrix}  \gamma_1^{\prime \prime} (t_0)  \\ \gamma_2^{\prime \prime} (t_0)   \end{pmatrix}$ tegak lurus terhadap $\begin{pmatrix}  \gamma_1'(t_0)  \\ \gamma_2'(t_0)    \end{pmatrix}$ dan karena itu bergantung linear pada vektor normal satuan
$\begin{pmatrix}  - \gamma_2'(t_0)  \\ \gamma_1'(t_0)    \end{pmatrix}$. Dengan demikian,

\mavergleichskettedisp
{\vergleichskette
{ \begin{pmatrix}  \gamma_1^{\prime \prime} (t_0)  \\ \gamma_2^{\prime \prime} (t_0)   \end{pmatrix}
}
{ =} { \left\langle  \begin{pmatrix}  - \gamma_2'(t_0)  \\ \gamma_1'(t_0)    \end{pmatrix}   ,  \begin{pmatrix}  \gamma_1^{\prime \prime} (t_0) \\ \gamma_2^{\prime \prime} (t_0)   \end{pmatrix}   \right\rangle  \begin{pmatrix}  - \gamma_2'(t_0)  \\ \gamma_1'(t_0)    \end{pmatrix}
}
{ =} { \left(  \gamma_1'(t_0) \gamma_2^{\prime \prime} (t_0) - \gamma_2'(t_0) \gamma_1^{\prime \prime} (t_0)  \right) \begin{pmatrix}  - \gamma_2'(t_0)  \\ \gamma_1'(t_0)    \end{pmatrix}
}
{ } {
}
{ } {
}
}
{}{}{}
dan karenanya

\mavergleichskettealign
{\vergleichskettealign
{ \delta^{\prime \prime} (t_0)
}
{ =} { - { \frac{   1  }{ r } }  \begin{pmatrix}  \cos \left( \alpha + { \frac{  t_0  }{  r  } }  \right)  \\ \sin  \left(  \alpha + { \frac{  t_0  }{  r  } }  \right)    \end{pmatrix}
}
{ =} { - \left(  \gamma_1'(t_0) \gamma_2^{\prime \prime} (t_0) - \gamma_2'(t_0) \gamma_1^{\prime \prime} (t_0)  \right) \begin{pmatrix}  \gamma_2'(t_0) \\ - \gamma_1'(t_0)   \end{pmatrix}
}
{ =} { \left(  \gamma_1'(t_0) \gamma_2^{\prime \prime} (t_0) - \gamma_2'(t_0) \gamma_1^{\prime \prime} (t_0)  \right) \begin{pmatrix}  - \gamma_2'(t_0) \\ \gamma_1'(t_0)   \end{pmatrix}
}
{ =} { \begin{pmatrix}  \gamma_1^{\prime \prime} (t_0) \\ \gamma_2^{\prime \prime} (t_0)   \end{pmatrix}
}
}
{
\vergleichskettefortsetzungalign
{ =} { \gamma^{\prime \prime} (t_0)
}
{ } {
}
{ } {}
{ } {}
}
{}{.}

}

Pernyataan berikut menunjukkan bahwa setiap profil kelengkungan yang diinginkan dapat direalisasikan oleh suatu kurva.


\inputfaktbeweis
{Ebene Kurve/Krümmung/Stammfunktion zu trigonometrischer Funktion/Fakt}
{Lemma}
{}
{

\faktsituation {Misalkan $I$ adalah suatu interval,

\mavergleichskette
{\vergleichskette
{ 0
}
{ \in }{ I
}
{  }{
}
{    }{
}
{   }{
}
}
{}{}{},
dan
\maabbdisp {\psi} { I } {\R
} {}
adalah suatu fungsi yang terdiferensialkan. Definisikan

\mavergleichskettedisp
{\vergleichskette
{ \gamma(t)
}
{ \defeq} { \int_0^t  \begin{pmatrix}  \cos  \psi(s)  \\ \sin  \psi(s)    \end{pmatrix} ds
}
{ } {
}
{ } {
}
{ } {
}
}
{}{}{}}
\faktuebergang {Maka pernyataan-pernyataan berikut berlaku.}
\faktfolgerung {\aufzaehlungfuenf{Kita mempunyai

\mavergleichskettedisp
{\vergleichskette
{ \gamma(0)
}
{ =} { \begin{pmatrix}  0  \\ 0   \end{pmatrix}
}
{ } {
}
{ } {
}
{ } {
}
}
{}{}{.}
}{Kita mempunyai

\mavergleichskettedisp
{\vergleichskette
{ \gamma'(t)
}
{ =} { \begin{pmatrix}  \cos  \psi(t)  \\ \sin  \psi(t)    \end{pmatrix}
}
{ } {
}
{ } {
}
{ } {
}
}
{}{}{.}
}{ \maabbdisp {} { I } { \R^2
} {}
adalah kurva bidang berparametrisasi panjang busur yang terdiferensialkan dua kali.
}{Kita mempunyai

\mavergleichskettedisp
{\vergleichskette
{ \gamma^{\prime \prime}(t)
}
{ =} { \psi' (t) \begin{pmatrix}  - \sin  \psi(t)  \\ \cos  \psi(t)    \end{pmatrix}
}
{ } {
}
{ } {
}
{ } {
}
}
{}{}{.}
}{\definitionsverweis {Kelengkungan}{}{}
adalah

\mavergleichskettedisp
{\vergleichskette
{ \kappa(t)
}
{ =} { \psi'(t)
}
{ } {
}
{ } {
}
{ } {
}
}
{}{}{.}
}}
\faktzusatz {}
\faktzusatz {}

}
{

(1) jelas. (2) langsung mengikuti teorema dasar kalkulus. Dengan demikian, (3) dan (4) juga jelas; parametrisasi panjang busur diperoleh dari

\mavergleichskette
{\vergleichskette
{ \cos^{ 2  }  u + \sin^{ 2  }  u
}
{ = }{ 1
}
{  }{
}
{    }{
}
{   }{
}
}
{}{}{.}
Untuk (5), kelengkungannya adalah

\mavergleichskettedisp
{\vergleichskette
{ \kappa(t)
}
{ =} { \gamma_1'(t) \gamma_2^{\prime \prime} (t) - \gamma_2'(t) \gamma_1^{\prime \prime} (t)
}
{ =} { \psi'(t) \cos  \psi(t)  \cos  \psi(t) + \psi'(t)  \sin  \psi(t)  \sin  \psi(t)
}
{ =} { \psi'(t)
}
{ } {
}
}
{}{}{.}

}



\inputfaktbeweis
{Ebene reguläre Kurve/Krümmung/Skalarprodukt/Fakt}
{Lemma}
{}
{

\faktsituation {Misalkan
\maabb {\gamma} { I } {\R^2
} {}
adalah suatu pemetaan yang dua kali
\definitionsverweis {terdiferensialkan secara kontinu}{}{},
\definitionsverweis {berparametrisasi panjang busur}{}{}
dan merupakan
\definitionsverweis {kurva}{}{.}}
\faktfolgerung {Maka
\definitionsverweis {kelengkungan}{}{}
$\gamma$ di $t$ adalah

\mavergleichskettedisp
{\vergleichskette
{ \kappa(t)
}
{ =} { \left\langle  \gamma^{\prime \prime}(t)  ,  N(t )    \right\rangle
}
{ =} { \pm \Vert { \gamma^{\prime \prime}(t) } \Vert
}
{ } {
}
{ } {
}
}
{}{}{,}
dengan

\mavergleichskettedisp
{\vergleichskette
{ N( t)
}
{ =} { \begin{pmatrix}  -\gamma_2'(t) \\ \gamma_1'(t)   \end{pmatrix}
}
{ } {
}
{ } {
}
{ } {
}
}
{}{}{}
merupakan vektor normal satuan di $\gamma(t)$.}
\faktzusatz {}
\faktzusatz {}

}
{

Pernyataan pertama langsung mengikuti definisi

\mavergleichskettedisp
{\vergleichskette
{ \kappa(t)
}
{ =} { \gamma_1'(t) \gamma_2^{\prime \prime} (t) - \gamma_2'(t) \gamma_1^{\prime \prime} (t)
}
{ } {
}
{ } {
}
{ } {
}
}
{}{}{.}
Dari

\mavergleichskettedisp
{\vergleichskette
{ \left(  \gamma_1'(t)  \right)^2 +  \left(  \gamma_2' (t)  \right)^2
}
{ =} { 1
}
{ } {
}
{ } {
}
{ } {
}
}
{}{}{}
diperoleh

\mavergleichskettedisp
{\vergleichskette
{ \gamma_1'(t) \gamma_1^{\prime \prime}(t) + \gamma_2'(t) \gamma_2^{\prime \prime} (t)
}
{ =} { 0
}
{ } {
}
{ } {
}
{ } {
}
}
{}{}{,}
sehingga $\gamma^{\prime \prime}(t)$ tegak lurus terhadap $\gamma'(t)$ dan bergantung linear pada
\mathl{N(t)}{.} Dengan demikian,

\mavergleichskettedisp
{\vergleichskette
{ \gamma^{\prime \prime}(t)
}
{ =} { \left\langle  \gamma^{\prime \prime}(t) ,  N(t)  \right\rangle  N(t)
}
{ =} { \kappa(t) N(t)
}
{ } {
}
{ } {
}
}
{}{}{}
dan karenanya

\mavergleichskettedisp
{\vergleichskette
{ \Vert { \gamma^{\prime \prime}(t) } \Vert
}
{ =} { \betrag {  \kappa(t) }
}
{ } {
}
{ } {
}
{ } {
}
}
{}{}{.}

}

Dalam pernyataan di atas, tandanya positif jika kelengkungan positif dan negatif jika kelengkungan negatif; apabila kelengkungannya nol, pilihan tanda tidak berpengaruh. Hal yang menentukan bukan deskripsi eksplisit medan normal satuan, melainkan apakah vektor singgung dan vektor normal satuan merepresentasikan orientasi standar atau tidak.

Kita catat pula definisi berikut.

\inputdefinition
{}
{

Misalkan
\maabb {\gamma} { I } { \R^2
} {}
adalah kurva yang terdiferensialkan dua kali secara kontinu dengan

\mavergleichskette
{\vergleichskette
{ \gamma'(t)
}
{ \neq }{ 0
}
{  }{
}
{    }{
}
{   }{
}
}
{}{}{}
untuk setiap

\mavergleichskette
{\vergleichskette
{ t
}
{ \in }{ I
}
{  }{
}
{    }{
}
{   }{
}
}
{}{}{.}
Pemetaan yang didefinisikan pada titik-titik berkelengkungan tak nol,
\maabbeledisp {} { \{t\in I\mid \kappa(t)\neq 0\} } { \R^2
} { t } { M(t)
} {,}
yang memetakan $t$ ke pusat
\definitionsverweis {lingkaran kelengkungan}{}{}
dari $\gamma$ di $t$, disebut
\definitionswort {evolut}{}
dari $\gamma$.

}

  \zwischenueberschrift{Kelengkungan Kurva Secara Umum}

Kita membahas kelengkungan kurva bidang yang tidak harus berparametrisasi panjang busur, serta kurva bidang yang diberikan secara implisit. Misalkan
\maabbdisp {\alpha} {I} { \R^n
} {}
adalah kurva yang terdiferensialkan dua kali secara kontinu dengan

\mavergleichskette
{\vergleichskette
{ \alpha'(t)
}
{ \neq }{ 0
}
{  }{
}
{    }{
}
{   }{
}
}
{}{}{}
untuk setiap

\mavergleichskette
{\vergleichskette
{ t
}
{ \in }{ I
}
{  }{
}
{    }{
}
{   }{
}
}
{}{}{}
dan suatu titik tetap

\mavergleichskette
{\vergleichskette
{ a
}
{ \in }{ I
}
{  }{
}
{    }{
}
{   }{
}
}
{}{}{.}
Besaran

\mavergleichskettedisp
{\vergleichskette
{ \ell(t)
}
{ =} { \int_a^t \Vert { \alpha'(x)} \Vert dx
}
{ } {
}
{ } {
}
{ } {
}
}
{}{}{}
menurut
Teorema 38.6 (Analisis (Osnabrück 2021-2023))
merupakan panjang busur $\alpha$ antara
\mathkor {} {a} {dan} {t} {.}
Pemetaan
\maabbeledisp {\ell} {I} { \R
} {t} { \ell(t)
} {,}
bersifat naik tegas dan terdiferensialkan secara kontinu. Misalkan $J$ adalah interval citranya,
\maabb {\beta} {J} {I
} {}
adalah pemetaan invers dari $\ell$, dan

\mavergleichskette
{\vergleichskette
{  \gamma (u)
}
{ = }{  \alpha( \beta(u))
}
{  }{
}
{    }{
}
{   }{
}
}
{}{}{.}
Dengan menggunakan
Teorema 18.10 (Analisis (Osnabrück 2021-2023))
dan
Teorema 24.3 (Analisis (Osnabrück 2021-2023)),
kita memperoleh

\mavergleichskettealign
{\vergleichskettealign
{ \gamma' (u)
}
{ =} { ( \alpha( \beta(u)))'
}
{ =} { \alpha'( \beta(u)) \beta'(u)
}
{ =} { \alpha'( \beta(u)) { \frac{  1 }{  \ell'( \beta(u))  } }
}
{ =} { \alpha'( \beta(u)) { \frac{  1 }{  \Vert { \alpha'( \beta (u)) } \Vert  } }
}
}
{}
{}{,}
yang berarti bahwa $\gamma$ berparametrisasi panjang busur. Peralihan dari $\alpha$ ke $\gamma$ disebut parametrisasi panjang busur. Peralihan ini tidak mengubah citra kurva, melainkan hanya kecepatan kurva ketika dilalui. Kita memperoleh diagram komutatif berikut.

\mathdisp {\begin{matrix}I  & \stackrel{ \alpha }{\longrightarrow} & \R^n   & \\ \beta \uparrow & \nearrow \gamma \!\!\! \!\! & \\ J  & & & & \!\!\!\!\! \!\!\!  \\ \end{matrix}} {  }



\inputfaktbeweis
{Ebene differenzierbare Kurve/Krümmung/Beschreibung/Fakt}
{Lemma}
{}
{

\faktsituation {Misalkan
\maabb {\alpha} { I } {\R^2
} {}
adalah suatu
\definitionsverweis {kurva yang terdiferensialkan secara kontinu}{}{}
sebanyak dua kali dengan

\mavergleichskette
{\vergleichskette
{ \alpha'(t)
}
{ \neq }{ 0
}
{  }{
}
{    }{
}
{   }{
}
}
{}{}{}
untuk setiap

\mavergleichskette
{\vergleichskette
{ t
}
{ \in }{ I
}
{  }{
}
{    }{
}
{   }{
}
}
{}{}{.}
Definisikan

\mavergleichskettedisp
{\vergleichskette
{ N(t)
}
{ \defeq} { { \frac{   \begin{pmatrix}  -\alpha_2'(t) \\ \alpha_1'(t)   \end{pmatrix}  }{  \Vert {\alpha'(t) } \Vert  } }
}
{ } {
}
{ } {
}
{ } {
}
}
{}{}{.}}
\faktfolgerung {Maka
\definitionsverweis {kelengkungan}{}{}
kurva berparametrisasi panjang busur yang bersesuaian adalah

\mavergleichskettedisp
{\vergleichskette
{ \kappa(t)
}
{ =} { { \frac{   \left\langle  \alpha^{\prime \prime}(t) ,  N(t)   \right\rangle  }{  \Vert {\alpha'(t) } \Vert^2  } }
}
{ =} { { \frac{   \alpha_1'(t) \alpha_2^{\prime\prime} (t) - \alpha_2' (t) \alpha_1^{\prime \prime} (t)  }{  \left(  \alpha_1'(t)^2 + \alpha_2'(t)^2  \right)^{3/2}  } }
}
{ } {
}
{ } {
}
}
{}{}{.}}
\faktzusatz {}
\faktzusatz {}

}
{

Karena

\mavergleichskettealignhandlinks
{\vergleichskettealignhandlinks
{ { \frac{   \left\langle  \alpha^{\prime \prime}(t) ,  N(t)   \right\rangle  }{  \Vert {\alpha'(t) } \Vert^2  } }
}
{ =} { { \frac{   \left\langle  \alpha^{\prime \prime}(t)  ,  { \frac{   1  }{  \Vert {\alpha'(t) } \Vert  } } \begin{pmatrix}  -\alpha_2'(t) \\ \alpha_1'(t)   \end{pmatrix}    \right\rangle  }{  \Vert {\alpha'(t) } \Vert^2  } }
}
{ =} { { \frac{   \left\langle  \begin{pmatrix} \alpha_1^{\prime \prime}(t)   \\\alpha_2^{\prime \prime}(t)    \end{pmatrix}  ,  \begin{pmatrix}  -\alpha_2'(t) \\ \alpha_1'(t)   \end{pmatrix}    \right\rangle  }{  \Vert {\alpha'(t) } \Vert \cdot  \Vert {\alpha'(t) } \Vert^2  } }
}
{ =} { { \frac{   \alpha_1'(t) \alpha_2^{\prime\prime} (t) - \alpha_2' (t) \alpha_1^{\prime \prime} (t)  }{  \left(  \alpha_1'(t)^2 + \alpha_2'(t)^2  \right)^{3/2}  } }
}
{ } {}
}
{}
{}{}
kedua suku tersebut sama. Misalkan

\mavergleichskettedisp
{\vergleichskette
{ \alpha(t)
}
{ =} { \gamma( \beta(t))
}
{ } {
}
{ } {
}
{ } {
}
}
{}{}{}
untuk suatu reparametrisasi meningkat $\beta$ dan kurva berparametrisasi panjang busur $\gamma$. Maka

\mavergleichskettealigndrucklinks
{\vergleichskettealigndrucklinks
{ { \frac{   \alpha_1'(t) \alpha_2^{\prime\prime} (t) - \alpha_2' (t) \alpha_1^{\prime \prime} (t)  }{  \left(  \alpha_1'(t)^2 + \alpha_2'(t)^2  \right)^{3/2}  } }
}
{ =} { { \frac{   \gamma_1'( \beta(t)) \beta'(t) \left(  \gamma_2^{\prime\prime} ( \beta(t)) \beta'(t)^2 + \gamma_2'(\beta(t)) \beta^{\prime \prime} (t)  \right) - \gamma_2'( \beta(t)) \beta'(t) \left(  \gamma_1^{\prime\prime} ( \beta(t)) \beta'(t)^2 + \gamma_1'(\beta(t)) \beta^{\prime \prime} (t)  \right)  }{  \left(  \gamma_1'( \beta(t))^2 \beta'(t)^2 + \gamma_2'( \beta(t))^2 \beta'(t)^2  \right)^{3/2}  } }
}
{ =} { { \frac{   \gamma_1'( \beta(t)) \left(  \gamma_2^{\prime\prime} ( \beta(t)) \beta'(t)^2 + \gamma_2'(\beta(t)) \beta^{\prime \prime} (t)  \right) - \gamma_2'( \beta(t))  \left(  \gamma_1^{\prime\prime} ( \beta(t)) \beta'(t)^2 + \gamma_1'(\beta(t)) \beta^{\prime \prime} (t)  \right)  }{  \beta'(t)^2        \left(  \gamma_1'( \beta(t))^2  + \gamma_2'( \beta(t))^2   \right)^{3/2}  } }
}
{ =} { { \frac{   \gamma_1'( \beta(t)) \gamma_2^{\prime\prime} ( \beta(t)) \beta'(t)^2 - \gamma_2'( \beta(t))  \gamma_1^{\prime\prime} ( \beta(t)) \beta'(t)^2  }{  \beta'(t)^2   } }
}
{ =} { \gamma_1'( \beta(t)) \gamma_2^{\prime\prime} ( \beta(t)) - \gamma_2'( \beta(t))  \gamma_1^{\prime\prime} ( \beta(t))
}
}
{
\vergleichskettefortsetzungalign
{ =} { \kappa (\beta(t))
}
{ } {}
{ } {}
{ } {}
}{}{.}

}

Untuk kurva berorientasi yang diberikan secara implisit

\mavergleichskette
{\vergleichskette
{ Y
}
{ \subseteq }{ \R^2
}
{  }{
}
{    }{
}
{   }{
}
}
{}{}{}
dan suatu titik

\mavergleichskette
{\vergleichskette
{ P
}
{ \in }{ Y
}
{  }{
}
{    }{
}
{   }{
}
}
{}{}{},
kita juga menulis $\kappa(P)$ sebagai pengganti $\kappa(t)$ apabila terdapat suatu parametrisasi panjang busur yang selaras dengan orientasi kurva tersebut dan memenuhi

\mavergleichskette
{\vergleichskette
{ \gamma(t)
}
{ = }{ P
}
{  }{
}
{    }{
}
{   }{
}
}
{}{}{.}



\inputfaktbeweis
{Ebene differenzierbare Kurve/Faser/Krümmung/Fakt}
{Lemma}
{}
{

\faktsituation {Misalkan

\mavergleichskette
{\vergleichskette
{ W
}
{ \subseteq }{ \R^2
}
{  }{
}
{    }{
}
{   }{
}
}
{}{}{}
\definitionsverweis {terbuka}{}{,}
\maabb {h} { W } { \R
} {}
adalah suatu
\definitionsverweis {fungsi yang terdiferensialkan secara kontinu}{}{}
sebanyak dua kali,
dan

\mavergleichskette
{\vergleichskette
{ Y
}
{ = }{ h^{-1}(c)
}
{  }{
}
{    }{
}
{   }{
}
}
{}{}{}
adalah
\definitionsverweis {serat}{}{}
di atas

\mavergleichskette
{\vergleichskette
{ c
}
{ \in }{ \R
}
{  }{
}
{    }{
}
{   }{
}
}
{}{}{,}
dengan $h$
\definitionsverweis {reguler}{}{}
di setiap titik pada $Y$. Misalkan
\maabbdisp {N} { U } {\R^2
} {}
adalah suatu medan yang didefinisikan pada lingkungan terbuka

\mavergleichskette
{\vergleichskette
{ Y
}
{ \subseteq }{ U
}
{  }{
}
{    }{
}
{   }{
}
}
{}{}{}
dan merupakan
\definitionsverweis {medan normal satuan}{}{}
bagi $Y$. Selanjutnya, misalkan

\mavergleichskette
{\vergleichskette
{ P
}
{ \in }{ Y
}
{  }{
}
{    }{
}
{   }{
}
}
{}{}{.}}
\faktfolgerung {Maka untuk setiap vektor singgung

\mavergleichskette
{\vergleichskette
{ v
}
{ \in }{ T_PY \setminus \{0\}
}
{  }{
}
{    }{
}
{   }{
}
}
{}{}{}

berlaku
\mavergleichskettedisp
{\vergleichskette
{ \left(  D N   \right)_{ P } \left(  v  \right)
}
{ =} { - \kappa(P) v
}
{ } {
}
{ } {
}
{ } {
}
}
{}{}{,}
dengan $\kappa(P)$ adalah
\definitionsverweis {kelengkungan}{}{}
dari suatu parametrisasi panjang busur dari $Y$ yang selaras dengan orientasi yang diberikan oleh $v, N(P)$. Untuk vektor nol, persamaan yang sama tetap berlaku karena linearitas.}
\faktzusatz {}
\faktzusatz {}

}
{

Misalkan
\maabb {\gamma} { I } { \R^2
} {}
adalah parametrisasi panjang busur kurva $Y$ pada suatu lingkungan dari $P$, dengan

\mavergleichskette
{\vergleichskette
{ \gamma (0)
}
{ = }{ P
}
{  }{
}
{    }{
}
{   }{
}
}
{}{}{,}
yang selaras dengan orientasi yang diberikan. Diferensial total
\maabbdisp {\left(D N \right)_{ P }} { \R^2} {\R^2
} {}
bersifat linear. Karena itu, cukup membuktikan pernyataan tersebut untuk vektor $\gamma'(0)$. Menurut
aturan rantai,

\mavergleichskettedisp
{\vergleichskette
{ \left(  D N   \right)_{ P } \left(   \gamma'(0)  \right)
}
{ =} { (N \circ \gamma )'(0)
}
{ } {
}
{ } {
}
{ } {
}
}
{}{}{.}
Vektor ini merupakan kelipatan dari $\gamma'(0)$ dan karena itu sama dengan

\mathdisp {\left\langle  (N \circ  \gamma)'(0) ,  \gamma' (0)   \right\rangle  \gamma'(0)} { . }

Berdasarkan syarat keortogonalan,

\mavergleichskettedisp
{\vergleichskette
{ \left\langle  (N \circ  \gamma)(t) ,  \gamma' (t)   \right\rangle
}
{ =} { 0
}
{ } {
}
{ } {
}
{ } {
}
}
{}{}{}
untuk setiap

\mavergleichskette
{\vergleichskette
{ t
}
{ \in }{ I
}
{  }{
}
{    }{
}
{   }{
}
}
{}{}{},
dan karenanya

\mavergleichskettedisp
{\vergleichskette
{ \left\langle  (N \circ  \gamma)'(t) ,  \gamma' (t)   \right\rangle + \left\langle  (N \circ  \gamma)(t) ,  \gamma^{\prime \prime} (t)   \right\rangle
}
{ =} { 0
}
{ } {
}
{ } {
}
{ } {
}
}
{}{}{.}
Jadi,

\mavergleichskettedisp
{\vergleichskette
{ \left(  D N   \right)_{ P } \left(   \gamma'(0)  \right)
}
{ =} { - \left\langle  (N \circ  \gamma)(0) ,  \gamma^{\prime \prime} (0)   \right\rangle  \gamma'(0)
}
{ =} { - \kappa( \gamma(0)) \gamma'(0)
}
{ } {
}
{ } {
}
}
{}{}{}
menurut
Lema 3.8.

}

\chapter{Lembar Kerja 3}
\setcounter{section}{3}

  \zwischenueberschrift{Soal latihan}

\inputaufgabe
{}
{

Misalkan
\maabb {\gamma} { I } {\R^2
} {}
sebuah
\definitionsverweis {kurva berparametrisasi panjang busur}{}{}
yang dua kali terdiferensialkan secara kontinu, dan misalkan
\maabbeledisp {\delta} { -I } {\R^2
} { s } { \delta(s) = \gamma( -s)
} {,}
yaitu kurva yang ditelusuri dalam arah sebaliknya. Buktikan bahwa
\definitionsverweis {kelengkungan}{}{}
$\delta$ di $s$ merupakan negatif dari kelengkungan
\mathl{\gamma}{} di $-s$.

}
{} {}

\inputaufgabe
{}
{

Misalkan
\maabb {\gamma} { I } {\R^2
} {}
sebuah
\definitionsverweis {kurva berparametrisasi panjang busur}{}{}
yang dua kali terdiferensialkan secara kontinu, dan
\maabbdisp {L} {\R^2} {\R^2
} {}
sebuah
\definitionsverweis {rotasi}{}{.}
Buktikan bahwa
\definitionsverweis {kelengkungan}{}{}
$\gamma$ sama dengan kelengkungan
\mathl{L\circ \gamma}{}.

}
{} {}

\inputaufgabe
{}
{

Buktikan
Lema 3.6
dalam kasus dengan orientasi nonstandar.

}
{} {}

\inputaufgabe
{}
{

Misalkan
\maabb {\gamma} { I } {\R^2
} {}
sebuah
\definitionsverweis {kurva berparametrisasi panjang busur}{}{}
yang dua kali terdiferensialkan secara kontinu, dan
\maabbdisp {L} {\R^2} {\R^2
} {}
sebuah
\definitionsverweis {homoteti skalar}{}{}
dengan faktor homoteti

\mavergleichskette
{\vergleichskette
{ s
}
{ \neq }{ 0
}
{  }{
}
{    }{
}
{   }{
}
}
{}{}{.}
Bagaimana hubungan antara
\definitionsverweis {kelengkungan}{}{}
$\gamma$ dan kelengkungan
\mathl{L\circ \gamma}{?}

}
{} {}

\inputaufgabe
{}
{

Misalkan
\maabbeledisp {\gamma} {\R} { \R^2
} { t } { \begin{pmatrix}  \cos \left(  \omega t \right)  \\ \sin  \left(  \omega t \right)    \end{pmatrix}
} {,}
dengan

\mavergleichskette
{\vergleichskette
{ \omega
}
{ > }{ 0
}
{  }{
}
{    }{
}
{   }{
}
}
{}{}{.}
Buktikan bahwa terdapat tak berhingga banyak gerak melingkar dengan norma kecepatan konstan yang memiliki nilai dan turunan pertama yang sama dengan $\gamma$ pada

\mavergleichskette
{\vergleichskette
{ t
}
{ = }{ 0
}
{  }{
}
{    }{
}
{   }{
}
}
{}{}{.}

}
{} {}

Soal berikut penting dalam pembuatan komidi putar. Kesenangan menaiki komidi putar berasal dari percepatan, bukan dari kecepatan.

\inputaufgabe
{}
{

Misalkan
\maabbeledisp {\gamma} {\R} { \R^2
} { t } { \begin{pmatrix}  \cos \left(  \omega t \right)  \\ \sin  \left(  \omega t \right)    \end{pmatrix}
} {,}
dengan

\mavergleichskette
{\vergleichskette
{ \omega
}
{ > }{ 0
}
{  }{
}
{    }{
}
{   }{
}
}
{}{}{.}
Buktikan bahwa terdapat tak berhingga banyak gerak melingkar dengan norma kecepatan konstan yang berimpit dengan $\gamma$ pada

\mavergleichskette
{\vergleichskette
{ t
}
{ = }{ 0
}
{  }{
}
{    }{
}
{   }{
}
}
{}{}{}
dan juga memiliki percepatan yang sama di sana.

}
{} {}

\inputaufgabegibtloesung
{}
{

Misalkan
\maabbeledisp {\gamma} {\R} { \R^2
} { t } { \begin{pmatrix}  \cos \left(  \omega t \right)  \\ \sin  \left(  \omega t \right)    \end{pmatrix}
} {,}
dengan

\mavergleichskette
{\vergleichskette
{ \omega
}
{ > }{ 0
}
{  }{
}
{    }{
}
{   }{
}
}
{}{}{.}
Buktikan bahwa tidak ada gerak melingkar lain dengan norma kecepatan konstan yang memiliki nilai serta turunan pertama dan kedua yang sama dengan $\gamma$ pada

\mavergleichskette
{\vergleichskette
{ t
}
{ = }{ 0
}
{  }{
}
{    }{
}
{   }{
}
}
{}{}{.}

}
{} {}

\inputaufgabe
{}
{

Misalkan

\mavergleichskette
{\vergleichskette
{ W
}
{ \subseteq }{ \R^2
}
{  }{
}
{    }{
}
{   }{
}
}
{}{}{}
\definitionsverweis {terbuka}{}{,}
\maabb {h} { W } { \R
} {}
sebuah fungsi yang dua kali
\definitionsverweis {terdiferensialkan secara kontinu}{}{,}
dan

\mavergleichskette
{\vergleichskette
{ Y
}
{ = }{ h^{-1}(c)
}
{  }{
}
{    }{
}
{   }{
}
}
{}{}{}
adalah
\definitionsverweis {serat}{}{}
di atas

\mavergleichskette
{\vergleichskette
{ c
}
{ \in }{ \R
}
{  }{
}
{    }{
}
{   }{
}
}
{}{}{,}
dengan $h$
\definitionsverweis {reguler}{}{}
di setiap titik pada $Y$. Buktikan bahwa
\maabbdisp {\gamma} { I } { Y
} {}
merupakan
\definitionsverweis {kurva geodesik}{}{}
jika dan hanya jika $\gamma$
\definitionsverweis {berparametrisasi panjang busur}{}{.}

}
{} {}

\inputaufgabe
{}
{

\bild{ \begin{center}
\includegraphics[width=5.5cm]{ \bildeinlesungsvg {Euler spiral} {svg} }
\end{center}
\bildtext {} }

\bildlizenz { Euler spiral.svg } {} {AdiJapan} {Commons} {CC-by-sa 3.0} {}

Kita tinjau kurva terdiferensialkan
\zusatzklammer {yang disebut \stichwort {klotoid} {}} {} {}
\maabbeledisp {} {\R} {\R^2
} { t } { \gamma(t)
} {,}
dengan

\mavergleichskettedisp
{\vergleichskette
{ \gamma(t)
}
{ =} { \sqrt{\pi} \int_0^t  \begin{pmatrix}  \cos  { \frac{   \pi s^2 }{  2 } }  \\ \sin  { \frac{   \pi s^2 }{  2 } }     \end{pmatrix}  ds
}
{ } {
}
{ } {
}
{ } {
}
}
{}{}{.}
Buktikan bahwa
\definitionsverweis {kelengkungan}{}{}
kurva ini memenuhi persamaan

\mavergleichskettedisp
{\vergleichskette
{ \kappa(t)
}
{ =} { \sqrt{\pi} t
}
{ } {
}
{ } {
}
{ } {
}
}
{}{}{.}

}
{} {}

\inputaufgabe
{}
{

Misalkan
\maabb {f} { I } {\R
} {}
sebuah fungsi yang dua kali
\definitionsverweis {terdiferensialkan secara kontinu}{}{.}
Tentukan
\definitionsverweis {kelengkungan}{}{}
dan lingkaran kelengkungan dari grafik yang terkait.

}
{} {}

\inputaufgabe
{}
{

Misalkan
\maabb {f} {\R} {\R
} {}
dua kali
\definitionsverweis {terdiferensialkan secara kontinu}{}{}
dan

\mavergleichskettedisp
{\vergleichskette
{ \alpha(t)
}
{ \defeq} { \begin{pmatrix}  \cos f(t)  \\ \sin f(t)   \end{pmatrix}
}
{ } {
}
{ } {
}
{ } {
}
}
{}{}{.}
Tentukan
\definitionsverweis {kelengkungan}{}{}
$\alpha$ dan
\definitionsverweis {lingkaran kelengkungan}{}{}
untuk

\mavergleichskette
{\vergleichskette
{ t
}
{ \in }{ \R
}
{  }{
}
{    }{
}
{   }{
}
}
{}{}{}
dengan syarat $f'(t) \neq 0$, menggunakan rumus-rumus dalam
Lema 3.10.

}
{} {}

\inputaufgabe
{}
{

Tentukan
\definitionsverweis {kelengkungan}{}{}
\definitionsverweis {grafik}{}{}
\maabbdisp {\sin} {\R} {\R
} {}
untuk

\mavergleichskette
{\vergleichskette
{ t
}
{ = }{ { \frac{   \pi }{  2 } }
}
{  }{
}
{    }{
}
{   }{
}
}
{}{}{.}

}
{} {}

\inputaufgabe
{}
{

Tentukan
\definitionsverweis {kelengkungan}{}{}
spiral logaritmik
\maabbeledisp {\gamma} {\R} { \R^2
} { t } { \gamma(t) = e^t \left(  \cos  t   , \,   \sin  t                   \right)
} {,}
untuk setiap

\mavergleichskette
{\vergleichskette
{ t
}
{ \in }{ \R
}
{  }{
}
{    }{
}
{   }{
}
}
{}{}{.}

}
{} {}

Sebuah fungsi
\maabb {f} { I } {\R
} {}
pada interval terbuka

\mavergleichskette
{\vergleichskette
{ I
}
{ \subseteq }{ \R
}
{  }{
}
{    }{
}
{   }{
}
}
{}{}{}
disebut
\definitionswort {analitik}{,}
jika pada setiap titik

\mavergleichskette
{\vergleichskette
{ a
}
{ \in }{ I
}
{  }{
}
{    }{
}
{   }{
}
}
{}{}{}
fungsi tersebut dapat dinyatakan oleh suatu
\definitionsverweis {deret pangkat}{}{.}

Sebuah kurva disebut analitik jika setiap fungsi komponennya analitik.

\inputaufgabe
{}
{

Misalkan
\maabb {\kappa} { I } {\R
} {}
sebuah
\definitionsverweis {fungsi analitik}{}{.}
Buktikan bahwa terdapat kurva analitik
\maabb {\gamma} { I } {\R^2
} {}
yang
\definitionsverweis {kelengkungan}{}{}
di titik $\gamma(t)$ sama dengan $\kappa(t)$.

}
{} {}

\inputaufgabe
{}
{

\bild{ \begin{center}
\includegraphics[width=5.5cm]{ \bildeinlesungsvg {Evolute-parab} {svg} }
\end{center}
\bildtext {} }

\bildlizenz { Evolute-parab.svg } {} {Ag2gaeh} {Commons} {CC-by-sa 4.0} {}

Tentukan
\definitionsverweis {evolut}{}{}
kurva
\maabbeledisp {} {\R} {\R^2
} { t } { \begin{pmatrix}  t  \\ t^2    \end{pmatrix}
} {.}

}
{} {}

\inputaufgabegibtloesung
{}
{

Tentukan
\definitionsverweis {kelengkungan}{}{}
di setiap titik lingkaran satuan yang diberikan secara implisit oleh

\mavergleichskettedisp
{\vergleichskette
{ h(x,y)
}
{ =} { x^2+y^2
}
{ =} { 1
}
{ } {
}
{ } {
}
}
{}{}{}
menggunakan
Lema 3.11
dan medan gradien $h$.

}
{} {}

  \zwischenueberschrift{Soal untuk dikumpulkan}

\inputaufgabe
{2}
{

Tentukan
\definitionsverweis {kelengkungan}{}{}
\definitionsverweis {grafik}{}{}
fungsi eksponensial
\maabbdisp {\exp} {\R} {\R
} {}
untuk setiap

\mavergleichskette
{\vergleichskette
{ t
}
{ \in }{ \R
}
{  }{
}
{    }{
}
{   }{
}
}
{}{}{.}

}
{} {}

\inputaufgabe
{2}
{

Tentukan
\definitionsverweis {kelengkungan}{}{}
spiral Archimedes
\maabbeledisp {\gamma} {\R} { \R^2
} { t } { \gamma(t) = t \left(  \cos  t   , \,   \sin  t                   \right)
} {,}
untuk setiap

\mavergleichskette
{\vergleichskette
{ t
}
{ \in }{ \R
}
{  }{
}
{    }{
}
{   }{
}
}
{}{}{.}

}
{} {}

\inputaufgabe
{4}
{

Berikan sebuah contoh kurva berparametrisasi panjang busur
\maabbdisp {\gamma} {\R} {\R^2
} {}
yang dua kali terdiferensialkan secara kontinu, sedemikian sehingga untuk dua waktu

\mavergleichskette
{\vergleichskette
{ t_0
}
{ \neq }{ t_1
}
{  }{
}
{    }{
}
{   }{
}
}
{}{}{}
berlaku kesamaan

\mavergleichskette
{\vergleichskette
{ \gamma (t_0)
}
{ = }{ \gamma(t_1)
}
{  }{
}
{    }{
}
{   }{
}
}
{}{}{}
dan
\definitionsverweis {lingkaran-lingkaran kelengkungan}{}{}
pada
\mathkor {} {t_0} {dan} {t_1} {}
tidak sama.

}
{} {}

\inputaufgabe
{4}
{

Tentukan
\definitionsverweis {evolut}{}{}
kurva
\maabbeledisp {} {\R} {\R^2
} { t } { \begin{pmatrix}  2 \cos  t  \\ 3 \sin  t    \end{pmatrix}
} {.}
Di titik-titik manakah turunan evolut mempunyai titik nol?

}
{} {}

\inputaufgabe
{4}
{

Tentukan
\definitionsverweis {kelengkungan}{}{}
lingkaran satuan yang diberikan secara implisit oleh

\mavergleichskettedisp
{\vergleichskette
{ h(x,y)
}
{ =} { x^4+y^4
}
{ =} { 1
}
{ } {
}
{ } {
}
}
{}{}{}
menggunakan
Lema 3.11
dan medan gradien $h$ di titik-titik berikut.
\aufzaehlungzwei {
\mavergleichskette
{\vergleichskette
{ P
}
{ = }{ \begin{pmatrix}  1  \\ 0   \end{pmatrix}
}
{  }{
}
{    }{
}
{   }{
}
}
{}{}{,}
} {
\mavergleichskette
{\vergleichskette
{ Q
}
{ = }{ \begin{pmatrix}  { \frac{   1  }{  \sqrt[4]{2}  } }  \\ { \frac{   1  }{  \sqrt[4]{2}  } }     \end{pmatrix}
}
{  }{
}
{    }{
}
{   }{
}
}
{}{}{.}
}

}
{} {}

\section*{Solusi yang disediakan oleh sumber}
\addcontentsline{toc}{section}{Solusi yang disediakan oleh sumber}
\subsection*{Solusi Soal 3.7}
% Authority: German Wikiversity pageid 151168, revid 1095913, 2026-06-15T07:25:24Z.
% Exact UTF-8 wikitext witness SHA-256: c771b7c1ed173ddc05c8c4219c3c34e7e8099fe42079740b292d991b2ce3fb25.
% Expanded German TeX source SHA-256: 7e3437274bf4f79b6a3fa719876b09fdbcee4e10523a215e9e6f4ecb566cdb85.

Nilai-nilainya adalah

\mavergleichskettedisp
{\vergleichskette
{ \gamma(0)
}
{ =} { \begin{pmatrix}  1  \\ 0   \end{pmatrix}
}
{ } {
}
{ } {
}
{ } {
}
}
{}{}{,}

\mavergleichskettedisp
{\vergleichskette
{ \gamma'(0)
}
{ =} { \begin{pmatrix}  0  \\ \omega   \end{pmatrix}
}
{ } {
}
{ } {
}
{ } {
}
}
{}{}{}
dan

\mavergleichskettedisp
{\vergleichskette
{ \gamma^{\prime \prime} (0)
}
{ =} { \begin{pmatrix}  - \omega^2  \\ 0   \end{pmatrix}
}
{ } {
}
{ } {
}
{ } {
}
}
{}{}{.}
Suatu gerak melingkar yang memiliki nilai serta turunan pertama dan kedua yang sama dengan gerak yang diberikan pada

\mavergleichskette
{\vergleichskette
{ t
}
{ = }{ 0
}
{  }{
}
{    }{
}
{   }{
}
}
{}{}{}
harus, khususnya, mempunyai garis singgung yang sama di titik ini. Karena itu, pusat lingkaran harus terletak pada sumbu $x$. Selain itu, pusat tersebut harus berada di sebelah kiri
\mathl{\begin{pmatrix}  1  \\ 0   \end{pmatrix}}{},
karena percepatan mengarah ke kiri. Oleh sebab itu, kita dapat menuliskan gerak melingkar beraturan yang dicari, dengan pusat
\mathl{\begin{pmatrix}  x  \\ 0   \end{pmatrix}}{}
\zusatzklammer {
\mavergleichskettek
{\vergleichskettek
{ x
}
{ < }{  1
}
{  }{
}
{    }{
}
{   }{
}
}
{}{}{}} {} {}
dan jari-jari
\mathl{1-x}{}, sebagai

\mavergleichskettedisp
{\vergleichskette
{  \delta (t)
}
{ =} {  \begin{pmatrix}  x  \\ 0   \end{pmatrix} + (1-x) \begin{pmatrix}  \cos \left(  u t \right)  \\ \sin  \left(  u t \right)     \end{pmatrix}
}
{ } {
}
{ } {
}
{ } {
}
}
{}{}{.}
Kita memperoleh

\mavergleichskettedisp
{\vergleichskette
{ \delta(0)
}
{ =} { \gamma(0)
}
{ } {
}
{ } {
}
{ } {
}
}
{}{}{.}
Selanjutnya,

\mavergleichskettedisp
{\vergleichskette
{ \delta'( 0)
}
{ =} {  (1-x)  \begin{pmatrix}  0  \\ u     \end{pmatrix}
}
{ } {
}
{ } {
}
{ } {
}
}
{}{}{}
dan

\mavergleichskettedisp
{\vergleichskette
{ \delta^{\prime \prime}( 0)
}
{ =} {  (1-x)  \begin{pmatrix}  -u^2  \\ 0     \end{pmatrix}
}
{ } {
}
{ } {
}
{ } {
}
}
{}{}{.}
Jadi, agar syarat-syarat yang diminta terpenuhi, harus berlaku

\mavergleichskettedisp
{\vergleichskette
{ \omega
}
{ =} { (1-x) u
}
{ } {
}
{ } {
}
{ } {
}
}
{}{}{}
dan

\mavergleichskettedisp
{\vergleichskette
{ \omega^2
}
{ =} { (1-x) u^2
}
{ } {
}
{ } {
}
{ } {
}
}
{}{}{.}
Dari sini diperoleh

\mavergleichskettedisp
{\vergleichskette
{ \omega^2
}
{ =} {  \omega u
}
{ } {
}
{ } {
}
{ } {
}
}
{}{}{}
dan dengan demikian

\mavergleichskettedisp
{\vergleichskette
{ u
}
{ =} { \omega
}
{ } {
}
{ } {
}
{ } {
}
}
{}{}{}
serta

\mavergleichskettedisp
{\vergleichskette
{ x
}
{ =} { 0
}
{ } {
}
{ } {
}
{ } {
}
}
{}{}{.}
\subsection*{Solusi Soal 3.16}
% Authority: German Wikiversity pageid 151178, revid 1095920, 2026-06-15T07:26:34Z.
% Exact UTF-8 wikitext witness SHA-256: 237b3b08df28c060d3c5c2fb6efdf22ea7ae5df6c1fbfad68d64380fe7b4f852.
% Expanded German TeX source SHA-256: ed49f1889840b8352f634ef01400301849f8303ec27a2a38b334b744ae5e5951.

\definitionsverweis {Medan gradien}{}{}
dari $h$ adalah $\begin{pmatrix}  2x \\ 2y   \end{pmatrix}$, sehingga sebuah medan normal satuan adalah

\mavergleichskettedisp
{\vergleichskette
{ N \begin{pmatrix}  x  \\ y   \end{pmatrix}
}
{ =} { { \frac{   1 }{  \sqrt{x^2+y^2}  } } \begin{pmatrix}  x  \\ y   \end{pmatrix}
}
{ } {
}
{ } {
}
{ } {
}
}
{}{}{.}
\definitionsverweis {Diferensial total}{}{}
$N$ di suatu titik

\mavergleichskettedisp
{\vergleichskette
{ P
}
{ =} { \begin{pmatrix}  x  \\ y   \end{pmatrix}
}
{ } {
}
{ } {
}
{ } {
}
}
{}{}{}
adalah

\mavergleichskettedisp
{\vergleichskette
{  \left(D N \right)_{P}
}
{ =} { \begin{pmatrix}  { \frac{   y^2 }{   \left(  x^2+y^2  \right)^{3/2}  } }   &   - { \frac{   xy }{   \left(  x^2+y^2  \right)^{3/2}  } }   \\    - { \frac{   xy }{   \left(  x^2+y^2  \right)^{3/2}  } }    &  { \frac{   x^2 }{   \left(  x^2+y^2  \right)^{3/2}  } }  \end{pmatrix}
}
{ } {
}
{ } {
}
{ } {
}
}
{}{}{.}
Penerapan pada vektor singgung
\mathl{\begin{pmatrix}  y  \\ -x   \end{pmatrix}}{}
memberikan

\mavergleichskettealign
{\vergleichskettealign
{ \left(D N \right)_{P} \begin{pmatrix}  y  \\ -x   \end{pmatrix}
}
{ =} { \begin{pmatrix}  y^2    &   - xy    \\  - xy   &  x^2  \end{pmatrix} \begin{pmatrix}  y  \\ -x   \end{pmatrix}
}
{ =} { \begin{pmatrix}  y \left(  y^2 +x^2  \right)  \\ -x \left(  y^2 +x^2  \right)    \end{pmatrix}
}
{ =} { \begin{pmatrix}  y   \\ -x    \end{pmatrix}
}
{ } {
}
}
{}
{}{.}
Oleh karena itu, menurut
fakta
kelengkungan kurva di $P$ sama dengan $-1$.
\zusatzklammer {Hal ini sesuai karena
\mathl{\begin{pmatrix}  y  \\ -x   \end{pmatrix}}{} dan medan gradien merepresentasikan orientasi standar, sedangkan
\mathl{\begin{pmatrix}  y  \\ -x   \end{pmatrix}}{} merupakan arah singgung untuk gerak searah jarum jam.} {} {}

\part{Unit 4}
\chapter[Kuliah 4]{Kuliah 4: Pemetaan Weingarten}
\setcounter{section}{4}

  \zwischenueberschrift{Pemetaan Weingarten}

Misalkan

\mavergleichskette
{\vergleichskette
{ Y
}
{ \subseteq }{ \R^n
}
{  }{
}
{    }{
}
{   }{
}
}
{}{}{}
adalah suatu
\definitionsverweis {hipermuka terdiferensialkan}{}{.}
Untuk suatu titik

\mavergleichskette
{\vergleichskette
{ P
}
{ \in }{ Y
}
{  }{
}
{    }{
}
{   }{
}
}
{}{}{}
kita mendefinisikan suatu pemetaan linear
\maabbdisp {L_P} { T_PY } { T_PY
} {}
pada ruang singgung $Y$ di P. Misalkan $N$ adalah suatu medan normal satuan terdiferensialkan pada $Y$ yang didefinisikan pada suatu lingkungan terbuka di sekitar $Y$. Gagasan pokoknya adalah merepresentasikan suatu vektor singgung

\mavergleichskette
{\vergleichskette
{ v
}
{ \in }{ T_PY
}
{  }{
}
{    }{
}
{   }{
}
}
{}{}{}
dengan suatu kurva terdiferensialkan
\maabbdisp {\gamma} { I } { Y
} {}
sehingga

\mavergleichskettedisp
{\vergleichskette
{ \gamma'(0)
}
{ =} { v
}
{ } {
}
{ } {
}
{ } {
}
}
{}{}{.}
Medan normal satuan itu kemudian mendefinisikan pemetaan
\maabbdisp {N \circ \gamma} { I } { \R^n
} {}
sepanjang $I$. Perubahan infinitesimal medan normal satuan sepanjang $\gamma$ diukur oleh limit

\mathdisp {\operatorname{lim}_{   t  \rightarrow  0   } \,  { \frac{   N(\gamma(t))-N(P)  }{  t  } }} {  }

jika limit ini ada. Menurut
Lemma 43.4 (Analysis (Osnabrück 2021-2023)),
limit ini sama dengan turunan berarah $\left(  D_{ v }  N   \right) \left(   P   \right)$ dan, pada gilirannya, sama dengan
\mathl{\left(  D N   \right)_{ P } \left(   v   \right)}{,} yaitu
\definitionsverweis {diferensial total}{}{}
yang dievaluasi pada vektor $v$. Akan ditunjukkan bahwa pemetaan
\maabbeledisp {} { T_PY } { \R^n
} {v } { \left(  D_{v}  N   \right) \left(   P   \right)
} {,}
ini mengambil nilai di $T_PY$.

\inputdefinition
{}
{

Misalkan

\mavergleichskette
{\vergleichskette
{ W
}
{ \subseteq }{ \R^n
}
{  }{
}
{    }{
}
{   }{
}
}
{}{}{}
\definitionsverweis {terbuka}{}{,}
\maabb {h} { W } { \R
} {}
suatu
\definitionsverweis {fungsi yang dua kali terdiferensialkan secara kontinu}{}{}
dan

\mavergleichskette
{\vergleichskette
{ Y
}
{ = }{ h^{-1}(c)
}
{  }{
}
{    }{
}
{   }{
}
}
{}{}{}
adalah
\definitionsverweis {serat}{}{}
di atas

\mavergleichskette
{\vergleichskette
{ c
}
{ \in }{ \R
}
{  }{
}
{    }{
}
{   }{
}
}
{}{}{,}
dengan $h$
\definitionsverweis {reguler}{}{}
pada setiap titik di $Y$. Misalkan $N$ adalah suatu
\definitionsverweis {medan normal satuan}{}{}
dan misalkan

\mavergleichskette
{\vergleichskette
{ P
}
{ \in }{ Y
}
{  }{
}
{    }{
}
{   }{
}
}
{}{}{.}
Pemetaan
\maabbeledisp {L_P} {T_P Y} { T_PY
} { v } { - \left(  D_{v}  N   \right) \left(   P  \right)
} {,}
disebut
\definitionswort {pemetaan Weingarten}{}
pada $T_PY$.

}

Perhatikan bahwa pemetaan Weingarten bergantung pada medan normal satuan, meskipun hal ini tidak selalu dinyatakan secara eksplisit. Jika $Y$ diberikan sebagai serat dari $h$, biasanya digunakan medan gradien ternormalisasi yang bersesuaian dengan $h$.


\inputfaktbeweis
{Hyperfläche/Weingartenabbildung/Tangentialraum/Fakt}
{Lemma}
{}
{

\faktsituation {Misalkan

\mavergleichskette
{\vergleichskette
{ W
}
{ \subseteq }{ \R^n
}
{  }{
}
{    }{
}
{   }{
}
}
{}{}{}
\definitionsverweis {terbuka}{}{,}
\maabb {h} { W } { \R
} {}
suatu
\definitionsverweis {fungsi yang dua kali terdiferensialkan secara kontinu}{}{}
dan

\mavergleichskette
{\vergleichskette
{ Y
}
{ = }{ h^{-1}(c)
}
{  }{
}
{    }{
}
{   }{
}
}
{}{}{}
adalah
\definitionsverweis {serat}{}{}
di atas

\mavergleichskette
{\vergleichskette
{ c
}
{ \in }{ \R
}
{  }{
}
{    }{
}
{   }{
}
}
{}{}{,}
dengan $h$
\definitionsverweis {reguler}{}{}
pada setiap titik di $Y$. Misalkan

\mavergleichskette
{\vergleichskette
{ P
}
{ \in }{ Y
}
{  }{
}
{    }{
}
{   }{
}
}
{}{}{.}}
\faktfolgerung {Maka
\definitionsverweis {pemetaan Weingarten}{}{}
$L_P$ adalah suatu endomorfisme linear dari
\definitionsverweis {ruang singgung}{}{}
$T_PY$.}
Catatan edisi: sumber memakai huruf kecil p pada ruang singgung ini, padahal titik yang dikuantifikasi adalah P.
\faktzusatz {}
\faktzusatz {}

}
{

Medan

\mavergleichskette
{\vergleichskette
{ N(P)
}
{ = }{ { \frac{
\operatorname{Grad} \,  h  (  P )  }{  \Vert {
\operatorname{Grad} \,  h  (  P ) } \Vert  } }
}
{  }{
}
{    }{
}
{   }{
}
}
{}{}{}
adalah suatu medan vektor yang terdiferensialkan secara kontinu pada suatu lingkungan terbuka

\mavergleichskette
{\vergleichskette
{ Y
}
{ \subseteq }{ U
}
{  }{
}
{    }{
}
{   }{
}
}
{}{}{}
tempat medan tersebut didefinisikan. Oleh karena itu, menurut
Proposition 46.1 (Analysis (Osnabrück 2021-2023)),

\mavergleichskettedisp
{\vergleichskette
{ \left(  D_{v}  N   \right) \left(   P  \right)
}
{ =} { \left(  D N   \right)_{ P } \left(  v  \right)
}
{ } {
}
{ } {
}
{ } {
}
}
{}{}{}
bersifat linear terhadap arah $v$. Dari syarat normal satuan diperoleh

\mavergleichskette
{\vergleichskette
{ \left\langle  N(P)  ,  N(P)   \right\rangle
}
{ = }{ 1
}
{  }{
}
{    }{
}
{   }{
}
}
{}{}{}
untuk setiap

\mavergleichskette
{\vergleichskette
{ P
}
{ \in }{ U
}
{  }{
}
{    }{
}
{   }{
}
}
{}{}{}
dan, dengan menggunakan
Soal 43.14 (Analysis (Osnabrück 2021-2023)), diperoleh

\mavergleichskettedisp
{\vergleichskette
{ 0
}
{ =} { \left(  D_{v}  \left\langle  N  ,  N  \right\rangle   \right) \left(   P  \right)
}
{ =} { 2 \left\langle  N(P)  ,  \left(  D_{v}  N   \right) \left(   P  \right)   \right\rangle
}
{ } {
}
{ } {
}
}
{}{}{.}
Jadi,
\mathl{\left(  D_{v}  N   \right) \left(   P  \right)}{} tegak lurus terhadap $N(P)$ dan, untuk

\mavergleichskette
{\vergleichskette
{ P
}
{ \in }{ Y
}
{  }{
}
{    }{
}
{   }{
}
}
{}{}{}
vektor tersebut termasuk dalam ruang singgung $T_PY$.

}

Dengan demikian, pemetaan Weingarten $L_P$ adalah negatif dari pembatasan diferensial total medan normal satuan pada ruang singgung $T_PY$, dengan kodomain ruang singgung $T_PY$. Jika diferensial total dihitung melalui
\definitionsverweis {matriks Jacobi}{}{}
dari $-N$, kita harus menentukan cara matriks itu bekerja pada suatu basis ruang singgung untuk memperoleh representasi matriks dari pemetaan Weingarten.



\inputfaktbeweis
{Ebene Kurve/Weingartenabbildung/Krümmung/Fakt}
{Lemma}
{}
{

\faktsituation {Misalkan

\mavergleichskette
{\vergleichskette
{ W
}
{ \subseteq }{ \R^2
}
{  }{
}
{    }{
}
{   }{
}
}
{}{}{}
\definitionsverweis {terbuka}{}{,}
\maabb {h} { W } { \R
} {}
suatu
\definitionsverweis {fungsi yang dua kali terdiferensialkan secara kontinu}{}{}
dan

\mavergleichskette
{\vergleichskette
{ C
}
{ = }{ h^{-1}(c)
}
{  }{
}
{    }{
}
{   }{
}
}
{}{}{}
adalah serat di atas

\mavergleichskette
{\vergleichskette
{ c
}
{ \in }{ \R
}
{  }{
}
{    }{
}
{   }{
}
}
{}{}{,}
dengan $h$
\definitionsverweis {reguler}{}{}
pada setiap titik di $C$, dan misalkan

\mavergleichskette
{\vergleichskette
{ P
}
{ \in }{ C
}
{  }{
}
{    }{
}
{   }{
}
}
{}{}{.}}
\faktfolgerung {Maka
\definitionsverweis {pemetaan Weingarten}{}{}
di $P$ adalah perkalian dengan
\definitionsverweis {kelengkungan}{}{}
$\kappa(P)$ dari $C$ di $P$.}
\faktzusatz {}
\faktzusatz {}

}
{

Pernyataan ini merupakan perumusan ulang dari
Lemma 3.11.

}

Pernyataan di atas tidak secara eksplisit merujuk pada orientasi; orientasi yang dimaksud harus turut diperhitungkan.

\inputbeispiel{}
{

Misalkan

\mavergleichskettedisp
{\vergleichskette
{ h(x,y,z)
}
{ =} { x^2+y^2+z^2
}
{ } {
}
{ } {
}
{ } {
}
}
{}{}{}
dan perhatikan serat $Y$ dari $h$ di atas $r^2$, yakni permukaan bola berjari-jari

\mavergleichskette
{\vergleichskette
{ r
}
{ > }{ 0
}
{  }{
}
{    }{
}
{   }{
}
}
{}{}{}
yang berpusat di titik asal. Misalkan

\mavergleichskette
{\vergleichskette
{ P
}
{ = }{ \left(  x_0   , \,   y_0  , \,  z_0                 \right)
}
{ \in }{ Y
}
{    }{
}
{   }{
}
}
{}{}{.}
Dengan suatu isometri, titik ini dapat dipetakan ke

\mavergleichskette
{\vergleichskette
{ Q
}
{ = }{ \left(  r   , \,   0  , \,   0                 \right)
}
{  }{
}
{    }{
}
{   }{
}
}
{}{}{}
tanpa mengubah
\definitionsverweis {pemetaan Weingarten}{}{}.
Suatu basis ruang singgungnya ialah
\mathkor {} {\begin{pmatrix}  0  \\ 1 \\  0   \end{pmatrix}} {dan} {\begin{pmatrix}  0  \\ 0\\  1   \end{pmatrix}} {.}
Medan normal satuan $N$ yang mengarah ke dalam adalah
\mathl{- { \frac{   1  }{  2r } } \begin{pmatrix}  2x \\ 2y\\  2z   \end{pmatrix}}{} dan karena itu, untuk

\mavergleichskettedisp
{\vergleichskette
{ v
}
{ =} { \begin{pmatrix}  0  \\ b \\ c   \end{pmatrix}
}
{ } {
}
{ } {
}
{ } {
}
}
{}{}{}

\mavergleichskettedisp
{\vergleichskette
{ \left(  D_{v}  N   \right) \left(   Q   \right)
}
{ =} { b { \frac{   \partial   N   }{  \partial   y  } } (Q) +  c { \frac{   \partial   N   }{  \partial  z  } }  (Q)
}
{ =} { -b { \frac{   1  }{  2r } } \begin{pmatrix}  0  \\ 2 \\  0   \end{pmatrix} - c { \frac{   1  }{  2r } } \begin{pmatrix}  0  \\ 0\\  2   \end{pmatrix}
}
{ =} { - { \frac{   1  }{ r } } \begin{pmatrix}  0  \\ b \\ c   \end{pmatrix}
}
{ } {
}
}
{}{}{.}
Jadi, pemetaan Weingarten adalah perkalian dengan kebalikan ${ \frac{   1  }{ r } }$ dari jari-jari.

}

\inputbeispiel{}
{

Kita melanjutkan
Contoh 2.5.
Matriks Jacobi dari medan normal satuan

\mavergleichskettedisphandlinks
{\vergleichskettedisphandlinks
{ N( \begin{pmatrix}  x  \\ y \\ z   \end{pmatrix} )
}
{ =} { \left(  x^2+y^2+z^2  \right)^{ - { \frac{   1  }{  2 } } } \begin{pmatrix}  x  \\ y \\  -z   \end{pmatrix}
}
{ } {
}
{ } {
}
{ } {
}
}
{}{}{}
adalah

\mathdisp {\left(  x^2+y^2+z^2  \right)^{- { \frac{   3  }{  2 } } } \begin{pmatrix}  y^2+z^2  &   -xy &  -xz \\  -xy &  x^2+z^2 &  -yz \\ xz &  yz &  -x^2-y^2  \end{pmatrix}} { . }

\definitionsverweis {Pemetaan Weingarten}{}{}
adalah negatif dari pembatasan pemetaan ini pada ruang singgung permukaan tersebut; menurut
Lemma 4.2,
hasilnya kembali berada di ruang singgung. Kita mengasumsikan

\mavergleichskette
{\vergleichskette
{ x
}
{ \neq }{ 0
}
{  }{
}
{    }{
}
{   }{
}
}
{}{}{}
dan menggunakan basis
\mathkor {} {\begin{pmatrix}  y  \\ -x\\  0   \end{pmatrix}} {dan} {\begin{pmatrix}  z  \\ 0 \\  x   \end{pmatrix}} {}
dari ruang singgung. Kita mempunyai

\mavergleichskettealignhandlinks
{\vergleichskettealignhandlinks
{ \begin{pmatrix}  y^2+z^2  &   -xy &  -xz \\  -xy &  x^2+z^2 &  -yz \\ xz &  yz &  -x^2-y^2  \end{pmatrix} \begin{pmatrix}  y  \\ -x\\  0   \end{pmatrix}
}
{ =} { \begin{pmatrix}  y \left(  y^2+z^2  \right) +x^2y \\ -xy^2 -x \left(  x^2+z^2  \right) \\  0   \end{pmatrix}
}
{ =} { \left(  x^2+y^2+z^2  \right) \begin{pmatrix}  y  \\ -x\\  0   \end{pmatrix}
}
{ } {
}
{ } {
}
}
{}
{}{,}
sehingga langsung diperoleh vektor eigen $\begin{pmatrix}  y  \\ -x\\  0   \end{pmatrix}$ dengan nilai eigen
\mathl{- \left(  x^2+y^2+z^2  \right)^{ - { \frac{   1  }{  2 } } }}{}.
Selain itu,


\mavergleichskettealigndrucklinks
{\vergleichskettealigndrucklinks
{ \begin{pmatrix}  y^2+z^2  &   -xy &  -xz \\  -xy &  x^2+z^2 &  -yz \\ xz &  yz &  -x^2-y^2  \end{pmatrix} \begin{pmatrix}  z  \\ 0 \\  x   \end{pmatrix}
}
{ =} { \begin{pmatrix} z \left(  y^2+z^2  \right) -x^2z \\ -2xyz\\  xz^2-x \left(  x^2+y^2  \right)   \end{pmatrix}
}
{ =} { 2yz \begin{pmatrix}  y  \\ -x\\  0   \end{pmatrix} + \left(  z^2-x^2-y^2  \right) \begin{pmatrix}  z  \\ 0 \\  x   \end{pmatrix}
}
{ =} { 2yz \begin{pmatrix}  y  \\ -x\\  0   \end{pmatrix}- \begin{pmatrix}  z  \\ 0 \\  x   \end{pmatrix}
}
{ } {
}
}
{}{}{.}

Dengan demikian, terhadap basis ini, pemetaan Weingarten direpresentasikan oleh matriks

\mathdisp {\begin{pmatrix}  - \left(  x^2+y^2+z^2  \right)^{ - { \frac{   1  }{  2 } } }   &   - 2yz \left(  x^2+y^2+z^2  \right)^{ - { \frac{   3  }{  2 } } }    \\  0  &  \left(  x^2+y^2+z^2  \right)^{ - { \frac{   3  }{  2 } } }   \end{pmatrix}} {  }

Untuk memperoleh vektor eigen kedua dalam ruang singgung,
\zusatzklammer {dengan memperhatikan
Korolari 4.8} {} {}
cara termudah ialah menentukan suatu vektor yang tegak lurus terhadap vektor eigen pertama dan vektor normal. Hasilnya adalah vektor
\mathl{\begin{pmatrix}  xz \\ yz\\  x^2+y^2   \end{pmatrix}}{,} dan memang,
\zusatzklammer {tanpa faktor skalar di depan} {} {}


\mavergleichskettealigndrucklinks
{\vergleichskettealigndrucklinks
{ \begin{pmatrix}  y^2+z^2  &   -xy &  -xz \\  -xy &  x^2+z^2 &  -yz \\ xz &  yz &  -x^2-y^2  \end{pmatrix} \begin{pmatrix}  xz \\ yz\\  x^2+y^2   \end{pmatrix}
}
{ =} { \begin{pmatrix}  xz \left(  y^2+z^2  \right) -xy^2z -xz \left(  x^2+y^2  \right)  \\ - x^2yz +yz \left(  x^2+z^2  \right) -yz \left(  x^2+y^2  \right) \\  x^2z^2 +y^2z^2 - \left(  x^2+y^2  \right)^2   \end{pmatrix}
}
{ =} { \left(  - x^2-y^2+z^2  \right) \begin{pmatrix}  xz \\ yz\\  x^2+y^2   \end{pmatrix}
}
{ =} { - \begin{pmatrix}  xz \\ yz\\  x^2+y^2   \end{pmatrix}
}
{ } {
}
}
{}
{}{.}

Jadi, nilai eigennya adalah
\mathl{\left(  x^2+y^2+z^2  \right)^{- { \frac{   3  }{  2 } } }}{}
\zusatzklammer {hal ini juga dapat langsung dibaca dari matriks segitiga yang merepresentasikannya} {} {.}

}




\inputfaktbeweis
{Hyperfläche/Weingartenabbildung/Zweite Ableitung/Skalarprodukt/Fakt}
{Lemma}
{}
{

\faktsituation {Misalkan

\mavergleichskette
{\vergleichskette
{ W
}
{ \subseteq }{ \R^n
}
{  }{
}
{    }{
}
{   }{
}
}
{}{}{}
\definitionsverweis {terbuka}{}{,}
\maabb {h} { W } { \R
} {}
suatu
\definitionsverweis {fungsi yang dua kali terdiferensialkan secara kontinu}{}{}
dan

\mavergleichskette
{\vergleichskette
{ Y
}
{ = }{ h^{-1}(c)
}
{  }{
}
{    }{
}
{   }{
}
}
{}{}{}
adalah
\definitionsverweis {serat}{}{}
di atas

\mavergleichskette
{\vergleichskette
{ c
}
{ \in }{ \R
}
{  }{
}
{    }{
}
{   }{
}
}
{}{}{,}
dengan $h$
\definitionsverweis {reguler}{}{}
pada setiap titik di $Y$. Misalkan $N$ adalah suatu
\definitionsverweis {medan normal satuan}{}{}
dan misalkan

\mavergleichskette
{\vergleichskette
{ P
}
{ \in }{ Y
}
{  }{
}
{    }{
}
{   }{
}
}
{}{}{.}
Misalkan $\gamma$ adalah suatu realisasi yang terdiferensialkan dua kali pada $Y$ dari vektor singgung

\mavergleichskette
{\vergleichskette
{ v
}
{ \in }{ T_PY
}
{  }{
}
{    }{
}
{   }{
}
}
{}{}{.}}
\faktfolgerung {Maka

\mavergleichskettedisp
{\vergleichskette
{ \left\langle  \gamma^{\prime \prime} (0)  ,  N(P)  \right\rangle
}
{ =} { \left\langle  L_P(v) , v  \right\rangle
}
{ } {
}
{ } {
}
{ } {
}
}
{}{}{.}}
\faktzusatz {}
\faktzusatz {}

}
{

Misalkan
\maabbeledisp {\gamma} { I } { Y
} { t } { \gamma(t)
} {,}
terdiferensialkan dua kali, dengan

\mavergleichskette
{\vergleichskette
{ \gamma(0)
}
{ = }{ P
}
{  }{
}
{    }{
}
{   }{
}
}
{}{}{}
dan

\mavergleichskette
{\vergleichskette
{ \gamma'(0)
}
{ = }{ v
}
{ \in }{ T_PY
}
{    }{
}
{   }{
}
}
{}{}{.}
Kita mempunyai

\mavergleichskettedisp
{\vergleichskette
{ \left\langle  \gamma'(t)  ,  N(\gamma(t))   \right\rangle
}
{ =} { 0
}
{ } {
}
{ } {
}
{ } {
}
}
{}{}{}
untuk semua

\mavergleichskette
{\vergleichskette
{ t
}
{ \in }{ I
}
{  }{
}
{    }{
}
{   }{
}
}
{}{}{}
sehingga, dengan
Soal 43.14 (Analysis (Osnabrück 2021-2023)), diperoleh

\mavergleichskettealign
{\vergleichskettealign
{ 0
}
{ =} { \left\langle  \gamma'(t)  ,  N(\gamma(t))   \right\rangle' (0)
}
{ =} { \left\langle  \gamma^{\prime \prime}(0)  ,  N(\gamma(0))  \right\rangle + \left\langle  \gamma'(0)  ,  \left(  D_{ \gamma'(0) }  N   \right) \left(   \gamma(0)   \right)   \right\rangle
}
{ =} { \left\langle  \gamma^{\prime \prime}(0) , N(P)  \right\rangle - \left\langle  L_P(v)  ,  v   \right\rangle
}
{ } {
}
}
{}
{}{.}

}

Dalam pernyataan di atas tidak diperlukan penetapan tambahan mengenai orientasi, sebab medan normal satuan muncul pada kedua ruas, dan pada ruas kanan muncul melalui pemetaan Weingarten. Perhatikan pula bahwa $\gamma$ tidak harus memiliki kecepatan konstan. Dengan demikian, terdapat banyak realisasi dan percepatan $\gamma^{\prime \prime} (0)$ tidak ditentukan secara tunggal. Namun, nilai
\mathl{\left\langle  \gamma^{\prime \prime}(0)  ,  N(P)  \right\rangle}{,} yang menggambarkan komponen normal percepatan, tidak berubah karena $\gamma$ bergerak pada hipermuka tersebut.
Catatan edisi: sumber kehilangan argumen N(P) dalam hasil kali dalam pada kalimat sebelumnya; identitas yang baru saja dibuktikan menentukannya secara unik.



\inputfaktbeweis
{Hyperfläche/Weingartenabbildung/Selbstadjungiert/Fakt}
{Satz}
{}
{

\faktsituation {Misalkan

\mavergleichskette
{\vergleichskette
{ W
}
{ \subseteq }{ \R^n
}
{  }{
}
{    }{
}
{   }{
}
}
{}{}{}
\definitionsverweis {terbuka}{}{,}
\maabb {h} { W } { \R
} {}
suatu
\definitionsverweis {fungsi yang dua kali terdiferensialkan secara kontinu}{}{}
dan

\mavergleichskette
{\vergleichskette
{ Y
}
{ = }{ h^{-1}(c)
}
{  }{
}
{    }{
}
{   }{
}
}
{}{}{}
adalah serat di atas

\mavergleichskette
{\vergleichskette
{ c
}
{ \in }{ \R
}
{  }{
}
{    }{
}
{   }{
}
}
{}{}{,}
dengan $h$
\definitionsverweis {reguler}{}{}
pada setiap titik di $Y$.}
\faktfolgerung {Maka
\definitionsverweis {pemetaan Weingarten}{}{}
$L_P$ pada setiap titik

\mavergleichskette
{\vergleichskette
{ P
}
{ \in }{ Y
}
{  }{
}
{    }{
}
{   }{
}
}
{}{}{}
bersifat
\definitionsverweis {adjoin-diri (self-adjoint)}{}{.}}
\faktzusatz {}
\faktzusatz {}

}
{

Untuk vektor

\mavergleichskette
{\vergleichskette
{ v,w
}
{ \in }{ T_PY
}
{  }{
}
{    }{
}
{   }{
}
}
{}{}{}
kita harus menunjukkan bahwa

\mavergleichskettedisp
{\vergleichskette
{ \left\langle  v  ,  L_P(w)  \right\rangle
}
{ =} { \left\langle L_P(v) ,  w  \right\rangle
}
{ } {
}
{ } {
}
{ } {
}
}
{}{}{}
tersebut. Dengan

\mavergleichskette
{\vergleichskette
{ N(P)
}
{ = }{ { \frac{
\operatorname{Grad} \,  h  (  P )  }{  \Vert {
\operatorname{Grad} \,  h  (  P ) } \Vert  } }
}
{  }{
}
{    }{
}
{   }{
}
}
{}{}{}
menurut
Lemma 45.10 (Analysis (Osnabrück 2021-2023)), kita mempunyai

\mavergleichskettealignhandlinks
{\vergleichskettealignhandlinks
{ \left\langle  v  ,  L_P(w)  \right\rangle
}
{ =} { -\left\langle  v  ,  \left(  D_{w}  N   \right) \left(   P   \right)   \right\rangle
}
{ =} { -\left\langle  v  ,  \left(  D_{w}  { \frac{
\operatorname{Grad} \,  h  }{  \Vert {
\operatorname{Grad} \,  h } \Vert  } }   \right) \left(   P   \right)     \right\rangle
}
{ =} { -\left\langle  v  ,  \left(  D_{w}  { \frac{   1  }{  \Vert {
\operatorname{Grad} \,  h } \Vert  } }   \right) \left(   P   \right) \cdot
\operatorname{Grad} \,  h  (  P ) + { \frac{   1  }{  \Vert {
\operatorname{Grad} \,  h  (  P ) } \Vert  } } \cdot  \left(  D_{w}
\operatorname{Grad} \,  h   \right) \left(   P   \right)   \right\rangle
}
{ =} { - { \frac{   1  }{  \Vert {
\operatorname{Grad} \,  h  (  P ) } \Vert  } } \left\langle  v  ,  \left(  D_{w}
\operatorname{Grad} \,  h   \right) \left(   P   \right)    \right\rangle
}
}
{}
{}{,}
karena suku pertama tegak lurus terhadap vektor singgung $v$. Dalam fungsi-fungsi koordinat,

\mavergleichskettealignhandlinks
{\vergleichskettealignhandlinks
{ \left(  D_{w}
\operatorname{Grad} \,  h   \right) \left(   P   \right)
}
{ =} { \left(  D_{w}  \left(  { \frac{   \partial   h   }{  \partial   x_1   } }   , \,   \ldots  , \,   { \frac{   \partial   h   }{  \partial   x_n   } }                 \right)    \right) \left(   P   \right)
}
{ =} { \sum_{j = 1}^n  w_j  { \frac{   \partial    }{  \partial   x_j   } }  \left(  { \frac{   \partial   h   }{  \partial   x_1   } }   , \,   \ldots  , \,   { \frac{   \partial   h   }{  \partial   x_n   } }                 \right) (P)
}
{ =} { \left(  \sum_{j = 1}^n  w_j  { \frac{   \partial    }{  \partial   x_j   } }  { \frac{   \partial   h   }{  \partial   x_1   } } (P)   , \,   \ldots  , \,   \sum_{j = 1}^n  w_j  { \frac{   \partial    }{  \partial   x_j   } }  { \frac{   \partial   h   }{  \partial   x_n   } } (P)                 \right)
}
{ } {
}
}
{}
{}{.}
Dengan demikian, ungkapan di atas sama dengan

\mavergleichskettealign
{\vergleichskettealign
{ \left\langle  v  ,  L_P(w)  \right\rangle
}
{ =} { - { \frac{   1  }{  \Vert {
\operatorname{Grad} \,  h  (  P ) } \Vert  } } \left\langle  v  ,  \left(  D_{w}
\operatorname{Grad} \,  h   \right) \left(   P   \right)   \right\rangle
}
{ =} { -  { \frac{   1  }{  \Vert {
\operatorname{Grad} \,  h  (  P ) } \Vert  } }  \sum_{i,j = 1}^n v_i  w_j  { \frac{   \partial    }{  \partial   x_j   } } { \frac{   \partial   h   }{  \partial   x_i   } }  (P)
}
{ } {
}
{ } {
}
}
{}
{}{.}
Menurut
Teorema 44.10 (Analysis (Osnabrück 2021-2023)),
urutan turunan parsial dapat dipertukarkan, sehingga peran
\mathkor {} {v} {dan} {w} {}
juga dapat dipertukarkan.

}



\inputfaktbeweis
{Hyperfläche/Weingartenabbildung/Selbstadjungiert/Eigenwerte/Fakt}
{Korollar}
{}
{

\faktsituation {Misalkan

\mavergleichskette
{\vergleichskette
{ W
}
{ \subseteq }{ \R^n
}
{  }{
}
{    }{
}
{   }{
}
}
{}{}{}
\definitionsverweis {terbuka}{}{,}
\maabb {h} { W } { \R
} {}
suatu
\definitionsverweis {fungsi yang dua kali terdiferensialkan secara kontinu}{}{}
dan

\mavergleichskette
{\vergleichskette
{ Y
}
{ = }{ h^{-1}(c)
}
{  }{
}
{    }{
}
{   }{
}
}
{}{}{}
adalah serat di atas

\mavergleichskette
{\vergleichskette
{ c
}
{ \in }{ \R
}
{  }{
}
{    }{
}
{   }{
}
}
{}{}{,}
dengan $h$
\definitionsverweis {reguler}{}{}
pada setiap titik di $Y$.}
\faktfolgerung {Maka
\definitionsverweis {pemetaan Weingarten}{}{}
$L_P$ pada setiap titik

\mavergleichskette
{\vergleichskette
{ P
}
{ \in }{ Y
}
{  }{
}
{    }{
}
{   }{
}
}
{}{}{}
\definitionsverweis {dapat didiagonalkan}{}{}
dengan
\definitionsverweis {nilai-nilai eigen}{}{}
real dan ruang-ruang eigen yang saling ortogonal.}
\faktzusatz {}
\faktzusatz {}

}
{

Hal ini mengikuti dari
Teorema 4.7
dan
Teorema 41.11 (Lineare Algebra (Osnabrück 2024-2025)).

}



\inputfaktbeweis
{Differenzierbare Funktion/Graph/Hyperfläche/Weingartenabbildung/Fakt}
{Lemma}
{}
{

\faktsituation {Misalkan

\mavergleichskette
{\vergleichskette
{ V
}
{ \subseteq }{ \R^{n-1}
}
{  }{
}
{    }{
}
{   }{
}
}
{}{}{}
suatu
\definitionsverweis {himpunan terbuka}{}{}
dan misalkan

\mavergleichskette
{\vergleichskette
{ Y
}
{ \subseteq }{ V \times \R
}
{  }{
}
{    }{
}
{   }{
}
}
{}{}{}
adalah
\definitionsverweis {grafik}{}{}
dari suatu fungsi yang
\definitionsverweis {dua kali terdiferensialkan secara kontinu}{}{}
\maabbdisp {f} { V } {\R
} {.}}
\faktfolgerung {Maka
\definitionsverweis {pemetaan Weingarten}{}{}
$L_P$ pada suatu titik

\mavergleichskette
{\vergleichskette
{ P
}
{ = }{ (x_1 , \ldots , x_{n-1}, f(x_1 , \ldots , x_{n-1}))
}
{ \in }{ Y
}
{    }{
}
{   }{
}
}
{}{}{}
\zusatzklammer {terhadap medan normal satuan yang mengarah ke atas} {} {}
dapat direpresentasikan sebagai berikut. Tuliskan \mathl{g=\operatorname{Grad} f}{} sebagai vektor kolom dan
\mathl{G=I_{n-1}+gg^{\mathsf T}}{,}
yakni matriks metrik dalam basis grafik. Matriks pemetaan Weingarten dalam basis tersebut adalah
\mathl{G^{-1}{ \frac{   1  }{  \sqrt{ 1+ \left(  { \frac{   \partial   f   }{  \partial   x_1   } }  \right)^2 + \cdots + \left(  { \frac{   \partial   f   }{  \partial   x_{n-1}   } }  \right)^2 }  } } \operatorname{Hess}_{   } \,  f}{}, dengan $\operatorname{Hess}_{   } \,  f$ menyatakan
\definitionsverweis {matriks Hesse}{}{}
dari $f$, apabila vektor-vektor dasar

\mavergleichskette
{\vergleichskette
{ v
}
{ \in }{ \R^{n-1}
}
{  }{
}
{    }{
}
{   }{
}
}
{}{}{}
diidentifikasi dengan vektor-vektor singgung dari $T_PY$ dalam pengertian
Contoh 1.2.}
\faktzusatz {}
\faktzusatz {}

}
{

Kita melanjutkan
Contoh 1.2.
Medan normal satuan yang mengarah ke atas diberikan oleh

\mavergleichskettedisphandlinks
{\vergleichskettedisphandlinks
{ N(x_1  , \ldots , x_{n-1}, f(x_1 , \ldots , x_{n-1}))
}
{ =} { { \frac{   1  }{  \sqrt{ 1+ \left(  { \frac{   \partial   f   }{  \partial   x_1   } }  \right)^2 + \cdots + \left(  { \frac{   \partial   f   }{  \partial   x_{n-1}   } }  \right)^2 }  } } \begin{pmatrix}  -{ \frac{   \partial   f   }{  \partial   x_{1}   } }  \\\vdots\\  -{ \frac{   \partial   f   }{  \partial   x_{n-1}   } }\\ 1   \end{pmatrix}
}
{ } {
}
{ } {
}
{ } {
}
}
{}{}{}
yang akan digunakan. Kita meninjau lintasan

\mavergleichskettedisp
{\vergleichskette
{ \gamma(t)
}
{ =} { \begin{pmatrix}  x_1+tv_1  \\ \vdots \\  x_{n-1} +tv_{n-1}\\f \begin{pmatrix}  x_1+tv_1  \\ \vdots \\  x_{n-1} +tv_{n-1}   \end{pmatrix}   \end{pmatrix}
}
{ } {
}
{ } {
}
{ } {
}
}
{}{}{}
pada grafik yang bersesuaian dengan vektor dasar $v$. Turunan keduanya adalah

\mavergleichskettedisp
{\vergleichskette
{ \gamma^{\prime \prime} (0)
}
{ =} { \begin{pmatrix}  0  \\ \vdots \\  0\\ \sum_{1 \leq i, j \leq n-1}  { \frac{   \partial    }{  \partial   x_i   } }  { \frac{   \partial   f   }{  \partial   x_j   } } v_iv_j   \end{pmatrix}
}
{ } {
}
{ } {
}
{ } {
}
}
{}{}{.}
Menurut
Lemma 4.6,

\mavergleichskettealigndrucklinks
{\vergleichskettealigndrucklinks
{ \left\langle  L_P( \gamma'(0) ) ,  \gamma'(0)  \right\rangle
}
{ =} { \left\langle  \gamma^{\prime \prime} (0) ,  N(x_1 , \ldots , x_{n-1}, f(x_1  , \ldots , x_{n-1}))   \right\rangle
}
{ =} { { \frac{   1  }{  \sqrt{ 1 + \left(  { \frac{   \partial   f   }{  \partial   x_1   } }  \right)^2 + \cdots + \left(  { \frac{   \partial   f  }{  \partial   x_{n-1}   } }  \right)^2 }  } } \left\langle  \begin{pmatrix}  0  \\\vdots\\  0\\ \sum_{ 1 \leq i, j \leq n-1 }  { \frac{   \partial    }{  \partial   x_i   } }  { \frac{   \partial   f   }{  \partial   x_j   } } v_iv_j   \end{pmatrix}  ,  \begin{pmatrix}  -{ \frac{   \partial   f   }{  \partial   x_{1}   } }  \\\vdots\\  -{ \frac{   \partial   f   }{  \partial   x_{n-1}   } }\\ 1   \end{pmatrix}   \right\rangle
}
{ =} { { \frac{   \sum_{ 1 \leq i, j \leq n-1 } { \frac{   \partial    }{  \partial   x_i   } }  { \frac{   \partial   f   }{  \partial   x_j   } } v_i v_j     }{  \sqrt{ 1+  \left(  { \frac{   \partial   f   }{  \partial   x_1   } }  \right)^2 + \cdots + \left(  { \frac{   \partial   f  }{  \partial   x_{n-1}   } }  \right)^2 }  } }
}
{ =} { { \frac{   1  }{  \sqrt{ 1+ \left(  { \frac{   \partial   f   }{  \partial   x_1   } }  \right)^2 + \cdots + \left(  { \frac{   \partial   f   }{  \partial   x_{n-1}   } }  \right)^2 }  } }  \left( v_1   , \,   \ldots  , \,   v_{n-1}                 \right)  \operatorname{Hess}_{   } \,  f  \begin{pmatrix} v_1  \\ \vdots \\  v_{n-1}   \end{pmatrix}
}
}
{}{}{.}
Artinya, menurut
Teorema 4.7,
\definitionsverweis {bentuk bilinear}{}{}
simetris

\mathl{\left\langle  L_P(-) ,  -  \right\rangle}{} pada ruang singgung $T_PY$ bersesuaian dengan bentuk bilinear yang diberikan oleh matriks Hesse yang diskalakan
\zusatzklammer {oleh faktor di depan} {} {}
pada $\R^{n-1}$, apabila vektor yang sama disubstitusikan pada kedua tempat.
Menurut
Lemma 38.10 (Lineare Algebra (Osnabrück 2024-2025)),
kedua bentuk bilinear tersebut berimpit untuk sembarang pasangan argumen. Akan tetapi, ini menentukan bentuk fundamental kedua, bukan langsung matriks operator dalam basis grafik. Jika \mathl{A}{} adalah matriks $L_P$ dalam basis grafik dan \mathl{G=I_{n-1}+gg^{\mathsf T}}{} adalah matriks metriknya, maka matriks bentuk bilinear di atas adalah \mathl{GA}{.} Oleh karena itu,

\mathdisp {GA={ \frac{   1  }{  \sqrt{1+\lVert g\rVert^2}  } }\operatorname{Hess}_{   }\, f} { , }

sehingga

\mathdisp {A=G^{-1}{ \frac{   1  }{  \sqrt{1+\lVert g\rVert^2}  } }\operatorname{Hess}_{   }\, f} { . }

Catatan edisi: sumber menghilangkan faktor invers metrik G dan dengan demikian menyamakan matriks bentuk fundamental kedua dengan matriks pemetaan Weingarten. Rumus dan simpulan di atas memperbaiki pembedaan tersebut.

}

\chapter{Lembar Kerja 4}
\setcounter{section}{4}

  \zwischenueberschrift{Soal latihan}

\inputaufgabe
{}
{

Misalkan

\mavergleichskette
{\vergleichskette
{ h(x_1 , \ldots , x_n)
}
{ = }{ \sum_{i = 1}^n a_ix_i
}
{  }{
}
{    }{
}
{   }{
}
}
{}{}{}
adalah suatu
\definitionsverweis {bentuk linear}{}{}
$\neq 0$ dan

\mavergleichskette
{\vergleichskette
{ Y
}
{ = }{ h^{-1}(c)
}
{  }{
}
{    }{
}
{   }{
}
}
{}{}{}
adalah hipermuka yang bersesuaian. Buktikan bahwa
\definitionsverweis {pemetaan Weingarten}{}{}
merupakan pemetaan nol.

}
{} {}

\inputaufgabe
{}
{

Misalkan

\mavergleichskettedisp
{\vergleichskette
{ Y
}
{ =} { \left\{  \left(  x   , \,   y  , \,  z                 \right)   \mid   x^2+y^2 = z^2 , \,  \left(  x   , \,   y  , \,  z                 \right)  \neq  \left(  0   , \,   0  , \,   0                 \right)       \right\}
}
{ \subseteq} { W
}
{ =} { \R^3 \setminus \{  \left(  0   , \,   0  , \,   0                 \right) \}
}
{ } {
}
}
{}{}{}
adalah kerucut standar, dan misalkan

\mavergleichskette
{\vergleichskette
{ P
}
{ = }{ \left(  x   , \,   y  , \,  z                 \right)
}
{ \in }{ Y
}
{    }{
}
{   }{
}
}
{}{}{.}
Misalkan

\mavergleichskettedisp
{\vergleichskette
{ \gamma(t)
}
{ =} { \left(  x\cos t-y\sin t   , \,   x\sin t+y\cos t  , \,  z                 \right)
}
{ } {
}
{ } {
}
{ } {
}
}
{}{}{.}

Catatan edisi: kurva pada sumber tidak selalu melalui P dan tidak selalu tetap berada pada kerucut. Rumus di atas menggantinya dengan rotasi standar pada bidang xy.
Tentukan

\mathdisp {\operatorname{lim}_{   t  \rightarrow  0   } \,  { \frac{  N( \gamma(t) ) -  N( \gamma(0) )  }{ t } }} { , }

dengan $N$ medan normal satuan yang diperoleh dari gradien $x^2+y^2-z^2$.

}
{} {}

\inputaufgabe
{}
{

Misalkan

\mavergleichskettedisp
{\vergleichskette
{ Y
}
{ =} { \left\{  \left(  x   , \,   y  , \,  z                 \right)   \mid   x^2+y^2 =1         \right\}
}
{ \subseteq} { \R^3
}
{ } {
}
{ } {
}
}
{}{}{}
adalah silinder standar, dan misalkan

\mavergleichskette
{\vergleichskette
{ P
}
{ = }{ \left(  1   , \,   0  , \,  z                 \right)
}
{ \in }{ Y
}
{    }{
}
{   }{
}
}
{}{}{.}
Misalkan

\mavergleichskettedisp
{\vergleichskette
{ \gamma(t)
}
{ =} { \left(  \cos  t   , \,   \sin  t  , \,  z                 \right)
}
{ } {
}
{ } {
}
{ } {
}
}
{}{}{.}
Tentukan

\mathdisp {\operatorname{lim}_{   t  \rightarrow  0   } \,  { \frac{  N( \gamma(t) ) -  N( \gamma(0) )  }{ t } }} { , }

dengan $N$ medan normal satuan yang diperoleh dari gradien $x^2+y^2-1$.

}
{} {}

\inputaufgabe
{}
{

Misalkan

\mavergleichskettedisp
{\vergleichskette
{ Y
}
{ =} { \left\{  \left(  x   , \,   y  , \,  z                 \right)   \mid   x^2+y^2 =1         \right\}
}
{ \subseteq} { \R^3
}
{ } {
}
{ } {
}
}
{}{}{}
dan misalkan

\mavergleichskette
{\vergleichskette
{ P
}
{ \in }{ Y
}
{  }{
}
{    }{
}
{   }{}
}
{}{}{}
adalah suatu titik. Buktikan bahwa
\definitionsverweis {pemetaan Weingarten}{}{}
$L_P$ tidak bijektif, tetapi juga bukan pemetaan nol.

}
{} {}

\inputaufgabe
{}
{

Buktikan bahwa suatu
\definitionsverweis {permukaan terdiferensialkan}{}{}

\mavergleichskette
{\vergleichskette
{ Y
}
{ \subseteq }{ \R^3
}
{  }{
}
{    }{
}
{   }{
}
}
{}{}{}
dapat memuat suatu garis

\mavergleichskette
{\vergleichskette
{ G
}
{ \subseteq }{ Y
}
{  }{
}
{    }{
}
{   }{
}
}
{}{}{}
yang vektor arahnya, jika dipandang sebagai vektor di
\definitionsverweis {ruang singgung}{}{}
$T_PY$ untuk suatu titik

\mavergleichskette
{\vergleichskette
{ P
}
{ \in }{ G
}
{  }{
}
{    }{
}
{   }{
}
}
{}{}{,}
tidak harus termasuk dalam
\definitionsverweis {kernel}{}{}
dari
\definitionsverweis {pemetaan Weingarten}{}{}
$L_P$.

}
{} {}

\inputaufgabe
{}
{

Misalkan

\mavergleichskette
{\vergleichskette
{ W
}
{ \subseteq }{ \R^n
}
{  }{
}
{    }{
}
{   }{
}
}
{}{}{}
\definitionsverweis {terbuka}{}{,}
\maabb {h} { W } { \R
} {}
adalah suatu
\definitionsverweis {fungsi yang terdiferensialkan secara kontinu sebanyak dua kali}{}{}
dan

\mavergleichskette
{\vergleichskette
{ Y
}
{ = }{ h^{-1}(c)
}
{  }{
}
{    }{
}
{   }{
}
}
{}{}{}
adalah
\definitionsverweis {serat}{}{}
di atas

\mavergleichskette
{\vergleichskette
{ c
}
{ \in }{ \R
}
{  }{
}
{    }{
}
{   }{
}
}
{}{}{,}
dengan $h$
\definitionsverweis {reguler}{}{}
di setiap titik pada $Y$.
Catatan edisi: sumber hanya mengasumsikan h kelas C1. Di sini hipotesis diperkuat menjadi h kelas C2 agar medan normal gradien dan pemetaan Weingarten yang ditanyakan terdefinisi sebagaimana pada kuliah.
Kita juga memandang $h$ sebagai pemetaan $\tilde{h}$ pada
\mathl{W \times \R^m}{}{,}
dengan $m$ variabel terakhir tidak muncul secara eksplisit dalam $\tilde{h}$. Misalkan

\mavergleichskettedisp
{\vergleichskette
{ Z
}
{ =} { \tilde{h}^{-1}(c)
}
{ =} { Y \times \R^m
}
{ \subseteq} { W \times \R^{m}
}
{ } {
}
}
{}{}{.}
Misalkan

\mavergleichskette
{\vergleichskette
{ Q
}
{ \in }{ Z
}
{  }{
}
{    }{
}
{   }{
}
}
{}{}{}
adalah suatu titik yang terletak di atas

\mavergleichskette
{\vergleichskette
{ P
}
{ \in }{ Y
}
{  }{
}
{    }{
}
{   }{
}
}
{}{}{.}
Bagaimana hubungan antara
\definitionsverweis {pemetaan-pemetaan Weingarten}{}{}
\mathkor {} {L_Q} {dan} {L_P} {?}

}
{} {}

\inputaufgabegibtloesung
{}
{

Misalkan
\mathkor {} {r} {dan} {R} {}
adalah bilangan real positif dengan

\mavergleichskette
{\vergleichskette
{ 0
}
{ < }{ r
}
{ < }{ R
}
{    }{
}
{   }{
}
}
{}{}{,}
dan kita tinjau torus
\zusatzklammer {tertanam} {} {}

\mavergleichskettedisp
{\vergleichskette
{ T
}
{ =} { \left\{ (x,y,z) \in \R^3  \mid   \left(  \sqrt{x^2+y^2} -R  \right)^2 +z^2  = r^2         \right\}
}
{ } {
}
{ } {
}
{ } {
}
}
{}{}{.}
\aufzaehlungdrei{Tentukan suatu
\definitionsverweis {medan normal satuan}{}{}
$N$ pada $T$.
}{Pada titik
\mathl{\left( R-r  , \,   0  , \,   0                 \right)}{} dan untuk vektor

\mavergleichskette
{\vergleichskette
{ v
}
{ = }{ \begin{pmatrix}  0  \\ 1 \\  0   \end{pmatrix}
}
{  }{
}
{    }{
}
{   }{
}
}
{}{}{}
tentukan
\definitionsverweis {turunan berarah}{}{}

\mathl{D_vN}{.}
}{Pada titik
\mathl{\left( R-r  , \,   0  , \,   0                 \right)}{} dan untuk vektor

\mavergleichskette
{\vergleichskette
{ w
}
{ = }{ \begin{pmatrix}  0  \\ 0\\  1   \end{pmatrix}
}
{  }{
}
{    }{
}
{   }{
}
}
{}{}{}
tentukan
\definitionsverweis {turunan berarah}{}{}

\mathl{D_w N}{.}
}

}
{} {}

\inputaufgabe
{}
{

Untuk grafik $Y$ dari fungsi

\mavergleichskettedisp
{\vergleichskette
{ f(x,y)
}
{ =} { xy
}
{ } {
}
{ } {
}
{ } {
}
}
{}{}{}
tentukan
\definitionsverweis {pemetaan Weingarten}{}{}
$L_P$ pada setiap titik

\mavergleichskettedisp
{\vergleichskette
{ P
}
{ =} { (x,y,f(x,y) )
}
{ } {
}
{ } {
}
{ } {
}
}
{}{}{.}
Tentukan nilai eigen, vektor eigen, dan suatu matriks diagonal untuk $L_P$.

}
{} {}

\inputaufgabe
{}
{

Misalkan

\mavergleichskette
{\vergleichskette
{ W
}
{ \subseteq }{ \R^n
}
{  }{
}
{    }{
}
{   }{
}
}
{}{}{}
\definitionsverweis {terbuka}{}{,}
\maabb {h} { W } { \R
} {}
adalah suatu
\definitionsverweis {fungsi yang terdiferensialkan secara kontinu sebanyak dua kali}{}{}
dan

\mavergleichskette
{\vergleichskette
{ Y
}
{ = }{ h^{-1}(c)
}
{  }{
}
{    }{
}
{   }{
}
}
{}{}{}
adalah
\definitionsverweis {serat}{}{}
di atas

\mavergleichskette
{\vergleichskette
{ c
}
{ \in }{ \R
}
{  }{
}
{    }{
}
{   }{
}
}
{}{}{,}
dengan $h$
\definitionsverweis {reguler}{}{}
di setiap titik pada $Y$.
Catatan edisi: sumber hanya mengasumsikan h kelas C1. Di sini hipotesis diperkuat menjadi h kelas C2 agar pemetaan Weingarten yang dibandingkan terdefinisi sebagaimana pada kuliah.
Misalkan
\maabbdisp {L} {\R^n } { \R^n
} {}
adalah suatu
\definitionsverweis {isometri linear}{}{,}

\mavergleichskette
{\vergleichskette
{ W'
}
{ = }{ L^{-1}(W)
}
{  }{
}
{    }{
}
{   }{
}
}
{}{}{}
dan

\mavergleichskettedisp
{\vergleichskette
{ Z
}
{ =} { L^{-1}(Y)
}
{ =} { (h \circ L)^{-1}(c)
}
{ } {
}
{ } {
}
}
{}{}{.}
Buktikan bahwa
\definitionsverweis {pemetaan Weingarten}{}{}
$L_P$ untuk

\mavergleichskette
{\vergleichskette
{ P
}
{ \in }{ Y
}
{  }{
}
{    }{
}
{   }{
}
}
{}{}{}
bersesuaian dengan pemetaan Weingarten untuk

\mavergleichskette
{\vergleichskette
{ Q
}
{ = }{ L^{-1}(P)
}
{ \in }{ Z
}
{    }{
}
{   }{
}
}
{}{}{,}
jika melalui $L$ ruang singgung $T_PY$ dan $T_QZ$ dihubungkan satu sama lain.

}
{} {}

\inputaufgabegibtloesung
{}
{

Untuk permukaan $Y$ yang diberikan oleh

\mavergleichskettedisp
{\vergleichskette
{ x^2+y^2+2z^2
}
{ =} { 4
}
{ } {
}
{ } {
}
{ } {
}
}
{}{}{}
dan titik

\mavergleichskettedisp
{\vergleichskette
{ P
}
{ =} { \left(  1  , \,   1  , \,   1                 \right)
}
{ \in} { Y
}
{ } {
}
{ } {
}
}
{}{}{}
tentukan suatu matriks diagonal untuk
\definitionsverweis {pemetaan Weingarten}{}{}
$L_P$.

}
{} {}

\inputaufgabe
{}
{

Misalkan

\mavergleichskette
{\vergleichskette
{ W
}
{ \subseteq }{ \R^n
}
{  }{
}
{    }{
}
{   }{
}
}
{}{}{}
\definitionsverweis {terbuka}{}{,}
\maabb {h} { W } { \R
} {}
adalah suatu
\definitionsverweis {fungsi yang terdiferensialkan secara kontinu sebanyak dua kali}{}{}
dan

\mavergleichskette
{\vergleichskette
{ Y
}
{ = }{ h^{-1}(c)
}
{  }{
}
{    }{
}
{   }{
}
}
{}{}{}
adalah
\definitionsverweis {serat}{}{}
di atas

\mavergleichskette
{\vergleichskette
{ c
}
{ \in }{ \R
}
{  }{
}
{    }{
}
{   }{
}
}
{}{}{,}
dengan $h$
\definitionsverweis {reguler}{}{}
di setiap titik pada $Y$.
Bagaimana
\definitionsverweis {pemetaan Weingarten}{}{}
dan
\definitionsverweis {ruang-ruang eigennya}{}{}
berubah jika kita beralih dari suatu
\definitionsverweis {orientasi}{}{}
pada $Y$ ke orientasi yang berlawanan?

}
{} {}

  \zwischenueberschrift{Soal untuk dikumpulkan}

\inputaufgabe
{4}
{

Dalam
Contoh 4.5,
tentukan matriks representasi
\definitionsverweis {pemetaan Weingarten}{}{}
untuk titik-titik berbentuk
\mathl{\begin{pmatrix}  0  \\ y \\ z   \end{pmatrix}}{} terhadap basis
\mathkor {} {\begin{pmatrix}  y  \\ 0 \\  0   \end{pmatrix}} {dan} {\begin{pmatrix}  0  \\ z \\  y   \end{pmatrix}} {}
dari ruang singgung. Tentukan nilai-nilai eigen dan suatu basis yang terdiri atas vektor-vektor eigen.

}
{} {}

\inputaufgabe
{5}
{

Untuk permukaan $Y$ yang diberikan oleh

\mavergleichskettedisp
{\vergleichskette
{ x^2+2y^2+5z^2
}
{ =} { 11
}
{ } {
}
{ } {
}
{ } {
}
}
{}{}{}
dan titik

\mavergleichskettedisp
{\vergleichskette
{ P
}
{ =} { \left(  2  , \,   1  , \,   1                 \right)
}
{ \in} { Y
}
{ } {
}
{ } {
}
}
{}{}{}
tentukan suatu matriks diagonal untuk
\definitionsverweis {pemetaan Weingarten}{}{}
$L_P$.

}
{} {}

\inputaufgabe
{4}
{

Untuk hipermuka yang diberikan oleh

\mavergleichskettedisp
{\vergleichskette
{ x^2+2y^2+3z^2+4w^2
}
{ =} { 10
}
{ } {
}
{ } {
}
{ } {
}
}
{}{}{}

\mavergleichskette
{\vergleichskette
{ Y
}
{ \subseteq }{ \R^4
}
{  }{
}
{    }{
}
{   }{
}
}
{}{}{}
dan titik

\mavergleichskettedisp
{\vergleichskette
{ P
}
{ =} { \left(  1  , \,   1  , \,   1 , \,   1                \right)
}
{ \in} { Y
}
{ } {
}
{ } {
}
}
{}{}{}
tentukan suatu matriks untuk
\definitionsverweis {pemetaan Weingarten}{}{}
$L_P$ terhadap suatu
\definitionsverweis {basis}{}{}
yang sesuai dari
\definitionsverweis {ruang singgung}{}{}
$T_PY$.

}
{} {}

\inputaufgabe
{6}
{

Untuk grafik $Y$ dari fungsi

\mavergleichskettedisp
{\vergleichskette
{ f(x,y)
}
{ =} { x^3y^5
}
{ } {
}
{ } {
}
{ } {
}
}
{}{}{}
tentukan
\definitionsverweis {pemetaan Weingarten}{}{}
$L_P$ pada titik-titik

\mavergleichskette
{\vergleichskette
{ P_1
}
{ = }{ (0,0,0)
}
{  }{
}
{    }{
}
{   }{
}
}
{}{}{,}

\mavergleichskette
{\vergleichskette
{ P_2
}
{ = }{ (0,1,0)
}
{  }{
}
{    }{
}
{   }{
}
}
{}{}{,}

\mavergleichskette
{\vergleichskette
{ P_3
}
{ = }{ (2,1,8)
}
{  }{
}
{    }{
}
{   }{
}
}
{}{}{.}
Untuk masing-masing titik, tentukan nilai eigen, vektor eigen, dan suatu matriks diagonal untuk $L_P$.

}
{} {}

\section*{Solusi yang disediakan oleh sumber}
\addcontentsline{toc}{section}{Solusi yang disediakan oleh sumber}
\subsection*{Solusi Soal 4.7}
\textit{Catatan edisi: solusi sumber memuat rumus turunan yang rusak dan berhenti sebelum menghitung kedua turunan berarah. Perhitungan berikut memperbaiki dan melengkapinya.}

\aufzaehlungdrei{Tuliskan

\mathdisp {\rho=\sqrt{x^2+y^2}} { }

dan

\mavergleichskettedisp
{\vergleichskette
{ A(x,y,z)
}
{ =} { \begin{pmatrix} x-xR\rho^{-1} \\ y-yR\rho^{-1} \\ z \end{pmatrix}
}
{ } {
}
{ } {
}
{ } {
}
}
{}{}{.}
Dengan demikian,

\mavergleichskettedisp
{\vergleichskette
{ h(x,y,z)
}
{ =} { x^2+y^2-2R\sqrt{x^2+y^2}+z^2+R^2
}
{ } {
}
{ } {
}
{ } {
}
}
{}{}{}
dan

\mavergleichskettedisp
{\vergleichskette
{ T
}
{ =} { h^{-1}(r^2)
}
{  }{
}
{    }{
}
{   }{
}
}
{}{}{.}
Gradien $h$ adalah $2A$. Pada $T$ berlaku $\lVert A\rVert=r$, sehingga suatu
\definitionsverweis {medan normal satuan}{}{}
adalah

\mavergleichskettedisp
{\vergleichskette
{ N(x,y,z)
}
{ =} { \frac{1}{r}\begin{pmatrix} x-xR\rho^{-1} \\ y-yR\rho^{-1} \\ z \end{pmatrix}
}
{ } {
}
{ } {
}
{ } {
}
}
{}{}{.}
}{Ambil

\mathdisp {P=(R-r,0,0)} { . }

Untuk menghitung turunan, kita dapat memakai perpanjangan gradien ternormalisasi $\widehat N=A/\lVert A\rVert$ pada suatu lingkungan $P$. Di titik ini,

\mavergleichskettedisp
{\vergleichskette
{ A(P)
}
{ =} { \begin{pmatrix}-r\\0\\0\end{pmatrix}
}
{ ,\quad }{ D_vA(P)
}
{ =} { \begin{pmatrix}0\\-\frac{r}{R-r}\\0\end{pmatrix}
}
{ } {
}
}
{}{}{,}
dan $\langle A(P),D_vA(P)\rangle=0$. Oleh karena itu,

\mavergleichskettedisp
{\vergleichskette
{ D_vN(P)
}
{ =} { D_v\widehat N(P)
}
{ =} { \frac{1}{r}D_vA(P)
}
{ =} { \begin{pmatrix}0\\-\frac{1}{R-r}\\0\end{pmatrix}
}
}
{}{}{.}
}{Demikian pula,

\mavergleichskettedisp
{\vergleichskette
{ D_wA(P)
}
{ =} { \begin{pmatrix}0\\0\\1\end{pmatrix}
}
{ ,\quad }{ \langle A(P),D_wA(P)\rangle
}
{ =} { 0
}
{ } {
}
}
{}{}{.}
Jadi,

\mavergleichskettedisp
{\vergleichskette
{ D_wN(P)
}
{ =} { D_w\widehat N(P)
}
{ =} { \frac{1}{r}D_wA(P)
}
{ =} { \begin{pmatrix}0\\0\\\frac{1}{r}\end{pmatrix}
}
}
{}{}{.}
}
\subsection*{Solusi Soal 4.10}
\textit{Catatan edisi: solusi sumber menghitung diferensial medan normal tetapi tidak menerapkan tanda minus dalam definisi pemetaan Weingarten. Solusi itu juga menyebut vektor basis kedua sebagai vektor eigen meskipun hasil yang ditampilkannya masih memiliki komponen pada vektor basis pertama. Tanda dan basis eigen diperbaiki di bawah ini.}

Medan gradien ternormalisasi adalah

\mavergleichskettealign
{\vergleichskettealign
{ N(x,y,z)
}
{ =} { { \frac{   1  }{  \sqrt{ 4x^2+4y^2+16 z^2}  } } \begin{pmatrix}  2x \\ 2y\\  4z   \end{pmatrix}
}
{ =} { { \frac{   1  }{  \sqrt{ x^2+y^2+4 z^2}  } } \begin{pmatrix}  x  \\ y \\  2z   \end{pmatrix}
}
{ } {
}
{ } {
}
}
{}
{}{.}
Kita bekerja dengan medan normal satuan

\mavergleichskettedisp
{\vergleichskette
{ M(x,y,z)
}
{ =} { { \frac{   1  }{  \sqrt{ 4+2 z^2}  } } \begin{pmatrix}  x  \\ y \\  2z   \end{pmatrix}
}
{ } {
}
{ } {
}
{ } {
}
}
{}{}{,}
yang berimpit dengan $N$ pada $Y$. Diferensial total $M$ adalah

\mathdisp {\begin{pmatrix}  { \frac{   1  }{  \sqrt{4+2z^2}  } }   &   0 &  { \frac{   -2 xz }{  \left(  4+2z^2  \right)^{3/2}  } }  \\  0  &  { \frac{   1  }{  \sqrt{4+2z^2}  } }   &  { \frac{   -2 yz }{  \left(  4+2z^2  \right)^{3/2}  } }   \\ 0   &  0  &  { \frac{   8 }{  \left(  4+2z^2  \right)^{3/2}  } }  \end{pmatrix}} { . }

Di titik $P$ yang diberikan, gradiennya adalah $\begin{pmatrix}  2  \\ 2\\  4   \end{pmatrix}$. Kedua vektor

\mavergleichskettedisp
{\vergleichskette
{ e_1
}
{ =} { \begin{pmatrix}1\\-1\\0\end{pmatrix}
}
{ ,\quad }{ e_2
}
{ =} { \begin{pmatrix}1\\1\\-1\end{pmatrix}
}
{ } {
}
}
{}{}{}
membentuk basis ortogonal ruang singgung $T_PY$. Diferensial total $M$ di titik tersebut adalah

\mathdisp {\begin{pmatrix}  { \frac{   1  }{  \sqrt{6}  } }   &   0 &  { \frac{   -1 }{  3 \sqrt{6}  } }  \\  0  &  { \frac{   1  }{  \sqrt{6}  } }   &  { \frac{   -1  }{  3 \sqrt{6}  } }   \\ 0   &  0  &  { \frac{   4  }{  3 \sqrt{6}  } }   \end{pmatrix}} { . }

Dengan menerapkannya pada kedua vektor basis, diperoleh

\mavergleichskettedisp
{\vergleichskette
{ (DM)_P(e_1)
}
{ =} { \frac{1}{\sqrt 6}e_1
}
{ ,\quad }{ (DM)_P(e_2)
}
{ =} { \frac{4}{3\sqrt 6}e_2
}
{ } {
}
}
{}{}{.}
Karena pemetaan Weingarten didefinisikan oleh $L_P=-(DM)_P|_{T_PY}$, maka

\mavergleichskettedisp
{\vergleichskette
{ L_P(e_1)
}
{ =} { -\frac{1}{\sqrt 6}e_1
}
{ ,\quad }{ L_P(e_2)
}
{ =} { -\frac{4}{3\sqrt 6}e_2
}
{ } {
}
}
{}{}{.}
Jadi, terhadap basis eigen $(e_1,e_2)$, suatu matriks diagonal yang merepresentasikan $L_P$ adalah

\mathdisp {\begin{pmatrix}  -{ \frac{   1  }{  \sqrt{6}  } }   &   0  \\  0  &  -{ \frac{   4  }{  3 \sqrt{6}  } }  \end{pmatrix}} { . }

\part{Unit 5}
\chapter[Kuliah 5]{Kuliah 5: Kelengkungan Utama}
\setcounter{section}{5}

  \zwischenueberschrift{Kelengkungan Utama}

Misalkan

\mavergleichskette
{\vergleichskette
{ W
}
{ \subseteq }{ \R^n
}
{  }{
}
{    }{
}
{   }{
}
}
{}{}{}
\definitionsverweis {terbuka}{}{,}
\maabb {h} { W } { \R
} {}
suatu
\definitionsverweis {fungsi yang dua kali terdiferensialkan secara kontinu}{}{}
dan

\mavergleichskette
{\vergleichskette
{ Y
}
{ = }{ h^{-1}(c)
}
{  }{
}
{    }{
}
{   }{
}
}
{}{}{}
adalah
\definitionsverweis {serat}{}{}
di atas

\mavergleichskette
{\vergleichskette
{ c
}
{ \in }{ \R
}
{  }{
}
{    }{
}
{   }{
}
}
{}{}{,}
dengan $h$
\definitionsverweis {reguler}{}{}
pada setiap titik di $Y$. Misalkan

\mavergleichskette
{\vergleichskette
{ P
}
{ \in }{ Y
}
{  }{
}
{    }{
}
{   }{
}
}
{}{}{.}
Kita mendefinisikan berbagai konsep kelengkungan dengan mengacu pada
\definitionsverweis {pemetaan Weingarten}{}{}
\maabbdisp {L_P} {T_PY} { T_PY
} {.}

Catatan edisi: sumber memakai hipotesis $C^1$ dalam pengantar dan dua definisi berikut, meskipun pemetaan Weingarten mendiferensialkan medan normal. Ketiga hipotesis diperkuat menjadi $C^2$ agar objek yang didefinisikan benar-benar tersedia.

\inputdefinition
{}
{

Misalkan

\mavergleichskette
{\vergleichskette
{ W
}
{ \subseteq }{ \R^n
}
{  }{
}
{    }{
}
{   }{
}
}
{}{}{}
\definitionsverweis {terbuka}{}{,}
\maabb {h} { W } { \R
} {}
suatu
\definitionsverweis {fungsi yang dua kali terdiferensialkan secara kontinu}{}{}
dan

\mavergleichskette
{\vergleichskette
{ Y
}
{ = }{ h^{-1}(c)
}
{  }{
}
{    }{
}
{   }{
}
}
{}{}{}
adalah
\definitionsverweis {serat}{}{}
di atas

\mavergleichskette
{\vergleichskette
{ c
}
{ \in }{ \R
}
{  }{
}
{    }{
}
{   }{
}
}
{}{}{,}
dengan $h$
\definitionsverweis {reguler}{}{}
pada setiap titik di $Y$. Misalkan

\mavergleichskette
{\vergleichskette
{ P
}
{ \in }{ Y
}
{  }{
}
{    }{
}
{   }{
}
}
{}{}{.}
Setiap
\definitionsverweis {nilai eigen}{}{}
dari
\definitionsverweis {pemetaan Weingarten}{}{}
\maabbeledisp {L_P} {T_P Y} { T_PY
} { v } { - \left(  D_{v}  N   \right) \left(   P  \right)
} {,}
disebut suatu
\definitionswort {kelengkungan utama}{}
dari $Y$ di $P$.

}

\inputdefinition
{}
{

Misalkan

\mavergleichskette
{\vergleichskette
{ W
}
{ \subseteq }{ \R^n
}
{  }{
}
{    }{
}
{   }{
}
}
{}{}{}
\definitionsverweis {terbuka}{}{,}
\maabb {h} { W } { \R
} {}
suatu
\definitionsverweis {fungsi yang dua kali terdiferensialkan secara kontinu}{}{}
dan

\mavergleichskette
{\vergleichskette
{ Y
}
{ = }{ h^{-1}(c)
}
{  }{
}
{    }{
}
{   }{
}
}
{}{}{}
adalah
\definitionsverweis {serat}{}{}
di atas

\mavergleichskette
{\vergleichskette
{ c
}
{ \in }{ \R
}
{  }{
}
{    }{
}
{   }{
}
}
{}{}{,}
dengan $h$
\definitionsverweis {reguler}{}{}
pada setiap titik di $Y$. Misalkan

\mavergleichskette
{\vergleichskette
{ P
}
{ \in }{ Y
}
{  }{
}
{    }{
}
{   }{
}
}
{}{}{.}
Setiap
\definitionsverweis {vektor eigen}{}{}
dari
\definitionsverweis {pemetaan Weingarten}{}{}
\maabbeledisp {L_P} {T_P Y} { T_PY
} { v } { - \left(  D_{v}  N   \right) \left(   P  \right)
} {,}
disebut suatu
\definitionswort {arah kelengkungan utama}{}
dari $Y$ di $P$.

}

Istilah \stichwort {vektor kelengkungan utama} {} juga digunakan sebagai pengganti arah kelengkungan utama; istilah arah kelengkungan utama khususnya digunakan untuk vektor eigen dengan norma $1$. Berdasarkan
Korolari 4.8,
pemetaan Weingarten memiliki suatu basis ortonormal yang terdiri atas vektor-vektor eigen, yakni arah-arah kelengkungan utama. Pada umumnya, dipilih suatu basis ortonormal yang terdiri atas arah-arah kelengkungan utama. Kelengkungan-kelengkungan utama merupakan bilangan real yang sering diurutkan menurut besarnya dan ditulis sebagai

\mavergleichskettedisp
{\vergleichskette
{ \kappa_1
}
{ \leq} { \kappa_2
}
{ \leq   \ldots \leq} { \kappa_{n-1}
}
{ } {
}
{ } {
}
}
{}{}{}
dan disebut \stichwort {tupel kelengkungan utama} {.} Setiap $\kappa$ dicantumkan sebanyak dimensi ruang eigen yang bersesuaian, yaitu sebanyak
\definitionsverweis {multiplisitas geometrik}{}{nya.}

Menurut
Teorema 23.2 (Aljabar Linear (Osnabrück 2024-2025)),
kelengkungan-kelengkungan utama diperoleh dengan menentukan akar-akar
\definitionsverweis {polinom karakteristik}{}{}
dari
\definitionsverweis {pemetaan Weingarten}{}{,}
sedangkan
\definitionsverweis {ruang-ruang eigen}{}{}
yang bersesuaian diperoleh menggunakan
Lemma 22.1 (Aljabar Linear (Osnabrück 2024-2025)).

\inputbeispiel{}
{

Misalkan
\mathl{a_1 , \ldots , a_{n-1}}{} bilangan-bilangan real. Kita meninjau fungsi
\maabbeledisp {f} {\R^{n-1}} {\R
} { \left(  x_1   , \,   \ldots  , \,   x_{n-1}                 \right) } { a_1x_1^2  + \cdots + a_{n-1} x_{n-1}^2
} {.}
Menurut
Lemma 4.9,
\definitionsverweis {pemetaan Weingarten}{}{}
dari
\definitionsverweis {grafik}{}{}
$f$ di titik nol dijelaskan oleh
\definitionsverweis {matriks Hesse}{}{}
dari $f$, yaitu oleh

\mathdisp {\begin{pmatrix}  2a_1  & 0 & 0 & 0 \\ 0 &  2a_2  & 0  & 0 \\ 0 & 0 &   \ddots & 0 \\ 0 & 0 & 0 &  2a_{n-1}  \end{pmatrix}} { . }

Jadi,
\definitionsverweis {kelengkungan-kelengkungan utama}{}{}
grafik tersebut di titik nol adalah $2a_1  , \ldots , 2a_{n-1}$ dan
\definitionsverweis {arah-arah kelengkungan utama}{}{}
diberikan oleh vektor-vektor standar. Secara khusus, setiap tupel bilangan real muncul sebagai tupel kelengkungan utama suatu hipermuka terdiferensialkan.

}

  \zwischenueberschrift{Kelengkungan pada Permukaan}

Sekarang kita menelaah lebih saksama konsep-konsep kelengkungan pada permukaan terdiferensialkan di ruang.

\inputdefinition
{}
{

Misalkan

\mavergleichskette
{\vergleichskette
{ W
}
{ \subseteq }{ \R^3
}
{  }{
}
{    }{
}
{   }{
}
}
{}{}{}
\definitionsverweis {terbuka}{}{,}
\maabb {h} { W } { \R
} {}
suatu
\definitionsverweis {fungsi yang dua kali terdiferensialkan secara kontinu}{}{}
dan

\mavergleichskette
{\vergleichskette
{ Y
}
{ = }{ h^{-1}(c)
}
{  }{
}
{    }{
}
{   }{
}
}
{}{}{}
adalah serat di atas

\mavergleichskette
{\vergleichskette
{ c
}
{ \in }{ \R
}
{  }{
}
{    }{
}
{   }{
}
}
{}{}{,}
dengan $h$
\definitionsverweis {reguler}{}{}
pada setiap titik di $Y$, dan misalkan

\mavergleichskette
{\vergleichskette
{ P
}
{ \in }{ Y
}
{  }{
}
{    }{
}
{   }{
}
}
{}{}{.}
\definitionsverweis {Rata-rata aritmetika}{}{}

\mathl{{ \frac{   \kappa_1+\kappa_2  }{  2 } }}{} dari kedua
\definitionsverweis {kelengkungan utama}{}{}
di $P$ disebut
\definitionswort {kelengkungan rata-rata}{}
permukaan $Y$ di $P$.

}

\inputdefinition
{}
{

Misalkan

\mavergleichskette
{\vergleichskette
{ W
}
{ \subseteq }{ \R^3
}
{  }{
}
{    }{
}
{   }{
}
}
{}{}{}
\definitionsverweis {terbuka}{}{,}
\maabb {h} { W } { \R
} {}
suatu
\definitionsverweis {fungsi yang dua kali terdiferensialkan secara kontinu}{}{}
dan

\mavergleichskette
{\vergleichskette
{ Y
}
{ = }{ h^{-1}(c)
}
{  }{
}
{    }{
}
{   }{
}
}
{}{}{}
adalah serat di atas

\mavergleichskette
{\vergleichskette
{ c
}
{ \in }{ \R
}
{  }{
}
{    }{
}
{   }{
}
}
{}{}{,}
dengan $h$
\definitionsverweis {reguler}{}{}
pada setiap titik di $Y$, dan misalkan

\mavergleichskette
{\vergleichskette
{ P
}
{ \in }{ Y
}
{  }{
}
{    }{
}
{   }{
}
}
{}{}{.}
Hasil kali
\mathl{\kappa_1 \cdot \kappa_2}{} dari kedua
\definitionsverweis {kelengkungan utama}{}{}
di $P$ disebut
\definitionswort {kelengkungan Gauss}{}
permukaan $Y$ di $P$.

}

Kelengkungan Gauss adalah
\definitionsverweis {determinan}{}{}
dari pemetaan Weingarten.

\inputbeispiel{}
{

Pada permukaan bola berjari-jari $r$,
\definitionsverweis {pemetaan Weingarten}{}{}
\zusatzklammer {terhadap medan normal satuan yang mengarah ke dalam} {} {}
di setiap titik, menurut
Contoh 4.4,
adalah homoteti skalar dengan faktor ${ \frac{   1  }{ r } }$. Oleh karena itu, ${ \frac{   1  }{ r } }$ merupakan satu-satunya
\definitionsverweis {kelengkungan utama}{}{}
dengan multiplisitas $2$, dan
\definitionsverweis {kelengkungan Gauss}{}{}
adalah
\mathl{{ \frac{   1  }{ r^2 } }}{.}

}

\inputbeispiel{}
{

Untuk hiperboloid satu lembar

\mavergleichskettedisp
{\vergleichskette
{ Y
}
{ =} { \left\{ (x,y,z)  \mid   x^2+y^2-z^2  = 1         \right\}
}
{ \subseteq} { \R^3
}
{ } {
}
{ } {
}
}
{}{}{}
berdasarkan perhitungan dalam
Contoh 4.5,
pada suatu titik

\mavergleichskette
{\vergleichskette
{ (x,y,z)
}
{ \in }{ Y
}
{  }{
}
{    }{
}
{   }{
}
}
{}{}{}
kedua
\definitionsverweis {kelengkungan utama}{}{}
adalah
\mathkor {} {- \left(  x^2+y^2+z^2  \right)^{- { \frac{   1  }{  2 } } }} {dan} {\left(  x^2+y^2+z^2  \right)^{- { \frac{   3  }{  2 } } }} {,}
sehingga
\definitionsverweis {kelengkungan Gauss}{}{}
adalah

\mavergleichskettedisphandlinks
{\vergleichskettedisphandlinks
{ - \left(  x^2+y^2+z^2  \right)^{- { \frac{   1  }{  2 } } }  \cdot \left(  x^2+y^2+z^2  \right)^{- { \frac{   3  }{  2 } } }
}
{ =} { - \left(  x^2+y^2+z^2  \right)^{- 2 }
}
{ =} { -  { \frac{   1  }{  \left(  x^2+y^2+z^2  \right)^{ 2 }   } }
}
{ } {
}
{ } {
}
}
{}{}{.}

}

\inputbeispiel{}
{

Misalkan $Y$ adalah
\definitionsverweis {grafik}{}{}
suatu fungsi yang dua kali terdiferensialkan secara kontinu
\maabbele {f} { V } { \R
} {(x,y)} { f(x,y)
} {,}
pada suatu himpunan terbuka

\mavergleichskette
{\vergleichskette
{ V
}
{ \subseteq }{ \R^2
}
{  }{
}
{    }{
}
{   }{
}
}
{}{}{.}
Menurut
Lemma 4.9,
\definitionsverweis {pemetaan Weingarten}{}{}
\zusatzklammer {terhadap medan normal satuan yang mengarah ke atas} {} {}
di setiap titik

\mavergleichskette
{\vergleichskette
{ P
}
{ = }{ (Q,f(Q))
}
{  }{
}
{    }{
}
{   }{
}
}
{}{}{}
memiliki bentuk berikut. Tuliskan
\mathl{g=\operatorname{Grad} f(Q)}{,}
\mathl{G=I_2+gg^{\mathsf T}}{,}
\mathl{H=\operatorname{Hess} f(Q)}{,}
dan
\mathl{\omega=\sqrt{1+\lVert g\rVert^2}}{.}
Dalam basis koordinat grafik, $G$ adalah matriks metrik dan matriks $A$ dari pemetaan Weingarten adalah

\mathdisp {A=G^{-1}{ \frac{ 1 }{ \omega } }H} { . }

Jadi, koordinat $u$ dari suatu arah kelengkungan utama memenuhi masalah nilai eigen umum

\mathdisp {Hu=\kappa\,\omega Gu} { . }

Hanya apabila $g=0$ masalah ini mereduksi menjadi masalah nilai eigen biasa untuk matriks Hesse. Untuk permukaan ini,
\definitionsverweis {kelengkungan Gauss}{}{}
adalah

\mathdisp {K={ \frac{ \det H }{ \omega^4 } }={ \frac{ \det\!\left(\operatorname{Hess}f(Q)\right) }{ \left(1+\lVert\operatorname{Grad}f(Q)\rVert^2\right)^2 } }} { . }

Catatan edisi: sumber menunjuk contoh bola yang tidak terkait dan menghilangkan faktor invers metrik $G^{-1}$. Akibatnya, pernyataan sumber tentang arah utama dan penyebut kelengkungan Gauss tidak benar secara umum; pernyataan itu tereduksi ke rumus yang benar pada titik kritis $f$. Rumus di atas menerapkan Lemma 4.9 yang telah diperbaiki pada Unit 4.

}

\inputbeispiel{}
{

Misalkan $I$ suatu
\definitionsverweis {interval terbuka}{}{}
dan
\maabb {f} { I } { \R_+
} {}
suatu
\definitionsverweis {fungsi yang dua kali terdiferensialkan secara kontinu}{}{,}
misalkan $Y$ adalah
\definitionsverweis {permukaan putar}{}{}
yang bersesuaian
\zusatzklammer {mengelilingi sumbu $x$} {} {}
dengan
\definitionsverweis {grafik}{}{}
$f$, dan misalkan

\mavergleichskettedisp
{\vergleichskette
{ P
}
{ =} { \begin{pmatrix}  x_0  \\f(x_0)\\  0   \end{pmatrix}
}
{ \in} { Y
}
{ } {
}
{ } {
}
}
{}{}{}
\zusatzklammer {yang bukan merupakan pembatasan berarti} {} {.}
Kita meninjau bagian permukaan putar yang, pada suatu lingkungan $(x_0,0)$ yang termuat dalam
\mathl{\left\{ (x,z)   \mid   z^2 < f(x)^2        \right\}=\left\{(x,z)\mid |z|<f(x)\right\}}{,} dapat ditulis sebagai grafik dari

\mavergleichskettedisp
{\vergleichskette
{ g(x,z)
}
{ =} { \sqrt{ f(x)^2-z^2 }
}
{ =} { \left(  f(x)^2-z^2   \right)^{1/2}
}
{ } {
}
{ } {
}
}
{}{}{.}
Matriks Jacobi dari $g$ adalah

\mathdisp {\left(  \left(  f(x)^2-z^2  \right)^{-1/2}  f(x) f'(x)   , \,   - z \left(  f(x)^2-z^2  \right)^{-1/2}                   \right)} {  }

Untuk meringkas tampilan matriks berikut, tulis $s=f(x)^2-z^2>0$. Matriks Hesse dari $g$ adalah

\mathdisp {\begin{pmatrix}  -s^{-3/2} f(x)^2 f'(x)^2 + s^{-1/2} \left(  f'(x)^2 + f(x) f^{\prime \prime} (x)   \right)   &   z s^{-3/2} f(x) f'(x)   \\  z s^{-3/2} f(x) f'(x)  &  -s^{-1/2}  -z^2 s^{-3/2}  \end{pmatrix}} {  }

atau merupakan hasil kali skalar
\mathl{\left(  f(x)^2-z^2  \right)^{-3/2}}{} dengan matriks

\mavergleichskettealigndrucklinks
{\vergleichskettealigndrucklinks
{ \begin{pmatrix}  -  f(x)^2 f'(x)^2 + \left(  f(x)^2-z^2  \right) \left(  f'(x)^2 + f(x) f^{\prime \prime} (x)   \right)   &   z  f(x) f'(x)   \\  z f(x) f'(x)  &  - \left(  f(x)^2-z^2  \right)  -z^2   \end{pmatrix}
}
{ =} { \begin{pmatrix}  f(x)^3 f^{\prime \prime} (x)  -z^2 f'(x)^2    -z^2 f(x) f^{\prime \prime} (x)    &   z  f(x) f'(x)   \\  z f(x) f'(x)  &  - f(x)^2  \end{pmatrix}
}
{ } {
}
{ } {
}
{ } {
}
}
{}{}{.}
Untuk titik $P$ yang dipilih, hasil ini adalah

\mavergleichskettedisp
{\vergleichskette
{ { \frac{   1  }{  f(x_0)^3 } } \begin{pmatrix}  f(x_0)^3 f^{\prime \prime} (x_0)   &   0   \\  0  &  - f(x_0)^2  \end{pmatrix}
}
{ =} { \begin{pmatrix}  f^{\prime \prime} (x_0)   &   0   \\  0  &  -  { \frac{   1  }{ f(x_0) } }   \end{pmatrix}
}
{ } {
}
{ } {
}
{ } {
}
}
{}{}{.}
Menurut
Lemma 4.9,
dengan
\mathl{\omega=\sqrt{1+f'(x_0)^2}}{,}
faktor invers metrik grafik juga harus diperhitungkan. Pemetaan Weingarten di $P$ adalah

\mathdisp {\begin{pmatrix}  { \frac{ f^{\prime \prime}(x_0) }{ \omega^3 } } & 0 \\ 0 & -{ \frac{ 1 }{ f(x_0)\omega } } \end{pmatrix}} { . }

Dengan demikian, kedua
\definitionsverweis {kelengkungan utama}{}{}
adalah
\mathkor {} {{ \frac{   f^{\prime \prime} (x_0)  }{  \left(1+f'(x_0)^2\right)^{3/2}  } }} {dan} {- { \frac{   1  }{  f(x_0) \sqrt{ 1 +f'(x_0)^2 }  } }} {,}
sedangkan arah-arah kelengkungan utama adalah vektor-vektor standar atau citranya
\mathkor {} {\begin{pmatrix}  1  \\f'(x_0)\\  0   \end{pmatrix}} {dan} {\begin{pmatrix}  0  \\ 0\\  1   \end{pmatrix}} {}
di ruang singgung $T_PY$. Jadi, arah-arah kelengkungan utama membentang sepanjang grafik $f$ dan sepanjang gerak melingkar rotasi.
\definitionsverweis {Kelengkungan Gauss}{}{}
adalah

\mathdisp {{ \frac{   - f^{\prime \prime} (x_0)  }{  f(x_0)  \left(  1 +f'(x_0)^2  \right)^2    } }} { . }

Catatan edisi: sumber hanya mengasumsikan keterdiferensialan meskipun memakai $f''$, menuliskan daerah grafik yang keliru, dan kembali menghilangkan faktor invers metrik. Hipotesis, daerah, dan tiga rumus kelengkungan di atas telah diperbaiki.

}

  \zwischenueberschrift{Kelengkungan Normal}

\inputdefinition
{}
{

Misalkan

\mavergleichskette
{\vergleichskette
{ W
}
{ \subseteq }{ \R^n
}
{  }{
}
{    }{
}
{   }{
}
}
{}{}{}
\definitionsverweis {terbuka}{}{,}
\maabb {h} { W } { \R
} {}
suatu
\definitionsverweis {fungsi yang dua kali terdiferensialkan secara kontinu}{}{}
dan

\mavergleichskette
{\vergleichskette
{ Y
}
{ = }{ h^{-1}(c)
}
{  }{
}
{    }{
}
{   }{
}
}
{}{}{}
adalah
\definitionsverweis {serat}{}{}
di atas

\mavergleichskette
{\vergleichskette
{ c
}
{ \in }{ \R
}
{  }{
}
{    }{
}
{   }{
}
}
{}{}{,}
dengan $h$
\definitionsverweis {reguler}{}{}
pada setiap titik di $Y$. Misalkan

\mavergleichskette
{\vergleichskette
{ P
}
{ \in }{ Y
}
{  }{
}
{    }{
}
{   }{
}
}
{}{}{.}
Untuk suatu vektor singgung ternormalisasi

\mavergleichskette
{\vergleichskette
{ v
}
{ \in }{ T_PY
}
{  }{
}
{    }{
}
{   }{
}
}
{}{}{}
nilai
\mathl{\left\langle  v  ,  L_P(v)  \right\rangle}{} disebut
\definitionswort {kelengkungan normal}{}
dari $Y$ di $P$ dalam arah $v$. Nilai tersebut dilambangkan dengan

\mavergleichskette
{\vergleichskette
{ \kappa(P,v)
}
{ = }{ \kappa(v)
}
{  }{
}
{    }{
}
{   }{
}
}
{}{}{}
.

}

Menurut
Lemma 4.6,
kelengkungan normal sama dengan

\mavergleichskettedisp
{\vergleichskette
{ \left\langle  v  ,  L_P(v)  \right\rangle
}
{ =} { \left\langle  \gamma^{\prime \prime} (0)  ,  N(P)  \right\rangle
}
{ } {
}
{ } {
}
{ } {
}
}
{}{}{,}
dengan $\gamma$ suatu realisasi kurva dari vektor singgung $v$. Jadi, kelengkungan normal mengukur komponen normal percepatan kurva tersebut.



\inputfaktbeweis
{Hyperfläche/Normalkrümmung/Berechnung/Fakt}
{Lemma}
{}
{

\faktsituation {Misalkan

\mavergleichskette
{\vergleichskette
{ W
}
{ \subseteq }{ \R^n
}
{  }{
}
{    }{
}
{   }{
}
}
{}{}{}
\definitionsverweis {terbuka}{}{,}
\maabb {h} { W } { \R
} {}
suatu
\definitionsverweis {fungsi yang dua kali terdiferensialkan secara kontinu}{}{}
dan

\mavergleichskette
{\vergleichskette
{ Y
}
{ = }{ h^{-1}(c)
}
{  }{
}
{    }{
}
{   }{
}
}
{}{}{}
adalah
\definitionsverweis {serat}{}{}
di atas

\mavergleichskette
{\vergleichskette
{ c
}
{ \in }{ \R
}
{  }{
}
{    }{
}
{   }{
}
}
{}{}{,}
dengan $h$
\definitionsverweis {reguler}{}{}
pada setiap titik di $Y$. Misalkan

\mavergleichskette
{\vergleichskette
{ P
}
{ \in }{ Y
}
{  }{
}
{    }{
}
{   }{
}
}
{}{}{.}
Misalkan
\mathl{\kappa_1 , \ldots , \kappa_{n-1}}{} adalah
\definitionsverweis {kelengkungan-kelengkungan utama}{}{}
dari $Y$ di $P$ dan
\mathl{v_1 , \ldots , v_{n-1}}{} adalah
\definitionsverweis {arah-arah kelengkungan utama}{}{}
\definitionsverweis {ternormalisasi}{}{ yang bersesuaian dan membentuk suatu basis ortonormal.}}
\faktfolgerung {Maka
\definitionsverweis {kelengkungan normal}{}{}
dari $Y$ di $P$ dalam arah

\mavergleichskette
{\vergleichskette
{ v
}
{ \in }{ T_PY
}
{  }{
}
{    }{
}
{   }{
}
}
{}{}{}
yang memenuhi $\lVert v\rVert=1$ adalah

\mavergleichskettedisp
{\vergleichskette
{ \kappa(v)
}
{ =} { \sum_{i = 1}^{n-1} \kappa_i \left\langle  v  , v_i   \right\rangle^2
}
{ } {
}
{ } {
}
{ } {
}
}
{}{}{.}}
Untuk vektor tak nol yang belum dinormalisasi, bentuk kuadratik pada ruas kanan memenuhi

\mathdisp {\sum_{i=1}^{n-1}\kappa_i\left\langle v,v_i\right\rangle^2=\lVert v\rVert^2\kappa\!\left({ \frac{v}{\lVert v\rVert} }\right)} { . }

Catatan edisi: sumber hanya menyebut vektor-vektor utama ``ternormalisasi'', padahal bukti memakai ortogonalitas, dan tidak mensyaratkan arah $v$ berpanjang satu. Kedua hipotesis itu dibuat eksplisit di sini.
\faktzusatz {}
\faktzusatz {}

}
{

Dengan

\mavergleichskette
{\vergleichskette
{ v
}
{ = }{ \sum_{i = 1}^{n-1} a_i v_i
}
{  }{
}
{    }{
}
{   }{
}
}
{}{}{}
dengan menggunakan ortogonalitas, diperoleh

\mavergleichskettealign
{\vergleichskettealign
{ \kappa(v)
}
{ =} { \left\langle  v  ,  L_P(v)  \right\rangle
}
{ =} { \left\langle  \sum_{i = 1}^{n-1} a_i v_i  ,  L_P \left(  \sum_{i = 1}^{n-1} a_i v_i  \right)   \right\rangle
}
{ =} { \left\langle  \sum_{i = 1}^{n-1} a_i v_i  ,  \sum_{i = 1}^{n-1} \kappa_i a_i v_i   \right\rangle
}
{ =} { \sum_{i = 1}^{n-1} \kappa_i a_i^2
}
}
{
\vergleichskettefortsetzungalign
{ =} { \sum_{i = 1}^{n-1} \kappa_i \left\langle  v  , v_i   \right\rangle^2
}
{ } {}
{ } {}
{ } {}
}
{}{.}

}

\bild{ \begin{center}
\includegraphics[width=5.5cm]{ \bildeinlesungsvg {Minimal surface curvature planes-de} {svg} }
\end{center}
\bildtext {Bidang-bidang kelengkungan utama, vektor normal, dan bidang singgung pada sebuah permukaan minimal. Label dalam berkas sumber dijelaskan seluruhnya oleh keterangan ini.} }

\bildlizenz { Minimal surface curvature planes-de.svg } {} {Eric Gaba} {Commons} {CC-by-sa 3.0} {}

\inputdefinition
{}
{

Misalkan

\mavergleichskette
{\vergleichskette
{ W
}
{ \subseteq }{ \R^n
}
{  }{
}
{    }{
}
{   }{
}
}
{}{}{}
\definitionsverweis {terbuka}{}{,}
\maabb {h} { W } { \R
} {}
suatu
\definitionsverweis {fungsi yang dua kali terdiferensialkan secara kontinu}{}{}
dan

\mavergleichskette
{\vergleichskette
{ Y
}
{ = }{ h^{-1}(c)
}
{  }{
}
{    }{
}
{   }{
}
}
{}{}{}
adalah
\definitionsverweis {serat}{}{}
di atas

\mavergleichskette
{\vergleichskette
{ c
}
{ \in }{ \R
}
{  }{
}
{    }{
}
{   }{
}
}
{}{}{,}
dengan $h$
\definitionsverweis {reguler}{}{}
pada setiap titik di $Y$. Misalkan

\mavergleichskette
{\vergleichskette
{ P
}
{ \in }{ Y
}
{  }{
}
{    }{
}
{   }{
}
}
{}{}{.}
Misalkan

\mavergleichskette
{\vergleichskette
{ v
}
{ \in }{ T_PY
}
{  }{
}
{    }{
}
{   }{
}
}
{}{}{}
suatu
\definitionsverweis {vektor singgung}{}{}
yang berbeda dari $0$, dan misalkan

\mavergleichskette
{\vergleichskette
{ u
}
{ \in }{ N_PY
}
{  }{
}
{    }{
}
{   }{
}
}
{}{}{}
suatu
\definitionsverweis {vektor normal}{}{.}
Vektor tersebut berbeda dari $0$. Bidang $\R u + \R v$ disebut
\definitionswort {bidang normal}{}
dari $Y$ melalui $P$.

}

Setiap vektor singgung

\mavergleichskette
{\vergleichskette
{ v
}
{ \neq }{ 0
}
{  }{
}
{    }{
}
{   }{
}
}
{}{}{}
menentukan sebuah bidang normal $E$ secara tunggal, yang biasanya dipandang sebagai
\mathl{P+E}{}.


\inputfaktbeweis
{Hyperfläche/Normalebene/Kurvenschnitt/Regularität/Fakt}
{Lemma}
{}
{

\faktsituation {Misalkan

\mavergleichskette
{\vergleichskette
{ W
}
{ \subseteq }{ \R^n
}
{  }{
}
{    }{
}
{   }{
}
}
{}{}{}
\definitionsverweis {terbuka}{}{,}
\maabb {h} { W } { \R
} {}
suatu
\definitionsverweis {fungsi yang dua kali terdiferensialkan secara kontinu}{}{}
dan

\mavergleichskette
{\vergleichskette
{ Y
}
{ = }{ h^{-1}(c)
}
{  }{
}
{    }{
}
{   }{
}
}
{}{}{}
adalah
\definitionsverweis {serat}{}{}
di atas

\mavergleichskette
{\vergleichskette
{ c
}
{ \in }{ \R
}
{  }{
}
{    }{
}
{   }{
}
}
{}{}{,}
dengan $h$
\definitionsverweis {reguler}{}{}
pada setiap titik di $Y$. Misalkan

\mavergleichskette
{\vergleichskette
{ P
}
{ \in }{ Y
}
{  }{
}
{    }{
}
{   }{
}
}
{}{}{.}
Misalkan

\mavergleichskette
{\vergleichskette
{ E
}
{ \subseteq }{ \R^n
}
{  }{
}
{    }{
}
{   }{
}
}
{}{}{}
suatu
\definitionsverweis {bidang normal}{}{}
melalui $P$ pada $Y$.}
\faktfolgerung {Maka irisan
\mathl{Y \cap (P+E)}{} merupakan suatu kurva bidang di $P+E$ yang reguler pada titik $P$.}
\faktzusatz {}
\faktzusatz {}

}
{

Setelah melakukan translasi, kita dapat mengasumsikan bahwa $P$ adalah titik nol ruang tersebut. Diferensial total
\maabbdisp {\left(D h \right)_{ 0}} {\R^n } {\R
} {}
bersifat surjektif dan kernelnya adalah ruang singgung $T_PY$. Bidang $E$ memuat vektor-vektor normal $\neq 0$ yang tidak berada dalam kernel. Oleh karena itu, diferensial total dari pemetaan komposit

\mathdisp {E \supseteq E \cap W \longrightarrow W \stackrel{h}{\longrightarrow} \R} {  }

bersifat surjektif di titik $0$. Dengan demikian, teorema fungsi implisit dapat diterapkan, dan diperoleh bahwa serat tersebut merupakan kurva yang reguler di titik $0$.

}

Perhatikan bahwa irisan kurva tersebut tidak harus reguler di setiap titik.


\inputfaktbeweis
{Hyperfläche/Normalkrümmung/Normalebene/Kurvenschnitt/Fakt}
{Satz}
{}
{

\faktsituation {Misalkan

\mavergleichskette
{\vergleichskette
{ W
}
{ \subseteq }{ \R^n
}
{  }{
}
{    }{
}
{   }{
}
}
{}{}{}
\definitionsverweis {terbuka}{}{,}
\maabb {h} { W } { \R
} {}
suatu
\definitionsverweis {fungsi yang dua kali terdiferensialkan secara kontinu}{}{}
dan

\mavergleichskette
{\vergleichskette
{ Y
}
{ = }{ h^{-1}(c)
}
{  }{
}
{    }{
}
{   }{
}
}
{}{}{}
adalah
\definitionsverweis {serat}{}{}
di atas

\mavergleichskette
{\vergleichskette
{ c
}
{ \in }{ \R
}
{  }{
}
{    }{
}
{   }{
}
}
{}{}{,}
dengan $h$
\definitionsverweis {reguler}{}{}
pada setiap titik di $Y$. Misalkan

\mavergleichskette
{\vergleichskette
{ P
}
{ \in }{ Y
}
{  }{
}
{    }{
}
{   }{
}
}
{}{}{.}
Misalkan

\mavergleichskette
{\vergleichskette
{ E
}
{ \subseteq }{ \R^n
}
{  }{
}
{    }{
}
{   }{
}
}
{}{}{}
suatu
\definitionsverweis {bidang normal}{}{}
melalui $P$ pada $Y$. Ambil suatu vektor satuan
\mathl{v\in E\cap T_PY}{,}
dan orientasikan bidang $P+E$ dengan menetapkan pasangan $(v,N(P))$ sebagai basis berorientasi positif. Dengan demikian, normal satuan kurva irisan di dalam bidang yang bersesuaian dengan arah singgung $v$ adalah $N(P)$.}
\faktfolgerung {Maka
\definitionsverweis {kelengkungan normal}{}{}
dari $Y$ di $P$ dalam arah $v$ sama dengan kelengkungan bertanda
\definitionsverweis {kelengkungan}{}{}
kurva bidang

\mavergleichskette
{\vergleichskette
{ Y \cap (P+E)
}
{ \subseteq }{ P+E
}
{ \cong }{ \R^2
}
{    }{
}
{   }{
}
}
{}{}{}
di titik $P$. Tanpa pilihan orientasi yang kompatibel, hanya nilai mutlak kedua kelengkungan yang ditentukan secara invarian.}
\faktzusatz {}
\faktzusatz {}

}
{

Untuk menyederhanakan notasi, kita mentranslasikan $P$ ke titik nol. Menurut
Lemma 5.13,
irisan $Y \cap E$ merupakan kurva yang reguler di $0$. Misalkan
\maabbdisp {\gamma} { I } { Y \cap E
} {}
suatu realisasi yang dua kali terdiferensialkan dan berparametrisasi panjang busur dengan
\mathl{\gamma(0)=0}{,}
\mathl{\gamma'(0)=v}{,}
dan normal satuan kurva $\gamma$ di dalam bidang yang kompatibel sama dengan $N(0)$. Menurut
Lemma 4.6,
kelengkungan normal sama dengan
\mathl{\left\langle  \gamma^{\prime \prime}(0) , N(0)  \right\rangle}{.} Karena bidang normal memuat vektor normal satuan $N(0)$, seluruh perhitungan berlangsung di bidang $E$. Jadi, pernyataan tersebut mengikuti dari
Lemma 3.8.

Catatan edisi: sumber tidak menetapkan orientasi irisan bidang, sehingga kesamaan kelengkungan bertanda bersifat ambigu. Konvensi orientasi di atas menghilangkan ambiguitas tersebut.

}

\chapter{Lembar Kerja 5}
\setcounter{section}{5}

  \zwischenueberschrift{Soal latihan}

\inputaufgabegibtloesung
{}
{

Misalkan

\mavergleichskettedisp
{\vergleichskette
{ f(x,y)
}
{ =} { ax^2+by^2+cxy
}
{ } {
}
{ } {
}
{ } {
}
}
{}{}{.}
Tentukan
\definitionsverweis {kelengkungan utama}{}{}
dan
\definitionsverweis {arah kelengkungan utama}{}{}
pada
\definitionsverweis {grafik}{}{}
fungsi $f$ di titik nol.

}
{} {}

\inputaufgabe
{}
{

Tentukan
\definitionsverweis {kelengkungan utama}{}{}
dan
\definitionsverweis {arah kelengkungan utama}{}{}
pada setiap titik
\definitionsverweis {grafik}{}{}
fungsi
\maabbeledisp {f} {\R^{n-1}} {\R
} { \left(  x_1   , \,   \ldots  , \,   x_{n-1}                 \right) } { a_1x_1^2  + \cdots + a_{n-1} x_{n-1}^2
} {.}

}
{} {}

\inputaufgabe
{}
{

Misalkan

\mavergleichskettedisp
{\vergleichskette
{ f(x,y,z)
}
{ =} { ax^2 +by^2 +c z^2 +dxy
}
{ } {
}
{ } {
}
{ } {
}
}
{}{}{.}
Hitung
\definitionsverweis {polinom karakteristik}{}{}
dari
\definitionsverweis {pemetaan Weingarten}{}{}
pada
\definitionsverweis {grafik}{}{}
fungsi $f$ di titik nol. Hitung pula
\definitionsverweis {kelengkungan utama}{}{}
dan
\definitionsverweis {arah kelengkungan utama}{}{.}

}
{} {}

\inputaufgabe
{}
{

Misalkan

\mavergleichskettedisp
{\vergleichskette
{ f(x,y,z)
}
{ =} { ax^2 +by^2 +c z^2 +dxy + e yz
}
{ } {
}
{ } {
}
{ } {
}
}
{}{}{.}
Hitung
\definitionsverweis {polinom karakteristik}{}{}
dari
\definitionsverweis {pemetaan Weingarten}{}{}
pada
\definitionsverweis {grafik}{}{}
fungsi $f$ di titik nol. Dalam kasus ini, seberapa sulitkah menentukan
\definitionsverweis {kelengkungan utama}{}{?}

}
{} {}

\inputaufgabe
{}
{

Misalkan

\mavergleichskettedisp
{\vergleichskette
{ f(x,y)
}
{ =} { x^2+y^3
}
{ } {
}
{ } {
}
{ } {
}
}
{}{}{.}
Tentukan
\definitionsverweis {kelengkungan utama}{}{,}
\definitionsverweis {arah kelengkungan utama}{}{}
dan
\definitionsverweis {kelengkungan Gauss}{}{}
pada
\definitionsverweis {grafik}{}{}
fungsi $f$ di setiap titik

\mavergleichskette
{\vergleichskette
{ P
}
{ = }{ (Q,f(Q))
}
{  }{
}
{    }{
}
{   }{
}
}
{}{}{,}

\mavergleichskette
{\vergleichskette
{ Q
}
{ \in }{ \R^2
}
{  }{
}
{    }{
}
{   }{
}
}
{}{}{.}

}
{} {}

\inputaufgabe
{}
{

Misalkan

\mavergleichskette
{\vergleichskette
{ V
}
{ \subseteq }{ \R^{n-1}
}
{  }{
}
{    }{
}
{   }{
}
}
{}{}{}
\definitionsverweis {terbuka}{}{}
dan misalkan
\maabb {f} { V } {\R
} {}
adalah suatu
\definitionsverweis {fungsi yang dua kali terdiferensialkan secara kontinu}{}{.}
Buktikan bahwa
\definitionsverweis {kelengkungan utama}{}{}
dan
\definitionsverweis {arah kelengkungan utama}{}{}
pada
\definitionsverweis {grafik}{}{}
fungsi $f$ di suatu titik

\mavergleichskette
{\vergleichskette
{ P
}
{ = }{ (Q,f(Q))
}
{  }{
}
{    }{
}
{   }{
}
}
{}{}{}
hanya bergantung pada turunan-turunan parsial pertama dan kedua dari $f$ di $Q$.

}
{} {}

\inputaufgabe
{}
{

Misalkan

\mavergleichskette
{\vergleichskette
{ W
}
{ \subseteq }{ \R^n
}
{  }{
}
{    }{
}
{   }{
}
}
{}{}{}
\definitionsverweis {terbuka}{}{,}
\maabb {h} { W } { \R
} {}
adalah suatu
\definitionsverweis {fungsi yang dua kali terdiferensialkan secara kontinu}{}{}
dan

\mavergleichskette
{\vergleichskette
{ Y
}
{ = }{ h^{-1}(c)
}
{  }{
}
{    }{
}
{   }{
}
}
{}{}{}
adalah
\definitionsverweis {serat}{}{}
di atas

\mavergleichskette
{\vergleichskette
{ c
}
{ \in }{ \R
}
{  }{
}
{    }{
}
{   }{
}
}
{}{}{,}
dengan $h$
\definitionsverweis {reguler}{}{}
di setiap titik pada $Y$.
Catatan edisi: sumber hanya mengasumsikan $h$ kelas $C^1$, padahal kelengkungan utama dan arah utamanya memerlukan pemetaan Weingarten. Hipotesis di atas diperkuat menjadi $C^2$ agar semua objek yang ditanyakan terdefinisi sebagaimana pada kuliah.
Kita juga memandang $h$ sebagai suatu pemetaan $\tilde{h}$ pada
\mathl{W \times \R^m}{},
dengan $m$ variabel terakhir tidak muncul secara eksplisit dalam $\tilde{h}$. Misalkan

\mavergleichskettedisp
{\vergleichskette
{ Z
}
{ =} { \tilde{h}^{-1}(c)
}
{ =} { Y \times \R^m
}
{ \subseteq} { W \times \R^{m}
}
{ } {
}
}
{}{}{.}
Misalkan

\mavergleichskette
{\vergleichskette
{ Q
}
{ \in }{ Z
}
{  }{
}
{    }{
}
{   }{
}
}
{}{}{}
adalah suatu titik yang terletak di atas

\mavergleichskette
{\vergleichskette
{ P
}
{ \in }{ Y
}
{  }{
}
{    }{
}
{   }{
}
}
{}{}{.}
Bagaimana hubungan antara
\definitionsverweis {kelengkungan utama}{}{}
dan
\definitionsverweis {arah kelengkungan utama}{}{}
dari $Z$ di $Q$ dengan yang dari $Y$ di $P$?

}
{} {}

Misalkan

\mavergleichskette
{\vergleichskette
{ W
}
{ \subseteq }{ \R^n
}
{  }{
}
{    }{
}
{   }{
}
}
{}{}{}
\definitionsverweis {terbuka}{}{,}
\maabb {h} { W } { \R
} {}
adalah suatu
\definitionsverweis {fungsi yang dua kali terdiferensialkan secara kontinu}{}{}
dan

\mavergleichskette
{\vergleichskette
{ Y
}
{ = }{ h^{-1}(c)
}
{  }{
}
{    }{
}
{   }{
}
}
{}{}{}
adalah
\definitionsverweis {serat}{}{}
di atas

\mavergleichskette
{\vergleichskette
{ c
}
{ \in }{ \R
}
{  }{
}
{    }{
}
{   }{
}
}
{}{}{,}
dengan $h$
\definitionsverweis {reguler}{}{}
di setiap titik pada $Y$. Untuk suatu titik

\mavergleichskette
{\vergleichskette
{ P
}
{ \in }{ Y
}
{  }{
}
{    }{
}
{   }{
}
}
{}{}{}
\definitionsverweis {determinan}{}{}
dari
\definitionsverweis {pemetaan Weingarten}{}{}
$L_P$ disebut
\definitionswort {kelengkungan Gauss--Kronecker}{}
dari $Y$ di $P$.

Kelengkungan Gauss--Kronecker merupakan perumuman langsung dari kelengkungan Gauss.

\inputaufgabe
{}
{

Misalkan

\mavergleichskette
{\vergleichskette
{ W
}
{ \subseteq }{ \R^n
}
{  }{
}
{    }{
}
{   }{
}
}
{}{}{}
\definitionsverweis {terbuka}{}{,}
\maabb {h} { W } { \R
} {}
adalah suatu
\definitionsverweis {fungsi yang dua kali terdiferensialkan secara kontinu}{}{}
dan

\mavergleichskette
{\vergleichskette
{ Y
}
{ = }{ h^{-1}(c)
}
{  }{
}
{    }{
}
{   }{
}
}
{}{}{}
adalah
\definitionsverweis {serat}{}{}
di atas

\mavergleichskette
{\vergleichskette
{ c
}
{ \in }{ \R
}
{  }{
}
{    }{
}
{   }{
}
}
{}{}{,}
dengan $h$
\definitionsverweis {reguler}{}{}
di setiap titik pada $Y$. Misalkan

\mavergleichskette
{\vergleichskette
{ P
}
{ \in }{ Y
}
{  }{
}
{    }{
}
{   }{
}
}
{}{}{.}
Buktikan bahwa hasil kali
\definitionsverweis {kelengkungan utama}{}{}
dari $Y$ di $P$ sama dengan
\definitionsverweis {kelengkungan Gauss--Kronecker}{}{}
dari $Y$ di $P$.

}
{} {}

\inputaufgabe
{}
{

Misalkan

\mavergleichskettedisp
{\vergleichskette
{ h(x,y,z,w)
}
{ =} { x^2+y^3+z^4+w^5
}
{ } {
}
{ } {
}
{ } {
}
}
{}{}{}
dan

\mavergleichskette
{\vergleichskette
{ Y
}
{ = }{ h^{-1}(1)
}
{  }{
}
{    }{
}
{   }{
}
}
{}{}{.}
Untuk titik

\mavergleichskettedisp
{\vergleichskette
{ P
}
{ =} { \left(  1   , \,   -1  , \,   1 , \,   0                \right)
}
{ \in} { Y
}
{ } {
}
{ } {
}
}
{}{}{}
tentukan
\definitionsverweis {pemetaan Weingarten}{}{}
$L_P$,
\definitionsverweis {polinom karakteristik}{}{}
darinya, serta
\definitionsverweis {kelengkungan Gauss--Kronecker}{}{}
dari $Y$ di $P$.

}
{} {}

\inputaufgabe
{}
{

Misalkan

\mavergleichskettedisp
{\vergleichskette
{ h(x,y,z)
}
{ =} { x^2+y^3+z^5
}
{ } {
}
{ } {
}
{ } {
}
}
{}{}{}
pada

\mavergleichskette
{\vergleichskette
{ W
}
{ = }{ \R^3 \setminus \{0\}
}
{  }{
}
{    }{
}
{   }{
}
}
{}{}{}
dan

\mavergleichskette
{\vergleichskette
{ Y
}
{ = }{ h^{-1}(0)
}
{  }{
}
{    }{
}
{   }{
}
}
{}{}{.}
Tentukan
\definitionsverweis {kelengkungan normal}{}{}
di titik

\mavergleichskette
{\vergleichskette
{ P
}
{ = }{ \begin{pmatrix}  1  \\ -1\\  0   \end{pmatrix}
}
{  }{
}
{    }{
}
{   }{
}
}
{}{}{}
dalam arah vektor singgung ternormalisasi
\mathl{{ \frac{   1  }{  \sqrt{29}  } }  \begin{pmatrix}  3  \\ -2\\  4   \end{pmatrix}}{.}

}
{} {}

  \zwischenueberschrift{Soal untuk dikumpulkan}

\inputaufgabe
{4}
{

Misalkan

\mavergleichskettedisp
{\vergleichskette
{ f(x,y)
}
{ =} { x^3 +x^2y+3xy^2-y^4
}
{ } {
}
{ } {
}
{ } {
}
}
{}{}{.}
Tentukan
\definitionsverweis {kelengkungan utama}{}{,}
\definitionsverweis {arah kelengkungan utama}{}{}
dan
\definitionsverweis {kelengkungan Gauss}{}{}
pada
\definitionsverweis {grafik}{}{}
fungsi $f$ di setiap titik

\mavergleichskette
{\vergleichskette
{ P
}
{ = }{ (Q,f(Q))
}
{  }{
}
{    }{
}
{   }{
}
}
{}{}{,}

\mavergleichskette
{\vergleichskette
{ Q
}
{ \in }{ \R^2
}
{  }{
}
{    }{
}
{   }{
}
}
{}{}{.}

}
{} {}

\inputaufgabe
{4}
{

Misalkan

\mavergleichskettedisp
{\vergleichskette
{ f(x,y,z)
}
{ =} { -7x^2 +3y^2 -5z^2+ 6xy
}
{ } {
}
{ } {
}
{ } {
}
}
{}{}{.}
Tentukan
\definitionsverweis {kelengkungan utama}{}{}
dan
\definitionsverweis {arah kelengkungan utama}{}{}
pada
\definitionsverweis {grafik}{}{}
fungsi $f$ di titik nol.

}
{} {}

\inputaufgabe
{6}
{

Misalkan

\mavergleichskettedisp
{\vergleichskette
{ h(x,y,z,w)
}
{ =} { x^2+yz+z^3+xyzw
}
{ } {
}
{ } {
}
{ } {
}
}
{}{}{}
dan

\mavergleichskette
{\vergleichskette
{ Y
}
{ = }{ h^{-1}(2)
}
{  }{
}
{    }{
}
{   }{
}
}
{}{}{.}
Untuk titik

\mavergleichskettedisp
{\vergleichskette
{ P
}
{ =} { \left(  1   , \,   1  , \,   1 , \,   -1                \right)
}
{ \in} { Y
}
{ } {
}
{ } {
}
}
{}{}{}
tentukan
\definitionsverweis {pemetaan Weingarten}{}{}
$L_P$,
\definitionsverweis {polinom karakteristik}{}{}
darinya, serta
\definitionsverweis {kelengkungan Gauss--Kronecker}{}{}
dari $Y$ di $P$.

}
{} {Petunjuk: Selesaikan persamaan terhadap suatu variabel yang sesuai dan perlakukan hipermuka tersebut sebagai grafik dalam variabel-variabel lainnya.}

\inputaufgabe
{6 (2+2+2)}
{

Misalkan

\mavergleichskettedisp
{\vergleichskette
{ h(x,y,z)
}
{ =} { x^3+y^3+z^3
}
{ } {
}
{ } {
}
{ } {
}
}
{}{}{}
pada

\mavergleichskette
{\vergleichskette
{ W
}
{ = }{ \R^3 \setminus \{0\}
}
{  }{
}
{    }{
}
{   }{
}
}
{}{}{}
dan

\mavergleichskette
{\vergleichskette
{ Y
}
{ = }{ h^{-1}(0)
}
{  }{
}
{    }{
}
{   }{
}
}
{}{}{.}
Misalkan

\mavergleichskette
{\vergleichskette
{ P
}
{ = }{ \begin{pmatrix}  1  \\ 1\\  - \sqrt[3]{2}   \end{pmatrix}
}
{ \in }{ Y
}
{    }{
}
{   }{
}
}
{}{}{}
dan

\mavergleichskettedisp
{\vergleichskette
{ v
}
{ =} { \begin{pmatrix}  1  \\ -1\\  0   \end{pmatrix}
}
{ \in} { T_PY
}
{ } {
}
{ } {
}
}
{}{}{.}
\aufzaehlungdrei{Tentukan suatu persamaan bidang untuk
\definitionsverweis {bidang normal}{}{}
yang bersesuaian dengan $v$.
}{Tentukan
\definitionsverweis {kelengkungan normal}{}{}
di $P$ dalam arah vektor singgung ternormalisasi
\mathl{{ \frac{   v  }{  \Vert { v } \Vert  } }}{.}
}{Tentukan suatu kurva berparametrisasi panjang busur
\maabbdisp {\gamma} { I } { Y
} {}
dengan

\mavergleichskette
{\vergleichskette
{ \gamma (0)
}
{ = }{ P
}
{  }{
}
{    }{
}
{   }{
}
}
{}{}{,}

\mavergleichskette
{\vergleichskette
{ \gamma' (0)
}
{ = }{ { \frac{   v  }{  \Vert { v } \Vert  } }
}
{  }{
}
{    }{
}
{   }{
}
}
{}{}{.}
Verifikasikan berlakunya
Teorema 5.14
dalam situasi ini.
}

Catatan edisi: pada butir ketiga, orientasikan bidang normal dengan menetapkan pasangan $(v/\lVert v\rVert,N(P))$ sebagai basis positif, sehingga normal satuan kurva irisan di dalam bidang yang bersesuaian dengan arah singgung $v/\lVert v\rVert$ adalah $N(P)$. Tanpa konvensi ini, verifikasi kelengkungan bertanda hanya menentukan kesamaan nilai mutlak.

}
{} {}

\inputaufgabe
{2}
{

Buatlah sketsa suatu permukaan terdiferensialkan $Y$ di ruang dan suatu bidang normal $E$ yang irisannya dengan permukaan tersebut menghasilkan suatu kurva yang tidak reguler pada setidaknya satu titik.

}
{} {}

\section*{Solusi yang disediakan oleh sumber}
\addcontentsline{toc}{section}{Solusi yang disediakan oleh sumber}
\subsection*{Solusi Soal 5.1}
\textit{Catatan edisi: solusi sumber memakai reduksi ke matriks Hesse tanpa menyebut matriks metrik grafik dan memberikan vektor eigen nol pada kasus sah $c=0$ dan $a<b$. Alasan metrik dan pemisahan kasus berikut memperbaiki kedua kekurangan itu.}

Kita menggunakan
Fakta,
dan di titik nol berlaku $\operatorname{Grad}f(0)=0$, sehingga faktor normalisasi $\omega=1$ dan matriks metrik grafik $G=I$. Oleh karena itu,
\definitionsverweis {pemetaan Weingarten}{}{}
langsung dinyatakan oleh
\definitionsverweis {matriks Hesse}{}{.}
Matriks Hesse dari fungsi tersebut adalah

\mathdisp {\begin{pmatrix}  2a  &   c   \\ c &  2b \end{pmatrix}} { , }

sedangkan
\definitionsverweis {polinom karakteristik}{}{}
darinya adalah

\mavergleichskettealign
{\vergleichskettealign
{ (t-2a)(t-2b)-c^2
}
{ =} { t^2 -2(a+b)t +4ab -c^2
}
{ =} { \left(  t-(a+b)  \right)^2 - (a+b)^2 +4ab -c^2
}
{ =} { \left(  t-(a+b)  \right)^2  - a^2 -b^2 +2ab -c^2
}
{ =} { \left(  t-(a+b)  \right)^2  - (a-b)^2 -c^2
}
}
{}
{}{.}
Jadi, nilai-nilai eigennya
\zusatzklammer {yang sekaligus merupakan kelengkungan utama} {} {}
adalah

\mavergleichskettedisp
{\vergleichskette
{ \lambda_{1,2}
}
{ =} { \pm \sqrt{ (a-b)^2+c^2 } +a+b
}
{ } {
}
{ } {
}
{ } {
}
}
{}{}{.}
Jika $c=0$, matriks tersebut diagonal. Apabila $a>b$, arah eigen
\mathkor {} {\begin{pmatrix}1\\0\end{pmatrix}} {dan} {\begin{pmatrix}0\\1\end{pmatrix}} {}
masing-masing bersesuaian dengan nilai eigen $2a$ dan $2b$. Apabila $a<b$, arah eigen
\mathkor {} {\begin{pmatrix}0\\1\end{pmatrix}} {dan} {\begin{pmatrix}1\\0\end{pmatrix}} {}
masing-masing bersesuaian dengan nilai eigen $2b$ dan $2a$. Jika

\mavergleichskette
{\vergleichskette
{ a
}
{ = }{ b
}
{  }{
}
{    }{
}
{   }{
}
}
{}{}{}
dan

\mavergleichskette
{\vergleichskette
{ c
}
{ = }{ 0
}
{  }{
}
{    }{
}
{   }{
}
}
{}{}{}
maka pemetaan Weingarten adalah perkalian dengan $2a$ dan setiap vektor tak nol merupakan vektor eigen. Dengan demikian, seluruh kasus $c=0$ telah selesai. Mulai sekarang, anggap $c\neq 0$. Untuk

\mavergleichskettedisp
{\vergleichskette
{ \lambda_{1}
}
{ =} { \sqrt{ (a-b)^2+c^2 } +a+b
}
{ } {
}
{ } {
}
{ } {
}
}
{}{}{}
kita perlu menentukan kernel matriks

\mathdisp {\begin{pmatrix}  \sqrt{(a-b)^2 +c^2} + b-a  &   -c  \\  -c &  \sqrt{(a-b)^2 +c^2} + a-b \end{pmatrix}} {  }

yang diberikan oleh

\mathdisp {\R \begin{pmatrix}  \sqrt{(a-b)^2 +c^2} + a-b  \\c   \end{pmatrix}} { . }

Dengan demikian, salah satu arah kelengkungan utama adalah
\mathl{\begin{pmatrix}  \sqrt{(a-b)^2 +c^2} + a-b  \\c   \end{pmatrix}}{} dengan kelengkungan utama
\mathl{\sqrt{ (a-b)^2+c^2 } +a+b}{.}

Dengan cara yang sama, diperoleh arah kelengkungan utama
\mathl{\begin{pmatrix}  -c \\ \sqrt{(a-b)^2 +c^2} + a-b    \end{pmatrix}}{} dengan kelengkungan utama
\mathl{-\sqrt{ (a-b)^2+c^2 } +a+b}{.}

\part{Unit 6}
\chapter[Kuliah 6]{Kuliah 6: Turunan Kovarian, Medan Vektor Paralel, dan Transpor Paralel}
\setcounter{section}{6}

  \zwischenueberschrift{Turunan Kovarian}

Untuk suatu hipermuka terdiferensialkan

\mavergleichskette
{\vergleichskette
{ Y
}
{ \subseteq }{ \R^n
}
{  }{
}
{    }{
}
{   }{
}
}
{}{}{}
kita ingin mendefinisikan konsep-konsep turunan yang mudah didefinisikan di ruang sekitar $\R^n$, tetapi memiliki kaitan alami dengan $Y$. Konsep ini disebut turunan kovarian; alat bantu terpentingnya adalah proyeksi ortogonal
\maabbdisp {\pi_P} {\R^n } { T_PY
} {}
di suatu titik

\mavergleichskette
{\vergleichskette
{ P
}
{ \in }{ Y
}
{  }{
}
{    }{
}
{   }{
}
}
{}{}{.}

\emph{Catatan edisi: Relasi pembuka diperbaiki dari $Y=\R^n$ menjadi $Y\subseteq\R^n$; sebuah hipermuka tidak sama dengan seluruh ruang sekitarnya.}

\inputdefinition
{}
{

Misalkan

\mavergleichskette
{\vergleichskette
{ Y
}
{ \subseteq }{ W
}
{ \subseteq }{ \R^n
}
{    }{
}
{   }{
}
}
{}{}{,}
$W$
\definitionsverweis {terbuka}{}{,}
dan $Y$ suatu
\definitionsverweis {hipermuka terdiferensialkan}{}{,}
misalkan

\mavergleichskette
{\vergleichskette
{ P
}
{ \in }{ Y
}
{  }{
}
{    }{
}
{   }{
}
}
{}{}{}
suatu titik dan

\mavergleichskette
{\vergleichskette
{ v
}
{ \in }{ T_PY
}
{  }{
}
{    }{
}
{   }{
}
}
{}{}{}
suatu
\definitionsverweis {vektor singgung}{}{.}
Misalkan
\maabbdisp {F} { U} { \R^n
} {}
suatu medan vektor terdiferensialkan yang didefinisikan pada lingkungan terbuka

\mavergleichskette
{\vergleichskette
{ P
}
{ \in }{ U
}
{ \subseteq }{ \R^n
}
{    }{
}
{   }{
}
}
{}{}{}
dan pada $Y \cap U$ tangensial terhadap $Y$. Maka

\mavergleichskettedisp
{\vergleichskette
{ \nabla_v F
}
{ \defeq} { \pi_P   \left(  \left(  D F   \right)_{ P } \left(  v  \right)  \right)
}
{ } {
}
{ } {
}
{ } {
}
}
{}{}{,}
dengan
\maabbdisp {\pi_P} {\R^n } { T_PY
} {}
menyatakan
\definitionsverweis {proyeksi ortogonal}{}{,}
disebut
\definitionswort {turunan kovarian}{}
dari $F$ dalam arah $v$.

}

Secara keseluruhan, situasinya adalah

\mathdisp {T_PY \longrightarrow \R^n \stackrel{  \left(D F \right)_{ P} }{\longrightarrow} \R^n\stackrel{ \pi_P }{\longrightarrow} T_PY} {  }

dan $\nabla_v F$ merupakan hasil ketika vektor singgung

\mavergleichskette
{\vergleichskette
{ v
}
{ \in }{ T_PY
}
{  }{
}
{    }{
}
{   }{
}
}
{}{}{}
dimasukkan pada awal rangkaian tersebut. Jika $N$ menyatakan suatu medan normal satuan, maka turunan kovarian sama dengan

\mavergleichskettedisp
{\vergleichskette
{ \nabla_v F
}
{ =} { \left(  D F   \right)_{ P } \left(  v  \right) -  \left\langle  \left(  D F   \right)_{ P } \left(  v  \right)  ,  N(P)  \right\rangle N(P)
}
{ } {
}
{ } {
}
{ } {
}
}
{}{}{,}
karena dengan cara inilah komplemen ortogonal dihitung.

\inputdefinition
{}
{

Misalkan

\mavergleichskette
{\vergleichskette
{ Y
}
{ \subseteq }{ W
}
{ \subseteq }{ \R^n
}
{    }{
}
{   }{
}
}
{}{}{,}
$W$
\definitionsverweis {terbuka}{}{,}
dan $Y$ suatu
\definitionsverweis {hipermuka terdiferensialkan}{}{}
dan misalkan
\maabbdisp {G} { Y } { TY
} {}
suatu medan vektor tangensial pada $Y$. Misalkan
\maabbdisp {F} { U} { \R^n
} {}
suatu medan vektor terdiferensialkan yang didefinisikan pada lingkungan terbuka

\mavergleichskette
{\vergleichskette
{ P
}
{ \in }{ U
}
{ \subseteq }{ \R^n
}
{    }{
}
{   }{
}
}
{}{}{}
dan pada $Y \cap U$ tangensial terhadap $Y$. Maka medan vektor tangensial
\maabbdisp {\nabla_G F} { Y \cap U } { T(Y \cap U)
} {,}
yang didefinisikan oleh

\mavergleichskettedisp
{\vergleichskette
{ \left(  \nabla_G F  \right) (P)
}
{ \defeq} { \nabla_{G(P)} F }
{ } {
}
{ } {
}
{ } {
}
}
{}{}{}
disebut
\definitionswort {turunan kovarian}{}
dari $F$ dalam arah $G$.

}



\inputfaktbeweis
{Hyperfläche/Tangentiales Vektorfeld/Bezüglich Tangentenvektor/Kovariante Ableitung/Eigenschaften/Fakt}
{Lemma}
{}
{

\faktsituation {Misalkan

\mavergleichskette
{\vergleichskette
{ Y
}
{ \subseteq }{ W
}
{ \subseteq }{\R^n
}
{    }{
}
{   }{
}
}
{}{}{,}
$W$
\definitionsverweis {terbuka}{}{,}
dan $Y$ suatu
\definitionsverweis {hipermuka terdiferensialkan}{}{,}
dan misalkan

\mavergleichskette
{\vergleichskette
{ P
}
{ \in }{ Y
}
{  }{
}
{    }{
}
{   }{
}
}
{}{}{}
suatu titik.}
\faktuebergang {Maka pernyataan-pernyataan berikut berlaku.}
\faktfolgerung {\aufzaehlungdrei{Untuk suatu
\definitionsverweis {medan vektor tangensial}{}{}
terdiferensialkan $F$ yang tetap, pemetaan
\maabbeledisp {} {T_PY} { T_PY
} { v } { \nabla_v F
} {,}
bersifat
\definitionsverweis {linear}{}{.}
}{Untuk suatu vektor singgung tetap

\mavergleichskette
{\vergleichskette
{ v
}
{ \in }{ T_PY
}
{  }{
}
{    }{
}
{   }{
}
}
{}{}{}
pemetaan
\maabbdisp {} { F } { \nabla_v F
} {,}
yang memetakan sebuah medan vektor terdiferensialkan ke turunan kovariannya sepanjang $v$, bersifat linear.
}{Untuk suatu fungsi terdiferensialkan
\maabb {g} { U } {\R
} {}
dan suatu medan vektor tangensial terdiferensialkan $F$, berlaku

\mavergleichskettedisp
{\vergleichskette
{ \nabla_v (gF)
}
{ =} { g (P) \nabla_v F + \left(  D_{v}  g   \right) \left(   P   \right) F(P)
}
{ } {
}
{ } {
}
{ } {
}
}
{}{}{.}
\emph{Catatan edisi: Evaluasi berlebih $(P)$ setelah $\nabla_vF$ dihapus karena $v\in T_PY$ sudah menentukan titiknya.}
}}
\faktzusatz {}
\faktzusatz {}

}
{

\aufzaehlungdrei{Hal ini langsung mengikuti representasi

\mavergleichskettedisp
{\vergleichskette
{ \nabla_v F
}
{ =} { \left(  D F   \right)_{ P } \left(  v  \right) - \left\langle  \left(  D F   \right)_{ P } \left(  v  \right)  ,  N(P)  \right\rangle N(P)
}
{ } {
}
{ } {
}
{ } {
}
}
{}{}{,}
karena kedua suku tersebut linear dalam $v$.
}{Hal ini juga mengikuti representasi

\mavergleichskettedisp
{\vergleichskette
{ \nabla_v F
}
{ =} { \left(  D F   \right)_{ P } \left(  v  \right) - \left\langle  \left(  D F   \right)_{ P } \left(  v  \right)  ,  N(P)  \right\rangle N(P)
}
{ } {
}
{ } {
}
{ } {
}
}
{}{}{,}
karena untuk vektor tetap $v$, menurut
Lemma 43.6 (Analisis (Osnabrück 2021--2023)),
kedua suku bergantung secara linear pada $F$.
}{Menurut
Lemma 45.10 (Analisis (Osnabrück 2021--2023)),
berlaku

\mavergleichskettedisp
{\vergleichskette
{ \left(D( g F)\right)_{ P }
}
{ =} { g(P) \cdot \left(D F \right)_{ P } + \left(D g \right)_{ P } \cdot F (P)
}
{ } {
}
{ } {
}
{ } {
}
}
{}{}{.}
Oleh karena itu,

\mavergleichskettealignhandlinks
{\vergleichskettealignhandlinks
{ \nabla_v (gF)
}
{ =} { \left(  D(gF)  \right)_{ P } \left(  v  \right) - \left\langle  \left(  D(gF)  \right)_{ P } \left(  v  \right)  ,  N(P)  \right\rangle N(P)
}
{ =} { \begin{aligned}[t]
&g(P) \cdot \left(  D F   \right)_{ P } \left(  v  \right) + \left(  D g   \right)_{ P } \left(  v  \right) \cdot F (P) \\
&\quad - \left\langle  g(P) \cdot \left(  D F   \right)_{ P } \left(  v  \right) + \left(  D g   \right)_{ P } \left(  v  \right) \cdot F (P)  ,  N(P)  \right\rangle N(P)
\end{aligned}
}
{ =} { \begin{aligned}[t]
&g(P) \cdot \left(  D F   \right)_{ P } \left(  v  \right) - \left\langle  g(P) \cdot \left(  D F   \right)_{ P } \left(  v  \right)  ,  N(P)  \right\rangle N(P) \\
&\quad + \left(  D g   \right)_{ P } \left(  v  \right) \cdot F (P) - \left\langle  \left(  D g   \right)_{ P } \left(  v  \right) \cdot F (P)  ,  N(P)  \right\rangle N(P)
\end{aligned}
}
{ =} { g (P) \nabla_v F + \left(  D_{v}  g   \right) \left(   P   \right) \cdot F(P)
}
}
{}
{}{,}
karena medan vektor $F$ tegak lurus terhadap $N$.
}

}



\inputfaktbeweis
{Hyperfläche/Tangentiales Vektorfeld/Bezüglich tangentialem Vektorfeld/Kovariante Ableitung/Eigenschaften/Fakt}
{Lemma}
{}
{

\faktsituation {Misalkan

\mavergleichskette
{\vergleichskette
{ Y
}
{ \subseteq }{ W
}
{ \subseteq }{\R^n
}
{    }{
}
{   }{
}
}
{}{}{,}
$W$
\definitionsverweis {terbuka}{}{,}
dan $Y$ suatu
\definitionsverweis {hipermuka terdiferensialkan}{}{.}}
\faktuebergang {Maka pernyataan-pernyataan berikut berlaku bagi medan-medan vektor tangensial pada $Y$.}
\faktfolgerung {\aufzaehlungdrei{Untuk suatu medan vektor tangensial terdiferensialkan $F$ yang tetap, pemetaan
\mathl{G \mapsto \nabla_G F}{}
bersifat
\definitionsverweis {linear}{}{}
dalam $G$. Selanjutnya, untuk suatu fungsi kontinu
\maabb {g} { Y } {\R
} {}

\mavergleichskettedisp
{\vergleichskette
{ \nabla_{gG} F
}
{ =} { g \nabla_G F
}
{ } {
}
{ } {
}
{ } {
}
}
{}{}{.}
}{Untuk suatu medan vektor tangensial $G$ yang tetap, pemetaan
\mathl{F \mapsto \nabla_G F}{,} yang memetakan sebuah medan vektor tangensial terdiferensialkan ke turunan kovariannya sepanjang $G$, bersifat linear.
}{Untuk suatu fungsi terdiferensialkan
\maabb {g} { U } {\R
} {}
dan suatu medan vektor tangensial terdiferensialkan $F$, berlaku

\mavergleichskettedisp
{\vergleichskette
{ (\nabla_G (gF))(P)
}
{ =} { g (P) ( \nabla_G F) (P) + (\nabla_G g )(P) F(P)
}
{ } {
}
{ } {
}
{ } {
}
}
{}{}{.}
\emph{Catatan edisi: Ruas kiri dievaluasi di $P$ agar identitas ini konsisten dengan ruas kanan yang juga dievaluasi di $P$.}
}}
\faktzusatz {}
\faktzusatz {}

}
{

Semua pernyataan mengikuti dari
Lemma 6.3.

}

Pada persamaan terakhir,

\mavergleichskettedisp
{\vergleichskette
{ \left(  \nabla_G g  \right) (P)
}
{ =} { \left(  D_{ G(P)}  g   \right) \left(   P   \right)
}
{ } {
}
{ } {
}
{ } {
}
}
{}{}{,}
yang berarti turunan berarah fungsi $g$ diambil dalam arah $G(P)$.

\inputdefinition
{}
{

Misalkan

\mavergleichskette
{\vergleichskette
{ Y
}
{ \subseteq }{ W
}
{ \subseteq }{\R^n
}
{    }{
}
{   }{
}
}
{}{}{,} $W$
\definitionsverweis {terbuka}{}{,}
dan $Y$ suatu
\definitionsverweis {hipermuka terdiferensialkan}{}{,}
misalkan
\maabbdisp {\gamma} {I} {Y
} {}
suatu
\definitionsverweis {kurva terdiferensialkan}{}{}
dan misalkan
\maabbdisp {F} { I } { \R^n
} {}
suatu medan vektor terdiferensialkan sepanjang $\gamma$ yang tangensial terhadap $Y$. Maka

\mavergleichskettedisp
{\vergleichskette
{ (\nabla_\gamma F) (t)
}
{ \defeq} { \pi_{\gamma(t)} F' (t)
}
{ } {
}
{ } {
}
{ } {
}
}
{}{}{,}
disebut
\definitionswort {turunan kovarian}{}
dari $F$ sepanjang $\gamma$.

\emph{Catatan edisi: Domain $\gamma$ dan $F$ diseragamkan menjadi interval $I$, dan $F$ dinyatakan sebagai medan sepanjang $\gamma$; sumber memakai $[a,b]$ dan kemudian interval $I$ yang belum didefinisikan.}

}

Medan vektor $F$ merupakan kurva bernilai vektor di $\R^n$ dengan

\mavergleichskette
{\vergleichskette
{ F(t)
}
{ \in }{ T_{\gamma(t)}Y
}
{  }{
}
{    }{
}
{   }{
}
}
{}{}{.}
Turunan kovariannya diperoleh dengan menurunkan kurva ini lalu mengambil proyeksi ortogonal hasilnya ke ruang singgung. Kadang-kadang ditulis pula $\nabla_{\gamma'} F$ atau $\nabla_{\partial} F$\zusatzfussnote {Untuk
Definisi 6.2
dan
Definisi 6.5
terdapat generalisasi berikut. Misalkan $Y$ adalah hipermuka terdiferensialkan
\zusatzklammer {atau suatu manifold diferensiabel dengan koneksi pada bundel tangen} {} {,}
misalkan $Z$ suatu manifold diferensiabel, misalkan
\maabb {\varphi} { Z } { Y
} {}
suatu pemetaan terdiferensialkan, misalkan
\maabbdisp {F} { Z } { \R^n
} {}
suatu pemetaan terdiferensialkan dengan $F(Q)\in T_{\varphi(Q)}Y$ untuk setiap $Q\in Z$, yaitu suatu medan vektor tangensial sepanjang $\varphi$, dan $G$ suatu medan vektor pada $Z$. Maka
\zusatzklammer {untuk

\mavergleichskette
{\vergleichskette
{ Q
}
{ \in }{ Z
}
{  }{
}
{    }{
}
{   }{
}
}
{}{}{}
dan

\mavergleichskette
{\vergleichskette
{ P
}
{ = }{ \varphi(Q)
}
{  }{
}
{    }{
}
{   }{
}
}
{}{}{}} {} {}
turunan kovarian $F$ terhadap $G$
\zusatzklammer {sepanjang $\varphi$} {} {}
didefinisikan oleh

\mavergleichskettedisp
{\vergleichskette
{ ( \nabla_G F )(Q)
}
{ =} { \pi_{ P  } \left(D F \right)_{ Q}  (G(Q))
}
{ } {
}
{ } {
}
{ } {
}
}
{}{}{}
dan $\nabla_GF$ kembali merupakan suatu pemetaan
\maabb {} { Z } { TY
}{.}
\emph{Catatan edisi: Kodomain $F$ diperbaiki menjadi $\R^n$ dengan syarat nilai tangensial, sehingga diferensial $DF$ yang diproyeksikan pada rumus di atas bertipe benar.}
Dalam definisi pertama,

\mavergleichskette
{\vergleichskette
{ Z
}
{ = }{ Y
}
{  }{
}
{    }{
}
{   }{
}
}
{}{}{}
dan $\varphi$ adalah identitas; dalam definisi kedua,

\mavergleichskette
{\vergleichskette
{ Z
}
{ = }{ I
}
{  }{
}
{    }{
}
{   }{
}
}
{}{}{}
adalah suatu interval,

\mavergleichskette
{\vergleichskette
{ \varphi
}
{ = }{ \gamma
}
{  }{
}
{    }{
}
{   }{
}
}
{}{}{}
adalah lintasan tersebut, dan

\mavergleichskette
{\vergleichskette
{ G
}
{ = }{ \partial
}
{  }{
}
{    }{
}
{   }{
}
}
{}{}{}
adalah medan vektor konstan pada interval $I$ dengan arah $1$.} {} {.}

  \zwischenueberschrift{Medan Vektor Paralel}

\inputdefinition
{}
{

Misalkan

\mavergleichskette
{\vergleichskette
{ Y
}
{ \subseteq }{ W
}
{ \subseteq }{ \R^n
}
{    }{
}
{   }{
}
}
{}{}{,}
$W$
\definitionsverweis {terbuka}{}{,}
dan $Y$ suatu
\definitionsverweis {hipermuka terdiferensialkan}{}{}
dan misalkan
\maabb {\gamma} { I } { Y
} {}
suatu
\definitionsverweis {kurva terdiferensialkan}{}{.}
Suatu medan vektor tangensial terdiferensialkan yang didefinisikan sepanjang $\gamma$,
\maabb {F} { I } {\R^n
} {}
disebut
\definitionswort {paralel sepanjang}{}
$\gamma$ jika

\mavergleichskettedisp
{\vergleichskette
{ (\nabla_{\gamma} F)(t)
}
{ =} { 0
}
{ } {
}
{ } {
}
{ } {
}
}
{}{}{}
untuk semua

\mavergleichskette
{\vergleichskette
{ t
}
{ \in }{ I
}
{  }{
}
{    }{
}
{   }{
}
}
{}{}{.}

}

Suatu medan vektor paralel sepanjang $\gamma$ tepat ketika
\mathl{F'(t)}{} selalu tegak lurus terhadap ruang singgung
\mathl{T_{\gamma(t)}Y}{}, yakni ketika turunannya berada pada garis normal.


\inputfaktbeweis
{Hyperfläche/Kurve/Paralleles Vektorfeld/Eigenschaften/Fakt}
{Lemma}
{}
{

\faktsituation {Misalkan

\mavergleichskette
{\vergleichskette
{ Y
}
{ \subseteq }{ W
}
{ \subseteq }{ \R^n
}
{    }{
}
{   }{
}
}
{}{}{,}
$W$
\definitionsverweis {terbuka}{}{,}
dan $Y$ suatu
\definitionsverweis {hipermuka terdiferensialkan}{}{}
dan misalkan
\maabb {\gamma} { I } { Y
} {}
suatu
\definitionsverweis {kurva terdiferensialkan}{}{.}}
\faktuebergang {Maka pernyataan-pernyataan berikut berlaku.}
\faktfolgerung {\aufzaehlungzwei {Jika $F$ adalah medan vektor
\definitionsverweis {paralel}{}{} sepanjang $\gamma$ dan

\mavergleichskette
{\vergleichskette
{ a
}
{ \in }{ \R
}
{  }{
}
{    }{
}
{   }{
}
}
{}{}{,}
maka $aF$ juga merupakan medan vektor paralel.
} {Jika $F$ dan $G$ adalah medan-medan vektor paralel sepanjang $\gamma$, maka
\mathl{F +G}{} juga merupakan medan vektor paralel.
}}
\faktzusatz {}
\faktzusatz {}

 }
{ Lihat Soal 6.3. }

Untuk suatu kurva terdiferensialkan
\maabbdisp {\gamma} { I } { Y
} {}
medan kecepatan $\gamma'$ khususnya merupakan suatu medan vektor sepanjang $\gamma$; karena itu, konsep apakah medan ini paralel dapat diterapkan.



\inputfaktbeweis
{Differenzierbare Hyperfläche/Differenzierbare Kurve/Ableitung parallel/Geodätische/Fakt}
{Lemma}
{}
{

\faktsituation {Misalkan

\mavergleichskette
{\vergleichskette
{ Y
}
{ \subseteq }{ W
}
{ \subseteq }{ \R^n
}
{    }{
}
{   }{
}
}
{}{}{,}
$W$
\definitionsverweis {terbuka}{}{,}
dan $Y$ suatu
\definitionsverweis {hipermuka terdiferensialkan}{}{}
dan misalkan
\maabb {\gamma} { I } { Y
} {}
suatu kurva yang
\definitionsverweis {dua kali terdiferensialkan secara kontinu}{}{.}}
\faktfolgerung {Maka $\gamma$ merupakan suatu
\definitionsverweis {kurva geodesik}{}{}
tepat ketika medan turunannya
\zusatzklammer {medan kecepatannya} {} {} $\gamma'$ merupakan suatu
\definitionsverweis {medan vektor paralel}{}{}
sepanjang $\gamma$.}
\faktzusatz {}
\faktzusatz {}

}
{

Misalkan
\maabbdisp {\gamma^{\prime \prime}} { I } { \R^n
} {}
adalah turunan kedua. Berdasarkan definisi, $\gamma$ merupakan suatu kurva geodesik jika
\mathl{\gamma^{\prime \prime} (t)}{} selalu tegak lurus terhadap ruang singgung $T_{\gamma(t)} Y$, yang berlaku tepat ketika proyeksi ortogonal $\gamma^{\prime \prime}(t)$ pada ruang singgung sama dengan $0$. Hal ini ekuivalen dengan

\mavergleichskettedisp
{\vergleichskette
{ (\nabla_\gamma \gamma')(t)
}
{ =} { \pi_{\gamma(t)} \left(  \gamma^{\prime \prime} (t)  \right)
}
{ =} { 0
}
{ } {
}
{ } {
}
}
{}{}{}
untuk semua

\mavergleichskette
{\vergleichskette
{ t
}
{ \in }{ I
}
{  }{
}
{    }{
}
{   }{
}
}
{}{}{,}
yang berarti bahwa $\gamma'$ merupakan suatu
\definitionsverweis {medan vektor paralel}{}{}
sepanjang $\gamma$.

}

Pernyataan berikut menjadi dasar bagi keberadaan transpor paralel.


\inputfaktbeweis
{Hyperfläche/Kurve/Anfangstangentenvektor/Paralleles Vektorfeld/Existenz/Fakt}
{Teorema}
{}
{

\faktsituation {Misalkan

\mavergleichskette
{\vergleichskette
{ Y
}
{ \subseteq }{ W
}
{ \subseteq }{\R^n
}
{    }{
}
{   }{
}
}
{}{}{,}
$W$
\definitionsverweis {terbuka}{}{,}
dan $Y$ suatu
\definitionsverweis {hipermuka terdiferensialkan}{}{} kelas $C^2$
dan misalkan
\maabb {\gamma} { I } { Y
} {}
suatu
\definitionsverweis {kurva terdiferensialkan}{}{.}
Misalkan

\mavergleichskette
{\vergleichskette
{ t_0
}
{ \in }{ I
}
{  }{
}
{    }{
}
{   }{
}
}
{}{}{,}

\mavergleichskette
{\vergleichskette
{ \gamma(t_0)
}
{ = }{ P
}
{  }{
}
{    }{
}
{   }{
}
}
{}{}{}
dan misalkan

\mavergleichskette
{\vergleichskette
{ v
}
{ \in }{ T_PY
}
{  }{
}
{    }{
}
{   }{
}
}
{}{}{}
suatu
\definitionsverweis {vektor singgung}{}{.}}
\faktfolgerung {Maka terdapat suatu
\definitionsverweis {medan vektor tangensial}{}{}
$F$ yang ditentukan secara unik sepanjang $\gamma$, yang bersifat
\definitionsverweis {paralel}{}{}
dan memenuhi

\mavergleichskettedisp
{\vergleichskette
{ F(t_0)
}
{ =} { v
}
{ } {
}
{ } {
}
{ } {
}
}
{}{}{.}
}
\faktzusatz {}
\faktzusatz {}

}
{

Pilih suatu medan normal satuan terdiferensialkan $N$ sepanjang $\gamma$; regularitas $C^2$ pada $Y$ menjamin pilihan lokal tersebut dapat diteruskan sepanjang interval $I$. Kita mencari suatu pemetaan terdiferensialkan
\maabbdisp {F} { I } { \R^n
} {,}
yang memenuhi ketiga syarat
\aufzaehlungdrei{
\mavergleichskettedisp
{\vergleichskette
{ F(t)
}
{ \in} { T_{\gamma(t)} Y
}
{ } {
}
{ } {
}
{ } {
}
}
{}{}{}
untuk semua

\mavergleichskette
{\vergleichskette
{ t
}
{ \in }{ I
}
{  }{
}
{    }{
}
{   }{
}
}
{}{}{,}
}{
\mavergleichskettedisp
{\vergleichskette
{ F'(t)
}
{ \in} { N_{\gamma(t)} Y
}
{ } {
}
{ } {
}
{ } {
}
}
{}{}{}
untuk semua

\mavergleichskette
{\vergleichskette
{ t
}
{ \in }{ I
}
{  }{
}
{    }{
}
{   }{
}
}
{}{}{,}
}{
\mavergleichskettedisp
{\vergleichskette
{ F (t_0)
}
{ =} { v
}
{ } {
}
{ } {
}
{ } {
}
}
{}{}{.}
}
Syarat pertama berarti bahwa $F(t)$ tegak lurus terhadap vektor normal satuan $N (\gamma(t))$, sehingga

\mavergleichskettedisp
{\vergleichskette
{ \left\langle F( t) ,  N( \gamma(t))   \right\rangle
}
{ =} { 0
}
{ } {
}
{ } {
}
{ } {
}
}
{}{}{.}
Oleh karena itu, berlaku pula

\mavergleichskettedisp
{\vergleichskette
{ 0
}
{ =} { \left\langle F(t)  ,  N( \gamma(t))   \right\rangle'
}
{ =} { \left\langle F'(t)  ,  N( \gamma(t))   \right\rangle +  \left\langle F(t)  ,  (N( \gamma(t)))'   \right\rangle
}
{ } {
}
{ } {
}
}
{}{}{.}
Syarat kedua, bahwa $F'(t)$ berada di ruang normal, berarti

\mavergleichskettedisp
{\vergleichskette
{ F'(t) - \left\langle F'( t) ,  N( \gamma(t))   \right\rangle N(\gamma(t) )
}
{ =} { 0
}
{ } {
}
{ } {
}
{ } {
}
}
{}{}{}
untuk semua

\mavergleichskette
{\vergleichskette
{ t
}
{ \in }{ I
}
{  }{
}
{    }{
}
{   }{
}
}
{}{}{,}
yang dengan bantuan persamaan sebelumnya dapat kita ubah menjadi

\mavergleichskettedisp
{\vergleichskette
{ F'(t) + \left\langle F( t) ,  (N( \gamma(t)))'   \right\rangle N(\gamma(t) )
}
{ =} { 0
}
{ } {
}
{ } {
}
{ } {
}
}
{}{}{}
untuk semua

\mavergleichskette
{\vergleichskette
{ t
}
{ \in }{ I
}
{  }{
}
{    }{
}
{   }{
}
}
{}{}{}
atau menjadi

\mavergleichskettedisp
{\vergleichskette
{ F'(t)
}
{ =} { - \left\langle F( t) ,  (N( \gamma(t)))'   \right\rangle N(\gamma(t) )
}
{ } {
}
{ } {
}
{ } {
}
}
{}{}{}
untuk semua

\mavergleichskette
{\vergleichskette
{ t
}
{ \in }{ I
}
{  }{
}
{    }{
}
{   }{
}
}
{}{}{.}
Di sini terdapat suatu persamaan diferensial linear biasa orde pertama untuk $F$ yang, bersama dengan syarat awal

\mavergleichskette
{\vergleichskette
{ F(t_0)
}
{ = }{ v
}
{  }{
}
{    }{
}
{   }{
}
}
{}{}{,}
menurut
Teorema 56.2 (Analisis (Osnabrück 2021--2023)),
memiliki solusi lokal tunggal. Karena persamaan ini linear dengan koefisien kontinu pada $I$, solusi tersebut dapat dilanjutkan secara unik ke seluruh interval $I$. Karena $v\in T_PY$, nilai awal memenuhi $\langle v,N(P)\rangle=0$. Sepanjang solusi persamaan diferensial tersebut,
\[
\frac{d}{dt}\left\langle F(t),N(\gamma(t))\right\rangle
=-\left\langle F(t),(N\circ\gamma)'(t)\right\rangle
+\left\langle F(t),(N\circ\gamma)'(t)\right\rangle=0.
\]
Jadi $F(t)$ tetap tangensial dan persamaan diferensial membuat $F'(t)$ tetap normal. Dengan demikian, solusi tersebut memang merupakan medan paralel yang diminta dan keunikan persamaan diferensial memberi keunikannya.

\emph{Catatan edisi: Sumber hanya mengasumsikan hipermuka terdiferensialkan, menerapkan teorema eksistensi lokal langsung pada seluruh interval $I$, dan tidak membuktikan bahwa solusi persamaan diferensial tetap tangensial. Hipotesis diperkuat menjadi $C^2$, kelanjutan global untuk sistem linear dinyatakan, medan normal diperkenalkan secara eksplisit, dan perhitungan invarian yang hilang dilengkapi.}

}

\inputbemerkung
{}
{

Dalam situasi
Teorema 6.9,
persamaan diferensial yang harus diselesaikan adalah sistem

\mavergleichskettedisp
{\vergleichskette
{ F_i'(t)
}
{ =} { - \left(  \sum_{j = 1}^n  F_j (t) (N ( \gamma(t)))'_j  \right)   N(\gamma(t))_i
}
{ } {
}
{ } {
}
{ } {
}
}
{}{}{}
untuk

\mavergleichskette
{\vergleichskette
{ i
}
{ = }{ 1 , \ldots , n
}
{  }{
}
{    }{
}
{   }{
}
}
{}{}{}
atau, dalam notasi matriks,

\mavergleichskettedisp
{\vergleichskette
{ \begin{pmatrix} F_1' \\ \vdots\\ F_n'   \end{pmatrix}
}
{ =} { - A \begin{pmatrix} F_1  \\ \vdots\\ F_n   \end{pmatrix}
}
{ } {
}
{ } {
}
{ } {
}
}
{}{}{}
dengan entri-entri

\mavergleichskettedisp
{\vergleichskette
{ A_{ij} (t)
}
{ =} { N(\gamma(t))_i (N(\gamma(t)))'_j
}
{ } {
}
{ } {
}
{ } {
}
}
{}{}{.}

\emph{Catatan edisi: Satu tanda kurung tutup berlebih pada faktor $(N(\gamma(t)))'_j$ dihapus, sehingga koefisiennya adalah $A_{ij}=N_iN'_j$.}

}

\inputbeispiel{}
{

Kita meninjau seperempat lingkaran besar
\maabbeledisp {} {[0 , \pi/2] } { \R^3
} { t } { \left(  \cos  t   , \,   \sin  t  , \,   0                 \right)
} {,}
pada permukaan bola satuan $S$. Berlaku

\mavergleichskette
{\vergleichskette
{ \gamma(0)
}
{ = }{ \left(  1   , \,   0  , \,   0                 \right)
}
{ = }{ P
}
{    }{
}
{   }{
}
}
{}{}{}
dan

\mavergleichskette
{\vergleichskette
{ \gamma(\pi/2)
}
{ = }{ \left(  0   , \,   1  , \,   0                 \right)
}
{ = }{ Q
}
{    }{
}
{   }{
}
}
{}{}{}
dengan
\definitionsverweis {ruang-ruang singgung}{}{}

\mavergleichskette
{\vergleichskette
{ T_PS
}
{ = }{ \R e_2 +\R e_3
}
{  }{
}
{    }{
}
{   }{
}
}
{}{}{}
dan

\mavergleichskette
{\vergleichskette
{ T_QS
}
{ = }{ \R e_1 +\R e_3
}
{  }{
}
{    }{
}
{   }{
}
}
{}{}{.}
Vektor normal satuan yang mengarah ke dalam sepanjang lintasan tersebut adalah

\mavergleichskette
{\vergleichskette
{ N(t)
}
{ = }{ -\gamma(t)
}
{  }{
}
{    }{
}
{   }{
}
}
{}{}{.}
Matriks dari
Catatan 6.10,
yang mendeskripsikan persamaan diferensial bagi suatu medan vektor paralel, adalah

\mathdisp {\begin{pmatrix}  -\sin  t \cos  t   &   \cos  t \cos  t  &  0  \\  - \sin  t  \sin  t  &  \sin  t  \cos  t  &  0  \\ 0 &  0  &  0 \end{pmatrix}} { . }

Fungsi konstan

\mavergleichskettedisp
{\vergleichskette
{ F(t)
}
{ =} { e_3
}
{ =} { \begin{pmatrix}  0  \\ 0\\  1   \end{pmatrix}
}
{ } {
}
{ } {
}
}
{}{}{}
merupakan suatu solusi. Selanjutnya,

\mavergleichskettedisp
{\vergleichskette
{ G(t)
}
{ =} { \begin{pmatrix}  - \sin  t  \\ \cos  t \\  0   \end{pmatrix}
}
{ } {
}
{ } {
}
{ } {
}
}
{}{}{}
merupakan suatu solusi; di satu pihak,

\mavergleichskettedisp
{\vergleichskette
{ G'(t)
}
{ =} { \begin{pmatrix}  - \cos  t  \\ - \sin  t \\  0   \end{pmatrix}
}
{ } {
}
{ } {
}
{ } {
}
}
{}{}{}
dan, di pihak lain,

\mavergleichskettedisphandlinks
{\vergleichskettedisphandlinks
{ - \begin{pmatrix}  -\sin  t \cos  t   &   \cos  t \cos  t  &  0  \\  - \sin  t  \sin  t  &  \sin  t  \cos  t  &  0  \\ 0 &  0  &  0 \end{pmatrix} \begin{pmatrix}  - \sin  t  \\ \cos  t \\  0   \end{pmatrix}
}
{ =} { - \begin{pmatrix}  \sin^{ 2  }  t \cos  t + \cos^{ 3  }  t  \\ \sin^{ 3  }  t +  \sin  t \cos^{ 2  }  t \\  0   \end{pmatrix}
}
{ =} { - \begin{pmatrix}  \cos  t  \\ \sin  t \\  0   \end{pmatrix}
}
{ } {
}
{ } {
}
}
{}{}{.}

}

  \zwischenueberschrift{Transpor Paralel}

Misalkan

\mavergleichskette
{\vergleichskette
{ Y
}
{ \subseteq }{ W
}
{ \subseteq }{ \R^n
}
{    }{
}
{   }{
}
}
{}{}{,}
$W$
\definitionsverweis {terbuka}{}{,}
dan $Y$ suatu
\definitionsverweis {hipermuka terdiferensialkan}{}{} kelas $C^2$.
Untuk titik-titik

\mavergleichskette
{\vergleichskette
{ P,Q
}
{ \in }{ Y
}
{  }{
}
{    }{
}
{   }{
}
}
{}{}{}
tidak terdapat hubungan alami antara
\definitionsverweis {ruang singgung}{}{}
$T_PY$ dan ruang singgung $T_QY$. Misalkan
\maabbdisp {\gamma} {[a,b]} {Y
} {}
suatu
\definitionsverweis {kurva terdiferensialkan}{}{}
di $Y$ dengan

\mavergleichskette
{\vergleichskette
{ \gamma(a)
}
{ = }{ P
}
{  }{
}
{    }{
}
{   }{
}
}
{}{}{}
dan

\mavergleichskette
{\vergleichskette
{ \gamma(b)
}
{ = }{ Q
}
{  }{
}
{    }{
}
{   }{
}
}
{}{}{.}
Kita dapat bertanya apakah sepanjang lintasan ini terdapat hubungan yang bermakna antara ruang-ruang singgung tersebut.

Untuk suatu kurva terdiferensialkan tetap $\gamma$ di $Y$ yang menghubungkan titik-titik
\mathkor {} {P =\gamma(a)} {dan} {Q = \gamma(b)} {,}
berdasarkan
Teorema 6.9,
ditentukan suatu pemetaan
\maabbdisp {} {T_PY} {T_QY
} {.}
Pemetaan ini memetakan vektor awal

\mavergleichskette
{\vergleichskette
{ v
}
{ \in }{ T_PY
}
{  }{
}
{    }{
}
{   }{
}
}
{}{}{}
ke vektor

\mavergleichskette
{\vergleichskette
{ F(b)
}
{ \in }{ T_QY
}
{  }{
}
{    }{
}
{   }{
}
}
{}{}{,}
dengan $F$ adalah medan vektor paralel yang ditentukan secara unik sepanjang $\gamma$. Pemetaan ini disebut \stichwort {transpor paralel} {} sepanjang $\gamma$. Pemetaan tersebut dilambangkan dengan
\maabbdisp {\Psi_\gamma} { T_{\gamma(a)}Y} {T_{\gamma(b)}Y
} {.}

\bild{ \begin{center}
\includegraphics[width=5.5cm]{ \bildeinlesungsvg {Parallel transport sphere2} {svg} }
\end{center}
\bildtext {} }

\bildlizenz { Parallel transport sphere2.svg } {} {Silly rabbit} {en Wikipedia} {CC-by-sa 3.0} {}



\inputfaktbeweis
{Hyperfläche/Kurve/Paralleltransport/Isometrie/Fakt}
{Teorema}
{}
{

\faktsituation {Misalkan

\mavergleichskette
{\vergleichskette
{ Y
}
{ \subseteq }{ W
}
{ \subseteq }{ \R^n
}
{    }{
}
{   }{
}
}
{}{}{,}
$W$
\definitionsverweis {terbuka}{}{,}
dan $Y$ suatu
\definitionsverweis {hipermuka terdiferensialkan}{}{} kelas $C^2$
dan misalkan
\maabb {\gamma} {[a,b]} {Y
} {}
suatu
\definitionsverweis {kurva terdiferensialkan}{}{}
dengan

\mavergleichskette
{\vergleichskette
{ \gamma(a)
}
{ = }{ P
}
{  }{
}
{    }{
}
{   }{
}
}
{}{}{}
dan

\mavergleichskette
{\vergleichskette
{ \gamma(b)
}
{ = }{ Q
}
{  }{
}
{    }{
}
{   }{
}
}
{}{}{.}}
\faktfolgerung {Maka
\definitionsverweis {transpor paralel}{}{}
sepanjang $\gamma$ merupakan suatu
\definitionsverweis {isometri}{}{}
\maabbdisp {} {T_PY} {T_QY
} {.}}
\faktzusatz {}
\faktzusatz {}

}
{

Untuk membuktikan linearitas, misalkan

\mavergleichskette
{\vergleichskette
{ v,w
}
{ \in }{ T_PY
}
{  }{
}
{    }{
}
{   }{
}
}
{}{}{}
dan

\mavergleichskette
{\vergleichskette
{ r,s
}
{ \in }{ \R
}
{  }{
}
{    }{
}
{   }{
}
}
{}{}{}
diberikan. Misalkan
\mathkor {} {F} {dan} {G} {}
adalah medan-medan vektor paralel sepanjang $\gamma$ yang ditentukan secara unik menurut
Teorema 6.9
dengan

\mavergleichskette
{\vergleichskette
{ F(a)
}
{ = }{ v
}
{  }{
}
{    }{
}
{   }{
}
}
{}{}{}
dan

\mavergleichskette
{\vergleichskette
{ G(a)
}
{ = }{ w
}
{  }{
}
{    }{
}
{   }{
}
}
{}{}{.}
Menurut
Lemma 6.7,
\mathl{rF+sG}{} merupakan suatu medan vektor paralel dengan

\mavergleichskette
{\vergleichskette
{ (rF+sG)(a)
}
{ = }{ rv+sw
}
{  }{
}
{    }{
}
{   }{
}
}
{}{}{.}
\emph{Catatan edisi: Nilai awal kombinasi linear diperbaiki dari $v+w$ menjadi $rv+sw$.}
Berdasarkan keunikan dari
Teorema 6.9,
\mathl{rF+sG}{} dengan demikian merupakan medan vektor paralel untuk vektor singgung
\mathl{rv+sw}{.} Oleh karena itu,

\mavergleichskettedisp
{\vergleichskette
{ \Psi_{\gamma} (rv+sw)
}
{ =} { (rF+sG)(b)
}
{ =} { r F(b) +s G(b)
}
{ =} { r \Psi_{\gamma} (v) + s \Psi_{\gamma} (w)
}
{ } {
}
}
{}{}{.}

Untuk membuktikan kompatibilitas dengan hasil kali dalam, misalkan kembali

\mavergleichskette
{\vergleichskette
{ v,w
}
{ \in }{ T_PY
}
{  }{
}
{    }{
}
{   }{
}
}
{}{}{}
diberikan, dan misalkan $F,G$ adalah medan-medan vektor paralel yang bersesuaian. Berlaku

\mavergleichskettedisp
{\vergleichskette
{ \left\langle  F(t) , G(t)  \right\rangle'
}
{ =} { \left\langle  F'(t) , G(t)  \right\rangle  +  \left\langle  F(t) , G'(t)  \right\rangle
}
{ =} { 0
}
{ } {
}
{ } {
}
}
{}{}{,}
karena $F,G$ bersifat tangensial, sedangkan $F',G'$ ortogonal terhadap ruang singgung. Jadi,
\mathl{\left\langle  F(t) , G(t)  \right\rangle}{} konstan sepanjang lintasan. Oleh karena itu,

\mavergleichskettedisp
{\vergleichskette
{ \left\langle \Psi_\gamma(v)  , \Psi_\gamma(w)  \right\rangle
}
{ =} { \left\langle  F(b)  ,  G(b)   \right\rangle
}
{ =} { \left\langle  F(a)  ,  G(a)   \right\rangle
}
{ =} { \left\langle  v  , w  \right\rangle
}
{ } {
}
}
{}{}{.}

\emph{Catatan edisi: Faktor kedua pada hasil kali dalam di titik awal diperbaiki dari $F(a)$ menjadi $G(a)$.}
Bijektivitasnya juga jelas dari sini.

}

\inputbeispiel{}
{

Kita meninjau seperempat lingkaran besar
\maabbeledisp {} {[0 , \pi/2] } { \R^3
} { t } { \left(  \cos  t   , \,   \sin  t  , \,   0                 \right)
} {,}
pada permukaan bola satuan $S$ dan melanjutkan pembahasan dari
Contoh 6.11.
Berlaku

\mavergleichskette
{\vergleichskette
{ \gamma(0)
}
{ = }{ \left(  1   , \,   0  , \,   0                 \right)
}
{ = }{ P
}
{    }{
}
{   }{
}
}
{}{}{}
dan

\mavergleichskette
{\vergleichskette
{ \gamma(\pi/2)
}
{ = }{ \left(  0   , \,   1  , \,   0                 \right)
}
{ = }{ Q
}
{    }{
}
{   }{
}
}
{}{}{}
dengan
\definitionsverweis {ruang-ruang singgung}{}{}

\mavergleichskette
{\vergleichskette
{ T_PS
}
{ = }{ \R e_2 +\R e_3
}
{  }{
}
{    }{
}
{   }{
}
}
{}{}{}
dan

\mavergleichskette
{\vergleichskette
{ T_QS
}
{ = }{ \R e_1 +\R e_3
}
{  }{
}
{    }{
}
{   }{
}
}
{}{}{.}
Berdasarkan perhitungan dalam contoh yang dirujuk tersebut, untuk
\definitionsverweis {transpor paralel}{}{}
berlaku

\mavergleichskette
{\vergleichskette
{ \Psi_\gamma(e_3)
}
{ = }{ e_3
}
{  }{
}
{    }{
}
{   }{
}
}
{}{}{}
dan

\mavergleichskette
{\vergleichskette
{ \Psi_\gamma(e_2)
}
{ = }{ - e_1
}
{  }{
}
{    }{
}
{   }{
}
}
{}{}{.}

}

Jika
\maabbdisp {\gamma} { I } { Y
} {}
merupakan suatu kurva kontinu yang terdiferensialkan sepotong-sepotong, transpor paralel sepanjang $\gamma$ didefinisikan dengan melaksanakan secara berurutan transpor-transpor paralel pada setiap potongan kurva yang terdiferensialkan.

\chapter{Lembar Kerja 6}
\setcounter{section}{6}

  \zwischenueberschrift{Soal latihan}

\inputaufgabe
{}
{

Misalkan $Y$ adalah
\definitionsverweis {grafik}{}{}
fungsi

\mavergleichskette
{\vergleichskette
{ f(x,y)
}
{ = }{ x^2-y^2
}
{  }{
}
{    }{
}
{   }{
}
}
{}{}{.}
\aufzaehlungdrei{Buktikan bahwa

\mavergleichskettedisp
{\vergleichskette
{ F(x,y,z)
}
{ =} { \begin{pmatrix}  y  \\ x \\  0   \end{pmatrix}
}
{ } {
}
{ } {
}
{ } {
}
}
{}{}{}
memberikan suatu
\definitionsverweis {medan vektor tangensial}{}{}
pada $Y$.
}{Misalkan

\mavergleichskette
{\vergleichskette
{ P
}
{ = }{ \begin{pmatrix}  2  \\ 3 \\  -5   \end{pmatrix}
}
{ \in }{ Y
}
{    }{
}
{   }{
}
}
{}{}{.}
Hitung
\mathl{\nabla_v F}{} untuk

\mavergleichskette
{\vergleichskette
{ v
}
{ = }{ \begin{pmatrix}  2  \\ 1 \\  2   \end{pmatrix}
}
{ \in }{ T_PY
}
{    }{
}
{   }{
}
}
{}{}{.}
\emph{Catatan edisi: Tanda sama dengan pada sumber diperbaiki menjadi tanda keanggotaan, karena $v$ adalah vektor di $T_PY$.}
}{Verifikasikan bahwa $\nabla_v F$ tangensial pada $Y$ di $P$.
}

}
{} {}

\inputaufgabegibtloesung
{}
{

Pada
\definitionsverweis {permukaan bola satuan}{}{}
$Y$ yang berdimensi dua, kita tinjau
\definitionsverweis {medan vektor tangensial}{}{}

\mavergleichskettedisp
{\vergleichskette
{ F \begin{pmatrix}  x  \\ y \\ z   \end{pmatrix}
}
{ =} { \begin{pmatrix}  y  \\ -x\\  0   \end{pmatrix}
}
{ } {
}
{ } {
}
{ } {
}
}
{}{}{.}
\aufzaehlungdrei{Tentukan
\definitionsverweis {matriks Jacobi}{}{}
dari $F$ pada $\R^3$.
}{Untuk titik

\mavergleichskette
{\vergleichskette
{ P
}
{ = }{ \begin{pmatrix}  0  \\ 1 \\  0   \end{pmatrix}
}
{ \in }{ Y
}
{    }{
}
{   }{
}
}
{}{}{}
tentukan suatu matriks bagi pemetaan
\maabbeledisp {} {T_PY} { T_PY
} { v } { \nabla_v F
} {,}
terhadap suatu
\definitionsverweis {basis}{}{}
yang sesuai.
}{Untuk titik

\mavergleichskette
{\vergleichskette
{ Q
}
{ = }{ \begin{pmatrix}  0  \\ 0\\  1   \end{pmatrix}
}
{ \in }{ Y
}
{    }{
}
{   }{
}
}
{}{}{}
tentukan suatu matriks bagi pemetaan
\maabbeledisp {} {T_QY} { T_QY
} { v } { \nabla_v F
} {,}
terhadap suatu
\definitionsverweis {basis}{}{}
yang sesuai.
}

}
{} {}

\inputaufgabe
{}
{

Misalkan

\mavergleichskette
{\vergleichskette
{ Y
}
{ \subseteq }{ W
}
{ \subseteq }{ \R^n
}
{    }{
}
{   }{
}
}
{}{}{,}
dengan $W$
\definitionsverweis {terbuka}{}{,}
adalah suatu
\definitionsverweis {hipermuka terdiferensialkan}{}{,}
dan misalkan
\maabb {\gamma} { I } { Y
} {}
adalah suatu
\definitionsverweis {kurva terdiferensialkan}{}{.}
Buktikan pernyataan-pernyataan berikut.
\aufzaehlungzwei {Jika $F$ adalah suatu
\definitionsverweis {medan vektor paralel}{}{}
sepanjang $\gamma$ dan

\mavergleichskette
{\vergleichskette
{ a
}
{ \in }{ \R
}
{  }{
}
{    }{
}
{   }{
}
}
{}{}{,}
maka \mathl{aF}{} juga merupakan medan vektor paralel.
} {Jika $F$ dan $G$ adalah medan-medan vektor paralel sepanjang $\gamma$, maka
\mathl{F +G}{} juga merupakan medan vektor paralel.
}

}
{} {}

\inputaufgabe
{}
{

Misalkan

\mavergleichskette
{\vergleichskette
{ Y
}
{ \subseteq }{ \R^n
}
{  }{
}
{    }{
}
{   }{
}
}
{}{}{}
adalah suatu
\definitionsverweis {hiperbidang}{}{.}
Buktikan bahwa setiap
\definitionsverweis {medan vektor tangensial}{}{}
$F$ yang konstan
\zusatzklammer {yang bersesuaian dengan suatu medan vektor pada

\mavergleichskette
{\vergleichskette
{ Y
}
{ \cong }{\R^{n-1}
}
{  }{
}
{    }{
}
{   }{
}
}
{}{}{}} {} {}
bersifat
\definitionsverweis {paralel}{}{}
\zusatzklammer {sepanjang setiap kurva} {}{}{.}

\emph{Catatan edisi: Syarat ``konstan'' ditambahkan; medan tangensial umum pada hiperbidang tidak harus paralel.}

}
{} {}

\inputaufgabe
{}
{

Misalkan

\mavergleichskettedisp
{\vergleichskette
{ Y
}
{ =} { \left\{  \begin{pmatrix}  x  \\ y \\ z   \end{pmatrix}   \mid   x^2+y^2+z^2 = 1        \right\}
}
{ } {
}
{ } {
}
{ } {
}
}
{}{}{}
adalah
\definitionsverweis {permukaan bola satuan}{}{}
berdimensi dua dan
\maabbele {\gamma} {[0,2\pi]} {Y
} { t } { \begin{pmatrix}  \cos  t  \\ 0 \\  \sin  t   \end{pmatrix}
} {.}
Tentukan mana di antara medan-medan vektor berikut
\maabbdisp {F} {[0,2\pi]} { \R^3
} {}
sepanjang $\gamma$ yang tangensial pada $Y$ dan mana yang
\definitionsverweis {paralel}{}{.}
\aufzaehlungfuenf{
\mavergleichskette
{\vergleichskette
{ F (t)
}
{ = }{ \begin{pmatrix}  2  \\ 0 \\  1   \end{pmatrix}
}
{  }{
}
{    }{
}
{   }{
}
}
{}{}{.}
}{
\mavergleichskette
{\vergleichskette
{ F (t)
}
{ = }{ \begin{pmatrix}  0  \\ 5 \\  0   \end{pmatrix}
}
{  }{
}
{    }{
}
{   }{
}
}
{}{}{.}
}{
\mavergleichskette
{\vergleichskette
{ F (t)
}
{ = }{ \begin{pmatrix}  \cos  t  \\ 0 \\  \sin  t   \end{pmatrix}
}
{  }{
}
{    }{
}
{   }{
}
}
{}{}{.}
}{
\mavergleichskette
{\vergleichskette
{ F (t)
}
{ = }{ \begin{pmatrix}  - \sin  t  \\ 0 \\  \cos  t   \end{pmatrix}
}
{  }{
}
{    }{
}
{   }{
}
}
{}{}{.}
}{
\mavergleichskette
{\vergleichskette
{ F (t)
}
{ = }{ \begin{pmatrix}  t  \\e^t\\  \sin  t   \end{pmatrix}
}
{  }{
}
{    }{
}
{   }{
}
}
{}{}{.}
}

}
{} {}

\inputaufgabegibtloesung
{}
{

Misalkan

\mavergleichskette
{\vergleichskette
{ Y
}
{ \subseteq }{ W
}
{ \subseteq }{ \R^3
}
{    }{
}
{   }{
}
}
{}{}{,}
dengan $W$
\definitionsverweis {terbuka}{}{,}
adalah suatu
\definitionsverweis {permukaan terdiferensialkan}{}{}
yang dilengkapi dengan
\definitionsverweis {medan normal satuan}{}{}
$N$. Misalkan
\maabb {\gamma} { I } { Y
} {}
adalah suatu kurva yang terdiferensialkan secara kontinu dan
\maabbdisp {F} { I } { \R^3
} {}
adalah suatu
\definitionsverweis {medan vektor tangensial}{}{}
terdiferensialkan sepanjang $\gamma$ yang
\definitionsverweis {paralel}{}{.}
Buktikan bahwa medan vektor
\maabbeledisp {} { I } { \R^3
} { t } { N(\gamma(t)) \times F(t)
} {,}
juga paralel, dengan $\times$ menyatakan
\definitionsverweis {hasil kali silang}{}{}
di $\R^3$.

}
{} {}

\inputaufgabe
{}
{

Misalkan

\mavergleichskettedisp
{\vergleichskette
{ h(x,y,z)
}
{ =} { x^3-2y^4-3z^5
}
{ } {
}
{ } {
}
{ } {
}
}
{}{}{}
dan

\mavergleichskette
{\vergleichskette
{ Y
}
{ = }{ h^{-1}(1)
}
{  }{
}
{    }{
}
{   }{
}
}
{}{}{.}
Misalkan
\maabbeledisp {\gamma} { (-3^{-1/5},\infty)} {Y
} { t } { \begin{pmatrix}  \sqrt[3]{ 1+3t^5  } \\ 0 \\ t   \end{pmatrix}
} {,}
adalah suatu lintasan terdiferensialkan pada $Y$. Tentukan persamaan diferensial biasa dalam bentuk
Catatan 6.10
untuk
\definitionsverweis {medan-medan vektor paralel}{}{}
sepanjang $\gamma$.

\emph{Catatan edisi: Domain sumber $\R$ dibatasi pada interval di atas karena komponen akar pangkat tiga tidak terdiferensialkan di $t=-3^{-1/5}$.}

}
{} {}

\inputaufgabe
{}
{

Misalkan

\mavergleichskette
{\vergleichskette
{ Y
}
{ \subseteq }{ \R^n
}
{  }{
}
{    }{
}
{   }{
}
}
{}{}{}
adalah suatu
\definitionsverweis {hiperbidang}{}{.}
Buktikan bahwa
\definitionsverweis {transpor paralel}{}{}
sepanjang setiap kurva pada $Y$ adalah identitas setelah semua
\definitionsverweis {ruang singgung}{}{}
$T_PY$ diidentifikasi secara kanonis dengan ruang arah linear $V=Y-P$ dari hiperbidang tersebut.

\emph{Catatan edisi: Sumber mengidentifikasi ruang singgung hiperbidang afin langsung dengan $Y$. Identifikasi kanonis sebenarnya adalah dengan ruang arah linear $V=Y-P$; pernyataan diperjelas sesuai identifikasi tersebut.}

}
{} {}

\inputaufgabegibtloesung
{}
{

Misalkan

\mavergleichskette
{\vergleichskette
{ Y
}
{ \subseteq }{ W
}
{ \subseteq }{ \R^n
}
{    }{
}
{   }{
}
}
{}{}{,}
dengan $W$
\definitionsverweis {terbuka}{}{,}
adalah suatu
\definitionsverweis {hipermuka terdiferensialkan}{}{,}
dan misalkan
\maabb {\gamma} {I =[a,b]} {Y
} {}
adalah suatu
\definitionsverweis {kurva geodesik}{}{.}
Buktikan bahwa
\definitionsverweis {transpor paralel}{}{}
sepanjang $\gamma$ membawa vektor singgung

\mavergleichskette
{\vergleichskette
{ \gamma'(a)
}
{ \in }{ T_{\gamma(a)}Y
}
{  }{
}
{    }{
}
{   }{
}
}
{}{}{}
ke vektor singgung

\mavergleichskette
{\vergleichskette
{ \gamma'(b)
}
{ \in }{ T_{\gamma(b)}Y
}
{  }{
}
{    }{
}
{   }{
}
}
{}{}{.}

}
{} {}

\inputaufgabe
{}
{

Misalkan $Y$ adalah
\definitionsverweis {permukaan bola satuan}{}{}
berdimensi dua dan
\maabb {\gamma} {[0,2\pi]} {Y
} {}
adalah
\definitionsverweis {parametrisasi panjang busur}{}{}
dari suatu
\definitionsverweis {lingkaran besar}{}{}
dengan

\mavergleichskette
{\vergleichskette
{ \gamma(0)
}
{ = }{ \gamma(2 \pi)
}
{ = }{ P
}
{ \in   }{ Y
}
{   }{
}
}
{}{}{.}
Buktikan bahwa
\definitionsverweis {transpor paralel}{}{}
yang bersesuaian adalah identitas pada $T_PY$.

}
{} {}

\inputaufgabe
{}
{

Misalkan

\mavergleichskettedisp
{\vergleichskette
{ Y
}
{ =} { \left\{  \begin{pmatrix}  x  \\ y \\ z   \end{pmatrix}   \mid   x^2+y^2+z^2 = 1        \right\}
}
{ } {
}
{ } {
}
{ } {
}
}
{}{}{}
adalah
\definitionsverweis {permukaan bola satuan}{}{}
berdimensi dua,

\mavergleichskette
{\vergleichskette
{ P
}
{ = }{ \begin{pmatrix}  { \frac{   1  }{  2 } }  \\ 0 \\  { \frac{   \sqrt{3}  }{  2 } }   \end{pmatrix}
}
{  }{
}
{    }{
}
{   }{
}
}
{}{}{}
dan
\maabbele {\gamma} {[0,2\pi]} {Y
} { t } { \begin{pmatrix}  { \frac{   1  }{  2  } } \cos  t  \\ { \frac{   1  }{  2  } }  \sin  t \\  { \frac{   \sqrt{3}  }{  2  } }   \end{pmatrix}
} {.}
Deskripsikan
\definitionsverweis {transpor paralel}{}{}
\maabb {\Psi_\gamma} { T_PY} {T_PY
} {}
terhadap suatu basis yang sesuai.

}
{} {}

\inputaufgabe
{}
{

Misalkan

\mavergleichskette
{\vergleichskette
{ Y
}
{ \subseteq }{ W
}
{ \subseteq }{ \R^n
}
{    }{
}
{   }{
}
}
{}{}{,}
dengan $W$
\definitionsverweis {terbuka}{}{,}
adalah suatu
\definitionsverweis {hipermuka terdiferensialkan}{}{,}
dan misalkan

\mavergleichskette
{\vergleichskette
{ P,Q
}
{ \in }{ Y
}
{  }{
}
{    }{
}
{   }{
}
}
{}{}{}
adalah titik-titik dengan
\definitionsverweis {ruang-ruang singgung}{}{}
\mathkor {} {T_PY} {dan} {T_QY} {.}
Kita tinjau pemetaan
\maabbeledisp {\varphi_{P,Q}} {T_PY} { T_QY
} { v } { \pi_Q(v)
} {,}
yang memandang suatu vektor singgung di $P$ sebagai vektor di ruang sekitar $\R^n$, lalu memproyeksikannya secara ortogonal ke $T_QY$.
\aufzaehlungvier{Apakah
\mathl{\varphi_{P,Q}}{} linear?
}{Apakah
\mathl{\varphi_{P,Q}}{} bijektif?
}{Apakah
\mathl{\varphi_{P,Q}}{} suatu isometri?
}{Apakah terdapat hubungan antara
\mathl{\varphi_{P,Q}}{} dan
\definitionsverweis {transpor paralel}{}{} $\Psi_\gamma$ sepanjang suatu kurva terdiferensialkan
\maabb {\gamma} {[a,b]} { Y
} {}
dengan

\mavergleichskette
{\vergleichskette
{ \gamma(a)
}
{ = }{ P
}
{  }{
}
{    }{
}
{   }{
}
}
{}{}{}
dan

\mavergleichskette
{\vergleichskette
{ \gamma(b)
}
{ = }{ Q
}
{  }{
}
{    }{
}
{   }{
}
}
{}{}{?}
}

}
{} {}

Misalkan

\mavergleichskette
{\vergleichskette
{ W
}
{ \subseteq }{ \R^n
}
{  }{
}
{    }{
}
{   }{
}
}
{}{}{}
\definitionsverweis {terbuka}{}{,}
\maabb {h} { W } { \R
} {}
adalah suatu
\definitionsverweis {fungsi yang dua kali terdiferensialkan secara kontinu}{}{}
dan

\mavergleichskette
{\vergleichskette
{ Y
}
{ = }{ h^{-1}(c)
}
{  }{
}
{    }{
}
{   }{
}
}
{}{}{}
adalah
\definitionsverweis {serat}{}{}
di atas

\mavergleichskette
{\vergleichskette
{ c
}
{ \in }{ \R
}
{  }{
}
{    }{
}
{   }{
}
}
{}{}{,}
dengan $h$
\definitionsverweis {reguler}{}{}
di setiap titik pada $Y$. Untuk suatu titik

\mavergleichskette
{\vergleichskette
{ P
}
{ \in }{ Y
}
{  }{
}
{    }{
}
{   }{
}
}
{}{}{,}
\definitionsverweis {subgrup}{}{}

\mavergleichskette
{\vergleichskette
{ H
}
{ \subseteq }{ \operatorname{Iso} \left(  T_PY  \right)
}
{  }{
}
{    }{
}
{   }{
}
}
{}{}{,}
yang terdiri atas semua
\definitionsverweis {transpor paralel}{}{}
\maabbdisp {\Psi_\gamma} {T_PY} {T_PY
} {}
yang diinduksi oleh
\definitionsverweis {lintasan tertutup}{}{}
$\gamma$ yang terdiferensialkan sepotong-sepotong pada $Y$ dari $P$ ke $P$, disebut
\definitionswort {grup holonomi}{}
di $P$.

\emph{Catatan edisi: Regularitas sumber diperkuat dari $C^1$ menjadi $C^2$ agar medan normal dapat diturunkan dan transpor paralel yang dipakai dalam definisi ini terjamin.}

\inputaufgabe
{}
{

Misalkan

\mavergleichskette
{\vergleichskette
{ W
}
{ \subseteq }{ \R^n
}
{  }{
}
{    }{
}
{   }{
}
}
{}{}{}
\definitionsverweis {terbuka}{}{,}
\maabb {h} { W } { \R
} {}
adalah suatu
\definitionsverweis {fungsi yang dua kali terdiferensialkan secara kontinu}{}{}
dan

\mavergleichskette
{\vergleichskette
{ Y
}
{ = }{ h^{-1}(c)
}
{  }{
}
{    }{
}
{   }{
}
}
{}{}{}
adalah
\definitionsverweis {serat}{}{}
di atas

\mavergleichskette
{\vergleichskette
{ c
}
{ \in }{ \R
}
{  }{
}
{    }{
}
{   }{
}
}
{}{}{,}
dengan $h$
\definitionsverweis {reguler}{}{}
di setiap titik pada $Y$. Buktikan bahwa
\definitionsverweis {grup holonomi}{}{}
di suatu titik

\mavergleichskette
{\vergleichskette
{ P
}
{ \in }{ Y
}
{  }{
}
{    }{
}
{   }{
}
}
{}{}{}
memang merupakan suatu
\definitionsverweis {subgrup}{}{}
dari
\definitionsverweis {grup isometri}{}{.}

\emph{Catatan edisi: Hipotesis $C^1$ pada sumber diperkuat menjadi $C^2$, sesuai dengan regularitas yang diperlukan oleh transpor paralel.}

}
{} {}

\inputaufgabe
{}
{

Misalkan

\mavergleichskette
{\vergleichskette
{ Y
}
{ \subseteq }{ \R^n
}
{  }{
}
{    }{
}
{   }{
}
}
{}{}{}
adalah suatu
\definitionsverweis {hiperbidang}{}{.}
Tentukan
\definitionsverweis {grup holonomi}{}{}
di

\mavergleichskette
{\vergleichskette
{ P
}
{ \in }{ Y
}
{  }{
}
{    }{
}
{   }{
}
}
{}{}{.}

}
{} {}

  \zwischenueberschrift{Soal untuk dikumpulkan}

\inputaufgabe
{2}
{

Pada
\definitionsverweis {permukaan bola satuan}{}{}
$Y$ yang berdimensi dua, kita tinjau
\definitionsverweis {medan vektor tangensial}{}{}

\mavergleichskettedisp
{\vergleichskette
{ F \begin{pmatrix}  x  \\ y \\ z   \end{pmatrix}
}
{ =} { \begin{pmatrix}  y  \\ -x\\  0   \end{pmatrix}
}
{ } {
}
{ } {
}
{ } {
}
}
{}{}{.}
Untuk titik

\mavergleichskette
{\vergleichskette
{ R
}
{ = }{ \begin{pmatrix}  { \frac{   1  }{  2 } }  \\ { \frac{   1  }{  2 } } \\  { \frac{   1  }{  \sqrt{2}  } }   \end{pmatrix}
}
{ \in }{ Y
}
{    }{
}
{   }{
}
}
{}{}{}
tentukan suatu matriks bagi pemetaan
\maabbeledisp {} {T_RY} { T_RY
} { v } { \nabla_v F
} {,}
terhadap suatu
\definitionsverweis {basis}{}{}
yang sesuai.

}
{} {}

\inputaufgabe
{4}
{

Misalkan

\mavergleichskettedisp
{\vergleichskette
{ h(x,y,z)
}
{ =} { x^2+y^2-2z^7
}
{ } {
}
{ } {
}
{ } {
}
}
{}{}{}
dan

\mavergleichskette
{\vergleichskette
{ Y
}
{ = }{ h^{-1}(3)
}
{  }{
}
{    }{
}
{   }{
}
}
{}{}{.}
Misalkan
\maabbeledisp {\gamma} { [1,5] } { Y
} { t } { \begin{pmatrix}  \sqrt{3  -4t^2+2t^7  } \\ 2t\\ t   \end{pmatrix}
} {,}
adalah suatu lintasan terdiferensialkan pada $Y$. Tentukan persamaan diferensial biasa dalam bentuk
Catatan 6.10
untuk
\definitionsverweis {medan-medan vektor paralel}{}{}
sepanjang $\gamma$.

}
{} {}

\inputaufgabe
{4}
{

Buktikan bahwa
\definitionsverweis {grup holonomi}{}{}
\definitionsverweis {permukaan bola satuan}{}{}
di setiap titik $P$ adalah grup isometri linear yang mempertahankan orientasi,
\mathl{\operatorname{SO}(T_PY) \cong \operatorname{SO}(2)}{.}

\emph{Catatan edisi: Sumber menyatakan grup isometri penuh $\operatorname{O}(2)$; transpor paralel mempertahankan orientasi, sehingga grup holonominya adalah $\operatorname{SO}(2)$.}

}
{} {}

\inputaufgabe
{4}
{

Misalkan

\mavergleichskette
{\vergleichskette
{ W
}
{ \subseteq }{ \R^n
}
{  }{
}
{    }{
}
{   }{
}
}
{}{}{}
\definitionsverweis {terbuka}{}{,}
\maabb {h} { W } { \R
} {}
adalah suatu
\definitionsverweis {fungsi yang dua kali terdiferensialkan secara kontinu}{}{}
dan

\mavergleichskette
{\vergleichskette
{ Y
}
{ = }{ h^{-1}(c)
}
{  }{
}
{    }{
}
{   }{
}
}
{}{}{}
adalah
\definitionsverweis {serat}{}{}
di atas

\mavergleichskette
{\vergleichskette
{ c
}
{ \in }{ \R
}
{  }{
}
{    }{
}
{   }{
}
}
{}{}{,}
dengan $h$
\definitionsverweis {reguler}{}{}
di setiap titik pada $Y$. Misalkan
\maabbdisp {L} {\R^n } { \R^n
} {}
adalah suatu
\definitionsverweis {isometri linear}{}{,}

\mavergleichskette
{\vergleichskette
{ W'
}
{ = }{ L^{-1}(W)
}
{  }{
}
{    }{
}
{   }{
}
}
{}{}{}
dan

\mavergleichskettedisp
{\vergleichskette
{ Z
}
{ =} { L^{-1}(Y)
}
{ =} { (h \circ L)^{-1}(c)
}
{ } {
}
{ } {
}
}
{}{}{.}
Jelaskan bagaimana konsep-konsep
\definitionsverweis {medan vektor paralel}{}{,}
\definitionsverweis {transpor paralel}{}{}
dan
\definitionsverweis {grup holonomi}{}{}
berperilaku di bawah $L$.

\emph{Catatan edisi: Hipotesis $C^1$ pada sumber diperkuat menjadi $C^2$, sesuai dengan regularitas yang diperlukan oleh transpor paralel dan holonomi.}

}
{} {}

\section*{Solusi yang disediakan oleh sumber}
\addcontentsline{toc}{section}{Solusi yang disediakan oleh sumber}
\subsection*{Solusi Soal 6.2}
\aufzaehlungdrei{Matriks Jacobi tersebut adalah

\mathdisp {\begin{pmatrix}  0   &   1  &  0  \\  -1 &  0  &  0 \\ 0  &  0 &  0  \end{pmatrix}} { . }

}{Sebagai medan normal, kita ambil $\begin{pmatrix}  x  \\ y \\ z   \end{pmatrix}$. Di

\mavergleichskette
{\vergleichskette
{ P
}
{ = }{ \begin{pmatrix}  0  \\ 1 \\  0   \end{pmatrix}
}
{  }{
}
{    }{
}
{   }{
}
}
{}{}{,}
vektor-vektor
\mathkor {} {\begin{pmatrix}  1  \\ 0 \\  0   \end{pmatrix}} {dan} {\begin{pmatrix}  0  \\ 0\\  1   \end{pmatrix}} {}
membentuk suatu basis ruang singgung. Kita mempunyai

\mavergleichskettedisp
{\vergleichskette
{ \begin{pmatrix}  0   &   1  &  0  \\  -1 &  0  &  0 \\ 0  &  0 &  0  \end{pmatrix} \begin{pmatrix}  1  \\ 0 \\  0   \end{pmatrix}
}
{ =} { \begin{pmatrix}  0  \\ -1\\  0   \end{pmatrix}
}
{ } {
}
{ } {
}
{ } {
}
}
{}{}{}
dan

\mavergleichskettedisp
{\vergleichskette
{  \begin{pmatrix}  0   &   1  &  0  \\  -1 &  0  &  0 \\ 0  &  0 &  0  \end{pmatrix} \begin{pmatrix}  0  \\ 0\\  1   \end{pmatrix}
}
{ =} { \begin{pmatrix}  0  \\ 0\\  0   \end{pmatrix}
}
{ } {
}
{ } {
}
{ } {
}
}
{}{}{.}
Di sini,
\mathl{\begin{pmatrix}  0  \\ -1\\  0   \end{pmatrix}}{} merupakan kelipatan vektor normal, sehingga proyeksi ortogonalnya ke ruang singgung juga merupakan vektor nol. Dengan demikian, pemetaan
\mathl{v \mapsto \nabla_vF}{} adalah pemetaan nol.
}{Di

\mavergleichskette
{\vergleichskette
{ Q
}
{ = }{ \begin{pmatrix}  0  \\ 0\\  1   \end{pmatrix}
}
{  }{
}
{    }{
}
{   }{
}
}
{}{}{,}
vektor-vektor
\mathkor {} {\begin{pmatrix}  1  \\ 0 \\  0   \end{pmatrix}} {dan} {\begin{pmatrix}  0  \\ 1 \\  0   \end{pmatrix}} {}
membentuk suatu basis ruang singgung. Kita mempunyai

\mavergleichskettedisp
{\vergleichskette
{ \begin{pmatrix}  0   &   1  &  0  \\  -1 &  0  &  0 \\ 0  &  0 &  0  \end{pmatrix} \begin{pmatrix}  1  \\ 0 \\  0   \end{pmatrix}
}
{ =} { \begin{pmatrix}  0  \\ -1\\  0   \end{pmatrix}
}
{ } {
}
{ } {
}
{ } {
}
}
{}{}{}
dan

\mavergleichskettedisp
{\vergleichskette
{ \begin{pmatrix}  0   &   1  &  0  \\  -1 &  0  &  0 \\ 0  &  0 &  0  \end{pmatrix} \begin{pmatrix}  0  \\ 1 \\  0   \end{pmatrix}
}
{ =} { \begin{pmatrix}  1  \\ 0 \\  0   \end{pmatrix}
}
{ } {
}
{ } {
}
{ } {
}
}
{}{}{.}
Kedua vektor citra ini sudah tangensial pada $Y$ di $Q$. Suatu matriks yang mendeskripsikan
\mathl{v \mapsto \nabla_vF}{} adalah

\mathdisp {\begin{pmatrix}  0   &   1  \\  -1  &  0 \end{pmatrix}} { . }

\emph{Catatan edisi: Tanda pada kedua kolom diperbaiki agar matriks ini sesuai dengan citra kedua vektor basis yang dihitung tepat di atasnya.}

}
\subsection*{Solusi Soal 6.6}
Kita mempunyai

\mavergleichskettealign
{\vergleichskettealign
{ (N (\gamma(t)) \times F(t))'
}
{ =} { (N (\gamma(t)))' \times F(t) +  N (\gamma(t)) \times F'(t)
}
{ =} { \left(  D N   \right)_{ \gamma(t)} \left(   \gamma'(t)   \right)  \times F(t) +  N (\gamma(t)) \times F'(t)
}
{ } {
}
{ } {
}
}
{}
{}{.}
Kita perlu membuktikan bahwa vektor ini selalu tegak lurus terhadap ruang singgung $T_{\gamma(t)}Y$. Karena $F'(t)$ tegak lurus terhadap ruang singgung, turunan tersebut merupakan kelipatan skalar vektor normal $N(\gamma(t))$; oleh karena itu, suku kedua sama dengan $0$ menurut
Fakta (3).
Pada suku di sebelah kiri, $F(t)$ bersifat tangensial menurut asumsi, sedangkan $\left(  D N   \right)_{ \gamma(t)} \left(   \gamma'(t)   \right)$ bersifat tangensial menurut
Fakta.
Oleh karena itu, menurut
Fakta (6),
suku di sebelah kiri tegak lurus terhadap ruang singgung. Jadi, medan vektor tersebut secara keseluruhan bersifat paralel.
\subsection*{Solusi Soal 6.9}
Menurut
Fakta,
medan kecepatan $\gamma'$ adalah suatu
\definitionsverweis {medan vektor paralel}{}{}
sepanjang $\gamma$. Oleh karena itu,
\definitionsverweis {transpor paralel}{}{}
$\Psi_\gamma$ sepanjang $\gamma$ membawa vektor

\mavergleichskette
{\vergleichskette
{ \gamma'(a)
}
{ \in }{  T_{\gamma(a)} Y
}
{  }{
}
{    }{
}
{   }{
}
}
{}{}{}
ke vektor

\mavergleichskette
{\vergleichskette
{ \gamma'(b)
}
{ \in }{  T_{\gamma(b)} Y
}
{  }{
}
{    }{
}
{   }{
}
}
{}{}{.}

\part{Unit 7}
\chapter[Kuliah 7]{Kuliah 7: Konsep sebuah manifold}
\setcounter{section}{7}

  \zwischenueberschrift{Konsep sebuah manifold}

Dalam geometri diferensial, gagasan tentang manifold berada di pusat pembahasan. Sebagai contoh, kita meninjau Bumi
\zusatzklammer {permukaannya} {} {,}
yang dalam sejarah ilmu pengetahuan lama dianggap sebagai sebuah cakram, dan ada alasan yang baik untuk itu. Secara lokal Bumi tampak seperti sebuah bidang. Hal ini juga tercermin dalam peta yang dibuat tentangnya. Sebuah peta adalah sebuah \anfuehrung{lembaran}{,} yang titik-titiknya berada dalam bijeksi dengan suatu bagian permukaan Bumi. Khususnya untuk bagian-bagian kecil hal ini tidak bermasalah, tetapi pada peta yang harus menggambarkan bagian besar atau bahkan seluruh Bumi segera muncul pertanyaan tentang apa yang dipertahankan oleh peta dan apa yang tidak: pertanyaan tentang pelestarian panjang, pelestarian luas, pelestarian sudut, titik-titik yang hilang atau muncul berkali-kali, pertanyaan perpanjangan, pertanyaan kelengkungan, dan seterusnya.

\bild{ \begin{center}
\includegraphics[width=5.5cm]{ \bildeinlesungpng {Stereographic_projection_in_3D} {png} }
\end{center}
\bildtext {Proyeksi stereografis ketika bidang tidak diletakkan melalui khatulistiwa, melainkan melalui Kutub Selatan.} }

\bildlizenz { Stereographic projection in 3D.png } {} {Mark.Howison} {en.Wikipedia} {PD} {}

Mula-mula kita membahas \stichwort {proyeksi stereografis} {} dari permukaan bola.

\inputbeispiel{
}
{

Kita meninjau permukaan bola

\mavergleichskettedisp
{\vergleichskette
{ K
}
{ =} { \left\{ (x,y,z)  \mid   x^2+y^2+z^2 = 1        \right\}
}
{ } {
}
{ } {
}
{ } {
}
}
{}{}{}
dan menyebut titik

\mavergleichskette
{\vergleichskette
{ N
}
{ = }{ (0,0,1)
}
{  }{
}
{    }{
}
{   }{
}
}
{}{}{}
sebagai Kutub Utara dan titik

\mavergleichskette
{\vergleichskette
{ S
}
{ = }{ (0,0,-1)
}
{  }{
}
{    }{
}
{   }{
}
}
{}{}{}
sebagai Kutub Selatan. Sebuah titik

\mathbed {P=(x,y,z) \in K} {}
 {P \neq N} {}
 {} {} {} {,}
bersama Kutub Utara menentukan sebuah garis tunggal di ruang, yang diparameterkan oleh

\mathdisp {(tx,ty ,1+t(z-1))=(0,0,1)  + t ((x,y,z) -(0,0,1)), \, t \in \R} { , }

Vektor \mathl{(x,y,z) -(0,0,1)}{,} yang menentukan garis ini, tidak sejajar dengan bidang $x-y$; artinya garis tersebut mempunyai tepat satu titik potong dengan bidang ini. Titik itu diperoleh untuk parameter

\mavergleichskettedisp
{\vergleichskette
{ t
}
{ =} { { \frac{   1  }{  1-z } }
}
{ } {
}
{ } {
}
{ } {
}
}
{}{}{,}
yaitu titik

\mathdisp {\left(  { \frac{   x  }{  1-z  } }, { \frac{   y  }{  1-z  } } ,0  \right)} { . }

Kita rangkum konstruksi ini sebagai pemetaan
\maabbeledisp {\alpha} {K \setminus \{N\}} { \R^2
} {(x,y,z)} { \left(  { \frac{   x  }{  1-z  } }, { \frac{   y  }{  1-z  } }  \right)
} {}
Pemetaan ini secara intuitif jelas merupakan bijeksi, dan hal tersebut juga mudah dihitung dari rumus-rumusnya. Pemetaan invers diperoleh dengan menghubungkan sebuah titik \mathl{(u,v,0)}{} di bidang dengan Kutub Utara, lalu menghitung titik tembusnya $\neq N$ dengan permukaan bola. Hal ini menghasilkan syarat

\mavergleichskettealignhandlinks
{\vergleichskettealignhandlinks
{ \Vert {(0,0,1) +a((u,v,0) - (0,0,1) ) } \Vert^2
}
{ =} { \Vert {(au,av,1-a) } \Vert^2
}
{ =} { a^2u^2+a^2v^2+(1-a)^2
}
{ =} { 1
}
{ } {
}
}
{}
{}{,}
yang berujung pada

\mavergleichskettedisp
{\vergleichskette
{ a \left(  au^2+av^2 -2 +a  \right)
}
{ =} { 0
}
{ } {
}
{ } {
}
{ } {
}
}
{}{}{}
Solusinya adalah

\mavergleichskette
{\vergleichskette
{ a
}
{ = }{ 0
}
{  }{
}
{    }{
}
{   }{
}
}
{}{}{}
yang bersesuaian dengan Kutub Utara dan tidak kita minati, sehingga kita memperoleh

\mavergleichskette
{\vergleichskette
{ a
}
{ = }{ { \frac{   2  }{ u^2+v^2+1 } }
}
{  }{
}
{    }{
}
{   }{
}
}
{}{}{}
dan dengan demikian pemetaan
\maabbeledisp {} { \R^2 } { K \setminus \{N\}
} { (u,v) } { \left(  { \frac{   2u }{ u^2+v^2+1 } }, { \frac{   2v }{ u^2+v^2+1 } } ,1- { \frac{   2  }{ u^2+v^2+1 } }  \right)
} {.}
Jadi bidang riil adalah \definitionsverweis {homeomorf}{}{} dengan permukaan bola yang \anfuehrung{dilubangi}{} pada satu titik.

Pertimbangan yang sama dapat dilakukan untuk \mathl{K \setminus \{S\}}{}. Hal ini menghasilkan pemetaan
\maabbeledisp {\beta} { K \setminus \{S\} } { \R^2
} { (x,y,z) } { \left(  { \frac{   x  }{ z+1 } }, { \frac{   y  }{ z+1 } }  \right)
} {}
dengan pemetaan invers
\maabbeledisp {} { \R^2 } { K \setminus \{S\}
} { (s,t) } { \left(  { \frac{   2s }{ s^2+t^2+1 } }, { \frac{   2t }{ s^2+t^2+1 } }, -1+ { \frac{   2  }{ s^2+t^2+1 } }  \right)
} {.}
Kutub Selatan dipetakan ke pusat bidang Euklides oleh pemetaan pertama, sedangkan Kutub Utara dipetakan ke pusat bidang Euklides oleh pemetaan kedua. Kita menyebut kedua pemetaan, maupun pemetaan inversnya, sebagai peta. Kedua peta tersebut bersama-sama menutupi seluruh permukaan bola. Karena merupakan homeomorfisme, keduanya mempertahankan sifat-sifat topologis utama bola. Keduanya tidak cocok untuk geografi Bumi karena peta membutuhkan seluruh bidang dan sangat mendistorsi panjang.

Kedua peta itu sama baiknya. Titik-titik \zusatzklammer {dan secara umum gambar-gambar lain} {} {} pada satu peta dapat dengan mudah diubah ke peta yang lain. Namun, perlu diperhatikan bahwa kedua kutub hanya muncul pada salah satu peta. Titik-titik dari himpunan \mathl{K \setminus \{N,S\}}{} terdapat pada kedua peta, dan melalui keduanya berada dalam bijeksi dengan bidang yang dilubangi di pusatnya, \mathl{\R^2 \setminus \{0\}}{.} Pemetaan \stichwort {transisi} {} \zusatzklammer {atau \stichwort {pergantian peta} {}} {} {} diberikan oleh

\mavergleichskette
{\vergleichskette
{ \psi
}
{ = }{ \beta \circ \left(  \alpha^{-1}\!\mid_{\R^2 \setminus \{0\} }  \right)
}
{  }{
}
{    }{
}
{   }{
}
}
{}{}{}
Selanjutnya,

\mavergleichskettealignhandlinks
{\vergleichskettealignhandlinks
{ \psi(u,v)
}
{ =} { \beta \left(  { \frac{   2u }{ u^2+v^2+1 } }, { \frac{   2v }{ u^2+v^2+1 } } ,1- { \frac{   2  }{ u^2+v^2+1 } }  \right)
}
{ =} { \left(  { \frac{   { \frac{   2u }{ u^2+v^2+1 } }  }{  2 - { \frac{   2  }{ u^2+v^2+1 } }  } } , { \frac{   { \frac{   2v }{ u^2+v^2+1 } }  }{  2 - { \frac{   2  }{ u^2+v^2+1 } }  } }  \right)
}
{ =} { \left(  { \frac{   { \frac{   u  }{ u^2+v^2+1 } }  }{  { \frac{  (u^2+v^2+1) -1  }{ u^2+v^2+1 } }  } } , { \frac{   { \frac{   v  }{ u^2+v^2+1 } }  }{  { \frac{  (u^2+v^2+1) -1  }{ u^2+v^2+1 } }  } }  \right)
}
{ =} { \left(  { \frac{   u  }{ u^2+v^2  } } , { \frac{   v  }{ u^2+v^2 } }  \right)
}
}
{}
{}{.}
Pemetaan transisi ini tidak hanya menginduksi homeomorfisme antara \mathl{\R^2 \setminus \{0\}}{} dengan \mathl{\R^2 \setminus \{0\}}{\zusatzfussnote {Sebaiknya di sini tidak dikatakan \anfuehrung{dengan dirinya sendiri}{} karena kedua bidang peta sebaiknya dibayangkan sebagai bidang yang independen. Hubungan di antara keduanya muncul semata-mata karena keduanya menggambarkan permukaan bola yang sama} {.} {,}} sebagaimana langsung mengikuti dari fakta bahwa kedua pemetaan peta \mathkor {} {\alpha} {dan} {\beta} {} adalah homeomorfisme, melainkan bahkan sebuah \definitionsverweis {difeomorfisme}{}{.} Ini langsung terlihat dari rumusnya; tetapi tidak masuk akal mengatakan bahwa pemetaan peta adalah difeomorfisme, karena permukaan bola bukan himpunan terbuka di dalam $\R^3$. Yang masih kurang adalah sebuah \anfuehrung{struktur diferensiabel}{} pada permukaan ini agar kita dapat berbicara tentang difeomorfisme.

}
Sebuah \stichwort {manifold} {} adalah suatu bangun geometri yang secara \anfuehrung{lokal}{} tampak seperti ruang Euklides $\R^n$. Kita menganggap bangun ini sebagai sebuah
\definitionsverweis {ruang topologis}{}{}
dan memperjelas kata lokal dengan menyatakan bahwa terdapat tutupan oleh himpunan-himpunan terbuka yang homeomorfik dengan himpunan terbuka di $\R^n$. Walaupun selanjutnya kita bekerja dengan ruang topologis, kandungan intuitif pembahasan tidak berkurang jika ruang topologis dibayangkan sekadar sebagai ruang metrik.

  \zwischenueberschrift{Manifold diferensiabel}

\inputdefinition
{}
{

Sebuah
\definitionsverweis {ruang topologis}{}{}
\definitionsverweis {Hausdorff}{}{}
$M$ disebut \definitionswort {manifold topologis}{} berdimensi \definitionswort {dimensi}{} $n$, jika terdapat suatu
\definitionsverweis {tutupan terbuka}{}{}

\mavergleichskette
{\vergleichskette
{ M
}
{ = }{  \bigcup_{i \in I} U_i
}
{  }{
}
{    }{
}
{   }{
}
}
{}{}{}
sedemikian sehingga setiap $U_i$
\definitionsverweis {homeomorf}{}{}
dengan suatu
\definitionsverweis {himpunan bagian terbuka}{}{}
dari $\R^n$.

}

Untuk setiap titik

\mavergleichskette
{\vergleichskette
{P
}
{ \in }{ M
}
{  }{
}
{    }{
}
{   }{
}
}
{}{}{}
terdapat lingkungan terbuka

\mavergleichskette
{\vergleichskette
{P
}
{ \in }{ U
}
{ \subseteq }{ M
}
{    }{
}
{   }{
}
}
{}{}{,}
yang homeomorfik dengan suatu himpunan terbuka

\mavergleichskette
{\vergleichskette
{V
}
{ \subseteq }{ \R^n
}
{  }{
}
{    }{
}
{   }{
}
}
{}{}{}.
Misalkan
\maabbdisp {\varphi} {U} {V
} {}
sebuah homeomorfisme dan misalkan

\mavergleichskette
{\vergleichskette
{ Q
}
{ = }{ \varphi(P)
}
{  }{
}
{    }{
}
{   }{
}
}
{}{}{.}
Maka lingkungan bola terbuka

\mavergleichskette
{\vergleichskette
{Q
}
{ \in }{ U \left(  Q, \epsilon  \right)
}
{ \subseteq }{ V
}
{    }{
}
{   }{
}
}
{}{}{}
berpadanan dengan lingkungan terbuka

\mavergleichskette
{\vergleichskette
{ U'
}
{ = }{  \varphi^{-1}(U \left(  Q, \epsilon  \right))
}
{  }{
}
{    }{
}
{   }{
}
}
{}{}{}
dengan

\mavergleichskette
{\vergleichskette
{P
}
{ \in }{ U'
}
{ \subseteq }{ U
}
{    }{
}
{   }{
}
}
{}{}{,}
yang menurut konstruksi homeomorfik dengan bola terbuka. Karena itu manifold topologis juga dapat dicirikan sebagai ruang Hausdorff topologis yang \stichwort {secara lokal Euklides} {}.

\inputdefinition
{}
{

Misalkan $M$ adalah sebuah
\definitionsverweis {manifold topologis}{}{.}
Setiap
\definitionsverweis {homeomorfisme}{}{}
\maabbdisp {\varphi} { U } { V
} {,}
dengan

\mavergleichskette
{\vergleichskette
{ U
}
{ \subseteq }{ M
}
{  }{
}
{    }{
}
{   }{
}
}
{}{}{}
dan

\mavergleichskette
{\vergleichskette
{ V
}
{ \subseteq }{ \R^n
}
{  }{
}
{    }{
}
{   }{
}
}
{}{}{}
\definitionsverweis {terbuka}{}{}
adalah sebuah
\zusatzklammer {topologis} {} {}
\definitionswort {peta}{} untuk $M$.

}

Himpunan terbuka

\mavergleichskette
{\vergleichskette
{ U
}
{ \subseteq }{ M
}
{  }{
}
{    }{
}
{   }{
}
}
{}{}{}
kadang-kadang disebut \stichwort {daerah peta} {}, sedangkan

\mavergleichskette
{\vergleichskette
{ V
}
{ \subseteq }{ \R^n
}
{  }{
}
{    }{
}
{   }{
}
}
{}{}{}
kadang-kadang disebut \stichwort {citra peta} {.} Untuk sebuah peta
\maabbdisp {\varphi} {U} {V
} {}
dan himpunan bagian terbuka

\mavergleichskette
{\vergleichskette
{ U'
}
{ \subseteq }{ U
}
{  }{
}
{    }{
}
{   }{
}
}
{}{}{}
pemetaan terinduksi
\maabbdisp {\varphi \vert_{U'}} {U'} {\varphi(U')
} {}
juga merupakan peta. Kadang-kadang pemetaan invers juga disebut peta. Selain peta, kita juga berbicara tentang \stichwort {sistem koordinat lokal} {.} Melalui peta
\maabb {\varphi} {U} {V
} {}
koordinat pada

\mavergleichskette
{\vergleichskette
{V
}
{ \subseteq }{ \R^n
}
{  }{
}
{    }{
}
{   }{
}
}
{}{}{}
dipindahkan ke $U$. Koordinat ke-$j$
\zusatzklammer {yakni proyeksi ke-$j$} {} {}
\maabb {x_j} {V} {\R
} {}
menginduksi \zusatzklammer {koordinat lokal} {} {-}fungsi
\maabbdisp {x_j \circ \varphi} {U} {\R
} {}
\zusatzklammer {yang sering kembali dilambangkan dengan $x_j$} {} {,}
dan sebuah titik

\mavergleichskette
{\vergleichskette
{ Q
}
{ = }{ \left(  x_1   , \,   , \ldots ,  , \,   x_n                 \right)
}
{ \in }{ V
}
{    }{
}
{   }{
}
}
{}{}{}
berpadanan dengan titik

\mavergleichskette
{\vergleichskette
{P
}
{ = }{\varphi^{-1}(Q)
}
{  }{
}
{    }{
}
{   }{
}
}
{}{}{.}

\bild{ \begin{center}
\includegraphics[width=5.5cm]{ \bildeinlesungpng {Manifold_zahyou3} {png} }
\end{center}
\bildtext {} }

% TeX-safe display form for the moving list-of-figures argument.  The exact
% Unicode creator string remains in source/unit_media.json and the rights
% manifest; this ASCII rendering avoids an engine-dependent Unicode failure.
\bildlizenz { Manifold zahyou3.png } {} {132ninme} {ja. Wikipedia} {CC-by-sa 3.0} {}

\inputdefinition
{}
{

Misalkan $M$ adalah sebuah
\definitionsverweis {manifold topologis}{}{,}
misalkan

\mavergleichskette
{\vergleichskette
{  U_1, U_2
}
{ \subseteq }{ M
}
{  }{
}
{    }{
}
{   }{
}
}
{}{}{}
merupakan
\definitionsverweis {himpunan bagian terbuka}{}{}
dan
\maabb {\alpha_1} { U_1 } { V_1
} {} dan \maabb {\alpha_2} { U_2 } { V_2
} {}
merupakan
\definitionsverweis {peta}{}{}
\zusatzklammer {dengan

\mavergleichskettek
{\vergleichskettek
{  V_1,V_2
}
{ \subseteq }{\R^n
}
{  }{
}
{    }{
}
{   }{
}
}
{}{}{}
terbuka} {} {.}
Maka
\definitionsverweis {pemetaan}{}{}
\maabbdisp {\alpha_2  \circ \alpha_1^{-1}} { \alpha_1(U_1 \cap U_2) } { \alpha_2(U_1 \cap U_2)
} {} disebut \definitionswort {pemetaan transisi}{} untuk peta-peta tersebut.

}

Irisan \mathl{U_1 \cap U_2}{} adalah himpunan bagian terbuka tempat kedua peta didefinisikan dan tempat keduanya dapat dibandingkan. Secara lebih teliti, pada definisi tersebut kita harus berbicara tentang pembatasan $\alpha_1^{-1}$ pada himpunan bagian terbuka

\mavergleichskette
{\vergleichskette
{ \alpha_1(U_1 \cap U_2)
}
{ \subseteq }{ V_1
}
{  }{
}
{    }{
}
{   }{
}
}
{}{}{}.

\inputdefinition
{}
{

Misalkan

\mavergleichskette
{\vergleichskette
{ n
}
{ \in }{ \N
}
{  }{
}
{    }{
}
{   }{
}
}
{}{}{}
dan

\mavergleichskette
{\vergleichskette
{ k
}
{ \in }{ \overline{ \N }_+
}
{  }{
}
{    }{
}
{   }{
}
}
{}{}{.}
Sebuah
\definitionsverweis {ruang topologis}{}{}
\definitionsverweis {Hausdorff}{}{}
$M$ bersama suatu
\definitionsverweis {tutupan terbuka}{}{}

\mavergleichskette
{\vergleichskette
{ M
}
{ = }{ \bigcup_{i \in I} U_i
}
{  }{
}
{    }{
}
{   }{
}
}
{}{}{}
dan
\definitionsverweis {peta}{}{}
\maabbdisp {\alpha_i} {U_i } { V_i
} {}
dengan

\mavergleichskette
{\vergleichskette
{ V_i
}
{ \subseteq }{ \R^n
}
{  }{
}
{    }{
}
{   }{
}
}
{}{}{}
terbuka, sedemikian sehingga
\definitionsverweis {pemetaan transisi}{}{}
\maabbdisp {\alpha_j \circ (\alpha_i)^{-1}} {V_i \cap \alpha_i(U_i \cap U_j) } { V_j \cap \alpha_j(U_i \cap U_j)
} {}
merupakan
$C^k$-\definitionsverweis {difeomorfisme}{}{}
untuk semua

\mavergleichskette
{\vergleichskette
{ i,j
}
{ \in }{ I
}
{  }{
}
{    }{
}
{   }{
}
}
{}{}{}
dinamakan
\definitionswortpraemath {C^k}{ manifold }{}
atau \definitionswort {manifold diferensiabel}{}
\zusatzklammer {berdimensi \definitionswort {dimensi}{} $n$ dan berderajat diferensiabilitas $k$} {} {.} Kumpulan peta

\mathbed {(U_i,\alpha_i)} {}
 {i \in I} {}
 {} {} {} {,}
dinamakan juga
\definitionswortpraemath {C^k}{ atlas }{}
manifold tersebut.

}

Manifold diferensiabel khususnya merupakan sebuah
\definitionsverweis {manifold topologis}{}{.}
Menurut definisi kita, atlas merupakan bagian integral dari konsep manifold. Namun, yang lebih penting daripada atlas adalah \stichwort {struktur diferensiabel} {} yang ditentukan oleh atlas tersebut. Hal ini akan menjadi lebih jelas setelah kita memiliki konsep pemetaan diferensiabel antara dua manifold dan dapat berbicara tentang manifold yang difeomorf.

\inputdefinition
{}
{

Misalkan $M$ sebuah
\definitionsverweis {manifold diferensiabel}{}{.}
Sebuah
\definitionsverweis {himpunan bagian terbuka}{}{}

\mavergleichskette
{\vergleichskette
{ U
}
{ \subseteq }{ M
}
{  }{
}
{    }{
}
{   }{
}
}
{}{}{,}
yang dilengkapi dengan
\definitionsverweis {peta}{}{}
\definitionsverweis {terbatas}{}{}
disebut \definitionswort {submanifold terbuka}{.}

}

\inputbeispiel{}
{

Setiap himpunan bagian terbuka

\mavergleichskette
{\vergleichskette
{ V
}
{ \subseteq }{ \R^n
}
{  }{
}
{    }{
}
{   }{
}
}
{}{}{}
merupakan
$C^\infty$-\definitionsverweis {manifold}{}{,}
jika kita menggunakan
\definitionsverweis {identitas}{}{}
\maabb {\operatorname{Id}} { V } { V
} {}
sebagai
\definitionsverweis {peta}{}{}.
Satu-satunya
\definitionsverweis {pemetaan transisi}{}{}
adalah identitas itu sendiri, yang merupakan sebuah $C^\infty$-\definitionsverweis {difeomorfisme}{}{}.
Ini adalah manifold dengan sebuah
\definitionsverweis {atlas}{}{,}
yang terdiri atas satu peta. Namun kita juga dapat menggunakan atlas yang terdiri atas semua himpunan bagian terbuka

\mavergleichskette
{\vergleichskette
{ U
}
{ \subseteq }{ V
}
{  }{
}
{    }{
}
{   }{
}
}
{}{}{}
beserta peta identitas terkait $\varphi_U$. Pemetaan transisinya adalah identitas pada \mathl{U_1 \cap U_2}{.}

}

Kita telah membahas ruang metrik yang
\definitionsverweis {terhubung}{}{}
dalam konteks teorema nilai antara. Definisi yang sama juga kita gunakan untuk ruang topologis.

\inputdefinition
{}
{

Sebuah
\definitionsverweis {ruang topologis}{}{}
$X$ disebut \definitionswort {terhubung}{,} jika di dalam $X$ tepat ada dua himpunan
\zusatzklammer {yaitu $\emptyset$ dan seluruh ruang \mathlk{X \neq \emptyset}{}} {} {,}
yang sekaligus
\definitionsverweis {terbuka}{}{}
dan
\definitionsverweis {tertutup}{}{}.

}

Sering kali kita hanya tertarik pada manifold yang terhubung, terutama karena pada kasus tak terhubung komponen-komponen \anfuehrung{keterhubungan}{} dapat diteliti secara terpisah. Sekarang kita membahas secara singkat manifold berdimensi rendah.

\inputbeispiel{}
{

Pada sebuah
\definitionsverweis {manifold}{}{}
$M$ berdimensi nol, untuk setiap titik

\mavergleichskette
{\vergleichskette
{ P
}
{ \in }{ M
}
{  }{
}
{    }{
}
{   }{
}
}
{}{}{}
terdapat lingkungan terbuka

\mavergleichskette
{\vergleichskette
{ P
}
{ \in }{ U
}
{  }{
}
{    }{
}
{   }{
}
}
{}{}{,}
yang
\definitionsverweis {homeomorf}{}{}
dengan suatu himpunan terbuka dari

\mavergleichskette
{\vergleichskette
{ \R^0
}
{ = }{ \{0\}
}
{  }{
}
{    }{
}
{   }{
}
}
{}{}{}.
Artinya, himpunan satu titik \mathl{\{P\}}{} harus terbuka, sehingga $M$ harus membawa
\definitionsverweis {topologi diskret}{}{},
yaitu setiap himpunan bagiannya terbuka. Karena itu satu-satunya manifold berdimensi nol yang
\definitionsverweis {terhubung}{}{}
adalah himpunan satu titik.

}

\bild{ \begin{center}
\includegraphics[width=5.5cm]{ \bildeinlesungsvg {Circle_-_black_simple} {svg} }
\end{center}
\bildtext {Sebuah garis lingkaran adalah manifold kompak berdimensi satu} }

\bildlizenz { Circle - black simple.svg } {} {Dakdada} {Commons} {PD} {}

\inputbeispiel{}
{

Pada manifold berdimensi satu
\definitionsverweis {manifold}{}{}
mula-mula terdapat himpunan bagian terbuka dari $\R^1$. Himpunan-himpunan ini merupakan gabungan interval terbuka, dan \definitionsverweis {terhubung}{}{,} tepat ketika merupakan sebuah interval terbuka. Setiap interval terbuka, baik terbatas maupun tak terbatas, adalah
\definitionsverweis {homeomorf}{}{}
dan juga
\definitionsverweis {difeomorf}{}{}
dengan
\definitionsverweis {interval satuan terbuka}{}{}
\mathl{]0,1[}{} dan dengan bilangan-bilangan riil $\R$ sendiri. Interval tertutup

\mathbed {[a,b]} {dengan}
 {a < b} {}
 {} {} {} {}
bukan manifold, karena untuk titik-titik batasnya
\zusatzklammer {ujung-ujung interval} {} {}
tidak ada lingkungan terbuka yang homeomorfik dengan interval terbuka
\zusatzklammer {meskipun interval itu disebut \stichwort {manifold dengan batas} {}} {} {.}

Selain itu ada \stichwort {lingkaran} {}
\zusatzklammer {atau \stichwort {sfera} {}} {} {}
$S^1$ sebagai manifold terhubung berdimensi satu lainnya. Berlaku

\mavergleichskettedisp
{\vergleichskette
{ S^1
}
{ =} { \left\{ (x,y) \in \R^2  \mid   x^2+y^2 = 1         \right\}
}
{ } {
}
{ } {
}
{ } {
}
}
{}{}{.}
Untuk setiap titik

\mavergleichskette
{\vergleichskette
{ P
}
{ \in }{ S^1
}
{  }{
}
{    }{
}
{   }{
}
}
{}{}{}
berlaku
\mathl{S^1 \setminus \{P \}}{} homeomorfik dengan $\R$
\zusatzklammer {melalui proyeksi stereografis} {} {.}
Lingkaran tidak homeomorfik dengan $\R$, sebab lingkaran
\definitionsverweis {kompak}{}{\zusatzfussnote {Sebelumnya kita mendefinisikan kekompakan hanya untuk himpunan bagian di dalam $\R^n$; sebentar lagi akan terlihat bahwa ini merupakan konsep absolut yang tidak bergantung pada penyematan. Jadi, tidak mungkin menyematkan $\R$ secara homeomorfik ke dalam $\R^n$ sedemikian sehingga citranya kompak} {.} {}}
sedangkan bilangan riil tidak kompak. Selain \mathkor {} {S^1} {dan} {\R} {}, tidak ada manifold terhubung berdimensi satu lainnya yang memiliki
\definitionsverweis {basis topologi terhitung}{}{}
\zusatzklammer {hal ini dinyatakan tanpa bukti di sini} {} {.}

}
Mulai dari dimensi dua, tanpa syarat tambahan yang kuat, mustahil membuat gambaran menyeluruh tentang semua manifold.

\chapter{Lembar Kerja 7}
\setcounter{section}{7}

  \zwischenueberschrift{Soal latihan}

\inputaufgabe
{}
{

Tunjukkan bahwa untuk dua titik
\mathl{A,B \neq N,S}{} pada bola satuan $S^2$, selisih jarak yang diambil dalam dua peta standar stereografis, yaitu
\mathkor {} {d( \alpha_1(A),\alpha_1(B))} {dan} {d( \alpha_2(A),\alpha_2(B))} {,}
dapat dibuat sekecil dan sebesar apa pun.

}
{} {}
Soal di atas menunjukkan bahwa secara umum kita tidak dapat memperoleh pengertian jarak yang bermakna pada suatu manifold melalui peta. Metrik alami pada bola satuan diperoleh dari metrik terinduksi pada $\R^3$ atau dari jarak geodesik; lihat
Aufgabe 7.16.

\inputaufgabe
{}
{

Tunjukkan bahwa suatu
\definitionsverweis {interval setengah terbuka}{}{}
\mathl{[a,b[}{} bukan
\definitionsverweis {manifold topologis}{}{}.

}
{} {}

\inputaufgabe
{}
{

Misalkan
\mathl{I=[0,1[}{} adalah
\zusatzklammer {setengah terbuka ke atas} {} {}
interval satuan dan $S^1$ adalah
\definitionsverweis {lingkaran satuan}{}{.}
Tunjukkan bahwa terdapat suatu
\definitionsverweis {pemetaan bijektif}{}{}
\definitionsverweis {kontinu}{}{}
\maabbdisp {f} { I } { S^1
} {}
tetapi
\mathkor {} {I} {dan} {S^1} {}
tidak
\definitionsverweis {homeomorfik}{}{}
.

}
{} {}

\inputaufgabegibtloesung
{}
{

Tunjukkan bahwa ketiga manifold berdimensi satu
\mathlistdisp {S^1 = \left\{ (x,y) \in \R^2  \mid   \betrag { (x,y) } = 1         \right\}} {} {\R} {dan} {\R \setminus \{0\}} {}
tidak homeomorfik satu sama lain secara berpasangan.

}
{} {}

\inputaufgabe
{}
{

Misalkan $M$ adalah
\definitionsverweis {grafik}{}{}
dari fungsi

\mavergleichskettedisp
{\vergleichskette
{ f(x)
}
{ =} {\begin{cases} \sin \frac{1}{x} \text{ untuk } x \neq 0\, , \\ 0 \text{ selain itu }\, , \end{cases}
}
{ } {
}
{ } {
}
{ } {
}
}
{}{}{}
yang diberikan oleh pemetaan
\maabb {f} {\R} {\R
} {.}
Tunjukkan bahwa $M$ bukan
\definitionsverweis {manifold topologis}{}{}.

}
{} {}

\inputaufgabe
{}
{

Tunjukkan bahwa sebuah
\definitionsverweis {bola terbuka}{}{}

\mavergleichskette
{\vergleichskette
{ U \left(   P ,r  \right)
}
{ \subseteq }{ \R^m
}
{  }{
}
{    }{
}
{   }{
}
}
{}{}{}
adalah $C^\infty$-\definitionsverweis {difeomorfik}{}{}
dengan $\R^m$.

}
{} {}

\inputaufgabegibtloesung
{}
{

Misalkan

\mathdisp {S= \left\{ P \in \R^3  \mid   \Vert { P } \Vert = 1         \right\}} {  }

adalah bola satuan. Untuk
\mathl{v=(a,b,c) \neq 0}{} berlaku

\mathdisp {E_v = \left\{ (x,y,z) \in \R^3  \mid  ax+by+cz = 0         \right\}} {  }

sebuah bidang melalui titik nol, yang mendefinisikan sebuah lingkaran besar
\zusatzklammer {sebuah \anfuehrung{khatulistiwa}{}} {} {}
dan dua hemisfer terbuka pada $S$.

a) Jelaskan, untuk
\mathl{v=(a,b,c) \neq 0}{}, lingkaran besar terkait dan kedua hemisfer dengan persamaan atau pertidaksamaan.

b) Tunjukkan bahwa $S$ tidak dapat ditutupi oleh tiga hemisfer terbuka.

}
{} {}

\inputaufgabe
{}
{

Tunjukkan bahwa permukaan silinder terbuka
\mathl{S^1 \times {]0,1[}}{} homeomorfik dengan
\mathl{S^1 \times \R}{,} dengan bidang yang ditusuk
\mathl{\R^2 \setminus \{(0,0)\}}{} dan dengan
\mathl{S^2 \setminus \{N,S\}}{}
\definitionsverweis {homeomorfik}{}{.}

}
{} {}

\inputaufgabe
{}
{

Misalkan
\mathl{]a,b[}{} sebuah
\definitionsverweis {interval terbuka}{}{}
dan
\maabbdisp {f} {]a,b[} {\R_{\geq 0}
} {}
sebuah
\definitionsverweis {fungsi kontinu}{}{.}
Misalkan $M$ adalah permukaan luar benda hasil rotasi terkait. Tunjukkan bahwa ketika

\mavergleichskette
{\vergleichskette
{ f
}
{ > }{ 0
}
{  }{
}
{    }{
}
{   }{
}
}
{}{}{}
ia merupakan manifold yang homeomorfik dengan selubung silinder terbuka. Tunjukkan lebih lanjut bahwa tidak ada manifold jika $f$ memiliki baik titik nol maupun nilai positif.

}
{} {}

\inputaufgabe
{}
{

Tunjukkan bahwa sebuah
\definitionsverweis {manifold topologis}{}{}
\definitionsverweis {terhubung}{}{}
tepat jika dan hanya jika
\definitionsverweis {terhubung lintasan}{}{}
.

}
{} {}

\inputaufgabe
{}
{

Misalkan $M$ sebuah
\definitionsverweis {manifold topologis}{}{}
dan misalkan

\mathl{M= \bigcup_{i \in I} U_i}{} sebuah
\definitionsverweis {tutupan terbuka}{}{}
dengan
\definitionsverweis {peta}{}{}
\maabbdisp {\alpha_i} {U_i } { V_i
} {}
dengan
\mathl{V_i \subseteq \R^n}{.} Untuk
\mathl{i,j \in I}{} definisikan
\maabbdisp {\varphi_{ij} = \alpha_j \circ (\alpha_i)^{-1}} {V_i \cap \alpha_i(U_i \cap U_j)  } { V_j \cap \alpha_j(U_i \cap U_j)
} {}
yang disebut
\definitionsverweis {pemetaan transisi}{}{.}
Tunjukkan bahwa untuk
\mathl{i,j,k \in I}{} apa yang disebut \stichwort {syarat koklus} {}
\zusatzklammer {pada himpunan bagian mana} {?} {}

\mavergleichskettedisp
{\vergleichskette
{ \varphi_{ij}
}
{ =} {\varphi_{kj} \circ \varphi_{ik}
}
{ } {
}
{ } {
}
{ } {
}
}
{}{}{}
berlaku.

}
{} {}

\inputaufgabe
{}
{

Periksa syarat koklus
\zusatzklammer {lihat
Aufgabe 7.11} {} {}
untuk bola satuan $M=S^2$ dan tiga proyeksi stereografis dari Kutub Utara, Kutub Selatan, dan
\mathl{P=(1,0,0)}{.}

}
{} {}

\inputaufgabegibtloesung
{}
{

Buat sebuah animasi yang menampilkan objek-objek geometri dari
Aufgabe 7.17.

}
{} {}

\inputaufgabe
{}
{

Perhatikan kertas kado. Dengan cara apa kertas tersebut dapat dipotong dan/atau direkatkan sehingga menghasilkan manifold berdimensi dua? Apakah tepi kertas harus menjadi bagian darinya atau tidak? Manifold mana yang dihasilkan yang terhubung, mana yang kompak? Bagaimana pita Möbius dibuat? Kemungkinan apa yang ada jika kita mengizinkan sejumlah hingga titik pengecualian, tempat tidak ada struktur manifold?

Terapkan teori ini untuk merancang kemasan yang seorisinal mungkin. Ikatlah kemasan tersebut dengan manifold berdimensi satu yang sesuai.

}
{} {}

  \zwischenueberschrift{Soal untuk dikumpulkan}

\inputaufgabe
{3}
{

Tentukan
\definitionsverweis {citra}{}{} dari
\definitionsverweis {lingkaran-lingkaran besar}{}{}
melalui kedua kutub pada
\definitionsverweis {bola satuan}{}{} di bawah
\definitionsverweis {proyeksi stereografis}{}{}
dari Kutub Utara.

}
{} {}

\inputaufgabe
{5}
{

Tunjukkan bahwa pada
\definitionsverweis {bola satuan}{}{}

\mavergleichskette
{\vergleichskette
{ K
}
{ \subset }{ \R^3
}
{  }{
}
{    }{
}
{   }{
}
}
{}{}{}
melalui korespondensi berikut ditentukan sebuah
\definitionsverweis {metrik}{}{} . Untuk

\mavergleichskette
{\vergleichskette
{ P,Q
}
{ \in }{ K
}
{  }{
}
{    }{
}
{   }{
}
}
{}{}{}
\mathl{d(P,Q)}{} adalah panjang lintasan penghubung
\zusatzklammer {yang lebih pendek} {} {} dari $P$ ke $Q$ pada
\definitionsverweis {lingkaran besar}{}{}
\zusatzklammer {perhatikan juga kasus

\mavergleichskette
{\vergleichskette
{ P
}
{ = }{ Q
}
{  }{
}
{    }{
}
{   }{
}
}
{}{}{}
dan $P,Q$
\definitionsverweis {antipodal}{}{}} {} {.}

}
{} {}

\inputaufgabe
{8}
{

Kita tetapkan dua titik $N=(0,0,1)$ dan $P=(1,0,0)$ pada
\definitionsverweis {bola satuan}{}{} $K$. Misalkan $G$ adalah garis penghubung dan $H$ adalah bidang tegak lurus terhadap $G$ melalui $N$. Pada $H$, buat lingkaran satuan berparameter $E$ dengan $N$ sebagai pusat. Tentukan, untuk $S \in E$, panjang lintasan
\zusatzklammer {yang lebih pendek} {} {} dari $N$ ke $P$ pada lingkaran yang ditentukan oleh irisan $K$ dengan bidang yang melalui
\mathkor {} {N,P} {dan} {S} {},
tersebut.

}
{} {}

\inputaufgabe
{6}
{

Berikan suatu
\definitionsverweis {pemetaan injektif}{}{}
\definitionsverweis {kontinu}{}{}
\maabbdisp {\varphi} {\R} {S^2
} {,}
yang
\zusatzklammer {sebagai pemetaan ke $\R^3$} {} {}
\definitionsverweis {dapat direktifikasi}{}{}
memiliki
\definitionsverweis {panjang}{}{}
tak hingga, dan memenuhi
\mathkor {} {\operatorname{lim}_{   t  \rightarrow  \infty  } \, \varphi (t) = N} {dan} {\operatorname{lim}_{   t  \rightarrow  - \infty  } \, \varphi (t) = S} {}.

}
{} {}

\inputaufgabe
{4}
{

Tunjukkan bahwa sumbu-sumbu koordinat bukan
\definitionsverweis {manifold topologis}{}{}.

}
{} {}

\section*{Solusi yang disediakan oleh sumber}
\addcontentsline{toc}{section}{Solusi yang disediakan oleh sumber}
\subsection*{Solusi Soal 7.4}
Sebagai himpunan bagian tertutup dan terbatas dari $\R^2$, lingkaran satuan kompak
\zusatzklammer {Satz von Heine-Borel} {} {.} Garis riil $\R$ dan $\R \setminus \{0\}$, sebaliknya, tidak kompak karena keduanya merupakan himpunan bagian tak terbatas dari $\R$. Jadi
\mathl{S^1 \not \cong \R,\, \R \setminus \{0\}}{.}

Garis riil terhubung, sebagaimana mengikuti dari
Zwischenwertsatz.
Sebaliknya, $\R \setminus \{0\}$ tidak terhubung, karena kita dapat menuliskannya sebagai

\mavergleichskettedisp
{\vergleichskette
{\R_+
}
{ =} { (\R \setminus \{0\}) \cap [0, \infty]
}
{ =} { (\R \setminus \{0\}) \cap \, ]0, \infty]
}
{ } {
}
{ } {
}
}
{}{}{}
sehingga diperoleh suatu himpunan bagian tak kosong yang sekaligus terbuka dan tertutup. Jadi
\mathl{\R \,\not \cong \R \setminus \{0\}}{.}
\subsection*{Solusi Soal 7.7}
a) Lingkaran besar adalah

\mathdisp {G= \left\{ (x,y,z) \in S  \mid  ax+by+cz = 0        \right\}} {  }

dan kedua hemisfer terbuka adalah

\mathdisp {U_+ = \left\{ (x,y,z) \in S  \mid  ax+by+cz > 0        \right\}} {  }

dan

\mathdisp {U_- = \left\{ (x,y,z) \in S  \mid  ax+by+cz < 0        \right\}} { . }

b) Untuk setiap hemisfer terbuka $U$ dan lingkaran besar $G$ yang terkait berlaku
\mathl{U \cap G = \emptyset}{.} Jarak maksimum antara dua titik
\mathl{P,Q \in S}{} adalah
\mathl{2}{,} dan hal ini terjadi tepat ketika kedua titik tersebut saling berhadapan
\zusatzklammer {antipodal} {} {}, yaitu ketika garis penghubungnya melalui pusat bola
\zusatzklammer {yakni pada
\mathl{Q=-P}{}} {} {}. Pasangan antipodal seperti itu tidak terletak pada satu hemisfer terbuka, sebab untuk
\mathl{P=(x,y,z)}{} dan
\mathl{P \in U_+}{} berlaku
\mathl{ax+by+cz>0}{} dan karena itu
\mathl{a(-x) +b(-y)+c(-z) <0}{} berlaku, sehingga
\mathl{-P \in U_-}{.}

Sekarang anggaplah bahwa terdapat suatu tutupan terbuka

\mathdisp {S \subseteq U_1 \cup U_2 \cup U_3} {  }

dengan tiga hemisfer terbuka
\mathl{U_1,U_2,U_3}{}
\zusatzklammer {bidang dan lingkaran besar yang bersesuaian masing-masing dilambangkan dengan
\mathkor {} {E_i} {atau} {G_i} {}} {} {.}
Karena
\mathl{U_1 \cap G_1 = \emptyset}{} diperoleh

\mathdisp {G_1 \subseteq U_2 \cup U_3} { . }

Irisan
\mathl{G_1 \cap E_2}{} memuat setidaknya dua titik antipodal
\mathkor {} {P} {dan} {-P} {}
dan
\mathl{P,-P \not\in U_2}{.} Karena, berdasarkan pengamatan sebelumnya, $U_3$ tidak memuat pasangan titik antipodal, salah satu dari kedua titik ini juga tidak termasuk dalam $U_3$, sehingga kita memperoleh kontradiksi.
\subsection*{Solusi Soal 7.13}
Animasi pertama memperlihatkan secara bertahap konstruksi objek-objek individual. Animasi lainnya memperlihatkan variasi S dan perubahan panjang lintasan yang bersesuaian.

Dari Lukas Freudenberg.
\href{https://commons.wikimedia.org/wiki/File:Aufgabe75.22.1.gif}{Konstruktion der Objekte}
\href{https://commons.wikimedia.org/wiki/File:Aufgabe75.22.2.gif}{Variation von S}

\part{Unit 8}
\chapter[Kuliah 8]{Kuliah 8: Teorema Pemetaan Implisit dan Manifold}
\setcounter{section}{8}

  \zwischenueberschrift{Teorema pemetaan implisit dan manifold}

Permukaan bola satuan, yang dalam kuliah sebelumnya kita bahas sebagai contoh pemotivasi sebuah manifold, merupakan serat di atas $1$ dari pemetaan diferensiabel
\maabbeledisp {} { \R^3 } { \R
} { (x,y,z) } { x^2+y^2+z^2
} {,}
Pemetaan ini
\definitionsverweis {reguler}{}{}
kecuali di titik nol. Dalam situasi ini, teorema pemetaan implisit memberikan pernyataan yang luas tentang bentuk lokal serat suatu pemetaan
\maabb {\varphi} {\R^n} { \R^m
} {,}
yakni bahwa secara lokal terdapat homeomorfisme antara serat di suatu titik reguler dan sebuah himpunan terbuka di $\R^k$, dengan $k$ selisih antara dimensi ruang asal dan dimensi ruang sasaran. Sebentar lagi akan kita lihat bahwa serat-serat semacam itu bukan hanya manifold topologis, melainkan juga manifold diferensiabel. Kita merumuskan teorema pemetaan implisit dalam suatu versi yang memperlihatkan bahwa serat-serat reguler merupakan manifold diferensiabel.



\inputfaktbeweis
{Satz über implizite Abbildungen/Globale Diffeomorphismen/Induzierte Diffeomorphismen zwischen Faser und Achsenräumen/Fakt}
{Teorema}
{}
{

\faktsituation {Misalkan

\mavergleichskette
{\vergleichskette
{ G
}
{ \subseteq }{ \R^n
}
{  }{
}
{    }{
}
{   }{
}
}
{}{}{}
\definitionsverweis {terbuka}{}{}
dan misalkan
\maabbdisp {\varphi} { G } {\R^m
} {}
suatu
\definitionsverweis {pemetaan terdiferensialkan secara kontinu}{}{.}
Misalkan

\mavergleichskette
{\vergleichskette
{ Z
}
{ = }{ \varphi^{-1}(0)
}
{  }{
}
{    }{
}
{   }{
}
}
{}{}{}
merupakan
\definitionsverweis {serat}{}{}
di atas

\mavergleichskette
{\vergleichskette
{ 0
}
{ \in }{ \R^m
}
{  }{
}
{    }{
}
{   }{
}
}
{}{}{,}
dan misalkan $\varphi$
\definitionsverweis {reguler}{}{}
di setiap titik serat tersebut.}
\faktfolgerung {Maka untuk setiap titik

\mavergleichskette
{\vergleichskette
{ P
}
{ \in }{ Z
}
{  }{
}
{    }{
}
{   }{
}
}
{}{}{}
terdapat suatu lingkungan terbuka

\mavergleichskette
{\vergleichskette
{ P
}
{ \in }{ W
}
{ \subseteq }{ G
}
{    }{
}
{   }{
}
}
{}{}{,}
himpunan-himpunan terbuka

\mavergleichskette
{\vergleichskette
{ V
}
{ \subseteq }{ \R^{n- m}
}
{  }{
}
{    }{
}
{   }{
}
}
{}{}{}
dan

\mavergleichskette
{\vergleichskette
{ V'
}
{ \subseteq }{ \R^{ m }
}
{  }{
}
{    }{
}
{   }{
}
}
{}{}{,}
serta suatu
$C^1$-\definitionsverweis {difeomorfisme}{}{}
\maabbdisp {\theta} { W } { V \times V'
} {}
dengan

\mavergleichskette
{\vergleichskette
{ \varphi\!\mid_{W}
}
{ = }{ p_2 \circ \theta
}
{  }{
}
{    }{
}
{   }{
}
}
{}{}{,}
yang menginduksi bijeksi antara
\mathkor {} {Z \cap W} {dan} {V \times \{0\}} {}
sedemikian sehingga
\definitionsverweis {diferensial total}{}{}

\mathl{\left(D\theta\right)_{Q}}{} untuk setiap

\mavergleichskette
{\vergleichskette
{ Q
}
{ \in }{ W
}
{  }{
}
{    }{
}
{   }{
}
}
{}{}{}
menginduksi bijeksi antara
\mathkor {} {\ker  \left(D\varphi\right)_{ Q }} {dan} {\R^{n-m }} {}
.}
\faktzusatz {}
\faktzusatz {}

}
{

Pernyataan ini telah dibuktikan sekaligus dalam bukti
teorema pemetaan implisit.
Pernyataan tambahan diperoleh dari

\mavergleichskettedisp
{\vergleichskette
{ \ker  \left(D\varphi\right)_{Q}
}
{ \cong} { \ker  \left(Dp_2 \right)_{ \theta(Q)}
}
{ =} { \R^{n- m } \times \{0\} \cong \R^{n-m}
}
{ } {
}
{ } {
}
}
{}{}{.}

Catatan edisi: Sumber mengidentifikasi kernel proyeksi kedua langsung dengan ruang berdimensi lebih rendah. Bentuk di atas menampilkan dahulu subruang yang sebenarnya, lalu isomorfisme kanonisnya.

}

Dari hasil ini, serat itu sendiri memperoleh struktur manifold. Dengan demikian, teorema pemetaan implisit menyediakan bagi kita suatu kelas manifold yang sangat besar. Manifold-manifold ini disebut \stichwort {submanifold tertutup} {,} yang segera akan kita bahas secara lebih sistematis setelah ruang singgung tersedia.



\inputfaktbeweis
{Satz über implizite Abbildungen/Faser ist Mannigfaltigkeit/Fakt}
{Teorema}
{}
{

\faktsituation {Misalkan

\mavergleichskette
{\vergleichskette
{ G
}
{ \subseteq }{ \R^n
}
{  }{
}
{    }{
}
{   }{
}
}
{}{}{}
\definitionsverweis {terbuka}{}{}
dan misalkan
\maabbdisp {\varphi} { G } { \R^{ m }
} {}
suatu
\definitionsverweis {pemetaan terdiferensialkan secara kontinu}{}{.}
Misalkan

\mavergleichskette
{\vergleichskette
{ Z
}
{ = }{ \varphi^{-1}( Q)
}
{  }{
}
{    }{
}
{   }{
}
}
{}{}{}
merupakan
\definitionsverweis {serat}{}{}
di atas suatu titik

\mavergleichskette
{\vergleichskette
{ Q
}
{ \in }{ \R^{ m }
}
{  }{
}
{    }{
}
{   }{
}
}
{}{}{.}}
\faktvoraussetzung {Misalkan
\definitionsverweis {diferensial total}{}{}

\mathl{\left(D\varphi\right)_{ P }}{}
\definitionsverweis {surjektif}{}{}
untuk setiap titik

\mavergleichskette
{\vergleichskette
{ P
}
{ \in }{ Z
}
{  }{
}
{    }{
}
{   }{
}
}
{}{}{.}}
\faktfolgerung {Maka $Z$ merupakan
\definitionsverweis {manifold diferensiabel}{}{}
kelas $C^1$ dengan
\definitionsverweis {dimensi}{}{}

\mathl{n - { m }}{.}}
\faktzusatz {}
\faktzusatz {}

}
{

Dengan mengurangkan nilai serat dari pemetaan, kita dapat mereduksi ke kasus

\mavergleichskette
{\vergleichskette
{ Q
}
{ = }{ 0
}
{  }{
}
{    }{
}
{   }{
}
}
{}{}{}
.
Catatan edisi: Sumber langsung menetapkan nilai serat sama dengan nol; kalimat sebelumnya menambahkan translasi kodomain yang diperlukan agar reduksi itu sah.
Berdasarkan
Teorema 8.1,
untuk setiap titik

\mavergleichskette
{\vergleichskette
{ P
}
{ \in }{ Z
}
{  }{
}
{    }{
}
{   }{
}
}
{}{}{}
terdapat suatu lingkungan terbuka

\mavergleichskette
{\vergleichskette
{ P
}
{ \in }{ W
}
{ \subseteq }{ G
}
{    }{
}
{   }{
}
}
{}{}{}
dan suatu
$C^1$-\definitionsverweis {difeomorfisme}{}{}
\maabbdisp {\theta} { W } { V \times V'
} {}
dengan himpunan-himpunan terbuka
\mathkor {} {V \subseteq \R^{n - m }} {dan} {V' \subseteq \R^m} {}
sedemikian sehingga $\theta$ menginduksi bijeksi antara
\mathl{Z \cap W}{} dan

\mavergleichskette
{\vergleichskette
{ V \times \{0\}
}
{ \cong }{ V
}
{  }{
}
{    }{
}
{   }{
}
}
{}{}{}
.
Pembatasan difeomorfisme ini masing-masing pada
\mathl{Z \cap W}{} dan $V$ kita ambil sebagai peta untuk $Z$.
\teilbeweis {}{}{}
{Untuk menunjukkan bahwa hal ini mendefinisikan struktur diferensiabel pada $Z$, misalkan diberikan lingkungan-lingkungan terbuka
\zusatzklammer {di $\R^n$} {} {}
\mathkor {} {W_1} {dan} {W_2} {}
dari

\mavergleichskette
{\vergleichskette
{ P
}
{ \in }{ Z
}
{  }{
}
{    }{
}
{   }{
}
}
{}{}{}
bersama difeomorfisme
\maabbdisp {\theta_1} {W_1 } { V_1 \times V_1'
} {} dan \maabbdisp {\theta_2} {W_2 } { V_2 \times V_2'
} {.}
Dengan beralih ke

\mavergleichskette
{\vergleichskette
{ W
}
{ = }{ W_1 \cap W_2
}
{  }{
}
{    }{
}
{   }{
}
}
{}{}{}
kita dapat mengasumsikan bahwa kedua himpunan terbuka tersebut sama. Pemetaan transisi
\mathl{\theta_2 \circ \theta_1^{-1}}{} merupakan difeomorfisme $C^1$ antara
\zusatzklammer {himpunan-himpunan bagian terbuka dari} {} {}
\mathkor {} {\theta_1(W)} {dan} {\theta_2(W)} {,}
yang memetakan
\mathl{\theta_1(W\cap Z)}{} ke
\mathl{\theta_2(W\cap Z)}{}. Oleh karena itu, menurut
Soal 52 (Analisis (Osnabrück 2021--2023)),
pembatasan pemetaan transisi pada himpunan-himpunan bagian tersebut juga merupakan difeomorfisme $C^1$
\zusatzklammer {antara himpunan-himpunan bagian terbuka di \mathlk{\R^{n - m }}{}} {} {.}}
{}

Catatan edisi: Sumber menyebut seluruh irisan sumbu pada kedua kodomain. Setelah lingkungan dipersempit menjadi irisan kedua lingkungan awal, domain yang tepat ialah kedua citra terbuka dari irisan tersebut sebagaimana ditulis di atas.

}

  \zwischenueberschrift{Pemetaan diferensiabel}

\inputdefinition
{}
{

Misalkan
\mathkor {} {L} {dan} {M} {}
dua
\definitionsverweis {manifold}{}{}
$C^k$ dengan
\definitionsverweis {atlas}{}{}
\mathkor {} {(U_i,U_i', \alpha_i, i \in I)} {dan} {(V_j,V_j', \beta_j, j \in J)} {.}
Misalkan

\mavergleichskette
{\vergleichskette
{ 1
}
{ \leq }{ \ell
}
{ \leq }{ k
}
{    }{
}
{   }{
}
}
{}{}{.}
Suatu
\definitionsverweis {pemetaan kontinu}{}{}
\maabbdisp {\varphi} { L } { M
} {}
disebut suatu $C^\ell$-\definitionswort {pemetaan diferensiabel}{,} jika untuk semua
\mathkor {} {i \in I} {dan semua} {j \in J} {}
pemetaan
\maabbdisp {\beta_j \circ  \varphi \circ (\alpha_i)^{-1}} {\alpha_i( \varphi^{-1}(V_j) \cap U_i) } { V_j'
} {}
bersifat
$C^\ell$-\definitionsverweis {diferensiabel}{}{.}

}

Karena
\mathl{\alpha_i \left(  \varphi^{-1}(V_j) \cap U_i  \right)}{} terbuka, definisi ini mereduksi konsep diferensiabilitas bagi pemetaan antara manifold pada konsep diferensiabilitas bagi pemetaan antara himpunan-himpunan terbuka di ruang vektor real. Karena manifold $C^k$ dapat dipandang sebagai manifold $C^\ell$ untuk

\mavergleichskette
{\vergleichskette
{ \ell
}
{ \leq }{ k
}
{  }{
}
{    }{
}
{   }{
}
}
{}{}{}
,
pada dasarnya cukup membicarakan pemetaan $C^k$ antara manifold-manifold $C^k$. Kasus yang khususnya penting adalah
\mathl{k=1,2, \infty}{.} Perhatikan bahwa untuk

\mavergleichskette
{\vergleichskette
{ k
}
{ = }{ 1
}
{  }{
}
{    }{
}
{   }{
}
}
{}{}{}
kita berbicara tentang pemetaan diferensiabel, meskipun
\zusatzklammer {sejauh ini} {} {}
belum ada suatu \anfuehrung{turunan}{.}



\inputfaktbeweis
{Differenzierbare Mannigfaltigkeit/Differenzierbare Abbildungen/Elementare funktorielle Eigenschaften/Fakt}
{Proposisi}
{}
{

\faktsituation {Misalkan
\mathkor {} {L,M} {dan} {N} {}
merupakan
\definitionsverweis {manifold}{}{}
$C^k$.}
\faktuebergang {Maka berlaku pernyataan-pernyataan berikut.}
\faktfolgerung {\aufzaehlungvier{\definitionsverweis {Identitas}{}{}
\maabbdisp {\operatorname{Id} \,} { L } { L
} {}
merupakan
\definitionsverweis {pemetaan}{}{}
$C^k$.
}{Setiap
\definitionsverweis {pemetaan konstan}{}{}
\maabbdisp {\varphi} { L } { M
} {}
merupakan
\definitionsverweis {pemetaan}{}{}
$C^k$.
}{Untuk setiap
\definitionsverweis {himpunan bagian terbuka}{}{}

\mavergleichskette
{\vergleichskette
{ U
}
{ \subseteq }{ L
}
{  }{
}
{    }{
}
{   }{
}
}
{}{}{}
\definitionsverweis {penyematan terbuka}{}{}
\maabb {} { U } { L
} {}
merupakan
\definitionsverweis {pemetaan}{}{}
$C^k$.
}{Misalkan
\maabbdisp {\varphi} { L } { M
} {}
dan
\maabbdisp {\psi} { M } { N
} {}
merupakan
\definitionsverweis {pemetaan}{}{}
$C^k$. Maka
\definitionsverweis {komposisi}{}{}
\maabbdisp {\psi \circ \varphi} { L } { N
} {}
juga merupakan
\definitionsverweis {pemetaan}{}{}
$C^k$.
}}
\faktzusatz {}
\faktzusatz {}

}
{

\teilbeweis {}{}{}
{(1). Pemetaan-pemetaan yang harus diperiksa tepatnya adalah pemetaan transisi
\mathl{\alpha_j \circ \alpha_i^{-1}}{,} yang menurut definisi suatu
\definitionsverweis {manifold diferensiabel}{}{}
$C^k$ merupakan
\definitionsverweis {difeomorfisme}{}{}
$C^k$.}
{}
\teilbeweis {}{}{}
{(2). Dalam setiap peta, pemetaan yang harus diperiksa bersifat konstan, sehingga diferensiabel hingga sebarang orde.}
{}
\teilbeweis {}{}{}
{(3). Pemetaan-pemetaan yang harus diperiksa untuk

\mavergleichskette
{\vergleichskette
{ i,j
}
{ \in }{ I
}
{  }{
}
{    }{
}
{   }{
}
}
{}{}{}
sama dengan

\mathdisp {\alpha_i \left(  U \cap U_i  \cap U_j  \right)  \subseteq  \alpha_i \left(  U_i  \cap U_j  \right) \stackrel{ \alpha_j \circ \alpha_i^{-1} }{  \longrightarrow} U'_j} { , }

sehingga merupakan penyematan terbuka yang diikuti oleh pemetaan transisi diferensiabel.}
{}
\teilbeweis {}{}{}
{(4). Misalkan
\maabbdisp {\gamma_{  \ell }} {W_{  \ell }} {W'_{  \ell }
} {}
merupakan peta-peta untuk $N$. Maka, untuk semua kombinasi indeks yang mungkin, komposisi
\zusatzklammer {yang dibatasi pada himpunan-himpunan bagian terbuka tertentu} {} {}

\mavergleichskettedisphandlinks
{\vergleichskettedisphandlinks
{ \gamma_{  \ell } \circ (\psi  \circ \varphi) \circ \alpha_i^{-1}
}
{ =} { \gamma_{  \ell } \circ  \psi  \circ \beta_j^{-1} \circ \beta_j \circ \varphi \circ \alpha_i^{-1}
}
{ =} { \left(  \gamma_{  \ell } \circ \psi \circ \beta_j^{-1}  \right)  \circ \left(  \beta_j \circ \varphi \circ \alpha_i^{-1}  \right)
}
{ } {
}
{ } {
}
}
{}{}{}
bersifat
\definitionsverweis {diferensiabel}{}{}
menurut aturan rantai.
Untuk

\mavergleichskette
{\vergleichskette
{ k
}
{ \geq }{ 2
}
{  }{
}
{    }{
}
{   }{
}
}
{}{}{}
kita menggunakan
Soal 46.10 (Analisis (Osnabrück 2021--2023)).}
{}

}

\inputdefinition
{}
{

Misalkan
\mathkor {} {L} {dan} {M} {}
dua
\definitionsverweis {manifold}{}{}
$C^k$. Suatu
\definitionsverweis {homeomorfisme}{}{}
\maabbdisp {\varphi} { L } { M
} {}
disebut
\definitionswortpraemath {C^k}{ difeomorfisme }{,}
jika
\mathkor {} {\varphi} {dan} {\varphi^{-1}} {}
keduanya merupakan
\definitionsverweis {pemetaan}{}{}
$C^k$.

}

\inputdefinition
{}
{

Dua
\definitionsverweis {manifold}{}{}
$C^k$
\mathkor {} {L} {dan} {M} {}
disebut
\definitionswortpraemath {C^k}{ difeomorfik }{,}
jika terdapat suatu
\definitionsverweis {difeomorfisme}{}{}
$C^k$ di antara keduanya.

}

\inputbemerkung
{}
{

Untuk suatu
\definitionsverweis {manifold}{}{}
$C^k$ $M$ dengan atlas $C^k$
\mathl{\left(  U_i,U_i',\alpha_i, i \in I  \right)}{} terdapat suatu \stichwort {atlas maksimal} {,} yang kompatibel dengan struktur diferensiabel yang diberikan oleh atlas tersebut. Atlas ini terdiri atas himpunan semua
\definitionsverweis {homeomorfisme}{}{}
\maabbdisp {\beta} { U } { V
} {}
dengan himpunan-himpunan terbuka

\mavergleichskette
{\vergleichskette
{ U
}
{ \subseteq }{ M
}
{  }{
}
{    }{
}
{   }{
}
}
{}{}{}
dan

\mavergleichskette
{\vergleichskette
{ V
}
{ \subseteq }{ \R^n
}
{  }{
}
{    }{
}
{   }{
}
}
{}{}{}
yang mempunyai sifat bahwa pemetaan ini beserta inversnya merupakan
\definitionsverweis {pemetaan}{}{}
$C^k$
\zusatzklammer {terhadap struktur yang diberikan oleh atlas tersebut} {} {.}
Atlas maksimal ini tentu memuat atlas awal, tetapi secara umum jauh lebih besar. Sebagai contoh, untuk setiap peta
\maabb {\beta} { U } { V
} {}
dan setiap himpunan bagian terbuka

\mavergleichskette
{\vergleichskette
{ U'
}
{ \subseteq }{ U
}
{  }{
}
{    }{
}
{   }{
}
}
{}{}{}
atlas tersebut juga memuat pembatasan peta itu pada $U'$. Hal yang penting adalah bahwa pemetaan identitas
\maabbdisp {\operatorname{Id}} {(M,A)} {(M,B)
} {,}
dengan $A$ menyatakan atlas awal dan $B$ atlas maksimal, merupakan
\definitionsverweis {difeomorfisme}{}{}
$C^k$ antara kedua manifold tersebut, sebagaimana langsung mengikuti dari definisi. Struktur diferensiabel pada manifold yang direpresentasikan oleh atlas lebih penting daripada atlas itu sendiri; struktur ini menentukan pemetaan mana yang diferensiabel dan pemetaan mana yang merupakan difeomorfisme. Peta-peta dalam atlas maksimal kadang-kadang juga disebut peta
\zusatzklammer {yang diperumum} {} {}
bagi manifold tersebut.

Catatan edisi: Sumber hanya mensyaratkan regularitas pemetaan maju. Regularitas invers ditambahkan karena sebuah homeomorfisme yang diferensiabel belum tentu merupakan difeomorfisme dan belum tentu kompatibel dengan atlas awal.

}

  \zwischenueberschrift{Fungsi diferensiabel}

Suatu
\definitionsverweis {pemetaan diferensiabel}{}{}
$C^k$
\maabbdisp {f} {M} {\R
} {}
dari suatu
\definitionsverweis {manifold diferensiabel}{}{}
ke bilangan real juga disebut \stichwort {fungsi diferensiabel} {} $C^k$. Menurut definisi, ini berarti bahwa untuk setiap peta
\maabbdisp {\alpha} {U} {V
} {}
fungsi komposit
\maabbdisp {f \circ \alpha^{-1}} {V} {\R
} {}
merupakan
\definitionsverweis {fungsi}{}{}
$C^k$. Himpunan semua fungsi $C^k$ pada $M$ dinotasikan dengan
\mathl{C^k(M,\R)}{}.


\inputfaktbeweis
{Mannigfaltigkeit/Differenzierbare Funktionen/Strukturelle Eigenschaften/Fakt}
{Lema}
{}
{

\faktsituation {Misalkan $M$ suatu
\definitionsverweis {manifold diferensiabel}{}{} dan
\maabbdisp {f,g} { M } {\R
} {}
merupakan
\definitionsverweis {fungsi diferensiabel}{}{}
pada $M$.}
\faktuebergang {Maka berlaku pernyataan-pernyataan berikut.}
\faktfolgerung {\aufzaehlungvier{Pemetaan
\maabbeledisp {f \times g} { M } {\R^2
} { x } {(f(x),g(x))
} {,} bersifat diferensiabel.
}{ $f+g$
\definitionsverweis {diferensiabel}{}{.}
}{ $f \cdot g$
\definitionsverweis {diferensiabel}{}{.}
}{Jika $f$ tidak mempunyai titik nol, maka
\mathl{f^{-1}}{} juga diferensiabel.
}}
\faktzusatz {}
\faktzusatz {}

 }
{ Lihat Soal 8.3. }

Secara khusus, fungsi-fungsi diferensiabel pada suatu manifold membentuk
\definitionsverweis {gelanggang komutatif}{}{.}

Jika
\maabbdisp {\alpha} {U} {V
} {}
merupakan suatu peta dengan

\mavergleichskette
{\vergleichskette
{ V
}
{ \subseteq }{ \R^n
}
{  }{
}
{    }{
}
{   }{
}
}
{}{}{}
terbuka, maka setiap proyeksi $x_i$ menghasilkan suatu fungsi diferensiabel
\maabbdisp {x_i \circ \alpha} {U} {\R
} {,}
yang biasanya kembali dinotasikan dengan $x_i$. Dalam hal ini dikatakan bahwa fungsi-fungsi
\mathl{x_1 , \ldots , x_n}{} membentuk \stichwort {koordinat diferensiabel} {} untuk

\mavergleichskette
{\vergleichskette
{ U
}
{ \subseteq }{ M
}
{  }{
}
{    }{
}
{   }{
}
}
{}{}{}
.
Bagi suatu fungsi yang terdiferensialkan secara kontinu
\maabbdisp {f} {U} {\R
} {}
menurut definisi fungsi
\maabbdisp {f \circ \alpha^{-1}} {V} {\R
} {}
\definitionsverweis {terdiferensialkan secara kontinu}{}{,}
artinya, untuk setiap $i$ terdapat
\definitionsverweis {turunan parsial}{}{}

\mathdisp {{ \frac{   \partial  (f \circ \alpha^{-1})  }{  \partial   x_i  } }} { , }

yang kembali merupakan fungsi
\zusatzklammer {kontinu} {} {}
pada $V$. Oleh karena itu,

\mathdisp {{ \frac{   \partial  (f \circ \alpha^{-1})   }{  \partial   x_i  } }  \circ \alpha} {  }

merupakan fungsi pada $U$. Fungsi ini secara umum kembali dinotasikan dengan
\mathl{{ \frac{   \partial  f  }{  \partial   x_i  } }}{}.

\chapter{Lembar Kerja 8}
\setcounter{section}{8}

  \zwischenueberschrift{Soal latihan}

\inputaufgabe
{}
{

Misalkan
\mathkor {} {M} {dan} {N} {}
adalah
\definitionsverweis {manifold diferensiabel}{}{}
dan
\maabbdisp {\varphi} { M } { N
} {}
sebuah pemetaan. Misalkan
\mathl{M= \bigcup_{i \in I} U_i}{} suatu
\definitionsverweis {tutupan terbuka}{}{}
dari $M$. Tunjukkan bahwa $\varphi$
\definitionsverweis {diferensiabel}{}{}
jika dan hanya jika semua pembatasan
\mathl{\varphi_i= \varphi \vert_{U_i }}{} diferensiabel.

}
{} {}

\inputaufgabe
{}
{

Misalkan $M$ sebuah
\definitionsverweis {manifold diferensiabel}{}{}
dan
\maabbdisp {\alpha} { U } { V
} {}
sebuah peta
\zusatzklammer {jadi
\mathl{U \subseteq M}{} dan
\mathl{V \subseteq \R^n}{} terbuka} {} {.}
Tunjukkan bahwa $\alpha$ adalah sebuah
\definitionsverweis {difeomorfisme}{}{}.

}
{} {}

\inputaufgabe
{}
{

Misalkan $M$ sebuah
\definitionsverweis {manifold diferensiabel}{}{} dan
\maabbdisp {f,g} { M } {\R
} {}
\definitionsverweis {fungsi-fungsi diferensiabel}{}{} pada $M$. Buktikan pernyataan-pernyataan berikut. \aufzaehlungvier{Pemetaan
\maabbeledisp {f \times g} { M } {\R^2
} { x } {(f(x),g(x))
} {,} diferensiabel.
}{ $f+g$
\definitionsverweis {diferensiabel}{}{.}
}{ $f \cdot g$
\definitionsverweis {diferensiabel}{}{.}
}{Jika $f$ tidak memiliki titik nol, maka
\mathl{f^{-1}}{} juga diferensiabel.
}

}
{} {}

\inputaufgabe
{}
{

Misalkan
\mathl{S^2 \subseteq \R^3}{} adalah
\definitionsverweis {sfera}{}{.}
Dengan menggunakan
\definitionsverweis {peta-peta stereografis}{}{,}
tunjukkan bahwa pembatasan koordinat ruang
\mathl{x,y,z}{} pada sfera merupakan
\definitionsverweis {fungsi-fungsi diferensiabel}{}{}.

}
{} {}

\inputaufgabe
{}
{

Misalkan
\mathkor {} {M} {dan} {N} {}
adalah
\definitionsverweis {manifold diferensiabel}{}{}
dan
\maabbdisp {\varphi} { M } { N
} {}
sebuah
\definitionsverweis {pemetaan diferensiabel}{}{.}
Tunjukkan bahwa pemetaan ini menginduksi sebuah
\definitionsverweis {homomorfisme gelanggang}{}{}
\maabbeledisp {\varphi^*} {C^1(N,\R)} {C^1(M,\R)
} { f } {f \circ \varphi
} {,}
melalui penarikan balik.

}
{} {}

\inputaufgabe
{}
{

Tunjukkan bahwa untuk
\mathl{m \leq n}{}, penyematan subruang
\mathl{\R^m}{} ke dalam $\R^n$ yang diberikan oleh
\mathl{(x_1 , \ldots , x_m) \mapsto (x_1 , \ldots , x_m,0 , \ldots , 0)}{}
\definitionsverweis {diferensiabel hingga orde berapa pun}{}{}.

}
{} {}

\inputaufgabe
{}
{

Tunjukkan bahwa permukaan silinder terbuka
\mathl{S^1 \times {]0,1[}}{},
\mathl{S^1 \times \R}{,} bidang yang ditusuk
\mathl{\R^2 \setminus \{(0,0)\}}{}, dan
\mathl{S^2 \setminus \{N,S\}}{}
saling \definitionsverweis {difeomorfik}{}{}.

}
{} {}

Soal berikut menggunakan definisi di bawah ini.

Sebuah fungsi
\maabbdisp {f} { {\mathbb K}^n } { {\mathbb K}
} {}
disebut
\definitionswort {homogen berderajat}{}
$d$, jika untuk setiap titik
\mathl{P \in {\mathbb K}^n}{} dan setiap
\mathl{\lambda \in {\mathbb K}}{} berlaku hubungan

\mavergleichskettedisp
{\vergleichskette
{f(\lambda P)
}
{ =} { \lambda^d f(P)
}
{ } {
}
{ } {
}
{ } {
}
}
{}{}{}
untuk semua pilihan tersebut.

\inputaufgabe
{}
{

Misalkan
\maabbdisp {f} {\R^n } {\R
} {}
sebuah
\definitionsverweis {fungsi homogen}{}{} yang
\definitionsverweis {terdiferensialkan secara kontinu}{}{,}
dan yang pada
\definitionsverweis {serat}{}{}
$F$ di atas
\mathl{a \neq 0}{}
bersifat \definitionsverweis {reguler}{}{.}
Tunjukkan bahwa setiap serat di atas
\mathl{b \neq 0,\ b=\lambda^d a}{}
untuk suatu skalar real tak nol,
\definitionsverweis {difeomorfik}{}{} dengan $F$ dan merupakan sebuah
\definitionsverweis {manifold}{}{}.

Catatan edisi: Syarat penskalaan real ditambahkan. Klaim sumber untuk semua nilai tak nol gagal pada derajat genap ketika tanda nilai serat berbeda; misalnya, fungsi kuadrat norma mempunyai serat positif berbentuk bola tetapi serat negatif kosong.

}
{} {}

\inputaufgabe
{}
{

Misalkan
\mathl{]a,b[}{} sebuah
\definitionsverweis {interval terbuka}{}{}
dan
\maabbdisp {f} {]a,b[} {\R_{+}
} {}
sebuah
\definitionsverweis {fungsi yang terdiferensialkan secara kontinu}{}{.}
Misalkan $M$ adalah permukaan benda putar terkait. Tunjukkan bahwa himpunan ini merupakan manifold yang difeomorfik dengan silinder terbuka.

}
{} {}

\inputaufgabe
{}
{

Tunjukkan bahwa permukaan elipsoid dan permukaan bola satuan
$C^\infty$-\definitionsverweis {difeomorfik}{}{}.

}
{} {}

\inputaufgabegibtloesung
{}
{

Berikan sebuah contoh
\definitionsverweis {manifold diferensiabel}{}{}
\definitionsverweis {berdimensi dua}{}{} yang
\definitionsverweis {terhubung}{}{},
sebutlah $M$, dan sebuah titik

\mavergleichskette
{\vergleichskette
{ P
}
{ \in }{ M
}
{  }{
}
{    }{
}
{   }{
}
}
{}{}{}
sedemikian sehingga
\mathkor {} {M} {dan} {M \setminus \{P\}} {}
\definitionsverweis {difeomorfik}{}{}
satu sama lain.

}
{} {}

\inputaufgabe
{}
{

Tunjukkan bahwa
\definitionsverweis {pemetaan antipodal}{}{}
\maabbeledisp {} {S^n } { S^n
} {(x_1 , \ldots , x_{n+1})} {(-x_1 , \ldots , -x_{n+1})
} {,}
adalah sebuah
\definitionsverweis {difeomorfisme}{}{} \definitionsverweis {tanpa titik tetap}{}{}
yang inversnya adalah dirinya sendiri.

}
{} {}

\inputaufgabegibtloesung
{}
{

Tunjukkan bahwa sebuah
\definitionsverweis {manifold diferensiabel}{}{}
$M$ berdimensi

\mavergleichskette
{\vergleichskette
{ n
}
{ \geq }{ 1
}
{  }{
}
{    }{
}
{   }{
}
}
{}{}{}
memiliki tak hingga banyak
\definitionsverweis {difeomorfisme}{}{}
\maabbdisp {\varphi} { M } { M
} {}
yang berbeda.

}
{} {}

Soal-soal berikut dimaksudkan untuk menjelaskan mengapa manifold dirumuskan dengan tutupan terbuka.

Misalkan
\mathl{\left(  x_n  \right)_{n \in  \N  }}{} sebuah
\definitionsverweis {barisan}{}{}
dalam suatu
\definitionsverweis {ruang topologis}{}{}
$X$. Kita mengatakan bahwa, untuk suatu

\mavergleichskette
{\vergleichskette
{ x
}
{ \in }{ X
}
{  }{
}
{    }{
}
{   }{
}
}
{}{}{}
barisan tersebut \definitionswort {konvergen ke titik itu}{,} jika sifat berikut terpenuhi.

Untuk setiap
\definitionsverweis {lingkungan terbuka}{}{}

\mavergleichskette
{\vergleichskette
{ U
}
{ \subseteq }{ X
}
{  }{
}
{    }{
}
{   }{
}
}
{}{}{}
yang memuat $x$, terdapat suatu

\mavergleichskette
{\vergleichskette
{ n_0
}
{ \in }{ \N
}
{  }{
}
{    }{
}
{   }{
}
}
{}{}{}
sedemikian sehingga untuk semua

\mavergleichskette
{\vergleichskette
{ n
}
{ \geq }{ n_0
}
{  }{
}
{    }{
}
{   }{
}
}
{}{}{}
suku barisan $x_n$ termasuk dalam $U$.

Dalam hal ini $x$ disebut \definitionswort {nilai batas}{} atau \definitionswort {limit}{} barisan tersebut. Kita juga menuliskannya sebagai

\mavergleichskettedisp
{\vergleichskette
{  \lim_{n \rightarrow \infty}  x_n
}
{ =} { x
}
{ } {
}
{ } {
}
{ } {
}
}
{}{}{.}
Jika barisan memiliki limit, kita juga mengatakan bahwa barisan itu \definitionswort {konvergen}{}
\zusatzklammer {tanpa menyebut nilai limit tertentu} {} {,}
dan jika tidak, bahwa barisan itu \definitionswort {divergen}{.}

\inputaufgabe
{}
{

Misalkan $M$ sebuah
\definitionsverweis {manifold topologis}{}{}
dan
\mathl{\left(   x_n  \right)_{n \in \N }}{} sebuah
\definitionsverweis {barisan}{}{}
dalam $M$. Tunjukkan bahwa barisan tersebut
\definitionsverweis {konvergen}{}{,}
jika dan hanya jika terdapat sebuah
\definitionsverweis {domain peta}{}{}
dari $M$ yang memuat hampir semua suku barisan dan sedemikian sehingga barisan citra yang bersesuaian konvergen di dalam citra peta.

}
{} {}

\inputaufgabe
{}
{

Misalkan $M$ sebuah
\definitionsverweis {manifold topologis}{}{}
dan
\maabbdisp {\alpha_i} {U_i } { V_i
} {}
sebuah keluarga
\definitionsverweis {peta}{}{}
dengan
\definitionsverweis {pemetaan transisi}{}{}
\maabbdisp {\varphi_{ij} = \alpha_j \circ (\alpha_i)^{-1}} {V_i \cap \alpha_i(U_i \cap U_j) } { V_j \cap \alpha_j(U_i \cap U_j)
} {.}
Tunjukkan bahwa manifold $M$ dapat direkonstruksi dari keluarga

\mathbed {V_i} {}
 {i \in I} {}
 {} {} {} {,}
himpunan-himpunan bagian
\mathl{V_{ij} \subseteq V_i}{} dan pemetaan transisi
\maabbdisp {\varphi_{ij}} {V_{ij}} { V_{ji}
} {}
yang bersesuaian.

a) Pada

\mavergleichskettedisp
{\vergleichskette
{N
}
{ \defeq} { \biguplus_{i \in I} V_i
}
{ } {
}
{ } {
}
{ } {
}
}
{}{}{}
perhatikan
\definitionsverweis {relasi ekuivalensi}{}{,}
yang menyamakan dua titik
\mathl{P\in V_i}{} dan
\mathl{Q \in V_j}{} jika
\mathl{\varphi_{ij}}{} memetakan titik pertama ke titik kedua.

b) Lengkapi himpunan hasil bagi
\mathl{N/ \sim}{} dengan sebuah topologi yang sesuai.

c) Definisikan peta-peta pada
\mathl{N/ \sim}{}.

d) Tunjukkan bahwa
\mathkor {} {M} {dan} {N/ \sim} {}
\definitionsverweis {homeomorfik}{}{}.

}
{} {}

Prosedur konstruksi manifold yang dijelaskan dalam soal sebelumnya berlaku untuk suatu keluarga himpunan bagian terbuka di dalam $\R^n$ dengan pemetaan transisi yang memenuhi syarat koklus dari
Soal 7.11.
Namun, ruang topologis yang dihasilkan tidak serta-merta merupakan ruang Hausdorff. Konstruksi ini disebut \stichwort {perekatan terbuka} {} ruang-ruang.

\inputaufgabe
{}
{

Kita mengambil dua salinan garis riil, yaitu

\mavergleichskette
{\vergleichskette
{G_1
}
{ = }{\R
}
{  }{
}
{    }{
}
{   }{
}
}
{}{}{}
dan

\mavergleichskette
{\vergleichskette
{G_2
}
{ = }{\R
}
{  }{
}
{    }{
}
{   }{
}
}
{}{}{}
beserta pemetaan perekatan
\maabbeledisp {\varphi} {G_1 \setminus \{0\} } { G_2 \setminus \{0\}
} { x } { x^{-1}
} {.}
Misalkan $M$ ruang topologis yang dihasilkan menurut konstruksi yang dijelaskan dalam
Soal 8.15.
Tunjukkan bahwa $M$
\definitionsverweis {homeomorfik}{}{}
dengan
$1$-\definitionsverweis {sfera}{}{}.

}
{} {}

\inputaufgabe
{}
{

Kita mengambil dua salinan garis riil, yaitu

\mavergleichskette
{\vergleichskette
{G_1
}
{ = }{\R
}
{  }{
}
{    }{
}
{   }{
}
}
{}{}{}
dan

\mavergleichskette
{\vergleichskette
{G_2
}
{ = }{\R
}
{  }{
}
{    }{
}
{   }{
}
}
{}{}{}
beserta pemetaan perekatan
\maabbdisp {\operatorname{Id}} {G_1 \setminus \{0\} } { G_2 \setminus \{0\}
} {.}
Misalkan $M$ ruang topologis yang dihasilkan menurut konstruksi yang dijelaskan dalam
Soal 8.15.
Tunjukkan bahwa $M$ bukan
\definitionsverweis {manifold}{}{}.

}
{} {}

  \zwischenueberschrift{Soal untuk dikumpulkan}

Dalam soal berikut, tafsirkan ${\mathbb C}$ sebagai $\R^2$.

\inputaufgabe
{4}
{

Perhatikan pemetaan
\maabbeledisp {} { {\mathbb C}^2} { {\mathbb C}
} {(z,w)} { zw
} {.}
Untuk titik manakah
\mathl{u \in {\mathbb C}}{} serat di atas $u$ merupakan manifold? Dalam setiap kasus, berikan deskripsi kelas difeomorfisme yang sesederhana mungkin.

}
{} {}

\inputaufgabe
{6}
{

Misalkan diberikan dua titik
\mathkor {} {P} {dan} {Q} {}
pada
\definitionsverweis {permukaan bola satuan}{}{.}
Tunjukkan bahwa terdapat sebuah
\definitionsverweis {difeomorfisme}{}{}
dari permukaan bola itu ke dirinya sendiri yang memetakan $P$ ke $Q$.

}
{} {}

\inputaufgabe
{4 (1+1+2)}
{

Misalkan $M$ sebuah
\definitionsverweis {manifold diferensiabel}{}{.}
Untuk setiap himpunan bagian terbuka

\mavergleichskette
{\vergleichskette
{ U
}
{ \subseteq }{ M
}
{  }{
}
{    }{
}
{   }{
}
}
{}{}{}
kita perhatikan himpunan
\mathl{C^1(U,\R)}{} yang terdiri atas
\definitionsverweis {fungsi-fungsi diferensiabel}{}{}
pada $U$. Misalkan

\mavergleichskette
{\vergleichskette
{ M
}
{ = }{ \bigcup_{i \in I} U_i
}
{  }{
}
{    }{
}
{   }{
}
}
{}{}{}
sebuah
\definitionsverweis {tutupan terbuka}{}{.}
\aufzaehlungdreiabc{Tunjukkan bahwa jika

\mavergleichskette
{\vergleichskette
{ V
}
{ \subseteq }{ U
}
{  }{
}
{    }{
}
{   }{
}
}
{}{}{}
terbuka dan

\mavergleichskette
{\vergleichskette
{ f
}
{ \in }{ C^1(U,\R)
}
{  }{
}
{    }{
}
{   }{
}
}
{}{}{}
maka
\definitionsverweis {pembatasan}{}{}
$f \vert_V$ juga termasuk dalam
\mathl{C^1(V,\R)}{}.
}{Misalkan

\mavergleichskette
{\vergleichskette
{ f
}
{ \in }{ C^1(M,\R)
}
{  }{
}
{    }{
}
{   }{
}
}
{}{}{.}
Tunjukkan bahwa

\mavergleichskette
{\vergleichskette
{ f
}
{ = }{ 0
}
{  }{
}
{    }{
}
{   }{
}
}
{}{}{}
jika dan hanya jika semua pembatasan

\mavergleichskette
{\vergleichskette
{ f \vert_{U_i }
}
{ = }{ 0
}
{  }{
}
{    }{
}
{   }{
}
}
{}{}{}
berlaku.
}{Misalkan diberikan suatu keluarga

\mavergleichskette
{\vergleichskette
{ f_i
}
{ \in }{ C^1(U_i,\R)
}
{  }{
}
{    }{
}
{   }{
}
}
{}{}{}
fungsi-fungsi yang memenuhi \anfuehrung{syarat kecocokan}{}

\mavergleichskettedisp
{\vergleichskette
{ f_i \vert_{U_i \cap U_j }
}
{ =} { f_j \vert_{U_i \cap U_j }
}
{ } {
}
{ } {
}
{ } {
}
}
{}{}{}
untuk semua $i,j$. Tunjukkan bahwa terdapat suatu

\mavergleichskette
{\vergleichskette
{ f
}
{ \in }{ C^1(M,\R)
}
{  }{
}
{    }{
}
{   }{
}
}
{}{}{}
yang memenuhi

\mavergleichskette
{\vergleichskette
{ f \vert_{U_i }
}
{ = }{ f_i
}
{  }{
}
{    }{
}
{   }{
}
}
{}{}{}
untuk semua $i$.
}

}
{} {}

\inputaufgabe
{5}
{

Misalkan $M$ sebuah
\definitionsverweis {manifold diferensiabel}{}{,} yang memiliki setidaknya dua titik. Tunjukkan bahwa terdapat
\definitionsverweis {fungsi-fungsi diferensiabel}{}{}
\maabbdisp {f,g} { M } {\R
} {}
dengan
\mathl{f,g \neq 0}{,} tetapi
\mathl{fg=0}{.}

}
{} {}

\section*{Solusi yang disediakan oleh sumber}
\addcontentsline{toc}{section}{Solusi yang disediakan oleh sumber}
\subsection*{Solusi Soal 8.11}
Misalkan

\mavergleichskettedisp
{\vergleichskette
{M
}
{ =} { \R^2 \setminus \N_+
}
{ =} { \R^2 \setminus \{ (1,0) , (2,0), (3,0), \ldots  \}
}
{ } {
}
{ } {
}
}
{}{}{,}
yakni dari $\R^2$ kita mengeluarkan bilangan-bilangan asli positif yang terletak pada sumbu $x$. Himpunan ini terbuka di dalam $\R^2$ dan karena itu merupakan manifold diferensiabel. Manifold ini terhubung lintasan: misalnya, setiap titik di $M$ dengan
\mathl{y \geq 0}{} dapat dihubungkan oleh lintasan lurus dengan
\mathl{(0,1)}{}, dan setiap titik di $M$ dengan
\mathl{y \leq 0}{} dapat dihubungkan oleh lintasan lurus dengan
\mathl{(0,-1)}{}; kedua titik tersebut juga dapat dihubungkan oleh lintasan lurus. Sekarang misalkan

\mavergleichskettedisp
{\vergleichskette
{P
}
{ =} { (0,0)
}
{ } {
}
{ } {
}
{ } {
}
}
{}{}{}
dan

\mavergleichskettedisp
{\vergleichskette
{M'
}
{ =} { M \setminus \{P\}
}
{ } {
}
{ } {
}
{ } {
}
}
{}{}{.}
Pemetaan
\maabbeledisp {} {\R^2} {\R^2
} {(x,y)} {(x-1,y)
} {,}
merupakan sebuah difeomorfisme
\zusatzklammer {sebagai translasi} {} {},
dan citra $M$ tepat sama dengan $M'$. Oleh karena itu,
\mathkor {} {M} {dan} {M'=M \setminus \{P\}} {}
difeomorfik satu sama lain.
\subsection*{Solusi Soal 8.13}
Kita perhatikan fungsi-fungsi berbentuk
\maabbeledisp {\psi_c} {[0,1]} {[0,1]
} { x } { x + \varphi_c (x)
} {,}
dengan

\mavergleichskettedisp
{\vergleichskette
{  \varphi_c(x)
}
{ =} { \begin{cases}  0 \text{ untuk } x \leq { \frac{   1  }{  3 } } \, , \\  c \left(  x- { \frac{   1  }{  3 } }   \right)^2 \left(  x- { \frac{   2 }{  3 } }   \right)^2 \text{ untuk } { \frac{   1  }{  3 } } < x < { \frac{   2 }{  3 } }  \, ,  \\  0 \text{ untuk } x \geq { \frac{   2 }{  3 } }  \, .  \end{cases}
}
{ } {
}
{ } {
}
{ } {
}
}
{}{}{}
dan parameter

\mavergleichskette
{\vergleichskette
{ c
}
{ \in }{  \R_{\geq 0}
}
{  }{
}
{    }{
}
{   }{
}
}
{}{}{.}
Fungsi $\varphi_c$ kontinu dan juga terdiferensialkan secara kontinu karena akar-akar polinom tersebut berlipat dua. Dengan demikian, $\psi_c$ juga terdiferensialkan secara kontinu. Untuk $c$ yang cukup kecil, $\psi_c$ naik ketat dan mendefinisikan pemetaan bijektif yang terdiferensialkan secara kontinu dari interval satuan ke dirinya sendiri, serta merupakan identitas pada sepertiga bawah dan sepertiga atas.

Dengan $\psi_c$, untuk setiap vektor tak nol di bola satuan terbuka kita mendefinisikan
\maabbeledisp {} { U \left(   0 , 1  \right) \setminus \{0\} } { U \left(   0 , 1  \right) \setminus \{0\}
} { \begin{pmatrix}  x_1  \\\vdots \\  x_n   \end{pmatrix} } { \frac{\psi_c \left(  \sqrt{  x_1^2 + \cdots + x_n^2 }  \right)}{\sqrt{  x_1^2 + \cdots + x_n^2 }} \begin{pmatrix}  x_1  \\\vdots \\  x_n   \end{pmatrix}
} {.}
Pemetaan ini diperluas dengan mengirim titik nol ke titik nol. Karena fungsi radial itu merupakan identitas pada suatu lingkungan titik nol, perluasan tersebut mulus dan merupakan difeomorfisme bola satuan terbuka. Jadi titik tersebut diskalakan bergantung pada jari-jarinya, sedangkan pada
\mathl{U \left(   0 , { \frac{   1  }{  3 } }   \right)}{} dan pada
\mathl{U \left(   0 , 1  \right) \setminus  U \left(   0 ,   { \frac{   2 }{  3 } }   \right)}{} pemetaan itu merupakan identitas.

Catatan edisi: Rumus sumber mengalikan vektor langsung dengan nilai radius baru sehingga bahkan bukan identitas ketika fungsi gangguannya nol. Pembagian dengan radius di atas memberikan penskalaan radial yang benar; definisi terpisah di titik nol sah karena faktor itu identik dengan satu di dekat nol.

Jika $M$ sebuah manifold berdimensi $\geq 1$, kita pilih sebuah peta
\maabbdisp {\alpha} { U } {V
} {}
dengan

\mavergleichskette
{\vergleichskette
{ V
}
{ \subseteq }{  \R^n
}
{  }{
}
{    }{
}
{   }{
}
}
{}{}{,}
dan kita dapat menganggap
\zusatzklammer {setelah memperkecilnya} {} {}
bahwa $V$
\zusatzklammer {merupakan sebuah bola terbuka dan kemudian} {} {}
adalah bola satuan terbuka. Difeomorfisme yang telah dikonstruksi pada $V$ menghasilkan difeomorfisme pada $U$. Karena difeomorfisme ini merupakan identitas pada bola bagian luar
\zusatzklammer {mulai dari jari-jari ${ \frac{   2 }{  3 } }$} {} {},
difeomorfisme tersebut dapat diperluas ke $M$ dengan menggunakan identitas pada $M \setminus \alpha^{-1} \left(  U \left(   0 , { \frac{   2 }{  3 } }   \right)  \right)$.

\part{Unit 9}
\chapter[Kuliah 9]{Kuliah 9: Ruang Singgung Manifold Diferensiabel}
\setcounter{section}{9}

  \zwischenueberschrift{Ruang singgung manifold diferensiabel}

Untuk suatu himpunan terbuka

\mavergleichskette
{\vergleichskette
{V
}
{ \subseteq }{ \R^n
}
{  }{
}
{    }{
}
{   }{
}
}
{}{}{}
dan suatu titik

\mavergleichskette
{\vergleichskette
{P
}
{ \in }{V
}
{  }{
}
{    }{
}
{   }{
}
}
{}{}{}
ruang vektor real ambien $\R^n$ merepresentasikan himpunan semua arah yang mungkin di titik tersebut. Vektor

\mavergleichskette
{\vergleichskette
{v
}
{ \in }{ \R^n
}
{  }{
}
{    }{
}
{   }{
}
}
{}{}{}
dapat diidentifikasi dengan lintasan
\maabbeledisp {} {I} {V
} {t} { P+t v
} {,}
dengan

\mavergleichskette
{\vergleichskette
{ 0
}
{ \in }{ I
}
{  }{
}
{    }{
}
{   }{
}
}
{}{}{}
suatu interval real sedemikian sehingga citra pemetaan tersebut terletak di dalam $V$. Ruang vektor ini disebut \stichwort {ruang singgung} {} di titik $P$ pada $V$.

Untuk serat suatu pemetaan diferensiabel
\maabb {\varphi} {G} {\R^m
} {,}

\mavergleichskette
{\vergleichskette
{G
}
{ \subseteq }{ \R^n
}
{  }{
}
{    }{
}
{   }{
}
}
{}{}{}
terbuka, di suatu titik reguler dalam arti submersif

\mavergleichskette
{\vergleichskette
{P
}
{ \in }{G
}
{  }{
}
{    }{
}
{   }{
}
}
{}{}{}, dengan $m\leq n$ dan diferensial total $(D\varphi)_P$ surjektif, kita telah mendefinisikan
\definitionsverweis {ruang singgung}{}{}
serat yang melalui $P$ sebagai kernel diferensial total
\maabb {\left(D\varphi\right)_{P}} {\R^n } { \R^m
} {}.
Dengan demikian, ruang singgung merupakan subruang vektor
\mathl{(n-m)}{-}dimensional dari ruang vektor ambien $\R^n$.

Catatan edisi: Sumber hanya menyebut titik ber-rank maksimal, tetapi dimensi kernel $n-m$ memerlukan kasus submersif di atas. Jika $n<m$, rank maksimal berarti injektif dan dimensi kernel ialah nol, bukan $n-m$.

Untuk konsep manifold abstrak kita, tidak ada ruang vektor ambien semacam itu yang menjadi tempat berlangsungnya segala sesuatu. Meskipun demikian, kita juga dapat mendefinisikan suatu ruang singgung di setiap titik suatu manifold. Ruang ini akan berupa ruang vektor
\zusatzklammer {yang dimensinya sama dengan dimensi manifold tersebut} {} {,}
dan bagi suatu pemetaan diferensiabel antara dua manifold, diferensial total di setiap titik akan berupa pemetaan linear antara ruang-ruang singgungnya.

Jika untuk suatu titik

\mavergleichskette
{\vergleichskette
{P
}
{ \in }{M
}
{  }{
}
{    }{
}
{   }{
}
}
{}{}{}
dari suatu manifold diferensiabel $M$ kita memilih suatu lingkungan terbuka

\mavergleichskette
{\vergleichskette
{ P
}
{ \in }{ U
}
{ \subseteq }{ M
}
{    }{
}
{   }{
}
}
{}{}{}
dan suatu peta
\maabbdisp {\alpha} {U} {V
} {}
dengan

\mavergleichskette
{\vergleichskette
{V
}
{ \subseteq }{ \R^k
}
{  }{
}
{    }{
}
{   }{
}
}
{}{}{}
maka wajar untuk memandang $\R^k$ ini sebagai ruang singgung. Memang, akan terdapat suatu isomorfisme semacam itu, tetapi sebagai definisi pendekatan ini tidak dapat digunakan karena bergantung pada peta yang dipilih. Sebagai gantinya, kita bekerja dengan kelas-kelas ekuivalensi kurva terdiferensialkan.

\bild{ \begin{center}
\includegraphics[width=5.5cm]{ \bildeinlesungsvg {Tangentialvektor} {svg} }
\end{center}
\bildtext {} }

\bildlizenz { Tangentialvektor.svg } {} {TN} {de Wikipedia} {PD} {}

\inputdefinition
{}
{

Misalkan $M$ suatu
\definitionsverweis {manifold diferensiabel}{}{}
dan

\mavergleichskette
{\vergleichskette
{ P
}
{ \in }{ M
}
{  }{
}
{    }{
}
{   }{
}
}
{}{}{}
suatu titik. Misalkan
\maabbdisp {\gamma_1} {I_1 } { M
} {}
dan
\maabbdisp {\gamma_2} {I_2 } { M
} {}
dua kurva yang didefinisikan pada
\definitionsverweis {interval-interval terbuka}{}{}

\mavergleichskette
{\vergleichskette
{ 0
}
{ \in }{ I_1,I_2
}
{ \subseteq }{ \R
}
{    }{
}
{   }{
}
}
{}{}{}
dan
\definitionsverweis {terdiferensialkan}{}{}
, dengan

\mavergleichskette
{\vergleichskette
{ \gamma_1(0)
}
{ = }{ P
}
{ = }{ \gamma_2(0)
}
{    }{
}
{   }{
}
}
{}{}{.}
Kurva
\mathkor {} {\gamma_1} {dan} {\gamma_2} {}
disebut
\definitionswort {ekuivalen secara tangensial }{} di $P$ jika terdapat suatu lingkungan terbuka

\mavergleichskette
{\vergleichskette
{P
}
{ \in }{ U
}
{  }{
}
{    }{
}
{   }{
}
}
{}{}{}
dan suatu
\definitionsverweis {peta}{}{}
\maabbdisp {\alpha} { U } { V
} {}
dengan

\mavergleichskette
{\vergleichskette
{V
}
{ \subseteq }{ \R^n
}
{  }{
}
{    }{
}
{   }{
}
}
{}{}{}
sedemikian sehingga berlaku

\mavergleichskettedisp
{\vergleichskette
{ \left(  \alpha \circ \left(  \gamma_1  \vert_{\gamma_1^{-1} (U)}  \right)  \right)'(0)
}
{ =} { \left(  \alpha \circ \left(  \gamma_2  \vert_{\gamma_2^{-1} (U)}  \right)  \right)'(0)
}
{ } {
}
{ } {
}
{ } {
}
}
{}{}{}
.

}

Kita memerlukan beberapa lema sederhana untuk menunjukkan bahwa konsep ini bermakna.


\inputfaktbeweis
{Differenzierbare Mannigfaltigkeit/Differenzierbare Kurve/Tangential äquivalent/Beliebige offene Umgebung/Fakt}
{Lema}
{}
{

\faktsituation {Misalkan $M$ suatu
\definitionsverweis {manifold diferensiabel}{}{}
dan

\mavergleichskette
{\vergleichskette
{ P
}
{ \in }{ M
}
{  }{
}
{    }{
}
{   }{
}
}
{}{}{}
suatu titik. Misalkan
\maabbdisp {\gamma_1} {I_1 } { M
} {}
dan
\maabbdisp {\gamma_2} {I_2 } { M
} {}
dua kurva yang didefinisikan pada
\definitionsverweis {interval-interval terbuka}{}{}

\mavergleichskette
{\vergleichskette
{ 0
}
{ \in }{ I_1,I_2
}
{ \subseteq }{ \R
}
{    }{
}
{   }{
}
}
{}{}{}
dan
\definitionsverweis {terdiferensialkan}{}{}
, dengan

\mavergleichskette
{\vergleichskette
{ \gamma_1(0)
}
{ = }{ P
}
{ = }{ \gamma_2(0)
}
{    }{
}
{   }{
}
}
{}{}{.}}
\faktfolgerung {Maka
\mathkor {} {\gamma_1} {dan} {\gamma_2} {}
\definitionsverweis {ekuivalen secara tangensial}{}{}
di $P$ jika dan hanya jika untuk setiap
\definitionsverweis {peta}{}{}
\maabbdisp {\alpha} { U } { V
} {}
dengan

\mavergleichskette
{\vergleichskette
{P
}
{ \in }{ U
}
{  }{
}
{    }{
}
{   }{
}
}
{}{}{}
dan

\mavergleichskette
{\vergleichskette
{V
}
{ \subseteq }{ \R^n
}
{  }{
}
{    }{
}
{   }{
}
}
{}{}{}
berlaku kesamaan

\mavergleichskettedisp
{\vergleichskette
{ \left(  \alpha \circ \left(  \gamma_1  \vert_{\gamma_1^{-1} (U)}  \right)  \right)'(0)
}
{ =} { \left(  \alpha \circ \left(  \gamma_2  \vert_{\gamma_2^{-1} (U)}  \right)  \right)'(0)
}
{ } {
}
{ } {
}
{ } {
}
}
{}{}{}
.}
\faktzusatz {}
\faktzusatz {}

}
{

Untuk suatu
\definitionsverweis {kurva terdiferensialkan}{}{}
\maabbdisp {\gamma} { I } { M
} {}
dengan

\mavergleichskette
{\vergleichskette
{ 0
}
{ \in }{ I
}
{  }{
}
{    }{
}
{   }{
}
}
{}{}{}
dan

\mavergleichskette
{\vergleichskette
{ \gamma(0)
}
{ = }{ P
}
{  }{
}
{    }{
}
{   }{
}
}
{}{}{}
serta suatu
\definitionsverweis {peta}{}{}
\maabbdisp {\alpha} { U } { V
} {}
\zusatzklammer {dengan

\mavergleichskettek
{\vergleichskettek
{P
}
{ \in }{ U
}
{  }{
}
{    }{
}
{   }{
}
}
{}{}{}
dan

\mavergleichskette
{\vergleichskette
{V
}
{ \subseteq }{\R^n
}
{  }{
}
{    }{
}
{   }{
}
}
{}{}{}} {} {}
nilai ungkapan

\mathdisp {\left(  \alpha \circ \left(  \gamma \vert_{\gamma^{-1} (U)}  \right)  \right)'(0)} {  }

tidak berubah apabila kita beralih ke interval terbuka yang lebih kecil

\mavergleichskette
{\vergleichskette
{ 0
}
{ \in }{ I'
}
{ \subseteq }{ I
}
{    }{
}
{   }{
}
}
{}{}{}
dan himpunan terbuka yang lebih kecil

\mavergleichskette
{\vergleichskette
{P
}
{ \in }{ U'
}
{ \subseteq }{ U
}
{    }{
}
{   }{
}
}
{}{}{}
\zusatzklammer {dengan peta terinduksi} {} {.}
Kita dengan demikian dapat mengasumsikan bahwa
\mathkor {} {\gamma_1} {dan} {\gamma_2} {}
didefinisikan pada interval yang sama dan citranya terletak di $U$, serta bahwa untuk $U$ ini terdapat dua peta
\maabbdisp {\alpha_1} { U } { V_1
} {}
dan
\maabbdisp {\alpha_2} { U } { V_2
} {}.
Selanjutnya, dari

\mavergleichskettedisp
{\vergleichskette
{ ( \alpha_1  \circ \gamma_1 )'(0)
}
{ =} { ( \alpha_1  \circ \gamma_2 )'(0)
}
{ } {
}
{ } {
}
{ } {
}
}
{}{}{}
menurut
aturan rantai,
dengan menggunakan diferensiabilitas
\definitionsverweis {pemetaan transisi}{}{}

\mathl{\alpha_2 \circ \alpha_1^{-1}}{}
, langsung diperoleh

\mavergleichskettealign
{\vergleichskettealign
{ \left(  \alpha_2 \circ \gamma_1  \right)'(0)
}
{ =} { \left(  \left(  \alpha_2 \circ \alpha_1^{-1}  \right) \circ \left(  \alpha_1 \circ \gamma_1  \right)  \right)'(0)
}
{ =} { \left(  D   \left(   \alpha_2 \circ \alpha_1^{-1}   \right)      \right)_{\alpha_1(P)} \left(   \left(  \alpha_1 \circ \gamma_1  \right)'(0)   \right)
}
{ =} { \left(  D   \left(   \alpha_2 \circ \alpha_1^{-1}   \right)      \right)_{\alpha_1(P)} \left(   \left(  \alpha_1 \circ \gamma_2  \right)'(0)   \right)
}
{ =} { \left(  \alpha_2 \circ \gamma_2  \right)'(0)
}
}
{}
{}{.}

}



\inputfaktbeweis
{Differenzierbare Mannigfaltigkeit/Differenzierbare Kurve/Tangential äquivalent/Äquivalenzrelation/Fakt}
{Lema}
{}
{

\faktsituation {Misalkan $M$ suatu
\definitionsverweis {manifold diferensiabel}{}{}
dan

\mavergleichskette
{\vergleichskette
{ P
}
{ \in }{ M
}
{  }{
}
{    }{
}
{   }{
}
}
{}{}{}
suatu titik.}
\faktfolgerung {Maka
\definitionsverweis {ekuivalensi tangensial}{}{}
dari
\definitionsverweis {kurva-kurva terdiferensialkan}{}{}
yang melalui $P$ merupakan suatu
\definitionsverweis {relasi ekuivalensi}{}{.}}
\faktzusatz {}
\faktzusatz {}

}
{

Refleksivitas dan simetri relasi tersebut langsung jelas. Untuk membuktikan transitivitas, misalkan diberikan tiga
\definitionsverweis {kurva terdiferensialkan}{}{}
\maabbdisp {\gamma_1,\gamma_2, \gamma_3} { I } { M
} {}
,
dan kita dapat langsung mengasumsikan bahwa ketiganya didefinisikan pada
\definitionsverweis {interval terbuka}{}{}
yang sama

\mavergleichskette
{\vergleichskette
{ 0
}
{ \in }{ I
}
{ \subseteq }{\R
}
{    }{
}
{   }{
}
}
{}{}{}
.
Misalkan

\mavergleichskette
{\vergleichskette
{P
}
{ \in }{ U_1,U_2
}
{  }{
}
{    }{
}
{   }{
}
}
{}{}{}
himpunan-himpunan terbuka yang dapat digunakan untuk membuktikan ekuivalensi tangensial antara
\mathkor {} {\gamma_1} {dan} {\gamma_2} {}
serta antara
\mathkor {} {\gamma_2} {dan} {\gamma_3} {}
,
secara berurutan. Maka, menurut
Lema 9.2,
dengan

\mavergleichskette
{\vergleichskette
{U
}
{ = }{ U_1 \cap U_2
}
{  }{
}
{    }{
}
{   }{
}
}
{}{}{}
kita dapat membuktikan ekuivalensi tangensial antara
\mathkor {} {\gamma_1} {dan} {\gamma_3} {}
.

}

Berdasarkan lema ini, definisi berikut bermakna.

\inputdefinition
{}
{

Misalkan $M$ suatu
\definitionsverweis {manifold diferensiabel}{}{}
dan

\mavergleichskette
{\vergleichskette
{ P
}
{ \in }{ M
}
{  }{
}
{    }{
}
{   }{
}
}
{}{}{}
suatu titik. Yang dimaksud dengan \definitionswort {vektor singgung}{} di $P$ adalah suatu
\definitionsverweis {kelas ekuivalensi}{}{}
dari
\definitionsverweis {kurva-kurva terdiferensialkan}{}{}
yang
\definitionsverweis {ekuivalen secara tangensial}{}{}
dan melalui $P$. Himpunan vektor-vektor singgung ini dinotasikan dengan

\mathdisp {T_PM} {  }

.

}



\inputfaktbeweis
{Differenzierbare Mannigfaltigkeit/Differenzierbare Kurve/Tangentialklassen und R^n unter Karte/Fakt}
{Lema}
{}
{

\faktsituation {Misalkan $M$ suatu
($C^1$)-\definitionsverweis {manifold diferensiabel}{}{,}

\mavergleichskette
{\vergleichskette
{ P
}
{ \in }{ M
}
{  }{
}
{    }{
}
{   }{
}
}
{}{}{}
suatu titik,

\mavergleichskette
{\vergleichskette
{ P
}
{ \in }{U
}
{ \subseteq }{M
}
{    }{
}
{   }{
}
}
{}{}{}
\definitionsverweis {terbuka}{}{}
dan
\maabbdisp {\alpha} {U} {V
} {}
suatu
\definitionsverweis {peta}{}{.}

Catatan edisi: Sumber menulis $U=M$, yang tanpa sengaja mensyaratkan peta global. Yang diperlukan hanyalah domain peta terbuka $U\subseteq M$ yang memuat $P$.}
\faktuebergang {Maka berlaku pernyataan-pernyataan berikut.}
\faktfolgerung {\aufzaehlungzwei {\definitionsverweis {Pemetaan}{}{}
\maabbeledisp {} {T_PM } { \R^n
} { [\gamma] } { \left(  \alpha \circ \gamma \vert_{\gamma^{-1}(U)}  \right)'(0)
} {,}
merupakan suatu
\definitionsverweis {bijeksi}{}{}
yang terdefinisi dengan baik.
} {Pada
\mathl{T_PM}{},
\definitionsverweis {struktur ruang vektor}{}{}
yang didefinisikan melalui pemetaan ini tidak bergantung pada peta yang dipilih.
}}
\faktzusatz {}
\faktzusatz {}

}
{

\teilbeweis {}{}{}
{(1). Bahwa pemetaan tersebut terdefinisi dengan baik jelas berdasarkan
Lema 9.2.
\definitionsverweis {Injektivitas}{}{}
langsung mengikuti dari
Definisi 9.1.
Untuk
\definitionsverweis {surjektivitas}{}{,}
misalkan

\mavergleichskette
{\vergleichskette
{ v
}
{ \in }{ \R^n
}
{  }{
}
{    }{
}
{   }{
}
}
{}{}{.}
Kita meninjau
\definitionsverweis {kurva afin-linear}{}{}
\maabbeledisp {\theta} { \R } { \R^n
} { t } { \theta(t) =  \alpha(P) + t v
} {,}
yang
\definitionsverweis {turunannya}{}{}
di $0$ tepat sama dengan $v$. Kita membatasi kurva ini pada suatu interval

\mavergleichskette
{\vergleichskette
{ 0
}
{ \in }{ I
}
{ \subseteq }{ \R
}
{    }{
}
{   }{
}
}
{}{}{}
sedemikian sehingga

\mavergleichskette
{\vergleichskette
{ \theta(I)
}
{ \subseteq }{ V
}
{  }{
}
{    }{
}
{   }{
}
}
{}{}{}
dan meninjau
\maabbdisp {\gamma = \alpha^{-1} \circ \theta} { I } { M
} {.}
Untuk kurva ini berlaku

\mavergleichskettedisp
{\vergleichskette
{ \gamma(0)
}
{ =} { \left(  \alpha^{-1} \circ \theta  \right) (0)
}
{ =} { \alpha^{-1} (\theta(0))
}
{ =} { \alpha^{-1} (\alpha(P))
}
{ =} { P
}
}
{}{}{}
dan

\mavergleichskettedisp
{\vergleichskette
{ \left(  \alpha \circ \gamma  \right)'(0)
}
{ =} { \left(  \alpha \circ \left(  \alpha^{-1} \circ \theta  \right)  \right)'(0)
}
{ =} { \theta'(0)
}
{ =} { v
}
{ } {
}
}
{}{}{.}}
{}
\teilbeweis {}{}{}
{(2). Dengan beralih ke himpunan-himpunan terbuka yang lebih kecil, kita dapat mengasumsikan bahwa terdapat dua peta
\maabbdisp {\alpha_1} { U } { V_1
} {}
dan
\maabbdisp {\alpha_2} { U } { V_2
} {}
.
\definitionsverweis {Pemetaan transisi}{}{}
\maabbdisp {\alpha_2 \circ \alpha_1^{-1}} { V_1 } { V_2
} {}
merupakan suatu
$C^1$-\definitionsverweis {difeomorfisme}{}{}
dan menurut
aturan rantai,
untuk
\definitionsverweis {diferensial total}{}{}
pemetaan tersebut di
\mathl{\alpha_1(P)}{}
berlaku hubungan

\mavergleichskettedisp
{\vergleichskette
{ \left(  D(\alpha_2 \circ \alpha_1^{-1})  \right)_{\alpha_1(P)} \left(   (\alpha_1 \circ \gamma)'(0)   \right)
}
{ =} { (\alpha_2 \circ \gamma)'(0)
}
{ } {
}
{ } {
}
{ } {
}
}
{}{}{.}
Ini berarti bahwa diagram

\mathdisp {\begin{matrix}T_PM  & \stackrel{  }{\longrightarrow} & \R^n   & \\ &  \searrow & \downarrow & \\ & & \R^n   & \!\!\!\!\! \!\!\! , \\ \end{matrix}} {  }

komutatif, dengan diferensial total dari
\mathl{\alpha_2 \circ \alpha_1^{-1}}{}
sebagai pemetaan vertikal. Karena diferensial total tersebut merupakan suatu
\definitionsverweis {pemetaan linear}{}{}
yang bijektif dalam situasi ini, tidak ada perbedaan apakah penjumlahan dan perkalian skalar pada
\mathl{T_PM}{}
didefinisikan dengan mengacu pada pemetaan horizontal atas atau bawah.}
{}

}

\inputdefinition
{}
{

Misalkan $M$ suatu
\definitionsverweis {manifold diferensiabel}{}{}
dan

\mavergleichskette
{\vergleichskette
{ P
}
{ \in }{ M
}
{  }{
}
{    }{
}
{   }{
}
}
{}{}{}
suatu titik. Yang dimaksud dengan \definitionswort {ruang singgung}{} di $P$, ditulis
\mathl{T_PM}{,} adalah himpunan
\definitionsverweis {vektor-vektor singgung}{}{}
di $P$ yang dilengkapi dengan
\definitionsverweis {struktur ruang vektor}{}{}
real yang diberikan oleh sebarang
\definitionsverweis {peta}{}{.}

}

Dimensi ruang singgung sama dengan dimensi manifold. Setiap peta menginduksi suatu isomorfisme antara
\mathl{T_P M}{} dan $\R^n$, tetapi isomorfisme-isomorfisme ini bergantung pada peta yang dipilih. Secara khusus, tidak terdapat basis standar pada ruang singgung.

\inputdefinition
{}
{

Misalkan $M$ suatu
\definitionsverweis {manifold diferensiabel}{}{}
dan

\mavergleichskette
{\vergleichskette
{ P
}
{ \in }{ M
}
{  }{
}
{    }{
}
{   }{
}
}
{}{}{}
suatu titik. \definitionsverweis {Ruang dual}{}{}
dari
\definitionsverweis {ruang singgung}{}{}

\mathl{T_PM}{} di $P$ disebut \definitionswort {ruang kotangen}{} di $P$. Ruang ini dinotasikan dengan
\mathdisp {T^*_PM} {  }
.

}


\inputfaktbeweis
{Differenzierbare Mannigfaltigkeiten/Differenzierbare Abbildung/Tangential äquivalent/Fakt}
{Lema}
{}
{

\faktsituation {Misalkan
\mathkor {} {M} {dan} {N} {}
\definitionsverweis {manifold-manifold diferensiabel}{}{}
dan misalkan
\maabbdisp {\varphi} { M } { N
} {}
suatu
\definitionsverweis {pemetaan diferensiabel}{}{.}
Misalkan
\mathkor {} {P \in M} {dan} {Q=\varphi(P)} {}
serta misalkan
\maabbdisp {\gamma_1,\gamma_2} { I } { M
} {}
dua
\definitionsverweis {kurva terdiferensialkan}{}{}
yang didefinisikan pada suatu interval terbuka

\mavergleichskette
{\vergleichskette
{ 0
}
{ \in }{ I
}
{  }{
}
{    }{
}
{   }{
}
}
{}{}{}
dan

\mavergleichskette
{\vergleichskette
{ \gamma_1 (0)
}
{ = }{ \gamma_2(0)
}
{ = }{ P
}
{    }{
}
{   }{
}
}
{}{}{.}}
\faktvoraussetzung {Misalkan
\mathkor {} {\gamma_1} {dan} {\gamma_2} {}
\definitionsverweis {ekuivalen secara tangensial}{}{}
di titik $P$.}
\faktfolgerung {Maka
\definitionsverweis {komposisi}{}{}
\mathkor {} {\varphi \circ \gamma_1} {dan} {\varphi \circ \gamma_2} {}
juga ekuivalen secara tangensial di $Q$.}
\faktzusatz {}
\faktzusatz {}

 }
{ Lihat Soal 9.3. }

Berdasarkan lema ini, konsep berikut terdefinisi dengan baik.

\inputdefinition
{}
{

Misalkan
\mathkor {} {M} {dan} {N} {}
\definitionsverweis {manifold-manifold diferensiabel}{}{}
dan misalkan
\maabbdisp {\varphi} { M } { N
} {}
suatu
\definitionsverweis {pemetaan diferensiabel}{}{.}
Misalkan
\mathkor {} {P \in M} {dan} {Q=\varphi(P)} {.}
Pemetaan
\maabbeledisp {} {T_PM} {T_{\varphi(P)}N
} { [\gamma] } { [\varphi \circ \gamma]
} {,}
disebut \definitionswort {pemetaan singgung di titik}{} $P$ yang bersesuaian. Pemetaan ini dinotasikan dengan
\mathl{T_P(\varphi)}{}.

}



\inputfaktbeweis
{Tangentialabbildung/Punktweise/Elementare Eigenschaften/Fakt}
{Lema}
{}
{

\faktsituation {Misalkan
\mathkor {} {M} {dan} {N} {}
\definitionsverweis {manifold-manifold diferensiabel}{}{}
dan misalkan
\maabbdisp {\varphi} { M } { N
} {}
suatu
\definitionsverweis {pemetaan diferensiabel}{}{.}
Misalkan
\mathkor {} {P \in M} {,} {Q=\varphi(P)} {}
dan misalkan
\maabbdisp {T_P(\varphi)} {T_PM} {T_QN
} {}
\definitionsverweis {pemetaan singgung}{}{}
yang terkait.}
\faktuebergang {Maka berlaku pernyataan-pernyataan berikut.}
\faktfolgerung {\aufzaehlungsechs{Jika
\mathkor {} {M \subseteq \R^m} {dan} {N \subseteq \R^n} {}
merupakan
\definitionsverweis {himpunan-himpunan bagian terbuka}{}{}
dan ruang-ruang singgung diidentifikasi dengan ruang-ruang Euklides ambien, maka pemetaan singgung sama dengan
\definitionsverweis {diferensial total}{}{}

\mathl{\left(D\varphi\right)_{P}}{.}
}{Jika
\maabbdisp {\alpha} { U } { V
} {}
dengan

\mavergleichskette
{\vergleichskette
{P
}
{ \in }{ U
}
{  }{
}
{    }{
}
{   }{
}
}
{}{}{}
dan

\mavergleichskette
{\vergleichskette
{V
}
{ \subseteq }{ \R^m
}
{  }{
}
{    }{
}
{   }{
}
}
{}{}{}
serta
\maabbdisp {\beta} {U'} {W
} {}
dengan

\mavergleichskette
{\vergleichskette
{Q
}
{ \in }{ U'
}
{  }{
}
{    }{
}
{   }{
}
}
{}{}{}
dan

\mavergleichskette
{\vergleichskette
{W
}
{ \subseteq }{ \R^n
}
{  }{
}
{    }{
}
{   }{
}
}
{}{}{}
merupakan
\definitionsverweis {peta-peta}{}{,}
maka diagram

\mathdisp {\begin{matrix} T_PM & \stackrel{ T_P(\varphi) }{\longrightarrow} & T_QN &  \\ \downarrow & & \downarrow  & \\  \R^m  & \stackrel{ \left(D \left(  \beta \circ \varphi \circ \alpha^{-1}  \right)  \right)_{ \alpha(P)} }{\longrightarrow} & \R^n  & \! \! \! \!\!\!\!\!   \\ \end{matrix}} {  }

komutatif, dengan pemetaan-pemetaan vertikal diberikan oleh isomorfisme
\mathl{[\gamma] \mapsto ( \alpha \circ \gamma)'(0)}{}
dan
\mathl{[\gamma] \mapsto ( \beta \circ \gamma)'(0)}{}
,
secara berurutan.
}{
\mathl{T_P(\varphi)}{}
bersifat
$\R$-\definitionsverweis {linear}{}{.}
}{Jika $L$ suatu
\definitionsverweis {manifold}{}{}
lain,

\mavergleichskette
{\vergleichskette
{R
}
{ \in }{ L
}
{  }{
}
{    }{
}
{   }{
}
}
{}{}{}
dan
\maabbdisp {\psi} { L } { M
} {}
suatu pemetaan diferensiabel lain dengan

\mavergleichskette
{\vergleichskette
{ \psi(R)
}
{ = }{ P
}
{  }{
}
{    }{
}
{   }{
}
}
{}{}{}
,
maka berlaku

\mavergleichskettedisp
{\vergleichskette
{ T_R( \varphi \circ \psi)
}
{ =} { T_P( \varphi) \circ T_R( \psi)
}
{ } {
}
{ } {
}
{ } {
}
}
{}{}{.}
}{Jika $\varphi$ suatu
\definitionsverweis {difeomorfisme}{}{,}
maka
\mathl{T_P(\varphi)}{}
merupakan suatu
\definitionsverweis {isomorfisme}{}{.}
}{Untuk suatu
\definitionsverweis {kurva terdiferensialkan}{}{}
\maabbdisp {\gamma} { I } { M
} {}
dengan suatu interval terbuka

\mavergleichskette
{\vergleichskette
{I
}
{ \subseteq }{\R
}
{  }{
}
{    }{
}
{   }{
}
}
{}{}{,}

\mavergleichskette
{\vergleichskette
{ 0
}
{ \in }{ I
}
{  }{
}
{    }{
}
{   }{
}
}
{}{}{}
dan

\mavergleichskette
{\vergleichskette
{ \gamma(0)
}
{ = }{ P
}
{  }{
}
{    }{
}
{   }{
}
}
{}{}{}
,
di dalam ruang singgung
\mathl{T_PM}{}
berlaku kesamaan

\mavergleichskettedisp
{\vergleichskette
{ [\gamma]
}
{ =} { (T_0(\gamma))(1)
}
{ } {
}
{ } {
}
{ } {
}
}
{}{}{.}
}}
\faktzusatz {}
\faktzusatz {}

}
{

\teilbeweis {}{}{}
{(1). Setiap vektor singgung direpresentasikan oleh suatu lintasan afin-linear pada interval terbuka yang cukup kecil di sekitar nol agar citranya tetap berada di domain terbuka,
\mathl{t \mapsto \gamma(t)=P +t v}{}
dengan suatu vektor

\mavergleichskette
{\vergleichskette
{v
}
{ \in }{\R^m
}
{  }{
}
{    }{
}
{   }{
}
}
{}{}{,}
sehingga kita dapat beralih bolak-balik antara vektor-vektor ini dan vektor-vektor singgung yang didefinisikannya. Untuk lintasan komposit
\mathl{\varphi \circ \gamma}{}
menurut
aturan rantai
berlaku

\mavergleichskettedisp
{\vergleichskette
{ (T_P\varphi )( [\gamma ] )
}
{ =} { [ \varphi \circ \gamma ]
}
{ =} { (\varphi \circ \gamma)'(0)
}
{ =} { (D \varphi)_P \left( \left(  D \gamma  \right)_{ 0 } \left(   1  \right) \right)
}
{ =} { \left(  D\varphi  \right)_{ P } \left(  v  \right)
}
}
{}{}{.}}
{}
\teilbeweis {}{}{}
{(2). Aturan rantai yang diterapkan pada
\zusatzklammer {dengan $I$ dan $V$ harus diganti oleh himpunan-himpunan terbuka yang lebih kecil} {} {}

\mathdisp {I \stackrel{ \alpha \circ \gamma}{ \longrightarrow} V \stackrel{ \beta \circ \varphi \circ \alpha^{-1} }{ \longrightarrow} W} {  }

menghasilkan

\mavergleichskettedisp
{\vergleichskette
{ \left(  D   \left(   \beta \circ \varphi \circ \alpha^{-1}    \right)      \right)_{ \alpha(P)} \left(   (\alpha \circ \gamma)' (0)    \right)
}
{ =} { ( \beta \circ ( \varphi \circ \gamma) )'(0)
}
{ } {
}
{ } {
}
{ } {
}
}
{}{}{,}
yang tepat menyatakan komutativitas diagram tersebut.}
{}
\teilbeweis {}{}{}
{(3). Pernyataan ini mengikuti dari (2) dan linearitas
\definitionsverweis {diferensial total}{}{.}}
{}
\teilbeweis {}{}{}
{(4). Dengan beralih ke peta-peta, pernyataan ini mengikuti dari (2) dan aturan rantai.}
{}
\teilbeweis {}{}{}
{(5) mengikuti dari (4) yang diterapkan pada pemetaan invers
\mathl{\varphi^{-1}}{.}}
{}
\teilbeweis {}{}{}
{(6). Elemen

\mavergleichskette
{\vergleichskette
{ 1
}
{ \in }{\R
}
{  }{
}
{    }{
}
{   }{
}
}
{}{}{}
sebagai vektor singgung di suatu titik

\mavergleichskette
{\vergleichskette
{a
}
{ \in }{ I
}
{  }{
}
{    }{
}
{   }{
}
}
{}{}{}
harus ditafsirkan sebagai lintasan
\mathl{s \mapsto a+s}{}
pada interval terbuka yang cukup kecil di sekitar nol sehingga $a+s\in I$.
Untuk

\mavergleichskette
{\vergleichskette
{a
}
{ = }{ 0
}
{  }{
}
{    }{
}
{   }{
}
}
{}{}{}
lintasan ini adalah lintasan identitas. Oleh karena itu,

\mavergleichskettedisp
{\vergleichskette
{ (T_0(\gamma))(1)
}
{ =} { (T_0(\gamma))(\operatorname{Id})
}
{ =} { [\gamma \circ \operatorname{Id}]
}
{ =} { [\gamma ]
}
{ } {
}
}
{}{}{.}}
{}

}

\inputdefinition
{}
{

Misalkan
\mathkor {} {M} {dan} {N} {}
\definitionsverweis {manifold-manifold diferensiabel}{}{}
dan
\maabbdisp {\varphi} { M } { N
} {}
suatu
\definitionsverweis {pemetaan diferensiabel}{}{.}
Misalkan
\mathkor {} {P \in M} {dan} {Q=\varphi(P)} {.}
Bagi
\definitionsverweis {pemetaan singgung}{}{}
\maabbdisp {T_P(\varphi)} { T_PM } { T_QN
} {}
,
\definitionsverweis {pemetaan dual}{}{}
\maabbeledisp {} { T^*_QN } { T^*_PM
} { h } { h \circ  T_P(\varphi)
} {,}
disebut \definitionswort {pemetaan kotangen}{} di titik $P$. Pemetaan ini dinotasikan dengan
\mathl{T^*_P(\varphi)}{}.

}

Jika dituliskan secara eksplisit, pemetaan tersebut adalah
\maabbeledisp {} {T^*_QN} {T^*_PM
} {h} {( [\gamma] \mapsto h([\varphi \circ \gamma]))
} {.}

\inputdefinition
{}
{

Misalkan
\mathkor {} {L} {dan} {M} {}
\definitionsverweis {manifold-manifold diferensiabel}{}{}
dan misalkan
\maabbdisp {\varphi} { L } { M
} {}
suatu
\definitionsverweis {pemetaan diferensiabel}{}{.}
Pemetaan $\varphi$ di titik

\mavergleichskette
{\vergleichskette
{ Q
}
{ \in }{ L
}
{  }{
}
{    }{
}
{   }{
}
}
{}{}{}
disebut
\definitionswort {reguler}{}
\zusatzklammer {dan $Q$ disebut \definitionswort {titik reguler}{} bagi $\varphi$} {} {,}
jika
\definitionsverweis {pemetaan singgung}{}{}
\maabbdisp {T_Q(\varphi)} { T_QL } { T_{\varphi(Q)}M
} {}
di titik $Q$ mempunyai
\definitionsverweis {rank maksimal}{}{.}

}

Definisi ini memperumum
definisi Euklides tentang titik reguler
yang bersesuaian dari himpunan-himpunan bagian Euklides ke manifold. Artinya, jika

\mavergleichskette
{\vergleichskette
{ \operatorname{dim} (L)
}
{ \geq }{  \operatorname{dim} (M)
}
{  }{
}
{    }{
}
{   }{
}
}
{}{}{}
maka pemetaan singgung di $Q$ harus surjektif, sedangkan jika

\mavergleichskette
{\vergleichskette
{ \operatorname{dim} (L)
}
{ \leq }{  \operatorname{dim} (M)
}
{  }{
}
{    }{
}
{   }{
}
}
{}{}{}
pemetaan tersebut harus injektif.

\chapter{Lembar Kerja 9}
\setcounter{section}{9}

  \zwischenueberschrift{Soal pemanasan}

\inputaufgabe
{}
{

Tunjukkan bahwa setiap
\definitionsverweis {vektor singgung}{}{}

\mavergleichskette
{\vergleichskette
{ v
}
{ \in }{ T_PS^2
}
{  }{
}
{    }{
}
{   }{
}
}
{}{}{}
pada suatu titik $P$ di
\definitionsverweis {permukaan bola satuan}{}{}
dapat direalisasikan oleh sebuah lintasan
\anfuehrung{berkecepatan tetap}{}
yang
\definitionsverweis {diferensiabel}{}{}
pada sebuah
\definitionsverweis {lingkaran besar}{}{.}

}
{} {}

\inputaufgabe
{}
{

Berikan realisasi sesederhana mungkin bagi
\definitionsverweis {vektor-vektor singgung}{}{}
pada suatu titik
\mathl{P=(s,t)}{} di permukaan silinder
\mathl{S^1 \times \R}{.}

}
{} {}

\inputaufgabe
{}
{

Misalkan
\mathkor {} {M} {dan} {N} {}
adalah
\definitionsverweis {manifold diferensiabel}{}{}
dan
\maabbdisp {\varphi} { M } { N
} {}
sebuah
\definitionsverweis {pemetaan diferensiabel}{}{.}
Misalkan $P \in M$ dan
\mathl{Q=\varphi(P)}{}, serta misalkan
\maabbdisp {\gamma_1,\gamma_2} { I } { M
} {}
dua
\definitionsverweis {kurva terdiferensialkan}{}{}
dengan interval terbuka
\mathl{0 \in I}{} dan
\mathl{\gamma_1 (0) = \gamma_2(0)=P}{.} Misalkan
\mathkor {} {\gamma_1} {dan} {\gamma_2} {}
\definitionsverweis {ekuivalen secara tangensial}{}{}
di titik $P$.
Tunjukkan bahwa
\definitionsverweis {komposisi}{}{}
\mathkor {} {\varphi \circ \gamma_1} {dan} {\varphi \circ \gamma_2} {}
juga ekuivalen secara tangensial di $Q$.

}
{} {}

\inputaufgabe
{}
{

Berikan sebuah contoh pemetaan yang
\definitionsverweis {injektif}{}{,}
tidak
\definitionsverweis {surjektif}{}{,}
dan
\definitionsverweis {diferensiabel}{}{},
\maabbdisp {\varphi} {\R} {\R
} {}
sedemikian sehingga
\definitionsverweis {pemetaan singgung}{}{}
yang bersesuaian,
\maabbdisp {T_P(\varphi)} {T_P \R} {T_{\varphi(P)} \R
} {}
bijektif pada setiap titik

\mavergleichskette
{\vergleichskette
{ P
}
{ \in }{ \R
}
{  }{
}
{    }{
}
{   }{
}
}
{}{}{.}

}
{} {}

\inputaufgabe
{}
{

Berikan sebuah contoh pemetaan yang
\definitionsverweis {surjektif}{}{,}
tidak
\definitionsverweis {injektif}{}{,}
dan
\definitionsverweis {diferensiabel}{}{},
\maabbdisp {\varphi} {S^1} {S^1
} {}
sedemikian sehingga
\definitionsverweis {pemetaan singgung}{}{}
yang bersesuaian,
\maabbdisp {T_P(\varphi)} {T_P S^1} {T_{\varphi(P)} S^1
} {}
bijektif pada setiap titik
\mathl{P \in S^1}{.}

}
{} {}

\inputaufgabe
{}
{

Tunjukkan bahwa pada sebuah
\definitionsverweis {manifold diferensiabel}{}{}, penjumlahan lintasan
\maabbdisp {\gamma_1, \gamma_2} { I } { M
} {}
dengan
\mathl{\gamma_1(0) =\gamma_2(0)=P \in M}{,} yang dapat diperoleh melalui sebuah peta
\maabbdisp {\alpha} { U } { V
} {}
dengan
\mathl{\alpha(P)=0}{} dari penjumlahan di $\R^n$, pada umumnya bergantung pada peta yang dipilih.

}
{} {}

\inputaufgabe
{}
{

Tunjukkan bahwa pemetaan singgung $T_P(\varphi)$ yang bersesuaian dengan
\maabbeledisp {\varphi} {\R^1} {S^1
} { t } {( \cos  t, \sin  t  )
} {,}
bijektif pada setiap titik
\mathl{P \in \R}{.}

}
{} {}

\inputaufgabe
{}
{

Berikan sebuah contoh pemetaan yang
\definitionsverweis {surjektif}{}{}
dan
\definitionsverweis {diferensiabel}{}{}
\maabbdisp {\varphi} { M } { N
} {}
antara dua
\definitionsverweis {manifold diferensiabel}{}{},
\mathkor {} {M} {dan} {N} {}
sedemikian sehingga
\definitionsverweis {pemetaan singgung}{}{}
yang bersesuaian,
\maabbdisp {T_P(\varphi)} {T_PM} {T_{\varphi(P)}N
} {}
tidak surjektif pada suatu titik
\mathl{P \in M}{.}

}
{} {}

\inputaufgabe
{}
{

Berikan sebuah contoh pemetaan yang
\definitionsverweis {injektif}{}{}
dan
\definitionsverweis {diferensiabel}{}{}
\maabbdisp {\varphi} { M } { N
} {}
antara dua
\definitionsverweis {manifold diferensiabel}{}{},
\mathkor {} {M} {dan} {N} {}
sedemikian sehingga
\definitionsverweis {pemetaan singgung}{}{}
yang bersesuaian,
\maabbdisp {T_P(\varphi)} {T_PM} {T_{\varphi(P)}N
} {}
tidak injektif pada suatu titik
\mathl{P \in M}{.}

}
{} {}

\inputaufgabe
{}
{

Misalkan $M$ sebuah
\definitionsverweis {manifold diferensiabel}{}{}
dan

\mavergleichskette
{\vergleichskette
{ P
}
{ \in }{ M
}
{  }{
}
{    }{
}
{   }{
}
}
{}{}{}
sebuah titik. Perhatikan himpunan berikut.

\mavergleichskettedisp
{\vergleichskette
{ T
}
{ =} { \left\{  (U,f)   \mid   U \subseteq M \text{ terbuka} , \,  P \in U   , \,  f \in C^1(U,\R)    \right\}
}
{ } {
}
{ } {
}
{ } {
}
}
{}{}{.}
Perhatikan
\definitionsverweis {relasi}{}{}

\mathdisp {(U,f) \sim (V,g) : \text{ terdapat himpunan terbuka } W \text{ dengan } P \in W \subseteq U \cap V \text{ dan } f  \vert_W = g \vert_W} { . }

\aufzaehlungzwei {Tunjukkan bahwa relasi ini merupakan sebuah
\definitionsverweis {relasi ekuivalensi}{}{}
pada $T$.
} {Tunjukkan bahwa terdapat suatu
\definitionsverweis {struktur gelanggang}{}{}
alami pada himpunan
\definitionsverweis {kelas-kelas ekuivalensi}{}{}
dari relasi ekuivalensi tersebut.
}

}
{} {}

\inputaufgabe
{}
{

Misalkan $X$ sebuah
\definitionsverweis {ruang topologis}{}{,}
$Y$ sebuah himpunan, dan
\maabbdisp {\varphi} { X } { Y
} {}
sebuah
\definitionsverweis {pemetaan}{}{.}
Tunjukkan bahwa
\definitionsverweis {sistem himpunan}{}{}

\mathdisp {{\mathcal T } = \left\{ V \subseteq Y  \mid  \varphi^{-1}(V) \text{ terbuka di } X        \right\}} {  }

mendefinisikan sebuah
\definitionsverweis {topologi}{}{}
pada $Y$ sehingga $\varphi$
\definitionsverweis {kontinu}{}{.}

}
{} {}

Topologi yang diperkenalkan dalam soal sebelumnya disebut \stichwort {topologi citra} {.}

\inputaufgabe
{}
{

Tunjukkan bahwa pada $\R^n$, hubungan

\mathdisp {P \sim Q, \text{ jika } P-Q \in \Z^n} {  }

mendefinisikan sebuah
\definitionsverweis {relasi ekuivalensi}{}{.}
Lengkapi
\definitionsverweis {himpunan hasil bagi}{}{}

\mavergleichskettedisp
{\vergleichskette
{Y
}
{ =} {\R^n /\!\!\sim
}
{ =} {\R^n/\Z^n
}
{ } {
}
{ } {
}
}
{}{}{}
dengan
\definitionsverweis {topologi citra}{}{}
yang diinduksi oleh
\definitionsverweis {pemetaan hasil bagi}{}{}
\maabb {\varphi} {\R^n } { Y
} {}
tersebut. Tunjukkan bahwa $Y$ merupakan sebuah
\definitionsverweis {ruang Hausdorff}{}{.}

}
{} {}

\inputaufgabe
{}
{

Uraikan garis besar suatu teori tentang \anfuehrung{manifold kompleks}{.} Apa manifold riil yang mendasarinya, berapakah
\zusatzklammer {kompleks/riil} {} {}
dimensinya, dan seperti apakah ruang singgungnya?

}
{} {}

  \zwischenueberschrift{Soal untuk dikumpulkan}

\inputaufgabe
{4}
{

Misalkan $M$ sebuah
\definitionsverweis {manifold diferensiabel}{}{} dan
\mathl{P \in M}{.} Tunjukkan bahwa untuk sebuah
\definitionsverweis {kurva terdiferensialkan}{}{}
\maabbdisp {\gamma} { I } { M
} {}
dengan
\mathl{\gamma(0)= P}{} dan
\mathl{a \in \R}{}, di dalam
\definitionsverweis {ruang singgung}{}{}
\mathl{T_PM}{} berlaku hubungan

\mavergleichskettedisp
{\vergleichskette
{a [\gamma]
}
{ =} { [  \lambda]
}
{ } {
}
{ } {
}
{ } {
}
}
{}{}{}. Definisikan interval terbuka
\mathl{J:=\{t\in\R\mid at\in I\}}{}
dan pemetaan $\lambda:J\to M$ oleh
\mathl{\lambda(t):= \gamma(at)}{}; tunjukkan hubungan yang sama untuk kelas kurva ini.

Catatan edisi: Sumber mendefinisikan $\lambda$ pada interval semula $I$, padahal $at$ tidak harus berada di $I$. Domain $J$ di atas adalah prapeta $I$ oleh perkalian dengan $a$, selalu merupakan lingkungan terbuka dari nol, dan membuat reparametrisasi terdefinisi dengan benar.

}
{} {}

\inputaufgabe
{6}
{

Misalkan $M$ sebuah
$C^k$-\definitionsverweis {manifold}{}{}
dan

\mavergleichskette
{\vergleichskette
{ P
}
{ \in }{ M
}
{  }{
}
{    }{
}
{   }{
}
}
{}{}{.}
Untuk
$C^k$-\definitionsverweis {kurva}{}{}
\maabbdisp {\gamma_1,\gamma_2} { I } { M
} {}
dengan

\mavergleichskette
{\vergleichskette
{ \gamma_1(0)
}
{ = }{ \gamma_2(0)
}
{ = }{ P
}
{    }{
}
{   }{
}
}
{}{}{}, definisikan sebuah relasi ekuivalensi yang, dalam
\zusatzklammer {setiap} {} {}
\definitionsverweis {peta}{}{,}
memperhitungkan
\definitionsverweis {turunan}{}{}
hingga orde $k$. Seperti apakah wakil-wakil sederhana dari relasi ekuivalensi ini? Untuk $k=1$, pulihkan struktur ruang vektor pada
\definitionsverweis {himpunan hasil bagi}{}{}
dan tentukan dimensinya. Untuk $k\geq 2$, periksa apakah konstruksi ruang vektor itu tidak bergantung pada peta. Setelah sebuah peta dipilih, identifikasikan himpunan hasil bagi dengan $(\R^n)^k$ dan tentukan dimensinya sebagai manifold; lalu tunjukkan melalui suatu pergantian peta nonlinier bahwa struktur ruang vektor tersebut pada umumnya tidak kanonis.

Catatan edisi: Perintah terakhir dalam sumber meminta struktur ruang vektor untuk semua $k$. Untuk $k\geq2$, turunan berorde lebih tinggi berubah secara nonlinier di bawah pergantian peta, sehingga ruang jet kurva ini berdimensi $kn$ sebagai manifold tetapi tidak mempunyai struktur ruang vektor kanonis tanpa pilihan tambahan seperti peta atau koneksi.

}
{} {}

\inputaufgabe
{5}
{

Misalkan $M$ sebuah
\definitionsverweis {manifold diferensiabel}{}{}
dan
\mathl{P \in M}{.} Kita mengatakan bahwa dua
\definitionsverweis {kurva}{}{}
\maabbdisp {\gamma_1, \gamma_2} { I } { M
} {}
dengan
\mathl{\gamma_1(0)= \gamma_2(0)= P}{} mendefinisikan \stichwort {germ kurva} {} yang sama jika terdapat
\mathl{\epsilon >0}{} sedemikian sehingga

\mavergleichskettedisp
{\vergleichskette
{ \gamma_1 \!\mid_{ [-\epsilon, \epsilon]}
}
{ =} { \gamma_2 \!\mid_{ [-\epsilon, \epsilon]}
}
{ } {
}
{ } {
}
{ } {
}
}
{}{}{.}

a) Tunjukkan bahwa ini mendefinisikan sebuah
\definitionsverweis {relasi ekuivalensi}{}{}
pada himpunan semua kurva
\maabb {\gamma} { I } { M
} {}
dengan $\gamma(0)= P$
\zusatzklammer {dan dengan interval terbuka berbeda-beda
\mathl{0 \in I}{}} {} {.}

b) Tunjukkan bahwa
\definitionsverweis {kurva-kurva terdiferensialkan}{}{,}
yang merepresentasikan germ kurva yang sama, juga merepresentasikan vektor singgung yang sama.

}
{} {}

\inputaufgabe
{8}
{

Misalkan
\definitionsverweis {ruang hasil bagi}{}{}

\mavergleichskette
{\vergleichskette
{ Y
}
{ = }{ \R^n/\Z^n
}
{  }{
}
{    }{
}
{   }{
}
}
{}{}{}
dilengkapi dengan
\definitionsverweis {topologi citra}{}{.}
Definisikan sebuah struktur manifold pada $Y$ dengan peta-peta yang sesuai. Tunjukkan bahwa
\definitionsverweis {pemetaan hasil bagi}{}{}
\maabbdisp {\varphi} {\R^n } { Y
} {}
merupakan
\definitionsverweis {pemetaan diferensiabel}{}{,}
dan bahwa
\definitionsverweis {pemetaan singgung}{}{}
merupakan
\definitionsverweis {isomorfisme}{}{}
di setiap titik.

}
{} {}

\part{Unit 10}
\chapter[Kuliah 10]{Kuliah 10: Submanifold Tertutup}
\setcounter{section}{10}

  \zwischenueberschrift{Submanifold tertutup}

\inputdefinition
{}
{

Misalkan $N$ suatu
\definitionsverweis {manifold diferensiabel}{}{}
dengan
\definitionsverweis {dimensi}{}{}
$n$ dan

\mavergleichskette
{\vergleichskette
{ M
}
{ \subseteq }{ N
}
{  }{
}
{    }{
}
{   }{
}
}
{}{}{}
suatu
\definitionsverweis {himpunan bagian tertutup}{}{.}
Maka $M$ disebut \definitionswort {submanifold tertutup}{} berdimensi $m$ dari $N$ jika untuk setiap titik

\mavergleichskette
{\vergleichskette
{ P
}
{ \in }{ M
}
{  }{
}
{    }{
}
{   }{
}
}
{}{}{}
terdapat suatu
\definitionsverweis {peta}{}{}\zusatzfussnote {Yang dimaksud dengan peta di sini adalah setiap peta yang kompatibel dengan atlas yang telah diberikan; peta itu sendiri tidak harus termasuk dalam atlas tersebut} {.} {}
\maabbdisp {\theta} { W } { W'
} {}
dengan

\mavergleichskette
{\vergleichskette
{ P
}
{ \in }{ W
}
{ \subseteq }{ N
}
{    }{
}
{   }{
}
}
{}{}{}
\definitionsverweis {terbuka}{}{,}

\mavergleichskette
{\vergleichskette
{ W'
}
{ \subseteq }{ \R^n
}
{  }{
}
{    }{
}
{   }{
}
}
{}{}{}
terbuka, dan

\mavergleichskettedisp
{\vergleichskette
{ M \cap W
}
{ =} { \theta^{-1} (( \R^m  \times\{0\}) \cap W' )
}
{ } {
}
{ } {
}
{ } {
}
}
{}{}{.}

}

Ini tepat merupakan sifat yang dimiliki serat suatu pemetaan diferensiabel antara ruang-ruang Euklides di suatu titik reguler berdasarkan
teorema pemetaan implisit.
Dengan kata lain, serat-serat semacam itu merupakan submanifold tertutup dari

\mavergleichskette
{\vergleichskette
{ N
}
{ = }{ \R^n
}
{  }{
}
{    }{
}
{   }{
}
}
{}{}{.}



\inputfaktbeweis
{Abgeschlossene Untermannigfaltigkeit/Ist Mannigfaltigkeit/Fakt}
{Teorema}
{}
{

\faktsituation {Misalkan $N$ suatu
\definitionsverweis {manifold diferensiabel}{}{} dengan
\definitionsverweis {dimensi}{}{}
$n$ dan

\mavergleichskette
{\vergleichskette
{ M
}
{ \subseteq }{ N
}
{  }{
}
{    }{
}
{   }{
}
}
{}{}{}
suatu
\definitionsverweis {submanifold tertutup}{}{}
berdimensi $m$ dari $N$.}
\faktfolgerung {Maka $M$ merupakan suatu
\definitionsverweis {manifold diferensiabel}{}{}
sedemikian sehingga inklusi
 \maabb {\iota} { M } { N
} {}
merupakan suatu pemetaan diferensiabel.}
\faktzusatz {}
\faktzusatz {}

}
{

Struktur diferensiabel pada $M$ diberikan dengan mengidentifikasi sumbu
$\R^m\times\{0\}$ dengan $\R^m$, melalui peta-peta
\maabbdisp {\beta=\operatorname{pr}_{\R^m}\circ\theta\!\mid_{M\cap W}} {M \cap W} {\operatorname{pr}_{\R^m}\bigl((\R^m \times \{0\}) \cap W'\bigr)
} {}
.
Bahwa sifat difeomorfisme dari pemetaan-pemetaan transisi diwariskan kepada pembatasannya diperoleh seperti dalam bukti
Teorema 8.2.
Bahwa pemetaan tersebut diferensiabel diperoleh dari fakta bahwa bagi suatu domain peta terbuka

\mavergleichskette
{\vergleichskette
{ W
}
{ \subseteq }{ N
}
{  }{
}
{    }{
}
{   }{
}
}
{}{}{}
terdapat diagram komutatif

\mathdisp {\begin{matrix} M\cap W & \stackrel{  }{\longrightarrow} &  W  &  \\ \downarrow & & \downarrow  & \\  M  & \stackrel{  }{\longrightarrow} &  N & \! \! \! \!\!\!\!\!   \\ \end{matrix}} {  }

dengan panah-panah vertikal merepresentasikan penyematan terbuka dan panah-panah horizontal merepresentasikan penyematan tertutup. Panah horizontal adalah pembatasan inklusi $\iota$ dan inklusi itu sendiri. Melalui peta $\beta$ pada sumber dan $\theta$ pada sasaran, panah atas bersesuaian dengan
\maabbdisp {j} {\operatorname{pr}_{\R^m}\bigl((\R^m \times \{0\}) \cap W'\bigr)} {W'
} {,}
dengan $j(x)=(x,0)$, yaitu penyematan tertutup suatu subruang koordinat, yang tentu saja diferensiabel.

}



\inputfaktbeweis
{Abgeschlossene Untermannigfaltigkeit/Punktweise/Tangentialraum als Unterraum/Fakt}
{Teorema}
{}
{

\faktsituation {Misalkan $N$ suatu
\definitionsverweis {manifold diferensiabel}{}{}
dengan
\definitionsverweis {dimensi}{}{}
$n$ dan misalkan

\mavergleichskette
{\vergleichskette
{ M
}
{ \subseteq }{ N
}
{  }{
}
{    }{
}
{   }{
}
}
{}{}{}
suatu
\definitionsverweis {submanifold tertutup}{}{}
berdimensi $m$, dengan inklusi
\maabb {\iota} {M} {N
} {.}}
\faktfolgerung {Maka untuk setiap titik

\mavergleichskette
{\vergleichskette
{ P
}
{ \in }{ M
}
{  }{
}
{    }{
}
{   }{
}
}
{}{}{}
\definitionsverweis {pemetaan singgung}{}{}
\maabbdisp {T_P\iota} { T_PM } { T_PN
} {}
bersifat
\definitionsverweis {injektif}{}{.}}
\faktzusatz {Dengan kata lain, ruang singgung
\mathl{T_PM}{} merupakan
\definitionsverweis {subruang vektor}{}{}
berdimensi $m$ dari
\mathl{T_PN}{.}}
\faktzusatz {}

}
{

Misalkan

\mavergleichskette
{\vergleichskette
{ P
}
{ \in }{ M
}
{  }{
}
{    }{
}
{   }{
}
}
{}{}{}
dan

\mavergleichskette
{\vergleichskette
{ P
}
{ \in }{ V
}
{ \subseteq }{ N
}
{    }{
}
{   }{
}
}
{}{}{}
suatu domain peta dengan peta
\maabbdisp {\alpha} { V } { V'
} {}
serta dengan peta submanifold
\maabbdisp {\beta=\operatorname{pr}_{\R^m}\circ\alpha\!\mid_U} {U = M \cap  V } { U' = \operatorname{pr}_{\R^m}\bigl((\R^m\times\{0\})\cap V'\bigr)
} {.}
Menurut
Lema 9.10  (2),
kita mempunyai diagram komutatif

\mathdisp {\begin{matrix} T_PM & \stackrel{ T_P\iota }{\longrightarrow} & T_PN &  \\ \downarrow & & \downarrow  & \\  \R^m  & \stackrel{ j }{\longrightarrow} & \R^n  & \! \! \! \!\!\!\!\! .  \\ \end{matrix}} {  }

Pemetaan horizontal bawah adalah diferensial total di $\beta(P)$ dari representasi koordinat inklusi,
\maabbdisp {\alpha\circ\iota\circ\beta^{-1}} {U'} {V'
} {,}
dan sama dengan inklusi linear

\mavergleichskette
{\vergleichskette
{ \R^m
}
{ \subseteq }{\R^n
}
{  }{
}
{    }{
}
{   }{
}
}
{}{}{,}
sehingga injektif. Karena pemetaan-pemetaan vertikal bersifat bijektif, pemetaan horizontal atas juga injektif.

}

Dari pernyataan terakhir juga diperoleh bahwa di suatu titik reguler $P$ dari serat $M$ suatu pemetaan diferensiabel
\maabb {\varphi} {G} {\R^k
} {,}

\mavergleichskette
{\vergleichskette
{G
}
{ \subseteq }{ \R^n
}
{  }{
}
{    }{
}
{   }{
}
}
{}{}{}
terbuka, ruang yang didefinisikan sebagai kernel diferensial total
\zusatzklammer {sebagai subruang vektor dari

\mavergleichskettek
{\vergleichskettek
{ \R^n
}
{ = }{ T_P \R^n
}
{  }{
}
{    }{
}
{   }{
}
}
{}{}{}} {} {}
dan disebut
\definitionsverweis {ruang singgung}{}{}
berimpit dengan
\definitionsverweis {ruang singgung}{}{}
serat tersebut sebagai manifold diferensiabel. Berdasarkan
Teorema 10.3,
ruang singgung
\zusatzklammer {abstrak} {} {}
\mathl{T_PM}{} merupakan subruang vektor dari

\mavergleichskette
{\vergleichskette
{ T_P\R^n
}
{ = }{ \R^n
}
{  }{
}
{    }{
}
{   }{
}
}
{}{}{}
berdimensi
\mathl{n-k}{.}
Kernel diferensial total yang surjektif
\maabb {\left(D\varphi\right)_{P}} { \R^n } { \R^k
} {}
juga merupakan subruang vektor
\mathl{(n-k)}{-}dimensional dari $\R^n$. Kesamaan kedua subruang vektor tersebut diperoleh dari fakta bahwa kurva-kurva terdiferensialkan yang mendefinisikan ruang singgung abstrak
\maabbdisp {\gamma} {I} {M
} {}
menjadi konstan setelah dikomposisikan dengan $\varphi$; lihat
Soal 10.12.

  \zwischenueberschrift{Bundel tangen}

Untuk setiap titik

\mavergleichskette
{\vergleichskette
{ P
}
{ \in }{ M
}
{  }{
}
{    }{
}
{   }{
}
}
{}{}{}
pada suatu manifold terdapat
\definitionsverweis {ruang singgung}{}{}

\mathl{T_PM}{.}
Ruang singgung merupakan ruang vektor berdimensi $n$, dengan $n$ dimensi manifold tersebut. Elemen-elemennya adalah vektor singgung, yakni \anfuehrung{arah infinitesimal}{} di titik tersebut. Arah-arah singgung di dua titik yang berbeda pada mulanya sama sekali tidak saling berkaitan, sebab definisi tepatnya masing-masing hanya bergantung pada lingkungan terbuka titik-titik itu yang sekecil apa pun, dan lingkungan tersebut dapat dipilih saling lepas berdasarkan
\definitionsverweis {sifat Hausdorff}{}{.}

Hal ini sangat bertentangan dengan gambaran yang diasosiasikan dengan suatu himpunan terbuka

\mavergleichskette
{\vergleichskette
{ V
}
{ \subseteq }{ \R^n
}
{  }{
}
{    }{
}
{   }{
}
}
{}{}{}
.
Di sana, untuk setiap titik

\mavergleichskette
{\vergleichskette
{ Q
}
{ \in }{ V
}
{  }{
}
{    }{
}
{   }{
}
}
{}{}{}
ruang singgung
\mathl{T_QV}{}
dapat diidentifikasi secara alami dengan ruang vektor ambien $\R^n$ dengan memetakan suatu vektor

\mavergleichskette
{\vergleichskette
{ v
}
{ \in }{ \R^n
}
{  }{
}
{    }{
}
{   }{
}
}
{}{}{}
ke vektor singgung yang didefinisikan oleh kurva linear
\mathl{t \mapsto Q+tv}{}
.
Karena identifikasi ini berlaku untuk setiap titik, terdapat kesejajaran langsung antara ruang-ruang singgung bagi

\mavergleichskette
{\vergleichskette
{ Q
}
{ \in }{ V
}
{ \subseteq }{ \R^n
}
{    }{
}
{   }{
}
}
{}{}{}
.

Karena suatu manifold ditutupi oleh himpunan-himpunan terbuka yang difeomorfik dengan himpunan-himpunan terbuka dalam ruang Euklides, wajar diduga bahwa berbagai ruang singgung tersebut tidak sepenuhnya terisolasi satu sama lain. Konsep \stichwort {bundel tangen} {} menyatukan semua ruang singgung dan memungkinkan keterhubungan lokal di antara ruang-ruang singgung itu dicerminkan.

\bild{ \begin{center}
\includegraphics[width=5.5cm]{ \bildeinlesungsvg {Tangent_bundle} {svg} }
\end{center}
\bildtext {Dua visualisasi bundel tangen sebuah lingkaran. Di bagian atas, ruang singgung lingkaran pada setiap titik $P$ ditempatkan \anfuehrung{secara tangensial}{} dan direalisasikan sebagai subruang afin berdimensi satu di dalam $\R^2$ ambien. Penyematan ini menimbulkan perpotongan yang sebenarnya tidak ada dalam bundel tangen, karena titik dasar $P$ juga harus diperhitungkan. Di bagian bawah, ruang-ruang singgung pada setiap titik lingkaran disusun secara paralel sehingga menghasilkan sebuah silinder.} }

\bildlizenz { Tangent bundle.svg } {} {Oleg Alexandrov} {Commons} {PD} {}

\inputdefinition
{}
{

Misalkan $M$ suatu
\definitionsverweis {manifold diferensiabel}{}{.}
Himpunan

\mavergleichskettedisp
{\vergleichskette
{TM
}
{ =} {\biguplus_{P \in M} T_PM
}
{ } {
}
{ } {
}
{ } {
}
}
{}{}{,}
yang dilengkapi dengan pemetaan proyeksi
\maabbeledisp {\pi} {TM} {M
} {(P,v)} {P
} {,}
disebut \definitionswort {bundel tangen}{} dari $M$.

}

Jadi, suatu titik

\mavergleichskette
{\vergleichskette
{ u
}
{ \in }{ TM
}
{  }{
}
{    }{
}
{   }{
}
}
{}{}{}
dalam bundel tangen selalu mempunyai suatu \stichwort {titik dasar} {}

\mavergleichskette
{\vergleichskette
{ P
}
{ \in }{ M
}
{  }{
}
{    }{
}
{   }{
}
}
{}{}{}
dan merupakan elemen ruang singgung
\mathl{T_PM}{.}
Titik semacam itu biasanya ditulis sebagai
\mathl{(P,v)}{} dengan

\mavergleichskette
{\vergleichskette
{ P
}
{ \in }{ M
}
{  }{
}
{    }{
}
{   }{
}
}
{}{}{}
dan

\mavergleichskette
{\vergleichskette
{ v
}
{ \in }{ T_PM
}
{  }{
}
{    }{
}
{   }{
}
}
{}{}{.}
Untuk suatu himpunan terbuka

\mavergleichskette
{\vergleichskette
{ V
}
{ \subseteq }{ \R^n
}
{  }{
}
{    }{
}
{   }{
}
}
{}{}{}
berlaku

\mavergleichskette
{\vergleichskette
{ TV
}
{ = }{ V \times \R^n
}
{  }{
}
{    }{
}
{   }{
}
}
{}{}{,}
sehingga merupakan ruang produk. Hal ini secara umum tidak berlaku bagi sebarang manifold. Pada mulanya, bundel tangen hanya menggabungkan berbagai ruang singgung secara saling lepas tanpa mengidentifikasi ruang-ruang singgung yang berbeda; namun, topologi yang sebentar lagi akan kita perkenalkan pada bundel tangen menghasilkan suatu \anfuehrung{struktur ketetanggaan}{} tambahan di antara ruang-ruang singgung tersebut.

\inputdefinition
{}
{

Misalkan
\mathkor {} {M} {dan} {N} {}
\definitionsverweis {manifold diferensiabel}{}{}
dan
\maabbdisp {\varphi} { M } { N
} {}
suatu
\definitionsverweis {pemetaan diferensiabel}{}{.}
Misalkan
\mathkor {} {TM} {dan} {TN} {}
\definitionsverweis {bundel-bundel tangen}{}{}
yang bersesuaian. Yang dimaksud dengan \definitionswort {pemetaan singgung}{}
\maabbdisp {T (\varphi)} {TM} {TN
} {}
adalah gabungan saling lepas dari
\definitionsverweis {pemetaan-pemetaan singgung}{}{}
di setiap titik, yaitu

\mavergleichskettedisp
{\vergleichskette
{ T (\varphi)
}
{ =} { \biguplus_{P \in M} T_P(\varphi)
}
{ } {
}
{ } {
}
{ } {
}
}
{}{}{.}

}

\inputbeispiel{}
{

Misalkan $M$ suatu
\definitionsverweis {manifold diferensiabel}{}{}
dan
\maabbdisp {\alpha} { U } { V
} {}
suatu
\definitionsverweis {peta}{}{}
dengan

\mavergleichskette
{\vergleichskette
{ V
}
{ \subseteq }{ \R^n
}
{  }{
}
{    }{
}
{   }{
}
}
{}{}{}
\definitionsverweis {terbuka}{}{.}
Maka peta tersebut menginduksi bijeksi alami
\maabbeledisp {T \left(  \alpha^{-1}  \right)} {TV =  V \times \R^n } { TU
} {(Q,v)} { \left(  \alpha^{-1}(Q) ,[s \mapsto \alpha^{-1}( Q+ sv) ]  \right)
} {.}
Di sini,

\mavergleichskette
{\vergleichskette
{ s
}
{ \in }{ I
}
{  }{
}
{    }{
}
{   }{
}
}
{}{}{}
bergerak dalam suatu
\definitionsverweis {interval real}{}{}
sedemikian sehingga

\mavergleichskette
{\vergleichskette
{ Q+sv
}
{ \in }{ V
}
{  }{
}
{    }{
}
{   }{
}
}
{}{}{}
\zusatzklammer {bandingkan
Lema 9.5} {} {.}
Karena
\mathl{V \times \R^n}{} merupakan
\definitionsverweis {produk ruang-ruang topologis}{}{,}
maka

\mavergleichskette
{\vergleichskette
{ T V
}
{ = }{ V \times \R^n
}
{  }{
}
{    }{
}
{   }{
}
}
{}{}{}
sendiri merupakan
\definitionsverweis {ruang topologis}{}{,}
dan wajar untuk memindahkan topologi ini ke $TU$ serta menggunakannya untuk mengonstruksi sebuah topologi pada
\definitionsverweis {bundel tangen}{}{}
$TM$ secara keseluruhan.

}

\inputdefinition
{}
{

Misalkan $M$ suatu
\definitionsverweis {manifold diferensiabel}{}{}
dengan
\definitionsverweis {dimensi}{}{}
$n$ dan

\mavergleichskettedisp
{\vergleichskette
{ TM
}
{ =} {\biguplus_{P \in M} T_P M
}
{ } {
}
{ } {
}
{ } {
}
}
{}{}{}
\definitionsverweis {bundel tangen}{}{,}
yang dilengkapi dengan pemetaan proyeksi
\maabbeledisp {\pi} { TM } { M
} { (P,v) } { P
} {.}
\definitionswort {Bundel tangen}{} dilengkapi dengan
\definitionsverweis {topologi}{}{}
yang mempunyai sifat bahwa suatu himpunan bagian

\mavergleichskette
{\vergleichskette
{ W
}
{ \subseteq }{ TM
}
{  }{
}
{    }{
}
{   }{
}
}
{}{}{}
terbuka jika dan hanya jika untuk setiap
\definitionsverweis {peta}{}{}
\maabbdisp {\alpha} { U } { V
} {}
himpunan
\mathl{(T(\alpha)) \left(  W \cap \pi^{-1}(U)  \right)}{}
terbuka di dalam
\mathl{V \times \R^n}{}
.

}

Secara khusus, untuk setiap himpunan terbuka

\mavergleichskette
{\vergleichskette
{ U
}
{ \subseteq }{ M
}
{  }{
}
{    }{
}
{   }{
}
}
{}{}{}
prapeta

\mavergleichskette
{\vergleichskette
{  \pi^{-1}(U)
}
{ = }{ TU
}
{ \subseteq }{ TM
}
{    }{
}
{   }{
}
}
{}{}{}
terbuka, yakni proyeksi $\pi$ kontinu.



\inputfaktbeweis
{Differenzierbare Mannigfaltigkeit/Abbildung/Tangentialabbildung/Eigenschaften/Fakt}
{Lema}
{}
{

\faktsituation {Misalkan
\mathkor {} {M} {dan} {N} {}
\definitionsverweis {manifold diferensiabel}{}{}
dan misalkan
\maabbdisp {\varphi} { M } { N
} {}
suatu
\definitionsverweis {pemetaan diferensiabel}{}{.}
Misalkan
\maabbdisp {T(\varphi)} {TM} {TN
} {}
\definitionsverweis {pemetaan singgung}{}{}
yang bersesuaian.}
\faktuebergang {Maka berlaku pernyataan-pernyataan berikut.}
\faktfolgerung {\aufzaehlungsechs{Terdapat suatu
\definitionsverweis {diagram komutatif}{}{}

\mathdisp {\begin{matrix} TM & \stackrel{ T(\varphi) }{\longrightarrow} & TN &  \\ \!\!\!\!\! \pi \downarrow & & \downarrow \pi \!\!\!\!\!  & \\  M  & \stackrel{ \varphi }{\longrightarrow} &  N  & \! \! \! \!\!\!\!\! .  \\ \end{matrix}} {  }

}{Untuk suatu peta
\maabbdisp {\alpha} { U } { V
} {}
dengan

\mavergleichskette
{\vergleichskette
{ U
}
{ \subseteq }{ M
}
{  }{
}
{    }{
}
{   }{
}
}
{}{}{}
terbuka dan

\mavergleichskette
{\vergleichskette
{ V
}
{ \subseteq }{ \R^m
}
{  }{
}
{    }{
}
{   }{
}
}
{}{}{}
terbuka, terdapat diagram komutatif

\mathdisp {\begin{matrix} TU & \stackrel{ T(\alpha) }{\longrightarrow} & TV = V \times \R^m  &  \\ \downarrow & & \downarrow  & \\  U  & \stackrel{ \alpha }{\longrightarrow} &  V  & \! \! \! \!\!\!\!\! .  \\ \end{matrix}} {  }

}{Jika
\mathkor {} {M \subseteq \R^m} {dan} {N \subseteq \R^n} {}
merupakan
\definitionsverweis {himpunan-himpunan bagian terbuka}{}{}
dan bundel-bundel tangen diidentifikasi berturut-turut dengan
\mathl{M \times \R^m}{} dan
\mathl{N \times \R^n}{}
,
maka pemetaan singgung sama dengan
\maabbeledisp {} {M \times \R^m } { N \times \R^n
} {(P,v)} { \left(  \varphi(P), \left(  D\varphi  \right)_{ P } \left(  v  \right)  \right)
} {.}
}{Jika $L$ suatu
\definitionsverweis {manifold}{}{}
lain dan
\maabbdisp {\psi} { L } { M
} {}
suatu pemetaan diferensiabel lain, maka berlaku

\mavergleichskettedisp
{\vergleichskette
{ T( \varphi \circ \psi)
}
{ =} { T( \varphi) \circ T( \psi)
}
{ } {
}
{ } {
}
{ } {
}
}
{}{}{.}
}{Pemetaan singgung
\mathl{T(\varphi)}{}
bersifat
\definitionsverweis {kontinu}{}{.}
}{Jika $\varphi$ suatu
\definitionsverweis {difeomorfisme}{}{,}
maka
\mathl{T(\varphi)}{}
merupakan suatu
\definitionsverweis {homeomorfisme}{}{.}
}}
\faktzusatz {}
\faktzusatz {}

}
{

\teilbeweis {}{}{}
{(1) langsung mengikuti dari definisi
\definitionsverweis {pemetaan singgung}{}{.}}
{}
\teilbeweis {}{}{}
{(2) mengikuti dari (1) dengan menggunakan identifikasi alami

\mavergleichskette
{\vergleichskette
{ TV
}
{ = }{ V \times \R^m
}
{  }{
}
{    }{
}
{   }{
}
}
{}{}{}
untuk suatu himpunan terbuka $V\subseteq\R^m$; dengan demikian $TV$ terbuka di dalam $\R^{2m}$.}
{}
\teilbeweis {}{}{}
{(3) mengikuti dari
Lema 9.10  (1).}
{}
\teilbeweis {}{}{}
{(4) mengikuti dari
Lema 9.10  (4).}
{}
\teilbeweis {}{}{}
{(5). Bagi suatu himpunan terbuka

\mavergleichskettedisp
{\vergleichskette
{ V
}
{ \subseteq} { N
}
{ } {
}
{ } {
}
{ } {
}
}
{}{}{}
berlaku

\mavergleichskette
{\vergleichskette
{ TV
}
{ = }{ \pi^{-1} V
}
{ \subseteq }{ TN
}
{    }{
}
{   }{
}
}
{}{}{}
terbuka, sehingga

\mavergleichskette
{\vergleichskette
{ (T(\varphi))^{-1} TV
}
{ = }{ T(\varphi^{-1}(V))
}
{ \subseteq }{ TM
}
{    }{
}
{   }{
}
}
{}{}{}
terbuka. Cukup dibuktikan kontinuitas
\maabbdisp {} {T(\varphi^{-1}(V))} {TV
} {}
.
Di sini, $V$ dapat diambil sebagai domain peta dan $\varphi^{-1}(V)$ dapat ditutupi oleh domain-domain peta. Maka cukup menunjukkan kontinuitas pembatasan
\maabbdisp {T(\varphi)\!\mid_{TU}} {TU} {TV
} {}
untuk domain-domain peta $U\subseteq M$ dan $V\subseteq N$ dengan $\varphi(U)\subseteq V$. Pilih peta
\maabb {\alpha} {U} {U'
} {}
dan
\maabbdisp {\beta} {V} {V'
} {,}
lalu tuliskan representasi koordinatnya sebagai
\maabbdisp {\widetilde\varphi=\beta\circ\varphi\circ\alpha^{-1}} {U'} {V'
} {.}
Selanjutnya terdapat diagram komutatif

\mathdisp {\begin{matrix} TU & \stackrel{ T(\varphi) }{\longrightarrow} & TV &  \\ \downarrow & & \downarrow  & \\ U' \times \R^m  & \stackrel{  }{\longrightarrow} &  V' \times \R^n  & \! \! \! \!\!\!\!\! ,  \\ \end{matrix}} {  }
 dengan pemetaan-pemetaan vertikal berupa homeomorfisme. Untuk pemetaan horizontal bawah kita berada dalam situasi yang dijelaskan pada (3). Jadi, kita harus membuktikan kontinuitas pemetaan
\maabbeledisp {\widetilde T} {U' \times \R^m } { V' \times \R^n
} {(x,v)} {(\widetilde\varphi(x), \left(  D\widetilde\varphi  \right)_{ x } \left(  v  \right) )
} {,}
dan kita hanya perlu meninjau komponen belakang, yaitu
\mathl{\left(  D\widetilde\varphi  \right)_{ x } \left(  v  \right)}{.}
Komponen ke-$j$ dari komponen tersebut adalah

\mathdisp {\sum_{i=1}^m v_i { \frac{   \partial  \widetilde\varphi_j   }{  \partial   x_i   } } ( x)} { , }

dan berdasarkan asumsi
$C^1$-\definitionsverweis {diferensiabilitas
}{}{-}, semua ini merupakan pemetaan kontinu.}
{}
\teilbeweis {}{}{}
{(6) mengikuti dari (5).}
{}

}

\bild{ \begin{center}
\includegraphics[width=5.5cm]{ \bildeinlesungjpg {Torus_vectors_oblique} {jpg} }
\end{center}
\bildtext {Suatu medan vektor pada torus. Setiap titik torus diberi suatu arah tangensial, yang ditunjukkan oleh panah-panah.} }

\bildlizenz { Torus vectors oblique.jpg } {} {RokerHRO} {Commons} {CC-by-sa 3.0} {}

\inputdefinition
{}
{

Misalkan $M$ suatu
\definitionsverweis {manifold diferensiabel}{}{.}
Suatu pemetaan
\maabbdisp {F} { M } { TM
} {}
dengan sifat bahwa

\mavergleichskette
{\vergleichskette
{ F(P)
}
{ \in }{  T_PM
}
{  }{
}
{    }{
}
{   }{
}
}
{}{}{}
untuk setiap titik

\mavergleichskette
{\vergleichskette
{ P
}
{ \in }{ M
}
{  }{
}
{    }{
}
{   }{
}
}
{}{}{}
disebut
\zusatzklammer {tak bergantung waktu} {} {}
\definitionswort {medan vektor}{.}

}
\emph{Catatan penyunting.} Definisi di atas mencakup sebarang penampang sebagai medan vektor. Untuk menerapkan teorema keberadaan atau ketunggalan bagi persamaan diferensial biasa, diperlukan syarat regularitas tambahan pada medan tersebut, sesuai dengan teorema yang digunakan.
Jadi, suatu medan vektor memberikan sebuah vektor arah di setiap titik. Secara ringkas, dikatakan pula bahwa medan vektor merupakan suatu \stichwort {penampang} {} dalam bundel tangen. Medan-medan vektor menghasilkan persamaan diferensial biasa pada manifold. Pada suatu himpunan terbuka

\mavergleichskette
{\vergleichskette
{ U
}
{ \subseteq }{ \R^n
}
{  }{
}
{    }{
}
{   }{
}
}
{}{}{}
terdapat \stichwort {medan-medan vektor standar} {}

\mathbed {\partial_i} {}
 {1 \leq i \leq n} {}
 {} {} {} {,}
yang memetakan setiap titik

\mavergleichskette
{\vergleichskette
{ P
}
{ \in }{ U
}
{  }{
}
{    }{
}
{   }{
}
}
{}{}{}
ke
\definitionsverweis {vektor standar}{}{}

\mavergleichskette
{\vergleichskette
{ e_i
}
{ \in }{ \R^n
}
{  }{
}
{    }{
}
{   }{
}
}
{}{}{,}
yang dipandang sebagai vektor singgung di $P$. Sebarang medan vektor $F$ pada $U$ kemudian dapat dituliskan sebagai

\mavergleichskette
{\vergleichskette
{ F
}
{ = }{ \sum_{i = 1}^n  f_i \partial_i
}
{  }{
}
{    }{
}
{   }{
}
}
{}{}{}
dengan fungsi-fungsi $f_i$.

\inputdefinition
{}
{

Misalkan $M$ suatu
\definitionsverweis {manifold diferensiabel}{}{.}
Himpunan

\mavergleichskettedisp
{\vergleichskette
{ T^*M
}
{ =} {\biguplus_{P \in M} T^*_PM
}
{ } {
}
{ } {
}
{ } {
}
}
{}{}{,}
yang dilengkapi dengan pemetaan proyeksi
\maabbeledisp {\pi} { T^*M } { M
} { (P,u) } { P
} {,}
dan dengan
\definitionsverweis {topologi}{}{}
yang mempunyai sifat bahwa suatu himpunan bagian

\mavergleichskette
{\vergleichskette
{ W
}
{ \subseteq }{ T^*M
}
{  }{
}
{    }{
}
{   }{
}
}
{}{}{}
\definitionsverweis {terbuka}{}{}
jika dan hanya jika untuk setiap
\definitionsverweis {peta}{}{}
\maabbdisp {\alpha} { U } { V
} {}
himpunan
\mathl{(T^*(\alpha))^{-1} \left(  W \cap \pi^{-1}(U)  \right)}{}
terbuka di dalam
\mathl{V \times (\R^n)^*}{}
,
disebut \definitionswort {bundel kotangen}{} dari $M$.

}
Penampang-penampang dalam bundel kotangen disebut bentuk diferensial-$1$. Kita akan kembali membahasnya secara terperinci.

\chapter{Lembar Kerja 10}
\setcounter{section}{10}

  \zwischenueberschrift{Soal latihan}

\inputaufgabe
{}
{

Tunjukkan bahwa

\mavergleichskettedisp
{\vergleichskette
{M
}
{ =} { \left\{ (x,y,z,t) \in \R^4  \mid   x+x^2y+z^2+t^3 = 0         \right\}
}
{ } {
}
{ } {
}
{ } {
}
}
{}{}{}
merupakan sebuah
\definitionsverweis {submanifold tertutup}{}{}
dari $\R^4$.

}
{} {}

\inputaufgabe
{}
{

Misalkan
\maabbdisp {f} {\R} {\R_+
} {}
sebuah
\definitionsverweis {fungsi yang terdiferensialkan secara kontinu}{}{.}
Tunjukkan bahwa
\definitionsverweis {permukaan putar}{}{}
dari
\definitionsverweis {grafik}{}{}
$f$ merupakan sebuah
\definitionsverweis {submanifold tertutup}{}{}
dari $\R^{3}$.

}
{} {}

\inputaufgabe
{}
{

Tunjukkan bahwa himpunan semua
\mathl{n\times n}{} matriks riil dengan
\definitionsverweis {determinan}{}{}
$1$ memiliki dimensi
\mathl{(n^2-1)}{} dan merupakan sebuah
\definitionsverweis {submanifold tertutup}{}{}
dari $\R^{n^2}$.

}
{} {}

\inputaufgabe
{}
{

Berikan sebuah contoh
\definitionsverweis {himpunan bagian tertutup}{}{}
\mathl{M\subseteq \R}{,} yang bukan
\definitionsverweis {submanifold tertutup}{}{} dari $\R$.

}
{} {}

\inputaufgabe
{}
{

Tentukan semua
\definitionsverweis {submanifold tertutup}{}{}
dari $\R$.

}
{} {}

\inputaufgabe
{}
{

Tentukan semua
\definitionsverweis {submanifold tertutup}{}{}
dari $S^1$.

}
{} {}

\inputaufgabe
{}
{

Misalkan
\mathkor {} {M} {dan} {N} {}
dua himpunan yang
\definitionsverweis {saling lepas}{}{}
dan masing-masing merupakan
\definitionsverweis {submanifold tertutup}{}{}
dari $\R^n$ dengan dimensi yang sama. Tunjukkan bahwa
\definitionsverweis {gabungan}{}{}
keduanya,

\mathl{M \cup N}{}, juga merupakan sebuah submanifold tertutup, dan bahwa pernyataan ini tidak berlaku tanpa asumsi saling lepas.

}
{} {}

\inputaufgabe
{}
{

Misalkan
\maabb {f} {\R^n } {\R
} {}
sebuah
\definitionsverweis {fungsi yang terdiferensialkan secara kontinu}{}{.}
Tunjukkan bahwa
\definitionsverweis {grafik}{}{}

\mavergleichskette
{\vergleichskette
{ \Gamma_f
}
{ \subseteq }{ \R^{n+1}
}
{  }{
}
{    }{
}
{   }{
}
}
{}{}{}
merupakan sebuah
\definitionsverweis {submanifold tertutup}{}{}
dari $\R^{n+1}$.

}
{} {}

\inputaufgabegibtloesung
{}
{

Misalkan

\mavergleichskette
{\vergleichskette
{ M
}
{ \neq }{ \emptyset
}
{  }{
}
{    }{
}
{   }{
}
}
{}{}{}
sebuah
\definitionsverweis {manifold diferensiabel}{}{}
berdimensi $n$. Tunjukkan bahwa terdapat sebuah rantai
\definitionsverweis {submanifold tertutup}{}{}

\mavergleichskettedisp
{\vergleichskette
{ M_0
}
{ \subseteq} { M_1
}
{ \subseteq} { M_2
}
{ \subseteq  \ldots \subseteq} { M_{n-1}
}
{ \subseteq} { M_n
}
}
{
\vergleichskettefortsetzung
{ =} { M
}
{ } {}
{ } {}
{ } {}
}{}{}
sedemikian sehingga submanifold tertutup $M_i$ berdimensi $i$.

}
{} {}

\inputaufgabegibtloesung
{}
{

Perhatikan himpunan

\mavergleichskettedisp
{\vergleichskette
{N
}
{ =} { \left\{ A= \begin{pmatrix}  a   &   b   \\  c  & d \end{pmatrix}   \mid   A \text{ adalah nilpoten}         \right\}
}
{ \subseteq} {\R^4
}
{ } {
}
{ } {
}
}
{}{}{}
yang terdiri atas
$2\times 2$ \definitionsverweis {matriks}{}{}
riil nilpoten, serta himpunan

\mavergleichskettedisp
{\vergleichskette
{M
}
{ =} { N \setminus \{ \begin{pmatrix}  0   &   0  \\  0  &  0 \end{pmatrix}  \}
}
{ } {
}
{ } {
}
{ } {
}
}
{}{}{.}

a) Apakah $N$ terhubung?

b) Tunjukkan bahwa $M$ merupakan sebuah
\definitionsverweis {submanifold tertutup}{}{}
dari suatu himpunan bagian terbuka

\mavergleichskette
{\vergleichskette
{ G
}
{ \subseteq }{ \R^4
}
{  }{
}
{    }{
}
{   }{
}
}
{}{}{.}

c) Tentukan
\definitionsverweis {dimensi}{}{}
$M$.

d) Apakah $M$ terhubung?

e) Tutupi
\mathl{M}{} dengan peta-peta topologis eksplisit.

}
{} {}

\inputaufgabe
{}
{

Misalkan

\mavergleichskette
{\vergleichskette
{ M
}
{ \subseteq }{ \R^n
}
{  }{
}
{    }{
}
{   }{
}
}
{}{}{}
sebuah
\definitionsverweis {submanifold tertutup}{}{}
dan

\mavergleichskette
{\vergleichskette
{ P
}
{ \in }{ M
}
{  }{
}
{    }{
}
{   }{
}
}
{}{}{}
sebuah titik. Misalkan
\mathl{T_P(i)}{} adalah
\definitionsverweis {pemetaan singgung}{}{}
yang bersesuaian dengan inklusi
\maabbdisp {i} { M } {\R^n
} {}
dan $\gamma$ sebuah
\definitionsverweis {lintasan diferensiabel}{}{}
di $M$ dengan

\mavergleichskettedisp
{\vergleichskette
{ \gamma(0)
}
{ =} { P
}
{ } {
}
{ } {
}
{ } {
}
}
{}{}{.}
Tunjukkan bahwa, di dalam

\mavergleichskette
{\vergleichskette
{ \R^n
}
{ \cong }{ T_P \R^n
}
{  }{
}
{    }{
}
{   }{
}
}
{}{}{}, kesamaan

\mavergleichskettedisp
{\vergleichskette
{ T_P(i) ( [\gamma] )
}
{ =} { \left(  i \circ \gamma  \right) '(0)
}
{ } {
}
{ } {
}
{ } {
}
}
{}{}{}
berlaku.

}
{} {}

\inputaufgabe
{}
{

Misalkan

\mavergleichskette
{\vergleichskette
{ G
}
{ \subseteq }{ \R^n
}
{  }{
}
{    }{
}
{   }{
}
}
{}{}{}
\definitionsverweis {terbuka}{}{,}
\maabb {\varphi} { G } {\R^m
} {}
sebuah
\definitionsverweis {pemetaan diferensiabel}{}{}
dan $M$ adalah
\definitionsverweis {serat}{}{}
di atas

\mavergleichskette
{\vergleichskette
{ 0
}
{ \in }{ \R^m
}
{  }{
}
{    }{
}
{   }{
}
}
{}{}{.}
Andaikan
\definitionsverweis {diferensial total}{}{}
bersifat
\definitionsverweis {surjektif}{}{}
di setiap titik serat ini. Tunjukkan bahwa untuk

\mavergleichskette
{\vergleichskette
{ P
}
{ \in }{ M
}
{  }{
}
{    }{
}
{   }{
}
}
{}{}{}
\definitionsverweis {ruang singgung}{}{}
dalam pengertian
Definisi 53.7 (Analisis (Osnabrück 2021-2023)) berimpit dengan
\definitionsverweis {ruang singgung}{}{}
dari
\definitionsverweis {manifold diferensiabel}{}{}
$M$ di titik $P$.

}
{} {}

\inputaufgabe
{}
{

Misalkan $M$ sebuah
\definitionsverweis {manifold diferensiabel}{}{}
dan
\maabbdisp {\pi} {TM} {M
} {}
adalah
\definitionsverweis {bundel tangen}{}{.}
Tunjukkan bahwa pemetaan proyeksi ini
\definitionsverweis {kontinu}{}{.}

}
{} {}

\inputaufgabe
{}
{

Misalkan $TM$ adalah
\definitionsverweis {bundel tangen}{}{}
dari suatu
\definitionsverweis {manifold diferensiabel}{}{}
$M$. Tunjukkan bahwa himpunan-himpunan
\mathl{(T(\alpha))^{-1} (V \times W)}{} untuk semua peta
\maabbdisp {\alpha} { U } { V
} {}
dengan
\mathl{V,W \subseteq \R^n}{} terbuka membentuk suatu
\definitionsverweis {basis topologi}{}{}
pada bundel tangen tersebut.

}
{} {}

\inputaufgabegibtloesung
{}
{

Misalkan

\mavergleichskette
{\vergleichskette
{ W
}
{ \subseteq }{ \R^n
}
{  }{
}
{    }{
}
{   }{
}
}
{}{}{}
\definitionsverweis {terbuka}{}{,}
\maabb {h} { W } { \R
} {}
sebuah
fungsi kelas $C^2$
dan

\mavergleichskette
{\vergleichskette
{ Y
}
{ = }{ h^{-1}(c)
}
{  }{
}
{    }{
}
{   }{
}
}
{}{}{}
adalah
\definitionsverweis {serat}{}{}
di atas

\mavergleichskette
{\vergleichskette
{ c
}
{ \in }{ \R
}
{  }{
}
{    }{
}
{   }{
}
}
{}{}{,}
dengan $h$
\definitionsverweis {reguler}{}{}
di setiap titik pada $Y$. Realisasikan
\definitionsverweis {bundel tangen}{}{}
dari $Y$ sebagai sebuah
\definitionsverweis {submanifold tertutup}{}{}
dari $W \times \R^n$.

}
{} {}

\inputaufgabe
{}
{

Misalkan

\mavergleichskette
{\vergleichskette
{ Y
}
{ \subseteq }{ \R^n
}
{  }{
}
{    }{
}
{   }{
}
}
{}{}{}
sebuah
\definitionsverweis {submanifold tertutup}{}{}
dan
\maabb {\varphi} { V } { U \subseteq Y
} {}
sebuah parametrisasi difeomorfik dari suatu himpunan bagian terbuka

\mavergleichskette
{\vergleichskette
{ U
}
{ \subseteq }{ Y
}
{  }{
}
{    }{
}
{   }{
}
}
{}{}{}
oleh suatu himpunan bagian terbuka

\mavergleichskette
{\vergleichskette
{ V
}
{ \subseteq }{ \R^d
}
{  }{
}
{    }{
}
{   }{
}
}
{}{}{,}
dengan $\varphi$ juga dipandang sebagai pemetaan ke $\R^n$. Definisikan
\maabbdisp {X_i} {U} {TU
} {}
melalui $X_i(\varphi(x))=T_x(\varphi)(e_i)=\partial_i\varphi(x)$. Tunjukkan bahwa diagram

\mathdisp {\begin{matrix}  V  & \stackrel{ \partial_i }{\longrightarrow} &  TV &  \\ \!\!\!\!\! \,\, \, \varphi \downarrow & \partial_i \varphi \searrow & \downarrow T(\varphi) \!\!\!\!\!  & \\  U  & \stackrel{ X_i }{\longrightarrow} &  TU & \!\!\!\!\!   \\ \end{matrix}} {  }

tersebut komutatif.

}
{} {}

\inputaufgabe
{}
{

Misalkan
\mathkor {} {M} {dan} {N} {}
adalah
\definitionsverweis {manifold diferensiabel}{}{} dan
\maabb {\varphi} { M } { N
} {} sebuah
\definitionsverweis {pemetaan diferensiabel}{}{.} Tunjukkan bahwa
\definitionsverweis {pemetaan singgung}{}{}
yang bersesuaian,
\maabbdisp {T(\varphi)} {TM} {TN
} {}
\definitionsverweis {kontinu}{}{.}

}
{} {}

\inputaufgabe
{}
{

Hitung
\definitionsverweis {pemetaan singgung}{}{} $T\varphi$ untuk
\maabbeledisp {\varphi} {\R^3} {\R^2
} {(x,y,z)} {(x^2y-3xz^3+y^2,x \sin  y  -e^{yz})
} {} dengan menggunakan identifikasi
\mathkor {} {T\R^3 =\R^3 \times \R^3} {dan} {T\R^2 =\R^2 \times \R^2} {.}

}
{} {}

\inputaufgabe
{}
{

Berikan contoh sebuah kurva diferensiabel
\maabbdisp {\gamma} { [0,1[} { S^1
} {}
sedemikian sehingga
\definitionsverweis {limit}{}{}

\mathl{\operatorname{lim}_{   t  \rightarrow  1   } \,  \gamma(t)}{} ada, tetapi limit
\mathl{\operatorname{lim}_{   t  \rightarrow  1   } \, ( \gamma(t), T_t (\gamma ) (1))}{} di dalam
\mathl{TS^1}{} tidak ada.

}
{} {}

\inputaufgabe
{}
{

Misalkan

\mavergleichskette
{\vergleichskette
{ M
}
{ \subseteq }{ \R^n
}
{  }{
}
{    }{
}
{   }{
}
}
{}{}{}
sebuah
\definitionsverweis {submanifold tertutup}{}{.}
Tafsirkan komposisi

\mathdisp {TM \longrightarrow T\R^n = \R^n \times \R^n \stackrel{+}{\longrightarrow} \R^n} { . }

}
{} {}

\inputaufgabe
{}
{

Tunjukkan bahwa pemetaan
\maabbeledisp {} {TS^1} { \R^2
} {( (a,b),t (-b,a))} { (a,b) + t (-b,a)
} {,}
memiliki dua titik prabayang untuk setiap titik
\mathl{(x,y) \in \R^2}{} di luar cakram satuan, satu titik prabayang pada lingkaran satuan, dan tidak memiliki titik prabayang di dalam cakram satuan terbuka. Tafsirkan hasil ini secara geometris.

}
{} {}

\inputaufgabe
{}
{

Misalkan

\mavergleichskette
{\vergleichskette
{ M
}
{ \subseteq }{ N
}
{  }{
}
{    }{
}
{   }{
}
}
{}{}{}
sebuah
\definitionsverweis {submanifold tertutup}{}{}
dari
\definitionsverweis {manifold diferensiabel}{}{}
$N$. Tunjukkan bahwa

\mavergleichskette
{\vergleichskette
{ TM
}
{ \subseteq }{ TN
}
{  }{
}
{    }{
}
{   }{
}
}
{}{}{}
merupakan sebuah himpunan bagian tertutup.

}
{} {}

\inputaufgabe
{}
{

Berikan sebuah contoh pemetaan yang
\definitionsverweis {injektif}{}{}
dan
\definitionsverweis {diferensiabel}{}{},
\maabbdisp {\varphi} { M } { N
} {}
antara dua
\definitionsverweis {manifold diferensiabel}{}{},
\mathkor {} {M} {dan} {N} {}
sedemikian sehingga
\definitionsverweis {pemetaan singgung}{}{}
yang bersesuaian,
\maabbdisp {T(\varphi)} {TM} {TN
} {},
tidak injektif.

}
{} {}

\inputaufgabe
{}
{

Berikan sebuah contoh pemetaan yang
\definitionsverweis {surjektif}{}{}
dan
\definitionsverweis {diferensiabel}{}{},
\maabbdisp {\varphi} { M } { N
} {}
antara dua
\definitionsverweis {manifold diferensiabel}{}{},
\mathkor {} {M} {dan} {N} {}
sedemikian sehingga
\definitionsverweis {pemetaan singgung}{}{}
yang bersesuaian,
\maabbdisp {T(\varphi)} {TM} {TN
} {},
tidak surjektif.

}
{} {}

\inputaufgabegibtloesung
{}
{

Berikan sebuah
\definitionsverweis {medan vektor}{}{}
yang
\definitionsverweis {kontinu}{}{}
pada $S^2$ dan hanya memiliki satu titik nol.

}
{} {}

\inputaufgabe
{}
{

\definitionsverweis {Permukaan bola satuan}{}{}

\mathl{S^2}{} berputar satu putaran penuh dalam satu detik
\zusatzklammer {berlawanan arah jarum jam jika dilihat dari kutub utara} {} {}
mengelilingi sumbu kutub. Kurva diferensiabel dan vektor singgung manakah pada suatu titik
\mathl{P \in S^2}{} yang bersesuaian dengan gerak ini? Deskripsikan
\definitionsverweis {medan vektor}{}{}
yang terkait,
\maabbdisp {} {S^2} {TS^2 \subseteq T\R^3 = \R^6
} {.}

}
{} {}

Misalkan $M$ sebuah
\definitionsverweis {manifold diferensiabel}{}{}
dengan
\definitionsverweis {bundel tangen}{}{}
\maabb {p} {TM} {M
} {}
dan misalkan
\maabbdisp {\psi} { Z } { M
} {}
sebuah
\definitionsverweis {pemetaan}{}{,}
dengan $Z$ menyatakan suatu himpunan. Sebuah pemetaan
\maabbdisp {F} { Z } { TM
} {}
disebut
\definitionswort {medan vektor sepanjang}{}
$\psi$, jika

\mavergleichskette
{\vergleichskette
{ p \circ F
}
{ = }{ \psi
}
{  }{
}
{    }{
}
{   }{
}
}
{}{}{}
berlaku.

  \zwischenueberschrift{Soal untuk dikumpulkan}

\inputaufgabe
{8 (3+3+2)}
{

Misalkan

\mavergleichskette
{\vergleichskette
{ m,n
}
{ \in }{\N_+
}
{  }{
}
{    }{
}
{   }{
}
}
{}{}{.}

a) Tunjukkan bahwa himpunan

\mavergleichskettedisp
{\vergleichskette
{ M
}
{ =} { \left\{ (x,y,u,v) \in \R^4  \mid  ux^m+vy^n = 1         \right\}
}
{ } {
}
{ } {
}
{ } {
}
}
{}{}{}
merupakan sebuah
\definitionsverweis {submanifold tertutup}{}{}
dari $\R^4$.

b) Tunjukkan bahwa
\definitionsverweis {pemetaan}{}{}
\maabbeledisp {\varphi} { M } {\R^2
} {(x,y,u,v)} {(x,y)
} {,}
\definitionsverweis {diferensiabel}{}{}
dan di setiap titik

\mavergleichskette
{\vergleichskette
{ P
}
{ \in }{ M
}
{  }{
}
{    }{
}
{   }{
}
}
{}{}{}
\definitionsverweis {reguler}{}{.}

c) Deskripsikan
\definitionsverweis {serat-serat}{}{}
dari $\varphi$.

}
{} {}

\inputaufgabe
{10 (2+3+5)}
{

Perhatikan
\definitionsverweis {pemetaan}{}{}
\maabbeledisp {\varphi} {\R} {\R^2
} { x } {(x^2,x^3)
} {.}

a) Tunjukkan bahwa pemetaan ini
\definitionsverweis {diferensiabel}{}{}
dan
\definitionsverweis {injektif}{}{.}

b) Tunjukkan bahwa $\varphi$ tidak
\definitionsverweis {reguler}{}{}
pada setidaknya satu titik.

c) Tunjukkan bahwa
\definitionsverweis {citra}{}{}
$\varphi$
\definitionsverweis {tertutup}{}{}
di $\R^2$, tetapi bukan
\definitionsverweis {submanifold tertutup}{}{}
dari $\R^2$.

}
{} {}

\inputaufgabe
{4}
{

Tunjukkan bahwa terdapat sebuah
\definitionsverweis {homeomorfisme}{}{}
antara
\definitionsverweis {bundel tangen}{}{}
\mathl{TS^1}{} dari
$1$-\definitionsverweis {sfera}{}{} $S^1$ dan
\definitionsverweis {produk}{}{} $S^1 \times \R$.

}
{} {}

\inputaufgabe
{5}
{

Misalkan $M$ sebuah
\definitionsverweis {manifold diferensiabel}{}{}
dan
\maabbdisp {\pi} {TM} {M
} {}
adalah
\definitionsverweis {bundel tangen}{}{.}
Tunjukkan bahwa
\mathl{TM}{} sendiri secara alami merupakan sebuah
\definitionsverweis {manifold topologis}{}{.}

}
{} {}

Dalam soal berikut digunakan konsep modul atas $R$
\zusatzklammer {yang merupakan generalisasi langsung dari konsep ruang vektor} {} {.}

Misalkan $R$ sebuah
\definitionsverweis {gelanggang komutatif}{}{}
dan

\mavergleichskette
{\vergleichskette
{ M
}
{ = }{ (M,+,0)
}
{  }{
}
{    }{
}
{   }{
}
}
{}{}{}.
Dengan notasi \stichwort {aditif} {}, struktur ini merupakan sebuah
\definitionsverweis {grup komutatif}{}{.}
$M$ disebut sebuah
\definitionswortpraemath {R}{ modul }{,}
jika ditetapkan suatu operasi
\maabbeledisp {} {R \times M } { M
} {(r,v)} { rv = r\cdot v
} {,}
\zusatzklammer {yang disebut \stichwort {perkalian skalar} {}} {} {}
dan operasi itu memenuhi aksioma-aksioma berikut
\zusatzklammer {di sini

\mavergleichskettek
{\vergleichskettek
{ r,s
}
{ \in }{ R
}
{  }{
}
{    }{
}
{   }{
}
}
{}{}{}
dan

\mavergleichskettek
{\vergleichskettek
{ u,v
}
{ \in }{ M
}
{  }{
}
{    }{
}
{   }{
}
}
{}{}{}
dipilih sebarang} {} {:}
\aufzaehlungvier{
\mavergleichskette
{\vergleichskette
{ r(su)
}
{ = }{ (rs) u
}
{  }{
}
{    }{
}
{   }{
}
}
{}{}{,}
}{
\mavergleichskette
{\vergleichskette
{ r(u+v)
}
{ = }{ (ru) + (rv)
}
{  }{
}
{    }{
}
{   }{
}
}
{}{}{,}
}{
\mavergleichskette
{\vergleichskette
{ (r+s)u
}
{ = }{ (ru)+ (su)
}
{  }{
}
{    }{
}
{   }{
}
}
{}{}{,}
}{
\mavergleichskette
{\vergleichskette
{ 1 u
}
{ = }{  u
}
{  }{
}
{    }{
}
{   }{
}
}
{}{}{.}
}

\inputaufgabe
{4}
{

Misalkan $M$ sebuah
\definitionsverweis {manifold diferensiabel}{}{,}
misalkan

\mavergleichskette
{\vergleichskette
{ R
}
{ = }{ C^1(M,\R)
}
{  }{
}
{    }{
}
{   }{
}
}
{}{}{}
adalah
\definitionsverweis {gelanggang}{}{}
dari
\definitionsverweis {fungsi-fungsi diferensiabel}{}{}
pada $M$, dan misalkan $F$ adalah himpunan semua
\definitionsverweis {medan vektor}{}{}
pada $M$.

a) Definisikan suatu penjumlahan pada $F$ sedemikian sehingga $F$ menjadi sebuah
\definitionsverweis {grup komutatif}{}{.}

b) Definisikan suatu perkalian skalar
\maabbeledisp {} {R \times F } { F
} {(f,s)} {fs
} {,}
sedemikian sehingga $F$ menjadi sebuah
$R$-\definitionsverweis {modul}{}{.}

}
{} {}

\section*{Solusi yang disediakan oleh sumber}
\addcontentsline{toc}{section}{Solusi yang disediakan oleh sumber}
\subsection*{Solusi Soal 10.9}
Jika $n=0$, ambil saja $M_0=M$. Sekarang anggap $n\geq1$. Pilih suatu titik $P\in M$, sebuah lingkungan peta terbuka
\mathl{P \in U}{},
dan sebuah peta
\maabbdisp {\alpha} { U } {V \subseteq \R^n
} {,}
dengan $V=B(0,3)$ bola terbuka berpusat di $0$ berjari-jari $3$ dan
\mathl{\alpha(P)=0}{}.
Untuk
\mathl{i=0 , \ldots , n-1}{}, perhatikan

\mathdisp {B_i = \left\{  x \in V  \mid  (x_1-1)^2+x_2^2 + \cdots + x_{i+1}^2 = 1 , \, x_{i+2}=\cdots=x_n=0       \right\}} { , }

dengan syarat koordinat terakhir dipahami kosong jika $i=n-1$.
Jadi $B_{n-1}$ adalah permukaan bola berpusat di
\mathl{(1,0 , \ldots , 0)}{}
dan berjari-jari $1$, sedangkan $B_{n-2}$ adalah
\anfuehrung{ekuator}{}
di dalamnya yang didefinisikan oleh
\mathl{x_n=0}{}, dan seterusnya. Himpunan $B_i$ diperoleh dari
\mathl{B_{i+1}}{}
dengan menambahkan syarat
\mathl{x_{i+2}=0}{}.
Dengan demikian terdapat rantai menurun himpunan-himpunan tertutup

\mathdisp {B_{n-1} \supseteq B_{n-2} \supseteq  \ldots \supseteq B_1 \supseteq B_0} { , }

dengan $B_0$ terdiri atas dua titik
\mathl{(0,0 , \ldots , 0 )}{}
dan
\mathl{(2,0 , \ldots , 0 )}{.}
Kita memandang $B_i$ sebagai serat di atas titik nol dari pemetaan
\maabbeledisp {\varphi_i} { V } {\R^{n-i}
} {(x_1 , \ldots , x_n) } {  ((x_1-1)^2+x_2^2 + \cdots + x_n^2 - 1 , x_{i+2} , \ldots , x_n)
} {,}
dengan komponen koordinat terakhir kembali dipahami kosong jika $i=n-1$.
Matriks Jacobi pemetaan ini ialah

\mathdisp {\begin{pmatrix}  2x_1-2 &  2x_2  &  \ldots  &  2x_{i+1} &  2x_{i+2} &  \ldots &  2x_{n-1} &  2x_n
 \\  0 &  0 &  \ldots &  0 &  1 &  0 &  \ldots &  0
 \\ \vdots  & \vdots & \vdots &  \ddots &  \ddots &  \ddots  &  \ddots  & \vdots \\  0 &  0 &  \ldots &  \ldots &  \ldots &  0 &  1 &  0 \\  0 &  0 &  \ldots &  \ldots &  \ldots  &  \ldots &  0 &  1 \end{pmatrix}} { . }

Rank matriks ini lebih kecil daripada $n-i$ hanya jika
\mathl{x_1=1, \, x_2 = \cdots =x_{i+1}=0}{},
dan titik semacam itu tidak terletak di $B_i$. Jadi $\varphi_i$ reguler pada serat $B_i$, sehingga menurut teorema pemetaan implisit, $B_i$ merupakan submanifold tertutup dari $V$ berdimensi
\mathl{i}{.}

Sekarang tetapkan
\mathl{M_i= \alpha^{-1}(B_i)}{}
untuk
\mathl{i=0 , \ldots , n-1}{}
dan
\mathl{M_n= M}{.}
Karena setiap $B_i$ kompak, setiap $M_i$ merupakan himpunan tertutup di $M$. Karena kondisi submanifold tertutup bersifat lokal, himpunan-himpunan ini merupakan submanifold tertutup dari $M$, dan
\mathl{\operatorname{dim}M_i=i}{}.

Catatan edisi: Sumber memulai definisi $B_i$ pada $i=1$ tetapi kemudian memakai $B_0$ dan $M_0$ tanpa mendefinisikannya. Sumber juga menyatakan titik-titik $B_0$ sebagai $(\pm1,0,\ldots,0)$, padahal bola yang dipakai berpusat di $(1,0,\ldots,0)$; titik yang benar ialah $0$ dan $(2,0,\ldots,0)$. Rentang dan titik-titik itu diperbaiki di atas.
\subsection*{Solusi Soal 10.10}
a) Setiap matriks nilpoten
\mathl{\begin{pmatrix}  a   &   b   \\  c  & d \end{pmatrix}}{}
dapat dihubungkan dengan matriks nol melalui lintasan linear
\maabbeledisp {} {[0,1]} {N
} { t } {t \begin{pmatrix}  a   &   b   \\  c  & d \end{pmatrix}
} {,}
yang seluruhnya berada di dalam himpunan matriks nilpoten. Jadi $N$ terhubung lintasan dan, karena itu, terhubung.

b) Sebuah
$2 \times 2$-\definitionsverweis {matriks}{}{}
nilpoten tepat ketika jejak dan determinannya sama dengan $0$. Dengan demikian, himpunan $N$ semua matriks nilpoten dapat ditulis sebagai

\mathdisp {\left\{  (a,b,c,d)  \mid   a+d = 0 \text{ dan } ad-bc = 0         \right\}} { . }

Perhatikan pemetaan
\maabbeledisp {} {\R^4} {\R^2
} {(a,b,c,d)} { (a+d, ad-bc)
} {.}
Matriks Jacobi pemetaan ini ialah

\mathdisp {\begin{pmatrix}  1 &  0 &  0 &  1 \\  d  &  -c &  -b &  a  \end{pmatrix}} { . }

Pemetaan ini tidak
\definitionsverweis {reguler}{}{}
di titik nol, tetapi reguler di setiap titik lain pada serat $N$. Memang, jika
\mathl{P \in M}{}, maka dari syarat jejak dan determinan diperoleh

\mavergleichskettedisp
{\vergleichskette
{a^2
}
{ =} {-bc
}
{ } {
}
{ } {
}
{ } {
}
}
{}{}{}.
Karena itu,

\mavergleichskettedisp
{\vergleichskette
{b
}
{ =} { c
}
{ =} { 0
}
{ } {
}
{ } {
}
}
{}{}{}
langsung mengakibatkan

\mavergleichskettedisp
{\vergleichskette
{a
}
{ =} {b
}
{ =} { c
}
{ =} { d
}
{ =} { 0
}
}
{}{}{.}
Jadi matriks Jacobi mempunyai rank maksimal di semua titik $M$. Dengan

\mavergleichskettedisp
{\vergleichskette
{G
}
{ =} {\R^4 \setminus \{0\}
}
{ } {
}
{ } {
}
{ } {
}
}
{}{}{}
kita dapat memandang $M$ sebagai serat pemetaan yang dibatasi ke $G$, dan pemetaan itu reguler di seluruh serat. Oleh karena itu, berdasarkan
fakta tentang serat reguler,
$M$ merupakan submanifold tertutup dari $G$.

Catatan edisi: Sumber menulis $a^2=bc$ setelah memakai $d=-a$. Tanda yang benar dari $ad-bc=0$ adalah $a^2=-bc$; argumen rank di atas menggunakan identitas yang telah diperbaiki ini.

c) Berdasarkan
fakta tentang serat reguler,
dimensi $M$ ialah

\mavergleichskette
{\vergleichskette
{ 4-2
}
{ = }{ 2
}
{  }{
}
{    }{
}
{   }{
}
}
{}{}{.}

d) Tuliskan

\mavergleichskettedisp
{\vergleichskette
{M_+
}
{ =} { \left\{ (a,b,c,d) \in M  \mid   b>c        \right\}
}
{ } {
}
{ } {
}
{ } {
}
}
{}{}{}
dan

\mavergleichskettedisp
{\vergleichskette
{M_-
}
{ =} { \left\{ (a,b,c,d) \in M  \mid   b<c        \right\}
}
{ } {
}
{ } {
}
{ } {
}
}
{}{}{.}
Keduanya merupakan himpunan terbuka di $M$
\zusatzklammer {sebagai irisan $M$ dengan himpunan terbuka di $\R^4$ yang masing-masing ditentukan oleh $b>c$ dan $b<c$} {} {.}
Matriks
\mathkor {} {\begin{pmatrix}  0  &   1  \\  0 &  0 \end{pmatrix}} {dan} {\begin{pmatrix}  0  &   0  \\  1 &  0 \end{pmatrix}} {}
menunjukkan bahwa keduanya tidak kosong. Selain itu, keduanya menutupi seluruh $M$. Jika

\mavergleichskette
{\vergleichskette
{b
}
{ = }{ c
}
{  }{
}
{    }{
}
{   }{
}
}
{}{}{},
maka

\mavergleichskettedisp
{\vergleichskette
{ 0
}
{ =} { ad-bc
}
{ =} { -a^2 -b^2
}
{ } {
}
{ } {
}
}
{}{}{}
langsung memberikan

\mavergleichskette
{\vergleichskette
{a
}
{ = }{b
}
{ = }{ c
}
{ =   }{ d
}
{ =  }{ 0
}
}
{}{}{,}
dan titik ini tidak termasuk $M$. Jadi $M$ mempunyai tutupan oleh dua himpunan terbuka tak kosong yang saling lepas; akibatnya, $M$ tidak terhubung.

e) Untuk $M_+$, gunakan pemetaan
\maabbeledisp {\psi} {\R^2 \setminus \{(0,0)\}} { M_+
} {(a,u)} { (a, u + \sqrt{a^2+u^2}, u - \sqrt{a^2 +u^2},-a)
} {.}
Karena

\mavergleichskettealign
{\vergleichskettealign
{ad-bc
}
{ =} { -a^2 - \left(  u + \sqrt{a^2+u^2}  \right) \left( u - \sqrt{a^2 +u^2}  \right)
}
{ =} { -a^2 - \left(  u^2 - \left(  a^2+u^2 \right)   \right)
}
{ =} { 0
}
{ } {
}
}
{}
{}{}
syarat determinan terpenuhi, dan karena

\mavergleichskettedisp
{\vergleichskette
{b-c
}
{ =} { u + \sqrt{a^2+u^2} - \left(  u- \sqrt{a^2 +u^2})  \right)
}
{ =} { 2 \sqrt{a^2+u^2}
}
{ >} { 0
}
{ } {
}
}
{}{}{},
citranya terletak di $M_+$. Pemetaan
\maabbeledisp {\varphi} {M_+} { \R^2 \setminus \{(0,0)\}
} {(a,b,c,d)} { \left(  a  , \,   { \frac{  b+c }{  2 } }                   \right)
} {,}
merupakan invers dari pemetaan tersebut. Identitas

\mavergleichskettedisp
{\vergleichskette
{ \varphi \circ \psi
}
{ =} {
\operatorname{Id}_{ \R^2 \setminus \{(0,0)\}  }
}
{ } {
}
{ } {
}
{ } {
}
}
{}{}{}
jelas. Identitas lainnya mengikuti dari

\mavergleichskettealign
{\vergleichskettealign
{ { \frac{  b+c }{  2 } } + \sqrt{a^2 + \left( { \frac{  b+c }{  2 } }  \right)^2 }
}
{ =} { { \frac{  b+c }{  2 } } + { \frac{   1 }{  2 } }  \sqrt{4 a^2 + b^2 +c^2 +2bc }
}
{ =} { { \frac{  b+c }{  2 } } + { \frac{   1 }{  2 } }  \sqrt{ b^2 +c^2 -2bc }
}
{ =} { { \frac{  b+c }{  2 } } + { \frac{  b-c }{  2 } }
}
{ =} {b
}
}
{}
{}{}
dan

\mavergleichskettedisp
{\vergleichskette
{ { \frac{  b+c }{  2 } } - \sqrt{a^2 + \left(   { \frac{  b+c }{  2 } }   \right)^2 }
}
{ =} { { \frac{  b+c }{  2 } } - { \frac{  b-c }{  2 } }
}
{ =} { c
}
{ } {}
{ } {}
}
{}{}{.}
Kedua pemetaan kontinu, sehingga diperoleh suatu
\definitionsverweis {homeomorfisme}{}{.}
Untuk
\mathl{M_-}{}, tukarkan peran
\mathkor {} {b} {dan} {c} {.}
\subsection*{Solusi Soal 10.15}
Dengan regularitas $C^2$ yang dinyatakan dalam soal edisi ini, kita klaim

\mavergleichskettealignhandlinks
{\vergleichskettealignhandlinks
{ TY
}
{ =} { \left\{  (P,u)   \mid   h(P) = c , \,  \left(  D h   \right)_{ P } \left(  u  \right) = 0       \right\}
}
{ =} { \left\{  (P,u_1 , \ldots , u_n)   \mid   h(P) = c , \,  \sum_{i = 1}^n u_i { \frac{   \partial   h   }{  \partial   x_i   } } (P) = 0       \right\}
}
{ \subseteq} { W \times \R^n
}
{ } {
}
}
{}
{}{}
dengan proyeksi alami ke $Y$.

Definisikan pemetaan
\maabbeledisp {\Phi} {W\times\R^n} {\R^2
} {(P,u)} {\left(h(P),\left(Dh\right)_P(u)\right)
} {.}
Karena $h$ berkelas $C^2$, pemetaan $\Phi$ berkelas $C^1$, dan

\mathdisp {TY=\Phi^{-1}(c,0)} { . }

Kita tunjukkan bahwa $(c,0)$ merupakan nilai reguler. Untuk
\mathl{(P,u)\in TY}{}, diferensial $\Phi$ diberikan oleh

\mathdisp {(D\Phi)_{(P,u)}(v,w)=\left((Dh)_P(v),\,(D^2h)_P(v,u)+(Dh)_P(w)\right)} { . }

Karena $h$ reguler pada $Y$, bentuk linear $(Dh)_P$ surjektif. Untuk sebarang
\mathl{(r,s)\in\R^2}{}, pilih $v$ dengan $(Dh)_P(v)=r$, lalu pilih $w$ dengan

\mathdisp {(Dh)_P(w)=s-(D^2h)_P(v,u)} { . }

Dengan pilihan ini, $(D\Phi)_{(P,u)}(v,w)=(r,s)$, sehingga diferensial $\Phi$ surjektif di setiap titik serat. Teorema pemetaan implisit kini menunjukkan bahwa $TY$ merupakan submanifold tertutup dari $W\times\R^n$ berdimensi $2n-2$. Ketertutupannya juga langsung mengikuti dari kontinuitas $\Phi$ dan ketertutupannya titik $(c,0)$ di $\R^2$.

Untuk mencocokkan realisasi tertanam ini dengan struktur bundel, bagi setiap
\mathl{P\in Y}{}
pilih suatu parametrisasi lokal
\maabbdisp {\beta} {V} {U\subseteq Y
} {,}
dengan $V\subseteq\R^{n-1}$ terbuka. Pemetaan
\maabbeledisp {} {V \times \R^{n-1} } { U \times \R^n
} {(Q,z)} { ( \beta(Q), \left(  D\beta  \right)_{ Q } \left(  z  \right) )
} {,}
memberikan trivialisasi lokal yang citranya tepat terdiri atas pasangan titik dan vektor singgung pada serat. Jadi topologi serta proyeksi realisasi tertanam ini sama dengan topologi dan proyeksi bundel singgung yang telah didefinisikan.

Catatan edisi: Solusi sumber mengganti syarat $h(P)=c$ dengan $h(P)=0$, menulis turunan parsial dari fungsi $f$ yang tidak didefinisikan, dan memakai $\beta(q)$ alih-alih $\beta(Q)$. Lebih mendasar, asumsi $C^1$ sumber hanya membuat turunan pertama kontinu dan tidak cukup untuk menyimpulkan bahwa $TY$ adalah submanifold diferensiabel dari $W\times\R^n$. Edisi ini memperkuat asumsi soal menjadi $C^2$ dan memberikan pemeriksaan nilai reguler yang hilang.
\subsection*{Solusi Soal 10.25}
Misalkan $N=(0,0,1)$ dan gunakan invers proyeksi stereografis
\maabbeledisp {\sigma} {\R^2} {S^2\setminus\{N\}
} {(x,y)} {\left(\frac{2x}{1+x^2+y^2},\frac{2y}{1+x^2+y^2},\frac{x^2+y^2-1}{1+x^2+y^2}\right)
} {.}
Pada $S^2\setminus\{N\}$, dorong medan vektor konstan $e_1$ ke depan dan definisikan

\mathdisp {X(\sigma(x,y))=(D\sigma)_{(x,y)}(e_1)} { . }

Secara eksplisit,

\mathdisp {X(\sigma(x,y))=\frac{1}{(1+x^2+y^2)^2}\left(2(1-x^2+y^2),-4xy,4x\right)} { . }

Karena $\sigma$ merupakan difeomorfisme, diferensialnya injektif; jadi medan ini tidak mempunyai titik nol pada $S^2\setminus\{N\}$. Selain itu,

\mathdisp {\left\|X(\sigma(x,y))\right\|=\frac{2}{1+x^2+y^2}} { , }

yang menuju $0$ ketika $x^2+y^2\to\infty$, yakni ketika $\sigma(x,y)\to N$. Oleh sebab itu, medan tersebut dapat diperluas secara kontinu dengan menetapkan

\mavergleichskettedisp
{\vergleichskette
{X(N)
}
{ =} {0
}
{ \in }{T_NS^2
}
{ } {
}
{ } {
}
}
{}{}{.}
Dalam realisasi $TS^2\subseteq S^2\times\R^3$, rumus norma di atas langsung membuktikan kontinuitas pada $N$. Dengan demikian, $X$ adalah medan vektor kontinu pada $S^2$ dan satu-satunya titik nolnya ialah $N$.

Catatan edisi: Solusi sumber menyatakan medan $e_1/(x^2+y^2)$ pada seluruh $\R^2$, padahal rumus itu tidak terdefinisi di titik asal, lalu menyimpulkan kontinuitas di kutub hanya dari komponen koordinat. Konstruksi di atas memakai medan konstan yang terdefinisi di seluruh bidang dan memeriksa norma dorong-majunya di dalam realisasi ambien, sehingga perluasan di kutub benar-benar kontinu.

\part{Unit 11}
\chapter[Kuliah 11]{Kuliah 11: Produk Manifold}
\setcounter{section}{11}

  \zwischenueberschrift{Produk manifold}

\inputdefinition
{}
{

Misalkan
\mathkor {} {M} {dan} {N} {}
dua
\definitionsverweis {manifold diferensiabel}{}{}
dengan
\definitionsverweis {atlas-atlas}{}{}
\mathkor {} {(U_i,U_i',\alpha_i, i \in I)} {dan} {(V_j,V_j',\beta_j, j \in J)} {.}
Maka
\definitionsverweis {ruang produk}{}{}

\mathl{M \times N}{} dengan
\definitionsverweis {peta-peta}{}{}
\maabbdisp {\alpha_i \times \beta_j} {U_i \times V_j } { U_i' \times V_j'
} {}
\zusatzklammer {dengan

\mavergleichskettek
{\vergleichskettek
{  (i,j)
}
{ \in }{ I \times J
}
{  }{
}
{    }{
}
{   }{
}
}
{}{}{}
dan

\mavergleichskettek
{\vergleichskettek
{ U_i' \times V_j'
}
{ \subseteq }{ \R^m \times \R^n
}
{  }{
}
{    }{
}
{   }{
}
}
{}{}{}} {} {}
disebut \definitionswort {produk manifold}{}
\mathkor {} {M} {dan} {N} {.}

}

Objek ini memang merupakan suatu manifold; lihat
Soal 11.1.
Manifold produk berbentuk
\mathl{M \times \R}{} disebut juga \stichwort {silinder} {} di atas $M$. Jika

\mavergleichskette
{\vergleichskette
{ M
}
{ \subseteq }{ \R^\ell
}
{  }{
}
{    }{
}
{   }{
}
}
{}{}{}
dan

\mavergleichskette
{\vergleichskette
{ N
}
{ \subseteq }{ \R^k
}
{  }{
}
{    }{
}
{   }{
}
}
{}{}{}
merupakan submanifold-submanifold tertutup, maka manifold produk
\mathl{M \times N}{} merupakan submanifold tertutup dari

\mavergleichskette
{\vergleichskette
{ \R^\ell \times \R^k
}
{ = }{ \R^{\ell + k}
}
{  }{
}
{    }{
}
{   }{
}
}
{}{}{,}
lihat
Soal 11.2.



\inputfaktbeweis
{Produkt von Mannigfaltigkeiten/Abbildungseigenschaften/Fakt}
{Lema}
{}
{

\faktsituation {Misalkan
\mathkor {} {M} {dan} {N} {}
\definitionsverweis {manifold-manifold diferensiabel}{}{}
dan
\mathl{M \times N}{} merupakan
\definitionsverweis {produk}{}{} keduanya.}
\faktuebergang {Maka berlaku sifat-sifat berikut.}
\faktfolgerung {\aufzaehlungdrei{Kedua
\definitionsverweis {proyeksi}{}{}
\maabbeledisp {p_1} {M \times N} {M
} {(x,y)} { x
} {,}
dan
\maabbeledisp {p_2} {M \times N} {N
} {(x,y)} { y
} {,}
merupakan
\definitionsverweis {pemetaan diferensiabel}{}{.}
}{
\definitionsverweis {Ruang singgung}{}{}
di suatu titik

\mavergleichskette
{\vergleichskette
{ R
}
{ = }{ (P,Q)
}
{  }{
}
{    }{
}
{   }{
}
}
{}{}{}
secara kanonik teridentifikasi dengan

\mavergleichskette
{\vergleichskette
{ T_R(M \times N)
}
{ = }{ T_PM \times T_QN
}
{  }{
}
{    }{
}
{   }{
}
}
{}{}{.}
}{Misalkan $L$ suatu manifold diferensiabel lain. Maka suatu pemetaan
\maabbeledisp {\varphi \times \psi} { L } { M \times N
} { u } { (\varphi(u), \psi(u))
} {,}
diferensiabel jika dan hanya jika kedua
\definitionsverweis {pemetaan komponennya}{}{}
\mathkor {} {\varphi} {dan} {\psi} {}
diferensiabel.
}}
\faktzusatz {}
\faktzusatz {}

}
{

\teilbeweis {}{}{}
{(1). Dengan beralih ke peta, kita dapat mengasumsikan bahwa
\mathkor {} {M} {dan} {N} {}
merupakan
\definitionsverweis {himpunan-himpunan bagian terbuka}{}{}
masing-masing di
\mathkor {} {\R^m} {dan} {\R^n} {.}
Dalam hal ini, pemetaan tersebut merupakan suatu
\definitionsverweis {pembatasan}{}{}
dari
\definitionsverweis {proyeksi linear}{}{}
\maabb {} {\R^m \times \R^n } {\R^m
} {,}
yang
menurut Proposisi 45.3 (Analysis (Osnabrück 2021-2023))
\definitionsverweis {terdiferensialkan secara kontinu}{}{.}}
{}
\teilbeweis {}{}{}
{(2). Proyeksi-proyeksi diferensiabel
\maabb {p_1} {M \times N} {M
} {}
dan
\maabb {p_2} {M \times N} {N
} {}
menghasilkan
\definitionsverweis {pemetaan-pemetaan singgung}{}{}
linear
\maabb {T_R(p_1)} {T_R(M \times N)} {T_PM
} {}
dan
\maabb {T_R(p_2)} {T_R(M \times N)} {T_QN
} {}
sehingga secara keseluruhan menghasilkan pemetaan linear
\maabbdisp {T_R(p_1) \times T_R(p_2)} { T_R(M \times N)} {T_PM \times T_QN} {.}
Untuk membuktikan bijektivitas, kita dapat beralih ke peta dan mengasumsikan bahwa

\mavergleichskette
{\vergleichskette
{ M
}
{ \subseteq }{ \R^m
}
{  }{
}
{    }{
}
{   }{
}
}
{}{}{}
dan

\mavergleichskette
{\vergleichskette
{ N
}
{ \subseteq }{ \R^n
}
{  }{
}
{    }{
}
{   }{
}
}
{}{}{}
merupakan himpunan-himpunan bagian terbuka. Pemetaan ini kemudian menjadi bijeksi
\maabbdisp {p_1 \times p_2} {\R^{m+n} } { \R^m \times \R^n
} {.}}
{}
\teilbeweis {}{}{}
{(3). Untuk suatu titik tetap

\mavergleichskette
{\vergleichskette
{ A
}
{ \in }{ L
}
{  }{
}
{    }{
}
{   }{
}
}
{}{}{}
dengan menggunakan lingkungan peta dari $A$ dan dari
\mathkor {} {\varphi(A)} {dan} {\psi(A)} {}
kita dapat mereduksi keadaan ke kasus ketika ketiga manifold merupakan himpunan-himpunan terbuka dalam ruang-ruang Euklides. Jika kedua pemetaan terdiferensialkan secara kontinu, maka menurut
Soal 45.9 (Analysis (Osnabrück 2021-2023))
pemetaan keseluruhannya terdiferensialkan secara kontinu. Arah sebaliknya jelas.}
{}

}

\bild{ \begin{center}
\includegraphics[width=5.5cm]{ \bildeinlesungpng {Toroidal_coord} {png} }
\end{center}
\bildtext {} }

\bildlizenz { Toroidal coord.png } {} {DaveBurke} {Commons} {CC-by 2.5} {}

\inputbeispiel{}
{

\definitionsverweis {Produk}{}{}
dari
\definitionsverweis {lingkaran}{}{}
dengan dirinya sendiri, yaitu

\mavergleichskette
{\vergleichskette
{ M
}
{ = }{ S^1 \times S^1
}
{  }{
}
{    }{
}
{   }{
}
}
{}{}{,}
disebut \stichwort {torus} {.} Ini merupakan manifold berdimensi dua. Karena

\mavergleichskette
{\vergleichskette
{ S^1
}
{ \subset }{ \R^2
}
{  }{
}
{    }{
}
{   }{
}
}
{}{}{}
merupakan
\definitionsverweis {submanifold tertutup}{}{,}
torus dapat direalisasikan sebagai submanifold tertutup di

\mavergleichskette
{\vergleichskette
{ \R^2 \times \R^2
}
{ = }{ \R^4
}
{  }{
}
{    }{
}
{   }{
}
}
{}{}{.}
Namun, torus juga dapat direalisasikan sebagai submanifold tertutup di $\R^3$. Untuk itu, misalkan
\mathkor {} {r} {dan} {R} {}
bilangan-bilangan real positif dengan

\mavergleichskette
{\vergleichskette
{ 0
}
{ < }{ r
}
{ < }{ R
}
{    }{
}
{   }{
}
}
{}{}{.}
Maka himpunan

\mathdisp {\left\{ (x,y,z) \in \R^3  \mid   \left(  \sqrt{x^2+y^2} -R  \right)^2 +z^2 = r^2         \right\}} {  }

merupakan suatu torus. Dalam realisasi ini, torus merupakan permukaan
\zusatzklammer {yang digelembungkan} {} {}
\anfuehrung{ban dalam sepeda}{,} dengan \anfuehrung{jari-jari roda}{} sama dengan $R$ dan \anfuehrung{jari-jari tabung}{} sama dengan $r$
\zusatzklammer {roda terletak pada bidang \mathlk{x-y}{-}} {} {.}
Setelah setiap faktor lingkaran diidentifikasi melalui
\mathl{S^1 \cong \R/(2\pi\mathbb{Z})}{,}
hubungan dengan produk
\mathl{S^1 \times S^1}{} diperoleh dengan memetakan pasangan kelas sudut
\mathl{([ \varphi ],[ \psi ])}{} ke titik
\mathl{( (R +r \cos  \psi) \cos \varphi, (R+ r \cos  \psi ) \sin \varphi ,  r \sin  \psi )}{}.

}

  \zwischenueberschrift{Braket Lie medan vektor}

\inputdefinition
{}
{

Misalkan $M$ suatu
\definitionsverweis {manifold diferensiabel}{}{.}
Suatu pemetaan
\maabbdisp {D} {C^1(M,\R) } { C^0(M,\R)
} {}
disebut
\definitionswort {operator diferensial orde pertama}{,}
jika memenuhi sifat-sifat berikut.
\aufzaehlungzwei { $D$ bersifat
$\R$-\definitionsverweis {linear}{}{.}
} {Berlaku

\mavergleichskette
{\vergleichskette
{ D(fg)
}
{ = }{ fD(g)+gD(f)
}
{  }{
}
{    }{
}
{   }{
}
}
{}{}{.}
}

}

Operator diferensial orde pertama disebut juga \stichwort {derivasi} {.}

Untuk suatu himpunan terbuka

\mavergleichskette
{\vergleichskette
{ U
}
{ \subseteq }{ \R^n
}
{  }{
}
{    }{
}
{   }{
}
}
{}{}{,}
turunan-turunan parsial $\partial_i$ dan $\partial_j$, sebagai operasi pada fungsi-fungsi yang dua kali terdiferensialkan secara kontinu, dapat dipertukarkan; jadi

\mavergleichskettedisp
{\vergleichskette
{ \partial_i \circ \partial_j
}
{ =} { \partial_j \circ \partial_i
}
{ } {
}
{ } {
}
{ } {
}
}
{}{}{.}
Suatu kombinasi linear

\mavergleichskettedisp
{\vergleichskette
{ V
}
{ =} { \sum_{i = 1}^n f_i \partial_i
}
{ } {
}
{ } {
}
{ } {
}
}
{}{}{}
dengan koefisien-koefisien kontinu
\mathl{f_i \in C^0(U,\R)}{}
dapat dipandang sebagai
\definitionsverweis {medan vektor}{}{}
\maabbeledisp {} { U } { TU = U \times \R^n
} { P } { (P, \sum_{i = 1}^n f_i(P) \partial_i ) = (P,f_1(P) , \ldots , f_n(P) )
} {,}
dan juga sebagai operator turunan
\maabbeledisp {} {C^1(U,\R)} {C^0(U,\R)
} {h } { \sum_{i = 1}^n f_i  \partial_i  h
} {,}
\zusatzklammer {atau dari $C^\infty(U,\R)$ ke $C^\infty(U,\R)$ apabila
\mathl{f_i \in C^\infty(U,\R)}{} untuk setiap $i$} {} {.}
Dalam pengertian ini, medan vektor dan operator diferensial orde pertama saling bersesuaian. Lebih lanjut, operator-operator diferensial dapat dikomposisikan satu sama lain, dan ternyata urutannya penting.



\inputfaktbeweis
{Offene Menge/R^n/Vektorfelder/Hintereinanderschaltung/Explizit/Fakt}
{Lema}
{}
{

\faktsituation {Misalkan

\mavergleichskette
{\vergleichskette
{ U
}
{ \subseteq }{ \R^n
}
{  }{
}
{    }{
}
{   }{
}
}
{}{}{}
\definitionsverweis {terbuka}{}{}
dan misalkan $f_1 , \ldots , f_n, g_1 , \ldots , g_n,h$ merupakan fungsi-fungsi pada $U$ yang dua kali
\definitionsverweis {terdiferensialkan secara kontinu}{}{.}}
\faktfolgerung {Maka berlaku

\mavergleichskettedisphandlinks
{\vergleichskettedisphandlinks
{ \left(  \sum_{i = 1}^n f_i \partial_i  \right) \left(  \left(  \sum_{j = 1}^n g_j \partial_j   \right) (h)  \right)
}
{ =} { \sum_{j = 1}^n  \left(  \sum_{i = 1}^n  f_i \partial_i g_j  \right) \partial_j h + \sum_{j = 1}^n \sum_{i = 1}^n    f_i g_j \partial_i \partial_j h
}
{ } {
}
{ } {
}
{ } {
}
}
{}{}{.}}
\faktzusatz {}
\faktzusatz {}

}
{

Berlaku

\mavergleichskettealignhandlinks
{\vergleichskettealignhandlinks
{ \left(  \sum_{i = 1}^n f_i \partial_i  \right) \left(  \left(  \sum_{j = 1}^n g_j \partial_j   \right) (h)  \right)
}
{ =} { \left(  \sum_{i = 1}^n f_i \partial_i  \right) \left(  \sum_{j = 1}^n g_j \partial_j (h)  \right)
}
{ =} { \sum_{i = 1}^n f_i \partial_i \left(  \sum_{j = 1}^n g_j \partial_j (h)  \right)
}
{ =} { \sum_{i = 1}^n f_i \left(  \sum_{j = 1}^n \left(  \partial_i (g_j) \partial_j (h) +  g_j \partial_i \partial_j (h)  \right)  \right)
}
{ =} { \sum_{j = 1}^n \left(  \sum_{i = 1}^n f_i  \partial_i g_j  \right) \partial_j h +  \sum_{j = 1}^n \sum_{i = 1}^n    f_i g_j \partial_i \partial_j h
}
}
{}
{}{.}

}

Dari deskripsi eksplisit ini tampak bahwa kedua komposisi pada umumnya dapat berbeda; yakni, dapat terjadi

\mavergleichskettedisp
{\vergleichskette
{ \left(  \sum_{i = 1}^n f_i \partial_i  \right) \circ  \left(  \sum_{j = 1}^n g_j \partial_j   \right)
}
{ \neq} { \left(  \sum_{j = 1}^n g_j \partial_j   \right) \circ  \left(  \sum_{i = 1}^n f_i \partial_i  \right)
}
{ } {
}
{ } {
}
{ } {
}
}
{}{}{,}
karena meskipun suku kedua selalu sama, suku pertamanya pada umumnya dapat berbeda. Hal ini juga mempunyai konsekuensi berikut.


\inputfaktbeweis
{Offene Menge/R^n/Vektorfelder/Lie-Klammer/Vektorfeld/Fakt}
{Lema}
{}
{

\faktsituation {Misalkan

\mavergleichskette
{\vergleichskette
{ U
}
{ \subseteq }{ \R^n
}
{  }{
}
{    }{
}
{   }{
}
}
{}{}{}
\definitionsverweis {terbuka}{}{}
dan misalkan $f_1 , \ldots , f_n, g_1 , \ldots , g_n$ merupakan fungsi-fungsi pada $U$ yang dua kali
\definitionsverweis {terdiferensialkan secara kontinu}{}{}
dengan operator-operator diferensial terkait

\mavergleichskette
{\vergleichskette
{ D
}
{ = }{ \sum_{i = 1}^n f_i \partial_i
}
{  }{
}
{    }{
}
{   }{
}
}
{}{}{}
dan

\mavergleichskette
{\vergleichskette
{ E
}
{ = }{ \sum_{j = 1}^n g_j \partial_j
}
{  }{
}
{    }{
}
{   }{
}
}
{}{}{}}
\faktfolgerung {Maka berlaku

\mavergleichskettedisp
{\vergleichskette
{ D \circ E -E \circ D
}
{ =} { \sum_{j = 1}^n \left(  \sum_{i = 1}^n \left(  f_i  \partial_i g_j - g_i \partial_i f_j  \right)  \right)  \partial_j
}
{ } {
}
{ } {
}
{ } {
}
}
{}{}{.}
Secara khusus, ungkapan ini sendiri bersesuaian dengan suatu operator diferensial orde $1$, dan dengan demikian dengan suatu medan vektor.}
\faktzusatz {}
\faktzusatz {}

}
{

Menurut
Lema 11.5
berlaku

\mavergleichskettealignhandlinks
{\vergleichskettealignhandlinks
{D \circ E-E \circ D
}
{ =} { \left(  \sum_{i = 1}^n f_i \partial_i  \right) \circ  \left(  \sum_{j = 1}^n g_j \partial_j   \right) -  \left(  \sum_{j = 1}^n g_j \partial_j   \right) \circ  \left(  \sum_{i = 1}^n f_i \partial_i  \right)
}
{ =} { \sum_{j = 1}^n  \left(  \sum_{i = 1}^n  \left(  f_i  \partial_i g_j - g_i \partial_i f_j  \right)  \right)  \partial_j
}
{ } {
}
{ } {
}
}
{}
{}{.}

}

Pada suatu manifold, operator diferensial yang bersesuaian dengan medan vektor $V$ kita nyatakan dengan $D_V$. Untuk suatu fungsi diferensiabel
\maabb {f} { M } {\R
} {}
operator $D_V$ didefinisikan secara titik-demi-titik oleh

\mavergleichskette
{\vergleichskette
{ \left( D_V f \right)(P)
}
{ = }{ d f_P\left(V_P\right)
}
{ \in }{ \R
}
{    }{
}
{   }{
}
}
{}{}{,}
dengan $V_P=V(P)$ dan dengan identifikasi kanonik
\mathl{T_{f(P)}\R \cong \R}{.}
Dengan kata lain, untuk suatu titik

\mavergleichskette
{\vergleichskette
{ P
}
{ \in }{ M
}
{  }{
}
{    }{
}
{   }{
}
}
{}{}{}
pemetaan singgung penuh memenuhi

\mavergleichskettedisp
{\vergleichskette
{ \left( T(f) \circ V \right)(P)
}
{ =} { \left( f(P), \left(D_V f\right)(P) \right)
}
{ } {
}
{ } {
}
{ } {
}
}
{}{}{.}
Pernyataan yang bersesuaian berlaku pula untuk fungsi diferensiabel yang didefinisikan pada suatu himpunan bagian terbuka dari $M$.

\inputdefinition
{}
{

Misalkan
\mathkor {} {V} {dan} {W} {}
merupakan
\definitionsverweis {medan-medan vektor}{}{}
yang terdiferensialkan secara kontinu pada suatu
$C^2$-\definitionsverweis {manifold diferensiabel}{}{}
dengan operator-operator diferensial terkait $D_V$ dan $D_W$. Medan vektor yang bersesuaian dengan komposisi
\mathl{D_V \circ D_W -D_W \circ D_V}{} disebut
\definitionswort {braket Lie}{}
dari kedua medan vektor tersebut. Braket ini dinotasikan dengan $[V,W]$.

}

Urutan operator yang diberi tanda minus dalam definisi merupakan suatu konvensi.


\inputfaktbeweis
{Mannigfaltigkeit/Vektorfelder/Lie-Klammer/Eigenschaften/Fakt}
{Lema}
{}
{

\faktsituation {Misalkan $M$ suatu
$C^2$-\definitionsverweis {manifold diferensiabel}{}{.}}
\faktuebergang {Maka
\definitionsverweis {braket Lie}{}{}
dari
\definitionsverweis {medan-medan vektor}{}{}
memenuhi sifat-sifat berikut.}
\faktfolgerung {\aufzaehlungdrei{Berlaku

\mavergleichskette
{\vergleichskette
{ [V,W]
}
{ = }{ -[W,V]
}
{  }{
}
{    }{
}
{   }{
}
}
{}{}{.}
}{Berlaku

\mavergleichskette
{\vergleichskette
{ [V_1+V_2,W]
}
{ = }{ [V_1,W] + [V_2,W]
}
{  }{
}
{    }{
}
{   }{
}
}
{}{}{.}
}{Untuk suatu fungsi diferensiabel $f$ pada $M$ berlaku

\mavergleichskettedisp
{\vergleichskette
{ [fV,W]
}
{ =} { f [V,W] - \left(  D_Wf  \right)  V
}
{ } {
}
{ } {
}
{ } {
}
}
{}{}{.}
}}
\faktzusatz {}
\faktzusatz {}

}
{

(1) mengikuti dari definisi, dan (2) mengikuti dari hubungan

\mavergleichskettedisp
{\vergleichskette
{ D_{V_1 +V_2 }
}
{ =} { D_{V_1 } + D_{V_2 }
}
{ } {
}
{ } {
}
{ } {
}
}
{}{}{}
antara operator diferensial dan medan vektor. (3). Karena pernyataan tersebut bersifat lokal, kita dapat menetapkan

\mavergleichskettedisp
{\vergleichskette
{ V
}
{ =} { \sum_{i = 1}^n f_i \partial_i
}
{ } {
}
{ } {
}
{ } {
}
}
{}{}{}
dan

\mavergleichskette
{\vergleichskette
{ W
}
{ = }{ \sum_{j = 1}^n g_j \partial_j
}
{  }{
}
{    }{
}
{   }{
}
}
{}{}{.}
Maka

\mavergleichskettedisp
{\vergleichskette
{ fV
}
{ =} { f \sum_{i = 1}^n f_i \partial_i
}
{ =} { \sum_{i = 1}^n ff_i \partial_i
}
{ } {
}
{ } {
}
}
{}{}{}
dan menurut
Lema 11.6
berlaku

\mavergleichskettealignhandlinks
{\vergleichskettealignhandlinks
{ \, [fV,W]
}
{ =} { \sum_{j = 1}^n \left(  \sum_{i = 1}^n \left(  f f_i \partial_i g_j - g_i \partial_i (f f_j) \right)  \right) \partial_j
}
{ =} { \sum_{j = 1}^n \left(  \sum_{i = 1}^n \left(  f f_i \partial_i g_j - g_i \left(  f \partial_i ( f_j) +f_j \partial_i (f)  \right)  \right)  \right)  \partial_j
}
{ =} { f \sum_{j = 1}^n \left(  \sum_{i = 1}^n \left(  f_i \partial_i g_j - g_i \partial_i ( f_j)  \right)  \right) \partial_j - \sum_{j = 1}^n \left(  \sum_{i = 1}^n g_i f_j \partial_i (f)  \right)  \partial_j
}
{ =} { f[V,W] - \sum_{j = 1}^n \left(  \sum_{i = 1}^n g_i \partial_i (f)  \right) f_j \partial_j
}
}
{
\vergleichskettefortsetzungalign
{ =} { f[V,W] - \left(  \sum_{i = 1}^n g_i \partial_i (f)  \right) \left(  \sum_{j = 1}^n f_j \partial_j  \right)
}
{ =} { f[V,W] - D_W(f) V
}
{ } {}
{ } {}
}
{}{.}

}

\chapter{Lembar Kerja 11}
\setcounter{section}{11}

  \zwischenueberschrift{Soal latihan}

\inputaufgabe
{}
{

Tunjukkan bahwa
\definitionsverweis {produk}{}{}

\mathl{M \times N}{} dari dua
\definitionsverweis {manifold diferensiabel}{}{}
\mathkor {} {M} {dan} {N} {}
sendiri merupakan sebuah manifold diferensiabel.

}
{} {}

\inputaufgabe
{}
{

Misalkan
\mathkor {} {M_1 \subseteq N_1} {dan} {M_2 \subseteq N_2} {}
\definitionsverweis {submanifold tertutup}{}{.} Tunjukkan bahwa
\definitionsverweis {produk}{}{}
\mathl{M_1 \times M_2}{} merupakan sebuah submanifold tertutup dari
\mathl{N_1 \times N_2}{}.

}
{} {}

\inputaufgabe
{}
{

Deskripsikan peta-peta pada
\definitionsverweis {torus}{}{}

\mathl{S^1 \times S^1}{,} yang diperoleh dari proyeksi-proyeksi stereografik.

}
{} {}

\inputaufgabe
{}
{

Tunjukkan bahwa
\mathl{\R^2 \setminus \{0\}}{}
\definitionsverweis {difeomorfik}{}{}
dengan suatu produk dari
\definitionsverweis {manifold-manifold}{}{}
berdimensi satu.

}
{} {}

\inputaufgabe
{}
{

Tunjukkan bahwa
\definitionsverweis {produk}{}{}

\mathl{M \times N}{} dari dua
\definitionsverweis {manifold yang terhubung lintasan}{}{} dan
\definitionsverweis {manifold diferensiabel}{}{}
\mathkor {} {M} {dan} {N} {}
juga terhubung lintasan.

}
{} {}

\inputaufgabe
{}
{

Misalkan $M$ sebuah
\definitionsverweis {manifold diferensiabel}{}{}
dan
\maabbeledisp {\varphi} { M } { M \times M
} { x } {(x,x)
} {,}
merupakan
\definitionsverweis {pemetaan diagonal}{}{}
ke dalam
\definitionsverweis {produk}{}{}

\mathl{M \times M}{.} Tunjukkan bahwa
\definitionsverweis {diagonal}{}{}
$\varphi(M)$ merupakan sebuah
\definitionsverweis {submanifold tertutup}{}{}.

}
{} {}

\inputaufgabe
{}
{

Misalkan $M$ sebuah
\definitionsverweis {manifold diferensiabel}{}{.}
Tunjukkan bahwa pemetaan pertukaran
\maabbeledisp {} {M \times M} { M \times M
} {(P,Q)} {(Q,P)
} {,}
merupakan sebuah
\definitionsverweis {difeomorfisme}{}{}.

}
{} {}

\inputaufgabe
{}
{

Deskripsikan torus
\mathl{S^1 \times S^1}{} sebagai
\definitionsverweis {himpunan hasil rotasi}{}{}
di dalam $\R^3$.

}
{} {}

\inputaufgabe
{}
{

Misalkan

\mavergleichskette
{\vergleichskette
{ R
}
{ > }{ 0
}
{  }{
}
{    }{
}
{   }{
}
}
{}{}{}
dan perhatikan pemetaan
\maabbeledisp {} { \R^3} { \R
} { (x,y,z) } { \left(  \sqrt{x^2+y^2}- R  \right)^2 +z^2
} {.}
Pertama, tunjukkan bahwa pemetaan yang ditampilkan tidak diferensiabel pada seluruh sumbu $z$. Pada himpunan terbuka
$U=\{(x,y,z)\in\R^3\mid x^2+y^2>0\}$, tentukan
\definitionsverweis {titik-titik reguler}{}{}
dan lokus kritis dari pembatasannya, serta bentuk serat di atas

\mavergleichskette
{\vergleichskette
{ s
}
{ \in }{ \R
}
{  }{
}
{    }{
}
{   }{
}
}
{}{}{.}
Nyatakan pula apa yang terjadi untuk $s<0$. Bagaimana bentuk itu berubah ketika $s$ melewati $R^2$, yakni ketika beralih dari

\mavergleichskette
{\vergleichskette
{ \sqrt{s}
}
{ < }{ R
}
{  }{
}
{    }{
}
{   }{
}
}
{}{}{}
ke

\mavergleichskette
{\vergleichskette
{ \sqrt{s}
}
{ > }{ R
}
{  }{
}
{    }{
}
{   }{
}
}
{}{}{?}

}
{} {}

\inputaufgabegibtloesung
{}
{

Buatlah sebuah animasi yang mengilustrasikan
Soal 11.9.

}
{} {}

\inputaufgabe
{}
{

Definisikan pemetaan
\maabbdisp {} {S^1 \times S^1} { S^2
} {,}
yang, bagi suatu pasangan sudut
\mathl{(\alpha, \beta)}{}, menafsirkan komponen pertama sebagai titik pada khatulistiwa, lalu bergerak dari titik itu sepanjang lingkaran besar menuju utara sesuai dengan komponen kedua. Apakah pemetaan tersebut
\definitionsverweis {diferensiabel}{}{?}
Bagaimanakah bentuk serat-serat pemetaan itu?

}
{} {}

\inputaufgabe
{}
{

Berikan sebuah argumen heuristik bahwa permukaan bola satuan
\mathl{S^2}{} dan torus
\mathl{S^1 \times S^1}{} tidak
\definitionsverweis {homeomorfik}{}{}.

}
{} {}

\inputaufgabe
{}
{

Dengan
\definitionsverweis {manifold diferensiabel}{}{}
manakah
\mathl{S^1 \times S^1 \setminus \triangle}{,} yaitu torus tanpa
\definitionsverweis {diagonal}{}{,}
\definitionsverweis {difeomorfik}{}{?}

}
{} {}

\inputaufgabegibtloesung
{}
{

Misalkan $X$ sebuah
\definitionsverweis {torus}{}{.}
Berikan sebuah
\definitionsverweis {pemetaan diferensiabel}{}{}
surjektif
\maabbdisp {\varphi} {\R^2} {X
} {}
sedemikian sehingga
\definitionsverweis {pemetaan singgung}{}{}
\maabbdisp {T_P(\varphi)} { T_P\R^2 } { T_{\varphi(P)}X
} {}
juga surjektif di setiap titik

\mavergleichskette
{\vergleichskette
{ P
}
{ \in }{ \R^2
}
{  }{
}
{    }{
}
{   }{
}
}
{}{}{}.

}
{} {}

Dalam soal-soal berikut diperkenalkan produk relatif dari himpunan-himpunan dan ruang-ruang topologis.

\inputaufgabe
{}
{

Misalkan
\mathl{L_1,L_2,M}{} himpunan-himpunan dan
\maabbdisp {p_1} {L_1 } { M
} {}
serta
\maabbdisp {p_2} {L_2 } { M
} {}
pemetaan-pemetaan. Kita definisikan

\mavergleichskettedisp
{\vergleichskette
{L_1 \times_M L_2
}
{ \defeq} { \left\{ (x_1,x_2)  \mid   p_1(x_1) = p_2(x_2)         \right\}
}
{ \subseteq} { L_1 \times L_2
}
{ } {
}
{ } {
}
}
{}{}{.}
\aufzaehlungzwei {Tunjukkan bahwa terdapat sebuah diagram komutatif

\mathdisp {\begin{matrix} L_1 \times_M L_2  & \stackrel{  }{\longrightarrow} &  L_1  &  \\ \downarrow & & \downarrow  & \\  L_2  & \stackrel{  }{\longrightarrow} &  M & \! \! \! \!\!\!\!\!   \\ \end{matrix}} {  }

.
} {Misalkan $T$ sebuah himpunan lain dan
\maabb {\psi_1} { T } { L_1
} {}
serta
\maabb {\psi_2} { T } { L_2
} {}
pemetaan-pemetaan yang memenuhi

\mavergleichskettedisp
{\vergleichskette
{ p_1 \circ \psi_1
}
{ =} { p_2 \circ \psi_2
}
{ } {
}
{ } {
}
{ } {}
}
{}{}{.}
Tunjukkan bahwa terdapat tepat satu pemetaan
\maabbdisp {\psi} { T } { L_1 \times_M L_2
} {}
sedemikian sehingga proyeksinya ke
\mathkor {} {L_1} {dan} {L_2} {}
masing-masing berimpit dengan $\psi_1,\psi_2$.
}

}
{} {}

\inputaufgabe
{}
{

Misalkan
\mathl{L_1,L_2,M}{}
\definitionsverweis {ruang-ruang topologis}{}{}
dan
\maabbdisp {p_1} {L_1 } { M
} {}
serta
\maabbdisp {p_2} {L_2 } { M
} {}
\definitionsverweis {pemetaan-pemetaan kontinu}{}{.}
Kita definisikan

\mavergleichskettedisp
{\vergleichskette
{L_1 \times_M L_2
}
{ \defeq} { \left\{ (x_1,x_2)  \mid   p_1(x_1) = p_2(x_2)         \right\}
}
{ \subseteq} { L_1 \times L_2
}
{ } {
}
{ } {
}
}
{}{}{}
dengan
\definitionsverweis {topologi yang diinduksi dari ruang produk di atas}{}{.}
\aufzaehlungzwei {Tunjukkan bahwa terdapat sebuah diagram komutatif

\mathdisp {\begin{matrix} L_1 \times_M L_2  & \stackrel{  }{\longrightarrow} &  L_1  &  \\ \downarrow & & \downarrow  & \\  L_2  & \stackrel{  }{\longrightarrow} &  M & \! \! \! \!\!\!\!\!   \\ \end{matrix}} {  }

dengan pemetaan-pemetaan kontinu.
} {Misalkan $T$ sebuah ruang topologis lain dan
\maabb {\psi_1} { T } { L_1
} {}
serta
\maabb {\psi_2} { T } { L_2
} {}
pemetaan-pemetaan kontinu yang memenuhi

\mavergleichskettedisp
{\vergleichskette
{ p_1 \circ \psi_1
}
{ =} { p_2 \circ \psi_2
}
{ } {
}
{ } {
}
{ } {}
}
{}{}{.}
Tunjukkan bahwa terdapat tepat satu pemetaan kontinu
\maabbdisp {\psi} { T } { L_1 \times_M L_2
} {}
sedemikian sehingga proyeksinya ke
\mathkor {} {L_1} {dan} {L_2} {}
masing-masing berimpit dengan $\psi_1,\psi_2$.
}

}
{} {}

Dalam situasi kedua soal sebelumnya, $L_1 \times_M L_2$
\zusatzklammer {dengan proyeksi ke $L_2$} {} {}
juga dilambangkan dengan $p_2^*L_1$ atau dengan
\mathl{p_1^*L_2}{.} Gagasannya ialah bahwa objek
\maabb {p_1} {L_1} {M
} {}
ditarik balik sepanjang $p_2$ ke basis baru $L_2$.

\inputaufgabe
{}
{

Misalkan
\maabb {p} { Y } { X
} {}
sebuah
\definitionsverweis {pemetaan kontinu}{}{}
di antara
\definitionsverweis {ruang-ruang topologis}{}{}
\mathkor {} {X} {dan} {Y} {.}
Tunjukkan bahwa untuk
\definitionsverweis {produk relatif}{}{}
berlaku relasi

\mavergleichskettedisp
{\vergleichskette
{ X \times_X Y
}
{ \cong} { Y
}
{ } {
}
{ } {
}
{ } {
}
}
{}{}{}
\zusatzklammer {dengan
\maabb {} { X } { X
} {}
menyatakan identitas} {} {.}

}
{} {}

\inputaufgabe
{}
{

Misalkan
\maabb {p} { Y } { X
} {}
sebuah
\definitionsverweis {pemetaan kontinu}{}{}
di antara
\definitionsverweis {ruang-ruang topologis}{}{}
\mathkor {} {X} {dan} {Y} {.}
Misalkan

\mavergleichskette
{\vergleichskette
{ x
}
{ \in }{ X
}
{  }{
}
{    }{
}
{   }{
}
}
{}{}{}
sebuah titik dengan pemetaan yang bersesuaian
\maabbdisp {\iota} {\{x\}} { X
} {.}
Tunjukkan bahwa untuk
\definitionsverweis {produk relatif}{}{}
berlaku relasi

\mavergleichskettedisp
{\vergleichskette
{ \{x\}  \times_X Y
}
{ \cong} { Y_x
}
{ } {
}
{ } {
}
{ } {
}
}
{}{}{,}
yaitu produk relatif tersebut berimpit dengan
\definitionsverweis {serat}{}{}.

}
{} {}

\inputaufgabe
{}
{

Misalkan
\mathkor {} {X,U} {dan} {V} {}
\definitionsverweis {ruang-ruang topologis}{}{}
dan misalkan
\maabb {p_1} {X \times U} {X
} {}
serta
\maabb {p_2} {X \times V} {X
} {}
adalah proyeksi kanonik ke $X$. Tunjukkan bahwa untuk
\definitionsverweis {produk relatif}{}{}
berlaku relasi

\mavergleichskettedisp
{\vergleichskette
{ \left(  X \times U  \right) \times_X \left(  X \times V  \right)
}
{ \cong} { X \times U \times V
}
{ } {
}
{ } {
}
{ } {
}
}
{}{}{.}

}
{} {}

\inputaufgabe
{}
{

Misalkan
\maabb {f,g} {\R} {\R
} {}
\definitionsverweis {fungsi-fungsi kontinu}{}{.}
\aufzaehlungdrei{Tunjukkan bahwa untuk
\definitionsverweis {produk relatif}{}{}
\zusatzklammer {fungsi-fungsi tersebut tidak muncul secara eksplisit dalam notasi} {} {}

\mavergleichskettedisp
{\vergleichskette
{ \R \times_\R \R
}
{ =} { \left\{ (x,y) \in \R^2   \mid   f(x) = g(y)         \right\}
}
{ } {
}
{ } {
}
{ } {
}
}
{}{}{}
berlaku.
}{Sketsalah
\mathl{\R \times_\R \R}{} untuk beberapa pilihan fungsi
\mathl{(f,g)}{.}
}{Tentukan pada (2) apakah
\mathl{\R \times_\R \R}{}
memiliki
\definitionsverweis {titik-titik singular}{}{}.
}

}
{} {}

\inputaufgabe
{}
{

Misalkan
\mathkor {} {X,Y} {dan} {Z} {}
\definitionsverweis {ruang-ruang topologis}{}{}
dan misalkan
\maabb {\varphi} { Y } { X
} {}
serta
\maabb {p} { Z } { X
} {}
\definitionsverweis {pemetaan-pemetaan kontinu}{}{.}
Misalkan
\maabbdisp {p_Y} {Y \times_XZ } { Y
} {.}
Tunjukkan bahwa sebuah
\definitionsverweis {penampang kontinu}{}{}
\maabb {s} { Y } { Y \times_XZ
} {}
setara dengan sebuah pemetaan kontinu
\maabb {t} { Y } { Z
} {}
yang memenuhi

\mavergleichskette
{\vergleichskette
{ p \circ t
}
{ = }{ \varphi
}
{  }{
}
{    }{
}
{   }{
}
}
{}{}{.}

}
{} {}

\inputaufgabe
{}
{

Misalkan
\mathl{X,Y,Z}{} dan $X'$
\definitionsverweis {ruang-ruang topologis}{}{}
dan diberikan sebuah diagram komutatif dari
\definitionsverweis {pemetaan-pemetaan kontinu}{}{}

\mathdisp {\begin{matrix}  & Y& \stackrel{}{\longrightarrow} &Z  \\  &  & \searrow& \downarrow   \\  & X' & \stackrel{\varphi}{  \longrightarrow } & X    \end{matrix}} {  }

. Misalkan

\mavergleichskette
{\vergleichskette
{ Y'
}
{ = }{ \varphi^*Y
}
{ = }{ Y \times_X X'
}
{    }{
}
{   }{
}
}
{}{}{}
dan

\mavergleichskette
{\vergleichskette
{ Z'
}
{ = }{ \varphi^*Z
}
{ = }{ Z \times_X X'
}
{    }{
}
{   }{
}
}
{}{}{.}
Tunjukkan bahwa terdapat diagram komutatif

\mathdisp {\begin{matrix}  & Y' & \stackrel{}{\longrightarrow} & Z'  \\  &  & \searrow& \downarrow   \\  &  &  & X'    \end{matrix}} {  }

.

}
{} {}

\inputaufgabe
{}
{

Misalkan
\mathkor {} {X,Y,Z} {dan} {W} {}
\definitionsverweis {ruang-ruang topologis}{}{}
dan misalkan
\maabb {} { W } { X
} {}
serta

\mathdisp {Z \stackrel{\varphi}{\longrightarrow} Y\stackrel{\psi}{\longrightarrow} X} {  }

merupakan
\definitionsverweis {pemetaan-pemetaan kontinu}{}{.}
Tunjukkan bahwa
\definitionsverweis {produk relatif}{}{}
\zusatzklammer {atau, secara ekuivalen, tarik balik} {} {}
memenuhi isomorfisme-isomorfisme kanonik

\mavergleichskettedisp
{\vergleichskette
{ \varphi^* \left(  \psi^*W  \right)
}
{ \cong} { Z \times_Y \left(  Y \times_X W  \right)
}
{ \cong} { Z \times_XW
}
{ \cong} { ( \psi \circ \varphi)^*W
}
{ } {
}
}
{}{}{,}
yang pada representasi produk diberikan oleh
\mathl{(z,(\varphi(z),w))\longmapsto(z,w)}{.}

}
{} {}

\inputaufgabe
{}
{

Perhatikan
\definitionsverweis {lingkaran}{}{}
$S^1$. Definisikan sebuah \stichwort {struktur grup diferensiabel} {} pada $S^1$, yaitu sebuah unsur identitas
\mathl{P \in S^1}{,} sebuah
\definitionsverweis {pemetaan diferensiabel}{}{}
\maabbeledisp {\alpha} {S^1} {S^1
} { x } { \alpha (x)
} {,}
dan sebuah pemetaan diferensiabel
\maabbeledisp {} {T = S^1 \times S^1} {S^1
} {(x,y)} { \varphi(x,y)
} {,}
sedemikian sehingga $S^1$ menjadi sebuah
\definitionsverweis {grup komutatif}{}{}
dengan data tersebut.

}
{} {}

\inputaufgabe
{}
{

Perhatikan
\definitionsverweis {grup linear umum}{}{}

\mavergleichskette
{\vergleichskette
{G
}
{ = }{ \operatorname{GL}_{  n   } \! \left(  \R  \right)
}
{ \subseteq }{ \R^{n^2}
}
{    }{
}
{   }{
}
}
{}{}{}
sebagai
\definitionsverweis {submanifold terbuka}{}{}
dari $\R^{n^2}$. Definisikan sebuah \stichwort {struktur grup diferensiabel} {} pada $G$, yaitu sebuah unsur identitas
\mathl{P \in G}{,} sebuah
\definitionsverweis {pemetaan diferensiabel}{}{}
\maabbeledisp {\alpha} { G } { G
} { x } { \alpha (x)
} {,}
dan sebuah pemetaan diferensiabel
\maabbeledisp {\varphi} {G \times G} {G
} {(x,y)} { \varphi(x,y)
} {,}
sedemikian sehingga $G$ menjadi sebuah
\definitionsverweis {grup}{}{}
dengan data tersebut.

}
{} {}

\inputaufgabe
{}
{

Tunjukkan bahwa
\definitionsverweis {pemetaan}{}{}
\maabbeledisp {\varphi} {S^1 \times \R } { S^1 \times \R
} {(x_1 , x_2;t)} {(-x_1 , -x_2; -t)
} {,}
merupakan sebuah
\definitionsverweis {difeomorfisme}{}{}
\definitionsverweis {tanpa titik tetap}{}{} yang inversnya adalah dirinya sendiri.

}
{} {}

\inputaufgabe
{}
{

Misalkan

\mavergleichskette
{\vergleichskette
{ D
}
{ = }{ x^2 \partial_x +  \left(  x-y^2  \right) \partial_y
}
{  }{
}
{    }{
}
{   }{
}
}
{}{}{.}
Hitung

\mathdisp {D  \left(  4x^2-xy+5y^3  \right)} { . }

}
{} {}

\inputaufgabe
{}
{

Misalkan

\mavergleichskette
{\vergleichskette
{ D
}
{ = }{ \sum_{j = 1}^n g_j \partial_j
}
{  }{
}
{    }{
}
{   }{
}
}
{}{}{}
sebuah operator diferensial pada suatu himpunan terbuka

\mavergleichskette
{\vergleichskette
{ U
}
{ \subseteq }{ \R^n
}
{  }{
}
{    }{
}
{   }{
}
}
{}{}{.}
Penerapan $D$ pada fungsi-fungsi diferensiabel mana yang menghasilkan fungsi-fungsi koefisien $g_j$?

}
{} {}

\inputaufgabe
{}
{

Untuk suatu himpunan terbuka $U\subseteq\R^n$, misalkan fungsi-fungsi $g_j\in C^0(U)$ diberikan, dan definisikan

\mavergleichskette
{\vergleichskette
{ D
}
{ = }{ \sum_{j = 1}^n g_j \partial_j
}
{  }{
}
{    }{
}
{   }{
}
}
{}{}{}
dengan

\mavergleichskette
{\vergleichskette
{ U
}
{ \subseteq }{ \R^n
}
{  }{
}
{    }{
}
{   }{
}
}
{}{}{.}
Pada setiap fungsi diferensiabel, $D$ bekerja menurut rumus di atas. Tunjukkan bahwa $D$ merupakan sebuah operator diferensial orde pertama dalam pengertian
Definisi 11.4.

}
{} {}

\inputaufgabe
{}
{

Misalkan $M$ sebuah
\definitionsverweis {manifold diferensiabel}{}{}
dan $D$ sebuah
\definitionsverweis {medan vektor}{}{}
diferensiabel pada $M$. Tunjukkan bahwa

\mavergleichskette
{\vergleichskette
{ [ D,D]
}
{ =} { 0
}
{  }{
}
{    }{
}
{   }{
}
}
{}{}{.}

}
{} {}

\inputaufgabe
{}
{

Misalkan $M$ sebuah
$C^3$-\definitionsverweis {manifold diferensiabel}{}{.}
Tunjukkan bahwa untuk
$C^2$-\definitionsverweis {medan-medan vektor}{}{}
$D,E,F$ pada $M$ berlaku apa yang disebut \stichwort {identitas Jacobi} {},

\mavergleichskettedisp
{\vergleichskette
{ [ [D,E],F] + [ [E,F],D] + [ [F,D],E]
}
{ =} { 0
}
{ } {
}
{ } {
}
{ } {
}
}
{}{}{.}

}
{} {}

Misalkan $V$ sebuah
\definitionsverweis {ruang vektor}{}{}
atas suatu
\definitionsverweis {lapangan}{}{}
$K$, dan misalkan
\maabbeledisp {} {V \times V} { V
} {(f,g)} { [f,g]
} {,}
sebuah
\definitionsverweis {operasi biner}{}{}
pada $V$. Pasangan
\mathl{(V, [-,-])}{} disebut sebuah
\definitionswort {aljabar Lie}{,}
jika ketiga syarat berikut dipenuhi.
\aufzaehlungdrei{ $[-,-]$ bersifat
\definitionsverweis {bilinear}{}{.}
}{ $[-,-]$ memenuhi apa yang disebut \stichwort {identitas Jacobi} {,} yaitu untuk

\mavergleichskette
{\vergleichskette
{ u,v,w
}
{ \in }{ V
}
{  }{
}
{    }{
}
{   }{
}
}
{}{}{}
berlaku

\mavergleichskettedisp
{\vergleichskette
{ [ [u,v],w] + [ [v,w],u] + [ [w,u],v]
}
{ =} { 0
}
{ } {
}
{ } {
}
{ } {
}
}
{}{}{.}
}{Berlaku

\mavergleichskette
{\vergleichskette
{ [v,v]
}
{ = }{ 0
}
{  }{
}
{    }{
}
{   }{
}
}
{}{}{}
untuk setiap

\mavergleichskette
{\vergleichskette
{ v
}
{ \in }{ V
}
{  }{
}
{    }{
}
{   }{
}
}
{}{}{.}
}

\inputaufgabe
{}
{

Misalkan $M$ sebuah
$C^\infty$-\definitionsverweis {manifold}{}{}
dan misalkan $V$ ruang semua
$C^\infty$-\definitionsverweis {medan vektor}{}{}
pada $M$. Tunjukkan bahwa $V$, dengan
\definitionsverweis {braket Lie}{}{}
medan-medan vektor, merupakan sebuah
\definitionsverweis {aljabar Lie}{}{}.

}
{} {}

\inputaufgabe
{}
{

Misalkan $V$ sebuah
$K$-\definitionsverweis {ruang vektor}{}{}
atas suatu
\definitionsverweis {lapangan}{}{}
$K$ dan misalkan

\mavergleichskettedisp
{\vergleichskette
{ L
}
{ =} { \operatorname{End}_{  } \, ( V )
}
{ } {
}
{ } {
}
{ } {
}
}
{}{}{}
adalah himpunan semua
$K$-\definitionsverweis {endomorfisme}{}{}
linear dari $V$ ke $V$. Tunjukkan bahwa $L$, yang dilengkapi dengan

\mavergleichskettedisp
{\vergleichskette
{ [f,g]
}
{ \defeq} { f \circ g -g \circ f
}
{ } {
}
{ } {
}
{ } {
}
}
{}{}{,}
merupakan sebuah
\definitionsverweis {aljabar Lie}{}{}.

}
{} {}

  \zwischenueberschrift{Soal untuk dikumpulkan}

\inputaufgabe
{5}
{

Misalkan

\mavergleichskette
{\vergleichskette
{ 0
}
{ < }{ r
}
{ < }{ R
}
{    }{
}
{   }{
}
}
{}{}{}
Identifikasikan $S^1 \cong \R/(2\pi\Z)$ melalui kelas-kelas sudut. Dengan menuliskan unsur-unsurnya sebagai $[\varphi]$ dan $[\psi]$, misalkan

\mavergleichskettedisp
{\vergleichskette
{ T
}
{ =} { \left\{ (x,y,z) \in \R^3  \mid   \left(  \sqrt{x^2+y^2} -R  \right)^2 +z^2 = r^2         \right\}
}
{ } {
}
{ } {
}
{ } {
}
}
{}{}{.}
Tunjukkan bahwa pemetaan
\maabbeledisp {} { S^1 \times S^1 } { T
} { ( [\varphi], [\psi]) } {( (R +r \cos  \psi) \cos \varphi, (R+ r \cos  \psi ) \sin \varphi ,  r \sin  \psi )
} {}
merupakan sebuah
\definitionsverweis {bijeksi}{}{}.

}
{} {}

\inputaufgabe
{6}
{

Misalkan $T$ sebuah
\definitionsverweis {torus}{}{}
dan misalkan
\mathl{P,Q \in T}{} dua titik. Tunjukkan bahwa terdapat suatu lingkungan koordinat bersama
\mathl{P,Q \in U \subseteq T}{} sedemikian sehingga pemetaan koordinat
\maabbdisp {\alpha} { U } { V
} {}
merupakan sebuah
\definitionsverweis {homeomorfisme}{}{}
dengan
\mathl{V= {]0,1[} \times  {]0,1[}}{}.

}
{} {}

\inputaufgabe
{3}
{

Tunjukkan bahwa
\definitionsverweis {produk relatif}{}{}
$\R \times_\R \R$ untuk fungsi-fungsi yang terdiferensialkan secara kontinu,
\maabbdisp {f,g} {\R} {\R
} {}
tidak harus merupakan sebuah
\definitionsverweis {manifold topologis}{}{}.

}
{} {}

\inputaufgabe
{2}
{

Misalkan

\mavergleichskettedisp
{\vergleichskette
{ D
}
{ =} { \left(  xy^2-  \sin z  \right)  \partial_x +  \left(  e^{x+y+z}  \right) \partial_y+  \left(  xy^3z  \right) \partial_z
}
{ } {
}
{ } {
}
{ } {
}
}
{}{}{.}
Hitung

\mathdisp {D  \left(  7e^{xy}-  \cos \left( z^3+y^5 \right) -x^4y^3 + y^2z^3e^{z}  \right)} { . }

}
{} {}

\inputaufgabe
{2}
{

Misalkan

\mavergleichskette
{\vergleichskette
{ W
}
{ \subseteq }{ \R^n
}
{  }{
}
{    }{
}
{   }{
}
}
{}{}{}
\definitionsverweis {terbuka}{}{,}
\maabb {h} { W } { \R
} {}
sebuah
\definitionsverweis {fungsi yang terdiferensialkan secara kontinu}{}{}
dan

\mavergleichskette
{\vergleichskette
{ Y
}
{ = }{ h^{-1}(0)
}
{  }{
}
{    }{
}
{   }{
}
}
{}{}{}
adalah
\definitionsverweis {serat}{}{}
di atas $0$, dengan $h$
\definitionsverweis {reguler}{}{}
di setiap titik $Y$ dan karena itu mendefinisikan sebuah manifold diferensiabel. Misalkan $D$ sebuah operator diferensial orde pertama yang bersesuaian dengan suatu medan vektor kontinu pada $Y$. Tunjukkan bahwa

\mavergleichskettedisp
{\vergleichskette
{ D(gh)
}
{ =} { 0
}
{ } {
}
{ } {
}
{ } {
}
}
{}{}{}
untuk setiap fungsi diferensiabel $g$ pada $Y$.

}
{} {}

\inputaufgabe
{4 (1+1+1+1)}
{

Pada $\R^3$, perhatikan
\definitionsverweis {medan-medan vektor}{}{}
berikut,

\mavergleichskettedisp
{\vergleichskette
{ D
}
{ =} { -y \partial_x  + z\partial_y -x \partial_z
}
{ } {
}
{ } {
}
{ } {
}
}
{}{}{,}

\mavergleichskettedisp
{\vergleichskette
{ E
}
{ =} { 5 \partial_x - 3\partial_y +  \cos \left(  xy \right)  \partial_z
}
{ } {
}
{ } {
}
{ } {
}
}
{}{}{,}
dan

\mavergleichskettedisp
{\vergleichskette
{ F
}
{ =} { 4e^{xyz} \partial_y -z^3 \partial_z
}
{ } {
}
{ } {
}
{ } {
}
}
{}{}{.}
\aufzaehlungvier{Hitung
\mathl{[D,E]}{.}
}{Hitung
\mathl{[D,F]}{.}
}{Hitung
\mathl{[E,F]}{.}
}{Hitung
\mathl{[ [D,E], F]}{.}
}

}
{} {}

\section*{Solusi yang disediakan oleh sumber}
\addcontentsline{toc}{section}{Solusi yang disediakan oleh sumber}
\subsection*{Solusi Soal 11.10}
Tuliskan $\rho=\sqrt{x^2+y^2}$. Fungsi
$F(x,y,z)=(\rho-R)^2+z^2$ tidak diferensiabel ketika $\rho=0$, yaitu pada seluruh sumbu $z$. Pada himpunan terbuka $U=\{\rho>0\}$ berlaku
\mathdisp {dF=2(\rho-R)\left(\frac{x}{\rho}\,dx+\frac{y}{\rho}\,dy\right)+2z\,dz} { . }
Jadi, lokus kritis pembatasan $F\vert_U$ adalah lingkaran yang diberikan oleh $\rho=R$ dan $z=0$; semua titik lain di $U$ merupakan titik reguler. Titik-titik pada sumbu $z$ bukanlah titik kritis, sebab $F$ sama sekali tidak diferensiabel di sana.

Untuk $s<0$, seratnya kosong. Untuk $s=0$, seratnya adalah lingkaran kritis tersebut. Jika $0<s<R^2$, seratnya berupa torus cincin yang mulus. Pada $s=R^2$, lubangnya menutup dan diperoleh torus tanduk yang bertemu sumbu $z$ di titik asal, tempat $F$ tidak diferensiabel. Jika $s>R^2$, seratnya berupa permukaan berbentuk torus gelendong yang bertemu sumbu $z$ di dua titik; pada kedua titik itu $F$ juga tidak diferensiabel.

Animasi di samping memperlihatkan, sebagai contoh untuk $R=2$, serat-serat bagi berbagai nilai $s$ dalam soal ini. Peralihan terjadi pada $s=4=R^2$, atau secara ekuivalen $\sqrt{s}=R$.
\href{https://commons.wikimedia.org/wiki/File:Aufgabe79.27.gif}{Animasi yang memperlihatkan serat-serat fungsi torus pada Soal 11.10 untuk berbagai nilai s.}
Oleh Lukas Freudenberg.
\subsection*{Solusi Soal 11.14}
Kita tuliskan torus sebagai $X=S^1 \times S^1$ dengan lingkaran satuan
\mathl{S^1= \left\{ (x,y) \in \R^2  \mid   \betrag { (x,y) } = 1        \right\}}{.}
Pemetaan
\maabbeledisp {\varphi} {\R^2} {S^1 \times S^1
} {(u,v)} {( \cos  u, \sin  u , \cos  v , \sin  v )
} {,}
bersifat surjektif karena kedua komponennya merupakan parametrisasi standar lingkaran satuan. Kesurjektifan pemetaan singgung juga dapat diperiksa per komponen karena ruang singgung suatu manifold produk merupakan produk dari ruang-ruang singgungnya. Turunan pemetaan
\maabbeledisp {} {\R} { \R^2
} { t } { ( \cos  t, \sin  t)
} {,}
adalah
\mathl{(- \sin  t, \cos  t) \neq (0,0)}{,} sehingga vektor citra ini selalu merentang ruang singgung satu dimensi dari lingkaran satuan pada titik citranya.

\part{Unit 12}
\chapter[Kuliah 12]{Kuliah 12: Bundel Vektor}
\setcounter{section}{12}

  \zwischenueberschrift{Bundel vektor}

Misalkan $M$ suatu
\definitionsverweis {manifold diferensiabel}{}{}
dengan
\definitionsverweis {bundel tangen}{}{}
\maabb {p} {TM} {M
} {.}
Untuk suatu domain peta

\mavergleichskette
{\vergleichskette
{ U
}
{ \subseteq }{ M
}
{  }{
}
{    }{
}
{   }{
}
}
{}{}{}
dan suatu peta
\maabbdisp {\alpha} { U } { V \subseteq \R^n
} {}
berdasarkan
Lema 10.8
terdapat suatu
\definitionsverweis {homeomorfisme}{}{}
\maabbdisp {T(\alpha)} {TM\vert_U \cong TU} { TV \cong V \times \R^n
} {,}
yang linear pada komponen kedua untuk setiap titik dasar

\mavergleichskette
{\vergleichskette
{ P
}
{ \in }{ U
}
{  }{
}
{    }{
}
{   }{
}
}
{}{}{}
. Untuk peta lain
\maabbdisp {\alpha'} {U'} {V'
} {}
pemetaan transisinya
\maabbdisp {} {TV \vert_{\alpha(U \cap U') }  \cong \alpha(U \cap U') \times \R^n } { TV' \vert_{\alpha'(U \cap U') }  \cong \alpha'(U \cap U') \times \R^n
} {}
linear pada serat-seratnya. Melalui homeomorfisme

\mavergleichskette
{\vergleichskette
{ U
}
{ \cong }{ V
}
{  }{
}
{    }{
}
{   }{
}
}
{}{}{}
kita juga langsung memperoleh homeomorfisme

\mavergleichskettedisp
{\vergleichskette
{ TM \vert_U \cong TU
}
{ \cong} { U \times \R^n
}
{ } {
}
{ } {
}
{ } {
}
}
{}{}{}
di atas $U$ beserta pemetaan transisi terkait
\maabbdisp {} {(U \cap U') \times \R^n } { (U \cap U') \times \R^n
} {,}
yang linear terhadap $\R^n$.

Bundel tangen merupakan contoh penting dari bundel vektor.

\inputdefinition
{}
{

Misalkan $X$ suatu
\definitionsverweis {ruang topologis}{}{}
dan

\mavergleichskette
{\vergleichskette
{r
}
{ \in }{\N
}
{  }{
}
{    }{
}
{   }{
}
}
{}{}{.}
Sebuah
\definitionswort {bundel vektor riil}{}
$V$ ber-
\definitionswort {rank}{}
$r$ adalah suatu ruang topologis yang dilengkapi dengan
\definitionsverweis {pemetaan kontinu}{}{}
\maabb {p} { V } { X
} {}
sedemikian sehingga setiap
\definitionsverweis {serat}{}{}
$p^{-1}(x)$ merupakan
\definitionsverweis {ruang vektor riil}{}{}
yang
\definitionsverweis {berdimensi}{}{}
$r$ dan terdapat suatu
\definitionsverweis {penutup terbuka}{}{}

\mavergleichskette
{\vergleichskette
{X
}
{ = }{ \bigcup_{i \in I} U_i
}
{  }{
}
{    }{
}
{   }{
}
}
{}{}{}
beserta
\definitionsverweis {homeomorfisme-homeomorfisme}{}{}
\maabbdisp {\varphi_i} {p^{-1} \left(  U_i  \right) } { U_i \times \R^r
} {}
di atas $U_i$ yang pada setiap serat menginduksi suatu
\definitionsverweis {isomorfisme linear}{}{}
\maabbdisp {\left(  \varphi_i  \right)_x} {p^{-1} (x) } { \R^r
} {}
.

}

Dalam hal ini, $V$ juga disebut \stichwort {ruang total} {} dan $X$ disebut \stichwort {ruang dasar} {} dari bundel vektor tersebut. Untuk serat sering pula ditulis

\mavergleichskette
{\vergleichskette
{ V_x
}
{ = }{ p^{-1} (x)
}
{  }{
}
{    }{
}
{   }{
}
}
{}{}{.}

Dalam homeomorfisme
\maabb {} { p^{-1}(U) } { U \times \R^r
} {}
sisi kanan dilengkapi dengan
\definitionsverweis {topologi produk}{}{}
dan $\R^r$ dengan
\definitionsverweis {topologi Euklides}{}{}
alaminya, sedangkan
\mathl{p^{-1}(U)}{} dilengkapi dengan topologi yang diinduksi dari $V$. Dengan demikian, semua serat $V_x$ menyandang topologi alami suatu ruang vektor riil berdimensi hingga. Yang dimaksud dengan homeomorfisme di atas $U$ adalah bahwa diagram

\mathdisp {\begin{matrix} p^{-1}(U)   & \stackrel{  }{\longrightarrow} &  U \times \R^r   & \\ &  \searrow & \downarrow & \\ & &  U   & \!\!\!\!\! \!\!\!  \\ \end{matrix}} {  }

komutatif. Produk
\mathl{X \times \R^r}{} merupakan suatu bundel vektor yang disebut \stichwort {bundel vektor trivial} {.} Suatu homeomorfisme linear
\maabb {} { p^{-1}(U) } { U \times \R^r
} {}
disebut
\zusatzklammer {lokal} {} {}
\stichwort {trivialisasi} {.} Pembatasan suatu bundel vektor pada $U_i$ bersifat trivial; jadi, secara lokal setiap bundel vektor bersifat trivial. Bundel vektor ber-rank $1$ juga disebut \stichwort {bundel garis} {.}

\inputbemerkung
{}
{

Suatu
\definitionsverweis {bundel vektor riil}{}{}
\maabb {p} { V } { X
} {}
di atas
\definitionsverweis {ruang topologis}{}{}
$X$ juga dapat didefinisikan hanya dengan mengacu pada suatu penutup terbuka

\mavergleichskettedisp
{\vergleichskette
{ X
}
{ =} { \bigcup_{i \in I} U_i
}
{ } {
}
{ } {
}
{ } {
}
}
{}{}{}
dan trivialisasi-trivialisasi
\maabbdisp {\varphi_i} {p^{-1} \left(  U_i  \right) } { U_i \times \R^r
} {}
dengan mensyaratkan agar pemetaan transisi
\zusatzklammer {yang diperoleh di atas $p^{-1} (U_i \cap U_j)$} {} {}
\maabbdisp {\varphi_j \circ \varphi_i^{-1}} {U_i \cap U_j \times \R^r } { U_i \cap U_j \times \R^r
} {}
linear pada komponen ruang vektornya. Hal ini menghasilkan struktur ruang vektor pada setiap serat, dengan menggunakan sebarang trivialisasi.

}



\inputfaktbeweis
{Topologischer Raum/Reelles Vektorbündel/Einschränkung/Fakt}
{Lema}
{}
{

\faktsituation {Misalkan
\maabb {p} { V } { X
} {}
suatu
\definitionsverweis {bundel vektor riil}{}{}
di atas
\definitionsverweis {ruang topologis}{}{}
$X$.}
\faktfolgerung {Maka untuk setiap
\definitionsverweis {himpunan terbuka}{}{}

\mavergleichskette
{\vergleichskette
{ W
}
{ \subseteq }{ X
}
{  }{
}
{    }{
}
{   }{
}
}
{}{}{}
\definitionsverweis {pembatasan}{}{}
\maabbdisp {} {p^{-1} (W)} {W
} {}
juga merupakan suatu bundel vektor.}
\faktzusatz {}
\faktzusatz {}

}
{

Hal ini jelas; kita cukup membatasi homeomorfisme-homeomorfisme yang linear pada serat
\maabbdisp {\varphi_i} {p^{-1} \left(  U_i  \right) } { U_i \times \R^r
} {}
menjadi
\maabbdisp {\varphi_i \vert_{W \cap U_i }} {p^{-1} \left(  W \cap U_i  \right) } { \left(  W \cap U_i  \right)  \times \R^r
} {}
.

}

\inputdefinition
{}
{

Misalkan
\mathkor {} {E} {dan} {F} {}
\definitionsverweis {bundel-bundel vektor riil}{}{}
di atas suatu
\definitionsverweis {ruang topologis}{}{}
$X$. Suatu
\definitionswort {homomorfisme bundel vektor}{}
\maabb {\varphi} { E } { F
} {}
adalah
\definitionsverweis {pemetaan kontinu}{}{}
di atas $X$ sedemikian sehingga untuk setiap titik

\mavergleichskette
{\vergleichskette
{ x
}
{ \in }{ X
}
{  }{
}
{    }{
}
{   }{
}
}
{}{}{}
pemetaan terinduksi
\maabbdisp {\varphi_x} { E_x} {F_x
} {}
bersifat
$\R$-\definitionsverweis {linear}{}{}
.

}

Pengamatan berikut memungkinkan konstruksi berbagai macam bundel vektor.

\inputbemerkung
{}
{

Misalkan $X$ suatu
\definitionsverweis {ruang topologis}{}{}
dan misalkan

\mavergleichskette
{\vergleichskette
{ f_1 , \ldots , f_n
}
{ \in }{ C^0(X,\R)
}
{  }{
}
{    }{
}
{   }{
}
}
{}{}{}
merupakan
\definitionsverweis {fungsi-fungsi kontinu}{}{}
bernilai riil pada $X$. Fungsi-fungsi ini mendefinisikan suatu pemetaan
\maabbeledisp {} {X \times \R^n } { X \times \R
} {(P,t_1 , \ldots , t_n) } { (P, \sum_{i= 1}^n f_i (P) t_i)
} {,}
antara
\definitionsverweis {bundel vektor}{}{}
trivial ber-rank $n$ di atas $X$ dan bundel vektor trivial ber-rank $1$ di atas $X$. Pemetaan ini kontinu, dan di atas setiap titik

\mavergleichskette
{\vergleichskette
{ P
}
{ \in }{ X
}
{  }{
}
{    }{
}
{   }{
}
}
{}{}{}
terdapat
\definitionsverweis {bentuk linear}{}{}
\maabb {} {\R^n } { \R
} {}
dengan matriks baris
\mathl{\left( f_1(P)  , \,   \ldots  , \,  f_n(P)                 \right)}{}; oleh karena itu, pemetaan ini merupakan suatu
\definitionsverweis {homomorfisme}{}{}
bundel vektor. Melalui kernel tersebut, pada setiap serat $\R^n$ diperoleh suatu subruang vektor berdimensi $n-1$, kecuali apabila semua fungsi $f_i$ lenyap secara serentak di titik $P$. Misalkan

\mavergleichskette
{\vergleichskette
{ D(f_i)
}
{ = }{ \left\{ P \in X  \mid  f_i(P) \neq 0        \right\}
}
{  }{
}
{    }{
}
{   }{
}
}
{}{}{,}
maka

\mavergleichskette
{\vergleichskette
{ U
}
{ = }{ \bigcup_{i = 1}^n D(f_i)
}
{ \subseteq }{ X
}
{    }{
}
{   }{
}
}
{}{}{}
adalah himpunan bagian terbuka tempat sekurang-kurangnya satu $f_i$ tidak lenyap, sehingga kernel-kernel tersebut selalu berdimensi $n-1$. \stichwort {Bundel kernel} {} $E$ yang terkait dengan fungsi-fungsi $f_i$ adalah bundel vektor di atas $U$ yang ditentukan serat demi serat oleh kernel-kernel itu, yaitu

\mavergleichskettedisp
{\vergleichskette
{ E
}
{ =} { \left\{  (P; t_1 , \ldots , t_n)  \mid   P \in U , \,  \sum_{i= 1}^n f_i(P)t_i = 0       \right\}
}
{ \subseteq} { U \times \R^n
}
{ } {
}
{ } {
}
}
{}{}{.}
Berdasarkan
Soal 12.26
terdapat trivialisasi-trivialisasi di atas $D(f_i)$, dan objek ini memang merupakan suatu bundel vektor.

}

\inputbeispiel{}
{

Di atas $\R^2$, kita meninjau pemetaan
\maabbeledisp {} {\R^2 \times \R^2} { \R^2 \times \R
} {(x,y;s,t)} { (x,y;xs+yt)
} {,}
sebagai suatu
\definitionsverweis {homomorfisme bundel vektor}{}{}
dalam pengertian
Catatan 12.5.
Satu-satunya nol bersama kedua fungsi koordinat $x,y$ adalah titik asal $(0,0)$, sehingga

\mavergleichskette
{\vergleichskette
{ U
}
{ = }{ \R^2 \setminus \{ (0,0) \}
}
{  }{
}
{    }{
}
{   }{
}
}
{}{}{}
dan bundel kernel terkait adalah

\mavergleichskettedisp
{\vergleichskette
{ E
}
{ =} { \left\{  (P; s ,t)   \mid   P \in U  , \,  xs + yt = 0       \right\}
}
{ \subseteq} { U \times \R^2
}
{ } {
}
{ } {
}
}
{}{}{.}
Bundel ini terdiri atas suatu keluarga garis di dalam suatu bidang, dengan garis-garis tersebut bergerak bersama titik dasar

\mavergleichskette
{\vergleichskette
{ P
}
{ = }{ (x,y)
}
{  }{
}
{    }{
}
{   }{
}
}
{}{}{}
.

}

\inputbeispiel{}
{

Misalkan

\mavergleichskette
{\vergleichskette
{ W
}
{ \subseteq }{ \R^n
}
{  }{
}
{    }{
}
{   }{
}
}
{}{}{}
\definitionsverweis {terbuka}{}{,}
\maabb {h} { W } { \R
} {}
suatu fungsi yang dua kali
\definitionsverweis {terdiferensialkan secara kontinu}{}{}
dan

\mavergleichskette
{\vergleichskette
{ Y
}
{ = }{ h^{-1}(c)
}
{  }{
}
{    }{
}
{   }{
}
}
{}{}{}
merupakan
\definitionsverweis {serat}{}{}
di atas

\mavergleichskette
{\vergleichskette
{ c
}
{ \in }{ \R
}
{  }{
}
{    }{
}
{   }{
}
}
{}{}{,}
dengan $h$
\definitionsverweis {reguler}{}{}
di setiap titik dari $Y$. \definitionsverweis {Ruang singgung}{}{}
$T_PY$ pada

\mavergleichskette
{\vergleichskette
{ P
}
{ \in }{ Y
}
{  }{
}
{    }{
}
{   }{
}
}
{}{}{}
adalah kernel
\definitionsverweis {diferensial total}{}{}
$\left(D h \right)_{ P}$, yang diberikan oleh turunan-turunan parsial ${ \frac{   \partial   h   }{  \partial   x_1   } } ( P) , \ldots , { \frac{   \partial   h   }{  \partial   x_n   } } ( P)$ dan karenanya merupakan subruang vektor berdimensi $(n-1)$ dari $\R^n$; lihat juga
Teorema 10.3.
Kita meninjau pemetaan
\maabbeledisp {} {Y \times \R^n } { Y \times \R
} {(P;t_1 , \ldots , t_n) } { (P; \sum_{i= 1}^n { \frac{   \partial   h   }{  \partial   x_i   } } ( P)  t_i )
} {}
dan bundel kernel terkait

\mavergleichskettedisp
{\vergleichskette
{ E
}
{ =} { \left\{  (P; t_1 , \ldots , t_n)  \mid   P \in Y , \,  \sum_{i= 1}^n { \frac{   \partial   h   }{  \partial   x_i   } } ( P) t_i = 0       \right\}
}
{ \subseteq} { Y \times \R^n
}
{ } {
}
{ } {
}
}
{}{}{.}
dalam pengertian
Catatan 12.5.
Bundel ini secara langsung berimpit, serat demi serat, dengan ruang-ruang singgung pada $Y$. Namun, bundel ini memang merupakan bundel tangen, sebagaimana diperoleh dengan membandingkannya dengan
Soal 10.22.

}

\inputdefinition
{}
{

Misalkan
\mathkor {} {E} {dan} {F} {}
\definitionsverweis {bundel-bundel vektor riil}{}{}
di atas suatu
\definitionsverweis {ruang topologis}{}{}
$X$. Suatu
\definitionsverweis {homomorfisme bundel vektor}{}{}
\maabb {\varphi} { E } { F
} {}
disebut
\definitionsverweis {isomorfisme}{}{,}
apabila terdapat suatu homomorfisme
\maabb {\psi} { F } { E
} {}
yang, ketika dikomposisikan dengan $\varphi$
\zusatzklammer {dalam kedua urutan} {} {}, menghasilkan identitas.

}

\inputdefinition
{}
{

Misalkan
\mathkor {} {X} {dan} {Y} {}
\definitionsverweis {ruang-ruang topologis}{}{}
dan misalkan
\maabbdisp {p} { Y } { X
} {}
suatu
\definitionsverweis {pemetaan kontinu}{}{.}
Yang dimaksud dengan
\definitionswort {penampang kontinu}{}
dari $p$ adalah suatu pemetaan kontinu
\maabb {s} { X } { Y
} {}
dengan

\mavergleichskettedisp
{\vergleichskette
{  p \circ s
}
{ =} {
\operatorname{Id}_{  X  }
}
{ } {
}
{ } {
}
{ } {
}
}
{}{}{.}

}
Sebagai contoh, $Y$ dapat dipandang sebagai suatu bundel vektor di atas $X$. Suatu penampang hanya mungkin ada apabila $p$ surjektif, dan ini selalu berlaku bagi suatu bundel vektor. Kadang-kadang suatu penampang diidentifikasi dengan citranya; hal ini tidak menimbulkan masalah karena suatu penampang selalu injektif. Peran khusus dimainkan oleh \stichwort {penampang nol} {,} yang memetakan setiap titik dasar $P$ ke titik nol dalam ruang vektor $V_P$.
Untuk
\definitionsverweis {bundel tangen}{}{}
penampang-penampang kontinu tepat merupakan
\definitionsverweis {medan-medan vektor}{}{}
kontinu. Untuk bundel vektor trivial
\maabbdisp {} {X \times \R^n } { X
} {}
terdapat penampang-penampang konstan

\mathbed {e_i} {}
 {i = 1 , \ldots , n} {}
 {} {} {} {,}
yang memetakan setiap titik ke vektor

\mavergleichskette
{\vergleichskette
{ e_i
}
{ \in }{ \R^n
}
{  }{
}
{    }{
}
{   }{
}
}
{}{}{}
. Keluarga penampang ini membentuk suatu basis pada serat di atas setiap titik.

Misalkan
\maabb {p} { V } { X
} {}
suatu
\definitionsverweis {bundel vektor riil}{}{}
ber-
\definitionsverweis {rank}{}{}
$r$ di atas suatu
\definitionsverweis {ruang topologis}{}{}
$X$, dan misalkan

\mavergleichskette
{\vergleichskette
{ U
}
{ \subseteq }{ X
}
{  }{
}
{    }{
}
{   }{
}
}
{}{}{}
suatu
\definitionsverweis {himpunan bagian terbuka}{}{.}
Suatu keluarga
\definitionsverweis {penampang}{}{}
kontinu
\maabbdisp {s_1 , \ldots , s_r} { U} {V  \vert_U
} {}
disebut keluarga
\definitionswort {penampang basis}{,}
apabila untuk setiap titik

\mavergleichskette
{\vergleichskette
{ P
}
{ \in }{ U
}
{  }{
}
{    }{
}
{   }{
}
}
{}{}{}
vektor-vektor

\mavergleichskette
{\vergleichskette
{ s_1(P)  , \ldots , s_r(P)
}
{ \in }{ V_P
}
{  }{
}
{    }{
}
{   }{
}
}
{}{}{}
membentuk suatu
\definitionsverweis {basis}{}{}
dari $V_P$.

Keberadaan penampang-penampang basis
\zusatzklammer {kontinu} {} {}
pada suatu himpunan terbuka $U$ ekuivalen dengan keberadaan suatu trivialisasi
\zusatzklammer {kontinu} {} {}
dari bundel vektor di atas $U$; lihat
Soal 12.14.

  \zwischenueberschrift{Data perekatan}

Gagasan untuk merekatkan data lokal guna mendeskripsikan objek global telah kita jumpai dalam
Soal 8.15.

\inputdefinition
{}
{

Yang dimaksud dengan suatu
\definitionswort {data perekatan}{}
untuk
\definitionsverweis {bundel vektor riil}{}{}
ber-rank $r$ di atas suatu
\definitionsverweis {ruang topologis}{}{}
$X$ adalah kumpulan data berikut.
\aufzaehlungvier{Suatu
\definitionsverweis {penutup terbuka}{}{}

\mavergleichskettedisp
{\vergleichskette
{X
}
{ =} { \bigcup_{i \in I} U_i
}
{ } {
}
{ } {
}
{ } {
}
}
{}{}{}
}{Suatu keluarga

\mathbed {E_i \rightarrow U_i} {}
 {i \in I} {}
 {} {} {} {,}
bundel vektor riil ber-rank $r$.
}{Untuk setiap pasangan
\mathl{(i,j)}{} suatu
\definitionsverweis {isomorfisme bundel vektor}{}{}
\maabbdisp {\varphi_{ji}} { E_i {{|}}_{U_i \cap U_j }
} { E_{j} {{|}}_{U_i \cap U_j }
} {}
di atas
\mathl{U_i \cap U_j}{.}
}{Untuk indeks-indeks

\mavergleichskette
{\vergleichskette
{ i,j,k
}
{ \in }{ I
}
{  }{
}
{    }{
}
{   }{
}
}
{}{}{}
berlaku
\definitionswort {syarat koklus}{}

\mavergleichskettedisp
{\vergleichskette
{ \varphi_{kj} \circ \varphi_{ji}
}
{ =} { \varphi_{ki}
}
{ } {
}
{ } {
}
{ } {
}
}
{}{}{}
sebagai pemetaan dari
\mathl{E_{i} {{|}}_{U_i \cap U_j \cap U_k }}{} ke
\mathl{E_{k}  {{|}}_{U_i \cap U_j \cap U_k }}{}.
}

}

\inputbemerkung
{}
{

Secara khas, dalam
Definisi 12.10
bundel-bundel vektor pada (2) adalah bundel-bundel vektor trivial di atas $U_i$, yakni

\mavergleichskette
{\vergleichskette
{E_i
}
{ = }{ \R^r \times U_i
}
{  }{
}
{    }{
}
{   }{
}
}
{}{}{.}
Isomorfisme-isomorfisme pada (3) kemudian hanya diberikan oleh pemetaan linear bijektif
\maabb {\varphi_{ji}} { \R^r} { \R^r
} {}
yang bergantung secara kontinu pada titik dasar dari
\mathl{U_i \cap U_j}{}. Pemetaan-pemetaan ini dapat dideskripsikan secara ringkas oleh pemetaan kontinu
\maabbdisp {\varphi_{ji}} {U_i \cap U_j } { \operatorname{GL}_{  r   } \! \left(  \R  \right)
} {}
ke dalam
\definitionsverweis {grup linear umum}{}{}
. Jadi, setiap titik dasar dipetakan secara kontinu ke suatu objek
\definitionsverweis {invertibel}{}{}
berupa $r \times r$-
\definitionsverweis {matriks}{}{,}
di mana kontinuitas berarti bahwa semua entri matriks merupakan fungsi kontinu. Ini disebut \stichwort {deskripsi matriks} {} dari bundel tersebut. Syarat koklus tetap berlaku.

}


\inputfaktbeweis
{Topologisches reelles Vektorbündel/Verklebungsdatum/Existenz/Fakt}
{Lema}
{}
{

\faktsituation {Misalkan diberikan suatu
\definitionsverweis {data perekatan}{}{}

\mathbed {E_i} {}
 {i \in I} {}
 {} {} {} {,}
di atas suatu
\definitionsverweis {ruang topologis}{}{}

\mavergleichskettedisp
{\vergleichskette
{X
}
{ =} { \bigcup_{i \in I} U_i
}
{ } {
}
{ } {
}
{ } {
}
}
{}{}{}
.}
\faktfolgerung {Maka terdapat suatu bundel vektor riil yang ditentukan secara tunggal
\maabb {} { E } { X
} {}
beserta
\definitionsverweis {isomorfisme-isomorfisme}{}{}
\maabb {\psi_i} {E_i } { E \vert_{U_i }
} {}
sedemikian sehingga

\mavergleichskettedisp
{\vergleichskette
{ \psi_i \vert_{ E_i {{ |}}_{ U_{i} \cap U_j } }
}
{ =} { \psi_j \vert_{ E_j {{ |}}_{ U_{i} \cap U_j } }  \circ \varphi_{ji}
}
{ } {
}
{ } {
}
{ } {
}
}
{}{}{}
berlaku.}
\faktzusatz {}
\faktzusatz {}

 }
{ Lihat Soal 12.12. }

\bild{ \begin{center}
\includegraphics[width=5.5cm]{ \bildeinlesunggif {Fiddler_crab_mobius_strip} {gif} }
\end{center}
\bildtext {Animasi seekor kepiting biola yang mengelilingi pita Möbius dan kembali dengan orientasi tercermin. Edisi PDF menampilkan bingkai pertama yang statis; animasi asli dipertahankan untuk edisi HTML dan unduhan.} }

\bildlizenz { Fiddler_crab_mobius_strip.gif } {} {Hamishtodd1} {Commons} {CC-by-sa 4.0} {}

\inputbeispiel{}
{

Pada lingkaran satu dimensi

\mavergleichskettedisp
{\vergleichskette
{S^1
}
{ =} { \left\{ (x,y) \in \R^2  \mid   x^2+y^2 = 1        \right\}
}
{ } {
}
{ } {
}
{ } {
}
}
{}{}{}
kita meninjau
\definitionsverweis {penutup terbuka}{}{}

\mavergleichskette
{\vergleichskette
{S^1
}
{ = }{ U \cup V
}
{  }{
}
{    }{
}
{   }{
}
}
{}{}{}
dengan

\mavergleichskette
{\vergleichskette
{U
}
{ = }{S^1 \setminus \{(0,1)\}
}
{  }{
}
{    }{
}
{   }{
}
}
{}{}{}
dan

\mavergleichskette
{\vergleichskette
{V
}
{ = }{S^1 \setminus \{(0,-1)\}
}
{  }{
}
{    }{
}
{   }{
}
}
{}{}{.}
Di atasnya kita mendeskripsikan suatu
\definitionsverweis {data perekatan}{}{}
untuk bundel vektor riil ber-rank $1$. Kedua himpunan terbuka tersebut
\definitionsverweis {homeomorfik}{}{}
dengan garis riil. Kita mempunyai

\mavergleichskettedisp
{\vergleichskette
{U \cap V
}
{ =} { S^1 \setminus \{(0,1), (0,-1)  \}
}
{ =} { \left\{ (x,y) \in  S^1   \mid   x \neq 0        \right\}
}
{ } {
}
{ } {
}
}
{}{}{,}
dan himpunan ini tidak
\definitionsverweis {terhubung}{}{,}
melainkan homeomorfik dengan dua setengah garis riil terbuka yang saling lepas
\zusatzklammer {atau, ekuivalen, dengan garis-garis} {} {.}
Kita tetapkan

\mavergleichskette
{\vergleichskette
{L
}
{ = }{ U \times \R
}
{  }{
}
{    }{
}
{   }{
}
}
{}{}{}
dan

\mavergleichskette
{\vergleichskette
{M
}
{ = }{ V \times \R
}
{  }{
}
{    }{
}
{   }{
}
}
{}{}{.}
Kita mendefinisikan suatu
\definitionsverweis {isomorfisme}{}{}
\maabbdisp {\varphi} {L \vert_{U \cap V} } { M\vert_{U \cap V}
} {}
melalui

\mavergleichskettedisp
{\vergleichskette
{ \varphi(x,y,t)
}
{ \defeq} { \begin{cases}  (x,y,t) \text{ untuk } x > 0 \, , \\  (x,y,-t) \text{ untuk } x < 0  \,  \end{cases}
}
{ } {
}
{ } {
}
{ } {
}
}
{}{}{}
. Perhatikan bahwa $\varphi$ kontinu karena kedua ekspresi fungsional itu berlaku pada himpunan bagian terbuka yang saling lepas. Pada separuh yang satu pemetaannya adalah identitas, sedangkan pada separuh yang lain pemetaannya membalik. Dalam pengertian
Catatan 12.11
kita memperoleh deskripsi matriks
\zusatzklammer {konstan} {} {}
yang kontinu

\mavergleichskettedisp
{\vergleichskette
{ \psi (x,y)
}
{ \defeq} { \begin{cases}  (1) \text{ untuk } x > 0 \, , \\  (-1) \text{ untuk } x < 0  \, , \end{cases}
}
{ } {
}
{ } {
}
{ } {
}
}
{}{}{}
pada
\mathl{U \cap V}{}. Karena hanya terdapat dua himpunan terbuka, syarat koklus terpenuhi secara otomatis. Berdasarkan
Lema 12.12
data perekatan ini mendefinisikan suatu bundel vektor riil ber-rank $1$ pada lingkaran tersebut, yang disebut \stichwort {pita Möbius} {}.

}

\chapter{Lembar Kerja 12}
\setcounter{section}{12}

\inputaufgabe
{}
{

Tunjukkan bahwa sebuah
\definitionsverweis {bundel vektor riil}{}{}
di atas satu titik
\zusatzklammer {yakni sebuah ruang topologis bertitik tunggal} {} {}
sama dengan sebuah ruang vektor riil berdimensi hingga.

}
{} {}

\inputaufgabe
{}
{

Misalkan
\maabb {p} { V } { X
} {}
sebuah
\definitionsverweis {bundel vektor riil}{}{}
di atas sebuah
\definitionsverweis {ruang topologis}{}{}
$X$. Tunjukkan bahwa $V$ merupakan sebuah
\definitionsverweis {ruang Hausdorff}{}{}
jika dan hanya jika $X$ merupakan ruang Hausdorff.

}
{} {}

\inputaufgabe
{}
{

Misalkan $X$ sebuah
\definitionsverweis {ruang topologis}{}{.}
Tunjukkan bahwa
\definitionsverweis {identitas}{}{}
\maabb {\operatorname{Id}_{  X  }} { X } { X
} {}
dapat dipandang sebagai sebuah
\definitionsverweis {bundel vektor riil}{}{}
ber-
\definitionsverweis {rank}{}{}
$0$.

}
{} {}

\inputaufgabe
{}
{

Tunjukkan bahwa bundel kernel dari
Contoh 12.6
adalah trivial.

}
{} {}

\inputaufgabe
{}
{

Untuk
\definitionsverweis {bundel vektor}{}{}

\mavergleichskettealign
{\vergleichskettealign
{L
}
{ =} { \left\{ (r,s,t,u,v,w)  \mid  ru+sv+tw = 0 , \, (r,s,t) \neq (0,0,0)       \right\}
}
{ \subseteq} { (\R^3 \setminus \{(0,0,0)\}) \times \R^3
}
{ \longrightarrow} { \R^3 \setminus \{(0,0,0)\}
}
{ } {
}
}
{}{}{}
tentukan trivialisasi-trivialisasi linear di atas
\mathl{D(r)}{,}
\mathl{D(s)}{} dan
\mathl{D(t)}{,} yakni basis-basis yang bergantung pada $r,s,t$ di atas $D(r)$ dan seterusnya. Tentukan pemetaan-pemetaan perubahan basis pada

\mavergleichskette
{\vergleichskette
{D(rs)
}
{ = }{ D(r) \cap D(s)
}
{  }{
}
{    }{
}
{   }{
}
}
{}{}{.}

}
{} {}

\inputaufgabe
{}
{

Untuk
\definitionsverweis {bundel vektor}{}{}

\mavergleichskettealign
{\vergleichskettealign
{L
}
{ =} { \left\{ (r,s,t,u,v,w)  \mid  ru+sv+tw = 0 , \, (r,s,t) \neq (0,0,0)       \right\}
}
{ \subseteq} { (\R^3 \setminus \{(0,0,0)\}) \times \R^3
}
{ \longrightarrow} { \R^3 \setminus \{(0,0,0)\}
}
{ } {
}
}
{}{}{}
tentukan parameter-parameter
\mathl{(r,s,t)}{,} yang untuknya vektor $\left(  3   , \,   7  , \,   4                 \right)$ termasuk dalam serat
\mathl{L_{(r,s,t)}}{}.

}
{} {}

\inputaufgabe
{}
{

Deskripsikan
\definitionsverweis {bundel tangen}{}{}
dari
\definitionsverweis {sfera}{}{}
berdimensi $(n-1 )$

\mavergleichskettedisp
{\vergleichskette
{ S^{n-1}
}
{ =} { \left\{  \left(  x_1   , \,   \ldots  , \,   x_n          \right)   \mid   x_1^2 + \cdots + x_n^2 = 1         \right\}
}
{ \subset} { \R^n
}
{ } {
}
{ } {
}
}
{}{}{}
sebagai sebuah
\definitionsverweis {bundel kernel}{}{}
dalam arti
Catatan 12.5,
bandingkan pula dengan
Contoh 12.7.

}
{} {}

\inputaufgabe
{}
{

Misalkan $X$ sebuah
\definitionsverweis {ruang topologis}{}{.}
Tunjukkan bahwa sebuah
\definitionsverweis {homomorfisme}{}{}
antara bundel-bundel vektor
\zusatzklammer {trivial} {} {}
\maabbdisp {\varphi} {X \times \R^n } { X \times \R^m
} {}
sama dengan sebuah susunan berukuran
$m \times n$, yaitu sebuah
\definitionsverweis {matriks}{}{,}
yang entri-entrinya merupakan fungsi-fungsi kontinu dari $X$ ke $\R$.

}
{} {}

\inputaufgabe
{}
{

Misalkan

\mavergleichskette
{\vergleichskette
{U
}
{ \subseteq }{ \R^n
}
{  }{
}
{    }{
}
{   }{
}
}
{}{}{}
\definitionsverweis {terbuka}{}{}
dan
\maabb {f} { U } { \R^m
} {}
sebuah
\definitionsverweis {pemetaan terdiferensialkan kontinu}{}{.}
Tunjukkan bahwa
\definitionsverweis {diferensial total}{}{}
dalam bentuk
\maabbeledisp {} { U \times \R^n } { U \times \R^m
} { (x,v)} { (x, \left(D f \right)_{ x} (v) )
} {,}
memberikan sebuah homomorfisme dari bundel vektor
\mathl{U \times \R^n}{} ke bundel vektor
\mathl{U \times \R^m}{}.

}
{} {}

\inputaufgabe
{}
{

Perhatikan
\definitionsverweis {ruang topologis}{}{}

\mavergleichskettedisp
{\vergleichskette
{Y
}
{ \defeq} { \left\{ (s,t,u,v) \in \R^4  \mid  su+tv = 1        \right\}
}
{ } {
}
{ } {
}
{ } {
}
}
{}{}{}
dengan proyeksi
\maabbeledisp {p} { Y } { \R^2 \setminus \{ (0,0) \} = X
} {(s,t,u,v)} { (s,t)
} {.}
\aufzaehlungvier{Tunjukkan bahwa setiap serat $p$
\definitionsverweis {homeomorfik}{}{}
dengan sebuah garis riil.
}{Tunjukkan bahwa

\mavergleichskettedisp
{\vergleichskette
{ \varphi(s,t)
}
{ =} {(s,t,u(s,t),v(s,t))
}
{ =} { \left( s,t, { \frac{   s  }{ s^2+t^2 } } , { \frac{   t  }{ s^2+t^2 } }  \right)
}
{ } {
}
{ } {
}
}
{}{}{}
mendefinisikan sebuah pemetaan kontinu
\maabb {\varphi} { X} { Y
} {}
dengan

\mavergleichskettedisp
{\vergleichskette
{ p \circ \varphi
}
{ =} {
\operatorname{Id}_{  X  }
}
{ } {
}
{ } {
}
{ } {
}
}
{}{}{}.
}{Definisikan sebuah
\definitionsverweis {homeomorfisme}{}{}
antara
\mathkor {} {Y} {dan} {X \times \R} {.}
}{Tunjukkan bahwa tidak ada pemetaan polinomial
\maabb {\psi} { X } { Y
} {}
dengan

\mavergleichskettedisp
{\vergleichskette
{ p \circ \psi
}
{ =} {
\operatorname{Id}_{  X  }
}
{ } {
}
{ } {
}
{ } {
}
}
{}{}{}.
}

}
{} {}

\bild{ \begin{center}
\includegraphics[width=5.5cm]{ \bildeinlesungsvg {Inclusion-exclusion} {svg} }
\end{center}
\bildtext {Diagram Venn tiga himpunan A, B, dan C yang memperlihatkan daerah-daerah tumpang tindih untuk prinsip inklusi--eksklusi.} }

\bildlizenz { Inclusion-exclusion.svg } {Pembuat tidak diketahui} {} {} {domain publik} {}

Bukan hanya bundel vektor, tetapi secara lebih umum ruang-ruang topologis juga dapat dideskripsikan melalui data perekatan.

Yang dimaksud dengan sebuah
\definitionswort {data perekatan}{}
untuk
\definitionsverweis {ruang-ruang topologis}{}{}
adalah kumpulan data berikut.
\aufzaehlungvier{Sebuah keluarga ruang topologis

\mathbed {U_i} {}
 {i \in I} {}
 {} {} {} {,}
}{Untuk setiap pasangan
\mathl{(i,j)}{} sebuah
\definitionsverweis {himpunan bagian terbuka}{}{}

\mavergleichskette
{\vergleichskette
{U_{ij}
}
{ \subseteq }{ U_i
}
{  }{
}
{    }{
}
{   }{
}
}
{}{}{}
\zusatzklammer {dengan

\mavergleichskettek
{\vergleichskettek
{U_{ii}
}
{ = }{ U_i
}
{  }{
}
{    }{
}
{   }{
}
}
{}{}{}} {} {.}
}{Untuk setiap pasangan
\mathl{(i,j)}{} sebuah
\definitionsverweis {homeomorfisme}{}{}
\maabbdisp {\varphi_{ji}} {U_{ij}} { U_{ji}
} {}
\zusatzklammer {dengan

\mavergleichskettek
{\vergleichskettek
{ \varphi_{ii}
}
{ = }{
\operatorname{Id}_{ U_i  }
}
{  }{
}
{    }{
}
{   }{
}
}
{}{}{}} {} {.}
}{Untuk indeks-indeks

\mavergleichskette
{\vergleichskette
{ i,j,k
}
{ \in }{ I
}
{  }{
}
{    }{
}
{   }{
}
}
{}{}{}
terpenuhi
\definitionswort {syarat koklus}{}

\mavergleichskettedisp
{\vergleichskette
{ \varphi_{kj} \circ \varphi_{ji}
}
{ =} { \varphi_{ki}
}
{ } {
}
{ } {
}
{ } {
}
}
{}{}{}
sebagai pemetaan dari
\mathl{U_{ik} \cap U_{ij}}{} ke
\mathl{U_{k}}{}.
}

\inputaufgabegibtloesung
{}
{

Misalkan diberikan sebuah
\definitionsverweis {data perekatan}{}{}

\mathbed {U_i} {}
 {i \in I} {}
 {} {} {} {,}
untuk
\definitionsverweis {ruang-ruang topologis}{}{.}
Tunjukkan bahwa terdapat sebuah ruang topologis $X$ yang ditentukan secara tunggal, sebuah penutup terbuka

\mavergleichskette
{\vergleichskette
{ X
}
{ = }{ \bigcup_{i \in I} V_i
}
{  }{
}
{    }{
}
{   }{
}
}
{}{}{}
dan
\definitionsverweis {homeomorfisme-homeomorfisme}{}{}
\maabb {\psi_i} {U_i } { V_i
} {}
sedemikian sehingga

\mavergleichskettedisp
{\vergleichskette
{ \psi_i  \left(  U_{ij}  \right)
}
{ =} { V_i \cap V_j
}
{ } {
}
{ } {
}
{ } {
}
}
{}{}{}
dan

\mavergleichskettedisp
{\vergleichskette
{ \psi_i \vert_{U_{ij} }
}
{ =} { \psi_j \vert_{U_{ji} }  \circ   \varphi_{ji}
}
{ } {
}
{ } {
}
{ } {
}
}
{}{}{}.

}
{} {}

\inputaufgabegibtloesung
{}
{

Misalkan diberikan sebuah
\definitionsverweis {data perekatan}{}{}

\mathbed {E_i} {}
 {i \in I} {}
 {} {} {} {,}
di atas sebuah
\definitionsverweis {ruang topologis}{}{}

\mavergleichskettedisp
{\vergleichskette
{ X
}
{ =} { \bigcup_{i \in I} U_i
}
{ } {
}
{ } {
}
{ } {
}
}
{}{}{}.
Tunjukkan bahwa terdapat sebuah bundel vektor riil yang ditentukan secara tunggal
\maabb {} { E } { X
} {}
dan
\definitionsverweis {isomorfisme-isomorfisme}{}{}
\maabb {\psi_i} {E_i } { E \vert_{U_i }
} {}
sedemikian sehingga

\mavergleichskettedisp
{\vergleichskette
{ \psi_i \vert_{ E_i {{ |}}_{ U_{i} \cap U_j } }
}
{ =} { \psi_j \vert_{ E_j {{ |}}_{ U_{i} \cap U_j } }  \circ \varphi_{ji}
}
{ } {
}
{ } {
}
{ } {
}
}
{}{}{}.

}
{} {}

\inputaufgabe
{}
{

Misalkan
\maabb {s} { X } { V
} {}
sebuah
\definitionsverweis {penampang kontinu}{}{}
dari sebuah
\definitionsverweis {bundel vektor riil}{}{}
\maabb {p} { V } { X
} {}
di atas sebuah
\definitionsverweis {ruang topologis}{}{}
$X$. Tunjukkan bahwa
\definitionsverweis {citra}{}{}

\mavergleichskette
{\vergleichskette
{s(X)
}
{ \subseteq }{ V
}
{  }{
}
{    }{
}
{   }{
}
}
{}{}{}
merupakan sebuah
\definitionsverweis {himpunan bagian tertutup}{}{}
yang
\definitionsverweis {homeomorfik}{}{}
dengan $X$.

}
{} {}

\inputaufgabe
{}
{

Misalkan
\maabb {p} { V } { X
} {}
sebuah
\definitionsverweis {bundel vektor riil}{}{}
ber-
\definitionsverweis {rank}{}{}
$r$ di atas sebuah
\definitionsverweis {ruang topologis}{}{}
$X$, dan misalkan

\mavergleichskette
{\vergleichskette
{ U
}
{ \subseteq }{ X
}
{  }{
}
{    }{
}
{   }{
}
}
{}{}{}
sebuah
\definitionsverweis {himpunan bagian terbuka}{}{.}
Tunjukkan bahwa pada $U$ terdapat sebuah
\definitionsverweis {trivialisasi}{}{}
kontinu dari $V \vert_U$ jika dan hanya jika terdapat sebuah keluarga
\definitionsverweis {penampang-penampang basis}{}{}
kontinu
\maabbdisp {s_1 , \ldots , s_r} { U} {V  \vert_U
} {}.

}
{} {}

\inputaufgabe
{}
{

Misalkan sebuah
\definitionsverweis {bundel vektor}{}{}
\maabb {} { V } { X
} {}
ber-rank $m$ di atas sebuah
\definitionsverweis {ruang topologis}{}{}
$X$ diberikan oleh pemetaan-pemetaan matriks kontinu
\maabbdisp {\varphi_{ij}} { U_i \cap U_j } { \operatorname{GL}_{  m   } \! \left(  \R  \right)
} {}
untuk sebuah penutup terbuka

\mavergleichskette
{\vergleichskette
{X
}
{ = }{ \bigcup_{i \in I} U_i
}
{  }{
}
{    }{
}
{   }{
}
}
{}{}{}.
Tunjukkan bahwa sebuah
\definitionsverweis {penampang kontinu}{}{}
\maabbdisp {s} { X } { V
} {}
sama dengan sebuah keluarga pemetaan kontinu
\maabb {t_i} {U_i } { \R^m
} {,}
yang memenuhi syarat

\mavergleichskettedisp
{\vergleichskette
{t_j
}
{ =} { \varphi_{ij} t_i
}
{ } {
}
{ } {
}
{ } {
}
}
{}{}{}
untuk semua $i,j$.

}
{} {}

\inputaufgabe
{}
{

Misalkan

\mavergleichskette
{\vergleichskette
{X
}
{ = }{ \bigcup_{i \in I} U_i
}
{  }{
}
{    }{
}
{   }{
}
}
{}{}{}
sebuah
\definitionsverweis {penutup terbuka}{}{}
dari sebuah
\definitionsverweis {ruang topologis}{}{}
$X$. Misalkan
\maabbdisp {\varphi_{ij}} {U_i \cap U_j } { \operatorname{GL}_{  r   } \! \left(  \R  \right)
} {}
dan
\maabbdisp {\psi_{ij}} {U_i \cap U_j } { \operatorname{GL}_{  r   } \! \left(  \R  \right)
} {}
merupakan
\definitionsverweis {deskripsi-deskripsi matriks}{}{}
yang secara berturut-turut menghasilkan bundel vektor
\mathkor {} {E} {dan} {F} {.}
Tunjukkan bahwa bundel-bundel ini
\definitionsverweis {isomorfik}{}{}
jika dan hanya jika terdapat pemetaan-pemetaan kontinu
\maabbdisp {\alpha_i} { U_i } { \operatorname{GL}_{  r   } \! \left(  \R  \right)
} {}
sedemikian sehingga
\zusatzklammer {untuk pembatasan-pembatasan yang bermakna} {} {}

\mavergleichskettedisp
{\vergleichskette
{ \alpha_i \circ  \varphi_{ij}  \circ \alpha_j^{-1}
}
{ =} { \psi_{ij}
}
{ } {
}
{ } {
}
{ } {
}
}
{}{}{}
untuk semua $i,j$.

}
{} {}

\inputaufgabe
{}
{

Tunjukkan bahwa komplemen penampang nol dalam sebuah
\definitionsverweis {bundel garis}{}{}
trivial tidak
\definitionsverweis {terhubung}{}{.}

}
{} {}

\inputaufgabe
{}
{

Ambillah sebuah pita persegi panjang yang tipis, puntir pita itu satu putaran penuh terhadap sumbu panjangnya, lalu rekatkan kedua sisi pendeknya. Sekarang potong pita itu memanjang tepat di tengah dengan gunting. Apakah objek yang dihasilkan
\definitionsverweis {terhubung}{}{?}
Bagaimanakah bentuknya jika dibuat $n$ setengah putaran?

}
{} {}

Misalkan
\maabb {} { V } { X
} {}
sebuah
\definitionsverweis {bundel vektor riil}{}{}
di atas sebuah
\definitionsverweis {ruang topologis}{}{.}
Sebuah bundel vektor riil
\maabb {} { U } { X
} {}
disebut
\definitionswort {subbundel}{}
dari $V$ jika diberikan sebuah
\definitionsverweis {homomorfisme bundel vektor}{}{}
\definitionsverweis {injektif}{}{}
\maabb {} { U } { V
} {}.

\inputaufgabe
{}
{

Misalkan $X$ sebuah
\definitionsverweis {ruang topologis}{}{}
dan
\maabb {\varphi} { V } { W
} {}
sebuah
\definitionsverweis {homomorfisme}{}{}
antara
\definitionsverweis {bundel-bundel vektor riil}{}{}
di atas $X$. Misalkan

\mavergleichskettedisp
{\vergleichskette
{ Y
}
{ =} { \left\{ P \in X  \mid  \varphi_P \text{ surjektif }         \right\}
}
{ } {
}
{ } {
}
{ } {
}
}
{}{}{}.
Tunjukkan bahwa
\zusatzklammer {keluarga} {} {}
\definitionsverweis {subruang-subruang vektor}{}{}

\mavergleichskette
{\vergleichskette
{ \ker \varphi_P
}
{ \subseteq }{ V_P
}
{  }{
}
{    }{
}
{   }{
}
}
{}{}{}
membentuk sebuah
\definitionsverweis {subbundel vektor}{}{}
dari $V\vert_Y$.

}
{} {}

\inputaufgabe
{}
{

Misalkan $X$ sebuah
\definitionsverweis {ruang topologis}{}{}
dan

\mavergleichskette
{\vergleichskette
{ U
}
{ \subseteq }{ V
}
{  }{
}
{    }{
}
{   }{
}
}
{}{}{}
sebuah
\definitionsverweis {subbundel vektor}{}{}
dari
\definitionsverweis {bundel vektor riil}{}{}
\maabb {} { V } { X
} {.}
Tunjukkan bahwa
\zusatzklammer {keluarga} {} {}
\definitionsverweis {ruang-ruang hasil bagi}{}{}

\mathl{V_P/U_P}{} membentuk sebuah bundel vektor di atas $X$ dan bahwa terdapat sebuah
\definitionsverweis {homomorfisme bundel}{}{}
\definitionsverweis {surjektif}{}{}
\maabb {} { V } { V/U
} {}.

}
{} {}

Kita perkenalkan dua istilah lagi.

Misalkan $K$ sebuah
\definitionsverweis {lapangan}{}{}
dan
\mathl{U,V,W}{}
\definitionsverweis {ruang-ruang vektor}{}{}
di atas $K$. Sebuah diagram berbentuk

\mathdisp {0 \, \stackrel{  }{\longrightarrow} \, U  \, \stackrel{  }{ \longrightarrow}   \, V  \, \stackrel{  }{\longrightarrow} \, W \,\stackrel{  } {\longrightarrow}  \, 0} {  }

disebut sebuah \definitionswort {barisan eksak pendek}{} ruang-ruang vektor-$K$ jika $U$ merupakan sebuah
\definitionsverweis {subruang vektor}{}{}
dari $V$ dan jika $W$
\definitionsverweis {isomorfik}{}{}
dengan
\definitionsverweis {ruang hasil bagi}{}{}
$V/U$.

Misalkan $X$ sebuah
\definitionsverweis {ruang topologis}{}{,}
misalkan
\maabbdisp {} {U,V,W} {X
} {}
\definitionsverweis {bundel-bundel vektor}{}{}
di atas $X$, serta misalkan
\maabb {\varphi} { U } { V
} {}
dan
\maabb {\psi} { V } { W
} {}
\definitionsverweis {homomorfisme-homomorfisme}{}{.}
Kita katakan bahwa terdapat sebuah
\definitionswort {barisan eksak pendek}{}

\mathdisp {0 \, \stackrel{  }{\longrightarrow} \,  U   \, \stackrel{ \varphi }{ \longrightarrow}   \,  V   \, \stackrel{ \psi }{\longrightarrow} \,  W  \,\stackrel{  } {\longrightarrow}  \, 0} {  }

jika untuk setiap titik

\mavergleichskette
{\vergleichskette
{ P
}
{ \in }{ X
}
{  }{
}
{    }{
}
{   }{
}
}
{}{}{}
terdapat sebuah
\definitionsverweis {barisan eksak pendek}{}{}

\mathdisp {0 \, \stackrel{  }{\longrightarrow} \, U_P   \, \stackrel{ \varphi_P }{ \longrightarrow}   \,  V_P   \, \stackrel{ \psi_P }{\longrightarrow} \,  W_P  \,\stackrel{  } {\longrightarrow}  \, 0} {  }

pada serat-seratnya.

\inputaufgabe
{}
{

Misalkan $X$ sebuah
\definitionsverweis {ruang topologis}{}{}
dan

\mavergleichskette
{\vergleichskette
{ U
}
{ \subseteq }{ V
}
{  }{
}
{    }{
}
{   }{
}
}
{}{}{}
sebuah
\definitionsverweis {subbundel vektor}{}{}
dari
\definitionsverweis {bundel vektor riil}{}{}
\maabb {} { V } { X
} {.}
Tunjukkan bahwa terdapat sebuah
\definitionsverweis {barisan eksak pendek}{}{}

\mathdisp {0 \, \stackrel{  }{\longrightarrow} \,  U   \, \stackrel{  }{ \longrightarrow}   \,  V   \, \stackrel{  }{\longrightarrow} \,  V/U \,\stackrel{  } {\longrightarrow}  \, 0} {  }.

}
{} {}

\inputaufgabe
{}
{

Misalkan
\mathkor {} {X} {dan} {Y} {}
\definitionsverweis {ruang-ruang topologis}{}{}
dan
\maabb {\varphi} { Y } { X
} {}
sebuah
\definitionsverweis {pemetaan kontinu}{}{.}
Misalkan
\maabb {p} { V } { X
} {}
sebuah
\definitionsverweis {bundel vektor}{}{}
di atas $X$. Tunjukkan bahwa
\mathl{Y \times_X V}{}
\zusatzklammer {lihat
Soal 11.16} {} {}
merupakan sebuah bundel vektor di atas $Y$.

}
{} {}

\inputaufgabe
{}
{

Misalkan
\mathkor {} {X} {dan} {Y} {}
\definitionsverweis {ruang-ruang topologis}{}{}
dan
\maabb {\varphi} {Y} {X
} {}
sebuah
\definitionsverweis {pemetaan kontinu}{}{.}
Misalkan
\maabb {p} {V} {X
} {}
sebuah
\definitionsverweis {bundel vektor}{}{}
di atas $X$. Deskripsikan
\definitionsverweis {tarik balik}{}{}
sebuah bundel vektor dengan bantuan
\definitionsverweis {data perekatan}{}{.}

}
{} {}

\inputaufgabe
{}
{

Misalkan $X$ sebuah
\definitionsverweis {ruang topologis}{}{}
dan
\maabb {p} { V } { X
} {}
serta
\maabb {q} { W } { X
} {}
\definitionsverweis {bundel-bundel vektor}{}{}
di atas $X$. Tunjukkan bahwa
\mathl{V \times_X W}{}
\zusatzklammer {lihat
Soal 11.16} {} {}
merupakan sebuah bundel vektor di atas $X$ yang sama dengan
\definitionsverweis {jumlah langsung bundel-bundel vektor}{}{}
di atas $X$.

}
{} {}

\inputaufgabe
{}
{

Misalkan
\maabb {\varphi} { X } { Y
} {}
sebuah
\definitionsverweis {pemetaan diferensiabel}{}{}
antara
\definitionsverweis {manifold-manifold diferensiabel}{}{}
\mathkor {} {X} {dan} {Y} {.}
Tunjukkan bahwa pemetaan ini menghasilkan
\definitionsverweis {homomorfisme bundel vektor}{}{}
alami
\maabbdisp {T\varphi} {TX} { \varphi^*TY
} {}
antara
\definitionsverweis {bundel-bundel vektor}{}{}
di atas $X$.

}
{} {}

  \zwischenueberschrift{Soal untuk dikumpulkan}

\inputaufgabe
{4}
{

Misalkan $X$ sebuah
\definitionsverweis {ruang topologis}{}{}
dan misalkan

\mavergleichskette
{\vergleichskette
{ f_1 , \ldots , f_n
}
{ \in }{ C^0(X,\R)
}
{  }{
}
{    }{
}
{   }{
}
}
{}{}{}
merupakan fungsi-fungsi bernilai riil
\definitionsverweis {kontinu}{}{}
pada $X$. Misalkan

\mavergleichskette
{\vergleichskette
{ U
}
{ = }{ \bigcup_{i = 1}^n D(f_i)
}
{ \subseteq }{ X
}
{    }{
}
{   }{
}
}
{}{}{}
dan

\mavergleichskettedisp
{\vergleichskette
{ S
}
{ =} { \left\{ (P; t_1 , \ldots , t_n)  \mid  P \in U , \,  \sum_{i = 1}^n f_i(P)t_i = 0       \right\}
}
{ \subseteq} { U \times \R^n
}
{ } {
}
{ } {
}
}
{}{}{,}
bandingkan dengan
Catatan 12.5.
Berikan secara eksplisit trivialisasi-trivialisasi untuk $S$ di atas $D(f_i)$ dan tunjukkan bahwa $S$ merupakan sebuah bundel vektor di atas $U$ ber-rank $n-1$.

}
{} {}

\inputaufgabe
{4}
{

Tunjukkan bahwa sebuah
\definitionsverweis {bundel garis riil}{}{}
\maabb {} { L } { X
} {}
di atas sebuah
\definitionsverweis {ruang topologis}{}{}
$X$ adalah trivial jika dan hanya jika bundel itu memiliki sebuah
\definitionsverweis {penampang}{}{}
kontinu yang tidak pernah nol.

}
{} {}

\inputaufgabe
{3}
{

Tunjukkan bahwa sebuah data perekatan untuk sebuah bundel garis terhadap sebuah penutup terbuka

\mavergleichskette
{\vergleichskette
{X
}
{ = }{ \bigcup_{i \in I} U_i
}
{  }{
}
{    }{
}
{   }{
}
}
{}{}{}
sama dengan sebuah keluarga fungsi kontinu yang tidak pernah nol
\maabb {f_{ij}} {U_i \cap U_j } { \R
} {,}
yang pada
\mathl{U_i \cap U_j \cap U_k}{} masing-masing memenuhi syarat

\mavergleichskette
{\vergleichskette
{ f_{ki}
}
{ = }{ f_{kj} \cdot f_{ji}
}
{  }{
}
{    }{
}
{   }{
}
}
{}{}{.}

}
{} {}

\inputaufgabe
{4}
{

Tunjukkan bahwa
\definitionsverweis {pita Möbius}{}{}
merupakan sebuah
\definitionsverweis {manifold}{}{}
berdimensi dua.

}
{} {}

\section*{Solusi yang disediakan oleh sumber}
\addcontentsline{toc}{section}{Solusi yang disediakan oleh sumber}
\subsection*{Solusi Soal 12.11}
Misalkan $Y$ merupakan gabungan saling lepas dari semua $U_i$. Pada $Y$ kita definisikan sebuah
\definitionsverweis {relasi ekuivalensi}{}{}
$\sim$ dengan menyatakan titik-titik

\mavergleichskette
{\vergleichskette
{ x_i
}
{ \in }{ U_i
}
{  }{
}
{    }{
}
{   }{
}
}
{}{}{}
dan

\mavergleichskette
{\vergleichskette
{ x_j
}
{ \in }{ U_j
}
{  }{
}
{    }{
}
{   }{
}
}
{}{}{}
ekuivalen jika

\mavergleichskette
{\vergleichskette
{ x_i
}
{ \in }{ U_{ij}
}
{  }{
}
{    }{
}
{   }{
}
}
{}{}{,}

\mavergleichskette
{\vergleichskette
{ x_j
}
{ \in }{ U_{ji}
}
{  }{
}
{    }{
}
{   }{
}
}
{}{}{}
dan

\mavergleichskette
{\vergleichskette
{ \varphi_{ji} (x_i)
}
{ = }{ x_j
}
{  }{
}
{    }{
}
{   }{
}
}
{}{}{}.
Sifat-sifat relasi ekuivalensi dijamin oleh syarat koklus; lihat
soal tentang relasi ekuivalensi.
Kita tetapkan

\mavergleichskettedisp
{\vergleichskette
{X
}
{ \defeq} { Y/ \sim
}
{ } {
}
{ } {
}
{ } {
}
}
{}{}{}
dan melengkapi $X$ dengan
\definitionsverweis {topologi hasil bagi}{}{.}
Komposisi-komposisi

\mathdisp {U_i \longrightarrow Y \longrightarrow X} {  }

adalah $\psi_i$, dan $V_i$ merupakan citra pemetaan-pemetaan ini. Jadi diperoleh homeomorfisme-homeomorfisme
\maabb {\psi_i} {U_i } { V_i
} {}.
Untuk

\mavergleichskette
{\vergleichskette
{ x
}
{ \in }{ U_i
}
{  }{
}
{    }{
}
{   }{
}
}
{}{}{}

\mavergleichskettedisp
{\vergleichskette
{ \psi_i (x)
}
{ \in} { V_j
}
{ } {
}
{ } {
}
{ } {
}
}
{}{}{}
berlaku jika dan hanya jika

\mavergleichskette
{\vergleichskette
{ x
}
{ \in }{ U_{ij}
}
{  }{
}
{    }{
}
{   }{
}
}
{}{}{},
sebab tepat dalam kasus ini $x$ diidentifikasi dengan
\mathl{\varphi_{ji}(x)}{}.
Oleh karena itu,

\mavergleichskette
{\vergleichskette
{ \psi_i(U_{ij})
}
{ = }{ V_i \cap V_j
}
{  }{
}
{    }{
}
{   }{
}
}
{}{}{.}
Komutativitas diagram

\mathdisp {\begin{matrix} U_{ij}   & \stackrel{ \varphi_{ji} }{\longrightarrow} &  U_{ji}   & \\ & \!\!\! \!\! \psi_i \searrow & \downarrow \psi_j \!\!\! \!\! & \\ & &  V_i \cap V_j   & \!\!\!\! \!\!\!\!  \\ \end{matrix}} {  }

mengikuti dengan cara yang sama.
\subsection*{Solusi Soal 12.12}
Keberadaan sebuah ruang topologis $E$ dengan sifat-sifat tersebut mengikuti dari
fakta tentang keberadaan ruang hasil perekatan
\zusatzklammer {himpunan-himpunan terbuka yang harus direkatkan adalah

\mavergleichskettek
{\vergleichskettek
{W_{ij}
}
{ \defeq }{ E_i \vert_{U_i \cap U_j }
}
{  }{
}
{    }{
}
{   }{
}
}
{}{}{}} {} {}
dan keberadaan pemetaan kontinu ke $X$ mengikuti dari
fakta tentang pemetaan kontinu dari ruang hasil perekatan.
Pada setiap serat $E_x$ terdapat struktur ruang vektor yang terdefinisi dengan baik, yang berasal dari $E_i$ untuk sebarang lingkungan terbuka

\mavergleichskette
{\vergleichskette
{ x
}
{ \in }{ U_i
}
{  }{
}
{    }{
}
{   }{
}
}
{}{}{}.
Ketaktergantungan terhadap pilihan tersebut didasarkan pada kenyataan bahwa untuk

\mavergleichskette
{\vergleichskette
{ x
}
{ \in }{ U_i \cap U_j
}
{  }{
}
{    }{
}
{   }{
}
}
{}{}{}
menurut asumsi terdapat sebuah isomorfisme bundel vektor
\maabbdisp {\varphi_{ji}} {E_i \vert_{U_i \cap U_j } } { E_j \vert_{U_i \cap U_j }
} {}
yang menginduksi sebuah
\definitionsverweis {isomorfisme ruang vektor}{}{}
\maabbdisp {} { \left(  E_i  \right)_x } { \left(  E_j  \right)_x
} {}.

\part{Unit 13}
\chapter[Kuliah 13]{Kuliah 13: Konstruksi Bundel Vektor}
\setcounter{section}{13}

  \zwischenueberschrift{Konstruksi bundel vektor}

Semua konstruksi yang dilakukan pada ruang vektor dalam aljabar linear dapat diterapkan pada bundel vektor. Pada setiap serat, kita menggunakan konstruksi untuk ruang vektor. Namun, karena trivialisasi juga harus diperhitungkan, cara termudah adalah bekerja dengan data perekatan.

\inputdefinition
{}
{

Untuk
\definitionsverweis {bundel-bundel vektor riil}{}{}
\mathkor {} {E} {dan} {F} {}
di atas suatu
\definitionsverweis {ruang topologis}{}{}
$X$ dengan trivialisasi
\maabbdisp {\alpha_i} {E {{|}}_{U_i } } { U_i \times \R^m
} {}
dan
\maabbdisp {\beta_i} {F{{|}}_{U_i } } { U_i \times \R^n
} {}
bundel vektor yang terkait dengan
\definitionsverweis {data perekatan}{}{}

\mavergleichskettedisp
{\vergleichskette
{G_i
}
{ =} { U_i \times \R^m \times \R^n
}
{ } {
}
{ } {
}
{ } {
}
}
{}{}{}
dan
\maabbdisp {\varphi_{ij}} {G_i {{|}}_{U_i \cap U_j }  } { G_j {{|}}_{U_i \cap U_j }
} {}
dengan

\mavergleichskettedisp
{\vergleichskette
{ \varphi_{ij} (x,v,w)
}
{ \defeq} { \left(  x   , \,   \alpha_j (\alpha_i^{-1} (x,v))  , \,   \beta_j (\beta_i^{-1} (x,w))                 \right)
}
{ } {
}
{ } {
}
{ } {
}
}
{}{}{}
disebut
\definitionswort {jumlah langsung}{}
dari bundel-bundel vektor
\mathkor {} {E} {dan} {F} {.}
Bundel ini dilambangkan dengan
\mathl{E \oplus F}{}.

}

Jika $E$ diberikan oleh deskripsi matriks
\maabbdisp {\varphi_{ij}} { U_i \cap U_j } { \operatorname{GL}_{ m  } \! \left(  \R  \right)
} {}
dan $F$ oleh deskripsi matriks
\maabbdisp {\psi_{ij}} { U_i \cap U_j } { \operatorname{GL}_{ n  } \! \left(  \R  \right)
} {,}
maka deskripsi matriks dari
\mathl{E \oplus F}{} diperoleh dengan menyusun kedua matriks secara diagonal menjadi suatu matriks
\mathl{(m+n) \times (m+n)}{} dan mengisi entri-entri lainnya dengan nol.

\inputdefinition
{}
{

Untuk
\definitionsverweis {bundel-bundel vektor riil}{}{}
\mathkor {} {E} {dan} {F} {}
di atas suatu
\definitionsverweis {ruang topologis}{}{}
$X$ dengan trivialisasi
\maabbdisp {\alpha_i} {E {{|}}_{U_i } } { U_i \times \R^m
} {}
dan
\maabbdisp {\beta_i} {F{{|}}_{U_i } } { U_i \times \R^n
} {}
bundel vektor yang terkait dengan
\definitionsverweis {data perekatan}{}{}

\mavergleichskettedisp
{\vergleichskette
{G_i
}
{ =} { U_i \times \left(  \R^m  \otimes  \R^n  \right)
}
{ } {
}
{ } {
}
{ } {
}
}
{}{}{}
dan
\maabbdisp {\varphi_{ij}} {G_i {{|}}_{U_i \cap U_j }  } { G_j {{|}}_{U_i \cap U_j }
} {}
dengan

\mavergleichskettedisp
{\vergleichskette
{ \varphi_{ij}
}
{ \defeq} { \left(  \alpha_j \circ \alpha_i^{-1}  \right)   \otimes \left(  \beta_j  \circ \beta_i^{-1}  \right)
}
{ } {
}
{ } {
}
{ } {
}
}
{}{}{}
\zusatzklammer {di sini, pada setiap titik dasar diambil
\definitionsverweis {produk tensor pemetaan linear}{}{}} {} {}
disebut
\definitionswort {produk tensor}{}
dari bundel-bundel vektor
\mathkor {} {E} {dan} {F} {.}
Bundel ini dilambangkan dengan
\mathl{E \otimes F}{}.

}

Jika deskripsi-deskripsi matriksnya diberikan, deskripsi matriks produk tensor diperoleh melalui apa yang disebut
\definitionsverweis {produk Kronecker}{}{.}
Dalam proses ini, setiap entri matriks yang satu dikalikan dengan setiap entri matriks yang lain.

\inputdefinition
{}
{

Untuk suatu
\definitionsverweis {bundel vektor riil}{}{}
$E$ ber-rank $m$ di atas suatu
\definitionsverweis {ruang topologis}{}{}
$X$ dengan trivialisasi
\maabbdisp {\alpha_i} {E {{|}}_{U_i } } { U_i \times \R^m
} {}
dan

\mavergleichskette
{\vergleichskette
{r
}
{ \in }{\N
}
{  }{
}
{    }{
}
{   }{
}
}
{}{}{}
bundel vektor yang terkait dengan
\definitionsverweis {data perekatan}{}{}

\mavergleichskettedisp
{\vergleichskette
{G_i
}
{ =} { U_i \times \left(   \bigwedge^r \R^m   \right)
}
{ } {
}
{ } {
}
{ } {
}
}
{}{}{}
dan
\maabbdisp {\varphi_{ij}} {G_i {{|}}_{U_i \cap U_j }  } { G_j {{|}}_{U_i \cap U_j }
} {}
dengan

\mavergleichskettedisp
{\vergleichskette
{ \varphi_{ij}
}
{ \defeq} {  \bigwedge^r \left(  \alpha_j \circ \alpha_i^{-1}  \right)
}
{ } {
}
{ } {
}
{ } {
}
}
{}{}{}
\zusatzklammer {di sini, pada setiap titik dasar diambil
\definitionsverweis {produk eksterior ke-$r$ dari pemetaan linear}{}{}} {} {}
disebut
\definitionswort {produk eksterior ke-$r$}{}
dari bundel vektor $E$. Bundel ini dilambangkan dengan
\mathl{\bigwedge^r E}{}.

}
Jika deskripsi matriks $E$ diberikan, deskripsi matriks produk eksterior ke-$r$ diperoleh dengan menghimpun semua
\definitionsverweis {determinan}{}{}
dari
$r \times r$-\definitionsverweis {submatriks}{}{}
ke dalam suatu matriks.

\inputdefinition
{}
{

Untuk suatu
\definitionsverweis {bundel vektor riil}{}{}
$E$ ber-rank $m$ di atas suatu
\definitionsverweis {ruang topologis}{}{}
$X$, objek yang disebut \definitionsverweis {produk eksterior ke-$m$}{}{}, yaitu

\mathl{\bigwedge^m E}{} disebut
\definitionswort {bundel determinan}{}
dari $E$. Bundel ini dilambangkan dengan
\mathl{\det  E}{}.

}
Bundel determinan merupakan bundel garis. Deskripsi matriksnya diberikan oleh determinan.

\inputdefinition
{}
{

Untuk
\definitionsverweis {bundel-bundel vektor riil}{}{}
\mathkor {} {E} {dan} {F} {}
di atas suatu
\definitionsverweis {ruang topologis}{}{}
$X$ dengan trivialisasi
\maabbdisp {\alpha_i} {E {{|}}_{U_i } } { U_i \times \R^m
} {}
dan
\maabbdisp {\beta_i} {F{{|}}_{U_i } } { U_i \times \R^n
} {}
bundel vektor yang terkait dengan
\definitionsverweis {data perekatan}{}{}

\mavergleichskettedisp
{\vergleichskette
{G_i
}
{ =} { U_i \times  \operatorname{Hom}_{ \R } \left(   \R^m , \R^n   \right)
}
{ } {
}
{ } {
}
{ } {
}
}
{}{}{}
dan
\maabbdisp {\varphi_{ij}} {G_i {{|}}_{U_i \cap U_j }  } { G_j {{|}}_{U_i \cap U_j }
} {}
dengan

\mavergleichskettedisp
{\vergleichskette
{ \varphi_{ij} (\theta)
}
{ \defeq} { \left(  \beta_j  \circ \beta_i^{-1}  \right)    \circ \theta \circ \left(  \alpha_i \circ \alpha_j^{-1}  \right)
}
{ } {
}
{ } {
}
{ } {
}
}
{}{}{}
disebut
\definitionswort {bundel homomorfisme}{}
dari bundel-bundel vektor
\mathkor {} {E} {dan} {F} {.}
Bundel ini dilambangkan dengan
\mathl{\operatorname{Hom} \left(    E ,  F  \right)}{}.

}

\inputdefinition
{}
{

Untuk suatu
\definitionsverweis {bundel vektor riil}{}{}
$E$ di atas suatu
\definitionsverweis {ruang topologis}{}{}
$X$,
\definitionsverweis {bundel homomorfisme}{}{}

\mathl{\operatorname{Hom} \left(    E ,  X \times \R  \right)}{} disebut
\definitionswort {bundel dual}{}
dari $E$. Bundel ini dilambangkan dengan
\mathl{E ^*}{}.

}

  \zwischenueberschrift{Bundel vektor diferensiabel}

\inputdefinition
{}
{

Misalkan $X$ suatu
\definitionsverweis {manifold diferensiabel}{}{}
dan

\mavergleichskette
{\vergleichskette
{r
}
{ \in }{\N
}
{  }{
}
{    }{
}
{   }{
}
}
{}{}{.}
Sebuah
\definitionswort {bundel vektor diferensiabel}{}
$V$ ber-
\definitionswort {rank}{}
$r$ adalah sebuah manifold yang dilengkapi dengan suatu
\definitionsverweis {pemetaan diferensiabel}{}{}
\maabb {p} { V } { X
} {}
sedemikian sehingga setiap
\definitionsverweis {serat}{}{}
$p^{-1}(x)$ merupakan
\definitionsverweis {ruang vektor riil}{}{}
yang
\definitionsverweis {berdimensi}{}{}
$r$ dan terdapat suatu
\definitionsverweis {tutupan terbuka}{}{}

\mavergleichskette
{\vergleichskette
{X
}
{ = }{ \bigcup_{i \in I} U_i
}
{  }{
}
{    }{
}
{   }{
}
}
{}{}{}
beserta
\definitionsverweis {difeomorfisme-difeomorfisme}{}{}
\maabbdisp {\varphi_i} {p^{-1} \left(  U_i  \right) } { U_i \times \R^r
} {}
di atas $U_i$ yang pada setiap serat menginduksi suatu
\definitionsverweis {isomorfisme linear}{}{}
\maabbdisp {\left(  \varphi_i  \right)_x} {p^{-1} (x) } { \R^r
} {}
.

}

  \zwischenueberschrift{Orientasi pada manifold}

\inputdefinition
{}
{

Misalkan $M$ sebuah
\definitionsverweis {manifold diferensiabel}{}{.}
Sebuah
\definitionsverweis {peta}{}{}
\maabbdisp {\alpha} { U } { V
} {}
dengan

\mavergleichskette
{\vergleichskette
{ U
}
{ \subseteq }{ M
}
{  }{
}
{    }{
}
{   }{
}
}
{}{}{}
dan

\mavergleichskette
{\vergleichskette
{ V
}
{ \subseteq }{ \R^n
}
{  }{
}
{    }{
}
{   }{
}
}
{}{}{}
yang terbuka disebut \definitionswort {berorientasi}{,} jika $\R^n$
\definitionsverweis {berorientasi}{}{.}

}

Jika tersedia suatu atlas yang terdiri atas peta-peta berorientasi
\mathl{(U_i,V_i, \alpha_i)}{}, maka orientasi-orientasi pada ruang bilangan sekitar $\R^n$ yang memuat citra terbuka $V_i$ dari peta-peta tersebut pada awalnya tidak saling berkaitan
\zusatzklammer {meskipun kita selalu menuliskan $\R^n$} {} {.}
Kaitan antarorientasi baru dapat dirumuskan melalui dua konsep berikut.

\inputdefinition
{}
{

Misalkan $M$ sebuah
\definitionsverweis {manifold diferensiabel}{}{}
dan misalkan
\mathkor {} {(U_1,V_1, \alpha_1)} {dan} {(U_2,V_2, \alpha_2)} {}
\definitionsverweis {peta-peta berorientasi}{}{.}
Maka
\definitionsverweis {pergantian peta}{}{}
yang terkait
\maabbdisp {\psi=\alpha_2 \circ \alpha_1^{-1}} {  \alpha_1 (U_1 \cap U_2)} {  \alpha_2 (U_1 \cap U_2 )
} {}
disebut
\definitionswort {mempertahankan orientasi}{,} jika untuk setiap titik

\mavergleichskette
{\vergleichskette
{ Q
}
{ \in }{ \alpha_1 (U_1 \cap U_2)
}
{  }{
}
{    }{
}
{   }{
}
}
{}{}{}
\definitionsverweis {diferensial total}{}{}
\maabbdisp {\left(D \psi\right)_{Q}} {\R^n } {\R^n
} {}
\definitionsverweis {mempertahankan orientasi}{}{}
.

}

\inputdefinition
{}
{

Sebuah
\definitionsverweis {manifold diferensiabel}{}{}
$M$ dengan suatu
\definitionsverweis {atlas}{}{}

\mathl{(U_i, V_i, \alpha_i)}{} disebut \definitionswort {berorientasi}{,} jika setiap peta
\definitionsverweis {berorientasi}{}{}
dan semua pergantian peta
\definitionsverweis {mempertahankan orientasi}{}{.}

}

\bild{ \begin{center}
\includegraphics[width=5.5cm]{ \bildeinlesungjpg {Möbius strip} {jpg} }
\end{center}
\bildtext {Pita Möbius adalah contoh khas manifold yang tidak dapat diorientasikan. Agar menjadi manifold, tepinya tidak boleh termasuk di dalamnya; tetapi dengan demikian pita ini juga bukan submanifold tertutup dari $\R^3$, sebab submanifold tertutup seperti itu selalu dapat diorientasikan.} }

\bildlizenz { Möbius strip.jpg } {} {Dbenbenn} {Commons} {CC-by-sa 3.0} {}

Pada sebuah manifold berorientasi, setiap ruang singgung
\mathl{T_PM}{}  memiliki suatu orientasi. Kita cukup memilih sebarang lingkungan koordinat

\mavergleichskette
{\vergleichskette
{ P
}
{ \in }{ U
}
{  }{
}
{    }{
}
{   }{
}
}
{}{}{}
\zusatzklammer {dari atlas berorientasi} {} {}
dan mentranspor orientasi pada

\mavergleichskettedisp
{\vergleichskette
{ V
}
{ \subseteq} { \R^n
}
{ } {
}
{ } {
}
{ } {
}
}
{}{}{}
melalui
\mathl{T_P(\alpha^{-1})}{} ke
\mathl{T_PM}{}. Karena pergantian peta mempertahankan orientasi, orientasi ini tidak bergantung pada lingkungan koordinat yang dipilih.

Pada sebuah manifold berorientasi, kita juga dapat menentukan apakah dua basis dalam ruang-ruang singgung pada dua titik yang berbeda merepresentasikan orientasi yang sama. Hal ini terjadi jika kedua basis merepresentasikan orientasi manifold tersebut atau jika keduanya sama-sama tidak merepresentasikannya.

Sebuah manifold disebut \stichwort {dapat diorientasikan} {,} jika manifold tersebut difeomorfik dengan suatu manifold berorientasi. Artinya, terdapat sebuah atlas yang mendefinisikan struktur diferensiabel yang sama dan juga dapat diberi orientasi.

\inputbemerkung
{}
{

Misalkan

\mavergleichskette
{\vergleichskette
{ W
}
{ \subseteq }{ \R^n
}
{  }{
}
{    }{
}
{   }{
}
}
{}{}{}
\definitionsverweis {terbuka}{}{,}
\maabb {h} { W } { \R
} {}
suatu
\definitionsverweis {fungsi yang terdiferensialkan secara kontinu}{}{}
dan

\mavergleichskette
{\vergleichskette
{ Y
}
{ = }{ h^{-1}(c)
}
{  }{
}
{    }{
}
{   }{
}
}
{}{}{}
merupakan
\definitionsverweis {serat}{}{}
di atas

\mavergleichskette
{\vergleichskette
{ c
}
{ \in }{ \R
}
{  }{
}
{    }{
}
{   }{
}
}
{}{}{,}
dengan $h$
\definitionsverweis {reguler}{}{}
di setiap titik dari $Y$. Pada ruang sekitar $\R^n$, kita tetapkan suatu
\definitionsverweis {orientasi}{}{}
\zusatzklammer {misalnya
\definitionsverweis {orientasi standar}{}{}} {} {}
dan kita tetapkan suatu
\definitionsverweis {medan normal satuan}{}{}
$N$ pada $Y$. Pilihan ini menetapkan suatu
\definitionsverweis {orientasi}{}{}
pada $Y$
\zusatzklammer {sebagai manifold} {} {}
dengan menetapkan pada setiap ruang singgung $T_PY$ bahwa suatu
\definitionsverweis {basis}{}{}

\mavergleichskette
{\vergleichskette
{ v_1 , \ldots , v_{n-1}
}
{ \in }{ T_PY
}
{  }{
}
{    }{
}
{   }{
}
}
{}{}{}
merepresentasikan orientasi jika

\mavergleichskette
{\vergleichskette
{ N(P) , v_1 , \ldots , v_{n-1}
}
{ \in }{ \R^n
}
{  }{
}
{    }{
}
{   }{
}
}
{}{}{}
merepresentasikan orientasi $\R^n$.

}

  \zwischenueberschrift{Kekompakan}

Himpunan bagian dari suatu ruang Euklides yang sekaligus tertutup dan terbatas disebut kompak. Pada ruang topologis yang tidak dilengkapi metrik, kita tidak dapat berbicara tentang keterbatasan; bahkan pada ruang metrik yang bukan himpunan bagian dari suatu $\R^n$, kedua sifat tertutup dan terbatas tidak banyak membantu. Konsep berikut lebih berdaya guna.

\inputdefinition
{}
{

Suatu
\definitionsverweis {ruang topologis}{}{}
$X$ disebut \definitionswort {kompak}{}
\zusatzklammer {atau \definitionswort {kompak terhadap tutupan}{}} {} {,}
jika untuk setiap tutupan terbuka

\mathdisp {X= \bigcup_{i \in I} U_i \, \, \, \text{ dengan } U_i \text{ terbuka dan untuk suatu himpunan indeks sebarang }I} {  }

terdapat suatu himpunan bagian hingga

\mavergleichskette
{\vergleichskette
{J
}
{ \subseteq }{ I
}
{  }{
}
{    }{
}
{   }{
}
}
{}{}{}
sedemikian sehingga

\mavergleichskettedisp
{\vergleichskette
{ X
}
{ =} { \bigcup_{i \in J} U_i
}
{ } {
}
{ } {
}
{ } {
}
}
{}{}{}
.

}

Sifat ini kadang-kadang juga disebut \stichwort {kompak terhadap tutupan} {.} Sering kali sifat Hausdorff juga disertakan dalam pengertian kompak. Perlu ditekankan bahwa sifat ini
\betonung{tidak}{} menyatakan adanya suatu tutupan hingga yang terdiri atas himpunan-himpunan terbuka
\zusatzklammer {tutupan terbuka trivial yang hanya terdiri atas seluruh ruang selalu ada} {} {,}
melainkan menyatakan bahwa dari sebarang tutupan terbuka, dengan sebarang himpunan indeks, cukup dipilih suatu himpunan bagian hingga dari himpunan indeks itu untuk tetap menutupi ruang tersebut.



\inputfaktbeweis
{Topologischer Raum/Abzählbare Basis/Überdeckungskompakt und folgenkompakt/Fakt}
{Lema}
{}
{

\faktsituation {Misalkan $X$ suatu
\definitionsverweis {ruang topologis}{}{}
dengan suatu
\definitionsverweis {basis terhitung}{}{.}}
\faktfolgerung {Maka $X$
\definitionsverweis {kompak}{}{}
jika dan hanya jika setiap barisan
\mathl{\left(  x_n  \right)_{n \in  \N  }}{} di $X$ memiliki suatu
\definitionsverweis {titik akumulasi}{}{}
\zusatzklammer {di $X$} {} {}.}
\faktzusatz {}
\faktzusatz {}

}
{

\teilbeweis {}{}{}
{Misalkan $X$ kompak dan diberikan suatu barisan
\mathl{\left(  x_n  \right)_{n \in  \N  }}{}.
Anggap bahwa barisan ini tidak memiliki titik akumulasi. Artinya, untuk setiap

\mavergleichskette
{\vergleichskette
{ y
}
{ \in }{ X
}
{  }{
}
{    }{
}
{   }{
}
}
{}{}{}
terdapat suatu lingkungan terbuka

\mavergleichskette
{\vergleichskette
{ y
}
{ \in }{ U_y
}
{  }{
}
{    }{
}
{   }{
}
}
{}{}{}
yang hanya memuat sejumlah hingga suku barisan. Karena

\mavergleichskettedisp
{\vergleichskette
{ X
}
{ =} { \bigcup_{y \in X} U_y
}
{ } {
}
{ } {
}
{ } {
}
}
{}{}{}
menurut asumsi terdapat suatu subtutupan hingga

\mavergleichskettedisp
{\vergleichskette
{ X
}
{ =} { \bigcup_{i = 1}^n U_{y_i }
}
{ } {
}
{ } {
}
{ } {
}
}
{}{}{.} Tutupan ini di satu pihak memuat semua suku barisan, tetapi di pihak lain hanya memuat sejumlah hingga suku barisan, suatu kontradiksi.}
{}

\teilbeweis {}{}{}
{Misalkan sifat barisan tersebut terpenuhi dan misalkan

\mavergleichskette
{\vergleichskette
{ X
}
{ = }{ \bigcup_{i \in I} U_i
}
{  }{
}
{    }{
}
{   }{
}
}
{}{}{}
suatu tutupan oleh himpunan-himpunan terbuka. Karena $X$ memiliki suatu
\definitionsverweis {basis terhitung}{}{,}
berdasarkan
Soal 2.8 (Teori Ukuran dan Integrasi (Osnabrück 2022-2023))
terdapat suatu himpunan bagian terhitung

\mavergleichskette
{\vergleichskette
{ J
}
{ \subseteq }{ I
}
{  }{
}
{    }{
}
{   }{
}
}
{}{}{}
dengan

\mavergleichskettedisp
{\vergleichskette
{ X
}
{ =} { \bigcup_{i \in J} U_i
}
{ } {
}
{ } {
}
{ } {
}
}
{}{}{.}
Kita dapat mengambil

\mavergleichskette
{\vergleichskette
{ J
}
{ = }{ \N
}
{  }{
}
{    }{
}
{   }{
}
}
{}{}{}
sebagai asumsi. Anggap bahwa tutupan

\mavergleichskette
{\vergleichskette
{ X
}
{ = }{ \bigcup_{i \in \N} U_i
}
{  }{
}
{    }{
}
{   }{
}
}
{}{}{}
tidak memiliki subtutupan hingga. Maka, khususnya,

\mavergleichskette
{\vergleichskette
{ \bigcup_{i = 0}^n  U_i
}
{ \neq }{ X
}
{  }{
}
{    }{
}
{   }{
}
}
{}{}{}
untuk setiap

\mavergleichskette
{\vergleichskette
{ n
}
{ \in }{ \N
}
{  }{
}
{    }{
}
{   }{
}
}
{}{}{}
dan oleh karena itu untuk setiap

\mavergleichskette
{\vergleichskette
{ n
}
{ \in }{ \N
}
{  }{
}
{    }{
}
{   }{
}
}
{}{}{}
terdapat suatu

\mathbed {x_n \in X} {dengan}
 {x_n \not\in \bigcup_{i =0}^n  U_i} {}
 {} {} {} {.}
Menurut asumsi, barisan ini memiliki suatu titik akumulasi $x$. Karena

\mavergleichskette
{\vergleichskette
{ X
}
{ = }{ \bigcup_{i \in \N} U_i
}
{  }{
}
{    }{
}
{   }{
}
}
{}{}{}
merupakan suatu tutupan, terdapat suatu

\mavergleichskette
{\vergleichskette
{ k
}
{ \in }{ \N
}
{  }{
}
{    }{
}
{   }{
}
}
{}{}{}
dengan

\mavergleichskette
{\vergleichskette
{ x
}
{ \in }{ U_k
}
{  }{
}
{    }{
}
{   }{
}
}
{}{}{.}
Karena $x$ merupakan titik akumulasi, terdapat tak hingga banyak suku barisan di $U_k$. Ini merupakan kontradiksi, sebab menurut konstruksi, untuk

\mavergleichskette
{\vergleichskette
{ n
}
{ \geq }{ k
}
{  }{
}
{    }{
}
{   }{
}
}
{}{}{}
suku-suku barisan $x_n$ tidak termasuk dalam $U_k$.}
{}

}

Teorema berikut disebut \stichwort {teorema Heine--Borel} {.}


\inputfaktbeweis
{Kompaktheit/Satz von Heine-Borel/Fakt}
{Teorema}
{}
{

\faktsituation {Misalkan

\mavergleichskette
{\vergleichskette
{T
}
{ \subseteq }{ \R^n
}
{  }{
}
{    }{
}
{   }{
}
}
{}{}{}
suatu himpunan bagian.}
\faktuebergang {Maka pernyataan-pernyataan berikut ekuivalen.}
\faktfolgerung {\aufzaehlungvier{ $T$
\definitionsverweis {kompak terhadap tutupan}{}{.}
}{Setiap
\definitionsverweis {barisan}{}{}

\mathl{\left(  x_n  \right)_{n \in  \N  }}{} di $T$ memiliki suatu
\definitionsverweis {titik akumulasi}{}{}
di $T$.
}{Setiap
\definitionsverweis {barisan}{}{}

\mathl{\left(  x_n  \right)_{n \in  \N  }}{} di $T$ memiliki suatu
\definitionsverweis {barisan bagian}{}{} yang \definitionsverweis {konvergen}{}{} di $T$.
}{ $T$
\definitionsverweis {tertutup}{}{}
dan
\definitionsverweis {terbatas}{}{.}
}}
\faktzusatz {}
\faktzusatz {}

}
{

\teilbeweis {}{}{}
{Ekuivalensi (1) dan (2) telah dibuktikan dalam bentuk yang lebih umum pada
Lema 13.13;
untuk keberadaan basis terhitung, lihat
Soal 13.17.}
{}
\teilbeweis {}{}{}
{Ekuivalensi (2) dan (3) jelas.}
{}
\teilbeweis {}{}{}
{Ekuivalensi (3) dan (4) telah ditunjukkan dalam
Teorema 36.9 (Analisis (Osnabrück 2021-2023)).}
{}

}

  \zwischenueberschrift{Ukuran pada manifold}

Misalkan $M$ sebuah manifold. Adakah pengertian volume yang bermakna untuk
\zusatzklammer {himpunan-himpunan bagian dari} {} {}
$M$, dan kapan suatu fungsi yang didefinisikan pada $M$ dapat diintegralkan secara bermakna? Jika teori ukuran dijadikan kerangka umum, diperoleh gambaran berikut. Anggap bahwa $M$ memiliki suatu atlas terhitung
\mathl{(U_i,V_i,\alpha_i, i \in I)}{}. Suatu ukuran $\mu$ pada himpunan-himpunan Borel
\mathl{{\mathcal B }(M)}{} ditentukan secara tunggal oleh pembatasan-pembatasan

\mavergleichskette
{\vergleichskette
{ \mu_i
}
{ = }{ \mu \vert_{U_i }
}
{  }{
}
{    }{
}
{   }{
}
}
{}{}{}
dari ukuran tersebut pada himpunan-himpunan bagian terbuka $U_i$. Untuk setiap

\mavergleichskette
{\vergleichskette
{ i
}
{ \in }{ I
}
{  }{
}
{    }{
}
{   }{
}
}
{}{}{}
homeomorfisme
\maabbdisp {\alpha_i} {U_i } { V_i
} {}
mendefinisikan ukuran citra

\mavergleichskette
{\vergleichskette
{ \nu_i
}
{ = }{ {\alpha_i }_* \mu_i
}
{  }{
}
{    }{
}
{   }{
}
}
{}{}{}
pada

\mavergleichskette
{\vergleichskette
{ V_i
}
{ \subseteq }{ \R^n
}
{  }{
}
{    }{
}
{   }{
}
}
{}{}{.}
Ukuran-ukuran citra

\mathbed {\nu_i} {}
 {i \in I} {}
 {} {} {} {,}
saling berhubungan melalui

\mavergleichskettedisp
{\vergleichskette
{\nu_i ( \alpha_i(T))
}
{ =} { \mu(T)
}
{ =} {\nu_j( \alpha_j(T))
}
{ } {
}
{ } {
}
}
{}{}{}
untuk setiap himpunan bagian terukur

\mavergleichskette
{\vergleichskette
{ T
}
{ \subseteq }{ U_i \cap U_j
}
{  }{
}
{    }{
}
{   }{
}
}
{}{}{.}
Dengan pergantian peta

\mavergleichskette
{\vergleichskette
{ \psi_{ij}
}
{ = }{ \alpha_j \circ \alpha_i^{-1}
}
{  }{
}
{    }{
}
{   }{
}
}
{}{}{}
hal ini berarti

\mavergleichskettedisp
{\vergleichskette
{ \nu_i(S)
}
{ =} { \nu_j ( \psi_{ij} (S))
}
{ } {
}
{ } {
}
{ } {
}
}
{}{}{}
untuk setiap himpunan terukur

\mavergleichskette
{\vergleichskette
{ S
}
{ \subseteq }{ V_i
}
{  }{
}
{    }{
}
{   }{
}
}
{}{}{,}
yang seluruhnya terletak di dalam domain pemetaan transisi.

Sekarang anggap bahwa setiap ukuran citra $\nu_i$ dapat dituliskan menggunakan suatu
\definitionsverweis {kerapatan}{}{}
terhadap
\definitionsverweis {ukuran Borel--Lebesgue}{}{}
$\lambda^n$, misalnya

\mavergleichskettedisp
{\vergleichskette
{ \nu_i
}
{ =} { g_i d \lambda^n
}
{ } {
}
{ } {
}
{ } {
}
}
{}{}{,}
dengan
\definitionsverweis {fungsi-fungsi yang dapat diintegralkan}{}{}
yang didefinisikan pada $V_i$
dan dinyatakan oleh
\maabb {g_i} {V_i } {\R
} {.}
Untuk suatu himpunan bagian terukur

\mavergleichskette
{\vergleichskette
{ T
}
{ \subseteq }{ U_i
}
{  }{
}
{    }{
}
{   }{
}
}
{}{}{}
berlaku

\mavergleichskettedisp
{\vergleichskette
{ \mu(T)
}
{ =} { \nu_i(\alpha_i(T))
}
{ =} { \int_{ \alpha_i(T)  }  g_i    \,  d  \lambda^n
}
{ } {
}
{ } {
}
}
{}{}{.}
Untuk suatu himpunan bagian terukur

\mavergleichskette
{\vergleichskette
{ T
}
{ \subseteq }{ U_i \cap U_j
}
{  }{
}
{    }{
}
{   }{
}
}
{}{}{}
maka, menurut
rumus transformasi,
yang diterapkan pada pemetaan transisi difeomorfik
\maabbdisp {\psi_{ij}} { \alpha_i(U_i \cap U_j) } { \alpha_j(U_i \cap U_j)
} {,}
yang membawa
\mathl{\alpha_i(T)}{} ke
\mathl{\alpha_j(T)}{}, berlaku kesamaan

\mavergleichskettealign
{\vergleichskettealign
{ \int_{ \alpha_i(T)  }  g_i    \,  d  \lambda^n
}
{ =} { \int_{ \alpha_j(T)  }  g_j    \,  d  \lambda^n
}
{ =} { \int_{ \alpha_i(T)  }   \betrag {  \det  \left(  D  \psi_{ij}  \right)  } \cdot  ( g_j \circ  \psi_{ij} )   \,  d  \lambda^n
}
{ } {
}
{ } {
}
}
{}
{}{.}
Hal ini menyarankan, bagi fungsi-fungsi kerapatan

\mathbed {g_i} {}
 {i \in I} {}
 {} {} {} {,}
perilaku transformasi

\mavergleichskettedisp
{\vergleichskette
{ g_i
}
{ =} { \betrag {  \det  \left(  D  \psi_{ij}  \right)  } \cdot ( g_j \circ \psi_{ij} )
}
{ } {
}
{ } {
}
{ } {
}
}
{}{}{}
\zusatzklammer {meskipun perilaku ini tidak dipaksakan, karena suatu kerapatan tidak ditentukan secara tunggal oleh ukurannya} {} {.}
Kita akan membangun teori integrasi pada manifold berdasarkan konsep bentuk diferensial-$n$, yang secara alami memiliki perilaku transformasi ini
\zusatzklammer {tanpa nilai mutlak} {} {}.

\chapter{Lembar Kerja 13}
\setcounter{section}{13}

  \zwischenueberschrift{Soal latihan}

\definitionswort {Produk Kronecker}{}
dari matriks-matriks

\mavergleichskettedisp
{\vergleichskette
{ A
}
{ =} { \left(  a_{ij}  \right)_{1 \leq i \leq m,\, 1 \leq j \leq n}
}
{ } {
}
{ } {
}
{ } {
}
}
{}{}{}
dan

\mavergleichskettedisp
{\vergleichskette
{ B
}
{ =} { \left(  b_{k \ell }  \right)_{1 \leq k \leq p,\, 1 \leq \ell \leq r}
}
{ } {
}
{ } {
}
{ } {
}
}
{}{}{}
diberikan oleh

\mathdisp {\left(  a_{ij} \cdot b_{k \ell}  \right)_{1 \leq i \leq m,\,     1 \leq  k \leq p;\,   1 \leq j \leq n,     \, 1 \leq \ell \leq r   }} {  }

\inputaufgabegibtloesung
{}
{

Hitunglah
\definitionsverweis {produk Kronecker}{}{}
dari kedua matriks
\mathkor {} {\begin{pmatrix}  3   &   -4  \\  5  &  -2 \end{pmatrix}} {dan} {\begin{pmatrix}  -2  &   7   \\  6  &  3 \end{pmatrix}} {.}

}
{} {}

\inputaufgabe
{}
{

Misalkan $K$ sebuah
\definitionsverweis {lapangan}{}{}
dan

\mavergleichskettedisp
{\vergleichskette
{ A
}
{ =} { \left(  a_{ij}  \right)_{1 \leq i \leq m,\, 1 \leq j \leq n}
}
{ } {
}
{ } {
}
{ } {
}
}
{}{}{}
serta

\mavergleichskettedisp
{\vergleichskette
{ B
}
{ =} { \left(  b_{k \ell }  \right)_{1 \leq  k \leq p,\, 1 \leq \ell \leq r}
}
{ } {
}
{ } {
}
{ } {
}
}
{}{}{}
\definitionsverweis {matriks-matriks}{}{}
dengan
\definitionsverweis {pemetaan-pemetaan linear}{}{}
yang bersesuaian
\maabb {A} {K^n  } { K^m
} {}
dan
\maabb {B} {K^r } { K^p
} {.}
Tunjukkan bahwa
\definitionsverweis {produk tensor}{}{}
dari pemetaan-pemetaan linear tersebut, terhadap basis-basis

\mathbed {e_j \otimes e_\ell} {}
 {1 \leq j \leq n,\, 1 \leq \ell \leq r} {}
 {} {} {} {,}
dari
\mathl{K^n \otimes K^r}{} dan

\mathbed {e_i \otimes e_k} {}
 {1 \leq i \leq m,\, 1 \leq  k \leq p} {}
 {} {} {} {,}
dari
\mathl{K^m \otimes K^p}{,}
dideskripsikan oleh
\definitionsverweis {produk Kronecker}{}{}
dari $A$ dan $B$.

}
{} {}

\inputaufgabe
{}
{

Tunjukkan bahwa
\definitionsverweis {produk tensor}{}{}
dari
\definitionsverweis {pita Möbius}{}{}
dengan dirinya sendiri merupakan sebuah bundel garis trivial.

}
{} {}

\inputaufgabe
{}
{

Misalkan
\mathkor {} {M_1} {dan} {M_2} {}
\definitionsverweis {manifold-manifold berorientasi}{}{.}
Tunjukkan bahwa
\definitionsverweis {produk}{}{}

\mathl{M_1 \times M_2}{} merupakan sebuah manifold berorientasi
\zusatzklammer {dengan orientasi yang bergantung pada urutan pada \mathlk{\{1,2\}}{}} {} {.}

}
{} {}

\inputaufgabe
{}
{

Tunjukkan bahwa sfera berdimensi $1$, yaitu $S^1$, merupakan sebuah
\definitionsverweis {manifold diferensiabel}{}{}
yang
\definitionsverweis {dapat diorientasikan}{}{.}

}
{} {}

Misalkan
\mathkor {} {M} {dan} {N} {}
\definitionsverweis {manifold-manifold berorientasi}{}{}
dan
\maabbdisp {\varphi} { M } { N
} {}
sebuah
\definitionsverweis {pemetaan diferensiabel}{}{.}
Pemetaan $\varphi$ disebut \definitionswort {mempertahankan orientasi}{,} jika untuk setiap titik

\mavergleichskette
{\vergleichskette
{ P
}
{ \in }{ M
}
{  }{
}
{    }{
}
{   }{
}
}
{}{}{}
\definitionsverweis {pemetaan singgung}{}{}
\maabbdisp {T\varphi} { T_PM } { T_{\varphi(P)} N
} {}
bijektif dan
\definitionsverweis {mempertahankan orientasi}{}{.}

\inputaufgabe
{}
{

Tunjukkan bahwa
\definitionsverweis {pemetaan antipodal}{}{}
\maabbeledisp {\varphi} {S^1} {S^1
} { P } { -P
} {,}
\definitionsverweis {mempertahankan orientasi}{}{.}

}
{} {}

\inputaufgabe
{}
{

Misalkan $M$ sebuah
\definitionsverweis {manifold berorientasi}{}{.}
Tunjukkan bahwa pemetaan pertukaran
\maabbeledisp {\varphi} {M \times M} {M \times M
} {(P,Q)} {(Q,P)
} {,}
terhadap masing-masing
\definitionsverweis {orientasi produk}{}{,}
tidak harus
\definitionsverweis {mempertahankan orientasi}{}{.}

}
{} {}

\inputaufgabe
{}
{

Misalkan $X$ sebuah
\definitionsverweis {ruang topologis}{}{}
yang hanya terdiri atas
\definitionsverweis {berhingga banyak}{}{}
elemen. Tunjukkan bahwa $X$
\definitionsverweis {kompak}{}{.}

}
{} {}

\inputaufgabe
{}
{

Misalkan $X$ sebuah
\definitionsverweis {ruang topologis}{}{}
dan misalkan

\mavergleichskette
{\vergleichskette
{Y_1 , \ldots , Y_n
}
{ \subseteq }{ X
}
{  }{
}
{    }{
}
{   }{
}
}
{}{}{}
\definitionsverweis {himpunan-himpunan bagian kompak}{}{.}
Tunjukkan bahwa
\definitionsverweis {gabungan}{}{}

\mavergleichskette
{\vergleichskette
{ Y
}
{ = }{ \bigcup_{i = 1}^n Y_i
}
{  }{
}
{    }{
}
{   }{
}
}
{}{}{}
juga kompak.

}
{} {}

\inputaufgabegibtloesung
{}
{

Misalkan $X$ sebuah
\definitionsverweis {ruang kompak}{}{}
dan misalkan

\mavergleichskette
{\vergleichskette
{Y
}
{ \subseteq }{ X
}
{  }{
}
{    }{
}
{   }{
}
}
{}{}{}
sebuah
\definitionsverweis {himpunan bagian tertutup}{}{,}
yang dilengkapi dengan
\definitionsverweis {topologi terinduksi}{}{.}
Tunjukkan bahwa $Y$ juga kompak.

}
{} {}

\inputaufgabegibtloesung
{}
{

Misalkan $X$ sebuah
\definitionsverweis {ruang topologis}{}{}
\definitionsverweis {kompak}{}{}
yang tak kosong dan misalkan
\maabbdisp {f} { X } {\R
} {}
sebuah
\definitionsverweis {fungsi kontinu}{}{.}
Tunjukkan bahwa terdapat suatu

\mavergleichskette
{\vergleichskette
{ x
}
{ \in }{ X
}
{  }{
}
{    }{
}
{   }{
}
}
{}{}{}
dengan

\mathdisp {f(x) \geq f(x')  \text { untuk semua } x' \in X} {  }
.

}
{} {}

\inputaufgabe
{}
{

Misalkan
\maabbdisp {f} {\R} {S^1
} {}
sebuah
\definitionsverweis {pemetaan kontinu}{}{.}
Tunjukkan bahwa
\definitionsverweis {citra}{}{}
dari $f$
\definitionsverweis {homeomorfik}{}{}
dengan suatu
\definitionsverweis {interval}{}{}
terbuka, setengah terbuka, atau tertutup, atau dengan $S^1$.

}
{} {}

\inputaufgabe
{}
{

Misalkan
\maabbdisp {f} {S^1} {\R
} {}
sebuah
\definitionsverweis {pemetaan kontinu}{}{.}
Tunjukkan bahwa
\definitionsverweis {citra}{}{}
dari $f$
\definitionsverweis {homeomorfik}{}{}
dengan suatu
\definitionsverweis {interval tertutup}{}{.}

}
{} {}

\inputaufgabe
{}
{

Tunjukkan bahwa himpunan
\definitionsverweis {bilangan-bilangan riil}{}{} $\R$ tidak
\definitionsverweis {kompak terhadap tutupan}{}{.}

}
{} {}

\inputaufgabe
{}
{

Kita meninjau
\definitionsverweis {bilangan-bilangan asli}{}{}
$\N$ dan melengkapinya dengan
\definitionsverweis {metrik diskret}{}{.}
Tunjukkan bahwa $\N$
\definitionsverweis {
tertutup}{}{}
dan
\definitionsverweis {terbatas}{}{,}
tetapi tidak
\definitionsverweis {
kompak terhadap tutupan}{}{.}

}
{} {}

\inputaufgabegibtloesung
{}
{

Misalkan $X$ sebuah
\definitionsverweis {ruang metrik}{}{}
\definitionsverweis {kompak}{}{.}
Tunjukkan bahwa $X$
\definitionsverweis {
lengkap}{}{.}

}
{} {}

\inputaufgabe
{}
{

Tunjukkan bahwa terdapat sebuah
\definitionsverweis {keluarga terhitung}{}{}
dari
\definitionsverweis {bola-bola terbuka}{}{}
di $\R^n$ yang membentuk sebuah
\definitionsverweis {basis topologi}{}{.}

}
{} {}

\inputaufgabegibtloesung
{}
{

Tunjukkan bahwa himpunan

\mavergleichskettedisp
{\vergleichskette
{ M
}
{ =} { \left\{ (x,y,z) \in \R^3  \mid   x^2+y^4+z^6 = 1        \right\}
}
{ } {
}
{ } {
}
{ } {
}
}
{}{}{}
merupakan sebuah manifold diferensiabel kompak berdimensi dua.

}
{} {}

\inputaufgabegibtloesung
{}
{

Misalkan $M$ sebuah manifold topologis kompak berdimensi $d$,

\mavergleichskette
{\vergleichskette
{ d
}
{ \geq }{ 1
}
{  }{
}
{    }{
}
{   }{
}
}
{}{}{.}
Tunjukkan bahwa terdapat sebuah himpunan bagian terbuka terbatas

\mavergleichskette
{\vergleichskette
{ U
}
{ \subseteq }{ \R^d
}
{  }{
}
{    }{
}
{   }{
}
}
{}{}{}
dan sebuah pemetaan kontinu surjektif
\maabbdisp {\varphi} { U } { M
} {}.

}
{} {}

  \zwischenueberschrift{Soal untuk dikumpulkan}

\inputaufgabe
{4}
{

Tunjukkan bahwa sfera berdimensi $2$, yaitu $S^2$, merupakan sebuah
\definitionsverweis {manifold diferensiabel}{}{}
yang
\definitionsverweis {dapat diorientasikan}{}{.}

}
{} {}

\inputaufgabegibtloesung
{4}
{

Tunjukkan bahwa
\definitionsverweis {pemetaan antipodal}{}{}
\maabbeledisp {} {S^2} {S^2
} {(x,y,z)} {(-x,-y,-z)
} {,}
tidak
\definitionsverweis {mempertahankan orientasi}{}{.}

}
{} {}

\inputaufgabegibtloesung
{4}
{

Misalkan $X$ sebuah
\definitionsverweis {ruang Hausdorff}{}{}
dan misalkan

\mavergleichskette
{\vergleichskette
{Y
}
{ \subseteq }{ X
}
{  }{
}
{    }{
}
{   }{
}
}
{}{}{}
sebuah himpunan bagian yang dilengkapi dengan
\definitionsverweis {topologi terinduksi}{}{.}
Misalkan $Y$
\definitionsverweis {kompak}{}{.}
Tunjukkan bahwa $Y$
\definitionsverweis {tertutup}{}{}
di $X$.

}
{} {}

\inputaufgabe
{4}
{

Misalkan
\mathkor {} {X} {dan} {Y} {}
\definitionsverweis {ruang-ruang topologis}{}{}
dan misalkan
\maabbdisp {\varphi} { X } { Y
} {}
sebuah
\definitionsverweis {pemetaan kontinu}{}{.}
Misalkan $X$
\definitionsverweis {kompak}{}{.}
Tunjukkan bahwa
\definitionsverweis {citra}{}{}

\mavergleichskette
{\vergleichskette
{ \varphi(X)
}
{ \subseteq }{ Y
}
{  }{
}
{    }{
}
{   }{
}
}
{}{}{}
juga kompak.

}
{} {}

\inputaufgabe
{4}
{

Misalkan $V$ sebuah
\definitionsverweis {ruang vektor Euklides}{}{}
ber-
\definitionsverweis {dimensi}{}{}
$n$ dan misalkan
\definitionsverweis {produk}{}{}

\mavergleichskette
{\vergleichskette
{ V^n
}
{ = }{ V \times \cdots \times V
}
{  }{
}
{    }{
}
{   }{
}
}
{}{}{}
dilengkapi dengan
\definitionsverweis {topologi produk}{}{.}
Misalkan $I$ sebuah
\definitionsverweis {interval riil}{}{}
dan
\maabbdisp {\varphi} { I } { V^n
} {}
sebuah
\definitionsverweis {pemetaan kontinu}{}{}
dengan sifat bahwa

\mavergleichskettedisp
{\vergleichskette
{ \varphi(t)
}
{ =} { (\varphi_1(t), \varphi_2(t) , \ldots , \varphi_n(t))
}
{ } {
}
{ } {
}
{ } {
}
}
{}{}{}
untuk setiap

\mavergleichskette
{\vergleichskette
{ t
}
{ \in }{ I
}
{  }{
}
{    }{
}
{   }{
}
}
{}{}{}
merupakan sebuah
\definitionsverweis {basis}{}{}
dari $V$. Tunjukkan bahwa semua basis

\mathbed {\varphi(t)} {}
 {t \in I} {}
 {} {} {} {,}
merepresentasikan
\definitionsverweis {orientasi}{}{}
yang sama pada $V$.

}
{} {}

\section*{Solusi yang disediakan oleh sumber}
\addcontentsline{toc}{section}{Solusi yang disediakan oleh sumber}
\subsection*{Solusi Soal 13.1}
Produk Kronecker dari kedua matriks tersebut adalah

\mavergleichskettealign
{\vergleichskettealign
{ \begin{pmatrix}  3  &   -4  \\  5 &  -2 \end{pmatrix} \otimes \begin{pmatrix}  -2  &   7  \\  6 &  3 \end{pmatrix}
}
{ =} { \begin{pmatrix}  3 \begin{pmatrix}  -2  &   7  \\  6 &  3 \end{pmatrix}   &   -4\begin{pmatrix}  -2  &   7  \\  6 &  3 \end{pmatrix}   \\  5\begin{pmatrix}  -2  &   7  \\  6 &  3 \end{pmatrix}  &  -2\begin{pmatrix}  -2  &   7  \\  6 &  3 \end{pmatrix} \end{pmatrix}
}
{ =} { \begin{pmatrix}  -6 &  21 &  8 &  -28 \\  18 &  9 &  -24 &  -12 \\  -10 &  35 &  4 &  -14 \\  30 &  15 &  -12 &  -6 \end{pmatrix}
}
{ } {
}
{ } {
}
}
{}
{}{.}
\subsection*{Solusi Soal 13.10}
Misalkan

\mavergleichskettedisp
{\vergleichskette
{ Y
}
{ =} { \bigcup_{i \in I} V_i
}
{ } {
}
{ } {
}
{ } {
}
}
{}{}{}
sebuah tutupan oleh himpunan-himpunan bagian terbuka

\mavergleichskette
{\vergleichskette
{ V_i
}
{ \subseteq }{ Y
}
{  }{
}
{    }{
}
{   }{
}
}
{}{}{.}
Ini berarti terdapat himpunan-himpunan terbuka

\mavergleichskette
{\vergleichskette
{ U_i
}
{ \subseteq }{ X
}
{  }{
}
{    }{
}
{   }{
}
}
{}{}{}
sedemikian sehingga

\mavergleichskette
{\vergleichskette
{ V_i
}
{ = }{ Y \cap U_i
}
{  }{
}
{    }{
}
{   }{
}
}
{}{}{.}
Karena itu,

\mavergleichskettedisp
{\vergleichskette
{ Y
}
{ \subseteq} { \bigcup_{i \in I} U_i
}
{ } {
}
{ } {
}
{ } {
}
}
{}{}{.}
Karena $Y$ tertutup di $X$, maka $X \setminus Y$ terbuka, sehingga

\mavergleichskettedisp
{\vergleichskette
{ X
}
{ =} { Y \cup (X \setminus Y)
}
{ =} { \bigcup_{i \in I} U_i  \cup  (X \setminus Y)
}
{ } {
}
{ } {
}
}
{}{}{}
merupakan sebuah tutupan terbuka dari $X$. Karena $X$ kompak, terdapat sebuah subtutupan hingga

\mavergleichskettedisp
{\vergleichskette
{ X
}
{ =} { \bigcup_{i \in J} U_i  \cup  (X \setminus Y)
}
{ } {
}
{ } {
}
{ } {}
}
{}{}{.}
Hal ini pada gilirannya berarti

\mavergleichskettedisp
{\vergleichskette
{ Y
}
{ =} { \bigcup_{i \in J} V_i
}
{ =} { Y \cap \bigcup_{i \in J} U_i
}
{ } {
}
{ } {}
}
{}{}{.}
\subsection*{Solusi Soal 13.11}
Berdasarkan
Fakta,

\mavergleichskette
{\vergleichskette
{ f(X)
}
{ \subseteq }{ \R
}
{  }{
}
{    }{
}
{   }{
}
}
{}{}{}
kompak, sehingga menurut
Fakta
himpunan ini
\definitionsverweis {tertutup}{}{}
dan
\definitionsverweis {terbatas}{}{.}
Secara khusus,

\mavergleichskette
{\vergleichskette
{ f(X)
}
{ \subseteq }{ (-\infty,M]
}
{  }{
}
{    }{
}
{   }{
}
}
{}{}{}
untuk suatu bilangan riil $M$. Karena

\mavergleichskette
{\vergleichskette
{ X
}
{ \neq }{ \emptyset
}
{  }{
}
{    }{
}
{   }{
}
}
{}{}{}
menurut
Fakta,
$f(X)$ memiliki sebuah
\definitionsverweis {supremum}{}{}
$s$ di $\R$. Karena himpunan itu tertutup, menurut Fakta,
supremum tersebut termasuk dalam $f(X)$, sehingga merupakan maksimum dari $f(X)$. Oleh karena itu, terdapat pula suatu

\mavergleichskette
{\vergleichskette
{ x
}
{ \in }{ X
}
{  }{
}
{    }{
}
{   }{
}
}
{}{}{}
dengan

\mavergleichskette
{\vergleichskette
{ f(x)
}
{ = }{ s
}
{  }{
}
{    }{
}
{   }{
}
}
{}{}{.}
\subsection*{Solusi Soal 13.16}
Misalkan
\mathl{\left(   x _n  \right)_{n \in \N }}{} sebuah
\definitionsverweis {barisan Cauchy}{}{}
di $X$. Andaikan barisan ini tidak konvergen. Menurut
Soal,
barisan tersebut juga tidak memiliki titik akumulasi. Ini berarti bahwa untuk setiap titik

\mavergleichskette
{\vergleichskette
{ y
}
{ \in }{ X
}
{  }{
}
{    }{
}
{   }{
}
}
{}{}{}
terdapat sebuah lingkungan terbuka

\mavergleichskette
{\vergleichskette
{ y
}
{ \in }{ U_y
}
{  }{
}
{    }{
}
{   }{
}
}
{}{}{}
yang hanya memuat berhingga banyak suku barisan. Berdasarkan kekompakan, tutupan

\mavergleichskettedisp
{\vergleichskette
{ X
}
{ =} { \bigcup_{y \in X} U_y
}
{ } {
}
{ } {
}
{ } {
}
}
{}{}{}
memiliki sebuah subtutupan hingga, yakni

\mavergleichskettedisp
{\vergleichskette
{ X
}
{ =} { \bigcup_{ i = 1}^m  U_{y_i }
}
{ } {
}
{ } {
}
{ } {
}
}
{}{}{.}
Dengan demikian, mulai suatu $N$, semua suku barisan akan berada di luar himpunan ini, suatu kemustahilan.
\subsection*{Solusi Soal 13.18}
Kita meninjau pemetaan diferensiabel
\maabbeledisp {\varphi} {\R^3} {\R
} {(x,y,z)} { x^2+y^4+z^6-1
} {.}
Himpunan $M$ adalah serat $\varphi$ di atas $0$. Berlaku

\mathdisp {\left(D\varphi\right)_{(x,y,z)}  = (2x,4y^3,6z^5)} { . }

Turunan ini sama dengan $(0,0,0)$ hanya di
\mathl{(x,y,z)=(0,0,0)}{}, dan titik ini bukan anggota $M$, sehingga $\varphi$ reguler di setiap titik $M$. Karena itu, menurut teorema pemetaan implisit, himpunan tersebut merupakan sebuah manifold berdimensi dua.

Sebagai serat dari sebuah pemetaan kontinu, $M$ merupakan sebuah himpunan bagian tertutup dari $\R^3$. Selain itu, $M$ terbatas. Memang, untuk
\mathl{(x,y,z) \in M}{} berlaku
\mathl{\betrag {  x  } , \betrag {  y  } ,\betrag {  z  } \leq 1}{,} sebab jika tidak,
\mathl{x^2+y^4+z^6 >1}{}. Hal ini mengimplikasikan kekompakan.
\subsection*{Solusi Soal 13.19}
Untuk setiap titik
\mathl{P \in M}{} kita pilih sebuah lingkungan peta terbuka
\mathl{P \in U_P}{} dengan sebuah peta
\maabbdisp {\alpha_P} {U_P } {V_P
} {}
dan
\mathl{V_P \subseteq \R^d}{.} Dengan beralih ke sebuah lingkungan bola dari
\mathl{\alpha_P(P) \in V_P}{} beserta prabayangnya, kita dapat mengandaikan bahwa $V_P$ adalah bola-bola terbuka dengan jari-jari paling besar
\mathl{{ \frac{   1 }{  3 } }}{}. Keluarga

\mathbed {U_P} {}
 {P\in M} {}
 {} {} {} {,}
menutupi manifold tersebut. Karena $M$ kompak, terdapat berhingga banyak titik
\mathl{P_1 , \ldots , P_n}{} sedemikian sehingga

\mathbed {U_i=U_{P_i}} {}
 {i=1 , \ldots , n} {}
 {} {} {} {,}
juga menutupi manifold tersebut. Kita menempatkan bola-bola terbuka
\mathl{V_i=V_{P_i}}{} di dalam
\zusatzklammer {\anfuehrung{$\R^d$ baru}{}} {} {}
dengan titik-titik pusat

\mathdisp {(1,0 , \ldots , 0), (2,0 , \ldots , 0) , \ldots , (n,0 , \ldots , 0)} { . }

Berdasarkan syarat jari-jari, bola-bola ini saling lepas. Kita meninjau himpunan terbuka terbatas
\mathl{U= \bigcup_{i=1}^n V_i}{.} Peta-peta $\alpha_i$ menghasilkan pemetaan-pemetaan kontinu

\mathdisp {\varphi_i :V_i \stackrel{ \left(  \alpha_i  \right)^{-1} }{\longrightarrow} U_i \longrightarrow M} { . }

Karena himpunan-himpunan tersebut saling lepas, dengan menetapkan
\mathl{\varphi(z)= \varphi_i(z)}{,} jika
\mathl{z \in V_i}{} berlaku, kita memperoleh sebuah pemetaan kontinu
\maabbdisp {\varphi} { U } {M
} {.}
Karena sifat tutupan tersebut, pemetaan ini surjektif.
\subsection*{Solusi Soal 13.21}
Untuk pemetaan antipodal, kita memperoleh matriks
:$
\ dA = \begin{pmatrix}
-1 & 0 & 0 \\
0 & -1 & 0 \\
0 & 0 & -1 \end{pmatrix}
$ $ \in $ $\R^{3x3}$.

Jadi, determinan $dA$ tepat sama dengan -1, yakni
 $ \det(dA) = -1 <0$.
Dengan orientasi standar $S^2$ sebagai batas bola satuan, pemetaan antipodal tersebut membalik orientasi dan dengan demikian tidak mempertahankan orientasi.

Secara umum, derajat pemetaan antipodal pada $S^n$ adalah $(-1)^{n+1}$; karena itu, pemetaan tersebut membalik orientasi tepat ketika dimensinya genap.
\subsection*{Solusi Soal 13.22}
Kita tunjukkan bahwa untuk setiap titik

\mavergleichskette
{\vergleichskette
{ x
}
{ \in }{ U
}
{ = }{ X \setminus Y
}
{    }{
}
{   }{
}
}
{}{}{}
terdapat sebuah lingkungan terbuka

\mavergleichskette
{\vergleichskette
{ x
}
{ \in }{ V
}
{ \subseteq }{ U
}
{    }{
}
{   }{
}
}
{}{}{}.
Gabungan semua lingkungan tersebut kemudian sama dengan $U$. Jadi, misalkan

\mavergleichskette
{\vergleichskette
{ x
}
{ \notin }{ Y
}
{  }{
}
{    }{
}
{   }{
}
}
{}{}{}
diberikan. Untuk setiap titik

\mavergleichskette
{\vergleichskette
{ y
}
{ \in }{ Y
}
{  }{
}
{    }{
}
{   }{
}
}
{}{}{}
terdapat himpunan-himpunan terbuka

\mavergleichskette
{\vergleichskette
{ x
}
{ \in }{ V_y
}
{  }{
}
{    }{
}
{   }{
}
}
{}{}{}
dan

\mavergleichskette
{\vergleichskette
{ y
}
{ \in }{ W_y
}
{  }{
}
{    }{
}
{   }{
}
}
{}{}{,}
yang saling lepas. Himpunan-himpunan terbuka

\mathbed {W_y} {}
 {y \in Y} {}
 {} {} {} {,}
menutupi $Y$. Berdasarkan kekompakan, terdapat sebuah subtutupan hingga, katakanlah

\mavergleichskettedisp
{\vergleichskette
{ Y
}
{ \subseteq} { \bigcup_{i = 1}^n  W_i
}
{ } {
}
{ } {
}
{ } {
}
}
{}{}{.}
Maka

\mavergleichskette
{\vergleichskette
{ V
}
{ = }{ \bigcap_{i = 1}^n  V_i
}
{  }{
}
{    }{
}
{   }{
}
}
{}{}{}
terbuka sebagai irisan berhingga dari himpunan-himpunan terbuka, dan karena itu merupakan sebuah lingkungan terbuka dari $x$ yang saling lepas dengan $Y$.

\part{Unit 14}
\chapter[Kuliah 14]{Kuliah 14: Bentuk Diferensial pada Manifold}
\setcounter{section}{14}

  \zwischenueberschrift{Bentuk diferensial pada manifold}

\inputdefinition
{}
{

Misalkan $M$ suatu
\definitionsverweis {manifold diferensiabel}{}{.}
Suatu
\definitionswortpraemath {k}{ bentuk diferensial }{}
\zusatzklammer {atau $k$-\definitionswort {bentuk}{} atau \definitionswort {bentuk berderajat}{} $k$} {} {}
adalah suatu
\definitionsverweis {penampang}{}{}
dalam
\definitionsverweis {produk eksterior}{}{}
ke-$k$ dari
\definitionsverweis {bundel kotangen}{}{,}
yaitu suatu pemetaan
\maabbeledisp {\omega} { M } { \bigwedge^k T^*M
} { P } { \omega(P)
} {,}
dengan

\mavergleichskette
{\vergleichskette
{ \omega(P)
}
{ \in }{ \bigwedge^k T_P^*M
}
{  }{
}
{    }{
}
{   }{
}
}
{}{}{.}

}

Kita melambangkan himpunan semua $k$-bentuk pada $M$ dengan

\mathdisp {{ \mathcal E }^{ k } ( M  )} { . }

\inputbemerkung
{}
{

Jadi, suatu $k$-\definitionsverweis {bentuk}{}{}
memasangkan kepada setiap titik $P$ pada manifold sebuah elemen dari
\mathl{\bigwedge^k T^*_PM}{}. Berdasarkan
Korolari 57.9 (Aljabar Linear (Osnabrück 2024-2025))
dan Teorema 58.4 (Aljabar Linear (Osnabrück 2024-2025)),
hal ini sama dengan suatu
\definitionsverweis {pemetaan multilinear}{}{}
\definitionsverweis {selang-seling}{}{}
\maabbdisp {} {T_PM \times \cdots \times T_PM } { \R
} {.}
Pemetaan yang bersesuaian ini juga kita lambangkan dengan
\mathl{\omega(P)}{;} untuk $k$
\definitionsverweis {vektor singgung}{}{}

\mavergleichskette
{\vergleichskette
{ v_1 , \ldots , v_k
}
{ \in }{ T_PM
}
{  }{
}
{    }{
}
{   }{
}
}
{}{}{}
dengan demikian,

\mathdisp {\omega(P)(v_1 , \ldots , v_k)} {  }

adalah suatu bilangan riil. Di sini muncul dua jenis argumen yang sama sekali berbeda: di satu pihak titik pada manifold, dan di pihak lain elemen-elemen ruang singgung di titik tersebut. Ketergantungan pada vektor-vektor singgung relatif sederhana karena pemetaannya multilinear selang-seling, sedangkan ketergantungan pada titik manifold dapat serumit apa pun. Karena produk-produk eksterior bundel kotangen sendiri merupakan manifold menurut
Soal 14.2,
kita dapat langsung membicarakan bentuk diferensial yang kontinu atau
\zusatzklammer {jika $M$ merupakan manifold $C^2$} {} {}
diferensiabel.

Untuk

\mavergleichskette
{\vergleichskette
{ k
}
{ = }{ 0
}
{  }{
}
{    }{
}
{   }{
}
}
{}{}{}
ruang kotangen hanya muncul secara formal; suatu $0$-bentuk tidak lain adalah sebuah fungsi
\maabb {f} { M } {\R
} {.}
Suatu $1$-bentuk
\zusatzklammer {juga disebut
\definitionswortenp{bentuk Pfaff}{}} {} {}
memasangkan suatu bilangan riil kepada setiap titik $P$ dan setiap vektor singgung pada titik tersebut. Untuk

\mavergleichskette
{\vergleichskette
{k
}
{ > }{ n
}
{ = }{ \operatorname{dim} \, M
}
{    }{
}
{   }{
}
}
{}{}{}
produk eksterior ke-$k$ adalah ruang nol sehingga sama sekali tidak ada bentuk taktrivial berderajat ini. Kasus yang sangat penting adalah

\mavergleichskette
{\vergleichskette
{ k
}
{ = }{ n
}
{ = }{ \operatorname{dim} \, M
}
{    }{
}
{   }{
}
}
{}{}{.}
Dalam hal ini, produk eksterior ke-$n$ memiliki rank $1$
\zusatzklammer {artinya, dimensinya $1$ di setiap titik} {} {}
dan suatu penampang di dalamnya secara lokal dideskripsikan oleh satu fungsi saja. Gambaran yang berguna ialah bahwa, untuk $n$ vektor singgung, bilangan
\mathl{\omega(P)(v_1 , \ldots , v_n)}{} menyatakan volume
\zusatzklammer {\anfuehrung{berorientasi}{}} {} {}
dari paralelotop yang direntang oleh vektor-vektor tersebut di ruang singgung. Gambaran ini juga membantu untuk $k$ yang lebih kecil: dengan
\mathl{\omega(P) (v_1 , \ldots , v_k)}{} kita dapat menghitung volume berdimensi $k$ dari paralelotop yang dibangkitkan oleh $k$ vektor singgung. Gambaran ini akan dibuat tepat ketika kita mengintegralkan pada suatu
\definitionsverweis {submanifold tertutup}{}{}
berdimensi $k$.

}



\inputfaktbeweis
{Mannigfaltigkeit/Differentialform/Elementare Eigenschaften/Fakt}
{Lemma}
{}
{

\faktsituation {Misalkan $M$ suatu
\definitionsverweis {manifold diferensiabel}{}{}
dan, untuk

\mavergleichskette
{\vergleichskette
{ k
}
{ \in }{ \N
}
{  }{
}
{    }{
}
{   }{
}
}
{}{}{},
misalkan
\mathl{{ \mathcal E }^{ k } (  M   )}{} himpunan semua
$k$-\definitionsverweis {bentuk}{}{}
pada $M$.}
\faktuebergang {Maka berlaku sifat-sifat berikut.}
\faktfolgerung {\aufzaehlungfuenf{Dengan operasi-operasi alaminya,
\mathl{{ \mathcal E }^{ k } (  M   )}{} membentuk
\definitionsverweis {ruang vektor riil}{}{.}
}{Untuk suatu bentuk diferensial

\mavergleichskette
{\vergleichskette
{ \omega
}
{ \in }{
{ \mathcal E }^{ k } (  M   )
}
{  }{
}
{    }{
}
{   }{
}
}
{}{}{}
dan suatu fungsi
\maabbdisp {f} { M } {\R
} {},

\mavergleichskette
{\vergleichskette
{ f \omega
}
{ \in }{
{ \mathcal E }^{ k } (  M   )
}
{  }{
}
{    }{
}
{   }{
}
}
{}{}{,}
dengan $f \omega$ didefinisikan oleh

\mavergleichskettedisp
{\vergleichskette
{ (f \omega)(P)
}
{ \defeq} { f(P) \omega (P)
}
{ } {
}
{ } {
}
{ } {
}
}
{}{}{}.
}{Setiap
\definitionsverweis {fungsi diferensiabel}{}{}
$C^1$
\maabbdisp {f} { M } {\R
} {}
melalui
\definitionsverweis {pemetaan singgung}{}{}

\mathl{Tf}{} mendefinisikan suatu
$1$-\definitionsverweis {bentuk diferensial}{}{}
\maabbeledisp {df} { M } { T^*M
} { P } { T_Pf
} {,}
dengan ruang singgung dari $\R$ di
\mathl{f(P)}{} diidentifikasi dengan $\R$. Hal ini menghasilkan suatu pemetaan
\maabbeledisp {d} {C^1(M,\R)} {
{ \mathcal E }^{  1 } (  M   )
} { f } {df
} {.}
}{Jika

\mavergleichskette
{\vergleichskette
{ M
}
{ \subseteq }{ \R^m
}
{  }{
}
{    }{
}
{   }{
}
}
{}{}{}
merupakan himpunan bagian terbuka, maka di bawah identifikasi

\mavergleichskettedisp
{\vergleichskette
{ TM
}
{ \cong} { M \times \R^m
}
{ } {
}
{ } {
}
{ } {
}
}
{}{}{}
pemetaan pada (3) sama dengan
\maabbeledisp {} { M } { M \times ( \R^m)^*
} { P } { (P, \left(  D f   \right)_{ P } \left(   -  \right) )
} {.}
}{Pemetaan $d$ pada (3) adalah
$\R$-\definitionsverweis {linear}{}{.}
}}
\faktzusatz {}
\faktzusatz {}

}
{

\teilbeweis {}{}{}
{(1) dan (2) langsung mengikuti dari peninjauan titik demi titik.}
{} \teilbeweis {}{}{}
{(3). Untuk setiap titik

\mavergleichskette
{\vergleichskette
{ P
}
{ \in }{ M
}
{  }{
}
{    }{
}
{   }{
}
}
{}{}{},

\maabbdisp {T_Pf} { T_PM } { T_{f(P)} \R \cong \R
} {}
merupakan suatu
\definitionsverweis {pemetaan linear}{}{}
menurut Lemma 9.10  (3),
dan dengan demikian merupakan elemen dari
\mathl{T_P^* M}{,} yang kita lambangkan dengan
\mathl{(df)_P}{}. Oleh karena itu, pemasangan
\mathl{P \mapsto (df)_P}{} merupakan suatu
\definitionsverweis {bentuk diferensial}{}{.}}
{}
\teilbeweis {}{}{}
{(4) mengikuti dari
Lemma 9.10  (1).}
{}
\teilbeweis {}{}{}
{(5). Pemetaan pada (3) ditentukan untuk setiap titik

\mavergleichskette
{\vergleichskette
{ P
}
{ \in }{ M
}
{  }{
}
{    }{
}
{   }{
}
}
{}{}{}
pada setiap lingkungan terbuka. Karena itu, kita boleh menganggap bahwa

\mavergleichskette
{\vergleichskette
{ M
}
{ \subseteq }{ \R^m
}
{  }{
}
{    }{
}
{   }{
}
}
{}{}{}
merupakan himpunan terbuka, sehingga pernyataannya mengikuti dari (4) dan
Proposisi 45.6 (Analisis (Osnabrück 2021-2023)).}
{}

}

Untuk suatu himpunan terbuka

\mavergleichskette
{\vergleichskette
{ V
}
{ \subseteq }{ \R^n
}
{  }{
}
{    }{
}
{   }{
}
}
{}{}{},
tersedia fungsi-fungsi koordinat
\maabbdisp {x_j} {V} {\R
} {},
yang, jika diberikan suatu peta, dapat dipindahkan ke suatu himpunan bagian terbuka dari manifold. Di setiap titik

\mavergleichskette
{\vergleichskette
{ Q
}
{ \in }{ V
}
{  }{
}
{    }{
}
{   }{
}
}
{}{}{},
bentuk-bentuk

\mathbed {dx_j} {}
 {j=1 , \ldots , n} {}
 {} {} {} {,}
membentuk basis ruang kotangen di $Q$. Basis ini tidak lain adalah
\definitionsverweis {basis dual}{}{}
dari basis standar di ruang ambien $\R^n$, yang digunakan sebagai ruang singgung di seluruh $V$. Untuk suatu himpunan bagian beranggotakan $k$

\mavergleichskettedisp
{\vergleichskette
{ J
}
{ =} { \{j_1 , \ldots , j_k\}
}
{ \subseteq} { \{1 , \ldots , n\}
}
{ } {
}
{ } {
}
}
{}{}{},
kita tetapkan

\mavergleichskettedisp
{\vergleichskette
{  dx_J
}
{ =} {  dx_{j_1} \wedge \ldots \wedge dx_{j_k}
}
{ } {
}
{ } {
}
{ } {
}
}
{}{}{,}
yang merupakan $k$-bentuk yang sangat sederhana pada $V$. Untuk setiap titik

\mavergleichskette
{\vergleichskette
{ Q
}
{ \in }{ V
}
{  }{
}
{    }{
}
{   }{
}
}
{}{}{},
berlaku

\mavergleichskettedisp
{\vergleichskette
{ dx_J(Q)
}
{ =} { (dx_{j_1} \wedge \ldots \wedge dx_{j_k} )(Q)
}
{ =} { dx_{j_1}(Q) \wedge \ldots \wedge dx_{j_k} (Q)
}
{ =} { e_{j_1}^* \wedge \ldots \wedge e_{j_k}^*
}
{ } {
}
}
{}{}{.}
Aksi bentuk ini pada

\mavergleichskette
{\vergleichskette
{  v_1 \wedge \ldots \wedge v_k
}
{ \in }{ \bigwedge^k T_QV
}
{  }{
}
{    }{
}
{   }{
}
}
{}{}{}
, menurut
Teorema 58.4 (Aljabar Linear (Osnabrück 2024-2025))
, diberikan oleh

\mavergleichskettedisp
{\vergleichskette
{ \left(  e_{j_1}^* \wedge \ldots \wedge e_{j_k}^*  \right)  (v_1 \wedge \ldots \wedge v_k)
}
{ =} { \det  \left(  e_{j_i}^*(v_\ell)  \right)_{i \ell}
}
{ =} { \det  \left(  (v_\ell)_{j_i}  \right)_{i \ell}
}
{ } {
}
{ } {
}
}
{}{}{.}
Menurut
Teorema 58.1 (Aljabar Linear (Osnabrück 2024-2025)),
nilai-nilai bentuk diferensial
\zusatzklammer {dengan \mathlk{j_1 <   \ldots < j_k}{}} {} {}

\mavergleichskettedisp
{\vergleichskette
{ dx_J
}
{ =} { dx_{j_1} \wedge \ldots \wedge dx_{j_k}
}
{ } {
}
{ } {
}
{ } {
}
}
{}{}{}
membentuk suatu
\definitionsverweis {basis}{}{}
dari
\mathl{\bigwedge^k T_Q^*V}{}
untuk setiap titik $Q$. Oleh karena itu, setiap
\definitionsverweis {bentuk diferensial}{}{} pada $V$ yang berderajat $k$

\mavergleichskette
{\vergleichskette
{ \omega
}
{ \in }{
{ \mathcal E }^{ k } ( V  )
}
{  }{
}
{    }{
}
{   }{
}
}
{}{}{}
dapat ditulis secara tunggal sebagai

\mavergleichskettedisp
{\vergleichskette
{ \omega
}
{ =} { \sum_{J,\, { \# \left(  J \right) } = k} f_J dx_J
}
{ } {
}
{ } {
}
{ } {
}
}
{}{}{}
dengan fungsi-fungsi yang ditentukan secara tunggal
\maabbdisp {f_J} {V} {\R
} {.}



\inputfaktbeweis
{Mannigfaltigkeit/Differentialform/Lokale Beschreibung/Fakt}
{Lemma}
{}
{

\faktsituation {Misalkan $M$ suatu
\definitionsverweis {manifold diferensiabel}{}{}
dan

\mavergleichskette
{\vergleichskette
{ U
}
{ \subseteq }{ M
}
{  }{
}
{    }{
}
{   }{
}
}
{}{}{}
suatu
\definitionsverweis {himpunan bagian terbuka}{}{}
dengan peta
\maabbdisp {\alpha} { U } { V
} {},
dan misalkan

\mavergleichskette
{\vergleichskette
{ V
}
{ \subseteq }{ \R^n
}
{  }{
}
{    }{
}
{   }{
}
}
{}{}{}
terbuka. Misalkan
\maabbdisp {x_j} { U } {\R
} {}
fungsi-fungsi koordinat yang bersesuaian,

\mavergleichskette
{\vergleichskette
{ 1
}
{ \leq }{ j
}
{ \leq }{ n
}
{    }{
}
{   }{
}
}
{}{}{.}}
\faktfolgerung {Maka setiap
\definitionsverweis {bentuk diferensial}{}{} pada $U$ yang berderajat $k$

\mavergleichskette
{\vergleichskette
{ \omega
}
{ \in }{
{ \mathcal E }^{ k } (  U   )
}
{  }{
}
{    }{
}
{   }{
}
}
{}{}{}
dapat ditulis secara tunggal sebagai

\mavergleichskettedisp
{\vergleichskette
{ \omega
}
{ =} { \sum_{J,\, { \# \left(   J  \right) } = k}  f_J dx_J
}
{ } {
}
{ } {
}
{ } {
}
}
{}{}{}
dengan fungsi-fungsi yang ditentukan secara tunggal
\maabbdisp {f_J} { U } {\R
} {.}}
\faktzusatz {}
\faktzusatz {}

}
{

Hal ini mengikuti dari
Teorema 58.1 (Aljabar Linear (Osnabrück 2024-2025)).

}



\inputfaktbeweis
{Mannigfaltigkeit/Pfaffsche Differentialform zu Funktion/Beschreibung mit partiellen Ableitungen/Fakt}
{Korollar}
{}
{

\faktsituation {Misalkan $M$ suatu
\definitionsverweis {manifold diferensiabel}{}{}
dan

\mavergleichskette
{\vergleichskette
{U
}
{ \subseteq }{ M
}
{  }{
}
{    }{
}
{   }{
}
}
{}{}{}
suatu
\definitionsverweis {himpunan bagian terbuka}{}{}
dengan peta
\maabbdisp {\alpha} { U } { V
} {},
dan misalkan

\mavergleichskette
{\vergleichskette
{V
}
{ \subseteq }{\R^n
}
{  }{
}
{    }{
}
{   }{
}
}
{}{}{}
terbuka. Misalkan
\maabbdisp {x_j} { U } {\R
} {}
fungsi-fungsi koordinat yang bersesuaian,

\mavergleichskette
{\vergleichskette
{ 1
}
{ \leq }{j
}
{ \leq }{n
}
{    }{
}
{   }{
}
}
{}{}{.}
Misalkan
\maabbdisp {f} { U } {\R
} {}
suatu
\definitionsverweis {fungsi diferensiabel}{}{.}}
\faktfolgerung {Maka $1$-\definitionsverweis {bentuk diferensial}{}{}
$df$ yang bersesuaian mempunyai representasi\zusatzfussnote {Turunan-turunan
\mathl{{ \frac{   \partial  f  }{  \partial  x_j  } }}{} diperkenalkan dalam kuliah kedelapan} {.} {}

\mavergleichskettedisp
{\vergleichskette
{ df
}
{ =} { \sum_{j = 1}^n { \frac{   \partial   f   }{  \partial   x_j   } } dx_j
}
{ } {
}
{ } {
}
{ } {
}
}
{}{}{.}}
\faktzusatz {}
\faktzusatz {}

}
{

Kita dapat langsung menganggap bahwa semuanya berlangsung pada himpunan terbuka

\mavergleichskette
{\vergleichskette
{V
}
{ \subseteq }{\R^n
}
{  }{
}
{    }{
}
{   }{
}
}
{}{}{}.
Untuk setiap titik

\mavergleichskette
{\vergleichskette
{Q
}
{ \in }{ V
}
{  }{
}
{    }{
}
{   }{
}
}
{}{}{},
berlaku kesamaan berikut antara
\definitionsverweis {bentuk-bentuk linear}{}{}
pada $\R^n$,

\mavergleichskettealign
{\vergleichskettealign
{(df)_Q
}
{ =} { \left(D f \right)_{ Q }
}
{ =} { \left(  { \frac{   \partial   f   }{  \partial   x_1   } } ( Q)   , \,   \ldots  , \,   { \frac{   \partial   f   }{  \partial   x_n   } } ( Q)                 \right)
}
{ =} { { \frac{   \partial   f   }{  \partial   x_1   } } ( Q) dx_1 + \cdots + { \frac{   \partial   f   }{  \partial   x_n   } } ( Q) dx_n
}
{ } {
}
}
{}
{}{.}

}

  \zwischenueberschrift{Penarikan balik bentuk diferensial}

\inputdefinition
{}
{

Misalkan
\mathkor {} {L} {dan} {M} {}
\definitionsverweis {dua manifold diferensiabel}{}{}
dan misalkan
\maabbdisp {\varphi} { L } { M
} {}
suatu
\definitionsverweis {pemetaan diferensiabel}{}{.}
Misalkan $\omega$ suatu $k$-\definitionsverweis {bentuk diferensial}{}{}
pada $M$. Maka pemetaan yang menentukan suatu $k$-bentuk pada $L$ diberikan oleh

\mathdisp {(P ,v_1 , \ldots , v_k) \longmapsto   \omega(\varphi(P), T_P(\varphi)(v_1) , \ldots ,  T_P(\varphi)(v_k) )} {  }

Pemetaan ini merupakan
\definitionsverweis {pemetaan selang-seling}{}{}.
Bentuk yang bersesuaian dengannya disebut, terhadap $\varphi$, \definitionswort {bentuk tarik balik}{} berderajat $k$. Bentuk ini dilambangkan dengan

\mathdisp {\varphi^* \omega} {  }.

}



\inputfaktbeweis
{Mannigfaltigkeit/Differenzierbare Abbildung/Zurückziehen von Differentialformen/Elementare Eigenschaften/Fakt}
{Lemma}
{}
{

\faktsituation {Misalkan
\mathkor {} {L} {dan} {M} {}
\definitionsverweis {dua manifold diferensiabel}{}{}
dan
\maabbdisp {\varphi} { L } { M
} {}
suatu
\definitionsverweis {pemetaan diferensiabel}{}{.}}
\faktuebergang {Maka
\definitionsverweis {penarikan balik bentuk diferensial}{}{}
memenuhi sifat-sifat berikut.}
\faktfolgerung {\aufzaehlungvier{Untuk suatu fungsi

\mavergleichskette
{\vergleichskette
{ f
}
{ \in }{ \operatorname{Map} (M,\R)
}
{ = }{
{ \mathcal E }^{  0 } (  M   )
}
{    }{
}
{   }{
}
}
{}{}{},
berlaku

\mavergleichskette
{\vergleichskette
{ \varphi^*(f)
}
{ = }{ f \circ \varphi
}
{  }{
}
{    }{
}
{   }{
}
}
{}{}{.}
}{Pemetaan-pemetaan
\maabbeledisp {\varphi^*} {
{ \mathcal E }^{ k } (  M   ) } {
{ \mathcal E }^{ k } (  L   )
} { \omega} { \varphi^* \omega
} {,}
adalah
$\R$-\definitionsverweis {linear}{}{.}
}{Jika

\mavergleichskette
{\vergleichskette
{ L
}
{ \subseteq }{ M
}
{  }{
}
{    }{
}
{   }{
}
}
{}{}{}
merupakan suatu
\definitionsverweis {submanifold terbuka}{}{,}
dan pemetaan di atas adalah inklusi, maka penarikan balik suatu bentuk diferensial $\omega$ tidak lain adalah pembatasan
\mathl{\omega \vert_L}{} pada himpunan bagian ini.
}{Misalkan $N$ suatu manifold diferensiabel lain dan
\maabbdisp {\psi} { M } { N
} {}
suatu pemetaan diferensiabel lain. Maka berlaku

\mavergleichskettedisp
{\vergleichskette
{ ( \psi \circ \varphi )^*(\omega)
}
{ =} { \varphi^*(\psi^*(\omega))
}
{ } {
}
{ } {
}
{ } {
}
}
{}{}{}
untuk setiap bentuk diferensial

\mavergleichskette
{\vergleichskette
{ \omega
}
{ \in }{
{ \mathcal E }^{ k } (  N   )
}
{  }{
}
{    }{
}
{   }{
}
}
{}{}{.}
}}
\faktzusatz {}
\faktzusatz {}

}
{

\teilbeweis {}{}{}
{(1) langsung mengikuti dari
Definisi 14.6.}
{}
\teilbeweis {}{}{}
{(2). Untuk bentuk-bentuk diferensial
\mathkor {} {\omega_1} {dan} {\omega_2} {}
dan skalar-skalar

\mavergleichskette
{\vergleichskette
{ a,b
}
{ \in }{\R
}
{  }{
}
{    }{
}
{   }{
}
}
{}{}{},
kita harus menunjukkan bahwa

\mavergleichskette
{\vergleichskette
{ \varphi^*(a \omega_1+ b \omega_2)
}
{ = }{ a \varphi^* \omega_1 + b \varphi^* \omega_2
}
{  }{
}
{    }{
}
{   }{
}
}
{}{}{}.
Kesamaan bentuk diferensial semacam ini berarti bahwa kesamaan tersebut berlaku di setiap titik

\mavergleichskette
{\vergleichskette
{ P
}
{ \in }{ L
}
{  }{
}
{    }{
}
{   }{
}
}
{}{}{}
dan untuk setiap $k$-tupel vektor singgung

\mavergleichskette
{\vergleichskette
{ v_1 , \ldots , v_k
}
{ \in }{ T_PL
}
{  }{
}
{    }{
}
{   }{
}
}
{}{}{}.
Karena itu, pernyataannya mengikuti dari

\mavergleichskettealignhandlinks
{\vergleichskettealignhandlinks
{ \left(  \varphi^*(a \omega_1 + b \omega_2)  \right) \left(  P, v_1 , \ldots , v_k  \right)
}
{ =} { (a \omega_1 + b \omega_2) \left(  \varphi(P), T_P(\varphi)(v_1) , \ldots ,  T_P(\varphi)(v_k)  \right)
}
{ =} { a \omega_1 \left(  \varphi(P), T_P(\varphi)(v_1) , \ldots ,  T_P(\varphi)(v_k)  \right)
}
{ \,\,\,\, \,} {+ b \omega_2 \left(  \varphi(P), T_P(\varphi)(v_1) , \ldots ,  T_P(\varphi)(v_k)  \right)
}
{ =} { \left(  a  \varphi^* \omega_1  \right)(P, v_1 , \ldots , v_k)  + \left(  b \varphi^* \omega_2  \right) (P, v_1 , \ldots , v_k)
}
}
{}
{}{.}}
{}
\teilbeweis {}{}{}
{(3) langsung mengikuti dari definisi.}
{}
\teilbeweis {}{}{}
{(4). Misalkan

\mavergleichskette
{\vergleichskette
{ Q
}
{ \in }{ L
}
{  }{
}
{    }{
}
{   }{
}
}
{}{}{,}

\mavergleichskette
{\vergleichskette
{ u_1 , \ldots , u_k
}
{ \in }{ T_QL
}
{  }{
}
{    }{
}
{   }{
}
}
{}{}{},
dan $\omega$ suatu $k$-bentuk pada $N$. Maka, dengan menggunakan
Lemma 9.10  (4),
berlaku

\mavergleichskettealigndrucklinks
{\vergleichskettealigndrucklinks
{(( \psi \circ \varphi)^* (\omega))(Q, u_1 , \ldots , u_k)
}
{ =} { \omega \left(  ( \psi \circ \varphi) (Q), T_Q( \psi \circ \varphi)(u_1) , \ldots , T_Q(  \psi \circ \varphi) (u_k)  \right)
}
{ =} { \omega \left(  \psi ( \varphi (Q)), \left(  T_{\varphi(Q)}\psi  \right)  ((T_Q\varphi)(u_1)) , \ldots , \left(  T_{\varphi(Q)} \psi  \right) ((T_Q\varphi)(u_k))  \right)
}
{ =} { \left(  \psi^*(\omega)  \right) \left(  \varphi(Q),T_Q(\varphi)(u_1) , \ldots , T_Q(\varphi)(u_k)  \right)
}
{ =} { \left(  \varphi^* \left(  \psi^*(\omega)  \right)   \right) (Q,u_1 , \ldots , u_k)
}
}
{}{}{,}
dan inilah pernyataannya.}
{}

}



\inputfaktbeweis
{Offene Mengen/Differenzierbare Abbildung/Zurückziehen von Differentialformen/In Koordinaten/Fakt}
{Lemma}
{}
{

\faktsituation {Misalkan
\mathkor {} {U \subseteq \R^n} {dan} {V \subseteq \R^m} {}
\definitionsverweis {himpunan-himpunan bagian terbuka}{}{,}
dengan koordinat masing-masing dilambangkan
\mathl{x_1 , \ldots , x_n}{} dan
\mathl{y_1 , \ldots , y_m}{}. Misalkan
\maabbdisp {\varphi} { U } { V
} {}
suatu
\definitionsverweis {pemetaan diferensiabel}{}{}
dan misalkan $\omega$ suatu
$k$-\definitionsverweis {bentuk diferensial}{}{}
pada $V$ dengan representasi

\mavergleichskettedisp
{\vergleichskette
{ \omega
}
{ =} { \sum_{I,\, { \# \left(   I  \right) } = k}  f_I  dy_I
}
{ } {
}
{ } {
}
{ } {
}
}
{}{}{,}
dengan
\maabb {f_I} { V } {\R
} {}
merupakan
\definitionsverweis {fungsi-fungsi}{}{.}}
\faktfolgerung {Maka
\definitionsverweis {bentuk tarik balik}{}{}
mempunyai representasi

\mavergleichskettealignhandlinks
{\vergleichskettealignhandlinks
{ \varphi^* \omega
}
{ =} { \sum_{I,\, { \# \left(   I  \right) } = k}  \left(  f_I \circ \varphi  \right)  \left( \sum_{J,\, { \# \left(   J  \right) } = k} \det  \left(  \left(  { \frac{   \partial  \varphi_i   }{  \partial   x_j   } }  \right)_{ i \in I, j \in J}  \right)  dx_J \right)
}
{ =} { \sum_{J,\, { \# \left(   J  \right) } = k}   \left( \sum_{I,\, { \# \left(   I  \right) } = k} \left(  f_I \circ \varphi  \right)  \det  \left(  \left(  { \frac{   \partial  \varphi_i   }{  \partial   x_j   } }  \right)_{ i \in I, j \in J}  \right)  \right)  dx_J
}
{ } {
}
{ } {
}
}
{}{}{.}}
\faktzusatz {}
\faktzusatz {}

}
{

Kesamaan kedua didasarkan pada pengaturan ulang yang sederhana. Berdasarkan
Lemma 14.7 (2),
kita dapat membatasi diri pada kasus

\mavergleichskette
{\vergleichskette
{ \omega
}
{ = }{ f_Idy_I
}
{  }{
}
{    }{
}
{   }{
}
}
{}{}{}.
Kita tetapkan

\mavergleichskette
{\vergleichskette
{ f
}
{ = }{ f_I
}
{  }{
}
{    }{
}
{   }{
}
}
{}{}{}
dan boleh menganggap

\mavergleichskette
{\vergleichskette
{ I
}
{ = }{ \{ 1 , \ldots , k \}
}
{  }{
}
{    }{
}
{   }{
}
}
{}{}{}.
Kita menunjukkan kesamaan antara kedua $k$-bentuk pada $U$ dengan menunjukkan bahwa keduanya menghasilkan nilai yang sama untuk setiap titik

\mavergleichskette
{\vergleichskette
{ P
}
{ \in }{ U
}
{  }{
}
{    }{
}
{   }{
}
}
{}{}{}
dan setiap produk eksterior
\mathl{e_{j_1 } \wedge \ldots \wedge e_{j_k }}{} dengan

\mavergleichskette
{\vergleichskette
{ 1
}
{ \leq }{ j_1
}
{ < }{ \ldots
}
{ <   }{ j_k
}
{ \leq  }{ n
}
}
{}{}{}.
Di satu pihak,

\mavergleichskettealigndrucklinks
{\vergleichskettealigndrucklinks
{ \varphi^*( f dy_1 \wedge \ldots \wedge dy_k)(P,e_{j_1 } \wedge \ldots \wedge e_{j_k }  )
}
{ =} { ( f dy_1 \wedge \ldots \wedge dy_k)(\varphi(P),T_P(\varphi)(e_{j_1 }) \wedge \ldots \wedge T_P(\varphi)(e_{j_k }))
}
{ =} { f(\varphi(P)) (dy_1 \wedge \ldots \wedge dy_k)( \begin{pmatrix}  { \frac{   \partial  \varphi_1   }{  \partial   x_{j_1 }  } } (P )  \\\vdots\\  { \frac{   \partial  \varphi_m   }{  \partial   x_{j_1 }  } } (P )   \end{pmatrix} \wedge \ldots \wedge \begin{pmatrix}  { \frac{   \partial  \varphi_1   }{  \partial   x_{j_k }  } } (P )  \\\vdots\\  { \frac{   \partial  \varphi_m   }{  \partial   x_{j_k }  } } (P )   \end{pmatrix} )
}
{ =} { f(\varphi(P)) \cdot \det  \left(  (dy_i) \begin{pmatrix}  { \frac{   \partial  \varphi_1   }{  \partial   x_{j_\ell}  } } (P )  \\\vdots\\  { \frac{   \partial  \varphi_m   }{  \partial   x_{j_\ell}  } } (P )   \end{pmatrix}  \right)_{1 \leq i, \ell \leq k}
}
{ =} { f(\varphi(P)) \cdot \det  \left(  { \frac{   \partial  \varphi_i   }{  \partial   x_{j_\ell}  } } (P )  \right)_{1 \leq i, \ell \leq k}
}
}
{}{}{.}
Di pihak lain, jika jumlah tersebut diterapkan pada
\mathl{e_{j_1 } \wedge \ldots \wedge e_{j_k }}{},
maka

\mavergleichskette
{\vergleichskette
{ dx_J (e_{j_1 } \wedge \ldots \wedge e_{j_k })
}
{ = }{ 0
}
{  }{
}
{    }{
}
{   }{
}
}
{}{}{}
kecuali apabila

\mavergleichskette
{\vergleichskette
{ J
}
{ = }{ \{j_1 , \ldots , j_k\}
}
{  }{
}
{    }{
}
{   }{
}
}
{}{}{,}
dan dalam kasus itu nilainya $1$, sehingga diperoleh nilai yang sama.

}



\inputfaktbeweis
{Offene Mengen/Differenzierbare Abbildung/Zurückziehen von Volumenform auf gleichdimensionalem Raum/In Koordinaten/Fakt}
{Korollar}
{}
{

\faktsituation {Misalkan

\mavergleichskette
{\vergleichskette
{ U,V
}
{ \subseteq }{ \R^n
}
{  }{
}
{    }{
}
{   }{
}
}
{}{}{}
\definitionsverweis {himpunan-himpunan bagian terbuka}{}{,}
dengan koordinat masing-masing dilambangkan
\mathl{x_1 , \ldots , x_n}{} dan
\mathl{y_1 , \ldots , y_n}{}. Misalkan
\maabbdisp {\varphi} { U } { V
} {}
suatu
\definitionsverweis {pemetaan diferensiabel}{}{}
dan misalkan $\omega$ suatu
$n$-\definitionsverweis {bentuk diferensial}{}{}
pada $V$ dengan representasi

\mavergleichskettedisp
{\vergleichskette
{ \omega
}
{ =} { f dy_1 \wedge \ldots \wedge dy_n
}
{ } {
}
{ } {
}
{ } {
}
}
{}{}{.}}
\faktfolgerung {Maka
\definitionsverweis {bentuk tarik balik}{}{}
mempunyai representasi

\mavergleichskettedisp
{\vergleichskette
{ \varphi^* \omega
}
{ =} { (f \circ \varphi) \cdot \det  \left(  \left(  { \frac{   \partial  \varphi_i   }{  \partial   x_j   } }  \right)_{1 \leq  i, j \leq n}  \right) dx_1 \wedge \ldots \wedge dx_n
}
{ } {
}
{ } {
}
{ } {}
}
{}{}{.}}
\faktzusatz {}
\faktzusatz {}

}
{

Hal ini langsung mengikuti dari
Lemma 14.8.

}
Dengan demikian, fungsi-fungsi yang mendeskripsikan suatu bentuk diferensial mempunyai perilaku transformasi yang sama dengan kerapatan berorientasi pada suatu peta. Untuk kerapatan ukuran positif, faktor determinannya harus diganti dengan nilai mutlak determinan.



\inputfaktbeweis
{Differentialform/Lokal/Zurückziehen unter partiell konstanter Abbildung/Fakt}
{Korollar}
{}
{

\faktsituation {Misalkan

\mavergleichskette
{\vergleichskette
{ U
}
{ \subseteq }{ \R^n
}
{  }{
}
{    }{
}
{   }{
}
}
{}{}{}
dan

\mavergleichskette
{\vergleichskette
{ V
}
{ \subseteq }{ \R^m
}
{  }{
}
{    }{
}
{   }{
}
}
{}{}{}
\definitionsverweis {himpunan-himpunan bagian terbuka}{}{,}
dengan koordinat masing-masing dilambangkan
\mathl{x_1 , \ldots , x_n}{} dan
\mathl{y_1 , \ldots , y_m}{}. Misalkan
\maabbdisp {\varphi} { U } { V
} {}
suatu
\definitionsverweis {pemetaan diferensiabel}{}{}
dengan
\mathl{\varphi_{i_0 }}{} konstan untuk suatu

\mavergleichskette
{\vergleichskette
{ i_0
}
{ \in }{ \{1 , \ldots , m\}
}
{  }{
}
{    }{
}
{   }{
}
}
{}{}{},
dan misalkan $\omega$ suatu
$k$-\definitionsverweis {bentuk diferensial}{}{}
pada $V$ dengan representasi

\mavergleichskettedisp
{\vergleichskette
{ \omega
}
{ =} { f dy_I
}
{ } {
}
{ } {
}
{ } {
}
}
{}{}{}
dengan

\mavergleichskette
{\vergleichskette
{ i_0
}
{ \in }{ I
}
{  }{
}
{    }{
}
{   }{
}
}
{}{}{.}}
\faktfolgerung {Maka

\mavergleichskette
{\vergleichskette
{ \varphi^*\omega
}
{ = }{ 0
}
{  }{
}
{    }{
}
{   }{
}
}
{}{}{.}}
\faktzusatz {}
\faktzusatz {}

}
{

Menurut
Lemma 14.8,
berlaku

\mavergleichskettedisp
{\vergleichskette
{ \varphi^* \omega
}
{ =} { \sum_{J,\, { \# \left(   J  \right) } = k} (f \circ \varphi) \cdot \det  \left(  \left(  { \frac{   \partial  \varphi_i   }{  \partial   x_j   } }  \right) _{ i \in I, j \in J}  \right) dx_J
}
{ } {
}
{ } {
}
{ } {}
}
{}{}{.}
Karena

\mavergleichskettedisp
{\vergleichskette
{ { \frac{   \partial  \varphi_{i_0 }  }{  \partial   x_j   } }
}
{ =} { 0
}
{ } {
}
{ } {
}
{ } {
}
}
{}{}{}
untuk semua

\mavergleichskette
{\vergleichskette
{ j
}
{ \in }{ J
}
{  }{
}
{    }{
}
{   }{
}
}
{}{}{,}
untuk setiap $J$, matriks tersebut mempunyai satu baris bernilai $0$, sehingga semua determinannya bernilai $0$.

}

\chapter{Lembar Kerja 14}
\setcounter{section}{14}

  \zwischenueberschrift{Soal latihan}

\inputaufgabe
{}
{

Misalkan $M$ sebuah
\definitionsverweis {manifold diferensiabel}{}{}
dengan
\definitionsverweis {bundel kotangen}{}{}

\mathl{T^*M}{.} Tunjukkan bahwa pada
\mathl{\bigwedge^k T^*M}{} untuk setiap $k$ dapat didefinisikan suatu
\definitionsverweis {topologi}{}{}
sedemikian sehingga, untuk setiap peta
\maabb {\alpha} { U } { V
} {}
pemetaan
\maabbdisp {} {\bigwedge^k T^*U} { \bigwedge^k T^*V \cong V \times \bigwedge^k \R^{n*}
} {}
merupakan sebuah
\definitionsverweis {homeomorfisme}{}{.}

}
{} {}
Dengan demikian, kita dapat berbicara tentang bentuk diferensial kontinu maupun terukur.

\inputaufgabe
{}
{

Misalkan $M$ sebuah
$C^2$-\definitionsverweis {manifold diferensiabel}{}{}
dengan
\definitionsverweis {bundel kotangen}{}{}

\mathl{T^*M}{.} Tunjukkan bahwa
\mathl{\bigwedge^k T^*M}{} merupakan sebuah manifold diferensiabel untuk setiap $k$.

}
{} {}
Dengan demikian, kita juga dapat berbicara tentang bentuk diferensial yang terdiferensialkan secara kontinu. Secara lebih umum, jika $M$ sebuah
$C^m$-\definitionsverweis {manifold}{}{}
dan

\mavergleichskette
{\vergleichskette
{  \ell
}
{ < }{ m
}
{  }{
}
{    }{
}
{   }{
}
}
{}{}{}
maka
\mathl{{ \mathcal E }^{ k }_{\ell } ( M  )}{} menyatakan himpunan semua bentuk diferensial yang terdiferensialkan secara kontinu hingga orde $\ell$ dan berderajat $k$.

\inputaufgabe
{}
{

Misalkan $M$ sebuah
\definitionsverweis {manifold diferensiabel}{}{}
dan
\mathl{{ \mathcal E }^{ k } (  M   )}{} himpunan semua
$k$-\definitionsverweis {bentuk}{}{}
pada $M$. Tunjukkan bahwa ${ \mathcal E }^{ k } (  M   )$ merupakan sebuah
$R$-\definitionsverweis {modul}{}{}
untuk

\mavergleichskette
{\vergleichskette
{R
}
{ = }{ C^0(M,\R)
}
{  }{
}
{    }{
}
{   }{
}
}
{}{}{.}

}
{} {}

\inputaufgabe
{}
{

Misalkan $M$ sebuah
\definitionsverweis {manifold diferensiabel}{}{,}

\mavergleichskette
{\vergleichskette
{ P
}
{ \in }{ M
}
{  }{
}
{    }{
}
{   }{
}
}
{}{}{}
sebuah titik, dan

\mavergleichskette
{\vergleichskette
{ f
}
{ \in }{ C^1(M,\R)
}
{  }{
}
{    }{
}
{   }{
}
}
{}{}{}
sebuah
\definitionsverweis {fungsi yang terdiferensialkan secara kontinu}{}{.}
Misalkan

\mavergleichskette
{\vergleichskette
{ v
}
{ \in }{ T_PM
}
{  }{
}
{    }{
}
{   }{
}
}
{}{}{}
sebuah
\definitionsverweis {vektor singgung}{}{,}
yang direpresentasikan oleh suatu lintasan diferensiabel
\maabbdisp {\gamma} {{]{-\delta},\delta [}} {M
} {}
dengan

\mavergleichskette
{\vergleichskette
{ \gamma(0)
}
{ = }{ P
}
{  }{
}
{    }{
}
{   }{
}
}
{}{}{.}
Tunjukkan kesamaan

\mavergleichskettedisp
{\vergleichskette
{(df)(P,v)
}
{ =} { (f \circ \gamma)'(0)
}
{ } {
}
{ } {
}
{ } {
}
}
{}{}{.}

}
{} {}

\inputaufgabegibtloesung
{}
{

Misalkan $M$ sebuah
\definitionsverweis {manifold diferensiabel}{}{}
dan
\maabbdisp {f} { M } {\R
} {}
sebuah
\definitionsverweis {fungsi yang terdiferensialkan secara kontinu}{}{.}
Pada titik

\mavergleichskette
{\vergleichskette
{ P
}
{ \in }{ M
}
{  }{
}
{    }{
}
{   }{
}
}
{}{}{}
fungsi tersebut mempunyai suatu
\definitionsverweis {ekstremum lokal}{}{.}
Tunjukkan bahwa

\mavergleichskettedisp
{\vergleichskette
{ (df)_P
}
{ =} { 0
}
{ } {
}
{ } {
}
{ } {
}
}
{}{}{.}

}
{} {}

\inputaufgabegibtloesung
{}
{

Misalkan $M$ sebuah
\definitionsverweis {manifold diferensiabel}{}{}
\definitionsverweis {kompak}{}{}
dan
\maabbdisp {f} { M } {\R
} {}
sebuah
\definitionsverweis {fungsi yang terdiferensialkan secara kontinu}{}{.}
Tunjukkan bahwa terdapat sedikitnya dua titik pada $M$
\zusatzklammer {dengan syarat $M$ mempunyai sedikitnya dua titik} {} {}
di mana
$1$-\definitionsverweis {bentuk diferensial}{}{}
$df$ bernilai nol.

}
{} {}

\inputaufgabe
{}
{

Misalkan
\mathl{i: M\subseteq \R^n}{} sebuah
\definitionsverweis {submanifold tertutup}{}{.}
Tunjukkan bahwa untuk setiap
\definitionsverweis {fungsi diferensiabel}{}{}
\maabbdisp {f} {\R^n } {\R
} {}
berlaku hubungan

\mavergleichskettedisp
{\vergleichskette
{ i^* (df)
}
{ =} { d(f \circ i )
}
{ } {
}
{ } {
}
{ } {
}
}
{}{}{.}

}
{} {}

\inputaufgabe
{}
{

Misalkan
\maabbeledisp {f} {\R} {\R
} { t } {f(t)
} {,}
sebuah
\definitionsverweis {fungsi}{}{}
\definitionsverweis {yang terdiferensialkan secara kontinu}{}{,}
dan misalkan

\mavergleichskette
{\vergleichskette
{ \omega
}
{ = }{ g(s)ds
}
{  }{
}
{    }{
}
{   }{
}
}
{}{}{}
sebuah
$1$-\definitionsverweis {bentuk diferensial}{}{}
pada $\R$. Tentukan
\mathl{f^* \omega}{.}

}
{} {}

\inputaufgabegibtloesung
{}
{

Kita meninjau pemetaan
\maabbeledisp {\varphi} {\R^3} {\R^2
} {(x,y,z)} {(xyz^2,xy-z^3)
} {.}
Hitunglah matriks pemetaan
\maabbdisp {\bigwedge^2 T_P (\varphi)} { \bigwedge^2 T_P \R^3 } { \bigwedge^2 T_{\varphi(P)}  \R^2
} {}
di titik

\mavergleichskette
{\vergleichskette
{ P
}
{ = }{ (1,3,5)
}
{  }{
}
{    }{
}
{   }{
}
}
{}{}{}
terhadap suatu basis yang sesuai.

}
{} {}

\inputaufgabe
{}
{

Kita meninjau
\definitionsverweis {pemetaan terdiferensialkan kontinu}{}{}
\maabbeledisp {\varphi} {\R^2 \setminus \{(0,0)\}} {\R^2 \setminus \{(0,0)\}
} {(u,v)} {(u^2,v^3-u)
} {,}
dan
$2$-\definitionsverweis {bentuk diferensial}{}{}

\mavergleichskettedisp
{\vergleichskette
{ \omega
}
{ =} { { \frac{   1  }{  x^2+y^2 } } dx \wedge dy
}
{ } {
}
{ } {
}
{ } {
}
}
{}{}{.}
Tentukan
\definitionsverweis {bentuk diferensial tarik balik}{}{}

\mathl{\varphi^* \omega}{.}

}
{} {}

\inputaufgabegibtloesung
{}
{

Kita meninjau
$1$-\definitionsverweis {bentuk diferensial}{}{}

\mavergleichskettedisp
{\vergleichskette
{ \omega
}
{ =} { -vdu+udv
}
{ } {
}
{ } {
}
{ } {
}
}
{}{}{}
pada
$1$-\definitionsverweis {sfera}{}{}

\mavergleichskette
{\vergleichskette
{ S^1
}
{ \subset }{ \R^2
}
{  }{
}
{    }{
}
{   }{
}
}
{}{}{,}
dengan koordinat-koordinat pada ruang ambien $\R^2$ dinotasikan sebagai
\mathkor {} {u} {dan} {v} {.}
Tentukan
\mathl{\pi^* \omega}{} di bawah
\definitionsverweis {pemetaan terdiferensialkan kontinu}{}{}
\maabbeledisp {\pi} {\R^2 \setminus \{0\}} { S^{1}
} {(x , y) } { { \frac{   1  }{  \sqrt{x^2 + y^2}  } } (x , y)
} {.}

}
{} {}

\inputaufgabegibtloesung
{}
{

Hitunglah bentuk diferensial tarik balik
\mathl{\varphi^* \tau}{} untuk

\mavergleichskettedisp
{\vergleichskette
{ \tau
}
{ =} { dx \wedge dy \wedge dz - w dx \wedge dy \wedge dw + \cos (xy)  dx \wedge dz \wedge dw-yw dy \wedge dz \wedge dw
}
{ } {
}
{ } {
}
{ } {
}
}
{}{}{}
di bawah pemetaan
\maabbeledisp {\varphi} {\R^3} {\R^4
} {(r,s,t)} { (r^2s,t, \sin  r ,e^{st}) = (x,y,z,w)
} {.}

}
{} {}

\inputaufgabegibtloesung
{}
{

Misalkan
\maabbdisp {\psi} { L } { M
} {}
sebuah
\definitionsverweis {pemetaan diferensiabel}{}{}
antara
\definitionsverweis {dua manifold diferensiabel}{}{,}
misalkan
\maabbdisp {f} { M } {\R
} {}
sebuah fungsi pada $M$, dan $\omega$ sebuah
$k$-\definitionsverweis {bentuk}{}{}
pada $M$. Tunjukkan bahwa

\mavergleichskettedisp
{\vergleichskette
{ \psi^*(f \omega)
}
{ =} { \psi^*(f)  \psi^* \omega
}
{ } {
}
{ } {
}
{ } {
}
}
{}{}{.}

}
{} {}

\inputaufgabegibtloesung
{}
{

Misalkan

\mavergleichskette
{\vergleichskette
{ W
}
{ \subseteq }{ \R^m
}
{  }{
}
{    }{
}
{   }{
}
}
{}{}{}
dan

\mavergleichskette
{\vergleichskette
{ U
}
{ \subseteq }{ \R^n
}
{  }{
}
{    }{
}
{   }{
}
}
{}{}{}
merupakan
\definitionsverweis {himpunan-himpunan bagian terbuka}{}{}
dan misalkan
\maabbdisp {\psi} { W } { U
} {}
sebuah
\definitionsverweis {pemetaan}{}{} yang
\definitionsverweis {terdiferensialkan secara kontinu}{}{.}
Misalkan
\maabbdisp {f} { U } {\R
} {}
sebuah fungsi yang terdiferensialkan secara kontinu. Simpulkan dari
aturan rantai
bahwa

\mavergleichskettedisp
{\vergleichskette
{ d (\psi^*f)
}
{ =} { \psi^*(df)
}
{ } {
}
{ } {
}
{ } {
}
}
{}{}{,}
dengan $\psi^*$ menyatakan
\definitionsverweis {penarikan balik bentuk diferensial}{}{.}

}
{} {}

  \zwischenueberschrift{Soal untuk dikumpulkan}

\inputaufgabe
{6}
{

Misalkan $M$ sebuah
\definitionsverweis {manifold diferensiabel}{}{}
dengan
\definitionsverweis {bundel kotangen}{}{}

\mathl{T^*M}{.} Misalkan $\omega$ sebuah
$k$-\definitionsverweis {bentuk diferensial}{}{,} yaitu suatu pemetaan
\maabbdisp {\omega} { M } {\bigwedge^k T^*M
} {}
dengan
\mathl{\omega(P) \in \bigwedge^k T^*_P M}{} untuk setiap
\mathl{P \in M}{,} dengan produk eksterior ini dilengkapi topologi alaminya
\zusatzklammer {lihat
Soal 14.1} {} {.}
Tunjukkan bahwa pernyataan-pernyataan berikut ekuivalen:
\aufzaehlungdrei{ $\omega$
\definitionsverweis {
kontinu}{}{.}
}{Untuk setiap
\definitionsverweis {peta}{}{}
\maabb {\alpha} { U } { V
} {}
dengan
\mathl{V \subseteq \R^n}{} dan representasi lokal
\mathl{\alpha_* \omega =  \sum_{J,\, { \# \left(   J  \right) } =k} f_J dx_J}{,}
fungsi-fungsi $f_J$ kontinu.
}{Terdapat suatu tutupan terbuka
\mathl{M= \bigcup_{i \in I} U_i}{} oleh
\definitionsverweis {domain-domain peta}{}{} $U_i$ sedemikian sehingga, dalam representasi lokal
\mathl{\alpha_{i*} \omega =  \sum_{J,\, { \# \left(   J  \right) } =k} f_{i J} dx_J}{,}
fungsi-fungsi $f_{i J}$ kontinu.
}

}
{} {}

\inputaufgabe
{4}
{

Misalkan
\maabbdisp {\varphi} { L } { M
} {}
sebuah
\definitionsverweis {pemetaan diferensiabel}{}{}
antara dua
\definitionsverweis {manifold diferensiabel}{}{}
\mathkor {} {L} {dan} {M} {.}
Misalkan
\mathkor {} {\omega \in
{ \mathcal E }^{ k } (  M   )} {dan} {\tau \in
{ \mathcal E }^{  \ell } (  M   )} {}
\definitionsverweis {bentuk-bentuk diferensial}{}{} pada $M$. Tunjukkan kesamaan

\mavergleichskettedisp
{\vergleichskette
{ \varphi^* ( \omega \wedge \tau)
}
{ =} {\varphi^* \omega \wedge \varphi^* \tau
}
{ } {
}
{ } {
}
{ } {
}
}
{}{}{.}

}
{} {}

\inputaufgabe
{4}
{

Kita meninjau
\definitionsverweis {pemetaan terdiferensialkan kontinu}{}{}
\maabbeledisp {\varphi} {\R^3 \setminus \{(0,0,0)\}} {N = \left\{ (x,y,z) \in \R^3  \mid   z \neq 0        \right\}
} {(u,v,w)} {(uvw,u^2-vw^5,u^2+v^2+w^2)
} {,}
dan
$2$-\definitionsverweis {bentuk diferensial}{}{}

\mathdisp {\omega =  z^2dx \wedge dy+{ \frac{   xy }{ z } }dx \wedge dz+(xe^y-z)dy \wedge dz} {  }

pada $N$. Tentukan
\definitionsverweis {bentuk diferensial tarik balik}{}{}

\mathl{\varphi^* \omega}{.}

}
{} {}

\inputaufgabe
{6}
{

Kita meninjau
\definitionsverweis {permukaan bola satuan}{}{}

\mavergleichskette
{\vergleichskette
{ S^2
}
{ \subseteq }{ \R^3
}
{  }{
}
{    }{
}
{   }{
}
}
{}{}{,}
dengan koordinat-koordinat pada $\R^3$ dinotasikan sebagai $x,y,z$. Untuk titik-titik manakah

\mavergleichskette
{\vergleichskette
{ P
}
{ \in }{ S^2
}
{  }{
}
{    }{
}
{   }{
}
}
{}{}{}
pembatasan
\mathkor {} {dx} {dan} {dy} {}
pada $S^2$ membentuk suatu
\definitionsverweis {basis}{}{}
dari
\definitionsverweis {ruang kotangen}{}{}

\mathl{T^*_PS^2}{.}

}
{} {}

\section*{Solusi yang disediakan oleh sumber}
\addcontentsline{toc}{section}{Solusi yang disediakan oleh sumber}
\subsection*{Solusi Soal 14.5}
Misalkan

\mavergleichskette
{\vergleichskette
{ P
}
{ \in }{  U
}
{ \subseteq }{ M
}
{    }{
}
{   }{
}
}
{}{}{}
dengan U sebagai suatu
\definitionsverweis {domain peta}{}{}
dan dengan peta
\maabb {\alpha} { U } {V \subseteq \R^n
} {.}
Maka
\mathl{f \circ \alpha^{-1}}{} juga mempunyai suatu ekstremum lokal di titik $\alpha(P)$.
Oleh karena itu, berdasarkan
Fakta   (2),
\definitionsverweis {diferensial total}{}{}
\maabbdisp {\left(D (f \circ \alpha^{-1}) \right)_{\alpha(P) }} {\R^n } {\R
} {}
merupakan pemetaan nol. Dengan demikian,

\mavergleichskette
{\vergleichskette
{ (df)_P
}
{ = }{ 0
}
{  }{
}
{    }{
}
{   }{
}
}
{}{}{.}
\subsection*{Solusi Soal 14.6}
Jika $f$ konstan, maka

\mavergleichskette
{\vergleichskette
{ df
}
{ = }{ 0
}
{  }{
}
{    }{
}
{   }{
}
}
{}{}{}
dan pernyataan tersebut jelas. Jadi, misalkan $f$ tidak konstan. Berdasarkan
Fakta,
fungsi kontinu tersebut mencapai maksimum global dan minimum global, katakanlah masing-masing di titik
\mathkor {} {P} {dan} {Q} {.}
Berdasarkan
soal,
maka

\mavergleichskettedisp
{\vergleichskette
{ (df)_P
}
{ =} { (df)_Q
}
{ =} { 0
}
{ } {
}
{ } {
}
}
{}{}{.}
\subsection*{Solusi Soal 14.9}
Matriks Jacobi dari $\varphi$ secara umum adalah

\mathdisp {\begin{pmatrix}  yz^2 &  xz^2 &  2xyz \\  y  &  x  &  -3z^2 \end{pmatrix}} { . }

Oleh karena itu, di titik

\mavergleichskette
{\vergleichskette
{ P
}
{ = }{ (1,3,5)
}
{  }{
}
{    }{
}
{   }{
}
}
{}{}{}
matriks Jacobi tersebut adalah

\mathdisp {\begin{pmatrix}  75 &  25 &  30 \\  3 &  1 &  -75  \end{pmatrix}} {  }
.

Matriks ini mendeskripsikan pemetaan linear
\maabbdisp {L} {T_P\R^3 \cong \R^3} { T_{\varphi(P)} \R^2 \cong \R^2
} {}
terhadap basis-basis standar. Kita menentukan representasi matriks untuk pemetaan
\maabbdisp {\bigwedge^2 L} {\bigwedge^2 \R^3 } { \bigwedge^2 \R^2
} {}
terhadap basis-basis
$e_1 \wedge e_2,\, e_1 \wedge e_3$ , $e_2 \wedge e_3$
\zusatzklammer {di ruas kiri} {} {}
dan
\mathl{e_1 \wedge e_2}{}
\zusatzklammer {di ruas kanan} {} {.}
Untuk itu, kita perlu menghitung citra produk-produk eksterior tersebut. Kita peroleh

\mavergleichskettealign
{\vergleichskettealign
{ (\bigwedge^2 L ) (e_1 \wedge e_2)
}
{ =} { L(e_1) \wedge L(e_2)
}
{ =} { (75e_1 + 3e_2) \wedge (  25 e_1 + 1 e_2 )
}
{ =} { 75 e_1 \wedge e_2 + 75 e_2 \wedge  e_1
}
{ =} { 75 e_1 \wedge e_2 - 75 e_1 \wedge  e_2
}
}
{
\vergleichskettefortsetzungalign
{ =} { 0
}
{ } {}
{ } {}
{ } {}
}
{}{,}

\mavergleichskettealign
{\vergleichskettealign
{ (\bigwedge^2 L ) (e_1 \wedge e_3)
}
{ =} { L(e_1) \wedge L(e_3)
}
{ =} { (75e_1 + 3e_2) \wedge (  30 e_1 - 75 e_2 )
}
{ =} { - 75^2 e_1 \wedge e_2 + 90 e_2 \wedge  e_1
}
{ =} { (-5625 - 90) e_1 \wedge e_2
}
}
{
\vergleichskettefortsetzungalign
{ =} { -5715 e_1 \wedge e_2
}
{ } {}
{ } {}
{ } {}
}
{}{,}
dan

\mavergleichskettealign
{\vergleichskettealign
{ (\bigwedge^2 L ) (e_2 \wedge e_3)
}
{ =} { L(e_2) \wedge L(e_3)
}
{ =} { (  25 e_1 + 1 e_2 )   \wedge ( 30 e_1 - 75 e_2)
}
{ =} { - 1875 e_1 \wedge e_2 + 30 e_2 \wedge  e_1
}
{ =} { -  1905 e_1 \wedge  e_2
}
}
{}
{}{.}
Jadi, matriks yang mendeskripsikan pemetaan tersebut adalah

\mathdisp {\left(  0,-5715,-1905                     \right)} { . }
\subsection*{Solusi Soal 14.11}
Kita peroleh

\mavergleichskettealign
{\vergleichskettealign
{ \pi^*du
}
{ =} { d (u \circ \pi)
}
{ =} { d   \left(  { \frac{   x  }{  \sqrt{x^2 + y^2}  } }  \right)
}
{ =} { { \frac{   \partial   \left(  { \frac{   x  }{  \sqrt{x^2 + y^2}  } }  \right)    }{  \partial   x  } }   dx + { \frac{   \partial   \left(  { \frac{   x  }{  \sqrt{x^2 + y^2}  } }  \right)    }{  \partial   y  } } dy
}
{ =} { { \frac{   \sqrt{x^2+y^2}     -x^2 (x^2 + y^2)^{-1/2}  }{  x^2+y^2  } }  dx   - { \frac{   xy }{  \sqrt{x^2+y^2}^3  } }  dy
}
}
{
\vergleichskettefortsetzungalign
{ =} { { \frac{   x^2+y^2 -x^2    }{  \sqrt{ x^2+y^2}^3  } }  dx  - { \frac{   xy }{  \sqrt{x^2+y^2}^3  } }  dy
}
{ =} { { \frac{   y^2  }{  \sqrt{ x^2+y^2}^3  } }  dx  - { \frac{   xy }{  \sqrt{x^2+y^2}^3  } }  dy
}
{ } {}
{ } {}
}
{}{}
dan

\mavergleichskettealign
{\vergleichskettealign
{ \pi^*dv
}
{ =} { d (v \circ \pi)
}
{ =} { d \left(  { \frac{   y  }{  \sqrt{x^2 + y^2}  } }  \right)
}
{ =} { { \frac{   \partial   \left(  { \frac{   y  }{  \sqrt{x^2 + y^2}  } }  \right)   }{  \partial   x  } } dx + { \frac{   \partial   \left(  { \frac{   y  }{  \sqrt{x^2 + y^2}  } }  \right)    }{  \partial   y  } } dy
}
{ =} { { \frac{   -xy }{  \sqrt{x^2+y^2}^3  } } dx + { \frac{   \sqrt{x^2+y^2} -y^2 (x^2 + y^2)^{-1/2}  }{  x^2+y^2  } } dy
}
}
{
\vergleichskettefortsetzungalign
{ =} { { \frac{   -xy }{  \sqrt{x^2+y^2}^3  } } dx + { \frac{   x^2  }{  \sqrt{x^2+y^2}^3  } }  dy
}
{ } {
}
{ } {}
{ } {}
}
{}{.}
Dengan demikian,

\mavergleichskettealigndrucklinks
{\vergleichskettealigndrucklinks
{ \pi^*(-vdu+udv)
}
{ =} { - (v \circ \pi)  \pi^*du + (u \circ \pi) \pi^*dv
}
{ =} { - { \frac{   y  }{  \sqrt{x^2+y^2}  } } \left(  { \frac{   y^2  }{  \sqrt{ x^2+y^2}^3  } }  dx  - { \frac{   xy }{  \sqrt{x^2+y^2}^3  } }  dy   \right) + { \frac{   x  }{  \sqrt{x^2+y^2}  } }  \left(  { \frac{   -xy }{  \sqrt{x^2+y^2}^3  } } dx + { \frac{   x^2  }{  \sqrt{x^2+y^2}^3  } } dy  \right)
}
{ =} { { \frac{   1 }{  \left(  x^2+y^2  \right)^2 } }  \left(  \left(  -y^3-x^2y  \right) dx + \left(  xy^2+x^3  \right) dy \right)
}
{ =} { { \frac{   1 }{  x^2+y^2  } }  \left(  -  y dx + x dy \right)
}
}
{}{}{.}
\subsection*{Solusi Soal 14.12}
Kita peroleh

\mavergleichskettedisp
{\vergleichskette
{dx
}
{ =} { d(r^2s)
}
{ =} { 2rs dr  +r^2ds
}
{ } {
}
{ } {
}
}
{}{}{,}

\mavergleichskettedisp
{\vergleichskette
{dy
}
{ =} { dt
}
{ } {
}
{ } {
}
{ } {
}
}
{}{}{,}

\mavergleichskettedisp
{\vergleichskette
{dz
}
{ =} { d (\sin  r )
}
{ =} { \cos  r  dr
}
{ } {
}
{ } {
}
}
{}{}{}
dan

\mavergleichskettedisp
{\vergleichskette
{dw
}
{ =} { d(e^{st})
}
{ =} { se^{st} dt+te^{st} ds
}
{ } {
}
{ } {
}
}
{}{}{.}
Dengan demikian,

\mavergleichskettealign
{\vergleichskettealign
{ \varphi^* \tau
}
{ =} { d(r^2s) \wedge dt \wedge d( \sin  r )  -e^{st} d(r^2s) \wedge dt \wedge d(e^{st})
}
{ \, \, \, \,} {+ \cos (r^2st)  d (r^2s) \wedge d ( \sin  r ) \wedge d(e^{st}) - te^{st} dt \wedge d( \sin  r ) \wedge d(e^{st})
}
{ =} { r^2 \cos  r  ds \wedge dt \wedge dr  - 2rs t e^{st} e^{st} dr \wedge dt \wedge ds
}
{ \, \, \, \,} { + r^2 se^{st} \cos (r^2st) \cos  r  ds \wedge dr \wedge dt  -  t^2 e^{st}  e^{st} \cos  r dt \wedge dr \wedge ds
}
}
{
\vergleichskettefortsetzungalign
{ =} { (r^2 \cos  r +2 rst e^{2st} - r^2 s e ^{st} \cos (r^2st) \cos  r - t^2e^{ 2st} \cos  r     )  dr \wedge ds \wedge dt
}
{ } {}
{ } {}
{ } {}
}
{}{}
\subsection*{Solusi Soal 14.13}
Misalkan
\mathl{P \in L}{} dan
\mathl{v_1 , \ldots , v_k \in T_PL}{.} Maka

\mavergleichskettealign
{\vergleichskettealign
{ \psi^* \left(  f \omega) (P,v_1 , \ldots , v_k  \right)
}
{ =} { (f \omega) \left(  \psi(P), T_P\psi (v_1) , \ldots , T_P\psi (v_k )       \right)
}
{ =} { f(\psi(P))   \cdot \omega \left(  \psi(P), T_P\psi (v_1) , \ldots , T_P\psi (v_k )       \right)
}
{ =} { \left(  \psi^* (f)  \right) (P)  \cdot \omega \left(  \psi(P), T_P\psi (v_1) , \ldots , T_P\psi (v_k )       \right)
}
{ =} { \left(  \psi^* (f)  \right) (P)   \cdot \left(  \psi^*( \omega)  \right)  \left(  P,v_1 , \ldots , v_k  \right)
}
}
{
\vergleichskettefortsetzungalign
{ =} { \left(  \psi^* (f)  \psi^*( \omega)  \right) \left(  P,v_1 , \ldots , v_k  \right)
}
{ } {}
{ } {}
{ } {}
}
{}{.}
\subsection*{Solusi Soal 14.14}
Misalkan

\mavergleichskette
{\vergleichskette
{ P
}
{ \in }{  W
}
{  }{
}
{    }{
}
{   }{
}
}
{}{}{}
dan

\mavergleichskette
{\vergleichskette
{ v
}
{ \in }{  \R^m
}
{  }{
}
{    }{
}
{   }{
}
}
{}{}{.}
Di satu pihak,

\mavergleichskettedisp
{\vergleichskette
{ d( \psi^* f) (P,v)
}
{ =} { D_P( f \circ \psi) (v)
}
{ =} { \left(  \left(  D_{\psi(P)} f  \right)  \circ \left(  D_P  \psi  \right)  \right) (v)
}
{ =} { D_{\psi(P)} (f) \left(  D_P \psi(v)  \right)
}
{ } {
}
}
{}{}{.}
Di pihak lain, kita juga memperoleh

\mavergleichskettedisp
{\vergleichskette
{ \left(  \psi^*(df)  \right) (P,v)
}
{ =} { (df) \left(  \psi(P), (D_P \psi )(v)  \right)
}
{ =} { D_{\psi(P)} (f) \left(  D_P \psi (v)  \right)
}
{ } {
}
{ } {
}
}
{}{}{.}

\part{Unit 15}
\chapter[Kuliah 15]{Kuliah 15: Integrasi pada Manifold}
\setcounter{section}{15}

Kini kita beralih ke teori integrasi pada manifold. Titik pangkalnya ialah tersedianya suatu $n$-bentuk pada manifold berdimensi $n$. Untuk suatu himpunan bagian terbuka

\mavergleichskette
{\vergleichskette
{ V
}
{ \subseteq }{ \R^n
}
{  }{
}
{    }{
}
{   }{
}
}
{}{}{}
dengan koordinat
\mathl{x_1 , \ldots , x_n}{,} integrasi terhadap bentuk
\mathl{dx_1 \wedge \ldots \wedge dx_n}{} bersesuaian dengan integrasi terhadap ukuran Lebesgue. Pada suatu manifold, bentuk tersebut dan ukuran yang terkait harus \anfuehrung{direkatkan bersama}{.}

  \zwischenueberschrift{Bentuk volume positif pada manifold}

Dalam definisi berikut, untuk suatu peta
\maabb {\alpha} { U } { V
} {}
dan suatu bentuk diferensial $\omega$ pada $U$, kita melambangkan bentuk diferensial yang ditranspor ke $V$ dengan
\mathl{\alpha_* \omega}{.} Bentuk ini sama dengan bentuk tarik balik
\mathl{\alpha^{-1 *} \omega}{.}

\inputdefinition
{}
{

Misalkan $M$ suatu
\definitionsverweis {manifold diferensiabel}{}{}
\definitionsverweis {berdimensi}{}{} $n$
dan $\omega$ suatu
\definitionsverweis {bentuk diferensial}{}{}
\definitionsverweis {terukur}{}{}
berderajat $n$ pada $M$. Bentuk $\omega$ disebut \definitionswort {bentuk volume positif}{,} jika untuk setiap
\definitionsverweis {peta}{}{}
\zusatzklammer {dari suatu atlas yang diberikan} {} {}
\maabbdisp {\alpha} { U } { V
} {}
\zusatzklammer {dengan

\mavergleichskettek
{\vergleichskettek
{ V
}
{ \subseteq }{ \R^n
}
{  }{
}
{    }{
}
{   }{
}
}
{}{}{}
dan fungsi koordinat \mathlk{x_1 , \ldots , x_n}{}} {} {}
dalam representasi lokal bentuk diferensial tersebut,

\mavergleichskettedisp
{\vergleichskette
{ \alpha_* \omega
}
{ =} { f dx_1 \wedge \ldots \wedge dx_n
}
{ } {
}
{ } {
}
{ } {
}
}
{}{}{}
fungsi $f$ bernilai positif di mana-mana.\zusatzfussnote {Fungsi yang terkait dengan setiap peta $U$, yang di sini dilambangkan $f$ dan dikalikan dengan bentuk standar berderajat $n$, bersesuaian dengan kerapatan yang disebutkan pada akhir Kuliah 13, yang dengannya suatu ukuran pada manifold dapat dideskripsikan} {.} {}

}
Di sini fungsi $f$ ditentukan oleh

\mavergleichskettedisphandlinks
{\vergleichskettedisphandlinks
{ f(Q)
}
{ =} { \omega \left(  \alpha^{-1} (Q) ,T_Q \left(  \alpha^{-1}  \right)(e_1) \wedge \ldots \wedge T_Q \left(  \alpha^{-1}  \right)( e_n)  \right)
}
{ } {
}
{ } {
}
{ } {
}
}
{}{}{}
secara tunggal. Bentuk volume positif semacam itu hanya dapat ada jika manifold tersebut
\definitionsverweis {dapat diorientasikan}{}{}
\zusatzklammer {lihat
Lema 15.5
di bawah} {} {.}



\inputfaktbeweis
{Mannigfaltigkeit/Abzählbar/Positive Volumenform/Zugehöriges Maß/Vorbereitende Eigenschaften/Fakt}
{Lema}
{}
{

\faktsituation {Misalkan $M$ suatu
\definitionsverweis {manifold diferensiabel}{}{}
\definitionsverweis {berdimensi}{}{} $n$
dengan
\definitionsverweis {basis topologi terhitung}{}{}
dan misalkan $\omega$ suatu
\definitionsverweis {bentuk volume positif}{}{}
pada $M$. Untuk suatu peta
\maabbdisp {\alpha} { U } { V
} {}
dengan

\mavergleichskette
{\vergleichskette
{ \alpha_* ( \omega \vert_U)
}
{ = }{ fdx_1 \wedge \ldots \wedge dx_n
}
{  }{
}
{    }{
}
{   }{
}
}
{}{}{}
dan suatu
\definitionsverweis {himpunan bagian terukur}{}{}

\mavergleichskette
{\vergleichskette
{ T
}
{ \subseteq }{ U
}
{  }{
}
{    }{
}
{   }{
}
}
{}{}{}
kita tetapkan

\mavergleichskettedisp
{\vergleichskette
{ \nu ( \alpha,T )
}
{ \defeq} { \int_{ \alpha(T)  }   f    \,  d  \lambda^n
}
{ } {
}
{ } {
}
{ } {
}
}
{}{}{} }
\faktuebergang {Maka berlaku sifat-sifat berikut.}
\faktfolgerung {\aufzaehlungdrei{Jika

\mavergleichskette
{\vergleichskette
{ T
}
{ \subseteq }{ U_1 \cap U_2
}
{  }{
}
{    }{
}
{   }{
}
}
{}{}{}
dan kedua himpunan di ruas kanan merupakan lingkungan koordinat, maka

\mavergleichskette
{\vergleichskette
{ \nu(\alpha_1,T)
}
{ = }{ \nu(\alpha_2,T)
}
{  }{
}
{    }{
}
{   }{
}
}
{}{}{.}
}{Untuk suatu himpunan bagian terukur

\mavergleichskette
{\vergleichskette
{ T
}
{ \subseteq }{ M
}
{  }{
}
{    }{
}
{   }{
}
}
{}{}{}
terdapat suatu
\definitionsverweis {gabungan saling lepas}{}{}
\definitionsverweis {terhitung}{}{}

\mavergleichskette
{\vergleichskette
{ T
}
{ = }{ \biguplus_{i \in I} T_i
}
{  }{
}
{    }{
}
{   }{
}
}
{}{}{}
sedemikian sehingga setiap $T_i$ seluruhnya termuat dalam suatu lingkungan koordinat $U_i$.
}{Jumlah
\mathl{\sum_{i \in I} \nu( \alpha_i ,T_i)}{} tidak bergantung pada dekomposisi saling lepas terhitung yang dipilih pada (2).
}}
\faktzusatz {}
\faktzusatz {}

}
{

\teilbeweis {}{}{}
{(1). Karena

\mavergleichskette
{\vergleichskette
{ T
}
{ \subseteq }{ U_1 \cap U_2
}
{  }{
}
{    }{
}
{   }{
}
}
{}{}{}
kita dapat mengasumsikan

\mavergleichskette
{\vergleichskette
{ U
}
{ = }{ U_1
}
{ = }{ U_2
}
{    }{
}
{   }{
}
}
{}{}{}
\zusatzklammer {tetapi dengan dua pemetaan koordinat yang berbeda,
\mathkor {} {\alpha_1} {dan} {\alpha_2} {}
masing-masing menuju
\mathkor {} {V_1} {dan} {V_2} {}} {} {.}
Misalkan
\maabbdisp {\psi=\alpha_2 \circ \alpha_1^{-1}} { V_1 } { V_2
} {}
merupakan
\definitionsverweis {pergantian koordinat difeomorfik}{}{.}
Maka berdasarkan
Teorema 14.3 (Teori Ukuran dan Integrasi (Osnabrück 2022-2023))
dan
Korolari 14.9,
serta karena dari kepositifan

\mavergleichskette
{\vergleichskette
{ f_1
}
{ = }{ \left(  f_2 \circ \psi  \right) \det  \left(  D  \psi   \right)
}
{  }{
}
{    }{
}
{   }{
}
}
{}{}{}
dan $f_2$ diperoleh bahwa determinannya juga positif, berlaku kesamaan

\mavergleichskettealign
{\vergleichskettealign
{ \nu \left(  \alpha_2,T  \right)
}
{ =} { \int_{  \alpha_2(T)   }   f_2    \,  d  \lambda^n
}
{ =} { \int_{  \alpha_1(T)   }   \left(  f_2 \circ \psi  \right) \mid  \det  \left(  D  \psi   \right)  \mid      \,  d  \lambda^n
}
{ =} { \int_{  \alpha_1(T)   }   \left(  f_2 \circ \psi  \right) \cdot \det  \left(  D  \psi   \right)     \,  d  \lambda^n
}
{ =} { \int_{  \alpha_1(T)   }   f_1    \,  d  \lambda^n
}
}
{
\vergleichskettefortsetzungalign
{ =} { \nu(\alpha_1,T)
}
{ } {}
{ } {}
{ } {}
}
{}{.}}
{}
\teilbeweis {}{}{}
{(2). Misalkan

\mathbed {U_n} {}
 {n \in \N} {}
 {} {} {} {,}
suatu atlas terhitung. Maka kita dapat mengambil himpunan

\mavergleichskette
{\vergleichskette
{ T_n
}
{ = }{ T \cap (U_n \setminus \bigcup_{m < n} U_m )
}
{  }{
}
{    }{
}
{   }{
}
}
{}{}{.}}
{}
\teilbeweis {}{}{}
{(3). Misalkan

\mavergleichskette
{\vergleichskette
{ T
}
{ = }{ \biguplus_{i \in I} T_i
}
{ = }{ \biguplus_{j \in J} S_j
}
{    }{
}
{   }{
}
}
{}{}{}
dua dekomposisi terukur saling lepas terhitung yang setiap anggotanya termuat dalam suatu daerah koordinat. Daerah beserta pemetaan koordinatnya masing-masing adalah
\mathl{(U_i, \alpha_i)}{} dengan fungsi $f_i$ yang mendeskripsikan bentuk tersebut dan
\mathl{(V_j, \beta_j)}{} dengan fungsi $g_j$ yang mendeskripsikannya. Kita tinjau dekomposisi terhitung yang juga diberikan oleh himpunan

\mathbed {T_i \cap S_j} {}
 {(i,j) \in I \times J} {}
 {} {} {} {,}
Berdasarkan
Lema 10.1 (Teori Ukuran dan Integrasi (Osnabrück 2022-2023))
\zusatzklammer {yang diterapkan pada masing-masing citra peta} {} {}
dan dengan menggunakan bagian (1), diperoleh

\mavergleichskettealign
{\vergleichskettealign
{ \sum_{i \in I } \left(  \int_{  \alpha_i \left(  T_i  \right)   }   f_i    \,  d  \lambda^n  \right)
}
{ =} { \sum_{i \in I } \left(  \sum_{j \in J} \int_{ \alpha_i \left(  T_i \cap S_j  \right)   }   f_i    \,  d  \lambda^n  \right)
}
{ =} { \sum_{i \in I } \left(  \sum_{j \in J} \int_{ \beta_j \left(  T_i \cap S_j  \right)   }   g_j    \,  d  \lambda^n  \right)
}
{ =} { \sum_{j \in J } \left(  \sum_{i \in I} \int_{ \beta_j \left(  T_i \cap S_j  \right)   }   g_j    \,  d  \lambda^n  \right)
}
{ =} { \sum_{j \in J } \left(  \int_{  \beta_j(S_j)   }   g_j    \,  d  \lambda^n  \right)
}
}
{}
{}{.}}
{}

}

\inputdefinition
{}
{

Misalkan $M$ suatu
\definitionsverweis {manifold diferensiabel}{}{}
\definitionsverweis {berdimensi}{}{} $n$
dengan suatu
\definitionsverweis {basis topologi terhitung}{}{}
dan misalkan $\omega$ suatu
\definitionsverweis {bentuk volume positif}{}{}
pada $M$. Untuk setiap
\definitionsverweis {himpunan Borel}{}{}

\mavergleichskette
{\vergleichskette
{ T
}
{ \subseteq }{ M
}
{  }{
}
{    }{
}
{   }{
}
}
{}{}{}
pilih suatu
\definitionsverweis {dekomposisi}{}{}
\definitionsverweis {terhitung}{}{}

\mavergleichskette
{\vergleichskette
{ T
}
{ = }{  \biguplus_{i \in I} T_i
}
{  }{
}
{    }{
}
{   }{
}
}
{}{}{}
\zusatzklammer {dengan suatu daerah koordinat terbuka yang memenuhi

\mavergleichskettek
{\vergleichskettek
{ T_i
}
{ \subseteq }{ U_i
}
{  }{
}
{    }{
}
{   }{
}
}
{}{}{}
dan

\mavergleichskettek
{\vergleichskettek
{  \alpha_{i*} \omega
}
{ = }{ f_i dx_1 \wedge \ldots \wedge dx_n
}
{  }{
}
{    }{
}
{   }{
}
}
{}{}{}} {} {.}
Bilangan yang didefinisikan oleh

\mavergleichskettedisp
{\vergleichskette
{
\int_{ T  }   \omega
}
{ =} { \sum_{i \in I} \int_{ \alpha_i(T_i)  }  f_i    \,  d  \lambda^n
}
{ } {
}
{ } {
}
{ } {
}
}
{}{}{}
\zusatzklammer {sebagai elemen $\overline{ \R }_{\geq 0}$} {} {}
disebut \definitionswort {ukuran}{} dari $T$ terhadap $\omega$, atau \definitionswort {integral}{} $\omega$ atas $T$.

}
Berdasarkan lema di atas, ukuran volume ini terdefinisi dengan baik dan merupakan suatu ukuran
$\sigma$-\definitionsverweis {hingga}{}{.}
Pada setiap daerah koordinat, citranya dapat ditutupi secara terhitung oleh irisan kotak terbatas dengan himpunan tempat fungsi kerapatan dibatasi; atlas terhitung kemudian memberikan tutupan terhitung bagi manifold tersebut. Untuk bentuk volume kontinu, kesimpulan ini juga mengikuti dari
Soal 15.2.
Untuk suatu himpunan terbuka

\mavergleichskette
{\vergleichskette
{ M
}
{ \subseteq }{ \R^n
}
{  }{
}
{    }{
}
{   }{
}
}
{}{}{,}
suatu himpunan bagian terukur

\mavergleichskette
{\vergleichskette
{ T
}
{ \subseteq }{ M
}
{  }{
}
{    }{
}
{   }{
}
}
{}{}{}
dan suatu $n$-bentuk positif

\mavergleichskette
{\vergleichskette
{ \omega
}
{ = }{ f dx_1 \wedge \ldots \wedge dx_n
}
{  }{
}
{    }{
}
{   }{
}
}
{}{}{}
berlaku langsung

\mavergleichskettedisp
{\vergleichskette
{ \int_T \omega
}
{ =} { \int_T f d \lambda^n
}
{ } {
}
{ } {
}
{ } {
}
}
{}{}{.}


\inputfaktbeweis
{Volumenform/Integration/Eigenschaften/Fakt}
{Lema}
{}
{

\faktsituation {Misalkan $M$ suatu
\definitionsverweis {manifold diferensiabel}{}{}
\definitionsverweis {berdimensi}{}{} $n$ dengan suatu
\definitionsverweis {basis topologi terhitung}{}{}
dan misalkan
\mathkor {} {\omega_1} {dan} {\omega_2} {}
\definitionsverweis {bentuk volume positif}{}{}
pada $M$.}
\faktfolgerung {Maka untuk setiap
\definitionsverweis {himpunan bagian terukur}{}{}

\mavergleichskette
{\vergleichskette
{T
}
{ \subseteq }{ M
}
{  }{
}
{    }{
}
{   }{
}
}
{}{}{}
dan

\mavergleichskette
{\vergleichskette
{a,b
}
{ \in }{\R_+
}
{  }{
}
{    }{
}
{   }{
}
}
{}{}{}
berlaku hubungan

\mavergleichskettedisp
{\vergleichskette
{
\int_{ T  }  (a \omega_1 + b \omega_2)
}
{ =} { a
\int_{  T  }   \omega_1 + b
\int_{  T  }   \omega_2
}
{ } {
}
{ } {
}
{ } {
}
}
{}{}{.}}
\faktzusatz {}
\faktzusatz {}

 }
{ Lihat Soal 15.5. }

  \zwischenueberschrift{Bentuk volume dan orientasi}

Keberadaan suatu bentuk volume kontinu yang tidak pernah nol pada manifold berkaitan erat dengan keterorientasiannya. Kebalikan dari pernyataan berikut juga akan kita buktikan dalam
Teorema 22.11.



\inputfaktbeweis
{Nullstellenfreie Volumenform/Impliziert orientierte (Karten) Mannigfaltigkeit/Fakt}
{Lema}
{}
{

\faktsituation {Misalkan $M$ suatu
\definitionsverweis {manifold diferensiabel}{}{}
\definitionsverweis {berdimensi}{}{} $m$
dan}
\faktvoraussetzung {$\omega$ suatu
\definitionsverweis {bentuk volume}{}{}
\definitionsverweis {kontinu}{}{}
yang tidak pernah nol pada $M$.}
\faktfolgerung {Maka terdapat suatu
\zusatzklammer {yang ekuivalen secara difeomorfik} {} {}
\definitionsverweis {atlas berorientasi}{}{}
untuk $M$ sedemikian sehingga $\omega$ menjadi suatu
\definitionsverweis {bentuk volume positif}{}{}
terhadap atlas tersebut.}
\faktzusatz {Khususnya, $M$
\definitionsverweis {dapat diorientasikan}{}{.}}
\faktzusatz {}

}
{

Untuk

\mavergleichskette
{\vergleichskette
{ P
}
{ \in }{ M
}
{  }{
}
{    }{
}
{   }{
}
}
{}{}{}
kita tinjau
\definitionsverweis {daerah koordinat}{}{}

\mavergleichskette
{\vergleichskette
{ P
}
{ \in }{ U
}
{  }{
}
{    }{
}
{   }{
}
}
{}{}{}
yang mempunyai sifat bahwa $U$ homeomorfik dengan suatu bola terbuka

\mavergleichskette
{\vergleichskette
{ V
}
{ \subseteq }{ \R^m
}
{  }{
}
{    }{
}
{   }{
}
}
{}{}{.}
Berlaku

\mavergleichskettedisp
{\vergleichskette
{ \bigwedge^m TU
}
{ \cong} { \bigwedge^m TV
}
{ \cong} { V \times \bigwedge^m \R^m
}
{ \cong} { V \times \R
}
{ } {
}
}
{}{}{}
melalui
\mathl{\bigwedge^m T (\alpha)}{.} Isomorfisme terakhir diberikan oleh basis standar dengan koordinat
\mathl{x_1 , \ldots , x_m}{.} Misalkan

\mavergleichskette
{\vergleichskette
{ \alpha_* \omega
}
{ = }{ fdx_1 \wedge \ldots \wedge dx_m
}
{  }{
}
{    }{
}
{   }{
}
}
{}{}{}
suatu $m$-bentuk diferensial yang bersesuaian pada $V$. Bentuk ini tidak pernah nol, dan karena $V$
\definitionsverweis {terhubung}{}{,}
menurut
teorema nilai antara,
fungsi $f$ bernilai positif seluruhnya atau negatif seluruhnya. Dalam kasus negatif, kita mengomposisikan peta dengan pencerminan koordinat yang membalik orientasi, setara dengan mengganti tanda satu vektor basis. Dengan cara ini, untuk setiap titik kita memperoleh suatu lingkungan koordinat tempat bentuk tersebut positif. Untuk dua peta
\maabb {\alpha} { U } { V
} {}
dan
\maabb {\beta} { U } { V'
} {}
dengan pemetaan transisi

\mavergleichskette
{\vergleichskette
{ \psi
}
{ = }{ \beta \circ \alpha^{-1}
}
{  }{
}
{    }{
}
{   }{
}
}
{}{}{}
dan deskripsi lokal
\mathkor {} {\alpha_*\omega =f dx_1 \wedge \ldots \wedge dx_m} {dan} {\beta_*\omega =g dy_1 \wedge \ldots \wedge dy_m} {}
maka, karena

\mavergleichskette
{\vergleichskette
{ \psi^* (\beta_*\omega)
}
{ = }{ \alpha_* \omega
}
{  }{
}
{    }{
}
{   }{
}
}
{}{}{}
menurut Korolari 14.9 berlaku hubungan

\mavergleichskette
{\vergleichskette
{ f
}
{ = }{ g \circ \psi \det  \left(  D  \psi  \right)
}
{  }{
}
{    }{
}
{   }{
}
}
{}{}{.}
Karena
\mathkor {} {f} {dan} {g} {}
positif, determinan tersebut juga harus positif, sehingga pemetaan transisinya
\definitionsverweis {mempertahankan orientasi}{}{}
dan dengan demikian manifoldnya
\definitionsverweis {berorientasi}{}{.}

}

  \zwischenueberschrift{Bentuk volume pada serat}

Untuk pernyataan berikut, bandingkan dengan
Catatan 13.11.


\inputfaktbeweis
{Differenzierbare reguläre Abbildung/R^n/Faser besitzt Volumenform über Gradienten/Fakt}
{Korollar}
{}
{

\faktsituation {Misalkan

\mavergleichskette
{\vergleichskette
{ G
}
{ \subseteq }{ \R^n
}
{  }{
}
{    }{
}
{   }{
}
}
{}{}{}
terbuka dan misalkan
\maabbdisp {\varphi} { G } {\R^\ell
} {}
\zusatzklammer {dengan

\mavergleichskettek
{\vergleichskettek
{ m
}
{ = }{ n - \ell
}
{ \geq }{ 0
}
{    }{
}
{   }{
}
}
{}{}{}} {} {}
suatu
\definitionsverweis {pemetaan yang terdiferensialkan secara kontinu}{}{,}}
\faktvoraussetzung {yang di setiap titik
\definitionsverweis {serat}{}{}
$M$ di atas

\mavergleichskette
{\vergleichskette
{ 0
}
{ \in }{ \R^\ell
}
{  }{
}
{    }{
}
{   }{
}
}
{}{}{}
\definitionsverweis {reguler}{}{.}}
\faktfolgerung {Maka pemetaan

\maabbeledisp {} {\bigwedge^{m} T_PM } { \R
} {v_1 \wedge \ldots \wedge v_m } { \det  (
\operatorname{Grad} \, \varphi_1  (  P ) , \ldots ,
\operatorname{Grad} \, \varphi_\ell ( P ),  v_1 , \ldots , v_m)
} {,}
di setiap titik

\mavergleichskette
{\vergleichskette
{ P
}
{ \in }{ M
}
{  }{
}
{    }{
}
{   }{
}
}
{}{}{}
merupakan suatu isomorfisme, sehingga memberikan suatu bentuk volume kontinu yang tidak pernah nol pada $M$.}
\faktzusatz {}
\faktzusatz {}

}
{

Dalam pemetaan ini, vektor

\mavergleichskette
{\vergleichskette
{ v_i
}
{ \in }{ T_PM
}
{ \subseteq }{ T_P \R^n
}
{ =   }{ \R^n
}
{   }{
}
}
{}{}{}
dan gradien dipandang sebagai vektor dalam $\R^n$. Menurut
Soal 12.13
dan
Korolari 57.6 (Aljabar Linear (Osnabrück 2024-2025)),
pemetaan tersebut terdefinisi dengan baik dan linear. Domain pemetaan itu berdimensi satu, sama seperti kodomainnya. Misalkan
\mathl{v_1 , \ldots , v_m}{} suatu
\definitionsverweis {basis}{}{}
dari

\mavergleichskettedisp
{\vergleichskette
{ T_PM
}
{ =} { \ker  \left(D\varphi\right)_{P}
}
{ =} { (
\operatorname{Grad} \, \varphi_1  (  P ) , \ldots ,
\operatorname{Grad} \, \varphi_\ell ( P ) ) ^{ { \perp  } }
}
{ } {
}
{ } {
}
}
{}{}{.}
Karena pemetaan tersebut reguler, gradien itu
\definitionsverweis {bebas linear}{}{}
satu sama lain; akibat relasi ortogonalitas, gradien tersebut terlebih lagi bebas linear dari $v_i$. Dengan demikian, keseluruhannya merupakan suatu basis dari $\R^n$, sehingga menurut
Teorema 16.11 (Aljabar Linear (Osnabrück 2024-2025))
determinannya $\neq 0$. Ketergantungan bentuk volume ini pada titik pangkal $P$ bersifat kontinu, sebagaimana mengikuti dari kontinuitas gradien, kontinuitas determinan, dan
teorema pemetaan implisit.

}

Walaupun dalam situasi penting ini hasil di atas menjamin keberadaan suatu ukuran positif, hasil itu belum memberikan bentuk volume kanonik yang akan kita perkenalkan melalui metrik Riemann pada kuliah berikutnya. Untuk hubungan antara kedua konsep ini, lihat
Soal 16.10
dan
Soal 16.11.



\inputfaktbeweis
{Differenzierbare reguläre Funktion/R^n/Volumenform über Gradienten/Als Differentialform/Fakt}
{Korollar}
{}
{

\faktsituation {Misalkan

\mavergleichskette
{\vergleichskette
{ G
}
{ \subseteq }{ \R^n
}
{  }{
}
{    }{
}
{   }{
}
}
{}{}{}
terbuka dan misalkan
\maabbdisp {\varphi} { G } {\R
} {}
suatu
\definitionsverweis {pemetaan yang terdiferensialkan secara kontinu}{}{,}}
\faktvoraussetzung {yang di setiap titik
\definitionsverweis {serat}{}{}
$M$ di atas

\mavergleichskette
{\vergleichskette
{ 0
}
{ \in }{\R
}
{  }{
}
{    }{
}
{   }{
}
}
{}{}{}
\definitionsverweis {reguler}{}{.}}
\faktfolgerung {Maka
\definitionsverweis {bentuk volume}{}{}
dari
Korolari 15.6
sama dengan

\mathdisp {\sum_{i = 1}^n (-1)^{i+1} { \frac{   \partial  \varphi  }{  \partial   x_i   } } dx_1 \wedge \ldots \wedge dx_{i-1} \wedge dx_{i+1} \wedge \ldots \wedge dx_n} { , }

di mana
\mathl{x_1 , \ldots , x_n}{} menyatakan koordinat pada $\R^n$ dan
\mathl{dx_i}{} adalah $1$-bentuk yang bersesuaian pada $M$.}
\faktzusatz {}
\faktzusatz {}

}
{

Berdasarkan
Korolari 15.6,
bentuk volume tersebut diberikan dengan memasangkan, di suatu titik

\mavergleichskette
{\vergleichskette
{ P
}
{ \in }{ M
}
{  }{
}
{    }{
}
{   }{
}
}
{}{}{}
dan pada suatu tupel beranggotakan $n-1$

\mavergleichskette
{\vergleichskette
{ v_1 , \ldots , v_{n-1}
}
{ \in }{ T_PM
}
{ \subset }{ \R^n
}
{    }{
}
{   }{
}
}
{}{}{}
dengan

\mavergleichskette
{\vergleichskette
{ v_s
}
{ = }{ \begin{pmatrix} v_{1s} \\\vdots\\ v_{ns}   \end{pmatrix}
}
{  }{
}
{    }{
}
{   }{
}
}
{}{}{,}
bilangan

\mavergleichskettedisphandlinks
{\vergleichskettedisphandlinks
{ \det  \left(
\operatorname{Grad} \, \varphi ( P )   , \,  v_1  , \,   \ldots , \,  v_{n-1}                \right)
}
{ =} { \det  \begin{pmatrix}  { \frac{   \partial  \varphi  }{  \partial   x_1   } } ( P)  & v_{11}   &  \ldots & v_{1 \, n-1 } \\  { \frac{   \partial  \varphi  }{  \partial   x_2   } } ( P)  & v_{21}   &  \ldots & v_{2 \, n-1 } \\ \vdots & \vdots & \vdots & \vdots \\  { \frac{   \partial  \varphi  }{  \partial   x_n   } } ( P)  & v_{n1} &  \ldots &  v_{n\, n-1}  \end{pmatrix}
}
{ } {
}
{ } {
}
{ } {}
}
{}{}{.}
Kita harus menunjukkan bahwa bentuk yang dinyatakan dalam pernyataan memang sama dengan bentuk ini. Jika bentuk tersebut dievaluasi di titik $P$ pada vektor itu, maka menurut
Teorema 58.4 (Aljabar Linear (Osnabrück 2024-2025))
diperoleh

\mavergleichskettealignhandlinks
{\vergleichskettealignhandlinks
{ \sum_{i = 1}^n (-1)^{i+1} { \frac{   \partial  \varphi  }{  \partial   x_i   } } (P) dx_1 \wedge \ldots \wedge dx_{i-1} \wedge dx_{i+1} \wedge \ldots \wedge dx_n  (v_1 , \ldots , v_{n-1})
}
{ } {
}
{ } {
}
{ } {
}
{ } {}
}
{}{}{}
\mavergleichskettealignhandlinks
{\vergleichskettealignhandlinks
{ }
{ =} { \sum_{i = 1}^n (-1)^{i+1} { \frac{   \partial  \varphi  }{  \partial   x_i   } } (P) \det  \left(  dx_r  (v_{s} )  \right)_{\substack{1\leq r\leq n,\, r\neq i\\1\leq s\leq n-1}}
}
{ =} { \sum_{i = 1}^n (-1)^{i+1} { \frac{   \partial  \varphi  }{  \partial   x_i   } } (P) \det  \left(  v_{r s }  \right)_{\substack{1\leq r\leq n,\, r\neq i\\1\leq s\leq n-1}}
}
{ } {
}
{ } {
}
}
{}
{}{.}
Ini tepat merupakan ekspansi determinan di atas menurut kolom pertama.

}

\inputbemerkung
{}
{

Dalam situasi
Korolari 15.6,
kita tidak hanya memperoleh suatu bentuk volume yang tidak pernah nol, tetapi juga suatu orientasi pada setiap ruang singgung, bahkan suatu manifold berorientasi. Orientasi pada
\mathl{T_PM}{} didefinisikan dengan menetapkan bahwa suatu basis
\mathl{v_1 , \ldots , v_m}{} merepresentasikan orientasi tersebut jika basis yang diperluas
\mathl{\operatorname{Grad} \, \varphi_1  (  P ) , \ldots ,
\operatorname{Grad} \, \varphi_\ell ( P ), v_1 , \ldots , v_m}{} merepresentasikan orientasi standar $\R^n$. Penetapan ini bergantung pada $\varphi$. Sebagai contoh, jika $\varphi$ diganti dengan $- \varphi$, maka untuk $\ell$ ganjil orientasinya berubah.

}

\inputbeispiel{}
{

Kita memandang
\definitionsverweis {permukaan bola berdimensi}{}{} $2$, yaitu
$S^2$, sebagai
\definitionsverweis {serat}{}{}
di atas $0$ dari
\definitionsverweis {pemetaan diferensiabel}{}{}
\maabbeledisp {\varphi} {\R^3} {\R
} {(x,y,z)} { x^2+y^2+z^2-1
} {.}
Kita dapat menerapkan
Korolari 15.6
dan melalui
\maabbeledisp {} {\bigwedge^2 T_P S^2} {\R
} {v_1 \wedge v_2 } { \det  (
\operatorname{Grad} \, \varphi(P),v_1,v_2)
} {,}
\zusatzklammer {dengan vektor singgung
\mathkor {} {v_1} {dan} {v_2} {}
karena

\mavergleichskette
{\vergleichskette
{ T_PS^2
}
{ \subseteq }{ T_P\R^3
}
{ = }{ \R^3
}
{    }{
}
{   }{
}
}
{}{}{}
dapat langsung dipandang sebagai vektor dalam $\R^3$} {.} {,}
kita memperoleh suatu
\definitionsverweis {bentuk luas}{}{}
kontinu yang tidak pernah nol, yaitu $\omega$. Bentuk ini menghasilkan suatu
\definitionsverweis {bentuk luas positif}{}{}
dan suatu
\definitionsverweis {orientasi}{}{}
pada $S^2$. Dua vektor singgung yang bebas linear,
\mathkor {} {v_1} {dan} {v_2} {,}
merepresentasikan orientasi tersebut jika

\mavergleichskette
{\vergleichskette
{ \omega(v_1,v_2)
}
{ > }{ 0
}
{  }{
}
{    }{
}
{   }{
}
}
{}{}{,}
dan ini terjadi tepat ketika ketiga vektor
\mathl{\operatorname{Grad} \, \varphi(P),\, v_1,\, v_2}{} merepresentasikan
\definitionsverweis {orientasi standar}{}{}
$\R^3$. Menurut
Korolari 15.7,
bentuk luas ini juga dapat ditulis sebagai dua kali

\mathdisp {x dy \wedge dz -y dx \wedge dz + z dx \wedge dy} {  }

}

\inputbeispiel{}
{

Misalkan

\mavergleichskette
{\vergleichskette
{ V
}
{ \subseteq }{ \R^n
}
{  }{
}
{    }{
}
{   }{
}
}
{}{}{}
suatu himpunan bagian terbuka,
\maabbdisp {h} { V } {\R
} {}
suatu
\definitionsverweis {fungsi yang terdiferensialkan secara kontinu}{}{}
dan

\mavergleichskette
{\vergleichskette
{M
}
{ \subset }{ \R^n \times \R
}
{ = }{\R^{n+1}
}
{    }{
}
{   }{
}
}
{}{}{}
grafik yang bersesuaian. Kita memandang grafik ini sebagai suatu
\definitionsverweis {manifold diferensiabel}{}{}
berdimensi $n$ yang difeomorfik dengan $V$ melalui
\mathl{(x_1 , \ldots , x_n) \mapsto (x_1 , \ldots , x_n, h (x_1 , \ldots , x_n) )}{}
\definitionsverweis {difeomorfik}{}{.}
Manifold ini sekaligus merupakan serat di atas $0$ dari pemetaan
\maabbeledisp {\varphi} {V \times \R} {\R
} {(x_1 , \ldots , x_n,y)} { h( x_1 , \ldots , x_n)-y
} {.}
Gradien pemetaan ini adalah

\mathdisp {\left( \frac{ \partial h } { \partial x_1 } (P)  , \,   \ldots  , \,  \frac{ \partial h } { \partial x_n } (P) , \,   -1                \right)} {  }

Oleh karena itu, menurut
Korolari 15.6,
pemasangan

\mathdisp {(v_1 , \ldots , v_n) \mapsto \det  \begin{pmatrix}  \frac{ \partial h } { \partial x_1 }(P)  & v_{11} &  \ldots & v_{1n} \\  \vdots & \vdots &  \ddots & \vdots \\  \frac{ \partial h } { \partial x_n } (P) & v_{n1} &  \ldots & v_{nn} \\  -1 & v_{n+1,1} &  \ldots & v_{n+1,n}  \end{pmatrix}} {  }

memberikan suatu $n$-bentuk kontinu $\omega$ yang tidak pernah nol pada $M$. Jika bentuk ini ditarik balik ke $V$ melalui peta
\zusatzklammer {satu-satunya} {} {}
yang dideskripsikan di atas, maka

\mavergleichskette
{\vergleichskette
{ \alpha_*\omega
}
{ = }{ f dx_1 \wedge \ldots \wedge dx_n
}
{  }{
}
{    }{
}
{   }{
}
}
{}{}{,}
di mana
\mathl{f(Q)}{} diperoleh sebagai nilai bentuk $\omega$ di titik

\mavergleichskette
{\vergleichskette
{ \alpha^{-1}(Q)
}
{ \in }{ M
}
{  }{
}
{    }{
}
{   }{
}
}
{}{}{}
pada vektor
\mathl{T_Q (\alpha^{-1})(e_1) , \ldots , T_Q (\alpha^{-1})(e_n)}{.}
Karena

\mavergleichskette
{\vergleichskette
{ T_Q (\alpha^{-1})(e_i)
}
{ = }{ ( e_i , \frac{ \partial h } { \partial x_i } (Q))
}
{  }{
}
{    }{
}
{   }{
}
}
{}{}{,}
nilai tersebut adalah

\mavergleichskettealignhandlinks
{\vergleichskettealignhandlinks
{ f(Q)
}
{ =} { \det  \begin{pmatrix}  \frac{ \partial h } { \partial x_1 }(Q)   &   1  &  0  &  \ldots &  0  \\  \frac{ \partial h } { \partial x_2 }(Q)  &  0  &  1  &   \ldots  &  0   \\  \vdots & \vdots  & \vdots &  \ddots  & \vdots \\   \frac{ \partial h } { \partial x_n } (Q) &  0  &  0 &   \ldots  &  1  \\  -1 &  \frac{ \partial h } { \partial x_1 }(Q)  &  \frac{ \partial h } { \partial x_2 }(Q) &  \ldots &  \frac{ \partial h } { \partial x_n }(Q)   \end{pmatrix}
}
{ =} { \pm \left(  \left( \frac{ \partial h } { \partial x_1 } (Q)   \right)^2 + \cdots + \left(  \frac{ \partial h } { \partial x_n } (Q)   \right)^2 +1  \right)
}
{ } {
}
{ } {
}
}
{}{}{,}
dengan tanda yang bergantung pada $n$.

}

\zwischenueberschrift{Integrasi sepanjang suatu pemetaan diferensiabel}

Pada suatu manifold berdimensi $n$, yaitu $M$, hanya $n$-bentuk pada $M$ yang secara wajar dapat diintegralkan. Namun, kita juga ingin dapat mengintegralkan $k$-bentuk
\zusatzklammer {
\mavergleichskettek
{\vergleichskettek
{ 1
}
{ \leq }{k
}
{ \leq }{n
}
{    }{
}
{   }{
}
}
{}{}{}} {} {}
atas subobjek berdimensi $k$ tertentu. Konsep yang sesuai ialah integrasi sepanjang suatu pemetaan diferensiabel
\maabbdisp {\varphi} {L} {M
} {}
dari suatu manifold berdimensi $k$, yaitu $L$. Pada tahap ini, atas $L$ kita mengintegralkan
\definitionsverweis {bentuk diferensial tarik balik}{}{}
\mathl{\varphi^* \omega}{,} yaitu bentuk tarik balik oleh $\varphi$ dari suatu bentuk

\mavergleichskette
{\vergleichskette
{ \omega
}
{ \in }{
{ \mathcal E }^{ k } ( M  )
}
{  }{
}
{    }{
}
{   }{
}
}
{}{}{.}
Pada $L$, dimensi manifold dan derajat bentuk tersebut bersesuaian. Kasus khusus yang penting ialah $1$-bentuk dan kurva diferensiabel
\maabbdisp {\gamma} {I} {M
} {,}
integral yang dihasilkan disebut \stichwort {integral lintasan} {.}

\inputdefinition
{}
{

Misalkan $M$ suatu
\definitionsverweis {manifold diferensiabel}{}{}
dan

\mavergleichskette
{\vergleichskette
{ \omega
}
{ \in }{
{ \mathcal E }^{  1 } (  M   )
}
{  }{
}
{    }{
}
{   }{
}
}
{}{}{}
suatu
$1$-\definitionsverweis {bentuk diferensial}{}{.}
Misalkan
\maabbdisp {\gamma} { [a,b] } { M
} {}
suatu
\definitionsverweis {kurva yang terdiferensialkan secara kontinu}{}{.}
Jika integran di bawah terintegralkan pada interval tersebut, maka

\mavergleichskettedisp
{\vergleichskette
{ \int_\gamma \omega
}
{ \defeq} {
\int_{  [a,b]  }   \gamma^* \omega
}
{ =} { \int_{  a  }^{  b  } \omega ( \gamma(t); \gamma'(t)) \, d  t
}
{ } {
}
{ } {
}
}
{}{}{}
disebut \definitionswort {integral lintasan}{} dari $\omega$ sepanjang $\gamma$.

}
Di sini

\mavergleichskette
{\vergleichskette
{  \gamma'(t)
}
{ = }{ (T_t \gamma)(1)
}
{ \in }{ T_{\gamma (t)} M
}
{    }{
}
{   }{
}
}
{}{}{.}

\inputbemerkung
{}
{

Dalam konteks fisika, suatu
$1$-\definitionsverweis {bentuk diferensial}{}{}
bernilai riil
\zusatzklammer {atau versi dualnya, yaitu suatu medan vektor} {} {}
mendeskripsikan suatu gaya; maka
\definitionsverweis {integral lintasan}{}{}
adalah \stichwort {usaha} {} atau \stichwort {energi} {,} yang dibutuhkan atau dilepaskan ketika sebuah partikel bergerak sepanjang lintasan tersebut.

}

Kita akan sering meninjau bentuk diferensial pada suatu submanifold tertutup

\mavergleichskette
{\vergleichskette
{M
}
{ \subseteq }{G
}
{  }{
}
{    }{
}
{   }{
}
}
{}{}{,}
dengan $G$ terbuka dalam $\R^n$, yang bahkan terdefinisi pada $G$ dan karena itu berbentuk

\mavergleichskette
{\vergleichskette
{ \omega
}
{ = }{ \sum_{i = 1}^n g_i d x_i
}
{  }{
}
{    }{
}
{   }{
}
}
{}{}{,}
di mana $x_i$ adalah koordinat $\R^n$ dan $g_i$ merupakan fungsi yang terdefinisi pada $G$. Untuk suatu lintasan dalam $M$, menurut
Soal 15.9,
tidak ada perbedaan apakah integral lintasannya ditinjau terhadap
\mathkor {} {G} {dan} {\omega} {}
atau terhadap
\mathkor {} {M} {dan bentuk diferensial yang dibatasi} {\omega \vert_M} {.}

\inputbemerkung
{}
{

Suatu
\definitionsverweis {integral lintasan}{}{}
dihitung sebagai berikut. Misalkan $\omega$ suatu $1$-bentuk pada suatu himpunan terbuka

\mavergleichskette
{\vergleichskette
{ G
}
{ \subseteq }{ \R^n
}
{  }{
}
{    }{
}
{   }{
}
}
{}{}{}
yang dideskripsikan oleh

\mavergleichskettedisp
{\vergleichskette
{ \omega
}
{ =} { g_1dx_1 + \cdots + g_ndx_n
}
{ } {
}
{ } {
}
{ } {
}
}
{}{}{,}
di mana
\maabb {g_j} { G } {\R
} {}
merupakan fungsi terukur. Misalkan diberikan suatu
\definitionsverweis {kurva yang terdiferensialkan secara kontinu}{}{}
\maabb {\gamma} { [a,b] } { G
} {}
dengan
\definitionsverweis {fungsi komponen}{}{}
$\gamma_j$
\zusatzklammer {yang masing-masing terdiferensialkan secara kontinu} {} {.}
Menurut Lema 37.4 (Analisis (Osnabrück 2021-2023)),
\definitionsverweis {turunan}{}{}
di suatu titik

\mavergleichskette
{\vergleichskette
{ t
}
{ \in }{ [a,b]
}
{  }{
}
{    }{
}
{   }{
}
}
{}{}{}
dideskripsikan oleh vektor
\mathl{\left(  \gamma_1'(t)  , \,   \ldots  , \,   \gamma_n'(t)                 \right)}{.}
Di titik $t$ dan pada arah $1$, nilai
\definitionsverweis {bentuk diferensial tarik balik}{}{}
\mathl{\gamma^* \omega}{} adalah

\mavergleichskettealignhandlinks
{\vergleichskettealignhandlinks
{ \omega(\gamma(t); \gamma'(t))
}
{ =} { \left(  g_1(\gamma(t))dx_1 + \cdots + g_n(\gamma(t))dx_n  \right) \begin{pmatrix}  \gamma_1'(t)  \\\vdots \\  \gamma_n'(t)   \end{pmatrix}
}
{ =} { g_1(\gamma(t)) \gamma_1'(t) + \cdots + g_n(\gamma(t)) \gamma_n'(t)
}
{ } {
}
{ } {
}
}
{}
{}{.}
Pada ekspresi di tengah, suatu
\definitionsverweis {bentuk linear}{}{}
diterapkan pada sebuah vektor. Jadi, dalam
\mathl{g_j(x_1 , \ldots , x_n)}{}
\mathkor {} {x_i} {diganti oleh} {\gamma_i(t)} {}
dan
\mathkor {} {dx_i} {diganti oleh} {\gamma_i'(t) dt} {.}
Hasil keseluruhannya adalah suatu $1$-bentuk terukur pada
\mathl{[a,b]}{} atau, secara ekuivalen, suatu fungsi terukur dari
\mathl{[a,b]}{} ke $\R$. Integral lintasan ada jika fungsi koefisien pada ruas kanan terintegralkan pada interval tersebut; keterukuran saja tidak cukup.

}

\inputbeispiel{}
{

Kita meninjau bentuk diferensial

\mavergleichskettedisp
{\vergleichskette
{ \omega
}
{ =} { \left(  xy+z^2  \right) dx+zdy+x^3dz
}
{ } {
}
{ } {
}
{ } {
}
}
{}{}{}
pada $\R^3$ dan
\definitionsverweis {lintasan afin-linear}{}{}
\maabbeledisp {\gamma} { [0,1] } {\R^3
} { t } {(1,2,0) +t(3,0,2) = (1+3t,2,2t)
} {.}
Melalui $\gamma$,
\definitionsverweis {bentuk diferensial tarik balik}{}{}
$\gamma^* \omega$ adalah

\mavergleichskettealignhandlinks
{\vergleichskettealignhandlinks
{ \left(  \left(  (1+3t)2 + (2t)^2  \right) dx + 2tdy + (1+3t)^3 dz  \right) \begin{pmatrix}  3  \\ 0 \\  2   \end{pmatrix}
}
{ =} { \left(  3  ((1+3t)2  + (2t)^2) + 2(1+3t)^3  \right) dt
}
{ =} { \left(  12t^2+18t+6+54t^3+54t^2+18t+2  \right) dt
}
{ =} { \left(  54t^3+66t^2+36t+8  \right) dt
}
{ } {
}
}
{}
{}{.}
Untuk integral pada interval satuan, kita memperoleh

\mavergleichskettedisphandlinks
{\vergleichskettedisphandlinks
{ \int_{  0  }^{  1  } 54t^3+66t^2+36t+8 \, d t
}
{ =} { \left(  { \frac{   27 }{  2 } } t^4 + 22 t^3 + 18t^2 +8t \right) \vert_{  0  }^{  1  }
}
{ =} { { \frac{   123 }{  2  } }
}
{ } {
}
{ } {
}
}
{}{}{.}

}

\chapter{Lembar Kerja 15}
\setcounter{section}{15}

  \zwischenueberschrift{Soal latihan}

\inputaufgabegibtloesung
{}
{

Misalkan $M$ sebuah manifold diferensiabel kompak dengan suatu bentuk volume positif yang kontinu $\omega$. Tunjukkan bahwa

\mavergleichskettedisp
{\vergleichskette
{ \int_M \omega
}
{ <} { \infty
}
{ } {
}
{ } {
}
{ } {
}
}
{}{}{}
berlaku.

}
{} {}

\inputaufgabe
{}
{

Tunjukkan bahwa ukuran volume yang diperkenalkan dalam
Definisi 15.3
untuk suatu
\definitionsverweis {bentuk volume positif}{}{}
\definitionsverweis {kontinu}{}{}
pada sebuah
\definitionsverweis {manifold diferensiabel}{}{}
merupakan suatu
\definitionsverweis {ukuran}{}{}
$\sigma$-hingga.

}
{} {}

\inputaufgabe
{}
{

Misalkan

\mavergleichskettedisp
{\vergleichskette
{ \omega
}
{ =} { dx_1 \wedge \ldots \wedge dx_n
}
{ =} { e_1^* \wedge \ldots \wedge e_n^*
}
{ } {
}
{ } {
}
}
{}{}{}
merupakan bentuk volume standar pada $\R^n$. Tunjukkan bahwa untuk setiap
\definitionsverweis {himpunan bagian terukur}{}{}

\mathl{T \subseteq \R^n}{} berlaku kesamaan

\mavergleichskettedisp
{\vergleichskette
{
\int_{ T  }   \omega
}
{ =} { \int_{  T  }     \,  d  \lambda^n
}
{ =} { \lambda^n(T)
}
{ } {
}
{ } {
}
}
{}{}{.}

}
{} {}

\inputaufgabe
{}
{

Misalkan $M$ sebuah
\definitionsverweis {manifold diferensiabel}{}{}
dengan suatu
\definitionsverweis {bentuk volume positif}{}{}
$\omega$. Misalkan
\mathl{T \subseteq M}{}
\definitionsverweis {terukur}{}{}
dan
\mathl{N \subseteq M}{} suatu
\definitionsverweis {himpunan nol}{}{.}
Tunjukkan bahwa

\mavergleichskettedisp
{\vergleichskette
{
\int_{ T  }   \omega
}
{ =} {
\int_{ T \setminus N  }   \omega
}
{ } {
}
{ } {
}
{ } {
}
}
{}{}{}
berlaku.

}
{} {}

\inputaufgabe
{}
{

Misalkan $M$ sebuah
\definitionsverweis {manifold diferensiabel}{}{}
\definitionsverweis {berdimensi}{}{} $n$
dengan suatu
\definitionsverweis {basis topologi terhitung}{}{}
dan misalkan
\mathkor {} {\omega_1} {dan} {\omega_2} {}
merupakan
\definitionsverweis {bentuk volume positif}{}{}
pada $M$. Tunjukkan bahwa untuk setiap
\definitionsverweis {himpunan bagian terukur}{}{}

\mathl{T \subseteq M}{} dan
\mathl{a,b \in \R_+}{} berlaku hubungan

\mavergleichskettedisp
{\vergleichskette
{
\int_{ T  }  (a \omega_1 + b \omega_2)
}
{ =} { a
\int_{  T  }   \omega_1+ b
\int_{  T  }   \omega_2
}
{ } {
}
{ } {
}
{ } {
}
}
{}{}{.}

}
{} {}

\inputaufgabe
{}
{

Misalkan $M$ sebuah
\definitionsverweis {manifold diferensiabel}{}{}
dengan
\definitionsverweis {basis topologi terhitung}{}{.}
Jelaskan bagaimana mendefinisikan \anfuehrung{himpunan nol}{} pada $M$ dengan mengacu pada peta, tanpa diberikan suatu
\definitionsverweis {ukuran}{}{.}
Tunjukkan pula bahwa jika suatu
\definitionsverweis {bentuk volume positif}{}{}
diberikan, maka himpunan-himpunan nol tersebut juga merupakan
\definitionsverweis {himpunan nol}{}{}
dalam pengertian teori ukuran.

}
{} {}

\inputaufgabe
{}
{

Nyatakan lingkaran satuan
\mathl{S^1 \subset \R^2}{} sebagai serat suatu pemetaan
\maabbdisp {} {\R^2} {\R
} {}
sedemikian sehingga bentuk volume $\omega$ yang diberikan menurut
Korolari 15.6
bersifat positif. Hitunglah
\mathl{\int_{S^1} \omega}{.}

}
{} {}

\inputaufgabe
{}
{

Nyatakan permukaan bola satuan
\mathl{S^2 \subset \R^3}{} sebagai serat suatu pemetaan
\maabbdisp {} {\R^3} {\R
} {}
sedemikian sehingga bentuk volume $\omega$ yang diberikan menurut
Korolari 15.6
bersifat positif. Hitunglah
\mathl{\int_{S^2} \omega}{.}

}
{} {}

\inputaufgabe
{}
{

Misalkan
\mathkor {} {L} {dan} {M} {}
merupakan dua
\definitionsverweis {manifold diferensiabel}{}{}
dan
\maabbdisp {\varphi} { L } { M
} {}
suatu
\definitionsverweis {pemetaan diferensiabel}{}{.}
Misalkan

\mavergleichskette
{\vergleichskette
{ \omega
}
{ \in }{
{ \mathcal E }^{  1 } (  M   )
}
{  }{
}
{    }{
}
{   }{
}
}
{}{}{}
sebuah bentuk diferensial terukur dengan
\definitionsverweis {bentuk diferensial tarik balik}{}{}

\mavergleichskette
{\vergleichskette
{ \varphi^*\omega
}
{ \in }{
{ \mathcal E }^{  1 } (  L   )
}
{  }{
}
{    }{
}
{   }{
}
}
{}{}{}
dan misalkan
\maabbdisp {\gamma} { I } { L
} {}
suatu
\definitionsverweis {kurva yang terdiferensialkan secara kontinu}{}{}
\zusatzklammer {$I$ suatu
\definitionsverweis {interval riil}{}{}} {} {.}
Tunjukkan bahwa untuk
\definitionsverweis {integral lintasan}{}{}
berlaku kesamaan

\mavergleichskettedisp
{\vergleichskette
{ \int_\gamma \varphi^* \omega
}
{ =} { \int_{\varphi \circ \gamma} \omega
}
{ } {
}
{ } {
}
{ } {
}
}
{}{}{.}

}
{} {}

\inputaufgabe
{}
{

Diberikan
\maabbeledisp {\gamma} { [0,2\pi] } {\R^2
} { t } {( \cos  t , \sin  t )
} {,}
hitunglah
\definitionsverweis {integral lintasan}{}{}
sepanjang lintasan ini untuk
\definitionsverweis {bentuk diferensial}{}{}
berikut:

a)
\mathl{xdx +ydy}{,}

b)
\mathl{xdx -ydy}{,}

c)
\mathl{ydx +xdy}{,}

d)
\mathl{ydx -xdy}{.}

}
{} {}

\inputaufgabegibtloesung
{}
{

Hitunglah integral lintasan
\mathl{\int_\gamma \omega}{} untuk
\maabbeledisp {\gamma} { [-1,0] } {\R^3
} { t } {(-t^2,t^3-1,t+2)
} {,}
dengan bentuk diferensial berderajat $1$

\mavergleichskettedisp
{\vergleichskette
{ \omega
}
{ =} { x^3dx -yzdy +xz^2 dz
}
{ } {
}
{ } {
}
{ } {
}
}
{}{}{}
pada $\R^3$.

}
{} {}

\inputaufgabegibtloesung
{}
{

Kita meninjau dua pemetaan diferensiabel
\maabbeledisp {\gamma} { [1,2] } {\R^2
} { t } {(t,t^{-1})
} {,}
dan
\maabbeledisp {\varphi} {\R^2} {\R^3
} {(u,v)} {(u^2,uv,-u+v^2)
} {,}
serta bentuk diferensial

\mavergleichskettedisp
{\vergleichskette
{ \omega
}
{ =} { xdx-zdy+dz
}
{ } {
}
{ } {
}
{ } {
}
}
{}{}{}
pada $\R^3$.

a) Hitunglah bentuk diferensial tarik balik
\mathl{\varphi^* \omega}{} pada $\R^2$.

b) Hitunglah integral lintasan dari bentuk diferensial
\mathl{\varphi^* \omega}{} sepanjang lintasan $\gamma$.

c) Hitunglah
\zusatzklammer {tanpa mengacu pada b)} {} {}
integral lintasan dari bentuk diferensial
\mathl{\omega}{} sepanjang lintasan $\varphi \circ \gamma$.

}
{} {}

\inputaufgabegibtloesung
{}
{

Kita meninjau dua pemetaan diferensiabel
\maabbeledisp {\gamma} { [1,c] } {\R^2
} { t } {(t,t^{3})
} {,}
\zusatzklammer {dengan \mathlk{c \geq 1}{}} {} {}
dan
\maabbeledisp {\varphi} {\R_+ \times \R_+ } {\R^3
} {(u,v)} {(u^3,u^2+v^2,u^{-1}v^{-1} )
} {,}
serta bentuk diferensial

\mavergleichskettedisp
{\vergleichskette
{ \omega
}
{ =} { (x-y)dx-z^2dy+dz
}
{ } {
}
{ } {
}
{ } {
}
}
{}{}{}
pada $\R^3$.

a) Hitunglah bentuk diferensial tarik balik
\mathl{\varphi^* \omega}{} pada $\R_+ \times \R_+$.

b) Hitunglah integral lintasan dari bentuk diferensial
\mathl{\varphi^* \omega}{} sepanjang lintasan $\gamma$ sebagai fungsi $c$.

}
{} {}

  \zwischenueberschrift{Soal untuk dikumpulkan}

\inputaufgabe
{4}
{

Misalkan $M$ sebuah
\definitionsverweis {manifold diferensiabel}{}{}
\definitionsverweis {berdimensi}{}{} $n$
dengan
\definitionsverweis {basis topologi terhitung}{}{.}
Misalkan $\omega$ suatu
\definitionsverweis {bentuk volume positif}{}{} pada $M$ dan misalkan $\mu$ suatu
\definitionsverweis {ukuran}{}{} pada $M$ yang didefinisikan oleh bentuk volume ini. Tunjukkan bahwa setiap
\definitionsverweis {submanifold tertutup}{}{} berdimensi
\mathl{\leq n-1}{} merupakan suatu
\definitionsverweis {himpunan nol}{}{.}

}
{} {}

\inputaufgabe
{4}
{

Misalkan

\mavergleichskette
{\vergleichskette
{a,b,c,d,r,s
}
{ \geq }{ 1
}
{  }{
}
{    }{
}
{   }{
}
}
{}{}{}
merupakan
\definitionsverweis {bilangan asli}{}{.}
Kita meninjau
\definitionsverweis {kurva yang terdiferensialkan secara kontinu}{}{}
\maabbeledisp {} { [0,1] } {\R^2
} { t } {(t^r,t^s)
} {.}
Hitunglah
\definitionsverweis {integral lintasan}{}{}
sepanjang lintasan ini untuk
\definitionsverweis {bentuk diferensial}{}{}

\mavergleichskette
{\vergleichskette
{ \omega
}
{ = }{ x^ay^b dx +x^c y^d dy
}
{  }{
}
{    }{
}
{   }{
}
}
{}{}{.}

}
{} {}

\inputaufgabe
{5}
{

Diberikan
\maabbeledisp {\gamma} { [0,2\pi] } {\R^3
} { t } {( \cos  t , \sin  t, t )
} {,}
hitunglah
\definitionsverweis {integral lintasan}{}{}
sepanjang lintasan ini untuk
\definitionsverweis {bentuk diferensial}{}{}

\mavergleichskette
{\vergleichskette
{ \omega
}
{ = }{ (y-z^3) dx +x^2dy -xzdz
}
{  }{
}
{    }{
}
{   }{
}
}
{}{}{.}

}
{} {}

\section*{Solusi yang disediakan oleh sumber}
\addcontentsline{toc}{section}{Solusi yang disediakan oleh sumber}
\subsection*{Solusi Soal 15.1}
Untuk setiap titik

\mavergleichskette
{\vergleichskette
{ P
}
{ \in }{ M
}
{  }{
}
{    }{
}
{   }{
}
}
{}{}{}
terdapat suatu lingkungan koordinat terbuka

\mavergleichskette
{\vergleichskette
{ P
}
{ \in }{ U
}
{  }{
}
{    }{
}
{   }{
}
}
{}{}{}
dan suatu pemetaan koordinat
\maabbdisp {\alpha} { U } { V
} {}
dengan

\mavergleichskette
{\vergleichskette
{ V
}
{ \subseteq }{ \R^n
}
{  }{
}
{    }{
}
{   }{
}
}
{}{}{}
terbuka, sedemikian sehingga

\mavergleichskette
{\vergleichskette
{ \alpha^{-1 *} \omega
}
{ = }{ fdx_1 \wedge \ldots \wedge dx_n
}
{  }{
}
{    }{
}
{   }{
}
}
{}{}{}
dengan
\mathl{f}{} kontinu dan positif. Kita juga dapat menemukan suatu lingkungan terbuka

\mavergleichskette
{\vergleichskette
{ P
}
{ \in }{ U'
}
{ \subseteq }{ U
}
{    }{
}
{   }{
}
}
{}{}{,}
yang homeomorfik dengan suatu bola terbuka

\mavergleichskette
{\vergleichskette
{ U'
}
{ \cong }{ B'
}
{ \subseteq }{ V
}
{    }{
}
{   }{
}
}
{}{}{}
dan kita dapat pula mengandaikan bahwa penutupan bola tersebut seluruhnya termuat dalam $V$. Penutupan $B'$ bersifat tertutup dan terbatas; oleh karena itu, fungsi kontinu $f$ terbatas pada penutupan $B'$ dan dengan demikian juga pada $B'$. Akibatnya,
\mathl{\int_{U'} \omega}{} berhingga, dengan $U'$ sebagai suatu lingkungan terbuka dari $P$.

Himpunan-himpunan terbuka

\mavergleichskette
{\vergleichskette
{ U'
}
{ = }{ U'(P)
}
{  }{
}
{    }{
}
{   }{
}
}
{}{}{}
ini menutupi $M$. Berdasarkan kekompakan, terdapat suatu tutupan hingga

\mavergleichskettedisp
{\vergleichskette
{ M
}
{ =} { \bigcup_{i = 1} ^N U_i
}
{ } {
}
{ } {
}
{ } {
}
}
{}{}{}
dengan

\mavergleichskettedisp
{\vergleichskette
{ \int_{U_i } \omega
}
{ <} { \infty
}
{ } {
}
{ } {
}
{ } {
}
}
{}{}{.}
Karena sifat positifnya, dengan demikian berlaku

\mavergleichskettedisp
{\vergleichskette
{ \int_{M} \omega
}
{ \leq} { \sum_{i = 1}^N  \left(  \int_{U_i } \omega  \right)
}
{ <} { \infty
}
{ } {
}
{ } {
}
}
{}{}{.}
\subsection*{Solusi Soal 15.11}
Kita peroleh

\mavergleichskettealign
{\vergleichskettealign
{ \gamma^* \omega
}
{ =} { (-t^2)^3 d(-t^2) - (t^3-1) (t+2)d (t^3-1)  - t^2(t+2)^2d(t+2)
}
{ =} { 2 t^7 dt     - 3t^2 (t^4 +2t^3 -t -2) dt - t^2(  t^2  +4 t +4) dt
}
{ =} {  (2t^7 -3t^6 -6t^5   +3t^3 +6t^2 -t^4  -4t^3 -4t^2       )dt
}
{ =} { (2t^7 -3t^6 -6t^5   -t^4  -t^3 +2 t^2       )dt
}
}
{}
{}{.}
Oleh karena itu,

\mavergleichskettealign
{\vergleichskettealign
{ \int_\gamma \omega
}
{ =} { \int_{-1}^0    (2t^7 -3t^6 -6t^5   -t^4  -t^3 +2 t^2       )dt
}
{ =} { ({ \frac{   1 }{  4 } }  t^8 - { \frac{   3 }{  7 } }  t^7 - t^6- { \frac{   1 }{  5 } } t^5 - { \frac{   1 }{  4 } } t^4 + { \frac{   2 }{  3 } } t^3) \vert_{  -1 }^{  0 }
}
{ =} { -( { \frac{   1 }{  4 } }  + { \frac{   3 }{  7 } }  -1 + { \frac{   1 }{  5 } }  - { \frac{   1 }{  4 } }  - { \frac{   2 }{  3 } } )
}
{ =} { - { \frac{   3 }{  7 } }  +1 - { \frac{   1 }{  5 } } + { \frac{   2 }{  3 } }
}
}
{
\vergleichskettefortsetzungalign
{ =} { { \frac{   -45 +105-21+70 }{  105 } }
}
{ =} { { \frac{   109 }{  105 } }
}
{ } {}
{ } {}
}
{}{.}
\subsection*{Solusi Soal 15.12}
a) Bentuk diferensial tarik baliknya adalah

\mavergleichskettealign
{\vergleichskettealign
{\varphi^* ( xdx -zdy +dz )
}
{ =} { u^2 d(u^2)    -(-u+v^2) d(uv) +d (-u+v^2)
}
{ =} { 2u^3 du + (u-v^2) (vdu +udv) -du +2vdv
}
{ =} { (2u^3 +uv -v^3 -1) du    + (u^2 -uv^2 +2v )dv
}
{ } {
}
}
{}
{}{.}

b) Integral lintasannya adalah

\mavergleichskettealign
{\vergleichskettealign
{ \int_1^2 ( 2t^3 +1 -t^{-3} -1 )dt  +(t^2 -t^{-1} +2t^{-1} ) d t^{-1}
}
{ =} { \int_1^2 ( 2t^3 -t^{-3} ) dt  +(t^2 + t^{-1} ) d t^{-1}
}
{ =} { \int_1^2 ( 2t^3 -t^{-3} ) dt  -(1 + t^{-3} ) dt
}
{ =} { \int_1^2 ( 2t^3 - 2t^{-3} -1 ) dt
}
{ =} { \left(  { \frac{   1 }{  2 } }  t^4  + t^{-2}  -t  \right) \vert_{  1 }^{  2 }
}
}
{
\vergleichskettefortsetzungalign
{ =} { 8 + { \frac{   1 }{  4 } }  -2 - { \frac{   1 }{  2 } } -1 +1
}
{ =} { 5 + { \frac{   3 }{  4 } }
}
{ } {}
{ } {}
}
{}{.}

c) Lintasan komposisinya adalah

\mavergleichskettedisp
{\vergleichskette
{ \psi (t)
}
{ =} { (t^2,1,-t +t^{-2} )
}
{ } {
}
{ } {
}
{ } {
}
}
{}{}{.}
Dengan demikian,

\mavergleichskettealign
{\vergleichskettealign
{ \int_1^2 \psi^* (xdx -zdy +dz )
}
{ =} { \int_1^2 t^2d(t^2) + d (-t + t^{-2} )
}
{ =} {  \int_1^2 (2 t^3-1 -2t^{-3} ) dt
}
{ =} { 5 + { \frac{   3 }{  4 } }
}
{ } {
}
}
{}
{}{.}
\subsection*{Solusi Soal 15.13}
a) Bentuk diferensial tarik baliknya adalah

\mavergleichskettealigndrucklinks
{\vergleichskettealigndrucklinks
{\varphi^* (  (x-y)dx-z^2dy+dz   )
}
{ =} { (u^3-u^2-v^2 ) d(u^3)  -u^{-2}v^{-2} d(u^2+v^2) +d (u^{-1}v^{-1})
}
{ =} { 3 (u^5-u^4-u^2v^2 ) du  -2u^{-1}v^{-2} du -2u^{-2}v^{-1} dv - u^{-2} v^{-1}du - u^{-1} v^{-2}dv
}
{ =} { (3 u^5-3u^4-3u^2v^2 -2u^{-1}v^{-2}- u^{-2} v^{-1} )  du + (  -2u^{-2}v^{-1} - u^{-1} v^{-2})dv
}
{ } {
}
}
{}{}{.}

b) Integral lintasannya adalah

\mavergleichskettealigndrucklinks
{\vergleichskettealigndrucklinks
{ \int_1^c (3 t^5-3t^4-3t^2t^6 -2t^{-1}t^{-6}- t^{-2} t^{-3} )dt  +(  -2t^{-2}t^{-3} - t^{-1} t^{-6}) d t^3
}
{ =} { \int_1^c   (3 t^5-3t^4-3t^8-2t^{-7} - t^{-5})dt + ( -2t^{-5} - t^{-7})3t^2 dt
}
{ =} { \int_1^c (3 t^5-3t^4-3t^8-2t^{-7} - t^{-5}   -6t^{-3} - 3 t^{-5}      ) dt
}
{ =} { \int_1^c (-3t^8 + 3 t^5-3t^4 -6t^{-3} - 4 t^{-5}  -2t^{-7}       )  dt
}
{ =} { (-{ \frac{   1 }{  3 } } t^9 + { \frac{   1 }{  2 } } t^{6}  - { \frac{   3 }{  5 } } t^5  +  3 t^{-2} +  t^{-4} + { \frac{   1 }{  3 } }  t^{-6}     ) \vert_{  1 }^{  c  }
}
}
{
\vergleichskettefortsetzungalign
{ =} { -{ \frac{   1 }{  3 } } c^9 + { \frac{   1 }{  2 } } c^{6}  - { \frac{   3 }{  5 } } c^5  +  3c^{-2} +   c^{-4} + { \frac{   1 }{  3 } }  c^{-6}
-(-{ \frac{   1 }{  3 } } + { \frac{   1 }{  2 } } - { \frac{   3 }{  5 } } +3 + 1 + { \frac{   1 }{  3 } }  )
}
{ =} { -{ \frac{   1 }{  3 } } c^9 + { \frac{   1 }{  2 } } c^{6}  - { \frac{   3 }{  5 } } c^5 + 3 c^{-2} +  c^{-4} + { \frac{   1 }{  3 } } c^{-6} - { \frac{   39 }{  10 } }  }
{ } {}
{ } {}
}{}{.}

\part{Unit 16}
\chapter[Kuliah 16]{Kuliah 16: Manifold Riemann}
\setcounter{section}{16}

\bild{ \begin{center}
\includegraphics[width=5.5cm]{ \bildeinlesungjpeg {Georg_Friedrich_Bernhard_Riemann} {jpeg} }
\end{center}
\bildtext {Georg Friedrich Bernhard Riemann (1826-1866)} }

\bildlizenz { Georg Friedrich Bernhard Riemann.jpeg } {} {Ævar Arnfjörð Bjarmason} {Commons} {PD} {http://www.sil.si.edu/digitalcollections/hst/scientific-identity/explore.htm}

  \zwischenueberschrift{Manifold Riemann}

Permukaan sebuah bola berjari-jari $r$ mempunyai luas
\mathl{4 \pi r^2}{.} Ini merupakan hasil klasik, tetapi bagaimana kita dapat merumuskan secara tepat luas suatu manifold dua dimensi seperti ini? Agar dapat menerapkan teori ukuran dan integrasi dari kuliah-kuliah sebelumnya, kita memerlukan suatu $2$-bentuk pada permukaan yang sealami mungkin. Melalui konsep metrik Riemann, kita akan menunjukkan bahwa permukaan yang tertanam dalam ruang Euklides tiga dimensi mempunyai ukuran permukaan alami yang dapat digunakan untuk menghitung luasnya.

\bild{ \begin{center}
\includegraphics[width=5.5cm]{ \bildeinlesungpng {Sphere_with_three_handles} {png} }
\end{center}
\bildtext {Permukaan hijau mewarisi hasil kali dalam dari ruang Euklides di sekelilingnya. Dengan demikian, luas permukaan tersebut dapat diukur secara bermakna.} }

\bildlizenz { Sphere with three handles.png } {} {Oleg Alexandrov} {Commons} {PD} {}

\inputdefinition
{}
{

Suatu
\definitionsverweis {manifold diferensiabel}{}{}
$M$ disebut \definitionswort {manifold Riemann}{,} jika pada setiap
\definitionsverweis {ruang singgung}{}{}

\mathbed {T_PM} {}
 {P \in M} {}
 {} {} {} {,}
diberikan suatu
\definitionsverweis {hasil kali dalam}{}{}

\mathl{\left\langle  - ,  -  \right\rangle_P}{} sedemikian sehingga untuk setiap peta
\maabbdisp {\alpha} { U } { V
} {}
dengan

\mavergleichskette
{\vergleichskette
{ V
}
{ \subseteq }{  \R^n
}
{  }{
}
{    }{
}
{   }{
}
}
{}{}{}
fungsi-fungsi
\zusatzklammer {untuk

\mavergleichskettek
{\vergleichskettek
{ 1
}
{ \leq }{ i,j
}
{ \leq }{ n
}
{    }{
}
{   }{
}
}
{}{}{}} {} {}
\maabbeledisp {g_{ij}} { V } {\R
} { Q } {g_{ij}(Q) = \left\langle  T(\alpha^{-1})(e_i)  ,  T(\alpha^{-1}) (e_j)   \right\rangle_{ \alpha^{-1}(Q) }
} {,}
\definitionsverweis {terdiferensialkan}{}{} kelas $C^1$
\zusatzfussnote {Banyak penulis mensyaratkan bahwa manifold Riemann dan fungsi-fungsi ini berkelas $C^\infty$} {.} {.}

}

Fungsi-fungsi
\mathl{g_{ij}}{} yang didefinisikan pada peta-peta disebut
\zusatzklammer {metrik atau Riemann} {} {}
\stichwort {fungsi fundamental} {.} Fungsi-fungsi tersebut dihimpun menjadi matriks
\mathl{\left(  g_{ij}  \right)_{1 \leq i,j \leq n}}{} yang juga disebut \stichwort {matriks fundamental metrik} {}
\zusatzklammer {atau \stichwort {matriks fundamental pertama} {} atau \stichwort {tensor fundamental metrik} {}} {} {.}
Pada setiap titik

\mavergleichskette
{\vergleichskette
{ Q
}
{ \in }{ V
}
{  }{
}
{    }{
}
{   }{
}
}
{}{}{}
matriks ini simetris dan definit positif. Determinannya juga penting, yaitu

\mavergleichskettedisp
{\vergleichskette
{ g
}
{ =} { \det  \left(  g_{ij}  \right)_{1 \leq i,j \leq n}
}
{ } {
}
{ } {
}
{ } {
}
}
{}{}{,}
yang juga terdiferensialkan secara kontinu dan, menurut
Korolari 39.5 (Aljabar Linear (Osnabrück 2024-2025)),
bernilai positif di mana-mana.

Contoh paling sederhana dari manifold Riemann ialah ruang Euklides
\mathl{\R^n}{} dengan hasil kali dalam standar pada setiap titik
\zusatzklammer {dan, secara lebih umum, sembarang ruang Euklides} {} {},
serta setiap himpunan bagian terbukanya. Yang lebih penting, setiap
\definitionsverweis {submanifold tertutup}{}{}
dari suatu manifold Riemann juga merupakan manifold Riemann. Hal ini menghasilkan banyak contoh tak trivial, misalnya permukaan di $\R^3$ seperti sfer atau torus.



\inputfaktbeweis
{Riemannsche Mannigfaltigkeit/Abgeschlossene Untermannigfaltigkeit/Riemannsch/Fakt}
{Teorema}
{}
{

\faktsituation {Misalkan $N$ suatu
\definitionsverweis {manifold Riemann}{}{}
dan

\mavergleichskette
{\vergleichskette
{ M
}
{ \subseteq }{ N
}
{  }{
}
{    }{
}
{   }{
}
}
{}{}{}
suatu
\definitionsverweis {submanifold tertutup}{}{.}}
\faktfolgerung {Maka $M$ juga merupakan manifold Riemann.}
\faktzusatz {}
\faktzusatz {}

}
{

Untuk setiap titik

\mavergleichskette
{\vergleichskette
{ P
}
{ \in }{ M
}
{  }{
}
{    }{
}
{   }{
}
}
{}{}{}
ruang

\mavergleichskette
{\vergleichskette
{ T_PM
}
{ \subseteq }{ T_PN
}
{  }{
}
{    }{
}
{   }{
}
}
{}{}{}
merupakan suatu
\definitionsverweis {subruang vektor}{}{}
menurut
Teorema 10.3.
Karena itu,
\definitionsverweis {hasil kali dalam}{}{}

\mathl{\left\langle  - ,  -  \right\rangle_P}{} pada
\mathl{T_PN}{} menginduksi hasil kali dalam pada
\mathl{T_PM}{.} Untuk membuktikan keterdiferensialan kontinu hasil kali dalam tersebut, misalkan
\maabbdisp {\theta} { W } { W'
} {}
suatu peta untuk $N$ dengan

\mavergleichskette
{\vergleichskette
{ W'
}
{ \subseteq }{\R^n
}
{  }{
}
{    }{
}
{   }{
}
}
{}{}{,}
yang menginduksi suatu bijeksi $\alpha$ antara
\mathl{M \cap W}{} dan
\mathl{\R^m \times \{0\} \cap W'}{}
\zusatzklammer {dengan

\mavergleichskettek
{\vergleichskettek
{ 0
}
{ \in }{ \R^{n-m}
}
{  }{
}
{    }{
}
{   }{
}
}
{}{}{}} {} {.}
Di bawah identifikasi ini berlaku

\mavergleichskette
{\vergleichskette
{ T_PM
}
{ \cong }{ \R^m
}
{ \subseteq }{ \R^n
}
{    }{
}
{   }{
}
}
{}{}{}
dengan vektor-vektor basis

\mathbed {e_i} {}
 {i \leq m} {}
 {} {} {} {.}
Untuk pasangan

\mathbed {e_i, e_j} {}
 {1 \leq i,j \leq m} {}
 {} {} {} {,}
dari vektor-vektor tersebut dan untuk

\mavergleichskette
{\vergleichskette
{ Q
}
{ \in }{ \R^m \times \{0\} \cap W'
}
{  }{
}
{    }{
}
{   }{
}
}
{}{}{}
berlaku kesamaan

\mavergleichskettealign
{\vergleichskettealign
{ h_{ij}(Q)
}
{ \defeq} { \left\langle  T(\alpha^{-1})(e_i)  ,  T(\alpha^{-1}) (e_j)   \right\rangle_{ \alpha^{-1}(Q) }
}
{ =} { \left\langle  T(\theta^{-1})(e_i)  ,  T(\theta^{-1}) (e_j)   \right\rangle_{ \theta^{-1}(Q) }
}
{ =} { g_{ij}(Q)
}
{ } {
}
}
{}
{}{,}
sebab hasil kali dalam pada
\mathl{T_{\alpha^{-1}(Q)}M}{} hanyalah pembatasan hasil kali dalam pada
\mathl{T_{\alpha^{-1}(Q)}N}{}, dan $\alpha$ merupakan pembatasan $\theta$.

}

Contoh paling sederhana ialah submanifold tertutup

\mavergleichskette
{\vergleichskette
{M
}
{ \subseteq }{ \R^n
}
{  }{
}
{    }{
}
{   }{
}
}
{}{}{,}
karena hasil kali dalam standar dapat langsung dibatasi ke $M$.

  \zwischenueberschrift{Medan vektor dan $1$-bentuk pada manifold Riemann}

Orang yang gemar menyindir mengatakan bahwa fisikawan tidak mengenal perbedaan antara medan vektor dan $1$-bentuk. Pada manifold Riemann, kedua objek ini memang saling bersesuaian.


\inputfaktbeweis
{Riemannsche Mannigfaltigkeit/Vektorfelder und 1-Formen/Fakt}
{Lema}
{}
{

\faktsituation {Misalkan $M$ suatu
\definitionsverweis {manifold Riemann}{}{.}}
\faktfolgerung {Maka
\definitionsverweis {pemetaan}{}{}

\maabbeledisp {} {
{ \mathcal V } ( M )  } {
{ \mathcal E }^{  1 } (  M   )
} { F } { \omega_F
} {,}
dengan

\mavergleichskettedisp
{\vergleichskette
{ \left(  \omega_F(P)  \right) (v)
}
{ \defeq} { \left\langle F(P) , v  \right\rangle_P
}
{ } {
}
{ } {
}
{ } {
}
}
{}{}{,}
di mana

\mavergleichskette
{\vergleichskette
{ P
}
{ \in }{ M
}
{  }{
}
{    }{
}
{   }{
}
}
{}{}{}
dan $v$ menyatakan suatu
\definitionsverweis {vektor singgung}{}{}
dari
\mathl{T_PM}{}, merupakan suatu
\definitionsverweis {isomorfisme}{}{}
antara
\definitionsverweis {medan vektor}{}{}
pada $M$ dan
$1$-\definitionsverweis {bentuk
}{}{}
pada $M$.}
\faktzusatz {}
\faktzusatz {}

}
{

Untuk setiap titik

\mavergleichskette
{\vergleichskette
{P
}
{ \in }{ M
}
{  }{
}
{    }{
}
{   }{
}
}
{}{}{}
pemetaan
\maabbeledisp {} {T_PM} {T^*_PM
} { u } { \left\langle  u  ,  -  \right\rangle_P
} {,}
menurut
Lema 38.5 (Aljabar Linear (Osnabrück 2024-2025)),
merupakan suatu
\definitionsverweis {isomorfisme}{}{.}
Dari sini langsung diperoleh bahwa korespondensi global tersebut merupakan bijeksi. Untuk linearitas korespondensi itu, lihat
Soal 16.3.

}

\inputbemerkung
{}
{

Pada suatu
\definitionsverweis {ruang vektor Euklides}{}{},
\definitionsverweis {medan vektor}{}{}
dan
$1$-\definitionsverweis {bentuk diferensial}{}{}
saling bersesuaian menurut
Lema 16.3.
Hal yang sama berlaku untuk suatu
\definitionsverweis {submanifold tertutup}{}{}

\mavergleichskette
{\vergleichskette
{ M
}
{ \subseteq }{ \R^n
}
{  }{
}
{    }{
}
{   }{
}
}
{}{}{,}
dan bentuk diferensial pada $\R^n$ dapat dibatasi ke $M$. Karena itu, suatu medan vektor $F$ pada $\R^n$ juga dapat ditarik balik menjadi medan vektor pada $M$: kita meninjau bentuk diferensial terkait pada $\R^n$, bentuk diferensial tarik balik pada $M$, lalu medan vektor yang bersesuaian dengannya pada $M$. Namun, secara geometrik, suatu titik

\mavergleichskette
{\vergleichskette
{ P
}
{ \in }{ M
}
{  }{
}
{    }{
}
{   }{
}
}
{}{}{}
tidak diberi arah
\mathl{F(P)}{}, sebab pada umumnya vektor ini bahkan tidak termasuk dalam
\definitionsverweis {ruang singgung}{}{}

\mavergleichskette
{\vergleichskette
{ T_PM
}
{ \subseteq }{ T_P\R^n
}
{ = }{ \R^n
}
{    }{
}
{   }{
}
}
{}{}{.}
Sebagai gantinya, kita harus mengambil
\definitionsverweis {proyeksi ortogonal}{}{}
dari
\mathl{F(P)}{} ke
\mathl{T_PM}{}
\zusatzklammer {jadi, struktur Euklides digunakan di sini} {} {.}

}

\inputbeispiel{}
{

Sebagai contoh untuk
Catatan 16.4,
kita meninjau
\definitionsverweis {lingkaran satuan}{}{}

\mavergleichskette
{\vergleichskette
{ S^1
}
{ \subseteq }{ \R^2
}
{  }{
}
{    }{
}
{   }{
}
}
{}{}{}
sebagai
\definitionsverweis {submanifold tertutup}{}{},
serta
\definitionsverweis {medan vektor}{}{}
konstan $e_1$ pada $\R^2$, yang memasangkan setiap titik

\mavergleichskette
{\vergleichskette
{ P
}
{ \in }{ \R^2
}
{  }{
}
{    }{
}
{   }{
}
}
{}{}{}
dengan vektor standar pertama sebagai arahnya. Bentuk diferensial terkait ialah $dx$, yang memetakan $e_1$ ke $1$ dan $e_2$ ke $0$. Bentuk diferensial yang
\definitionsverweis {ditarik balik}{}{}
ke $S^1$ juga dilambangkan dengan $dx$ dan bekerja dengan cara yang sama, tetapi dibatasi pada
\definitionsverweis {ruang singgung}{}{}

\mavergleichskette
{\vergleichskette
{ T_PS^1
}
{ \subseteq }{ \R^2
}
{  }{
}
{    }{
}
{   }{
}
}
{}{}{.}
Medan vektor pada $S^1$ yang bersesuaian dengan bentuk diferensial ini, yang kita nyatakan dengan $H$, dihitung menurut
Lema 16.3
sebagai berikut: Untuk setiap titik

\mavergleichskette
{\vergleichskette
{ P
}
{ = }{ \begin{pmatrix}  a  \\b   \end{pmatrix}
}
{ \in }{ S^1
}
{    }{
}
{   }{
}
}
{}{}{}
dan setiap vektor

\mavergleichskette
{\vergleichskette
{ v
}
{ = }{ \begin{pmatrix} v_1  \\v_2    \end{pmatrix}
}
{ \in }{ T_PS^1
}
{    }{
}
{   }{
}
}
{}{}{}
harus berlaku

\mavergleichskettedisp
{\vergleichskette
{ \left\langle H(P) , v  \right\rangle
}
{ =} { dx(P)(v)
}
{ =} { v_1
}
{ } {
}
{ } {
}
}
{}{}{}
dengan syarat

\mavergleichskette
{\vergleichskette
{ H(P)
}
{ \in }{ T_PS^1
}
{  }{
}
{    }{
}
{   }{
}
}
{}{}{.}
Ruang singgung tersebut berdimensi satu dan direntang oleh
\mathl{\begin{pmatrix}  -b \\a   \end{pmatrix}}{.} Karena itu,

\mavergleichskette
{\vergleichskette
{ v
}
{ = }{ d \begin{pmatrix}  -b \\a   \end{pmatrix}
}
{  }{
}
{    }{
}
{   }{
}
}
{}{}{}
dan

\mavergleichskettedisp
{\vergleichskette
{H(P)
}
{ =} { c \begin{pmatrix}  -b \\a   \end{pmatrix}
}
{ } {
}
{ } {
}
{ } {
}
}
{}{}{}
untuk suatu pasangan

\mavergleichskette
{\vergleichskette
{ d,c
}
{ \in }{ \R
}
{  }{
}
{    }{
}
{   }{
}
}
{}{}{.}
Dari syarat

\mavergleichskettedisp
{\vergleichskette
{ \left\langle H(P) , v  \right\rangle
}
{ =} { \left\langle c \begin{pmatrix}  -b \\a   \end{pmatrix}  , d \begin{pmatrix}  -b \\a   \end{pmatrix}   \right\rangle
}
{ =} { c d
}
{ =} { -bd
}
{ } {
}
}
{}{}{}
langsung diperoleh

\mavergleichskette
{\vergleichskette
{ c
}
{ = }{ -b
}
{  }{
}
{    }{
}
{   }{
}
}
{}{}{.}
Jadi, medan vektor tarik balik tersebut adalah

\mavergleichskettedisp
{\vergleichskette
{ H \left(  \begin{pmatrix}  a  \\b   \end{pmatrix}  \right)
}
{ =} { -b \begin{pmatrix}  -b \\a   \end{pmatrix}
}
{ } {
}
{ } {
}
{ } {
}
}
{}{}{.}

}

  \zwischenueberschrift{Bentuk volume kanonik pada manifold Riemann berorientasi}

\inputdefinition
{}
{

Misalkan $M$ suatu
\definitionsverweis {manifold Riemann}{}{}
\definitionsverweis {berorientasi}{}{}
dengan
\definitionsverweis {dimensi}{}{}
$n$. Untuk

\mavergleichskette
{\vergleichskette
{ P
}
{ \in }{ M
}
{  }{
}
{    }{
}
{   }{
}
}
{}{}{}
misalkan $\omega_P$ merupakan
\definitionsverweis {bentuk selang-seling}{}{} pada
\mathl{T_PM}{}
\zusatzklammer {atau, ekuivalen, elemen yang bersesuaian dari \mathlk{\bigwedge^n T^*_PM}{}} {} {,}
yang memberi nilai $1$ pada setiap
\definitionsverweis {basis ortonormal}{}{}
yang merepresentasikan
\definitionsverweis {orientasi}{}{.}
Maka
$n$-\definitionsverweis {bentuk diferensial
}{}{}
\maabbeledisp {} { M } { \bigwedge^n T^*M
} {P } { \omega_P
} {,}
disebut \definitionswort {bentuk volume kanonik}{} pada $M$.

}

\definitionsverweis {Ukuran terkait}{}{}
dengan bentuk positif ini disebut \stichwort {ukuran kanonik} {} pada $M$. Kita melambangkannya dengan $\lambda_M$. Dengan demikian,
\mathl{\lambda_M(M)}{} adalah ukuran total
\zusatzklammer {luas permukaan atau volume} {} {}
dari $M$.



\inputfaktbeweis
{Riemannsche Mannigfaltigkeit/Orientiert/Kanonische Volumenform/Lokale Berechnung/Fakt}
{Lema}
{}
{

\faktsituation {Misalkan $M$ suatu
\definitionsverweis {manifold Riemann}{}{}
\definitionsverweis {berorientasi}{}{}
dan $\omega$ merupakan
\definitionsverweis {bentuk volume kanonik}{}{.}
Misalkan
\maabbdisp {\alpha} { U } { V
} {}
suatu
\definitionsverweis {peta berorientasi}{}{}
dengan

\mavergleichskette
{\vergleichskette
{V
}
{ \subseteq }{\R^n
}
{  }{
}
{    }{
}
{   }{
}
}
{}{}{}
terbuka, dengan koordinat
\mathl{x_1 , \ldots , x_n}{}, dengan
\definitionsverweis {matriks fundamental metrik}{}{}

\mavergleichskette
{\vergleichskette
{ G
}
{ = }{ (g_{ij})_{1 \leq i,j \leq n}
}
{  }{
}
{    }{
}
{   }{
}
}
{}{}{}
dan

\mavergleichskette
{\vergleichskette
{ g
}
{ = }{ \det  G
}
{  }{
}
{    }{
}
{   }{
}
}
{}{}{.}}
\faktfolgerung {Maka

\mavergleichskettedisp
{\vergleichskette
{ \alpha_* \omega
}
{ =} { \sqrt{ g} dx_1 \wedge \ldots \wedge dx_n
}
{ } {
}
{ } {
}
{ } {
}
}
{}{}{.}}
\faktzusatz {Untuk suatu
\definitionsverweis {himpunan bagian terukur}{}{}

\mavergleichskette
{\vergleichskette
{T
}
{ \subseteq }{ U
}
{  }{
}
{    }{
}
{   }{
}
}
{}{}{}
berlaku

\mavergleichskettedisp
{\vergleichskette
{
\int_{  T   }   \omega
}
{ =} {
\int_{ \alpha(T)  }   \sqrt{g} dx_1 \wedge \ldots \wedge dx_n
}
{ =} { \int_{ \alpha(T)  }   \sqrt{g}    \,  d  \lambda^n
}
{ } {
}
{ } {
}
}
{}{}{.}}
\faktzusatz {}

}
{

Menurut
Definisi 15.3,
kita harus menghitung bentuk diferensial
\mathl{\left(  \alpha^{-1}  \right)^*\omega}{} untuk setiap titik

\mavergleichskette
{\vergleichskette
{ Q
}
{ \in }{ V
}
{  }{
}
{    }{
}
{   }{
}
}
{}{}{.}
Bentuk ini mempunyai wujud
\mathl{c_Q dx_1 \wedge \ldots \wedge dx_n}{} dan ditentukan oleh nilainya pada
\mathl{e_1 \wedge \ldots \wedge e_n}{.} Berlaku

\mavergleichskettedisphandlinks
{\vergleichskettedisphandlinks
{ \left(  \alpha^{-1}  \right)^*\omega \left(  Q,  e_1 \wedge \ldots \wedge e_n  \right)
}
{ =} { \omega \left(  \alpha^{-1} (Q) ,T_Q \left(  \alpha^{-1}  \right)(e_1) \wedge \ldots \wedge T_Q \left(  \alpha^{-1}  \right)( e_n)  \right)
}
{ } {
}
{ } {
}
{ } {
}
}
{}{}{.}
Menurut definisi matriks fundamental metrik,

\mavergleichskettedisp
{\vergleichskette
{ g_{ij}(Q)
}
{ =} { \left\langle  T_Q \left(  \alpha^{-1}  \right) (e_i)   ,  T_Q \left(  \alpha^{-1}  \right) (e_j)   \right\rangle_{\alpha^{-1}(Q)}
}
{ } {
}
{ } {
}
{ } {
}
}
{}{}{.}
Menurut
Teorema 7.8 (Teori Ukuran dan Integrasi (Osnabrück 2022-2023)),
berlaku

\mavergleichskettealigndrucklinks
{\vergleichskettealigndrucklinks
{ \omega \left(  \alpha^{-1}(Q), T_Q \left(  \alpha^{-1}  \right) (e_1) \wedge \ldots \wedge T_Q \left(  \alpha^{-1}  \right) ( e_n)  \right)
}
{ =} { \left(  \det  \left(  \left\langle T_Q \left(  \alpha^{-1}  \right)  (e_i)  ,  T_Q \left(  \alpha^{-1}  \right) (e_j)   \right\rangle  \right)_{1 \leq i,j \leq n}  \right)^{1/2}
}
{ =} { \left(  \det  \left(  g_{ij}(Q)   \right)_{1 \leq i,j \leq n}  \right)^{1/2}
}
{ =} { \sqrt{g(Q)}
}
{ } {
}
}
{}{}{.}

}



\inputfaktbeweis
{Riemannsche Untermannigfaltigkeit des R^n/Einbettung/Volumenform/Fakt}
{Teorema}
{}
{

\faktsituation {Misalkan

\mavergleichskette
{\vergleichskette
{G
}
{ \subseteq }{ \R^m
}
{  }{
}
{    }{
}
{   }{
}
}
{}{}{}
suatu himpunan
\definitionsverweis {terbuka}{}{}
dan misalkan

\mavergleichskette
{\vergleichskette
{M
}
{ \subseteq }{ G
}
{  }{
}
{    }{
}
{   }{
}
}
{}{}{}
suatu
\definitionsverweis {submanifold tertutup}{}{}
\definitionsverweis {berdimensi}{}{} $n$,
yang
\definitionsverweis {berorientasi}{}{}
dan dilengkapi dengan struktur Riemann terinduksi serta
\definitionsverweis {bentuk volume kanonik}{}{}
$\omega$. Misalkan

\mavergleichskette
{\vergleichskette
{W
}
{ \subseteq }{\R^n
}
{  }{
}
{    }{
}
{   }{
}
}
{}{}{}
terbuka dan misalkan
\maabbdisp {\varphi} { W } { U
} {}
suatu
\definitionsverweis {difeomorfisme}{}{}
yang mempertahankan orientasi menuju himpunan terbuka

\mavergleichskette
{\vergleichskette
{ U
}
{ \subseteq }{ M
}
{  }{
}
{    }{
}
{   }{
}
}
{}{}{\zusatzfussnote {Kita juga mengatakan bahwa $\varphi$ merupakan
\zusatzklammer {difeomorfik} {} {}
\stichwort {parametrisasi} {} untuk $U$} {.} {.}}}
\faktfolgerung {Maka $\varphi^{-1}$ merupakan suatu peta untuk $M$, dan pada $W$ berlaku

\mavergleichskettealignhandlinks
{\vergleichskettealignhandlinks
{\varphi^* \left(  \omega \vert_U  \right)
}
{ =} { \left( \det  \left( \left\langle  \begin{pmatrix}  { \frac{   \partial  \varphi_1   }{  \partial   x_i   } }  \\ \vdots \\  { \frac{   \partial  \varphi_m   }{  \partial   x_i   } }   \end{pmatrix}  ,  \begin{pmatrix}  { \frac{   \partial  \varphi_1   }{  \partial   x_j   } }  \\ \vdots \\  { \frac{   \partial  \varphi_m   }{  \partial   x_j   } }   \end{pmatrix}   \right\rangle \right)_{1 \leq i,j \leq n} \right)^{1/2}  dx_1 \wedge \ldots \wedge dx_n
}
{ =} { \left( \det  \left(  { \frac{   \partial  \varphi_1   }{  \partial   x_i   } } { \frac{   \partial  \varphi_1   }{  \partial   x_j   } } + \cdots + { \frac{   \partial  \varphi_m   }{  \partial   x_i   } } { \frac{   \partial  \varphi_m   }{  \partial   x_j   } }  \right)_{1 \leq i,j \leq n} \right)^{1/2}  dx_1 \wedge \ldots \wedge dx_n
}
{ } {
}
{ } {
}
}
{}
{}{.}}
\faktzusatz {}
\faktzusatz {}

}
{

Kesamaan kedua langsung diperoleh dengan mengevaluasi hasil kali dalam standar pada $\R^m$. Menurut definisi
\definitionsverweis {matriks fundamental metrik}{}{},
untuk

\mavergleichskette
{\vergleichskette
{ Q
}
{ \in }{ W
}
{  }{
}
{    }{
}
{   }{
}
}
{}{}{}
berlaku

\mavergleichskettealign
{\vergleichskettealign
{ g_{ij} (Q)
}
{ =} { \left\langle  T_Q(\varphi)(e_i)  ,  T_Q(\varphi)(e_j)   \right\rangle_{\varphi(Q)}
}
{ =} { \left\langle  T_Q(\varphi)(e_i)  ,  T_Q(\varphi)(e_j)   \right\rangle
}
{ =} { \left\langle  \left(  D \varphi  \right)_{ Q } \left(   e_i   \right)  ,  \left(  D\varphi  \right)_{ Q } \left(   e_j   \right)   \right\rangle
}
{ =} { \left\langle  \begin{pmatrix}  { \frac{   \partial  \varphi_1   }{  \partial   x_i   } } ( Q)  \\ \vdots \\  { \frac{   \partial  \varphi_m   }{  \partial   x_i   } } ( Q)   \end{pmatrix}  ,  \begin{pmatrix}  { \frac{   \partial   \varphi_1   }{  \partial   x_j   } } ( Q )  \\ \vdots \\  { \frac{   \partial  \varphi_m   }{  \partial   x_j   } } ( Q)   \end{pmatrix}   \right\rangle
}
}
{}
{}{,}
sebab ruang singgung
\mathl{T_{\varphi(Q)}M}{} membawa hasil kali dalam yang diinduksi dari $\R^m$, pemetaan singgung dalam situasi lokal sama dengan diferensial total, dan entri-entri diferensial total tersebut dapat dinyatakan melalui turunan parsial. Karena itu, pernyataan teorema diperoleh dari
Lema 16.7.

}

\inputbeispiel{}
{

Misalkan

\mavergleichskette
{\vergleichskette
{ I
}
{ \subseteq }{ \R
}
{  }{
}
{    }{
}
{   }{
}
}
{}{}{}
suatu
\definitionsverweis {interval terbuka}{}{}
dan
\maabbdisp {\varphi} { I } {\R^m
} {}
suatu
\definitionsverweis {kurva terdiferensialkan}{}{}
\definitionsverweis {reguler}{}{,}
yakni

\mavergleichskette
{\vergleichskette
{ \varphi'(t)
}
{ \neq }{ 0
}
{  }{
}
{    }{
}
{   }{
}
}
{}{}{}
di mana-mana. Misalkan pula $\varphi$ injektif dan citra

\mavergleichskette
{\vergleichskette
{ M
}
{ = }{ \varphi(I)
}
{  }{
}
{    }{
}
{   }{
}
}
{}{}{}
dari $I$ merupakan suatu
\definitionsverweis {submanifold tertutup}{}{}
berdimensi satu dari suatu himpunan bagian terbuka

\mavergleichskette
{\vergleichskette
{ G
}
{ \subseteq }{ \R^m
}
{  }{
}
{    }{
}
{   }{
}
}
{}{}{.}
Kita memberi submanifold ini orientasi yang diinduksi oleh parametrisasinya sehingga parametrisasi tersebut mempertahankan orientasi. Menurut
Teorema 16.8,
untuk
\definitionsverweis {bentuk kanonik}{}{}
$\omega$ pada $M$
\zusatzklammer {atau
\definitionsverweis {ukuran kanonik}{}{,}
yang dalam hal ini merupakan ukuran panjang} {} {}
berlaku hubungan

\mavergleichskettealign
{\vergleichskettealign
{ \varphi^* \omega
}
{ =} { \left(  \left\langle  \begin{pmatrix} \varphi'_1(t) \\\vdots\\ \varphi'_m(t)   \end{pmatrix}  ,  \begin{pmatrix} \varphi'_1(t) \\\vdots\\ \varphi'_m(t)   \end{pmatrix}   \right\rangle  \right)^{1/2}  dt
}
{ =} { \sqrt{ (\varphi'_1(t))^2 + \cdots + (\varphi'_m(t))^2  }   dt
}
{ =} { \Vert {\varphi'(t)} \Vert dt
}
{ } {
}
}
{}
{}{.}
Dengan demikian, jika

\mavergleichskette
{\vergleichskette
{ I
}
{ = }{ {]a,b[}
}
{  }{
}
{    }{
}
{   }{
}
}
{}{}{}
maka untuk ukuran
\zusatzklammer {yaitu panjang} {} {}
dari $M$ berlaku rumus

\mavergleichskettedisp
{\vergleichskette
{ \lambda_M (M)
}
{ =} { \int_{  a  }^{  b  } \sqrt{ (\varphi'_1(t))^2 + \cdots + (\varphi'_m(t))^2  } \, d t
}
{ } {
}
{ } {
}
{ } {
}
}
{}{}{.}
Rumus ini bersesuaian dengan rumus yang diperoleh dalam
Teorema 38.6 (Analisis (Osnabrück 2021-2023))
melalui teori
\definitionsverweis {kurva terektifikasi}{}{.}

}

\inputbeispiel{}
{

Misalkan

\mavergleichskette
{\vergleichskette
{ I
}
{ \subseteq }{ \R
}
{  }{
}
{    }{
}
{   }{
}
}
{}{}{}
suatu
\definitionsverweis {interval riil}{}{}
dan misalkan
\maabbdisp {g} { I } {\R_+
} {}
suatu fungsi positif yang terdiferensialkan, yang kita tafsirkan sebagai
\definitionsverweis {metrik Riemann}{}{}
pada $I$. Makna $g(t)$ ialah mendefinisikan suatu hasil kali dalam pada ruang singgung satu dimensi dari interval tersebut di titik $t$. Jadi,

\mavergleichskettedisp
{\vergleichskette
{ \left\langle  1  ,  1  \right\rangle_t
}
{ =} { g(t)
}
{ } {
}
{ } {
}
{ } {
}
}
{}{}{}
dan karena itu

\mavergleichskettedisp
{\vergleichskette
{ \Vert { 1 } \Vert_t
}
{ =} { \sqrt{ g(t) }
}
{ } {
}
{ } {
}
{ } {
}
}
{}{}{.}
Jika
\maabb {\gamma} { I } { \R^n
} {}
suatu
\definitionsverweis {kurva yang terdiferensialkan secara kontinu}{}{}
\definitionsverweis {reguler}{}{}
dan $\R^n$ dilengkapi dengan metrik Euklides, maka melalui pemetaan tersebut $I$ mewarisi metrik seperti ini, yaitu

\mavergleichskettedisp
{\vergleichskette
{ g(t)
}
{ \defeq} { \left\langle  \gamma'(t) ,  \gamma'(t)  \right\rangle
}
{ } {
}
{ } {
}
{ } {
}
}
{}{}{.}
Menurut
Teorema 38.6 (Analisis (Osnabrück 2021-2023)),
panjang kurva $\gamma$ pada

\mavergleichskette
{\vergleichskette
{ [a,b]
}
{ \subseteq }{ I
}
{  }{
}
{    }{
}
{   }{
}
}
{}{}{}
adalah

\mavergleichskettedisp
{\vergleichskette
{ \int_a^b \Vert { \gamma'(t) } \Vert  dt
}
{ =} { \int_a^b  \sqrt{ g(t) }  dt
}
{ } {
}
{ } {
}
{ } {
}
}
{}{}{.}
Jadi, panjang kurva pada interval dapat dihitung jika pengukuran panjang infinitesimal pada interval tersebut diketahui. Sebaliknya, setiap metrik positif $g$ pada $I$ dapat diperoleh melalui fungsi
\maabbdisp {\gamma} { I } { \R
} {}
dengan

\mavergleichskettedisp
{\vergleichskette
{ \gamma(t)
}
{ \defeq} { \int_0^t \sqrt{ g(s) } ds
}
{ } {
}
{ } {
}
{ } {
}
}
{}{}{.}
Dengan demikian, metrik ini merupakan metrik tarik balik dari metrik standar pada $\R$.

}

\chapter{Lembar Kerja 16}
\setcounter{section}{16}

  \zwischenueberschrift{Soal latihan}

\inputaufgabegibtloesung
{}
{

Kita meninjau parabola standar

\mavergleichskette
{\vergleichskette
{ Y
}
{ \subseteq }{ \R^2
}
{  }{
}
{    }{
}
{   }{
}
}
{}{}{}
dengan struktur Riemann terinduksi dan parametrisasi
\maabbeledisp {\varphi} {\R} { Y
} { t } { \left(  t   , \,  t^2                   \right)
} {,}
yang kita pandang sebagai suatu peta
\zusatzklammer {invers} {} {.}
Tentukan
\definitionsverweis {fungsi fundamental}{}{}
Riemann

\mavergleichskette
{\vergleichskette
{ g(t)
}
{ = }{ g_{11} (t)
}
{  }{
}
{    }{
}
{   }{
}
}
{}{}{.}

}
{} {}

\inputaufgabe
{}
{

Misalkan $M$ suatu
\definitionsverweis {manifold diferensiabel}{}{}
sedemikian sehingga
\definitionsverweis {bundel tangen}{}{}
$TM$ trivial. Tunjukkan bahwa melalui suatu trivialisasi diferensiabel

\mavergleichskette
{\vergleichskette
{ TM
}
{ \cong }{ M \times \R^n
}
{  }{
}
{    }{
}
{   }{
}
}
{}{}{}
\zusatzklammer {dengan bantuan
\definitionsverweis {hasil kali dalam standar}{}{}
pada $\R^n$} {} {}
$M$ dapat dijadikan suatu
\definitionsverweis {manifold Riemann}{}{.}

}
{} {}

\inputaufgabe
{}
{

Kita meninjau lingkaran

\mavergleichskette
{\vergleichskette
{ S^1
}
{ \subseteq }{ \R^2
}
{  }{
}
{    }{
}
{   }{
}
}
{}{}{}
dengan
\definitionsverweis {struktur Riemann terinduksi}{}{.}
Trivialisasi
\definitionsverweis {bundel tangen}{}{}
\maabbeledisp {} { S^1 \times \R } { T S^1
} {( a,b, u )} { (a,b, -ub,ua )
} {,}
melalui hasil kali dalam standar pada $\R$, juga menghasilkan suatu struktur Riemann pada $S^1$. Tunjukkan bahwa kedua struktur Riemann tersebut berimpit.

}
{} {}

\inputaufgabe
{}
{

Misalkan
\mathkor {} {L} {dan} {M} {}
merupakan
\definitionsverweis {manifold Riemann}{}{.}
Tunjukkan bahwa $L \times M$ secara alami juga merupakan suatu manifold Riemann.

}
{} {}

\inputaufgabe
{}
{

Kita meninjau suatu himpunan terbuka
\mathl{V \subseteq \R^n}{} sebagai
\definitionsverweis {manifold Riemann}{}{.} Apa
\definitionsverweis {bentuk volume kanonik}{}{} pada $V$?

}
{} {}

\inputaufgabe
{}
{

Misalkan $M$ suatu
\definitionsverweis {manifold Riemann}{}{.}
Tunjukkan bahwa
\definitionsverweis {pemetaan}{}{}
\maabbeledisp {} {
{ \mathcal V } ( M )  } {
{ \mathcal E }^{  1 } (  M   )
} { F } { \omega_F
} {,}
dengan

\mavergleichskettedisp
{\vergleichskette
{ \left(  \omega_F(P)  \right) (v)
}
{ \defeq} { \left\langle F(P) , v  \right\rangle_P
}
{ } {
}
{ } {
}
{ } {
}
}
{}{}{,}
bersifat
\definitionsverweis {linear}{}{.}

}
{} {}

\inputaufgabe
{}
{

Kita meninjau suatu himpunan terbuka
\mathl{V \subseteq \R^n}{} sebagai
\definitionsverweis {manifold Riemann}{}{.} Apa yang dinyatakan oleh korespondensi antara
\definitionsverweis {medan vektor}{}{}
dan
$1$-\definitionsverweis {bentuk diferensial}{}{}
yang diuraikan dalam
Lema 16.3
untuk situasi ini?

}
{} {}

\inputaufgabe
{}
{

Berikan alasan untuk
Catatan 16.4.

}
{} {}

\inputaufgabe
{}
{

Misalkan $M$ suatu
\definitionsverweis {manifold Riemann}{}{}
\definitionsverweis {terorientasi}{}{.} Tunjukkan bahwa
\definitionsverweis {bentuk volume kanonik}{}{} $\omega$ ditentukan oleh syarat bahwa pada setiap titik, nilainya pada setiap basis ortonormal yang merepresentasikan orientasi adalah $1$.

}
{} {}

\inputaufgabe
{}
{

Misalkan $V$ suatu ruang vektor riil terorientasi berdimensi $n$ dan $\lambda$ suatu
\definitionsverweis {ukuran yang invarian terhadap translasi}{}{}
pada $V$. Tunjukkan bahwa pemetaan
\maabbeledisp {} {V \times \cdots \times V} {\R
} {(v_1 , \ldots , v_n) } { \pm \lambda (P(v_1 , \ldots , v_n))
} {,}
dengan tanda positif dipilih jika vektor tersebut merepresentasikan orientasi, merupakan suatu pemetaan multilinear selang-seling.

}
{} {}

\inputaufgabe
{}
{

Tunjukkan bahwa pada suatu
\definitionsverweis {manifold Riemann}{}{},
\definitionsverweis {pemetaan koordinat}{}{}
\maabbdisp {\alpha} { U } { V
} {}
secara umum tidak menginduksi suatu
\definitionsverweis {isometri}{}{}
\maabbdisp {T_P(\alpha)} {T_PU} {T_{\alpha(P)} V
} {}
\zusatzklammer {jika
\mathl{T_PU}{} dilengkapi dengan
\mathl{\left\langle  - ,  -  \right\rangle_P}{} dan
\mathl{T_{\alpha(P)} V= \R^n}{} dilengkapi dengan hasil kali dalam standar} {} {.}

}
{} {}

\inputaufgabegibtloesung
{}
{

Kita meninjau grafik $M$ dari pemetaan
\maabbeledisp {\varphi} {\R^2} {\R
} {(u,v)} { u^2+uv-v^3
} {,} sebagai suatu submanifold tertutup berdimensi dua dari $\R^3$, yaitu

\mavergleichskettedisp
{\vergleichskette
{ M
}
{ =} { \left\{ (u,v,u^2+uv-v^3)  \mid  (u,v) \in \R^2         \right\}
}
{ } {
}
{ } {
}
{ } {
}
}
{}{}{,}
dengan metrik Riemann yang diinduksi dari $\R^3$. Pemetaan
\maabbeledisp {\psi} {\R^2} { M
} {(u,v)} { (u,v,u^2+uv-v^3)
} {,}
merupakan difeomorfisme yang terkait.

a) Tentukan diferensial total dari $\psi$ serta vektor citra
\mathl{T_P(\psi) (e_1)}{} dan
\mathl{T_P(\psi) (e_2)}{} dalam
\mathl{T_{\psi(P)}M}{.}

b) Untuk setiap titik berbentuk

\mavergleichskette
{\vergleichskette
{ P
}
{ = }{ (u,0)
}
{  }{
}
{    }{
}
{   }{
}
}
{}{}{}
tentukan luas jajaran genjang yang direntang oleh
\mathl{T_P(\psi) (e_1)}{} dan
\mathl{T_P(\psi) (e_2)}{} dalam
\mathl{T_{\psi(P)}M}{.}

c) Untuk setiap titik berbentuk

\mavergleichskette
{\vergleichskette
{ P
}
{ = }{ (0,v)
}
{  }{
}
{    }{
}
{   }{
}
}
{}{}{}
tentukan luas jajaran genjang yang direntang oleh
\mathl{T_P(\psi) (e_1)}{} dan
\mathl{T_P(\psi) (e_2)}{} dalam
\mathl{T_{\psi(P)}M}{.}

}
{} {}

\inputaufgabe
{}
{

Misalkan
\maabbdisp {\varphi} {\R^n } {\R
} {}
\zusatzklammer {dengan \mathlk{m=n -1 \geq 0}{}} {} {}
suatu
\definitionsverweis {fungsi yang terdiferensialkan secara kontinu}{}{,}
yang pada setiap titik dari
\definitionsverweis {serat}{}{}
$M$ di atas
\mathl{0 \in \R}{}, bersifat
\definitionsverweis {reguler}{}{.}
Kita memandang $M$ sebagai suatu manifold Riemann terorientasi. Tunjukkan bahwa antara bentuk volume $\tau$ dari
Korolari 15.6
dan
\definitionsverweis {bentuk volume kanonik}{}{}
$\omega$ berlaku hubungan

\mavergleichskettedisp
{\vergleichskette
{\tau(P,v_1 , \ldots , v_m)
}
{ =} { \pm \Vert {
\operatorname{Grad} \, \varphi ( P ) } \Vert  \omega(P, v_1 , \ldots , v_m)
}
{ } {
}
{ } {
}
{ } {
}
}
{}{}{.}

}
{} {}

\inputaufgabe
{}
{

Misalkan
\maabbdisp {\varphi} {\R^n } {\R^\ell
} {}
\zusatzklammer {dengan \mathlk{m=n - \ell \geq 0}{}} {} {}
suatu
\definitionsverweis {pemetaan yang terdiferensialkan secara kontinu}{}{,}
yang pada setiap titik dari
\definitionsverweis {serat}{}{}
$M$ di atas
\mathl{0 \in \R^\ell}{}, bersifat
\definitionsverweis {reguler}{}{.}
Kita memandang $M$ sebagai suatu manifold Riemann terorientasi. Andaikan bahwa gradien

\mathdisp {\operatorname{Grad} \, \varphi_1  (  P ) , \ldots ,
\operatorname{Grad} \, \varphi_\ell ( P )} {  }

saling tegak lurus pada setiap titik
\mathl{P\in M}{.} Tunjukkan bahwa antara bentuk volume $\tau$ dari
Korolari 15.6
dan
\definitionsverweis {bentuk volume kanonik}{}{}
$\omega$ berlaku hubungan

\mavergleichskettedisp
{\vergleichskette
{\tau(P,v_1 , \ldots , v_m)
}
{ =} { \pm \Vert {
\operatorname{Grad} \, \varphi_1  (  P ) } \Vert  \cdots \Vert {
\operatorname{Grad} \, \varphi_\ell ( P ) } \Vert \omega(P, v_1 , \ldots , v_m)
}
{ } {
}
{ } {
}
{ } {
}
}
{}{}{.}

}
{} {}

\inputaufgabe
{}
{

Tunjukkan bahwa
\definitionsverweis {bentuk luas kanonik}{}{}
pada
\definitionsverweis {permukaan bola satuan}{}{}

\mavergleichskette
{\vergleichskette
{S^2
}
{ \subset }{ \R^3
}
{  }{
}
{    }{
}
{   }{
}
}
{}{}{}
adalah

\mathdisp {x dy \wedge dz  -y dx \wedge dz +  z dx \wedge dy} {  }

dengan $x,y,z$ menyatakan koordinat pada $\R^3$.

}
{} {}

Suatu
\definitionsverweis {bundel vektor}{}{}
\maabb {p} { E } { X
} {}
di atas suatu
\definitionsverweis {manifold diferensiabel}{}{}
$X$ disebut
\definitionswort {bundel vektor Riemann}{,}
jika pada setiap serat $E_x$ diberikan suatu
\definitionsverweis {hasil kali dalam}{}{}
dengan sifat bahwa terdapat
\definitionsverweis {peta}{}{}
yang mentrivialkan
\maabbdisp {\alpha} { U } { V \subseteq \R^n
} {}
bersama
\maabbdisp {\theta} {E\vert_U \cong U \times \R^r } { V \times \R^r
} {}
sedemikian sehingga fungsi

\mavergleichskettedisp
{\vergleichskette
{ h_{ij} (Q )
}
{ \defeq} { \left\langle  \theta^{-1}(Q,e_i)  , \theta^{-1}(Q,e_j)    \right\rangle_{\alpha^{-1}(Q)}
}
{ } {
}
{ } {
}
{ } {
}
}
{}{}{}
yang terdiferensialkan secara kontinu.

Jadi, pada suatu
\definitionsverweis {manifold Riemann}{}{},
\definitionsverweis {bundel tangen}{}{}
merupakan suatu bundel vektor Riemann.

\inputaufgabe
{}
{

Misalkan
\maabb {p} { E } { M
} {}
suatu
\definitionsverweis {bundel vektor Riemann}{}{}
di atas suatu
\definitionsverweis {manifold diferensiabel}{}{}
$M$, dan misalkan
\maabb {\varphi} { L } { M
} {}
suatu
\definitionsverweis {pemetaan diferensiabel}{}{}
dengan
\definitionsverweis {bundel vektor tarik balik}{}{}
\maabb {} {\varphi^*E} {L
} {.}
Tunjukkan bahwa $\varphi^*E$ secara alami juga merupakan suatu bundel vektor Riemann.

}
{} {}

  \zwischenueberschrift{Soal untuk dikumpulkan}

Pada suatu manifold Riemann $M$, untuk suatu vektor singgung

\mavergleichskette
{\vergleichskette
{ v
}
{ \in }{ T_PM
}
{  }{
}
{    }{
}
{   }{
}
}
{}{}{}
normanya didefinisikan oleh

\mavergleichskette
{\vergleichskette
{ \Vert {v} \Vert
}
{ = }{ \sqrt{ \left\langle v , v  \right\rangle_P }
}
{  }{
}
{    }{
}
{   }{
}
}
{}{}{.}

\inputaufgabe
{4}
{

Misalkan $M$ suatu
\definitionsverweis {manifold Riemann}{}{.}
Tunjukkan bahwa pemetaan
\maabbeledisp {} {TM} {\R
} { v } { \Vert { v } \Vert
} {,}
bersifat
\definitionsverweis {kontinu}{}{.}

}
{} {}

\inputaufgabe
{3}
{

Tunjukkan bahwa
\mathl{\R \times \R_+}{} dengan bentuk bilinear yang diberikan oleh
\definitionsverweis {bentuk Hesse}{}{} dari
\definitionsverweis {fungsi}{}{}
\maabbeledisp {f} {\R \times \R_+} {\R
} {(x,y)} { x^2+y^4
} {,}
merupakan suatu
\definitionsverweis {manifold Riemann}{}{.}

}
{} {}

\inputaufgabe
{4}
{

Untuk setiap titik
\mathl{P=(x,y,z)}{} pada
\definitionsverweis {permukaan bola satuan}{}{}
$K$, berikan suatu
\definitionsverweis {basis ortonormal}{}{}
dalam
\mathl{T_PK \subset \R^3}{}
\zusatzklammer {terhadap
\definitionsverweis {struktur Riemann terinduksi}{}{}} {} {.}

}
{} {}

\inputaufgabe
{6}
{

Dalam $\R^3$, diberikan
\definitionsverweis {elipsoid}{}{}

\mathdisp {E= \left\{ (x,y,z)  \mid   x^2+y^2+3z^2 \leq 5        \right\}} {  }
dan bidang

\mathdisp {M= \left\{ (x,y,z)  \mid   7x-3y-2z =  2         \right\}} {  }.
Hitunglah luas irisan
\mathl{M \cap E}{.}

}
{} {}

\inputaufgabe
{6}
{

Buatlah suatu grafik komputer yang memvisualisasikan situasi yang diuraikan dalam
Catatan 16.4
dengan menggunakan suatu permukaan di $\R^3$.

}
{} {}

\section*{Solusi yang disediakan oleh sumber}
\addcontentsline{toc}{section}{Solusi yang disediakan oleh sumber}
\subsection*{Solusi Soal 16.1}
Berlaku

\mavergleichskettealign
{\vergleichskettealign
{ g(t)
}
{ =} { \left\langle  T_t \varphi(1)   , T_t \varphi(1)   \right\rangle_{ Y, \varphi(t) }
}
{ =} { \left\langle    \varphi' (t)   ,   \varphi'(t)   \right\rangle
}
{ =} { \left\langle  \begin{pmatrix}  1 \\ 2t   \end{pmatrix}     ,  \begin{pmatrix}  1 \\ 2t   \end{pmatrix}   \right\rangle
}
{ =} { 1+4t^2
}
}
{}
{}{.}
\subsection*{Solusi Soal 16.12}
a) Diferensial total dari $\psi$ pada titik
\mathl{P=(u,v)}{} adalah

\mathdisp {\left(D \psi\right)_{P} = \begin{pmatrix}  1 &  0 \\  0 &  1 \\   2u+v & u-3v^2 \end{pmatrix}} {  }

dan berlaku

\mathdisp {T_P(\psi) (e_1) = \left(  D \psi  \right)_{ P } \left(  e_1  \right) = \begin{pmatrix}  1 \\ 0\\  2u+v   \end{pmatrix} \text{   dan } T_P(\psi) (e_2)   = \left(  D \psi  \right)_{ P } \left(  e_2  \right) = \begin{pmatrix}  0 \\ 1\\ u-3v^2   \end{pmatrix}} { . }

b) dan c) Untuk menentukan luas, pertama-tama kita menghitung ketiga hasil kali dalam yang dibentuk oleh kedua vektor
\mathkor {} {\begin{pmatrix}  1 \\ 0\\  2u+v   \end{pmatrix}} {dan} {\begin{pmatrix}  0 \\ 1\\ u-3v^2   \end{pmatrix}} {.}
Berlaku

\mathdisp {\left\langle \begin{pmatrix}  1 \\ 0\\  2u+v   \end{pmatrix}  , \begin{pmatrix}  1 \\ 0\\  2u+v   \end{pmatrix}   \right\rangle = 1+(2u+v)^2 = 1 +4u^2+4uv +v^2} { , }

\mathdisp {\left\langle \begin{pmatrix}  1 \\ 0\\  2u+v   \end{pmatrix}  , \begin{pmatrix}  0 \\ 1\\ u-3v^2   \end{pmatrix}    \right\rangle = (2u+v)(u-3v^2) = 2u^2 -6uv^2 +uv -3v^3} {  }

dan

\mathdisp {\left\langle \begin{pmatrix}  0 \\ 1\\ u-3v^2   \end{pmatrix}   , \begin{pmatrix}  0 \\ 1\\ u-3v^2   \end{pmatrix}    \right\rangle = 1+(u-3v^2)^2 = 1 +u^2 -6uv^2  + 9v^4} { . }

b) Untuk
\mathl{P=(u,0)}{}, luas jajaran genjang yang direntang oleh
\mathl{T_P(\psi) (e_1)}{} dan
\mathl{T_P(\psi) (e_2)}{} dalam
\mathl{T_{\psi(P)}M}{} adalah

\mavergleichskettedisp
{\vergleichskette
{ \sqrt { \det  \begin{pmatrix}  1+4u^2  &   2u^2   \\  2 u^2 &  1+ u^2  \end{pmatrix} }
}
{ =} { \sqrt{  (1+4u^2)(1+u^2) - 4u^4 }
}
{ =} { \sqrt{  1+5u^2 }
}
{ } {
}
{ } {
}
}
{}{}{.}

c) Untuk
\mathl{P=(0,v)}{}, luas jajaran genjang yang direntang oleh
\mathl{T_P(\psi) (e_1)}{} dan
\mathl{T_P(\psi) (e_2)}{} dalam
\mathl{T_{\psi(P)}M}{} adalah

\mavergleichskettedisp
{\vergleichskette
{ \sqrt { \det  \begin{pmatrix}  1+v^2  &   -3v^3   \\  -3  v^3 &  1+ 9v^4  \end{pmatrix} }
}
{ =} { \sqrt{  (1+v^2)(1+9v^4) - 9v^6 }
}
{ =} { \sqrt{  1+v^2 +9v^4  }
}
{ } {
}
{ } {
}
}
{}{}{.}

\part{Unit 17}
\chapter[Kuliah 17]{Kuliah 17: Perhitungan pada Manifold Riemann}
\setcounter{section}{17}

Dalam kuliah ini kita melanjutkan teori manifold Riemann dan, khususnya, menghitung beberapa luas permukaan.

  \zwischenueberschrift{Perhitungan pada manifold Riemann}



\inputfaktbeweis
{Graph einer Funktion/Riemannsche Untermannigfaltigkeit des R^n/Volumenform/Fakt}
{Korollar}
{}
{

\faktsituation {Misalkan

\mavergleichskette
{\vergleichskette
{W
}
{ \subseteq }{\R^n
}
{  }{
}
{    }{
}
{   }{
}
}
{}{}{}
suatu
\definitionsverweis {himpunan bagian terbuka}{}{}
dan
\maabbdisp {\psi} { W } {\R
} {}
suatu
\definitionsverweis {fungsi terdiferensialkan}{}{.}
Misalkan

\mavergleichskettedisphandlinks
{\vergleichskettedisphandlinks
{M
}
{ =} { \left\{  \left(  x_1 , \ldots , x_n , \, \psi(x_1 , \ldots , x_n) \right)   \mid  (x_1 , \ldots , x_n) \in W         \right\}
}
{ \subseteq} { W \times \R
}
{ } {
}
{ } {
}
}
{}{}{}
merupakan
\definitionsverweis {grafik}{}{}
dari $\psi$.}
\faktfolgerung {Maka $M$ merupakan suatu
\definitionsverweis {manifold Riemann}{}{} yang \definitionsverweis {berorientasi}{}{,}
dan untuk
\definitionsverweis {bentuk volume kanonik}{}{}
$\omega$ pada $M$ berlaku

\mavergleichskettedisp
{\vergleichskette
{ (\operatorname{Id} \times \psi)^*(\omega)
}
{ =} { \left(  1+ \sum_{i = 1}^n \left(  { \frac{   \partial   \psi  }{  \partial   x_i   } }  \right)^2  \right)^{1/2} dx_1 \wedge \ldots \wedge dx_n
}
{ } {
}
{ } {
}
{ } {
}
}
{}{}{.}}
\faktzusatz {}
\faktzusatz {}

}
{

Pemetaan
\maabbeledisp {\varphi = \operatorname{id} \times \psi} { W } { W \times \R
} { x } {(x, \psi(x))
} {,}
merupakan suatu
\definitionsverweis {difeomorfisme}{}{}
antara $W$ dan grafik $M$. Grafik tersebut merupakan suatu
\definitionsverweis {submanifold tertutup}{}{}
dari
\mathl{W \times \R}{} sehingga membawa
\definitionsverweis {struktur Riemann terinduksi}{}{}
dan
\zusatzklammer {karena orientasi $W$ dipindahkan ke $M$} {} {}
suatu
\definitionsverweis {bentuk volume kanonik}{}{}
$\omega$. Pada situasi ini kita dapat menerapkan
Teorema 16.8.
\definitionsverweis {Turunan parsial}{}{}
$\varphi$ terhadap variabel ke-$i$ adalah

\mavergleichskettedisp
{\vergleichskette
{ { \frac{   \partial  \varphi  }{  \partial   x_i   } }
}
{ =} { \begin{pmatrix}  0  \\\vdots\\  0\\ 1\\  0\\ \vdots\\  0\\  { \frac{   \partial   \psi  }{  \partial   x_i   } }   \end{pmatrix}
}
{ } {
}
{ } {
}
{ } {
}
}
{}{}{.}
Misalkan

\mavergleichskette
{\vergleichskette
{Q
}
{ \in }{ W
}
{  }{
}
{    }{
}
{   }{
}
}
{}{}{}
suatu titik yang selanjutnya kita substitusikan ke dalam fungsi-fungsi tersebut, sehingga di seluruh perhitungan kita bekerja dengan bilangan real. Hasil kali dalam yang membentuk entri-entri
\mathl{b_{ij}}{} dari matriks $B$
\zusatzklammer {yang determinannya perlu kita ambil akarnya} {} {}
ialah

\mavergleichskettedisp
{\vergleichskette
{ b_{ij}
}
{ =} { \left\langle  { \frac{   \partial  \varphi  }{  \partial   x_i   } }(Q)  ,  { \frac{   \partial  \varphi  }{  \partial   x_j   } }(Q)   \right\rangle
}
{ } {}
{ } {}
{ } {}
}
{}{}{}
Kita tulis

\mavergleichskette
{\vergleichskette
{B
}
{ = }{ E_n+A
}
{  }{
}
{    }{
}
{   }{
}
}
{}{}{}
dengan

\mavergleichskette
{\vergleichskette
{ A
}
{ = }{ \left(  a_{ij}  \right)_{1 \leq i,j \leq n}
}
{  }{
}
{    }{
}
{   }{
}
}
{}{}{.}
Dengan

\mavergleichskette
{\vergleichskette
{ c_i
}
{ = }{ { \frac{   \partial   \psi  }{  \partial   x_i   } } ( Q)
}
{  }{
}
{    }{
}
{   }{
}
}
{}{}{}
kita dapat menuliskan

\mavergleichskette
{\vergleichskette
{ a_{ij}
}
{ = }{ c_ic_j
}
{  }{
}
{    }{
}
{   }{
}
}
{}{}{}
dan seluruh matriks $A$ sebagai

\mavergleichskettedisp
{\vergleichskette
{A
}
{ =} { \begin{pmatrix} c_1  \\\vdots\\ c_n   \end{pmatrix} \left( c_1   , \,   \ldots  , \,  c_n                 \right)
}
{ } {
}
{ } {
}
{ } {
}
}
{}{}{}
Dengan demikian, $A$ mendeskripsikan suatu pemetaan linear dari $\R^n$ ke $\R^n$ yang
\definitionsverweis {terfaktorkan}{}{} melalui $\R$,
sehingga mempunyai suatu
\definitionsverweis {kernel}{}{}
yang berdimensi sekurang-kurangnya $(n-1)$. Sebut kernel itu $K$. Jika dimensinya $n$, maka

\mavergleichskette
{\vergleichskette
{A
}
{ = }{ 0
}
{  }{
}
{    }{
}
{   }{
}
}
{}{}{}
dan $B$ adalah identitas, sehingga pernyataannya benar. Jadi, misalkan

\mavergleichskette
{\vergleichskette
{A
}
{ \neq }{ 0
}
{  }{
}
{    }{
}
{   }{
}
}
{}{}{.}
Maka

\mavergleichskette
{\vergleichskette
{v
}
{ = }{ \begin{pmatrix} c_1  \\\vdots\\ c_n   \end{pmatrix}
}
{  }{
}
{    }{
}
{   }{
}
}
{}{}{}
merupakan suatu
\definitionsverweis {vektor eigen}{}{}
dari $A$ untuk
\definitionsverweis {nilai eigen}{}{}

\mavergleichskette
{\vergleichskette
{ c_1^2 + \cdots + c_n^2
}
{ \neq }{ 0
}
{  }{
}
{    }{
}
{   }{
}
}
{}{}{.}
Vektor ini merupakan vektor eigen $B$ untuk nilai eigen
\mathl{1+c_1^2 + \cdots + c_n^2}{,} sedangkan $K$ merupakan
\definitionsverweis {ruang eigen}{}{}
bagi $B$ untuk nilai eigen $1$ dan berdimensi
\mathl{(n-1)}{.} Secara keseluruhan, $B$
\definitionsverweis {dapat didiagonalkan}{}{}
dan
\definitionsverweis {determinan}{}{}
matriks tersebut adalah hasil kali nilai-nilai eigennya, yaitu
\mathl{1+c_1^2 + \cdots + c_n^2}{.}

}

Dengan pendekatan ini kita, misalnya, dapat menghitung luas permukaan sfer satuan; lihat
Soal 17.2.



\inputfaktbeweis
{Flächenstück im Raum/Einbettung/Flächenform/Fakt}
{Korollar}
{}
{

\faktsituation {Misalkan

\mavergleichskette
{\vergleichskette
{ M
}
{ \subseteq }{ G
}
{  }{
}
{    }{
}
{   }{
}
}
{}{}{}
suatu
\definitionsverweis {submanifold tertutup terbenam}{}{\zusatzfussnote {Permukaan adalah manifold dua dimensi} {.} {}}
berdimensi dua dan berorientasi di dalam suatu
\definitionsverweis {himpunan terbuka}{}{}

\mavergleichskette
{\vergleichskette
{ G
}
{ \subseteq }{ \R^3
}
{  }{
}
{    }{
}
{   }{
}
}
{}{}{,}
yang dilengkapi dengan struktur Riemann terinduksi dan
\definitionsverweis {bentuk luas kanonik}{}{}
$\omega$. Misalkan

\mavergleichskette
{\vergleichskette
{ W
}
{ \subseteq }{ \R^2
}
{  }{
}
{    }{
}
{   }{
}
}
{}{}{}
terbuka dan misalkan
\maabbdisp {\varphi} { W } { U
} {}
suatu
\definitionsverweis {difeomorfisme}{}{}
yang mempertahankan orientasi ke himpunan terbuka

\mavergleichskette
{\vergleichskette
{ U
}
{ \subseteq }{ M
}
{  }{
}
{    }{
}
{   }{
}
}
{}{}{.}
Koordinat-koordinat pada $\R^2$ adalah
\mathkor {} {u} {dan} {v} {}
dan kita tetapkan
\zusatzfussnote {Notasi ini sudah digunakan oleh Carl Friedrich Gauss} {.} {}

\mathdisp {E = \left\langle  \begin{pmatrix}  { \frac{   \partial  \varphi_1   }{  \partial  u  } }  \\ { \frac{   \partial  \varphi_2   }{  \partial  u  } } \\  { \frac{   \partial  \varphi_3   }{  \partial  u  } }   \end{pmatrix}  ,  \begin{pmatrix}  { \frac{   \partial  \varphi_1   }{  \partial  u  } }  \\ { \frac{   \partial  \varphi_2   }{  \partial  u  } } \\  { \frac{   \partial  \varphi_3   }{  \partial  u  } }   \end{pmatrix}   \right\rangle ,\, F = \left\langle  \begin{pmatrix}  { \frac{   \partial  \varphi_1   }{  \partial  u  } }  \\ { \frac{   \partial  \varphi_2   }{  \partial  u  } } \\  { \frac{   \partial  \varphi_3   }{  \partial  u  } }   \end{pmatrix}   ,  \begin{pmatrix}  { \frac{   \partial  \varphi_1   }{  \partial  v  } }  \\ { \frac{   \partial  \varphi_2   }{  \partial  v  } } \\  { \frac{   \partial  \varphi_3   }{  \partial  v  } }   \end{pmatrix}   \right\rangle \text{ dan } G = \left\langle  \begin{pmatrix}  { \frac{   \partial  \varphi_1   }{  \partial  v  } }  \\ { \frac{   \partial  \varphi_2   }{  \partial  v  } } \\  { \frac{   \partial  \varphi_3   }{  \partial  v  } }   \end{pmatrix}   ,  \begin{pmatrix}  { \frac{   \partial  \varphi_1   }{  \partial  v  } }  \\ { \frac{   \partial  \varphi_2   }{  \partial  v  } } \\  { \frac{   \partial  \varphi_3   }{  \partial  v  } }   \end{pmatrix}   \right\rangle} { . }
}
\faktfolgerung {Maka pada $W$ berlaku

\mavergleichskettealign
{\vergleichskettealign
{ \varphi^*\left(  \omega \vert_U  \right)
}
{ =} { \sqrt{EG-F^2}  du \wedge dv
}
{ } {
}
{ } {
}
{ } {
}
}
{}
{}{.}}
\faktzusatz {}
\faktzusatz {}

}
{

Hal ini langsung mengikuti dari
Teorema 16.8.

}

\inputbemerkung
{}
{

Misalkan

\mavergleichskettedisp
{\vergleichskette
{ M
}
{ =} { \left\{ (u,v, \psi(u,v))  \mid  (u,v) \in W         \right\}
}
{ \subseteq} { W \times \R
}
{ } {
}
{ } {
}
}
{}{}{}
merupakan
\definitionsverweis {grafik}{}{}
dari
\maabbdisp {\psi} { W } {\R
} {,}
dengan

\mavergleichskette
{\vergleichskette
{W
}
{ \subseteq }{ \R^2
}
{  }{
}
{    }{
}
{   }{
}
}
{}{}{}
suatu himpunan bagian terbuka. Dalam kasus ini,
Korolari 17.1
dan
Korolari 17.2
berkaitan sebagai berikut. Turunan-turunan parsialnya adalah
\mathkor {} {\begin{pmatrix}  1  \\ 0 \\  { \frac{   \partial \psi }{  \partial u } }   \end{pmatrix}} {dan} {\begin{pmatrix}  0  \\ 1 \\  { \frac{   \partial \psi }{  \partial v } }   \end{pmatrix}} {.}
Oleh karena itu,

\mavergleichskette
{\vergleichskette
{ E
}
{ = }{ 1 + \left(  { \frac{   \partial \psi }{  \partial u } }  \right)^2
}
{  }{
}
{    }{
}
{   }{
}
}
{}{}{,}

\mavergleichskette
{\vergleichskette
{ F
}
{ = }{ { \frac{   \partial \psi }{  \partial u } } { \frac{   \partial \psi }{  \partial v } }
}
{  }{
}
{    }{
}
{   }{
}
}
{}{}{,}
dan

\mavergleichskette
{\vergleichskette
{ G
}
{ = }{ 1 + \left(  { \frac{   \partial \psi }{  \partial v } }  \right)^2
}
{  }{
}
{    }{
}
{   }{
}
}
{}{}{.}
Dengan demikian,

\mavergleichskettedisphandlinks
{\vergleichskettedisphandlinks
{EG-F^2
}
{ =} { \left(  1 + \left(  { \frac{   \partial \psi }{  \partial u } }  \right)^2  \right) \left(  1 + \left(  { \frac{   \partial \psi }{  \partial v } }  \right)^2  \right) - \left(  { \frac{   \partial \psi }{  \partial u } }  \right)^2 \left(  { \frac{   \partial \psi }{  \partial v } }  \right)^2
}
{ =} { 1  + \left(  { \frac{   \partial \psi }{  \partial u } }  \right)^2  + \left(  { \frac{   \partial \psi }{  \partial v } }  \right)^2
}
{ } {
}
{ } {
}
}
{}{}{.}

}

\inputbemerkung
{}
{

Kita melanjutkan dengan notasi dari
Korolari 17.2.
Jika himpunan bagian terbuka yang dicakup oleh
\mathkor {} {W} {dan} {\varphi} {}

\mavergleichskette
{\vergleichskette
{ U
}
{ \subseteq }{ M
}
{  }{
}
{    }{
}
{   }{
}
}
{}{}{}
mempunyai sifat bahwa
\definitionsverweis {komplemen}{}{}

\mathl{M \setminus U}{} merupakan suatu
\definitionsverweis {himpunan nol}{}{}
terhadap
\definitionsverweis {ukuran kanonik}{}{}
pada $M$, maka luas permukaan $M$ dapat dihitung hanya dengan memakai rumus untuk
\mathl{\varphi^*\left(  \omega \vert_U  \right)}{.} Hal ini, misalnya, berlaku jika
\mathl{M \setminus U}{} merupakan suatu
\definitionsverweis {submanifold tertutup}{}{}
dari $M$ yang berdimensi
\mathl{\leq 1}{;} lihat
Soal 15.14.
Dalam perhitungan, himpunan nol sering diabaikan secara tersirat.

}

  \zwischenueberschrift{Permukaan putar}

Misalkan diberikan suatu
\definitionsverweis {kurva terdiferensialkan}{}{}
\maabbeledisp {\gamma} {]a,b[} {\R^2
} {t} {(x(t),y(t))
} {,}
dengan

\mavergleichskette
{\vergleichskette
{ y(t)
}
{ \geq }{ 0
}
{  }{
}
{    }{
}
{   }{
}
}
{}{}{.}
Kita tertarik pada
\definitionsverweis {permukaan putar}{}{}
yang bersesuaian, yakni himpunan bagian

\mathdisp {\left\{  (x(t), y(t) \cos \alpha , y(t) \sin \alpha)   \mid   t \in ]a,b[  , \,  \alpha \in [0,2\pi]       \right\}} {  }

dari $\R^3$ yang diperoleh dengan memutar lintasan kurva tersebut mengelilingi sumbu $x$. Kita juga mensyaratkan bahwa $\gamma$ menghasilkan suatu difeomorfisme ke citranya

\mavergleichskette
{\vergleichskette
{ M
}
{ = }{ \gamma \left(  ]a,b[  \right)
}
{  }{
}
{    }{
}
{   }{
}
}
{}{}{}
dan bahwa $M$ merupakan submanifold tertutup dalam suatu himpunan terbuka

\mavergleichskettedisp
{\vergleichskette
{ G
}
{ \subseteq} {\R \times \R_+
}
{ } {
}
{ } {
}
{ } {
}
}
{}{}{}

\zusatzklammer {jadi juga disyaratkan bahwa komponen kedua $\gamma$ bernilai positif di mana-mana} {} {.}
Permukaan putarnya kemudian merupakan suatu submanifold tertutup dua dimensi dari $\R^3$ yang tidak berpotongan dengan sumbu $x$, sehingga kita memperoleh suatu manifold Riemann. Luas permukaannya dapat dihitung sebagai berikut.



\inputfaktbeweis
{Rotationsfläche zu ebener Kurve/Flächenberechnung/Fakt}
{Teorema}
{}
{

\faktsituation {Misalkan
\maabbeledisp {\gamma} {]a,b[} {\R^2
} { t } {(x(t),y(t))
} {,}
suatu
\definitionsverweis {kurva terdiferensialkan}{}{}
dengan

\mavergleichskette
{\vergleichskette
{ y(t)
}
{ > }{ 0
}
{  }{
}
{    }{
}
{   }{
}
}
{}{}{,}}
\faktvoraussetzung {yang menginduksi suatu
\definitionsverweis {difeomorfisme}{}{}
ke

\mavergleichskette
{\vergleichskette
{M
}
{ = }{ \gamma(]a,b[)
}
{  }{
}
{    }{
}
{   }{
}
}
{}{}{,}
dengan

\mavergleichskette
{\vergleichskette
{M
}
{ \subseteq }{ G
}
{  }{
}
{    }{
}
{   }{
}
}
{}{}{}
suatu
\definitionsverweis {submanifold tertutup}{}{}
satu dimensi dalam suatu himpunan terbuka

\mavergleichskette
{\vergleichskette
{G
}
{ \subseteq }{\R \times \R_+
}
{  }{
}
{    }{
}
{   }{
}
}
{}{}{.}}
\faktfolgerung {Maka
\definitionsverweis {permukaan putar}{}{}
yang bersesuaian merupakan suatu submanifold tertutup dari
\mathl{\R^3}{} yang tidak berpotongan dengan sumbu $x$, dan luas permukaannya sama dengan

\mathdisp {2 \pi \int_{  a  }^{  b  } \sqrt{ (x'(t))^2+ (y'(t))^2 } \cdot y(t) \, d t} { . }
}
\faktzusatz {}
\faktzusatz {}

}
{

Misalkan $S$ permukaan putar tersebut, yang merupakan submanifold tertutup dua dimensi dalam suatu himpunan terbuka di $\R^3$. Kita menerapkan
Korolari 17.2
pada parametrisasi
\maabbeledisp {} { ]a,b[ \times ]0, 2 \pi[} {S
} { (t, \alpha) } { (x(t),  y(t) \cos \alpha  , y(t) \sin \alpha  )
} {,}
Turunan-turunan parsialnya adalah

\mathdisp {\begin{pmatrix}  x'(t) \\ y'(t) \cos \alpha \\  y'(t) \sin \alpha   \end{pmatrix} \text{   dan } \begin{pmatrix}  0  \\ - y(t) \sin \alpha \\  y(t) \cos \alpha   \end{pmatrix}} {  }

dan karena itu

\mathdisp {E(t, \alpha) = (x'(t))^2 + (y'(t))^2,\, F(t, \alpha)=0,\, G(t, \alpha) = (y(t))^2} { . }

Dengan demikian, luas permukaannya sama dengan

\mavergleichskettedisphandlinks
{\vergleichskettedisphandlinks
{ \int_{  a  }^{  b  } \int_{  0  }^{  2 \pi } \sqrt{(x'(t))^2+(y'(t))^2} \cdot   y(t) \, d \alpha \, d t
}
{ =} { 2 \pi \int_{  a  }^{  b  } \sqrt{(x'(t))^2+(y'(t))^2} \cdot   y(t) \, d t
}
{ } {
}
{ } {
}
{ } {
}
}
{}{}{.}

}

  \zwischenueberschrift{Kartografi}

Kartografi
\zusatzklammer {abstrak} {} {}
mempelajari peta untuk permukaan suatu sfer.

\bild{ \begin{center}
\includegraphics[width=5.5cm]{ \bildeinlesungjpg {Cilinderprojectie-constructie} {jpg} }
\end{center}
\bildtext {} }

\bildlizenz { Cilinderprojectie-constructie.jpg } {} {KoenB} {Commons} {PD} {}

\inputbeispiel{}
{

Kita meninjau pemetaan
\maabbeledisp {\varphi} { [0, 2 \pi] \times [-1,1] } {\R^3
} {(u,v)} { \left(  \sqrt{1-v^2} \cos  u   , \,   \sqrt{1-v^2} \sin  u  , \,   v                 \right)
} {,}
yang
\definitionsverweis {citranya}{}{}
terletak pada \stichwort {sfer satuan} {.} Pemetaan ini kontinu pada persegi panjang tertutup tersebut. Pemetaan ini dapat dibayangkan dengan mula-mula membentuk persegi panjang itu menjadi permukaan tabung, kemudian memproyeksikan lingkaran-lingkaran tabung ke lingkaran-lingkaran horizontal pada sebuah sfer yang mempunyai tinggi sama. Pada interior persegi panjang, pemetaan ini
\definitionsverweis {terdiferensialkan}{}{}
dengan
\definitionsverweis {turunan parsial}{}{}

\mathdisp {{ \frac{   \partial  \varphi  }{  \partial  u  } } = \begin{pmatrix}  - \sqrt{1-v^2} \sin  u  \\ \sqrt{1-v^2} \cos  u \\  0   \end{pmatrix} \text{   dan } { \frac{   \partial  \varphi  }{  \partial  v  } } = \begin{pmatrix}  { \frac{   -v }{  \sqrt{1-v^2}  } } \cos  u   \\ { \frac{   -v }{  \sqrt{1-v^2}  } } \sin  u \\  1   \end{pmatrix}} { . }

\definitionsverweis {Pembatasan}{}{}
pemetaan ini pada persegi panjang terbuka bersifat
\definitionsverweis {injektif}{}{,}
dan citranya adalah sfer satuan kecuali sebuah busur setengah lingkaran bujur. Dengan demikian, kita dapat menghitung luas permukaan sfer memakai koordinat-koordinat ini. Dengan notasi yang digunakan dalam
Korolari 17.2,
kita memperoleh

\mavergleichskettedisp
{\vergleichskette
{E
}
{ =} { \left(  1-v^2  \right) \sin^{ 2  }  u  + \left(  1-v^2  \right) \cos^{ 2  }  u
}
{ =} { 1-v^2
}
{ } {
}
{ } {
}
}
{}{}{,}

\mavergleichskettedisp
{\vergleichskette
{F
}
{ =} { v \sin  u\cos  u  - v \sin  u\cos  u
}
{ =} { 0
}
{ } {
}
{ } {
}
}
{}{}{}
dan

\mavergleichskettedisp
{\vergleichskette
{G
}
{ =} { { \frac{  v^2 }{  1-v^2 } } \cos^{ 2  }  u + { \frac{  v^2 }{  1-v^2 } } \sin^{ 2  }  u +1
}
{ =} { { \frac{  v^2 }{  1-v^2 } }  +1
}
{ =} { { \frac{   1  }{  1-v^2 } }
}
{ } {
}
}
{}{}{.}
Oleh karena itu,

\mavergleichskettedisp
{\vergleichskette
{ \sqrt{EG-F^2 }
}
{ =} { 1
}
{ } {
}
{ } {
}
{ } {
}
}
{}{}{,}
artinya, pemetaan kartografis ini \stichwort {mempertahankan luas} {\zusatzfussnote {Akan tetapi, pemetaan ini tidak mempertahankan panjang. Ruas horizontal pada persegi panjang sangat dipendekkan ketika mendekati kutub. Sebagai gantinya, ruas vertikal semakin diregangkan ketika mendekati kutub, dan kedua gejala ini saling menetralkan} {.} {,}}
sehingga luas permukaan sfer sama dengan

\mavergleichskettealign
{\vergleichskettealign
{A
}
{ =} {
\int_{  ]0, 2 \pi[ \times ]- 1 ,  1 [   }   1  du \wedge dv
}
{ =} { \int_{  [0, 2 \pi] \times [- 1 ,  1 ]   }   1     \,  d  \lambda^2
}
{ =} { 2 \cdot 2 \pi
}
{ =} { 4 \pi
}
}
{}
{}{.}

}

\inputbeispiel{}
{

Kita meninjau pemetaan
\maabbeledisp {\varphi} { [0, 2 \pi] \times \R} {\R^3
} {(u,v)} { { \frac{   1  }{  \sqrt{1+v^2}  } }  ( \cos  u  , \sin  u  , v )
} {,}
yang
\definitionsverweis {citranya}{}{}
terletak pada \stichwort {sfer satuan} {.} Pemetaan ini dapat dibayangkan dengan mula-mula membentuk persegi panjang
\zusatzklammer {yang tak terbatas dalam satu arah} {} {}
\mathl{[0, 2 \pi] \times \R}{} menjadi selimut tabung tak hingga di atas lingkaran satuan, lalu memproyeksikan setiap titik selimut tabung itu ke sfer melalui garis yang menghubungkannya dengan pusat sfer. Di bawah pemetaan ini, semua titik permukaan sfer tercapai kecuali kutub utara dan kutub selatan. Selain itu, pemetaan tersebut bersifat injektif jika titik-titik ujung interval dibuang
\zusatzklammer {maka setengah lingkaran bujur tidak termasuk dalam citra} {} {.}
Pemetaan ini
\definitionsverweis {terdiferensialkan}{}{}
dengan
\definitionsverweis {turunan parsial}{}{}

\mathdisp {{ \frac{   \partial  \varphi  }{  \partial  u  } } = \begin{pmatrix}  { \frac{   - \sin  u  }{  \sqrt{1+v^2}  } }   \\ { \frac{   \cos  u  }{  \sqrt{1+v^2}  } }  \\  0   \end{pmatrix} = { \frac{   1  }{  \sqrt{1+v^2}  } } \begin{pmatrix}  - \sin  u  \\ \cos  u  \\  0   \end{pmatrix} \text{   dan } { \frac{   \partial  \varphi  }{  \partial  v  } } = \begin{pmatrix}  { \frac{   -v \cos  u  }{  \sqrt{1+v^2}^3  } }    \\ { \frac{   -v \sin  u  }{  \sqrt{1+v^2}^3  } }    \\  { \frac{   \sqrt{1+v^2 } - v^2 \left(  1+v^2  \right)^{- 1/2}   }{  1+v^2  } }   \end{pmatrix} = { \frac{   1  }{  \sqrt{1+v^2}^3  } } \begin{pmatrix}  -v \cos  u  \\ -v \sin  u  \\  1   \end{pmatrix}} { . }

Kita dapat menghitung luas permukaan sfer memakai koordinat-koordinat ini. Dengan notasi yang digunakan dalam
Korolari 17.2,
kita memperoleh

\mavergleichskettedisp
{\vergleichskette
{E
}
{ =} { { \frac{   1  }{  1+v^2 } }
}
{ } {
}
{ } {
}
{ } {
}
}
{}{}{,}

\mavergleichskettedisp
{\vergleichskette
{F
}
{ =} { 0
}
{ } {
}
{ } {
}
{ } {
}
}
{}{}{}
dan

\mavergleichskettedisp
{\vergleichskette
{G
}
{ =} { { \frac{   1  }{  \left(  1+v^2  \right)^3 } } \left(  v^2 \cos^{ 2  }  u + v^2 \sin^{ 2  }  u +1  \right)
}
{ =} { { \frac{   1  }{  \left(  1+v^2  \right)^2 } }
}
{ } {
}
{ } {
}
}
{}{}{.}
Oleh karena itu,

\mavergleichskettedisp
{\vergleichskette
{ \sqrt{EG-F^2 }
}
{ =} { \sqrt{ { \frac{   1  }{  \left(  1+v^2  \right)^3  } }  }
}
{ =} { { \frac{   1  }{  \sqrt{1+v^2}^3  } }
}
{ } {
}
{ } {
}
}
{}{}{.}
Dengan menggunakan
Teorema 12.10 (Teori Ukuran dan Integrasi (Osnabrück 2022-2023)),
luas permukaan sfer dengan demikian sama dengan

\mavergleichskettealign
{\vergleichskettealign
{A
}
{ =} {
\int_{  ]0, 2 \pi[ \times \R   }   { \frac{   1  }{  \sqrt{1+v^2}^3  } }  du \wedge dv
}
{ =} { \int_{ ]0, 2 \pi[ \times  \R   }   { \frac{   1  }{  \sqrt{1+v^2}^3  } }    \,  d  \lambda^2
}
{ =} { 2 \pi \int_{  - \infty }^{ + \infty } { \frac{   1  }{  \sqrt{1+v^2}^3  } } \, d  v
}
{ } {
}
}
{}
{}{.}
Menurut
Contoh 31.7 (Analisis (Osnabrück 2021-2023)),
integral tersebut sama dengan $2$, sehingga diperoleh luas permukaan
\mathl{4\pi}{.}

}
\stichwort {Proyeksi Mercator} {} berangkat dari proyeksi yang disebut terakhir dan mereparametrisasi koordinat vertikal sebagai $v=\sinh w,\;w\in\R$. Dengan demikian, koordinat lintang yang terbatas diwakili oleh koordinat vertikal tak terbatas, dan peta yang dihasilkan mempertahankan sudut.

\inputbeispiel{}
{

Kita meninjau pemetaan
\maabbeledisp {\varphi} {\R^2} {\R^3
} {(u,v)} {(\cos  u \cos  v  , \cos  u \sin  v  , \sin  u )
} {,}
yang
\definitionsverweis {citranya}{}{}
terletak pada \stichwort {sfer satuan} {.} Dalam istilah geografi, $u$ menyatakan \stichwort {lingkaran lintang} {} dan $v$ menyatakan \stichwort {lingkaran bujur} {} dari titik yang bersesuaian pada bumi satuan
\zusatzklammer {dalam \stichwort {koordinat geosentrik} {;} koordinat yang dipakai dalam geografi sedikit berbeda karena bumi bukan benar-benar sebuah sfer} {} {.}
Pemetaan ini
\definitionsverweis {terdiferensialkan}{}{}
dengan
\definitionsverweis {turunan parsial}{}{}

\mathdisp {{ \frac{   \partial  \varphi  }{  \partial  u  } } = \begin{pmatrix}  - \sin  u \cos  v  \\ -\sin  u \sin  v \\  \cos  u   \end{pmatrix} \text{   dan } { \frac{   \partial  \varphi  }{  \partial  v  } } = \begin{pmatrix}  - \cos  u \sin  v  \\ \cos  u \cos  v \\  0   \end{pmatrix}} { . }

\definitionsverweis {Pembatasan}{}{}
pemetaan ini pada persegi panjang terbuka

\mathdisp {{]{- { \frac{   \pi }{  2 } }}, { \frac{   \pi }{  2 } } [} \times {]{-\pi}, \pi[}} {  }

bersifat
\definitionsverweis {injektif}{}{,}
dan citranya adalah sfer satuan kecuali satu lingkaran bujur. Dengan demikian, kita dapat menghitung luas permukaan sfer memakai koordinat-koordinat ini. Dengan notasi yang digunakan dalam
Korolari 17.2,
kita memperoleh

\mavergleichskettedisp
{\vergleichskette
{E
}
{ =} { \sin^{ 2  }  u \cos^{ 2  }  v + \sin^{ 2  }  u \sin^{ 2  }  v + \cos^{ 2  }  u
}
{ =} { \sin^{ 2  }  u + \cos^{ 2  }  u
}
{ =} { 1
}
{ } {
}
}
{}{}{,}

\mavergleichskettedisp
{\vergleichskette
{F
}
{ =} { \sin  u\cos  u \sin  v \cos  v  -\sin  u\cos  u \sin  v \cos  v
}
{ =} { 0
}
{ } {
}
{ } {
}
}
{}{}{}
dan

\mavergleichskettedisp
{\vergleichskette
{G
}
{ =} { \cos^{ 2  }  u \sin^{ 2  }  v + \cos^{ 2  }  u \cos^{ 2  }  v
}
{ =} { \cos^{ 2  }  u
}
{ } {
}
{ } {
}
}
{}{}{.}
Oleh karena itu,

\mavergleichskettedisp
{\vergleichskette
{ \sqrt{EG-F^2 }
}
{ =} { \sqrt{ 1 \cdot \cos^{ 2  }  u  }
}
{ =} { \cos  u
}
{ } {
}
{ } {
}
}
{}{}{.}
Jadi, menurut
teorema Fubini,
luas permukaan sfer sama dengan

\mavergleichskettealign
{\vergleichskettealign
{A
}
{ =} {
\int_{  ]- { \frac{   \pi }{  2 } } , { \frac{   \pi }{  2 } } [ \times ]-\pi, \pi[   }   \cos  u  \,du \wedge dv
}
{ =} { \int_{  [- { \frac{   \pi }{  2 } }, { \frac{   \pi }{  2 } }] \times [-\pi, \pi]   }   \cos  u    \,  d  \lambda^2
}
{ =} { \int_{  -\pi }^{  \pi } \int_{  - { \frac{   \pi }{  2 } }  }^{  { \frac{   \pi }{  2 } }  } \cos  u \, d  u \, d v
}
{ =} { \int_{  -\pi }^{  \pi } 2 \, d v
}
}
{
\vergleichskettefortsetzungalign
{ =} { 4 \pi
}
{ } {}
{ } {}
{ } {}
}
{}{.}

}

\chapter{Lembar Kerja 17}
\setcounter{section}{17}

  \zwischenueberschrift{Soal latihan}

\inputaufgabe
{}
{

Misalkan
\maabbeledisp {\varphi} {\R^n } {\R
} { \left(  x_1   , \,   \ldots  , \,   x_n                 \right) } {a_1x_1 + \cdots + a_nx_n
} {,}
suatu
\definitionsverweis {bentuk linear}{}{.}
Misalkan $M$ adalah
\definitionsverweis {grafik}{}{}
fungsi ini, yang kita pandang sebagai
\definitionsverweis {manifold Riemann}{}{.}
Tunjukkan bahwa terdapat hubungan konstan antara volume himpunan-himpunan bagian yang saling bersesuaian di $\R^n$ dan pada grafik tersebut.

}
{} {}

\inputaufgabegibtloesung
{}
{

Hitung luas permukaan bola dengan menggunakan
Korolari 17.1.

}
{} {}

\inputaufgabe
{}
{

Bahas
\definitionsverweis {permukaan putar}{}{} $S$ yang diperoleh dengan memutar

\mathdisp {M= \left\{ ( \sin  { \frac{   1  }{  y } } ,y)  \mid   y> 0        \right\}} {  }

mengelilingi sumbu $x$, yaitu $A$. Apakah $S$ merupakan
\definitionsverweis {submanifold tertutup}{}{} dari
\mathl{\R^3 \setminus A}{?} Apakah himpunan $S$
\definitionsverweis {tertutup}{}{} di $\R^3$? Apakah
\definitionsverweis {penutupan}{}{} $S$ di $\R^3$ merupakan suatu manifold?

}
{} {}

\inputaufgabegibtloesung
{}
{

Tunjukkan bahwa luas permukaan putar yang diperoleh dengan memutar grafik

\mathdisp {\Gamma = \left\{ (x,e^x)  \mid   x \leq 0        \right\}} {  }

mengelilingi sumbu $x$ lebih kecil daripada $10$.

}
{} {}

\inputaufgabe
{}
{

Pastikan bahwa
\definitionsverweis {pemetaan-pemetaan}{}{}
yang diberikan dalam
Contoh 17.6,
Contoh 17.7,
dan
Contoh 17.8
memiliki
\definitionsverweis {citra}{}{}
yang terletak pada
\definitionsverweis {sfer satuan}{}{}
dan bersifat surjektif kecuali pada suatu
\definitionsverweis {himpunan nol}{}{.}

}
{} {}

\inputaufgabe
{}
{

Tentukan
\zusatzklammer {yang terdefinisi secara parsial} {} {}
\definitionsverweis {pemetaan-pemetaan invers}{}{}
dari
\definitionsverweis {pemetaan-pemetaan}{}{}
yang diberikan dalam
Contoh 17.6,
Contoh 17.7,
dan
Contoh 17.8.

}
{} {}

\inputaufgabe
{}
{

Tunjukkan bahwa
\definitionsverweis {lingkaran bujur}{}{}
dan
\definitionsverweis {lingkaran lintang}{}{}
pada bola bumi saling
\definitionsverweis {tegak lurus}{}{.}

}
{} {}

\inputaufgabe
{}
{

Berapakah panjang
\definitionsverweis {lingkaran lintang}{}{} pada lintang $30^\circ$ di Bumi
\zusatzklammer {gunakan $6370$ km sebagai jari-jari bumi} {} {?}

}
{} {}

\inputaufgabe
{}
{

Berapa persen permukaan Bumi yang pada suatu saat disinari Matahari
\zusatzklammer {karena jaraknya sangat jauh, sinar Matahari dianggap saling sejajar} {} {?}

}
{} {}

\inputaufgabe
{}
{

Tinjau elipsoid

\mathdisp {M= \left\{ (x,y,z)\in \R^3  \mid   \ \frac{x^2}{a^2}+\frac{y^2}{b^2}+\frac{z^2}{c^2} =1         \right\}} { , }

untuk $a,b,c\in\mathbb{R}_{+}$. Hitung luas permukaan M untuk
\mathl{a= b}{.} Apa yang terjadi ketika
\mathl{a\rightarrow c}{?}

}
{} {}

\inputaufgabe
{}
{

Tentukan
\definitionsverweis {infimum}{}{}
dan
\definitionsverweis {supremum}{}{}
dari
\definitionsverweis {panjang}{}{}
citra
\definitionsverweis {lingkaran-lingkaran besar}{}{}
pada peta yang diuraikan dalam
Contoh 17.6.

}
{} {}

\inputaufgabe
{}
{

Hubungkan
Korolari 17.1
dengan
Contoh 15.10
dengan menggunakan
Latihan 16.10.

}
{} {}

\inputaufgabe
{}
{

Misalkan
\maabbdisp {\alpha} {S^2 \setminus \{N\}} { \R^2
} {}
adalah
\definitionsverweis {proyeksi stereografis}{}{.}

a) Tentukan
\mathl{\alpha_* \omega}{} untuk
\definitionsverweis {bentuk luas kanonik}{}{}
pada $S^2$.

b) Gunakan
\mathl{\alpha_* \omega}{} untuk menghitung luas permukaan bola.

}
{} {}

  \zwischenueberschrift{Soal untuk dikumpulkan}

\inputaufgabe
{5}
{

Kita meninjau
\definitionsverweis {grafik}{}{}
$M$ dari fungsi
\maabbeledisp {\psi} {\R^2} {\R
} {(x,y)} { y+x^2
} {,}
sebagai
\definitionsverweis {manifold Riemann}{}{.}
Hitung luas permukaan grafik di atas persegi
\mathl{[-1,1] \times [-1,1]}{.}

}
{} {}

\inputaufgabe
{5}
{

Misalkan

\mathdisp {M= \left\{ (x,x^2)  \mid   x \in \R        \right\}  \subseteq \R^2} {  }

adalah
\definitionsverweis {parabola}{}{,}
yaitu
\definitionsverweis {grafik}{}{}
dari
\definitionsverweis {fungsi}{}{}
\maabbeledisp {} {\R} {\R
} { x } { x^2
} {.}
Tunjukkan bahwa
\definitionsverweis {permukaan putar}{}{}
terkait yang diperoleh dengan memutarnya mengelilingi sumbu $x$ bukan suatu
\definitionsverweis {manifold}{}{.}

}
{} {}

\inputaufgabe
{4}
{

Nyatakan suatu
\definitionsverweis {permukaan bola}{}{} sebagai
\definitionsverweis {permukaan putar}{}{} dan gunakan representasi ini untuk menghitung luas permukaan bola tersebut.

}
{} {}

\inputaufgabe
{4}
{

Nyatakan suatu
\definitionsverweis {torus}{}{} sebagai
\definitionsverweis {permukaan putar}{}{} dan gunakan representasi ini untuk menghitung luas permukaannya.

}
{} {}

\inputaufgabe
{6}
{

Tentukan \anfuehrung{jarak}{} antara Osnabrück dan Bangalore
\zusatzklammer {gunakan $6370$ km sebagai jari-jari bumi} {} {}
dalam kedua pengertian berikut.

a) Menyusuri permukaan bumi
\zusatzklammer {jarak garis udara} {} {.}

b) Menembus bumi
\zusatzklammer {jarak garis tikus tanah} {} {.}

}
{} {}

\inputaufgabe
{6}
{

Berapakah panjang citra
\definitionsverweis {lingkaran lintang}{}{} pada lintang $30^\circ$
pada
\definitionsverweis {peta-peta}{}{}
yang diuraikan dalam
Contoh 17.6,
Contoh 17.7,
dan
Contoh 17.8
\zusatzklammer {gunakan $6370$ km sebagai jari-jari bumi} {} {?}

}
{} {}

\section*{Solusi yang disediakan oleh sumber}
\addcontentsline{toc}{section}{Solusi yang disediakan oleh sumber}
\subsection*{Solusi Soal 17.2}
Hemisfer atas adalah grafik dari fungsi
\maabbeledisp {\psi} { B  \left(  0, 1 \right) } {\R
} {(u,v)} { \sqrt{1- u^2-v^2}
} {.}
Turunan parsial dari parametrisasi grafik ini adalah
\mathkor {} {\begin{pmatrix}  1  \\ 0\\  { \frac{   -u }{  \sqrt{1-u^2 -v^2}  } }   \end{pmatrix}} {dan} {\begin{pmatrix}  0  \\ 1\\  { \frac{   -v }{  \sqrt{1-u^2 -v^2}  } }   \end{pmatrix}} {}
dan karena itu, menurut
Korolari 17.1,
akar kuadrat dari

\mavergleichskettedisp
{\vergleichskette
{ 1 + { \frac{  u^2 }{  1-u^2-v^2 } }  + { \frac{  v^2 }{  1-u^2-v^2 } }
}
{ =} { { \frac{   1 }{  1-u^2-v^2 } }
}
{ } {
}
{ } {
}
{ } {
}
}
{}{}{}
merupakan faktor yang menentukan. Dengan demikian, luas hemisfer satuan adalah

\mathdisp {\int_B { \frac{   1 }{  \sqrt{1- u^2 -v^2}  } } d \lambda^2} { . }

Dengan substitusi
\mathkor {} {u= \sqrt{w} \cos \alpha} {dan} {v= \sqrt{w} \sin \alpha} {}
serta
rumus transformasi,
dengan menggunakan
latihan,
diperoleh

\mavergleichskettedisp
{\vergleichskette
{ \int_B { \frac{   1 }{  \sqrt{1- u^2-v^2}  } } d \lambda^2
}
{ =} { \int_{[0,2 \pi] \times [0,1] } { \frac{   1 }{  2  } }  \cdot { \frac{   1 }{  \sqrt{1- w}  } }  d \lambda^2
}
{ =} { \pi \cdot \int_0^1 { \frac{   1 }{  \sqrt{1- w}  } } dw
}
{ =} { \pi \cdot (  -2\sqrt{1-w})\vert_0^1
}
{ =} { 2 \pi
}
}
{}{}{.}
Jadi, luas permukaan bola adalah $4 \pi$.
\subsection*{Solusi Soal 17.4}
Kita meninjau permukaan untuk $s \leq x \leq 0$ dengan suatu bilangan negatif sebarang
\mathl{s}{.} Menurut rumus permukaan putar, luas permukaan tersebut sama dengan

\mathdisp {2 \pi \int_s^0    \sqrt{1 + e ^{2t} } e^t dt} { . }

Karena $t$ berada di daerah negatif, berlaku
\mathl{e^{2t} \leq 1}{} sehingga integrannya tidak melebihi $\sqrt{2} e^t$. Oleh karena itu, integral ini kurang dari atau sama dengan

\mavergleichskettedisp
{\vergleichskette
{ 2 \pi \sqrt{2}  \int_s^0 e^t dt
}
{ =} { 2 \pi \sqrt{2} (1- e^s )
}
{ \leq} { 2 \pi \sqrt{2}
}
{ \leq} { 2 \cdot 3{,}2 \cdot 1{,}5
}
{ =} { 2 \cdot 4{,}8
}
}
{
\vergleichskettefortsetzung
{ <} { 10
}
{ } {}
{ } {}
{ } {}
}{}{.}
Taksiran ini juga berlaku ketika $s \rightarrow - \infty$.

\part{Unit 18}
\chapter[Kuliah 18]{Kuliah 18: Isometri dan Permukaan Hiperbolik}
\setcounter{section}{18}

  \zwischenueberschrift{Rumus transformasi untuk bentuk volume}



\inputfaktbeweis
{Differenzierbare Mannigfaltigkeit/Positive Volumenform/Diffeomorphismus/Rückzug/Maß/Fakt}
{Lema}
{}
{

\faktsituation {Misalkan $M$ suatu
\definitionsverweis {manifold diferensiabel}{}{}
dengan
\definitionsverweis {basis topologi terhitung}{}{}
dan misalkan $\omega$ suatu
\definitionsverweis {bentuk volume positif}{}{}
pada $M$. Misalkan
\maabbdisp {\varphi} { L } { M
} {}
suatu
\definitionsverweis {difeomorfisme}{}{}
dengan manifold $L$, yang diberi orientasi tarikan balik melalui pemetaan tersebut, dan

\mavergleichskette
{\vergleichskette
{ S
}
{ \subseteq }{ L
}
{  }{
}
{    }{
}
{   }{
}
}
{}{}{}
suatu
\definitionsverweis {himpunan bagian terukur}{}{.}}
\faktfolgerung {Maka

\mavergleichskettedisp
{\vergleichskette
{ \int_S \varphi^* \omega
}
{ =} { \int_{ \varphi(S)} \omega
}
{ } {
}
{ } {
}
{ } {
}
}
{}{}{.}}
\faktzusatz {}
\faktzusatz {}

}
{

Menurut
Lema 15.2,
kita dapat mengasumsikan bahwa
\mathkor {} {S} {dan} {\varphi(S)} {}
masing-masing seluruhnya terletak di dalam sebuah peta berorientasi positif pada
\mathkor {} {L} {dan} {M} {;}
katakanlah

\mavergleichskette
{\vergleichskette
{ S
}
{ \subseteq }{ U
}
{  }{
}
{    }{
}
{   }{
}
}
{}{}{}
dengan peta
\maabbdisp {\alpha} { U } { U' \subseteq \R^n
} {}
dan

\mavergleichskette
{\vergleichskette
{ \varphi(S)
}
{ \subseteq }{ V
}
{  }{
}
{    }{
}
{   }{
}
}
{}{}{}
dengan peta
\maabbdisp {\beta} { V } { V' \subseteq \R^n
} {.}
Dengan memperkecil peta-peta itu lebih lanjut, kita dapat mengasumsikan bahwa
\mathkor {} {U} {dan} {V} {}
saling difeomorfik melalui $\varphi$. Dengan demikian, terdapat diagram komutatif difeomorfisme

\mathdisp {\begin{matrix}  U  & \stackrel{ \varphi }{\longrightarrow} &  V  &  \\ \!\!\!\!\! \alpha \downarrow & & \downarrow \beta \!\!\!\!\!  & \\  U'  & \stackrel{ \beta \circ \varphi \circ \alpha^{-1} }{\longrightarrow} &  V' & \! \! \! \!\!\!\!\!   \\ \end{matrix}} {  }

yang menginduksi diagram komutatif himpunan terukur

\mathdisp {\begin{matrix}  S  & \stackrel{ \varphi }{\longrightarrow} &  \varphi (S)  &  \\ \!\!\!\!\! \alpha \downarrow & & \downarrow \beta \!\!\!\!\!  & \\  \alpha (S) & \stackrel{ \beta \circ \varphi \circ \alpha^{-1} }{\longrightarrow} &  \beta ( \varphi (S) ) & \! \! \! \!\!\!\!\!   \\ \end{matrix}} {  }

Bagi bentuk volume $\omega$ pada $V$, berlaku

\mavergleichskettedisp
{\vergleichskette
{ \alpha^{ {-1}*} ( \varphi^* \omega)
}
{ =} { \left(  \beta \circ \varphi \circ \alpha^{-1}  \right)^*  ( \beta^{ {-1}*} \omega)
}
{ } {
}
{ } {
}
{ } {
}
}
{}{}{.}
menurut
Lema 14.7 (4).
Bentuk volume $\beta^{ {-1}*} \omega$ pada $V'$ mempunyai representasi
\mathl{g dy_1  \wedge \ldots \wedge dy_n}{} dengan suatu fungsi terukur
\maabb {g} {V'} {\R
} {}
dan, menurut definisi, $\int_{\varphi(S)} \omega$ sama dengan $\int_{\beta( \varphi(S))} g d \lambda^n$. Demikian pula,
\mathl{\alpha^{ {-1}*} ( \varphi^* \omega)}{} pada $U'$ mempunyai representasi berbentuk
\mathl{f dx_1 \wedge \ldots \wedge dx_n}{.} Bentuk volume pada $U'$ ini berimpit dengan penarikan balik
\mathl{g dy_1  \wedge \ldots \wedge dy_n}{} dan, menurut
Korolari 14.9,
penarikan balik tersebut sama dengan

\mathdisp {(g \circ \psi) \cdot \det  \left(  \left(  { \frac{   \partial   \psi_i   }{  \partial   x_j   } }  \right)_{1 \leq  i, j \leq n}  \right) dx_1 \wedge \ldots \wedge dx_n} {  }

\zusatzklammer {dengan

\mavergleichskettek
{\vergleichskettek
{ \psi
}
{ = }{ \beta \circ \varphi \circ \alpha^{-1}
}
{  }{
}
{    }{
}
{   }{
}
}
{}{}{}} {} {.}
Dengan demikian, pernyataan tersebut mengikuti dari
Teorema 14.3 (Teori Ukuran dan Integrasi (Osnabrück 2022-2023)).

}

  \zwischenueberschrift{Isometri antara manifold Riemann}

\inputdefinition
{}
{

Misalkan
\mathkor {} {L} {dan} {M} {}
merupakan
\definitionsverweis {manifold Riemann}{}{.}
Suatu
\definitionsverweis {pemetaan terdiferensialkan}{}{}
\maabb {\varphi} { L } { M
} {}
disebut
\definitionswort {isometri lokal}{,}
jika untuk setiap titik

\mavergleichskette
{\vergleichskette
{ P
}
{ \in }{ L
}
{  }{
}
{    }{
}
{   }{
}
}
{}{}{}
\definitionsverweis {pemetaan singgung}{}{}
\maabbdisp {T_P \varphi} {T_PL} { T_{\varphi(P) } M
} {}
merupakan suatu
\definitionsverweis {isometri}{}{}
terhadap hasil kali dalam yang diberikan.

}

Pemetaan
\maabbeledisp {} {\R} { \R^2
} { t } { \begin{pmatrix}  \cos  t  \\ \sin  t    \end{pmatrix}
} {,}
sudah menunjukkan bahwa suatu isometri lokal tidak harus injektif.

\inputdefinition
{}
{

Misalkan
\mathkor {} {L} {dan} {M} {}
merupakan
\definitionsverweis {manifold Riemann}{}{.}
Suatu
\definitionsverweis {pemetaan terdiferensialkan}{}{}
\maabb {\varphi} { L } { M
} {}
disebut
\definitionswort {isometri}{,}
jika pemetaan tersebut merupakan suatu
\definitionsverweis {difeomorfisme}{}{}
dan secara lokal merupakan isometri.

}

Manifold-manifold Riemann yang isometrik dianggap sama dalam geometri Riemann. Secara khusus, panjang, sudut, dan volume dipertahankan oleh isometri; lihat
Lema 18.6.
Akan tetapi, sama sekali tidak mudah menentukan sifat mana dari suatu manifold Riemann yang tertanam di dalam $\R^n$ yang hanya bergantung pada struktur Riemann dan bukan pada penanamannya. Sifat yang hanya bergantung pada struktur Riemann juga disebut intrinsik
\zusatzklammer {sedangkan sifat yang bergantung pada penanaman disebut ekstrinsik} {} {.}
Untuk sifat intrinsik, gambaran yang sesuai ialah menanyakan apakah sifat itu dapat dikenali dan dideskripsikan oleh suatu makhluk yang hanya boleh bergerak pada manifold tersebut dan tidak dapat meninggalkannya.

\inputbeispiel{}
{

Misalkan

\mavergleichskette
{\vergleichskette
{ M
}
{ \subseteq }{ \R^n
}
{  }{
}
{    }{
}
{   }{
}
}
{}{}{}
suatu submanifold Riemann tertutup. Maka setiap
\definitionsverweis {isometri linear}{}{}
dari $\R^n$ menginduksi suatu
\definitionsverweis {isometri}{}{}
dari $M$ ke $\varphi(M)$; simbol pada kodomain menyatakan isometri linear yang bersangkutan.

}

\inputbeispiel{}
{

Misalkan
\maabb {\gamma} {]a,b[} { \R^2
} {}
suatu kurva injektif
\definitionsverweis {terdiferensialkan}{}{}
dan
\definitionsverweis {berparametrisasi panjang busur}{}{}
dengan kurva citra

\mavergleichskettedisp
{\vergleichskette
{ C
}
{ \subseteq} { V
}
{ \subseteq} { \R^2
}
{ } {
}
{ } {
}
}
{}{}{,}
yang tertutup di dalam himpunan bagian terbuka $V$. Maka pemetaan
\maabbdisp {\varphi = \gamma \times
\operatorname{Id}_{ \R }} { ]a,b[ \times \R} { C \times \R
} {}
merupakan suatu
\definitionsverweis {isometri}{}{,}
dengan
\mathl{]a,b[ \times \R}{} membawa struktur Euklides alami dan $C \times \R$ membawa struktur submanifold Riemann

\mavergleichskette
{\vergleichskette
{ C \times \R
}
{ \subseteq }{ V \times \R
}
{ \subseteq }{ \R^3
}
{    }{
}
{   }{
}
}
{}{}{}
yang terinduksi. Pada dasarnya, pemetaan ini berarti bahwa selembar kertas dibentangkan mengikuti sebuah kurva dasar. Serat-serat vertikal tidak berubah, sedangkan serat-serat horizontal merupakan salinan kurva dasar. Contoh tak trivial yang paling sederhana ialah membengkokkan lembar tersebut menjadi selimut tabung bercelah. Secara intuitif jelas bahwa panjang kurva pada permukaan dan luas permukaan tidak berubah. Kita dapat menunjukkan bahwa pemetaan ini merupakan isometri sebagai berikut. Untuk itu, misalkan

\mavergleichskette
{\vergleichskette
{ P
}
{ = }{ (s,t)
}
{ \in }{ ]a,b[ \times \R
}
{    }{
}
{   }{
}
}
{}{}{.}
Maka

\mavergleichskettedisp
{\vergleichskette
{ \left(  D\varphi  \right)_{ P } \left(  e_1   \right)
}
{ =} { \begin{pmatrix}  \gamma_1'(s) \\ \gamma_2'(s)\\  0   \end{pmatrix}
}
{ } {
}
{ } {
}
{ } {
}
}
{}{}{}
dan

\mavergleichskettedisp
{\vergleichskette
{ \left(  D\varphi  \right)_{ P } \left(  e_2   \right)
}
{ =} { \begin{pmatrix}  0  \\ 0\\  1   \end{pmatrix}
}
{ } {
}
{ } {
}
{ } {
}
}
{}{}{,}
jadi kedua vektor citra dari vektor-vektor standar mempunyai panjang $1$ dan saling tegak lurus. Dengan demikian,
\maabbdisp {} {T_P(  ]a,b[ \times \R)} { T_{\varphi(P)} (C \times \R)
} {}
merupakan suatu isometri.

}



\inputfaktbeweis
{Riemannsche Mannigfaltigkeit/Isometrie/Grundlegende Eigenschaften/Fakt}
{Lema}
{}
{

\faktsituation {Misalkan $L,M$ merupakan
\definitionsverweis {manifold Riemann}{}{}
yang
\definitionsverweis {berorientasi}{}{}
dan misalkan
\maabb {\varphi} { L } { M
} {}
suatu
\definitionsverweis {isometri}{}{}
yang
\definitionsverweis {mempertahankan orientasi}{}{.}}
\faktuebergang {Maka pernyataan-pernyataan berikut berlaku.}
\faktfolgerung {\aufzaehlungvier{Bentuk volume kanonik
\definitionsverweis {ditarik balik}{}{}
dari
\definitionsverweis {bentuk volume kanonik}{}{}
pada $M$ menjadi bentuk volume kanonik pada $L$.
}{Untuk setiap kurva terdiferensialkan
\maabbdisp {\gamma} {[a,b]} { L
} {}
\definitionsverweis {panjang kurva}{}{}
$\gamma$ sama dengan panjang kurva $\varphi \circ \gamma$.
}{ $\varphi$
\definitionsverweis {mempertahankan ukuran}{}{.}
}{ $\varphi$
\definitionsverweis {konformal}{}{.}
}}
\faktzusatz {}
\faktzusatz {}

}
{

\aufzaehlungvier{Misalkan $\omega$ bentuk volume kanonik pada $M$ dan

\mavergleichskette
{\vergleichskette
{ P
}
{ \in }{ L
}
{  }{
}
{    }{
}
{   }{
}
}
{}{}{.}
Misalkan

\mavergleichskette
{\vergleichskette
{ u_1 , \ldots , u_n
}
{ \in }{ T_PL
}
{  }{
}
{    }{
}
{   }{
}
}
{}{}{}
suatu
\definitionsverweis {basis ortonormal}{}{}
dari $T_PL$ yang merepresentasikan orientasi $L$. Maka

\mavergleichskettedisp
{\vergleichskette
{ \left(  \varphi^* \omega  \right) (P, u_1 , \ldots , u_n)
}
{ =} { \omega( \varphi(P), T_P\varphi(u_1) , \ldots , T_P\varphi(u_n))
}
{ } {
}
{ } {
}
{ } {
}
}
{}{}{.}
Karena sifat isometri,
\mathl{T_P\varphi(u_1) , \ldots , T_P\varphi(u_n)}{} merupakan suatu basis ortonormal dari $T_{\varphi(P)} M$. Karena pemetaan tersebut mempertahankan orientasi, basis itu merepresentasikan orientasi pada $M$, sehingga nilainya sama dengan $1$. Sifat ini mengarakterisasi bentuk volume kanonik pada $L$.
}{Panjang kurva $\gamma$ adalah integral dari
\mathl{\Vert { \gamma'(t) } \Vert}{,} dengan norma diambil dalam
\mathl{T_{\gamma(t)}L}{.} Demikian pula, panjang kurva $\varphi \circ \gamma$ adalah integral dari
\mathl{\Vert { \left(  \varphi \circ \gamma  \right)'(t) } \Vert}{,} yang sama karena sifat isometri.
}{Untuk suatu himpunan bagian

\mavergleichskette
{\vergleichskette
{ S
}
{ \subseteq }{ L
}
{  }{
}
{    }{
}
{   }{
}
}
{}{}{}
ukuran yang berasal dari bentuk volume kanonik pada $L$, jika diterapkan pada $S$, menurut bagian (1) dan
Lema 18.1,
sama dengan

\mavergleichskettedisp
{\vergleichskette
{ \int_S \omega_L
}
{ =} { \int_S \varphi^*(\omega_M)
}
{ =} { \int_{\varphi(S)} \omega_M
}
{ } {
}
{ } {
}
}
{}{}{.}
}{Hal ini langsung mengikuti dari sifat isometri.
}

}



\inputfaktbeweis
{Riemannsche Mannigfaltigkeit/Karte/Isometrie/Fakt}
{Lema}
{}
{

\faktsituation {Misalkan $M$ suatu
\definitionsverweis {manifold Riemann}{}{}
dan misalkan
\maabbdisp {\alpha} { U } { V \subseteq \R^n
} {}
suatu
\definitionsverweis {peta}{}{}
pada $M$.}
\faktfolgerung {Maka $\alpha$ merupakan suatu
\definitionsverweis {isometri}{}{}
antara
\mathkor {} {U} {dan} {V} {,}
jika $V$ dilengkapi dengan
\definitionsverweis {hasil kali dalam}{}{}
yang ditentukan oleh
\definitionsverweis {fungsi-fungsi fundamental metrik}{}{}
$g_{ij}$.}
\faktzusatz {}
\faktzusatz {}

}
{

Hal ini langsung mengikuti dari definisi $g_{ij}$ melalui

\mavergleichskettedisp
{\vergleichskette
{ g_{ij}(Q)
}
{ =} { \left\langle  T(\alpha^{-1})(e_i)  ,  T(\alpha^{-1}) (e_j)    \right\rangle_{ \alpha^{-1}(Q) }
}
{ } {
}
{ } {
}
{ } {
}
}
{}{}{.}

}

  \zwischenueberschrift{Permukaan hiperbolik}

Kita membahas tiga model untuk apa yang disebut permukaan hiperbolik
\zusatzklammer {atau bidang hiperbolik} {} {}
dan memberikan
\definitionsverweis {isometri}{}{}
di antara model-model tersebut. Dalam
Contoh 29.10,
kita akan melihat bahwa ini merupakan permukaan dengan
\definitionsverweis {kelengkungan seksional}{}{}
konstan $-1$. Model berikut disebut setengah bidang Poincaré.

\inputbeispiel{}
{

Misalkan

\mavergleichskettedisp
{\vergleichskette
{ H = {\mathbb H}
}
{ =} { \R \times \R_+
}
{ =} { \left\{ (x,y)  \mid   y > 0        \right\}
}
{ } {
}
{ } {
}
}
{}{}{}
merupakan setengah bidang atas yang dilengkapi dengan
\definitionsverweis {metrik Riemann}{}{,}
yang diberikan oleh

\mavergleichskettedisp
{\vergleichskette
{ g_{11}
}
{ =} { g_{22}
}
{ =} { { \frac{   1  }{  y^2 } }
}
{ } {
}
{ } {
}
}
{}{}{}
dan

\mavergleichskette
{\vergleichskette
{ g_{12}
}
{ = }{ 0
}
{  }{
}
{    }{
}
{   }{
}
}
{}{}{.}
Jadi, hasil kali dalam standar diskalakan dengan fungsi ${ \frac{   1  }{  y^2 } }$. Norma suatu vektor singgung pada titik
\mathl{(x,y)}{} adalah norma standar yang diskalakan dengan faktor ${ \frac{   1  }{  \betrag {  y  } } }$. Dengan demikian, bentuk luas diberikan oleh kerapatan ${ \frac{   1  }{  y^2 } }$.

}

Model berikut disebut \stichwort {cakram hiperbolik} {.}

\inputbeispiel{}
{

Pada cakram terbuka

\mavergleichskettedisp
{\vergleichskette
{ V
}
{ =} { \left\{ (x,y)  \mid   x^2+y^2 < 1        \right\}
}
{ } {
}
{ } {
}
{ } {
}
}
{}{}{}
kita mendefinisikan suatu
\definitionsverweis {struktur Riemann}{}{}
melalui bentuk bilinear

\mavergleichskettedisp
{\vergleichskette
{ g(Q) (u,v)
}
{ =} { { \frac{   4  }{  \left(  1-  \Vert { Q } \Vert^2   \right)^2  } }   \left\langle  u  , v  \right\rangle
}
{ } {
}
{ } {
}
{ } {
}
}
{}{}{,}
jadi hasil kali dalam standar diskalakan dengan fungsi
\mathl{{ \frac{   4  }{  \left(  1 - \Vert { Q } \Vert^2   \right)^2  } }}{,} dan matriks fundamental pertamanya adalah

\mathdisp {{ \frac{   4  }{  \left(  1- \Vert { Q } \Vert^2   \right)^2  } } \begin{pmatrix}  1   &   0   \\  0 &  1 \end{pmatrix}} { . }

}

\bild{ \begin{center}
\includegraphics[width=5.5cm]{ \bildeinlesunggif {Poincarehalfplaneconform} {gif} }
\end{center}
\bildtext {Animasi transformasi konformal setengah bidang Poincaré menjadi cakram satuan melalui inversi terhadap lingkaran yang berpusat di titik bergerak I. Edisi PDF menampilkan bingkai pertama yang statis; animasi asli dipertahankan untuk edisi HTML dan unduhan.} }

\bildlizenz { Poincarehalfplaneconform.gif } {} {HB} {Commons} {CC-by-sa 3.0} {}



\inputfaktbeweis
{Hyperbolische Fläche/Kreisscheibe/Halbebene/Isometrie/Fakt}
{Lema}
{}
{

\faktsituation {Pemetaan
\maabbeledisp {\psi} { V } { H
} { w } { { \frac{   \left(  w +1  \right)  { \mathrm i}  }{  1 - w } }
} {,}}
\faktfolgerung {mendefinisikan suatu
\definitionsverweis {isometri}{}{}
antara
\definitionsverweis {cakram hiperbolik}{}{}
dan
\definitionsverweis {setengah bidang Poincaré}{}{.}}
\faktzusatz {}
\faktzusatz {}

}
{

Menurut
Soal 18.10,
pemetaan ini bijektif dan terdiferensialkan kompleks dengan invers yang terdiferensialkan kompleks, sehingga khususnya terdapat suatu difeomorfisme antara manifold-manifold diferensiabel. Menurut
Soal 18.12,
pemetaan ini dalam koordinat real diberikan oleh

\mavergleichskettedisp
{\vergleichskette
{ \psi \begin{pmatrix}  a  \\b   \end{pmatrix}
}
{ =} { { \frac{   1  }{  -2a+1+a^2+b^2 } }  \begin{pmatrix}  -2b \\ 1-a^2-b^2   \end{pmatrix}
}
{ } {
}
{ } {
}
{ } {
}
}
{}{}{}
Menurut
Soal 18.13,
turunan-turunan parsialnya adalah

\mavergleichskettedisp
{\vergleichskette
{ { \frac{   \partial   \psi  }{  \partial  a  } }
}
{ =} { { \frac{   1  }{  \left(  -2a+1+a^2+b^2  \right)^2  } } \begin{pmatrix}  4ab -4b  \\ 2a^2 -4a+2 -2 b^2     \end{pmatrix}
}
{ } {
}
{ } {}
{ } {}
}
{}{}{}
dan

\mavergleichskettedisp
{\vergleichskette
{ { \frac{   \partial   \psi  }{  \partial  b  } }
}
{ =} { { \frac{   1  }{  \left(  -2a+1+a^2+b^2  \right)^2  } } \begin{pmatrix}  4a  -2-2a^2+2b^2   \\ 4ab -4b     \end{pmatrix}
}
{ } {
}
{ } {
}
{ } {}
}
{}{}{.}
Kita harus menunjukkan bahwa

\mavergleichskettedisp
{\vergleichskette
{ \left\langle e_i  , e_j   \right\rangle_{V, (a,b)}
}
{ =} { \left\langle  \left(  D \psi  \right)_{(a,b)} \left(  e_i   \right)   ,  \left(  D \psi  \right)_{(a,b)} \left(  e_j   \right)   \right\rangle_{H, \psi (a,b)}
}
{ } {
}
{ } {
}
{ } {
}
}
{}{}{.}
Karena kedua turunan parsial saling tegak lurus dan karena kedua metrik Riemann mewarisi ortogonalitas hasil kali dalam standar, hal ini hanya perlu ditunjukkan untuk

\mavergleichskette
{\vergleichskette
{ i
}
{ = }{ j
}
{  }{
}
{    }{
}
{   }{
}
}
{}{}{.}
Jadi, kita harus menunjukkan bahwa diferensial total mempertahankan panjang. Kita memperoleh

\mavergleichskettealigndrucklinks
{\vergleichskettealigndrucklinks
{ \Vert { { \frac{   \partial   \psi  }{  \partial   a   } } ( \begin{pmatrix}  a  \\b   \end{pmatrix} ) } \Vert_{H, \psi(a,b) }
}
{ =} { \Vert { { \frac{   1  }{  \left(  -2a+1+a^2+b^2  \right)^2  } } \begin{pmatrix}  4ab -4b  \\ 2a^2 -4a+2 -2 b^2    \end{pmatrix} } \Vert_{H, \psi(a,b) }
}
{ =} { { \frac{   \left(  -2a+1+a^2+b^2  \right)  }{  \left(  1-a^2-b^2  \right)  } } \Vert { { \frac{   2  }{  \left(  -2a+1+a^2+b^2  \right)^2  } } \begin{pmatrix}  2(a-1)b  \\ \left(  a-1  \right)^2 - b^2    \end{pmatrix} } \Vert
}
{ =} { { \frac{   2  }{  \left(  1-a^2-b^2  \right) \left(  \left(  a-1  \right)^2 +b^2  \right)  } } \sqrt{ 4(a-1)^2b^2 + \left(  a-1  \right)^4 + b^4 -2\left(  a-1  \right)^2  b^2 }
}
{ =} { { \frac{   2  }{  \left(  1-a^2-b^2  \right) \left(  \left(  a-1  \right)^2 +b^2  \right)  } } \left(  \left(  a-1  \right)^2 + b^2   \right)
}
}
{
\vergleichskettefortsetzungalign
{ =} { { \frac{   2  }{  \left(  1- \left(  a^2+b^2  \right)  \right)  } }
}
{ =} { \Vert {e_1 } \Vert_{V, (a,b)}
}
{ } {}
{ } {}
}{}{.}
Dengan cara yang sama, pernyataan tersebut diperoleh untuk turunan parsial kedua.

}

Metrik Euklides standar pada $\R^n$ menginduksi suatu metrik Riemann pada setiap submanifold tertutup. Akan tetapi, terdapat pula situasi ketika $\R^n$ dilengkapi dengan suatu
\definitionsverweis {bentuk bilinear}{}{}
tak terdegenerasi yang berbeda, yaitu tidak
\definitionsverweis {definit positif}{}{,}
namun pembatasannya pada suatu submanifold tertutup tetap menghasilkan manifold Riemann. Untuk contoh berikut, yaitu hiperboloid dua lembar, lihat juga Kurs:Lineare Algebra (Osnabrück 2017-2018)/Teil II/Vorlesung 40.

\inputbeispiel{}
{

\bild{ \begin{center}
\includegraphics[width=5.5cm]{ \bildeinlesungpng {Hyperboloid2} {png} }
\end{center}
\bildtext {Yang dibahas ialah lembar atas hiperboloid dua lembar,\\
yang dilengkapi dengan bentuk bilinear terinduksi dari bentuk Minkowski.} }

\bildlizenz { Hyperboloid2.png } {} {RokerHRO} {Commons} {gemeinfrei} {}

Kita melengkapi $\R^3$ dengan
\definitionsverweis {bentuk Minkowski standar}{}{,}
yakni bentuk bilinear yang bersesuaian dengan bentuk kuadratik $x^2+y^2-z^2$, dan kita meninjau himpunan bagian

\mavergleichskettedisp
{\vergleichskette
{ Y
}
{ =} { \left\{ (x,y,z)  \mid   x^2 +y^2-z^2 = -1  , \,  z > 0       \right\}
}
{ } {
}
{ } {
}
{ } {
}
}
{}{}{.}
Jika $\R^3$ dengan bentuk Minkowski dipandang sebagai model tiga dimensi untuk relativitas khusus, maka ini merupakan himpunan vektor pengamat berarah masa depan. Untuk suatu titik

\mavergleichskette
{\vergleichskette
{ P
}
{ = }{ (x,y,z)
}
{ \in }{ Y
}
{    }{
}
{   }{
}
}
{}{}{}
menurut
Lema 40.4 (Aljabar Linear (Osnabrück 2024-2025)),
pembatasan bentuk tersebut pada ruang singgung bersifat definit positif. Vektor-vektor
\mathkor {} {\begin{pmatrix}  z  \\ 0\\  x   \end{pmatrix}} {dan} {\begin{pmatrix}  0  \\ z \\  y   \end{pmatrix}} {}
membentuk suatu basis dari
\definitionsverweis {ruang singgung}{}{}
$T_PY$. Untuk sembarang vektor singgung

\mavergleichskettedisp
{\vergleichskette
{ v
}
{ =} { a \begin{pmatrix}  z  \\ 0\\  x   \end{pmatrix} + b \begin{pmatrix}  0  \\ z\\  y   \end{pmatrix}
}
{ =} { \begin{pmatrix} az \\ bz\\ ax+by   \end{pmatrix}
}
{ } {
}
{ } {
}
}
{}{}{}
berlaku

\mavergleichskettealignhandlinks
{\vergleichskettealignhandlinks
{ \left\langle  v  , v  \right\rangle_{Y,P}
}
{ =} { {  \begin{pmatrix} az \\ bz\\ ax+by \end{pmatrix} ^{ \text{tr} } } \begin{pmatrix}  1   &   0  &  0 \\  0  &  1  &  0  \\ 0 &  0  &  -1 \end{pmatrix} \begin{pmatrix} az \\ bz\\ ax+by \end{pmatrix}
}
{ =} { {  \begin{pmatrix} az \\ bz\\ ax+by \end{pmatrix} ^{ \text{tr} } }\begin{pmatrix} az \\ bz\\  -ax-by \end{pmatrix}
}
{ =} { a^2z^2+b^2z^2-\left(  ax+by  \right)^2
}
{ =} { a^2+b^2+\left(  ay-bx  \right)^2
}
}
{}
{}{,}
yang positif setiap kali kedua koefisien tidak sekaligus nol.

}



\inputfaktbeweis
{Hyperbolische Fläche/Kreisscheibe/Zweischaliges Hyperboloid/Isometrie/Fakt}
{Lema}
{}
{

\faktsituation {Pemetaan
\maabbeledisp {\theta} { V } { Y
} { \begin{pmatrix}  a  \\b   \end{pmatrix} } { \begin{pmatrix}  { \frac{   -2a }{ a^2+b^2-1  } }  \\ { \frac{   -2b }{ a^2+b^2-1  } } \\  { \frac{   -a^2-b^2-1  }{ a^2+b^2-1  } }   \end{pmatrix}
} {,}}
\faktfolgerung {mendefinisikan suatu
\definitionsverweis {isometri}{}{}
antara
\definitionsverweis {cakram hiperbolik}{}{}
$V$ dan lembar atas
\definitionsverweis {hiperboloid dua lembar}{}{}
dengan
\definitionsverweis {metrik Minkowski}{}{}
terinduksi.}
\faktzusatz {}
\faktzusatz {}

}
{

Turunan parsial dari $\theta$ adalah

\mavergleichskettealignhandlinks
{\vergleichskettealignhandlinks
{ { \frac{   \partial  \theta  }{  \partial   a   } } ( \begin{pmatrix}  a  \\b   \end{pmatrix} )
}
{ =} { { \frac{   1  }{  \left(  a^2+b^2-1  \right)^2  } } \begin{pmatrix}  -2 \left(  a^2+b^2-1  \right)+ 2a (2a)  \\ 4ab\\  -2a \left(  a^2+b^2-1  \right) +2a  \left(  a^2+b^2+1  \right)   \end{pmatrix}
}
{ =} { { \frac{   1  }{  \left(  a^2+b^2-1  \right)^2  } } \begin{pmatrix}  2a^2-2b^2+2  \\ 4ab\\  4a   \end{pmatrix}
}
{ } {
}
{ } {
}
}
{}
{}{}
dan

\mavergleichskettedisp
{\vergleichskette
{ { \frac{   \partial  \theta  }{  \partial   b   } } ( \begin{pmatrix}  a  \\b   \end{pmatrix} )
}
{ =} { { \frac{   1  }{  \left(  a^2+b^2-1  \right)^2  } } \begin{pmatrix}  4ab \\ - 2a^2+2b^2+2 \\  4b   \end{pmatrix}
}
{ } {
}
{ } {
}
{ } {}
}
{}{}{.}
Kita memperoleh

\mavergleichskettealignhandlinks
{\vergleichskettealignhandlinks
{ \left\langle  \begin{pmatrix}  2a^2-2b^2+2  \\ 4ab\\  4a   \end{pmatrix}   ,  \begin{pmatrix}  4ab \\ - 2a^2+2b^2+2 \\  4b   \end{pmatrix}     \right\rangle_{Y}
}
{ =} { 4 ab \left(  2a^2-2b^2+2   \right) +  4 ab \left(  - 2a^2 +2b^2+2   \right)      -16ab
}
{ =} { 0
}
{ } {
}
{ } {
}
}
{}
{}{,}

\mavergleichskettealignhandlinks
{\vergleichskettealignhandlinks
{ \left\langle  \begin{pmatrix}  2a^2-2b^2+2  \\ 4ab\\  4a   \end{pmatrix}   ,  \begin{pmatrix}  2a^2-2b^2+2  \\ 4ab\\  4a   \end{pmatrix}      \right\rangle_{Y}
}
{ =} { \left(  2a^2-2b^2+2   \right)^2  +  16 a^2b^2 -16a^2
}
{ =} { 4  \left(  \left(  a^2-b^2+1   \right)^2  +  4 a^2b^2 -4a^2  \right)
}
{ =} { 4 \left(  a^4+b^4+1 -2a^2+2a^2b^2 -2b^2  \right)
}
{ =} { 4  \left(  a^2 +b^2-1  \right)^2
}
}
{}
{}{,}
dan

\mavergleichskettedisphandlinks
{\vergleichskettedisphandlinks
{ \left\langle  \begin{pmatrix}  4ab \\ - 2a^2+2b^2+2 \\  4b   \end{pmatrix}   ,  \begin{pmatrix}  4ab \\ - 2a^2+2b^2+2 \\  4b   \end{pmatrix}     \right\rangle_{Y}
}
{ =} { 4  \left(  a^2 +b^2-1  \right)^2
}
{ } {}
{ } {}
{ } {}
}
{}{}{.}
Dengan memperhitungkan faktor awal ${ \frac{   1  }{  \left(  a^2+b^2-1  \right)^2  } }$, nilai hasil kali dalam pada $Y$ berimpit dengan hasil kali dalam pada cakram.

}

\chapter{Lembar Kerja 18}
\setcounter{section}{18}

  \zwischenueberschrift{Soal latihan}

\inputaufgabe
{}
{

Misalkan $L,M,N$ adalah
\definitionsverweis {manifold-manifold Riemann}{}{.}
Tunjukkan sifat-sifat berikut.
\aufzaehlungdrei{Identitas
\maabbdisp {\operatorname{Id}_{  M  }} { M } { M
} {}
adalah suatu
\definitionsverweis {isometri}{}{.}
}{Jika
\maabb {\varphi} { L } { M
} {}
adalah suatu isometri, maka
\definitionsverweis {pemetaan invers}{}{}
\maabb {\varphi^{-1}} { M } { L
} {}
juga merupakan suatu isometri.
}{Jika
\maabb {\varphi} { L } { M
} {}
dan
\maabb {\psi} { M } { N
} {}
adalah isometri, maka $\psi \circ \varphi$ juga merupakan suatu isometri.
}

}
{} {}

\inputaufgabe
{}
{

Misalkan $M$ adalah suatu
\definitionsverweis {manifold Riemann}{}{.}
Tunjukkan bahwa himpunan semua
\definitionsverweis {isometri}{}{}
pada $M$ membentuk suatu
\definitionsverweis {grup}{}{.}

}
{} {}

\inputaufgabe
{}
{

Misalkan
\maabbeledisp {\varphi} {\R} {S^1 \subseteq \R^2
} { t } { \begin{pmatrix}  \cos  t  \\ \sin  t    \end{pmatrix}
} {.}
Tunjukkan bahwa $\varphi$ adalah suatu
\definitionsverweis {isometri lokal}{}{}
antara
\definitionsverweis {manifold-manifold Riemann}{}{}
jika $\R$ dan $\R^2$ dilengkapi dengan
\definitionsverweis {struktur Euklides}{}{}
dan $S^1$ dilengkapi dengan struktur Riemann terinduksi.

}
{} {}

\inputaufgabe
{}
{

Misalkan
\maabb {\varphi} {\R^n } { \R^n
} {}
adalah suatu
\definitionsverweis {pemetaan linear}{}{,}
dan misalkan $\R^n$ dilengkapi dengan
\definitionsverweis {struktur Euklides}{}{}
beserta
\definitionsverweis {struktur Riemann}{}{}
yang diinduksikannya. Tunjukkan bahwa $\varphi$ merupakan suatu
\definitionsverweis {isometri linear}{}{}
jika dan hanya jika $\varphi$ merupakan suatu
\definitionsverweis {isometri}{}{}
terhadap struktur Riemann tersebut.

}
{} {}

\inputaufgabe
{}
{

Misalkan

\mavergleichskette
{\vergleichskette
{ I
}
{ \subseteq }{ \R
}
{  }{
}
{    }{
}
{   }{
}
}
{}{}{}
adalah suatu
\definitionsverweis {interval terbuka}{}{}
dan
\maabbdisp {\gamma} { I } { \R^n
} {}
adalah suatu
\definitionsverweis {kurva yang terdiferensialkan secara kontinu}{}{}
yang injektif. Kita memandang $I$ maupun $\R^n$ sebagai
\definitionsverweis {manifold Riemann}{}{}
melalui
\definitionsverweis {struktur Euklides}{}{}
masing-masing. Andaikan bahwa

\mavergleichskette
{\vergleichskette
{ \gamma(I)
}
{ \subseteq }{ U
}
{ \subseteq }{ \R^n
}
{    }{
}
{   }{
}
}
{}{}{}
dengan $U$
\definitionsverweis {terbuka}{}{}
merupakan suatu
\definitionsverweis {submanifold tertutup}{}{.}
Tunjukkan bahwa pemetaan

\mavergleichskettedisp
{\vergleichskette
{ \gamma
}
{ :} { I
}
{ \longrightarrow} { \gamma(I) \subseteq \R^n
}
{ } {
}
{ } {
}
}
{}{}{}
merupakan suatu
\definitionsverweis {isometri}{}{}
jika dan hanya jika $\gamma$
\definitionsverweis {berparametrisasi panjang busur}{}{.}

}
{} {}

\inputaufgabe
{}
{

Misalkan $L_1,L_2,M_1,M_2$ adalah
\definitionsverweis {manifold-manifold Riemann}{}{}
dan misalkan
\maabb {\varphi_1} {L_1 } { M_1
} {}
serta
\maabb {\varphi_2} {L_2 } { M_2
} {}
adalah
\definitionsverweis {isometri}{}{.}
Tunjukkan bahwa
\maabbdisp {\varphi_1 \times \varphi_2} {L_1 \times L_2 } { M_1 \times M_2
} {}
juga merupakan suatu isometri.

}
{} {}

\inputaufgabe
{}
{

Kita meninjau
\definitionsverweis {difeomorfisme}{}{}
\maabbeledisp {\varphi} {S^1 \times \R_+} { \Complex \setminus \{0\}  =  \R^2 \setminus \{(0,0)\}
} {(s,t) } { st
} {,}
dengan silinder $S^1 \times \R_+$ dilengkapi
\definitionsverweis {struktur produk Riemann}{}{.}
Uraikan struktur Riemann pada
\mathl{\R^2 \setminus \{(0,0)\}}{}
yang membuat $\varphi$ menjadi suatu
\definitionsverweis {isometri}{}{.}

}
{} {}

\inputaufgabegibtloesung
{}
{

Misalkan
\mathkor {} {r} {dan} {R} {}
adalah bilangan real positif dengan

\mavergleichskette
{\vergleichskette
{ 0
}
{ < }{ r
}
{ < }{ R
}
{    }{
}
{   }{
}
}
{}{}{,}
dan kita meninjau torus
\zusatzklammer {tertanam} {} {}

\mavergleichskettedisp
{\vergleichskette
{ T
}
{ =} { \left\{ (x,y,z) \in \R^3  \mid   \left(  \sqrt{x^2+y^2} -R  \right)^2 +z^2  = r^2         \right\}
}
{ \subseteq} { \R^3
}
{ } {
}
{ } {
}
}
{}{}{}
dengan
\definitionsverweis {struktur Riemann terinduksi}{}{.}
Tunjukkan bahwa untuk setiap

\mavergleichskette
{\vergleichskette
{ \alpha
}
{ \in }{ [0,2 \pi [
}
{  }{
}
{    }{
}
{   }{
}
}
{}{}{}
pemetaan
\maabbeledisp {\Psi_\alpha} {\R^3} {\R^3
} { \left(  x   , \,   y  , \,  z                 \right) } { \left(  x \cos \alpha - y \sin  \alpha   , \,   x \sin \alpha + y \cos  \alpha  , \,  z                 \right)
} {,}
menginduksi suatu
\definitionsverweis {isometri}{}{}
dari $T$ ke dirinya sendiri.

}
{} {}

\inputaufgabe
{}
{

Hitung panjang lintasan linear dari
\mathl{\begin{pmatrix}  0  \\ 1   \end{pmatrix}}{} ke
\mathl{\begin{pmatrix}  1  \\ 1   \end{pmatrix}}{} dalam
\definitionsverweis {setengah bidang Poincaré}{}{.}

}
{} {}

\inputaufgabe
{}
{

Hitung luas persegi panjang dalam
\definitionsverweis {setengah bidang Poincaré}{}{}
yang diberikan oleh titik-titik sudut
\mathl{(1,1),\, (5,1),\, (1,3),\, (5,3)}{}{.}

}
{} {}

\inputaufgabegibtloesung
{}
{

Tunjukkan bahwa
\definitionsverweis {pemetaan-pemetaan}{}{}
\maabbeledisp {\varphi} { {\mathbb H} } { U \left(  0,1  \right)
} { z} { { \frac{  z-  { \mathrm i}  }{ z + { \mathrm i}  } }
} {,}
dan
\maabbeledisp {\psi} {  U \left(  0,1  \right)  } { {\mathbb H}
} {w} { { \frac{   (w+1)  { \mathrm i}  }{ 1-w  } }
} {,}
memberikan suatu
\definitionsverweis {pemetaan biholomorfik}{}{}
antara
\definitionsverweis {setengah bidang atas}{}{}
${\mathbb H}$ dan cakram satuan terbuka $U \left(  0,1  \right)$.

}
{} {}

\inputaufgabe
{}
{

Tentukan
\definitionsverweis {turunan}{}{}
dari
\definitionsverweis {pemetaan-pemetaan}{}{}
\maabbeledisp {\varphi} { {\mathbb H} } { U \left(  0,1  \right)
} { z} { { \frac{  z-  { \mathrm i}  }{ z + { \mathrm i}  } }
} {,}
dan
\maabbeledisp {\psi} {  U \left(  0,1  \right)  } { {\mathbb H}
} {w} { { \frac{   (w+1)  { \mathrm i}  }{ 1-w  } }
} {.}

}
{} {}

\inputaufgabegibtloesung
{}
{

Uraikan
\definitionsverweis {pemetaan}{}{}
\maabbeledisp {\psi} { \Complex \setminus \{1\} } { \Complex
} {w} { { \frac{   (w+1)  { \mathrm i}  }{ 1-w  } }
} {}
dalam koordinat real.

}
{} {}

\inputaufgabegibtloesung
{}
{

Tentukan
\definitionsverweis {turunan parsial}{}{}
dari
\definitionsverweis {pemetaan}{}{}
\maabbeledisp {\psi} { \R^2 \setminus \{  (1,0 )\} } { \R^2
} {  \begin{pmatrix} a \\b   \end{pmatrix} } {  { \frac{  1 }{ -2a+1+a^2+b^2 } }  \begin{pmatrix} -2b \\1-a^2-b^2   \end{pmatrix}
} {.}

}
{} {}

\inputaufgabe
{}
{

Tentukan titik-titik potong garis yang melalui
\mathl{\begin{pmatrix}  0  \\ 0\\  -1   \end{pmatrix}}{} dan suatu titik
\mathl{\begin{pmatrix}  a  \\ b \\  0   \end{pmatrix}}{} dengan

\mavergleichskette
{\vergleichskette
{ a^2+b^2
}
{ < }{ 1
}
{  }{
}
{    }{
}
{   }{
}
}
{}{}{}
dengan hiperboloid dua lembar; lihat
Contoh 18.11.

}
{} {}

Misalkan
\mathkor {} {L} {dan} {M} {}
adalah dua
$C^k$-\definitionsverweis {manifold}{}{.}
Suatu
\definitionsverweis {pemetaan kontinu}{}{}
\maabbdisp {\varphi} { L } { M
} {}
disebut suatu
\definitionswortpraemath {C^k}{ difeomorfisme lokal }{,}
jika untuk setiap titik

\mavergleichskette
{\vergleichskette
{ P
}
{ \in }{ L
}
{  }{
}
{    }{
}
{   }{
}
}
{}{}{}
terdapat suatu
\definitionsverweis {lingkungan terbuka}{}{}

\mavergleichskette
{\vergleichskette
{ P
}
{ \in }{ U
}
{ \subseteq }{ L
}
{    }{
}
{   }{
}
}
{}{}{}
sedemikian sehingga $\varphi  \vert_U$ menginduksi suatu
\definitionsverweis {difeomorfisme}{}{}
antara $U$ dan $\varphi(U)$.

  \zwischenueberschrift{Soal untuk dikumpulkan}

\inputaufgabe
{}
{

Kita meninjau permukaan elipsoid

\mavergleichskettedisp
{\vergleichskette
{ E
}
{ =} { \left\{ (x,y,z)  \mid   2x^2+3y^2+5z^2 = 1        \right\}
}
{ \subset} { \R^3
}
{ } {
}
{ } {
}
}
{}{}{}
dengan
\definitionsverweis {struktur Riemann}{}{}
yang diinduksi oleh ruang sekitarnya. Tentukan
\definitionsverweis {isometri linear}{}{}
pada $\R^3$ yang menginduksi suatu
\definitionsverweis {isometri}{}{}
pada $E$.

}
{} {}

\inputaufgabe
{4}
{

Misalkan $L,M$ adalah
\definitionsverweis {manifold terdiferensialkan}{}{,}
misalkan
\maabbdisp {\varphi} { L } { M
} {}
adalah suatu
\definitionsverweis {difeomorfisme lokal}{}{,}
dan misalkan $M$ adalah suatu
\definitionsverweis {manifold Riemann}{}{.}
Tunjukkan bahwa terdapat tepat satu struktur Riemann pada $L$ yang membuat $\varphi$ menjadi suatu
\definitionsverweis {isometri lokal}{}{.}

}
{} {}

\inputaufgabe
{3}
{

Berikan suatu contoh
\definitionsverweis {manifold Riemann}{}{}
$M$ dan suatu
\definitionsverweis {difeomorfisme}{}{}
\maabbdisp {\varphi} { M } { M
} {,}
yang bukan suatu
\definitionsverweis {isometri}{}{.}

}
{} {}

\inputaufgabe
{4}
{

Hitung panjang busur pada setengah lingkaran atas berjari-jari $\sqrt{2}$ dari
\mathl{\begin{pmatrix}  1  \\ 1   \end{pmatrix}}{} ke
\mathl{\begin{pmatrix}  -1 \\ 1   \end{pmatrix}}{} dalam
\definitionsverweis {setengah bidang Poincaré}{}{.}

}
{} {}

\inputaufgabe
{4}
{

Hitung luas cakram berpusat di
\mathl{\begin{pmatrix}  2  \\ 2   \end{pmatrix}}{} dan berjari-jari $1$ dalam
\definitionsverweis {setengah bidang Poincaré}{}{.}

}
{} {}

\inputaufgabe
{4}
{

Uraikan suatu
\definitionsverweis {isometri}{}{}
eksplisit antara
\definitionsverweis {setengah bidang Poincaré}{}{}
dan lembar atas hiperboloid dua lembar.

}
{} {}

\section*{Solusi yang disediakan oleh sumber}
\addcontentsline{toc}{section}{Solusi yang disediakan oleh sumber}
\subsection*{Solusi Soal 18.8}
Pada

\mavergleichskettedisp
{\vergleichskette
{ \R^3
}
{ =} { \R^2 \times \R
}
{ } {
}
{ } {
}
{ } {
}
}
{}{}{}
pemetaan tersebut merupakan hasil kali suatu
\definitionsverweis {rotasi bidang}{}{}
dengan identitas. Karena itu, pemetaan tersebut adalah suatu
\definitionsverweis {isometri}{}{}
Euklides pada $\R^3$. Misalkan

\mavergleichskette
{\vergleichskette
{ (x,y,z)
}
{ \in }{  T
}
{  }{
}
{    }{
}
{   }{
}
}
{}{}{}
dan

\mavergleichskettedisp
{\vergleichskette
{ \Psi_\alpha(x,y,z)
}
{ =} { ( u,v,z)
}
{ } {
}
{ } {
}
{ } {
}
}
{}{}{.}
Maka

\mavergleichskettealigndrucklinks
{\vergleichskettealigndrucklinks
{ \left(  \sqrt{u^2+v^2} -R  \right)^2 +z^2
}
{ =} { \left(  \sqrt{ \left(  x \cos \alpha - y \sin  \alpha   \right)^2+ \left(  x \sin \alpha + y \cos  \alpha  \right)^2} -R  \right)^2 +z^2
}
{ =} { \left(  \sqrt{  x^2 \cos^{ 2  } \alpha + y^2 \sin^{ 2  }  \alpha  -2xy \cos \alpha \sin  \alpha    + x^2 \sin^{ 2  } \alpha + y^2 \cos^{ 2  }  \alpha  +2xy \cos \alpha \sin  \alpha     } -R  \right)^2 +z^2
}
{ =} { \left(  \sqrt{  x^2   + y^2 } -R  \right)^2 +z^2
}
{ =} { r^2
}
}
{}{}{,}
sehingga

\mavergleichskette
{\vergleichskette
{ \Psi_\alpha(x,y,z)
}
{ \in }{  T
}
{  }{
}
{    }{
}
{   }{
}
}
{}{}{.}
Jadi,
\maabbdisp {\Psi_\alpha} { T} {T
} {}
adalah suatu difeomorfisme sekaligus isometri karena $T$ menyandang struktur Riemann terinduksi.
\subsection*{Solusi Soal 18.11}
$\varphi$ terdefinisi karena

\mavergleichskette
{\vergleichskette
{ - { \mathrm i}
}
{ \notin }{ {\mathbb H}
}
{  }{
}
{    }{
}
{   }{
}
}
{}{}{.}
Dengan demikian, pemetaan tersebut terdefinisi pada ${\mathbb H}$. Kita akan menunjukkan bahwa

\mavergleichskettedisp
{\vergleichskette
{ \betrag {  { \frac{  z-  { \mathrm i}  }{ z + { \mathrm i}  } }  }
}
{ <} { 1
}
{ } {
}
{ } {
}
{ } {
}
}
{}{}{,}
yang ekuivalen dengan

\mavergleichskettedisp
{\vergleichskette
{ \betrag {  z- { \mathrm i}  }
}
{ <} { \betrag { z + { \mathrm i}  }
}
{ } {
}
{ } {
}
{ } {
}
}
{}{}{.}
Dengan

\mavergleichskette
{\vergleichskette
{z
}
{ = }{ a+b { \mathrm i}
}
{  }{
}
{    }{
}
{   }{
}
}
{}{}{}
ketaksamaan ini ekuivalen dengan

\mavergleichskettedisp
{\vergleichskette
{ a^2 +(b+1)^2
}
{ >} { a^2 +(b-1)^2
}
{ } {
}
{ } {
}
{ } {
}
}
{}{}{,}
yang selanjutnya ekuivalen dengan

\mavergleichskettedisp
{\vergleichskette
{ b^2+2b+1
}
{ >} { b^2 -2b +1
}
{ } {
}
{ } {
}
{ } {
}
}
{}{}{.}
Pernyataan terakhir benar karena

\mavergleichskette
{\vergleichskette
{ b
}
{ > }{ 0
}
{  }{
}
{    }{
}
{   }{
}
}
{}{}{.}

Untuk

\mavergleichskette
{\vergleichskette
{ \betrag {  w  }
}
{ < }{ 1
}
{  }{
}
{    }{
}
{   }{
}
}
{}{}{}
langsung berlaku

\mavergleichskette
{\vergleichskette
{ 1-w
}
{ \neq }{ 0
}
{  }{
}
{    }{
}
{   }{
}
}
{}{}{,}
sehingga $\psi$ terdefinisi. Misalkan

\mavergleichskette
{\vergleichskette
{ w
}
{ = }{ c+d  { \mathrm i}
}
{  }{
}
{    }{
}
{   }{
}
}
{}{}{}
dengan

\mavergleichskette
{\vergleichskette
{ c^2+d^2
}
{ < }{ 1
}
{  }{
}
{    }{
}
{   }{
}
}
{}{}{.}
Maka

\mavergleichskettealignhandlinks
{\vergleichskettealignhandlinks
{ { \frac{   (w+1) { \mathrm i}  }{  1-w  } }
}
{ =} { { \frac{   (c+1+d { \mathrm i} ) { \mathrm i}  }{  1-c -d { \mathrm i}  } }
}
{ =} { { \frac{   -d +(c+1) { \mathrm i}  }{  1-c -d { \mathrm i}  } }
}
{ =} { { \frac{   -d +(c+1) { \mathrm i}  }{  1-c -d { \mathrm i}  } } \cdot { \frac{   1-c+d { \mathrm i}  }{  1-c +d { \mathrm i}  } }
}
{ =} { { \frac{   -d (1-c) -d(c+1) + ( (1+c)(1-c)-d^2 ) { \mathrm i}  }{  (1-c)^2 +d^2  } }
}
}
{}
{}{.}
Karena penyebutnya merupakan bilangan real positif, bagian imajiner bilangan ini positif berdasarkan

\mavergleichskettedisp
{\vergleichskette
{ 1-c^2-d^2
}
{ >0} {
}
{ } {
}
{ } {
}
{ } {
}
}
{}{}{.}
Jadi, citra $\psi$ terletak di ${\mathbb H}$.

Sekarang kita menghitung kedua komposisi. Berlaku

\mavergleichskettealign
{\vergleichskettealign
{ \varphi( \psi(w))
}
{ =} { \varphi \left( \frac{   (w+1) { \mathrm i}  }{  1-w  } \right)
}
{ =} { { \frac{   { \frac{  (w+1) { \mathrm i}  }{  1-w  } } - { \mathrm i}  }{  { \frac{   (w+1) { \mathrm i}  }{  1-w  } } + { \mathrm i}  } }
}
{ =} { { \frac{   (w+1) { \mathrm i} - (1-w) { \mathrm i}  }{  (w+1) { \mathrm i} + (1-w) { \mathrm i}  } }
}
{ =} { { \frac{   2 w { \mathrm i}  }{  2 { \mathrm i}  } }
}
}
{
\vergleichskettefortsetzungalign
{ =} { w
}
{ } {}
{ } {}
{ } {}
}
{}{}
dan

\mavergleichskettealign
{\vergleichskettealign
{ \psi( \varphi(z))
}
{ =} { \psi \left(  { \frac{   z- { \mathrm i}  }{  z + { \mathrm i}  } }  \right)
}
{ =} { { \frac{   \left(  { \frac{   z- { \mathrm i}  }{ z + { \mathrm i}  } } +1  \right)  { \mathrm i}  }{  1- { \frac{   z- { \mathrm i}  }{ z + { \mathrm i}  } }  } }
}
{ =} { { \frac{   ( z- { \mathrm i} + z+ { \mathrm i} ) { \mathrm i}  }{  z+ { \mathrm i} - (z- { \mathrm i} )  } }
}
{ =} { { \frac{   2 z { \mathrm i}  }{  2 { \mathrm i}  } }
}
}
{
\vergleichskettefortsetzungalign
{ =} { z
}
{ } {}
{ } {}
{ } {}
}
{}{.}
\subsection*{Solusi Soal 18.13}
Kita tulis

\mavergleichskette
{\vergleichskette
{ w
}
{ = }{ a+b { \mathrm i}
}
{  }{
}
{    }{
}
{   }{
}
}
{}{}{.}
Maka

\mavergleichskettealign
{\vergleichskettealign
{  \psi ( a+b { \mathrm i}   )
}
{ =} { { \frac{   \left(  a+b { \mathrm i} +1  \right) { \mathrm i}      }{  1-  a-b { \mathrm i}      } }
}
{ =} { { \frac{     a { \mathrm i} +{ \mathrm i}  -b  }{  1-  a-b { \mathrm i}      } }
}
{ =} { { \frac{     \left(  a { \mathrm i} +{ \mathrm i}  -b  \right) \left(  1-  a+b { \mathrm i}  \right)  }{  \left(  1-  a-b { \mathrm i}  \right) \left(  1-  a+b { \mathrm i}  \right)      } }
}
{ =} { { \frac{     \left(  a+1 \right) \left(  1-a  \right) { \mathrm i}  -b^2 { \mathrm i} -b \left(  1-  a \right) -b\left(  a+1  \right)  }{  \left(  1-  a \right)^2 +b^2  } }
}
}
{
\vergleichskettefortsetzungalign
{ =} { { \frac{     \left(  1-a ^2-b^2  \right) { \mathrm i}   -2 b  }{  \left(  1-  a \right)^2 +b^2  } }
}
{ } {
}
{ } {}
{ } {}
}
{}{.}
Jadi, dalam koordinat real,

\mavergleichskettedisp
{\vergleichskette
{ \psi \begin{pmatrix}  a  \\b   \end{pmatrix}
}
{ =} { { \frac{   1  }{  -2a+1+a^2+b^2 } } \begin{pmatrix}  -2b \\ 1-a^2-b^2   \end{pmatrix}
}
{ } {
}
{ } {
}
{ } {
}
}
{}{}{.}
\subsection*{Solusi Soal 18.14}
Turunan parsial dari

\mavergleichskettedisp
{\vergleichskette
{ \psi \begin{pmatrix}  a  \\b   \end{pmatrix}
}
{ =} { { \frac{   1 }{  -2a+1+a^2+b^2 } } \begin{pmatrix}  -2b \\ 1-a^2-b^2   \end{pmatrix}
}
{ } {
}
{ } {
}
{ } {
}
}
{}{}{}
adalah

\mavergleichskettealign
{\vergleichskettealign
{ { \frac{   \partial   \psi  }{  \partial  a  } }
}
{ =} { { \frac{   1 }{  \left(  -2a+1+a^2+b^2  \right)^2  } } \begin{pmatrix}  2b \left(  2a-2  \right)  \\ -2a \left(  -2a+1+a^2+b^2  \right) - \left(  1-a^2-b^2  \right) \left(  -2+2a  \right)    \end{pmatrix}
}
{ =} { { \frac{   1 }{  \left(  -2a+1+a^2+b^2  \right)^2  } } \begin{pmatrix}  2b \left(  2a-2  \right)  \\ 4a^2 -4a+2 \left(  1-a^2-b^2  \right)    \end{pmatrix}
}
{ =} { { \frac{   1 }{  \left(  -2a+1+a^2+b^2  \right)^2  } } \begin{pmatrix}  4ab -4b  \\ 2a^2 -4a+2 -2 b^2     \end{pmatrix}
}
{ } {
}
}
{}
{}{}
dan

\mavergleichskettealign
{\vergleichskettealign
{ { \frac{   \partial   \psi  }{  \partial  b  } }
}
{ =} { { \frac{   1 }{  \left(  -2a+1+a^2+b^2  \right)^2  } } \begin{pmatrix}  -2 \left(  -2a+1+a^2+b^2  \right) +4b^2   \\ -2b \left(  -2a+1+a^2+b^2  \right) - 2b \left(  1-a^2-b^2  \right)     \end{pmatrix}
}
{ =} { { \frac{   1 }{  \left(  -2a+1+a^2+b^2  \right)^2  } } \begin{pmatrix}  4a  -2-2a^2+2b^2   \\ 4ab -4b    \end{pmatrix}
}
{ } {
}
{ } {
}
}
{}
{}{.}

\part{Unit 19}
\chapter[Kuliah 19]{Kuliah 19: Pemetaan Weingarten dan Kelengkungan Gauss}
\setcounter{section}{19}

  \zwischenueberschrift{Pemetaan Weingarten untuk manifold terparametrisasi}

Misalkan

\mavergleichskette
{\vergleichskette
{ W
}
{ \subseteq }{ \R^n
}
{  }{
}
{    }{
}
{   }{
}
}
{}{}{}
\definitionsverweis {terbuka}{}{,}
\maabb {f} { W } { \R
} {}
suatu
\definitionsverweis {fungsi terdiferensialkan kontinu dua kali}{}{}
dan

\mavergleichskette
{\vergleichskette
{ Y
}
{ = }{ f^{-1}(c)
}
{  }{
}
{    }{
}
{   }{
}
}
{}{}{}
\definitionsverweis {serat}{}{}
di atas

\mavergleichskette
{\vergleichskette
{ c
}
{ \in }{ \R
}
{  }{
}
{    }{
}
{   }{
}
}
{}{}{,}
dengan $f$
\definitionsverweis {reguler}{}{}
pada setiap titik di $Y$. Misalkan
\maabbdisp {\varphi} { V } { Y
} {}
dengan

\mavergleichskette
{\vergleichskette
{ V
}
{ \subseteq }{ \R^{n-1}
}
{  }{
}
{    }{
}
{   }{
}
}
{}{}{}
terbuka. Misalkan pemetaan tersebut merupakan suatu parametrisasi difeomorfik pada himpunan terbuka

\mavergleichskette
{\vergleichskette
{ U
}
{ \subseteq }{ Y
}
{  }{
}
{    }{
}
{   }{}
}
{}{}{,}
Dengan demikian, $\varphi^{-1}$ dapat dipandang sebagai suatu
\definitionsverweis {peta}{}{}
untuk $Y$. Untuk

\mavergleichskette
{\vergleichskette
{ Q
}
{ \in }{ V
}
{  }{
}
{    }{
}
{   }{
}
}
{}{}{}
dengan

\mavergleichskette
{\vergleichskette
{ \varphi(Q)
}
{ = }{ P
}
{  }{
}
{    }{
}
{   }{
}
}
{}{}{}

\mavergleichskettedisp
{\vergleichskette
{ T_{P}Y
}
{ =} { \operatorname{im}  \left(D\varphi\right)_{Q}
}
{ } {
}
{ } {
}
{ } {
}
}
{}{}{,}
Dengan demikian, terdapat isomorfisme linear
\maabbdisp {\left(D\varphi\right)_{Q}} {\R^{n-1}} { T_{P} Y
} {}
yang mendeskripsikan ruang singgung
\mathl{T_{P} Y}{.} Vektor-vektor basis standar $e_i$ dipetakan ke

\mavergleichskette
{\vergleichskette
{ \partial_i \varphi(Q)
}
{ = }{ \begin{pmatrix}  { \frac{   \partial  \varphi_1   }{  \partial  u_i   } } ( Q)  \\\vdots\\  { \frac{   \partial  \varphi_n   }{  \partial  u_i   } } ( Q)   \end{pmatrix}
}
{  }{
}
{    }{
}
{   }{
}
}
{}{}{}
dan membentuk suatu basis ruang singgung. Jika

\mavergleichskette
{\vergleichskette
{ Y
}
{ \subseteq }{ \R^n
}
{  }{
}
{    }{
}
{   }{
}
}
{}{}{}
dilengkapi dengan struktur Euklides pada $\R^n$, struktur ini menginduksi
\definitionsverweis {struktur Riemann}{}{}
dan menghasilkan fungsi-fungsi

\mavergleichskettedisp
{\vergleichskette
{ g_{ij}(Q)
}
{ =} { \left\langle  \partial_i \varphi (Q)   ,  \partial_j \varphi (Q)   \right\rangle
}
{ } {
}
{ } {
}
{ } {
}
}
{}{}{,}
yang dihimpun menjadi matriks fundamental pertama
\zusatzklammer {matriks fundamental metrik} {} {}

\mavergleichskettedisp
{\vergleichskette
{ G
}
{ =} { \left(  g_{ij}  \right)
}
{ } {
}
{ } {
}
{ } {
}
}
{}{}{}
Matriks fundamental pertama sepenuhnya ditentukan oleh struktur Riemann pada $Y$.

\definitionsverweis {Pemetaan Weingarten}{}{}
\maabbdisp {L_{P}} { T_{P}Y } { T_{P}Y
} {}
kini dapat dideskripsikan terhadap basis
\mathl{\partial_i \varphi (Q)}{.} Cara terbaik untuk melakukannya ialah dengan memperkenalkan apa yang disebut matriks fundamental kedua. Untuk itu, orientasikan $Y$ dengan suatu
\definitionsverweis {medan normal satuan}{}{}
$N$
\definitionsverweis {yang menentukan orientasi}{}{.}
Hal ini dapat dilakukan dengan menggunakan fungsi pendeskripsi $f$. Medan normal satuan yang dibatasi pada $U$ dapat langsung dipandang sebagai medan pada $V$. Karena $f$ dan $\varphi$ diasumsikan terdiferensialkan kontinu dua kali, medan normal satuan pada $V$ terdiferensialkan. Kita mendefinisikan hal berikut.

\inputdefinition
{}
{

Misalkan

\mavergleichskette
{\vergleichskette
{ W
}
{ \subseteq }{ \R^n
}
{  }{
}
{    }{
}
{   }{
}
}
{}{}{}
\definitionsverweis {terbuka}{}{,}
\maabb {f} { W } { \R
} {}
suatu
\definitionsverweis {fungsi terdiferensialkan kontinu dua kali}{}{}
dan

\mavergleichskette
{\vergleichskette
{ Y
}
{ = }{ f^{-1}(c)
}
{  }{
}
{    }{
}
{   }{
}
}
{}{}{}
\definitionsverweis {serat}{}{}
di atas

\mavergleichskette
{\vergleichskette
{ c
}
{ \in }{ \R
}
{  }{
}
{    }{
}
{   }{
}
}
{}{}{,}
dengan $f$
\definitionsverweis {reguler}{}{}
pada setiap titik di $Y$. Misalkan
\maabbdisp {\varphi} { V } { Y
} {,}

\mavergleichskette
{\vergleichskette
{ V
}
{ \subseteq }{ \R^{n-1}
}
{  }{
}
{    }{
}
{   }{
}
}
{}{}{,}
suatu parametrisasi terdiferensialkan dua kali pada suatu himpunan terbuka

\mavergleichskette
{\vergleichskette
{ U
}
{ \subseteq }{ Y
}
{  }{
}
{    }{
}
{   }{
}
}
{}{}{}
dengan parameter
\mathl{u_1 , \ldots , u_{n-1}}{.} Orientasikan $Y$ dengan suatu
\definitionsverweis {medan normal satuan}{}{}
$N$
\definitionsverweis {yang menentukan orientasi}{}{,}
dan pandang $N$ sebagai medan pada $V$. Tetapkan

\mavergleichskettedisp
{\vergleichskette
{ h_{i j}
}
{ =} { \left\langle  \partial_i \partial_j \varphi , N  \right\rangle
}
{ } {
}
{ } {
}
{ } {
}
}
{}{}{.}
\definitionswort {Matriks fundamental kedua}{}
untuk $\varphi$ adalah matriks
\zusatzklammer {yang bergantung pada

\mavergleichskettek
{\vergleichskettek
{ Q
}
{ \in }{ V
}
{  }{
}
{    }{
}
{   }{
}
}
{}{}{}} {} {}

\mavergleichskettedisp
{\vergleichskette
{ H
}
{ =} { \left(  h_{ij}   \right)_{1 \leq i,j \leq n-1}
}
{ } {
}
{ } {
}
{ } {
}
}
{}{}{.}

}

Definisi ini menggunakan medan normal satuan dan, melalui hasil kali skalar, struktur Riemann pada $Y$. Menurut
teorema Schwarz,
matriks fundamental kedua bersifat simetris.



\inputfaktbeweis
{Hyperfläche/R^n/Parametrisierung/Zweite Fundamentalform/Variante/Fakt}
{Lema}
{}
{

\faktsituation {Misalkan

\mavergleichskette
{\vergleichskette
{ W
}
{ \subseteq }{ \R^n
}
{  }{
}
{    }{
}
{   }{
}
}
{}{}{}
\definitionsverweis {terbuka}{}{,}
\maabb {f} { W } { \R
} {}
suatu
\definitionsverweis {fungsi terdiferensialkan kontinu dua kali}{}{}
dan

\mavergleichskette
{\vergleichskette
{ Y
}
{ = }{ f^{-1}(c)
}
{  }{
}
{    }{
}
{   }{
}
}
{}{}{}
\definitionsverweis {serat}{}{}
di atas

\mavergleichskette
{\vergleichskette
{ c
}
{ \in }{ \R
}
{  }{
}
{    }{
}
{   }{
}
}
{}{}{,}
dengan $f$
\definitionsverweis {reguler}{}{}
pada setiap titik di $Y$. Misalkan
\maabbdisp {\varphi} { V } { Y
} {,}

\mavergleichskette
{\vergleichskette
{ V
}
{ \subseteq }{ \R^{n-1}
}
{  }{
}
{    }{
}
{   }{
}
}
{}{}{,}
suatu parametrisasi terdiferensialkan dua kali pada suatu himpunan terbuka

\mavergleichskette
{\vergleichskette
{ U
}
{ \subseteq }{ Y
}
{  }{
}
{    }{
}
{   }{
}
}
{}{}{}
dengan parameter
\mathl{u_1 , \ldots , u_{n-1}}{.} Orientasikan $Y$ dengan suatu
\definitionsverweis {medan normal satuan}{}{}
$N$
\definitionsverweis {yang menentukan orientasi}{}{,}
dan pandang $N$ sebagai medan pada $V$.}
\faktfolgerung {Maka berlaku

\mavergleichskettedisp
{\vergleichskette
{ h_{i j}
}
{ =} { - \left\langle  \partial_i  \varphi ,  \partial_j N  \right\rangle
}
{ =} { \left\langle  \partial_i  \varphi ,  L_{} ( \partial_j \varphi)  \right\rangle
}
{ } {
}
{ } {
}
}
{}{}{.}}
\faktzusatz {}
\faktzusatz {}

}
{

Karena
\definitionsverweis {medan normal satuan}{}{}
tegak lurus terhadap
\definitionsverweis {ruang-ruang singgung}{}{}
pada $Y$, berlaku

\mavergleichskettedisp
{\vergleichskette
{ \left\langle  \partial_i \varphi , N  \right\rangle
}
{ =} { 0
}
{ } {
}
{ } {
}
{ } {
}
}
{}{}{.}
Dari sini diperoleh

\mavergleichskettedisp
{\vergleichskette
{ 0
}
{ =} { \left\langle  \partial_j \partial_i \varphi , N  \right\rangle + \left\langle  \partial_i \varphi ,  \partial_j N  \right\rangle
}
{ } {
}
{ } {
}
{ } {
}
}
{}{}{,}
yang memberikan persamaan pertama. Persamaan kedua langsung mengikuti definisi pemetaan Weingarten.

}

Sekarang kita membatasi pembahasan pada kasus

\mavergleichskette
{\vergleichskette
{ n
}
{ = }{ 3
}
{  }{
}
{    }{
}
{   }{
}
}
{}{}{.}
Matriks fundamental pertama adalah

\mavergleichskettedisp
{\vergleichskette
{ \begin{pmatrix} g_{11}  &  g_{12}  \\ g_{21} & g_{22}  \end{pmatrix}
}
{ =} { \begin{pmatrix}  E   &   F   \\  F & G  \end{pmatrix}
}
{ } {
}
{ } {
}
{ } {
}
}
{}{}{}
dengan

\mavergleichskettedisp
{\vergleichskette
{ g_{ij}
}
{ =} { \left\langle  \partial_i \varphi ,  \partial_j \varphi  \right\rangle
}
{ } {
}
{ } {
}
{ } {
}
}
{}{}{}
Dalam konteks ini kita memilih notasi pertama. Matriks fundamental kedua adalah

\mathdisp {\begin{pmatrix} h_{11}  &  h_{12}  \\ h_{21} & h_{22}  \end{pmatrix}} {  }

Untuk menyingkat hipotesis dalam pernyataan berikut, kita juga mengatakan bahwa terdapat suatu \stichwort {permukaan berorientasi} {} atau suatu permukaan berorientasi yang terparametrisasi sepotong demi sepotong di $\R^3$.


\inputfaktbeweis
{Hyperfläche/Raum/Parametrisierung/Weingartenabbildung/Fundamentalformen/Fakt}
{Lema}
{}
{

\faktsituation {Misalkan

\mavergleichskette
{\vergleichskette
{ T
}
{ \subseteq }{ \R^3
}
{  }{
}
{    }{
}
{   }{
}
}
{}{}{}
\definitionsverweis {terbuka}{}{,}
\maabb {f} { T } { \R
} {}
suatu
\definitionsverweis {fungsi terdiferensialkan kontinu dua kali}{}{}
dan

\mavergleichskette
{\vergleichskette
{ Y
}
{ = }{ f^{-1}(c)
}
{  }{
}
{    }{
}
{   }{
}
}
{}{}{}
\definitionsverweis {serat}{}{}
di atas

\mavergleichskette
{\vergleichskette
{ c
}
{ \in }{ \R
}
{  }{
}
{    }{
}
{   }{
}
}
{}{}{,}
dengan $f$
\definitionsverweis {reguler}{}{}
di setiap titik $Y$. Misalkan
\maabbdisp {\varphi} { V } { Y
} {,}

\mavergleichskette
{\vergleichskette
{ V
}
{ \subseteq }{ \R^{2}
}
{  }{
}
{    }{
}
{   }{
}
}
{}{}{,}
suatu parametrisasi terdiferensialkan dua kali pada suatu himpunan terbuka

\mavergleichskette
{\vergleichskette
{ U
}
{ \subseteq }{ Y
}
{  }{
}
{    }{
}
{   }{
}
}
{}{}{}
dengan parameter $u_1,u_2$. Orientasikan $Y$ dengan suatu
\definitionsverweis {medan normal satuan}{}{}
$N$
\definitionsverweis {yang menentukan orientasi}{}{,}
dan pandang $N$ sebagai medan pada $V$. Misalkan $G$ adalah
\definitionsverweis {matriks fundamental pertama}{}{}
dan $H$ adalah
\definitionsverweis {matriks fundamental kedua}{}{}
untuk $\varphi$. Untuk

\mavergleichskette
{\vergleichskette
{ P
}
{ = }{ \varphi(Q)
}
{  }{
}
{    }{
}
{   }{
}
}
{}{}{}
misalkan $W$ adalah matriks yang merepresentasikan
\definitionsverweis {pemetaan Weingarten}{}{}
\maabbdisp {L_P} {T_PY} {T_PY
} {}
terhadap basis
\mathkor {} {\partial_1 \varphi (Q)} {dan} {\partial_2 \varphi (Q)} {}
dari $T_PY$.}
\faktuebergang {Maka berlaku pernyataan-pernyataan berikut.}
\faktfolgerung {\aufzaehlungvier{Berlaku

\mavergleichskettedisp
{\vergleichskette
{ H
}
{ =} { G \cdot W
}
{ } {
}
{ } {
}
{ } {
}
}
{}{}{.}
}{Berlaku

\mavergleichskettedisp
{\vergleichskette
{ W
}
{ =} { G^{-1} H
}
{ } {
}
{ } {
}
{ } {
}
}
{}{}{.}
}{Berlaku

\mavergleichskettedisphandlinks
{\vergleichskettedisphandlinks
{ \partial_1 N
}
{ =} { -  { \frac{   1  }{ g_{11} g_{22} - g_{12}^2  } } \left(  \left(  g_{22} h_{11}  -g_{12}h_{12}  \right) \partial_1 \varphi + \left(  - g_{12} h_{11} + g_{11}h_{12}  \right) \partial_2 \varphi \right)
}
{ } {
}
{ } {
}
{ } {
}
}
{}{}{}
dan

\mavergleichskettedisphandlinks
{\vergleichskettedisphandlinks
{ \partial_2 N
}
{ =} { - { \frac{   1  }{  g_{11} g_{22} - g_{12}^2  } }  \left(  \left(  g_{22} h_{12} -g_{12} h_{22}  \right)  \partial_1 \varphi + \left(  -g_{12} h_{12} + g_{11} h_{22}   \right) \partial_2 \varphi  \right)
}
{ } {
}
{ } {
}
{ } {
}
}
{}{}{}
}{Untuk
\definitionsverweis {kelengkungan Gauss}{}{}
pada $Y$, berlaku

\mavergleichskettedisp
{\vergleichskette
{ K
}
{ =} { { \frac{   \det  H  }{  \det  G  } }
}
{ } {
}
{ } {
}
{ } {
}
}
{}{}{.}
}}
\faktzusatz {}
\faktzusatz {}

}
{

\aufzaehlungvier{Menurut definisi $W$, berlaku

\mavergleichskettedisp
{\vergleichskette
{ L_P ( \partial_1 \varphi)
}
{ =} { w_{11}  \partial_1 \varphi + w_{21}\partial_2 \varphi
}
{ } {
}
{ } {
}
{ } {
}
}
{}{}{}
dan

\mavergleichskettedisp
{\vergleichskette
{ L_P ( \partial_2 \varphi)
}
{ =} { w_{12}  \partial_1 \varphi + w_{22}\partial_2 \varphi
}
{ } {
}
{ } {
}
{ } {
}
}
{}{}{.}
Karena itu, menurut
Lema 19.2,

\mavergleichskettedisp
{\vergleichskette
{ h_{ij}
}
{ =} { \left\langle  \partial_i \varphi  ,  L_P ( \partial_j \varphi)  \right\rangle
}
{ =} { \left\langle  \partial_i \varphi  ,  w_{1j}  \partial_1 \varphi + w_{2j} \partial_2 \varphi    \right\rangle
}
{ =} { w_{1j} g_{i1 } + w_{2j} g_{i2}
}
{ } {
}
}
{}{}{.}
}{Hal ini langsung mengikuti (1) dengan mengalikan dari kiri oleh $G^{-1}$.
}{Berlaku

\mavergleichskettedisp
{\vergleichskette
{ \partial_i N
}
{ =} { - L( \partial_i \varphi)
}
{ =} { -w_{1i} \partial_1 \varphi - w_{2i} \partial_2 \varphi
}
{ } {
}
{ } {
}
}
{}{}{}
dan dengan demikian pernyataan tersebut mengikuti (2), dengan memperhatikan bahwa

\mavergleichskettedisp
{\vergleichskette
{ G^{-1}
}
{ =} { { \frac{   1  }{  g_{11} g_{22} - g_{12}^2  } } \begin{pmatrix} g_{22}  &   - g_{12}   \\  -g_{12}  & g_{11}  \end{pmatrix}
}
{ } {
}
{ } {
}
{ } {
}
}
{}{}{.}
Oleh karena itu,

\mavergleichskettedisp
{\vergleichskette
{ W
}
{ =} { { \frac{   1  }{  g_{11} g_{22} - g_{12}^2  } } \begin{pmatrix} g_{22}  &   - g_{12}   \\  -g_{12}  & g_{11}  \end{pmatrix}  \begin{pmatrix} h_{11}  &   h_{12}   \\  h_{12}  & h_{22}  \end{pmatrix}
}
{ } {
}
{ } {
}
{ } {
}
}
{}{}{}
dan dengan demikian

\mavergleichskettealignhandlinks
{\vergleichskettealignhandlinks
{ \partial_1 N
}
{ =} { -w_{11} \partial_1 \varphi - w_{21} \partial_2 \varphi
}
{ =} { -{ \frac{   1  }{  g_{11} g_{22} - g_{12}^2  } }  \left(  \left(  g_{22} h_{11} -g_{12} h_{12}   \right)  \partial_1 \varphi + \left(  -g_{12} h_{11} + g_{11} h_{12}   \right)  \partial_2 \varphi  \right)
}
{ } {
}
{ } {}
}
{}
{}{}
dan

\mavergleichskettealignhandlinks
{\vergleichskettealignhandlinks
{ \partial_2 N
}
{ =} { -w_{12} \partial_1 \varphi - w_{22} \partial_2 \varphi
}
{ =} { - { \frac{   1  }{  g_{11} g_{22} - g_{12}^2  } }  \left(  \left(  g_{22} h_{12} -g_{12} h_{22}  \right)  \partial_1 \varphi + \left(  -g_{12} h_{12} + g_{11} h_{22}   \right) \partial_2 \varphi  \right)
}
{ } {
}
{ } {}
}
{}
{}{}
}{Kelengkungan Gauss adalah determinan pemetaan Weingarten. Karena itu, pernyataan ini mengikuti (2) bersama
teorema perkalian determinan.
}

}

  \zwischenueberschrift{Kelengkungan Gauss sebagai konsep intrinsik}

Contoh 18.5
menunjukkan bahwa terdapat isometri di antara permukaan-permukaan Riemann
\zusatzklammer {misalnya antara suatu bagian bidang dan selimut silinder} {} {}
yang tidak mempertahankan pemetaan Weingarten maupun kelengkungan utama. Pada bidang, pemetaan Weingarten adalah pemetaan nol; pada silinder, nilai eigennya ialah $0$ dan ${ \pm \frac{  1 }{ r } }$, dengan $r$ menyatakan jari-jari dan tandanya ditentukan oleh pilihan medan normal satuan. Namun, hasil kali nilai-nilai eigen tersebut, yakni determinan pemetaan Weingarten dan sekaligus kelengkungan Gauss, sama dengan $0$ untuk kedua permukaan. Kita akan membuktikan bahwa kelengkungan Gauss suatu permukaan di $\R^3$ memang sepenuhnya ditentukan oleh struktur Riemannnya.

\inputdefinition
{}
{

Misalkan

\mavergleichskette
{\vergleichskette
{ Y
}
{ \subseteq }{ \R^3
}
{  }{
}
{    }{
}
{   }{
}
}
{}{}{}
suatu
\definitionsverweis {permukaan berorientasi}{}{}
yang terdiferensialkan kontinu dua kali, dan misalkan
\maabbdisp {\varphi} { V } { Y
} {,}

\mavergleichskette
{\vergleichskette
{ V
}
{ \subseteq }{ \R^2
}
{  }{
}
{    }{
}
{   }{
}
}
{}{}{}
\definitionsverweis {terbuka}{}{,}
suatu parametrisasi lokal dari $Y$ yang terdiferensialkan kontinu dua kali, dengan parameter $u,v$. Misalkan $N$ adalah
\definitionsverweis {medan normal satuan}{}{}
yang dipandang sebagai medan pada $V$. Misalkan

\mavergleichskette
{\vergleichskette
{ H
}
{ = }{ \left(  h_{ij}  \right)
}
{  }{
}
{    }{
}
{   }{
}
}
{}{}{}
adalah
\definitionsverweis {matriks fundamental kedua}{}{}
untuk $\varphi$. Fungsi-fungsi yang pada setiap titik didefinisikan melalui sistem persamaan linear

\mavergleichskettedisp
{\vergleichskette
{ { \frac{   \partial^2  \varphi  }{  \partial u^2  } }
}
{ =} { \Gamma^1_{11}  { \frac{   \partial  \varphi  }{  \partial  u  } } + \Gamma^2_{11} { \frac{   \partial  \varphi  }{  \partial  v  } } +h_{11} N
}
{ } {
}
{ } {
}
{ } {
}
}
{}{}{,}

\mavergleichskettedisp
{\vergleichskette
{ { \frac{   \partial    }{  \partial   u  } } { \frac{   \partial  \varphi  }{  \partial  v  } }
}
{ =} { \Gamma^1_{12} { \frac{   \partial  \varphi  }{  \partial  u  } } + \Gamma^2_{12}  { \frac{   \partial  \varphi  }{  \partial  v  } } +h_{12} N
}
{ } {
}
{ } {
}
{ } {
}
}
{}{}{,}

\mavergleichskettedisp
{\vergleichskette
{ { \frac{   \partial    }{  \partial   v  } } { \frac{   \partial  \varphi  }{  \partial  u  } }
}
{ =} { \Gamma^1_{21} { \frac{   \partial  \varphi  }{  \partial  u  } } + \Gamma^2_{21}   { \frac{   \partial  \varphi  }{  \partial  v  } } +h_{21} N
}
{ } {
}
{ } {
}
{ } {
}
}
{}{}{,}

\mavergleichskettedisp
{\vergleichskette
{ { \frac{   \partial^2  \varphi  }{  \partial v^2  } }
}
{ =} { \Gamma^1_{22} { \frac{   \partial  \varphi  }{  \partial  u  } } + \Gamma^2_{22}   { \frac{   \partial  \varphi  }{  \partial  v  } } +h_{22} N
}
{ } {
}
{ } {
}
{ } {
}
}
{}{}{,}
yakni
\maabbdisp {\Gamma^{k}_{ij}} { V } {\R
} {}
disebut
\definitionswort {simbol Christoffel}{}
untuk $\varphi$.

}

Perhatikan bahwa ketiga vektor
\mathl{{ \frac{   \partial  \varphi  }{  \partial   u   } } ( Q), { \frac{   \partial  \varphi  }{  \partial   v   } } ( Q), N(Q)}{}
pada setiap titik

\mavergleichskette
{\vergleichskette
{ Q
}
{ \in }{ V
}
{  }{
}
{    }{
}
{   }{
}
}
{}{}{}
membentuk suatu basis $\R^3$;
karena itu, setiap vektor dapat ditulis secara tunggal sebagai kombinasi linear terhadap basis ini. Akan tetapi, basis tersebut berubah bersama $Q$, sehingga simbol-simbol Christoffel adalah fungsi. Karena $N$ ternormalisasi dan tegak lurus terhadap kedua vektor basis lainnya, koefisien pada $N$ diperoleh sebagai hasil kali skalar

\mavergleichskette
{\vergleichskette
{ \left\langle  \partial_i \partial_j \varphi , N  \right\rangle
}
{ = }{ h_{ij}
}
{  }{
}
{    }{
}
{   }{
}
}
{}{}{.}



\inputfaktbeweis
{Hyperfläche/Raum/Parametrisierung/Christoffelsymbole/Bezug zu riemannscher Metrik/Fakt}
{Lema}
{}
{

\faktsituation {Misalkan

\mavergleichskette
{\vergleichskette
{ Y
}
{ \subseteq }{ \R^3
}
{  }{
}
{    }{
}
{   }{
}
}
{}{}{}
suatu
\definitionsverweis {permukaan berorientasi}{}{}
yang terdiferensialkan kontinu dua kali, dan misalkan
\maabbdisp {\varphi} { V } { Y
} {,}

\mavergleichskette
{\vergleichskette
{ V
}
{ \subseteq }{ \R^2
}
{  }{
}
{    }{
}
{   }{
}
}
{}{}{,}
suatu parametrisasi lokal dari $Y$ yang terdiferensialkan dua kali, dengan parameter $u,v$. Misalkan
\mathl{\left(  g_{ij}  \right)}{} adalah
\definitionsverweis {matriks fundamental pertama}{}{}
pada $V$, dan misalkan
\mathl{\left(  g^{kl}  \right)}{} adalah
\definitionsverweis {matriks invers}{}{}
dari
\mathl{g_{ij}}{.}}
\faktfolgerung {Maka, untuk
\definitionsverweis {simbol-simbol Christoffel}{}{,}
berlaku

\mavergleichskettedisphandlinks
{\vergleichskettedisphandlinks
{ \Gamma^k_{ij}
}
{ =} { { \frac{   1  }{  2 } } g^{1k}  \left(  \partial_i g_{j1}  + \partial_j g_{1i} - \partial_1 g_{ij}  \right)  +  { \frac{   1  }{  2 } } g^{2k}  \left(  \partial_i g_{j2}  + \partial_j g_{2i} -  \partial_2 g_{ij}  \right)
}
{ } {
}
{ } {
}
{ } {
}
}
{}{}{.}}
\faktzusatz {Secara khusus, simbol-simbol Christoffel dapat dinyatakan melalui data matriks fundamental pertama.}
\faktzusatz {}

}
{

Berlaku

\mavergleichskettealignhandlinks
{\vergleichskettealignhandlinks
{ \partial_k g_{ij}
}
{ =} { \partial_k  \left\langle  \partial_i \varphi  ,  \partial_j \varphi  \right\rangle
}
{ =} { \left\langle  \partial_k \partial_i  \varphi ,  \partial_j \varphi  \right\rangle +  \left\langle  \partial_i  \varphi ,  \partial_k \partial_j \varphi  \right\rangle
}
{ =} { \left\langle  \Gamma_{ki}^1 \partial_1 \varphi + \Gamma^2_{ki} \partial_2 \varphi + h_{ki} N ,  \partial_j \varphi  \right\rangle +  \left\langle  \partial_i  \varphi ,  \Gamma_{kj}^1 \partial_1 \varphi + \Gamma^2_{kj} \partial_2 \varphi + h_{kj} N   \right\rangle
}
{ =} { \left\langle  \Gamma_{ki}^1 \partial_1 \varphi + \Gamma^2_{ki} \partial_2 \varphi ,  \partial_j \varphi  \right\rangle +  \left\langle  \partial_i  \varphi ,  \Gamma_{kj}^1 \partial_1 \varphi + \Gamma^2_{kj} \partial_2 \varphi  \right\rangle
}
}
{
\vergleichskettefortsetzungalign
{ =} { \Gamma^1_{ki} g_{1j} +  \Gamma^2_{ki} g_{2j} + \Gamma^1_{kj} g_{i1} +  \Gamma^2_{kj} g_{i2}
}
{ } {
}
{ } {}
{ } {}
}
{}{.}
Dengan demikian,

\mavergleichskettealigndrucklinks
{\vergleichskettealigndrucklinks
{ \partial_i g_{jk} +\partial_j g_{ki}   -\partial_k g_{ij}
}
{ =} { \begin{aligned}[t]
& \Gamma^1_{ij} g_{1k} +  \Gamma^2_{ij} g_{2k} + \Gamma^1_{ik} g_{j1}  +  \Gamma^2_{ik} g_{j2}
\\
& {}+ \Gamma^1_{jk} g_{1i} +  \Gamma^2_{jk} g_{2i} + \Gamma^1_{ji} g_{k1} +  \Gamma^2_{ji} g_{k2}
\\
& {}-  \left(  \Gamma^1_{ki} g_{1j} +  \Gamma^2_{ki} g_{2j} + \Gamma^1_{kj} g_{i1} +  \Gamma^2_{kj} g_{i2} \right)
\end{aligned}
}
{ =} { 2  \left(  \Gamma_{ij}^1 g_{1k}  + \Gamma_{ij}^2 g_{2k}  \right)
}
{ } {
}
{ } {
}
}
{}{}{.}
Dengan demikian, berlaku relasi matriks

\mavergleichskettedisp
{\vergleichskette
{ \begin{pmatrix}  \partial_i g_{j1} +\partial_j g_{1i}  -\partial_1 g_{ij} \\ \partial_i g_{j2} +\partial_j g_{2i}   -\partial_2 g_{ij}    \end{pmatrix}
}
{ =} { 2  \begin{pmatrix} g_{11}  &  g_{12}  \\ g_{12} & g_{22}  \end{pmatrix}  \begin{pmatrix} \Gamma_{ij}^1 \\\Gamma_{ij}^2   \end{pmatrix}
}
{ } {
}
{ } {
}
{ } {
}
}
{}{}{}
Dengan mengalikan relasi ini dengan
\mathl{{ \frac{   1  }{  2 } } G^{-1}}{,} diperoleh pernyataan yang dimaksud.

}

\inputbemerkung
{}
{

Dalam situasi
Lema 19.5,
berlaku relasi-relasi

\mavergleichskettedisp
{\vergleichskette
{ \begin{pmatrix}  \partial_1 g_{11} \\ 2 \partial_1 g_{12} - \partial_2 g_{11}    \end{pmatrix}
}
{ =} { 2  \begin{pmatrix} g_{11}  &  g_{12}  \\ g_{12} & g_{22}  \end{pmatrix}  \begin{pmatrix} \Gamma^1_{11}  \\\Gamma^2_{11}    \end{pmatrix}
}
{ } {
}
{ } {
}
{ } {
}
}
{}{}{,}

\mavergleichskettedisp
{\vergleichskette
{ \begin{pmatrix}  2 \partial_2 g_{12} - \partial_1 g_{22}  \\ \partial_2 g_{22}   \end{pmatrix}
}
{ =} { 2  \begin{pmatrix} g_{11}  &  g_{12}  \\ g_{12} & g_{22}  \end{pmatrix}  \begin{pmatrix} \Gamma^1_{22}  \\\Gamma^2_{22}    \end{pmatrix}
}
{ } {
}
{ } {
}
{ } {
}
}
{}{}{}
dan

\mavergleichskettedisp
{\vergleichskette
{ \begin{pmatrix}  \partial_2 g_{11} \\ \partial_1 g_{22}    \end{pmatrix}
}
{ =} { 2  \begin{pmatrix} g_{11}  &  g_{12}  \\ g_{12} & g_{22}  \end{pmatrix}  \begin{pmatrix} \Gamma^1_{12}  \\\Gamma^2_{12}    \end{pmatrix}
}
{ } {
}
{ } {
}
{ } {
}
}
{}{}{.}

}



\inputfaktbeweis
{Hyperfläche/Raum/Parametrisierung/Gaußkrümmung/Christoffelsymbole/Fakt}
{Teorema}
{}
{

\faktsituation {Misalkan

\mavergleichskette
{\vergleichskette
{ Y
}
{ \subseteq }{ \R^3
}
{  }{
}
{    }{
}
{   }{
}
}
{}{}{}
suatu
\definitionsverweis {permukaan berorientasi}{}{,}
dan misalkan
\maabbdisp {\varphi} { V } { Y
} {,}

\mavergleichskette
{\vergleichskette
{ V
}
{ \subseteq }{ \R^2
}
{  }{
}
{    }{
}
{   }{
}
}
{}{}{,}
suatu parametrisasi lokal dari $Y$ yang terdiferensialkan kontinu tiga kali, dengan parameter $u,v$. Misalkan
\mathl{\left(  g_{ij}  \right)}{} adalah
\definitionsverweis {matriks fundamental pertama}{}{}
pada $V$.}
\faktfolgerung {Dengan menggunakan
\definitionsverweis {simbol-simbol Christoffel}{}{,}
\definitionsverweis {kelengkungan Gauss}{}{}
memenuhi relasi

\mavergleichskettedisphandlinks
{\vergleichskettedisphandlinks
{ K
}
{ =} { { \frac{   1  }{ g_{11}  } }  \left(  \partial_2 \Gamma_{11}^2 - \partial_1 \Gamma_{12}^2 + \Gamma^1_{11} \Gamma^2_{12}  + \Gamma_{11}^2 \Gamma_{22}^2 - \Gamma_{12}^1 \Gamma_{11}^2-  \left(  \Gamma_{12}^2  \right)^2  \right)
}
{ } {
}
{ } {
}
{ } {
}
}
{}{}{.}}
\faktzusatz {}
\faktzusatz {}

}
{

Kita mendiferensialkan persamaan penentu pertama bagi simbol-simbol Christoffel, yakni

\mavergleichskettedisp
{\vergleichskette
{ \partial_1 \partial_1 \varphi
}
{ =} { \Gamma^1_{11}  \partial_1 \varphi  + \Gamma^2_{11}  \partial_2 \varphi +h_{11} N
}
{ } {
}
{ } {
}
{ } {
}
}
{}{}{,}
terhadap variabel kedua, dan persamaan penentu kedua, yakni

\mavergleichskettedisp
{\vergleichskette
{ \partial_1 \partial_2 \varphi
}
{ =} { \Gamma^1_{12}  \partial_1 \varphi + \Gamma^2_{12}  \partial_2  \varphi + h_{12} N
}
{ } {
}
{ } {
}
{ } {
}
}
{}{}{,}
terhadap variabel pertama. Menurut teorema Schwarz, kita memperoleh

\mavergleichskettealigndrucklinks
{\vergleichskettealigndrucklinks
{ \left(  \partial_2 \Gamma^1_{11}  \right) \partial_1 \varphi +  \left(  \partial_2 \Gamma^2_{11}  \right)  \partial_2 \varphi+ \left(  \partial_2 h_{11}  \right)  N + \Gamma^1_{11} \partial_2 \partial_1 \varphi + \Gamma^2_{11}  \partial_2 \partial_2 \varphi  +h_{11} \partial_2 N
}
{ =} { \left(  \partial_1 \Gamma^1_{12}  \right)  \partial_1  \varphi+ \left(  \partial_1 \Gamma^2_{12}  \right)\partial_2 \varphi + \left(  \partial_1 h_{12}  \right)  N+ \Gamma^1_{12}  \partial_1 \partial_1 \varphi + \Gamma^2_{12}  \partial_1 \partial_2 \varphi + h_{12} \partial_1 N
}
{ } {
}
{ } {
}
{ } {
}
}
{}{}{.}
Selisih kedua ruas ini adalah $0$. Kita menentukan koefisien vektor basis $\partial_2 \varphi$. Dengan menggunakan persamaan-persamaan penentu simbol Christoffel dan
Lema 19.3 (3),
kita memperoleh

\mavergleichskettedisphandlinks
{\vergleichskettedisphandlinks
{ 0
}
{ =} { \partial_2 \Gamma^2_{11} - \partial_1 \Gamma^2_{12} +\Gamma^1_{11} \Gamma^2_{12} - \Gamma^1_{12} \Gamma^2_{11} + \Gamma_{11}^2 \Gamma_{22}^2- \left(  \Gamma_{12}^2  \right)^2 + h_{11}  { \frac{   g_{12} h_{12} - g_{11} h_{22}   }{  g_{11}g_{22} -g_{12}^2 } } - h_{12}  { \frac{   g_{12} h_{11} - g_{11}h_{12}  }{  g_{11}g_{22} -g_{12}^2 } }
}
{ } {
}
{ } {
}
{ } {
}
}
{}{}{.}
Dengan identitas

\mavergleichskettedisphandlinks
{\vergleichskettedisphandlinks
{ - g_{11} \left(  h_{11}h_{22} -h_{12}^2  \right)
}
{ =} { h_{11} \left(  g_{12} h_{12}- g_{11} h_{22}  \right) - h_{12} \left(  g_{12} h_{11} - g_{11}h_{12}  \right)
}
{ } {
}
{ } {
}
{ } {
}
}
{}{}{}
kita dapat mengganti dua suku terakhir. Kemudian, dengan
Lema 19.3 (4),
diperoleh

\mavergleichskettealignhandlinks
{\vergleichskettealignhandlinks
{ g_{11} K
}
{ =} { g_{11} { \frac{   h_{11}h_{22} -h_{12}^2  }{ g_{11}g_{22} -g_{12}^2 } }
}
{ =} { \left(  \partial_2 \Gamma_{11}^2 - \partial_1 \Gamma_{12}^2 + \Gamma^1_{11} \Gamma^2_{12}  + \Gamma_{11}^2 \Gamma_{22}^2 - \Gamma_{12}^1 \Gamma_{11}^2-  \left(  \Gamma_{12}^2  \right)^2  \right)
}
{ } {
}
{ } {
}
}
{}
{}{.}

}


\inputfaktbeweis
{Hyperfläche/Raum/Parametrisierung/Gaußkrümmung/Bezug zu riemannscher Metrik/Brioschi/Fakt}
{Korollar}
{}
{

\faktsituation {Misalkan

\mavergleichskette
{\vergleichskette
{ Y
}
{ \subseteq }{ \R^3
}
{  }{
}
{    }{
}
{   }{
}
}
{}{}{}
suatu
\definitionsverweis {permukaan berorientasi}{}{,}
dan misalkan
\maabbdisp {\varphi} { V } { Y
} {,}

\mavergleichskette
{\vergleichskette
{ V
}
{ \subseteq }{ \R^2
}
{  }{
}
{    }{
}
{   }{
}
}
{}{}{,}
suatu parametrisasi lokal dari $Y$ yang terdiferensialkan kontinu tiga kali, dengan parameter $u,v$. Misalkan
\mathl{\left(  g_{ij}  \right)}{} adalah
\definitionsverweis {matriks fundamental pertama}{}{}
pada $V$.}
\faktfolgerung {Maka, untuk
\definitionsverweis {kelengkungan Gauss}{}{}
berlaku

\mavergleichskettedisp
{\vergleichskette
{ K
}
{ =} { \begin{aligned}[t]
& { \frac{   1  }{  \left(  g_{11} g_{22} -g_{12}^2  \right)^2 } }  \Biggl(
\\
& \det  \begin{pmatrix}  - { \frac{   1  }{  2 } } \partial_2^2  g_{11}  + \partial_1 \partial_2 g_{12} - { \frac{   1  }{  2 } }  \partial_1^2  g_{22}   &   { \frac{   1  }{  2 } } \partial_1 g_{11}  &  \partial_1 g_{12} - { \frac{   1  }{  2 } } \partial_2 g_{11}  \\  \partial_2 g_{12} -{ \frac{   1  }{  2 } } \partial_1 g_{22} &  g_{11}  &  g_{12}  \\ { \frac{   1  }{  2 } } \partial_2 g_{22}  &  g_{12}  &  g_{22}  \end{pmatrix}
\\
& {}- \det  \begin{pmatrix}  0   &   { \frac{   1  }{  2 } }  \partial_2 g_{11}  &  { \frac{   1  }{  2 } }  \partial_1 g_{22}   \\  { \frac{   1  }{  2 } } \partial_2 g_{11} & g_{11} & g_{12} \\ { \frac{   1  }{  2 } } \partial_1 g_{22}  &  g_{12} & g_{22} \end{pmatrix} \Biggr)
\end{aligned}
}
{ } {
}
{ } {
}
{ } {
}
}
{}{}{.}}
\faktzusatz {}
\faktzusatz {}

 }
{ Lihat Soal 19.9. }

Dua pernyataan di atas merupakan varian-varian Theorema Egregium. Terlepas dari bentuk rumusnya, isinya menyatakan bahwa kelengkungan Gauss pada suatu permukaan tertanam dapat dideskripsikan hanya melalui pengukuran pada permukaan itu sendiri, tanpa merujuk pada ruang sekelilingnya. Pengukuran pada permukaan tersebut dikodekan oleh matriks fundamental metrik. Jadi, kita perlu mengetahui hasil kali skalar pada ruang-ruang singgung permukaan, yang persis berarti memandang permukaan itu sebagai manifold Riemann. Kelengkungan-kelengkungan utama tidak dapat ditentukan secara intrinsik saja; lihat
Contoh 18.5
dan
Soal 19.12.

\chapter{Lembar Kerja 19}
\setcounter{section}{19}

  \zwischenueberschrift{Soal latihan}

\inputaufgabe
{}
{

Tunjukkan bahwa
\definitionsverweis {matriks fundamental kedua}{}{}
bersifat
\definitionsverweis {simetris}{}{.}

}
{} {}

\inputaufgabe
{}
{

Misalkan
\maabbeledisp {\psi} {\R^2} {\R
} {(x,y)} { x^2-xy-y^3
} {,}
dan misalkan $Y$ adalah
\definitionsverweis {grafik}{}{}
dari $\psi$, yang dilengkapi dengan
\definitionsverweis {medan normal satuan}{}{}
berarah ke atas, serta
\maabbdisp {\varphi  =
\operatorname{Id}  \times \psi} {\R^2} {Y
} {}
adalah parametrisasi grafik tersebut. Tentukan
\definitionsverweis {matriks fundamental kedua}{}{}
untuk $\varphi$.

}
{} {}

\inputaufgabe
{}
{

Misalkan
\maabbeledisp {\psi} {\R^2} {\R
} {(x,y)} { e^{xy} -x^2y^3  \sin  \left(  x-y^2 \right)
} {,}
dan misalkan $Y$ adalah
\definitionsverweis {grafik}{}{}
dari $\psi$, yang dilengkapi dengan
\definitionsverweis {medan normal satuan}{}{}
berarah ke atas, serta
\maabbdisp {\varphi  =
\operatorname{Id}  \times \psi} {\R^2} {Y
} {}
adalah parametrisasi grafik tersebut. Tentukan
\definitionsverweis {matriks fundamental kedua}{}{}
untuk $\varphi$.

}
{} {}

\inputaufgabe
{}
{

Kita meninjau sfer satuan dengan
\definitionsverweis {medan normal satuan}{}{}
berarah ke luar dan parametrisasi
\maabbeledisp {\varphi} {[0, 2 \pi] \times [-1,1]} {\R^3
} {(u,v)} { \left(  \sqrt{1-v^2}  \cos  u   , \,   \sqrt{1-v^2}  \sin  u  , \,   v                 \right)
} {,}
lihat
Contoh 17.6.
Tentukan
\definitionsverweis {matriks fundamental kedua}{}{}
untuk $\varphi$ pada interior domain parameter.

}
{} {}

\inputaufgabe
{}
{

Kita meninjau sfer satuan dengan
\definitionsverweis {medan normal satuan}{}{}
berarah ke luar dan parametrisasi
\maabbeledisp {\varphi} {[0, 2 \pi] \times \R} {\R^3
} {(u,v)} { { \frac{   1  }{  \sqrt{1+v^2}  } }  ( \cos  u  ,  \sin  u  , v )
} {,}
lihat
Contoh 17.7.
Tentukan
\definitionsverweis {matriks fundamental kedua}{}{}
untuk $\varphi$.

}
{} {}

\inputaufgabe
{}
{

Tentukan
\definitionsverweis {pemetaan Weingarten}{}{}
terhadap basis $\partial_1 \varphi, \partial_2 \varphi$ untuk berbagai parametrisasi $\varphi$ dari sfer satuan dalam
Contoh 17.6,
Contoh 17.7,
dan
Contoh 17.8.

}
{} {}

\inputaufgabe
{}
{

Tentukan
\definitionsverweis {kelengkungan Gauss}{}{}
sfer satuan dengan memakai berbagai parametrisasi $\varphi$ dari
Contoh 17.6,
Contoh 17.7,
dan
Contoh 17.8,
dengan menggunakan
Lema 19.3  (4).

}
{} {}

\inputaufgabe
{}
{

Misalkan $Y$ adalah suatu
\definitionsverweis {kurva}{}{}
di $\R^2$ yang dilengkapi dengan suatu medan normal satuan. Hubungkan
\definitionsverweis {kelengkungan}{}{}
kurva tersebut dengan
\definitionsverweis {matriks fundamental kedua}{}{.}

}
{} {}

\inputaufgabe
{}
{

Misalkan

\mavergleichskette
{\vergleichskette
{ Y
}
{ \subseteq }{ \R^3
}
{  }{
}
{    }{
}
{   }{
}
}
{}{}{}
adalah suatu
\definitionsverweis {permukaan berorientasi}{}{}
dan misalkan
\maabbdisp {\varphi} { V } { Y
} {,}

\mavergleichskette
{\vergleichskette
{ V
}
{ \subseteq }{ \R^2
}
{  }{
}
{    }{
}
{   }{
}
}
{}{}{,}
adalah parametrisasi lokal $Y$ yang terdiferensialkan tiga kali, dengan parameter $u,v$. Misalkan
\mathl{\left(  g_{ij}  \right)}{} adalah
\definitionsverweis {matriks fundamental pertama}{}{}
pada $V$. Tunjukkan bahwa
\definitionsverweis {kelengkungan Gauss}{}{}
memenuhi hubungan

\mavergleichskettedisp
{\vergleichskette
{ K
}
{ =} { \begin{aligned}[t]
& { \frac{   1  }{  \left(  g_{11} g_{22} -g_{12}^2  \right)^2 } }  \Biggl(
\\
& \det  \begin{pmatrix}  - { \frac{   1  }{  2 } } \partial_2^2  g_{11}  + \partial_1 \partial_2 g_{12} - { \frac{   1  }{  2 } }  \partial_1^2  g_{22}   &   { \frac{   1  }{  2 } } \partial_1 g_{11}  &  \partial_1 g_{12} - { \frac{   1  }{  2 } } \partial_2 g_{11}  \\  \partial_2 g_{12} -{ \frac{   1  }{  2 } } \partial_1 g_{22} &  g_{11}  &  g_{12}  \\ { \frac{   1  }{  2 } } \partial_2 g_{22}  &  g_{12}  &  g_{22}  \end{pmatrix}
\\
& {}- \det  \begin{pmatrix}  0   &   { \frac{   1  }{  2 } }  \partial_2 g_{11}  &  { \frac{   1  }{  2 } }  \partial_1 g_{22}   \\  { \frac{   1  }{  2 } } \partial_2 g_{11} & g_{11} & g_{12} \\ { \frac{   1  }{  2 } } \partial_1 g_{22}  &  g_{12} & g_{22} \end{pmatrix} \Biggr)
\end{aligned}
}
{ } {
}
{ } {
}
{ } {
}
}
{}{}{}
tersebut.

}
{} {}

  \zwischenueberschrift{Soal untuk dikumpulkan}

\inputaufgabe
{3}
{

Misalkan
\maabbeledisp {\psi} {\R^2} {\R
} {(x,y)} { x^3-2x^2y-7xy^3+x^2y^2+5y^4
} {,}
dan misalkan $Y$ adalah
\definitionsverweis {grafik}{}{}
dari $\psi$, yang dilengkapi dengan
\definitionsverweis {medan normal satuan}{}{}
berarah ke atas, serta
\maabbdisp {\varphi  =
\operatorname{Id}  \times \psi} {\R^2} {Y
} {}
adalah parametrisasi grafik tersebut. Tentukan
\definitionsverweis {matriks fundamental kedua}{}{}
untuk $\varphi$.

}
{} {}

\inputaufgabe
{4}
{

Kita meninjau sfer satuan dengan
\definitionsverweis {medan normal satuan}{}{}
berarah ke luar dan parametrisasi
\maabbeledisp {\varphi} {{]{- { \frac{   \pi }{  2 } }},  { \frac{   \pi }{  2 } } [} \times {]{-\pi}, \pi[} } {\R^3
} {(u,v)} {(\cos  u  \cos  v  , \cos  u  \sin  v  , \sin  u )
} {,}
lihat
Contoh 17.8.
Tentukan
\definitionsverweis {matriks fundamental kedua}{}{}
untuk $\varphi$.

}
{} {}

\inputaufgabe
{5 (2+1+1+1)}
{

Misalkan
\maabb {\gamma} {]a,b[} { \R^2
} {}
adalah suatu kurva injektif yang
\definitionsverweis {terdiferensialkan dua kali}{}{}
dan
\definitionsverweis {berparameter panjang busur}{}{,}
dengan kurva citra

\mavergleichskettedisp
{\vergleichskette
{ C
}
{ \subseteq} { V
}
{ \subseteq} { \R^2
}
{ } {
}
{ } {
}
}
{}{}{,}
yang tertutup dalam himpunan terbuka $V$. Kita meninjau silinder di atas kurva ini, yakni pemetaan
\maabbdisp {\varphi = \gamma \times
\operatorname{Id}_{ \R }} {]a,b[ \times \R } { C \times \R
} {.}
Lengkapi silinder tersebut dengan salah satu dari kedua medan normal satuannya.
\aufzaehlungvier{Hitung
\definitionsverweis {matriks fundamental pertama}{}{}
dan
\definitionsverweis {matriks fundamental kedua}{}{}
dari $\varphi$.
}{Hitung
\definitionsverweis {pemetaan Weingarten}{}{}
terhadap basis $\partial_1 \varphi,  \partial_2 \varphi$.
}{Hitung
\definitionsverweis {kelengkungan-kelengkungan utama}{}{}
dari $C \times \R$.
}{Hitung
\definitionsverweis {kelengkungan Gauss}{}{}
dari $C \times \R$.
}

}
{} {}

\part{Unit 20}
\chapter[Kuliah 20]{Kuliah 20: Turunan Eksterior dan Setengah Ruang Euklides}
\setcounter{section}{20}

  \zwischenueberschrift{Turunan eksterior}

Pada kuliah ini kita akan mengenal suatu objek matematika jenis baru, yang disebut turunan eksterior. Turunan ini merupakan konsep turunan yang mengubah bentuk diferensial berderajat $k$ menjadi bentuk diferensial berderajat
\mathl{k+1}{.} Untuk bentuk diferensial berderajat $0$, yaitu suatu fungsi $f$, turunan eksterior yang bersesuaian hanyalah bentuk diferensial berderajat $1$, yaitu $df$, yang memasangkan pada setiap titik $P$
\zusatzklammer {untuk ruang Euklides} {} {}
diferensial total
\maabbdisp {\left(Df\right)_{P}} { \R^n } { \R
} {}
atau
\zusatzklammer {untuk suatu manifold $M$} {} {}
pemetaan singgung
\maabbdisp {T_P(f)} {T_PM} {\R
} {.}

Dalam kalkulus diferensial satu dimensi, fungsi dan turunannya atau antiturunannya merupakan objek sejenis
\zusatzklammer {hal ini masih berlaku untuk kurva terdiferensialkan} {} {,}
tetapi sejak diperkenalkannya diferensial total bagi suatu fungsi beberapa variabel, turunannya sudah merupakan objek yang secara mendasar berbeda dari fungsi itu sendiri. Turunan berarah tingkat tinggi memang dapat didefinisikan sepanjang arah-arah yang telah ditentukan dan kembali berupa fungsi, tetapi masing-masing hanya menangkap satu aspek parsial dari turunan fungsi, sedangkan diferensial total memuat seluruh informasinya.

Perbedaan mendasar antara fungsi dan turunannya ini juga berkaitan dengan kenyataan bahwa, dalam dimensi lebih tinggi, kita belum membahas pertanyaan sebaliknya: turunan mana yang memiliki antiturunan. Suatu fungsi beberapa variabel tidak dapat memiliki antiturunan dalam pengertian ini; pertanyaan tersebut hanya bermakna bagi suatu bentuk diferensial-$1$
\zusatzklammer {atau medan vektor yang bersesuaian} {} {.}
Teorema Schwarz tentang dapat dipertukarkannya turunan berarah sudah memberikan suatu kriteria perlu yang penting bagi keberadaan antiturunan dari suatu bentuk diferensial-$1$.

Pertanyaan tentang antiturunan atau bentuk primitif memperoleh kerangka alaminya dalam teori turunan eksterior. Selain itu, teori ini memungkinkan teorema Stokes dirumuskan secara ringkas. Turunan eksterior juga dapat digunakan untuk mengarakterisasi sifat-sifat topologis penting suatu manifold, walaupun hal tersebut jauh melampaui kuliah ini.

\inputdefinition
{}
{

Misalkan

\mavergleichskette
{\vergleichskette
{ U
}
{ \subseteq }{ \R^n
}
{  }{
}
{    }{
}
{   }{
}
}
{}{}{}
\definitionsverweis {terbuka}{}{}
dan misalkan

\mavergleichskette
{\vergleichskette
{ \omega
}
{ \in }{
{ \mathcal E }^{ k } (  U   )
}
{  }{
}
{    }{
}
{   }{
}
}
{}{}{}
suatu
$k$-\definitionsverweis {bentuk diferensial}{}{}
yang \definitionsverweis {terdiferensialkan secara kontinu}{}{}
dengan penyajian

\mavergleichskettedisp
{\vergleichskette
{ \omega
}
{ =} { \sum_{I \subseteq \{1 , \ldots , n \},\, { \# \left(   I  \right) } = k } f_I dx_I
}
{ } {
}
{ } {
}
{ } {
}
}
{}{}{}
dengan
\definitionsverweis {fungsi-fungsi yang terdiferensialkan secara kontinu}{}{}
\maabbdisp {f_I} { U } {\R
} {.}
Maka bentuk diferensial berderajat $(k+1)$

\mavergleichskettedisp
{\vergleichskette
{d \omega
}
{ =} { \sum_{I \subseteq \{1 , \ldots , n \},\,  { \# \left(   I  \right) } = k } df_I \wedge dx_I
}
{ =} { \sum_{I \subseteq \{1 , \ldots , n \},\,  { \# \left(   I  \right) } = k } \left(  \sum_{j = 1}^n { \frac{   \partial  f_I  }{  \partial   x_j   } } dx_j  \right)  \wedge dx_I
}
{ } {
}
{ } {
}
}
{}{}{}
disebut \definitionswort {turunan eksterior}{} dari $\omega$.

}

Kadang-kadang kita menyebutnya secara lebih tepat sebagai turunan eksterior ke-$k$. Derajat keterdiferensialan bentuk diferensial tersebut berkurang sebesar $1$, sebagaimana dapat langsung dilihat dari fungsi-fungsi koefisiennya.

Turunan eksterior relevan untuk

\mavergleichskette
{\vergleichskette
{ k
}
{ = }{ 0 , \ldots , n
}
{  }{
}
{    }{
}
{   }{
}
}
{}{}{}
sedangkan untuk

\mavergleichskette
{\vergleichskette
{ k
}
{ \geq }{ n
}
{  }{
}
{    }{
}
{   }{
}
}
{}{}{}
pemetaan ini merupakan pemetaan nol. Jika kita membatasi diri pada bentuk-bentuk diferensial mulus, diperoleh suatu barisan turunan eksterior

\mathdisp {C^\infty(U,\R) =
{ \mathcal E }^{  0  }_{ \infty } (  U   ) \stackrel{d}{\longrightarrow}
{ \mathcal E }^{  1  }_{ \infty } (  U   ) \stackrel{d}{\longrightarrow}
{ \mathcal E }^{  2  }_{ \infty } (  U   ) \stackrel{d}{\longrightarrow} \ldots \stackrel{d}{\longrightarrow}
{ \mathcal E }^{ n-1 }_{ \infty } (  U   ) \stackrel{d}{\longrightarrow}
{ \mathcal E }^{  n  }_{ \infty } (  U   ) \stackrel{d}{\longrightarrow} 0} { . }

Pada posisi pertama hanya terdapat turunan suatu fungsi
\zusatzklammer {satu-satunya himpunan indeks dengan $0$ elemen adalah himpunan kosong} {} {,}
yakni pemetaan
\mathl{f \mapsto df}{.}

Sifat-sifat terpenting turunan eksterior dirangkum sebagai berikut.


\inputfaktbeweis
{Differentialform/Offene Menge/Äußere Ableitung/Eigenschaften/Fakt}
{Lemma}
{}
{

\faktsituation {Misalkan

\mavergleichskette
{\vergleichskette
{ U
}
{ \subseteq }{ \R^n
}
{  }{
}
{    }{
}
{   }{
}
}
{}{}{}
\definitionsverweis {terbuka}{}{,}

\mavergleichskette
{\vergleichskette
{ k
}
{ \in }{ \N
}
{  }{
}
{    }{
}
{   }{
}
}
{}{}{}
dan misalkan
\maabbeledisp {d} {
{ \mathcal E }^{  k  }_{ 1 } (  U   ) } {
{ \mathcal E }^{ k+1 }_{ 0 } (  U   )
} { \omega} { d \omega
} {,}
adalah
\definitionsverweis {turunan eksterior}{}{.}}
\faktuebergang {Maka berlaku sifat-sifat berikut.}
\faktfolgerung {\aufzaehlungfuenf{Turunan eksterior
\maabbdisp {d} {
{ \mathcal E }^{  0  }_{ 1 } (  U   ) } {
{ \mathcal E }^{  1  }_{ 0 } (  U   )
} {,}
adalah
\definitionsverweis {diferensial total}{}{.}
}{Turunan eksterior bersifat
$\R$-\definitionsverweis {linear}{}{.}
}{Untuk

\mavergleichskette
{\vergleichskette
{ \omega
}
{ \in }{
{ \mathcal E }^{  k  }_{ 1 } (  U   )
}
{  }{
}
{    }{
}
{   }{
}
}
{}{}{}
dan

\mavergleichskette
{\vergleichskette
{ \tau
}
{ \in }{
{ \mathcal E }^{  \ell }_{ 1 } (  U   )
}
{  }{
}
{    }{
}
{   }{
}
}
{}{}{}
berlaku aturan hasil kali

\mavergleichskettedisp
{\vergleichskette
{ d( \omega \wedge \tau)
}
{ =} {( d \omega) \wedge \tau  + (-1)^k \omega  \wedge  (d \tau)
}
{ } {
}
{ } {
}
{ } {
}
}
{}{}{.}
Untuk

\mavergleichskette
{\vergleichskette
{k
}
{ = }{ 0
}
{  }{
}
{    }{
}
{   }{
}
}
{}{}{}
hal ini dibaca sebagai

\mavergleichskettedisp
{\vergleichskette
{ d( f \tau)
}
{ =} {( d f) \wedge \tau  +  f  d \tau
}
{ } {
}
{ } {
}
{ } {
}
}
{}{}{}
}{Untuk setiap bentuk diferensial yang dua kali terdiferensialkan secara kontinu, $\omega$,

\mavergleichskette
{\vergleichskette
{ d(d \omega)
}
{ = }{ 0
}
{  }{
}
{    }{
}
{   }{
}
}
{}{}{.}
}{Untuk suatu
\definitionsverweis {pemetaan yang dua kali terdiferensialkan secara kontinu}{}{}
\zusatzklammer {dengan

\mavergleichskettek
{\vergleichskettek
{ W
}
{ \subseteq \R^m }{}
{  }{}
{    }{}
{   }{}
}
{}{}{}
terbuka} {} {}
\maabbdisp {\psi} { W } { U
} {}
dan setiap

\mavergleichskette
{\vergleichskette
{ \omega
}
{ \in }{
{ \mathcal E }^{  k  }_{ 1 } (  U   )
}
{  }{
}
{    }{
}
{   }{
}
}
{}{}{}
berlaku bagi
\definitionsverweis {bentuk diferensial tarik balik}{}{}

\mavergleichskettedisp
{\vergleichskette
{d (\psi^* \omega)
}
{ =} { \psi^* (d \omega)
}
{ } {
}
{ } {
}
{ } {
}
}
{}{}{.}
}}
\faktzusatz {}
\faktzusatz {}

}
{

\teilbeweis {}{}{}
{(1) langsung mengikuti definisi
\zusatzklammer {himpunan kosong adalah satu-satunya himpunan indeks yang relevan} {} {.}}
{}
\teilbeweis {}{}{}
{(2). Linearitas langsung mengikuti definisi,
linearitas
dari
\definitionsverweis {diferensial total}{}{}
dan
\definitionsverweis {multilinearitas}{}{}
dari
\definitionsverweis {produk eksterior}{}{.}}
{}

\teilbeweis {}{}{}
{(3). Misalkan
\mathl{x_1 , \ldots , x_n}{} koordinat-koordinat pada $\R^n$. Berkat linearitas $d$ dan
multilinearitas produk eksterior,
kedua bentuk diferensial dapat ditulis sebagai
\mathkor {} {\omega=f dx_I} {dan} {\tau=g dx_J} {}
dengan himpunan indeks
\mathkor {} {I=\{i_1 , \ldots , i_k\}} {dan} {J=\{ j_1 , \ldots , j_\ell \}} {.}
Maka

\mavergleichskettealignhandlinks
{\vergleichskettealignhandlinks
{ d( \omega \wedge \tau)
}
{ =} { d \left(  fdx_I \wedge gdx_J  \right)
}
{ =} { d \left(  (f g) dx_I \wedge dx_J  \right)
}
{ =} { \sum_{s = 1}^n { \frac{   \partial  (fg)  }{  \partial   x_s  } }  dx_s  \wedge  dx_I  \wedge dx_J
}
{ =} { \sum_{s = 1}^n \left(  g { \frac{   \partial   f   }{  \partial   x_s  } } + f { \frac{   \partial   g   }{  \partial   x_s  } }  \right)  dx_s \wedge  dx_I  \wedge dx_J
}
}
{
\vergleichskettefortsetzungalign
{ =} { \sum_{s = 1}^n  g { \frac{   \partial   f   }{  \partial   x_s  } } dx_s \wedge  dx_I \wedge dx_J + \sum_{s = 1}^n  f { \frac{   \partial   g   }{  \partial   x_s  } }  dx_s \wedge  dx_I  \wedge dx_J
}
{ =} { \sum_{s = 1}^n  { \frac{   \partial   f   }{  \partial   x_s  } } dx_s \wedge  dx_I \wedge gdx_J + \sum_{s = 1}^n  { \frac{   \partial   g   }{  \partial   x_s  } }  dx_s \wedge f dx_I \wedge dx_J
}
{ =} { d(fdx_I ) \wedge g dx_J  + \sum_{s = 1}^n (-1)^k  f dx_I \wedge { \frac{   \partial   g   }{  \partial   x_s  } }  dx_s \wedge dx_J
}
{ =} { d(fdx_I ) \wedge g dx_J + (-1)^k  f dx_I \wedge  d( g dx_J)
}
}
{}{.}}
{}

\teilbeweis {}{}{}
{(4). Untuk suatu bentuk-$1$

\mavergleichskette
{\vergleichskette
{ \omega
}
{ = }{ \sum_{j = 1}^n g_jdx_j
}
{  }{
}
{    }{
}
{   }{
}
}
{}{}{}
dengan menggunakan

\mavergleichskette
{\vergleichskette
{ dx_i \wedge dx_j
}
{ = }{ -dx_j \wedge dx_i
}
{  }{
}
{    }{
}
{   }{
}
}
{}{}{}
diperoleh

\mavergleichskettealign
{\vergleichskettealign
{d \omega
}
{ =} { \sum_{j = 1 }^n  d(g_j dx_j)
}
{ =} { \sum_{j = 1 }^n   \left(  \sum_{i = 1}^n { \frac{   \partial  g_j   }{  \partial   x_i   } } dx_i \wedge dx_j  \right)
}
{ =} { \sum _{1 \leq i <j\leq n} \left(  { \frac{   \partial  g_j   }{  \partial   x_i   } } - { \frac{   \partial  g_i   }{  \partial   x_j   } }  \right) dx_i \wedge dx_j
}
{ } {
}
}
{}
{}{.}
Untuk fungsi $f$ yang dua kali terdiferensialkan secara kontinu,

\mavergleichskette
{\vergleichskette
{ df
}
{ = }{ \sum_{j = 1}^n g_j dx_j
}
{  }{
}
{    }{
}
{   }{
}
}
{}{}{}
dengan
\definitionsverweis {turunan parsial}{}{}

\mavergleichskette
{\vergleichskette
{ g_j
}
{ = }{ { \frac{   \partial   f   }{  \partial   x_j   } }
}
{  }{
}
{    }{
}
{   }{
}
}
{}{}{,}
sehingga

\mavergleichskette
{\vergleichskette
{ d(df)
}
{ = }{ 0
}
{  }{
}
{    }{
}
{   }{
}
}
{}{}{}
menurut
teorema Schwarz.

Untuk bentuk diferensial berderajat $k$, kita mengambil

\mavergleichskette
{\vergleichskette
{ \omega
}
{ = }{ f dx_{i_1 } \wedge \ldots \wedge dx_{i_k }
}
{  }{
}
{    }{
}
{   }{
}
}
{}{}{}
dan memperoleh

\mavergleichskettedisp
{\vergleichskette
{ d(d \omega)
}
{ =} { d (df \wedge  dx_{i_1 } \wedge \ldots \wedge dx_{i_k })
}
{ } {
}
{ } {
}
{ } {
}
}
{}{}{.}
Menurut aturan hasil kali (3), ungkapan ini merupakan jumlah dari
\mathl{k+1}{} produk eksterior; pada setiap produk itu, salah satu \anfuehrung{faktor produk-eksterior}{} berbentuk

\mavergleichskette
{\vergleichskette
{ d(dg)
}
{ = }{ 0
}
{  }{
}
{    }{
}
{   }{
}
}
{}{}{.}}
{}

\teilbeweis {}{}{}
{(5). Kita tulis

\mavergleichskettedisp
{\vergleichskette
{ \psi_i
}
{ =} { x_i \circ \psi
}
{ } {
}
{ } {
}
{ } {
}
}
{}{}{.}
Berkat linearitas turunan eksterior (2) dan
linearitas penarikan balik bentuk diferensial,
kita dapat mengambil

\mavergleichskette
{\vergleichskette
{ \omega
}
{ = }{ f dx_I
}
{  }{
}
{    }{
}
{   }{
}
}
{}{}{}
dengan

\mavergleichskette
{\vergleichskette
{ I
}
{ = }{ \{i_1 , \ldots , i_k\}
}
{  }{
}
{    }{
}
{   }{
}
}
{}{}{.}
Menurut
Soal 14.13
dan
Soal 14.16,
penarikan balik serasi dengan perkalian oleh fungsi skalar dan dengan produk eksterior. Dengan menggunakan aturan hasil kali (3), aturan (4), dan
aturan rantai
\zusatzklammer {dalam pengertian
Soal 14.14} {} {,}
diperoleh

\mavergleichskettealignhandlinks
{\vergleichskettealignhandlinks
{ d \left(  \psi^* \omega  \right)
}
{ =} { d \left(  \psi^* \left(  f dx_{i_1 } \wedge \ldots \wedge dx_{i_k }  \right)   \right)
}
{ =} { d \left(  \psi^*(f)  \psi^* \left(  dx_{i_1 } \wedge \ldots \wedge dx_{i_k }  \right)  \right)
}
{ =} { d \left(  \psi^*(f) \left(  \psi^* dx_{i_1 }  \right) \wedge \ldots \wedge \left(  \psi^* dx_{i_k }  \right)  \right)
}
{ =} { d \left(  \psi^*(f)  d \psi_{i_1 } \wedge \ldots \wedge d\psi_{i_k }  \right)
}
}
{
\vergleichskettefortsetzungalign
{ =} { d \left(  \psi^*(f)  \right) \wedge \left(  d \psi_{i_1 } \wedge \ldots \wedge d\psi_{i_k }   \right) + (\psi^*(f))  d \left(  d \psi_{i_1 } \wedge \ldots \wedge d\psi_{i_k }  \right)
}
{ =} { d \left(  \psi^*(f)  \right)  \wedge \left(  d \psi_{i_1 } \wedge \ldots \wedge d\psi_{i_k }  \right)
}
{ =} { \psi^* (df)  \wedge  d \psi_{i_1 } \wedge \ldots \wedge d\psi_{i_k }
}
{ =} { \psi^* ( df) \wedge \psi^*\left(  d x_{i_1 }  \right) \wedge \ldots \wedge \psi^*\left(  d x_{i_k }  \right)
}
}
{
\vergleichskettefortsetzungalign
{ =} { \psi^* \left(  df \wedge  d x_{i_1 } \wedge \ldots \wedge d x_{i_k }  \right)
}
{ =} { \psi^* \left(  d \left(  f d x_{i_1 } \wedge \ldots \wedge d x_{i_k }  \right)   \right)
}
{ } {}
{ } {}
}{.}}
{}

}

\inputdefinition
{}
{

Misalkan $M$ suatu
\definitionsverweis {manifold diferensiabel}{}{.}
Untuk suatu
\definitionsverweis {bentuk diferensial}{}{}
\definitionsverweis {terdiferensialkan}{}{}

\mavergleichskette
{\vergleichskette
{ \omega
}
{ \in }{
{ \mathcal E }^{  k  }_{ 1 } (  M   )
}
{  }{
}
{    }{
}
{   }{
}
}
{}{}{}
kita definisikan \definitionswort {turunan eksterior}{}

\mavergleichskette
{\vergleichskette
{ d\omega
}
{ \in }{
{ \mathcal E }^{ k+1 }_{ 0 } (  M   )
}
{  }{
}
{    }{
}
{   }{
}
}
{}{}{}
dengan mengacu pada
\definitionsverweis {kasus lokal}{}{}
dan peta-peta
\maabbdisp {\alpha} { U } { V
} {}
\zusatzklammer {dengan
\mavergleichskettek
{\vergleichskettek
{ U
}
{ \subseteq }{ M
}
{  }{
}
{    }{
}
{   }{
}
}
{}{}{}
dan

\mavergleichskettek
{\vergleichskettek
{ V
}
{ \subseteq }{ \R^n
}
{  }{
}
{    }{
}
{   }{
}
}
{}{}{}
terbuka} {} {}
melalui

\mavergleichskettedisp
{\vergleichskette
{ (d \omega){{|}}_U
}
{ =} { d \left(  \omega{{|}}_U  \right)
}
{ =} { \alpha^* \left(  d \left(  \alpha^{-1 *} \left(  \omega {{|}}_U  \right)  \right)  \right)
}
{ } {
}
{ } {
}
}
{}{}{.}

}

Jadi, bentuk diferensial yang dibatasi pada $U$ ditarik balik ke $V$, turunan eksteriornya diambil di sana sesuai aturan lokal, lalu hasilnya ditarik balik ke $U$. Kita perlu memastikan bahwa prosedur ini menghasilkan bentuk diferensial yang terdefinisi dengan baik pada $M$: untuk suatu titik

\mavergleichskette
{\vergleichskette
{ P
}
{ \in }{ M
}
{  }{
}
{    }{
}
{   }{
}
}
{}{}{}
hasil turunan eksterior tidak bergantung pada lingkungan koordinat yang digunakan. Misalkan diberikan dua peta untuk $P$; kita dapat langsung menganggap bahwa keduanya mempunyai domain definisi yang sama, yaitu $U$. Misalkan peta-peta itu
\maabbdisp {\alpha} { U } { V
} {}
dan
\maabbdisp {\beta} { U } { W
} {}
dan kita tetapkan

\mavergleichskette
{\vergleichskette
{\tau
}
{ = }{ \omega  \vert_U
}
{  }{
}
{    }{
}
{   }{
}
}
{}{}{.}
Dengan menerapkan
Lema 20.2 (5)
pada
\mathl{\alpha \circ \beta^{-1}}{} dan
\mathl{\alpha^{-1*} \tau}{,}
kita memperoleh

\mavergleichskettealignhandlinks
{\vergleichskettealignhandlinks
{ \beta^* \left(  d \left(  \beta^{-1*}(\tau)  \right)   \right)
}
{ =} { \left(  \beta \circ \alpha^{-1} \circ \alpha  \right)^* \left(  d \left(  \left(  \alpha^{-1} \circ \alpha  \circ \beta^{-1}  \right)^* (\tau )  \right)   \right)
}
{ =} {\alpha^* \left(  \beta \circ \alpha^{-1}  \right)^* \left(  d \left(  \left(  \alpha  \circ  \beta^{-1}  \right)^*  \alpha^{-1*} (\tau)  \right)   \right)
}
{ =} {\alpha^* \left(  d \left(  \alpha^{-1*}( \tau)  \right)   \right)
}
{ } {}
}
{}
{}{.}

Sifat-sifat dasar di atas juga terbawa ke manifold.


\inputfaktbeweis
{Differentialform/Mannigfaltigkeit/Äußere Ableitung/Eigenschaften/Fakt}
{Satz}
{}
{

\faktsituation {Misalkan $M$ suatu
\definitionsverweis {manifold diferensiabel}{}{,}

\mavergleichskette
{\vergleichskette
{ k
}
{ \in }{ \N
}
{  }{
}
{    }{
}
{   }{
}
}
{}{}{}
dan misalkan
\maabbeledisp {d} {
{ \mathcal E }^{  k  }_{ 1 } (  M   ) } {
{ \mathcal E }^{ k+1 }_{ 0 } (  M   )
} { \omega} { d \omega
} {,}
adalah
\definitionsverweis {turunan eksterior}{}{.}}
\faktuebergang {Maka berlaku sifat-sifat berikut.}
\faktfolgerung {\aufzaehlungfuenf{Turunan eksterior
\maabbdisp {d} {
{ \mathcal E }^{  0  }_{ 1 } (  M   ) } {
{ \mathcal E }^{  1  }_{ 0 } (  M   )
} {,}
adalah
\definitionsverweis {pemetaan singgung}{}{.}
}{Turunan eksterior bersifat
$\R$-\definitionsverweis {linear}{}{.}
}{Untuk

\mavergleichskette
{\vergleichskette
{ \omega
}
{ \in }{
{ \mathcal E }^{  k  }_{ 1 } (  M   )
}
{  }{
}
{    }{
}
{   }{
}
}
{}{}{}
dan

\mavergleichskette
{\vergleichskette
{ \tau
}
{ \in }{
{ \mathcal E }^{  \ell }_{ 1 } (  M   )
}
{  }{
}
{    }{
}
{   }{
}
}
{}{}{}
berlaku aturan hasil kali

\mavergleichskettedisp
{\vergleichskette
{ d( \omega \wedge \tau)
}
{ =} { ( d \omega) \wedge \tau  + (-1)^k \omega \wedge d \tau
}
{ } {
}
{ } {
}
{ } {
}
}
{}{}{}
}{Untuk setiap bentuk diferensial yang dua kali terdiferensialkan secara kontinu, $\omega$,

\mavergleichskette
{\vergleichskette
{ d(d \omega)
}
{ = }{ 0
}
{  }{
}
{    }{
}
{   }{
}
}
{}{}{.}
}{Misalkan $L$ suatu manifold diferensiabel lainnya. Untuk suatu
\definitionsverweis {pemetaan yang dua kali terdiferensialkan secara kontinu}{}{}
\maabbdisp {\psi} { L } { M
} {}
dan setiap

\mavergleichskette
{\vergleichskette
{ \omega
}
{ \in }{
{ \mathcal E }^{  k  }_{ 1 } (  M   )
}
{  }{
}
{    }{
}
{   }{
}
}
{}{}{}
berlaku bagi
\definitionsverweis {bentuk diferensial tarik balik}{}{}

\mavergleichskettedisp
{\vergleichskette
{ d (\psi^* \omega)
}
{ =} { \psi^* (d \omega)
}
{ } {
}
{ } {
}
{ } {
}
}
{}{}{.}
}}
\faktzusatz {}
\faktzusatz {}

}
{

Semua pernyataan ini bersifat lokal, sehingga langsung diperoleh dari
Lema 20.2.

}

\inputdefinition
{}
{

Misalkan $M$ suatu
\definitionsverweis {manifold diferensiabel}{}{.}
Suatu
\definitionsverweis {bentuk diferensial}{}{}
\definitionsverweis {terdiferensialkan}{}{}
$\omega$ pada $M$ disebut \definitionswort {tertutup}{,} jika
\definitionsverweis {turunan eksterior}{}{}

\mavergleichskette
{\vergleichskette
{ d \omega
}
{ = }{ 0
}
{  }{
}
{    }{
}
{   }{
}
}
{}{}{.}

}

\inputdefinition
{}
{

Misalkan $M$ suatu
\definitionsverweis {manifold diferensiabel}{}{.}
Untuk derajat positif, suatu
$k$-\definitionsverweis {bentuk diferensial}{}{}
$\omega$ pada $M$, disebut \definitionswort {eksak}{,} jika terdapat suatu bentuk diferensial yang
\definitionsverweis {terdiferensialkan}{}{}
dan berderajat \mathl{(k-1)}{}, yang kita nyatakan dengan $\sigma$, pada $M$ sedemikian sehingga

\mavergleichskette
{\vergleichskette
{ d \sigma
}
{ = }{ \omega
}
{  }{
}
{    }{
}
{   }{
}
}
{}{}{.}

}

Jadi, bentuk diferensial eksak adalah bentuk diferensial berderajat positif yang memiliki suatu \stichwort {bentuk primitif} {} $\sigma$. Dengan istilah-istilah ini, pernyataan di atas

\mavergleichskette
{\vergleichskette
{ dd
}
{ = }{ 0
}
{  }{
}
{    }{
}
{   }{
}
}
{}{}{}
dapat dirumuskan dengan mengatakan bahwa setiap bentuk eksak adalah tertutup. Jadi, sifat tertutup merupakan syarat perlu bagi keberadaan suatu bentuk primitif. Tanpa bukti, kita catat bahwa untuk bentuk berderajat positif kriteria perlu ini juga cukup pada $\R^n$. Namun, ekuivalensi tersebut sama sekali tidak berlaku pada setiap manifold.

  \zwischenueberschrift{Setengah ruang Euklides}

\inputdefinition
{}
{

Dalam dimensi positif, yang dimaksud dengan \definitionswort {setengah ruang Euklides}{} dengan
\definitionswort {dimensi}{} $n$ adalah himpunan

\mavergleichskettedisp
{\vergleichskette
{H
}
{ =} { \left\{  x \in \R^n   \mid   x_1 \geq 0        \right\}
}
{ } {
}
{ } {
}
{ } {
}
}
{}{}{}
dengan
\definitionsverweis {topologi terinduksi}{}{.}

}
Untuk

\mavergleichskette
{\vergleichskette
{ n
}
{ = }{ 0
}
{  }{
}
{    }{
}
{   }{
}
}
{}{}{}
kita gunakan konvensi bahwa setengah ruang Euklides berupa satu titik; untuk

\mavergleichskette
{\vergleichskette
{ n
}
{ = }{ 1
}
{  }{
}
{    }{
}
{   }{
}
}
{}{}{}
himpunan ini berupa interval
\mathl{[0, \infty)}{;} untuk

\mavergleichskette
{\vergleichskette
{ n
}
{ = }{ 2
}
{  }{
}
{    }{
}
{   }{
}
}
{}{}{}
himpunan ini merupakan suatu \stichwort {setengah bidang} {;} dan untuk

\mavergleichskette
{\vergleichskette
{ n
}
{ = }{ 3
}
{  }{
}
{    }{
}
{   }{
}
}
{}{}{}
himpunan ini merupakan suatu setengah ruang. Jika indeks koordinat $1$ diganti dengan indeks lain, atau \anfuehrung{$\leq$}{} digunakan sebagai ganti \anfuehrung{$\geq$ }{,} objek-objek ini juga disebut setengah ruang. Karena suatu setengah ruang $H$ tertutup dalam $\R^n$, suatu himpunan bagian

\mavergleichskette
{\vergleichskette
{ T
}
{ \subseteq }{ H
}
{  }{
}
{    }{
}
{   }{
}
}
{}{}{}
tertutup di $H$ jika dan hanya jika ia tertutup di $\R^n$. Ekuivalensi ini
\betonung{tidak}{} berlaku untuk himpunan terbuka. Sebagai contoh, seluruh ruang $H$ terbuka di $H$, tetapi tidak terbuka di $\R^n$. Himpunan

\mavergleichskettedisp
{\vergleichskette
{ \partial H
}
{ =} { \left\{  x \in \R^n   \mid   x_1 = 0        \right\}
}
{ } {
}
{ } {
}
{ } {
}
}
{}{}{}
merupakan bagian dari $H$ dan disebut \stichwort {batas} {} dari $H$. Batas ini homeomorfik dengan $\R^{n-1}$
\zusatzklammer {yang untuk $n=0$ harus ditafsirkan sebagai himpunan kosong} {} {.}
Dengan $H_+$ kita nyatakan bagian positif

\mavergleichskette
{\vergleichskette
{ H_+
}
{ = }{ \left\{  x \in \R^n   \mid   x_1 >0        \right\}
}
{  }{
}
{    }{
}
{   }{
}
}
{}{}{,}
yang merupakan himpunan bagian terbuka dari $\R^n$.

Setengah ruang merupakan model standar bagi manifold berbatas yang akan kita perkenalkan sekarang. Konsep ini memperumum konsep manifold. Contoh khas manifold berbatas adalah bola pejal tertutup; batasnya adalah sfer. Suatu titik di interior bola mempunyai lingkungan berupa bola terbuka yang lebih kecil, sehingga di sekitar titik itu bentuknya \anfuehrung{secara lokal}{} seperti $\R^3$. Titik pada batas bola tidak mempunyai lingkungan seperti itu: batasnya hadir dalam setiap lingkungan terbuka titik tersebut. Secara lokal, titik batas demikian tampak seperti titik batas suatu setengah ruang Euklides.

Citra koordinat peta-peta suatu manifold berbatas berupa himpunan-himpunan terbuka di dalam setengah ruang. Karena itu, untuk membahas pemetaan transisi, kita harus dapat berbicara tentang pemetaan diferensiabel yang didefinisikan pada setengah ruang. Definisi berikut memungkinkan hal tersebut.

\inputdefinition
{}
{

Misalkan

\mavergleichskette
{\vergleichskette
{ U
}
{ \subseteq }{ H
}
{  }{
}
{    }{
}
{   }{
}
}
{}{}{}
suatu
\definitionsverweis {himpunan bagian terbuka}{}{}
di dalam
\definitionsverweis {setengah ruang Euklides}{}{}

\mavergleichskette
{\vergleichskette
{ H
}
{ \subseteq }{ \R^n
}
{  }{
}
{    }{
}
{   }{
}
}
{}{}{,}

\mavergleichskette
{\vergleichskette
{ P
}
{ \in }{ U
}
{  }{
}
{    }{
}
{   }{
}
}
{}{}{}
suatu titik, dan misalkan
\maabbdisp {\varphi} { U } {\R^m
} {}
suatu
\definitionsverweis {pemetaan}{}{.}
Maka $\varphi$ disebut \definitionswort {terdiferensialkan secara kontinu}{} di $P$ jika terdapat suatu lingkungan terbuka

\mavergleichskette
{\vergleichskette
{ P
}
{ \in }{ V
}
{ \subseteq }{ \R^n
}
{    }{
}
{   }{
}
}
{}{}{}
dan suatu fungsi
\definitionsverweis {terdiferensialkan secara kontinu}{}{}
\maabbdisp {\tilde{\varphi}} { V } {\R^m
} {}
dengan

\mavergleichskette
{\vergleichskette
{ \tilde{\varphi} {{|}}_{U \cap V}
}
{ = }{\varphi {{|}}_{U \cap V}
}
{  }{
}
{    }{
}
{   }{
}
}
{}{}{.}

}
Jadi, konsep keterdiferensialan yang baru direduksi ke konsep lama. Untuk suatu himpunan terbuka

\mavergleichskette
{\vergleichskette
{ U
}
{ \subseteq }{ H
}
{  }{
}
{    }{
}
{   }{
}
}
{}{}{,}
yang tidak memotong batas $H$, definisi ini ekuivalen dengan definisi bagi suatu himpunan terbuka di $\R^n$.

Melalui strategi mendefinisikan konsep-konsep pada titik batas dengan mensyaratkan adanya lingkungan terbuka dan objek yang diperluas, banyak konsep penting terbawa ke situasi baru yang lebih umum ini; kita tidak akan selalu menguraikannya satu per satu. Sebagai contoh, sudah jelas apa yang dimaksud dengan \stichwort {difeomorfisme} {} antara himpunan-himpunan terbuka di dalam setengah ruang, serta apa yang dimaksud dengan diferensial total suatu pemetaan diferensiabel. Dengan latar ini, definisi manifold berbatas juga tidak mengejutkan.

\chapter{Lembar Kerja 20}
\setcounter{section}{20}

  \zwischenueberschrift{Soal latihan}

\inputaufgabe
{}
{

Tentukan
\definitionsverweis {turunan eksterior}{}{}
dari
$1$-\definitionsverweis {bentuk diferensial}{}{}

\mavergleichskettedisp
{\vergleichskette
{ \omega
}
{ =} { \left(  x^2-y^3  \right) dx+x^3y^2dy
}
{ } {
}
{ } {
}
{ } {
}
}
{}{}{}
pada $\R^2$.

}
{} {}

\inputaufgabe
{}
{

Tentukan
\definitionsverweis {turunan eksterior}{}{}
dari
$1$-\definitionsverweis {bentuk diferensial}{}{}

\mavergleichskettedisp
{\vergleichskette
{ \omega
}
{ =} { xy^2dx+yzdy+x^3dz
}
{ } {
}
{ } {
}
{ } {
}
}
{}{}{}
pada $\R^3$.

}
{} {}

\inputaufgabegibtloesung
{}
{

Hitung
\definitionsverweis {turunan eksterior}{}{}
$d \omega$ dari
$1$-\definitionsverweis {bentuk diferensial}{}{}

\mavergleichskettedisp
{\vergleichskette
{ \omega
}
{ =} { { \frac{   x^2 }{  y } } dx- { \frac{   x  }{  y^2 } } dy
}
{ } {
}
{ } {
}
{ } {
}
}
{}{}{}
pada

\mavergleichskette
{\vergleichskette
{ U
}
{ = }{ \left\{ (x,y) \in \R^2  \mid   y \neq 0         \right\}
}
{  }{
}
{    }{
}
{   }{
}
}
{}{}{.}

}
{} {}

\inputaufgabegibtloesung
{}
{

Hitung turunan eksterior $d \omega$ dari bentuk diferensial

\mavergleichskettedisp
{\vergleichskette
{ \omega
}
{ =} {  e^{xz} dx \wedge dy - xyz dx \wedge dz + \left(  \sin ( \cos (xy) )  +y^{10}z^{100}  \right) dy \wedge dz
}
{ } {
}
{ } {
}
{ } {
}
}
{}{}{}
pada
\mathl{\R^3}{.}

}
{} {}

\inputaufgabegibtloesung
{}
{

Misalkan
\maabbeledisp {f} {\R} {\R
} { t } {f(t)
} {,}
diberikan oleh

\mavergleichskettedisp
{\vergleichskette
{ f(t)
}
{ =} { { \frac{   \sin^{ 3  } (t^4)  }{  1+t^2 } }
}
{ } {
}
{ } {
}
{ } {
}
}
{}{}{}

a) Hitung turunan eksterior dari $f$.

b) Hitung turunan eksterior dari $fdt$.

}
{} {}

\inputaufgabe
{}
{

Tentukan
\definitionsverweis {turunan eksterior}{}{} dari
$2$-\definitionsverweis {bentuk diferensial}{}{}

\mathdisp {\omega = xdx \wedge dy+xy^2zdy \wedge dz+xe^ydx\wedge dz} {  }
 pada $\R^3$.

}
{} {}

\inputaufgabegibtloesung
{}
{

Tunjukkan bahwa
$1$-\definitionsverweis {bentuk diferensial}{}{}

\mavergleichskette
{\vergleichskette
{ \omega
}
{ = }{ (2x- \sin  y )dx-x\cos  ydy
}
{  }{
}
{    }{
}
{   }{
}
}
{}{}{}
pada $\R^2$
\definitionsverweis {tertutup}{}{}
dan juga
\definitionsverweis {eksak}{}{.}

}
{} {}

\inputaufgabe
{}
{

Misalkan

\mavergleichskette
{\vergleichskette
{ U
}
{ \subseteq }{ \R^n
}
{  }{
}
{    }{
}
{   }{
}
}
{}{}{}
suatu
\definitionsverweis {himpunan bagian terbuka}{}{}
dan

\mavergleichskette
{\vergleichskette
{ \omega
}
{ = }{ \sum_{i = 1}^n f_i dx_i
}
{  }{
}
{    }{
}
{   }{
}
}
{}{}{}
suatu
$1$-\definitionsverweis {bentuk diferensial}{}{,}
yang dinyatakan dengan
\definitionsverweis {fungsi-fungsi yang terdiferensialkan secara kontinu}{}{}
\maabb {f_i} { U } { \R
} {}.
Tunjukkan bahwa $\omega$
\definitionsverweis {tertutup}{}{}
jika dan hanya jika

\mavergleichskettedisp
{\vergleichskette
{ \partial_i f_j
}
{ =} { \partial_j f_i
}
{ } {
}
{ } {
}
{ } {
}
}
{}{}{}
untuk semua $i,j$.

}
{} {}

\inputaufgabe
{}
{

Misalkan

\mavergleichskette
{\vergleichskette
{ U
}
{ \subseteq }{ \R^n
}
{  }{
}
{    }{
}
{   }{
}
}
{}{}{}
suatu
\definitionsverweis {himpunan bagian terbuka}{}{}
dan $\omega$ suatu
\definitionsverweis {bentuk diferensial}{}{} berderajat $1$ yang
\definitionsverweis {terdiferensialkan secara kontinu}{}{}
pada $U$, dengan
\definitionsverweis {medan vektor}{}{}
$F$ pada $U$ yang bersesuaian menurut
Lema 16.3.
Tunjukkan bahwa $\omega$
\definitionsverweis {tertutup}{}{}
jika dan hanya jika $F$ memenuhi
\definitionsverweis {syarat keterintegralan}{}{.}

}
{} {}

\inputaufgabe
{}
{

Misalkan
\mathl{U \subseteq \R^n}{} suatu
\definitionsverweis {himpunan bagian terbuka}{}{}
dan $\omega$ suatu
\definitionsverweis {bentuk diferensial}{}{} berderajat $1$ yang
\definitionsverweis {terdiferensialkan secara kontinu}{}{}
pada $U$, dengan
\definitionsverweis {medan vektor}{}{}
$F$ pada $U$ yang bersesuaian menurut
Lema 16.3.
Tunjukkan bahwa $\omega$
\definitionsverweis {eksak}{}{}
jika dan hanya jika $F$ merupakan
\definitionsverweis {medan gradien}{}{.}

}
{} {}

\inputaufgabe
{}
{

Misalkan $\omega$ suatu
\definitionsverweis {bentuk diferensial}{}{} \definitionsverweis {terdiferensialkan}{}{}
pada manifold $M$, dan
\maabb {\varphi} { L } { M
} {}
suatu pemetaan yang dua kali terdiferensialkan secara kontinu dan didefinisikan pada manifold $L$.

\aufzaehlungzwei {Misalkan $\omega$
\definitionsverweis {eksak}{}{.}
Tunjukkan bahwa bentuk diferensial tarik balik $\varphi^*\omega$ juga eksak.
} {Misalkan $\omega$
\definitionsverweis {tertutup}{}{.}
Tunjukkan bahwa
\definitionsverweis {bentuk diferensial tarik balik}{}{}
$\varphi^*\omega$ juga tertutup.
}

}
{} {}

Untuk soal-soal berikut, bandingkan
Contoh 58.7 (Analisis (Osnabrück 2021–2023))
dan
Contoh 58.10 (Analisis (Osnabrück 2021–2023)).

\inputaufgabe
{}
{

Tunjukkan bahwa
$1$-\definitionsverweis {bentuk diferensial}{}{}

\mathdisp {-ydx+xdy} {  }

pada $\R^2$ tidak
\definitionsverweis {tertutup}{}{.}

}
{} {}

\inputaufgabe
{}
{

Tunjukkan bahwa
$1$-\definitionsverweis {bentuk diferensial}{}{}

\mathdisp {- { \frac{   y  }{  x^2+y^2 } }  dx+ { \frac{   x  }{  x^2+y^2 } } dy} {  }

pada $\R^2 \setminus \{0\}$
\definitionsverweis {tertutup}{}{,}
tetapi tidak
\definitionsverweis {eksak}{}{.}

}
{} {}

\inputaufgabegibtloesung
{}
{

Dalam dimensi positif, tunjukkan bahwa setiap
$n$-\definitionsverweis {bentuk diferensial}{}{}
kontinu $\omega$ pada suatu himpunan terbuka

\mavergleichskette
{\vergleichskette
{U
}
{ \subseteq }{\R^n
}
{  }{
}
{    }{
}
{   }{
}
}
{}{}{}
merupakan bentuk
\definitionsverweis {eksak}{}{.}

}
{} {}

\inputaufgabe
{}
{

Misalkan $M$ suatu
\definitionsverweis {manifold diferensiabel}{}{}
yang
\definitionsverweis {terhubung}{}{}
dan $\omega$ suatu
$1$-\definitionsverweis {bentuk diferensial}{}{}
\definitionsverweis {terdiferensialkan}{}{}
pada $M$. Tunjukkan bahwa $\omega$
\definitionsverweis {eksak}{}{}
jika dan hanya jika, untuk setiap lintasan
\definitionsverweis {terdiferensialkan secara kontinu}{}{}
\maabbdisp {\gamma} { [a,b] } { M
} {}
\definitionsverweis {integral lintasan}{}{}

\mathl{\int_\gamma \omega}{} hanya bergantung pada
\mathl{\gamma(a)}{} dan
\mathl{\gamma(b)}{}.

}
{} {}

\inputaufgabe
{}
{

Misalkan $M$ dan $N$ merupakan
\definitionsverweis {manifold diferensiabel}{}{}
dan
\maabbdisp {\varphi} { M } { N
} {}
suatu
\definitionsverweis {pemetaan yang dua kali terdiferensialkan secara kontinu}{}{.}
Misalkan $\omega$ suatu
$n$-\definitionsverweis {bentuk diferensial}{}{}
\definitionsverweis {terdiferensialkan}{}{}
pada $N$, dengan $n$ adalah dimensi $N$. Tunjukkan bahwa $\varphi^* \omega$ merupakan
\definitionsverweis {bentuk diferensial}{}{}
\definitionsverweis {tertutup}{}{}
pada $M$.

}
{} {}

\inputaufgabe
{}
{

Manakah di antara fungsi-fungsi berikut
\maabbdisp {} {\R_+} {\R
} {}
yang dapat diperluas menjadi fungsi
\definitionsverweis {diferensiabel}{}{}
pada
\definitionsverweis {titik batas}{}{}
$0$?
\aufzaehlungsechs{
\mathl{x^3+ \sin^{ 3  }  x -e^{-x}}{,}
}{
\mathl{{ \frac{   1  }{  x  } }}{,}
}{
\mathl{\sin  { \frac{   1  }{  x  } }}{,}
}{
\mathl{x \sin  { \frac{   1  }{  x  } }}{,}
}{
\mathl{{ \frac{   e^{ { \frac{   1  }{  x  } } }  }{  x  } }}{,}
}{
\mathl{x^2 \sin  { \frac{   1  }{  x  } }}{.}
}

}
{} {}

\inputaufgabe
{}
{

Misalkan
\mathl{H \subset \R^n}{} suatu
\definitionsverweis {setengah ruang}{}{.}
Misalkan
\mathl{Q\in H}{} suatu titik dan $Q \in U \subseteq H$, dengan $U$ suatu
\definitionsverweis {himpunan bagian terbuka}{}{}
dari $\R^n$. Tunjukkan bahwa $Q$ bukan titik batas $H$.

}
{} {}

\inputaufgabe
{}
{

Definisikan istilah \stichwort {difeomorfisme} {,} \stichwort {diferensial total} {} dan \stichwort {turunan tingkat tinggi} {} untuk
\definitionsverweis {setengah ruang}{}{}
\zusatzklammer {atau himpunan bagian terbuka darinya} {} {.}

}
{} {}

  \zwischenueberschrift{Soal untuk dikumpulkan}

\inputaufgabe
{3}
{

Tentukan
\definitionsverweis {turunan eksterior}{}{} dari
$1$-\definitionsverweis {bentuk diferensial}{}{}

\mathdisp {\omega = xy^2z^3dx+xyzdy+x^3yz^4dz} {  }
 pada $\R^3$.

}
{} {}

\inputaufgabe
{3}
{

Tentukan
\definitionsverweis {turunan eksterior}{}{} dari
$2$-\definitionsverweis {bentuk diferensial}{}{}

\mathdisp {\omega = xy^2dx \wedge dy+(x^3-y^2z^4)dy \wedge dz+ \sin (xy) dx \wedge dz} {  }

pada $\R^3$.

}
{} {}

\inputaufgabe
{5}
{

Misalkan

\mavergleichskette
{\vergleichskette
{ U
}
{ \subseteq }{ \R^n
}
{  }{
}
{    }{
}
{   }{
}
}
{}{}{}
\definitionsverweis {terbuka}{}{}
dan misalkan
\mathl{\omega_1 , \ldots , \omega_r}{} bentuk-bentuk diferensial pada $U$, dengan $\omega_i$ suatu
$k_i$-\definitionsverweis {bentuk diferensial}{}{.}
Temukan dan buktikan suatu rumus untuk

\mathdisp {d( \omega_1 \wedge \ldots \wedge \omega_r)} { . }

}
{} {}

\inputaufgabe
{5}
{

Tunjukkan bahwa
\definitionsverweis {bentuk diferensial}{}{}

\mavergleichskettedisp
{\vergleichskette
{ \omega
}
{ =} { \left(  2xy+3x^2-y e^{xy}  \right) dx + \left(  x^2-xe^{xy} +8y  \right)  dy
}
{ } {
}
{ } {
}
{ } {
}
}
{}{}{}
pada $\R^2$
\definitionsverweis {tertutup}{}{}
dan juga
\definitionsverweis {eksak}{}{.}

}
{} {}

\inputaufgabe
{6}
{

Tunjukkan bahwa
\definitionsverweis {setengah bidang}{}{}
$\R_{\geq 0} \times \R$ dan
\definitionsverweis {kuadran}{}{}

\mathl{\R_{\geq 0} \times \R_{\geq 0}}{}
\definitionsverweis {homeomorfik}{}{.}

}
{} {}

\section*{Solusi yang disediakan oleh sumber}
\addcontentsline{toc}{section}{Solusi yang disediakan oleh sumber}
\subsection*{Solusi Soal 20.3}
Kita peroleh

\mavergleichskettealign
{\vergleichskettealign
{d \omega
}
{ =} { d { \frac{   x^2 }{  y } } \wedge dx - d { \frac{   x  }{  y^2 } } \wedge dy
}
{ =} { - { \frac{   x^2 }{  y^2 } } dy  \wedge dx  -   { \frac{   1 }{  y^2 } }  dx \wedge dy
}
{ =} { { \frac{   x^2 -1 }{  y^2 } }  dx \wedge dy
}
{ } {
}
}
{}
{}{.}
\subsection*{Solusi Soal 20.4}
Kita peroleh

\mavergleichskettealign
{\vergleichskettealign
{ d \omega
}
{ =} { xe^{xz} dz \wedge d x \wedge dy - xzdy \wedge dx \wedge dz  -y \sin (xy)  \cos ( \cos (xy) )  dx \wedge dy \wedge dz
}
{ =} {  \left(  xe^{xz} +xz  -y \sin (xy)  \cos ( \cos (xy) )  \right) dx \wedge dy \wedge dz
}
{ } {
}
{ } {}
}
{}
{}{.}
\subsection*{Solusi Soal 20.5}
a) Berlaku
\mathl{df=f'dt}{,} dan, menurut aturan hasil bagi,

\mavergleichskettealign
{\vergleichskettealign
{f'
}
{ =} { { \frac{   (1+t^2)  3\sin^{ 2  } (t^4) \cos (t^4) 4 t^3  - 2t \sin^{ 3  } (t^4)  }{ (1+t^2)^2 } }
}
{ =} { { \frac{   12 (t^3+t^5) \sin^{ 2  } (t^4) \cos (t^4)   - 2t \sin^{ 3  } (t^4)  }{  1+2t^2+t^4 } }
}
{ } {
}
{ } {
}
}
{}
{}{.}

b) Turunan eksterior dari $fdt$ adalah
\mathl{f'dt \wedge dt =0}{.}
\subsection*{Solusi Soal 20.7}
\textit{Catatan edisi: solusi sumber ditulis dengan notasi ASCII yang rusak
dan menyebut $d\omega$ sebagai bentuk yang tertutup. Perhitungan berikut
menormalkan notasi, membuktikan bahwa $\omega$ tertutup, dan menampilkan
bentuk primitifnya secara eksplisit.}

Pertama,

\mavergleichskettealign
{\vergleichskettealign
{d\omega
}
{ =} { -\cos(y)\,dy\wedge dx-\cos(y)\,dx\wedge dy
}
{ =} { \cos(y)\,dx\wedge dy-\cos(y)\,dx\wedge dy
}
{ =} {0
}
{ } {}
}
{}
{}{.}

Jadi, $\omega$ tertutup. Selain itu,

\mavergleichskettedisp
{\vergleichskette
{\omega
}
{ =} {d\left(x^2-x\sin(y)\right)
}
{ } {
}
{ } {
}
{ } {
}
}
{}{}{,}

sehingga $\omega$ eksak.
\subsection*{Solusi Soal 20.14 (fragmen sumber; belum selesai)}
\textit{Catatan edisi: halaman sumber secara eksplisit berstatus belum
selesai. Halaman itu berhenti setelah menampilkan kandidat bentuk
berderajat $n-1$ dan sebuah rantai persamaan kosong. Jika setiap
$\partial_iH_i=h$, turunan kandidat yang ditampilkan sumber justru sama
dengan $-n\omega$, bukan $\omega$; sumber juga belum membuktikan bahwa
pilihan antiturunan pada irisan-irisan koordinat menghasilkan fungsi global
yang sesuai pada himpunan terbuka sembarang. Fragmen berikut dipertahankan
sebagai fragmen solusi sumber, bukan sebagai bukti lengkap.}

Kita dapat menuliskan bentuk tersebut sebagai

\mavergleichskettedisp
{\vergleichskette
{ \omega
}
{ =} { h(x_1 , \ldots , x_n) dx_1 \wedge \ldots \wedge dx_n
}
{ } {
}
{ } {
}
{ } {
}
}
{}{}{}
dengan suatu fungsi kontinu
\maabbdisp {h} { U } {\R
} {}
dengan
\mathl{x_1 , \ldots , x_n}{} menyatakan koordinat-koordinat pada $\R^n$.
Untuk setiap tupel tetap

\mathdisp {(x_1 , \ldots , x_{i-1},x_{i+1} , \ldots , x_n)} { , }

pembatasan fungsi $h$
pada satu komponen interval dari irisan koordinat hanya bergantung pada
$x_i$; kita nyatakan pembatasan ini dengan
\mathl{h_i(x_1 , \ldots , x_{i-1},-,x_{i+1} , \ldots , x_n)}{.}
Karena $h_i$ kontinu, pada komponen interval tersebut terdapat suatu
antiturunan
\mathl{H_i(x_1 , \ldots , x_{i-1},-,x_{i+1} , \ldots , x_n)}{.}
Turunan parsialnya terhadap $x_i$ sama dengan $h_i$. Sumber kemudian
menampilkan kandidat

\mathdisp {\sum_{i = 1}^n (-1)^{i} H_i dx_1 \wedge \ldots \wedge  dx_{i-1} \wedge d x_{i+1} \wedge \ldots \wedge  dx_n} { . }

\textit{Fragmen sumber berakhir di sini. Rantai persamaan kosong pada
halaman sumber tidak dimasukkan sebagai isi matematis.}

\part{Unit 21}
\chapter[Kuliah 21]{Kuliah 21: Manifold Berbatas dan Orientasi Batas}
\setcounter{section}{21}

  \zwischenueberschrift{Manifold berbatas}

\bild{ \begin{center}
\includegraphics[width=5.5cm]{ \bildeinlesungsvg {Runge_theorem} {svg} }
\end{center}
\bildtext {Suatu manifold berdimensi dua dengan batas. Batasnya terdiri atas empat busur tertutup.} }

\bildlizenz { Runge theorem.svg } {} {Oleg Alexandrov} {Commons} {PD} {}

\inputdefinition
{}
{

Misalkan

\mavergleichskette
{\vergleichskette
{ n
}
{ \in }{ \N
}
{  }{
}
{    }{
}
{   }{
}
}
{}{}{}
dan

\mavergleichskette
{\vergleichskette
{ k
}
{ \in }{ \overline{ \N }_+
}
{  }{
}
{    }{
}
{   }{
}
}
{}{}{.}
Suatu
\definitionsverweis {ruang topologis}{}{}
\definitionsverweis {Hausdorff}{}{}
$M$ bersama dengan suatu
\definitionsverweis {tutupan terbuka}{}{}

\mavergleichskette
{\vergleichskette
{ M
}
{ = }{ \bigcup_{i \in I} U_i
}
{  }{
}
{    }{
}
{   }{
}
}
{}{}{}
dan
\definitionsverweis {peta-peta}{}{}
\maabbdisp {\alpha_i} {U_i } { V_i
} {,}
dengan

\mavergleichskette
{\vergleichskette
{ V_i
}
{ \subseteq }{ H
}
{ \subset }{ \R^n
}
{    }{
}
{   }{
}
}
{}{}{}
merupakan himpunan-himpunan terbuka di dalam
\definitionsverweis {setengah ruang Euklides}{}{}
$H$ berdimensi $n$, serta dengan sifat bahwa
\definitionsverweis {pemetaan-pemetaan transisi}{}{}
\maabbdisp {\alpha_j \circ \alpha_i^{-1}} {V_i \cap \alpha_i(U_i \cap U_j) } { V_j \cap \alpha_j(U_i \cap U_j)
} {}
merupakan
$C^k$-\definitionsverweis {difeomorfisme}{}{}{,}
disebut \definitionswort {manifold-$C^k$ berbatas}{} atau \definitionswort {manifold diferensiabel berbatas}{}
\zusatzklammer {berderajat $k$} {} {,}
atau \definitionswort {manifold dengan batas}{.} Himpunan peta

\mathbed {(U_i,\alpha_i)} {}
 {i \in I} {}
 {} {} {} {,}
juga disebut \definitionswort {atlas-$C^k$}{} dari manifold berbatas tersebut.

}
Karena himpunan-himpunan terbuka di dalam setengah ruang $H$ yang sama sekali tidak memotong batas

\mavergleichskette
{\vergleichskette
{ \partial H
}
{ = }{ 0 \times \R^{n-1}
}
{  }{
}
{    }{
}
{   }{
}
}
{}{}{}
juga diperbolehkan, konsep ini mencakup konsep manifold diferensiabel. Batas suatu manifold berbatas
\zusatzklammer {yang segera akan kita definisikan dengan cara yang wajar} {} {}
dapat saja kosong. Hal ini terjadi tepat pada manifold diferensiabel \anfuehrung{biasa}{}.

\inputdefinition
{}
{

Misalkan $M$ suatu
\definitionsverweis {manifold berbatas}{}{.}
Maka \definitionswort {batas}{} dari $M$, yang ditulis
\mathl{\partial M}{,} didefinisikan oleh

\mavergleichskettedisp
{\vergleichskette
{ \partial M
}
{ =} { \left\{  x \in M  \mid   \alpha_i(x) \in  \partial H \cap V_i  \text{ untuk } \text{suatu } i \in I         \right\}
}
{ } {
}
{ } {
}
{ } {
}
}
{}{}{}
dengan $\alpha_i$ adalah peta-peta.

}
Pada definisi ini, peta mana pun dapat digunakan untuk menguji apakah suatu titik yang diberikan merupakan titik batas.


\inputfaktbeweis
{Mannigfaltigkeit mit Rand/Rand/Elementare Eigenschaften/Fakt}
{Lema}
{}
{

\faktsituation {Misalkan $M$ suatu
\definitionsverweis {manifold diferensiabel berbatas}{}{}

\mathl{\partial M}{.}}
\faktuebergang {Maka pernyataan-pernyataan berikut berlaku.}
\faktfolgerung {\aufzaehlungdrei{Suatu titik

\mavergleichskette
{\vergleichskette
{ P
}
{ \in }{ M
}
{  }{
}
{    }{
}
{   }{
}
}
{}{}{}
merupakan titik batas jika dan hanya jika hal ini berlaku bagi setiap
\definitionsverweis {peta}{}{,}
yang domain petanya memuat titik tersebut\zusatzfussnote {Pada manifold topologis berbatas, yaitu ketika pemetaan-pemetaan transisinya hanya berupa homeomorfisme, pernyataan ini juga berlaku, tetapi dengan pembuktian yang berbeda.} {} {.}
}{Batas $\partial M$ merupakan himpunan
\definitionsverweis {tertutup}{}{.}
}{Komplemen
\mathl{M \setminus \partial M}{} merupakan suatu
\definitionsverweis {manifold diferensiabel}{}{}
\zusatzklammer {tanpa batas} {} {.}
}}
\faktzusatz {}
\faktzusatz {}

}
{

\teilbeweis {}{}{}
{(1). Misalkan

\mavergleichskette
{\vergleichskette
{ P
}
{ \in }{ M
}
{  }{
}
{    }{
}
{   }{
}
}
{}{}{}
dan

\mavergleichskette
{\vergleichskette
{ P
}
{ \in }{ U
}
{  }{
}
{    }{
}
{   }{
}
}
{}{}{}
suatu domain peta dengan dua peta
\maabbdisp {\alpha_1} { U } { V_1
} {}
dan
\maabbdisp {\alpha_2} { U } { V_2
} {}
dengan
\mathkor {} {V_1 \subseteq H_1} {dan} {V_2 \subseteq H_2} {}
terbuka di dalam
\definitionsverweis {setengah ruang Euklides}{}{}

\mavergleichskette
{\vergleichskette
{ H_1, H_2
}
{ \cong }{ \R_{\geq 0} \times \R^{n-1}
}
{  }{
}
{    }{
}
{   }{
}
}
{}{}{.}
Pemetaan pergantian peta

\mavergleichskette
{\vergleichskette
{ \varphi
}
{ = }{ \alpha_2 \circ \alpha_1^{-1}
}
{  }{
}
{    }{
}
{   }{
}
}
{}{}{}
merupakan suatu
\definitionsverweis {difeomorfisme}{}{,}
dan menurut Soal 21.14 ini berarti bahwa bagi setiap titik

\mavergleichskette
{\vergleichskette
{ Q
}
{ \in }{ V_1
}
{  }{
}
{    }{
}
{   }{
}
}
{}{}{,}
terdapat lingkungan-lingkungan terbuka
\mathkor {} {W_1} {dan} {W_2} {}
di $\R^n$ dengan

\mavergleichskette
{\vergleichskette
{ Q
}
{ \in }{ W_1
}
{  }{
}
{    }{
}
{   }{
}
}
{}{}{}
serta suatu perluasan
\definitionsverweis {difeomorfik}{}{}
\maabbdisp {\tilde{\varphi}} {W_1 } { W_2
} {}
dari
\mathl{\varphi\!\mid_{W_1 \cap H_1 }}{.} Oleh karena itu,
\mathl{\tilde{\varphi} \left(  H_1^+ \cap W_1  \right)}{} terbuka di $W_2$.

Sekarang, tetapkan

\mavergleichskette
{\vergleichskette
{ Q_1
}
{ = }{ \alpha_1(P)
}
{  }{
}
{    }{
}
{   }{
}
}
{}{}{}
dan pilih
\mathl{W_1,W_2, \tilde{\varphi}}{} dengan sifat-sifat yang baru saja disebutkan. Jika $Q_1$ bukan titik batas dalam peta pertama, maka

\mavergleichskette
{\vergleichskette
{ Q_1
}
{ \in }{ H_1^+ \cap W_1
}
{  }{
}
{    }{
}
{   }{
}
}
{}{}{}
merupakan suatu lingkungan terbuka dan karenanya

\mavergleichskette
{\vergleichskette
{ \alpha_2 (P)
}
{ \in }{ \tilde{\varphi} \left(  H_1^+ \cap W_1  \right)
}
{  }{
}
{    }{
}
{   }{
}
}
{}{}{}
merupakan suatu lingkungan terbuka di

\mavergleichskette
{\vergleichskette
{ W_2
}
{ \subseteq }{ \R^n
}
{  }{
}
{    }{
}
{   }{
}
}
{}{}{.}
Selain itu,

\mavergleichskette
{\vergleichskette
{ \tilde{\varphi} \left(  H_1^+ \cap W_1  \right)
}
{ \subseteq }{ H_2
}
{  }{
}
{    }{
}
{   }{
}
}
{}{}{.}
Dengan kata lain,

\mavergleichskette
{\vergleichskette
{ \alpha_2(P)
}
{ \in }{ H_2
}
{  }{
}
{    }{
}
{   }{
}
}
{}{}{}
memiliki suatu lingkungan yang terbuka di $\R^n$ dan termuat di dalam $H_2$; maka menurut Soal 21.13, titik ini juga tidak mungkin menjadi titik batas dalam peta kedua.}
{}
\teilbeweis {}{}{}
{(2). Misalkan

\mavergleichskette
{\vergleichskette
{ P
}
{ \notin }{ \partial M
}
{  }{
}
{    }{
}
{   }{
}
}
{}{}{}
dan misalkan

\mavergleichskette
{\vergleichskette
{ P
}
{ \in }{ U
}
{  }{
}
{    }{
}
{   }{
}
}
{}{}{}
suatu domain peta dengan
\definitionsverweis {homeomorfisme}{}{}
\maabbdisp {\alpha} { U } { V
} {}
dengan

\mavergleichskette
{\vergleichskette
{ V
}
{ \subseteq }{ \R_{\geq 0} \times \R^{n-1}
}
{ = }{ H
}
{    }{
}
{   }{
}
}
{}{}{}
terbuka. Karena $P$ bukan titik batas, komponen pertama
\mathl{\alpha(P)}{} positif, sehingga terdapat suatu himpunan terbuka

\mavergleichskette
{\vergleichskette
{ \alpha(P)
}
{ \in }{ V'
}
{ \subseteq }{ \R_{+} \times \R^{n-1}
}
{ =   }{ H_+
}
{   }{
}
}
{}{}{.}
Dengan demikian,

\mavergleichskette
{\vergleichskette
{ U'
}
{ = }{ \alpha^{-1}(V')
}
{  }{
}
{    }{
}
{   }{
}
}
{}{}{}
merupakan suatu lingkungan terbuka dari $P$ yang
\zusatzklammer {menurut Bagian 1} {} {}
tidak memotong batas.}
{}
\teilbeweis {}{}{}
{(3). Bagi setiap titik

\mavergleichskette
{\vergleichskette
{ P
}
{ \notin }{ \partial M
}
{  }{
}
{    }{
}
{   }{
}
}
{}{}{}
kita dapat menentukan, seperti pada (2), suatu domain peta yang terpisah dari batas dan citra petanya merupakan himpunan terbuka di $\R^n$. Jadi, kita memperoleh suatu
\definitionsverweis {manifold}{}{}
\zusatzklammer {tanpa batas} {} {.}}
{}

}

Konsep pemetaan diferensiabel, difeomorfisme, dan ruang singgung juga terbawa ke manifold berbatas; lihat
Soal 21.2.



\inputfaktbeweis
{Mannigfaltigkeit mit Rand/Rand ist Mannigfaltigkeit/Fakt}
{Teorema}
{}
{

\faktsituation {Misalkan $M$ suatu
\definitionsverweis {manifold diferensiabel berbatas}{}{} berdimensi $n$.}
\faktfolgerung {Maka batas
\mathl{\partial M}{} merupakan suatu
\definitionsverweis {manifold diferensiabel}{}{}
\zusatzklammer {tanpa batas} {} {} berdimensi
\mathl{n-1}{.}}
\faktzusatz {}
\faktzusatz {}

}
{

Pertama-tama,

\mavergleichskette
{\vergleichskette
{ L
}
{ = }{ \partial M
}
{  }{
}
{    }{
}
{   }{
}
}
{}{}{}
dengan
\definitionsverweis {topologi terinduksi}{}{}
merupakan suatu
\definitionsverweis {ruang Hausdorff}{}{.}
Misalkan

\mavergleichskette
{\vergleichskette
{ P
}
{ \in }{ L
}
{  }{
}
{    }{
}
{   }{
}
}
{}{}{}
dan misalkan
\maabbdisp {\alpha} { U } { V
} {}
suatu
\definitionsverweis {peta}{}{}
dengan

\mavergleichskette
{\vergleichskette
{ V
}
{ \subseteq }{ H
}
{  }{
}
{    }{
}
{   }{
}
}
{}{}{}
terbuka dan

\mavergleichskette
{\vergleichskette
{ \alpha(P)
}
{ \in }{ \partial H
}
{  }{
}
{    }{
}
{   }{
}
}
{}{}{.}
Karena $\alpha$ merupakan suatu
\definitionsverweis {homeomorfisme}{}{}
dan menurut
Lema 21.3 (1)
titik batas berkorespondensi dengan titik batas pada setiap peta, peta ini menginduksi suatu homeomorfisme
\maabbdisp {} {U \cap L} {V \cap \partial H
} {.}
Di sini,
\mathl{U \cap L}{} merupakan suatu lingkungan terbuka dari $P$ di $L$, sehingga himpunan-himpunan ini dapat kita gunakan sebagai domain peta. Sekarang tinjau suatu pergantian peta. Kita boleh langsung memulai dengan satu domain peta

\mavergleichskette
{\vergleichskette
{ U
}
{ \subseteq }{ M
}
{  }{
}
{    }{
}
{   }{
}
}
{}{}{}
dan dua peta
\maabb {\alpha_1} { U } { V_1
} {} dan \maabb {\alpha_2} { U } { V_2
} {}
dengan

\mavergleichskette
{\vergleichskette
{ V_1,V_2
}
{ \subseteq }{ H
}
{  }{
}
{    }{
}
{   }{
}
}
{}{}{}
terbuka. Dengan demikian, terdapat suatu
$C^1$-\definitionsverweis {difeomorfisme
}{}{}
\maabbdisp {\varphi=\alpha_2 \circ \alpha_1^{-1}} {V_1 } { V_2
} {.}
Pertama, hal ini berarti terdapat suatu homeomorfisme
\maabb {} { \partial H \cap V_1 } { \partial H \cap V_2
} {.}
Sifat difeomorfisme dari $\varphi$ berarti bahwa bagi setiap titik

\mavergleichskette
{\vergleichskette
{ P
}
{ \in }{ U \cap L
}
{  }{
}
{    }{
}
{   }{
}
}
{}{}{,}
terdapat lingkungan-lingkungan terbuka

\mavergleichskette
{\vergleichskette
{ \alpha_1(P)
}
{ \in }{ W_1
}
{ \subseteq }{ \R^n
}
{    }{
}
{   }{
}
}
{}{}{}
dan

\mavergleichskette
{\vergleichskette
{ \alpha_2(P)
}
{ \in }{ W_2
}
{ \subseteq }{ \R^n
}
{    }{
}
{   }{
}
}
{}{}{}
serta suatu
\definitionsverweis {perluasan}{}{}
\maabbdisp {\tilde{\varphi}} { W_1 } { W_2
} {}
yang difeomorfik dari $\varphi$, dari
\mathkor {} {V_1 \cap W_1} {ke} {V_2 \cap W_2} {.}
Menurut Soal 21.19, perluasan ini juga menginduksi suatu
$C^1$-\definitionsverweis {difeomorfisme}{}{}
antara batas-batas
\mathl{\partial H \cap W_1}{} dan
\mathl{\partial H \cap W_2}{,} sehingga secara keseluruhan terdapat suatu difeomorfisme
\maabbdisp {\varphi \vert_{\partial H}} { \partial H \cap V_1 } { \partial H \cap V_2
} {.}

}

Kita sudah mengetahui bahwa serat suatu pemetaan reguler diferensiabel antara manifold-manifold diferensiabel memiliki struktur submanifold tertutup. Teorema berikut didasarkan pada argumen serupa dan menjamin keberadaan sangat banyak manifold berbatas.


\inputfaktbeweis
{Reguläre Funktion auf Mannigfaltigkeit/Urbild halbseitiger Intervalle/Mannigfaltigkeit mit Rand/Fakt}
{Teorema}
{}
{

\faktsituation {Misalkan $L$ suatu
\definitionsverweis {manifold diferensiabel}{}{}
dan misalkan
\maabbdisp {f} { L } {\R
} {}
suatu
\definitionsverweis {fungsi yang terdiferensialkan secara kontinu}{}{.}
Misalkan

\mavergleichskette
{\vergleichskette
{ a
}
{ \in }{ \R
}
{  }{
}
{    }{
}
{   }{
}
}
{}{}{.}}
\faktvoraussetzung {Setiap titik dalam
\definitionsverweis {serat}{}{}

\mavergleichskette
{\vergleichskette
{ F
}
{ = }{ f^{-1}(a)
}
{  }{
}
{    }{
}
{   }{
}
}
{}{}{}
di atas $a$ merupakan titik
\definitionsverweis {reguler}{}{.}}
\faktfolgerung {Maka himpunan-himpunan bagian

\mathdisp {M= \left\{  x \in L  \mid   f(x) \leq a        \right\} \text{   dan } N= \left\{  x \in L  \mid   f(x) \geq a        \right\}} {  }

merupakan
\definitionsverweis {manifold diferensiabel berbatas}{}{,}
dan \definitionsverweis {batas}{}{} masing-masing sama dengan $F$.}
\faktzusatz {}
\faktzusatz {}

}
{

Untuk menyederhanakan notasi, misalkan

\mavergleichskette
{\vergleichskette
{ a
}
{ = }{ 0
}
{  }{
}
{    }{
}
{   }{
}
}
{}{}{,}
dan kita batasi perhatian pada $M$. Kita mempunyai

\mavergleichskettedisp
{\vergleichskette
{ M
}
{ =} { \left\{  x\in L  \mid  f(x)<0        \right\}  \uplus F
}
{ } {
}
{ } {
}
{ } {
}
}
{}{}{,}
dengan himpunan di sebelah kiri merupakan suatu
\definitionsverweis {himpunan terbuka}{}{}
di $L$ dan dengan demikian merupakan suatu submanifold terbuka. Jadi, hal yang menentukan ialah menunjukkan bahwa bagi titik-titik dari serat $F$ terdapat peta-peta, dan dengan demikian suatu struktur manifold. Misalkan

\mavergleichskette
{\vergleichskette
{ P
}
{ \in }{ F
}
{  }{
}
{    }{
}
{   }{
}
}
{}{}{.}
Menurut pembuktian
teorema pemetaan implisit,
bagi

\mavergleichskette
{\vergleichskette
{ P
}
{ \in }{ F
}
{  }{
}
{    }{
}
{   }{
}
}
{}{}{}
terdapat suatu lingkungan
\zusatzklammer {koordinat} {} {}
terbuka

\mavergleichskette
{\vergleichskette
{ P
}
{ \in }{ U
}
{ \subseteq }{ L
}
{    }{
}
{   }{
}
}
{}{}{}
dan suatu peta
\maabbdisp {\alpha} { U } { V
} {,}

\mavergleichskette
{\vergleichskette
{ V
}
{ \subseteq }{ \R^n
}
{  }{
}
{    }{
}
{   }{
}
}
{}{}{,}
sedemikian sehingga

\mavergleichskette
{\vergleichskette
{ f
}
{ = }{ x_1 \circ \alpha
}
{  }{
}
{    }{
}
{   }{
}
}
{}{}{.}
Dalam hal ini,
\mathl{F \cap U}{} berkorespondensi dengan
\mathl{\left\{  x \in V  \mid   x_1 = 0         \right\}}{} dan
\mathl{M \cap U}{} berkorespondensi dengan
\mathl{H_{\leq 0} \cap V}{,} sehingga pembatasan $\alpha$ pada $M$ menghasilkan suatu peta untuk $M$ di $P$. Pergantian-pergantian petanya
$C^1$-\definitionsverweis {difeomorfik}{}{,}
karena hal ini berlaku bagi peta-peta
\zusatzklammer {penuh} {} {}
pada $L$.

}

\inputbeispiel{}
{

Untuk $r>0$, \definitionsverweis {bola tertutup}{}{}

\mavergleichskettedisp
{\vergleichskette
{ B  \left(  0 ,r \right)
}
{ =} { \left\{  x \in \R^n   \mid   \sqrt{x_1^2 + \cdots + x_n^2} \leq r         \right\}
}
{ } {
}
{ } {
}
{ } {
}
}
{}{}{}
merupakan suatu
$C^\infty$-\definitionsverweis {manifold diferensiabel}{}{}
dengan \stichwort {sfer} {}

\mavergleichskettedisp
{\vergleichskette
{ S \left(   0 , r  \right)
}
{ =} { \left\{  x \in \R^n   \mid   \sqrt{x_1^2 + \cdots + x_n^2} =  r         \right\}
}
{ } {
}
{ } {
}
{ } {
}
}
{}{}{}
sebagai
\definitionsverweis {batas}{}{.}
Hal ini langsung mengikuti
Teorema 21.5
yang diterapkan pada fungsi mulus
\maabbeledisp {g} {\R^n } {\R
} {(x_1 , \ldots , x_n) } { x_1^2 + \cdots + x_n^2
} {,}
dengan nilai reguler $r^2$; setiap titik $x\in g^{-1}(r^2)$ memenuhi
$x\neq 0$, sehingga $g$ bersifat
\definitionsverweis {reguler}{}{} pada seluruh serat itu.

}

\inputbeispiel{}
{

Suatu
\definitionsverweis {balok}{}{} \definitionsverweis {tertutup}{}{}

\mathdisp {[a_1,b_1] \times [a_2,b_2] \times \cdots \times [a_n,b_n] \subseteq \R^n} {  }

untuk
\mathl{n \geq 2}{} bukan merupakan suatu
\definitionsverweis {manifold diferensiabel berbatas}{}{,}
karena balok itu tidak hanya mempunyai sisi, tetapi juga
\zusatzklammer {bergantung pada dimensinya} {} {}
sudut dan rusuk. Suatu persegi panjang mempunyai empat titik sudut yang tidak dapat diberi struktur manifold diferensiabel berbatas
\zusatzklammer {yang kompatibel dengan ruang sekitar} {} {}
\zusatzklammer {karena persegi panjang tertutup
\definitionsverweis {homeomorfik}{}{}
dengan cakram tertutup, suatu struktur manifold diferensiabel berbatas dapat didefinisikan padanya; tetapi dengan struktur ini, penyematan alami persegi panjang tersebut ke dalam $\R^2$ tidak diferensiabel} {} {,}
sedangkan suatu balok berdimensi tiga mempunyai dua belas rusuk dan delapan sudut yang tidak memiliki struktur manifold alami.

Namun, jika sudut, rusuk, dan seterusnya tersebut dihapus dan hanya \anfuehrung{sisi-sisi berkodimensi satu}{} dipertahankan, kita memperoleh suatu
\zusatzklammer {tak kompak} {} {}
manifold berbatas. Batasnya merupakan gabungan saling lepas dari muka-muka hiper tersebut. Seperti pada setiap manifold berbatas, batas ini tertutup di dalam manifold, tetapi tidak tertutup di dalam ruang Euklides sekitar.

}

  \zwischenueberschrift{Orientasi pada manifold berbatas}

Misalkan pada $\R^n$ vektor-vektor standar
\mathl{e_1 , \ldots , e_n}{} memberikan suatu
\definitionsverweis {orientasi}{}{;}
selanjutnya pilih setengah ruang

\mavergleichskettedisp
{\vergleichskette
{ H_{\leq 0}
}
{ =} { \left\{  x \in \R^n   \mid   x_1 \leq 0         \right\}
}
{ } {
}
{ } {
}
{ } {
}
}
{}{}{}
sebagai \anfuehrung{setengah ruang dalam}{}. Orientasi pada hiperbidang
\zusatzklammer {yaitu batas manifold berbatas \mathlk{H_{\leq 0}}{}} {} {}

\mavergleichskettedisp
{\vergleichskette
{ E
}
{ =} { \left\{  x \in \R^n   \mid   x_1 = 0         \right\}
}
{ } {
}
{ } {
}
{ } {
}
}
{}{}{}
yang didefinisikan oleh basis
\mathl{e_2 , \ldots , e_n}{} disebut \stichwort {orientasi oleh normal luar} {.} Suatu basis sembarang
\mathl{v_2 , \ldots , v_n}{} dari $E$ merepresentasikan orientasi ini jika dan hanya jika, bagi suatu vektor sembarang

\mavergleichskette
{\vergleichskette
{ v
}
{ \in }{ H_+
}
{  }{
}
{    }{
}
{   }{
}
}
{}{}{}
\zusatzklammer {artinya, mengarah ke \anfuehrung{luar}{,} yakni keluar dari setengah ruang} {} {}
basis
\mathl{v,v_2 , \ldots , v_n}{}
\zusatzklammer {jadi $v$ ditempatkan pertama} {} {}
dari $\R^n$ merepresentasikan orientasi semula
\zusatzfussnote {Untuk suatu setengah garis

\mavergleichskette
{\vergleichskette
{H
}
{ = }{ \R_{\geq 0}
}
{ \subseteq }{ \R
}
{    }{
}
{   }{
}
}
{}{}{}
dengan satu-satunya titik batas $\{0\}$, hal ini harus ditafsirkan sebagai berikut. Kedua orientasi pada
\mathl{\{0\}}{} adalah
\mathkor {} {+} {dan} {-} {,}
dan $-$ merepresentasikan orientasi oleh normal luar, sebab bagi suatu vektor yang mengarah ke luar

\mavergleichskette
{\vergleichskette
{ w
}
{ \in }{ \R_-
}
{  }{
}
{    }{
}
{   }{
}
}
{}{}{}
vektor lawannya $-w$ merepresentasikan orientasi standar dari $\R$. Sebaliknya, untuk setengah ruang negatif
\mathl{\R_{\leq 0}}{} di titik nol, $+$ merepresentasikan orientasi oleh normal luar} {.} {.}

Hubungan antara orientasi pada ruang vektor real dan orientasi pada batas suatu setengah ruang ini terbawa ke manifold berbatas. Hal yang penting di sini ialah bahwa ruang singgung
\mathl{T_PM}{} di suatu titik batas $P$ memuat suatu hiperbidang kanonik, yaitu ruang singgung
\mathl{T_P (\partial M)}{} dari batas. Dalam hal ini manifold tersebut menentukan suatu \anfuehrung{separuh dalam}{} dan suatu \anfuehrung{separuh luar}{} dari ruang singgung.



\inputfaktbeweis
{Mannigfaltigkeit mit Rand/Orientierung/Randorientierung/Fakt}
{Teorema}
{}
{

\faktsituation {Misalkan $M$ suatu
\definitionsverweis {manifold diferensiabel berbatas}{}{,}}
\faktvoraussetzung {yang dilengkapi dengan suatu
\definitionsverweis {orientasi}{}{.}}
\faktfolgerung {Maka
\definitionsverweis {manifold batas}{}{} juga dilengkapi dengan suatu orientasi kanonik, yaitu orientasi yang pada setiap
\definitionsverweis {peta}{}{} ditentukan oleh
\definitionsverweis {normal luar}{}{.}}
\faktzusatz {}
\faktzusatz {}

}
{

Untuk setiap peta
\zusatzklammer {berorientasi} {} {}
\maabbdisp {\alpha} { U } { V
} {}
dengan

\mavergleichskette
{\vergleichskette
{ U
}
{ \subseteq }{ M
}
{  }{
}
{    }{
}
{   }{
}
}
{}{}{}
terbuka, peta terinduksi
\maabbdisp {} {U \cap \partial M } { V \cap \partial H
} {}
dilengkapi dengan
\definitionsverweis {orientasi oleh normal luar}{}{}
pada
\mathl{\partial H}{.} Menurut asumsi, semua pergantian peta pada $M$ mempunyai
\definitionsverweis {determinan fundamental}{}{}
positif di setiap titik terhadap basis-basis yang merepresentasikan orientasi; kita harus menunjukkan bahwa hal ini juga berlaku bagi pergantian peta yang terinduksi. Kita dapat memulai dengan suatu domain peta terbuka

\mavergleichskette
{\vergleichskette
{ P
}
{ \in }{ U
}
{ \subseteq }{ M
}
{    }{
}
{   }{
}
}
{}{}{}
dan dua peta
\maabbdisp {\alpha_1} { U } { V_1
} {}
dan
\maabbdisp {\alpha_2} { U } { V_2
} {}
serta meninjau pemetaan transisi
\maabbdisp {\varphi = \alpha_2 \circ \alpha_1^{-1}} { V_1 } { V_2
} {}
dengan himpunan-himpunan terbuka

\mavergleichskette
{\vergleichskette
{ V_1
}
{ \subseteq }{ H_1
}
{ \subset }{ \R^n
}
{    }{
}
{   }{
}
}
{}{}{}
dan

\mavergleichskette
{\vergleichskette
{ V_2
}
{ \subseteq }{ H_2
}
{ \subset }{ \R^n
}
{    }{
}
{   }{
}
}
{}{}{.}
Misalkan
\mathl{v_2 , \ldots , v_n}{} suatu basis dari
\mathl{\partial H_1}{} dan

\mavergleichskette
{\vergleichskette
{ v_1
}
{ \in }{ \R^n
}
{  }{
}
{    }{
}
{   }{
}
}
{}{}{}
sedemikian sehingga $v_1$ merepresentasikan normal luar dari $H_1$, yakni
\mathl{v_1, v_2 , \ldots , v_n}{} merepresentasikan orientasi $\R^n$
\zusatzklammer {misalkan $x_1 , \ldots , x_n$ adalah koordinat-koordinat yang bersesuaian} {} {;}
demikian pula,

\mavergleichskette
{\vergleichskette
{ w_1, w_2 , \ldots , w_n
}
{ \in }{ \R^n
}
{  }{
}
{    }{
}
{   }{
}
}
{}{}{}
harus memenuhi sifat-sifat yang bersesuaian. Kita tuliskan matriks fundamental dari $\varphi$ terhadap basis-basis ini, yakni matriks dengan entri-entri

\mathdisp {\left(  { \frac{   \partial  \varphi_i   }{  \partial   x_j   } } ( Q)  \right)_{1 \leq i,j \leq n}} { . }

Menurut asumsi, determinannya positif. Karena

\mavergleichskette
{\vergleichskette
{ \varphi ( \partial H_1)
}
{ \subseteq }{ \partial H_2
}
{  }{
}
{    }{
}
{   }{
}
}
{}{}{}
maka untuk suatu titik

\mavergleichskette
{\vergleichskette
{ Q
}
{ \in }{ \partial H_1
}
{  }{
}
{    }{
}
{   }{
}
}
{}{}{}
berlaku hubungan

\mavergleichskettedisp
{\vergleichskette
{ { \frac{   \partial  \varphi_1   }{  \partial   x_j   } } ( Q)
}
{ =} { 0
}
{ } {
}
{ } {
}
{ } {
}
}
{}{}{}
untuk

\mavergleichskette
{\vergleichskette
{ j
}
{ = }{ 2 , \ldots , n
}
{  }{
}
{    }{
}
{   }{
}
}
{}{}{.}
Menurut
teorema ekspansi,
oleh karena itu tanda determinan matriks

\mathdisp {\left(  { \frac{   \partial  \varphi_i   }{  \partial   x_j   } } ( Q)  \right)_{ 2 \leq i,j \leq n}} { , }

yang merupakan matriks fundamental dari pemetaan transisi peta-peta batas
\maabbdisp {} {V_1 \cap \partial H_1 } { V_2 \cap \partial H_2
} {}
\zusatzklammer {terhadap basis
\mathl{v_2 , \ldots , v_n}{} dan
\mathl{w_2 , \ldots , w_n}{}} {} {}
di titik $Q$, hanya bergantung pada
\mathl{{ \frac{   \partial  \varphi_1   }{  \partial   x_1   } } ( Q)}{.} Dengan

\mavergleichskette
{\vergleichskette
{ Q
}
{ = }{ (0,a_2 , \ldots , a_n)
}
{  }{
}
{    }{
}
{   }{
}
}
{}{}{}
menurut
Lema 44.3 (Analisis (Osnabrück 2021--2023)),
kita memperoleh hubungan

\mavergleichskettedisp
{\vergleichskette
{ { \frac{   \partial  \varphi_1   }{  \partial   x_1   } } ( Q)
}
{ =} { \operatorname{lim}_{   \epsilon \rightarrow  0   } \,  { \frac{   \varphi_1 (\epsilon, a_2 , \ldots , a_n)  }{  \epsilon } }
}
{ } {
}
{ } {
}
{ } {
}
}
{}{}{.}
Karena $v_1$ merepresentasikan normal luar, untuk nilai negatif
\zusatzklammer {dengan nilai mutlak yang cukup kecil} {} {}
$\epsilon$, vektor dengan koordinat

\mavergleichskette
{\vergleichskette
{ (\epsilon,a_2 , \ldots , a_n)
}
{ \in }{ H_1
}
{  }{
}
{    }{
}
{   }{
}
}
{}{}{.}
Karena itu, vektor citranya harus termasuk dalam $H_2$, dan dengan demikian

\mavergleichskette
{\vergleichskette
{ \varphi_1(\epsilon,a_2 , \ldots , a_n)
}
{ < }{ 0
}
{  }{
}
{    }{
}
{   }{
}
}
{}{}{.}
Jadi, hasil bagi ini
\mathl{\geq 0}{,} dan hal yang sama berlaku bagi limitnya. Karena terdapat suatu difeomorfisme, limit tersebut bahkan harus positif, dan pernyataan pun terbukti.

}

\chapter{Lembar Kerja 21}
\setcounter{section}{21}

  \zwischenueberschrift{Soal latihan}

\inputaufgabe
{}
{

Deskripsikan berbagai jenis pakaian sebagai
\definitionsverweis {manifold berbatas}{}{}
berdimensi dua.

}
{} {}

\inputaufgabe
{}
{

Definisikan konsep pemetaan diferensiabel, difeomorfisme, ruang singgung, bundel tangen, dan seterusnya untuk
\definitionsverweis {manifold berbatas}{}{.}

}
{} {}

\inputaufgabe
{}
{

Deskripsikan
\definitionsverweis {interval riil}{}{}
sebagai
\definitionsverweis {manifold berbatas}{}{.}
Apa
\definitionsverweis {batas}{}{}
masing-masing interval tersebut, dan interval mana saja yang saling
\definitionsverweis {difeomorfik}{}{?}

}
{} {}

\inputaufgabe
{}
{

$\,$

\bild{ \begin{center}
\includegraphics[width=5.5cm]{ \bildeinlesungsvg {Circle_on_sphere_wireframe_10deg_6r} {svg} }
\end{center}
\bildtext {} }

\bildlizenz { Circle on sphere wireframe 10deg 6r.svg } {} {Itai} {Commons} {CC BY 3.0} {}

Tunjukkan bahwa potongan permukaan berwarna biru maupun merah, masing-masing bersama garis pembatasnya, merupakan
\definitionsverweis {manifold berbatas}{}{.}
Apa batasnya? Apakah kedua manifold tersebut
\definitionsverweis {difeomorfik}{}{?}
Adakah manifold yang lebih sederhana dan difeomorfik dengannya?

}
{} {}

\inputaufgabe
{}
{

Misalkan $M$ suatu
\definitionsverweis {manifold diferensiabel}{}{}
\zusatzklammer {tanpa batas} {} {}
dan $N$ suatu
\definitionsverweis {manifold diferensiabel berbatas}{}{.}
Apa yang dapat dikatakan tentang
\definitionsverweis {produk}{}{}
\mathl{M \times N}{?}

}
{} {}

\inputaufgabe
{}
{

Tunjukkan bahwa anulus tertutup

\mathdisp {\left\{ P \in \R^2  \mid   1 \leq \Vert { P } \Vert \leq 2         \right\}} {  }

dan silinder tertutup

\mathdisp {S^1 \times [0,1]} {  }

merupakan
\definitionsverweis {manifold berbatas}{}{}
yang saling
\definitionsverweis {difeomorfik}{}{.}

}
{} {}

\inputaufgabegibtloesung
{}
{

Jelaskan mengapa $\R^2$ tidak dapat diberi struktur
\definitionsverweis {manifold berbatas}{}{}
sedemikian rupa sehingga himpunan bagian $T$ yang diberikan menjadi
\definitionsverweis {batas}{}{}
dari struktur tersebut.
\aufzaehlungvierabc{
\mavergleichskette
{\vergleichskette
{ T
}
{ = }{ {]0,1[}
}
{  }{
}
{    }{
}
{   }{
}
}
{}{}{,}
dengan interval tersebut terletak pada sumbu $x$.
}{
\mavergleichskette
{\vergleichskette
{ T
}
{ = }{ [0,1]
}
{  }{
}
{    }{
}
{   }{
}
}
{}{}{,}
dengan interval tersebut terletak pada sumbu $x$.
}{
\mavergleichskette
{\vergleichskette
{ T
}
{ = }{ \R
}
{  }{
}
{    }{
}
{   }{
}
}
{}{}{,}
dengan $\R$ menyatakan sumbu $x$.
}{
\mavergleichskette
{\vergleichskette
{ T
}
{ = }{ S^1
}
{  }{
}
{    }{
}
{   }{
}
}
{}{}{.}
}

}
{} {}

\inputaufgabe
{}
{

Tunjukkan bahwa kedua
\definitionsverweis {difeomorfisme}{}{}
yang saling invers,
\maabbeledisp {\varphi} { {\mathbb H} } { U \left(   0 , 1  \right)
} { z} { { \frac{  z-  { \mathrm i}  }{ z + { \mathrm i}  } }
} {,}
dan
\maabbeledisp {\psi} { U \left(   0 , 1  \right)  } { {\mathbb H}
} { w } { { \frac{   (w+1)  { \mathrm i}  }{  1-w  } }
} {,}
\zusatzklammer {lihat
Soal 18.10} {} {,}
dapat diperluas menjadi difeomorfisme
\zusatzklammer {antara
\definitionsverweis {manifold berbatas}{}{}} {} {}
antara setengah bidang tertutup dan cakram tertutup berlubang
\mathl{B  \left(  0 , 1 \right) \setminus \{  \left(  1   , \,   0                   \right) \}}{}.

}
{} {}

\inputaufgabe
{}
{

Misalkan $M$ suatu
\definitionsverweis {manifold diferensiabel berbatas}{}{.}
Yang dimaksud dengan \stichwort {lintasan-separuh diferensiabel} {} ialah setiap
\definitionsverweis {pemetaan diferensiabel}{}{}
\maabbdisp {\gamma} { [0, \epsilon [ } { M
} {}
atau
\maabbdisp {\gamma} {] -\epsilon, 0] } { M
} {}
\zusatzklammer {dengan \mathlk{\epsilon >0}{;} masing-masing disebut mengarah ke dalam dan mengarah ke luar} {} {.}
Definisikan kapan dua lintasan-separuh dengan
\mathl{\gamma_1(0) = \gamma_2(0) =P \in M}{} disebut \stichwort {ekuivalen secara tangensial} {,} lalu tunjukkan bahwa definisi ini memberikan suatu relasi ekuivalensi dan bahwa
\definitionsverweis {himpunan hasil bagi}{}{}
yang diperoleh merupakan ruang vektor riil yang disebut \stichwort {ruang singgung} {} di $P$. Karakterisasikan kelas-kelas ekuivalensi yang dapat direpresentasikan baik oleh lintasan-separuh yang mengarah ke dalam maupun yang mengarah ke luar.

}
{} {}

\inputaufgabegibtloesung
{}
{

Misalkan $M$ suatu
\definitionsverweis {manifold berbatas}{}{}
dan

\mavergleichskette
{\vergleichskette
{ P
}
{ \in }{ \partial M
}
{  }{
}
{    }{
}
{   }{
}
}
{}{}{}
suatu titik batas. Tunjukkan bahwa untuk suatu vektor singgung

\mavergleichskette
{\vergleichskette
{ v
}
{ \in }{ T_PM
}
{  }{
}
{    }{
}
{   }{
}
}
{}{}{}
kedua sifat berikut ekuivalen.
\aufzaehlungzwei {Terdapat lintasan yang terdiferensialkan secara kontinu
\maabb {\gamma} { [- \epsilon, 0]} { M
} {}
dengan

\mavergleichskette
{\vergleichskette
{ \gamma(0)
}
{ = }{ P
}
{  }{
}
{    }{
}
{   }{
}
}
{}{}{}
dan

\mavergleichskette
{\vergleichskette
{ \gamma'(0)
}
{ = }{ v
}
{  }{
}
{    }{
}
{   }{
}
}
{}{}{.}
} {Pada setiap peta dengan setengah ruang negatif, vektor $v$ dipetakan oleh pemetaan singgung ke setengah ruang kanan.
}

}
{} {}

\inputaufgabe
{}
{

Misalkan $M$ suatu
\definitionsverweis {manifold diferensiabel berbatas}{}{}
dan
\maabbdisp {\gamma} {{]{-\epsilon}, \epsilon[}} {M
} {}
suatu
\definitionsverweis {lintasan diferensiabel}{}{}
dengan

\mavergleichskettedisp
{\vergleichskette
{ \gamma (0)
}
{ =} { P
}
{ \in} { \partial M
}
{ } {
}
{ } {
}
}
{}{}{.}
Tunjukkan bahwa
\mathl{\gamma'(0) \in T_P \partial M}{.}

}
{} {}

\inputaufgabegibtloesung
{}
{

Misalkan $M$ suatu
\definitionsverweis {manifold berbatas}{}{.}
Tunjukkan bahwa setiap himpunan bagian terbuka

\mavergleichskette
{\vergleichskette
{ N
}
{ \subseteq }{ M
}
{  }{
}
{    }{
}
{   }{
}
}
{}{}{}
juga merupakan manifold dengan batas
\zusatzklammer {yang mungkin kosong} {} {,}
dan bahwa

\mavergleichskettedisp
{\vergleichskette
{ \partial N
}
{ =} { \partial M \cap N
}
{ } {
}
{ } {
}
{ } {
}
}
{}{}{}
berlaku.

}
{} {}

\inputaufgabegibtloesung
{}
{

Misalkan

\mavergleichskette
{\vergleichskette
{ H
}
{ \subseteq }{ \R^n
}
{  }{
}
{    }{
}
{   }{
}
}
{}{}{}
suatu
\definitionsverweis {setengah ruang Euklides}{}{}
dan

\mavergleichskette
{\vergleichskette
{ P
}
{ \in }{ H
}
{  }{
}
{    }{
}
{   }{
}
}
{}{}{.}
Misalkan terdapat himpunan yang terbuka di $\R^n$, yaitu $V$, dengan

\mavergleichskette
{\vergleichskette
{ P
}
{ \in }{ V
}
{ \subseteq }{ H
}
{    }{
}
{   }{
}
}
{}{}{.}
Tunjukkan bahwa
\mathl{P}{} bukan titik batas $H$.

}
{} {}

\inputaufgabe
{}
{

Misalkan
\mathl{U_1 \subseteq H_1 \subset \R^n}{} dan
\mathl{U_2 \subseteq H_2 \subset \R^n}{}
merupakan
\definitionsverweis {himpunan bagian terbuka}{}{}
di dalam
\definitionsverweis {setengah ruang Euklides}{}{}
\mathkor {} {H_1} {dan} {H_2} {,}
serta
\maabbdisp {\varphi} {U_1 } { U_2
} {}
suatu
\definitionsverweis {difeomorfisme}{}{.}
Tunjukkan bahwa untuk setiap titik
\mathl{P\in U_1}{} terdapat lingkungan terbuka
\mathl{P\in V_1 \subset \R^n}{} dan
\mathl{\varphi(P) \in V_2 \subset \R^n}{,}
beserta suatu perluasan
\maabbdisp {\tilde{\varphi}} {V_1 } { V_2
} {}
yang merupakan difeomorfisme.

}
{} {}

\inputaufgabe
{}
{

Misalkan cakram tertutup
\mathl{B  \left(  0 , 1 \right)}{} dilengkapi dengan
\definitionsverweis {orientasi standar}{}{}
dari $\R^2$. Apakah orientasi yang ditentukan oleh
\definitionsverweis {normal luar}{}{}
pada batas
\zusatzklammer {yaitu pada lingkaran satuan} {} {}
searah atau berlawanan arah jarum jam?

}
{} {}

  \zwischenueberschrift{Soal untuk dikumpulkan}

\bild{ \begin{center}
\includegraphics[width=5.5cm]{ \bildeinlesungsvg {Not-star-shaped} {svg} }
\end{center}
\bildtext {} }

\bildlizenz { Not-star-shaped.svg } {} {WillT.Net} {Commons} {PD} {}

\inputaufgabe
{4}
{

Untuk anulus

\mavergleichskettedisp
{\vergleichskette
{M
}
{ =} { \left\{  x \in \R^2  \mid   1 \leq \Vert { x } \Vert \leq 2         \right\}
}
{ } {
}
{ } {
}
{ } {
}
}
{}{}{}
berikan secara eksplisit
\definitionsverweis {peta}{}{}
yang menunjukkan bahwa $M$ merupakan
\definitionsverweis {manifold berbatas}{}{.}

}
{} {}

\inputaufgabe
{6}
{

Tunjukkan bahwa
\definitionsverweis {setengah bidang}{}{} $\R_{\geq 0} \times \R$ dan
\definitionsverweis {kuadran}{}{}
\mathl{\R_{\geq 0} \times \R_{\geq 0}}{} tidak
\definitionsverweis {difeomorfik}{}{.}

}
{(Apa pengertian difeomorfisme yang sesuai dalam konteks ini?)} {}

\inputaufgabe
{6}
{

Misalkan
\mathl{M= B  \left(  0 , 1 \right) \setminus \{ (0,1), (0,-1)\} \subseteq \R^2}{,}
yaitu
\definitionsverweis {cakram tertutup}{}{}
yang dua titik batasnya telah dibuang. Misalkan
\mathl{N={]{-1},1[} \times  [-1,1]}{} merupakan produk sebuah interval terbuka dan sebuah interval tertutup. Tunjukkan bahwa
\mathkor {} {M} {dan} {N} {}
merupakan
\definitionsverweis {manifold berbatas}{}{}
yang saling
\definitionsverweis {difeomorfik}{}{.}

}
{} {}

\inputaufgabe
{4}
{

Misalkan
\mathl{V_1,V_2 \subseteq \R^n}{}
merupakan
\definitionsverweis {himpunan bagian terbuka}{}{}
dan
\maabbdisp {\varphi} {V_1 } { V_2
} {}
suatu
\definitionsverweis {difeomorfisme}{}{}
yang menginduksi
\definitionsverweis {homeomorfisme}{}{}
antara
\mathkor {} {V_1 \cap H} {dan} {V_2 \cap H} {}
dan, dengan demikian, juga antara
\mathkor {} {V_1 \cap \partial H} {dan} {V_2 \cap \partial H} {}
\zusatzklammer {$H$ menyatakan
\definitionsverweis {setengah ruang}{}{}
dan $\partial H$ batasnya} {} {.}
Tunjukkan bahwa pembatasan pada batas juga merupakan
\definitionsverweis {difeomorfisme}{}{.}

}
{} {}

\inputaufgabe
{4}
{

Misalkan bola satuan tertutup
\mathl{B  \left(  0 , 1 \right)}{} dilengkapi dengan
\definitionsverweis {orientasi standar}{}{} dari $\R^3$. Tentukan apakah kedua vektor singgung
\mathkor {} {(2,1,0)} {dan} {(3,-1,0)} {}
di kutub utara
\mathl{(0,0,1)}{} merepresentasikan orientasi pada batas yang diinduksi oleh
\definitionsverweis {normal luar}{}{}
\zusatzklammer {yaitu pada permukaan bola satuan} {} {,}
atau tidak.

}
{} {}

\section*{Solusi yang disediakan oleh sumber}
\addcontentsline{toc}{section}{Solusi yang disediakan oleh sumber}
\subsection*{Solusi Soal 21.7}
a) Batas
\mathl{\partial M}{} suatu manifold berbatas, menurut
Fakta (2),
\definitionsverweis {tertutup}{}{,}
sedangkan interval terbuka sebagai himpunan bagian dari $\R^2$ tidak tertutup.

b) Batas
\mathl{\partial M}{} suatu manifold berbatas, menurut
Fakta,
merupakan manifold tanpa batas, sedangkan interval tertutup memiliki titik-titik batas.

c) Andaikan bahwa $\R^2$ merupakan manifold dengan batas
\mathl{\R = \R \times \{0\}}{.} Untuk
\mathl{P \in \R}{} terdapat suatu lingkungan terbuka
\mathl{P\in U}{} dan suatu homeomorfisme
\maabbdisp {\alpha} { U } {V
} {}
dengan suatu himpunan terbuka $V \subseteq H$, dengan $H$ menyatakan suatu setengah bidang. Dalam hal ini, kita dapat mengganti $V$ dengan suatu lingkungan setengah bola terbuka yang lebih kecil di sekitar $P$. Menurut
Fakta (2),
pada peta seperti itu titik-titik batas dipetakan ke titik-titik batas, yaitu

\mavergleichskettedisp
{\vergleichskette
{\alpha ( U \cap \R  )
}
{ =} { V \cap \partial H
}
{ } {
}
{ } {
}
{ } {
}
}
{}{}{.}
Dengan demikian, kita juga memperoleh suatu homeomorfisme antara komplemen-komplemennya, yakni antara
\mathkor {} {U \setminus U \cap \R} {dan} {V \setminus V \cap \partial H} {.}
Namun, lingkungan setengah bola di ruas kanan
\definitionsverweis {terhubung lintasan}{}{,}
sedangkan himpunan di ruas kiri merupakan gabungan saling lepas dari irisannya dengan separuh positif dan separuh negatif. Kedua irisan tersebut terbuka dan juga tidak kosong karena $U$ memuat suatu lingkungan bola dari $P$. Oleh karena itu, himpunan di ruas kiri tidak terhubung, sehingga tidak mungkin terdapat suatu homeomorfisme.

d) Sama seperti c).
\clearpage
\subsection*{Solusi Soal 21.10}
Misalkan $\gamma$ suatu lintasan-separuh di $M$ seperti yang diberikan, dan misalkan

\mavergleichskette
{\vergleichskette
{ P
}
{ \in }{  U
}
{ \subseteq }{ M
}
{    }{
}
{   }{
}
}
{}{}{}
merupakan suatu daerah peta dengan peta
\maabbdisp {\alpha} { U } {V \subseteq  H_{\leq 0}
} {}
dengan

\mavergleichskettedisp
{\vergleichskette
{  H_{\leq 0}
}
{ =} { \left\{ (x_1,x_2 , \ldots, x_n)  \mid   x_1 \leq 0        \right\}
}
{ } {
}
{ } {
}
{ } {
}
}
{}{}{}
dan

\mavergleichskettedisp
{\vergleichskette
{ Q
}
{ =} { \alpha(P)
}
{ =} { \left(  0   , \,  a_2  , \,   \ldots , \,  a_n                \right)
}
{ } {
}
{ } {
}
}
{}{}{.}
Di bawah pemetaan singgung,

\mavergleichskette
{\vergleichskette
{ v
}
{ = }{ \gamma'(0)
}
{  }{
}
{    }{
}
{   }{
}
}
{}{}{}
dipetakan ke vektor turunan dari

\mavergleichskette
{\vergleichskette
{ \delta
}
{ = }{ \alpha \circ \gamma
}
{  }{
}
{    }{
}
{   }{
}
}
{}{}{}
di titik nol. Karena $\gamma$ bernilai di $M$, $\delta$ bernilai seluruhnya di $H_{\leq 0}$. Oleh karena itu, dalam hasil bagi selisih

\mathdisp {{ \frac{   \left(  \delta_1(t)-\delta_1(0)   , \,   \delta_2(t)-\delta_2(0)  , \,   \ldots , \,   \delta_n(t)-\delta_n(0)                \right)  }{ t } }} {  }

nilai $t$ negatif dan nilai $\delta_1(t)$ nonpositif, sehingga

\mavergleichskettedisp
{\vergleichskette
{  \delta_1'(0)
}
{ \geq} { 0
}
{ } {
}
{ } {
}
{ } {
}
}
{}{}{,}
dan dengan demikian termasuk separuh positif.

Sebaliknya, andaikan
\mathl{T_P(\alpha)(v)}{} termasuk separuh positif. Maka garis
\maabbeledisp {\delta} {\R} {\R^n
} { t } { Q+tT_P(\alpha)(v)
} {}
untuk

\mavergleichskette
{\vergleichskette
{ t
}
{ \leq }{  0
}
{  }{
}
{    }{
}
{   }{
}
}
{}{}{}
terletak seluruhnya di dalam separuh negatif. Lintasan
\mathl{\alpha^{-1} \circ \delta}{} terdefinisi pada suatu interval di sisi negatif, bernilai di $M$, dan merupakan realisasi lintasan-separuh dari $v$.
\subsection*{Solusi Soal 21.12}
Misalkan
\mathl{P \in N}{} dan
\mathl{P \in U}{} suatu daerah peta dari $M$. Misalkan
\maabbdisp {\alpha} { U } {V
} {}
suatu peta dengan himpunan terbuka
\mathl{V \subseteq H}{,}
\mathl{H=\R_{\geq 0} \times \R^{n-1}}{.} Peta ini menginduksi suatu homeomorfisme
\maabbdisp {} {N \cap U} { \alpha (N \cap U)
} {.}
Peta-peta inilah yang kita gunakan untuk $N$. Sifat difeomorfisme dari pergantian peta untuk $M$ langsung diwariskan oleh pergantian peta untuk $N$. Misalkan
\mathl{P \in N \cap \partial M}{.} Di bawah suatu peta, $P$ kemudian dipetakan ke suatu titik batas dari $H$. Karena peta-peta untuk $N$ merupakan pembatasan dari peta-peta semacam itu, sifat ini tetap berlaku. Sebaliknya, jika
\mathl{P \in \partial N}{}, maka tentu saja, di satu pihak,
\mathl{P \in N \subseteq M}{.} Di pihak lain, di sekitar $P$ terdapat suatu peta untuk $N$ yang diperoleh dengan membatasi suatu peta untuk $M$, dan dalam peta tersebut $P$ dipetakan ke suatu titik batas dari $H$.
\subsection*{Solusi Soal 21.13}
Misalkan

\mavergleichskettedisp
{\vergleichskette
{ P
}
{ =} { \begin{pmatrix}  x_1  \\ x_2 \\  \vdots\\ x_n   \end{pmatrix}
}
{ } {
}
{ } {
}
{ } {
}
}
{}{}{.}
Kita dapat mengganti lingkungan yang terbuka di dalam $\R^n$, yaitu $V$, dengan suatu bola terbuka
\zusatzklammer {di dalam $\R^n$} {} {}
\mathl{U \left(   P , \epsilon  \right)}{} yang memenuhi

\mavergleichskette
{\vergleichskette
{ P
}
{ \in }{   U \left(   P , \epsilon  \right)
}
{ \subseteq }{  V
}
{ \subseteq   }{ H
}
{   }{
}
}
{}{}{}.
Jika $P$ merupakan suatu titik batas, maka

\mavergleichskette
{\vergleichskette
{ x_1
}
{ = }{ 0
}
{  }{
}
{    }{
}
{   }{
}
}
{}{}{.}
Namun, dalam hal itu,

\mavergleichskette
{\vergleichskette
{ \begin{pmatrix}  - { \frac{   \epsilon }{  2 } }  \\ x_2 \\  \vdots\\ x_n   \end{pmatrix}
}
{ \in }{ U \left(   P , \epsilon  \right)
}
{  }{
}
{    }{
}
{   }{
}
}
{}{}{,}
padahal titik ini tidak termasuk $H$.

\part{Unit 22}
\chapter[Kuliah 22]{Kuliah 22: Partisi Satuan}
\setcounter{section}{22}

Sekarang kita membahas suatu teknik bantu analitik penting yang disebut \stichwort {partisi satuan} {.} Kita akan menggunakannya dalam pembuktian pernyataan bahwa manifold yang dapat diorientasikan mempunyai bentuk volume positif, dan dalam pembuktian teorema Stokes. Dalam kuliah ini kita akan mengonstruksi partisi satuan; untuk itu mula-mula kita memerlukan beberapa konsep topologis.

  \zwischenueberschrift{Penghabisan kompak}

\bild{ \begin{center}
\includegraphics[width=5.5cm]{ \bildeinlesungpng {Inner_point} {png} }
\end{center}
\bildtext {Interior terbuka adalah gabungan semua titik interior, yaitu titik-titik dari $T$
\zusatzklammer {pada gambar, $E$} {} {,}
yang memiliki suatu lingkungan terbuka yang seluruhnya termuat di dalam $T$.} }

\bildlizenz { Inner point.png } {} {Zasdfgbnm} {Commons} {PD} {}

\inputdefinition
{}
{

Misalkan $X$ suatu
\definitionsverweis {ruang topologis}{}{}
dan

\mavergleichskette
{\vergleichskette
{T
}
{ \subseteq }{ X
}
{  }{
}
{    }{
}
{   }{
}
}
{}{}{}
suatu himpunan bagian. Maka

\mavergleichskettedisp
{\vergleichskette
{ \overline{  T  }
}
{ =} {  \bigcap_{A \text{ tertutup}, \, T \subseteq A} A
}
{ } {
}
{ } {
}
{ } {
}
}
{}{}{}
disebut \definitionswort {penutupan}{}
\zusatzklammer {atau \definitionswort {penutupan topologis}{}} {} {}
dari $T$.

}
Untuk ruang metrik, kita sebelumnya telah memperkenalkan penutupan sebagai himpunan semua
\definitionsverweis {titik pelekatan}{}{}{;}
lihat khususnya
Soal 35.3 (Analisis (Osnabrück 2021--2023)).

\inputdefinition
{}
{

Misalkan $X$ suatu
\definitionsverweis {ruang topologis}{}{}
dan

\mavergleichskette
{\vergleichskette
{ T
}
{ \subseteq }{ X
}
{  }{
}
{    }{
}
{   }{
}
}
{}{}{}
suatu himpunan bagian. Maka

\mavergleichskettedisp
{\vergleichskette
{ T ^{o}
}
{ =} { \bigcup_{U \text{ terbuka}, \, U \subseteq T}   U
}
{ } {
}
{ } {
}
{ } {
}
}
{}{}{}
disebut \definitionswort {interior terbuka}{}
\zusatzklammer {atau \stichwort {interior} {}} {} {}
dari $T$.

}
Kedua konsep ini berhubungan melalui

\mavergleichskettedisp
{\vergleichskette
{ \overline{ T }
}
{ =} { X \setminus (X \setminus T)^{o}
}
{ } {
}
{ } {
}
{ } {
}
}
{}{}{.}
Konsep batas juga terbawa dari situasi metrik ke ruang topologis sembarang.

\inputdefinition
{
}
{

Misalkan $X$ suatu
\definitionsverweis {ruang topologis}{}{}
dan

\mavergleichskette
{\vergleichskette
{ T
}
{ \subseteq }{ X
}
{  }{
}
{    }{
}
{   }{
}
}
{}{}{}
suatu himpunan bagian. Yang dimaksud dengan
\definitionswort {batas}{}
dari $T$ adalah himpunan

\mavergleichskettedisp
{\vergleichskette
{ \operatorname{Batas} \, ( T )
}
{ =} { \overline{  T  }  \setminus T ^{o}
}
{ } {
}
{ } {
}
{ } {
}
}
{}{}{.}

}
Perhatikan bahwa batas topologis ini merupakan konsep yang berbeda dari batas suatu manifold berbatas. Meskipun demikian, terdapat hubungan yang dibahas dalam
Soal 22.1.

\inputdefinition
{}
{

Misalkan $X$ suatu
\definitionsverweis {ruang topologis}{}{}
dan
\maabbdisp {f} { X } {\R
} {}
suatu
\definitionsverweis {fungsi}{}{.}
Maka
\definitionsverweis {penutupan topologis}{}{}

\mathdisp {\overline{  \left\{  x \in X  \mid  f(x) \neq 0        \right\}   }} {  }

disebut \definitionswort {tumpuan}{} dari $f$.

}

\inputdefinition
{}
{

Misalkan $X$ suatu
\definitionsverweis {ruang topologis}{}{.}
Suatu \definitionswort {penghabisan kompak}{}

\mathbed {A_n} {}
 {n \in \N} {}
 {} {} {} {,}
dari $X$ adalah suatu
\definitionsverweis {barisan}{}{}
\definitionsverweis {himpunan-himpunan bagian kompak}{}{}

\mavergleichskette
{\vergleichskette
{ A_n
}
{ \subseteq }{ X
}
{  }{
}
{    }{
}
{   }{
}
}
{}{}{}
dengan

\mathdisp {A_n \subseteq A_{n+1}^{o} \text{ dan } \bigcup_{n=0}^\infty A_n = X} { . }

}

Ruang $\R^d$ dihabiskan secara kompak oleh bola-bola tertutup

\mathbed {B  \left( 0,n \right)} {}
 {n \in \N} {}
 {} {} {} {.}
Lihat
Soal 22.6.

  \zwischenueberschrift{Partisi satuan}



\inputfaktbeweis
{Mannigfaltigkeit/Abzählbare Basis/Kompakte Ausschöpfung/Fakt}
{Lema}
{}
{

\faktsituation {Misalkan $M$ suatu
\definitionsverweis {manifold}{}{}}
\faktvoraussetzung {dengan suatu
\definitionsverweis {basis terhitung bagi topologinya}{}{.}}
\faktfolgerung {Maka $M$ mempunyai suatu
\definitionsverweis {penghabisan kompak}{}{.}}
\faktzusatz {}
\faktzusatz {}

}
{

Untuk setiap titik

\mavergleichskette
{\vergleichskette
{ P
}
{ \in }{ M
}
{  }{
}
{    }{
}
{   }{
}
}
{}{}{}
terdapat suatu lingkungan koordinat terbuka

\mavergleichskette
{\vergleichskette
{ P
}
{ \in }{ U
}
{  }{
}
{    }{
}
{   }{
}
}
{}{}{,}
\maabbdisp {\alpha} { U } { V
} {}
serta lingkungan-lingkungan bola

\mavergleichskettedisp
{\vergleichskette
{ U \left(  \alpha(P), \epsilon  \right)
}
{ \subseteq} { B  \left( \alpha(P), \epsilon \right)
}
{ \subset} { V
}
{ } {
}
{ } {
}
}
{}{}{.}
Karena pemetaan koordinat tersebut merupakan
\definitionsverweis {homeomorfisme}{}{}
dan bola tertutup bersifat kompak, maka

\mavergleichskette
{\vergleichskette
{ B_P
}
{ = }{ \alpha^{-1} \left(  B  \left( \alpha(P), \epsilon \right)  \right)
}
{  }{
}
{    }{
}
{   }{
}
}
{}{}{}
merupakan suatu
\definitionsverweis {himpunan bagian kompak}{}{}
dari $M$ yang memuat
\definitionsverweis {lingkungan terbuka}{}{}

\mavergleichskette
{\vergleichskette
{ U_P
}
{ = }{ \alpha^{-1} \left(  U \left(  \alpha(P), \epsilon  \right)  \right)
}
{  }{
}
{    }{
}
{   }{
}
}
{}{}{}
dari $P$. Himpunan-himpunan

\mathbed {U_P} {}
 {P \in M} {}
 {} {} {} {,}
membentuk suatu
\definitionsverweis {tutupan terbuka}{}{}
dari $M$, sehingga menurut
Soal 2.8 (Teori Ukuran dan Integrasi (Osnabrück 2022--2023))
terdapat suatu subtutupan terhitung. Nyatakan subtutupan ini dengan

\mathbed {U_n} {}
 {n \in \N} {}
 {} {} {} {,}
\zusatzklammer {dengan $U_n$ termuat di dalam himpunan kompak $B_n$} {} {.}
Sekarang kita definisikan secara
\definitionsverweis {rekursif}{}{}
suatu pemetaan
\definitionsverweis {naik secara monoton}{}{}
\maabbeledisp {} {\N} {\N
} { k } {n_k
} {,}
sedemikian sehingga

\mathbeddisp {A_k = \bigcup_{n=0}^{n_k } B_n} {}
 {k \in \N} {}
 {} {} {} {,}
merupakan suatu
\definitionsverweis {penghabisan kompak}{}{}
dari $M$. Sebagai gabungan berhingga dari himpunan-himpunan kompak, setiap $A_k$ bersifat kompak. Kita mulai dengan

\mavergleichskette
{\vergleichskette
{ n_0
}
{ = }{ 0
}
{  }{
}
{    }{}
{   }{}
}
{}{}{.}
Misalkan $n_k$ sudah dikonstruksi. Himpunan

\mathdisp {A_k \cup B_{n_k +1}} {  }

bersifat
\definitionsverweis {kompak}{}{}
dan karena itu ditutupi oleh berhingga banyak himpunan terbuka
\mathl{\bigcup_{n=0}^{n_{k+1} } U_n}{,} dengan memilih

\mavergleichskette
{\vergleichskette
{ n_{k+1}
}
{ \geq }{ n_k +1
}
{  }{
}
{    }{
}
{   }{
}
}
{}{}{.}
Dengan pilihan ini,

\mavergleichskettedisp
{\vergleichskette
{ A_k
}
{ \subseteq} { \bigcup_{n = 0}^{n_{k+1} } U_n
}
{ \subseteq} { \bigcup_{n = 0}^{n_{k+1} } B_n
}
{ =} { A_{k+1}
}
{ } {
}
}
{}{}{,}
dan barisan ini membentuk suatu penghabisan karena himpunan-himpunan

\mathbed {U_n} {}
 {n \in \N} {}
 {} {} {} {,}
membentuk suatu tutupan terbuka dari $M$.

}

\bild{ \begin{center}
\includegraphics[width=5.5cm]{ \bildeinlesungsvg {Partition_of_unity_illustration} {svg} }
\end{center}
\bildtext {Ilustrasi beberapa fungsi taknegatif yang jumlahnya selalu 1.} }

\bildlizenz { Partition of unity illustration.svg } {} {Oleg Alexandrov} {Commons} {PD} {}

\inputdefinition
{}
{

Misalkan $X$ suatu
\definitionsverweis {ruang topologis}{}{.}
Suatu keluarga fungsi
\maabbdisp {h_j} { X } {\R
} {}
dengan

\mavergleichskette
{\vergleichskette
{ j
}
{ \in }{ J
}
{  }{
}
{    }{
}
{   }{
}
}
{}{}{}
disebut suatu \definitionswort {partisi satuan}{,} jika sifat-sifat berikut berlaku.
\aufzaehlungdrei{Berlaku

\mavergleichskette
{\vergleichskette
{ h_j(X)
}
{ \subseteq }{ [0,1]
}
{  }{
}
{    }{
}
{   }{
}
}
{}{}{}
untuk setiap

\mavergleichskette
{\vergleichskette
{ j
}
{ \in }{ J
}
{  }{
}
{    }{
}
{   }{
}
}
{}{}{.}
}{Setiap titik

\mavergleichskette
{\vergleichskette
{ P
}
{ \in }{ X
}
{  }{
}
{    }{
}
{   }{
}
}
{}{}{}
mempunyai suatu
\definitionsverweis {lingkungan terbuka}{}{}

\mavergleichskette
{\vergleichskette
{ P
}
{ \in }{ U
}
{  }{
}
{    }{
}
{   }{
}
}
{}{}{}
sedemikian sehingga
\definitionsverweis {fungsi-fungsi hasil pembatasan}{}{}

\mathl{h_j {{|}}_{ U }}{} semuanya, kecuali berhingga banyak, merupakan fungsi nol.
}{Berlaku

\mavergleichskette
{\vergleichskette
{  \sum_{j \in J} h_j
}
{ = }{ 1
}
{  }{
}
{    }{
}
{   }{
}
}
{}{}{.}
} Jika semua $h_j$ \definitionsverweis {kontinu}{}{,} partisi tersebut disebut suatu \definitionswort {partisi satuan kontinu}{.}

}

Sifat kedua memastikan bahwa jumlah dalam (3) terdefinisi, sebab untuk setiap titik

\mavergleichskette
{\vergleichskette
{ P
}
{ \in }{ X
}
{  }{
}
{    }{
}
{   }{
}
}
{}{}{}
dan bagi semua indeks

\mavergleichskette
{\vergleichskette
{ j
}
{ \in }{ J
}
{  }{
}
{    }{
}
{   }{
}
}
{}{}{}
kecuali berhingga banyak, berlaku kesamaan

\mavergleichskette
{\vergleichskette
{  h_j(P)
}
{ = }{ 0
}
{  }{
}
{    }{
}
{   }{
}
}
{}{}{.}
Pada suatu manifold, partisi seperti itu disebut \stichwort {diferensiabel} {,} jika semua $h_j$ merupakan
\definitionsverweis {fungsi-fungsi diferensiabel}{}{.}

\inputdefinition
{}
{

Misalkan

\mavergleichskette
{\vergleichskette
{ X
}
{ = }{ \bigcup_{i \in I} W_i
}
{  }{
}
{    }{
}
{   }{
}
}
{}{}{}
suatu
\definitionsverweis {tutupan terbuka}{}{}
dari suatu
\definitionsverweis {ruang topologis}{}{}
$X$. Suatu
\definitionsverweis {partisi satuan}{}{}
\maabbdisp {h_j} { X } {\R
} {}
dengan

\mavergleichskette
{\vergleichskette
{ j
}
{ \in }{ J
}
{  }{
}
{    }{
}
{   }{
}
}
{}{}{}
disebut suatu \definitionswort {partisi satuan subordinat terhadap tutupan}{,} jika untuk setiap

\mavergleichskette
{\vergleichskette
{ j
}
{ \in }{ J
}
{  }{
}
{    }{
}
{   }{
}
}
{}{}{}
terdapat suatu himpunan terbuka
\mathl{W_{i(j)}}{} dari tutupan tersebut sedemikian sehingga
\definitionsverweis {tumpuan}{}{}
dari $h_j$ termuat di dalam
\mathl{W_{i(j)}}{}{.}

}



\inputfaktbeweis
{Mannigfaltigkeit/Abzählbare Basis/Lokal endlicher Atlas/Fakt}
{Lema}
{}
{

\faktsituation {Misalkan $M$ suatu
\definitionsverweis {manifold}{}{}
dengan suatu
\definitionsverweis {basis terhitung}{}{}
bagi topologinya. Misalkan

\mavergleichskette
{\vergleichskette
{ M
}
{ = }{ \bigcup_{i \in I} W _i
}
{  }{
}
{    }{
}
{   }{
}
}
{}{}{}
suatu
\definitionsverweis {tutupan terbuka}{}{}
dari $M$.}
\faktfolgerung {Maka terdapat suatu
\definitionsverweis {atlas kompatibel}{}{}
yang terhitung

\mathbed {\left(  U_j, \alpha_j, V_j  \right)} {}
 {j \in J} {}
 {} {} {} {,}
dengan lingkungan-lingkungan bola

\mavergleichskettedisp
{\vergleichskette
{ U \left(   0 , \delta_j   \right)
}
{ \subset} { B  \left(  0 , \epsilon_j  \right)
}
{ \subset} { V_j
}
{ } {
}
{ } {
}
}
{}{}{}
\zusatzklammer {di sini,

\mavergleichskettek
{\vergleichskettek
{ 0
}
{ \in }{ V_j
}
{ \subseteq }{ \R^n
}
{    }{
}
{   }{
}
}
{}{}{}
dan

\mavergleichskettek
{\vergleichskettek
{ \delta_j
}
{ < }{ \epsilon_j
}
{  }{
}
{    }{
}
{   }{
}
}
{}{}{}} {} {}
sedemikian sehingga untuk setiap

\mavergleichskette
{\vergleichskette
{ j
}
{ \in }{ J
}
{  }{
}
{    }{
}
{   }{
}
}
{}{}{}
terdapat suatu
\mathl{W_{i(j)}}{} dengan

\mavergleichskette
{\vergleichskette
{ U_j
}
{ \subseteq }{ W_{i(j)}
}
{  }{
}
{    }{
}
{   }{
}
}
{}{}{}; selain itu, $M$ ditutupi oleh

\mathbed {\alpha_j^{-1} \left(  U \left(   0 , \delta_j   \right)  \right)} {}
 {j \in J} {}
 {} {} {} {,}
dan setiap titik

\mavergleichskette
{\vergleichskette
{ P
}
{ \in }{ M
}
{  }{
}
{    }{
}
{   }{
}
}
{}{}{}
terletak hanya di dalam berhingga banyak himpunan $U_j$.}
\faktzusatz {}
\faktzusatz {}

}
{

\teilbeweis {}{}{}
{Misalkan tutupan terbuka

\mathbed {W_i} {}
 {i \in I} {}
 {} {} {} {,}
diberikan. Selanjutnya, misalkan

\mathbed {A_n} {}
 {n \in \N} {}
 {} {} {} {,}
suatu
\definitionsverweis {penghabisan kompak}{}{}
dari $M$, yang keberadaannya dijamin oleh
Lema 22.6.
Himpunan-himpunan terbuka
\mathl{A_{n+2}^{o}  \setminus A_{n-1}}{} juga membentuk suatu tutupan terbuka, sebab untuk setiap titik

\mavergleichskette
{\vergleichskette
{ P
}
{ \in }{ M
}
{  }{
}
{    }{
}
{   }{
}
}
{}{}{}
terdapat suatu nilai minimal

\mavergleichskette
{\vergleichskette
{ n
}
{ \in }{ \N
}
{  }{
}
{    }{
}
{   }{
}
}
{}{}{}
dengan

\mavergleichskette
{\vergleichskette
{ P
}
{ \in }{ A_n
}
{  }{
}
{    }{
}
{   }{
}
}
{}{}{}
\zusatzklammer {tetapkan

\mavergleichskettek
{\vergleichskettek
{ A_{-1}
}
{ = }{ \emptyset
}
{  }{
}
{    }{
}
{   }{
}
}
{}{}{}} {} {.}
Untuk $n$ ini berlaku
\mathkor {} {P \not\in A_{n-1}} {dan} {P \in A_n \subseteq A_{n+1} ^{o}} {.}
Dengan meninjau irisan-irisan
\mathl{W_i \cap \left(  A_{n+2}^{o}  \setminus A_{n-1}  \right)}{}{,} kita dapat menganggap bahwa setiap himpunan dari tutupan termuat di dalam suatu
\mathl{A_{n+2}^{o} \setminus A_{n-1}}{}{.}}
{}
\teilbeweis {}{}{}
{Untuk setiap titik

\mavergleichskette
{\vergleichskette
{ P
}
{ \in }{ M
}
{  }{
}
{    }{
}
{   }{
}
}
{}{}{}
terdapat suatu lingkungan koordinat terbuka
\zusatzklammer {yang kompatibel} {} {}

\mavergleichskette
{\vergleichskette
{ P
}
{ \in }{ U_P
}
{  }{
}
{    }{
}
{   }{
}
}
{}{}{,}
yang termuat di dalam salah satu $W_i$ dan yang mempunyai lingkungan-lingkungan bola

\mavergleichskettedisp
{\vergleichskette
{ U \left(   0 , \delta_P   \right)
}
{ \subset} { B  \left(  0 , \epsilon_P  \right)
}
{ \subset} { V_P
}
{ } {
}
{ } {
}
}
{}{}{}
dengan

\mavergleichskette
{\vergleichskette
{ P
}
{ \in }{ \alpha_P^{-1} \left(  U \left(   0 , \delta_P   \right)  \right)
}
{  }{
}
{    }{
}
{   }{
}
}
{}{}{}
dan

\mavergleichskette
{\vergleichskette
{ \delta_P
}
{ < }{ \epsilon_P
}
{  }{
}
{    }{
}
{   }{
}
}
{}{}{.}
Himpunan-himpunan

\mathbed {\alpha_P^{-1} \left(  U \left(   0 , \delta_P   \right)  \right)} {}
 {P \in M} {}
 {} {} {} {,}
kemudian juga membentuk suatu tutupan terbuka dari $M$. Menurut
Soal 2.8 (Teori Ukuran dan Integrasi (Osnabrück 2022--2023))
kita dapat beralih ke suatu subtutupan terhitung. Jadi, kita dapat menganggap bahwa suatu sistem domain peta

\mathbed {U_j} {}
 {j \in \N} {}
 {} {} {} {,}
bersama dengan lingkungan-lingkungan bola

\mavergleichskettedisp
{\vergleichskette
{ U \left(   0 , \delta_j   \right)
}
{ \subset} { B  \left(  0 , \epsilon_j  \right)
}
{ \subset} { V_j
}
{ } {
}
{ } {
}
}
{}{}{}
diberikan sedemikian sehingga

\mathbed {\alpha_j^{-1} \left(  U \left(   0 , \delta_j   \right)  \right)} {}
 {j \in \N} {}
 {} {} {} {,}
juga merupakan suatu tutupan terbuka dari $M$, setiap $U_j$ termuat di dalam suatu
\mathl{W_{i(j)}}{}{,} dan hubungan dengan penghabisan kompak yang diuraikan di atas berlaku.}
{}
\teilbeweis {}{Kita akan mendefinisikan suatu himpunan bagian

\mavergleichskette
{\vergleichskette
{ J
}
{ \subseteq }{ \N
}
{  }{
}
{    }{
}
{   }{
}
}
{}{}{}
sedemikian sehingga keluarga

\mathbed {U_j} {}
 {j \in J} {}
 {} {} {} {,}
juga memenuhi sifat keberhinggaan tersebut.\leerzeichen{}}{}
{Mula-mula, kekompakan $A_0$ memberikan suatu himpunan indeks berhingga
\mathl{J_{-1}}{} sedemikian sehingga himpunan-himpunan

\mathbed {\alpha_j^{-1} \left(  U \left(   0 , \delta_j   \right)  \right)} {}
 {j \in J_{-1}} {}
 {} {} {} {,}
menutupi $A_0$. Untuk

\mavergleichskette
{\vergleichskette
{ n
}
{ \in }{ \N
}
{  }{
}
{    }{
}
{   }{
}
}
{}{}{}
tinjau himpunan kompak
\mathl{A_{n+1} \setminus A_n ^{o}}{.} Himpunan ini ditutupi oleh berhingga banyak himpunan

\mathbed {\alpha_j^{-1} \left(  U \left(   0 , \delta_j   \right)  \right)} {}
 {j \in \N} {}
 {} {} {} {,}
dan untuk itu hanya diperlukan indeks-indeks $j$ sedemikian sehingga $U_j$ termuat di dalam
\mathl{A_{n+2}^{o}  \setminus A_{n-1}}{}{.} Nyatakan himpunan indeks berhingga yang bersangkutan dengan $J_n$, dan tetapkan

\mavergleichskette
{\vergleichskette
{ J
}
{ = }{ J_{-1} \cup \bigcup_{n \in \N} J_n
}
{  }{
}
{    }{
}
{   }{
}
}
{}{}{.}
Maka setiap $A_n$ hanya memotong berhingga banyak himpunan

\mathbed {U_j} {}
 {j \in J} {}
 {} {} {} {.}}
{}

}




\inputfaktbeweis
{Mannigfaltigkeit/Abzählbare Basis/Überdeckung/Partition/Fakt}
{Teorema}
{}
{

\faktsituation {Misalkan $M$ suatu
\definitionsverweis {manifold diferensiabel}{}{}
dengan suatu
\definitionsverweis {basis terhitung}{}{}
bagi topologinya.}
\faktfolgerung {Maka untuk setiap tutupan terbuka terdapat suatu
\definitionsverweis {partisi satuan}{}{}
yang \definitionsverweis {terdiferensialkan secara kontinu}{}{}
dan subordinat terhadap tutupan tersebut.}
\faktzusatz {}
\faktzusatz {}

}
{

Menurut
Lema 22.9
kita dapat menganggap bahwa suatu tutupan terbuka terdiri atas domain-domain peta

\mathbed {U_j} {}
 {j \in J} {}
 {} {} {} {,}
\zusatzklammer {$J$ terhitung} {} {}
dengan
\maabbdisp {\alpha_j} {U_j } { V_j
} {}
dan dengan lingkungan-lingkungan bola

\mavergleichskettedisp
{\vergleichskette
{ U \left(   0 , \delta_j   \right)
}
{ \subset} { B  \left(  0 , \epsilon_j  \right)
}
{ \subset} { V_j
}
{ } {
}
{ } {
}
}
{}{}{}
\zusatzklammer {dengan

\mavergleichskettek
{\vergleichskettek
{ \delta_j
}
{ < }{ \epsilon_j
}
{  }{
}
{    }{
}
{   }{
}
}
{}{}{}} {} {}
sedemikian sehingga himpunan-himpunan
\mathl{\alpha_j^{-1} \left(  U \left(   0 , \delta_j   \right)  \right)}{} juga membentuk suatu tutupan dari $M$ dan setiap titik

\mavergleichskette
{\vergleichskette
{ P
}
{ \in }{ M
}
{  }{
}
{    }{
}
{   }{
}
}
{}{}{}
termuat hanya dalam berhingga banyak himpunan $U_j$ dan, khususnya, hanya dalam berhingga banyak himpunan
\mathl{\alpha_j^{-1} \left(  U \left(   0 , \delta_j   \right)  \right)}{}{.}
Pada $V_j$, tinjau fungsi $g_j$ yang didefinisikan oleh

\mavergleichskettedisp
{\vergleichskette
{ g_j (v)
}
{ =} { \begin{cases}  e^{ - { \frac{   1  }{  \left(  \delta_j^2- \Vert { v } \Vert^2  \right)^2 } } }  \text{ untuk } \Vert { v } \Vert < \delta_j \, , \\ 0 \text{ selainnya} \, , \end{cases}
}
{ } {
}
{ } {
}
{ } {
}
}
{}{}{}
Fungsi ini bernilai positif tepat pada
\mathl{U \left(   0 , \delta_j   \right)}{} dan
\definitionsverweis {tumpuan}{}{}
fungsi ini adalah
\mathl{B  \left(  0 , \delta_j  \right)}{.}
Peninjauan pada kedua himpunan bagian terbuka
\zusatzklammer {yang menutupi $V_j$} {} {}
\mathkor {} {U \left(   0 , \epsilon_j   \right)} {dan} {V_j \setminus B  \left(  0 , \delta_j  \right)} {}
menunjukkan bahwa $g_j$ terdiferensialkan tak hingga kali. Kita mendefinisikan suatu fungsi
\maabbdisp {\tilde{g}_j} { M } {\R
} {}
melalui

\mavergleichskettedisp
{\vergleichskette
{ \tilde{g}_j(x)
}
{ =} { \begin{cases}  g_j( \alpha_j(x)) \text{ untuk }  x \in \alpha_j^{-1}( U \left(   0 , \delta_j   \right) ) \, , \\  0 \text{ selainnya} \, . \end{cases}
}
{ } {
}
{ } {
}
{ } {
}
}
{}{}{}
Fungsi ini
\definitionsverweis {terdiferensialkan secara kontinu}{}{}
pada $M$, sebab \anfuehrung{pita}{}
\mathl{B  \left(  0 , \epsilon_j  \right) \setminus U \left(   0 , \delta_j   \right)}{} memungkinkan suatu transisi mulus. Kita tetapkan

\mavergleichskettedisp
{\vergleichskette
{ \tilde{g}(x)
}
{ \defeq} { \sum_{j \in J} \tilde{g}_j(x)
}
{ } {
}
{ } {
}
{ } {
}
}
{}{}{,}
yang merupakan jumlah berhingga untuk setiap titik, sebab
\definitionsverweis {tumpuan}{}{}
dari $\tilde{g}_j$ termuat di dalam

\mavergleichskettedisp
{\vergleichskette
{ \alpha_j^{-1} \left(  B  \left(  0 , \delta_j  \right)  \right)
}
{ \subset} { U_j
}
{ } {
}
{ } {
}
{ } {
}
}
{}{}{.}
Fungsi ini
\definitionsverweis {terdiferensialkan secara kontinu}{}{}
pada $M$ dan positif di setiap titik, sebab
\mathl{\tilde{g}_j(x)}{} masing-masing bernilai positif pada himpunan yang bersesuaian dari tutupan
\mathl{\alpha_j^{-1} \left(  U \left(   0 , \delta_j   \right)  \right)}{}{.}
Dengan demikian, fungsi-fungsi

\mavergleichskettedisp
{\vergleichskette
{ h_j
}
{ =} { { \frac{   \tilde{g}_j  }{ \tilde{g}  } }
}
{ } {
}
{ } {
}
{ } {
}
}
{}{}{}
membentuk partisi satuan yang dicari.

}

  \zwischenueberschrift{Orientasi pada manifold dan bentuk volume}

Dengan bantuan partisi satuan, sekarang kita dapat membuktikan kebalikan dari
Lema 15.5.



\inputfaktbeweis
{Nullstellenfreie Volumenform/Orientierte (Karten) Mannigfaltigkeit/Äquivalenz/Fakt}
{Teorema}
{}
{

\faktsituation {Misalkan $M$ suatu
\definitionsverweis {manifold diferensiabel
}{}{}
dengan suatu
\definitionsverweis {basis terhitung bagi topologinya}{}{.}}
\faktfolgerung {Maka terdapat suatu
\definitionsverweis {bentuk volume}{}{}
\definitionsverweis {kontinu}{}{}
tanpa titik nol pada $M$ jika dan hanya jika $M$
\definitionsverweis {dapat diorientasikan}{}{.}}
\faktzusatz {Bentuk volume ini juga dapat dipilih agar terdiferensialkan secara kontinu dan positif.}
\faktzusatz {}

}
{

\teilbeweis {}{}{}
{Salah satu arah sudah dibuktikan
dalam Lema 15.5.}
{}
\teilbeweis {}{}{}
{Sebaliknya, misalkan $M$
\definitionsverweis {dapat diorientasikan}{}{}
dan diberikan suatu
\definitionsverweis {atlas berorientasi}{}{}
yang
\definitionsverweis {terhitung}{}{}

\mathbed {\left(  U_i,\alpha_i,V_i  \right)} {}
 {i \in I} {}
 {} {} {} {,}
dari $M$. Di sini,

\mavergleichskette
{\vergleichskette
{ V_i
}
{ \subseteq }{ \R^n
}
{  }{
}
{    }{
}
{   }{
}
}
{}{}{}
merupakan himpunan
\definitionsverweis {terbuka}{}{,}
dan
\definitionsverweis {koordinat-koordinat}{}{}

\mathl{x_1 , \ldots , x_n}{} mendefinisikan suatu
\definitionsverweis {bentuk volume}{}{}
kontinu tanpa titik nol
\zusatzklammer {bahkan terdiferensialkan tak hingga kali} {} {}

\mathl{dx_1 \wedge \ldots \wedge dx_n}{} pada $V_i$. Kita tetapkan

\mavergleichskettedisp
{\vergleichskette
{ \omega_i
}
{ =} { \alpha_i^* \left( dx_1 \wedge \ldots \wedge dx_n \right)
}
{ } {
}
{ } {
}
{ } {
}
}
{}{}{}
dan dengan demikian memperoleh suatu bentuk volume tanpa titik nol pada $U_i$, yang kita perluas dengan $0$ di luar $U_i$\zusatzfussnote {Perluasan ini tentu saja tidak kontinu, tetapi hal itu tidak menjadi masalah untuk pembahasan berikut} {.} {.}

Sekarang, misalkan

\mathbed {h_j} {}
 {j \in J} {}
 {} {} {} {,}
suatu keluarga fungsi yang subordinat terhadap tutupan

\mathbed {U_i} {}
 {i \in I} {}
 {} {} {} {,}
dan terdiferensialkan secara kontinu, serta membentuk suatu
\definitionsverweis {partisi satuan}{}{,}
yang keberadaannya dijamin oleh
Teorema 22.10.
Jadi, khususnya, untuk setiap $j$ terdapat suatu
\mathl{i(j)}{} sedemikian sehingga
\definitionsverweis {tumpuan}{}{}
dari $h_j$ termuat di dalam
\mathl{U_{i(j)}}{}{.}
Oleh karena itu,
\mathl{h_j \omega_{i(j)}}{} merupakan suatu
\definitionsverweis {bentuk diferensial}{}{}
kontinu berderajat $n$ pada $M$. Kita tetapkan

\mavergleichskettedisp
{\vergleichskette
{ \omega
}
{ =} { \sum_{j \in J} h_j \omega_{i(j)}
}
{ } {
}
{ } {
}
{ } {
}
}
{}{}{.}
Di setiap titik

\mavergleichskette
{\vergleichskette
{ P
}
{ \in }{ M
}
{  }{
}
{    }{
}
{   }{
}
}
{}{}{}
jumlah ini berhingga dan karena itu mendefinisikan suatu bentuk diferensial kontinu berderajat $n$ pada $M$. Untuk suatu titik

\mavergleichskette
{\vergleichskette
{ P
}
{ \in }{ M
}
{  }{
}
{    }{
}
{   }{
}
}
{}{}{}
dan suatu
\definitionsverweis {basis}{}{}

\mathl{v_1 , \ldots , v_n}{} dari
\mathl{T_PM}{} yang merepresentasikan orientasi, berlaku

\mavergleichskettedisp
{\vergleichskette
{ \omega(P; v_1 , \ldots , v_n)
}
{ =} { \sum_{j \in J} h_j(P) \omega_{i(j)} \left(  P; v_1 , \ldots , v_n  \right)
}
{ } {
}
{ } {
}
{ } {
}
}
{}{}{.}
Terdapat suatu $j$ dengan

\mavergleichskette
{\vergleichskette
{ h_j(P)
}
{ > }{ 0
}
{  }{
}
{    }{
}
{   }{
}
}
{}{}{,}
dan untuk $j$ ini juga berlaku

\mavergleichskette
{\vergleichskette
{ \omega_{i(j)}(P; v_1 , \ldots , v_n)
}
{ > }{ 0
}
{  }{
}
{    }{
}
{   }{
}
}
{}{}{,}
sebab

\mavergleichskette
{\vergleichskette
{ P
}
{ \in }{ U_{i(j)}
}
{  }{
}
{    }{
}
{   }{
}
}
{}{}{}
berlaku; dengan demikian, bentuk ini positif di setiap titik.}
{}

}

\chapter{Lembar Kerja 22}
\setcounter{section}{22}

  \zwischenueberschrift{Soal latihan}

\inputaufgabe
{}
{

Misalkan $M$ suatu
\definitionsverweis {manifold berbatas}{}{}
dan
\mathl{\partial M}{} merupakan batasnya. Tunjukkan bahwa
\definitionsverweis {batas topologis}{}{}
dari
\mathl{M \setminus \partial M}{} sama dengan $\partial M$.

}
{} {}

\inputaufgabe
{}
{

Tentukan
\definitionsverweis {tumpuan}{}{} fungsi-fungsi berikut dari $\R$ ke $\R$.
\aufzaehlungsieben{Suatu fungsi polinomial.
}{Fungsi sinus.
}{Fungsi eksponensial.
}{Fungsi indikator $e_{ \Z }$.
}{Fungsi indikator $e_{ \Q }$.
}{Fungsi indikator $e_{  [a,b]  }$.
}{Fungsi indikator $e_{ ]a,b[ }$.
}

}
{} {}

\inputaufgabe
{}
{

Misalkan
\maabbdisp {f} { X } {\R
} {}
suatu fungsi pada
\definitionsverweis {ruang topologis}{}{}
$X$. Andaikan terdapat suatu
\definitionsverweis {himpunan bagian terbuka}{}{}

\mavergleichskette
{\vergleichskette
{U
}
{ \subseteq }{ X
}
{  }{
}
{    }{
}
{   }{
}
}
{}{}{,}
yang memuat
\definitionsverweis {tumpuan}{}{}
dari $f$. Tunjukkan bahwa $f$
\definitionsverweis {kontinu}{}{}
jika dan hanya jika pembatasan
\mathl{f \vert_U}{} kontinu.

}
{} {}

\inputaufgabe
{}
{

Misalkan $X$ suatu
\definitionsverweis {ruang topologis}{}{}
dan

\mavergleichskette
{\vergleichskette
{ T
}
{ \subseteq }{ X
}
{  }{
}
{    }{
}
{   }{
}
}
{}{}{}
suatu himpunan bagian. Tunjukkan bahwa
\definitionsverweis {penutupan}{}{}
dari $T$ sama dengan
\definitionsverweis {tumpuan}{}{}
dari
\definitionsverweis {fungsi indikator}{}{}

\mathl{e_{  T  }}{}.

}
{} {}

\inputaufgabe
{}
{

Berikan suatu
\definitionsverweis {penghabisan kompak}{}{} untuk
\definitionsverweis {bilangan riil}{}{} $\R$.

}
{} {}

\inputaufgabegibtloesung
{}
{

Berikan suatu
\definitionsverweis {penghabisan kompak}{}{}
untuk $\R^d$.

}
{} {}

\inputaufgabe
{}
{

Misalkan $X$ suatu
\definitionsverweis {ruang topologis}{}{.}
Untuk tutupan atas $X$ yang hanya terdiri atas $X$, tentukan suatu
\definitionsverweis {partisi satuan subordinat terhadap tutupan tersebut}{}{.}

}
{} {}

\inputaufgabe
{}
{

Misalkan $X$ suatu
\definitionsverweis {ruang topologis}{}{}
dan

\mavergleichskette
{\vergleichskette
{ X
}
{ = }{ \bigcup_{i \in I} U_i
}
{  }{
}
{    }{
}
{   }{
}
}
{}{}{}
suatu
\definitionsverweis {tutupan terbuka}{}{.}
Kita tinjau keluarga
\definitionsverweis {fungsi indikator}{}{}

\mathbeddisp {e_{  P  }} {}
 {P \in X} {}
 {} {} {} {.}
Untuk tutupan yang dimaksud
\zusatzklammer {tutupan ini} {}{}{,}
manakah sifat-sifat suatu
\definitionsverweis {partisi satuan subordinat}{}{}
terhadap tutupan tersebut yang dipenuhi oleh keluarga ini?

}
{} {}

\inputaufgabe
{}
{

Kita tinjau
\definitionsverweis {penghabisan kompak}{}{}

\mathbed {A_n= [-n,n]} {}
 {n \in \N} {}
 {} {} {} {,} dari
\definitionsverweis {bilangan riil}{}{}
beserta tutupan terbuka

\mathbed {W_n= A_{n+1}^{o}  \setminus A_{n-1}} {}
 {n \in \N} {}
 {} {} {} {,}
\zusatzklammer {dengan
\mathl{A_{-1}= \emptyset}{}} {} {.}
Carilah suatu tutupan $\R$ dengan
\definitionsverweis {interval terbuka}{}{,}
yang memenuhi sifat-sifat
dari Lemma 22.9
\zusatzklammer {beserta pembuktiannya} {} {}.

}
{} {}

\inputaufgabe
{}
{

Tunjukkan bahwa tidak terdapat barisan
\definitionsverweis {fungsi kontinu}{}{}
\maabbdisp {f_n} { [0,1] } { [0,1]
} {}
untuk

\mavergleichskette
{\vergleichskette
{ n
}
{ \in }{ \N
}
{  }{
}
{    }{
}
{   }{
}
}
{}{}{,}
yang membentuk suatu
\definitionsverweis {partisi satuan}{}{}
dan memiliki sifat bahwa $0$ termasuk dalam
\definitionsverweis {tumpuan}{}{}
setiap fungsi $f_n$.

}
{} {}

\inputaufgabe
{}
{

Misalkan

\mavergleichskette
{\vergleichskette
{ V
}
{ \subseteq }{ H
}
{  }{
}
{    }{
}
{   }{
}
}
{}{}{}
suatu
\definitionsverweis {himpunan terbuka}{}{}
di dalam
\definitionsverweis {setengah ruang}{}{}

\mavergleichskette
{\vergleichskette
{ H
}
{ \subset }{ \R^n
}
{  }{
}
{    }{
}
{   }{
}
}
{}{}{}
dan misalkan
\maabbdisp {f} { V } {\R
} {}
suatu
\definitionsverweis {fungsi yang terdiferensialkan secara kontinu}{}{.}
Tunjukkan bahwa terdapat suatu himpunan terbuka

\mavergleichskette
{\vergleichskette
{ \tilde{V}
}
{ \subseteq }{ \R^n
}
{  }{
}
{    }{
}
{   }{
}
}
{}{}{}
dan suatu
\definitionsverweis {fungsi yang terdiferensialkan secara kontinu}{}{}
\maabbdisp {\tilde{f}} {\tilde{V}} { \R
} {}
sedemikian sehingga

\mavergleichskette
{\vergleichskette
{ \tilde{V} \cap H
}
{ = }{ V
}
{  }{
}
{    }{
}
{   }{
}
}
{}{}{}
dan

\mavergleichskette
{\vergleichskette
{ \tilde{f} \vert_V
}
{ = }{ f
}
{  }{
}
{    }{
}
{   }{
}
}
{}{}{}
berlaku.

}
{} {}

\inputaufgabe
{}
{

Misalkan
\maabbdisp {f} {\R} {\R
} {}
suatu
\definitionsverweis {fungsi yang terdiferensialkan secara kontinu}{}{}
dengan
\definitionsverweis {tumpuan kompak}{}{.}
Tunjukkan bahwa
\definitionsverweis {turunan}{}{}
$f'$ juga memiliki tumpuan kompak dan bahwa

\mavergleichskette
{\vergleichskette
{ \int_\R f' d \lambda^1
}
{ = }{ 0
}
{  }{
}
{    }{
}
{   }{
}
}
{}{}{}
berlaku.

}
{} {}

\inputaufgabe
{}
{

Misalkan
\maabbdisp {f} {\R_{\geq 0}} {\R
} {}
suatu
\definitionsverweis {fungsi yang terdiferensialkan secara kontinu}{}{}
dengan
\definitionsverweis {tumpuan kompak}{}{.}
Tunjukkan bahwa
\definitionsverweis {turunan}{}{}
$f'$ juga memiliki tumpuan kompak dan bahwa

\mavergleichskette
{\vergleichskette
{ \int_{\R_{\geq 0} } f' d \lambda^1
}
{ = }{ -f(0)
}
{  }{
}
{    }{
}
{   }{
}
}
{}{}{}
berlaku. Tunjukkan bahwa pernyataan ini tidak harus berlaku jika $f$ tidak memiliki tumpuan kompak.

}
{} {}

\inputaufgabe
{}
{

Misalkan $X$ suatu
\definitionsverweis {manifold diferensiabel}{}{}
dengan suatu
\definitionsverweis {basis terhitung bagi topologinya}{}{.}
Tunjukkan bahwa terdapat suatu
\definitionsverweis {struktur Riemann}{}{}
pada $X$.

}
{} {}

\inputaufgabe
{}
{

Misalkan $X$ suatu
\definitionsverweis {manifold diferensiabel}{}{}
dengan suatu
\definitionsverweis {basis terhitung bagi topologinya}{}{}
dan misalkan
\maabb {p} { E } { X
} {}
suatu
\definitionsverweis {bundel vektor diferensiabel}{}{}
di atas $X$. Tunjukkan bahwa $E$ dapat dilengkapi dengan struktur
\definitionsverweis {bundel vektor Riemann}{}{.}

}
{} {}

  \zwischenueberschrift{Soal untuk dikumpulkan}

\inputaufgabe
{3}
{

Misalkan

\mathbed {A_n} {}
 {n \in \N} {}
 {} {} {} {,}
suatu
\definitionsverweis {penghabisan kompak}{}{} dari
\definitionsverweis {ruang topologis}{}{} $X$. Dengan konvensi
\mathl{A_{-1}= \emptyset}{,}
tunjukkan bahwa relasi

\mathdisp {A_{n+1} \setminus A_n ^{o}  \subseteq A_{n+2}^{o} \setminus A_{n-1}} {  }
 berlaku.

}
{} {}

\inputaufgabe
{4}
{

Untuk
\definitionsverweis {tutupan terbuka}{}{}

\mavergleichskettedisp
{\vergleichskette
{ \R
}
{ =} { \bigcup_{n \in \Z} ]n,n+3[
}
{ } {
}
{ } {
}
{ } {
}
}
{}{}{}
berikan suatu
\definitionsverweis {partisi satuan}{}{}
kontinu yang subordinat terhadap tutupan tersebut.

}
{} {}

\inputaufgabe
{6}
{

Kita tinjau
\definitionsverweis {penghabisan kompak}{}{}

\mathbeddisp {A_n = B  \left(  0 ,n \right)} {}
 {n \in \N} {}
 {} {} {} {,}
dari $\R^2$ beserta
\definitionsverweis {tutupan terbuka}{}{}

\mathbeddisp {W_n= A_{n+1}^{o} \setminus A_{n-1}} {}
 {n \in \N} {}
 {} {} {} {,}
\zusatzklammer {dengan
\mathl{A_{-1}= \emptyset}{}} {} {.}
Carilah suatu tutupan $\R^2$ dengan
\definitionsverweis {cakram terbuka}{}{,}
yang memenuhi sifat-sifat
dari Lemma 22.9
\zusatzklammer {beserta pembuktiannya} {} {}.

}
{} {}

\inputaufgabe
{3}
{

Berikan suatu contoh
\definitionsverweis {ruang topologis}{}{,}
yang tidak memiliki
\definitionsverweis {penghabisan kompak}{}{.}

}
{} {}

\section*{Solusi yang disediakan oleh sumber}
\addcontentsline{toc}{section}{Solusi yang disediakan oleh sumber}
\subsection*{Solusi Soal 22.6}
Kita meninjau
\definitionsverweis {bola-bola tertutup}{}{}

\mathbed {B  \left(  0 , n  \right)} {}
 {n \in \N_+} {}
 {} {} {} {,}
yang
\definitionsverweis {kompak}{}{}
dan menutupi $\R^d$. Interior terbuka dari
\mathl{B  \left(  0 , n+1  \right)}{}
adalah
\definitionsverweis {bola terbuka}{}{}

\mathl{U \left(   0,n+1  \right)}{} dan karena

\mavergleichskettedisp
{\vergleichskette
{ B  \left(  0,n \right)
}
{ \subseteq} { U \left(   0,n+1  \right)
}
{ } {
}
{ } {
}
{ } {
}
}
{}{}{}
suatu
\definitionsverweis {penghabisan kompak}{}{}
pun diperoleh.

% O011-COMPLETE-EXTENSION-BEGIN
\part{Kelanjutan Edisi Lengkap}
\chapter*{Catatan provenans bagian tambahan}
\addcontentsline{toc}{chapter}{Catatan provenans bagian tambahan}
Unit 23--29 melanjutkan sumber Holger Brenner dari Wikiversity berbahasa Jerman. Hanya solusi inti yang benar-benar disediakan oleh sumber yang dimuat pada bagian unit. Jembatan Grup Lie dan de Rham serta enam perbaikan kekosongan ujian merupakan materi asli edisi ini dan tidak dinisbahkan kepada Brenner atau Wikiversity. Formulir solusi resmi mempertahankan kekosongan solusi sumber; perbaikan asli disajikan kemudian sebagai bagian yang terpisah dan berlabel jelas.
\newtheorem{Korolari}[fakt]{Korolari}
\newtheorem{Akibat}[fakt]{Akibat}

\part{Unit 23}
\chapter[Kuliah 23]{Kuliah 23: Teorema Stokes dan Teorema Titik Tetap Brouwer}
\setcounter{section}{23}

\bild{ \begin{center}
\includegraphics[width=5.5cm]{ \bildeinlesungjpg {SS-stokes} {jpg} }
\end{center}
\bildtext {George Stokes (1819 -1903)} }

\bildlizenz { SS-stokes.jpg } {} {Kelson} {Commons} {PD} {}

Teorema Stokes termasuk salah satu teorema terpenting dalam matematika. Teorema ini membentuk hubungan langsung antara integral suatu bentuk diferensial di atas batas manifold berbatas dan integral turunan eksterior bentuk tersebut di atas seluruh manifold. Dengan demikian, teorema ini merupakan perumuman yang sangat luas dari teorema dasar kalkulus, yang menyatakan bahwa integral tentu suatu fungsi yang didefinisikan pada sebuah interval dapat dinyatakan melalui antiturunannya hanya dengan menggunakan nilai-nilai pada ujung interval.

  \zwischenueberschrift{Teorema Stokes versi balok}

Sebelum kita merumuskan dan membuktikan teorema Stokes secara umum, kita berikan dahulu versi baloknya. Dalam versi ini, domain definisi bentuk diferensial adalah sebuah balok yang batasnya terdiri atas muka-mukanya. Agar objek geometris ini menjadi manifold berbatas, kita harus membuang \anfuehrung{rusuk-rusuknya}{} dari balok tersebut. Namun, rusuk-rusuk itu merupakan himpunan nol di dalam setiap muka
\zusatzklammer {dan demikian pula muka-muka merupakan himpunan nol di dalam keseluruhan balok} {} {,}
sehingga himpunan-himpunan bagian tersebut dapat diabaikan dalam pengintegralan.



\inputfaktbeweis
{Satz von Stokes/Quaderversion/Fakt}
{Teorema}
{}
{

\faktsituation {Misalkan

\mavergleichskette
{\vergleichskette
{ Q
}
{ \subseteq }{ \R^n
}
{  }{
}
{    }{
}
{   }{
}
}
{}{}{}
sebuah balok berdimensi $n$ yang sejajar sumbu
\zusatzklammer {beserta muka-mukanya, tetapi tanpa rusuk-rusuknya} {} {}\zusatzfussnote {Syarat-syarat ini menjamin bahwa yang diperoleh adalah manifold berbatas, dengan batas berupa gabungan saling lepas dari muka-muka terbuka. Namun, dalam pembuktian kita juga akan menggunakan balok tertutup} {.} {}
dengan
\definitionsverweis {batas}{}{}

\mathl{\partial Q}{,} dan misalkan $\omega$ didefinisikan pada $Q$ sebagai suatu
\definitionsverweis {bentuk diferensial}{}{}
yang
\definitionsverweis {terdiferensialkan secara kontinu}{}{}
dan berderajat $(n-1)$.}
\faktfolgerung {Maka

\mavergleichskettedisp
{\vergleichskette
{
\int_{  Q  }  d \omega
}
{ =} {
\int_{  \partial Q  }   \omega
}
{ } {
}
{ } {
}
{ } {
}
}
{}{}{.}}
\faktzusatz {}
\faktzusatz {}

}
{

Karena kedua ruas persamaan ini linear terhadap $\omega$, kita dapat menganggap bahwa $\omega$ berbentuk

\mavergleichskettedisp
{\vergleichskette
{ \omega
}
{ =} { f dx_2 \wedge \ldots \wedge dx_n
}
{ } {
}
{ } {
}
{ } {
}
}
{}{}{}
dengan suatu fungsi yang terdiferensialkan secara kontinu dan didefinisikan pada suatu lingkungan terbuka dari $Q$, yaitu $f$. Integral pada ruas kiri dan ruas kanan adalah
\definitionsverweis {integral Lebesgue}{}{}
dari fungsi-fungsi kontinu pada himpunan bagian
\mathkor {} {\R^n} {dan} {\R^{n-1}} {.}
Karena itu, pada kedua ruas kita dapat beralih ke
\definitionsverweis {penutupan topologis}{}{}
sebab hal tersebut hanya mengubah daerah integrasi yang bersangkutan sebesar suatu himpunan nol, sehingga tidak mengubah integralnya.

Kita tuliskan balok tertutup itu sebagai

\mavergleichskettedisp
{\vergleichskette
{ \overline{Q}
}
{ =} { [a,b] \times \tilde{Q}
}
{ } {
}
{ } {
}
{ } {
}
}
{}{}{.}
Kita menerapkan
Korolari 14.10
pada setiap muka $S$, kecuali
\mathkor {} {a \times \tilde{Q}} {dan} {b \times \tilde{Q}} {,}
dan memperoleh pada muka tersebut

\mavergleichskettedisp
{\vergleichskette
{
\int_{ S  }   \omega
}
{ =} {
\int_{  S  }  f dx_2 \wedge \ldots \wedge dx_n
}
{ =} {
\int_{ S  }   0
}
{ =} { 0
}
{ } {
}
}
{}{}{,}
karena pada setiap muka ini salah satu variabel
\mathl{x_2 , \ldots , x_n}{} konstan. Berdasarkan
teorema Fubini
dan
teorema dasar kalkulus
\zusatzklammer {yang diterapkan untuk setiap \mathlk{(x_2 , \ldots , x_n)}{} yang tetap} {} {,}
berlaku

\mavergleichskettealignhandlinks
{\vergleichskettealignhandlinks
{
\int_{ Q  }   d\omega
}
{ =} {
\int_{  Q  }   df \wedge  dx_2 \wedge \ldots \wedge dx_n
}
{ =} {
\int_{  Q  }   \left(  \sum_{j = 1}^n { \frac{   \partial   f   }{  \partial   x_j   } } dx_j  \right) \wedge  dx_2 \wedge \ldots \wedge dx_n
}
{ =} {
\int_{  Q  }   { \frac{   \partial   f   }{  \partial   x_1   } } dx_1 \wedge  dx_2 \wedge \ldots \wedge dx_n
}
{ =} {
\int_{ \tilde{Q}   }   \left(  \int_{[a,b] } { \frac{   \partial   f   }{  \partial   x_1   } } dx_1  \right)  dx_2 \wedge \ldots \wedge dx_n
}
}
{
\vergleichskettefortsetzungalign
{ =} {
\int_{  \tilde{Q}   }   \left(  f(b,x_2 , \ldots , x_n ) - f(a,x_2 , \ldots , x_n )   \right) dx_2 \wedge \ldots \wedge dx_n
}
{ =} {
\int_{ b \times \tilde{Q}   }   f dx_2 \wedge \ldots \wedge dx_n  -
\int_{  a \times \tilde{Q}   }  f dx_2 \wedge \ldots \wedge dx_n
}
{ =} { \sum_{S \text{ muka berorientasi dari } Q}
\int_{  S  }  f dx_2 \wedge \ldots \wedge dx_n
}
{ =} {
\int_{  \partial Q  }   f dx_2 \wedge \ldots \wedge dx_n
}
}
{
\vergleichskettefortsetzungalign
{ =} {
\int_{  \partial Q  }   \omega
}
{ } {}
{ } {}
{ } {}
}{.}

}

  \zwischenueberschrift{Teorema Stokes}



\inputfaktbeweis
{Mannigfaltigkeit mit Rand/Satz von Stokes/Fakt}
{Teorema}
{}
{

\faktsituation {Misalkan $M$ suatu
\definitionsverweis {manifold diferensiabel berbatas}{}{}
yang
\definitionsverweis {berdimensi}{}{} $n$,
\definitionsverweis {berorientasi}{}{},
memiliki batas $\partial M$, serta memiliki
\definitionsverweis {basis terhitung bagi topologinya}{}{.}
Misalkan pula $\omega$ suatu
\definitionsverweis {bentuk diferensial}{}{}
berderajat $(n-1)$ yang
\definitionsverweis {terdiferensialkan secara kontinu}{}{}
dan memiliki
\definitionsverweis {tumpuan}{}{}
\definitionsverweis {kompak}{}{}\zusatzfussnote {Tumpuan suatu bentuk diferensial adalah penutupan topologis dari himpunan titik tempat bentuk tersebut \mathlk{\neq 0}{}} {.} {} pada $M$.}
\faktfolgerung {Maka

\mavergleichskettedisp
{\vergleichskette
{
\int_{ M  }  d\omega
}
{ =} {
\int_{  \partial M  }   \omega
}
{ } {
}
{ } {
}
{ } {
}
}
{}{}{.}}
\faktzusatz {}
\faktzusatz {}

}
{

\teilbeweis {}{}{}
{Misalkan

\mathbed {U_i} {}
 {i \in I} {}
 {} {} {} {,}
suatu
\definitionsverweis {tutupan terbuka}{}{}
dari $M$ oleh
\definitionsverweis {peta-peta berorientasi}{}{,}
dan misalkan

\mathbed {h_j} {}
 {j \in J} {}
 {} {} {} {,}
suatu
\definitionsverweis {partisi satuan yang terdiferensialkan secara kontinu dan subordinat terhadap tutupan ini}{}{,}
yang ada menurut
Teorema 22.10.
Untuk setiap

\mavergleichskette
{\vergleichskette
{ P
}
{ \in }{ M
}
{  }{
}
{    }{
}
{   }{
}
}
{}{}{}
terdapat suatu lingkungan terbuka

\mavergleichskette
{\vergleichskette
{ P
}
{ \in }{ V_P
}
{ \subseteq }{ M
}
{    }{
}
{   }{
}
}
{}{}{}
sedemikian rupa sehingga

\mavergleichskette
{\vergleichskette
{ h_j \vert_{V_P }
}
{ = }{ 0
}
{  }{
}
{    }{
}
{   }{
}
}
{}{}{}
selain untuk sejumlah berhingga pengecualian. Misalkan $Y$ adalah tumpuan $\omega$. Tutupan

\mavergleichskette
{\vergleichskette
{ Y
}
{ \subseteq }{ \bigcup_{P \in M} V_P
}
{  }{
}
{    }{
}
{   }{
}
}
{}{}{}
memiliki suatu subtutupan berhingga karena
\definitionsverweis {kekompakan}{}{}
yang diasumsikan; sebutlah

\mavergleichskettedisp
{\vergleichskette
{ Y
}
{ \subseteq} { V_{1} \cup \ldots \cup V_{r}
}
{ =} { V
}
{ } {
}
{ } {
}
}
{}{}{.}
Karena itu, hanya sejumlah berhingga dari $h_j$ pada $V$ yang berbeda dari $0$. Kita definisikan

\mavergleichskette
{\vergleichskette
{ \omega_j
}
{ = }{ h_j \omega
}
{  }{
}
{    }{
}
{   }{
}
}
{}{}{;}
bentuk-bentuk diferensial ini juga terdiferensialkan secara kontinu. Tumpuan $h_j$ adalah suatu
\definitionsverweis {himpunan bagian tertutup}{}{}
dari $M$ yang termuat dalam
\mathl{U_{i(j)}}{.} Oleh karena itu, tumpuan $\omega_j$ termuat dalam
\mathl{Y \cap U_{i(j)}}{} dan dengan sendirinya kompak menurut
Soal 13.10.
Berlaku

\mavergleichskettedisp
{\vergleichskette
{ \omega
}
{ =} { \sum_{j \in J} \omega_j
}
{ } {
}
{ } {
}
{ } {
}
}
{}{}{,}
dan hanya sejumlah berhingga dari bentuk-bentuk diferensial
\mathl{\omega_j}{} tersebut yang berbeda dari $0$, sebab

\mavergleichskette
{\vergleichskette
{ \omega_j \vert_{M \setminus Y}
}
{ = }{ 0
}
{  }{
}
{    }{
}
{   }{
}
}
{}{}{}
untuk semua $j$, sedangkan

\mavergleichskettedisp
{\vergleichskette
{ \omega_j \vert_V
}
{ =} { 0
}
{ } {
}
{ } {
}
{ } {
}
}
{}{}{}
untuk semua $j$ selain sejumlah berhingga pengecualian. Berdasarkan
aditivitas integral bentuk diferensial
dan
aditivitas turunan eksterior,
pernyataan tersebut dapat dibuktikan secara terpisah untuk setiap $\omega_j$. Jadi, kita dapat menganggap bahwa diberikan suatu bentuk diferensial berderajat $(n-1)$ yang terdiferensialkan secara kontinu, mempunyai tumpuan kompak, dan tumpuan itu seluruhnya termuat dalam suatu lingkungan peta $U$.}
{}
\teilbeweis {}{}{}
{Terdapat suatu
\definitionsverweis {difeomorfisme}{}{}
\maabb {\alpha} { U } { V
} {}
dengan

\mavergleichskette
{\vergleichskette
{ V
}
{ \subseteq }{ H
}
{  }{
}
{    }{
}
{   }{
}
}
{}{}{}
terbuka, yang sekaligus menginduksi suatu difeomorfisme antara batas-batas

\mavergleichskette
{\vergleichskette
{ \partial U
}
{ = }{ \partial M \cap U
}
{  }{
}
{    }{
}
{   }{
}
}
{}{}{}
dan

\mavergleichskette
{\vergleichskette
{ \partial V
}
{ = }{ \partial H \cap V
}
{  }{
}
{    }{
}
{   }{
}
}
{}{}{.}
Dalam hal ini berlaku

\mavergleichskettedisp
{\vergleichskette
{
\int_{  \partial U  }   \omega
}
{ =} {
\int_{  \partial V  }  \alpha^{-1 *}\omega
}
{ } {
}
{ } {
}
{ } {
}
}
{}{}{}
dan

\mavergleichskettedisp
{\vergleichskette
{
\int_{ U  }  d \omega
}
{ =} {
\int_{ V  }  \alpha^{-1 *} (  d \omega)
}
{ =} {
\int_{  V  }  d \left(  \alpha^{-1 *}\omega  \right)
}
{ } {
}
{ } {
}
}
{}{}{}
menurut
Lema 20.2 (5).}
{\leerzeichen{}Dengan demikian, kita dapat bertolak dari suatu bentuk diferensial yang terdiferensialkan secara kontinu, didefinisikan pada

\mavergleichskette
{\vergleichskette
{ V
}
{ \subseteq }{ H
}
{  }{
}
{    }{
}
{   }{
}
}
{}{}{,}
serta mempunyai tumpuan kompak. Di luar tumpuannya, bentuk ini dapat kita perluas ke seluruh $H$ sebagai bentuk nol.}
\teilbeweis {}{}{}
{Karena tumpuannya kompak, terdapat suatu balok yang cukup besar

\mavergleichskette
{\vergleichskette
{ Q
}
{ \subseteq }{ H
}
{  }{
}
{    }{
}
{   }{
}
}
{}{}{,}
yang salah satu mukanya, $S$, terletak pada $\partial H$ dan yang hanya berpotongan dengan tumpuan $\omega$ di $S$. Pada semua muka lain dari $Q$, $\omega$
\zusatzklammer {dan karena itu juga $d\omega$} {} {}
adalah bentuk nol. Oleh sebab itu, pada satu pihak berlaku

\mavergleichskettedisp
{\vergleichskette
{
\int_{ H  }  d\omega
}
{ =} {
\int_{ Q  }  d\omega
}
{ } {
}
{ } {
}
{ } {
}
}
{}{}{}
dan pada pihak lain

\mavergleichskettedisp
{\vergleichskette
{
\int_{  \partial H  }   \omega
}
{ =} {
\int_{  S  }   \omega
}
{ =} {
\int_{  \partial Q  }   \omega
}
{ } {
}
{ } {
}
}
{}{}{.}
Dengan demikian, pernyataan tersebut mengikuti dari
teorema Stokes versi balok.}
{}

}



\inputfaktbeweis
{Mannigfaltigkeiten ohne Rand/Satz von Stokes/Fakt}
{Korolari}
{}
{

\faktsituation {Misalkan $M$ suatu manifold
\definitionsverweis {berdimensi}{}{} $n$ yang
\definitionsverweis {berorientasi}{}{} dan
\definitionsverweis {diferensiabel}{}{}
\zusatzklammer {tanpa batas} {} {}
serta memiliki
\definitionsverweis {basis terhitung bagi topologinya}{}{.}
Misalkan pula $\omega$ suatu
\definitionsverweis {bentuk diferensial}{}{}
berderajat $(n-1)$ yang
\definitionsverweis {terdiferensialkan secara kontinu}{}{}
dan memiliki
\definitionsverweis {tumpuan}{}{}
\definitionsverweis {kompak}{}{} pada $M$.}
\faktfolgerung {Maka

\mavergleichskettedisp
{\vergleichskette
{
\int_{ M  }  d\omega
}
{ =} { 0
}
{ } {
}
{ } {
}
{ } {
}
}
{}{}{.}}
\faktzusatz {}
\faktzusatz {}

}
{

Hal ini segera mengikuti dari
Teorema 23.2
karena

\mavergleichskette
{\vergleichskette
{ \partial M
}
{ = }{ \emptyset
}
{  }{
}
{    }{
}
{   }{
}
}
{}{}{.}

}



\inputfaktbeweis
{Mannigfaltigkeiten mit Rand/Satz von Stokes/Differentialform ist null auf Rand/Fakt}
{Korolari}
{}
{

\faktsituation {Misalkan $M$ suatu
\definitionsverweis {manifold diferensiabel berbatas}{}{}
yang
\definitionsverweis {berdimensi}{}{} $n$,
\definitionsverweis {berorientasi}{}{},
dan memiliki
\definitionsverweis {basis terhitung bagi topologinya}{}{.}
Misalkan pula $\omega$ suatu
\definitionsverweis {bentuk diferensial}{}{}
berderajat $(n-1)$ yang
\definitionsverweis {terdiferensialkan secara kontinu}{}{}
dan memiliki
\definitionsverweis {tumpuan}{}{}
\definitionsverweis {kompak}{}{}
pada $M$,}
\faktvoraussetzung {serta pada batas $\partial M$ bernilai tetap $0$.}
\faktfolgerung {Maka

\mavergleichskettedisp
{\vergleichskette
{
\int_{ M  }  d\omega
}
{ =} { 0
}
{ } {
}
{ } {
}
{ } {
}
}
{}{}{.}}
\faktzusatz {}
\faktzusatz {}

}
{

Hal ini segera mengikuti dari
Teorema 23.2.

}

Hal penting dalam pernyataan di atas adalah bahwa $\omega$ bernilai $0$ pada batas; tidak cukup bahwa turunan eksterior $d\omega$ bernilai $0$ pada batas, sebagaimana sudah ditunjukkan oleh situasi berdimensi satu.

Ada banyak cara untuk merealisasikan bentuk volume

\mavergleichskette
{\vergleichskette
{ \tau
}
{ = }{ dx_1 \wedge \ldots \wedge dx_n
}
{  }{
}
{    }{
}
{   }{
}
}
{}{}{}
dari $\R^n$ sebagai turunan eksterior suatu bentuk berderajat
\mathl{(n-1)}{,} misalnya dengan

\mavergleichskette
{\vergleichskette
{  \omega
}
{ = }{ x_1dx_2 \wedge \ldots \wedge dx_n
}
{  }{
}
{    }{
}
{   }{
}
}
{}{}{.}
Dengan demikian, perhitungan volume suatu benda berbatas dapat dikembalikan pada perhitungan integral di atas batasnya. Dalam kasus bidang, pernyataan ini juga disebut \stichwort {teorema Green} {.}



\inputfaktbeweis
{Integration auf ebener Mannigfaltigkeit mit Rand/Satz von Green/Fakt}
{Teorema}
{}
{

\faktsituation {Misalkan

\mavergleichskette
{\vergleichskette
{M
}
{ \subseteq }{ \R^2
}
{  }{
}
{    }{
}
{   }{
}
}
{}{}{}
suatu manifold
\definitionsverweis {kompak}{}{}
\definitionsverweis {dengan batas}{}{}\zusatzfussnote {Bidang real sekitar hanya berperan dalam menetapkan koordinat dan bentuk diferensial pada $M$} {.} {}
$\partial M$, dan misalkan
\maabbdisp {f,g} { M } {\R
} {}
\definitionsverweis {fungsi-fungsi yang terdiferensialkan secara kontinu}{}{.}}
\faktfolgerung {Maka

\mavergleichskettealignhandlinks
{\vergleichskettealignhandlinks
{
\int_{  M  }   \left(  - { \frac{   \partial   f   }{  \partial   y   } } ( x,y) + { \frac{   \partial   g   }{  \partial   x   } } ( x,y)  \right) dx \wedge dy
}
{ =} {
\int_{  \partial M  }  f(x,y)dx +g(x,y)dy
}
{ } {
}
{ } {
}
{ } {
}
}
{}
{}{}}
\faktzusatz {}
\faktzusatz {}

}
{

Hal ini mengikuti dari
Teorema 23.2,
yang diterapkan pada suatu bentuk-$1$ yang
\definitionsverweis {terdiferensialkan secara kontinu}{}{},
yakni suatu
\definitionsverweis {bentuk diferensial}{}{}

\mavergleichskette
{\vergleichskette
{ \omega
}
{ = }{ f(x,y)dx + g(x,y)dy
}
{  }{
}
{    }{
}
{   }{
}
}
{}{}{.}

}



\inputfaktbeweis
{Integration auf ebener Mannigfaltigkeit mit Rand/Satz von Green/Flächenversion/Fakt}
{Korolari}
{}
{

\faktsituation {Misalkan

\mavergleichskette
{\vergleichskette
{ M
}
{ \subseteq }{ \R^2
}
{  }{
}
{    }{
}
{   }{
}
}
{}{}{}
suatu manifold
\definitionsverweis {kompak}{}{}
\definitionsverweis {dengan batas}{}{}

\mathl{\partial M}{.}}
\faktfolgerung {Maka luas $M$ sama dengan

\mavergleichskettedisp
{\vergleichskette
{ \lambda^2(M)
}
{ =} {
\int_{  \partial M  }   x dy
}
{ =} { -
\int_{  \partial M  }   y dx
}
{ } {
}
{ } {
}
}
{}{}{.}}
\faktzusatz {}
\faktzusatz {}

}
{

Hal ini mengikuti karena

\mavergleichskette
{\vergleichskette
{ \lambda^2(M)
}
{ = }{
\int_{ M  }  dx \wedge dy
}
{  }{
}
{    }{
}
{   }{
}
}
{}{}{}
dari
Teorema 23.5,
yang diterapkan pada
\mathkor {} {f=0,\, g=x} {atau} {f=-y,\, g=0} {.}

}

\inputbemerkung
{}
{

Teorema Green juga dapat diterapkan pada daerah kompak

\mavergleichskette
{\vergleichskette
{ M
}
{ \subseteq }{ \R^2
}
{  }{
}
{    }{
}
{   }{
}
}
{}{}{}
yang batas
\zusatzklammer {topologisnya} {} {}
terhubung dan tersusun atas sejumlah berhingga kurva mulus yang tidak harus mulus pada titik-titik sambungannya
\zusatzklammer {misalnya sebuah segitiga} {} {.}
Integral batas kemudian merupakan jumlah integral lintasan di atas potongan-potongan kurva mulus tersebut. Situasi yang sedikit lebih umum ini dapat dikembalikan pada
Teorema 23.5
dengan salah satu dari dua cara: melicinkan titik-titik sambungan untuk memperoleh batas mulus, sambil membuat penyimpangan integral sekecil yang diinginkan; atau melicinkan bentuk diferensial hingga bernilai $0$ dalam lingkungan-lingkungan kecil dari titik-titik sambungan.

}

\inputbeispiel{}
{

Salah satu kasus khusus
teorema Stokes
adalah apa yang disebut \stichwort {teorema divergensi} {} atau \stichwort {teorema Gauss} {.} Untuk suatu manifold kompak berdimensi tiga dengan batas

\mavergleichskette
{\vergleichskette
{ M
}
{ \subset }{ U
}
{ \subseteq }{ \R^3
}
{    }{
}
{   }{
}
}
{}{}{}
\zusatzklammer {$U$ suatu himpunan bagian terbuka} {} {}
dan suatu bentuk diferensial-$2$ $\omega$ pada $U$, teorema ini menyatakan persamaan

\mavergleichskettedisp
{\vergleichskette
{ \int_{\partial M} \omega
}
{ =} { \int_M d \omega
}
{ } {
}
{ } {
}
{ } {
}
}
{}{}{.}
Teorema ini berkaitan dengan situasi fisis suatu aliran. Bentuk $\omega$, untuk suatu titik

\mavergleichskette
{\vergleichskette
{ P
}
{ \in }{ U
}
{  }{
}
{    }{
}
{   }{
}
}
{}{}{}
dan sebuah jajaran genjang infinitesimal pada titik itu, menyatakan berapa banyak fluida yang mengalir menembus potongan tersebut per satuan waktu. Dalam deskripsi ekuivalen situasi ini dengan suatu medan vektor,
\mathl{F(P)}{} menyatakan arah aliran beserta kekuatannya. Turunan $d \omega$ adalah suatu bentuk volume yang disebut \stichwort {divergensi} {.} Pada suatu titik, bentuk ini menyatakan pertambahan infinitesimal bahan aliran
\zusatzklammer {\stichwort {sumber} {} atau \stichwort {serapan} {}} {} {}
di titik tersebut. Syarat

\mavergleichskette
{\vergleichskette
{d \omega
}
{ = }{ 0
}
{  }{
}
{    }{
}
{   }{
}
}
{}{}{}
sering dipenuhi dan berarti bahwa fluida tidak mengalami penambahan bahan di $U$. Teorema Gauss kemudian menyatakan bahwa aliran total yang menembus batas
\zusatzklammer {dengan orientasi menentukan aliran mana yang dianggap mengarah ke luar atau ke dalam} {} {}
sama dengan $0$. Jadi, apa pun yang mengalir masuk ke $M$ pada suatu tempat akan mengalir keluar lagi di tempat lain.

}

  \zwischenueberschrift{Teorema titik tetap Brouwer}



\inputfaktbeweis
{Mannigfaltigkeit mit Rand/Stokes/Nichtexistenz von Retraktionen/Fakt}
{Teorema}
{}
{

\faktsituation {Misalkan $M$ suatu
\definitionsverweis {manifold diferensiabel berbatas}{}{}
yang
\definitionsverweis {kompak}{}{},
\definitionsverweis {berorientasi}{}{},
memiliki batas
\mathl{\partial M}{,} dan memiliki
\definitionsverweis {basis terhitung bagi topologinya}{}{.}}
\faktfolgerung {Maka tidak ada
\definitionsverweis {pemetaan yang terdiferensialkan secara kontinu}{}{}
\maabbdisp {\varphi} { M } { \partial M
} {,}
yang
\definitionsverweis {pembatasannya}{}{}
pada $\partial M$ adalah identitas.}
\faktzusatz {}
\faktzusatz {}

}
{

Menurut
Teorema 21.8,
batas
\mathl{\partial M}{} adalah suatu
\definitionsverweis {manifold diferensiabel}{}{}
\definitionsverweis {berorientasi}{}{}
\zusatzklammer {tanpa batas} {} {.}
Oleh karena itu, menurut
Teorema 22.11,
terdapat suatu
\definitionsverweis {bentuk volume positif}{}{}
yang
\definitionsverweis {terdiferensialkan secara kontinu}{}{},
yaitu $\tau$, pada
\mathl{\partial M}{.} Berlaku

\mavergleichskette
{\vergleichskette
{
\int_{  \partial M  }  \tau
}
{ > }{ 0
}
{  }{
}
{    }{
}
{   }{
}
}
{}{}{.}
\definitionsverweis {Turunan eksterior}{}{}
dari bentuk volume $\tau$ adalah $0$.
Andaikan terdapat suatu
\definitionsverweis {pemetaan yang terdiferensialkan secara kontinu}{}{}
\maabbdisp {\varphi} { M } { \partial M
} {}
dengan

\mavergleichskette
{\vergleichskette
{ \varphi \vert_{\partial M}
}
{ = }{ \operatorname{Id}_{\partial M}
}
{  }{
}
{    }{
}
{   }{
}
}
{}{}{.}
Maka
\definitionsverweis {bentuk tarik balik}{}{}

\mathl{\varphi^* \tau}{} adalah suatu
\definitionsverweis {bentuk diferensial}{}{}
berderajat $(n-1)$ pada $M$, yang pembatasannya pada batas sama dengan $\tau$. Dengan menggunakan
Teorema 23.2
dan
Teorema 20.4 (5),
kita memperoleh

\mavergleichskettealign
{\vergleichskettealign
{
\int_{  \partial M  }   \tau
}
{ =} {
\int_{  \partial  M  }   \varphi^* \tau
}
{ =} {
\int_{  M  }   d(\varphi^* \tau)
}
{ =} {
\int_{  M  }   \varphi^* (d\tau)
}
{ =} {
\int_{  M  }   \varphi^* (0)
}
}
{
\vergleichskettefortsetzungalign
{ =} { 0
}
{ } {}
{ } {}
{ } {}
}
{}{.}
Ini sebuah kontradiksi.

}

Pernyataan ini juga dirumuskan dengan mengatakan bahwa tidak ada
\zusatzklammer {yang terdiferensialkan secara kontinu} {} {} \stichwort {retraksi} {} ke batas.

\bild{ \begin{center}
\includegraphics[width=5.5cm]{ \bildeinlesungjpg {Théorème-de-Brouwer-(cond-2)} {jpg} }
\end{center}
\bildtext {Ilustrasi cakram dengan lintasan-lintasan berarah yang berpilin menuju titik pusat.} }

\bildlizenz { Théorème-de-Brouwer-(cond-2).jpg } {} {Jean-Luc W} {Commons} {CC-by-sa 3.0} {}

\bild{ \begin{center}
\includegraphics[width=5.5cm]{ \bildeinlesungjpg {Théorème-de-Brouwer-(cond-1)} {jpg} }
\end{center}
\bildtext {Ilustrasi lintasan-lintasan melingkar berarah pada anulus dan pada pita persegi panjang padanannya.} }

\bildlizenz { Théorème-de-Brouwer-(cond-1).jpg } {} {Jean-Luc W} {Commons} {CC-by-sa 3.0} {}

Teorema berikut disebut \stichwort {teorema titik tetap Brouwer} {.}


\inputfaktbeweis
{Brouwersche Fixpunktsatz/Stetig differenzierbar/Retraktion/Fakt}
{Teorema}
{}
{

\faktsituation {Misalkan
\maabbdisp {\psi} { B  \left(  0 ,r \right) } { B  \left(  0 ,r \right)
} {}
suatu
\definitionsverweis {pemetaan yang terdiferensialkan secara kontinu}{}{}
dari
\definitionsverweis {bola tertutup}{}{}
di $\R^n$ ke dirinya sendiri.}
\faktfolgerung {Maka $\psi$ memiliki suatu
\definitionsverweis {titik tetap}{}{.}}
\faktzusatz {}
\faktzusatz {}

}
{

Untuk menyederhanakan notasi, misalkan

\mavergleichskette
{\vergleichskette
{ r
}
{ = }{ 1
}
{  }{
}
{    }{
}
{   }{
}
}
{}{}{.}
Andaikan ada suatu pemetaan $\psi$ yang terdiferensialkan secara kontinu dan tidak memiliki titik tetap. Maka selalu berlaku

\mavergleichskettedisp
{\vergleichskette
{ x
}
{ \neq} { \psi(x)
}
{ } {
}
{ } {
}
{ } {
}
}
{}{}{,}
sehingga kedua titik itu menentukan sebuah garis. Gagasannya adalah menggunakan garis ini untuk mengambil salah satu
\zusatzklammer {dari kedua} {} {}
titik potong dengan sfer sebagai titik citra suatu retraksi ke batas. Dengan fungsi bantu

\mavergleichskettedisp
{\vergleichskette
{ h(x)
}
{ =} { { \frac{   \psi(x)-x }{  \Vert { \psi (x)-x} \Vert  } }
}
{ } {
}
{ } {
}
{ } {
}
}
{}{}{,}
kita mendefinisikan suatu pemetaan
\maabbdisp {\varphi} { B  \left(  0 , 1 \right) } { \R^n
} {}
melalui

\mavergleichskettedisphandlinks
{\vergleichskettedisphandlinks
{ \varphi(x)
}
{ =} { x + \left(  - \left\langle  x  , h(x)  \right\rangle + \sqrt{1 + \left\langle  x  , h(x)  \right\rangle^2 - \Vert { x} \Vert^2  }  \right) \cdot h(x)
}
{ } {
}
{ } {
}
{ } {
}
}
{}{}{.}
Ekspresi di bawah akar kuadrat bernilai positif. Hal ini jelas jika

\mavergleichskette
{\vergleichskette
{ \Vert { x } \Vert
}
{ < }{ 1
}
{  }{
}
{    }{
}
{   }{
}
}
{}{}{,}
sedangkan jika

\mavergleichskette
{\vergleichskette
{ \Vert { x } \Vert
}
{ = }{ 1
}
{  }{
}
{    }{
}
{   }{
}
}
{}{}{,}
terdapat suatu titik pada sfer yang garis penghubungnya dengan titik bola
\mathl{\psi(x)}{} tidak tegak lurus terhadap $x$
\zusatzklammer {ruang singgung afin pada suatu titik sfer hanya berpotongan dengan bola di satu titik} {} {,}
sehingga

\mavergleichskettedisp
{\vergleichskette
{ \left\langle  x  , h(x)  \right\rangle
}
{ \neq} { 0
}
{ } {
}
{ } {
}
{ } {
}
}
{}{}{.}
Karena akar kuadrat dan nilai mutlak
\definitionsverweis {terdiferensialkan secara kontinu}{}{}
di luar titik nol, baik
\mathkor {} {h(x)} {maupun} {\varphi(x)} {}
merupakan pemetaan yang terdiferensialkan secara kontinu. Menurut
Soal 23.23,
pemetaan $\varphi$ memetakan bola ke sfer, dan pembatasannya pada sfer adalah identitas. Dengan demikian, terdapat suatu
\definitionsverweis {retraksi}{}{}
yang terdiferensialkan secara kontinu dari bola pejal tertutup ke batasnya, yang mustahil menurut
Teorema 23.9.

}

\chapter{Lembar Kerja 23}
\setcounter{section}{23}

  \zwischenueberschrift{Soal latihan}

\inputaufgabe
{}
{

Bahaslah
rumus Newton--Leibniz
sebagai kasus khusus dari
teorema Stokes.

}
{} {}

\inputaufgabe
{}
{

Buktikan secara langsung
versi balok dari teorema Stokes
untuk balok
\mathl{[a_1,b_1] \times \cdots \times [a_n,b_n]}{} dan suatu
$n-1$-\definitionsverweis {bentuk diferensial}{}{}
$\omega$ berbentuk

\mavergleichskettedisp
{\vergleichskette
{ \omega
}
{ =} { x_1^{m_1 } \cdots x_n^{m_n }  dx_2 \wedge \ldots \wedge dx_n
}
{ } {
}
{ } {
}
{ } {
}
}
{}{}{.}

}
{} {}

\inputaufgabe
{}
{

Misalkan $M$ suatu
$n$-\definitionsverweis {dimensional}{}{}
\definitionsverweis {berorientasi}{}{}
\definitionsverweis {manifold diferensiabel berbatas}{}{}
dengan
\definitionsverweis {basis topologi terhitung}{}{},
dan misalkan $\omega$ suatu $(n-1)$-\definitionsverweis {bentuk diferensial}{}{}
yang
\definitionsverweis {terdiferensialkan secara kontinu}{}{},
\definitionsverweis {tertutup}{}{}, dan mempunyai
\definitionsverweis {tumpuan kompak}{}{}
pada $M$. Tunjukkan bahwa

\mavergleichskettedisp
{\vergleichskette
{ \int_{\partial M} \omega
}
{ =} { 0
}
{ } {
}
{ } {
}
{ } {
}
}
{}{}{.}

}
{} {}

\inputaufgabe
{}
{

Misalkan $M$ suatu
\definitionsverweis {manifold diferensiabel}{}{}
$n$-\definitionsverweis {dimensional}{}{},
\definitionsverweis {kompak}{}{}, dan
\definitionsverweis {berorientasi}{}{}
\zusatzklammer {tanpa batas} {} {},
dengan
\definitionsverweis {basis topologi terhitung}{}{},
dan misalkan $\omega$ suatu
$(n-1)$-\definitionsverweis {bentuk diferensial}{}{}
yang
\definitionsverweis {terdiferensialkan secara kontinu}{}{}
pada $M$. Tunjukkan bahwa

\mavergleichskettedisp
{\vergleichskette
{
\int_{ M  }  d\omega
}
{ =} { 0
}
{ } {
}
{ } {
}
{ } {
}
}
{}{}{.}

}
{Apa arti pernyataan ini untuk $S^1$? Bagaimana pernyataan tersebut dapat dibuktikan dalam kasus ini melalui suatu integral lintasan?} {}

\inputaufgabe
{}
{

Misalkan $M$ suatu
\definitionsverweis {manifold diferensiabel}{}{}
\definitionsverweis {kompak}{}{} dan
\definitionsverweis {berorientasi}{}{}
\zusatzklammer {tanpa batas} {} {},
dengan
\definitionsverweis {basis topologi terhitung}{}{},
dan misalkan $\tau$ suatu
\definitionsverweis {bentuk volume positif}{}{}
pada $M$. Tunjukkan bahwa $\tau$ tidak
\definitionsverweis {eksak}{}{.}

}
{Bagaimana keadaannya tanpa asumsi kekompakan?} {}

\inputaufgabegibtloesung
{}
{

Misalkan $M$ suatu
\definitionsverweis {manifold diferensiabel}{}{}
dan $\omega$ suatu
\definitionsverweis {eksak}{}{} \definitionsverweis {bentuk diferensial}{}{}
pada $M$. Misalkan $S$ suatu
\definitionsverweis {manifold diferensiabel}{}{}
\definitionsverweis {kompak}{}{} dan
\definitionsverweis {berorientasi}{}{}
\zusatzklammer {tanpa batas} {} {},
dengan
\definitionsverweis {basis topologi terhitung}{}{},
serta misalkan
\maabbdisp {\varphi} { S } { M
} {}
suatu pemetaan yang
\definitionsverweis {terdiferensialkan secara kontinu}{}{.}
Tunjukkan bahwa

\mavergleichskettedisp
{\vergleichskette
{
\int_{  S  }   \varphi^*\omega
}
{ =} { 0
}
{ } {
}
{ } {
}
{ } {
}
}
{}{}{.}

}
{} {}

\inputaufgabe
{}
{

Kita tinjau
\definitionsverweis {pemetaan yang terdiferensialkan secara kontinu}{}{}
\zusatzklammer {\mathlk{n \geq 2}{}} {} {}
\maabbeledisp {\pi} {\R^n \setminus \{0\}} { S^{n-1}
} {(x_1 , \ldots , x_n) } { { \frac{   1  }{  \sqrt{x_1^2 + \cdots + x_n^2}  } } (x_1 , \ldots , x_n)
} {.}
Misalkan $\omega$ adalah
\definitionsverweis {bentuk volume kanonik}{}{}
pada $S^{n-1}$. Tunjukkan bahwa
\mathl{\pi^* \omega}{} pada $\R^n \setminus \{0\}$ merupakan
\definitionsverweis {tertutup}{}{},
tetapi tidak
\definitionsverweis {eksak}{}{}
$n-1$-\definitionsverweis {bentuk diferensial}{}{.}

}
{} {}

\inputaufgabe
{}
{

Misalkan

\mavergleichskette
{\vergleichskette
{ M
}
{ = }{ B  \left(  0 , 1 \right) \setminus \{(0,0)\}
}
{ \subset }{ \R^2
}
{    }{
}
{   }{
}
}
{}{}{}
dengan batas

\mavergleichskette
{\vergleichskette
{ \partial M
}
{ = }{ S^1
}
{  }{
}
{    }{
}
{   }{
}
}
{}{}{.}
Kita tinjau kedua
$1$-\definitionsverweis {bentuk diferensial}{}{}
yang
\definitionsverweis {terdiferensialkan secara kontinu}{}{}

\mathdisp {\omega= ydx-xdy \text{ dan } \tau= { \frac{   1  }{  x^2+y^2 } } \left(  ydx-xdy  \right)} {  }

pada $M$. Tunjukkan bahwa pembatasan kedua bentuk tersebut pada batas saling berimpit, dan khususnya

\mavergleichskette
{\vergleichskette
{ \int_{\partial M} \omega
}
{ =} { \int_{\partial M} \tau
}
{  }{
}
{    }{
}
{   }{
}
}
{}{}{}
berlaku. Bandingkan
\mathkor {} {\int_M d \omega} {dan} {\int_M d \tau} {.}

}
{} {}

\inputaufgabe
{}
{

Yakinkan diri bahwa
teorema Green
tidak menyatakan bahwa luas suatu daerah berbatas di $\R^2$ hanya bergantung pada panjang batasnya.

}
{} {}

\inputaufgabe
{}
{

Misalkan $D$ segitiga dengan titik sudut
$(0,2),\, (1,-1)$, dan $(-2,-1)$,
serta

\mavergleichskette
{\vergleichskette
{ \tau
}
{ = }{ x^2ydx \wedge dy
}
{  }{
}
{    }{
}
{   }{
}
}
{}{}{}
suatu
$2$-\definitionsverweis {bentuk diferensial}{}{} pada $D$. Temukan suatu
\definitionsverweis {bentuk primitif}{}{}
untuk $\tau$, lalu gunakan bentuk tersebut untuk menghitung
\mathl{\int_{ D  }  \tau}{} melalui suatu integral pada batas segitiga.

}
{} {}

\inputaufgabe
{}
{

Verifikasikan
teorema Green
melalui perhitungan eksplisit untuk persegi satuan

\mavergleichskette
{\vergleichskette
{ T
}
{ = }{ [0,1] \times [0,1]
}
{  }{
}
{    }{
}
{   }{
}
}
{}{}{}
dan
\definitionsverweis {bentuk-bentuk diferensial}{}{}

\mavergleichskettedisp
{\vergleichskette
{ \omega
}
{ =} { x^ay^b dx + x^cy^d dy
}
{ } {
}
{ } {
}
{ } {
}
}
{}{}{}
dengan

\mavergleichskette
{\vergleichskette
{ a,b,c,d
}
{ \in }{ \N
}
{  }{
}
{    }{
}
{   }{
}
}
{}{}{.}

}
{} {}

\inputaufgabe
{}
{

Verifikasikan
teorema Green
melalui perhitungan eksplisit untuk himpunan
\mathl{T=[-2,2] \times [-2,2] \setminus U \left(   0 , 1  \right)}{}
\zusatzklammer {yakni persegi berpusat di asal dengan panjang sisi $4$, tanpa lingkaran satuan terbuka} {} {}
dan
\definitionsverweis {bentuk diferensial}{}{}

\mavergleichskette
{\vergleichskette
{ \omega
}
{ = }{(x-y)dx + xy dy
}
{  }{
}
{    }{
}
{   }{
}
}
{}{}{.}

}
{} {}

\inputaufgabegibtloesung
{}
{

Buktikan teorema Green untuk segitiga $D$ dengan titik-titik sudut

\mavergleichskette
{\vergleichskette
{ P_1,P_2,P_3
}
{ \in }{ \R^2
}
{  }{
}
{    }{
}
{   }{
}
}
{}{}{}
dan untuk bentuk diferensial
\mathl{xdy}{.}

}
{} {}

\inputaufgabe
{}
{

Misalkan
\maabbdisp {f} { [a,b] } { \R_{+}
} {}
suatu
\definitionsverweis {fungsi yang terdiferensialkan secara kontinu}{}{.}
Kita pandang
\definitionsverweis {subgraf}{}{}
sebagai suatu
\definitionsverweis {manifold berbatas}{}{},
dengan batas yang terdiri atas grafik, interval dasar, dan kedua sisi tegak, tetapi tanpa keempat titik sudutnya. Hitung luas subgraf dengan menggunakan kedua
\definitionsverweis {bentuk diferensial}{}{}

\mathl{xdy}{} dan
\mathl{ydx}{}
pada batas.

}
{} {}

\inputaufgabe
{}
{

Tentukan volume
\definitionsverweis {bola satuan tertutup}{}{}
tiga dimensi melalui suatu integral permukaan yang sesuai pada
\definitionsverweis {permukaan bola satuan}{}{}
$S^2$.

}
{} {}

\inputaufgabegibtloesung
{}
{

Kita tinjau
$2$-\definitionsverweis {bentuk diferensial}{}{}

\mavergleichskettedisp
{\vergleichskette
{ \omega
}
{ =} { x dy \wedge dz -y dx \wedge dz + z dx \wedge dy
}
{ } {
}
{ } {
}
{ } {
}
}
{}{}{}
pada bola satuan
\mathl{B  \left(  0 , 1  \right)}{.}
\aufzaehlungdreiabc{Tunjukkan bahwa $d \omega$ adalah tiga kali bentuk volume standar pada bola tersebut.
}{Tunjukkan bahwa
\mathl{\omega\vert_{S^2}}{} adalah bentuk luas standar pada permukaan bola satuan.
}{Hitung luas permukaan bola dari
volume bola
dengan menggunakan
teorema Stokes.
}

}
{} {}

\inputaufgabegibtloesung
{}
{

Misalkan $\omega$ suatu
$n-1$-\definitionsverweis {bentuk diferensial}{}{}
yang terdiferensialkan secara kontinu pada suatu
\definitionsverweis {himpunan terbuka}{}{}

\mavergleichskette
{\vergleichskette
{ U
}
{ \subseteq }{ \R^n
}
{  }{
}
{    }{
}
{   }{
}
}
{}{}{},
dan misalkan
\definitionsverweis {turunan eksterior}{}{}
diberikan oleh

\mavergleichskettedisp
{\vergleichskette
{d \omega
}
{ =} { fdx_1 \wedge \ldots \wedge dx_n
}
{ } {
}
{ } {
}
{ } {
}
}
{}{}{}
dengan suatu fungsi
\maabb {f} { U } {\R
} {.}
Tunjukkan untuk

\mavergleichskette
{\vergleichskette
{ P
}
{ \in }{ U
}
{  }{
}
{    }{
}
{   }{
}
}
{}{}{}
kesamaan

\mavergleichskettedisp
{\vergleichskette
{ f(P)
}
{ =} { { \frac{   1  }{  \beta_n  } } \cdot \operatorname{lim}_{   \epsilon \rightarrow  0  } \, { \frac{   1  }{  \epsilon^n } } \int_{S^{n-1}(P, \epsilon)} \omega
}
{ } {
}
{ } {
}
{ } {
}
}
{}{}{,}
dengan $\beta_n$ menyatakan
volume
bola satuan $n$-dimensional.

}
{} {}

Untuk suatu kasus khusus dari pernyataan di atas, lihat
Soal 58.23 (Analisis (Osnabrück 2021--2023)).

\inputaufgabe
{}
{

Misalkan $M$ suatu
\definitionsverweis {manifold diferensiabel}{}{}
dengan batas tak kosong. Tunjukkan bahwa terdapat suatu
\definitionsverweis {pemetaan diferensiabel}{}{}
\maabb {} { M } { \partial M
} {}
ada.

}
{} {}

\inputaufgabe
{}
{

Misalkan

\mavergleichskette
{\vergleichskette
{ H
}
{ \subseteq }{ \R^n
}
{  }{
}
{    }{
}
{   }{
}
}
{}{}{}
\zusatzklammer {
\mavergleichskette
{\vergleichskette
{ n
}
{ \geq }{ 1
}
{  }{
}
{    }{
}
{   }{
}
}
{}{}{}} {} {}
suatu
\definitionsverweis {setengah ruang}{}{.}
Tunjukkan bahwa terdapat suatu
\definitionsverweis {pemetaan diferensiabel}{}{}
\maabbdisp {\varphi} { H } { \partial H
} {}
yang
\definitionsverweis {pembatasannya}{}{}
pada
\mathl{\partial H}{} adalah
\definitionsverweis {identitas}{}{.}

}
{Bagaimana keadaannya untuk

\mavergleichskette
{\vergleichskette
{ n
}
{ = }{ 0
}
{  }{
}
{    }{
}
{   }{
}
}
{}{}{}?} {}

\inputaufgabe
{}
{

Kita tinjau
\definitionsverweis {manifold berbatas}{}{}

\mavergleichskette
{\vergleichskette
{ M
}
{ = }{ \R^n \setminus  U \left(   0 , 1  \right)
}
{  }{
}
{    }{
}
{   }{
}
}
{}{}{.}
Tunjukkan bahwa terdapat suatu pemetaan yang terdiferensialkan secara kontinu dari $M$ ke batas
\mathl{S \left(   0 ,  1  \right)}{}
yang merupakan identitas pada batas tersebut.

}
{} {}

Untuk soal berikut, Anda dapat menggunakan
Soal 21.8.

\inputaufgabe
{}
{

Kita tinjau
\definitionsverweis {manifold berbatas}{}{}

\mavergleichskettedisp
{\vergleichskette
{ M
}
{ =} { B  \left(  0 , 1 \right) \setminus \{  \left(  1   , \,   0                   \right) \}
}
{ \subseteq} { \R^2
}
{ } {
}
{ } {
}
}
{}{}{.}
Tunjukkan bahwa terdapat suatu
\definitionsverweis {pemetaan yang terdiferensialkan secara kontinu}{}{}
\maabbdisp {\varphi} { M } { \partial M
} {}
yang merupakan identitas pada batas.

}
{} {}

\inputaufgabe
{}
{

Tunjukkan bahwa pada suatu
\definitionsverweis {anulus}{}{}
terdapat pemetaan
\definitionsverweis {bijektif}{}{} yang
\definitionsverweis {terdiferensialkan secara kontinu}{}{}
dan tidak mempunyai
\definitionsverweis {titik tetap}{}{.}

}
{} {}

\inputaufgabe
{}
{

Tunjukkan bahwa pemetaan

\mavergleichskettedisp
{\vergleichskette
{ \varphi(x)
}
{ =} { x + \left( - \left\langle  x  , h(x)  \right\rangle + \sqrt{1 + \left\langle  x  , h(x)  \right\rangle^2 - \Vert { x} \Vert^2  }\right) \cdot h(x)
}
{ } {
}
{ } {
}
{ } {
}
}
{}{}{}
\zusatzklammer {dengan

\mavergleichskette
{\vergleichskette
{ h(x)
}
{ = }{ { \frac{   \psi(x)-x }{  \Vert { \psi (x)-x} \Vert  } }
}
{  }{
}
{    }{
}
{   }{
}
}
{}{}{}} {} {}
yang digunakan dalam pembuktian
teorema titik tetap Brouwer
memetakan bola satuan ke permukaan bola satuan.

}
{} {}

  \zwischenueberschrift{Tugas untuk dikumpulkan}

\inputaufgabe
{4}
{

Misalkan $D$ segitiga dengan titik sudut
$(0,2),\, (1,-1)$, dan $(-2,-1)$,
serta
\mathl{\tau =( 3x^2y^5-x \sin  y ) dx \wedge dy}{} suatu
$2$-\definitionsverweis {bentuk diferensial}{}{} pada $D$. Temukan suatu
\definitionsverweis {bentuk primitif}{}{} untuk $\tau$, lalu gunakan bentuk tersebut untuk menghitung
\mathl{\int_{ D  }  \tau}{} melalui suatu integral pada batas segitiga.

}
{} {}

\inputaufgabe
{6}
{

Kita tinjau kubus

\mavergleichskettedisp
{\vergleichskette
{ Q
}
{ =} { [-1,1]^3
}
{ \subseteq} { \R^3
}
{ } {
}
{ } {
}
}
{}{}{}
dan
$2$-\definitionsverweis {bentuk diferensial}{}{}

\mavergleichskettedisp
{\vergleichskette
{ \omega
}
{ =} { dx \wedge dy +y dx \wedge dz +x^2y^2z^2 dy \wedge dz
}
{ } {
}
{ } {
}
{ } {
}
}
{}{}{.}
Hitung
\mathl{d \omega}{} dan kedua integral
\mathkor {} {\int_{  \partial Q  }   \omega} {dan} {\int_{  Q  }  d \omega} {}
\zusatzklammer {secara terpisah} {} {.}

}
{} {}

\inputaufgabe
{4}
{

Hitung luas
\definitionsverweis {cakram satuan tertutup}{}{}
melalui suatu
\definitionsverweis {integral lintasan}{}{}
yang sesuai.

}
{} {}

\inputaufgabe
{2}
{

Tunjukkan bahwa pada suatu
\definitionsverweis {torus}{}{}
terdapat pemetaan
\definitionsverweis {bijektif}{}{} yang
\definitionsverweis {terdiferensialkan secara kontinu}{}{}
dan tidak mempunyai
\definitionsverweis {titik tetap}{}{.}

}
{} {}

\inputaufgabe
{3}
{

Misalkan $B$ adalah
\definitionsverweis {cakram satuan tertutup}{}{}
dan $K$ busur setengah lingkaran atas. Tunjukkan bahwa terdapat suatu
\definitionsverweis {pemetaan diferensiabel}{}{}
\maabbdisp {\varphi} { B } { K
} {}
yang
\definitionsverweis {pembatasannya}{}{}
pada
\mathl{K}{} adalah
\definitionsverweis {identitas}{}{.}

}
{} {}

\inputaufgabe
{3}
{

Misalkan

\mavergleichskette
{\vergleichskette
{ v,w
}
{ \in }{ \R^n
}
{  }{
}
{    }{
}
{   }{
}
}
{}{}{}
dengan
\mathkor {} {\Vert { v } \Vert \leq 1} {dan} {\Vert { w } \Vert = 1} {.}
Tunjukkan bahwa terdapat

\mavergleichskette
{\vergleichskette
{ a
}
{ \in }{ \R
}
{  }{
}
{    }{
}
{   }{
}
}
{}{}{}
sedemikian sehingga

\mavergleichskette
{\vergleichskette
{ \Vert {v+aw} \Vert
}
{ = }{ 1
}
{  }{
}
{    }{
}
{   }{
}
}
{}{}{.}

}
{} {}

\inputaufgabe
{4}
{

Kita tinjau
\definitionsverweis {manifold berbatas}{}{}

\mavergleichskettedisp
{\vergleichskette
{ M
}
{ =} { B  \left(  0 , 1 \right) \setminus \{  \left(  1   , \,   0                   \right) \}
}
{ \subseteq} { \R^2
}
{ } {
}
{ } {
}
}
{}{}{.}
Tunjukkan bahwa terdapat suatu pemetaan bebas titik tetap
\definitionsverweis {yang terdiferensialkan secara kontinu}{}{}
\maabbdisp {\varphi} { M } { M
} {}
ada.

}
{} {}

\section*{Solusi yang disediakan oleh sumber}
\addcontentsline{toc}{section}{Solusi yang disediakan oleh sumber}
\subsection*{Solusi Soal 23.6}
Karena menurut
Fakta (5)
penarikan balik bentuk kompatibel dengan turunan eksterior, maka
\mathl{\varphi^*\omega}{} juga eksak. Misalkan $\tau$ suatu bentuk diferensiabel pada $S$ dengan

\mavergleichskettedisp
{\vergleichskette
{d \tau
}
{ =} { \varphi^* \omega
}
{ } {
}
{ } {
}
{ } {
}
}
{}{}{.}
Menurut
teorema Stokes,

\mavergleichskettedisp
{\vergleichskette
{ \int_S \varphi^* \omega
}
{ =} { \int_S d \tau
}
{ =} { \int_{ \partial S } \tau
}
{ =} { \int_\emptyset \tau
}
{ =} { 0
}
}
{}{}{.}
\subsection*{Solusi Soal 23.13}
Misalkan

\mathdisp {P_1 = \begin{pmatrix} a_1  \\b_1   \end{pmatrix} ,\, P_2 = \begin{pmatrix} a_2  \\b_2   \end{pmatrix} \text{ dan } P_3 = \begin{pmatrix} a_3  \\b_3   \end{pmatrix}} {  }

adalah ketiga titik sudut. Karena
\mathl{d (xdy )= dx \wedge dy}{}, integral bentuk luas ini pada $D$ sama dengan luas segitiga tersebut. Dari $P_1$, segitiga itu direntang oleh kedua vektor
\mathkor {} {\begin{pmatrix} a_2-a_1  \\b_2-b_1   \end{pmatrix}} {dan} {\begin{pmatrix} a_3-a_1  \\b_3-b_1   \end{pmatrix}} {.}
Jadi, menurut rumus determinan untuk sebuah paralelotop, luasnya sama dengan

\mavergleichskettedisp
{\vergleichskette
{ { \frac{   1 }{  2 } } \betrag { (a_2-a_1)(b_3-b_1) -(b_2-b_1)  (a_3-a_1)   }
}
{ =} { { \frac{   1 }{  2 } } \betrag {  a_2b_3-a_2b_1 -a_1b_3 -a_3b_2 +a_1b_2 +a_3b_1   }
}
{ } {
}
{ } {
}
{ } {
}
}
{}{}{.}
Jika kedua vektor tersebut merepresentasikan orientasi standar, sebagaimana akan kita asumsikan mulai sekarang, tanda nilai mutlak dapat dihilangkan.

Sekarang kita hitung integral lintasan dari
\mathl{xdy}{} di sepanjang batas segitiga yang ditelusuri berlawanan arah jarum jam. Lintasan bergerak dari $P_1$ ke $P_2$, lalu ke $P_3$, dan kembali ke $P_1$
\zusatzklammer {ini sesuai dengan arah berlawanan jarum jam untuk orientasi yang telah ditetapkan} {} {.}
Lintasan-lintasan linear ini adalah
\zusatzklammer {masing-masing didefinisikan pada interval satuan} {} {}

\mathdisp {\gamma_1(t)=P_1+t(P_2-P_1) = \begin{pmatrix} a_1  \\b_1   \end{pmatrix} + t \begin{pmatrix} a_2-a_1  \\b_2-b_1   \end{pmatrix}} { , }

\mathdisp {\gamma_2(t) = P_2+t(P_3-P_2) = \begin{pmatrix} a_2  \\b_2   \end{pmatrix} + t \begin{pmatrix} a_3-a_2  \\b_3-b_2   \end{pmatrix}} {  }

dan

\mathdisp {\gamma_3(t) = P_3+t(P_1-P_3)= \begin{pmatrix} a_3  \\b_3   \end{pmatrix} + t \begin{pmatrix} a_1-a_3  \\b_1-b_3   \end{pmatrix}} { . }

Berlaku

\mavergleichskettealign
{\vergleichskettealign
{ \int_{\gamma_1}  xdy
}
{ =} { \int_0^1  (a_1 +t (a_2-a_1)) (b_2-b_1) dt
}
{ =} { (  a_1 (b_2-b_1)t + { \frac{   1 }{  2 } } (a_2-a_1)(b_2-b_1) t^2) \vert_{  0 }^{  1 }
}
{ =} { a_1 (b_2-b_1)  + { \frac{   1 }{  2 } } (a_2-a_1)(b_2-b_1)
}
{ =} { { \frac{   1 }{  2 } } (a_2+a_1)(b_2-b_1)
}
}
{}
{}{.}
Demikian pula,

\mavergleichskettedisp
{\vergleichskette
{ \int_{\gamma_2}  xdy
}
{ =} { { \frac{   1 }{  2 } } (a_3+a_2)(b_3-b_2)
}
{ } {
}
{ } {
}
{ } {
}
}
{}{}{}
dan

\mavergleichskettedisp
{\vergleichskette
{ \int_{\gamma_3}  xdy
}
{ =} { { \frac{   1 }{  2 } } (a_1+a_3)(b_1-b_3)
}
{ } {
}
{ } {
}
{ } {
}
}
{}{}{.}
Jumlah ketiga integral lintasan ini adalah separuh dari

\mavergleichskettealigndrucklinks
{\vergleichskettealigndrucklinks
{(a_2+a_1)(b_2-b_1) + (a_3+a_2)(b_3-b_2) + (a_1+a_3)(b_1-b_3)
}
{ =} { a_2b_2 -a_2b_1+a_1b_2-a_1b_1 +a_3b_3-a_3b_2 +a_2b_3 -a_2b_2+a_1b_1-a_1b_3+a_3b_1-a_3b_3
}
{ =} { -a_2b_1+a_1b_2 -a_3b_2 +a_2b_3 -a_1b_3+a_3b_1
}
{ } {
}
{ } {
}
}
{}{}{,}
sehingga kedua integral itu berimpit.
\subsection*{Solusi Soal 23.16}
a) Berlaku

\mavergleichskettealign
{\vergleichskettealign
{d \omega
}
{ =} { d \left(  x dy \wedge dz  -y dx \wedge dz +  z dx \wedge dy \right)
}
{ =} { dx \wedge dy \wedge dz - dy \wedge dx \wedge dz + d z \wedge dx \wedge dy
}
{ =} { dx \wedge dy \wedge dz + dx \wedge dy \wedge dz + d x \wedge dy \wedge dz
}
{ =} { 3 d x \wedge dy \wedge dz
}
}
{}
{}{,}
di sini kita telah menggunakan fakta bahwa pertukaran faktor-faktor dalam hasil kali baji mengubah tandanya.

b) Untuk suatu titik
\mathl{P = \begin{pmatrix}  x  \\ y \\  z   \end{pmatrix}  \in S^2}{} dan dua vektor singgung ortonormal yang
\zusatzklammer {bersama $P$} {} {}
merepresentasikan orientasi, yaitu
\mathl{u= \begin{pmatrix} u_1  \\u_2 \\ u_3   \end{pmatrix}}{} dan
\mathl{v= \begin{pmatrix} v_1  \\v_2 \\ v_3   \end{pmatrix}}{}, menurut
fakta tentang hasil kali baji,

\mavergleichskettealign
{\vergleichskettealign
{ \omega \vert_{S^2}  (u,v)
}
{ =} { \omega (u,v)
}
{ =} { z \det  \begin{pmatrix} u_1   &  v_1   \\ u_2  & v_2  \end{pmatrix} - y \det  \begin{pmatrix} u_1   &  v_1   \\ u_3  & v_3  \end{pmatrix} + x \det  \begin{pmatrix} u_2   &  v_2   \\ u_3  & v_3  \end{pmatrix}
}
{ =} { \det  \begin{pmatrix}  x   &  u_1  & v_1  \\  y  & u_2  & v_2  \\ z  & u_3  & v_3  \end{pmatrix}
}
{ } {
}
}
{}
{}{.}
Namun, determinan suatu sistem ortonormal yang merepresentasikan orientasi adalah $1$. Jadi, bentuk diferensial yang dibatasi tersebut memenuhi sifat yang mencirikan
\definitionsverweis {bentuk volume standar}{}{}.

c) Dengan menggunakan a), b), teorema Stokes, dan volume bola, luas permukaan bola adalah

\mavergleichskettedisp
{\vergleichskette
{ \int_{S^2} \omega
}
{ =} { \int_{ B  \left(  0, 1 \right)  }  d \omega
}
{ =} { \int_{ B  \left(  0, 1 \right)  } 3 d x \wedge dy \wedge dz
}
{ =} { 3 \lambda^3 (   B  \left(  0, 1 \right)     )
}
{ =} { 4 \pi
}
}
{}{}{.}
\subsection*{Solusi Soal 23.17}
Kita terapkan
teorema Stokes
pada bola-bola
\mathl{B^n(P, \epsilon)}{} dengan batas
\mathl{S^{n-1}(P, \epsilon)}{}
\zusatzklammer {berpusat di $P$ dan berjari-jari $\epsilon$} {} {},
sehingga diperoleh

\mavergleichskettealign
{\vergleichskettealign
{ \operatorname{lim}_{   \epsilon \rightarrow  0  } \, { \frac{   1 }{  \epsilon^n } }     \int_{S^{n-1}(P, \epsilon)} \omega
}
{ =} { \operatorname{lim}_{   \epsilon \rightarrow  0  } \, { \frac{   1 }{  \epsilon^n } }     \int_{ \partial B^{n}(P, \epsilon)}   \omega
}
{ =} { \operatorname{lim}_{   \epsilon \rightarrow  0  } \, { \frac{   1 }{  \epsilon^n } }     \int_{B^{n}(P, \epsilon)} d \omega
}
{ =} { \operatorname{lim}_{   \epsilon \rightarrow  0  } \, { \frac{   1 }{  \epsilon^n } }     \int_{B^{n}(P, \epsilon)} f d \lambda^n
}
{ =} {\beta_n \cdot f(P)
}
}
{}
{}{,}
di mana pada kesamaan terakhir kita menggunakan kekontinuan $f$.

\part{Unit 24}
\chapter[Kuliah 24]{Kuliah 24: Koneksi dan Turunan Vertikal}
\setcounter{section}{24}

Misalkan

\mavergleichskette
{\vergleichskette
{ Y
}
{ \subseteq }{ \R^n
}
{  }{
}
{    }{
}
{   }{
}
}
{}{}{ }
adalah suatu hipermuka terdiferensialkan. Bagi setiap titik

\mavergleichskette
{\vergleichskette
{ P
}
{ \in }{ Y
}
{  }{
}
{    }{
}
{   }{
}
}
{}{}{,}
\definitionsverweis {ruang singgung}{}{}
merupakan subruang vektor

\mavergleichskette
{\vergleichskette
{ T_PY
}
{ \subseteq }{ \R^n
}
{  }{
}
{    }{
}
{   }{
}
}
{}{}{}
di ruang ambien. Dengan menggunakan hasil kali dalam standar pada $\R^n$, kita memperoleh
\definitionsverweis {proyeksi ortogonal}{}{}
\maabbdisp {\pi_P} {\R^n } { T_PY
} {,}
yang memungkinkan data di ruang ambien ditafsirkan sebagai data tangensial pada hipermuka. Dari konstruksi ini muncul berbagai konsep penting pada $Y$, seperti
\definitionsverweis {percepatan tangensial}{}{,}
\definitionsverweis {kurva geodesik}{}{,}
\definitionsverweis {kelengkungan}{}{,}
\definitionsverweis {turunan kovarian suatu medan vektor}{}{,}
\definitionsverweis {medan vektor paralel}{}{}
serta
\definitionsverweis {transpor paralel}{}{.}
Di sisi lain, hasil kali dalam standar tersebut juga menginduksi
\definitionsverweis {struktur Riemann}{}{}
pada $Y$. Pertanyaannya kemudian ialah apakah konsep-konsep tadi dapat dibangun secara intrinsik hanya dari $Y$ dan struktur Riemann-nya
\zusatzklammer {misalnya untuk
\definitionsverweis {permukaan hiperbolik}{}{}} {} {.}
Untuk menjawabnya, diperlukan perangkat baru, yaitu koneksi pada bundel vektor, terutama pada bundel tangen suatu
\definitionsverweis {manifold diferensiabel}{}{.}
Pada manifold Riemann, struktur ini menghasilkan secara kanonik \stichwort {koneksi Levi-Civita} {}
\zusatzklammer {lihat
Teorema 26.12} {} {,}
yang memungkinkan seluruh konsep di atas dirumuskan secara intrinsik.

  \zwischenueberschrift{Koneksi}

Misalkan $X$ adalah
\definitionsverweis {manifold diferensiabel}{}{}
riil, dan
\maabbdisp {p} { E } { X
} {}
adalah bundel vektor di atas $X$ yang ruang totalnya juga diperlakukan sebagai manifold diferensiabel. Adapun
\definitionsverweis {pemetaan singgung}{}{}
\maabb {T(p)} {TE} {TX
} {,}
berada dalam diagram komutatif

\mathdisp {\begin{matrix} TE & \stackrel{ T(p) }{\longrightarrow} & TX &  \\ \downarrow & & \downarrow  & \\  E  & \stackrel{ p }{\longrightarrow} &  X  & \! \! \! \!\!\!\!\!   \\ \end{matrix}} {  }

dan, bagi setiap titik

\mavergleichskette
{\vergleichskette
{ e
}
{ \in }{ E
}
{  }{
}
{    }{
}
{   }{
}
}
{}{}{}
dengan titik dasar

\mavergleichskette
{\vergleichskette
{ x
}
{ = }{ p(e)
}
{  }{
}
{    }{
}
{   }{
}
}
{}{}{,}
\definitionsverweis {pemetaan singgung}{}{}
\definitionsverweis {linear}{}{}
\maabbdisp {T_e(p)} {T_e E} { T_{x}X
} {}
bersifat surjektif. Salah satu alasannya ialah bahwa secara lokal, melalui setiap titik

\mavergleichskette
{\vergleichskette
{ e
}
{ \in }{ E
}
{  }{
}
{    }{
}
{   }{
}
}
{}{}{,}
terdapat suatu penampang bundel vektor $E$. Agar domain dan kodomain memiliki ruang dasar yang sama, kita memandang $T(p)$ sebagai
\definitionsverweis {pemetaan bundel}{}{}
\maabbdisp {Tp} {TE} { p^* TX
} {}
antara bundel vektor di atas $E$. Dalam notasi ini,

\mavergleichskettedisp
{\vergleichskette
{ p^*TX
}
{ =} { E \oplus TX
}
{ =} { E \times_X TX
}
{ } {
}
{ } {
}
}
{}{}{}
merupakan
\definitionsverweis {jumlah langsung}{}{}
bundel vektor di atas $X$; lihat
Soal 12.24.
Di sini, $p^*TX$ dipandang sebagai bundel vektor di atas $E$. Serat $p^*TX$ di atas titik

\mavergleichskette
{\vergleichskette
{ e
}
{ \in }{ E
}
{  }{
}
{    }{
}
{   }{
}
}
{}{}{}
dengan titik dasar

\mavergleichskette
{\vergleichskette
{ x
}
{ = }{ p(e)
}
{  }{
}
{    }{
}
{   }{
}
}
{}{}{}
hanyalah $T_{x}X$, sedangkan pemetaan pada serat tersebut tetaplah
\maabbdisp {T_e (p)} { T_e E } { T_{x} X
} {,}
sehingga di dalam $p^*TX$ terdapat satu salinan $T_{x} X$ untuk setiap

\mavergleichskette
{\vergleichskette
{ e
}
{ \in }{ E
}
{  }{
}
{    }{
}
{   }{
}
}
{}{}{.}

\inputbemerkung
{}
{

Misalkan
\maabb {p} { E } { X
} {}
adalah
\definitionsverweis {bundel vektor}{}{}
diferensiabel di atas
\definitionsverweis {manifold diferensiabel}{}{}
$X$. Ambil himpunan terbuka

\mavergleichskette
{\vergleichskette
{ U
}
{ \subseteq }{ X
}
{  }{
}
{    }{
}
{   }{
}
}
{}{}{,}
tempat $E$ dapat ditrivialkan, sehingga

\mavergleichskette
{\vergleichskette
{ E  \vert_U
}
{ \cong }{ U \times W
}
{  }{
}
{    }{
}
{   }{
}
}
{}{}{}
dengan ruang vektor atas $\R$ yang kita namakan $W$. Dengan demikian,

\mavergleichskettedisp
{\vergleichskette
{ T  \left(  E \vert_U  \right)
}
{ =} { T(U \times  W )
}
{ =} { T U \times T W
}
{ =} { TU \times W \times W
}
{ } {
}
}
{}{}{,}
Sementara itu,
\definitionsverweis {pemetaan singgung}{}{}
yang berkaitan dengan
\maabb {p_U} {E \vert_U } { U
} {}
merupakan proyeksi ke $TU$. Untuk menuliskannya sebagai homomorfisme bundel di atas $U\times W$, kodomainnya diidentifikasi sebagai

\mavergleichskettedisp
{\vergleichskette
{ p^*TU
}
{ =} { TU \times W
}
{ =} { TU \times_U (U \times W)
}
{ } {
}
{ } {
}
}
{}{}{}
dan pemetaannya ialah
\maabbdisp {} {T  \left(  E \vert_U  \right)  = TU \times W \times W  } { TU \times W
} {}
yang mengambil proyeksi ke $TU$ sekaligus proyeksi kedua ke $W$.

}

\inputdefinition
{}
{

Misalkan
\maabb {p} { E } { X
} {}
adalah
\definitionsverweis {bundel vektor diferensiabel}{}{}
di atas
\definitionsverweis {manifold diferensiabel}{}{}
$X$.
\definitionsverweis {Bundel kernel}{}{}
dari
\definitionsverweis {homomorfisme bundel}{}{}
\definitionsverweis {surjektif}{}{}
\maabbdisp {T(p)} {TE} {p^*TX
} {}
di atas $E$ disebut \definitionswort {bundel vertikal}{.} Notasinya adalah $V$.

}

Menurut Soal 12.19, bundel vertikal adalah bundel vektor di atas $E$; lebih khusus lagi, ia merupakan
\definitionsverweis {subbundel vektor}{}{}
dari $TE$.



\inputfaktbeweis
{Mannigfaltigkeit/Vektorbündel/R/Vertikalbündel/Eigenschaften/Fakt}
{Lema}
{}
{

\faktsituation {Misalkan
\maabb {p} { E } { X
} {}
adalah
\definitionsverweis {bundel vektor diferensiabel}{}{}
di atas
$C^2$-\definitionsverweis {manifold diferensiabel}{}{}
$X$.}
\faktuebergang {Pernyataan berikut berlaku.}
\faktfolgerung {\aufzaehlungdrei{Terdapat
\definitionsverweis {barisan eksak pendek}{}{}

\mathdisp {0 \, \stackrel{  }{\longrightarrow} \,  V   \, \stackrel{  }{ \longrightarrow}   \,  TE  \, \stackrel{  }{\longrightarrow} \, p^* TX \,\stackrel{  } {\longrightarrow}  \, 0} {  }

yang terdiri dari
\definitionsverweis {homomorfisme}{}{}
bundel vektor di atas $E$.
}{Untuk
\definitionsverweis {bundel vertikal}{}{}
terdapat
\definitionsverweis {isomorfisme bundel}{}{}

\mavergleichskette
{\vergleichskette
{ V \cong p^*E
}
{ = }{ E \oplus E
}
{  }{
}
{    }{
}
{   }{
}
}
{}{}{.}
}{Bagi setiap titik

\mavergleichskette
{\vergleichskette
{ e
}
{ \in }{ E
}
{  }{
}
{    }{
}
{   }{
}
}
{}{}{,}
diperoleh barisan eksak pendek ruang vektor berikut:

\mathdisp {0 \, \stackrel{  }{\longrightarrow} \, E_{p(e)}   \, \stackrel{  }{ \longrightarrow}   \,  T_eE  \, \stackrel{  }{\longrightarrow} \, T_{p(e)}X \,\stackrel{  } {\longrightarrow}  \, 0} {  }

}}
\faktzusatz {}
\faktzusatz {}

}
{

\aufzaehlungzwei {Pernyataan ini langsung mengikuti definisi.
} {Terdapat homomorfisme injektif bundel vektor
\maabbdisp {} {p^*E = E \oplus E} { TE
} {}
di atas $E$ yang memetakan
\mathl{(P,u)}{,}

\mavergleichskette
{\vergleichskette
{ u
}
{ \in }{ E_P
}
{  }{
}
{    }{
}
{   }{
}
}
{}{}{,}
ke vektor singgung yang diwakili oleh kurva diferensiabel
\zusatzklammer {vertikal} {} {}
\maabbeledisp {} {\R} {E
} { t } {(P,tu)
} {,}
ini. Selanjutnya, oleh pemetaan
\maabbdisp {} {TE} { p^*TX
} {}
vektor singgung tersebut dikirim ke $0$, sebab kurva wakilnya dikirim ke $P$ oleh
\maabb {p} { E } { X
} {}
.
Oleh karena itu,

\mavergleichskette
{\vergleichskette
{ E
}
{ \subseteq }{ V
}
{  }{
}
{    }{
}
{   }{
}
}
{}{}{,}
dan perbandingan rank menunjukkan bahwa inklusi ini sebenarnya suatu isomorfisme.
}

}

Kasus yang sering dijumpai adalah ketika $E$ merupakan bundel tangen dari $X$. Dalam kasus ini diperoleh barisan eksak pendek

\mathdisp {0 \, \stackrel{  }{\longrightarrow} \, p^* TX  \, \stackrel{  }{ \longrightarrow}   \, TTX  \, \stackrel{  }{\longrightarrow} \, p^* TX \,\stackrel{  } {\longrightarrow}  \, 0} {  }

Oleh sebab itu, kedua salinan bundel tangen di ruas kiri dan kanan harus selalu dibedakan dengan cermat.

\inputbemerkung
{}
{

Misalkan
\maabbdisp {f} {\R^n } { \R
} {}
adalah fungsi yang dua kali
\definitionsverweis {terdiferensialkan secara kontinu}{}{,}
dan

\mavergleichskette
{\vergleichskette
{ Y
}
{ = }{ f^{-1} (0)
}
{  }{
}
{    }{
}
{   }{
}
}
{}{}{}
adalah
\definitionsverweis {hipermuka terdiferensialkan}{}{}
terkait yang
\definitionsverweis {reguler}{}{}
pada setiap titik.
\definitionsverweis {Bundel tangen}{}{}
$TY$ dapat pula dinyatakan sebagai serat bersama dari fungsi-fungsi, yakni

\mavergleichskettealignhandlinks
{\vergleichskettealignhandlinks
{ TY
}
{ =} { \left\{  (P,u)   \mid   f(P) = 0 , \,  \left(  D f   \right)_{ P } \left(  u  \right) = 0       \right\}
}
{ =} { \left\{  (P,u_1 , \ldots , u_n)   \mid   f(P) = 0 , \,  \sum_{i = 1}^n u_i { \frac{   \partial   f   }{  \partial   x_i   } } ( P) = 0       \right\}
}
{ \subseteq} { \R^n \times \R^n
}
{ } {
}
}
{}
{}{,}
sebagaimana dalam
Soal 10.15.
Dalam rumusan ini, $f$ dipandang sebagai fungsi pada $\R^{2n}$ yang hanya bergantung pada $n$ variabel pertama, dan kita definisikan

\mavergleichskettedisp
{\vergleichskette
{ g(P,u)
}
{ =} { \sum_{i = 1}^n u_i { \frac{   \partial   f   }{  \partial   x_i   } } ( P)
}
{ } {
}
{ } {
}
{ } {
}
}
{}{}{.}
Bundel tangen kedua $TTY$, yaitu bundel tangen dari $TY$, selanjutnya dapat dideskripsikan dengan
\definitionsverweis {diferensial total}{}{}
dari pemetaan
\maabbdisp {\varphi = (f,g)} {\R^{2n}} { \R^2
} {}
sebagai berikut:

\mavergleichskettealignhandlinks
{\vergleichskettealignhandlinks
{ TTY
}
{ =} { \left\{  (P,u, v,w)  \mid   (P,u) \in TY , \,  \left(  D\varphi  \right)_{(P,u)} \left(  (v,w)  \right) = 0       \right\}
}
{ \subseteq} { \R^{4n}
}
{ } {
}
{ } {}
}
{}
{}{.}
\definitionsverweis {Matriks Jacobi}{}{,}
yang merepresentasikan diferensial total $D\varphi$, ialah

\mathdisp {\begin{pmatrix}   { \frac{   \partial   f   }{  \partial   x_1   } } ( P)   &   \ldots &  { \frac{   \partial   f   }{  \partial   x_n   } } ( P)  &   0 &  \ldots &  0  \\   \sum_{i = 1}^n u_i   { \frac{   \partial    }{  \partial   x_1   } } { \frac{   \partial   f   }{  \partial   x_i   } } ( P)  &  \ldots  &   \sum_{i = 1}^n u_i   { \frac{   \partial    }{  \partial   x_n   } } { \frac{   \partial   f   }{  \partial   x_i   } } ( P)   &  { \frac{   \partial   f   }{  \partial   x_1   } } ( P)  &  \ldots &  { \frac{   \partial   f   }{  \partial   x_n   } } ( P)    \end{pmatrix}} { . }

Di samping dua syarat yang mendefinisikan $TY$ dan hanya menyangkut
\mathkor {} {P} {dan} {u} {,}
diperoleh dua syarat tambahan yang melibatkan seluruh variabel, yaitu

\mavergleichskettedisp
{\vergleichskette
{ \sum_{j = 1}^n v_j { \frac{   \partial   f   }{  \partial   x_j   } } ( P)
}
{ =} { 0
}
{ } {
}
{ } {
}
{ } {
}
}
{}{}{}
dan

\mavergleichskettedisp
{\vergleichskette
{ \sum_{j = 1}^n v_j  \left(  \sum_{i = 1}^n u_i { \frac{   \partial    }{  \partial   x_i   } }  { \frac{   \partial   f   }{  \partial   x_j   } } ( P)  \right)      + \sum_{k = 1}^n w_k { \frac{   \partial   f   }{  \partial   x_k   } } ( P)
}
{ =} { 0
}
{ } {
}
{ } {
}
{ } {
}
}
{}{}{.}

}

\inputbemerkung
{}
{

Misalkan
\maabbdisp {f} {\R^n } { \R
} {}
adalah fungsi yang dua kali
\definitionsverweis {terdiferensialkan secara kontinu}{}{,}
dan

\mavergleichskette
{\vergleichskette
{ Y
}
{ = }{ f^{-1} (0)
}
{  }{
}
{    }{
}
{   }{
}
}
{}{}{}
adalah
\definitionsverweis {hipermuka terdiferensialkan}{}{}
terkait yang
\definitionsverweis {reguler}{}{}
pada setiap titik. Kita gunakan deskripsi bundel tangen kedua $TTY$ dari
Catatan 24.4
beserta notasi di sana. Bundel tangen tarik balik $\pi^*TY$ terhadap
\maabbdisp {\pi} {TY} {Y
} {,}
yang juga merupakan
\definitionsverweis {bundel vertikal}{}{,}
adalah

\mavergleichskettealign
{\vergleichskettealign
{ V
}
{ \cong} { \pi^*TY
}
{ =} { TY \times_Y TY
}
{ =} { \left\{ (P,u,z)  \mid  f(P) = 0 , \,  \sum_{i = 1}^n u_i { \frac{   \partial   f   }{  \partial   x_i   } } ( P) = 0   , \,  \sum_{k = 1}^n z_k { \frac{   \partial   f   }{  \partial   x_k   } } ( P) = 0    \right\}
}
{ } {
}
}
{}
{}{.}
Sekarang kita tentukan bentuk pemetaan singgung
\maabb {T(\pi)} {TTY } { \pi^* TY
} {}
serta pembenaman bundel vertikal ke dalam $TTY$. Berdasarkan
Soal 24.2,
diperoleh
\maabbeledisp {T(\pi)} {TTY} { \pi^* TY
} {(P,u,v,w)} { (P,u,v)
} {.}
Kernel pemetaan ini adalah

\mavergleichskettedisp
{\vergleichskette
{ V
}
{ =} { \left\{  (P,u,0,w)   \mid   (P,u,0,w) \in TTY          \right\}
}
{ \subset} { TTY
}
{ } {
}
{ } {
}
}
{}{}{,}
dan langsung terlihat bahwa
\mathl{(P,u,w)}{} memenuhi syarat bagi
\mathl{\pi^*TY}{}.
Oleh sebab itu, kedua homomorfisme bundel dalam barisan eksak pendek

\mathdisp {0 \, \stackrel{  }{\longrightarrow} \,  \pi^*TY  \, \stackrel{  }{ \longrightarrow}   \, TTY  \, \stackrel{  }{\longrightarrow} \,  \pi^*TY \,\stackrel{  } {\longrightarrow}  \, 0} {  }

dapat ditulis dalam koordinat sebagai

\mathl{(P,u,w) \mapsto  (P,u,0,w)}{} dan
\mathl{(P,u,v,w) \mapsto (P,u,v)}{}.

}

Jika $E$ merupakan bundel trivial, yaitu

\mavergleichskette
{\vergleichskette
{ E
}
{ = }{ X \times \R^r
}
{  }{
}
{    }{
}
{   }{
}
}
{}{}{,}
maka terdapat penguraian langsung
\zusatzklammer {di atas $E$} {} {}

\mavergleichskette
{\vergleichskette
{ TE
}
{ = }{ V \oplus p^* TX
}
{  }{
}
{    }{
}
{   }{
}
}
{}{}{.}
Di setiap titik

\mavergleichskette
{\vergleichskette
{ e
}
{ \in }{ E
}
{  }{
}
{    }{
}
{   }{
}
}
{}{}{,}
vektor singgung

\mavergleichskette
{\vergleichskette
{ t
}
{ \in }{ V
}
{  }{
}
{    }{
}
{   }{
}
}
{}{}{}
mewakili arah vertikal, sedangkan

\mavergleichskette
{\vergleichskette
{ t
}
{ \in }{ p^* T_xX
}
{  }{
}
{    }{
}
{   }{
}
}
{}{}{}
mewakili \anfuehrung{arah horizontal}{.}
Pada bundel vektor sembarang, suatu trivialisasi lokal menguraikan ruang singgung
\mathl{T_eE}{} sebagai jumlah langsung ruang vertikal dan ruang horizontal. Namun, subruang horizontal yang dihasilkan sangat bergantung pada trivialisasi yang dipilih. Tidak seperti subruang vertikal
\zusatzklammer {yang bersama-sama membentuk bundel vertikal} {} {,}
subruang horizontal ini belum tentu menyatu menjadi suatu subbundel. Keberadaan subbundel horizontal semacam inilah yang dirumuskan oleh konsep koneksi.

\inputdefinition
{}
{

Misalkan $X$ adalah
\definitionsverweis {manifold diferensiabel}{}{}
dan $E$ adalah
\definitionsverweis {bundel vektor diferensiabel}{}{}
di atas $X$. Suatu
\definitionswort {koneksi}{}
pada $E$ didefinisikan sebagai penguraian jumlah langsung
\definitionsverweis {bundel tangen}{}{}

\mavergleichskette
{\vergleichskette
{ TE
}
{ = }{ V \oplus H
}
{  }{
}
{    }{
}
{   }{
}
}
{}{}{}
menjadi dua
\definitionsverweis {subbundel vektor}{}{}
\mathkor {} {V} {dan} {H} {,}
dengan ketentuan bahwa

\mavergleichskette
{\vergleichskette
{ V
}
{ \subseteq }{ TE
}
{  }{
}
{    }{
}
{   }{
}
}
{}{}{}
merupakan
\definitionsverweis {bundel vertikal}{}{.}
Subbundel

\mavergleichskette
{\vergleichskette
{ H
}
{ \subseteq }{ TE
}
{  }{
}
{    }{
}
{   }{
}
}
{}{}{}
disebut
\definitionswort {bundel horizontal}{}
koneksi tersebut.

}

Koneksi biasanya dilambangkan dengan $\nabla$, khususnya ketika proses diferensiasi yang ditentukannya hendak ditekankan. Jadi, menurut definisi, koneksi hanyalah subbundel dari $TE$ yang komplementer terhadap bundel vertikal. Ada beberapa rumusan lain yang ekuivalen. Memberikan suatu koneksi sama artinya dengan memberikan homomorfisme bundel
\zusatzklammer {disebut \stichwort {proyeksi vertikal} {} koneksi} {} {}
\maabbdisp {\pi_{\text{vert} }} {TE} {V
} {}
dengan

\mavergleichskettedisp
{\vergleichskette
{ \pi_{\text{vert} }  \circ \iota
}
{ =} {
\operatorname{Id}_{  V  }
}
{ } {
}
{ } {
}
{ } {
}
}
{}{}{,}
di mana $\iota$ adalah pembenaman $V$ ke dalam $TE$. Dalam rumusan ini, bundel horizontal ialah kernel $\pi_{\text{vert}}$. Sebaliknya, penguraian jumlah langsung menyediakan proyeksi ke masing-masing suku. Penguraian seperti itu juga disebut \stichwort {pemisahan} {} dari barisan eksak pendek. Setiap koneksi dapat dibatasi ke himpunan terbuka

\mavergleichskette
{\vergleichskette
{ U
}
{ \subseteq }{ X
}
{  }{
}
{    }{
}
{   }{
}
}
{}{}{.}

\inputbeispiel{}
{

Pertimbangkan bundel trivial
\maabbdisp {p} {X \times W = W_X } { X
} {,}
dengan $W$ ruang vektor riil
\definitionsverweis {berdimensi}{}{}
$r$ dan $X$ sebuah
\definitionsverweis {manifold diferensiabel}{}{,}
kita memperoleh barisan eksak pendek

\mathdisp {0 \, \stackrel{  }{\longrightarrow} \, p^*(X \times W)  =  X \times W \times W   \, \stackrel{  }{ \longrightarrow}   \, p^* T(X \times W )  = TX \times W \times W   \, \stackrel{  }{\longrightarrow} \, p^* TX  =  TX \times W   \,\stackrel{  } {\longrightarrow}  \, 0} {  }

bundel vektor di atas
\mathl{X \times W}{.}
Semua produk di sini dipahami sebagai produk Kartesius biasa, setelah dilakukan identifikasi

\mavergleichskette
{\vergleichskette
{ (X \times W) \times_X (X \times W)
}
{ = }{ X \times W \times W
}
{  }{
}
{    }{
}
{   }{
}
}
{}{}{}
di sebelah kiri serta identifikasi

\mavergleichskette
{\vergleichskette
{ TX \times_X (X \times W)
}
{ = }{ TX \times W
}
{  }{
}
{    }{
}
{   }{
}
}
{}{}{}
di sebelah kanan. Pada titik

\mavergleichskette
{\vergleichskette
{ (x,w)
}
{ \in }{ X \times W
}
{  }{
}
{    }{
}
{   }{
}
}
{}{}{,}
barisan eksak pada serat berbentuk

\mathdisp {0 \, \stackrel{  }{\longrightarrow} \,  W   \, \stackrel{  }{ \longrightarrow}   \,  T_xX \times W   \, \stackrel{  }{\longrightarrow} \,  T_xX \,\stackrel{  } {\longrightarrow}  \, 0} { . }

Pemisahan baku dari barisan eksak pendek bundel vektor ini disebut \stichwort {koneksi trivial} {} pada bundel trivial tersebut.

}

\inputbemerkung
{}
{

Misalkan
\maabbdisp {f} {\R^n } { \R
} {}
adalah fungsi yang dua kali
\definitionsverweis {terdiferensialkan secara kontinu}{}{,}
dan

\mavergleichskette
{\vergleichskette
{ Y
}
{ = }{ f^{-1} (0)
}
{  }{
}
{    }{
}
{   }{
}
}
{}{}{}
adalah
\definitionsverweis {hipermuka terdiferensialkan}{}{}
terkait yang
\definitionsverweis {reguler}{}{}
pada setiap titik. Kita kembali memakai deskripsi bundel tangen kedua $TTY$ dari
Catatan 24.4
dan homomorfisme bundel dari
Catatan 24.5.
Melalui proyeksi ortogonal
\maabb {} {\R^n } { T_PY
} {}
diperoleh proyeksi vertikal
\maabbeledisp {} {TTY} { V = p^*TY
} {(P,u,v,w)} { (P,u, \pi_P (w) )
} {,}
pada
\definitionsverweis {bundel vertikal}{}{.}
Memang, $\pi_P(w)$ berada dalam ruang singgung $Y$. Apabila kita mulai dari vektor

\mavergleichskette
{\vergleichskette
{ (P,u,w)
}
{ \in }{ V
}
{  }{
}
{    }{
}
{   }{
}
}
{}{}{,}
komponen terakhir $w$ pada
\mathl{(P,u,0,w)}{} sudah otomatis berada dalam ruang singgung. Karena itu, proyeksi ortogonal mengirim unsur ini kembali ke $(P,u,w)$. Jadi, konstruksi tersebut mendefinisikan
\definitionsverweis {koneksi}{}{}
pada $TY$.

}

  \zwischenueberschrift{Turunan vertikal}

\inputdefinition
{}
{

Misalkan $X$ adalah
\definitionsverweis {manifold diferensiabel}{}{}
dan $E$ adalah
\definitionsverweis {bundel vektor diferensiabel}{}{}
di atas $X$ yang dilengkapi dengan suatu
\definitionsverweis {koneksi}{}{.}
\definitionswort {turunan vertikal}{}
$\nabla$ yang berkaitan dengan koneksi tersebut adalah pemetaan
\maabbeledisp {\nabla} {C^1(E)} { C^0(E \otimes T^* X )
} { s } { \nabla(s) = \pi_{\text{vert} }  \circ T(s)
} {.}

}
Dengan kata lain, setiap penampang diferensiabel $s$ dari $E$
\zusatzklammer {di atas $X$ ataupun di atas himpunan terbuka sembarang

\mavergleichskettek
{\vergleichskettek
{ U
}
{ \subseteq }{ X
}
{  }{
}
{    }{
}
{   }{
}
}
{}{}{}} {} {}
dipetakan ke komposisi

\mathdisp {TX \stackrel{T(s)}{ \longrightarrow } TE \stackrel{ \pi_{\text{vert} } }{\longrightarrow} V \cong p^*E \longrightarrow E} {  }

dengan identifikasi

\mavergleichskette
{\vergleichskette
{ \operatorname{Hom}(TX,E)
}
{ \cong }{ E \otimes T^*X
}
{  }{
}
{    }{
}
{   }{
}
}
{}{}{.}

Jika diberikan pula
\definitionsverweis {medan vektor}{}{}
$F$ pada $X$, yaitu penampang
\maabb {F} { X } { TX
} {}
pada
\definitionsverweis {bundel tangen}{}{}
$TX$, maka diperoleh pemetaan
\maabbeledisp {\nabla_F} {C^1(E)} {C^0(E)
} {s } {\nabla_F(s) = \nabla (s) \circ F
} {,}
yang disebut \stichwort {turunan vertikal} {} dalam arah $F$. Ketergantungan terhadap $F$ hanya bersifat titik demi titik. Sebaliknya, ketergantungan terhadap $s$ bersifat infinitesimal karena melibatkan pemetaan singgung
\mathl{T(s)}{}.



\inputfaktbeweis
{Mannigfaltigkeit/Vektorbündel/R/Zusammenhang/Ableitung bezüglich Vektorfeld/Eigenschaften/Fakt}
{Lema}
{}
{

\faktsituation {Misalkan $X$ adalah
\definitionsverweis {manifold diferensiabel}{}{}
dan $E$ adalah
\definitionsverweis {bundel vektor diferensiabel}{}{}
di atas $X$ yang dilengkapi dengan suatu
\definitionsverweis {koneksi}{}{.}
Tuliskan
\maabbdisp {\nabla_F} {C^1(E)} {C^0(E)
} {}
untuk operator turunan yang bersesuaian dengan
\definitionsverweis {medan vektor}{}{}
$F$.}
\faktuebergang {Sifat-sifat berikut berlaku.}
\faktfolgerung {\aufzaehlungzwei {
\mathl{\left(  \nabla_F(s)  \right)  (P)}{} hanya bergantung pada $F(P)$
\zusatzklammer {serta pada $s$} {} {.}
} {Berlaku

\mavergleichskettedisp
{\vergleichskette
{ \nabla_{g_1F_1+g_2F_2 }
}
{ =} { g_1 \nabla_{F_1 } +g_2 \nabla_{F_2 }
}
{ } {
}
{ } {
}
{ } {
}
}
{}{}{}
untuk medan vektor $F_1,F_2$ dan fungsi
\maabb {g_1,g_2} { X} { \R
} {.}
}}
\faktzusatz {}
\faktzusatz {}

}
{

\aufzaehlungzwei {Tetapkan titik

\mavergleichskette
{\vergleichskette
{ P
}
{ \in }{ X
}
{  }{
}
{    }{
}
{   }{
}
}
{}{}{}
dan penampang

\mavergleichskette
{\vergleichskette
{ s
}
{ \in }{ C^1(U,E)
}
{  }{
}
{    }{
}
{   }{
}
}
{}{}{}
pada himpunan terbuka

\mavergleichskette
{\vergleichskette
{ P
}
{ \in }{ U
}
{ \subseteq }{ X
}
{    }{
}
{   }{
}
}
{}{}{,}
maka

\mavergleichskettedisp
{\vergleichskette
{ \nabla_F (s)(P)
}
{ =} { ( \pi_{\text{vert} } \circ  T_P(s) ) (F(P))
}
{ } {
}
{ } {
}
{ } {
}
}
{}{}{.}
} {Pemetaan singgung dan proyeksi vertikal keduanya linear. Menurut (1), ketergantungan pada $F$ juga linear dan selaras dengan perkalian oleh fungsi.
}

}



\inputfaktbeweis
{Differenzierbare Hyperfläche/Zweites Tangentialbündel/Induzierter Zusammenhang/Vertikale Ableitung/Fakt}
{Lema}
{}
{

\faktsituation {Misalkan

\mavergleichskette
{\vergleichskette
{ Y
}
{ \subseteq }{ W
}
{ \subseteq }{ \R^n
}
{    }{
}
{   }{
}
}
{}{}{,}
dengan $W$
\definitionsverweis {terbuka}{}{,}
dan himpunan pertama adalah
\definitionsverweis {hipermuka terdiferensialkan}{}{.}
Misalkan pula
\maabbdisp {F} { Y} { \R^n
} {}
adalah medan vektor tangensial.}
\faktfolgerung {Maka
\definitionsverweis {turunan vertikal}{}{}
$\nabla_F$ yang ditentukan oleh
\definitionsverweis {koneksi}{}{}
pada
Catatan 24.8
berimpit dengan
\definitionsverweis {turunan kovarian}{}{.}}
\faktzusatz {}
\faktzusatz {}

}
{

Untuk penampang diferensiabel
\maabb {s} { Y } { TY
} {}
turunan vertikal koneksi tersebut searah $F$ diberikan oleh komposisi

\mathdisp {Y \stackrel{F}{\longrightarrow} TY \stackrel{T(s)}{\longrightarrow} TTY  \stackrel{  \pi_{\text{vert} } }{\longrightarrow} TY} { . }

Pemetaan $T(s)$ mengirim titik
\mathl{(P,u)}{} ke
\mathl{(P,u, s(P), \left(  D s   \right)_{ P } \left(  u   \right))}{} dalam notasi
Catatan 24.8,
kemudian proyeksi vertikal mengirimnya ke
\mathl{(P, \pi_P  \left(  \left(  D s   \right)_{ P } \left(  u   \right)  \right) )}{.}
Inilah tepatnya rumus dalam definisi turunan kovarian.

}

\chapter{Lembar Kerja 24}
\setcounter{section}{24}

  \zwischenueberschrift{Soal latihan}

\inputaufgabe
{}
{

Misalkan $X$ suatu
\definitionsverweis {manifold diferensiabel}{}{}
dengan
\definitionsverweis {basis topologi terhitung}{}{}
dan misalkan

\mathdisp {0 \, \stackrel{  }{\longrightarrow} \,  U   \, \stackrel{  }{ \longrightarrow}   \,  V   \, \stackrel{  }{\longrightarrow} \,  W \,\stackrel{  } {\longrightarrow}  \, 0} {  }

suatu
\definitionsverweis {barisan eksak pendek}{}{}
dari
\definitionsverweis {bundel-bundel vektor diferensiabel}{}{}
di atas $X$. Dengan menggunakan
Teorema 22.10,
tunjukkan bahwa barisan tersebut memiliki suatu pemisahan, yaitu suatu isomorfisme bundel vektor

\mavergleichskette
{\vergleichskette
{ V
}
{ \cong }{ U \oplus W
}
{  }{
}
{    }{
}
{   }{
}
}
{}{}{,}
yang kompatibel dengan homomorfisme-homomorfisme bundel.

}
{} {}

\inputaufgabe
{}
{

Misalkan

\mavergleichskette
{\vergleichskette
{ W
}
{ \subseteq }{ \R^n
}
{  }{
}
{    }{
}
{   }{
}
}
{}{}{}
\definitionsverweis {terbuka}{}{,}
\maabb {h} { W } { \R
} {}
suatu
\definitionsverweis {fungsi terdiferensialkan secara kontinu}{}{}
dan

\mavergleichskette
{\vergleichskette
{ Y
}
{ = }{ h^{-1}(c)
}
{  }{
}
{    }{
}
{   }{
}
}
{}{}{}
merupakan
\definitionsverweis {serat}{}{}
di atas

\mavergleichskette
{\vergleichskette
{ c
}
{ \in }{ \R
}
{  }{
}
{    }{
}
{   }{
}
}
{}{}{,}
dan anggap $h$
\definitionsverweis {reguler}{}{}
di setiap titik $Y$. Realisasikan bundel tangen pertama dan kedua
\definitionsverweis {bundel tangen}{}{}
dari $Y$ sebagai
\definitionsverweis {submanifold tertutup}{}{}
di $W \times \R^n$ dan, secara berturut-turut, di
\mathl{W \times \R^n \times \R^n \times \R^n}{} dalam pengertian
Catatan 24.4
bersama dengan proyeksi bundel
\maabbeledisp {\pi} {TY} {Y
} {(P,u)} {P
} {,}
dan
\maabbeledisp {q} {TTY} {TY
} {(P,u,v,w)} {(P,u)
} {,}
Tunjukkan pernyataan-pernyataan berikut.
\aufzaehlungdrei{Suatu
\definitionsverweis {kurva diferensiabel}{}{}
\maabbeledisp {\epsilon} { I } { TY
} { t } { (P(t), u(t))
} {,}
pada interval terbuka

\mavergleichskette
{\vergleichskette
{ I
}
{ \subseteq }{ \R
}
{  }{
}
{    }{
}
{   }{
}
}
{}{}{}
dengan

\mavergleichskette
{\vergleichskette
{ 0
}
{ \in }{ I
}
{  }{
}
{    }{
}
{   }{
}
}
{}{}{}
mendefinisikan
\definitionsverweis {vektor singgung}{}{}

\mavergleichskette
{\vergleichskette
{ [\epsilon]
}
{ = }{ (P(0), u(0), P'(0), u'(0))
}
{ \in }{ TTY
}
{    }{
}
{   }{
}
}
{}{}{.}
}{Di bawah
\definitionsverweis {pemetaan tangen}{}{}
\maabbdisp {T_{(P,u)} (\pi)} {T_{ (P,u)} TY} { T_PY
} {}
yang berkaitan dengan $\pi$, kelas
\mathl{[\epsilon]}{} dari (1) dipetakan ke

\mavergleichskettedisp
{\vergleichskette
{ [ \pi \circ \epsilon ]
}
{ =} { [P'(0)]
}
{ } {
}
{ } {
}
{ } {
}
}
{}{}{}
.
}{Di bawah
\definitionsverweis {pemetaan tangen}{}{}
\maabb {T(\pi)} {TTY} {TY
} {}
yang berkaitan dengan $\pi$, elemen
\mathl{(P,u,v,w)}{} dipetakan ke
\mathl{(P,v)}{}.
}

}
{} {}

\inputaufgabe
{}
{

Untuk lingkaran satuan

\mavergleichskettedisp
{\vergleichskette
{ Y
}
{ =} { \left\{ (x_1,x_2)  \mid   x_1^2 + x_2^2 = 1         \right\}
}
{ \subset} { \R^2
}
{ } {
}
{ } {
}
}
{}{}{}
tentukan deskripsi implisit
\definitionsverweis {bundel tangen}{}{}
$TY$ dan bundel tangen kedua $TTY$ dalam pengertian
Catatan 24.5.

}
{} {}

\inputaufgabe
{}
{

Misalkan

\mavergleichskette
{\vergleichskette
{ W
}
{ \subseteq }{ \R^n
}
{  }{
}
{    }{
}
{   }{
}
}
{}{}{}
\definitionsverweis {terbuka}{}{,}
\maabb {\varphi} { W } { \R^k
} {}
suatu pemetaan
$C^2$-\definitionsverweis {terdiferensialkan secara kontinu}{}{},

\mavergleichskette
{\vergleichskette
{ c
}
{ \in }{ \R^k
}
{  }{
}
{    }{
}
{   }{
}
}
{}{}{}
dan

\mavergleichskette
{\vergleichskette
{ Y
}
{ = }{ \varphi^{-1} (c)
}
{  }{
}
{    }{
}
{   }{
}
}
{}{}{}
merupakan
\definitionsverweis {serat}{}{} di atas $c$. Anggap serat ini
\definitionsverweis {reguler}{}{}
di setiap titiknya. Kembangkan suatu deskripsi implisit
\zusatzklammer {yakni dengan persamaan} {} {}
dari
\definitionsverweis {bundel tangen}{}{}
$TY$ dan bundel tangen kedua $TTY$, serupa dengan
Catatan 24.5.

}
{} {}

\inputaufgabe
{}
{

Misalkan $X$ suatu
$C^2$-\definitionsverweis {manifold diferensiabel}{}{.}
Misalkan

\mavergleichskette
{\vergleichskette
{ I,J
}
{ \subseteq }{ \R
}
{  }{
}
{    }{
}
{   }{
}
}
{}{}{}
\definitionsverweis {interval terbuka}{}{}
dengan

\mavergleichskette
{\vergleichskette
{ 0
}
{ \in }{ I,J
}
{  }{
}
{    }{
}
{   }{
}
}
{}{}{}
dan misalkan
\maabbeledisp {\varphi} {I \times J} {X
} {(s,t)} { \varphi(s,t)
} {,}
suatu pemetaan dua kali
\definitionsverweis {terdiferensialkan secara kontinu}{}{}
dengan

\mavergleichskette
{\vergleichskette
{ \varphi(0,0)
}
{ = }{ P
}
{ \in }{ X
}
{    }{
}
{   }{
}
}
{}{}{.}
\aufzaehlungvier{Tunjukkan bahwa untuk setiap

\mavergleichskette
{\vergleichskette
{ s
}
{ \in }{ I
}
{  }{
}
{    }{
}
{   }{
}
}
{}{}{}
pembatasan

\mathl{\varphi \vert_{ \{s\} \times J }}{} mendefinisikan suatu kurva diferensiabel di $X$.
}{Tunjukkan bahwa (1) mendefinisikan suatu kurva diferensiabel
\maabbeledisp {} { I } { TX
} { s } { [ \varphi \vert_{ \{s\} \times J } ]
} {,}
di mana
\mathl{[ \varphi \vert_{ \{s\} \times J } ]}{} menyatakan vektor singgung kurva tersebut pada titik $(s,0)$.
}{Tunjukkan bahwa (2) menentukan suatu elemen dalam
\mathl{TTY}{}.
}{Tunjukkan bahwa setiap elemen
\mathl{TTY}{} dapat direalisasikan oleh suatu pemetaan $\varphi$ seperti di atas.
}

}
{} {}

\inputaufgabe
{}
{

Misalkan $X$ suatu
$C^2$-\definitionsverweis {manifold diferensiabel}{}{.}
Kita memandang pemetaan-pemetaan yang dua kali
\definitionsverweis {terdiferensialkan secara kontinu}{}{}

\maabbeledisp {\varphi} {I \times J} {X
} {(s,t)} { \varphi(s,t)
} {,}
sebagai elemen dalam
\mathl{TTX}{} dalam arti
Soal 24.5.
Bagaimana realisasi ini berperilaku berdasarkan barisan eksak pendek

\mathdisp {0 \, \stackrel{  }{\longrightarrow} \, p^*TX  \, \stackrel{  }{ \longrightarrow}   \, TTX  \, \stackrel{  }{\longrightarrow} \, p^*TX \,\stackrel{  } {\longrightarrow}  \, 0} { . }

}
{} {}

\inputaufgabe
{}
{

Misalkan
\maabb {} { E } { X
} {}

suatu
\definitionsverweis {bundel vektor diferensiabel}{}{}
di atas suatu
\definitionsverweis {manifold diferensiabel}{}{}
$X$. Tentukan
\definitionsverweis {rank}{}{}
bundel-bundel vektor di atas $E$ dalam barisan eksak pendek

\mathdisp {0 \, \stackrel{  }{\longrightarrow} \,  V   \, \stackrel{  }{ \longrightarrow}   \,  TE  \, \stackrel{  }{\longrightarrow} \, p^* TX \,\stackrel{  } {\longrightarrow}  \, 0} { . }

}
{} {}

\inputaufgabe
{}
{

Tentukan
\definitionsverweis {koneksi}{}{}
pada bundel nol di atas suatu
\definitionsverweis {manifold diferensiabel}{}{}
$X$.

}
{} {}

\inputaufgabe
{}
{

Tentukan
\definitionsverweis {koneksi}{}{}
pada suatu
\definitionsverweis {bundel vektor}{}{}
di atas
\definitionsverweis {manifold}{}{}
satu titik $X$.

}
{} {}

\inputaufgabe
{}
{

Misalkan $X$ suatu
\definitionsverweis {manifold diferensiabel}{}{}
dan $E$ suatu
\definitionsverweis {bundel vektor diferensiabel}{}{}
di atas $X$ yang dilengkapi dengan
\definitionsverweis {koneksi}{}{}
. Tunjukkan bahwa konsep turunan vertikal juga dapat didefinisikan sepanjang suatu pemetaan diferensiabel
\maabbdisp {\psi} { Z } { X
} {}
; bandingkan pula dengan
Definisi 6.5.

}
{} {}

  \zwischenueberschrift{Soal untuk dikumpulkan}

\inputaufgabe
{3}
{

Misalkan $X$ suatu
\definitionsverweis {ruang topologis}{}{}
dan misalkan

\mathdisp {0 \, \stackrel{  }{\longrightarrow} \, U  \, \stackrel{  }{ \longrightarrow}   \, V  \, \stackrel{  }{\longrightarrow} \, W \,\stackrel{  } {\longrightarrow}  \, 0} {  }

suatu
\definitionsverweis {barisan eksak pendek}{}{}
dari
\definitionsverweis {bundel vektor}{}{}
di atas $X$. Misalkan
\maabb {\varphi} {Y} {X
} {}
suatu
\definitionsverweis {pemetaan kontinu}{}{}
dari suatu ruang topologis lain $Y$. Tunjukkan bahwa penarikan balik menghasilkan barisan eksak pendek

\mathdisp {0 \, \stackrel{  }{\longrightarrow} \, \varphi^*U  \, \stackrel{  }{ \longrightarrow}   \, \varphi^*V  \, \stackrel{  }{\longrightarrow} \, \varphi^*W \,\stackrel{  } {\longrightarrow}  \, 0} {  }

bundel vektor di atas $Y$.

}
{} {}

\inputaufgabe
{3 (1+2)}
{

Misalkan

\mavergleichskette
{\vergleichskette
{ I
}
{ \subseteq }{ \R
}
{  }{
}
{    }{
}
{   }{
}
}
{}{}{}

\definitionsverweis {interval terbuka}{}{,}
kita tinjau
\definitionsverweis {barisan eksak pendek}{}{}
\zusatzklammer {dari bundel vektor di atas $I \times \R$} {} {}

\mathdisp {0 \, \stackrel{  }{\longrightarrow} \, I \times \R \times \R   \, \stackrel{  }{ \longrightarrow}   \,  I \times \R \times \R \times \R   \, \stackrel{  }{\longrightarrow} \,  I \times \R \times \R  \,\stackrel{  } {\longrightarrow}  \, 0} { , }

di mana pemetaan pertama diberikan oleh
\mathl{(P,u,w) \mapsto (P,u,0,w)}{} dan pemetaan kedua oleh
\mathl{(P,u,v,w) \mapsto (P,u,v)}{}. Misalkan
\maabbdisp {g,h} {I \times \R } {\R
} {}
fungsi-fungsi terdiferensialkan secara kontinu.
\aufzaehlungdrei{Tunjukkan bahwa
\maabbeledisp {} {I \times \R \times \R \times \R  } { I \times \R \times \R
} {(P,u,v,w)} { (P,u,   v g(P,u) +w h(P,u) )
} {,}
mendefinisikan suatu homomorfisme bundel vektor di atas $I \times \R$.
}{Misalkan

\mavergleichskette
{\vergleichskette
{ h
}
{ = }{ 1
}
{  }{
}
{    }{
}
{   }{
}
}
{}{}{}
konstan. Tunjukkan bahwa pemetaan pada (1) mendefinisikan suatu pemisahan untuk bundel di tengah.
}{
}

}
{} {}

\inputaufgabe
{4}
{

Untuk kurva bidang

\mavergleichskettedisp
{\vergleichskette
{ Y
}
{ =} { \left\{ (x_1,x_2)  \mid   x_1^3 -5x_1x_2 + x_2^4 = 1         \right\}
}
{ \subset} { \R^2
}
{ } {
}
{ } {
}
}
{}{}{}
tentukan deskripsi implisit
\definitionsverweis {bundel tangen}{}{}
$TY$ dan bundel tangen kedua $TTY$ dalam pengertian
Catatan 24.5.

}
{} {}

\inputaufgabe
{5 (3+2)}
{

Kita tinjau permukaan diferensiabel

\mavergleichskettedisp
{\vergleichskette
{ Y
}
{ =} { \left\{ (x_1,x_2,x_3)  \mid   x_1^2+x_2^2 -x_3^2 = 0  , \, (x,y,z) \neq 0       \right\}
}
{ \subseteq} { W
}
{ =} { \R^3 \setminus \{(0,0,0) \}
}
{ \subset} { \R^3
}
}
{}{}{.}
\aufzaehlungzwei {Tentukan deskripsi implisit dari
\definitionsverweis {bundel tangen}{}{}
$TY$ dan bundel tangen kedua $TTY$ dalam pengertian
Catatan 24.5.
} {Hitung untuk kurva diferensiabel
\maabbeledisp {} { I } { Y
} { t } { \left(  \cos  t   , \,   \sin  t  , \,   1                 \right)
} {}
kurva terkait di bundel tangen dan di bundel tangen kedua.
}

}
{} {}

\inputaufgabe
{3}
{

Misalkan

\mavergleichskette
{\vergleichskette
{ U
}
{ \subseteq }{ \R^n
}
{  }{
}
{    }{
}
{   }{
}
}
{}{}{}
\definitionsverweis {terbuka}{}{}
dan misalkan

\mavergleichskette
{\vergleichskette
{ E
}
{ = }{ U \times \R
}
{  }{
}
{    }{
}
{   }{
}
}
{}{}{,}
di mana fungsi-fungsi
\maabb {f} { U } {\R
} {}
kita pandang sebagai
\definitionsverweis {penampang}{}{}
di $E$. Lengkapi $E$ dengan
\definitionsverweis {koneksi trivial}{}{.} Tunjukkan bahwa

\mavergleichskette
{\vergleichskette
{ \nabla (f)
}
{ = }{ df
}
{  }{
}
{    }{
}
{   }{
}
}
{}{}{}
untuk fungsi-fungsi terdiferensialkan secara kontinu.

}
{} {}

\part{Unit 25}
\chapter[Kuliah 25]{Kuliah 25: Penampang Horizontal dan Koneksi Linear}
\setcounter{section}{25}

  \zwischenueberschrift{Penampang horizontal}

\inputdefinition
{}
{

Misalkan $X$ suatu
\definitionsverweis {manifold diferensiabel}{}{}
dan $E$ suatu
\definitionsverweis {bundel vektor diferensiabel}{}{}
di atas $X$ yang dilengkapi dengan suatu
\definitionsverweis {koneksi}{}{}.
Misalkan pula $Z$ suatu
\definitionsverweis {manifold diferensiabel}{}{}
dan
\maabbdisp {\psi} { Z } { X
} {}
suatu
\definitionsverweis {pemetaan diferensiabel}{}{.}
Suatu
\definitionsverweis {penampang diferensiabel}{}{}
\maabbdisp {s} { Z } { E
} {}
\zusatzklammer {di atas $\psi$} {} {}
disebut
\definitionswort {horizontal}{,}
jika untuk setiap titik

\mavergleichskette
{\vergleichskette
{ z
}
{ \in }{ Z
}
{  }{
}
{    }{
}
{   }{
}
}
{}{}{}
citra
\mathl{T(s)(T_zZ)}{} seluruhnya termuat dalam bundel horizontal

\mavergleichskette
{\vergleichskette
{ H_{s(z)}
}
{ \subseteq }{ T_{s(z)}E
}
{  }{
}
{    }{
}
{   }{
}
}
{}{}{}.

}

Dengan demikian, terdapat diagram komutatif

\mathdisp {\begin{matrix}  &   &   &  E   \\ & & s \nearrow & \downarrow  \\  &  Z  & \stackrel{ \psi }{\longrightarrow}  &  X \! \\ \end{matrix}} {  }.

Selain penampang horizontal, istilah \stichwort {pengangkatan horizontal} {} juga digunakan. Konsep ini khususnya dapat diterapkan pada penampang
\maabb {s} { U } { E\vert_U
} {}
di suatu
\definitionsverweis {himpunan terbuka}{}{}

\mavergleichskette
{\vergleichskette
{ U
}
{ \subseteq }{ X
}
{  }{
}
{    }{
}
{   }{
}
}
{}{}{}.


\inputfaktbeweis
{Mannigfaltigkeit/Vektorbündel/R/Zusammenhang/Horizontaler Schnitt/Vertikale Ableitung/Fakt}
{Lema}
{}
{

\faktsituation {Misalkan $X$ suatu
\definitionsverweis {manifold diferensiabel}{}{}
dan $E$ suatu
\definitionsverweis {bundel vektor diferensiabel}{}{}
di atas $X$ yang dilengkapi dengan suatu
\definitionsverweis {koneksi}{}{}.}
\faktfolgerung {Maka suatu penampang yang terdiferensialkan secara kontinu
\maabb {s} { U } { E \vert_U
} {}
di suatu himpunan terbuka

\mavergleichskette
{\vergleichskette
{ U
}
{ \subseteq }{ X
}
{  }{
}
{    }{
}
{   }{
}
}
{}{}{}
bersifat
\definitionsverweis {horizontal}{}{}
jika dan hanya jika
\definitionsverweis {turunan vertikal}{}{}-nya memenuhi

\mavergleichskette
{\vergleichskette
{ \nabla s
}
{ = }{ 0
}
{  }{
}
{    }{
}
{   }{
}
}
{}{}{.}}
\faktzusatz {}
\faktzusatz {}

 }
{ Lihat Soal 25.4. }

Pernyataan ini memiliki versi yang bersesuaian untuk penampang horizontal sepanjang suatu peta; lihat
Soal 25.5.

\inputdefinition
{}
{

Misalkan $X$ suatu
\definitionsverweis {manifold diferensiabel}{}{}
dan $E$ suatu
\definitionsverweis {bundel vektor diferensiabel}{}{}
di atas $X$ yang dilengkapi dengan suatu
\definitionsverweis {koneksi}{}{}.
Koneksi tersebut disebut
\definitionswort {terintegralkan secara lokal}{,}
jika untuk setiap titik

\mavergleichskette
{\vergleichskette
{ e
}
{ \in }{ E
}
{  }{
}
{    }{
}
{   }{
}
}
{}{}{}
terdapat suatu lingkungan terbuka

\mavergleichskette
{\vergleichskette
{ p(e)
}
{ \in }{ U
}
{ \subseteq }{ X
}
{    }{
}
{   }{
}
}
{}{}{}
serta suatu
\definitionsverweis {penampang horizontal}{}{}
\maabbdisp {s} { U } { E  \vert_U
} {}
yang terdefinisi di sana dan melalui $e$.

}
Koneksi tersebut disebut \stichwort {terintegralkan secara global} {,} jika untuk setiap

\mavergleichskette
{\vergleichskette
{ e
}
{ \in }{ E
}
{  }{
}
{    }{
}
{   }{
}
}
{}{}{}
terdapat suatu penampang horizontal pada seluruh $X$.

  \zwischenueberschrift{Koneksi linear}

\inputdefinition
{}
{

Misalkan $X$ suatu
\definitionsverweis {manifold diferensiabel}{}{}
dan $E$ suatu
\definitionsverweis {bundel vektor diferensiabel}{}{}
di atas $X$ yang dilengkapi dengan suatu
\definitionsverweis {koneksi}{}{}.
Koneksi tersebut disebut
\definitionswort {linear}{,}
jika
\definitionsverweis {turunan vertikal}{}{}
\maabbeledisp {\nabla} {C^1(E)} { C^0(E \otimes T^* X )
} { s } { \nabla(s)
} {,}
yang terkait dengannya bersifat
$\R$-\definitionsverweis {linear}{}{}
\zusatzklammer {pada setiap himpunan terbuka} {} {.}

}

Perhatikan bahwa hal ini tidak berarti peta tersebut berasal dari suatu homomorfisme bundel atau merupakan homomorfisme modul ${\mathcal O}_X$. Peta itu tidak kompatibel dengan perkalian penampang oleh fungsi, melainkan hanya dengan penjumlahan dan perkalian skalar oleh konstanta. Sebagai gantinya, berlaku \stichwort {aturan Leibniz} {} berikut.



\inputfaktbeweis
{Mannigfaltigkeit/Vektorbündel/R/Linearer Zusammenhang/Leibnizregel/Fakt}
{Teorema}
{}
{

\faktsituation {Misalkan $X$ suatu
\definitionsverweis {manifold diferensiabel}{}{}
dan $E$ suatu
\definitionsverweis {bundel vektor diferensiabel}{}{} di atas $X$ yang dilengkapi dengan suatu
\definitionsverweis {koneksi linear}{}{}.}
\faktfolgerung {Maka
\definitionsverweis {turunan vertikal}{}{}
\maabbeledisp {\nabla} {C^1(E)} { C^0(E \otimes T^* X )
} { s } { \nabla(s)
} {,}
yang terkait dengannya memenuhi aturan Leibniz; yaitu, untuk setiap
\definitionsverweis {penampang diferensiabel}{}{}
$s$ di $E$ dan setiap
\definitionsverweis {fungsi diferensiabel}{}{}
$f$
\zusatzklammer {yang keduanya terdefinisi pada suatu himpunan terbuka

\mavergleichskette
{\vergleichskette
{ U
}
{ \subseteq }{ X
}
{  }{
}
{    }{
}
{   }{
}
}
{}{}{}} {} {}
berlaku

\mavergleichskettedisp
{\vergleichskette
{ \nabla (fs)
}
{ =} { s \otimes df + f \nabla (s)
}
{ } {
}
{ } {
}
{ } {
}
}
{}{}{.}}
\faktzusatz {}
\faktzusatz {}

}
{

Kedua ruas bersifat lokal di $X$, sehingga kita dapat mengandaikan bahwa

\mavergleichskette
{\vergleichskette
{ E
}
{ = }{ X \times W
}
{  }{
}
{    }{
}
{   }{
}
}
{}{}{}
adalah bundel trivial
\zusatzklammer {dengan suatu ruang vektor-$\R$ yang dinotasikan $W$} {} {}.
Misalkan
\mathl{s_1 , \ldots , s_r}{} penampang basis konstan dari $E$; artinya, $s_i$ secara konstan sama dengan suatu vektor

\mavergleichskette
{\vergleichskette
{ w_i
}
{ \in }{ W
}
{  }{
}
{    }{
}
{   }{
}
}
{}{}{,}
dan vektor-vektor ini membentuk suatu basis dari $W$. Jika kesamaan tersebut telah dibuktikan untuk $s_i$ ini
\zusatzklammer {dan $f$ sembarang} {} {,}
maka kasus umum mengikutinya. Memang, setiap penampang $s$ di $E$ dapat dituliskan secara tunggal sebagai

\mavergleichskette
{\vergleichskette
{s
}
{ = }{ \sum_{i = 1}^r g_is_i
}
{  }{
}
{    }{
}
{   }{
}
}
{}{}{}
dengan fungsi-fungsi koefisien $g_i$. Dengan menggunakan linearitas-$\R$ kedua ruas, aturan Leibniz untuk penampang konstan
\zusatzklammer {lihat di bawah} {} {,}
dan
aturan hasil kali,
kita memperoleh

\mavergleichskettealign
{\vergleichskettealign
{ \nabla (fs)
}
{ =} { \nabla \left(  f \sum_{i = 1}^r g_is_i  \right)
}
{ =} { \nabla \left(  \sum_{i = 1}^r \left(  f g_i  \right) s_i  \right)
}
{ =} { \sum_{i = 1}^r \nabla \left(  \left(  f g_i  \right) s_i  \right)
}
{ =} { \sum_{i = 1}^r \left(  s_i \otimes d(fg_i) + f g_i \nabla s_i  \right)
}
}
{
\vergleichskettefortsetzungalign
{ =} { \sum_{i = 1}^r s_i \otimes \left(  g_i df + fdg_i  \right) + \sum_{i = 1}^r f g_i \nabla s_i
}
{ =} { \sum_{i = 1}^r \left(  g_i s_i \otimes df  \right)  + \sum_{i = 1}^r \left(  s_i \otimes f dg_i + f g_i \nabla s_i  \right)
}
{ =} { \left(  \sum_{i = 1}^r g_i s_i  \right)  \otimes df + f  \sum_{i = 1}^r \left(  s_i \otimes dg_i + g_i \nabla s_i  \right)
}
{ =} { s\otimes df + f \sum_{i = 1}^r \nabla \left(  g_i s_i  \right)
}
}
{
\vergleichskettefortsetzungalign
{ =} { s\otimes df + f \nabla s
}
{ } {}
{ } {}
{ } {}
}{.}
Sekarang misalkan $s$ suatu penampang yang secara konstan sama dengan $w$. Ambil

\mavergleichskette
{\vergleichskette
{ x
}
{ \in }{ X
}
{  }{
}
{    }{
}
{   }{
}
}
{}{}{.}
Karena

\mavergleichskette
{\vergleichskette
{ \nabla(fs)
}
{ = }{ \nabla ((f -f(x))  s ) + f(x) \nabla s
}
{  }{
}
{    }{
}
{   }{
}
}
{}{}{,}
kita selanjutnya dapat mengandaikan bahwa

\mavergleichskette
{\vergleichskette
{ f(x)
}
{ = }{ 0
}
{  }{
}
{    }{
}
{   }{
}
}
{}{}{}.
Dengan demikian,

\mavergleichskette
{\vergleichskette
{ (f \nabla s)(x)
}
{ = }{ 0
}
{  }{
}
{    }{
}
{   }{
}
}
{}{}{,}
dan kita hanya perlu meninjau suku pertama. Ruang singgung dari

\mavergleichskette
{\vergleichskette
{ E
}
{ = }{ X \times W
}
{  }{
}
{    }{
}
{   }{
}
}
{}{}{}
di
\mathl{(x,0)}{} adalah
\mathl{(T_x X \times 0) \oplus (0 \times W)}{.}
Karena penampang nol bersifat horizontal untuk suatu koneksi linear menurut
Soal 25.9,
penguraian ini juga merupakan penguraian menjadi ruang horizontal dan ruang vertikal. Turunan vertikal
\mathl{\nabla (fs)}{} di titik $x$ adalah peta komposit

\mathdisp {T_x X \stackrel{T(fs)}{ \longrightarrow} T_{(x,0)} (X \times W) \stackrel{\pi }{ \longrightarrow} W} { , }
dan karena $s$ konstan, peta ini sama dengan
\mathl{s \otimes df}{.}

}

Sebaliknya, suatu koneksi linear dapat ditentukan melalui aturan pendiferensialan yang memenuhi aturan Leibniz. Dengan demikian, formalisme turunan semacam itu dan penguraian jumlah langsung linear dari bundel tangen bundel vektor pada hakikatnya merupakan konsep-konsep yang ekuivalen.



\inputfaktbeweis
{Differenzierbare Mannigfaltigkeit/R/Vektorbündel/Linearer Zusammenhang/Existenz/Fakt}
{Lema}
{}
{

\faktsituation {Misalkan $X$ suatu
\definitionsverweis {manifold diferensiabel}{}{}
dengan
\definitionsverweis {basis topologi terhitung}{}{}
dan $E$ suatu
\definitionsverweis {bundel vektor diferensiabel}{}{}
di atas $X$.}
\faktfolgerung {Maka terdapat suatu
\definitionsverweis {koneksi linear}{}{}
pada $E$.}
\faktzusatz {}
\faktzusatz {}

}
{

Misalkan

\mavergleichskette
{\vergleichskette
{ X
}
{ = }{ \bigcup_{i \in I} U_i
}
{  }{
}
{    }{
}
{   }{
}
}
{}{}{}
suatu
\definitionsverweis {tutupan terbuka}{}{}
yang mentrivialkan $E$. Terdapat
\definitionsverweis {koneksi trivial}{}{}
$\nabla_i$ pada $E\vert_{U_i }$, yang linear menurut
Soal 25.10.
Menurut
Teorema 22.10,
terdapat suatu partisi satuan yang tersubordinasi pada tutupan tersebut,

\mathbed {h_j} {}
 {j \in J} {}
 {} {} {} {.}
Maka
\mathl{\sum_{j \in J} h_j \nabla_{i(j)}}{} merupakan suatu koneksi linear yang terdefinisi dengan baik pada $X$ menurut
Soal 25.16.

}

  \zwischenueberschrift{Simbol Christoffel untuk suatu koneksi linear}

\inputdefinition
{}
{

Misalkan

\mavergleichskette
{\vergleichskette
{ U
}
{ \subseteq }{ \R^n
}
{  }{
}
{    }{
}
{   }{
}
}
{}{}{}
\definitionsverweis {terbuka}{}{}
dan $\nabla$ suatu
\definitionsverweis {koneksi linear}{}{}
pada
\definitionsverweis {bundel vektor}{}{}
trivial

\mavergleichskette
{\vergleichskette
{ E
}
{ = }{ U \times \R^r
}
{  }{
}
{    }{
}
{   }{
}
}
{}{}{.}
Misalkan $x_i$ koordinat-koordinat pada $U$ dan $e_j$ penampang-penampang standar di $E$. Fungsi-fungsi yang ditentukan oleh persamaan

\mavergleichskettedisp
{\vergleichskette
{ \nabla_{\partial_i } e_j
}
{ =} { \sum_{k = 1}^r \Gamma^k_{i j}  e_k
}
{ } {
}
{ } {
}
{ } {}
}
{}{}{}
dan dipetakan sebagai
\maabbdisp {\Gamma^k_{i j}} { U } {\R
} {}
disebut
\definitionswort {simbol Christoffel}{}
dari koneksi tersebut terhadap $x_i$ dan $e_j$.

}

Dengan demikian, simbol Christoffel mendeskripsikan peta-peta

\mathdisp {U \stackrel{ \partial_i }{\longrightarrow} TU \stackrel{ T(e_j) }{\longrightarrow} TE \stackrel{ \nabla}{\longrightarrow} E} {  }

terhadap penampang basis $e_k$ dari $E$. Linearitas koneksi tidak diperlukan untuk definisi ini. Namun, daya guna simbol Christoffel baru muncul bersama
aturan Leibniz,
yang mengandaikan linearitas; lihat
Lema 25.10.

\inputbemerkung
{}
{

Untuk $E$ trivial di

\mavergleichskette
{\vergleichskette
{ U
}
{ \subseteq }{ \R^n
}
{  }{
}
{    }{
}
{   }{
}
}
{}{}{}
dengan rank $r$, berlaku

\mavergleichskettedisp
{\vergleichskette
{ T (U \times \R^r)
}
{ =} { U \times \R^r \times \R^n \times \R^r
}
{ } {
}
{ } {
}
{ } {
}
}
{}{}{.}
Jika
\definitionsverweis {bundel vertikal}{}{}

\mavergleichskettedisp
{\vergleichskette
{ V E
}
{ \cong} { U \times \R^r \times \R^r
}
{ \subseteq} { T E
}
{ \cong} { U \times \R^r \times \R^n \times \R^r
}
{ } {
}
}
{}{}{}
dan fungsi-fungsi
\maabb {\Gamma^k_{ij}} { U } {\R
} {}
diberikan, maka suatu proyeksi vertikal
\zusatzklammer {yang ekuivalen dengan suatu pemisahan bundel tangen $TE$, dan karenanya dengan suatu
\definitionsverweis {koneksi}{}{}} {} {}
diberikan oleh
\maabbeledisp {} {U \times \R^r \times \R^n \times \R^r } { U \times \R^r \times  \R^r
} {(P,u, v, w)} {(P,u,\sum_{k = 1}^r    \left(  \sum_{i,j} v_iu_j \Gamma^k_{ij} (P)  \right)  e_k +w)
} {,}
dan secara khusus oleh
\maabbeledisp {} {U \times \R^r \times \R^n \times \R^r } { U \times \R^r \times  \R^r
} {(P,e_j, \partial_i, w)} {(P,e_j,\sum_{k = 1}^r \Gamma^k_{ij} (P) e_k +w)
} {.}
\definitionsverweis {Turunan vertikal}{}{}
yang terkait,

\mathdisp {U \stackrel{ \partial_i }{\longrightarrow}  TU  \stackrel{T(e_j)}{\longrightarrow} TE  \stackrel{\nabla}{\longrightarrow}  E} {  },

memetakan

\mathdisp {P \longmapsto (P, \partial_i) \longmapsto (P , e_j, \partial_i, 0)  \longmapsto  \sum_{k = 1}^r \Gamma^k_{ij} (P) e_k} { . }

Jadi kita memperoleh kembali sifat yang mendefinisikan
\definitionsverweis {simbol Christoffel}{}{}.

}

\inputdefinition
{}
{

Misalkan $X$ suatu
\definitionsverweis {manifold diferensiabel}{}{}
dan $E$ suatu
\definitionsverweis {bundel vektor diferensiabel}{}{}
di atas $X$ yang dilengkapi dengan suatu
\definitionsverweis {koneksi linear}{}{}
$\nabla$. Untuk setiap domain peta terbuka

\mavergleichskette
{\vergleichskette
{ U
}
{ \subseteq }{ X
}
{  }{
}
{    }{
}
{   }{
}
}
{}{}{}
dengan koordinat $x_i$
\zusatzklammer {dan medan turunan $\partial_i$ yang terkait} {} {,}
serta tempat $E$ ditrivialkan dengan
\definitionsverweis {penampang basis}{}{}
$s_j$, fungsi-fungsi yang ditentukan oleh persamaan

\mavergleichskettedisp
{\vergleichskette
{ \nabla_{\partial_i } s_j
}
{ =} { \sum_{k = 1}^r \Gamma^k_{i j} s_k
}
{ } {
}
{ } {
}
{ } {}
}
{}{}{}
dan dipetakan sebagai
\maabbdisp {\Gamma^k_{i j}} { U } {\R
} {}
disebut
\definitionswort {simbol Christoffel lokal}{}
dari koneksi tersebut terhadap trivialisasi-trivialisasi itu.

}



\inputfaktbeweis
{Differenzierbare Mannigfaltigkeit/R/Vektorbündel/Linearer Zusammenhang/Lokale Beschreibung/Fakt}
{Lema}
{}
{

\faktsituation {Misalkan $X$ suatu
\definitionsverweis {manifold diferensiabel}{}{}
dan $E$ suatu
\definitionsverweis {bundel vektor diferensiabel}{}{} di atas $X$ yang dilengkapi dengan suatu
\definitionsverweis {koneksi linear}{}{}.}
\faktfolgerung {Maka, secara lokal pada suatu domain peta dengan koordinat $x_i$, medan turunan $\partial_i$, dan penampang basis $s_j$ di $E$, untuk
\definitionsverweis {simbol Christoffel}{}{}
yang terkait dan fungsi-fungsi yang terdiferensialkan secara kontinu
\maabb {h_i,f_j} { U } {\R
} {}
berlaku hubungan

\mavergleichskettedisphandlinks
{\vergleichskettedisphandlinks
{ \nabla_{ \sum_{i = 1}^n h_i \partial_i }  \left(  \sum_{j =1}^r  f_js_j  \right)
}
{ =} { \sum_{k = 1}^r \left(  \sum_{i = 1}^n h_i  \partial_i f_k  +  \sum_{i ,j} h_i f_j \Gamma_{ij}^k  \right)  s_k
}
{ } {
}
{ } {
}
{ } {
}
}
{}{}{.}}
\faktzusatz {}
\faktzusatz {}

}
{

Menurut
Lema 24.10
dan
Teorema 25.5,
kita memperoleh

\mavergleichskettealignhandlinks
{\vergleichskettealignhandlinks
{ \nabla_{ \sum_{i = 1}^n h_i \partial_i }  \left(  \sum_{j =1}^r  f_js_j  \right)
}
{ =} { \sum_{i, j} h_i \nabla_{\partial_i } \left(  f_js_j  \right)
}
{ =} { \sum_{i, j} h_i  \left(  s_j \otimes \partial_i f_j  + f_j  \nabla_{\partial_i } (s_j)  \right)
}
{ =} { \sum_{i,j}  h_i s_j \partial_i f_j  +  \sum_{i ,j,k} h_i f_j \Gamma_{ij}^k s_k
}
{ =} { \sum_{k} \left(  \sum_{i = 1}^n h_i  \partial_i f_k  +  \sum_{i ,j} h_i f_j \Gamma_{ij}^k  \right)  s_k
}
}
{}
{}{.}

}

  \zwischenueberschrift{Penampang horizontal pada suatu koneksi linear}

Untuk suatu himpunan terbuka

\mavergleichskette
{\vergleichskette
{ U
}
{ \subseteq }{ X
}
{  }{
}
{    }{
}
{   }{
}
}
{}{}{,}
ruang penampang horizontal dari suatu bundel vektor $E$ dengan koneksi linear $\nabla$ kita nyatakan dengan
\mathl{H(\nabla,U)}{.}
Jadi,

\mavergleichskettedisp
{\vergleichskette
{ H(\nabla,U)
}
{ =} { \left\{ s \in C^1(U,E)  \mid  \nabla(s) =0         \right\}
}
{ } {
}
{ } {
}
{ } {
}
}
{}{}{.}
Ini merupakan pilihan penampang yang \anfuehrung{sangat kecil}{,} sebagaimana sudah terlihat dari
Soal 25.21
dan juga diperlihatkan oleh lema berikut. Untuk bundel vektor trivial ber-rank $1$ dengan koneksi trivial, penampang horizontal hanyalah penampang yang konstan secara lokal. Dalam konteks ini orang berbicara tentang berkas lokal konstan atau sistem lokal.


\inputfaktbeweis
{Mannigfaltigkeit/Vektorbündel/R/Linearer Zusammenhang/Horizontale Schnitte/Untervektorraum/Fakt}
{Lema}
{}
{

\faktsituation {Misalkan $X$ suatu
\definitionsverweis {manifold diferensiabel}{}{}
dan $E$ suatu
\definitionsverweis {bundel vektor diferensiabel}{}{} ber-rank $r$ di atas $X$ yang dilengkapi dengan suatu
\definitionsverweis {koneksi linear}{}{}.}
\faktfolgerung {Maka, untuk setiap himpunan bagian terbuka terhubung

\mavergleichskette
{\vergleichskette
{ U
}
{ \subseteq }{ X
}
{  }{
}
{    }{
}
{   }{
}
}
{}{}{,}
himpunan
\definitionsverweis {penampang horizontal}{}{}

\mathl{H(\nabla,U)}{} merupakan suatu
$\R$-\definitionsverweis {subruang vektor}{}{}
dari
\mathl{C^1(U,E)}{} yang berdimensi
\mathl{\leq r}{.}}
\faktzusatz {}
\faktzusatz {}

}
{

Misalkan
\mathl{s_1 , \ldots , s_r}{} penampang-penampang horizontal yang bebas linear pada $U$. Misalkan $s$ suatu penampang horizontal terdiferensiasi lainnya. Penampang $s$, menurut
Soal 25.6
\zusatzklammer {yang diterapkan pada $U$} {} {,}
dapat dituliskan secara tunggal sebagai

\mavergleichskette
{\vergleichskette
{ s
}
{ = }{ \sum_{i = 1}^r g_is_i
}
{  }{
}
{    }{
}
{   }{
}
}
{}{}{}
dengan fungsi-fungsi
\maabbdisp {g_i} { U} { \R
} {}.
Maka, menurut
Teorema 25.5,

\mavergleichskettedisp
{\vergleichskette
{ 0
}
{ =} { \nabla s
}
{ =} { \nabla \left(  \sum_{i = 1}^r g_is_i  \right)
}
{ =} { \sum_{i = 1}^r  \nabla \left(  g_is_i  \right)
}
{ =} { \sum_{i = 1}^r  s_i \otimes d g_i
}
}
{}{}{.}
Pada setiap himpunan terbuka yang lebih kecil dan difeomorfik dengan

\mavergleichskette
{\vergleichskette
{ U'
}
{ \subseteq }{ \R^n
}
{  }{
}
{    }{
}
{   }{
}
}
{}{}{,}
dengan memakai koordinat
\mathl{x_1 , \ldots , x_n}{} kita dapat menuliskannya sebagai

\mavergleichskettedisp
{\vergleichskette
{ \sum_{i = 1}^r  s_i \otimes \left(  \sum_{ j = 1}^n { \frac{   \partial  g_i   }{  \partial   x_j   } }  dx_j   \right)
}
{ =} { \sum_{i = 1}^r \sum_{ j = 1}^n { \frac{   \partial  g_i   }{  \partial   x_j   } } \left(  s_i \otimes d x_j  \right)
}
{ } {
}
{ } {
}
{ } {
}
}
{}{}{.}
Penampang-penampang
\mathl{s_i \otimes dx_j}{}
\zusatzklammer {di atas himpunan terbuka ini} {} {}
membentuk penampang basis dari
\mathl{E  \vert_{U'} \otimes T^* U'}{,}
sehingga harus berlaku

\mavergleichskette
{\vergleichskette
{ { \frac{   \partial  g_i   }{  \partial   x_j   } }
}
{ = }{ 0
}
{  }{
}
{    }{
}
{   }{
}
}
{}{}{.}
Jadi $g_i$ konstan
\zusatzklammer {karena keterhubungan, hal ini juga berlaku pada seluruh $U$} {} {,}
sehingga $s$ merupakan kombinasi linear-$\R$ dari $s_i$.

}

Jika secara global
\zusatzklammer {yaitu di seluruh $X$} {} {}
terdapat $r$ penampang horizontal yang bebas linear di dalam bundel vektor $E$, maka bundel vektor tersebut trivial.

Bagi suatu koneksi linear, keterintegralan lokal berarti bahwa untuk setiap titik

\mavergleichskette
{\vergleichskette
{ x
}
{ \in }{ X
}
{  }{
}
{    }{
}
{   }{
}
}
{}{}{}
dan

\mavergleichskette
{\vergleichskette
{ e
}
{ \in }{ E_x
}
{  }{
}
{    }{
}
{   }{
}
}
{}{}{,}
secara lokal
\zusatzklammer {di suatu lingkungan terbuka dari $x$} {} {}
tercapai banyaknya maksimum penampang horizontal yang bebas linear menurut
Lema 25.11
\zusatzklammer {banyaknya ini ditentukan oleh rank bundel} {} {.}
Keterintegralan lokal selalu terpenuhi jika manifold basis berdimensi satu.

\chapter{Lembar Kerja 25}
\setcounter{section}{25}

  \zwischenueberschrift{Soal latihan}

\inputaufgabegibtloesung
{}
{

Misalkan $X$ suatu
\definitionsverweis {manifold diferensiabel}{}{}
dan
\maabb {p} { E  =  X \times \R } { X
} {}
bundel vektor trivial di atas $X$ yang ber-
\definitionsverweis {rank}{}{}
$1$. Misalkan $\omega$ suatu
$1$-\definitionsverweis {bentuk diferensial}{}{}
pada $X$. Tunjukkan pernyataan-pernyataan berikut.
\aufzaehlungfuenf{Pemetaan
\maabbeledisp {} {TE} { p^*E
} {(P,u,v,w) } { (P, u, \omega(P,v) +w)
} {,}
mendefinisikan suatu
\definitionsverweis {koneksi}{}{.}
}{Untuk fungsi-fungsi yang terdiferensialkan secara kontinu
\maabb {f} { U } {\R
} {}
pada suatu himpunan terbuka

\mavergleichskette
{\vergleichskette
{ U
}
{ \subseteq }{ X
}
{  }{
}
{    }{
}
{   }{
}
}
{}{}{}
turunan vertikal dari koneksi tersebut diberikan oleh

\mavergleichskettedisp
{\vergleichskette
{ \nabla (f) (P)
}
{ =} { \omega (P) + (df)_P
}
{ } {
}
{ } {
}
{ } {
}
}
{}{}{}
\zusatzklammer {di kedua ruas kita mengidentifikasi fungsi bernilai riil $f$ dengan penampang
\maabbele {} { X } { X \times \R
} { x } { (x,f(x))
} {.}} {} {}
}{Suatu fungsi yang terdiferensialkan secara kontinu
\maabb {f} { U } {\R
} {}
merupakan
\definitionsverweis {penampang horizontal}{}{}
di atas $U$ terhadap $\nabla$ jika dan hanya jika $-f$ merupakan
\definitionsverweis {bentuk primitif}{}{}
dari $\omega$ pada $U$.
}{Bentuk $\omega$
\definitionsverweis {tertutup}{}{}
jika dan hanya jika koneksi tersebut
\definitionsverweis {terintegralkan secara lokal}{}{.}
}{Bentuk $\omega$
\definitionsverweis {eksak}{}{}
jika dan hanya jika koneksi tersebut
\definitionsverweis {terintegralkan secara global}{}{.}
}

}
{} {}

\inputaufgabe
{}
{

Kita tinjau bundel vektor trivial
\maabbdisp {p} {E= \R^2 \times \R} { \R^2
} {.}
Deskripsikan suatu
\definitionsverweis {koneksi}{}{}
pada $E$ sedemikian sehingga tidak terdapat
\definitionsverweis {penampang horizontal}{}{.}

}
{} {}

\inputaufgabe
{}
{

Misalkan $X$ suatu
\definitionsverweis {manifold diferensiabel}{}{}
dan
\maabb {p} { E  =  X \times \R } { X
} {}
bundel vektor trivial di atas $X$ yang ber-
\definitionsverweis {rank}{}{}
$1$. Misalkan $\omega$ suatu
$1$-\definitionsverweis {bentuk diferensial}{}{}
pada $X$. Tunjukkan pernyataan-pernyataan berikut.
\aufzaehlungdrei{Pemetaan
\maabbeledisp {} {TE} { p^*E
} {(P,u,v,w) } { (P, u, u\omega(P,v) +w)
} {,}
mendefinisikan suatu
\definitionsverweis {koneksi}{}{.}
}{Untuk fungsi-fungsi yang terdiferensialkan secara kontinu
\maabb {f} { U } {\R
} {}
pada suatu himpunan terbuka

\mavergleichskette
{\vergleichskette
{ U
}
{ \subseteq }{ X
}
{  }{
}
{    }{
}
{   }{
}
}
{}{}{}
turunan vertikal dari koneksi tersebut diberikan oleh

\mavergleichskettedisp
{\vergleichskette
{ \nabla (f) (P)
}
{ =} { f\omega (P) + (df)_P
}
{ } {
}
{ } {
}
{ } {
}
}
{}{}{}
\zusatzklammer {di kedua ruas kita mengidentifikasi fungsi bernilai riil $f$ dengan penampang
\maabbele {} { X } { X \times \R
} { x } { (x,f(x))
} {.}} {} {}
}{Suatu fungsi yang terdiferensialkan secara kontinu dan tidak pernah nol
\maabb {f} { U } {\R
} {}
merupakan
\definitionsverweis {penampang horizontal}{}{}
di atas $U$ terhadap $\nabla$ jika dan hanya jika $-   \ln  \betrag {  f  }$ merupakan
\definitionsverweis {bentuk primitif}{}{}
dari $\omega$ pada $U$.
}

}
{} {}

\inputaufgabe
{}
{

Misalkan $X$ suatu
\definitionsverweis {manifold diferensiabel}{}{}
dan $E$ suatu
\definitionsverweis {bundel vektor diferensiabel}{}{}
di atas $X$ yang dilengkapi dengan suatu
\definitionsverweis {koneksi}{}{.}
Tunjukkan bahwa suatu penampang yang terdiferensialkan secara kontinu
\maabb {s} { U } { E\vert_U
} {}
pada suatu himpunan terbuka

\mavergleichskette
{\vergleichskette
{ U
}
{ \subseteq }{ X
}
{  }{
}
{    }{
}
{   }{
}
}
{}{}{}
bersifat
\definitionsverweis {horizontal}{}{}
jika dan hanya jika
\definitionsverweis {turunan vertikal}{}{}-nya memenuhi

\mavergleichskette
{\vergleichskette
{ \nabla s
}
{ = }{ 0
}
{  }{
}
{    }{
}
{   }{
}
}
{}{}{.}

}
{} {}

\inputaufgabe
{}
{

Misalkan $X$ suatu
\definitionsverweis {manifold diferensiabel}{}{}
dan $E$ suatu
\definitionsverweis {bundel vektor diferensiabel}{}{}
di atas $X$ yang dilengkapi dengan suatu
\definitionsverweis {koneksi}{}{.}
Misalkan $Z$ suatu manifold diferensiabel lain dengan pemetaan diferensiabel
\maabb {\psi} { Z } { X
} {.}
Tunjukkan bahwa suatu penampang yang terdiferensialkan secara kontinu
\maabb {s} { Z } { E
} {}
di atas $\psi$ bersifat
\definitionsverweis {horizontal}{}{}
jika dan hanya jika
\definitionsverweis {turunan vertikal}{}{}-nya memenuhi

\mavergleichskette
{\vergleichskette
{ \nabla s
}
{ = }{ 0
}
{  }{
}
{    }{
}
{   }{
}
}
{}{}{.}

}
{} {}

Untuk soal berikut, gunakan
Soal 25.21.

\inputaufgabe
{}
{

Misalkan
\maabb {p} { E } { X
} {}
suatu
\definitionsverweis {bundel vektor}{}{}
ber-
\definitionsverweis {rank}{}{}
$r$ di atas suatu
\definitionsverweis {manifold diferensiabel}{}{}
$X$, yang dilengkapi dengan suatu
\definitionsverweis {koneksi yang terintegralkan secara global}{}{.}
Tunjukkan bahwa terdapat
\zusatzklammer {yang terdefinisi pada $X$} {} {}
\definitionsverweis {penampang-penampang horizontal}{}{}

\mathl{s_1  , \ldots , s_r}{} di $E$ sedemikian sehingga
\maabbeledisp {} {X \times \R^r} { E
} {(P,v_1  , \ldots , v_r)} { \sum_{i = 1}^r v_i s_i(P)
} {,}
merupakan suatu
\definitionsverweis {isomorfisme bundel vektor}{}{.}

}
{} {}

\inputaufgabegibtloesung
{}
{

Misalkan
\maabb {p} { E } { X
} {}
suatu
\definitionsverweis {bundel vektor}{}{}
ber-
\definitionsverweis {rank}{}{}
$r$ di atas suatu
\definitionsverweis {manifold diferensiabel}{}{}
$X$, yang dilengkapi dengan suatu koneksi linear yang
\definitionsverweis {terintegralkan secara global}{}{.}
Tunjukkan bahwa terdapat
\zusatzklammer {yang terdefinisi pada $X$} {} {}
\definitionsverweis {penampang-penampang horizontal}{}{}

\mathl{s_1 , \ldots , s_r}{} di $E$ sedemikian sehingga, di bawah
\definitionsverweis {isomorfisme bundel vektor}{}{}
\maabbeledisp {} {X \times \R^r} { E
} {(P,v_1 , \ldots , v_r)} { \sum_{i = 1}^r v_i s_i(P)
} {,}
\definitionsverweis {simbol-simbol Christoffel}{}{}
dari koneksi tersebut terhadap penampang basis
\mathl{s_1  , \ldots , s_r}{} semuanya bernilai nol.

}
{} {}

\inputaufgabegibtloesung
{}
{

Misalkan diberikan suatu
\definitionsverweis {medan vektor}{}{}
kontinu yang bergantung pada waktu,
\maabbeledisp {F} {I \times \R^n } { \R^n
} { \left(  t   , \,  u_1  , \,   \ldots , \,  u_n                \right) } { \left(  F_1 \left(  u_1 , \ldots ,  u_n  \right)    , \,   \ldots  , \,   F_n \left(  u_1 , \ldots , u_n  \right)                 \right)
} {,}
beserta
\definitionsverweis {sistem persamaan diferensial biasa}{}{}
yang terkait,

\mavergleichskettedisp
{\vergleichskette
{ u'
}
{ =} { F(t,u)
}
{ } {
}
{ } {
}
{ } {
}
}
{}{}{}
pada suatu
\definitionsverweis {interval terbuka}{}{}

\mavergleichskette
{\vergleichskette
{ I
}
{ \subseteq }{ \R
}
{  }{
}
{    }{
}
{   }{
}
}
{}{}{.}
\aufzaehlungdrei{Tunjukkan bahwa pemetaan
\maabbeledisp {} {I \times \R^n \times \R \times \R^n } { I \times \R^n  \times \R^n
} {(t,u,v,w)} { \left(  t   , \,  u  , \,   - vF \left(  t, u_1  , \ldots , u_n  \right)  +w                 \right)
} {,}
memberikan
\zusatzklammer {proyeksi vertikal dari suatu} {} {}
\definitionsverweis {koneksi}{}{}
pada bundel vektor trivial
\maabbdisp {} {I \times \R^n } { \R^n
} {}
diberikan.
}{Tentukan
\definitionsverweis {turunan vertikal}{}{}
$\nabla_\partial u$ dari suatu
\definitionsverweis {kurva diferensiabel}{}{}
\maabbdisp {u} { I } { \R^n
} {}
\zusatzklammer {yang dipandang sebagai penampang di $I \times \R^n$} {} {.}
}{Tunjukkan bahwa suatu kurva diferensiabel
\maabb {u} { I } { \R^n
} {}
merupakan solusi dari sistem persamaan diferensial jika dan hanya jika $u$ merupakan
\definitionsverweis {penampang horizontal}{}{}
terhadap koneksi tersebut.
}

}
{} {}

\inputaufgabe
{}
{

Misalkan $\nabla$ suatu
\definitionsverweis {koneksi linear}{}{}
pada suatu
\definitionsverweis {bundel vektor diferensiabel}{}{}
$E$ di atas suatu
\definitionsverweis {manifold diferensiabel}{}{}
$X$. Tunjukkan bahwa penampang nol bersifat
\definitionsverweis {horizontal}{}{.}

}
{} {}

\inputaufgabe
{}
{

Misalkan $\nabla$
\definitionsverweis {koneksi trivial}{}{}
pada bundel vektor trivial
\maabb {p} {X \times W } { X
} {}
yang terkait dengan suatu ruang vektor riil
\mathl{W}{} di atas suatu
\definitionsverweis {manifold diferensiabel}{}{}
$X$. Tunjukkan bahwa $\nabla$
\definitionsverweis {linear}{}{.}

}
{} {}

\inputaufgabegibtloesung
{}
{

Misalkan

\mavergleichskette
{\vergleichskette
{ U
}
{ \subseteq }{ \R^n
}
{  }{
}
{    }{
}
{   }{
}
}
{}{}{}
suatu
\definitionsverweis {himpunan bagian terbuka}{}{}
dan misalkan $\nabla$
\definitionsverweis {koneksi trivial}{}{}
pada
\mathl{U \times \R}{.} Tunjukkan bahwa
\definitionsverweis {simbol-simbol Christoffel}{}{}
dan
\definitionsverweis {turunan vertikal}{}{}
dari koneksi tersebut memenuhi sifat-sifat berikut.
\aufzaehlungvier{Simbol-simbol Christoffel memenuhi

\mavergleichskettedisp
{\vergleichskette
{ \Gamma_{ij}^k
}
{ =} { 0
}
{ } {
}
{ } {
}
{ } {
}
}
{}{}{}
untuk semua $i,j,k$.
}{Berlaku

\mavergleichskettedisp
{\vergleichskette
{ \nabla_{\partial_i }  \partial_j
}
{ =} { 0
}
{ } {
}
{ } {
}
{ } {
}
}
{}{}{}
untuk
\definitionsverweis {medan-medan vektor standar}{}{}
$\partial_i$.
}{Berlaku

\mavergleichskettedisp
{\vergleichskette
{ \nabla_{\partial_i }  \left(  \sum_{j = 1}^n f_j \partial_j  \right)
}
{ =} { \sum_{j = 1}^n \partial_i(f_j) \partial_j
}
{ } {
}
{ } {
}
{ } {
}
}
{}{}{}
untuk fungsi-fungsi diferensiabel $f_j$.
}{Untuk suatu medan vektor $V$ dan suatu medan vektor diferensiabel $W$ pada $U$ berlaku

\mavergleichskettedisp
{\vergleichskette
{ ( \nabla_V W)(P)
}
{ =} { \left(  D W   \right)_{ P } \left(   V(P)  \right)
}
{ } {
}
{ } {
}
{ } {
}
}
{}{}{.}
}

}
{} {}

\inputaufgabegibtloesung
{}
{

Pada
\definitionsverweis {bundel vektor trivial}{}{}

\mathl{\R \times \R^2}{} di atas $\R$, kita tinjau
\definitionsverweis {koneksi trivial}{}{}
$\nabla$.
\aufzaehlungdrei{Tunjukkan bahwa

\mavergleichskette
{\vergleichskette
{ f_1(t)
}
{ = }{ \left(  1+ t^2  , \,  t^3                   \right)
}
{  }{
}
{    }{
}
{   }{
}
}
{}{}{}
dan

\mavergleichskette
{\vergleichskette
{ f_2(t)
}
{ = }{ \left(  t  , \,   1+t^2                   \right)
}
{  }{
}
{    }{
}
{   }{
}
}
{}{}{}
merupakan
\definitionsverweis {penampang-penampang basis}{}{}
dari bundel vektor tersebut.
}{Nyatakan penampang basis standar $e_1,e_2$ sebagai kombinasi linear dari penampang basis $f_1,f_2$.
}{Tentukan
\definitionsverweis {simbol-simbol Christoffel}{}{}
dari $\nabla$ terhadap penampang basis tersebut.
}

}
{} {}

\inputaufgabe
{}
{

Pada
\definitionsverweis {bundel vektor trivial}{}{}

\mathl{\R^2 \times \R^2}{} di atas $\R^2$, kita tinjau
\definitionsverweis {koneksi trivial}{}{}
$\nabla$.
\aufzaehlungzwei {Tunjukkan bahwa

\mavergleichskette
{\vergleichskette
{ f_1(x,y)
}
{ = }{ \left(  x   , \,   y                   \right)
}
{  }{
}
{    }{
}
{   }{
}
}
{}{}{}
dan

\mavergleichskette
{\vergleichskette
{ f_2(x,y)
}
{ = }{ \left(  - y  , \,   x                   \right)
}
{  }{
}
{    }{
}
{   }{
}
}
{}{}{}
merupakan
\definitionsverweis {penampang-penampang basis}{}{}
dari bundel vektor pada

\mavergleichskette
{\vergleichskette
{ U
}
{ = }{ \R^2 \setminus \{ (0,0)\}
}
{  }{
}
{    }{
}
{   }{
}
}
{}{}{.}
} {Tentukan
\definitionsverweis {simbol-simbol Christoffel}{}{}
dari $\nabla$ terhadap penampang basis tersebut pada $U$.
}

}
{} {}

\inputaufgabegibtloesung
{}
{

Misalkan diberikan suatu
\definitionsverweis {sistem persamaan diferensial linear homogen}{}{}

\mavergleichskettedisp
{\vergleichskette
{ u'
}
{ =} { Mu
}
{ } {
}
{ } {
}
{ } {
}
}
{}{}{}
pada suatu
\definitionsverweis {interval terbuka}{}{}

\mavergleichskette
{\vergleichskette
{ I
}
{ \subseteq }{ \R
}
{  }{
}
{    }{
}
{   }{
}
}
{}{}{}
dengan

\mavergleichskettedisp
{\vergleichskette
{ M(t)
}
{ =} { \left(  a_{kj} (t)  \right)_{kj}
}
{ } {
}
{ } {
}
{ } {
}
}
{}{}{}
\zusatzklammer {di sini $k$ adalah indeks baris dan $j$ adalah indeks kolom} {} {}
serta fungsi-fungsi kontinu
\maabbdisp {a_{kj}} { I } {\R
} {}
diberikan.
Melalui persamaan

\mavergleichskettedisp
{\vergleichskette
{ \Gamma_{j}^k (t)
}
{ \defeq} { -a_{kj}
}
{ } {
}
{ } {
}
{ } {
}
}
{}{}{}
kita memandang fungsi-fungsi ini sebagai
\definitionsverweis {simbol-simbol Christoffel}{}{}
\zusatzklammer {dengan

\mavergleichskettek
{\vergleichskettek
{ i
}
{ = }{ 1
}
{  }{
}
{    }{
}
{   }{
}
}
{}{}{}} {} {}
dari
\definitionsverweis {koneksi linear}{}{}
$\nabla$ pada bundel vektor trivial
\maabb {p} {I \times \R^n } { I
} {}
dalam pengertian
Catatan 25.8.
Tunjukkan bahwa suatu
\definitionsverweis {kurva diferensiabel}{}{}
\maabbdisp {u} { I } { \R^n
} {}
merupakan solusi sistem persamaan diferensial jika dan hanya jika $u$ merupakan
\definitionsverweis {penampang horizontal}{}{}
terhadap koneksi tersebut.

}
{} {}

\inputaufgabe
{}
{

Misalkan

\mavergleichskette
{\vergleichskette
{ Y
}
{ \subseteq }{ W
}
{ \subseteq }{ \R^n
}
{    }{
}
{   }{
}
}
{}{}{,}
$W$
\definitionsverweis {terbuka}{}{,}
suatu
\definitionsverweis {hipermuka terdiferensialkan}{}{,}
dan
\maabb {\gamma} { I } { Y
} {}
suatu
\definitionsverweis {kurva diferensiabel}{}{.}
Misalkan
\maabbdisp {F} { I } { TY
} {}
suatu medan vektor diferensiabel sepanjang $\gamma$; dengan kata lain, terdapat diagram komutatif

\mathdisp {\begin{matrix}  &   &   & TY  \\ & & F \nearrow & \downarrow  \\  &  I  & \stackrel{ \gamma }{\longrightarrow}  &  Y \! \\ \end{matrix}} {  }

Tunjukkan bahwa $F$
\definitionsverweis {paralel sepanjang}{}{}
$\gamma$ jika dan hanya jika $F$ merupakan
\definitionsverweis {penampang horizontal}{}{}
terhadap
\zusatzklammer {yang didefinisikan menggunakan proyeksi ortogonal} {} {}
\definitionsverweis {koneksi}{}{}
pada $TY$.

}
{} {}

\inputaufgabe
{}
{

Misalkan

\mavergleichskette
{\vergleichskette
{ Y
}
{ \subseteq }{ W
}
{ \subseteq }{ \R^n
}
{    }{
}
{   }{
}
}
{}{}{,}
$W$
\definitionsverweis {terbuka}{}{,}
suatu
\definitionsverweis {hipermuka terdiferensialkan}{}{,}
dan misalkan $\nabla$
\zusatzklammer {yang didefinisikan menggunakan proyeksi ortogonal} {} {}
\definitionsverweis {koneksi}{}{}
pada $TY$. Tunjukkan bahwa $\nabla$
\definitionsverweis {linear}{}{.}

}
{} {}

\inputaufgabe
{}
{

Misalkan
\mathkor {} {\nabla_1} {dan} {\nabla_2} {}
\definitionsverweis {koneksi-koneksi linear}{}{}
pada suatu
\definitionsverweis {bundel vektor diferensiabel}{}{}
$E$ di atas suatu
\definitionsverweis {manifold diferensiabel}{}{}
$X$. Misalkan
\maabbdisp {g_1,g_2} { X } { \R
} {}
fungsi-fungsi yang
\definitionsverweis {terdiferensialkan secara kontinu}{}{}
dengan

\mavergleichskette
{\vergleichskette
{ g_1+g_2
}
{ = }{ 1
}
{  }{
}
{    }{
}
{   }{
}
}
{}{}{.}
Tunjukkan bahwa
\mathl{g_1 \nabla_1+ g_2 \nabla_2}{} juga merupakan suatu koneksi linear
\zusatzklammer {dalam pengertian yang telah didefinisikan} {?} {}
pada $E$.

}
{} {}

\inputaufgabe
{}
{

Misalkan

\mavergleichskette
{\vergleichskette
{ X
}
{ = }{ \bigcup_{i \in I} U_i
}
{  }{
}
{    }{
}
{   }{
}
}
{}{}{}
suatu
\definitionsverweis {tutupan terbuka}{}{}
dari suatu
\definitionsverweis {manifold diferensiabel}{}{}
$X$ yang mentrivialkan
\definitionsverweis {bundel vektor diferensiabel}{}{}
$E$. Misalkan $\nabla_i$
\definitionsverweis {koneksi trivial}{}{}
pada $E\vert_{U_i }$ dan misalkan
\mathkor {} {h_j} {} {j \in J} {,}
suatu
\definitionsverweis {partisi satuan}{}{}
yang subordinat terhadap tutupan tersebut,

\mathbed {h_j} {}
 {j \in J} {}
 {} {} {} {.}
Tunjukkan bahwa $\sum_{j \in J} h_j \nabla_{i(j)}$ merupakan suatu koneksi linear yang terdefinisi dengan baik pada $E$.

}
{} {}

\inputaufgabe
{}
{

Misalkan $X$ suatu
\definitionsverweis {manifold diferensiabel}{}{}
dan $\omega$ suatu
$1$-\definitionsverweis {bentuk diferensial}{}{}
kontinu pada $X$. Tunjukkan bahwa pemetaan
\maabbeledisp {} {C^1(X,\R)} { C^0(X,T^*X)
} { f } { \nabla_\omega (f) \defeq f \omega + df
} {,}
memenuhi aturan Leibniz dari
Teorema 25.5.

}
{} {}

\inputaufgabe
{}
{

Misalkan
\maabb {p} { E } { X
} {}
suatu
\definitionsverweis {bundel vektor}{}{}
diferensiabel di atas suatu
\definitionsverweis {manifold diferensiabel}{}{}
$X$. Misalkan $E$ sendiri merupakan suatu
\definitionsverweis {manifold Riemann}{}{.}
Tunjukkan bahwa suatu
\definitionsverweis {koneksi linear}{}{}
pada $E$ diperoleh dengan mengambil, untuk
\definitionsverweis {bundel vertikal}{}{}

\mavergleichskette
{\vergleichskette
{ V
}
{ \subseteq }{ T E
}
{  }{
}
{    }{
}
{   }{
}
}
{}{}{}
\definitionsverweis {komplemen ortogonal}{}{}
sebagai suatu
\definitionsverweis {bundel horizontal}{}{.}

}
{} {}

\inputaufgabe
{}
{

Misalkan
\maabb {p} { E } { X
} {}
suatu
\definitionsverweis {bundel vektor}{}{}
di atas
\definitionsverweis {manifold diferensiabel}{}{}
$X$ dengan
\definitionsverweis {basis topologi terhitung}{}{.}
Gunakan
Soal 22.15
dan
Soal 25.18
untuk menunjukkan bahwa terdapat suatu
\definitionsverweis {koneksi linear}{}{}
pada $E$.

}
{} {}

  \zwischenueberschrift{Soal yang dikumpulkan}

\inputaufgabe
{5 (3+2)}
{

Misalkan

\mavergleichskette
{\vergleichskette
{ X
}
{ = }{ S^1
}
{  }{
}
{    }{
}
{   }{
}
}
{}{}{}
\definitionsverweis {lingkaran satuan}{}{}
dengan
\definitionsverweis {bundel tangen}{}{}
\maabb {p} {TS^1 = S^1 \times \R} { S^1
} {.}
\aufzaehlungzwei {Deskripsikan suatu
\definitionsverweis {koneksi}{}{}
pada $TS^1$ sedemikian sehingga tidak terdapat
\definitionsverweis {penampang horizontal}{}{}
pada seluruh $S^1$.
} {Tentukan suatu penampang horizontal dari koneksi pada (1) sepanjang pemetaan
\maabbeledisp {} {\R} { S^1
} { t } { \left(  \cos  t   , \,   \sin  t                   \right)
} {.}
}

}
{} {}

\inputaufgabe
{4}
{

Misalkan $X$ suatu
\definitionsverweis {manifold diferensiabel}{}{}
dan $E$ suatu
\definitionsverweis {bundel vektor diferensiabel}{}{}
di atas $X$ yang dilengkapi dengan suatu
\definitionsverweis {koneksi}{}{.}
Misalkan $Z$ suatu manifold diferensiabel lain dengan pemetaan diferensiabel
\maabb {\psi} { Z } { X
} {.}
Misalkan $Z$
\definitionsverweis {terhubung}{}{}
dan misalkan
\maabb {s_1,s_2} { Z } { E
} {}
\definitionsverweis {penampang-penampang horizontal}{}{}
di atas $\psi$ dengan

\mavergleichskette
{\vergleichskette
{ s_1(z)
}
{ = }{ s_2(z)
}
{  }{
}
{    }{
}
{   }{
}
}
{}{}{}
untuk suatu titik

\mavergleichskette
{\vergleichskette
{ z
}
{ \in }{ Z
}
{  }{
}
{    }{
}
{   }{
}
}
{}{}{.}
Tunjukkan bahwa

\mavergleichskette
{\vergleichskette
{ s_1
}
{ = }{ s_2
}
{  }{
}
{    }{
}
{   }{
}
}
{}{}{.}

}
{} {}

\inputaufgabe
{3}
{

Misalkan $X$ suatu
\definitionsverweis {manifold diferensiabel}{}{}
dan $E$ suatu
\definitionsverweis {bundel vektor diferensiabel}{}{}
di atas $X$ yang dilengkapi dengan suatu
\definitionsverweis {koneksi linear}{}{,}
dan misalkan $G$ suatu
\definitionsverweis {medan vektor}{}{}
pada $X$. Tunjukkan bahwa untuk setiap
\definitionsverweis {penampang diferensiabel}{}{}
$s$ di $E$ dan setiap
\definitionsverweis {fungsi diferensiabel}{}{}
$f$
\zusatzklammer {yang keduanya terdefinisi pada suatu himpunan terbuka

\mavergleichskette
{\vergleichskette
{ U
}
{ \subseteq }{ X
}
{  }{
}
{    }{
}
{   }{
}
}
{}{}{}} {} {}
berlaku hubungan

\mavergleichskettedisp
{\vergleichskette
{ \nabla_G (fs)
}
{ =} { s \otimes D_G (f) + f \nabla_G (s)
}
{ } {
}
{ } {
}
{ } {
}
}
{}{}{.}

}
{} {}

\inputaufgabe
{3}
{

Berikan suatu contoh
\definitionsverweis {koneksi linear}{}{}
pada suatu
\definitionsverweis {bundel vektor}{}{}
\maabb {p} { E } { X
} {}
di atas suatu
\definitionsverweis {manifold diferensiabel}{}{}
$X$ sedemikian sehingga ruang vektor
\mathl{H(\nabla,X)}{} dari
\definitionsverweis {penampang-penampang horizontal}{}{}
memiliki
\definitionsverweis {dimensi ruang vektor}{}{}
yang lebih besar daripada
\definitionsverweis {rank}{}{}
$E$.

}
{} {}

\section*{Solusi yang disediakan oleh sumber}
\addcontentsline{toc}{section}{Solusi yang disediakan oleh sumber}
\subsection*{Solusi Soal 25.1}
\aufzaehlungfuenf{Pemetaan
\maabbeledisp {\pi_{\text{vert} }} {TE= TX \times \R \times \R } { p^*E  = X \times \R \times \R
} {(P,u,v,w)} { (P,u,  \omega(P,v) + w)
} {}
bersifat linear dalam
\mathkor {} {v} {dan} {w} {,}
untuk $(P,u)$ yang tetap, sehingga merupakan suatu homomorfisme bundel vektor. Ketergantungannya pada basis $E$ terdiferensialkan secara kontinu karena $\omega$ terdiferensialkan secara kontinu. Pemetaan ini merupakan suatu pemisahan bagi penyertaan alami
\maabb {} {p^*E} { TE
} {}
karena
\mathl{(P,u,w)}{} dipetakan ke
\mathl{(P,u,0,w)}{} dan kemudian ke

\mavergleichskette
{\vergleichskette
{ (P,u, \omega(P,0) + w)
}
{ = }{ (P,u,  w)
}
{  }{
}
{    }{
}
{   }{
}
}
{}{}{.}
}{Untuk
\maabb {f} { U } {\R
} {}
kita perlu meninjau komposisi

\mathdisp {TX \stackrel{T(f)}{  \longrightarrow } T E = TX \times \R \times \R \stackrel{ \pi_{\text{vert} } }{ \longrightarrow } X \times \R \longrightarrow \R} {  }

yang menghasilkan

\mathdisp {(P,v) \longmapsto  (P,f(P), v, (df)_P(v) ) \longmapsto (P, f(P), \omega(P,v) + (df)_P(v) ) \longmapsto  \omega(P,v) + (df)_P(v)} { . }

}{Suatu penampang horizontal ada jika dan hanya jika turunan vertikalnya lenyap; dalam hal ini, syarat tersebut berarti

\mavergleichskettedisp
{\vergleichskette
{ \omega(P,v) + (df)_P(v)
}
{ =} { 0
}
{ } {
}
{ } {
}
{ } {
}
}
{}{}{}
untuk semua

\mavergleichskette
{\vergleichskette
{ (P,v)
}
{ \in }{ TX
}
{  }{
}
{    }{
}
{   }{
}
}
{}{}{.}
}{Hal ini mengikuti dari kenyataan bahwa suatu bentuk eksak jika dan hanya jika secara lokal bentuk tersebut mempunyai suatu bentuk primitif.
}{Hal ini mengikuti dari hasil-hasil yang telah dibuktikan.
}
\subsection*{Solusi Soal 25.7}
Menurut
Soal,
terdapat suatu isomorfisme bundel vektor. Menurut
Lema,
untuk suatu penampang horizontal $s$ berlaku

\mavergleichskette
{\vergleichskette
{ \nabla s
}
{ = }{ 0
}
{  }{
}
{    }{
}
{   }{
}
}
{}{}{.}
Oleh karena itu, pada suatu himpunan terbuka

\mavergleichskette
{\vergleichskette
{ U
}
{ \cong }{ V
}
{ \subseteq }{  \R^n
}
{    }{
}
{   }{
}
}
{}{}{}
secara khusus berlaku

\mavergleichskettedisp
{\vergleichskette
{ \nabla_{\partial_i} (s)
}
{ =} { 0
}
{ } {
}
{ } {
}
{ } {
}
}
{}{}{}
untuk semua $i$. Ini berarti bahwa semua simbol Christoffel dari
\definitionsverweis {penampang basis}{}{}
$s_j$ terhadap medan vektor standar $\partial_i$ bernilai nol.
\subsection*{Solusi Soal 25.8}
\aufzaehlungdrei{Untuk setiap
\mathl{(t,u)}{} yang tetap, terdapat suatu pemetaan linear karena $v$ dan $w$ memasuki suku ketiga secara linear. Jika di sebelah kiri kita mulai dengan
\mathl{(t,u,w)}{}, titik ini dipetakan ke
\mathl{(t,u,0,w)}{} dan kemudian kembali ke
\mathl{(t,u,w)}{,} sehingga memang terdapat suatu pemisahan.
}{Pemetaan singgung dari
\maabbeledisp {u} { I } { I \times \R^n
} { t } { (t,u(t))
} {,}
adalah
\maabbeledisp {} {TI = I \times \R } { T \left( I \times \R^n  \right) = I \times \R^n \times \R \times \R^n
} { (t,a) } { (t,u(t), a, a u'(t) )
} {.}
Oleh karena itu,
\zusatzklammer {turunan vertikal ke arah $\partial$ bersesuaian dengan

\mavergleichskettek
{\vergleichskettek
{ a
}
{ = }{ 1
}
{  }{
}
{    }{
}
{   }{
}
}
{}{}{}} {} {}

\mavergleichskettedisp
{\vergleichskette
{  \nabla_\partial (u)
}
{ =} { \pi_{\text{vert} }  (t,u(t), 1,  u'(t) )
}
{ =} { (t,u(t), -  F \left(  t, u_1 (t) , \ldots , u_n(t)  \right)  +  u'(t)
}
{ } {
}
{ } {
}
}
{}{}{.}
}{Penampang $u$ merupakan
\definitionsverweis {penampang horizontal}{}{}
jika dan hanya jika turunan vertikalnya lenyap, yaitu jika

\mavergleichskettedisp
{\vergleichskette
{  -  F \left(  t, u_1 (t) , \ldots , u_n(t)  \right)  +  u'(t)
}
{ =} { 0
}
{ } {
}
{ } {
}
{ } {
}
}
{}{}{}
atau, secara ekuivalen,

\mavergleichskettedisp
{\vergleichskette
{ u'(t)
}
{ =} {  F \left(  t, u_1 (t) , \ldots , u_n(t)  \right)
}
{ } {
}
{ } {
}
{ } {
}
}
{}{}{}
untuk semua

\mavergleichskette
{\vergleichskette
{ t
}
{ \in }{  I
}
{  }{
}
{    }{
}
{   }{
}
}
{}{}{.}
Ini tepat berarti bahwa kurva tersebut merupakan solusi sistem persamaan diferensial.
}
\subsection*{Solusi Soal 25.11}
\aufzaehlungvier{Hal ini mengikuti dari deskripsi eksplisit turunan vertikal suatu koneksi linear; lihat
Catatan.
}{Jelas dari (1).
}{Hal ini mengikuti dari (2) dan
Lema.
}{Kita tuliskan

\mavergleichskette
{\vergleichskette
{ V
}
{ = }{ \sum_{i = 1}^n {h_i } \partial_i
}
{  }{
}
{    }{
}
{   }{
}
}
{}{}{}
dan

\mavergleichskette
{\vergleichskette
{ W
}
{ = }{ \sum_{j = 1}^n {f_j } \partial_j
}
{  }{
}
{    }{
}
{   }{
}
}
{}{}{.}
Menurut
Lema
dan bagian-bagian yang telah dibuktikan,

\mavergleichskettealign
{\vergleichskettealign
{ \nabla_V W
}
{ =} { \nabla_{\sum_{i = 1}^n {h_i } \partial_i }   \left(  \sum_{j = 1}^n {f_j } \partial_j   \right)
}
{ =} { \sum_{j = 1}^n \left(  \sum_{i = 1}^n h_i  \partial_i f_j   \right) \partial_j
}
{ } {
}
{ } {
}
}
{}
{}{.}
Pada suatu titik

\mavergleichskette
{\vergleichskette
{ P
}
{ \in }{ U
}
{  }{
}
{    }{
}
{   }{
}
}
{}{}{}
nilai ini adalah

\mathdisp {\sum_{j  = 1}^n \left(  \sum_{i = 1}^n h_i (P) ( \partial_i f_j)(P)   \right) \partial_j} { . }

Di sisi lain,

\mavergleichskettealign
{\vergleichskettealign
{ \left(  D W   \right)_{ P } \left(   V(P)  \right)
}
{ =} { \left(  (\partial_j f_i)(P)  \right)_{ji} \begin{pmatrix} h_1(P) \\ \vdots\\  h_n(P)   \end{pmatrix}
}
{ =} { \begin{pmatrix}  \sum_{i = 1}^n h_i (P) \left(  \partial_i f_1  \right)(P)  \\ \vdots\\  \sum_{i = 1}^n h_i (P)  \left(  \partial_i f_n  \right) (P)   \end{pmatrix}
}
{ } {
}
{ } {
}
}
{}
{}{.}
}
\subsection*{Solusi Soal 25.12}
\aufzaehlungdrei{Determinan

\mavergleichskettedisp
{\vergleichskette
{ \det \begin{pmatrix} 1+t^2 & t^3 \\ t & 1+t^2 \end{pmatrix}
}
{ =} { 1+2t^2+t^4 -t^4
}
{ =} { 1 +2t^2
}
{ >} {  0
}
{ } {
}
}
{}{}{}
selalu positif. Oleh karena itu, kedua vektor
\mathkor {} {f_1(t)} {dan} {f_2(t)} {}
untuk setiap

\mavergleichskette
{\vergleichskette
{ t
}
{ \in }{  \R
}
{  }{
}
{    }{
}
{   }{
}
}
{}{}{}
membentuk suatu basis.
}{Dari hubungan

\mavergleichskettedisp
{\vergleichskette
{ \begin{pmatrix} f_1 \\f_2 \end{pmatrix}
}
{ =} { \begin{pmatrix} 1+t^2 & t^3 \\ t & 1+t^2 \end{pmatrix} \begin{pmatrix} e_1 \\e_2 \end{pmatrix}
}
{ } {
}
{ } {
}
{ } {
}
}
{}{}{}
menghasilkan

\mavergleichskettedisp
{\vergleichskette
{ \begin{pmatrix} e_1 \\e_2 \end{pmatrix}
}
{ =} { { \frac{ 1 }{ 1+2t^2 } } \begin{pmatrix} 1+t^2 & -t^3 \\ -t & 1+t^2 \end{pmatrix} \begin{pmatrix} f_1 \\f_2 \end{pmatrix}
}
{ } {
}
{ } {
}
{ } {
}
}
{}{}{,}
sehingga

\mavergleichskettedisp
{\vergleichskette
{ e_1
}
{ =} { { \frac{   1+t^2 }{  1+2t^2 } }   f_1 - { \frac{  t^3 }{  1+2t^2 } } f_2
}
{ } {
}
{ } {
}
{ } {
}
}
{}{}{}
dan

\mavergleichskettedisp
{\vergleichskette
{ e_2
}
{ =} { -{ \frac{   t  }{  1+2t^2 } }   f_1 + { \frac{   1+t^2 }{  1+2t^2 } } f_2
}
{ } {
}
{ } {
}
{ } {
}
}
{}{}{.}
}{Untuk koneksi trivial berlaku

\mavergleichskettedisp
{\vergleichskette
{ \nabla_\partial (f)
}
{ =} { \partial f
}
{ } {
}
{ } {
}
{ } {
}
}
{}{}{.}
Dengan demikian,

\mavergleichskettealign
{\vergleichskettealign
{ \nabla_\partial (f_1)
}
{ =} { \nabla_\partial \left(  \left(  1+t^2  \right) e_1 + t^3e_2  \right)
}
{ =} { 2t e_1 + 3t^2 e_2
}
{ =} { 2t  \left(  { \frac{   1+t^2 }{  1+2t^2 } }   f_1 - { \frac{  t^3 }{  1+2t^2 } } f_2   \right)  + 3t^2  \left(  -{ \frac{   t  }{  1+2t^2 } }   f_1 + { \frac{   1+t^2 }{  1+2t^2 } } f_2   \right)
}
{ =} {  \left(   { \frac{   2t+2t^3 }{  1+2t^2 } } - { \frac{   3t^3 }{  1+2t^2 } }  \right) f_1  +  \left(  - { \frac{  t^3 }{  1+2t^2 } }   + { \frac{   3t^2+3t^4 }{  1+2t^2 } }     \right)   f_2
}
}
{
\vergleichskettefortsetzungalign
{ =} { { \frac{   2t - t^3 }{  1+2t^2 } }  f_1  +  { \frac{   3t^2-t^3+3t^4 }{  1+2t^2 } }      f_2
}
{ } {}
{ } {}
{ } {}
}
{}{}
dan

\mavergleichskettealign
{\vergleichskettealign
{ \nabla_\partial (f_2)
}
{ =} { \nabla_\partial \left(  t e_1 + \left(  1+t^2  \right) e_2  \right)
}
{ =} { e_1 + 2t  e_2
}
{ =} { { \frac{   1+t^2 }{  1+2t^2 } } f_1 - { \frac{  t^3 }{  1+2t^2 } } f_2 + 2t \left(  -{ \frac{   t  }{  1+2t^2 } } f_1 + { \frac{   1+t^2 }{  1+2t^2 } } f_2   \right)
}
{ =} { \left( { \frac{   1+t^2 }{  1+2t^2 } }  - { \frac{   2t^2 }{  1+2t^2 } }  \right) f_1  + \left(  - { \frac{   t^3  }{  1+2t^2  } } +  { \frac{   2t+2t^3  }{  1+2t^2 } }  \right) f_2
}
}
{
\vergleichskettefortsetzungalign
{ =} { { \frac{   1 - t^2  }{  1+2t^2 } } f_1 + { \frac{   2t   +t^3 }{  1+2t^2 } }      f_2
}
{ } {}
{ } {}
{ } {}
}
{}{.}
Fungsi-fungsi koefisien tersebut adalah simbol-simbol Christoffel.
}
\subsection*{Solusi Soal 25.14}
Pemetaan
\maabbdisp {u} { I } {\R^n
} {}
dinyatakan terhadap penampang basis
\mathl{e_1 , \ldots , e_n}{}
oleh
\mavergleichskettedisp
{\vergleichskette
{ u
}
{ =} { u_1e_1 + \cdots + u_ne_n
}
{ } {
}
{ } {
}
{ } {
}
}
{}{}{}
.
Menurut
\definitionsverweis {turunan vertikal}{}{}
dari koneksi tersebut, turunan penampang ini dalam satu-satunya arah turunan $\partial$, berdasarkan
Lema,
adalah

\mavergleichskettedisp
{\vergleichskette
{  \sum_{k = 1}^n \left(    \partial u_k  +  \sum_{j = 1}^n u_j \Gamma_{j}^k  \right) e_k
}
{ =} { \sum_{k = 1}^n \left(   u_k'  +  \sum_{j = 1}^n u_j \Gamma_{j}^k  \right) e_k
}
{ } {
}
{ } {
}
{ } {
}
}
{}{}{.}
Menurut
Lema,
$u$ merupakan penampang horizontal jika dan hanya jika turunan vertikalnya sama dengan $0$. Untuk setiap komponen, ini berarti

\mavergleichskettedisp
{\vergleichskette
{ u_k' + \sum_{j = 1}^n \Gamma_j^k u_j
}
{ =} { 0
}
{ } {
}
{ } {
}
{ } {
}
}
{}{}{}
untuk semua $k$. Kondisi ini ekuivalen dengan

\mavergleichskettedisp
{\vergleichskette
{ u_k'
}
{ =} { \sum_{j = 1}^n a_{kj} u_j
}
{ } {
}
{ } {
}
{ } {
}
}
{}{}{}
untuk semua $k$, yang tepat berarti bahwa $u$ memenuhi sistem matriks

\mavergleichskettedisp
{\vergleichskette
{ u'
}
{ =} { \left(  a_{kj}    \right) u
}
{ } {
}
{ } {
}
{ } {
}
}
{}{}{}
;
dengan demikian, pemetaan itu merupakan suatu solusi sistem persamaan diferensial linear tersebut.

\part{Unit 26}
\chapter[Kuliah 26]{Kuliah 26: Koneksi Levi--Civita}
\setcounter{section}{26}

  \zwischenueberschrift{Koneksi Levi--Civita pada himpunan terbuka}

Misalkan

\mavergleichskette
{\vergleichskette
{ U
}
{ \subseteq }{ \R^n
}
{  }{
}
{    }{
}
{   }{
}
}
{}{}{}
terbuka dan padanya diberikan suatu struktur Riemann melalui fungsi-fungsi $g_{ij}$. Kita mempunyai fungsi koordinat $x_i$ dan medan vektor konstan terkait $\partial_i$. Suatu koneksi linear pada bundel tangen
\mathl{U \times \R^n}{} dijelaskan oleh turunan

\mavergleichskettedisp
{\vergleichskette
{ \nabla_{\partial_i } \partial_j
}
{ =} { \sum_{k = 1}^n\Gamma_{ij}^k\,\partial_k
}
{ } {
}
{ } {
}
{ } {
}
}
{}{}{}
Berikut ini akan terlihat bahwa pilihan berikut menghasilkan suatu koneksi yang berkaitan erat dengan struktur Riemann.

\inputdefinition
{}
{

Misalkan

\mavergleichskette
{\vergleichskette
{ U
}
{ \subseteq }{ \R^n
}
{  }{
}
{    }{
}
{   }{
}
}
{}{}{}
terbuka dan padanya diberikan suatu struktur Riemann melalui fungsi-fungsi $\left(  g_{ij}  \right)$ dengan matriks invers
\mathl{\left(  g^{k l}  \right)}{.} Fungsi-fungsi bernilai riil yang didefinisikan pada $U$,

\mavergleichskettedisp
{\vergleichskette
{ \Gamma_{ij}^k
}
{ =} { { \frac{   1  }{  2 } } \sum_{l = 1}^n g^{kl}(\partial_ig_{jl}+\partial_jg_{il}-\partial_lg_{ij})
}
{ } {
}
{ } {
}
{ } {
}
}
{}{}{}
disebut
\definitionswort {simbol Christoffel}{}
untuk koneksi Levi--Civita.

}

\inputdefinition
{}
{

Misalkan

\mavergleichskette
{\vergleichskette
{ U
}
{ \subseteq }{ \R^n
}
{  }{
}
{    }{
}
{   }{
}
}
{}{}{}
terbuka dan padanya diberikan suatu struktur Riemann melalui fungsi-fungsi $\left(  g_{ij}  \right)$ dengan matriks invers
\mathl{\left(  g^{k l}  \right)}{.} Pada $U$,
\definitionsverweis {simbol Christoffel}{}{}

\mavergleichskettedisp
{\vergleichskette
{ \Gamma_{ij}^k
}
{ =} { { \frac{   1  }{  2 } } \sum_{l = 1}^n g^{kl}(\partial_ig_{jl}+\partial_jg_{il}-\partial_lg_{ij})
}
{ } {
}
{ } {
}
{ } {
}
}
{}{}{}
menentukan suatu
\definitionsverweis {koneksi linear}{}{}
pada

\mavergleichskette
{\vergleichskette
{ TU
}
{ = }{ U \times \R^n
}
{  }{
}
{    }{
}
{   }{
}
}
{}{}{}
yang disebut
\definitionswort {koneksi Levi--Civita}{}
pada $U$.

}



\inputfaktbeweis
{R^n/Offene Menge/Levi-Civita-Zusammenhang/Trivial/Fakt}
{Lemma}
{}
{

\faktsituation {Misalkan

\mavergleichskette
{\vergleichskette
{ U
}
{ \subseteq }{ \R^n
}
{  }{
}
{    }{
}
{   }{
}
}
{}{}{}
suatu
\definitionsverweis {himpunan bagian terbuka}{}{,}
yang dilengkapi dengan
\definitionsverweis {struktur Riemann}{}{}
terinduksi dari $\R^n$.}
\faktfolgerung {Maka
\definitionsverweis {koneksi Levi--Civita}{}{}
pada $U$
\definitionsverweis {trivial}{}{.}
Secara khusus, berlaku pernyataan-pernyataan berikut.
\aufzaehlungvier{Semua
\definitionsverweis {simbol Christoffel}{}{}
memenuhi

\mavergleichskettedisp
{\vergleichskette
{ \Gamma_{ij}^k
}
{ =} { 0
}
{ } {
}
{ } {
}
{ } {
}
}
{}{}{}
untuk setiap $i,j,k$.
}{Berlaku

\mavergleichskettedisp
{\vergleichskette
{ \nabla_{\partial_i }  \partial_j
}
{ =} { 0
}
{ } {
}
{ } {
}
{ } {
}
}
{}{}{}
untuk
\definitionsverweis {medan vektor standar}{}{}
$\partial_i$.
}{Untuk fungsi-fungsi terdiferensialkan $f_j$ berlaku

\mavergleichskettedisp
{\vergleichskette
{ \nabla_{\partial_i }  \left(  \sum_{j = 1}^n f_j \partial_j  \right)
}
{ =} { \sum_{j = 1}^n \partial_i(f_j) \partial_j
}
{ } {
}
{ } {
}
{ } {
}
}
{}{}{.}
}{Untuk suatu medan vektor $V$ dan suatu medan vektor terdiferensialkan $W$ pada $U$, berlaku

\mavergleichskettedisp
{\vergleichskette
{ ( \nabla_V W)(P)
}
{ =} { \left(  D W   \right)_{ P } \left(   V(P)  \right)
}
{ } {
}
{ } {
}
{ } {
}
}
{}{}{.}
}}
\faktzusatz {}
\faktzusatz {}

}
{

Fungsi-fungsi fundamental Riemann $g_{ij}$ semuanya konstan, sehingga turunan parsialnya yang muncul dalam definisi simbol Christoffel lenyap. Karena itu, menurut
Catatan 25.8,
koneksi tersebut adalah koneksi trivial. Sifat-sifat lainnya mengikuti dari
Latihan 25.11.

}

\inputbeispiel{}
{

Misalkan

\mavergleichskette
{\vergleichskette
{ I
}
{ \subseteq }{ \R
}
{  }{
}
{    }{
}
{   }{
}
}
{}{}{}
suatu
\definitionsverweis {interval riil}{}{}
dan misalkan
\maabbdisp {g} { I } {\R_+
} {}
suatu fungsi positif terdiferensialkan, yang kita tafsirkan sebagai
\definitionsverweis {metrik Riemann}{}{}
pada $I$; lihat
Contoh 16.10.
Satu-satunya
\definitionsverweis {simbol Christoffel}{}{}
untuk koneksi Levi--Civita ialah

\mavergleichskettedisp
{\vergleichskette
{ \Gamma
}
{ =} { { \frac{   1  }{  2 } }  g^{-1} g'
}
{ } {
}
{ } {
}
{ } {
}
}
{}{}{,}
yakni, selain faktor di depannya, merupakan
\definitionsverweis {turunan logaritmik}{}{}
dari $g$. Turunan vertikalnya adalah

\mavergleichskettedisp
{\vergleichskette
{ \nabla_\partial \partial
}
{ =} { { \frac{  g' }{  2g } } \partial
}
{ } {
}
{ } {
}
{ } {
}
}
{}{}{,}
dengan $\partial$ menyatakan turunan pada $I$. Jika diterapkan pada suatu fungsi yang terdiferensialkan secara kontinu
\maabb {f} { I } {\R
} {}
\zusatzklammer {atau, setara dengan itu, pada medan arah $f \partial$} {} {,}
maka menurut
Teorema 25.5
berlaku

\mavergleichskettedisp
{\vergleichskette
{ \nabla_\partial (f)
}
{ =} { \partial f + f  { \frac{  g' }{  2g } }
}
{ =} { f' + f { \frac{  g' }{  2g } }
}
{ } {
}
{ } {
}
}
{}{}{.}

}



\inputfaktbeweis
{Riemannsche Mannigfaltigkeit/Zweidimensional/Levi-Civita-Zusammenhang/Christoffelsymbole/Beziehung/Fakt}
{Lemma}
{}
{

\faktsituation {Misalkan

\mavergleichskette
{\vergleichskette
{ U
}
{ \subseteq }{ \R^2
}
{  }{
}
{    }{
}
{   }{
}
}
{}{}{}
\definitionsverweis {terbuka}{}{}
dan dilengkapi dengan suatu struktur Riemann yang diberikan oleh
\definitionsverweis {fungsi-fungsi terdiferensialkan}{}{}
\maabbdisp {g_{11},g_{12},g_{22}} { U } {\R
} {}
serta
\definitionsverweis {struktur Riemann}{}{}
tersebut mempunyai matriks invers $\left(  g^{kl}  \right)$.}
\faktfolgerung {Maka
\definitionsverweis {simbol Christoffel}{}{}
untuk koneksi Levi--Civita pada $U$ adalah

\mavergleichskettedisphandlinks
{\vergleichskettedisphandlinks
{ \Gamma_{ij}^k
}
{ =} { { \frac{   1  }{  2 } } \left(  g^{k1}(\partial_ig_{j1}+\partial_jg_{i1}-\partial_1 g_{ij})   +  g^{k2}(\partial_ig_{j2}+\partial_jg_{i2}-\partial_2 g_{ij})  \right)
}
{ } {
}
{ } {
}
{ } {
}
}
{}{}{.}
Jika

\mavergleichskette
{\vergleichskette
{ Y
}
{ \subseteq }{ \R^3
}
{  }{
}
{    }{
}
{   }{
}
}
{}{}{}
suatu permukaan berorientasi yang dua kali terdiferensialkan secara kontinu dan dilengkapi dengan struktur yang diinduksi oleh ruang sekitar $\R^3$ sebagai
\definitionsverweis {struktur Riemann}{}{,}
serta
\maabbdisp {\varphi} { V } { Y
} {,}

\mavergleichskette
{\vergleichskette
{ V
}
{ \subseteq }{ \R^2
}
{  }{
}
{    }{
}
{   }{
}
}
{}{}{}
terbuka, merupakan suatu parametrisasi lokal $Y$ yang dua kali terdiferensialkan, maka simbol-simbol Christoffel ini berimpit dengan
\definitionsverweis {simbol Christoffel}{}{}
yang didefinisikan dengan bantuan ruang sekitar.}
\faktzusatz {}
\faktzusatz {}

}
{

Pernyataan pertama langsung mengikuti dari
Definisi 26.1.
Pernyataan kedua mengikuti dari pernyataan pertama dan
Lemma 19.5.

}



\inputfaktbeweis
{Offene Menge/Riemannsche Struktur/Christoffelsymbole/Eigenschaften/Fakt}
{Lemma}
{}
{

\faktsituation {Misalkan

\mavergleichskette
{\vergleichskette
{ U
}
{ \subseteq }{ \R^n
}
{  }{
}
{    }{
}
{   }{
}
}
{}{}{}
terbuka dan padanya diberikan suatu
\definitionsverweis {struktur Riemann}{}{}
melalui fungsi-fungsi $\left(  g_{ij}  \right)$. Maka
\definitionsverweis {simbol Christoffel}{}{}
untuk koneksi Levi--Civita pada $U$ memenuhi

\mavergleichskettedisp
{\vergleichskette
{ \Gamma^k_{ij}
}
{ =} { \Gamma^k_{ji}
}
{ } {
}
{ } {
}
{ } {
}
}
{}{}{.}}
\faktfolgerung {}
\faktzusatz {}
\faktzusatz {}

}
{

Hal ini langsung mengikuti dari definisi dan kesimetrian
\mathl{\left(  g_{ij}  \right)}{.}

}

\inputbeispiel{}
{

Kita melanjutkan
Contoh 18.8,
yakni kita meninjau setengah bidang Poincaré $\mathbb H$ dengan fungsi-fungsi fundamental Riemann

\mavergleichskette
{\vergleichskette
{ g_{11}
}
{ = }{ g_{22}
}
{ = }{ { \frac{   1  }{  y^2 } }
}
{    }{
}
{   }{
}
}
{}{}{}
dan

\mavergleichskette
{\vergleichskette
{ g_{12}
}
{ = }{ 0
}
{  }{
}
{    }{
}
{   }{
}
}
{}{}{.}
Menurut
Lemma 26.5,
dari

\mathdisp {{ \frac{   1  }{  2 } } \left(  g^{k1}(\partial_ig_{j1}+\partial_jg_{i1}-\partial_1 g_{ij})   +  g^{k2}(\partial_ig_{j2}+\partial_jg_{i2}-\partial_2 g_{ij})  \right)} {  }

diperoleh

\mavergleichskettedisp
{\vergleichskette
{ \Gamma^1_{11}
}
{ =} { { \frac{   1  }{  2 } }  \left(  y^2  \partial_1  \left( \frac{   1  }{  y^2 } \right)     \right)
}
{ =} { 0
}
{ } {
}
{ } {}
}
{}{}{,}

\mavergleichskettedisp
{\vergleichskette
{ \Gamma^2_{11}
}
{ =} { { \frac{   1  }{  2 } }  \left(  y^2 \left(  - \partial_2  \left( \frac{   1  }{  y^2 } \right)  \right)      \right)
}
{ =} { { \frac{   1  }{  2 } }  \left(  y^2 \left(  -  \left( \frac{   -2 }{  y^3 } \right)  \right)      \right)
}
{ =} { { \frac{   1  }{  y } }
}
{ } {
}
}
{}{}{}
dan dengan cara serupa

\mavergleichskette
{\vergleichskette
{ \Gamma^1_{12}
}
{ = }{ - { \frac{   1  }{  y } }
}
{  }{
}
{    }{
}
{   }{
}
}
{}{}{,}

\mavergleichskette
{\vergleichskette
{ \Gamma^2_{12}
}
{ = }{ 0
}
{  }{
}
{    }{
}
{   }{
}
}
{}{}{,}

\mavergleichskette
{\vergleichskette
{ \Gamma^1_{22}
}
{ = }{ 0
}
{  }{
}
{    }{
}
{   }{
}
}
{}{}{,}

\mavergleichskette
{\vergleichskette
{ \Gamma^2_{22}
}
{ = }{ - { \frac{   1  }{  y } }
}
{  }{
}
{    }{
}
{   }{
}
}
{}{}{.}

}

  \zwischenueberschrift{Koneksi Levi--Civita}

Kita hendak memperluas konsep koneksi Levi--Civita dari suatu himpunan terbuka

\mavergleichskette
{\vergleichskette
{ U
}
{ \subseteq }{ \R^n
}
{  }{
}
{    }{
}
{   }{
}
}
{}{}{}
dengan struktur Riemann ke suatu manifold Riemann sembarang. Kesulitan utamanya adalah menunjukkan bahwa konstruksi langsung pada daerah-daerah peta terbuka saling konsisten pada irisan yang tumpang tindih. Untuk itu, kita memperkenalkan konsep-konsep berikut.

\inputdefinition
{}
{

Misalkan $M$ suatu
\definitionsverweis {manifold terdiferensialkan}{}{}
dan diberikan suatu
\definitionsverweis {koneksi linear}{}{}
pada
\definitionsverweis {bundel tangen}{}{}
$TM$. Koneksi tersebut disebut
\definitionswort {bebas torsi}{,}
jika

\mavergleichskettedisp
{\vergleichskette
{ \nabla_V W- \nabla_W V
}
{ =} { [V,W]
}
{ } {
}
{ } {
}
{ } {
}
}
{}{}{}
untuk setiap
\definitionsverweis {medan vektor}{}{}
$V,W$ yang
\definitionsverweis {terdiferensialkan secara kontinu}{}{}
pada setiap himpunan terbuka.

}

\inputdefinition
{}
{

Misalkan $M$ suatu
\definitionsverweis {manifold Riemann}{}{}
dan diberikan suatu
\definitionsverweis {koneksi linear}{}{}
pada
\definitionsverweis {bundel tangen}{}{}
$TM$. Koneksi tersebut disebut
\definitionswort {metrik}{,}
jika

\mavergleichskettedisp
{\vergleichskette
{ D_V \left\langle  W  ,  Z  \right\rangle
}
{ =} { \left\langle \nabla_V W , Z  \right\rangle  + \left\langle  W  ,  \nabla_VZ  \right\rangle
}
{ } {
}
{ } {
}
{ } {
}
}
{}{}{.}
untuk setiap
\definitionsverweis {medan vektor}{}{}
$V,W,Z$ yang
\definitionsverweis {terdiferensialkan secara kontinu}{}{}
pada setiap himpunan terbuka.

}



\inputfaktbeweis
{Riemannsche Mannigfaltigkeit/Tangentialbündel/Linearer Zusammenhang/Metrisch und torsionsfrei/Koszul-Eigenschaft/Eindeutigkeit/Fakt}
{Lemma}
{}
{

\faktsituation {Misalkan $M$ suatu
\definitionsverweis {manifold Riemann}{}{}
dan diberikan suatu
\definitionsverweis {koneksi linear}{}{}
$\nabla$ pada
\definitionsverweis {bundel tangen}{}{}
$TM$ yang
\definitionsverweis {metrik}{}{}
dan
\definitionsverweis {bebas torsi}{}{.}}
\faktfolgerung {Maka untuk setiap
\definitionsverweis {medan vektor}{}{}
$V,W,Z$ yang
\definitionsverweis {terdiferensialkan secara kontinu}{}{}
pada setiap himpunan terbuka berlaku apa yang disebut \stichwort {rumus Koszul} {}

\mavergleichskettedisphandlinks
{\vergleichskettedisphandlinks
{ \left\langle \nabla_V W , Z  \right\rangle
}
{ =} { { \frac{   1  }{  2 } }  \left(  D_V \left\langle  W  ,  Z  \right\rangle +  D_W \left\langle  Z  ,  V  \right\rangle -  D_Z \left\langle  V  ,  W  \right\rangle  - \left\langle  W  , [V,Z]   \right\rangle - \left\langle  Z  , [W,V]   \right\rangle + \left\langle  V  , [Z,W]   \right\rangle   \right)
}
{ } {
}
{ } {
}
{ } {
}
}
{}{}{.}}
\faktzusatz {Secara khusus, paling banyak ada satu koneksi dengan sifat-sifat ini.}
\faktzusatz {}

}
{

Karena koneksi bebas torsi, berlaku

\mavergleichskettedisp
{\vergleichskette
{ \nabla_VW- \nabla_WV
}
{ =} { [V,W]
}
{ } {
}
{ } {
}
{ } {
}
}
{}{}{,}

\mavergleichskettedisp
{\vergleichskette
{ \nabla_VZ- \nabla_ZV
}
{ =} { [V,Z]
}
{ } {
}
{ } {
}
{ } {
}
}
{}{}{}
dan

\mavergleichskettedisp
{\vergleichskette
{ \nabla_WZ- \nabla_ZW
}
{ =} { [W,Z]
}
{ } {
}
{ } {
}
{ } {
}
}
{}{}{,}
dan identitas pertama juga dapat ditulis sebagai

\mavergleichskettedisp
{\vergleichskette
{ \nabla_VW + \nabla_WV
}
{ =} { 2 \nabla_VW -[V,W]
}
{ } {
}
{ } {
}
{ } {
}
}
{}{}{.}
Karena koneksi metrik, berlaku identitas-identitas

\mavergleichskettedisp
{\vergleichskette
{ D_V \left\langle  W  ,  Z  \right\rangle
}
{ =} { \left\langle \nabla_V W , Z  \right\rangle  + \left\langle  W  ,  \nabla_V Z  \right\rangle
}
{ } {
}
{ } {
}
{ } {
}
}
{}{}{,}

\mavergleichskettedisp
{\vergleichskette
{ D_W \left\langle  Z  ,  V  \right\rangle
}
{ =} { \left\langle \nabla_WZ , V  \right\rangle  + \left\langle  Z  ,  \nabla_WV  \right\rangle
}
{ } {
}
{ } {
}
{ } {
}
}
{}{}{}
dan

\mavergleichskettedisp
{\vergleichskette
{ D_Z \left\langle  V  ,  W  \right\rangle
}
{ =} { \left\langle \nabla_Z V , W  \right\rangle  + \left\langle  V  ,  \nabla_Z W   \right\rangle
}
{ } {
}
{ } {
}
{ } {
}
}
{}{}{.}
Dengan menjumlahkan persamaan-persamaan ini, kita memperoleh

\mavergleichskettedisphandlinks
{\vergleichskettedisphandlinks
{ D_V \left\langle  W  ,  Z  \right\rangle+ D_W \left\langle  Z  ,  V  \right\rangle  - D_Z \left\langle  V  ,  W  \right\rangle
}
{ =} { \left\langle \nabla_V W+\nabla_W V  ,  Z  \right\rangle + \left\langle \nabla_W Z-\nabla_Z W   ,  V  \right\rangle+\left\langle  \nabla_V Z-\nabla_Z V   ,  W  \right\rangle
}
{ } {
}
{ } {
}
{ } {
}
}
{}{}{.}
Kita substitusikan ungkapan-ungkapan di atas ke ruas kanan dan memperoleh

\mavergleichskettealignhandlinks
{\vergleichskettealignhandlinks
{ D_V \left\langle  W  ,  Z  \right\rangle+ D_W \left\langle  Z  ,  V  \right\rangle  - D_Z \left\langle  V  ,  W  \right\rangle
}
{ =} { \left\langle  2 \nabla_VW -[V,W]  ,  Z  \right\rangle + \left\langle  [W,Z]  ,  V  \right\rangle + \left\langle  [V,Z]   ,  W  \right\rangle
}
{ =} { 2\left\langle  \nabla_VW   ,  Z  \right\rangle - \left\langle  [V,W]  ,  Z  \right\rangle + \left\langle  [W,Z]  ,  V  \right\rangle + \left\langle  [V,Z]   ,  W  \right\rangle
}
{ } {
}
{ } {
}
}
{}
{}{.}
Dengan menata ulang persamaan tersebut, diperoleh rumus yang dinyatakan.

Berdasarkan persamaan itu,
\mathl{\left\langle  \nabla_V W , Z  \right\rangle}{} untuk medan-medan vektor sembarang telah ditentukan oleh ruas kanan, yang sama sekali tidak memuat koneksi. Karena hal ini berlaku untuk setiap medan vektor $Z$, turunan vertikal $\nabla_V W$, dan dengan demikian koneksi linearnya, juga ditentukan secara unik.

}

\inputfaktbeweis
{R^n/Offene Menge/Riemannsche Struktur/Zusammenhang aus Christoffel-Symbole/Eigenschaften/Fakt}
{Lemma}
{}
{

\faktsituation {Misalkan

\mavergleichskette
{\vergleichskette
{ U
}
{ \subseteq }{ \R^n
}
{  }{
}
{    }{
}
{   }{
}
}
{}{}{}
\definitionsverweis {terbuka}{}{}
dan padanya diberikan suatu
\definitionsverweis {struktur Riemann}{}{}
melalui fungsi-fungsi $\left(  g_{ij}  \right)$ dengan
\definitionsverweis {matriks invers}{}{}

\mathl{\left(  g^{k l}  \right)}{.} Misalkan $\nabla$ adalah koneksi yang diberikan oleh
\definitionsverweis {simbol-simbol Christoffel}{}{}

\mavergleichskettedisp
{\vergleichskette
{ \Gamma_{ij}^k
}
{ =} { { \frac{   1  }{  2 } } \sum_{l = 1}^n g^{kl}(\partial_ig_{jl}+\partial_jg_{il}-\partial_lg_{ij})
}
{ } {
}
{ } {
}
{ } {
}
}
{}{}{}
sebagai
\definitionsverweis {koneksi Levi--Civita}{}{,}
yang dicirikan oleh

\mavergleichskettedisp
{\vergleichskette
{ \nabla_{\partial_i } \partial_j
}
{ =} { \sum_{k = 1}^n\Gamma_{ij}^k\,\partial_k
}
{ } {
}
{ } {
}
{ } {
}
}
{}{}{.}}
\faktuebergang {Maka $\nabla$ memenuhi sifat-sifat berikut.}
\faktfolgerung {\aufzaehlungvier{Berlaku

\mavergleichskettedisphandlinks
{\vergleichskettedisphandlinks
{ \left\langle  \nabla_{\partial_i }  \partial_j  ,  \partial_k   \right\rangle
}
{ =} { { \frac{   1  }{  2 } }  \left(  D_{\partial_i }  \left\langle  \partial_j  ,  \partial_k   \right\rangle   +D_{\partial_j }  \left\langle  \partial_i  ,  \partial_k   \right\rangle -D_{\partial_k }  \left\langle  \partial_i  ,  \partial_j   \right\rangle   \right)
}
{ } {}
{ } {}
{ } {}
}
{}{}{.}
}{Koneksi ini memenuhi rumus Koszul dari
Lemma 26.10.
}{Koneksi ini
\definitionsverweis {bebas torsi}{}{.}
}{Koneksi ini
\definitionsverweis {metrik}{}{.}
}}
\faktzusatz {}
\faktzusatz {}

}
{

\aufzaehlungvier{Berlaku

\mavergleichskettedisp
{\vergleichskette
{ D_{\partial_k }  \left\langle  \partial_i  ,  \partial_j   \right\rangle
}
{ =} { \partial_k  \left(  g_{ij}  \right)
}
{ } {
}
{ } {
}
{ } {
}
}
{}{}{.}
Selanjutnya,

\mavergleichskettealignhandlinks
{\vergleichskettealignhandlinks
{ \left\langle  \nabla_{\partial_i }  \partial_j  ,  \partial_k   \right\rangle
}
{ =} { \left\langle  \sum_{\ell = 1}^n\Gamma_{ij}^\ell \,\partial_\ell  ,  \partial_k   \right\rangle
}
{ =} { \sum_{\ell = 1}^n\Gamma_{ij}^\ell \left\langle  \,\partial_\ell  ,  \partial_k   \right\rangle
}
{ =} { \sum_{\ell = 1}^n\Gamma_{ij}^\ell g_{\ell k}
}
{ =} { { \frac{   1  }{  2 } } \sum_{\ell = 1}^n  \left(  \sum_{r = 1}^n g^{\ell r}  \left(  \partial_ig_{j r}+\partial_jg_{ir}-\partial_r g_{ij}  \right)   \right)  g_{\ell k}
}
}
{
\vergleichskettefortsetzungalign
{ =} { { \frac{   1  }{  2 } } \sum_{r = 1}^n  \left(  \partial_ig_{j r}+\partial_jg_{ir}-\partial_r g_{ij}  \right)  \left(  \sum_{\ell = 1}^n  g^{\ell r} g_{\ell k}   \right)
}
{ =} { { \frac{   1  }{  2 } }  \left(  \partial_ig_{j k}+\partial_jg_{i k}-\partial_k g_{ij}  \right)
}
{ =} { { \frac{   1  }{  2 } }  \left(  D_{\partial_i }  \left\langle  \partial_j  ,  \partial_k   \right\rangle   +D_{\partial_j }  \left\langle  \partial_i  ,  \partial_k   \right\rangle   -D_{\partial_k }  \left\langle  \partial_i  ,  \partial_j   \right\rangle   \right)
}
{ } {}
}
{}{.}
}{Karena kurung Lie
\mathl{[ \partial_i, \partial_j]}{} trivial pada $U$, menurut bagian (1) rumus Koszul berlaku untuk medan-medan basis $\partial_i$. Untuk kasus umum, lihat
Latihan 26.7.
}{Menurut
Lemma 26.6
dan karena

\mavergleichskette
{\vergleichskette
{ [\partial_i, \partial_j]
}
{ = }{ 0
}
{  }{
}
{    }{
}
{   }{
}
}
{}{}{}
berlaku

\mavergleichskettedisp
{\vergleichskette
{ \nabla_{\partial_i }  \partial_j
}
{ =} { \nabla_{\partial_j }  \partial_i
}
{ } {
}
{ } {
}
{ } {
}
}
{}{}{.}
Pernyataan itu sekarang mengikuti dari
Latihan 26.5.
}{Menurut bagian (1),

\mavergleichskettealignhandlinks
{\vergleichskettealignhandlinks
{ \left\langle  \nabla_{\partial_i }  \partial_j  ,  \partial_k   \right\rangle + \left\langle  \nabla_{\partial_i }  \partial_k  ,  \partial_j   \right\rangle
}
{ =} { { \frac{   1  }{  2 } }  \left(  D_{\partial_i }  \left\langle  \partial_j  ,  \partial_k   \right\rangle   +D_{\partial_j }  \left\langle  \partial_i  ,  \partial_k   \right\rangle   -D_{\partial_k }  \left\langle  \partial_i  ,  \partial_j   \right\rangle    +D_{\partial_i }  \left\langle  \partial_k  ,  \partial_j   \right\rangle +D_{\partial_k }  \left\langle  \partial_i  ,  \partial_j   \right\rangle  - D_{\partial_j }  \left\langle  \partial_i  ,  \partial_k   \right\rangle   \right)
}
{ =} { D_{\partial_i }  \left\langle  \partial_j  ,  \partial_k   \right\rangle
}
{ } {}
{ } {}
}
{}
{}{.}
Pernyataan itu sekarang mengikuti dari
Latihan 26.8.
}

}



\inputfaktbeweis
{Riemannsche Mannigfaltigkeit/Levi-Civita-Zusammenhang/Eindeutige Existenz/Fakt}
{Teorema}
{}
{

\faktsituation {Misalkan $X$ suatu
\definitionsverweis {manifold Riemann}{}{.}
Maka pada
\definitionsverweis {bundel tangen}{}{}
$TX$ terdapat suatu
\definitionsverweis {koneksi linear}{}{}
yang
\definitionsverweis {bebas torsi}{}{,}
\definitionsverweis {metrik}{}{,}
dan ditentukan secara unik.}
\faktfolgerung {}
\faktzusatz {}
\faktzusatz {}

}
{

Menurut
Lemma 26.10,
paling banyak ada satu koneksi semacam itu. Menurut
Lemma 26.11,
pada setiap daerah peta kita dapat memberikan secara eksplisit suatu koneksi yang mempunyai sifat-sifat yang disyaratkan. Karena keunikan yang baru saja disebutkan, koneksi-koneksi yang dikonstruksi dengan cara ini berimpit pada irisan daerah-daerah peta.

}

\inputdefinition
{}
{

Misalkan $X$ suatu
\definitionsverweis {manifold Riemann}{}{.}
\definitionsverweis {Koneksi}{}{}
pada
\definitionsverweis {bundel tangen}{}{}
$TX$ yang didefinisikan secara lokal oleh
\definitionsverweis {simbol-simbol Christoffel}{}{}
disebut
\definitionswort {koneksi Levi--Civita}{.}

}



\inputfaktbeweis
{Riemannsche Mannigfaltigkeit/Levi-Civita-Zusammenhang/Eigenschaften/Fakt}
{Teorema}
{}
{

\faktsituation {Misalkan $X$ suatu
\definitionsverweis {manifold Riemann}{}{.}
Maka
\definitionsverweis {koneksi Levi--Civita}{}{}
$\nabla$ memenuhi sifat-sifat berikut.}
\faktfolgerung {\aufzaehlungdrei{ $\nabla$
\definitionsverweis {linear}{}{.}
}{ $\nabla$ metrik, yakni

\mavergleichskettedisp
{\vergleichskette
{ D_V \left\langle  W  ,  Z  \right\rangle
}
{ =} { \left\langle \nabla_V W , Z  \right\rangle  + \left\langle  W  ,  \nabla_VZ  \right\rangle
}
{ } {
}
{ } {
}
{ } {
}
}
{}{}{.}
}{ $\nabla$ bebas torsi, yakni

\mavergleichskettedisp
{\vergleichskette
{ \nabla_V W- \nabla_W V
}
{ =} { [V,W]
}
{ } {
}
{ } {
}
{ } {
}
}
{}{}{.}
}}
\faktzusatz {}
\faktzusatz {}

}
{

Hal ini mengikuti dari
Lemma 26.11.

}



\inputfaktbeweis
{R^3/Riemannsche orientierte Fläche/Levi-Civita-Zusammenhang und induzierter Zusammenhang/Fakt}
{Lemma}
{}
{

\faktsituation {Misalkan

\mavergleichskette
{\vergleichskette
{ Y
}
{ \subseteq }{ \R^3
}
{  }{
}
{    }{
}
{   }{
}
}
{}{}{}
suatu permukaan yang dua kali terdiferensialkan dan dilengkapi oleh ruang sekitar $\R^3$ dengan
\definitionsverweis {struktur Riemann}{}{}
terinduksi.}
\faktfolgerung {Maka
\definitionsverweis {koneksi Levi--Civita}{}{}
pada $Y$ berimpit dengan apa yang
\zusatzklammer {melalui
\definitionsverweis {proyeksi ortogonal}{}{}
diberikan pada ruang-ruang tangen} {} {}
didefinisikan sebagai
\definitionsverweis {koneksi}{}{}
dari
Catatan 24.8.}
\faktzusatz {}
\faktzusatz {}

}
{

Kedua koneksi itu linear. Karena itu, cukup ditunjukkan secara lokal bahwa
\definitionsverweis {turunan vertikal}{}{}
terkait sama pada medan-medan basis. Misalkan

\mavergleichskette
{\vergleichskette
{ V
}
{ \subseteq }{ \R^2
}
{  }{
}
{    }{
}
{   }{
}
}
{}{}{}
\definitionsverweis {terbuka}{}{}
dan
\maabbdisp {\varphi} { V } { U
} {}
suatu parametrisasi lokal yang dua kali terdiferensialkan secara kontinu dari sebuah himpunan bagian terbuka

\mavergleichskette
{\vergleichskette
{ U
}
{ \subseteq }{ Y
}
{  }{
}
{    }{
}
{   }{
}
}
{}{}{.}
Misalkan $\partial_1,\partial_2$ adalah kedua medan arah konstan pada $V$. Untuk koneksi Levi--Civita berlaku

\mavergleichskettedisp
{\vergleichskette
{ \nabla_{\partial_i } \partial_j
}
{ =} { \sum_{k = 1}^2  \Gamma^k_{ij} \partial_k
}
{ } {
}
{ } {
}
{ } {
}
}
{}{}{,}
dengan simbol-simbol Christoffel yang didefinisikan berdasarkan

\mavergleichskettedisp
{\vergleichskette
{ g_{ij}
}
{ =} { \left\langle  \partial_i \varphi ,  \partial_j \varphi  \right\rangle
}
{ } {
}
{ } {
}
{ } {
}
}
{}{}{.}
Menurut
Lemma 26.5,
simbol-simbol Christoffel ini juga memenuhi

\mavergleichskettedisp
{\vergleichskette
{ \partial_i \partial_j \varphi
}
{ =} { \Gamma^1_{ij} \partial_1  \varphi +  \Gamma^2_{ij}  \partial_2  \varphi+   h_{ij} N
}
{ } {
}
{ } {
}
{ } {
}
}
{}{}{,}
dengan $N$ menyatakan
\definitionsverweis {medan normal satuan}{}{}
dan $h_{ij}$ merupakan entri-entri
\definitionsverweis {matriks fundamental kedua}{}{}
\zusatzklammer {semuanya dipandang pada $V$} {} {.}

Medan-medan vektor $\partial_i$ pada $V$ bersesuaian, pada

\mavergleichskette
{\vergleichskette
{ U
}
{ \subseteq }{ Y
}
{ \subseteq }{ \R^3
}
{    }{
}
{   }{
}
}
{}{}{}
dengan medan-medan vektor $\left(  \partial_i \varphi  \right)  \circ \varphi^{-1}$. Menurut definisi
\definitionsverweis {turunan vertikal}{}{}
untuk koneksi terbatas tersebut, kita harus meninjau

\mathdisp {U \stackrel{ \left(  \partial_i\varphi  \right) \circ \varphi^{-1} }{\longrightarrow} TU \stackrel{ T \left(  \left(  \partial_j \varphi  \right) \circ \varphi^{-1}    \right) }{\longrightarrow} TTU \stackrel{ \pi_{\text{vert} } } \longrightarrow TU} {  }

Ini menghasilkan proyeksi ortogonal
\mathl{\left(  \partial_i \partial_j \varphi  \right) \circ \varphi^{-1}}{}
ke bundel tangen pada $Y$, yang berimpit dengan
\mathl{\Gamma^1_{ij} \partial_1  \varphi +  \Gamma^2_{ij}  \partial_2  \varphi}{,}
dipandang di atas $U$
\zusatzklammer {lihat juga
Latihan 26.11} {} {.}

}

Pernyataan ini juga berlaku dengan cara yang sama untuk setiap submanifold tertutup

\mavergleichskette
{\vergleichskette
{ Y
}
{ \subseteq }{ \R^n
}
{  }{
}
{    }{
}
{   }{
}
}
{}{}{}
yang dilengkapi dengan struktur Riemann terinduksi.

\chapter{Lembar Kerja 26}
\setcounter{section}{26}

  \zwischenueberschrift{Soal latihan}

\inputaufgabe
{}
{

Misalkan bidang real dilengkapi dengan
\definitionsverweis {struktur Euklides}{}{}
dan
\definitionsverweis {koneksi Levi--Civita}{}{}
yang bersesuaian. Hitung

\mavergleichskette
{\vergleichskette
{ \nabla_FG
}
{ = }{
}
{  }{
}
{    }{
}
{   }{
}
}
{}{}{}
untuk medan-medan vektor

\mavergleichskettedisp
{\vergleichskette
{ F(x,y)
}
{ =} { \left(  x^2-y^3  , \,   xy+y^2                   \right)
}
{ } {
}
{ } {
}
{ } {
}
}
{}{}{}
dan

\mavergleichskettedisp
{\vergleichskette
{ G(x,y)
}
{ =} { \left( e^{xy^2}   , \,   x^3-x^2y^5                   \right)
}
{ } {
}
{ } {
}
{ } {
}
}
{}{}{.}

}
{} {}

\inputaufgabe
{}
{

Misalkan

\mavergleichskette
{\vergleichskette
{ I
}
{ \subseteq }{ \R
}
{  }{
}
{    }{
}
{   }{
}
}
{}{}{}
suatu
\definitionsverweis {interval real}{}{}
dan misalkan

\mavergleichskettedisp
{\vergleichskette
{ g(t)
}
{ =} { 3+t^2
}
{ } {
}
{ } {
}
{ } {
}
}
{}{}{,}
yang kita tafsirkan sebagai suatu
\definitionsverweis {metrik Riemann}{}{}
pada $I$. Tentukan
\definitionsverweis {turunan vertikal}{}{}
\zusatzklammer {terhadap
\definitionsverweis {koneksi Levi--Civita}{}{}
yang bersesuaian} {} {}

\mathl{\nabla_\partial  \left(  t^3-  \sin  \left( t^2 \right)   \right)}{.}

}
{} {}

\inputaufgabegibtloesung
{}
{

Misalkan

\mavergleichskette
{\vergleichskette
{ I
}
{ \subseteq }{ \R
}
{  }{
}
{    }{
}
{   }{
}
}
{}{}{}
suatu
\definitionsverweis {interval real}{}{}
dan misalkan
\maabbdisp {g} { I } {\R_+
} {}
suatu fungsi positif diferensiabel yang kita tafsirkan sebagai
\definitionsverweis {metrik Riemann}{}{}
pada $I$. Tentukan
\definitionsverweis {penampang horizontal}{}{}
terhadap
\definitionsverweis {koneksi Levi--Civita}{}{}
yang bersesuaian pada $I \times \R$
\zusatzklammer {lihat
Contoh 26.4} {} {.}

}
{} {}

\inputaufgabe
{}
{

Misalkan $M$ suatu
\definitionsverweis {manifold diferensiabel}{}{}
sedemikian sehingga
\definitionsverweis {bundel tangen}{}{}
$TM$ trivial, dan misalkan pada $M$, melalui suatu trivialisasi terdiferensialkan

\mavergleichskette
{\vergleichskette
{ TM
}
{ \cong }{ M \times \R^n
}
{  }{
}
{    }{
}
{   }{
}
}
{}{}{}
\zusatzklammer {dengan menggunakan
\definitionsverweis {hasil kali skalar standar}{}{}
pada $\R^n$} {} {}
diberikan suatu
\definitionsverweis {struktur Riemann}{}{}
\zusatzklammer {lihat
Soal 16.2} {} {.}
Tunjukkan bahwa
\definitionsverweis {koneksi Levi--Civita}{}{}
pada $M$ adalah
\definitionsverweis {koneksi trivial}{}{}
tersebut.

}
{} {}

\inputaufgabe
{}
{

Misalkan ${\mathbb H}$ adalah
\definitionsverweis {setengah bidang Poincaré}{}{}
dengan
\definitionsverweis {struktur Riemann}{}{}
dan
\definitionsverweis {koneksi Levi--Civita}{}{}
yang bersesuaian. Tentukan
\definitionsverweis {turunan vertikal}{}{}
\mathl{\nabla_{\partial_1 + \partial_2 } (G)}{} dari medan vektor

\mavergleichskettedisp
{\vergleichskette
{ G(x,y)
}
{ =} { \left(  x^5-y^3+2xy  , \,   4 x^2+ y^3                   \right)
}
{ } {
}
{ } {
}
{ } {
}
}
{}{}{.}

}
{} {}

\inputaufgabegibtloesung
{}
{

Misalkan diberikan suatu
\definitionsverweis {koneksi linear}{}{}
$\nabla$ pada
\definitionsverweis {bundel tangen}{}{}
$TX$ dari suatu
\definitionsverweis {manifold diferensiabel}{}{}
$X$. Andaikan terdapat
\definitionsverweis {penampang-penampang basis}{}{}

\mathl{V_1 , \ldots , V_n}{} dengan

\mavergleichskettedisp
{\vergleichskette
{ \nabla_{V_i } V_j -  \nabla_{V_j } V_i
}
{ =} { [V_i,V_j]
}
{ } {
}
{ } {
}
{ } {
}
}
{}{}{}
untuk semua $i,j$. Tunjukkan bahwa $\nabla$
\definitionsverweis {bebas torsi}{}{.}

}
{} {}

\inputaufgabe
{}
{

Tunjukkan bahwa suatu
\definitionsverweis {koneksi linear}{}{}
pada
\definitionsverweis {bundel tangen}{}{}
dari interval terbuka

\mavergleichskette
{\vergleichskette
{ I
}
{ \subseteq }{ \R
}
{  }{
}
{    }{
}
{   }{
}
}
{}{}{}
bersifat
\definitionsverweis {bebas torsi}{}{.}

}
{} {}

\inputaufgabe
{}
{

Misalkan

\mavergleichskette
{\vergleichskette
{ U
}
{ \subseteq }{ \R^n
}
{  }{
}
{    }{
}
{   }{
}
}
{}{}{}
\definitionsverweis {terbuka}{}{}
dan padanya diberikan suatu
\definitionsverweis {struktur Riemann}{}{}
melalui fungsi-fungsi $\left(  g_{ij}  \right)$, dengan
\definitionsverweis {matriks invers}{}{}

\mathl{\left(  g^{k l}  \right)}{.} Misalkan $\nabla$ adalah koneksi yang diberikan oleh
\definitionsverweis {simbol-simbol Christoffel}{}{}

\mavergleichskettedisp
{\vergleichskette
{ \Gamma_{ij}^k
}
{ =} { { \frac{   1  }{  2 } } \sum_{l = 1}^n g^{kl}(\partial_ig_{jl}+\partial_jg_{il}-\partial_lg_{ij})
}
{ } {
}
{ } {
}
{ } {
}
}
{}{}{}
sebagai
\definitionsverweis {koneksi linear}{}{,}
yang dicirikan oleh

\mavergleichskettedisp
{\vergleichskette
{ \nabla_{\partial_i } \partial_j
}
{ =} { \sum_{k = 1}^n\Gamma_{ij}^k\,\partial_k
}
{ } {
}
{ } {
}
{ } {
}
}
{}{}{.}
Tunjukkan dalam dua langkah bahwa $\nabla$ memenuhi rumus Koszul dari
Lemma 26.10.
\aufzaehlungzwei {Untuk medan-medan basis $\partial_i$.
} {Untuk medan-medan vektor sembarang $V,W,Z$.
}

}
{} {}

\inputaufgabegibtloesung
{}
{

Misalkan $M$ suatu
\definitionsverweis {manifold Riemann}{}{}
dan $\nabla$ suatu
\definitionsverweis {koneksi linear}{}{}
pada $TM$. Andaikan bahwa $\nabla$, secara lokal untuk medan-medan basis $\partial_1  , \ldots , \partial_n$, memenuhi sifat

\mavergleichskettedisp
{\vergleichskette
{ \partial_i \left\langle  \partial_j  ,  \partial_k   \right\rangle
}
{ =} { \left\langle \nabla_{\partial_i } \partial_j   ,  \partial_k   \right\rangle  + \left\langle  \partial_j  ,  \nabla_{\partial_i } \partial_k    \right\rangle
}
{ } {
}
{ } {
}
{ } {
}
}
{}{}{}
untuk semua $i,j,k$. Tunjukkan bahwa $\nabla$
\definitionsverweis {metrik}{}{.}

}
{} {}

Konsep koneksi metrik tidak hanya terdapat pada bundel tangen suatu manifold Riemann, tetapi secara lebih umum juga pada bundel vektor Riemann.

Misalkan
\maabb {p} { E } { M
} {}
suatu
\definitionsverweis {bundel vektor Riemann}{}{}
di atas suatu
\definitionsverweis {manifold diferensiabel}{}{}
$M$. Suatu
\definitionsverweis {koneksi linear}{}{}
pada $E$ disebut
\definitionswort {metrik}{,}
jika pada setiap himpunan terbuka berlaku

\mavergleichskettedisp
{\vergleichskette
{ D_V \left\langle  s  , t  \right\rangle
}
{ =} { \left\langle \nabla_V s , t  \right\rangle  + \left\langle  s  ,  \nabla_Vt  \right\rangle
}
{ } {
}
{ } {
}
{ } {
}
}
{}{}{}
untuk suatu
\definitionsverweis {medan vektor}{}{}
$V$ yang
\definitionsverweis {terdiferensialkan secara kontinu}{}{}
dan penampang-penampang $s,t$ dari $E$ yang terdiferensialkan secara kontinu.

\inputaufgabe
{}
{

Misalkan

\mavergleichskette
{\vergleichskette
{ I
}
{ \subseteq }{ \R
}
{  }{
}
{    }{
}
{   }{
}
}
{}{}{}
suatu
\definitionsverweis {interval terbuka}{}{}
dan misalkan
\maabbdisp {p} {E = I \times \R^n } { I
} {}
bundel vektor trivial di atas $I$, yang kita lengkapi dengan
\definitionsverweis {hasil kali skalar standar}{}{}
dan kita pandang sebagai
\definitionsverweis {bundel vektor Riemann}{}{.}
Tunjukkan bahwa suatu
\definitionsverweis {koneksi linear}{}{}
$\nabla$ pada $E$ bersifat
\definitionsverweis {metrik}{}{}
jika dan hanya jika
\definitionsverweis {simbol-simbol Christoffel}{}{}
memenuhi syarat
\zusatzklammer {semuanya mengacu pada satu-satunya arah turunan $\partial$} {} {}

\mavergleichskettedisp
{\vergleichskette
{ \Gamma_j^k
}
{ =} { - \Gamma_k^j
}
{ } {
}
{ } {
}
{ } {
}
}
{}{}{}
untuk

\mavergleichskette
{\vergleichskette
{ 1
}
{ \leq }{ j,k
}
{ \leq }{ n
}
{    }{
}
{   }{
}
}
{}{}{.}

}
{} {}

\inputaufgabe
{}
{

Misalkan
\maabb {p} { E } { M
} {}
suatu
\definitionsverweis {bundel vektor Riemann}{}{}
pada suatu
\definitionsverweis {manifold diferensiabel}{}{}
$M$, dan misalkan $\nabla$ suatu
\definitionsverweis {koneksi linear}{}{}
pada $E$. Andaikan bahwa $\nabla$, secara lokal untuk
\definitionsverweis {medan-medan vektor standar}{}{}
$\partial_1  , \ldots , \partial_n$ dan
\definitionsverweis {penampang-penampang basis}{}{}

\mathl{s_1  , \ldots , s_r}{} dari $E$, memenuhi sifat

\mavergleichskettedisp
{\vergleichskette
{ \partial_i \left\langle  s_j  , s_k   \right\rangle
}
{ =} { \left\langle \nabla_{\partial_i } s_j   ,  s_k   \right\rangle  + \left\langle s_j  ,  \nabla_{\partial_i } s_k    \right\rangle
}
{ } {
}
{ } {
}
{ } {
}
}
{}{}{}
untuk semua $i,j,k$. Tunjukkan bahwa $\nabla$
\definitionsverweis {metrik}{}{.}

}
{} {}

Untuk soal berikut, bandingkan dengan
Soal 10.16.

\inputaufgabe
{}
{

Misalkan

\mavergleichskette
{\vergleichskette
{ Y
}
{ \subseteq }{ \R^n
}
{  }{
}
{    }{
}
{   }{
}
}
{}{}{}
suatu
\definitionsverweis {submanifold tertutup}{}{}
dan
\maabb {\varphi} { V } { U \subseteq Y
} {}
suatu parametrisasi $C^2$-difeomorfik dari suatu himpunan bagian terbuka

\mavergleichskette
{\vergleichskette
{ U
}
{ \subseteq }{ Y
}
{  }{
}
{    }{
}
{   }{
}
}
{}{}{}
oleh suatu himpunan bagian terbuka

\mavergleichskette
{\vergleichskette
{ V
}
{ \subseteq }{ \R^d
}
{  }{
}
{    }{
}
{   }{
}
}
{}{}{,}
di mana $\varphi$ juga kita pandang sebagai pemetaan ke $\R^n$. Tunjukkan bahwa terdapat diagram komutatif

\mathdisp {\begin{matrix}  V  & \stackrel{ \partial_i }{\longrightarrow} &  TV & \stackrel{ T( \partial_j) }{\longrightarrow} & TTV & \\ \!\!\!\!\! \,\, \, \varphi \downarrow & & \!\!\!\!\! T(\varphi) \downarrow & &  \downarrow TT(\varphi) \!\!\!\!\! & & \\  U  & \stackrel{  }{\longrightarrow} &  TU & \stackrel{  }{\longrightarrow} & TTU & \!\!\!\!\!  \\ \end{matrix}} {  }

seperti di atas. Bagaimana pemetaan-pemetaan pada diagonal dapat dideskripsikan?

}
{} {}

  \zwischenueberschrift{Tugas untuk dikumpulkan}

\inputaufgabe
{2}
{

Misalkan

\mavergleichskette
{\vergleichskette
{ I
}
{ \subseteq }{ \R
}
{  }{
}
{    }{
}
{   }{
}
}
{}{}{}
suatu
\definitionsverweis {interval real}{}{}
dan misalkan

\mavergleichskettedisp
{\vergleichskette
{ g(t)
}
{ =} { 1+t^4
}
{ } {
}
{ } {
}
{ } {
}
}
{}{}{,}
yang kita tafsirkan sebagai suatu
\definitionsverweis {metrik Riemann}{}{}
pada $I$. Tentukan
\definitionsverweis {turunan vertikal}{}{}
\zusatzklammer {terhadap
\definitionsverweis {koneksi Levi--Civita}{}{}
yang bersesuaian} {} {}

\mathl{\nabla_\partial \left(  3e^{-t^2} + t  \cos \left(  t  \right) - t^5  \right)}{.}

}
{} {}

\inputaufgabe
{2}
{

Hitung $\Gamma^1_{12},  \Gamma^2_{12},  \Gamma^1_{22},  \Gamma^2_{22}$ dalam
Contoh 18.8.

}
{} {}

\inputaufgabe
{3}
{

Misalkan ${\mathbb H}$ adalah
\definitionsverweis {setengah bidang Poincaré}{}{}
dengan
\definitionsverweis {struktur Riemann}{}{}
dan
\definitionsverweis {koneksi Levi--Civita}{}{}
yang bersesuaian. Tentukan
\definitionsverweis {turunan vertikal}{}{}
\mathl{\nabla_{xy\partial_1 -2y^3 \partial_2 } (G)}{} dari medan vektor

\mavergleichskettedisp
{\vergleichskette
{ G(x,y)
}
{ =} { \left(  x^4-2y^2  , \,   3 x^3-  \ln  \left(  1+x^2y^2 \right)                   \right)
}
{ } {
}
{ } {
}
{ } {
}
}
{}{}{.}

}
{} {}

\inputaufgabe
{2}
{

Misalkan $\nabla$ suatu
\definitionsverweis {koneksi bebas torsi}{}{}
\definitionsverweis {linear}{}{}
pada suatu himpunan terbuka

\mavergleichskette
{\vergleichskette
{ U
}
{ \subseteq }{ \R^n
}
{  }{
}
{    }{
}
{   }{
}
}
{}{}{.}
Oleh berapa banyak
\definitionsverweis {simbol Christoffel}{}{}
yang saling bebas koneksi itu ditentukan
\zusatzklammer {sebagai fungsi dari $n$} {} {}?

}
{} {}

\section*{Solusi yang disediakan oleh sumber}
\addcontentsline{toc}{section}{Solusi yang disediakan oleh sumber}
\subsection*{Solusi Soal 26.3}
Penampang horizontal dicirikan oleh lenyapnya pemetaan vertikalnya. Dalam situasi ini, menurut
contoh tersebut,
turunan vertikal sama dengan

\mavergleichskettedisp
{\vergleichskette
{ \nabla_\partial (f)
}
{ =} { f' + f { \frac{  g' }{  2g } }
}
{ } {
}
{ } {
}
{ } {}
}
{}{}{.}
Pertama-tama,

\mavergleichskette
{\vergleichskette
{ f
}
{ = }{ 0
}
{  }{
}
{    }{
}
{   }{
}
}
{}{}{}
merupakan suatu solusi. Sekarang misalkan

\mavergleichskette
{\vergleichskette
{ f(t_0)
}
{ \neq }{ 0
}
{  }{
}
{    }{
}
{   }{
}
}
{}{}{}
dan karena itu sifat ini juga berlaku pada suatu lingkungan terbuka dari $t_0$. Pada lingkungan tersebut,

\mavergleichskettedisp
{\vergleichskette
{  { \frac{  f' }{ f } }
}
{ =} { { \frac{  g' }{  2g  } }
}
{ } {
}
{ } {
}
{ } {
}
}
{}{}{.}
Metode penyelesaian persamaan diferensial linear menghasilkan

\mavergleichskettedisp
{\vergleichskette
{   \ln  \betrag {  f  }
}
{ =} { c+ { \frac{   1  }{  2 } }  \ln   g
}
{ } {
}
{ } {
}
{ } {
}
}
{}{}{}
dan dengan demikian

\mavergleichskettedisp
{\vergleichskette
{ \betrag {  f  }
}
{ =} { e^c e^{ { \frac{   1  }{  2 } }  \ln  g }
}
{ =} { e^c  \sqrt{g}
}
{ } {
}
{ } {
}
}
{}{}{}
atau

\mavergleichskettedisp
{\vergleichskette
{ f
}
{ =} { d \sqrt{g}
}
{ } {
}
{ } {
}
{ } {
}
}
{}{}{}
\zusatzklammer {dengan

\mavergleichskette
{\vergleichskette
{ d
}
{ \in }{  \R
}
{  }{
}
{    }{
}
{   }{
}
}
{}{}{}} {} {,}
yang juga dapat diperiksa secara langsung.
\subsection*{Solusi Soal 26.6}
Pertama-tama, untuk suatu fungsi terdiferensialkan $g$ berlaku

\mavergleichskettealign
{\vergleichskettealign
{  \nabla_{ gV_i} V_j - \nabla_{V_j} \left(  g V_i  \right)
}
{ =} { g \nabla_{V_i} V_j -\left(  g \nabla_{V_j} \left(  V_i  \right)  + D_{V_j}(g)  V_i  \right)
}
{ =} { g [V_i,V_j] - D_{V_j}(g)  V_i
}
{ =} { [gV_i,V_j]
}
{ } {
}
}
{}
{}{}
menurut
Fakta (3).
Medan-medan vektor terdiferensialkan sembarang dapat ditulis sebagai

\mavergleichskettedisp
{\vergleichskette
{ V
}
{ =} { \sum_{i= 1}^n f_i V_i
}
{ } {
}
{ } {
}
{ } {
}
}
{}{}{}
dan

\mavergleichskettedisp
{\vergleichskette
{ W
}
{ =} { \sum_{i= 1}^n g_i V_i
}
{ } {
}
{ } {
}
{ } {
}
}
{}{}{.}
Dengan demikian,

\mavergleichskettealign
{\vergleichskettealign
{  \nabla_VW- \nabla_WV
}
{ =} { \nabla_{\sum_{i= 1}^n f_i V_i}  \left(  \sum_{j= 1}^n g_j V_j   \right) - \nabla_{\sum_{j= 1}^n g_j V_j}  \left(  \sum_{i= 1}^n f_i V_i   \right)
}
{ =} { \sum_{i,j}   f_i  \nabla_{V_i} \left(  g_jV_j  \right) - \sum_{i,j}   g_j  \nabla_{V_j} \left(  f_iV_i  \right)
}
{ =} { \sum_{i,j}   f_i  \left(  g_j \nabla_{V_i}  V_j - [V_i,g_jV_j]               \right)  - \sum_{i,j}   g_j  \left(  f_i \nabla_{V_j} V_i -    [ f_i V_i,V_j]  \right)
}
{ =} {  \sum_{i,j}   f_i   g_j \left(  \nabla_{V_i}  V_j -  \nabla_{V_j}  V_i              \right)  - \sum_{i} f_i   [ V_i, g_jV_j]  - \sum_{j} g_j   [ f_iV_i, V_j]
}
}
{
\vergleichskettefortsetzungalign
{ =} {  \sum_{i,j}   f_i   g_j [V_i,V_j]     - \sum_{i} f_i   [ V_i, g_jV_j]  - \sum_{j} g_j   [ f_iV_i, V_j]
}
{ } {}
{ } {}
{ } {}
}
{}{.}
\subsection*{Solusi Soal 26.9}
Pernyataan ini bersifat lokal; oleh karena itu, tanpa mengurangi keumuman, misalkan

\mavergleichskette
{\vergleichskette
{ U
}
{ \subseteq }{  \R^n
}
{  }{
}
{    }{
}
{   }{
}
}
{}{}{}
dilengkapi dengan suatu struktur Riemann melalui fungsi-fungsi
\mathl{g_{ij}}{.} Hipotesis menyatakan

\mavergleichskettedisp
{\vergleichskette
{ \left\langle \nabla_{\partial_i} \partial_j  ,  \partial_k   \right\rangle  + \left\langle  \partial_j  ,  \nabla_{\partial_i} \partial_k   \right\rangle
}
{ =} { \partial_i g_{jk}
}
{ } {
}
{ } {
}
{ } {}
}
{}{}{.}
Dalam bahasa simbol Christoffel,

\mavergleichskettedisp
{\vergleichskette
{  \left\langle \nabla_{\partial_i} \partial_j   ,  \partial_k   \right\rangle
}
{ =} { \left\langle  \sum_{\ell = 1}^n \Gamma^\ell_{ij}  \partial_\ell   ,  \partial_k   \right\rangle
}
{ =} {   \sum_{\ell = 1}^n \Gamma^\ell_{ij}   g_{ \ell k}
}
{ } {
}
{ } {
}
}
{}{}{}
dan

\mavergleichskettedisp
{\vergleichskette
{  \left\langle  \partial_j   ,  \nabla_{\partial_i} \partial_k   \right\rangle
}
{ =} { \left\langle  \partial_j   ,  \sum_{\ell = 1}^n \Gamma^\ell_{ik}  \partial_\ell     \right\rangle
}
{ =} {  \sum_{\ell = 1}^n \Gamma^\ell_{ik}   g_{ \ell j}
}
{ } {
}
{ } {
}
}
{}{}{,}
sehingga syarat tersebut menyatakan

\mavergleichskettedisp
{\vergleichskette
{   \sum_{\ell = 1}^n \Gamma^\ell_{ij}   g_{ \ell k} +\sum_{\ell = 1}^n \Gamma^\ell_{ik}   g_{ \ell j}
}
{ =} {   \partial_i g_{jk}
}
{ } {
}
{ } {
}
{ } {
}
}
{}{}{.}
Kita harus menunjukkan bahwa

\mavergleichskettedisp
{\vergleichskette
{ D_V \left\langle  W  , Z  \right\rangle
}
{ =} { \left\langle \nabla_V W , Z  \right\rangle  + \left\langle  W  ,  \nabla_VZ  \right\rangle
}
{ } {
}
{ } {
}
{ } {
}
}
{}{}{}
untuk medan-medan vektor $V,W,Z$ sembarang yang terdiferensialkan secara kontinu, yang masing-masing dapat ditulis sebagai
\mathl{\sum_{i = 1}^nf_i \partial_i}{} dengan fungsi-fungsi yang terdiferensialkan secara kontinu. Pertama, misalkan

\mavergleichskette
{\vergleichskette
{ V
}
{ = }{ \partial_i
}
{  }{
}
{    }{
}
{   }{
}
}
{}{}{}
dan

\mavergleichskettedisp
{\vergleichskette
{ W
}
{ =} { \sum_{j = 1}^nf_j \partial_j
}
{ } {
}
{ } {
}
{ } {
}
}
{}{}{}
dan

\mavergleichskettedisp
{\vergleichskette
{ Z
}
{ =} { \sum_{k = 1}^n h_k \partial_k
}
{ } {
}
{ } {
}
{ } {
}
}
{}{}{.}
Hasil kali skalar aditif dalam komponen-komponennya dan $\nabla_{\partial_i}$ juga aditif karena yang diberikan adalah koneksi linear. Demikian pula,

\mavergleichskette
{\vergleichskette
{ D_{\partial_i}
}
{ = }{ \partial_i
}
{  }{
}
{    }{
}
{   }{
}
}
{}{}{}
bersifat aditif. Karena itu, cukup kita batasi pada medan-medan berbentuk

\mavergleichskette
{\vergleichskette
{ W
}
{ = }{ f  \partial_j
}
{  }{
}
{    }{
}
{   }{
}
}
{}{}{}
dan

\mavergleichskette
{\vergleichskette
{ Z
}
{ = }{ h  \partial_k
}
{  }{
}
{    }{
}
{   }{
}
}
{}{}{}
dengan $j,k$ tetap. Untuk ruas kiri diperoleh

\mavergleichskettedisp
{\vergleichskette
{  \partial_i \left\langle f \partial_j  ,  h \partial_k   \right\rangle
}
{ =} { \partial_i \left(  fh g_{jk}  \right)
}
{ =} {  fg_{jk} \partial_i(h) + hg_{jk} \partial_i(f)   +  fh \partial_i \left(  g_{jk}   \right)
}
{ } {
}
{ } {
}
}
{}{}{.}
Menurut
fakta tentang deskripsi lokal koneksi linear,

\mavergleichskettedisp
{\vergleichskette
{ \nabla_{\partial_i} \left(  f \partial_j  \right)
}
{ =} {  \left(  \partial_i f  \right) \partial_j + f \nabla_{\partial_i }\partial_j
}
{ =} { \left(  \partial_i f  \right) \partial_j + f \left(  \sum_{ \ell = 1}^n \Gamma^\ell_{ij} \partial_\ell   \right)
}
{ } {
}
{ } {
}
}
{}{}{}
dan

\mavergleichskettedisp
{\vergleichskette
{ \nabla_{\partial_i} \left(  h \partial_k  \right)
}
{ =} { \left(  \partial_i h  \right) \partial_k + h \nabla_{\partial_i }\partial_k
}
{ =} { \left(  \partial_i h  \right) \partial_k  + h \left(  \sum_{ \ell = 1}^n \Gamma^\ell_{ik} \partial_\ell   \right)
}
{ } {
}
{ } {
}
}
{}{}{.}
Jadi,

\mavergleichskettealign
{\vergleichskettealign
{  \left\langle  \nabla_{\partial_i} \left(  f \partial_j  \right)   ,  h \partial_k  \right\rangle
}
{ =} {  \left\langle  \left(  \partial_i f  \right) \partial_j  + f \left(  \sum_{ \ell = 1}^n \Gamma^\ell_{ij} \partial_\ell   \right)   ,  h \partial_k  \right\rangle
}
{ =} { h \left(   \partial_i f  \right) \left\langle   \partial_j  ,  \partial_k   \right\rangle +fh \sum_{ \ell = 1}^n   \Gamma_{ij}^\ell \left\langle          \partial_\ell   ,  \partial_k   \right\rangle
}
{ =} { h \left(   \partial_i f  \right)  g_{j k} +fh \sum_{ \ell = 1}^n   \Gamma_{ij}^\ell  g_{\ell k }
}
{ } {
}
}
{}
{}{}
dan

\mavergleichskettealign
{\vergleichskettealign
{  \left\langle f \partial_j  ,  \nabla_{\partial_i} \left(  h \partial_k  \right)   \right\rangle
}
{ =} { \left\langle  f\partial_j  ,  \left(  \partial_i h \right) \partial_k  + h \left(  \sum_{ \ell = 1}^n \Gamma^\ell_{ik} \partial_\ell   \right)    \right\rangle
}
{ =} {   f \left(   \partial_i h  \right)  g_{j k} +fh \sum_{ \ell = 1}^n   \Gamma_{ik}^\ell  g_{\ell j }
}
{ } {
}
{ } {
}
}
{}
{}{.}
Jumlah kedua suku ini sama dengan ruas kiri.

Kedua ruas linear terhadap medan arah $V$ dan juga kompatibel dengan perkalian oleh fungsi.

\part{Unit 27}
\chapter[Kuliah 27]{Kuliah 27: Turunan Vertikal, Percepatan Tangensial, dan Geodesik}
\setcounter{section}{27}

  \zwischenueberschrift{Turunan vertikal sepanjang suatu pemetaan}

Misalkan
\maabb {p} { E } { M
} {}
suatu
\definitionsverweis {bundel vektor diferensiabel}{}{}
di atas suatu
\definitionsverweis {manifold diferensiabel}{}{}
$M$ yang dilengkapi dengan suatu
\definitionsverweis {koneksi linear}{}{.}
Koneksi ini menghasilkan suatu
\definitionsverweis {turunan vertikal}{}{.}
Untuk setiap medan vektor
\maabb {V} { M } { TM
} {}
dan setiap penampang yang terdiferensialkan secara kontinu
\maabbdisp {s} { M } { E
} {}
\mathl{\nabla_V s}{} adalah penampang kontinu di $E$ yang diberikan oleh

\mathdisp {M \stackrel{V}{ \longrightarrow  } TM  \stackrel{T(s) }{ \longrightarrow } TE \stackrel{  \pi_{\text{vert} }  }{ \longrightarrow} E} {  }

Pada suatu himpunan terbuka

\mavergleichskette
{\vergleichskette
{ U
}
{ \subseteq }{ M
}
{  }{
}
{    }{
}
{   }{
}
}
{}{}{}
dengan

\mavergleichskette
{\vergleichskette
{ U
}
{ \cong }{ V
}
{ \subseteq }{ \R^n
}
{    }{
}
{   }{
}
}
{}{}{}
serta turunan-turunan parsial
\mathl{\partial_1  , \ldots , \partial_n}{} dan suatu trivialisasi

\mavergleichskettedisp
{\vergleichskette
{ E \vert_U
}
{ \cong} { U \times \R^r
}
{ } {
}
{ } {
}
{ } {
}
}
{}{}{}
dengan penampang basis
\mathl{s_1 , \ldots , s_r}{,} turunan vertikal semacam itu dideskripsikan sepenuhnya oleh simbol-simbol Christoffel $\Gamma_{ij}^k$ melalui

\mavergleichskettedisp
{\vergleichskette
{ \nabla_{\partial_i } (s_j)
}
{ =} { \sum_{k= 1}^r  \Gamma_{ij}^k s_k
}
{ } {
}
{ } {
}
{ } {
}
}
{}{}{}
berdasarkan
Lema 25.10.
Kita ingin memiliki konsep ini bukan hanya bagi penampang di atas $M$, melainkan juga bagi penampang
\maabbdisp {s} { L } { E
} {}
sepanjang suatu
\definitionsverweis {pemetaan diferensiabel}{}{}
tetap
\maabbdisp {\psi} { L } { M
} {}
yaitu ketika terdapat diagram komutatif

\mathdisp {\begin{matrix}  &   &   &  E   \\ & & s \nearrow & \downarrow  \\  &  L  & \stackrel{ \psi }{\longrightarrow}  &  M \! \\ \end{matrix}} {  }

Situasi ini telah kita jumpai pada penampang horizontal. Perhatikan bahwa suatu penampang
\maabbdisp {s} { L } { E
} {}
sama dengan suatu penampang
\maabbdisp {\tilde{s}} { L  } { \psi^*(E) = E \times_XY
} {}
di dalam
\definitionsverweis {bundel vektor tarik balik}{}{}
$\psi^*(E)$. Untuk suatu medan vektor $F$ pada $L$ dan suatu penampang semacam itu, kita dapat meninjau komposisi pemetaan

\mathdisp {L \stackrel{F}{ \longrightarrow  } TL  \stackrel{T(s) }{ \longrightarrow } TE \stackrel{  \pi_{\text{vert} }  }{ \longrightarrow} E} {  }

yang kembali kita nyatakan dengan $\nabla_F s$. Koneksi yang ditentukan oleh turunan vertikal ini kita sebut
\stichwort {koneksi tarik balik} {} $\psi^*\nabla$. Koneksi tersebut mempunyai sifat-sifat berikut.



\inputfaktbeweis
{Linearer Zusammenhang/Vertikale Ableitung/Längs Abbildung/Eigenschaften/Fakt}
{Lema}
{}
{

\faktsituation {Misalkan
\maabb {p} { E } { M
} {}
suatu
\definitionsverweis {bundel vektor diferensiabel}{}{}
di atas suatu
\definitionsverweis {manifold diferensiabel}{}{}
$M$ yang dilengkapi dengan suatu
\definitionsverweis {koneksi linear}{}{}
$\nabla$. Misalkan
\maabb {\psi} { L } { M
} {}
suatu
\definitionsverweis {pemetaan diferensiabel}{}{}
dengan
\definitionsverweis {koneksi tarik balik}{}{}
$\psi^*\nabla$ pada
\definitionsverweis {bundel vektor tarik balik}{}{}
$\psi^*E$ di atas $L$. Maka berlaku sifat-sifat berikut.}
\faktfolgerung {\aufzaehlungzwei {Pemetaan
\maabbeledisp {\psi^*\nabla} {C^1(\psi^* E)} { C^0(\psi^* E \otimes T^* L )
} { s } { \left(  \psi^*\nabla  \right) (s)
} {,}
bersifat
$\R$-\definitionsverweis {linear}{}{.}
Secara khusus, $\psi^*\nabla$ merupakan suatu koneksi linear.
} {Untuk suatu pembatasan
\maabbdisp {\psi} {U'} {U
} {}
pada domain-domain peta terbuka dengan koordinat
\mathl{z_1 , \ldots , z_m}{} di $U'$ dan koordinat
\mathl{x_1 , \ldots , x_n}{}
di $U$, bagi
\definitionsverweis {simbol-simbol Christoffel}{}{}
$\tilde{\Gamma}_{\ell j}^k$ dari $\psi^*\nabla$ dan
\definitionsverweis {penampang basis}{}{}

\mathl{s_1  , \ldots , s_r}{} berlaku hubungan

\mavergleichskettedisp
{\vergleichskette
{ \tilde{\Gamma}_{j \ell }^k
}
{ =} { \sum_{i = 1}^n  \left(  \partial_j \psi_i  \right)   \Gamma_{i\ell }^k
}
{ } {
}
{ } {
}
{ } {
}
}
{}{}{.}
}}
\faktzusatz {}
\faktzusatz {}

}
{

\aufzaehlungzwei {Jelas.
} {Kita tinjau situasi lokal secara langsung. Suatu
\definitionsverweis {penampang basis}{}{}
$s_\ell$ sepanjang $L$ dapat dipandang langsung sebagai penampang basis di atas $M$. Pemetaan-pemetaan yang relevan adalah

\mathdisp {U' \stackrel{\partial_j }{\longrightarrow} TU'  \stackrel{T (\psi) }{\longrightarrow} T(U) \stackrel{ T(s_\ell ) }{\longrightarrow} T( U \times \R^r) \stackrel{  \pi_{\text{vert} } }{\longrightarrow} \R^r} { . }

Di sini,

\mavergleichskettedisp
{\vergleichskette
{ T(\psi) \circ \partial_j
}
{ =} { \partial_j \psi
}
{ =} { \sum_{i = 1}^n \left(  \partial_j \psi_i  \right)  \partial_i
}
{ } {
}
{ } {
}
}
{}{}{.}
Dengan demikian, untuk

\mavergleichskette
{\vergleichskette
{ Q
}
{ \in }{ U'
}
{  }{
}
{    }{
}
{   }{
}
}
{}{}{}
dan dengan menggunakan
Lema 25.10,

\mavergleichskettealignhandlinks
{\vergleichskettealignhandlinks
{ \left(  \left(  \psi^*\nabla  \right) _{\partial_j } (s_\ell)  \right) (Q)
}
{ =} { \left(  \psi^*\nabla  \right) _{\partial_j (Q) } (s_\ell) (Q)
}
{ =} { \nabla_{  \left(  T(\psi) \circ \partial_j  \right)  (Q)  }  (s_\ell) ( \psi(Q) )
}
{ =} { \nabla_{  \sum_{i = 1}^n  \left(  \partial_j \psi_i  \right)  (Q) \partial_i  }  (s_\ell) ( \psi(Q) )
}
{ =} { \sum_{i = 1}^n  \left(  \partial_j \psi_i  \right)  (Q)\nabla_{  \partial_i }  (s_\ell) ( \psi(Q) )
}
}
{
\vergleichskettefortsetzungalign
{ =} { \sum_{i = 1}^n  \left(  \partial_j \psi_i  \right)  (Q)  \left(  \sum_{k = 1}^r \Gamma_{i \ell }^k(\psi(Q)) s_k (\psi(Q))  \right)
}
{ =} { \sum_{k = 1}^r  \left(  \sum_{i = 1}^n  \left(  \partial_j \psi_i  \right)  (Q)  \Gamma_{i \ell }^k(\psi(Q))   \right) s_k (\psi(Q))
}
{ } {}
{ } {}
}
{}{.}
Dengan meninjau komponen ke-$k$, kita memperoleh

\mavergleichskettedisp
{\vergleichskette
{ \tilde{\Gamma}_{j \ell }^k (Q)
}
{ =} { \sum_{i = 1}^n \Gamma_{i\ell }^k (\psi(Q))  \left(  \partial_j \psi_i  \right)  (Q)
}
{ } {
}
{ } {
}
{ } {
}
}
{}{}{.}
}

}



\inputfaktbeweis
{Linearer Zusammenhang/Vertikale Ableitung/Längs Abbildung/Metrisch/Fakt}
{Lema}
{}
{

\faktsituation {Misalkan
\maabb {p} { E } { M
} {}
suatu
\definitionsverweis {bundel vektor diferensiabel}{}{}
di atas suatu
\definitionsverweis {manifold diferensiabel}{}{}
$M$ yang dilengkapi dengan suatu
\definitionsverweis {koneksi linear}{}{}
\definitionsverweis {metrik}{}{}
$\nabla$. Misalkan
\maabb {\psi} { L } { M
} {}
suatu
\definitionsverweis {pemetaan diferensiabel}{}{}
dengan
\definitionsverweis {koneksi tarik balik}{}{}
$\psi^*\nabla$ pada
\definitionsverweis {bundel vektor tarik balik}{}{}
$\psi^*E$ di atas $L$.}
\faktfolgerung {Maka
\mathl{\psi^*\nabla}{} juga metrik.}
\faktzusatz {}
\faktzusatz {}

}
{

Pertama, menurut
Soal 16.13,
$\psi^*E$ kembali merupakan suatu bundel Riemann. Kita tinjau situasi lokal
\maabbdisp {\psi} {U'} {U
} {}
dan

\mavergleichskette
{\vergleichskette
{ E
}
{ = }{ U \times \R^r
}
{  }{
}
{    }{
}
{   }{
}
}
{}{}{}
serta, bersesuaian dengan itu,

\mavergleichskette
{\vergleichskette
{ \psi^*E
}
{ = }{ U' \times \R^r
}
{  }{
}
{    }{
}
{   }{
}
}
{}{}{.}
Fungsi-fungsi yang mendeskripsikan struktur Riemann pada $E$,

\mavergleichskettedisp
{\vergleichskette
{ h_{ \ell m}
}
{ =} { \left\langle s_\ell , s_m   \right\rangle
}
{ } {
}
{ } {
}
{ } {
}
}
{}{}{,}
pada $U$ dan $U'$ berhubungan langsung melalui $\psi$. Misalkan $\partial_j$ suatu
\definitionsverweis {medan vektor standar}{}{}
pada $U'$ dan
\mathl{s_\ell, s_m}{} penampang basis di $E$. Maka

\mavergleichskettealign
{\vergleichskettealign
{ \partial_j  \left\langle s_\ell , s_m   \right\rangle
}
{ =} { \partial_j  \left(  h_{\ell , m} \circ\psi   \right)
}
{ =} { \sum_{i = 1}^n  \left(  \partial_j \psi_i  \right) \left(  \partial_i  h_{\ell , m}  \right)
}
{ =} { \sum_{i = 1}^n  \left(  \partial_j \psi_i  \right) \partial_i  \left\langle s_\ell , s_m   \right\rangle
}
{ } {}
}
{}
{}{.}
Menurut
Lema 27.1,

\mavergleichskettealignhandlinks
{\vergleichskettealignhandlinks
{ \left\langle  \left(  \psi^*\nabla  \right)_{\partial_j } s_\ell ,  s_m   \right\rangle
}
{ =} { \left\langle  \sum_{k = 1}^r \tilde{\Gamma}_{j \ell }^k s_ k ,  s_m   \right\rangle
}
{ =} { \sum_{k = 1}^r \left\langle  \tilde{\Gamma}_{j \ell }^k s_ k ,  s_m   \right\rangle
}
{ =} { \sum_{k = 1}^r \left\langle  \sum_{i = 1}^n \Gamma_{i\ell }^k  \left(  \partial_j \psi_i  \right)  s_k  ,  s_m   \right\rangle
}
{ =} { \sum_{i = 1}^n \left(  \partial_j \psi_i  \right) \left\langle  \sum_{k = 1}^r \Gamma_{i\ell }^k s_k  ,  s_m   \right\rangle
}
}
{
\vergleichskettefortsetzungalign
{ =} { \sum_{i = 1}^n \left(  \partial_j \psi_i  \right) \left\langle  \nabla_{\partial_i } s_\ell  ,  s_m   \right\rangle
}
{ } {}
{ } {}
{ } {}
}
{}{}
dan, secara bersesuaian,

\mavergleichskettedisp
{\vergleichskette
{ \left\langle  s_\ell ,  \left(  \psi^*\nabla  \right)_{\partial_j } s_m   \right\rangle
}
{ =} { \sum_{i = 1}^n  \left(  \partial_j \psi_i  \right)   \left\langle  s_\ell  , \nabla_{\partial_i } s_m   \right\rangle
}
{ } {
}
{ } {
}
{ } {
}
}
{}{}{.}
Karena $\nabla$ metrik, diperoleh

\mavergleichskettedisp
{\vergleichskette
{ \partial_j  \left\langle s_\ell , s_m   \right\rangle
}
{ =} { \left\langle  \left(  \psi^*\nabla  \right)_{\partial_j } s_\ell ,  s_m   \right\rangle    + \left\langle  s_\ell ,  \left(  \psi^*\nabla  \right)_{\partial_j } s_m   \right\rangle
}
{ } {
}
{ } {
}
{ } {
}
}
{}{}{}
dan dari sini, bersama
Soal 26.10,
diperoleh bahwa $\psi^*\nabla$ metrik.

}

  \zwischenueberschrift{Percepatan tangensial}

Misalkan $M$ suatu
\definitionsverweis {manifold}{}{.}
Untuk suatu
\definitionsverweis {kurva diferensiabel}{}{}
\maabbdisp {\gamma} { I } { M
} {}
berlaku bahwa
\maabbdisp {\gamma'} { I} {TM
} {}
merupakan kurva di dalam bundel tangen yang menjadi suatu pengangkatan dari $\gamma$. Jika kurva ini pada gilirannya terdiferensialkan, diperoleh suatu kurva
\maabbdisp {\gamma^{\prime \prime}} { I} {TTM
} {,}
namun kurva ini tidak dapat langsung dihubungkan dengan $\gamma'$ karena nilainya berada di ruang lain. Sebaliknya, apabila pada $TM$ diberikan suatu
\definitionsverweis {koneksi}{}{},
maka melalui turunan vertikal
\zusatzklammer {tarik balik} {} {}
turunan kedua dapat didefinisikan sebagai
\mathl{\nabla_{ \partial } \gamma'}{}.

\inputdefinition
{}
{

Misalkan $M$ suatu
\definitionsverweis {manifold Riemann}{}{}
yang dilengkapi dengan
\definitionsverweis {koneksi Levi--Civita}{}{.}
Untuk suatu
\definitionsverweis {kurva diferensiabel}{}{}
\maabbdisp {\gamma} { I } { M
} {}
yang dua kali terdiferensialkan secara kontinu, pemetaan
\mathl{\pi_{\text{vert} }  \circ TT(\gamma)}{} disebut
\definitionswort {percepatan tangensial}{}
dari $\gamma$.

}

Besaran ini juga ditulis sebagai
\mathl{\pi_{\text{vert} } \circ \gamma^{\prime \prime}}{} atau
\zusatzklammer {sebagai turunan vertikal sepanjang $\gamma$} {} {}

\mathl{\nabla_\partial \gamma'}{,} dengan
\maabbdisp {\gamma'} { I } { TM
} {}
dipandang sebagai penampang dalam bundel tangen sepanjang lintasan $\gamma$, dan dengan meninjau rantai pemetaan

\mathdisp {I \stackrel{\partial}{\longrightarrow} I \times \R = TI \stackrel{T(\gamma')}{\longrightarrow} TTM \stackrel{ \pi_{\text{vert} } }{\longrightarrow}  TM} {  }

.



\inputfaktbeweis
{Hyperfläche/Zweite Ableitung einer Kurve/Direkt und über Zusammenhang/Fakt}
{Lema}
{}
{

\faktsituation {Misalkan

\mavergleichskette
{\vergleichskette
{ W
}
{ \subseteq }{ \R^n
}
{  }{
}
{    }{
}
{   }{
}
}
{}{}{}
\definitionsverweis {terbuka}{}{,}
\maabb {h} { W } { \R
} {}
suatu
\definitionsverweis {fungsi yang terdiferensialkan secara kontinu}{}{}
dan

\mavergleichskette
{\vergleichskette
{ Y
}
{ = }{ h^{-1}(c)
}
{  }{
}
{    }{
}
{   }{
}
}
{}{}{}
\definitionsverweis {serat}{}{}
di atas

\mavergleichskette
{\vergleichskette
{ c
}
{ \in }{ \R
}
{  }{
}
{    }{
}
{   }{
}
}
{}{}{,}
dengan $h$ bersifat
\definitionsverweis {reguler}{}{}
di setiap titik $Y$. Misalkan
\maabbdisp {\gamma} { I } { Y
} {}
suatu kurva yang dua kali terdiferensialkan secara kontinu.}
\faktfolgerung {Maka percepatan tangensial $\gamma$ dalam pengertian
Definisi 1.7
sama dengan percepatan tangensial dalam pengertian
Definisi 27.3.}
\faktzusatz {}
\faktzusatz {}

}
{

Menurut
Lema 26.15,
\zusatzklammer {yang juga berlaku dalam dimensi lebih tinggi} {} {},
kedua koneksi tersebut sama.

}

  \zwischenueberschrift{Kurva geodesik}

\inputdefinition
{}
{

Misalkan $M$ suatu
\definitionsverweis {manifold Riemann}{}{}
yang dilengkapi dengan
\definitionsverweis {koneksi Levi--Civita}{}{.}
Suatu
\definitionsverweis {kurva diferensiabel}{}{}
\maabbdisp {\gamma} { I } { M
} {}
yang dua kali terdiferensialkan disebut
\definitionswort {kurva geodesik}{}
\zusatzklammer {atau
\definitionswort {geodesik}{}} {} {,}
apabila

\mavergleichskettedisp
{\vergleichskette
{ \nabla_{d/dt}  \gamma'
}
{ =} { 0
}
{ } {
}
{ } {
}
{ } {
}
}
{}{}{}
berlaku pada $I$.

}

Kadang-kadang digunakan pula istilah \stichwort {geodesik} {.} Perlu ditekankan bahwa kurva geodesik adalah sebuah kurva; dengan demikian, yang dimaksud bukan hanya
\definitionsverweis {citra}{}{}
kurva tersebut, melainkan juga proses geraknya.



\inputfaktbeweis
{Hyperfläche/Geodätische/Direkt und über Zusammenhang/Fakt}
{Lema}
{}
{

\faktsituation {Misalkan

\mavergleichskette
{\vergleichskette
{ W
}
{ \subseteq }{ \R^n
}
{  }{
}
{    }{
}
{   }{
}
}
{}{}{}
\definitionsverweis {terbuka}{}{,}
\maabb {h} { W } { \R
} {}
suatu
\definitionsverweis {fungsi yang terdiferensialkan secara kontinu}{}{}
dan

\mavergleichskette
{\vergleichskette
{ Y
}
{ = }{ h^{-1}(c)
}
{  }{
}
{    }{
}
{   }{
}
}
{}{}{}
\definitionsverweis {serat}{}{}
di atas

\mavergleichskette
{\vergleichskette
{ c
}
{ \in }{ \R
}
{  }{
}
{    }{
}
{   }{
}
}
{}{}{,}
dengan $h$ bersifat
\definitionsverweis {reguler}{}{}
di setiap titik $Y$.}
\faktfolgerung {Maka geodesik-geodesik di $Y$ dalam pengertian
Definisi 1.8
sama dengan geodesik-geodesik dalam pengertian
Definisi 27.5.}
\faktzusatz {}
\faktzusatz {}

}
{

Hal ini mengikuti dari
Lema 27.4.

}



\inputfaktbeweis
{Riemannsche Mannigfaltigkeit/Levi-Civita-Zusammenhang/Geodätische/Geschwindigkeitskonstanz/Fakt}
{Lema}
{}
{

\faktsituation {Misalkan
\maabb {\gamma} { I } { M
} {}
suatu
\definitionsverweis {kurva geodesik}{}{}
pada suatu
\definitionsverweis {manifold Riemann}{}{}
$M$ yang dilengkapi dengan
\definitionsverweis {koneksi Levi--Civita}{}{.}}
\faktfolgerung {Maka
\mathl{\Vert { \gamma'(t) } \Vert}{} konstan pada $I$.}
\faktzusatz {}
\faktzusatz {}

}
{

Pada suatu domain peta terbuka, dengan menggunakan
Teorema 26.14  (2)
dan
Lema 27.2,
diperoleh

\mavergleichskettealign
{\vergleichskettealign
{ { \frac{   \partial }{  \partial t } }  \left\langle  \gamma'(t) ,  \gamma'(t)  \right\rangle
}
{ =} { \left\langle \nabla_{ { \frac{   \partial }{  \partial t } }  }\gamma'(t) ,  \gamma'(t)  \right\rangle + \left\langle  \gamma'(t) , \nabla_{  { \frac{   \partial }{  \partial t } }  }\gamma'(t)  \right\rangle
}
{ =} { 2 \left\langle \nabla_{ { \frac{   \partial }{  \partial t } }  }\gamma'(t) ,  \gamma'(t)  \right\rangle
}
{ =} { 2 \left\langle  0  ,  \gamma'(t)  \right\rangle
}
{ =} { 0
}
}
{}
{}{,}
sehingga
\mathl{\left\langle  \gamma'(t) ,  \gamma'(t)  \right\rangle}{} konstan.

}



\inputfaktbeweis
{Riemannsche Mannigfaltigkeit/Levi-Civita-Zusammenhang/Geodätische/Differentialgleichung/Fakt}
{Teorema}
{}
{

\faktsituation {Suatu
\definitionsverweis {kurva diferensiabel}{}{}
\maabb {\gamma} { I } { M
} {}
yang dua kali terdiferensialkan pada suatu
\definitionsverweis {manifold Riemann}{}{}
$M$ adalah}
\faktfolgerung {suatu
\definitionsverweis {kurva geodesik}{}{}
\zusatzklammer {terhadap
\definitionsverweis {koneksi Levi--Civita}{}{}} {} {,}
jika dan hanya jika pada setiap peta kurva tersebut memenuhi
\definitionsverweis {sistem persamaan diferensial biasa orde kedua}{}{}

\mavergleichskettedisp
{\vergleichskette
{ \gamma_k^{\prime \prime} + \sum_{ij} \Gamma^k_{ij} \gamma_i' \gamma_j'
}
{ =} { 0
}
{ } {
}
{ } {
}
{ } {
}
}
{}{}{}
\zusatzklammer {untuk semua $k$} {} {}.}
\faktzusatz {}
\faktzusatz {}

}
{

Misalkan $n$ dimensi $M$. Kita meninjau situasinya secara langsung pada suatu citra peta terbuka

\mavergleichskette
{\vergleichskette
{ U
}
{ \subseteq }{ \R^n
}
{  }{
}
{    }{
}
{   }{
}
}
{}{}{.}
Menurut
Catatan 25.8,
turunan vertikal diberikan oleh
\maabbeledisp {} {T(U \times \R^n) = U \times \R^n \times \R^n \times \R^n } { U \times \R^n \times  \R^n
} {(P,e_j, \partial_i, w)} {(P,e_j,\sum_{k = 1}^n \Gamma^k_{ij} (P) e_k +w) \cong V(U \times \R^n)
} {},
sedangkan pemetaan ke $TU$ diperoleh dengan menghilangkan komponen tengah. Turunan kedua kurva mula-mula adalah pemetaan tangen kedua, yaitu
\zusatzklammer {dengan mengabaikan perkalian dengan arah berdimensi satu dari ruang tangen kurva} {} {}
\maabbeledisp {T(\gamma)} { I } { TU
} { t } { (\gamma(t), \gamma'(t))
} {,}
dan, sejalan dengan itu,
\maabbeledisp {T(T(\gamma))} { I } { T(TU)
} { t } { ((\gamma(t), \gamma'(t)), (\gamma'(t), \gamma^{\prime \prime} (t)))
} {}
\zusatzklammer {kedua komponen pemetaan tangen diturunkan} {} {.}
Dengan

\mavergleichskette
{\vergleichskette
{ \gamma'(t)
}
{ = }{ \sum_{j = 1}^n \gamma_j'(t) \partial_j
}
{  }{
}
{    }{
}
{   }{
}
}
{}{}{}
hal ini dipetakan oleh proyeksi vertikal ke

\mathdisp {\sum_{k = 1}^n  \left(  \sum_{i,j} \gamma_i'(t) \gamma_j'(t) \Gamma^k_{ij}(\gamma(t))  \right)  \partial_k  +\sum_{k = 1}^n \gamma_k^{\prime \prime}(t) \partial_k} {  }

Syarat bagi suatu geodesik bahwa turunan keduanya
\zusatzklammer {di dalam $TTM$} {} {}
selalu horizontal ekuivalen dengan syarat bahwa proyeksi vertikal yang dihitung di atas sama dengan $0$. Ini berarti setiap komponennya sama dengan $0$, yakni

\mavergleichskettedisp
{\vergleichskette
{ \sum_{i,j} \gamma_i'(t) \gamma_j'(t) \Gamma^k_{ij}(\gamma(t))  + \gamma_k^{\prime \prime}(t)
}
{ =} { 0
}
{ } {
}
{ } {
}
{ } {
}
}
{}{}{}
untuk

\mavergleichskette
{\vergleichskette
{ k
}
{ = }{ 1  , \ldots , n
}
{  }{
}
{    }{
}
{   }{
}
}
{}{}{.}

}

\inputbemerkung
{}
{

Dalam situasi
Teorema 27.8,
sistem persamaan diferensial

\mavergleichskettedisp
{\vergleichskette
{ \gamma_k^{\prime \prime}(t) + \sum_{ij} \Gamma^k_{ij}(\gamma(t)) \gamma_i'(t) \gamma_j'(t)
}
{ =} { 0
}
{ } {
}
{ } {
}
{ } {
}
}
{}{}{}
dengan

\mavergleichskette
{\vergleichskette
{ k
}
{ = }{ 1  , \ldots , n
}
{  }{
}
{    }{
}
{   }{
}
}
{}{}{,}
yang secara lokal mencirikan geodesik, disebut \stichwort {persamaan diferensial geodesik} {.} Ini adalah suatu
\definitionsverweis {sistem persamaan diferensial biasa}{}{}
orde kedua dengan $n$ persamaan. Sistem ini tidak linear dan secara umum sulit diselesaikan.

}

\inputbeispiel{}
{

\bild{ \begin{center}
\includegraphics[width=5.5cm]{ \bildeinlesungsvg {Parallel lines in Poincare's model of hyperbolic geometry} {svg} }
\end{center}
\bildtext {} }

\bildlizenz { Parallel lines in Poincare's model of hyperbolic geometry.svg } {} {Januszkaja} {Commons} {CC-by-sa 3.0} {}

Kita melanjutkan
Contoh 26.7.
Menurut
Teorema 27.8,
suatu geodesik harus memenuhi kedua syarat

\mavergleichskettedisp
{\vergleichskette
{ \gamma_1^{\prime \prime} (t)
}
{ =} { { \frac{   2  }{  \gamma_2(t) } } \gamma_1'(t) \gamma_2'(t)
}
{ } {
}
{ } {
}
{ } {
}
}
{}{}{}
dan

\mavergleichskettedisp
{\vergleichskette
{ \gamma_2^{\prime \prime} (t)
}
{ =} { -  { \frac{   1  }{  \gamma_2(t) } }  \gamma_1'(t) \gamma_1'(t) + { \frac{   1  }{  \gamma_2(t) } }    \gamma_2'(t) \gamma_2'(t)
}
{ } {
}
{ } {
}
{ } {
}
}
{}{}{}.
Untuk

\mavergleichskette
{\vergleichskette
{ \gamma_1
}
{ = }{ x_0
}
{  }{
}
{    }{
}
{   }{
}
}
{}{}{}
yang konstan, sistem ini menjadi

\mavergleichskettedisp
{\vergleichskette
{ \gamma_2^{\prime \prime} (t)
}
{ =} { { \frac{   1  }{  \gamma_2(t) } }  \gamma_2'(t) \gamma_2'(t)
}
{ } {
}
{ } {
}
{ } {
}
}
{}{}{}
dengan solusi komponen

\mavergleichskettedisp
{\vergleichskette
{ y_2(t)
}
{ =} { e^t
}
{ } {
}
{ } {
}
{ } {
}
}
{}{}{}
dan solusi keseluruhan

\mavergleichskettedisp
{\vergleichskette
{ y(t)
}
{ =} { \left(  x_0   , \,  e^t                   \right)
}
{ } {
}
{ } {
}
{ } {
}
}
{}{}{,}
yang bergerak pada suatu garis vertikal. Selain itu, menurut
Soal 27.13,

\mavergleichskettedisp
{\vergleichskette
{ \gamma(t)
}
{ =} { \left(  s   , \,   0                   \right) + r \left(  \tanh  t   , \,   { \frac{   1  }{  \cosh  t  } }                   \right)
}
{ } {
}
{ } {
}
{ } {
}
}
{}{}{}
untuk

\mavergleichskette
{\vergleichskette
{ s
}
{ \in }{ \R
}
{  }{
}
{    }{
}
{   }{
}
}
{}{}{,}

\mavergleichskette
{\vergleichskette
{ r
}
{ \in }{ \R_+
}
{  }{
}
{    }{
}
{   }{
}
}
{}{}{}
merupakan solusi yang, menurut
Soal 20.50 (Analysis (Osnabrück 2021-2023))  (3),
bergerak pada setengah lingkaran dengan pusat
\mathl{\left(  s   , \,   0                   \right)}{} dan jari-jari $r$.

}

\inputbemerkung
{}
{

Apa yang disebut \stichwort {aksioma paralel} {}
\zusatzklammer {atau postulat paralel} {} {}
dalam geometri Euklides bidang menyatakan bahwa untuk setiap garis dan setiap titik yang tidak terletak pada garis tersebut, ada tepat satu garis paralel yang melalui titik itu. Sebuah garis paralel ditentukan oleh sifat bahwa garis tersebut tidak memotong garis yang diberikan. Salah satu pertanyaan penting dalam sejarah adalah apakah postulat paralel dapat diturunkan dari aksioma dan postulat Euklides lainnya, sehingga postulat ini sebenarnya tidak diperlukan. Pada paruh pertama abad ke-19, pertanyaan ini dijawab secara negatif dengan memberikan model-model geometri
\zusatzklammer {yang khususnya memuat konsep titik dan garis} {} {}
yang memenuhi aksioma-aksioma geometri Euklides lainnya, tetapi tidak memenuhi aksioma paralel. Dalam kerangka geometri Riemann, khususnya dengan konsep
\definitionsverweis {geodesik}{}{},
model-model semacam itu muncul secara alami apabila geodesik-geodesik
\zusatzklammer {di sini yang dimaksud adalah citra suatu kurva geodesik} {} {}
dipandang sebagai garis. Perhatikan bahwa dalam pengembangan aksiomatik suatu geometri, garis tidak ditentukan oleh gambaran intuitif, melainkan oleh sifat-sifat yang dimilikinya dalam hubungannya dengan objek-objek lain. Sebagai contoh,
Contoh 27.10
bersama geodesik-geodesik yang dideskripsikan di sana memberikan suatu model geometri non-Euklides: bagi suatu geodesik yang diberikan dan suatu titik lain, selalu terdapat tak terhingga banyak geodesik melalui titik tersebut yang tidak berpotongan dengan geodesik yang diberikan.

}

\inputdefinition
{}
{

Misalkan $M$ suatu
\definitionsverweis {manifold Riemann}{}{}
yang
\definitionsverweis {terhubung}{}{.}
Untuk titik-titik

\mavergleichskette
{\vergleichskette
{ P,Q
}
{ \in }{ M
}
{  }{
}
{    }{
}
{   }{
}
}
{}{}{}
didefinisikan

\mavergleichskettedisp
{\vergleichskette
{ d(P,Q)
}
{ \defeq} { \operatorname{inf} ( L(\gamma), \gamma \text{ adalah lintasan penghubung kontinu yang terdiferensialkan secara kontinu sepotong-sepotong dari } P \text{ ke } Q  )
}
{ } {
}
{ } {
}
{ } {
}
}
{}{}{}
dan besaran ini disebut
\definitionswort {metrik Riemann}{}
pada $M$.

}

Metrik ini benar-benar menjadikan $M$ suatu ruang metrik yang topologinya sama dengan topologi yang diberikan. Selain itu, dapat dibuktikan bahwa realisasi jarak antara dua titik yang diparameterkan oleh panjang busur merupakan suatu geodesik.

\chapter{Lembar Kerja 27}
\setcounter{section}{27}

  \zwischenueberschrift{Soal latihan}

\inputaufgabe
{}
{

Misalkan
\maabb {p} { E } { M
} {}
suatu
\definitionsverweis {bundel vektor diferensiabel}{}{}
di atas
\definitionsverweis {manifold diferensiabel}{}{}
$M$, yang dilengkapi dengan
\definitionsverweis {koneksi linear}{}{}
$\nabla$. Misalkan $L$ manifold diferensiabel lain dan
\maabb {\psi} { L } { M
} {}
suatu
\definitionsverweis {pemetaan terdiferensialkan secara kontinu}{}{}
dengan
\definitionsverweis {koneksi tarikan balik}{}{}
$\psi^*\nabla$ pada
\definitionsverweis {bundel vektor tarikan balik}{}{}
$\psi^*E$ di atas $L$. Tunjukkan bahwa suatu penampang yang terdiferensialkan secara kontinu
\maabb {s} { L } { E
} {}
bersifat
\definitionsverweis {horizontal}{}{}
terhadap $\nabla$ jika dan hanya jika penampang terkait
\maabbdisp {s} { L } { \psi^*E
} {}
bersifat horizontal terhadap $\psi^* \nabla$.

}
{} {}

\inputaufgabe
{}
{

Misalkan
\maabb {p} {M \times \R^r } { M
} {}
adalah
\definitionsverweis {bundel vektor}{}{}
trivial di atas suatu
\definitionsverweis {manifold diferensiabel}{}{}
$M$, dilengkapi dengan
\definitionsverweis {koneksi trivial}{}{}
$\nabla$. Misalkan $L$ manifold diferensiabel lain dan
\maabb {\psi} { L } { M
} {}
suatu
\definitionsverweis {pemetaan terdiferensialkan secara kontinu}{}{.}
Tunjukkan bahwa
\definitionsverweis {koneksi tarikan balik}{}{}
$\psi^*\nabla$
\zusatzklammer {pada $L \times \R^r$} {} {}
juga trivial.

}
{} {}

\inputaufgabe
{}
{

Misalkan
\maabb {p} { E } { M
} {}
suatu
\definitionsverweis {bundel vektor diferensiabel}{}{}
di atas
\definitionsverweis {manifold diferensiabel}{}{}
$M$, yang dilengkapi dengan
\definitionsverweis {koneksi linear}{}{}
$\nabla$. Misalkan

\mavergleichskette
{\vergleichskette
{ I
}
{ \subseteq }{ \R
}
{  }{
}
{    }{
}
{   }{
}
}
{}{}{}
suatu
\definitionsverweis {interval terbuka}{}{}
dan
\maabb {\gamma} { I } { M
} {}
suatu
\definitionsverweis {kurva terdiferensialkan secara kontinu}{}{}
dengan
\definitionsverweis {koneksi tarikan balik}{}{}
$\gamma^*\nabla$ pada
\definitionsverweis {bundel vektor tarikan balik}{}{}
$\gamma^*E$ di atas $I$. Tentukan
\definitionsverweis {simbol-simbol Christoffel}{}{}
dari
\mathl{\gamma^* \nabla}{.}

}
{} {}

\inputaufgabegibtloesung
{}
{

Kita tinjau parametrisasi trigonometris
\maabbeledisp {\psi} {\R } { S^1
} { t } { \left(  \cos  t   , \,   \sin  t                   \right)
} {}
dari lingkaran satuan. Tetapkan

\mavergleichskette
{\vergleichskette
{ c
}
{ \in }{ [0, 2 \pi[
}
{  }{
}
{    }{
}
{   }{
}
}
{}{}{}
. Pada bundel vektor trivial
\maabbdisp {} {S^1 \times \R^2} { S^1
} {}
diberikan suatu
\definitionsverweis {koneksi linear}{}{}
$\nabla$ sedemikian sehingga sepanjang $\psi$
\definitionsverweis {koneksi tarikan balik}{}{}
$\psi^* \nabla$ diberikan oleh
\definitionsverweis {simbol-simbol Christoffel}{}{}
\zusatzklammer {indeks untuk satu-satunya arah penurunan tidak dituliskan} {} {}

\mavergleichskettedisp
{\vergleichskette
{ \begin{pmatrix} \tilde{\Gamma}_1^1 (t)   &  \tilde{\Gamma}_1^2 (t)   \\ \tilde{\Gamma}_2^1(t)  &  \tilde{\Gamma}_2^2(t)  \end{pmatrix}
}
{ =} { \begin{pmatrix}  0  &   - { \frac{  c  }{  2 \pi   } }   \\  { \frac{  c  }{  2 \pi  } }  &  0  \end{pmatrix}
}
{ } {
}
{ } {
}
{ } {
}
}
{}{}{}
. Untuk suatu titik

\mavergleichskette
{\vergleichskette
{ Q
}
{ = }{ \begin{pmatrix}  a  \\b   \end{pmatrix}
}
{ \in }{ \R^2
}
{    }{
}
{   }{
}
}
{}{}{}
kita tinjau pemetaan
\maabbeledisp {\Psi_{Q}} {\R} { S^1 \times \R^2
} { t } { \left(  \cos  t   , \,   \sin  t  , \,   M(t)  \begin{pmatrix}  a  \\b   \end{pmatrix}                 \right)
} {}
dengan

\mavergleichskettedisp
{\vergleichskette
{ M(t)
}
{ =} { \begin{pmatrix}  \cos  { \frac{  ct  }{  2 \pi   } }   &   -  \sin  { \frac{  ct  }{  2 \pi   } }    \\  \sin  { \frac{  ct  }{  2 \pi  } }   &  \cos  { \frac{  ct  }{  2 \pi   } }  \end{pmatrix}
}
{ } {
}
{ } {
}
{ } {
}
}
{}{}{.}
\aufzaehlungfuenf{Tentukan $\Psi_Q(0)$.
}{Tentukan $\Psi_Q(2 \pi)$.
}{Tunjukkan bahwa $\Psi_Q$ adalah suatu
\definitionsverweis {pengangkatan horizontal}{}{}
sepanjang $\psi$.
}{Tunjukkan bahwa
\maabbeledisp {} {\R^2} { \R^2
} { Q } { \Psi_Q(2 \pi)
} {,}
merupakan
\definitionsverweis {isometri linear}{}{}
\zusatzklammer {titik dasar $(1,0)$ tidak dituliskan di sini} {} {.}
}{Tunjukkan bahwa
\maabbeledisp {} {\Z} { \operatorname{SL}_{  2  } \! \left(  \R  \right)
} { k } { \left(  Q \mapsto \Psi_Q(2k \pi )  \right)
} {,}
merupakan
\definitionsverweis {homomorfisme grup}{}{.}
}

}
{} {}

\inputaufgabegibtloesung
{}
{

Untuk kurva
\maabbeledisp {\psi} {\R} { {\mathbb H}
} { t } { \left(  t   , \,  e^t                   \right)
} {,}
yang bernilai dalam
\definitionsverweis {setengah bidang Poincaré}{}{}
${\mathbb H}$, tentukan
\definitionsverweis {simbol-simbol Christoffel}{}{}
untuk
\definitionsverweis {koneksi tarikan balik}{}{}
\zusatzklammer {dari
\definitionsverweis {koneksi Levi-Civita}{}{}
pada ${\mathbb H}$} {} {}
di atas $\R$.

}
{} {}

Dalam soal berikut, perhatikan bahwa menurut
Soal 11.23,
tidak ada bedanya apakah tarikan balik bundel vektor ke $Y$ ditentukan langsung atau melalui $Z$.

\inputaufgabe
{}
{

Misalkan
\maabb {p} { E } { M
} {}
suatu
\definitionsverweis {bundel vektor diferensiabel}{}{}
di atas
\definitionsverweis {manifold diferensiabel}{}{}
$M$, yang dilengkapi dengan
\definitionsverweis {koneksi linear}{}{}
$\nabla$. Misalkan $Y,Z$ manifold-manifold diferensiabel lain dengan
\definitionsverweis {pemetaan-pemetaan terdiferensialkan secara kontinu}{}{}

\mathdisp {Y \stackrel{\varphi}{\longrightarrow} Z \stackrel{\psi}{\longrightarrow}  M} { . }

Tunjukkan bahwa untuk
\definitionsverweis {koneksi-koneksi tarikan balik}{}{}
pada
\mathl{\left(  \psi \circ \varphi  \right)^* E}{} berlaku hubungan

\mavergleichskettedisp
{\vergleichskette
{ \left(  \psi \circ \varphi  \right)^*(\nabla)
}
{ =} { \varphi^*  \left(  \psi^*\nabla  \right)
}
{ } {
}
{ } {
}
{ } {
}
}
{}{}{.}

}
{} {}

\inputaufgabe
{}
{

Misalkan $M$ suatu
\definitionsverweis {manifold Riemann}{}{}
yang dilengkapi dengan
\definitionsverweis {koneksi Levi-Civita}{}{}
dan misalkan
\maabbdisp {\gamma} { I } { M
} {}
suatu
\definitionsverweis {kurva}{}{}
yang dua kali terdiferensialkan secara kontinu. Tunjukkan bahwa $\gamma$ adalah
\definitionsverweis {kurva geodesik}{}{}
jika dan hanya jika kurva turunannya
\maabb {\gamma'} { I } { TM
} {}
merupakan
\definitionsverweis {penampang horizontal}{}{}
dalam
\definitionsverweis {bundel tangen}{}{}
$TM$.

}
{} {}

\inputaufgabe
{}
{

Misalkan

\mavergleichskette
{\vergleichskette
{ I
}
{ \subseteq }{ \R
}
{  }{
}
{    }{
}
{   }{
}
}
{}{}{}
suatu
\definitionsverweis {interval terbuka}{}{}
dengan fungsi positif yang terdiferensialkan secara kontinu
\maabbdisp {g} { I } { \R_+
} {,}
yang kita pandang sebagai
\definitionsverweis {metrik Riemann}{}{}
pada $I$. Karakterisasikan
\definitionsverweis {kurva-kurva geodesik}{}{}
\maabbdisp {\gamma} { J } { I
} {,}
dengan $J$ suatu interval lain.

}
{} {}

\inputaufgabegibtloesung
{}
{

Untuk kurva
\maabbeledisp {\psi} {\R} { {\mathbb H}
} { t } { \left(  t   , \,  e^t                   \right)
} {,}
yang bernilai dalam
\definitionsverweis {setengah bidang Poincaré}{}{}
${\mathbb H}$, tentukan
\definitionsverweis {percepatan tangensial}{}{}
untuk setiap $t$.

}
{} {}

\inputaufgabe
{}
{

Berikan contoh sebuah
\definitionsverweis {manifold Riemann}{}{}
yang
\definitionsverweis {terhubung}{}{}
$M$ dan dua titik

\mavergleichskette
{\vergleichskette
{ P,Q
}
{ \in }{ M
}
{  }{
}
{    }{
}
{   }{
}
}
{}{}{,}
yang tidak dapat dihubungkan oleh suatu
\definitionsverweis {kurva geodesik}{}{.}

}
{} {}

\inputaufgabe
{}
{

Berikan contoh sebuah
\definitionsverweis {manifold Riemann}{}{}
yang
\definitionsverweis {terhubung}{}{}
$M$ dan dua titik

\mavergleichskette
{\vergleichskette
{ P,Q
}
{ \in }{ M
}
{  }{
}
{    }{
}
{   }{
}
}
{}{}{,}
yang dapat dihubungkan oleh beberapa
\definitionsverweis {kurva geodesik}{}{}
yang semuanya mempunyai panjang sama.

}
{} {}

\inputaufgabe
{}
{

Berikan contoh sebuah
\definitionsverweis {manifold Riemann}{}{}
yang
\definitionsverweis {terhubung}{}{}
$M$ dan dua titik

\mavergleichskette
{\vergleichskette
{ P,Q
}
{ \in }{ M
}
{  }{
}
{    }{
}
{   }{
}
}
{}{}{,}
yang dapat dihubungkan oleh beberapa
\definitionsverweis {kurva geodesik}{}{}
yang panjangnya tidak semuanya sama.

}
{} {}

\inputaufgabegibtloesung
{}
{

Verifikasikan bahwa sistem persamaan diferensial biasa orde dua

\mavergleichskettedisp
{\vergleichskette
{ \gamma_1^{\prime \prime} (t)
}
{ =} { { \frac{   2  }{  \gamma_2(t) } } \gamma_1'(t) \gamma_2'(t)
}
{ } {
}
{ } {
}
{ } {
}
}
{}{}{}
dan

\mavergleichskettedisp
{\vergleichskette
{ \gamma_2^{\prime \prime} (t)
}
{ =} { -  { \frac{   1  }{  \gamma_2(t) } }  \gamma_1'(t) \gamma_1'(t) + { \frac{   1  }{  \gamma_2(t) } }    \gamma_2'(t) \gamma_2'(t)
}
{ } {
}
{ } {
}
{ } {
}
}
{}{}{}
diselesaikan oleh kurva-kurva

\mavergleichskettedisp
{\vergleichskette
{ \gamma(t)
}
{ =} { \begin{pmatrix}  s+r \tanh  t  \\ { \frac{   r  }{  \cosh  t  } }    \end{pmatrix}
}
{ } {
}
{ } {
}
{ } {
}
}
{}{}{}
untuk

\mavergleichskette
{\vergleichskette
{ s
}
{ \in }{ \R
}
{  }{
}
{    }{
}
{   }{
}
}
{}{}{,}

\mavergleichskette
{\vergleichskette
{ r
}
{ \in }{ \R_+
}
{  }{
}
{    }{
}
{   }{
}
}
{}{}{.}

}
{} {}

\inputaufgabe
{}
{

Kita tinjau lingkaran-lingkaran yang berpusat di
\mathl{(s,0)}{} dan berjari-jari

\mavergleichskette
{\vergleichskette
{ r
}
{ > }{ 0
}
{  }{
}
{    }{
}
{   }{
}
}
{}{}{.}
Karakterisasikan kapan dua lingkaran semacam ini berpotongan.

}
{} {}

\inputaufgabe
{}
{

Misalkan

\mavergleichskette
{\vergleichskette
{ P
}
{ = }{(a,b)
}
{ \in }{ {\mathbb H}
}
{    }{
}
{   }{
}
}
{}{}{}
suatu titik dalam
\definitionsverweis {setengah bidang Poincaré}{}{}
dan misalkan

\mavergleichskette
{\vergleichskette
{ v
}
{ = }{ (v_1,v_2)
}
{ \in }{ \R^2
}
{    }{
}
{   }{
}
}
{}{}{}
suatu
\definitionsverweis {vektor tangen}{}{}
di $P$.
\aufzaehlungzwei {Tentukan sebuah setengah lingkaran berpusat pada sumbu $x$
\zusatzklammer {dengan garis-garis vertikal sebagai kasus ekstrem} {} {,}
yang melalui $P$ dan menyinggung $\R v$.
} {Tentukan suatu
\definitionsverweis {kurva geodesik}{}{,}
yang merealisasikan vektor tangen ini.
}

}
{} {}

\inputaufgabe
{}
{

Misalkan
\maabb {\psi} { L } { M
} {}
suatu
\definitionsverweis {isometri}{}{}
antara
\definitionsverweis {manifold-manifold Riemann}{}{}
\mathkor {} {L} {dan} {M} {.}
\aufzaehlungdrei{Tunjukkan bahwa
\definitionsverweis {koneksi tarikan balik}{}{}
dari
\definitionsverweis {koneksi Levi-Civita}{}{}
pada $M$ sama dengan koneksi Levi-Civita pada $L$.
}{Tunjukkan bahwa
\definitionsverweis {percepatan-percepatan tangensial}{}{}
dari kurva yang dua kali terdiferensialkan
\maabb {\gamma} { I } { L
} {}
dan, berturut-turut,
\maabb {\psi \circ \gamma} { I } { M
} {}
saling bersesuaian.
}{Tunjukkan bahwa
\definitionsverweis {kurva-kurva geodesik}{}{}
di $L$ dan di $M$ saling bersesuaian.
}

}
{} {}

Misalkan $M$ suatu
\definitionsverweis {manifold Riemann}{}{,}
dengan
\definitionsverweis {bundel tangen}{}{}
$TM$ yang dilengkapi
\definitionsverweis {koneksi Levi-Civita}{}{.}
Misalkan
\maabb {\gamma} { I } { M
} {}
suatu
\definitionsverweis {kurva diferensiabel}{}{.}
Suatu objek diferensiabel yang didefinisikan sepanjang $\gamma$ dan merupakan
\definitionsverweis {medan vektor}{}{},
\maabb {F} { I } { TM
} {}
disebut
\definitionswort {paralel sepanjang}{}
$\gamma$ jika

\mavergleichskettedisp
{\vergleichskette
{ (\nabla_{\partial} F)(t)
}
{ =} { 0
}
{ } {
}
{ } {
}
{ } {
}
}
{}{}{}
untuk semua

\mavergleichskette
{\vergleichskette
{ t
}
{ \in }{ I
}
{  }{
}
{    }{
}
{   }{
}
}
{}{}{.}

Untuk soal berikut, bandingkan dengan
Soal 25.13.

\inputaufgabe
{}
{

Misalkan

\mavergleichskette
{\vergleichskette
{ Y
}
{ \subseteq }{ W
}
{ \subseteq }{ \R^n
}
{    }{
}
{   }{
}
}
{}{}{,}
dengan $W$
\definitionsverweis {terbuka}{}{,}
dan suatu
\definitionsverweis {hiperpermukaan diferensiabel}{}{}
yang juga kita pandang sebagai
\definitionsverweis {manifold Riemann}{}{}
melalui ruang sekitar $\R^n$ dengan
\definitionsverweis {hasil kali skalar standar}{}{.}
Misalkan
\maabb {\gamma} { I } { Y
} {}
suatu
\definitionsverweis {kurva diferensiabel}{}{}
dan
\maabb {F} { I } { TY \subseteq Y \times \R^n
} {}
suatu
\definitionsverweis {medan vektor}{}{}
diferensiabel sepanjang $\gamma$. Tunjukkan bahwa $F$
\definitionswort {paralel sepanjang}{}
$\gamma$ dalam arti
Definisi 6.6
jika dan hanya jika $F$
\definitionsverweis {paralel}{}{}
sepanjang $\gamma$ terhadap
\definitionsverweis {koneksi Levi-Civita}{}{}
pada $TY$.

}
{} {}

\inputaufgabe
{}
{

Misalkan $M$ suatu
\definitionsverweis {manifold Riemann}{}{,}
dengan
\definitionsverweis {bundel tangen}{}{}
$TM$ yang dilengkapi
\definitionsverweis {koneksi Levi-Civita}{}{.}
Misalkan
\maabb {\gamma} { I } { M
} {}
suatu kurva yang dua kali
\definitionsverweis {terdiferensialkan}{}{.}
Tunjukkan bahwa $\gamma$ adalah
\definitionsverweis {kurva geodesik}{}{}
jika dan hanya jika turunannya
\maabbdisp {\gamma'} { I } { TM
} {}
merupakan
\definitionsverweis {medan vektor paralel}{}{.}

}
{} {}

\inputaufgabe
{}
{

Tunjukkan bahwa pada
\definitionsverweis {ruang Euklides}{}{}
$\R^n$,
\definitionsverweis {metrik Riemann}{}{}
sama dengan
\definitionsverweis {jarak Euklides}{}{.}

}
{} {}

  \zwischenueberschrift{Tugas yang dikumpulkan}

\inputaufgabe
{4}
{

Kita tinjau $\R_+$ dengan struktur yang diberikan oleh

\mavergleichskette
{\vergleichskette
{ g(t)
}
{ = }{ t^2
}
{  }{
}
{    }{
}
{   }{
}
}
{}{}{,}
yaitu suatu
\definitionsverweis {metrik Riemann}{}{.}
Tentukan
\definitionsverweis {kurva-kurva geodesik}{}{}
di $\R_+$.

}
{} {}

\inputaufgabe
{6 (3+3)}
{

Kita tinjau kurva
\maabbeledisp {\psi} {\R} { {\mathbb H}
} { t } { \left(  t   , \,   1+t^2                   \right)
} {,}
yang bernilai dalam
\definitionsverweis {setengah bidang Poincaré}{}{}
${\mathbb H}$.
\aufzaehlungzwei {Tentukan
\definitionsverweis {simbol-simbol Christoffel}{}{}
untuk
\definitionsverweis {koneksi tarikan balik}{}{}
\zusatzklammer {dari
\definitionsverweis {koneksi Levi-Civita}{}{}
pada ${\mathbb H}$} {} {}
di atas $\R$.
} {Tentukan
\definitionsverweis {percepatan tangensial}{}{}
dari $\psi$ untuk setiap $t$.
}

}
{} {}

\inputaufgabe
{4}
{

Tentukan
\definitionsverweis {kurva-kurva geodesik}{}{}
pada
\definitionsverweis {cakram hiperbolik}{}{.}

}
{} {}

\inputaufgabe
{4}
{

Misalkan

\mavergleichskette
{\vergleichskette
{ n
}
{ \geq }{ 2
}
{  }{
}
{    }{
}
{   }{
}
}
{}{}{}
dan

\mavergleichskettedisp
{\vergleichskette
{ M
}
{ =} { \R^n \setminus \{P_1  , \ldots , P_m \}
}
{ \subseteq} { \R^n
}
{ } {
}
{ } {
}
}
{}{}{,}
yang dipandang sebagai
\definitionsverweis {manifold Riemann}{}{}
dengan
\definitionsverweis {hasil kali skalar standar}{}{}
terinduksi. Tunjukkan bahwa pada $M$,
\definitionsverweis {metrik Riemann}{}{}
sama dengan
\definitionsverweis {jarak Euklides}{}{.}

}
{} {}

\section*{Solusi yang disediakan oleh sumber}
\addcontentsline{toc}{section}{Solusi yang disediakan oleh sumber}
\subsection*{Solusi Soal 27.4}
\aufzaehlungfuenf{Berlaku

\mavergleichskettealign
{\vergleichskettealign
{ \Psi_Q(0)
}
{ =} { \left(  \cos  0   , \,   \sin  0  , \,   M(0) \begin{pmatrix}  a  \\b   \end{pmatrix}                 \right)
}
{ =} { \left(  1   , \,   0  , \,   \begin{pmatrix}   1   &   0   \\  0   &  1  \end{pmatrix}    \begin{pmatrix}  a  \\b   \end{pmatrix}                 \right)
}
{ =} { \left(  1   , \,   0  , \,     \begin{pmatrix}  a  \\b   \end{pmatrix}                 \right)
}
{ } {
}
}
{}
{}{.}
}{Berlaku

\mavergleichskettealign
{\vergleichskettealign
{ \Psi_Q(2 \pi )
}
{ =} { \left(  \cos  2 \pi   , \,   \sin  2 \pi  , \,   M(2 \pi) \begin{pmatrix}  a  \\b   \end{pmatrix}                 \right)
}
{ =} { \left(  1   , \,   0  , \,   \begin{pmatrix}  \cos  c   &   - \sin c    \\  \sin c   &   \cos  c     \end{pmatrix} \begin{pmatrix}  a  \\b   \end{pmatrix}                 \right)
}
{ =} { \left(  1   , \,   0  , \,     \begin{pmatrix} a \cos  c   - b \sin c   \\ a \sin  c   - b \cos c    \end{pmatrix}                 \right)
}
{ } {
}
}
{}
{}{.}
}{Berlaku

\mavergleichskettealign
{\vergleichskettealign
{  \Psi_Q(t)
}
{ =} { \left(  \cos  t    , \,   \sin  t  , \,   \begin{pmatrix}   \cos  { \frac{  ct  }{  2 \pi   } }   &   - \sin  { \frac{  ct  }{  2 \pi   } }    \\  \sin  { \frac{  ct  }{  2 \pi   } }   &   \cos  { \frac{  ct  }{  2 \pi   } }  \end{pmatrix}  \begin{pmatrix}  a  \\b   \end{pmatrix}                 \right)
}
{ =} { \left(  \cos  t   , \,   \sin  t  , \,   \begin{pmatrix}  a \cos  { \frac{  ct  }{  2 \pi  } }  - b \sin  { \frac{  ct  }{  2 \pi  } }   \\ a \sin  { \frac{  ct  }{  2 \pi  } }  +   b  \cos  { \frac{  ct  }{  2 \pi } }    \end{pmatrix}                 \right)
}
{ =} { \left(  \cos  t   , \,   \sin  t  , \,     \left(  a \cos  { \frac{  ct  }{  2 \pi  } }  - b \sin  { \frac{  ct  }{  2 \pi  } }  \right)   e_1 +  \left(  a \sin  { \frac{  ct  }{  2 \pi  } }  +   b  \cos  { \frac{  ct  }{  2 \pi } }  \right)   e_2                 \right)
}
{ } {
}
}
{}
{}{,}
khususnya, kita memperoleh suatu pengangkatan. Turunan vertikal penampang ini pada $\R^2$ di atas $\R$ dapat ditentukan dari simbol-simbol Christoffel yang diberikan; berlaku

\mavergleichskettealigndrucklinks
{\vergleichskettealigndrucklinks
{ \left(  \psi^* \nabla  \right)_\partial  \left(    \left(  a \cos  { \frac{  ct  }{  2 \pi  } }  - b \sin  { \frac{  ct  }{  2 \pi  } }  \right)   e_1 +  \left(  a \sin  { \frac{  ct  }{  2 \pi  } }  +   b  \cos  { \frac{  ct  }{  2 \pi } }  \right)   e_2  \right)
}
{ =} {   \left(  a \cos  { \frac{  ct  }{  2 \pi  } }  - b \sin  { \frac{  ct  }{  2 \pi  } }  \right)   \left(  \psi^* \nabla  \right)_\partial \left(  e_1  \right)  + \partial    \left(  a \cos  { \frac{  ct  }{  2 \pi  } }  - b \sin  { \frac{  ct  }{  2 \pi  } }  \right)   e_1    +   \left(  a \sin  { \frac{  ct  }{  2 \pi  } }  +   b  \cos  { \frac{  ct  }{  2 \pi } }  \right)      \left(  \psi^* \nabla  \right)_\partial \left(  e_2  \right) + \partial  \left(  a \sin  { \frac{  ct  }{  2 \pi  } }  +   b  \cos  { \frac{  ct  }{  2 \pi } }  \right)   e_2
}
{ =} { -  \left(  a \cos  { \frac{  ct  }{  2 \pi  } }  - b \sin  { \frac{  ct  }{  2 \pi  } }  \right)  { \frac{   c  }{  2 \pi } }  e_2  +   { \frac{   c  }{  2 \pi } } \left(  -a \sin  { \frac{  ct  }{  2 \pi  } }  - b \cos  { \frac{  ct  }{  2 \pi  } }  \right)   e_1    +   \left(  a \sin  { \frac{  ct  }{  2 \pi  } }  +   b  \cos  { \frac{  ct  }{  2 \pi } }  \right) { \frac{   c  }{  2 \pi } }   e_1    +  { \frac{   c  }{  2 \pi } }   \left(  a \cos  { \frac{  ct  }{  2 \pi  } } -  b  \sin  { \frac{  ct  }{  2 \pi } }  \right)   e_2
}
{ =} { 0
}
{ } {
}
}
{}{}{,}
jadi ini merupakan suatu
\definitionsverweis {pengangkatan horizontal}{}{}
.
}{Pemetaan yang dimaksud ialah
\maabbeledisp {} {\R^2} { \R^2
} { \begin{pmatrix}  a  \\b   \end{pmatrix} } {    \begin{pmatrix}  \cos  c   &   - \sin  c   \\  \sin  c  &  \cos  c  \end{pmatrix} \begin{pmatrix}  a  \\b   \end{pmatrix}
} {,}
yang linear. Matriks yang merepresentasikannya bersifat ortogonal
\zusatzklammer {sebuah matriks rotasi} {} {,}
sehingga pemetaan tersebut merupakan isometri.
}{Berlaku bahwa
\mathl{\Psi (2k \pi)}{} adalah pemetaan linear
\maabbeledisp {} {\R^2} { \R^2
} {\begin{pmatrix}  a  \\b   \end{pmatrix} } { \begin{pmatrix}  \cos kc   &   - \sin kc   \\  \sin kc  &  \cos kc  \end{pmatrix} \begin{pmatrix}  a  \\b   \end{pmatrix}
} {,}
yaitu rotasi dengan sudut $kc$. Matriks rotasi untuk jumlah dua sudut adalah hasil kali matriks kedua rotasi tersebut
\zusatzklammer {lihat
Fakta} {} {,}
sehingga kita memperoleh suatu
\definitionsverweis {homomorfisme grup}{}{}
.
}
\subsection*{Solusi Soal 27.5}
Satu-satunya turunan standar pada $\R$ adalah $\partial$, sehingga kita tidak menuliskan indeks bawah pada simbol-simbol Christoffel. Menurut
Fakta (2),
untuk

\mavergleichskette
{\vergleichskette
{ \ell,k
}
{ = }{ 1,2
}
{  }{
}
{    }{
}
{   }{
}
}
{}{}{}

berlaku

\mavergleichskettedisp
{\vergleichskette
{ \tilde{\Gamma}_\ell^k (t)
}
{ =} { \psi_1'(t) \Gamma_{1 \ell}^k (\psi(t)) + \psi'_2(t) \Gamma_{2 \ell}^k (\psi(t))
}
{ =} { \Gamma_{1 \ell}^k (t,e^t)   + e^t \Gamma_{2 \ell}^k (t,e^t)
}
{ } {}
{ } {}
}
{}{}{.}
Dari
contoh
kita memperoleh

\mavergleichskettedisp
{\vergleichskette
{ \tilde{\Gamma}_1^1 (t)
}
{ =} { \Gamma_{1 1}^1 (t,e^t) + e^t \Gamma_{2 1}^1 (t,e^t)
}
{ =} {  -e^t e^{-t}
}
{ =} { -1
}
{ } {
}
}
{}{}{,}

\mavergleichskettedisp
{\vergleichskette
{ \tilde{\Gamma}_1^2 (t)
}
{ =} { \Gamma_{1 1}^2 (t,e^t) + e^t \Gamma_{2 1}^2 (t,e^t)
}
{ =} { e^{-t}
}
{ } {
}
{ } {
}
}
{}{}{,}

\mavergleichskettedisp
{\vergleichskette
{ \tilde{\Gamma}_2^1 (t)
}
{ =} { \Gamma_{1 2}^1 (t,e^t) + e^t \Gamma_{2 2}^1 (t,e^t)
}
{ =} {  -e^{-t}
}
{ } {
}
{ } {
}
}
{}{}{}
dan

\mavergleichskettedisp
{\vergleichskette
{ \tilde{\Gamma}_2^2 (t)
}
{ =} { \Gamma_{1 2}^2 (t,e^t) + e^t \Gamma_{2 2}^2 (t,e^t)
}
{ =} { - e^t e^{-t}
}
{ =} { -1
}
{ } {
}
}
{}{}{.}
\subsection*{Solusi Soal 27.9}
Kita bekerja pada bundel vektor tarikan balik
\mathl{\R \times \R^2}{} di atas $\R$ dengan
\definitionsverweis {koneksi tarikan balik}{}{,}
yang
\definitionsverweis {simbol-simbol Christoffel}{}{}
$\tilde{\Gamma}_\ell^k$-nya telah ditentukan dalam
soal.
Misalkan $e_1,e_2$ adalah penampang-penampang basis. Maka, menurut
fakta,

\mavergleichskettealign
{\vergleichskettealign
{ \nabla_\partial \left(   \psi_1' e_1 + \psi_2' e_2  \right)
}
{ =} { D_\partial (\psi_1') e_1 + \psi_1' \nabla_\partial e_1 +  D_\partial (\psi_2') e_2 + \psi_2' \nabla_\partial e_2
}
{ =} { \psi_1^{\prime \prime} e_1 + \psi_1' \nabla_\partial e_1 +  \gamma_2^{\prime \prime} e_2 + \psi_2' \nabla_\partial e_2
}
{ =} { \psi_1^{\prime \prime} e_1 + \psi_1' \left(   \tilde{\Gamma}_1^1 e_1 +  \tilde{\Gamma}_1^2  e_2  \right)  + \psi_2^{\prime \prime} e_2 + \psi_2' \left(   \tilde{\Gamma}_2^1 e_1 +  \tilde{\Gamma}_2^2  e_2  \right)
}
{ =} { \left(   \psi_1^{\prime \prime} +  \psi_1'  \tilde{\Gamma}_1^1    +   \psi_2'  \tilde{\Gamma}_2^1     \right)  e_1 +  \left(    \psi_2^{\prime \prime} +  \psi_1'  \tilde{\Gamma}_1^2    +   \psi_2'  \tilde{\Gamma}_2^2   \right) e_2
}
}
{}
{}{.}
Dengan demikian,

\mavergleichskettealign
{\vergleichskettealign
{  \nabla_\partial \left(   \psi_1' e_1 + \psi_2' e_2  \right)  (t)
}
{ =} { \left(    \tilde{\Gamma}_1^1 (t)    +   e^t  \tilde{\Gamma}_2^1(t)  \right)  e_1 +  \left(   e^t +    \tilde{\Gamma}_1^2 (t) + e^t \tilde{\Gamma}_2^2(t)     \right) e_2
}
{ =} { \left(   -1  -   e^t  e^{-t}  \right)  e_1 +  \left(   e^t +   e^{-t} - e^t   \right) e_2
}
{ =} { -2e_1 +e^{-t} e_2
}
{ } {
}
}
{}
{}{.}
\subsection*{Solusi Soal 27.13}
Berlaku

\mavergleichskettealign
{\vergleichskettealign
{ \gamma'(t)
}
{ =} { \begin{pmatrix} s+r { \frac{   \sinh  t  }{  \cosh  t  } }   \\ { \frac{   r  }{  \cosh  t  } }    \end{pmatrix}'
}
{ =} { r \begin{pmatrix}  { \frac{   \cosh  t   -\sinh  t  }{  \cosh  t  } }   \\ - { \frac{  \sinh  t   }{  \cosh  t  } }    \end{pmatrix}
}
{ =} { r \begin{pmatrix}  { \frac{    1  }{  \cosh  t  } }   \\ - { \frac{  \sinh  t   }{  \cosh  t  } }    \end{pmatrix}
}
{ =} { { \frac{   r  }{  \cosh  t  } }\begin{pmatrix}  1   \\ - \sinh  t     \end{pmatrix}
}
}
{}
{}{}
dan

\mavergleichskettealign
{\vergleichskettealign
{ \gamma^{\prime \prime}(t)
}
{ =} { r \begin{pmatrix}  { \frac{    1  }{  \cosh  t  } }   \\ - { \frac{  \sinh  t   }{  \cosh  t  } }    \end{pmatrix}'
}
{ =} { r \begin{pmatrix}  { \frac{   -2 \sinh  t   }{  \cosh  t  } }   \\ { \frac{   - \cosh  t + 2 \sinh  t \cosh  t  }{  \cosh  t  } }    \end{pmatrix}
}
{ =} { r \begin{pmatrix}  { \frac{   -2 \sinh  t   }{  \cosh  t  } }   \\ { \frac{   - \cosh  t + 2 \sinh  t  }{  \cosh  t  } }    \end{pmatrix}
}
{ } {}
}
{}
{}{.}
Karena

\mavergleichskettedisp
{\vergleichskette
{  { \frac{   2 }{  { \frac{   r  }{  \cosh  t  } }  } } \cdot    { \frac{   r  }{ \cosh  t  } } \cdot (-1) { \frac{  r \sinh  t  }{ \cosh  t  } }
}
{ =} { - { \frac{   2 r \sinh  t   }{ \cosh  t  } }
}
{ =} { \gamma_1^{\prime \prime} (t)
}
{ } {
}
{ } {
}
}
{}{}{}
persamaan pertama terpenuhi. Karena

\mavergleichskettealign
{\vergleichskettealign
{ - { \frac{   1 }{  \gamma_2(t) } }  \gamma_1'(t) \gamma_1'(t) + { \frac{   1 }{  \gamma_2(t) } }  \gamma_2'(t) \gamma_2'(t)
}
{ =} { { \frac{   \cosh  t  }{ r } } \left(  - \left(  { \frac{   r  }{  \cosh  t  } }  \right)^2 + \left(  { \frac{  r \sinh  t  }{  \cosh  t  } }  \right)^2  \right)
}
{ =} { { \frac{   r \left(  -1 + \sinh  t  \right)  }{  \cosh  t  } }
}
{ =} { { \frac{   r \left(  - \cosh  t +2 \sinh  t  \right)  }{ \cosh  t  } }
}
{ =} { \gamma_2^{\prime \prime} (t)
}
}
{}
{}{}
maka persamaan kedua juga terpenuhi.

\part{Unit 28}
\chapter[Kuliah 28]{Kuliah 28: Kelengkungan dan Teorema Frobenius}
\setcounter{section}{28}

  \zwischenueberschrift{Kelengkungan}

Misalkan
\mathkor {} {V} {dan} {W} {}
adalah medan vektor yang dua kali terdiferensialkan secara kontinu pada suatu manifold diferensiabel $M$, dan misalkan $\nabla$ adalah
\definitionsverweis {koneksi linear}{}{}
pada
\definitionsverweis {bundel vektor}{}{}
\maabb {p} { E } { M
} {}
di atas $M$. Pemetaan
\maabbeledisp {\nabla_W} {C^2(E) } { C^1 (E)
} { s } { \nabla_W (s)
} {,}
adalah
\definitionsverweis {turunan vertikal}{}{.}
Pada hasil $\nabla_W (s)$ kita dapat sekali lagi menerapkan turunan vertikal dalam arah $V$, sehingga diperoleh penampang kontinu
\mathl{\nabla_V  \left(  \nabla_W (s)  \right)}{.} Jika semua objek berkelas $C^\infty$, derajat keterdiferensialannya tidak lagi perlu diperhatikan. Di sini kita meninjau ekspresi

\mavergleichskettedisp
{\vergleichskette
{ R(V,W)
}
{ =} { \nabla_V \nabla_W - \nabla_W \nabla_V- \nabla_{[V,W]}
}
{ } {
}
{ } {
}
{ } {
}
}
{}{}{,}
yang mengirim penampang
\zusatzklammer {yang cukup terdiferensialkan} {} {}
dari $E$ kembali ke penampang dari $E$. Di sini
\mathl{[V,W]}{} adalah
\definitionsverweis {braket Lie}{}{}
dari kedua medan vektor tersebut.

\inputdefinition
{}
{

Misalkan $M$ adalah
\definitionsverweis {manifold diferensiabel}{}{}
yang dua kali terdiferensialkan secara kontinu, dan
\maabb {} { E } { M
} {}
adalah
\definitionsverweis {bundel vektor diferensiabel}{}{}
yang dua kali terdiferensialkan secara kontinu di atas $M$, dilengkapi dengan
\definitionsverweis {koneksi linear}{}{}
$\nabla$. Untuk
\definitionsverweis {medan vektor}{}{}
diferensiabel
\mathkor {} {V} {dan} {W} {}
pada $M$, pemetaan
\maabbeledisp {} {C^2(E)} {C^0(E)
} { s } { \nabla_V \nabla_W (s) - \nabla_W \nabla_V (s) - \nabla_{[V,W]}(s)
} {,}
disebut
\definitionswort {operator kelengkungan}{}
untuk $V$ dan $W$. Operator ini dinotasikan dengan $R(V,W)$.

}



\inputfaktbeweis
{Reelles Vektorbündel/Trivial/Linearer Zusammenhang/Krümmungsoperator/Beschreibung/Fakt}
{Lema}
{}
{

\faktsituation {Misalkan

\mavergleichskette
{\vergleichskette
{ U
}
{ \subseteq }{ \R^d
}
{  }{
}
{    }{
}
{   }{
}
}
{}{}{}
adalah himpunan
\definitionsverweis {terbuka}{}{}
dan $\nabla$ adalah
\definitionsverweis {koneksi linear}{}{}
pada
\definitionsverweis {bundel vektor}{}{}
trivial

\mavergleichskette
{\vergleichskette
{ E
}
{ = }{ U \times \R^r
}
{  }{
}
{    }{
}
{   }{
}
}
{}{}{.}
Misalkan

\mathbed {\partial_i} {}
 {1 \leq i \leq d} {}
 {} {} {} {}
adalah
\definitionsverweis {medan vektor standar}{}{}
dan

\mathbed {s_j} {}
 {1 \leq j \leq r} {}
 {} {} {} {}
adalah penampang standar di $E$.}
\faktfolgerung {Untuk
\definitionsverweis {operator kelengkungan}{}{}
berlaku

\mavergleichskettedisp
{\vergleichskette
{ R( \partial_i , \partial_j)(s_k)
}
{ =} { \sum_{\ell = 1}^r  R^\ell_{ijk} s_\ell
}
{ } {
}
{ } {
}
{ } {
}
}
{}{}{}
dengan

\mavergleichskettedisp
{\vergleichskette
{ R^\ell_{ijk}
}
{ =} { {\partial_i } \Gamma^\ell_{jk} - {\partial_j } \Gamma^\ell_{ik}  +  \sum_{ a  = 1}^r  \Gamma^a_{jk} \Gamma_{i a}^\ell -\sum_{ a  = 1}^r   \Gamma^a_{ik}  \Gamma_{j a}^\ell
}
{ } {
}
{ } {
}
{ } {
}
}
{}{}{,}
di mana $\Gamma^\ell_{ik}$ menyatakan
\definitionsverweis {simbol Christoffel}{}{}
lokal dari koneksi tersebut.}
\faktzusatz {}
\faktzusatz {}

}
{

Kita mempunyai

\mavergleichskettealignhandlinks
{\vergleichskettealignhandlinks
{ R( \partial_i ,\partial_j)
}
{ =} { \nabla_{\partial_i } \circ \nabla_{\partial_j } -  \nabla_{\partial_j } \circ \nabla_{\partial_i } - \nabla_{[\partial_i, \partial_j]}
}
{ =} { \nabla_{\partial_i } \circ \nabla_{\partial_j } -  \nabla_{\partial_j } \circ \nabla_{\partial_i }
}
{ } {
}
{ } {
}
}
{}
{}{,}
sebab turunan parsial saling komutatif dan karena itu braket Lie-nya sama dengan $0$. Maka, untuk penampang standar $s_k$, dengan menggunakan
Lema 25.10,

\mavergleichskettealignhandlinks
{\vergleichskettealignhandlinks
{ R( \partial_i ,\partial_j) (s_k)
}
{ =} { \nabla_{\partial_i } ( \nabla_{\partial_j } (s_k) ) -  \nabla_{\partial_j } ( \nabla_{\partial_i } (s_k))
}
{ =} { \nabla_{\partial_i }  \left(  \sum_{\ell = 1}^r \Gamma^\ell_{jk} s_\ell  \right)  -  \nabla_{\partial_j }  \left(  \sum_{\ell = 1}^r \Gamma^\ell_{ik} s_\ell  \right)
}
{ =} { \sum_{\ell = 1}^r  \left(  \nabla_{\partial_i }  \left(  \Gamma^\ell_{jk} s_\ell  \right) - \nabla_{\partial_j }  \left(  \Gamma^\ell_{ik} s_\ell  \right)   \right)
}
{ =} { \sum_{\ell = 1}^r  \left(  {\partial_i } \Gamma^\ell_{jk} s_\ell +  \Gamma^\ell_{jk} \sum_{ \rho  = 1}^r \Gamma_{i \ell}^\rho  s_\rho  - {\partial_j } \Gamma^\ell_{ik} s_\ell - \Gamma^\ell_{ik} \sum_{ \rho  = 1}^r \Gamma_{j \ell}^\rho  s_\rho   \right)
}
}
{
\vergleichskettefortsetzungalign
{ =} { \sum_{\ell = 1}^r  \left(  {\partial_i } \Gamma^\ell_{jk} - {\partial_j } \Gamma^\ell_{ik}  \right) s_\ell  + \sum_{\ell = 1}^r \sum_{ \rho  = 1}^r  \left(  \Gamma^\ell_{jk} \Gamma_{i \ell}^\rho  s_\rho  - \Gamma^\ell_{ik}  \Gamma_{j \ell}^\rho  s_\rho  \right)
}
{ =} { \sum_{\ell = 1}^r  \left(  {\partial_i } \Gamma^\ell_{jk} - {\partial_j } \Gamma^\ell_{ik}  \right) s_\ell  + \sum_{\ell = 1}^r  \left(  \sum_{ a  = 1}^r  \Gamma^a_{jk} \Gamma_{i a}^\ell -\sum_{ a  = 1}^r   \Gamma^a_{ik}  \Gamma_{j a}^\ell  \right)  s_\ell
}
{ } {
}
{ } {}
}
{}{.}

}



\inputfaktbeweis
{Reelles Vektorbündel/Rang 1/Offene Menge/Krümmungsoperator/Fakt}
{Akibat}
{}
{

\faktsituation {Misalkan

\mavergleichskette
{\vergleichskette
{ U
}
{ \subseteq }{ \R^n
}
{  }{
}
{    }{
}
{   }{
}
}
{}{}{}
adalah
\definitionsverweis {himpunan terbuka}{}{}
dan $\nabla$ adalah
\definitionsverweis {koneksi linear}{}{}
pada bundel vektor trivial
\mathl{U \times \R}{} yang memiliki
\definitionsverweis {rank}{}{}
$1$, dengan
\definitionsverweis {simbol Christoffel}{}{}

\mathbed {\Gamma_i} {}
 {i = 1 , \ldots , n} {}
 {} {} {} {.}}
\faktuebergang {Pernyataan berikut berlaku bagi
\definitionsverweis {operator kelengkungan}{}{}
terhadap
\definitionsverweis {medan vektor standar}{}{}
$\partial_i$.}
\faktfolgerung {\aufzaehlungzwei {Berlaku

\mavergleichskettedisp
{\vergleichskette
{ R(\partial_i, \partial_j)
}
{ =} { \partial_i \Gamma_j - \partial_j \Gamma_i
}
{ } {
}
{ } {
}
{ } {
}
}
{}{}{.}
} {Persamaan

\mavergleichskette
{\vergleichskette
{ R(\partial_i, \partial_j)
}
{ = }{ 0
}
{  }{
}
{    }{
}
{   }{
}
}
{}{}{}
berlaku untuk semua $i,j$ jika dan hanya jika
$1$-\definitionsverweis {bentuk diferensial}{}{}
$\sum_{i = 1}^n \Gamma_i dx_i$
\definitionsverweis {tertutup}{}{.}
}}
\faktzusatz {}
\faktzusatz {}

}
{

\aufzaehlungzwei {Karena rank bundel sama dengan $1$, kita tuliskan

\mavergleichskette
{\vergleichskette
{ \Gamma_i
}
{ = }{ \Gamma_{i1}^1
}
{  }{
}
{    }{
}
{   }{
}
}
{}{}{}
dan, secara serupa,

\mavergleichskette
{\vergleichskette
{ R_{ij}
}
{ = }{ R_{ij1}^1
}
{  }{
}
{    }{
}
{   }{
}
}
{}{}{,}
sehingga pernyataan tersebut harus dipahami sebagai

\mavergleichskettedisp
{\vergleichskette
{ R(\partial_i, \partial_j) (e)
}
{ =} { \left(  \partial_i \Gamma_j - \partial_j \Gamma_i  \right) e
}
{ } {
}
{ } {
}
{ } {
}
}
{}{}{}
dengan $e$ penampang satuan konstan. Berdasarkan deskripsi eksplisit dalam
Lema 28.2,
kita memperoleh

\mavergleichskettedisp
{\vergleichskette
{ R( \partial_i , \partial_j)
}
{ =} { {\partial_i } \Gamma_{j} - {\partial_j } \Gamma_{i}  + \Gamma_{j} \Gamma_{i } - \Gamma_{i}  \Gamma_{j }
}
{ =} { {\partial_i } \Gamma_{j} - {\partial_j } \Gamma_{i}
}
{ } {
}
{ } {
}
}
{}{}{.}
} {Pernyataan ini mengikuti dari (1) dan
Soal 20.8.
}

}



\inputfaktbeweis
{Reelles Vektorbündel/Linearer Zusammenhang/Krümmungsoperator/Eigenschaften/Fakt}
{Lema}
{}
{

\faktsituation {Misalkan $M$ adalah
\definitionsverweis {manifold diferensiabel}{}{}
yang dua kali terdiferensialkan secara kontinu, dan
\maabb {} { E } { M
} {}
adalah
\definitionsverweis {bundel vektor diferensiabel}{}{}
yang dua kali terdiferensialkan secara kontinu di atas $M$, dilengkapi dengan
\definitionsverweis {koneksi linear}{}{}
$\nabla$.}
\faktuebergang {Maka
\definitionsverweis {operator kelengkungan}{}{}
memiliki sifat-sifat berikut.}
\faktfolgerung {\aufzaehlungdrei{Berlaku

\mavergleichskettedisp
{\vergleichskette
{ R(V,W)
}
{ =} { -R(W,V)
}
{ } {
}
{ } {
}
{ } {
}
}
{}{}{.}
}{ $R$ aditif dalam kedua komponennya.
}{Untuk setiap fungsi diferensiabel
\maabb {f} { M } { \R
} {}
berlaku

\mavergleichskettedisp
{\vergleichskette
{ R(fV,W)
}
{ =} { f R(V,W)
}
{ } {
}
{ } {
}
{ } {
}
}
{}{}{.}
}}
\faktzusatz {}
\faktzusatz {}

}
{

Sifat-sifat tersebut lokal pada manifold. Karena itu, cukup membuktikannya pada suatu tutupan terbuka yang setiap anggotanya difeomorfik dengan himpunan terbuka di $\R^d$ dan di atasnya bundel vektor tersebut trivial. Dengan demikian, (1) mengikuti dari
Lema 28.2.
Pernyataan (2) mengikuti dari
Lema 24.10
dan
Lema 11.8.
Untuk (3), tinjau situasi lokal dengan medan vektor $f \partial_i$ dan $g \partial_j$. Kita mempunyai

\mavergleichskettedisp
{\vergleichskette
{ R(f \partial_i, g \partial_j)
}
{ =} { \nabla_{f \partial_i } \circ \nabla_{g \partial_j } - \nabla_{g\partial_j } \circ \nabla_{f \partial_i } - \nabla_{[ f \partial_i, g \partial_j]}
}
{ } {
}
{ } {
}
{ } {
}
}
{}{}{.}
Menurut
Lema 11.8,

\mavergleichskettedisp
{\vergleichskette
{ [f \partial_i, g\partial_j]
}
{ =} { f [\partial_i, g\partial_j] - g \partial_j(f) \partial_i
}
{ } {
}
{ } {
}
{ } {
}
}
{}{}{.}
Oleh karena itu, untuk suatu penampang $s$, dengan menggunakan
Teorema 25.5,

\mavergleichskettealignhandlinks
{\vergleichskettealignhandlinks
{ R(f \partial_i, g\partial_j) (s)
}
{ =} { \nabla_{f \partial_i } \circ \nabla_{g \partial_j } (s)- \nabla_{g \partial_j } \circ \nabla_{f \partial_i } (s)- \nabla_{[ f \partial_i, g \partial_j]}(s)
}
{ =} { f  \nabla_{ \partial_i } ( \nabla_{g\partial_j } (s) ) - \nabla_{g\partial_j } ( \nabla_{f \partial_i } (s))  - \nabla_{ f [\partial_i, g\partial_j] - g \partial_j(f) \partial_i  } (s)
}
{ =} { f  \nabla_{ \partial_i } ( \nabla_{ g \partial_j } (s) )-\nabla_{ g \partial_j } (f \nabla_{ \partial_i } (s)) - \nabla_{f [\partial_i, g\partial_j] } (s)  + \nabla_{ g \partial_j (f) \partial_i }  (s)
}
{ =} { f  \nabla_{ \partial_i } ( \nabla_{ g \partial_j } (s) )- f \nabla_{ g \partial_j } ( \nabla_{ \partial_i } (s))- g \partial_j(f) \nabla_{\partial_i } (s)  - \nabla_{f [\partial_i, g\partial_j] } (s)          + g \partial_j( f)  \nabla_{\partial_i } (s)
}
}
{
\vergleichskettefortsetzungalign
{ =} { f \nabla_{ \partial_i } ( \nabla_{g \partial_j } (s) )- f \nabla_{ g \partial_j } ( \nabla_{ \partial_i } (s)) - f \nabla_{[\partial_i, g \partial_j]} (s)
}
{ =} { f R( \partial_i, \partial_j) (s)
}
{ } {
}
{ } {
}
}
{}{.}
Bersama (1) dan (2), perhitungan ini membuktikan pernyataan tersebut.

}



\inputfaktbeweis
{Eindimensionale Mannigfaltigkeit/Reelles Vektorbündel/Linearer Zusammenhang/Krümmungsoperator trivial/Fakt}
{Akibat}
{}
{

\faktsituation {Misalkan $M$ adalah
\definitionsverweis {manifold diferensiabel}{}{}
berdimensi satu dan berkelas $C^2$, serta
\maabb {} { E } { M
} {}
adalah
\definitionsverweis {bundel vektor diferensiabel}{}{}
yang dua kali terdiferensialkan secara kontinu di atas $M$, dilengkapi dengan
\definitionsverweis {koneksi linear}{}{}
$\nabla$.}
\faktfolgerung {Maka
\definitionsverweis {operator kelengkungan}{}{}
bersifat trivial untuk sebarang
\definitionsverweis {medan vektor}{}{}
$V,W$.}
\faktzusatz {}
\faktzusatz {}

}
{

Karena pernyataan bersifat lokal, kita dapat langsung mengambil

\mavergleichskette
{\vergleichskette
{ M
}
{ = }{ I
}
{  }{
}
{    }{
}
{   }{
}
}
{}{}{}
dengan suatu interval terbuka, serta

\mavergleichskette
{\vergleichskette
{ V
}
{ = }{ f \partial
}
{  }{
}
{    }{
}
{   }{
}
}
{}{}{,}

\mavergleichskette
{\vergleichskette
{ W
}
{ = }{ g \partial
}
{  }{
}
{    }{
}
{   }{
}
}
{}{}{}
dengan
\definitionsverweis {medan vektor standar}{}{}
$\partial$. Berdasarkan
Lema 28.4  (3),

\mavergleichskettedisp
{\vergleichskette
{ R(f \partial, g \partial)
}
{ =} { fg R( \partial, \partial)
}
{ =} { 0
}
{ } {
}
{ } {
}
}
{}{}{}
karena

\mavergleichskette
{\vergleichskette
{ [\partial, \partial ]
}
{ = }{ 0
}
{  }{
}
{    }{
}
{   }{
}
}
{}{}{.}

}



\inputfaktbeweis
{Reelles Vektorbündel/Trivialer Zusammenhang/Trivialer Krümmungsoperator/Fakt}
{Akibat}
{}
{

\faktsituation {Misalkan $M$ adalah
$C^2$-\definitionsverweis {manifold diferensiabel}{}{}
dan

\mavergleichskette
{\vergleichskette
{ E
}
{ = }{ M \times \R^r
}
{  }{
}
{    }{
}
{   }{
}
}
{}{}{}
adalah
\definitionsverweis {bundel vektor trivial}{}{}
dengan
\definitionsverweis {koneksi trivial}{}{.}}
\faktfolgerung {Maka, untuk sebarang
\definitionsverweis {medan vektor}{}{}
$V,W$,
\definitionsverweis {operator kelengkungan}{}{}
$R(V,W)$ adalah trivial.}
\faktzusatz {}
\faktzusatz {}

}
{

Pernyataan ini bersifat lokal, sehingga kita dapat mengambil

\mavergleichskette
{\vergleichskette
{ M
}
{ = }{ U
}
{ \subseteq }{ \R^n
}
{    }{
}
{   }{
}
}
{}{}{}
dan himpunan tersebut terbuka. Menurut
Lema 28.4  (3),
kita dapat mereduksi lebih lanjut ke kasus ketika medan-medannya adalah
\definitionsverweis {medan vektor standar}{}{}
$\partial_i$. Menurut
Soal 25.11,
semua simbol Christoffel koneksi trivial sama dengan nol. Karena itu, pernyataan mengikuti dari
Lema 28.2.

}

  \zwischenueberschrift{Teorema Frobenius}

Kita tidak dapat membuktikan teorema berikut secara lengkap. Perhatikan bahwa dari simbol Christoffel terhadap penampang basis yang diberikan
\zusatzklammer {secara lokal} {} {}
saja, tidak terlihat apakah simbol-simbol itu akan lenyap terhadap pilihan penampang basis lain. Keadaan ini tepat dicirikan oleh lenyapnya kelengkungan.


\inputfaktbeweis
{Differenzierbare Mannigfaltigkeit/Vektorbündel/Linearer Zusammenhang/Lokal integrabel/Krümmung/Fakt}
{Teorema}
{}
{

\faktsituation {Misalkan $M$ adalah
\definitionsverweis {manifold diferensiabel}{}{}
dan
\maabbdisp {p} { E } { M
} {}
adalah
\definitionsverweis {bundel vektor diferensiabel}{}{,}
yang dilengkapi dengan
\definitionsverweis {koneksi linear}{}{}
$\nabla$.}
\faktuebergang {Pernyataan-pernyataan berikut ekuivalen.}
\faktfolgerung {\aufzaehlungvier{Koneksi tersebut
\definitionsverweis {terintegralkan secara lokal}{}{.}
}{Secara lokal, terhadap
\definitionsverweis {penampang basis}{}{}
yang sesuai, koneksi tersebut isomorfik dengan
\definitionsverweis {koneksi trivial}{}{.}
}{Untuk setiap titik

\mavergleichskette
{\vergleichskette
{ P
}
{ \in }{ M
}
{  }{
}
{    }{
}
{   }{
}
}
{}{}{}
terdapat lingkungan peta terbuka

\mavergleichskette
{\vergleichskette
{ P
}
{ \in }{ U
}
{  }{
}
{    }{
}
{   }{
}
}
{}{}{}
sedemikian sehingga $E \vert_U$ trivial dan
\definitionsverweis {simbol Christoffel}{}{}
yang bersesuaian semuanya merupakan fungsi nol.
}{Untuk sebarang
\definitionsverweis {medan vektor}{}{},
\definitionsverweis {operator kelengkungan}{}{}
adalah trivial.
}}
\faktzusatz {}
\faktzusatz {}

}
{

Semua pernyataan bersifat lokal. Karena itu, kita dapat langsung mulai dengan himpunan terbuka

\mavergleichskette
{\vergleichskette
{ U
}
{ \subseteq }{ \R^n
}
{  }{
}
{    }{
}
{   }{
}
}
{}{}{}
dan bundel vektor trivial $U \times \R^r$. Andaikan (1) berlaku. Kita dapat pula mengandaikan bahwa pada $U$ diberikan keluarga yang terdiri dari $r$
\definitionsverweis {penampang horizontal}{}{}

\mathl{s_1 , \ldots , s_r}{} yang bebas linear. Menurut
Soal 25.6
dan
Soal 25.7,
koneksi itu kemudian dapat ditrivialkan secara lokal. Ini membuktikan (2). Ekuivalensi (2) dan (3) mengikuti dari
Soal 25.11
serta
Catatan 25.8.
Jika simbol Christoffel sama dengan nol, maka penampang standar bersifat horizontal dan secara lokal membentuk penampang basis horizontal. Jadi, ketiga rumusan pertama ekuivalen. Jika (2) berlaku, kita dapat menerapkan
Akibat 28.6
secara lokal dan memperoleh bahwa operator kelengkungan trivial.

Arah dari (4) ke (1) jauh lebih sulit.

}

Untuk kasus-kasus khusus dari arah (4) ke (1), lihat
Soal 28.4
dan
Soal 28.5.

\chapter{Lembar Kerja 28}
\setcounter{section}{28}

  \zwischenueberschrift{Soal latihan}

\inputaufgabe
{}
{

Kita tinjau
\definitionsverweis {koneksi trivial}{}{}
pada bundel vektor trivial
\maabbdisp {} {\R^2 \times \R } { \R^2
} {}
di atas $\R^2$. Misalkan diberikan
\definitionsverweis {medan-medan vektor}{}{}

\mavergleichskettedisp
{\vergleichskette
{ V
}
{ =} { y^3 \partial_1 + xy\partial_2
}
{ } {
}
{ } {
}
{ } {
}
}
{}{}{}
dan

\mavergleichskettedisp
{\vergleichskette
{ W
}
{ =} { \left(  x^2- y  \right)  \partial_1 + 4x^3y^2\partial_2
}
{ } {
}
{ } {
}
{ } {
}
}
{}{}{}
serta

\mavergleichskettedisp
{\vergleichskette
{ f
}
{ =} { x^3-xy^2+y^4
}
{ } {
}
{ } {
}
{ } {
}
}
{}{}{.}
Hitung fungsi atau medan vektor berikut.
\aufzaehlungsechs{
\mathl{\nabla_V f}{,}
}{
\mathl{\nabla_W f}{,}
}{
\mathl{[V,W]}{,}
}{
\mathl{\nabla_W  \nabla_V f}{,}
}{
\mathl{\nabla_V \nabla_W f}{,}
}{
\mathl{\nabla_{[V,W]} (f)}{.}
}

}
{} {}

\inputaufgabegibtloesung
{}
{

Pada
\definitionsverweis {bundel vektor trivial}{}{}
\maabbdisp {p} {\R^2 \times \R} { \R^2
} {}
ber-
\definitionsverweis {rank}{}{}
$1$ di atas $\R^2$, tinjau
\definitionsverweis {koneksi linear}{}{,}
yang diberikan oleh
\definitionsverweis {simbol-simbol Christoffel}{}{}

\mavergleichskette
{\vergleichskette
{ \Gamma_1 (x,y)
}
{ = }{ x^5-x^2y^2+3y^4
}
{  }{
}
{    }{
}
{   }{
}
}
{}{}{}
dan

\mavergleichskette
{\vergleichskette
{ \Gamma_2 (x,y)
}
{ = }{ 7x^3 +xy^2-4y^5
}
{  }{
}
{    }{
}
{   }{
}
}
{}{}{.}
Hitung
\definitionsverweis {operator kelengkungan}{}{}
$R \left(  \partial_1, \partial_2  \right)$.

}
{} {}

\inputaufgabe
{}
{

Pada
\definitionsverweis {bundel vektor trivial}{}{}
\maabbdisp {p} {\R^2 \times \R} { \R^2
} {}
ber-
\definitionsverweis {rank}{}{}
$1$ di atas $\R^2$, tinjau
\definitionsverweis {koneksi linear}{}{,}
yang diberikan oleh
\definitionsverweis {simbol-simbol Christoffel}{}{}

\mavergleichskette
{\vergleichskette
{ \Gamma_1 (x,y)
}
{ = }{ 2x^3-xy^2+  \sinh  \left(  xy^3  \right)
}
{  }{
}
{    }{
}
{   }{
}
}
{}{}{}
dan

\mavergleichskette
{\vergleichskette
{ \Gamma_2 (x,y)
}
{ = }{ x y^4 -   \exp  \left(  x^3-y^4 \right)
}
{  }{
}
{    }{
}
{   }{
}
}
{}{}{.}
Hitung
\definitionsverweis {operator kelengkungan}{}{}
$R \left(  \partial_1, \partial_2  \right)$.

}
{} {}

\inputaufgabe
{}
{

Misalkan $M$ suatu
$C^2$-\definitionsverweis {manifold diferensiabel}{}{}
berdimensi satu, dan misalkan
\maabb {} { E } { M
} {}
suatu
\definitionsverweis {bundel vektor diferensiabel}{}{}
dua kali terdiferensialkan secara kontinu di atas $M$, yang dilengkapi dengan
\definitionsverweis {koneksi linear}{}{}
$\nabla$.

}
{} {}

\inputaufgabegibtloesung
{}
{

Misalkan
\maabb {} { E } { M
} {}
suatu
\definitionsverweis {bundel vektor diferensiabel}{}{}
ber-
\definitionsverweis {rank}{}{}
$1$ di atas suatu
\definitionsverweis {manifold diferensiabel}{}{}
$M$. Misalkan $\nabla$ suatu
\definitionsverweis {koneksi linear}{}{}
pada $E$ dengan
\definitionsverweis {operator kelengkungan}{}{}
trivial. Tunjukkan bahwa $\nabla$
\definitionsverweis {terintegralkan secara lokal}{}{.}

}
{} {}

  \zwischenueberschrift{Soal untuk dikumpulkan}

\inputaufgabe
{3}
{

Kita tinjau
\definitionsverweis {koneksi trivial}{}{}
pada bundel vektor trivial
\maabbdisp {} {\R^2 \times \R } { \R^2
} {}
di atas $\R^2$. Misalkan diberikan
\definitionsverweis {medan-medan vektor}{}{}

\mavergleichskettedisp
{\vergleichskette
{ V
}
{ =} { 6xy^4 \partial_1 -3 y^2 e^{xy} \partial_2
}
{ } {
}
{ } {
}
{ } {
}
}
{}{}{}
dan

\mavergleichskettedisp
{\vergleichskette
{ W
}
{ =} { \left(  x+ y^3  \right)  \partial_1 + 2x^2y \partial_2
}
{ } {
}
{ } {
}
{ } {
}
}
{}{}{}
serta

\mavergleichskettedisp
{\vergleichskette
{ f
}
{ =} { 5x^2-2  \cos \left(  x^2y^3 \right) -3y^5
}
{ } {
}
{ } {
}
{ } {
}
}
{}{}{.}
Hitung fungsi atau medan vektor berikut.
\aufzaehlungsechs{
\mathl{\nabla_V f}{,}
}{
\mathl{\nabla_W f}{,}
}{
\mathl{[V,W]}{,}
}{
\mathl{\nabla_W  \nabla_V f}{,}
}{
\mathl{\nabla_V \nabla_W f}{,}
}{
\mathl{\nabla_{[V,W]} (f)}{.}
}

}
{} {}

\inputaufgabe
{2}
{

Misalkan

\mavergleichskette
{\vergleichskette
{ U
}
{ = }{ \R_+ \times \R_+
}
{  }{
}
{    }{
}
{   }{
}
}
{}{}{.}
Pada
\definitionsverweis {bundel vektor trivial}{}{}
\maabbdisp {p} {U \times \R} { U
} {}
ber-
\definitionsverweis {rank}{}{}
$1$ di atas $U$, tinjau
\definitionsverweis {koneksi linear}{}{,}
yang diberikan oleh
\definitionsverweis {simbol-simbol Christoffel}{}{}

\mavergleichskette
{\vergleichskette
{ \Gamma_1 (x,y)
}
{ = }{ \ln  \left(  xy \right)
}
{  }{
}
{    }{
}
{   }{
}
}
{}{}{}
dan

\mavergleichskette
{\vergleichskette
{ \Gamma_2 (x,y)
}
{ = }{ 5x^4 - \cos  x^2y^2 +  \ln  \left(  x+x^3 \right)
}
{  }{
}
{    }{
}
{   }{
}
}
{}{}{.}
Hitung
\definitionsverweis {operator kelengkungan}{}{}
$R \left(  \partial_1, \partial_2  \right)$.

}
{} {}

\section*{Solusi yang disediakan oleh sumber}
\addcontentsline{toc}{section}{Solusi yang disediakan oleh sumber}
\subsection*{Solusi Soal 28.2}
Menurut
Fakta,

\mavergleichskettealign
{\vergleichskettealign
{ R \left(  \partial_1, \partial_2  \right)
}
{ =} { \partial_1 \Gamma_2 - \partial_2 \Gamma_1
}
{ =} { \partial_1 \left(  7x^3 +xy^2-4y^5  \right) -  \partial_2 \left(   x^5-x^2y^2+3y^4   \right)
}
{ =} {  21x^2 +y^2 - 5x^4 -2xy^2
}
{ } {
}
}
{}
{}{.}
\subsection*{Solusi Soal 28.5}
Kita dapat bekerja secara lokal. Pada bundel vektor trivial
\maabbdisp {} {U \times \R} {U
} {}
terdapat suatu koneksi, dan kita boleh menganggap

\mavergleichskette
{\vergleichskette
{ U
}
{ \subseteq }{  \R^n
}
{  }{
}
{    }{
}
{   }{
}
}
{}{}{}
\definitionsverweis {berbentuk bintang}{}{.}
Koneksi tersebut dijelaskan oleh
\definitionsverweis {simbol-simbol Christoffel}{}{}
$\Gamma_1 , \ldots , \Gamma_n$. Karena koneksi diasumsikan bebas kelengkungan, menurut
Fakta,
\definitionsverweis {bentuk diferensial}{}{}
$1$ berikut,
$\sum_{i = 1}^n\Gamma_i dx_i$,
\definitionsverweis {tertutup}{}{.}
Menurut
Soal,
hal ini berarti bahwa medan vektor terkait
\maabbdisp {\left( \Gamma_1   , \,   \ldots  , \,  \Gamma_n                 \right)} { U } {\R^n
} {}
memenuhi
\definitionsverweis {syarat keterintegralan}{}{.}
Menurut
Fakta,
ini berarti bahwa pada himpunan terbuka berbentuk bintang $U$ terdapat suatu
\definitionsverweis {medan gradien}{}{.}
Jadi, terdapat fungsi yang terdiferensialkan secara kontinu
\maabbdisp {f} { U } { \R
} {}
dengan

\mavergleichskettedisp
{\vergleichskette
{ \Gamma_i
}
{ =} { \partial_i f
}
{ } {
}
{ } {
}
{ } {
}
}
{}{}{}
untuk semua $i$. Tetapkan

\mavergleichskettedisp
{\vergleichskette
{ g
}
{ =} { e^f
}
{ } {
}
{ } {
}
{ } {
}
}
{}{}{}
dan pandang $g$ sebagai suatu penampang pada bundel trivial $U \times \R$. Menurut
Fakta,

\mavergleichskettedisp
{\vergleichskette
{ \nabla_{\partial_i} g
}
{ =} { g \Gamma_i -  \partial_i g
}
{ =} { e^f \Gamma_i - \partial_i e^f
}
{ =} { e^f \Gamma_i - e^f \partial_i f
}
{ =} { 0
}
}
{}{}{,}
yakni $g$ merupakan
\definitionsverweis {penampang horizontal}{}{}
tanpa nol. Dengan demikian, penampang ini membuktikan bahwa $\nabla$ terintegralkan secara lokal.

\part{Unit 29}
\chapter[Kuliah 29]{Kuliah 29: Kelengkungan Riemann dan Kelengkungan Seksional}
\setcounter{section}{29}

  \zwischenueberschrift{Kelengkungan pada manifold Riemann}

Misalkan $M$ suatu
\definitionsverweis {manifold Riemann}{}{}
dan
\definitionsverweis {bundel tangen}{}{}
$TM$ dilengkapi dengan
\definitionsverweis {koneksi Levi--Civita}{}{.}
Kita hendak memahami
\definitionsverweis {operator kelengkungan}{}{}
$R$ dalam situasi ini secara lebih mendalam. Untuk medan-medan vektor $V,W$ yang diberikan, operator ini merupakan suatu pemetaan
\maabbdisp {R(V,W)} { C^2(M,TM) } { C^0(M,TM)
} {,}
yakni, ketika diterapkan pada suatu medan vektor $Z$
\zusatzklammer {yang dua kali terdiferensialkan} {} {,}
operator ini menghasilkan medan vektor

\mavergleichskettedisp
{\vergleichskette
{ R(V,W)(Z)
}
{ =} { \nabla_V \nabla_W Z
}
{ } {
}
{ } {
}
{ } {
}
}
{}{}{.}
Dengan sudut pandang ini, pada suatu manifold yang terdiferensialkan tak hingga, keseluruhan konstruksi merupakan pemetaan
\maabbdisp {} {C^\infty(TM) \times C^\infty(TM) \times C^\infty(TM)  } { C^\infty(TM)
} {,}
tetapi ketiga argumennya tidak mempunyai kedudukan yang setara.



\inputfaktbeweis
{R^n/Offene Menge/Riemannsche Metrik/Levi-Civita-Zusammenhang/Krümmungsoperator/Berechnung/Fakt}
{Lemma}
{}
{

\faktsituation {Misalkan

\mavergleichskette
{\vergleichskette
{ U
}
{ \subseteq }{ \R^n
}
{  }{
}
{    }{
}
{   }{
}
}
{}{}{}
\definitionsverweis {terbuka}{}{}
dan dilengkapi dengan suatu
\definitionsverweis {metrik Riemann}{}{}
$g_{ij}$. Misalkan $\nabla$ adalah
\definitionsverweis {koneksi Levi--Civita}{}{}
pada
\definitionsverweis {bundel tangen}{}{}

\mavergleichskette
{\vergleichskette
{ E
}
{ = }{ U \times \R^n
}
{  }{
}
{    }{
}
{   }{
}
}
{}{}{.}}
\faktfolgerung {Maka untuk
\definitionsverweis {operator kelengkungan}{}{}
berlaku

\mavergleichskettedisp
{\vergleichskette
{ R( \partial_i , \partial_j)( \partial_k)
}
{ =} { \sum_{\ell = 1}^n  R^\ell_{ijk} \partial_\ell
}
{ } {
}
{ } {
}
{ } {
}
}
{}{}{}
dengan

\mavergleichskettedisp
{\vergleichskette
{ R^\ell_{ijk}
}
{ =} { {\partial_i } \Gamma^\ell_{jk} - {\partial_j } \Gamma^\ell_{ik}  +  \sum_{ a  = 1}^n  \Gamma^a_{jk} \Gamma_{i a}^\ell -\sum_{ a  = 1}^n   \Gamma^a_{ik}  \Gamma_{j a}^\ell
}
{ } {
}
{ } {
}
{ } {
}
}
{}{}{,}
di mana

\mavergleichskettedisp
{\vergleichskette
{ \Gamma_{ij}^\ell
}
{ =} { { \frac{   1  }{  2 } } \sum_{\mu = 1}^n g^{\ell  \mu }(\partial_ig_{j\mu }+\partial_jg_{i\mu }-\partial_\mu g_{ij})
}
{ } {
}
{ } {
}
{ } {
}
}
{}{}{}
menyatakan
\definitionsverweis {simbol-simbol Christoffel}{}{}
dari koneksi tersebut.}
\faktzusatz {}
\faktzusatz {}

}
{

Ini merupakan kasus khusus dari
Lemma 28.2.

}



\inputfaktbeweis
{Riemannsche Mannigfaltigkeit/Levi-Civita-Zusammenhang/Krümmungsoperator/Eigenschaften/Fakt}
{Lemma}
{}
{

\faktsituation {Misalkan $M$ suatu
\definitionsverweis {manifold Riemann}{}{}
dan
\definitionsverweis {bundel tangen}{}{}
dilengkapi dengan
\definitionsverweis {koneksi Levi--Civita}{}{.}}
\faktuebergang {Maka
\definitionsverweis {operator kelengkungan}{}{}
mempunyai sifat-sifat berikut.}
\faktfolgerung {\aufzaehlungsechs{Ekspresi
\mathl{R(V,W)Z}{} linear dalam ketiga komponennya.
}{Untuk

\mavergleichskette
{\vergleichskette
{ P
}
{ \in }{ M
}
{  }{
}
{    }{
}
{   }{
}
}
{}{}{}
nilai
\mathl{(R(V,W)Z) (P)}{} hanya bergantung pada
\mathl{V(P),W(P),Z(P)}{}.
}{Berlaku

\mavergleichskettedisp
{\vergleichskette
{ R(V,W)
}
{ =} { -R(W,V)
}
{ } {
}
{ } {
}
{ } {
}
}
{}{}{}
}{Berlaku

\mavergleichskettedisp
{\vergleichskette
{ R(V,W)Z+ R(W,Z)V+R(Z,V)W
}
{ =} { 0
}
{ } {
}
{ } {
}
{ } {
}
}
{}{}{.}
}{Berlaku

\mavergleichskettedisp
{\vergleichskette
{ \left\langle R(V,W) T  ,  Z   \right\rangle
}
{ =} { - \left\langle R(V,W) Z  ,  T   \right\rangle
}
{ } {
}
{ } {
}
{ } {
}
}
{}{}{.}
}{Berlaku

\mavergleichskettedisp
{\vergleichskette
{ \left\langle R(V,W) T  ,  Z   \right\rangle
}
{ =} { \left\langle R(T,Z) V  ,  W   \right\rangle
}
{ } {
}
{ } {
}
{ } {
}
}
{}{}{.}
}}
\faktzusatz {}
\faktzusatz {}

}
{

\aufzaehlungsechs{Linearitas dalam $Z$ berasal dari linearitas koneksi. Linearitas dalam $V$ dan $W$ mengikuti dari
Lemma 24.10
dan
Lemma 11.8.
}{Untuk ketergantungan terhadap $V$ dan $W$, pernyataan tersebut mengikuti dari
Lemma 28.4.
Untuk menunjukkan bahwa ketergantungan terhadap $Z$ juga hanya bergantung pada $Z(P)$, kita dapat bekerja dalam situasi lokal pada

\mavergleichskette
{\vergleichskette
{ U
}
{ \subseteq }{ \R^n
}
{  }{
}
{    }{
}
{   }{
}
}
{}{}{}
dan mengambil

\mavergleichskette
{\vergleichskette
{ V
}
{ = }{ \partial_i
}
{  }{
}
{    }{
}
{   }{
}
}
{}{}{,}

\mavergleichskette
{\vergleichskette
{ W
}
{ = }{ \partial_j
}
{  }{
}
{    }{
}
{   }{
}
}
{}{}{}
dan

\mavergleichskette
{\vergleichskette
{ Z
}
{ = }{ h \partial_k
}
{  }{
}
{    }{
}
{   }{
}
}
{}{}{}
dengan suatu fungsi
\maabb {h} { U } {\R
} {}
yang dua kali terdiferensialkan secara kontinu. Menurut
Lemma 25.10,
Teorema 25.5,
dan
Teorema Schwarz,
kita memperoleh

\mavergleichskettealigndrucklinks
{\vergleichskettealigndrucklinks
{ R \left(  \partial_i, \partial_j  \right) \left(  h \partial_k  \right)
}
{ =} { \nabla_{\partial_i }  \left(  \nabla_{\partial_j }  \left(  h \partial_k  \right)  \right) - \nabla_{\partial_j } \left(  \nabla_{\partial_i }  \left(  h \partial_k  \right)  \right)
}
{ =} { \nabla_{\partial_i }  \left(  h \nabla_{\partial_j } \partial_k+ \left(  \partial_j h  \right)  \partial_k  \right) - \nabla_{\partial_j } \left( h \nabla_{\partial_i }  \partial_k  + \left(  \partial_i h  \right) \partial_k  \right)
}
{ =} { h \nabla_{\partial_i } \left(  \nabla_{\partial_j }  \partial_k   \right) +  \left(  \partial_i h  \right) \nabla_{\partial_j } \partial_k + \left(  \partial_j h  \right) \nabla_{\partial_i } \partial_k  +  \left(  \partial_i \partial_j h  \right) \partial_k
-h \nabla_{\partial_j } \left(  \nabla_{\partial_i } \partial_k   \right) -  \left(  \partial_j h  \right) \nabla_{\partial_i } \partial_k - \left(  \partial_i h  \right) \nabla_{\partial_j } \partial_k  -\left(  \partial_j  \partial_i h  \right) \partial_k
}
{ =} { h \nabla_{\partial_i } \left(  \nabla_{\partial_j }  \partial_k   \right)-h \nabla_{\partial_j } \left(  \nabla_{\partial_i } \partial_k   \right)
}
}
{}{}{,}
yang menunjukkan bahwa ekspresi ini hanya bergantung pada $h(P)$.
}{Hal ini langsung terlihat dari definisi dan
Lemma 11.8.
}{Karena
\definitionsverweis {koneksi bebas torsi}{}{}
\zusatzklammer {lihat
Teorema 26.14} {} {}
dan
identitas Jacobi,
berlaku

\mavergleichskettealigndrucklinks
{\vergleichskettealigndrucklinks
{ R(V,W)Z+ R(W,Z)V+R(Z,V)W
}
{ =} { \nabla_V  \left(  \nabla_W Z  \right) -\nabla_W  \left(  \nabla_V Z  \right) - \nabla_{[V,W]} Z +\nabla_W  \left(  \nabla_Z V  \right) -\nabla_Z  \left(  \nabla_W V  \right) - \nabla_{[W,Z]} V+\nabla_Z  \left(  \nabla_V W  \right) -\nabla_V  \left(  \nabla_Z W  \right) - \nabla_{[Z,V]} W
}
{ =} { \nabla_V  \left(  [W,Z ]  \right) + \nabla_W  \left(  [Z, V]  \right) + \nabla_Z  \left(  [V, W]  \right) - \nabla_{[V,W]} Z - \nabla_{[W,Z]} V  - \nabla_{[Z,V]} W
}
{ =} {\nabla_V  \left(  [W,Z ]  \right)  - \nabla_{[W,Z]} V + \nabla_W  \left(  [Z, V]  \right) - \nabla_{[Z,V]} W + \nabla_Z  \left(  [V, W]  \right) - \nabla_{[V,W]} Z
}
{ =} { [V, [W,Z]] + [W, [Z,V]] + [Z, [V,W]]
}
}
{
\vergleichskettefortsetzungalign
{ =} { 0
}
{ } {}
{ } {}
{ } {}
}{}{.}
}{Kita dapat membatasi diri pada medan-medan vektor $V,W$ dengan

\mavergleichskette
{\vergleichskette
{ [V,W]
}
{ = }{ 0
}
{  }{
}
{    }{
}
{   }{
}
}
{}{}{.}
Menurut
Teorema 26.14 (2)
dan
Teorema Schwarz,
kita kemudian memperoleh

\mavergleichskettealigndrucklinks
{\vergleichskettealigndrucklinks
{ \left\langle R(V,W)T , Z  \right\rangle
}
{ =} { \left\langle  \nabla_V \nabla_WT -\nabla_W \nabla_VT   ,  Z  \right\rangle
}
{ =} { \left\langle  \nabla_V \nabla_WT , Z  \right\rangle - \left\langle  \nabla_W \nabla_VT  ,  Z  \right\rangle
}
{ =} { - \left\langle  \nabla_WT , \nabla_V Z  \right\rangle + D_V\left\langle  \nabla_W T  ,  Z  \right\rangle + \left\langle  \nabla_VT  , \nabla_W  Z  \right\rangle -  D_W \left\langle  \nabla_V T  ,  Z  \right\rangle
}
{ =} { \left\langle  T ,  \nabla_W  \nabla_V Z  \right\rangle - D_W \left\langle  T ,  \nabla_VZ   \right\rangle  + D_V \left(  - \left\langle  T  ,  \nabla_WZ  \right\rangle +  D_W \left\langle T  ,  Z   \right\rangle  \right)  - \left\langle  T  ,  \nabla_V \nabla_W  Z  \right\rangle +  D_V \left\langle  T  ,  \nabla_W Z  \right\rangle - D_W  \left(  - \left\langle T  ,  \nabla_V Z   \right\rangle + D_V \left\langle T  ,  Z   \right\rangle  \right)
}
}
{
\vergleichskettefortsetzungalign
{ =} { \left\langle  T ,  \nabla_W  \nabla_V  Z  \right\rangle - \left\langle  T  ,  \nabla_V \nabla_W  Z  \right\rangle
}
{ =} { - \left\langle R(V,W)Z , T  \right\rangle
}
{ } {}
{ } {}
}{}{.}
}{Dengan menggunakan (3), (4), dan (5), diperoleh

\mavergleichskettealigndrucklinks
{\vergleichskettealigndrucklinks
{ \left\langle R(V,W)T , Z  \right\rangle
}
{ =} { - \left\langle R(T,V)W , Z  \right\rangle - \left\langle R(W,T)V , Z  \right\rangle
}
{ =} { \left\langle R(T,V) Z , W  \right\rangle + \left\langle R(W,T) Z ,  V  \right\rangle
}
{ =} { - \left\langle R(V,Z) T , W  \right\rangle - \left\langle R(Z,T) V  ,  W  \right\rangle - \left\langle R(T,Z) W ,  V  \right\rangle - \left\langle R(Z,W) T ,  V  \right\rangle
}
{ =} { 2 \left\langle R(T,Z) V  ,  W  \right\rangle - \left\langle R(V,Z) T , W  \right\rangle - \left\langle R(Z,W) T ,  V  \right\rangle
}
}
{
\vergleichskettefortsetzungalign
{ =} { 2 \left\langle R(T,Z) V  ,  W  \right\rangle + \left\langle R(V,Z) W , T  \right\rangle + \left\langle R(Z,W) V  ,  T  \right\rangle
}
{ =} { 2 \left\langle R(T,Z) V  ,  W  \right\rangle - \left\langle R(W,V) Z , T  \right\rangle
}
{ =} { 2 \left\langle R(T,Z) V  ,  W  \right\rangle - \left\langle R(V,W) T , Z  \right\rangle
}
{ } {}
}{}{.}
Ini membuktikan pernyataan tersebut.
}

}

Berdasarkan
Lemma 29.2 (2),
untuk vektor-vektor

\mavergleichskette
{\vergleichskette
{ v,w,z
}
{ \in }{ T_PM
}
{  }{
}
{    }{
}
{   }{
}
}
{}{}{}
di suatu titik

\mavergleichskette
{\vergleichskette
{ P
}
{ \in }{ M
}
{  }{
}
{    }{
}
{   }{
}
}
{}{}{}
terdefinisi dengan baik suatu vektor

\mavergleichskette
{\vergleichskette
{ R(v,w)(z)
}
{ \in }{ T_PM
}
{  }{
}
{    }{
}
{   }{
}
}
{}{}{,}
sebab vektor-vektor tersebut dapat direalisasikan oleh medan-medan vektor terdiferensialkan
\mathl{V,W,Z}{} dengan

\mavergleichskette
{\vergleichskette
{ V(P)
}
{ = }{ v
}
{  }{
}
{    }{
}
{   }{
}
}
{}{}{}
dan seterusnya
\zusatzklammer {secara lokal} {} {}
dan kemudian $R(V,W)(Z)$ dapat dievaluasi di titik tersebut.

  \zwischenueberschrift{Kelengkungan seksional}

\inputdefinition
{}
{

Misalkan $M$ suatu
\definitionsverweis {manifold Riemann}{}{,}
yang dilengkapi dengan
\definitionsverweis {koneksi Levi--Civita}{}{}
pada
\definitionsverweis {bundel tangen}{}{}
$TM$ dan
\definitionsverweis {operator kelengkungan}{}{}
terkait $R$. Untuk vektor-vektor
\definitionsverweis {bebas linear}{}{}

\mavergleichskette
{\vergleichskette
{ v,w
}
{ \in }{ T_PM
}
{  }{
}
{    }{
}
{   }{
}
}
{}{}{}
di atas suatu titik

\mavergleichskette
{\vergleichskette
{ P
}
{ \in }{ M
}
{  }{
}
{    }{
}
{   }{
}
}
{}{}{}
bilangan

\mavergleichskettedisp
{\vergleichskette
{ K(v,w)
}
{ =} { { \frac{   \left\langle R(v,w)w , v  \right\rangle  }{  \left\langle  v  , v  \right\rangle  \left\langle  w  , w  \right\rangle  -\left\langle  v  , w  \right\rangle^2   } }
}
{ } {
}
{ } {
}
{ } {
}
}
{}{}{}
disebut
\definitionswort {kelengkungan seksional}{}
untuk $v,w$ di $P$.

}

Kelengkungan seksional ini terdefinisi dengan baik karena, dalam kasus kebebasan linear, penyebutnya tidak sama dengan $0$.



\inputfaktbeweis
{Riemannsche Mannigfaltigkeit/Levi-Civita-Zusammenhang/Schnittkrümmung/Tangentiale Ebene/Fakt}
{Lemma}
{}
{

\faktsituation {Misalkan $M$ suatu
\definitionsverweis {manifold Riemann}{}{,}
yang dilengkapi dengan
\definitionsverweis {koneksi Levi--Civita}{}{}
pada
\definitionsverweis {bundel tangen}{}{}
$TM$ dan
\definitionsverweis {operator kelengkungan}{}{}
terkait $R$.}
\faktfolgerung {Maka
\definitionsverweis {kelengkungan seksional}{}{}
untuk dua vektor bebas linear

\mavergleichskette
{\vergleichskette
{ v,w
}
{ \in }{ T_PM
}
{  }{
}
{    }{
}
{   }{
}
}
{}{}{}
hanya bergantung pada bidang yang direntang oleh vektor-vektor tersebut,

\mavergleichskettedisp
{\vergleichskette
{ \R v + \R w
}
{ \subseteq} { T_PM
}
{ } {
}
{ } {
}
{ } {
}
}
{}{}{.}}
\faktzusatz {}
\faktzusatz {}

}
{

Kita meninjau ekspresi pendefinisi

\mavergleichskettedisp
{\vergleichskette
{ K(v,w)
}
{ =} { { \frac{   \left\langle R(v,w)w , v  \right\rangle  }{  \left\langle  v  , v  \right\rangle \left\langle  w  , w  \right\rangle  - \left\langle  v  , w  \right\rangle^2   } }
}
{ } {
}
{ } {
}
{ } {
}
}
{}{}{}
untuk kelengkungan seksional. Menurut
Lemma 29.2,
ekspresi
\mathl{R(v,w)z}{} linear dalam semua komponennya. Jika $v$ atau $w$ dikalikan dengan suatu skalar

\mavergleichskette
{\vergleichskette
{ a
}
{ \neq }{ 0
}
{  }{
}
{    }{
}
{   }{
}
}
{}{}{,}
faktor $a^2$ dapat dikeluarkan baik dari pembilang maupun penyebut. Selanjutnya,

\mavergleichskettealignhandlinks
{\vergleichskettealignhandlinks
{ \left\langle  R(v+w,w) w  , v +w   \right\rangle
}
{ =} { \left\langle  R(v,w)w +R(w,w)w  , v+w   \right\rangle
}
{ =} { \left\langle  R(v,w)w   , v   \right\rangle  +  \left\langle  R(w,w)w  ,  v   \right\rangle + \left\langle  R(v,w)w  ,  w   \right\rangle  +  \left\langle  R(w,w)w   , w   \right\rangle
}
{ } {
}
{ } {
}
}
{}
{}{.}
Menurut
Lemma 29.2 (4),
\mathkor {} {\left\langle  R(w,w)w  ,  v   \right\rangle} {dan} {\left\langle  R(w,w)w   , w   \right\rangle} {}
sama dengan $0$, dan menurut
Lemma 29.2 (5,6),
juga berlaku

\mavergleichskette
{\vergleichskette
{ \left\langle  R(v,w)w  ,  w   \right\rangle
}
{ = }{ 0
}
{  }{
}
{    }{
}
{   }{
}
}
{}{}{.}
Jadi pembilang kelengkungan seksional tidak berubah ketika $v$ diganti dengan $v+w$. Karena

\mavergleichskettealigndrucklinks
{\vergleichskettealigndrucklinks
{ \left\langle v+w , v+w  \right\rangle \left\langle  w  , w  \right\rangle -\left\langle v+w , w  \right\rangle^2
}
{ =} { \left\langle  v  , v  \right\rangle  \left\langle  w  , w  \right\rangle+\left\langle  w  , w  \right\rangle \left\langle  w  , w  \right\rangle +2\left\langle  v  , w  \right\rangle \left\langle  w  , w  \right\rangle -\left\langle v  , w  \right\rangle^2-\left\langle w  , w  \right\rangle^2 -2 \left\langle v  , w  \right\rangle \left\langle w  , w  \right\rangle
}
{ =} { \left\langle  v  , v  \right\rangle \left\langle  w  , w  \right\rangle -\left\langle v  , w  \right\rangle^2
}
{ } {
}
{ } {
}
}
{}{}{}
penyebutnya juga tidak berubah.

}

\inputbemerkung
{}
{

Pada suatu manifold Riemann berdimensi dua, satu-satunya subruang vektor berdimensi dua dalam setiap ruang tangen adalah ruang tangen itu sendiri. Ini berarti bahwa, dalam kasus berdimensi dua, menurut
Lemma 29.4,
argumen-argumen dalam
\definitionsverweis {kelengkungan seksional}{}{}
tidak diperlukan. Karena itu, dalam situasi ini kita memandang kelengkungan seksional sebagai suatu fungsi bernilai riil pada manifold tersebut.

}



\inputfaktbeweis
{Riemannsche Mannigfaltigkeit/Zweidimensional/Levi-Civita-Zusammenhang/Schnittkrümmung/Berechnung/Fakt}
{Lemma}
{}
{

\faktsituation {Untuk suatu
\definitionsverweis {manifold Riemann}{}{}
berdimensi dua}
\faktfolgerung {
\definitionsverweis {kelengkungan seksional}{}{}
pada setiap peta sama dengan
\mathl{-  { \frac{   R^2_{121} }{ g_{11}  } }}{.}}
\faktzusatz {}
\faktzusatz {}

}
{

Kita dapat langsung meninjau situasi pada suatu peta. Menurut
Lemma 29.4,
dalam kasus berdimensi dua kelengkungan seksional dapat dihitung menggunakan basis sembarang; kita memilih
\mathkor {} {\partial_1} {dan} {\partial_2} {.}
Dengan menggunakan
Lemma 29.2 (5),
untuk kelengkungan seksional diperoleh

\mavergleichskettealign
{\vergleichskettealign
{ K( \partial_1, \partial_2)
}
{ =} { { \frac{   \left\langle R(\partial_1 ,\partial_2)\partial_2  ,  \partial_1   \right\rangle  }{  \left\langle  \partial_1  ,  \partial_1   \right\rangle  \left\langle  \partial_2  ,  \partial_2   \right\rangle  -\left\langle  \partial_1  ,  \partial_2   \right\rangle^2   } }
}
{ =} { { \frac{   \left\langle R(\partial_1 ,\partial_2)\partial_2  ,  \partial_1   \right\rangle  }{  g_{11} g_{22}  -g_{12}^2   } }
}
{ =} { -  { \frac{   \left\langle R(\partial_1 ,\partial_2)\partial_1  ,  \partial_2   \right\rangle  }{  g_{11} g_{22}  -g_{12}^2   } }
}
{ } {
}
}
{}
{}{}
Menurut
Lemma 29.1,
berlaku

\mavergleichskettedisp
{\vergleichskette
{ R( \partial_1 , \partial_2)( \partial_1)
}
{ =} { R^1_{121} \partial_1+ R^2_{121} \partial_2
}
{ } {
}
{ } {
}
{ } {
}
}
{}{}{.}
Hal ini memberikan

\mavergleichskettealign
{\vergleichskettealign
{ K( \partial_1, \partial_2)
}
{ =} { - { \frac{   \left\langle R(\partial_1 ,\partial_2)\partial_1  ,  \partial_2   \right\rangle  }{  g_{11} g_{22}  -g_{12}^2   } }
}
{ =} { - { \frac{   R^1_{121} g_{12}+ R^2_{121} g_{22}  }{  g_{11} g_{22}  -g_{12}^2   } }
}
{ } {
}
{ } {
}
}
{}
{}{.}
Menurut
Lemma 29.2 (4),

\mavergleichskettedisp
{\vergleichskette
{ \left\langle  R(\partial_1, \partial_1) \partial_1  ,  \partial_2   \right\rangle
}
{ =} { R^1_{121}  g_{11}  + R_{121}^2 g_{12}
}
{ =} { 0
}
{ } {
}
{ } {
}
}
{}{}{.}
Kita mengalikan pembilang kelengkungan seksional dengan $g_{11}$ dan memperoleh

\mavergleichskettealignhandlinks
{\vergleichskettealignhandlinks
{ g_{11}  \left(  R^1_{121} g_{12}+ R^2_{121} g_{22}  \right)
}
{ =} { g_{11}  \left(  R^1_{121} g_{12}+ R^2_{121} g_{22}  \right) - g_{12}  \left(  R^1_{121}  g_{11}  + R_{121}^2 g_{12}   \right)
}
{ =} { g_{11} g_{22} R^2_{121} - g_{12}^2 R^2_{121}
}
{ =} { \left(  g_{11} g_{22}  - g_{12}^2  \right) R^2_{121}
}
{ } {
}
}
{}
{}{.}
Jadi,

\mavergleichskettedisp
{\vergleichskette
{ K( \partial_1, \partial_2)
}
{ =} { -  { \frac{   R^2_{121} }{ g_{11}  } }
}
{ } {
}
{ } {
}
{ } {
}
}
{}{}{.}

}



\inputfaktbeweis
{Fläche im Raum/Schnittkrümmung/Gaußkrümmung/Fakt}
{Lemma}
{}
{

\faktsituation {Misalkan

\mavergleichskette
{\vergleichskette
{ W
}
{ \subseteq }{ \R^3
}
{  }{
}
{    }{
}
{   }{
}
}
{}{}{}
\definitionsverweis {terbuka}{}{}
dan misalkan
\maabb {h} { W } { \R
} {}
suatu
\definitionsverweis {fungsi yang dua kali terdiferensialkan secara kontinu}{}{}
serta

\mavergleichskette
{\vergleichskette
{ Y
}
{ = }{ h^{-1}(c)
}
{  }{
}
{    }{
}
{   }{
}
}
{}{}{}
serat di atas

\mavergleichskette
{\vergleichskette
{ c
}
{ \in }{ \R
}
{  }{
}
{    }{
}
{   }{
}
}
{}{}{,}
dan andaikan $h$
\definitionsverweis {reguler}{}{}
di setiap titik $Y$.}
\faktfolgerung {Maka
\definitionsverweis {kelengkungan Gauss}{}{}
dari $Y$ sama dengan
\definitionsverweis {kelengkungan seksional}{}{}
dari $Y$.}
\faktzusatz {}
\faktzusatz {}

}
{

Menurut
Lemma 29.6,
kelengkungan seksional pada suatu peta sama dengan
\mathl{- { \frac{  R^2_{121}  }{ g_{11}  } }}{.} Karena

\mavergleichskettedisphandlinks
{\vergleichskettedisphandlinks
{ R^2_{121}
}
{ =} { {\partial_1 } \Gamma^2_{12} - {\partial_2 } \Gamma^2_{11} +  \Gamma^1_{12} \Gamma_{1 1}^2 + \Gamma^2_{12} \Gamma_{1 2}^2 - \Gamma^1_{11}  \Gamma_{12 }^2  - \Gamma^2_{11} \Gamma_{2 2}^2
}
{ =} { -g_{11} K
}
{ } {
}
{ } {
}
}
{}{}{.}
maka

\mavergleichskettedisphandlinks
{\vergleichskettedisphandlinks
{ K
}
{ =} { { \frac{   1  }{ g_{11}  } }  \left(  \partial_2 \Gamma_{11}^2 - \partial_1 \Gamma_{12}^2 + \Gamma^1_{11} \Gamma^2_{12}  + \Gamma_{11}^2 \Gamma_{22}^2 - \Gamma_{12}^1 \Gamma_{11}^2-  \left(  \Gamma_{12}^2  \right)^2  \right)
}
{ } {
}
{ } {
}
{ } {
}
}
{}{}{}
yang berimpit dengan kelengkungan Gauss menurut
Teorema 19.7.

}

\inputbeispiel{}
{

Untuk
\definitionsverweis {bidang Euklides}{}{}
$\R^2$ dengan
\definitionsverweis {hasil kali skalar standar}{}{,}
\definitionsverweis {kelengkungan seksional}{}{}
konstan sama dengan $0$.

}

\inputbeispiel{}
{

Untuk
\definitionsverweis {bola satuan}{}{}

\mavergleichskette
{\vergleichskette
{ S^2
}
{ \subseteq }{ \R^3
}
{  }{
}
{    }{
}
{   }{
}
}
{}{}{}
dengan
\definitionsverweis {struktur Riemann}{}{}
terinduksi,
\definitionsverweis {kelengkungan seksional}{}{}
konstan sama dengan $1$. Hal ini mengikuti dari
Contoh 5.6
bersama dengan
Lemma 29.7.

}

\inputbeispiel{}
{

\bild{ \begin{center}
\includegraphics[width=5.5cm]{ \bildeinlesungsvg {Parallel lines in Poincare's model of hyperbolic geometry} {svg} }
\end{center}
\bildtext {} }

\bildlizenz { Parallel lines in Poincare's model of hyperbolic geometry.svg } {} {Januszkaja} {Commons} {CC-by-sa 3.0} {}

Kita melanjutkan
Contoh 18.8
dan
Contoh 26.7.
Menurut
Lemma 29.1,

\mavergleichskettedisphandlinks
{\vergleichskettedisphandlinks
{ R^2_{121}
}
{ =} { {\partial_1 } \Gamma^2_{12} - {\partial_2 } \Gamma^2_{11} + \Gamma^1_{12} \Gamma_{1 1}^2 + \Gamma^2_{12} \Gamma_{1 2}^2 - \Gamma^1_{11} \Gamma_{12 }^2 - \Gamma^2_{11} \Gamma_{2 2}^2
}
{ =} { { \frac{   1  }{  y^2 } } - { \frac{   1  }{  y^2 } }+  { \frac{   1  }{  y^2 } }
}
{ =} { { \frac{   1  }{  y^2 } }
}
{ } {
}
}
{}{}{.}
Menurut
Lemma 29.6,
untuk
\definitionsverweis {kelengkungan seksional}{}{}
berlaku

\mavergleichskettedisp
{\vergleichskette
{ K
}
{ =} { - { \frac{   R^2_{121} }{ g_{11 } }}
}
{ =} { - { \frac{   { \frac{   1  }{  y^2 } }  }{  { \frac{   1  }{  y^2 } }  } }
}
{ =} { -1
}
{ } {
}
}
{}{}{,}
jadi kita memperoleh kelengkungan seksional konstan $-1$.

}

\chapter{Lembar Kerja 29}
\setcounter{section}{29}

  \zwischenueberschrift{Soal latihan}

\inputaufgabe
{}
{

Misalkan $M$ suatu
\definitionsverweis {manifold diferensiabel}{}{}
sedemikian sehingga
\definitionsverweis {bundel tangen}{}{}
$TM$ bersifat trivial. Melalui suatu trivialisasi diferensiabel

\mavergleichskette
{\vergleichskette
{ TM
}
{ \cong }{ M \times \R^n
}
{  }{
}
{    }{
}
{   }{
}
}
{}{}{}
lengkapi $M$
\zusatzklammer {dengan menggunakan
\definitionsverweis {hasil kali dalam standar}{}{}
pada $\R^n$} {} {}
dengan suatu
\definitionsverweis {struktur Riemann}{}{}
\zusatzklammer {lihat
Soal 16.2} {} {.}
Tunjukkan bahwa
\definitionsverweis {operator kelengkungan}{}{}
pada $M$ bersifat trivial.

}
{} {}

\inputaufgabegibtloesung
{}
{

Dalam situasi dua dimensi pada
Lemma 29.1,
hitung simbol-simbol Christoffel yang relevan untuk
\definitionsverweis {operator kelengkungan}{}{.}

}
{} {}

\inputaufgabe
{}
{

Kita tinjau
\definitionsverweis {torus}{}{}
$S^1 \times S^1$. Pertama, lengkapi torus tersebut dengan struktur yang diberikan oleh struktur produk
\zusatzklammer {bersama penyematan

\mavergleichskette
{\vergleichskette
{ S^1
}
{ \subseteq }{ \R^2
}
{  }{
}
{    }{
}
{   }{
}
}
{}{}{,}
lihat
Soal 16.3
dan
Soal 16.4} {} {}
dan
\definitionsverweis {struktur Riemann}{}{,}
kemudian bandingkan dengan realisasi tertanam berikut.

\mavergleichskettedisp
{\vergleichskette
{ T
}
{ =} { \left\{ (x,y,z) \in \R^3  \mid   \left(  \sqrt{x^2+y^2} -R  \right)^2 +z^2  = r^2         \right\}
}
{ \subseteq} { \R^3
}
{ } {
}
{ } {
}
}
{}{}{}
dengan struktur Riemann terinduksi. Bandingkan
\definitionsverweis {operator-operator kelengkungan}{}{}
untuk kedua struktur tersebut.

}
{} {}

  \zwischenueberschrift{Soal untuk dikumpulkan}

\section*{Solusi yang disediakan oleh sumber}
\addcontentsline{toc}{section}{Solusi yang disediakan oleh sumber}
\subsection*{Solusi Soal 29.2}
Berdasarkan berbagai identitas pada
Fakta,
hanya simbol-simbol Christoffel berikut yang mungkin relevan:

\mathdisp {R^1_{112}, R^2_{112},R^1_{122},R^2_{122},} {  }

Menurut
Fakta,

\mavergleichskettedisp
{\vergleichskette
{ R^1_{112}
}
{ =} { {\partial_1} \Gamma^1_{12} - {\partial_1} \Gamma^1_{12}  +  \Gamma^1_{12} \Gamma_{11}^1  +  \Gamma^2_{12} \Gamma_{12}^1        -     \Gamma^1_{12}  \Gamma_{11}^1        -     \Gamma^2_{12}  \Gamma_{12}^1
}
{ =} {  0
}
{ } {
}
{ } {
}
}
{}{}{,}

\mavergleichskettedisp
{\vergleichskette
{ R^2_{112}
}
{ =} { {\partial_1} \Gamma^2_{12} - {\partial_1} \Gamma^2_{12} +  \Gamma^1_{12} \Gamma_{11}^2 +\Gamma^2_{12} \Gamma_{12}^2 - \Gamma^1_{12}  \Gamma_{11}^2 - \Gamma^2_{12}  \Gamma_{12}^2
}
{ =} {  0
}
{ } {
}
{ } {
}
}
{}{}{,}

\mavergleichskettedisp
{\vergleichskette
{ R^1_{121}
}
{ =} { - {\partial_1} \Gamma^1_{12} - {\partial_2} \Gamma^1_{11} -  \Gamma^1_{12} \Gamma_{1 1}^1 - \Gamma^2_{12} \Gamma_{1 2}^1 +  \Gamma^1_{11}  \Gamma_{12 }^1  - \Gamma^2_{11} \Gamma_{2 2}^1
}
{ } {
}
{ } {
}
{ } {
}
}
{}{}{,}

\mavergleichskettedisp
{\vergleichskette
{ R^2_{121}
}
{ =} {  {\partial_1} \Gamma^2_{12} - {\partial_2} \Gamma^2_{11} +  \Gamma^1_{12} \Gamma_{1 1}^2 + \Gamma^2_{12} \Gamma_{1 2}^2 - \Gamma^1_{11}  \Gamma_{12 }^2  - \Gamma^2_{11} \Gamma_{2 2}^2
}
{ } {
}
{ } {
}
{ } {
}
}
{}{}{.}

\mavergleichskettedisp
{\vergleichskette
{ R^1_{122}
}
{ =} { {\partial_1} \Gamma^1_{22} - {\partial_2} \Gamma^1_{22}  +   \Gamma^1_{22} \Gamma_{1 1}^1   +   \Gamma^2_{22} \Gamma_{1 2}^1 -  \Gamma^1_{12}  \Gamma_{12 }^1  - \Gamma^2_{12} \Gamma_{2 2}^1
}
{ } {
}
{ } {
}
{ } {
}
}
{}{}{,}

\mavergleichskettedisp
{\vergleichskette
{ R^2_{122}
}
{ =} { {\partial_1} \Gamma^2_{22} - {\partial_2} \Gamma^2_{22}  +   \Gamma^1_{22} \Gamma_{1 1}^2  -  \Gamma^1_{12}  \Gamma_{12 }^2
}
{ } {
}
{ } {
}
{ } {
}
}
{}{}{.}

\part{Jembatan Asli}
\chapter{Jembatan Grup Lie dan Aljabar Lie}
\noindent\textit{Materi asli edisi ini; bukan solusi atau teks yang disediakan oleh sumber Brenner/Wikiversity.}
% O011-BL — modul asli, bukan terjemahan atau adaptasi dari teks pembanding.
% Judul: Jembatan Grup Lie dan Aljabar Lie
% Lisensi modul: CC BY-SA 4.0.
% Provenans: disusun oleh OpenAI Codex gpt-5.6-sol, Ultra, atas arahan pengguna.
% Prasyarat internal: Unit 7, 9, 10, 11, dan 28 edisi ini.

\setcounter{section}{30}
\section{Jembatan Grup Lie dan Aljabar Lie}
\label{o011-bridge-lie}

\subsection{Tujuan dan hubungan dengan materi sebelumnya}

Modul ini menghubungkan bahasa manifold dengan simetri kontinu. Unit 7
menyediakan manifold mulus, Unit 9 ruang singgung, Unit 10 pemetaan mulus dan
medan vektor, Unit 11 produk manifold, sedangkan Unit 28 menyediakan braket
Lie medan vektor. Semua hasil di bawah dibangun dari perangkat itu. Tujuan
utamanya ialah memahami bagaimana operasi grup yang mulus menghasilkan suatu
aljabar linear pada ruang singgung di unsur identitas, bagaimana eksponensial
menghubungkan kedua struktur tersebut, dan bagaimana aksi grup menghasilkan
orbit serta stabilisator.

Sepanjang modul, semua manifold dianggap berdimensi hingga, Hausdorff, dan
memiliki basis terhitung. Kata ``mulus'' berarti terdiferensialkan tak hingga.
Jika $f\colon M\to N$ mulus dan $p\in M$, diferensialnya ditulis
$T_pf\colon T_pM\to T_{f(p)}N$.

\subsection{Grup Lie dan contoh matriks}

\begin{definition}\label{o011-bl-def-lie-group}
Suatu \emph{grup Lie} ialah sebuah manifold mulus $G$ yang sekaligus merupakan
grup, sedemikian sehingga pemetaan perkalian dan invers
\[
  m\colon G\times G\longrightarrow G,\qquad m(g,h)=gh,
  \qquad
  \iota\colon G\longrightarrow G,\qquad \iota(g)=g^{-1},
\]
bersifat mulus. Unsur identitasnya ditulis $e$. Homomorfisme grup Lie ialah
homomorfisme grup yang juga merupakan pemetaan mulus.
\end{definition}

Syarat tersebut bukan sekadar hiasan topologis. Ia memungkinkan kita
mendiferensialkan identitas grup. Sebagai contoh, dari $gg^{-1}=e$ kita akan
memperoleh rumus diferensial untuk invers; dari keasosiatifan kita akan
memperoleh medan vektor invarian dan braket pada $T_eG$.

\begin{beispiel}\label{o011-bl-ex-vector-groups}
Ruang vektor nyata $\mathbb R^n$ dengan penjumlahan merupakan grup Lie. Unsur
identitasnya $0$, inversnya $x\mapsto -x$, dan kedua operasi itu polinomial.
Torus
\[
  \mathbb T^n=(S^1)^n
\]
dengan perkalian komponen demi komponen juga merupakan grup Lie. Pemetaan
\[
  \mathbb R^n\longrightarrow\mathbb T^n,
  \qquad (t_1,\ldots,t_n)\longmapsto
  (e^{it_1},\ldots,e^{it_n})
\]
ialah homomorfisme grup Lie surjektif dengan kernel $(2\pi\mathbb Z)^n$.
\end{beispiel}

\begin{beispiel}\label{o011-bl-ex-general-linear}
Di dalam ruang semua matriks nyata $M_n(\mathbb R)\cong\mathbb R^{n^2}$,
himpunan
\[
  \mathrm{GL}_n(\mathbb R)=\{A\in M_n(\mathbb R):\det A\ne0\}
\]
terbuka karena determinan kontinu. Jadi ia merupakan manifold. Perkalian
matriks bersifat polinomial, sedangkan
\[
  A^{-1}=\frac{1}{\det A}\operatorname{adj}(A)
\]
mulus pada himpunan tempat $\det A\ne0$. Dengan demikian
$\mathrm{GL}_n(\mathbb R)$ ialah grup Lie.
\end{beispiel}

\begin{beispiel}\label{o011-bl-ex-matrix-subgroups}
Contoh-contoh penting berikut merupakan subgrup tertutup
$\mathrm{GL}_n(\mathbb R)$ dan, dengan struktur manifold terinduksi yang
sesuai, merupakan grup Lie:
\[
\begin{aligned}
 \mathrm{SL}_n(\mathbb R)&=\{A:\det A=1\},\\
 \mathrm O(n)&=\{A:A^{\mathsf T}A=I\},\\
 \mathrm{SO}(n)&=\{A:A^{\mathsf T}A=I,\ \det A=1\}.
\end{aligned}
\]
Untuk $\mathrm{SL}_n$, nilai $1$ merupakan nilai regular determinan karena
$D(\det)_A(H)=\det(A)\operatorname{tr}(A^{-1}H)$ tidak nol sebagai fungsional.
Untuk $\mathrm O(n)$, persamaan $A^{\mathsf T}A=I$ dapat diperlakukan dengan
teorema nilai regular ke ruang matriks simetris. Kita tidak memerlukan teorema
umum bahwa setiap subgrup tertutup dari grup Lie adalah submanifold, meskipun
teorema itu memang benar.
\end{beispiel}

\begin{bemerkung}\label{o011-bl-rem-components}
$\mathrm{GL}_n(\mathbb R)$ mempunyai dua komponen terhubung, dibedakan oleh
tanda determinan. $\mathrm O(n)$ juga mempunyai dua komponen, sedangkan
$\mathrm{SO}(n)$ adalah komponen identitasnya. Aljabar Lie hanya merekam
geometri infinitesimal di sekitar $e$; karena itu ia tidak dengan sendirinya
merekam semua komponen terhubung suatu grup.
\end{bemerkung}

\subsection{Translasi kiri dan medan vektor invarian}

Untuk $g\in G$, definisikan translasi kiri dan kanan
\[
  L_g(h)=gh,
  \qquad
  R_g(h)=hg.
\]
Keduanya difeomorfisme, dengan invers masing-masing $L_{g^{-1}}$ dan
$R_{g^{-1}}$.

\begin{definition}\label{o011-bl-def-left-invariant}
Suatu medan vektor mulus $X$ pada $G$ disebut \emph{invarian-kiri} apabila
\[
  T_hL_g(X_h)=X_{gh}
\]
untuk semua $g,h\in G$.
\end{definition}

\begin{Proposition}\label{o011-bl-prop-left-invariant-correspondence}
Untuk setiap $v\in T_eG$ terdapat tepat satu medan vektor invarian-kiri
$X^v$, yaitu
\[
  X^v_g=T_eL_g(v).
\]
Pemetaan $v\mapsto X^v$ merupakan isomorfisme linear dari $T_eG$ ke ruang
medan vektor invarian-kiri.
\end{Proposition}

\begin{proof}
Keasosiatifan memberi $L_g\circ L_h=L_{gh}$. Maka
\[
 T_hL_g(X^v_h)
 =T_hL_g\bigl(T_eL_h(v)\bigr)
 =T_e(L_g\circ L_h)(v)
 =T_eL_{gh}(v)=X^v_{gh}.
\]
Sebaliknya, jika $X$ invarian-kiri, ambil $h=e$ untuk memperoleh
$X_g=T_eL_g(X_e)$. Jadi $X$ ditentukan secara unik oleh nilainya di $e$.
Linearitas jelas dari linearitas diferensial.
\end{proof}

\begin{Proposition}\label{o011-bl-prop-invariant-bracket}
Jika $X$ dan $Y$ invarian-kiri, maka braket medan vektor $[X,Y]$ juga
invarian-kiri.
\end{Proposition}

\begin{proof}
Untuk setiap difeomorfisme $F$, sifat natural braket menyatakan
$F_*[X,Y]=[F_*X,F_*Y]$. Terapkan ini pada $F=L_g$. Karena
$(L_g)_*X=X$ dan $(L_g)_*Y=Y$, diperoleh $(L_g)_*[X,Y]=[X,Y]$.
\end{proof}

\subsection{Aljabar Lie di unsur identitas}

\begin{definition}\label{o011-bl-def-lie-algebra}
\emph{Aljabar Lie} suatu grup Lie $G$ ialah ruang vektor
\[
  \mathfrak g=T_eG
\]
dengan operasi bilinear
\[
  [v,w]_{\mathfrak g}=[X^v,X^w]_e.
\]
Jika konteksnya jelas, subskrip $\mathfrak g$ dihilangkan.
\end{definition}

Proposisi sebelumnya menjamin bahwa $[X^v,X^w]$ kembali invarian-kiri,
sehingga nilainya di $e$ menentukan seluruh medan. Sifat braket medan vektor
langsung memberi
\[
 [v,w]=-[w,v],
 \qquad
 [u,[v,w]]+[v,[w,u]]+[w,[u,v]]=0.
\]
Identitas kedua disebut identitas Jacobi.

\begin{Proposition}\label{o011-bl-prop-homomorphism-differential}
Jika $F\colon G\to H$ homomorfisme grup Lie, maka
\[
  dF_e=T_eF\colon\mathfrak g\longrightarrow\mathfrak h
\]
mempertahankan braket Lie.
\end{Proposition}

\begin{proof}
Identitas $F\circ L_g=L_{F(g)}\circ F$ menunjukkan bahwa medan
invarian-kiri $X^v$ dan $X^{dF_e(v)}$ saling terkait oleh $F$. Naturalitas
braket untuk medan yang saling terkait menghasilkan
\[
  dF_e([v,w])=[dF_e(v),dF_e(w)].
\]
\end{proof}

\begin{beispiel}\label{o011-bl-ex-matrix-algebras}
Untuk grup matriks, kita mengidentifikasi ruang singgung di $I$ dengan
subruang matriks. Hasilnya ialah
\[
\begin{aligned}
 \mathfrak{gl}_n(\mathbb R)&=M_n(\mathbb R),\\
 \mathfrak{sl}_n(\mathbb R)&=\{X:\operatorname{tr}X=0\},\\
 \mathfrak o(n)=\mathfrak{so}(n)&=\{X:X^{\mathsf T}+X=0\}.
\end{aligned}
\]
Braketnya adalah komutator
\[
  [X,Y]=XY-YX.
\]
Sebagai contoh, diferensiasikan $A(t)^{\mathsf T}A(t)=I$ pada $t=0$ dengan
$A(0)=I$ untuk memperoleh $A'(0)^{\mathsf T}+A'(0)=0$.
\end{beispiel}

\begin{bemerkung}\label{o011-bl-rem-sign}
Konvensi braket di atas memakai medan invarian-kiri. Jika seseorang
mengidentifikasi $T_eG$ dengan medan invarian-kanan tanpa mengubah definisi,
tanda braket akan berbalik. Menyatakan konvensi ini mencegah kesalahan tanda
pada rumus adjoint.
\end{bemerkung}

\subsection{Subgrup satu-parameter dan pemetaan eksponensial}

\begin{definition}\label{o011-bl-def-one-parameter}
Subgrup satu-parameter dalam $G$ ialah homomorfisme grup Lie
\[
  \gamma\colon(\mathbb R,+)\longrightarrow G.
\]
\end{definition}

Jika $\gamma$ subgrup satu-parameter, maka
\[
  \gamma(t+s)=\gamma(t)\gamma(s)
\]
dan, setelah mendiferensialkan terhadap $s$ pada $0$,
\[
  \gamma'(t)=T_eL_{\gamma(t)}\bigl(\gamma'(0)\bigr).
\]
Jadi $\gamma$ adalah kurva integral medan invarian-kiri yang ditentukan oleh
$v=\gamma'(0)$.

\begin{Theorem}\label{o011-bl-thm-one-parameter-existence}
Untuk setiap $v\in\mathfrak g$ terdapat tepat satu subgrup satu-parameter
$\gamma_v$ dengan $\gamma_v'(0)=v$.
\end{Theorem}

\begin{proof}
Ambil kurva integral maksimal $\gamma$ dari medan invarian-kiri $X^v$ dengan
$\gamma(0)=e$. Untuk $s$ tetap, kurva $t\mapsto\gamma(s)\gamma(t)$ dan
$t\mapsto\gamma(s+t)$ memenuhi persamaan diferensial yang sama serta memiliki
nilai awal yang sama pada $t=0$. Ketunggalan solusi memberi
$\gamma(s+t)=\gamma(s)\gamma(t)$ selama kedua ruas terdefinisi.

Identitas ini memperpanjang kurva integral lokal langkah demi langkah ke
seluruh $\mathbb R$: pilih $\varepsilon>0$ dengan solusi pada
$(-\varepsilon,\varepsilon)$, lalu definisikan nilai pada selang yang lebih
panjang dengan hasil kali potongan-potongan kecil. Identitas lokal dan
ketunggalan menjamin bahwa definisi itu konsisten. Kurva global yang diperoleh
adalah homomorfisme. Ketunggalan subgrup satu-parameter juga mengikuti dari
ketunggalan kurva integral.
\end{proof}

\begin{definition}\label{o011-bl-def-exponential}
Pemetaan eksponensial grup Lie adalah
\[
  \exp_G\colon\mathfrak g\longrightarrow G,
  \qquad \exp_G(v)=\gamma_v(1).
\]
\end{definition}

Dari definisi dan penskalaan parameter diperoleh
\[
  \gamma_v(t)=\exp_G(tv),
  \qquad
  \exp_G((s+t)v)=\exp_G(sv)\exp_G(tv).
\]

\begin{Proposition}\label{o011-bl-prop-exp-local}
Pemetaan $\exp_G$ mulus, $\exp_G(0)=e$, dan diferensialnya di $0$ adalah
identitas pada $\mathfrak g$. Karena itu $\exp_G$ merupakan difeomorfisme dari
suatu lingkungan $0$ ke suatu lingkungan $e$.
\end{Proposition}

\begin{proof}
Kebergantungan mulus solusi persamaan diferensial pada nilai awal dan parameter
memberi kemulusan $\exp_G$. Kurva $t\mapsto\exp_G(tv)$ mempunyai kecepatan $v$
pada $0$, sehingga
\[
  T_0\exp_G(v)=\left.\frac{d}{dt}\right|_{0}\exp_G(tv)=v.
\]
Teorema fungsi invers menyelesaikan pernyataan terakhir.
\end{proof}

\begin{beispiel}\label{o011-bl-ex-matrix-exp}
Untuk $G=\mathrm{GL}_n(\mathbb R)$,
\[
  \exp_G(X)=e^X=\sum_{k=0}^{\infty}\frac{X^k}{k!}.
\]
Deret ini konvergen mutlak dalam setiap norma matriks. Turunan
$t\mapsto e^{tX}$ ialah $Xe^{tX}=e^{tX}X$, sehingga kurva tersebut merupakan
subgrup satu-parameter dengan kecepatan awal $X$. Jika $X$ dan $Y$ komutatif,
$e^{X+Y}=e^Xe^Y$; tanpa hipotesis komutatif, identitas ini umumnya salah.
\end{beispiel}

\begin{Proposition}\label{o011-bl-prop-exp-natural}
Untuk homomorfisme grup Lie $F\colon G\to H$ berlaku
\[
  F(\exp_G v)=\exp_H(dF_e v).
\]
\end{Proposition}

\begin{proof}
Kedua ruas, setelah $v$ diganti $tv$, merupakan subgrup satu-parameter di $H$
dengan kecepatan awal $dF_e v$. Gunakan ketunggalan pada Teorema
\ref{o011-bl-thm-one-parameter-existence} dan ambil $t=1$.
\end{proof}

\subsection{Konjugasi dan representasi adjoint}

Untuk $g\in G$, konjugasi oleh $g$ adalah difeomorfisme
\[
  C_g\colon G\longrightarrow G,
  \qquad C_g(h)=ghg^{-1}.
\]
Ia mempertahankan $e$.

\begin{definition}\label{o011-bl-def-adjoint}
Pemetaan adjoint pada tingkat grup adalah
\[
  \operatorname{Ad}\colon G\longrightarrow\mathrm{GL}(\mathfrak g),
  \qquad \operatorname{Ad}_g=T_eC_g.
\]
Diferensialnya di identitas adalah
\[
  \operatorname{ad}=T_e\operatorname{Ad}\colon
  \mathfrak g\longrightarrow\mathfrak{gl}(\mathfrak g).
\]
\end{definition}

Karena $C_{gh}=C_g\circ C_h$, pemetaan $\operatorname{Ad}$ merupakan
homomorfisme grup Lie. Untuk grup matriks,
\[
  \operatorname{Ad}_g(X)=gXg^{-1}.
\]

\begin{Theorem}\label{o011-bl-thm-ad-bracket}
Untuk $v,w\in\mathfrak g$ berlaku
\[
  \operatorname{ad}_v(w)=[v,w].
\]
Selain itu,
\[
  \operatorname{Ad}_g[v,w]
  =[\operatorname{Ad}_gv,\operatorname{Ad}_gw]
\]
dan
\[
  \operatorname{Ad}_{\exp(tv)}
  =\exp\bigl(t\operatorname{ad}_v\bigr).
\]
\end{Theorem}

\begin{proof}
Konjugasi membawa medan invarian-kiri $X^w$ ke medan invarian-kiri yang
bernilai $\operatorname{Ad}_g w$ di identitas. Naturalitas braket memberi
identitas kedua. Untuk identitas pertama, diferensiasikan keluarga
$C_{\exp(tv)}$ pada $t=0$; turunan aksi keluarga difeomorfisme ini pada medan
invarian-kiri adalah braket dengan $X^v$. Evaluasi di $e$ memberi
$\operatorname{ad}_v(w)=[v,w]$.

Terakhir, $t\mapsto\operatorname{Ad}_{\exp(tv)}$ ialah subgrup satu-parameter
di $\mathrm{GL}(\mathfrak g)$ dengan kecepatan awal
$\operatorname{ad}_v$. Contoh \ref{o011-bl-ex-matrix-exp} pada ruang
$\mathfrak g$ memberi identitas ketiga.
\end{proof}

\subsection{Aksi mulus, orbit, dan stabilisator}

\begin{definition}\label{o011-bl-def-action}
Aksi kiri mulus grup Lie $G$ pada manifold $M$ adalah pemetaan mulus
\[
  \Phi\colon G\times M\longrightarrow M,
  \qquad (g,p)\longmapsto g\cdot p,
\]
dengan $e\cdot p=p$ dan $(gh)\cdot p=g\cdot(h\cdot p)$.
Orbit dan stabilisator titik $p$ adalah
\[
  G\cdot p=\{g\cdot p:g\in G\},
  \qquad
  G_p=\{g\in G:g\cdot p=p\}.
\]
\end{definition}

Stabilisator $G_p$ adalah subgrup tertutup, sebab ia adalah prapeta
$\{p\}$ di bawah pemetaan kontinu $g\mapsto g\cdot p$. Pemetaan orbit
\[
  \Phi_p\colon G\longrightarrow M,
  \qquad \Phi_p(g)=g\cdot p
\]
memenuhi $\Phi_p(gh)=g\cdot\Phi_p(h)$.

\begin{definition}\label{o011-bl-def-fundamental-field}
Untuk $v\in\mathfrak g$, medan vektor fundamental yang dihasilkan oleh $v$
adalah
\[
  v_M(p)=\left.\frac{d}{dt}\right|_{0}\exp(tv)\cdot p
  =T_e\Phi_p(v).
\]
\end{definition}

Dengan konvensi aksi kiri dan medan invarian-kiri yang dipakai di sini,
pemetaan $v\mapsto v_M$ merupakan antihomomorfisme aljabar Lie:
\[
  [v_M,w_M]=-[v,w]_M.
\]
Tanda minus berasal dari fakta bahwa $p\mapsto\exp(tv)\cdot p$ membentuk
aksi kiri, sedangkan medan fundamentalnya terkait secara alami dengan medan
invarian-kanan pada $G$. Mengubah salah satu konvensi akan mengubah tanda.

\begin{Proposition}\label{o011-bl-prop-orbit-tangent}
Untuk setiap $p\in M$,
\[
  \operatorname{im}(T_e\Phi_p)
  =\{v_M(p):v\in\mathfrak g\}
\]
adalah ruang arah tangensial orbit di $p$. Kernel
\[
  \mathfrak g_p=\ker(T_e\Phi_p)
\]
adalah aljabar Lie stabilisator. Jika orbit mempunyai struktur manifold
homogen yang lazim, maka
\[
  T_p(G\cdot p)\cong\mathfrak g/\mathfrak g_p,
  \qquad
  \dim(G\cdot p)=\dim G-\dim G_p.
\]
\end{Proposition}

\begin{proof}
Identitas pertama adalah definisi diferensial pemetaan orbit. Kurva
$\exp(tv)$ terletak di $G_p$ tepat secara infinitesimal ketika
$v_M(p)=0$, sehingga kernelnya adalah ruang singgung stabilisator di $e$.
Teorema isomorfisme linear memberi
$\mathfrak g/\ker(T_e\Phi_p)\cong\operatorname{im}(T_e\Phi_p)$.
Apabila orbit direalisasikan sebagai manifold homogen $G/G_p$, diferensial
pemetaan hasil bagi mengidentifikasi ruang terakhir dengan $T_p(G\cdot p)$.
Rumus dimensi mengikuti.
\end{proof}

\begin{beispiel}\label{o011-bl-ex-sphere-action}
$\mathrm{SO}(n)$ bertindak pada $S^{n-1}$ melalui perkalian matriks. Aksi ini
transitif. Stabilisator $e_n$ terdiri atas matriks berbentuk
\[
  \begin{pmatrix} A&0\\0&1\end{pmatrix},
  \qquad A\in\mathrm{SO}(n-1),
\]
sehingga $S^{n-1}\cong\mathrm{SO}(n)/\mathrm{SO}(n-1)$. Rumus dimensi memberi
\[
 \frac{n(n-1)}2-\frac{(n-1)(n-2)}2=n-1,
\]
sesuai dengan dimensi bola satuan.
\end{beispiel}

\subsection{Ringkasan struktural}

Data global dan infinitesimal yang dibangun di atas tersusun sebagai berikut:
\[
\begin{array}{ccl}
G&\longmapsto&\mathfrak g=T_eG,\\
F\colon G\to H&\longmapsto&dF_e\colon\mathfrak g\to\mathfrak h,\\
v\in\mathfrak g&\longmapsto&X^v\text{ dan }\exp(tv),\\
g\in G&\longmapsto&\operatorname{Ad}_g,\\
v\in\mathfrak g&\longmapsto&\operatorname{ad}_v=[v,\cdot],\\
G\curvearrowright M&\longmapsto&v_M(p)=T_e\Phi_p(v).
\end{array}
\]
Eksponensial memberi koordinat lokal di sekitar identitas, tetapi tidak harus
surjektif secara global dan tidak mengubah persoalan global menjadi persoalan
linear. Orbit mengubah simetri menjadi geometri submanifold, sedangkan
stabilisator mengukur simetri yang tidak menggerakkan suatu titik.
% O011-BL-A — latihan, petunjuk, solusi, dan uji penguasaan asli.
% Lisensi: CC BY-SA 4.0.
% Provenans: disusun oleh OpenAI Codex gpt-5.6-sol, Ultra, atas arahan pengguna.

\subsection{Latihan, petunjuk bertahap, dan solusi}

Setiap latihan di bagian ini mempunyai identitas stabil yang tidak bergantung
pada bahasa. Petunjuk disusun bertahap: baca petunjuk kedua hanya jika petunjuk
pertama belum cukup. Solusi yang diberikan merupakan bagian asli modul ini,
bukan solusi yang dinisbahkan kepada sumber Brenner.

\subsubsection*{Latihan L1 — Grup perkalian positif (O011-BL-E01)}
\label{o011-bl-e01}

Tunjukkan bahwa $\mathbb R_{>0}$ dengan perkalian merupakan grup Lie
satu dimensi. Tentukan aljabar Lie dan pemetaan eksponensialnya. Tunjukkan pula
bahwa logaritma memberikan isomorfisme grup Lie
$\mathbb R_{>0}\cong(\mathbb R,+)$.

\paragraph{Petunjuk 1.}
Gunakan koordinat biasa pada himpunan terbuka $\mathbb R_{>0}\subset\mathbb R$.

\paragraph{Petunjuk 2.}
Subgrup satu-parameter berkecepatan awal $v$ harus memenuhi
$\gamma'(t)=v\gamma(t)$ dan $\gamma(0)=1$.

\paragraph{Solusi.}
Perkalian $(x,y)\mapsto xy$ dan invers $x\mapsto x^{-1}$ mulus pada
$\mathbb R_{>0}$. Ruang singgung di identitas $1$ dapat diidentifikasi dengan
$\mathbb R$. Karena grupnya abelian, braket Lie identik nol. Persamaan untuk
medan invarian-kiri yang bernilai $v$ di $1$ ialah
$\gamma'(t)=v\gamma(t)$, sehingga $\gamma_v(t)=e^{tv}$ dan
$\exp(v)=e^v$. Akhirnya $\log(xy)=\log x+\log y$; logaritma dan eksponensial
saling invers serta mulus.

\subsubsection*{Latihan L2 — Ruang singgung grup matriks (O011-BL-E02)}
\label{o011-bl-e02}

Buktikan langsung bahwa
\[
 T_I\mathrm{SL}_n(\mathbb R)=\{X:\operatorname{tr}X=0\}
 \quad\text{dan}\quad
 T_I\mathrm O(n)=\{X:X^{\mathsf T}+X=0\}.
\]
Hitung dimensinya.

\paragraph{Petunjuk 1.}
Diferensiasikan persamaan $\det A(t)=1$ dan
$A(t)^{\mathsf T}A(t)=I$ pada $t=0$.

\paragraph{Petunjuk 2.}
Untuk arah sebaliknya, gunakan $e^{tX}$ dan identitas
$\det(e^{tX})=e^{t\operatorname{tr}X}$; untuk matriks antisimetri gunakan
$(e^{tX})^{\mathsf T}e^{tX}=I$.

\paragraph{Solusi.}
Jika $A(0)=I$ dan $A'(0)=X$, rumus turunan determinan memberi
\[
 0=\left.\frac d{dt}\right|_0\det A(t)=\operatorname{tr}X.
\]
Sebaliknya, jika $\operatorname{tr}X=0$, kurva $e^{tX}$ berada di
$\mathrm{SL}_n$ dan berkecepatan $X$ pada nol. Jadi ruang singgung pertama
tepat seperti dinyatakan dan berdimensi $n^2-1$.

Untuk $A(t)\in\mathrm O(n)$,
\[
 0=\left.\frac d{dt}\right|_0 A(t)^{\mathsf T}A(t)=X^{\mathsf T}+X.
\]
Jika $X^{\mathsf T}=-X$, maka
$(e^{tX})^{\mathsf T}=e^{-tX}$, sehingga $e^{tX}\in\mathrm O(n)$. Matriks
antisimetri ditentukan oleh entri di atas diagonal; dimensinya
$n(n-1)/2$.

\subsubsection*{Latihan L3 — Medan invarian pada grup umum linear (O011-BL-E03)}
\label{o011-bl-e03}

Untuk $A\in M_n(\mathbb R)$, tentukan medan invarian-kiri $X^A$ dan medan
invarian-kanan $Y^A$ pada $\mathrm{GL}_n(\mathbb R)$. Hitung alirannya.

\paragraph{Petunjuk 1.}
Diferensial perkalian kiri $L_g(h)=gh$ pada $I$ adalah $A\mapsto gA$.

\paragraph{Petunjuk 2.}
Periksa kurva $g e^{tA}$ dan $e^{tA}g$.

\paragraph{Solusi.}
Kita memperoleh
\[
  X^A_g=gA,
  \qquad Y^A_g=Ag.
\]
Kurva integral $X^A$ yang berawal di $g$ ialah $t\mapsto ge^{tA}$, sebab
turunannya $ge^{tA}A$. Kurva integral $Y^A$ ialah
$t\mapsto e^{tA}g$. Keduanya lengkap untuk semua $t\in\mathbb R$.

\subsubsection*{Latihan L4 — Braket pada grup afin garis (O011-BL-E04)}
\label{o011-bl-e04}

Pertimbangkan
\[
 G=\left\{\begin{pmatrix}a&b\\0&1\end{pmatrix}:a>0,
 b\in\mathbb R\right\}.
\]
Tentukan aljabar Lie dan braketnya. Apakah aljabar Lie ini abelian?

\paragraph{Petunjuk 1.}
Tuliskan arah singgung di identitas sebagai
$X(x,y)=\left(\begin{smallmatrix}x&y\\0&0\end{smallmatrix}\right)$.

\paragraph{Petunjuk 2.}
Hitung $X(x,y)X(u,v)-X(u,v)X(x,y)$.

\paragraph{Solusi.}
Aljabar Lie terdiri atas semua
\[
 X(x,y)=\begin{pmatrix}x&y\\0&0\end{pmatrix}.
\]
Perkalian langsung memberi
\[
 [X(x,y),X(u,v)]
 =\begin{pmatrix}0&xv-uy\\0&0\end{pmatrix}
 =X(0,xv-uy).
\]
Aljabar ini tidak abelian; misalnya
$[X(1,0),X(0,1)]=X(0,1)$.

\subsubsection*{Latihan L5 — Eksponensial rotasi dan geseran (O011-BL-E05)}
\label{o011-bl-e05}

Hitung eksponensial dari
\[
 J=\begin{pmatrix}0&-1\\1&0\end{pmatrix},
 \qquad
 N=\begin{pmatrix}0&1\\0&0\end{pmatrix}.
\]
Jelaskan subgrup satu-parameter yang dihasilkan.

\paragraph{Petunjuk 1.}
Gunakan $J^2=-I$ dan $N^2=0$.

\paragraph{Petunjuk 2.}
Pisahkan deret pangkat $e^{tJ}$ menjadi suku genap dan ganjil.

\paragraph{Solusi.}
Karena $J^{2k}=(-1)^kI$ dan $J^{2k+1}=(-1)^kJ$,
\[
 e^{tJ}=\cos(t)I+\sin(t)J
 =\begin{pmatrix}\cos t&-\sin t\\\sin t&\cos t\end{pmatrix}.
\]
Ini adalah subgrup rotasi di $\mathrm{SO}(2)$. Karena $N^2=0$,
\[
 e^{tN}=I+tN=\begin{pmatrix}1&t\\0&1\end{pmatrix},
\]
yakni subgrup geseran unipoten.

\subsubsection*{Latihan L6 — Determinan eksponensial (O011-BL-E06)}
\label{o011-bl-e06}

Buktikan bahwa untuk setiap matriks nyata $X$,
\[
  \det(e^X)=e^{\operatorname{tr}X}.
\]
Gunakan hasil ini untuk menunjukkan bahwa eksponensial
$\mathfrak{sl}_n\to\mathrm{SL}_n$ terdefinisi dengan baik.

\paragraph{Petunjuk 1.}
Tinjau $f(t)=\det(e^{tX})$ dan gunakan rumus Jacobi untuk turunan determinan.

\paragraph{Petunjuk 2.}
Tunjukkan $f'(t)=\operatorname{tr}(X)f(t)$ dengan $f(0)=1$.

\paragraph{Solusi.}
Karena $(e^{tX})'=Xe^{tX}$ dan $e^{tX}$ invertibel,
\[
 f'(t)=f(t)\operatorname{tr}\bigl(e^{-tX}Xe^{tX}\bigr)
 =f(t)\operatorname{tr}X.
\]
Solusi persamaan skalar ini dengan $f(0)=1$ ialah
$f(t)=e^{t\operatorname{tr}X}$. Ambil $t=1$. Jika
$\operatorname{tr}X=0$, determinan $e^X$ sama dengan $1$.

\subsubsection*{Latihan L7 — Adjoint grup matriks (O011-BL-E07)}
\label{o011-bl-e07}

Untuk subgrup matriks $G\subseteq\mathrm{GL}_n(\mathbb R)$, buktikan
\[
 \operatorname{Ad}_gX=gXg^{-1},
 \qquad
 \operatorname{ad}_XY=XY-YX.
\]

\paragraph{Petunjuk 1.}
Diferensiasikan $g(I+tX+o(t))g^{-1}$.

\paragraph{Petunjuk 2.}
Kemudian diferensiasikan $e^{tX}Ye^{-tX}$ pada $t=0$.

\paragraph{Solusi.}
Kurva $A(t)$ melalui $I$ dengan $A'(0)=X$ memberi
\[
 \left.\frac d{dt}\right|_0gA(t)g^{-1}=gXg^{-1},
\]
yang membuktikan rumus $\operatorname{Ad}$. Selanjutnya
\[
 \left.\frac d{dt}\right|_0 e^{tX}Ye^{-tX}
 =XY-YX,
\]
sehingga $\operatorname{ad}_X(Y)=[X,Y]$.

\subsubsection*{Latihan L8 — Komutator infinitesimal (O011-BL-E08)}
\label{o011-bl-e08}

Untuk $X,Y\in M_n(\mathbb R)$, tetapkan
\[
 c(t)=e^{tX}e^{tY}e^{-tX}e^{-tY}.
\]
Tunjukkan bahwa
\[
 c(t)=I+t^2[X,Y]+O(t^3).
\]

\paragraph{Petunjuk 1.}
Gunakan $e^{tX}=I+tX+\tfrac12t^2X^2+O(t^3)$ untuk keempat faktor.

\paragraph{Petunjuk 2.}
Semua suku orde satu harus saling menghapus; kelompokkan suku orde dua.

\paragraph{Solusi.}
Kalikan dua faktor pertama dan dua faktor terakhir hingga orde dua:
\[
 e^{tX}e^{tY}
 =I+t(X+Y)+t^2\left(\tfrac12X^2+XY+\tfrac12Y^2\right)+O(t^3),
\]
\[
 e^{-tX}e^{-tY}
 =I-t(X+Y)+t^2\left(\tfrac12X^2+XY+\tfrac12Y^2\right)+O(t^3).
\]
Ketika kedua ekspansi dikalikan, koefisien orde satu hilang dan koefisien
orde dua menyederhana menjadi $XY-YX$. Jadi rumus yang diminta berlaku.
Rumus ini memperlihatkan bahwa braket mengukur kegagalan gerak kecil dalam
dua arah untuk saling komutatif.

\subsubsection*{Latihan L9 — Aksi lingkaran (O011-BL-E09)}
\label{o011-bl-e09}

Biarkan $\mathbb R$ bertindak pada $S^1\subset\mathbb C$ melalui
$t\cdot z=e^{it}z$. Tentukan orbit, stabilisator, medan fundamental untuk
$v\in\mathbb R$, dan kernel diferensial pemetaan orbit di $0$.

\paragraph{Petunjuk 1.}
Untuk $z$ tetap, selesaikan $e^{it}z=z$.

\paragraph{Petunjuk 2.}
Diferensiasikan $e^{itv}z$ pada $t=0$.

\paragraph{Solusi.}
Orbit setiap $z$ adalah seluruh $S^1$. Stabilisatornya
$2\pi\mathbb Z$, suatu subgrup diskret. Medan fundamental ialah
\[
 v_{S^1}(z)=ivz.
\]
Pemetaan $v\mapsto ivz$ dari $\mathbb R$ ke $T_zS^1$ adalah isomorfisme,
sehingga kernelnya nol. Ini sesuai dengan aljabar Lie stabilisator diskret,
yang juga nol, meskipun stabilisator grupnya sendiri tidak trivial.

\subsubsection*{Latihan L10 — Orbit rotasi di ruang Euklides (O011-BL-E10)}
\label{o011-bl-e10}

Tentukan orbit dan stabilisator aksi standar $\mathrm{SO}(3)$ pada
$\mathbb R^3$. Untuk $p\ne0$, buktikan
\[
  T_p(\mathrm{SO}(3)\cdot p)=p^\perp.
\]

\paragraph{Petunjuk 1.}
Rotasi mempertahankan norma dan bertindak transitif pada setiap bola berjari-jari
positif.

\paragraph{Petunjuk 2.}
Medan fundamental yang ditentukan $X\in\mathfrak{so}(3)$ bernilai $Xp$.
Gunakan antisimetri untuk menunjukkan $Xp\perp p$ dan hitung dimensi.

\paragraph{Solusi.}
Orbit $0$ adalah $\{0\}$ dan stabilisatornya seluruh $\mathrm{SO}(3)$. Jika
$p\ne0$, orbitnya bola $S^2_{\|p\|}$ dan stabilisatornya isomorfik dengan
$\mathrm{SO}(2)$, yaitu rotasi di bidang $p^\perp$. Karena
$\langle Xp,p\rangle=-\langle p,Xp\rangle$, setiap $Xp$ tegak lurus pada
$p$. Pemetaan $X\mapsto Xp$ berperingkat dua: kernelnya berdimensi satu,
yakni aljabar stabilisator. Jadi citranya adalah seluruh ruang dua dimensi
$p^\perp$.

\subsubsection*{Latihan L11 — Orbit konjugasi (O011-BL-E11)}
\label{o011-bl-e11}

$\mathrm{GL}_n(\mathbb R)$ bertindak pada $M_n(\mathbb R)$ melalui
$g\cdot A=gAg^{-1}$. Tentukan stabilisator $A$ dan ruang tangensial orbitnya
di $A$.

\paragraph{Petunjuk 1.}
Stabilisator adalah semua matriks invertibel yang komutatif dengan $A$.

\paragraph{Petunjuk 2.}
Diferensiasikan $e^{tX}Ae^{-tX}$.

\paragraph{Solusi.}
Stabilisatornya adalah sentralisator invertibel
\[
 (\mathrm{GL}_n)_A=\{g\in\mathrm{GL}_n:gA=Ag\}.
\]
Medan fundamental arah $X$ bernilai
\[
 \left.\frac d{dt}\right|_0e^{tX}Ae^{-tX}=XA-AX=[X,A].
\]
Oleh karena itu
\[
 T_A(\mathrm{GL}_n\cdot A)=\{[X,A]:X\in M_n(\mathbb R)\}.
\]
Kernelnya adalah aljabar semua matriks yang komutatif dengan $A$.

\subsubsection*{Latihan L12 — Kernel homomorfisme (O011-BL-E12)}
\label{o011-bl-e12}

Misalkan $F\colon G\to H$ homomorfisme grup Lie. Buktikan bahwa
$\ker(dF_e)$ adalah ideal dalam $\mathfrak g$ dan
$\operatorname{im}(dF_e)$ adalah subaljabar Lie dalam $\mathfrak h$.

\paragraph{Petunjuk 1.}
Gunakan $dF_e[v,w]=[dF_ev,dF_ew]$.

\paragraph{Petunjuk 2.}
Untuk sifat ideal, ambil $v$ di kernel dan $w$ sebarang.

\paragraph{Solusi.}
Jika $v\in\ker(dF_e)$ dan $w\in\mathfrak g$, maka
\[
 dF_e[v,w]=[dF_ev,dF_ew]=[0,dF_ew]=0.
\]
Jadi $[v,w]$ kembali berada dalam kernel; kernel adalah ideal. Jika
$a=dF_ev$ dan $b=dF_ew$ berada dalam citra, maka
\[
 [a,b]=[dF_ev,dF_ew]=dF_e[v,w]
\]
juga berada dalam citra. Jadi citra adalah subaljabar Lie.

\subsection{Uji penguasaan kumulatif}

Keempat soal berikut dimaksudkan sebagai penilaian setelah seluruh modul.
Setiap soal mempunyai rubrik, solusi lengkap, dan parameter alternatif untuk
menghasilkan bentuk ekuivalen tanpa mengubah kompetensi yang diuji.

\subsubsection*{Penguasaan M1 — Grup Heisenberg (O011-BL-M01)}
\label{o011-bl-m01}

Untuk
\[
 H=\left\{h(x,y,z)=
 \begin{pmatrix}1&x&z\\0&1&y\\0&0&1\end{pmatrix}:x,y,z\in\mathbb R\right\},
\]
(a) buktikan bahwa $H$ adalah grup Lie; (b) tentukan hukum grup dalam
koordinat $(x,y,z)$; (c) tentukan aljabar Lie dan semua braket basis
$X=E_{12}$, $Y=E_{23}$, $Z=E_{13}$; (d) hitung
$\exp(aX+bY+cZ)$.

\paragraph{Petunjuk.}
Semua matriks $N=aX+bY+cZ$ memenuhi $N^3=0$. Perhatikan $XY=Z$ tetapi
$YX=0$.

\paragraph{Solusi.}
Himpunan ini adalah submanifold afin tiga dimensi dari $M_3(\mathbb R)$.
Perkalian dan invers matriks membatas menjadi pemetaan mulus, dan
\[
 h(x,y,z)h(x',y',z')
 =h(x+x',y+y',z+z'+xy').
\]
Inversnya $h(x,y,z)^{-1}=h(-x,-y,-z+xy)$. Aljabar Lie ialah
$\operatorname{span}\{X,Y,Z\}$ dengan
\[
 [X,Y]=Z,
 \qquad [X,Z]=[Y,Z]=0.
\]
Untuk $N=aX+bY+cZ$, berlaku $N^2=abZ$ dan $N^3=0$, sehingga
\[
 \exp N=I+N+\tfrac12N^2
 =h\left(a,b,c+\tfrac12ab\right).
\]

\paragraph{Rubrik (12 poin).}
Dua poin untuk manifold dan kemulusan operasi; tiga poin untuk hukum grup dan
invers; tiga poin untuk ruang aljabar serta semua braket; empat poin untuk
perhitungan eksponensial beserta alasan deret berhenti.

\paragraph{Parameter alternatif.}
Ganti koordinat entri kanan atas dengan $\lambda z$ untuk suatu
$\lambda\ne0$. Peserta harus menyesuaikan hukum grup, basis pusat, dan
koefisien eksponensial; kompetensi yang diuji tetap sama.

\subsubsection*{Penguasaan M2 — Grup afin dan eksponensial (O011-BL-M02)}
\label{o011-bl-m02}

Untuk grup $G$ pada Latihan L4, (a) hitung
$\exp X(x,y)$; (b) tentukan $\operatorname{Ad}_{g(a,b)}X(x,y)$; (c) verifikasi
langsung $\operatorname{ad}_{X(x,y)}X(u,v)=[X(x,y),X(u,v)]$.

\paragraph{Petunjuk.}
Jika $x\ne0$, jumlahkan deret entri kanan atas; perlakukan $x=0$ sebagai
limit. Gunakan
$g(a,b)^{-1}=\left(\begin{smallmatrix}a^{-1}&-a^{-1}b\\0&1\end{smallmatrix}\right)$.

\paragraph{Solusi.}
Pangkat-pangkat $X=X(x,y)$ memenuhi, untuk $k\ge1$,
\[
 X^k=\begin{pmatrix}x^k&x^{k-1}y\\0&0\end{pmatrix}.
\]
Maka
\[
 \exp X(x,y)=
 \begin{cases}
 \left(\begin{smallmatrix}e^x&y(e^x-1)/x\\0&1\end{smallmatrix}\right),&x\ne0,\\[2mm]
 \left(\begin{smallmatrix}1&y\\0&1\end{smallmatrix}\right),&x=0.
 \end{cases}
\]
Perkalian matriks memberi
\[
 \operatorname{Ad}_{g(a,b)}X(x,y)=X(x,ay-bx).
\]
Diferensiasikan rumus ini pada $g=\exp(tX(x,y))$ dan $t=0$, atau hitung
komutator langsung, untuk memperoleh
\[
 \operatorname{ad}_{X(x,y)}X(u,v)=X(0,xv-uy).
\]

\paragraph{Rubrik (12 poin).}
Lima poin untuk eksponensial termasuk kasus $x=0$; empat poin untuk adjoint;
tiga poin untuk diferensiasi atau komutator yang memverifikasi rumus
$\operatorname{ad}$.

\paragraph{Parameter alternatif.}
Gunakan representasi afin dengan matriks
$\left(\begin{smallmatrix}a&0&b\\0&1&0\\0&0&1\end{smallmatrix}\right)$.
Peserta harus mengenali bahwa perhitungan esensial tetap dua dimensi.

\subsubsection*{Penguasaan M3 — Bola sebagai ruang homogen (O011-BL-M03)}
\label{o011-bl-m03}

Untuk aksi $\mathrm{SO}(n)$ pada $S^{n-1}$ dengan $n\ge3$, (a) buktikan
transitivitas; (b) tentukan stabilisator $e_n$; (c) hitung diferensial pemetaan
orbit di identitas; (d) deduksi
$S^{n-1}\cong\mathrm{SO}(n)/\mathrm{SO}(n-1)$ dan periksa dimensinya.

\paragraph{Petunjuk.}
Lengkapi setiap vektor satuan menjadi basis ortonormal berorientasi. Untuk
$X\in\mathfrak{so}(n)$, diferensial pemetaan orbit adalah $X\mapsto Xe_n$.

\paragraph{Solusi.}
Jika $p,q\in S^{n-1}$, lengkapi masing-masing menjadi basis ortonormal
berorientasi. Pemetaan yang membawa basis pertama ke basis kedua berada dalam
$\mathrm{SO}(n)$ dan membawa $p$ ke $q$, sehingga aksi transitif. Matriks yang
menetapkan $e_n$ mempunyai bentuk blok
$\operatorname{diag}(A,1)$ dengan $A\in\mathrm{SO}(n-1)$.

Pemetaan orbit $\Phi_{e_n}(g)=ge_n$ mempunyai diferensial
$T_I\Phi_{e_n}(X)=Xe_n$. Citranya adalah $e_n^\perp$, sedangkan kernelnya
terdiri atas matriks antisimetri dengan baris dan kolom terakhir nol, yaitu
$\mathfrak{so}(n-1)$. Karena itu ruang homogen adalah $S^{n-1}$ dan
\[
 \dim\mathrm{SO}(n)-\dim\mathrm{SO}(n-1)
 =\frac{n(n-1)-(n-1)(n-2)}2=n-1.
\]

\paragraph{Rubrik (12 poin).}
Tiga poin untuk transitivitas; dua poin untuk stabilisator; empat poin untuk
diferensial, kernel, dan citra; tiga poin untuk identifikasi hasil bagi dan
perhitungan dimensi.

\paragraph{Parameter alternatif.}
Ganti bola satuan dengan bola berjari-jari $r>0$ dan titik dasar $re_1$.
Stabilisator dan hasil bagi sama hingga pilihan titik dasar; diferensial
pemetaan orbit memperoleh faktor $r$.

\subsubsection*{Penguasaan M4 — Naturalisasi eksponensial dan kegagalan komutatif (O011-BL-M04)}
\label{o011-bl-m04}

Misalkan $F\colon G\to H$ homomorfisme grup Lie. (a) buktikan
$F(\exp_Gv)=\exp_H(dF_ev)$; (b) deduksi
$F(\exp_G(tv))=\exp_H(t,dF_ev)$ untuk semua $t$; (c) berikan matriks
$X,Y\in\mathfrak{gl}_2$ yang menunjukkan bahwa umumnya
$e^{X+Y}\ne e^Xe^Y$; (d) nyatakan hipotesis sederhana yang menjamin kesamaan.

\paragraph{Petunjuk.}
Untuk (a), bandingkan dua subgrup satu-parameter dengan kecepatan awal yang
sama. Untuk (c), gunakan kelipatan kecil matriks $E_{12}$ dan $E_{21}$ atau
bandingkan suku orde dua.

\paragraph{Solusi.}
Kurva $t\mapsto F(\exp_G(tv))$ adalah homomorfisme
$\mathbb R\to H$ dan turunannya pada nol ialah $dF_ev$. Berdasarkan ketunggalan
subgrup satu-parameter, kurva itu sama dengan
$t\mapsto\exp_H(t,dF_ev)$. Pernyataan (a) adalah kasus $t=1$, sedangkan (b)
adalah identitas kurva penuh.

Ambil $X=sE_{12}$ dan $Y=sE_{21}$ dengan $s\ne0$ cukup kecil. Ekspansi hingga
orde dua memberi
\[
 e^Xe^Y=I+X+Y+XY+O(s^3),
\]
sedangkan
\[
 e^{X+Y}=I+X+Y+\tfrac12(XY+YX)+O(s^3).
\]
Karena $XY\ne YX$, kedua matriks berbeda untuk $s$ cukup kecil. Jika
$[X,Y]=0$, deret pangkat dapat dikelompokkan seperti bilangan komutatif dan
$e^{X+Y}=e^Xe^Y$.

\paragraph{Rubrik (12 poin).}
Empat poin untuk argumen subgrup satu-parameter; dua poin untuk identitas
berparameter; empat poin untuk contoh tandingan yang benar beserta verifikasi;
dua poin untuk hipotesis komutatif dan alasannya.

\paragraph{Parameter alternatif.}
Ganti pasangan $E_{12},E_{21}$ dengan
$X=\left(\begin{smallmatrix}0&s\\0&0\end{smallmatrix}\right)$ dan
$Y=\left(\begin{smallmatrix}r&0\\0&0\end{smallmatrix}\right)$ untuk $rs\ne0$.
Peserta harus mendeteksi komutator tak nol dan membandingkan eksponensial.

\chapter{Jembatan Kohomologi de Rham dan Topologi Diferensial}
\noindent\textit{Materi asli edisi ini; bukan solusi atau teks yang disediakan oleh sumber Brenner/Wikiversity.}
% O011-BR — modul asli, bukan terjemahan atau adaptasi dari teks pembanding.
% Judul: Jembatan Kohomologi de Rham dan Topologi Diferensial
% Lisensi modul: CC BY-SA 4.0.
% Provenans: disusun oleh OpenAI Codex gpt-5.6-sol, Ultra, atas arahan pengguna.
% Prasyarat internal: Unit 14, 20, 22, dan 23 edisi ini.

\setcounter{section}{31}
\section{Jembatan Kohomologi de Rham dan Topologi Diferensial}
\label{o011-bridge-de-rham}

\subsection{Tujuan dan prasyarat}

Unit 14 memperkenalkan bentuk diferensial, Unit 20 turunan eksterior, Unit 22
partisi kesatuan, dan Unit 23 teorema Stokes. Modul ini menyatukan keempatnya
menjadi sebuah invariant topologis. Gagasan pokoknya sederhana: bentuk tertutup
yang bukan turunan eksterior bentuk lain mendeteksi ``lubang'' pada manifold.
Kita membuktikan hasil lokal, membangun invariansi homotopi, memakai barisan
Mayer--Vietoris untuk perhitungan global, lalu menandai jalan menuju derajat,
transversalitas, dan teori Morse tanpa berpura-pura membahas teori-teori itu
secara lengkap.

Semua manifold di bagian ini dianggap mulus, berdimensi hingga, Hausdorff, dan
memiliki basis terhitung. Ruang bentuk-$k$ mulus pada $M$ ditulis
$\Omega^k(M)$; untuk $k<0$ atau $k>\dim M$, ruang itu adalah nol.

\subsection{Kompleks de Rham}

\begin{definition}\label{o011-br-def-closed-exact}
Suatu bentuk $\omega\in\Omega^k(M)$ disebut \emph{tertutup} jika
$d\omega=0$, dan disebut \emph{eksak} jika terdapat
$\eta\in\Omega^{k-1}(M)$ dengan $\omega=d\eta$.
\end{definition}

Karena $d\circ d=0$, setiap bentuk eksak tertutup. Jadi turunan eksterior
membentuk kompleks rantai-kobalik
\[
  0\longrightarrow\Omega^0(M)\xrightarrow{d}\Omega^1(M)
  \xrightarrow{d}\cdots\xrightarrow{d}\Omega^n(M)
  \longrightarrow0.
\]

\begin{definition}\label{o011-br-def-cohomology}
Kohomologi de Rham derajat $k$ adalah ruang vektor
\[
  H^k_{\mathrm{dR}}(M)
  =\frac{\ker(d\colon\Omega^k(M)\to\Omega^{k+1}(M))}
  {\operatorname{im}(d\colon\Omega^{k-1}(M)\to\Omega^k(M))}.
\]
Kelas bentuk tertutup $\omega$ ditulis $[\omega]$.
\end{definition}

\begin{Proposition}\label{o011-br-prop-wedge-ring}
Produk wedge menurunkan produk bergradasi
\[
 H^p_{\mathrm{dR}}(M)\times H^q_{\mathrm{dR}}(M)
 \longrightarrow H^{p+q}_{\mathrm{dR}}(M),
 \qquad ([\alpha],[\beta])\longmapsto[\alpha\wedge\beta].
\]
Produk ini memenuhi
$[\alpha]\wedge[\beta]=(-1)^{pq}[\beta]\wedge[\alpha]$.
\end{Proposition}

\begin{proof}
Aturan Leibniz bergradasi memberi
\[
 d(\alpha\wedge\beta)=d\alpha\wedge\beta
 +(-1)^p\alpha\wedge d\beta.
\]
Jika $\alpha$ dan $\beta$ tertutup, wedge-nya tertutup. Jika
$\alpha$ diganti $\alpha+d\mu$, maka
\[
 (\alpha+d\mu)\wedge\beta-\alpha\wedge\beta
 =d(\mu\wedge\beta)
\]
karena $d\beta=0$. Argumen yang sama berlaku pada faktor kedua. Jadi produk
tidak bergantung pada wakil. Komutativitas bergradasi diwarisi dari bentuk.
\end{proof}

\begin{Proposition}\label{o011-br-prop-h0-components}
$H^0_{\mathrm{dR}}(M)$ adalah ruang fungsi yang konstan pada setiap komponen
terhubung. Khususnya, jika $M$ terhubung, $H^0_{\mathrm{dR}}(M)\cong\mathbb R$.
\end{Proposition}

\begin{proof}
Untuk fungsi mulus $f$, persamaan $df=0$ berarti semua turunan arah lokalnya
nol. Maka $f$ konstan pada setiap komponen terhubung. Tidak ada bentuk
berderajat $-1$, sehingga pada derajat nol tidak ada hasil bagi tambahan.
\end{proof}

\subsection{Tarikan balik dan funktorialitas}

Jika $f\colon M\to N$ mulus, tarikan balik memenuhi
\[
 f^*(\alpha\wedge\beta)=f^*\alpha\wedge f^*\beta,
 \qquad d(f^*\alpha)=f^*(d\alpha).
\]

\begin{Proposition}\label{o011-br-prop-pullback-cohomology}
Pemetaan $f$ menginduksi homomorfisme aljabar bergradasi
\[
 f^*\colon H^*_{\mathrm{dR}}(N)\longrightarrow H^*_{\mathrm{dR}}(M).
\]
Jika $g\colon N\to P$ mulus, maka $(g\circ f)^*=f^*\circ g^*$, dan
$\operatorname{id}_M^*$ adalah identitas.
\end{Proposition}

\begin{proof}
Komutasi dengan $d$ menunjukkan bahwa tarikan balik membawa bentuk tertutup
ke bentuk tertutup dan bentuk eksak ke bentuk eksak. Karena itu ia terdefinisi
pada kelas kohomologi. Semua identitas lain sudah berlaku pada tingkat bentuk.
\end{proof}

Perhatikan pembalikan arah: pemetaan $M\to N$ menghasilkan pemetaan
kohomologi dari $N$ ke $M$. Kohomologi de Rham adalah funktor kontravarian.

\subsection{Lema Poincar\'e pada himpunan berbentuk bintang}

Suatu himpunan terbuka $U\subseteq\mathbb R^n$ disebut berbentuk bintang
terhadap $0$ jika $tx\in U$ untuk semua $x\in U$ dan $t\in[0,1]$. Misalkan
$E=\sum_i x_i\partial_{x_i}$ adalah medan Euler dan $\iota_E$ kontraksi oleh
$E$.

Untuk $\omega\in\Omega^k(U)$ dengan $k\ge1$, definisikan operator homotopi oleh
\[
 (K\omega)_x(v_1,\ldots,v_{k-1})
 =\int_0^1 t^{k-1}\omega_{tx}(x,v_1,\ldots,v_{k-1})\,dt.
\]
Dengan medan Euler $E_y=y$, integran yang sama dapat ditulis
$t^{k-2}(\iota_E\omega)_{tx}$ untuk $t>0$; rumus pertama tetap teratur di
$t=0$ dan karena itu kita pakai sebagai definisi. Rumus yang lebih intrinsik
akan muncul pada homotopi umum di bawah.

\begin{Theorem}[Lema Poincar\'e]\label{o011-br-thm-poincare}
Jika $U\subseteq\mathbb R^n$ terbuka dan berbentuk bintang, setiap bentuk
tertutup $\omega\in\Omega^k(U)$ dengan $k\ge1$ adalah eksak. Dengan kata lain,
$H^k_{\mathrm{dR}}(U)=0$ untuk $k\ge1$.
\end{Theorem}

\begin{proof}
Ambil homotopi $F\colon U\times[0,1]\to U$, $F(x,t)=tx$. Perhitungan di
koordinat, atau rumus homotopi yang dibuktikan pada Teorema
\ref{o011-br-thm-chain-homotopy}, memberi
\[
 dK\omega+K(d\omega)=F_1^*\omega-F_0^*\omega.
\]
Untuk $k\ge1$, tarikan balik bentuk-$k$ oleh pemetaan konstan $F_0$ adalah nol,
sedangkan $F_1$ adalah identitas. Jika $d\omega=0$, maka
$\omega=d(K\omega)$.
\end{proof}

\begin{bemerkung}\label{o011-br-rem-local-global}
Lema Poincar\'e bersifat lokal: setiap titik manifold mempunyai lingkungan
koordinat yang dapat diperkecil menjadi bola, sehingga setiap bentuk tertutup
berderajat positif lokalnya eksak. Hambatan untuk memperoleh primitif global
adalah tepat informasi yang diukur kohomologi de Rham.
\end{bemerkung}

\subsection{Homotopi rantai dan invariansi homotopi}

Misalkan $F\colon M\times[0,1]\to N$ suatu homotopi mulus dan
$F_t(p)=F(p,t)$. Untuk $\omega\in\Omega^k(N)$, definisikan
\[
 K_F\omega
 =\int_0^1 \iota_t^*\bigl(\iota_{\partial_t}F^*\omega\bigr)\,dt
 \in\Omega^{k-1}(M),
\]
di mana $\iota_t(p)=(p,t)$.

\begin{Theorem}[Rumus homotopi rantai]\label{o011-br-thm-chain-homotopy}
Operator $K_F$ memenuhi
\[
 dK_F+K_Fd=F_1^*-F_0^*.
\]
\end{Theorem}

\begin{proof}
Pada $M\times[0,1]$, turunan terhadap $t$ dari tarikan balik sepanjang
$\iota_t$ memenuhi rumus Cartan
\[
 \frac d{dt}\iota_t^*(F^*\omega)
 =\iota_t^*\mathcal L_{\partial_t}(F^*\omega)
 =\iota_t^*\bigl(d\iota_{\partial_t}F^*\omega
 +\iota_{\partial_t}dF^*\omega\bigr).
\]
Integrasikan dari $0$ sampai $1$, gunakan bahwa tarikan balik komutatif dengan
$d$, dan kenali kedua suku sebagai $dK_F\omega$ serta $K_Fd\omega$.
\end{proof}

\begin{Korollar}\label{o011-br-cor-homotopy-invariance}
Pemetaan mulus yang homotopik menginduksi pemetaan kohomologi yang sama. Jika
$M$ dan $N$ ekuivalen-homotopi, maka
$H^*_{\mathrm{dR}}(M)\cong H^*_{\mathrm{dR}}(N)$ sebagai aljabar bergradasi.
\end{Korollar}

\begin{proof}
Jika $d\omega=0$, rumus homotopi memberi
$F_1^*\omega-F_0^*\omega=d(K_F\omega)$, sehingga kedua tarikan balik menentukan
kelas yang sama. Untuk ekuivalensi homotopi, terapkan funktorialitas pada
komposisi yang homotopik dengan identitas.
\end{proof}

\begin{Korollar}\label{o011-br-cor-contractible}
Jika $M$ kontraktibel, maka $H^0_{\mathrm{dR}}(M)\cong\mathbb R$ dan
$H^k_{\mathrm{dR}}(M)=0$ untuk $k>0$.
\end{Korollar}

\subsection{Barisan Mayer--Vietoris}

Misalkan $M=U\cup V$ dengan $U,V$ terbuka. Untuk setiap $k$ terdapat barisan
pendek kompleks
\[
0\longrightarrow\Omega^k(M)
\xrightarrow{r}(\Omega^k(U)\oplus\Omega^k(V))
\xrightarrow{\delta}\Omega^k(U\cap V)
\longrightarrow0,
\]
dengan
\[
 r(\omega)=(\omega|_U,\omega|_V),
 \qquad \delta(\alpha,\beta)=\alpha|_{U\cap V}-\beta|_{U\cap V}.
\]
Injektivitas dan ketepatan di tengah mengikuti sifat pengeleman bentuk. Untuk
surjektivitas $\delta$, pilih partisi kesatuan $\rho_U+\rho_V=1$ yang
tersubordinasi pada $U,V$; suatu bentuk $\gamma$ di $U\cap V$ diangkat oleh
perpanjangan nol yang dibangun dari $\rho_V\gamma$ pada $U$ dan
$-\rho_U\gamma$ pada $V$.

\begin{Theorem}[Mayer--Vietoris]\label{o011-br-thm-mv}
Penutup $M=U\cup V$ menghasilkan barisan eksak panjang
\[
\cdots\longrightarrow H^{k-1}_{\mathrm{dR}}(U\cap V)
\xrightarrow{\partial}H^k_{\mathrm{dR}}(M)
\xrightarrow{r^*}H^k_{\mathrm{dR}}(U)\oplus H^k_{\mathrm{dR}}(V)
\xrightarrow{\delta^*}H^k_{\mathrm{dR}}(U\cap V)
\xrightarrow{\partial}H^{k+1}_{\mathrm{dR}}(M)
\longrightarrow\cdots.
\]
\end{Theorem}

\begin{proof}
Barisan pendek kompleks di atas menghasilkan barisan eksak panjang kohomologi
melalui konstruksi standar pemetaan penghubung. Secara konkret, jika
$\gamma$ tertutup pada $U\cap V$, pilih $(\alpha,\beta)$ dengan
$\delta(\alpha,\beta)=\gamma$. Karena $d\gamma=0$, pasangan
$(d\alpha,d\beta)$ saling cocok pada irisan, sehingga menempel menjadi bentuk
tertutup $\omega$ pada $M$. Tetapkan $\partial[\gamma]=[\omega]$. Perubahan
pilihan hanya mengubah $\omega$ dengan bentuk eksak, dan pemeriksaan langsung
memberi ketepatan pada setiap suku.
\end{proof}

\subsection{Perhitungan dasar}

\begin{Proposition}\label{o011-br-prop-circle}
Kohomologi de Rham lingkaran adalah
\[
 H^0_{\mathrm{dR}}(S^1)\cong\mathbb R,
 \qquad H^1_{\mathrm{dR}}(S^1)\cong\mathbb R,
 \qquad H^k_{\mathrm{dR}}(S^1)=0\quad(k\ge2).
\]
Isomorfisme pada derajat satu dapat diberikan oleh periode
\[
 [\omega]\longmapsto\int_{S^1}\omega.
\]
\end{Proposition}

\begin{proof}
Derajat nol mengikuti keterhubungan, dan derajat di atas satu nol karena
dimensi. Tutupi $S^1$ oleh dua busur terbuka kontraktibel $U,V$ yang irisannya
mempunyai dua komponen. Bagian relevan Mayer--Vietoris adalah
\[
 0\to H^0(S^1)\to\mathbb R\oplus\mathbb R
 \to\mathbb R^2\to H^1(S^1)\to0.
\]
Pemetaan tengah mengirim $(a,b)$ ke $(a-b,a-b)$, sehingga cokernelnya satu
dimensi. Integral lenyap pada bentuk eksak dan tidak nol pada bentuk sudut
ternormalisasi; karena itu periode memberi isomorfisme yang dinyatakan.
\end{proof}

Pada $\mathbb R^2\setminus\{0\}$, bentuk
\[
 \eta=\frac{1}{2\pi}\frac{x\,dy-y\,dx}{x^2+y^2}
\]
tertutup dan integralnya pada lingkaran satuan adalah $1$. Jadi ia tidak eksak
dan mewakili generator kelas derajat satu.

\begin{Theorem}\label{o011-br-thm-sphere}
Untuk $n\ge1$,
\[
 H^k_{\mathrm{dR}}(S^n)\cong
 \begin{cases}
 \mathbb R,&k=0\text{ atau }k=n,\\
 0,&\text{selain itu}.
 \end{cases}
\]
\end{Theorem}

\begin{proof}
Kasus $n=1$ sudah dibuktikan. Untuk $n\ge2$, tutupi $S^n$ oleh lingkungan
terbuka kutub utara dan selatan yang masing-masing kontraktibel dan yang
irisannya ekuivalen-homotopi dengan $S^{n-1}$. Mayer--Vietoris, di luar
derajat nol, memberi isomorfisme penghubung
\[
 H^{k-1}_{\mathrm{dR}}(S^{n-1})\cong H^k_{\mathrm{dR}}(S^n)
 \qquad(1<k\le n).
\]
Pada derajat satu, keterhubungan penutup dan irisan memberi nol. Induksi pada
$n$ menghasilkan satu kelas di derajat teratas dan nol pada derajat antara.
Derajat nol kembali mengikuti keterhubungan.
\end{proof}

Jika $S^n$ diberi orientasi, kelas teratas dapat dinormalisasi oleh
$\int_{S^n}\omega=1$. Teorema Stokes memastikan bahwa integral hanya
bergantung pada kelas kohomologi.

\subsection{Teorema de Rham dan gerbang topologi diferensial}

\begin{Theorem}[Teorema de Rham — pernyataan]\label{o011-br-thm-de-rham}
Integrasi bentuk pada rantai singular mulus menginduksi isomorfisme alami
\[
 H^k_{\mathrm{dR}}(M)\xrightarrow{\ \cong\ }
 H^k_{\mathrm{sing}}(M;\mathbb R)
\]
untuk setiap manifold mulus $M$ dan setiap $k$.
\end{Theorem}

Teorema ini dinyatakan, bukan dibuktikan, karena pembuktian lengkap memerlukan
kompleks singular, subdivisi, dan argumen berkas atau resolusi halus yang berada
di luar jembatan ini. Isi konseptualnya ialah bahwa perhitungan dengan bentuk
diferensial menghasilkan invariant topologis yang sama dengan kohomologi
singular berkoefisien nyata.

\paragraph{Derajat.}
Jika $f\colon M\to N$ mulus antara manifold tertutup, terhubung, berorientasi,
dan berdimensi sama $n$, tindakan $f^*$ pada kohomologi teratas adalah
perkalian oleh suatu bilangan bulat $\deg f$. Secara ekuivalen,
\[
 \int_M f^*\omega=(\deg f)\int_N\omega
\]
untuk setiap bentuk volume tertutup $\omega$. Dengan memilih nilai regular,
derajat dapat dihitung sebagai jumlah bertanda titik-titik prapeta.

\paragraph{Transversalitas.}
Jika $f\colon M\to N$ transversal terhadap submanifold $Z\subseteq N$, maka
$f^{-1}(Z)$ adalah submanifold berdimensi
$\dim M-\operatorname{codim}Z$. Transversalitas memungkinkan kelas kohomologi,
bilangan perpotongan, dan derajat direpresentasikan oleh geometri prapeta yang
stabil di bawah perturbasi kecil. Modul ini hanya memakai pernyataan gerbang
tersebut; teorema aproksimasi transversal memerlukan pembahasan tersendiri.

\paragraph{Teori Morse.}
Fungsi Morse mempunyai titik kritis nondegenerat. Ketika nilai melewati titik
kritis berindeks $\lambda$, sublevel berubah melalui pemasangan pegangan
$\lambda$. Dari sini lahir kompleks Morse, ketaksamaan Morse, dan hubungan
antara jumlah titik kritis dengan kohomologi. Pembuktian memerlukan lema Morse,
aliran gradien, dan analisis kekompakan; tidak satu pun disamarkan sebagai
telah dibuktikan di sini.

\begin{bemerkung}\label{o011-br-rem-roadmap}
Jembatan ini menutup rantai logis yang diperlukan untuk kurikulum dasar:
bentuk dan $d$ $\to$ kohomologi $\to$ invariansi homotopi $\to$
Mayer--Vietoris dan contoh $\to$ teorema de Rham sebagai identifikasi global
$\to$ gerbang derajat, transversalitas, serta Morse. Teori lanjutan dapat
memulai tepat dari tiga gerbang terakhir tanpa mengulang bagian sebelumnya.
\end{bemerkung}
% O011-BR-A — latihan, petunjuk, solusi, dan uji penguasaan asli.
% Lisensi: CC BY-SA 4.0.
% Provenans: disusun oleh OpenAI Codex gpt-5.6-sol, Ultra, atas arahan pengguna.

\subsection{Latihan, petunjuk bertahap, dan solusi}

Latihan berikut memakai ID stabil yang tidak bergantung pada bahasa. Setiap
solusi adalah materi asli modul ini. Tidak ada solusi yang dinisbahkan kepada
sumber Brenner atau kepada teks pembanding.

\subsubsection*{Latihan R1 — Bentuk tertutup dan primitif (O011-BR-E01)}
\label{o011-br-e01}

Pada $\mathbb R^2$, tinjau
\[
 \omega=(2xy+1)\,dx+(x^2+3y^2)\,dy.
\]
Tentukan apakah $\omega$ tertutup dan, jika eksak, temukan suatu primitif.

\paragraph{Petunjuk 1.}
Untuk $P\,dx+Q\,dy$, hitung $d\omega=(Q_x-P_y)\,dx\wedge dy$.

\paragraph{Petunjuk 2.}
Integralkan $P$ terhadap $x$, lalu cocokkan turunan terhadap $y$.

\paragraph{Solusi.}
$P_y=2x$ dan $Q_x=2x$, sehingga $d\omega=0$. Jika $df=\omega$, integrasi
$f_x=2xy+1$ memberi
$f=x^2y+x+c(y)$. Persamaan $f_y=x^2+c'(y)=x^2+3y^2$ memberi
$c(y)=y^3+C$. Jadi
\[
 \omega=d(x^2y+x+y^3).
\]

\subsubsection*{Latihan R2 — Produk pada kelas kohomologi (O011-BR-E02)}
\label{o011-br-e02}

Misalkan $\alpha,\alpha'$ bentuk-$p$ tertutup dengan
$\alpha'-\alpha=d\mu$, dan $\beta,\beta'$ bentuk-$q$ tertutup dengan
$\beta'-\beta=d\nu$. Buktikan secara langsung bahwa
$\alpha'\wedge\beta'-\alpha\wedge\beta$ eksak.

\paragraph{Petunjuk 1.}
Ganti kedua wakil satu per satu.

\paragraph{Petunjuk 2.}
Gunakan $d(\mu\wedge\beta)=d\mu\wedge\beta$ dan
$d(\alpha'\wedge\nu)=(-1)^p\alpha'\wedge d\nu$.

\paragraph{Solusi.}
Kita tulis
\[
 \alpha'\wedge\beta'-\alpha\wedge\beta
 =d\mu\wedge\beta+\alpha'\wedge d\nu.
\]
Karena $d\beta=0$ dan $d\alpha'=0$,
\[
 d\mu\wedge\beta=d(\mu\wedge\beta),
 \qquad
 \alpha'\wedge d\nu=(-1)^p d(\alpha'\wedge\nu).
\]
Jadi selisihnya adalah
$d(\mu\wedge\beta+(-1)^p\alpha'\wedge\nu)$.

\subsubsection*{Latihan R3 — Funktorialitas konkret (O011-BR-E03)}
\label{o011-br-e03}

Untuk $f\colon\mathbb R^2\to\mathbb R^2$,
$f(u,v)=(u^2-v^2,2uv)$, dan bentuk $\omega=x\,dy-y\,dx$, hitung
$f^*\omega$ dan verifikasi $d(f^*\omega)=f^*(d\omega)$.

\paragraph{Petunjuk 1.}
Substitusikan $x=u^2-v^2$, $y=2uv$ beserta diferensialnya.

\paragraph{Petunjuk 2.}
$d\omega=2\,dx\wedge dy$ dan determinan Jacobian $f$ adalah
$4(u^2+v^2)$.

\paragraph{Solusi.}
Kita mempunyai
$dx=2u\,du-2v\,dv$ dan $dy=2v\,du+2u\,dv$. Maka
\[
 f^*\omega
 =(u^2-v^2)(2v\,du+2u\,dv)
 -2uv(2u\,du-2v\,dv)
 =-2v(u^2+v^2)\,du+2u(u^2+v^2)\,dv.
\]
Turunan eksteriornya adalah
$8(u^2+v^2)\,du\wedge dv$. Di sisi lain,
\[
 f^*(d\omega)=2f^*(dx\wedge dy)
 =2\det(Df)\,du\wedge dv
 =8(u^2+v^2)\,du\wedge dv.
\]

\subsubsection*{Latihan R4 — Kohomologi derajat nol (O011-BR-E04)}
\label{o011-br-e04}

Jika manifold $M$ mempunyai tepat $r<\infty$ komponen terhubung, buktikan
$H^0_{\mathrm{dR}}(M)\cong\mathbb R^r$. Apa yang berubah jika banyaknya
komponen tak hingga?

\paragraph{Petunjuk 1.}
Suatu fungsi dengan $df=0$ ditentukan oleh satu konstanta pada setiap
komponen.

\paragraph{Petunjuk 2.}
Untuk tak hingga banyak komponen, tidak ada syarat bahwa keluarga konstanta
harus berhingga dukungannya.

\paragraph{Solusi.}
Pemetaan yang mengirim fungsi lokal-konstan ke daftar nilainya pada
komponen-komponen adalah isomorfisme linear ke $\mathbb R^r$. Tidak ada bentuk
derajat $-1$, jadi tidak ada hasil bagi oleh bentuk eksak. Jika himpunan
komponen adalah $I$ tak hingga, ruangnya adalah produk semua keluarga
$\mathbb R^I$, bukan jumlah langsung $\mathbb R^{(I)}$.

\subsubsection*{Latihan R5 — Generator bidang berlubang (O011-BR-E05)}
\label{o011-br-e05}

Pada $\mathbb R^2\setminus\{0\}$, buktikan bahwa
\[
 \eta=\frac{1}{2\pi}\frac{x\,dy-y\,dx}{x^2+y^2}
\]
tertutup tetapi tidak eksak.

\paragraph{Petunjuk 1.}
Hitung langsung $d\eta$, atau gunakan koordinat polar lokal.

\paragraph{Petunjuk 2.}
Tarik balik sepanjang $c(t)=(\cos t,\sin t)$ dan integralkan dari $0$ sampai
$2\pi$.

\paragraph{Solusi.}
Perhitungan aturan hasil bagi memberi $d\eta=0$ di luar asal. Sepanjang
lingkaran satuan,
\[
 c^*\eta=\frac{1}{2\pi}\,dt,
 \qquad \int_{S^1}\eta=1.
\]
Jika $\eta=df$, integralnya pada kurva tertutup akan nol berdasarkan teorema
dasar kalkulus atau Stokes. Kontradiksi ini menunjukkan $\eta$ tidak eksak.

\subsubsection*{Latihan R6 — Operator Poincar\'e eksplisit (O011-BR-E06)}
\label{o011-br-e06}

Gunakan operator homotopi radial untuk bentuk
$\omega=dx\wedge dy$ pada $\mathbb R^2$. Hitung $K\omega$ dan verifikasi
$dK\omega=\omega$.

\paragraph{Petunjuk 1.}
Untuk $x=(x,y)$ dan vektor $v$, integran ialah
$t\,\omega_x((x,y),v)$.

\paragraph{Petunjuk 2.}
Kontraksi vektor radial dengan $dx\wedge dy$ adalah $x\,dy-y\,dx$.

\paragraph{Solusi.}
Karena bentuknya konstan dan berderajat dua,
\[
 K\omega=\int_0^1t(x\,dy-y\,dx)\,dt
 =\frac12(x\,dy-y\,dx).
\]
Dengan demikian
\[
 dK\omega=\frac12(dx\wedge dy-dy\wedge dx)=dx\wedge dy=\omega.
\]

\subsubsection*{Latihan R7 — Pemetaan homotopik (O011-BR-E07)}
\label{o011-br-e07}

Misalkan $F\colon M\times[0,1]\to N$ homotopi dari $f_0$ ke $f_1$ dan
$\omega$ bentuk tertutup pada $N$. Tunjukkan secara eksplisit bahwa
$f_1^*\omega-f_0^*\omega$ eksak. Apa primitifnya?

\paragraph{Petunjuk 1.}
Gunakan rumus homotopi rantai.

\paragraph{Petunjuk 2.}
Suku $K_F(d\omega)$ lenyap.

\paragraph{Solusi.}
Rumus $dK_F+K_Fd=f_1^*-f_0^*$ memberi
\[
 f_1^*\omega-f_0^*\omega=d(K_F\omega)+K_F(d\omega)
 =d(K_F\omega).
\]
Jadi primitif eksplisitnya adalah
\[
 K_F\omega=\int_0^1\iota_t^*(\iota_{\partial_t}F^*\omega)\,dt.
\]

\subsubsection*{Latihan R8 — Retraksi deformasi (O011-BR-E08)}
\label{o011-br-e08}

Buktikan bahwa penyertaan $S^{n-1}\hookrightarrow\mathbb R^n\setminus\{0\}$
menginduksi isomorfisme kohomologi de Rham.

\paragraph{Petunjuk 1.}
Gunakan retraksi radial $r(x)=x/\|x\|$.

\paragraph{Petunjuk 2.}
Bangun homotopi dari identitas ke $i\circ r$ tanpa melewati asal.

\paragraph{Solusi.}
Penyertaan $i$ memenuhi $r\circ i=\operatorname{id}_{S^{n-1}}$. Homotopi
\[
 H(x,t)=\left((1-t)+\frac{t}{\|x\|}\right)x
\]
menghubungkan identitas pada $\mathbb R^n\setminus\{0\}$ dengan $i\circ r$;
koefisiennya selalu positif. Jadi $i$ dan $r$ adalah ekuivalensi homotopi.
Invariansi homotopi memberi isomorfisme yang diminta.

\subsubsection*{Latihan R9 — Mayer--Vietoris untuk lingkaran (O011-BR-E09)}
\label{o011-br-e09}

Tutupi $S^1$ oleh dua busur terbuka kontraktibel yang irisannya mempunyai dua
komponen. Tuliskan bagian barisan Mayer--Vietoris yang menghitung
$H^1_{\mathrm{dR}}(S^1)$ dan tentukan pemetaan linearnya.

\paragraph{Petunjuk 1.}
$H^0(U)=H^0(V)=\mathbb R$ dan
$H^0(U\cap V)=\mathbb R^2$.

\paragraph{Petunjuk 2.}
Pemetaan beda restriksi mengirim $(a,b)$ ke $(a-b,a-b)$.

\paragraph{Solusi.}
Bagian barisan yang relevan adalah
\[
0\to\mathbb R\to\mathbb R^2
\xrightarrow{D}\mathbb R^2
\to H^1(S^1)\to0,
\qquad D(a,b)=(a-b,a-b).
\]
Citra $D$ adalah diagonal satu dimensi, sehingga cokernelnya satu dimensi.
Karena itu $H^1_{\mathrm{dR}}(S^1)\cong\mathbb R$.

\subsubsection*{Latihan R10 — Derajat peta pangkat (O011-BR-E10)}
\label{o011-br-e10}

Untuk $f_m\colon S^1\to S^1$, $f_m(e^{i\theta})=e^{im\theta}$ dengan
$m\in\mathbb Z$, tentukan tindakan $f_m^*$ pada $H^1_{\mathrm{dR}}(S^1)$ dan
derajat $f_m$.

\paragraph{Petunjuk 1.}
Gunakan generator $d\theta/(2\pi)$ pada koordinat sudut lokal dan integral
globalnya.

\paragraph{Petunjuk 2.}
Tarikan baliknya adalah $m\,d\theta/(2\pi)$.

\paragraph{Solusi.}
Jika $[\eta]$ adalah generator ternormalisasi dengan integral $1$, maka
\[
 f_m^*[\eta]=m[\eta].
\]
Tindakan pada kohomologi teratas lingkaran adalah perkalian $m$, sehingga
$\deg(f_m)=m$. Ini juga benar untuk $m<0$, ketika orientasi dibalik.

\subsubsection*{Latihan R11 — Dua kelas pada torus (O011-BR-E11)}
\label{o011-br-e11}

Pada $\mathbb T^2=S^1\times S^1$, misalkan $\eta_1$ dan $\eta_2$ adalah
tarikan balik generator $H^1(S^1)$ melalui dua proyeksi. Buktikan bahwa
$[\eta_1]$ dan $[\eta_2]$ bebas linear.

\paragraph{Petunjuk 1.}
Integralkan pada dua lingkaran koordinat
$c_1(z)=(z,1)$ dan $c_2(z)=(1,z)$.

\paragraph{Petunjuk 2.}
Matriks periodenya adalah matriks identitas.

\paragraph{Solusi.}
Dengan normalisasi integral generator sama dengan $1$,
\[
 \int_{c_1}\eta_1=1,\quad \int_{c_1}\eta_2=0,
 \qquad
 \int_{c_2}\eta_1=0,\quad \int_{c_2}\eta_2=1.
\]
Jika $a[\eta_1]+b[\eta_2]=0$, bentuk terkait eksak dan semua periodenya nol.
Integrasi pada $c_1$ memberi $a=0$, dan pada $c_2$ memberi $b=0$.

\subsubsection*{Latihan R12 — Bentuk eksak pada manifold tertutup (O011-BR-E12)}
\label{o011-br-e12}

Misalkan $M$ manifold kompak berorientasi tanpa batas berdimensi $n$. Buktikan
bahwa untuk setiap $\eta\in\Omega^{n-1}(M)$,
\[
 \int_M d\eta=0.
\]
Jelaskan mengapa suatu bentuk-$n$ dengan integral tak nol tidak mungkin eksak.

\paragraph{Petunjuk 1.}
Gunakan teorema Stokes dan $\partial M=\varnothing$.

\paragraph{Petunjuk 2.}
Argumen kedua adalah kontraposisi langsung.

\paragraph{Solusi.}
Teorema Stokes memberi
\[
 \int_Md\eta=\int_{\partial M}\eta=0.
\]
Jika bentuk-$n$ $\omega$ eksak, $\omega=d\eta$ dan integralnya harus nol.
Jadi integral tak nol mendeteksi kelas kohomologi de Rham tak nol.

\subsection{Uji penguasaan kumulatif}

\subsubsection*{Penguasaan N1 — Bidang berlubang (O011-BR-M01)}
\label{o011-br-m01}

(a) Bangun retraksi deformasi $\mathbb R^2\setminus\{0\}\to S^1$; (b) hitung
semua grup kohomologi de Rham bidang berlubang; (c) buktikan bahwa bentuk
$\eta$ dari Latihan R5 mewakili generator ternormalisasi; (d) tentukan apakah
$2\eta+d(x/(x^2+y^2))$ eksak.

\paragraph{Petunjuk.}
Gabungkan Latihan R5 dan R8. Integral bentuk eksak pada lingkaran tertutup
adalah nol.

\paragraph{Solusi.}
Retraksi dan homotopinya adalah
\[
 r(x)=\frac{x}{\|x\|},
 \qquad H(x,t)=\left((1-t)+\frac t{\|x\|}\right)x.
\]
Maka kohomologi sama dengan kohomologi $S^1$:
\[
 H^0\cong\mathbb R,\qquad H^1\cong\mathbb R,
 \qquad H^k=0\ (k\ge2).
\]
Karena $\eta$ tertutup dan periodenya pada lingkaran positif adalah $1$, ia
mewakili generator ternormalisasi. Bentuk terakhir mempunyai periode $2$,
sebab suku diferensial eksak berperiode nol; jadi bentuk itu tidak eksak.

\paragraph{Rubrik (12 poin).}
Tiga poin untuk retraksi dan homotopi yang menghindari asal; tiga poin untuk
semua derajat kohomologi; tiga poin untuk ketertutupan/periode/generator;
tiga poin untuk keputusan terakhir beserta alasan periodenya.

\paragraph{Parameter alternatif.}
Ganti lubang asal dengan titik $p\in\mathbb R^2$ dan gunakan koordinat
$x-p$. Generator dan retraksi harus ditranslasikan, sementara struktur
kohomologi tetap sama.

\subsubsection*{Penguasaan N2 — Lema Poincar\'e terhitung (O011-BR-M02)}
\label{o011-br-m02}

Pada $\mathbb R^3$, misalkan
\[
 \omega=x\,dy\wedge dz+y\,dz\wedge dx+z\,dx\wedge dy.
\]
(a) Hitung $d\omega$; (b) tentukan apakah $\omega$ tertutup; (c) untuk bentuk
$\alpha=dx\wedge dy$, hitung $K\alpha$ dengan operator radial; (d) verifikasi
identitas homotopi pada $\alpha$.

\paragraph{Petunjuk.}
Tiga suku pada $d\omega$ semuanya merupakan $dx\wedge dy\wedge dz$. Untuk
$\alpha$, gunakan perhitungan Latihan R6.

\paragraph{Solusi.}
Kita memperoleh
\[
 d\omega=3\,dx\wedge dy\wedge dz,
\]
sehingga $\omega$ tidak tertutup. Untuk $\alpha$,
\[
 K\alpha=\frac12(x\,dy-y\,dx),
 \qquad dK\alpha=dx\wedge dy=\alpha.
\]
Karena $d\alpha=0$ dan tarikan balik bentuk derajat dua oleh pemetaan konstan
nol, identitas $dK\alpha+Kd\alpha=\alpha-F_0^*\alpha$ tepat menjadi
$dK\alpha=\alpha$.

\paragraph{Rubrik (12 poin).}
Tiga poin untuk tanda dan koefisien $d\omega$; dua poin untuk keputusan
ketertutupan; empat poin untuk operator radial; tiga poin untuk verifikasi
identitas lengkap termasuk suku pemetaan konstan.

\paragraph{Parameter alternatif.}
Ganti $\omega$ dengan
$a x\,dy\wedge dz+b y\,dz\wedge dx+c z\,dx\wedge dy$.
Ketertutupan terjadi tepat ketika $a+b+c=0$.

\subsubsection*{Penguasaan N3 — Kohomologi bola melalui Mayer--Vietoris (O011-BR-M03)}
\label{o011-br-m03}

Dengan penutup dua lingkungan kutub, buktikan secara induktif rumus lengkap
untuk $H^k_{\mathrm{dR}}(S^n)$, mulai dari $S^1$. Jelaskan secara terpisah apa
yang terjadi pada derajat nol, satu, antara, dan teratas.

\paragraph{Petunjuk.}
Kedua bagian kontraktibel dan irisannya ekuivalen-homotopi dengan
$S^{n-1}$. Jangan memakai isomorfisme penghubung pada derajat nol tanpa
memeriksa keterhubungan.

\paragraph{Solusi.}
Kasus dasar $S^1$ diperoleh dari barisan
$0\to\mathbb R\to\mathbb R^2\to\mathbb R^2\to H^1(S^1)\to0$,
yang memberi $H^0=H^1=\mathbb R$. Untuk $n\ge2$, pada derajat $k>1$ bagian
kontraktibel menyumbang nol dan Mayer--Vietoris memberi
\[
 H^k(S^n)\cong H^{k-1}(S^{n-1}).
\]
Pada derajat satu, keterhubungan kedua bagian dan irisan memberi $H^1(S^n)=0$.
Derajat nol adalah $\mathbb R$ karena $S^n$ terhubung. Induksi lalu memberi
\[
 H^k(S^n)=\begin{cases}\mathbb R,&k=0,n,\\0,&\text{lainnya}.
 \end{cases}
\]

\paragraph{Rubrik (12 poin).}
Tiga poin untuk kasus dasar; tiga poin untuk penutup dan tipe homotopi irisan;
empat poin untuk penggunaan barisan eksak pada rentang derajat yang benar;
dua poin untuk kesimpulan induktif lengkap.

\paragraph{Parameter alternatif.}
Mulai induksi dari $S^0$ sebagai dua titik dan minta peserta menelusuri
derajat nol secara eksplisit sebelum melanjutkan ke $S^1$.

\subsubsection*{Penguasaan N4 — Derajat, nilai regular, dan batas cakupan (O011-BR-M04)}
\label{o011-br-m04}

Untuk $f_m(e^{i\theta})=e^{im\theta}$ dengan $m\ne0$: (a) hitung derajat dari
tarikan balik bentuk generator; (b) pilih nilai regular $1\in S^1$ dan hitung
jumlah bertanda prapetanya; (c) cocokkan kedua jawaban; (d) jelaskan dengan
satu kalimat masing-masing peran transversalitas dan teori Morse dalam jalur
lanjutan, tanpa mengklaim teorema yang belum dibuktikan.

\paragraph{Petunjuk.}
Prapeta $1$ adalah $e^{2\pi ik/m}$ dengan interpretasi yang sesuai untuk tanda
$m$. Turunan dalam koordinat sudut adalah perkalian $m$.

\paragraph{Solusi.}
Untuk generator $\eta$ dengan integral $1$,
$f_m^*\eta=m\eta$, sehingga $\deg f_m=m$. Nilai $1$ regular karena turunan
lokalnya $m\ne0$. Ada $|m|$ titik prapeta; masing-masing bertanda
$\operatorname{sgn}(m)$, sehingga jumlah bertandanya $m$, sama dengan hasil
kohomologi.

Transversalitas menjamin prapeta submanifold yang ditemui secara generik
merupakan submanifold dengan dimensi yang diharapkan dan mendasari hitungan
perpotongan. Teori Morse menguraikan perubahan topologi sublevel melalui titik
kritis nondegenerat dan menghubungkannya dengan kohomologi; pembuktian kedua
teori itu berada di luar modul ini.

\paragraph{Rubrik (12 poin).}
Tiga poin untuk tarikan balik dan derajat; empat poin untuk prapeta, regularitas,
dan tanda lokal; dua poin untuk pencocokan; tiga poin untuk dua penanda jalur
yang benar dan tidak melebih-lebihkan cakupan.

\paragraph{Parameter alternatif.}
Gunakan $g_m(e^{i\theta})=e^{i(m\theta+\theta_0)}$ dengan konstanta
$\theta_0$. Derajat tetap $m$, tetapi himpunan prapeta nilai regular bergeser.

\part{Formulir Ujian}
\chapter*{Provenans formulir ujian}
\addcontentsline{toc}{chapter}{Provenans formulir ujian}
Sepuluh formulir peserta berikut mempertahankan teks peserta resmi. Sepuluh formulir solusi resmi sesudahnya hanya memuat solusi yang tersedia pada sumber resmi yang dibekukan. Kekosongan sumber tetap kosong dan tidak diisi pada permukaan resmi ini.
% Lokalisasi khusus permukaan pembaca untuk formulir ujian dalam edisi lengkap.
% Nama perintah sumber dipertahankan; hanya teks yang dicetak yang dilokalkan.

\renewcommand{\fachbereich}{Bidang studi}
\renewcommand{\dozent}{Pengajar}
\renewcommand{\klausurdatum}{Tanggal}
\renewcommand{\klausurgebiet}{Matematika}
\renewcommand{\klausurtyp}{Ujian}

\renewcommand{\klausurmodus}{%
Durasi: sepuluh menit untuk membaca dan dua jam penuh untuk mengerjakan.
Alat bantu tidak diperbolehkan. Semua jawaban harus diberi alasan.

Jumlah poin seluruhnya $64$. Berlaku aturan ambang dasar, yakni penilaian
untuk setiap soal atau bagian soal pada umumnya dimulai pada separuh jumlah
poin. Untuk lulus diperlukan $16$ poin; mulai $32$ poin diperoleh nilai satu.

Tuliskan nama dan nomor mahasiswa Anda dengan jelas pada lembar sampul dan
setiap lembar berikutnya. Semoga berhasil!}

\renewcommand{\testklausurmodus}{\klausurmodus}

\renewcommand{\klausurname}{%
\renewcommand{\arraystretch}{1.5}%
\begin{tabular}{@{}p{3.2cm}p{10cm}}
Nama lengkap: &
..................................................................................\\
Nomor mahasiswa: &
..................................................................................\\
\end{tabular}}

\renewcommand{\klausurvorspann}[5]{%
\begin{tabular}{@{}p{3.2cm}p{10cm}}
Bidang studi: & #1\\
Tanggal: & #2\\
Pengajar: & #3
\end{tabular}
\vspace{3ex}

\begin{center}
{\bfseries\large Ujian\\#4\ #5}
\end{center}

\vspace{2ex}
\klausurmodus
\vspace{2ex}
\klausurname\par}

\renewcommand{\testklausurvorspann}[5]{\klausurvorspann{#1}{#2}{#3}{#4}{#5}}
\renewcommand{\klausurmitloesungenvorspann}[5]{%
\klausurvorspann{#1}{#2}{#3}{#4}{#5}
\begin{center}\bfseries Formulir dengan solusi\end{center}}

\renewcommand{\klausurnote}{%
\begin{tabular}{@{}p{3.2cm}p{10cm}}Nilai: & \end{tabular}}

\renewcommand{\klausurtaeuschungwarnung}{%
\vspace{2ex}%
Saya memahami bahwa pelanggaran berat terhadap aturan ujian dapat membatalkan
seluruh komponen penilaian dalam mata kuliah ini. Penggunaan literatur atau
alat bantu elektronik termasuk pelanggaran berat.
\vspace{1ex}}

\renewcommand{\klausureinverstaendnis}{%
\vspace{2ex}
Dengan tanda tangan saya, saya menyetujui bahwa hasil ujian dapat diumumkan
bersama nomor mahasiswa saya.
\vspace{1ex}

\begin{tabular}{@{}p{3.2cm}p{10cm}}
Tanda tangan: &
.....................................................................................
\end{tabular}}

\renewcommand{\inputaufgabeklausurloesung}[3]{%
\par\newpage
\begin{aufgabe}
{\ifthenelse{\equal{#1}{}}{}{\ifthenelse{\equal{#1}{1}}{(1 poin)}{(#1 poin)}}}%
\par\bigskip #2
\end{aufgabe}
\underline{Solusi:}\par\bigskip #3}

\renewcommand{\inputaufgabeloesung}[2]{%
\begin{aufgabe}#1\end{aufgabe}\aufgabenachskip
\bigskip
Solusi
\bigskip #2}

\renewcommand{\inputaufgabeloesungvar}[2]{%
\aufgabevorskip Soal: #1\aufgabenachskip Solusi: #2}

\renewcommand{\inputaufgabepunkteloesung}[3]{%
\aufgabevorskip
\begin{aufgabe}
{\ifthenelse{\equal{#1}{}}{}{\ifthenelse{\equal{#1}{1}}{(1 poin)}{(#1 poin)}}}%
\aufgabepunktskip #2
\end{aufgabe}
\aufgabenachskip\loesung{#3}}

\renewcommand{\inputaufgabepunktenummervonhand}[3]{%
\aufgabevorskip{\bfseries Soal #1} (#2 poin)\aufgabepunktskip #3\aufgabenachskip}

\renewcommand{\loesung}[1]{\underline{Solusi:} #1}

% Delapan formulir sumber memakai huruf penampung K sebagai nilai penghitung
% bagian. Pertahankan byte formulir dan tafsirkan penampung itu sebagai awal
% penomoran hanya pada lapisan kompilasi pembaca.
\let\OExamOriginalSetCounter\setcounter
\renewcommand{\setcounter}[2]{%
\ifthenelse{\equal{#1}{bagian}}%
  {\OExamOriginalSetCounter{section}{#2}}%
  {\ifthenelse{\equal{#1}{section}}%
  {\ifthenelse{\equal{#2}{K}}{\OExamOriginalSetCounter{section}{0}}{\OExamOriginalSetCounter{#1}{#2}}}%
  {\OExamOriginalSetCounter{#1}{#2}}}}

\newcommand{\OExamPointTable}[4]{%
\vspace{3ex}
\begin{center}
\begingroup
\scriptsize
\renewcommand{\arraystretch}{1.45}%
\setlength{\tabcolsep}{0.35ex}%
\begin{tabular}{@{}|l|*{#1}{c|}c|}
\hline
Soal & #2 & Jumlah\\ \hline
Maks. & #3 & #4\\ \hline
Diperoleh & \multicolumn{#1}{c|}{} & \\ \hline
\end{tabular}
\endgroup
\end{center}
\vspace{2ex}}

\renewcommand{\punktetabelleelf}{%
\OExamPointTable{11}%
{1&2&3&4&5&6&7&8&9&10&11}%
{\aeins&\azwei&\adrei&\avier&\afuenf&\asechs&\asieben&\aacht&\aneun&\azehn&\aelf}%
{\azwoelf}}

\renewcommand{\punktetabelledreizehn}{%
\OExamPointTable{13}%
{1&2&3&4&5&6&7&8&9&10&11&12&13}%
{\aeins&\azwei&\adrei&\avier&\afuenf&\asechs&\asieben&\aacht&\aneun&\azehn&\aelf&\azwoelf&\adreizehn}%
{\avierzehn}}

\renewcommand{\punktetabellevierzehn}{%
\OExamPointTable{14}%
{1&2&3&4&5&6&7&8&9&10&11&12&13&14}%
{\aeins&\azwei&\adrei&\avier&\afuenf&\asechs&\asieben&\aacht&\aneun&\azehn&\aelf&\azwoelf&\adreizehn&\avierzehn}%
{\afuenfzehn}}

\renewcommand{\punktetabellesechzehn}{%
\OExamPointTable{16}%
{1&2&3&4&5&6&7&8&9&10&11&12&13&14&15&16}%
{\aeins&\azwei&\adrei&\avier&\afuenf&\asechs&\asieben&\aacht&\aneun&\azehn&\aelf&\azwoelf&\adreizehn&\avierzehn&\afuenfzehn&\asechzehn}%
{\asiebzehn}}

\renewcommand{\punktetabellesiebzehn}{%
\OExamPointTable{17}%
{1&2&3&4&5&6&7&8&9&10&11&12&13&14&15&16&17}%
{\aeins&\azwei&\adrei&\avier&\afuenf&\asechs&\asieben&\aacht&\aneun&\azehn&\aelf&\azwoelf&\adreizehn&\avierzehn&\afuenfzehn&\asechzehn&\asiebzehn}%
{\aachtzehn}}
\providecommand{\theHfakt}{}

\chapter{Formulir Peserta Ujian 1}
\noindent\textit{Formulir peserta dari sumber resmi; teks soal dipertahankan sesuai terjemahan yang telah diverifikasi.}
\begingroup
\renewcommand{\theHfakt}{o011.exam.learner.01.\arabic{fakt}}
%Data institusi

%\input{Dozentdaten}

%\renewcommand{\fachbereich}{Bidang studi}

%\renewcommand{\dozent}{Prof. Dr. . }

%Data ujian

\renewcommand{\klausurgebiet}{ }

\renewcommand{\klausurtyp}{ }

\renewcommand{\klausurdatum}{. 20}

\klausurvorspann {\fachbereich} {\klausurdatum} {\dozent} {\klausurgebiet} {\klausurtyp}

%Data untuk tabel poin berikut

\renewcommand{\aeins}{  3  }

\renewcommand{\azwei}{  3   }

\renewcommand{\adrei}{  6  }

\renewcommand{\avier}{  5  }

\renewcommand{\afuenf}{  5  }

\renewcommand{\asechs}{  2  }

\renewcommand{\asieben}{  5  }

\renewcommand{\aacht}{  5  }

\renewcommand{\aneun}{  7  }

\renewcommand{\azehn}{  3  }

\renewcommand{\aelf}{  2  }

\renewcommand{\azwoelf}{  5  }

\renewcommand{\adreizehn}{  9   }

\renewcommand{\avierzehn}{  5  }

\renewcommand{\afuenfzehn}{  65  }

\renewcommand{\asechzehn}{    }

\renewcommand{\asiebzehn}{    }

\renewcommand{\aachtzehn}{    }

\renewcommand{\aneunzehn}{    }

\renewcommand{\azwanzig}{    }

\renewcommand{\aeinundzwanzig}{    }

\renewcommand{\azweiundzwanzig}{    }

\renewcommand{\adreiundzwanzig}{    }

\renewcommand{\avierundzwanzig}{    }

\renewcommand{\afuenfundzwanzig}{    }

\renewcommand{\asechsundzwanzig}{    }

\punktetabellevierzehn

\klausurnote

\newpage

\setcounter{section}{0}

\inputaufgabegibtloesung
{3}
{

Definisikan istilah-istilah berikut
\zusatzklammer {yang dicetak miring} {} {}.
\aufzaehlungsechs{\stichwort {Lingkaran kelengkungan} {}
dari suatu
\definitionsverweis {kurva}{}{}
yang dua kali
\definitionsverweis {terdiferensialkan secara kontinu}{}{}
dan
\definitionsverweis {berparametrisasi panjang busur}{}{}.
Misalkan
\maabbdisp {\gamma} {I} {\R^2
} {}
dan ambil titik

\mavergleichskette
{\vergleichskette
{ t_0
}
{ \in }{ I
}
{  }{
}
{    }{
}
{   }{
}
}
{}{}{.}

}{Suatu \stichwort {submanifold tertutup} {}
\mathl{M \subseteq N}{} dari manifold diferensiabel $N$.

}{\stichwort {Ruang singgung} {} pada titik
\mathl{P\in M}{} dari manifold diferensiabel $M$.

}{Istilah
\stichwort {berorientasi} {}
untuk suatu
\definitionsverweis {manifold diferensiabel}{}{}
$M$.

}{\stichwort {Batas} {}
suatu
\definitionsverweis {manifold berbatas}{}{.}

}{\stichwort {Simbol Christoffel} {}

\mathl{\Gamma_{ij}^k}{} bagi koneksi Levi-Civita dari struktur Riemann $\left(  g_{ij}  \right)$ pada himpunan terbuka

\mavergleichskette
{\vergleichskette
{ U
}
{ \subseteq }{ \R^n
}
{  }{
}
{    }{
}
{   }{
}
}
{}{}{.}
}

}
{} {}

\inputaufgabegibtloesung
{3}
{

Nyatakan teorema-teorema berikut.
\aufzaehlungdrei{Teorema tentang citra
\definitionsverweis {pemetaan Gauss}{}{}
dari suatu hipermuka.}{Rumus luas permukaan dari permukaan putar yang dibentuk oleh
\definitionsverweis {kurva diferensiabel}{}{}
\maabbeledisp {\gamma} {]a,b[} {\R^2
} { t } {(x(t),y(t))
} {,}
dengan

\mavergleichskette
{\vergleichskette
{  y(t)
}
{ > }{  0
}
{  }{
}
{    }{
}
{   }{
}
}
{}{}{.}}{\stichwort {Aturan hasil kali} {}
untuk turunan eksterior.}

}
{} {}

\inputaufgabegibtloesung
{6 (1+2+3)}
{

Misalkan $Y$ adalah hipermuka yang diberikan oleh persamaan

\mavergleichskettedisp
{\vergleichskette
{ 2x^2 +3y^2+5z^2
}
{ =} { 10
}
{ } {
}
{ } {
}
{ } {
}
}
{}{}{}
di $\R^3$. Misalkan

\mavergleichskette
{\vergleichskette
{ P
}
{ = }{ \begin{pmatrix}  -1 \\ -1\\  1   \end{pmatrix}
}
{ \in }{ Y
}
{    }{
}
{   }{
}
}
{}{}{}
dan

\mavergleichskette
{\vergleichskette
{ w
}
{ = }{ \begin{pmatrix}  -3 \\ 2 \\  4   \end{pmatrix}
}
{ \in }{ \R^3
}
{    }{
}
{   }{
}
}
{}{}{.}
\aufzaehlungdrei{Tentukan
\definitionsverweis {ruang singgung}{}{}

\mathl{T_PY}{.}
}{Tentukan
\definitionsverweis {dekomposisi ortogonal}{}{}
$w$ menjadi komponen tangensial dan komponen ortogonal pada titik $P$.
}{Realisasikan komponen tangensial $w$ di $T_PY$ melalui suatu kurva diferensiabel yang seluruhnya terletak pada $Y$.
}

}
{} {}

\inputaufgabegibtloesung
{5}
{

Misalkan
\maabbeledisp {\gamma} {\R} { \R^2
} { t } { \begin{pmatrix}  \cos \left(  \omega t \right)  \\ \sin  \left(  \omega t \right)    \end{pmatrix}
} {,}
dengan

\mavergleichskette
{\vergleichskette
{ \omega
}
{ > }{ 0
}
{  }{
}
{    }{
}
{   }{
}
}
{}{}{.}
Tunjukkan bahwa tidak ada gerak melingkar lain dengan norma kecepatan konstan yang memiliki nilai dan turunan hingga orde kedua yang sama dengan $\gamma$ di titik parameter

\mavergleichskette
{\vergleichskette
{ t
}
{ = }{ 0
}
{  }{
}
{    }{
}
{   }{
}
}
{}{}{}
tersebut.

}
{} {}

\inputaufgabegibtloesung
{5}
{

Buktikan teorema tentang kelengkungan suatu kurva bidang yang diberikan secara implisit.

}
{} {}

\inputaufgabegibtloesung
{2}
{

Berikan contoh manifold diferensiabel
\mathkor {} {M} {dan} {N} {}
dengan

\mavergleichskette
{\vergleichskette
{ \operatorname{ dim}_{  } ^{  }  \left(   M   \right), \, \operatorname{ dim}_{  } ^{  }  \left(   N   \right)
}
{ \geq }{ 1
}
{  }{
}
{    }{
}
{   }{
}
}
{}{}{}
serta pemetaan diferensiabel
\maabb {\varphi} { M } { N
} {}
dan
\maabb {\psi} { N } { M
} {}
sedemikian sehingga

\mavergleichskette
{\vergleichskette
{ \psi \circ \varphi
}
{ = }{
\operatorname{Id}_{ M }
}
{  }{
}
{    }{
}
{   }{
}
}
{}{}{}
dan

\mavergleichskette
{\vergleichskette
{\varphi \circ  \psi
}
{ \neq }{
\operatorname{Id}_{ N }
}
{  }{
}
{    }{
}
{   }{
}
}
{}{}{}
berlaku.

}
{} {}

\inputaufgabegibtloesung
{5}
{

Buktikan bahwa ekuivalensi tangensial lintasan-lintasan pada suatu manifold di titik $P$ dapat diperiksa dengan menggunakan sebarang peta koordinat.

}
{} {}

\inputaufgabegibtloesung
{5}
{

Tunjukkan bahwa pemetaan singgung $T(\varphi)$ dari
\maabbeledisp {\varphi} {\R^1} {S^1
} { t } {( \cos  t, \sin  t  )
} {,}
bersifat surjektif.

}
{} {}

\inputaufgabegibtloesung
{7 (3+4)}
{

Tinjau pemetaan diferensiabel
\maabbeledisp {\gamma} { [1,c] } {\R^2
} { t } {(t,t^{3})
} {,}
\zusatzklammer {dengan \mathlk{c \geq 1}{}} {} {}
dan
\maabbeledisp {\varphi} {\R_+ \times \R_+ } {\R^3
} {(u,v)} {(u^3,u^2+v^2,u^{-1}v^{-1} )
} {,}
serta bentuk diferensial

\mavergleichskettedisp
{\vergleichskette
{ \omega
}
{ =} { (x-y)dx-z^2dy+dz
}
{ } {
}
{ } {
}
{ } {
}
}
{}{}{}
pada $\R^3$.

a) Hitung bentuk diferensial tarik balik
\mathl{\varphi^* \omega}{} pada $\R_+ \times \R_+$.

b) Hitung integral lintasan dari bentuk diferensial
\mathl{\varphi^* \omega}{} sepanjang lintasan $\gamma$ sebagai fungsi dari $c$.

}
{} {}

\inputaufgabegibtloesung
{3}
{

Misalkan
\maabbdisp {\psi} { L } { M
} {}
adalah
\definitionsverweis {pemetaan diferensiabel}{}{}
antara
\definitionsverweis {manifold diferensiabel}{}{.}
Selanjutnya, misalkan
\maabbdisp {f} { M } {\R
} {}
adalah fungsi pada $M$, dan $\omega$ adalah
$k$-\definitionsverweis {bentuk}{}{}
pada $M$. Tunjukkan bahwa

\mavergleichskettedisp
{\vergleichskette
{ \psi^*(f \omega)
}
{ =} { \psi^*(f)  \psi^* \omega
}
{ } {
}
{ } {
}
{ } {
}
}
{}{}{.}

}
{} {}

\inputaufgabe
{2}
{

Deskripsikan sebanyak mungkin
\definitionsverweis {manifold berbatas}{}{}
berdimensi satu yang
\definitionsverweis {terhubung}{}{.}

}
{} {}

\inputaufgabegibtloesung
{5}
{

Misalkan
\maabbdisp {f} { [a,b] } { \R_{+}
} {}
adalah
\definitionsverweis {fungsi yang terdiferensialkan secara kontinu}{}{.}
Kita pandang
\definitionsverweis {subgraf}{}{}
sebagai
\definitionsverweis {manifold berbatas}{}{}.
Batasnya terdiri dari grafik, interval dasar, dan kedua sisi tegak, tetapi tidak mencakup keempat titik sudut. Buktikan secara langsung berlakunya teorema Stokes untuk
\definitionsverweis {bentuk diferensial}{}{}

\mavergleichskette
{\vergleichskette
{ \omega
}
{ = }{ xdy
}
{  }{
}
{    }{
}
{   }{
}
}
{}{}{.}

}
{} {}

\inputaufgabegibtloesung
{9}
{

Buktikan teorema titik tetap Brouwer.

}
{} {}

\inputaufgabegibtloesung
{5 (3+1+1)}
{

Tinjau silinder

\mavergleichskette
{\vergleichskette
{ Z
}
{ = }{ S^1 \times \R
}
{ \subseteq }{ \R^2 \times \R
}
{    }{
}
{   }{
}
}
{}{}{}
dengan
\definitionsverweis {koneksi trivial}{}{}
pada
\definitionsverweis {bundel tangen}{}{.}
\aufzaehlungdrei{Tentukan, dalam suatu peta isometrik yang sesuai,
\definitionsverweis {persamaan diferensial geodesik}{}{}
untuk kurva geodesik $\psi$ pada $M$ yang memenuhi

\mavergleichskettedisp
{\vergleichskette
{ \psi(0)
}
{ =} { (1,0,3)
}
{ } {
}
{ } {
}
{ } {
}
}
{}{}{}
dan

\mavergleichskettedisp
{\vergleichskette
{ \psi'(0)
}
{ =} { (0,1,-4)
}
{ } {
}
{ } {
}
{ } {
}
}
{}{}{}
.
}{Selesaikan persamaan diferensial dari (1) dalam peta tersebut.
}{Carilah suatu kurva geodesik di $Z$ yang memenuhi syarat-syarat pada (1).
}

}
{} {}
\endgroup

\chapter{Formulir Peserta Ujian 2}
\noindent\textit{Formulir peserta dari sumber resmi; teks soal dipertahankan sesuai terjemahan yang telah diverifikasi.}
\begingroup
\renewcommand{\theHfakt}{o011.exam.learner.02.\arabic{fakt}}
%Data institusi

%\input{Dozentdaten}

%\renewcommand{\fachbereich}{Fachbereich}

%\renewcommand{\dozent}{Prof. Dr. . }

%Data ujian

\renewcommand{\klausurgebiet}{ }

\renewcommand{\klausurtyp}{ }

\renewcommand{\klausurdatum}{ . 20}

\klausurvorspann {\fachbereich} {\klausurdatum} {\dozent} {\klausurgebiet} {\klausurtyp}

%Data untuk tabel nilai berikut

\renewcommand{\aeins}{  3  }

\renewcommand{\azwei}{  3   }

\renewcommand{\adrei}{  8  }

\renewcommand{\avier}{  7  }

\renewcommand{\afuenf}{  8  }

\renewcommand{\asechs}{  3  }

\renewcommand{\asieben}{  5  }

\renewcommand{\aacht}{  5  }

\renewcommand{\aneun}{  6  }

\renewcommand{\azehn}{  4  }

\renewcommand{\aelf}{  12  }

\renewcommand{\azwoelf}{  64  }

\renewcommand{\adreizehn}{     }

\renewcommand{\avierzehn}{    }

\renewcommand{\afuenfzehn}{    }

\renewcommand{\asechzehn}{    }

\renewcommand{\asiebzehn}{    }

\renewcommand{\aachtzehn}{    }

\renewcommand{\aneunzehn}{    }

\renewcommand{\azwanzig}{    }

\renewcommand{\aeinundzwanzig}{    }

\renewcommand{\azweiundzwanzig}{    }

\renewcommand{\adreiundzwanzig}{    }

\renewcommand{\avierundzwanzig}{    }

\renewcommand{\afuenfundzwanzig}{    }

\renewcommand{\asechsundzwanzig}{    }

\punktetabelleelf

\klausurnote

\newpage

\setcounter{section}{0}

\inputaufgabegibtloesung
{3}
{

Definisikan istilah-istilah berikut
\zusatzklammer {yang dicetak miring} {} {.}
\aufzaehlungsechs{Suatu
\stichwort {medan normal} {}
$N$ pada hiperpermukaan diferensiabel

\mavergleichskette
{\vergleichskette
{ Y
}
{ \subseteq }{ \R^n
}
{  }{
}
{    }{
}
{   }{
}
}
{}{}{.}

}{Suatu ruang topologis yang
\stichwort {parakompak}{}
.

}{
\stichwort {Ruang kotangen} {}
di suatu titik

\mavergleichskette
{\vergleichskette
{  P
}
{ \in }{ M
}
{  }{}
{    }{}
{   }{}
}
{}{}{}
dari
\definitionsverweis {manifold diferensiabel}{}{}
$M$.

}{Suatu
\stichwort {titik reguler} {}

\mathl{Q \in L}{} bagi
\definitionsverweis {pemetaan diferensiabel}{}{}
\maabbdisp {\varphi} { L } { M
} {}
antara
\definitionsverweis {manifold-manifold diferensiabel}{}{}
\mathkor {} {L} {dan} {M} {.}

}{
\stichwort {Produk} {}
dari manifold-manifold diferensiabel
\mathkor {} {M} {dan} {N} {.}

}{Suatu
\stichwort {koneksi} {}
pada bundel vektor diferensiabel
\maabb {} {E} {M
} {}
di atas manifold diferensiabel $M$.
}

}
{} {}

\inputaufgabegibtloesung
{3}
{

Nyatakan teorema-teorema berikut.
\aufzaehlungdrei{
\stichwort {Teorema tentang hubungan antara kelengkungan dan vektor normal satuan suatu kurva bidang berparameter panjang busur} {.}}{
\stichwort {Theorema egregium} {.}}{
\stichwort {Teorema Green} {}
untuk luas.}

}
{} {}

\inputaufgabegibtloesung
{8}
{

Misalkan

\mavergleichskettedisp
{\vergleichskette
{ f(x,y)
}
{ =} { ax^2+by^2+cxy
}
{ } {
}
{ } {
}
{ } {
}
}
{}{}{.}
Tentukan
\definitionsverweis {kelengkungan-kelengkungan utama}{}{}
dan
\definitionsverweis {arah-arah kelengkungan utama}{}{}
dari
\definitionsverweis {grafik}{}{}
$f$ di titik asal.

}
{} {}

\inputaufgabegibtloesung
{7}
{

Buktikan teorema eksistensi untuk transpor paralel pada suatu hiperpermukaan

\mavergleichskette
{\vergleichskette
{ Y
}
{ \subseteq }{ \R^n
}
{  }{
}
{    }{
}
{   }{
}
}
{}{}{.}

}
{} {}

\inputaufgabegibtloesung
{8}
{

Tunjukkan bahwa suatu
\definitionsverweis {manifold diferensiabel}{}{}
$M$ berdimensi

\mavergleichskette
{\vergleichskette
{ n
}
{ \geq }{ 1
}
{  }{
}
{    }{
}
{   }{
}
}
{}{}{}
mempunyai tak hingga banyak
\definitionsverweis {difeomorfisme}{}{}
\maabbdisp {\varphi} { M } { M
} {}
.

}
{} {}

\inputaufgabegibtloesung
{3}
{

Misalkan $M$ suatu
\definitionsverweis {manifold diferensiabel}{}{}
dan
\maabbdisp {f} { M } {\R
} {}
suatu
\definitionsverweis {fungsi terdiferensialkan secara kontinu}{}{.}
Misalkan di titik

\mavergleichskette
{\vergleichskette
{ P
}
{ \in }{ M
}
{  }{
}
{    }{
}
{   }{
}
}
{}{}{}
fungsi tersebut mempunyai
\definitionsverweis {ekstremum lokal}{}{.}
Tunjukkan bahwa

\mavergleichskettedisp
{\vergleichskette
{ (df)_P
}
{ =} { 0
}
{ } {
}
{ } {
}
{ } {
}
}
{}{}{.}

}
{} {}

\inputaufgabegibtloesung
{5}
{

Kita tinjau
$1$-\definitionsverweis {bentuk diferensial}{}{}

\mavergleichskettedisp
{\vergleichskette
{ \omega
}
{ =} { -vdu+udv
}
{ } {
}
{ } {
}
{ } {
}
}
{}{}{}
pada
$1$-\definitionsverweis {bola}{}{}

\mavergleichskette
{\vergleichskette
{ S^1
}
{ \subset }{ \R^2
}
{  }{
}
{    }{
}
{   }{
}
}
{}{}{,}
dengan koordinat ruang sekitar $\R^2$ dinotasikan oleh
\mathkor {} {u} {dan} {v} {.}
Tentukan
\mathl{\pi^* \omega}{} di bawah
\definitionsverweis {pemetaan terdiferensialkan secara kontinu}{}{}
\maabbeledisp {\pi} {\R^2 \setminus \{0\}} { S^{1}
} {(x , y) } { { \frac{   1  }{  \sqrt{x^2 + y^2}  } } (x , y)
} {.}

}
{} {}

\inputaufgabegibtloesung
{5}
{

Buktikan teorema tentang perhitungan volume kanonik pada manifold Riemann $M$.

}
{} {}

\inputaufgabegibtloesung
{6 (1+1+3+1)}
{

Jelaskan mengapa $\R^2$ tidak membawa struktur
\definitionsverweis {manifold berbatas}{}{}
sedemikian sehingga subhimpunan $T$ yang diberikan merupakan
\definitionsverweis {batas}{}{.}
\aufzaehlungvierabc{
\mavergleichskette
{\vergleichskette
{ T
}
{ = }{ {]0,1[}
}
{  }{
}
{    }{
}
{   }{
}
}
{}{}{,}
dengan interval tersebut terletak pada sumbu $x$.
}{
\mavergleichskette
{\vergleichskette
{ T
}
{ = }{ [0,1]
}
{  }{
}
{    }{
}
{   }{
}
}
{}{}{,}
dengan interval tersebut terletak pada sumbu $x$.
}{
\mavergleichskette
{\vergleichskette
{ T
}
{ = }{ \R
}
{  }{
}
{    }{
}
{   }{
}
}
{}{}{,}
dengan $\R$ adalah sumbu $x$.
}{
\mavergleichskette
{\vergleichskette
{ T
}
{ = }{ S^1
}
{  }{
}
{    }{
}
{   }{
}
}
{}{}{.}
}

}
{} {}

\inputaufgabegibtloesung
{4}
{

Misalkan $\omega$ suatu
$n-1$-\definitionsverweis {bentuk diferensial}{}{}
yang terdiferensialkan secara kontinu pada
\definitionsverweis {himpunan terbuka}{}{}

\mavergleichskette
{\vergleichskette
{ U
}
{ \subseteq }{ \R^n
}
{  }{
}
{    }{
}
{   }{
}
}
{}{}{}
dan misalkan
\definitionsverweis {turunan eksterior}{}{}
sama dengan

\mavergleichskettedisp
{\vergleichskette
{d \omega
}
{ =} { fdx_1 \wedge \ldots \wedge dx_n
}
{ } {
}
{ } {
}
{ } {
}
}
{}{}{}
dengan suatu fungsi
\maabb {f} { U } {\R
} {.}
Tunjukkan untuk

\mavergleichskette
{\vergleichskette
{ P
}
{ \in }{ U
}
{  }{
}
{    }{
}
{   }{
}
}
{}{}{}
kesamaan

\mavergleichskettedisp
{\vergleichskette
{ f(P)
}
{ =} { { \frac{   1  }{  \beta_n  } } \cdot \operatorname{lim}_{   \epsilon \rightarrow  0  } \, { \frac{   1  }{  \epsilon^n } } \int_{S^{n-1}(P, \epsilon)} \omega
}
{ } {
}
{ } {
}
{ } {
}
}
{}{}{,}
dengan $\beta_n$ menyatakan
volume
bola satuan berdimensi $n$.

}
{} {}

\inputaufgabegibtloesung
{12}
{

Buktikan teorema Stokes untuk manifold berbatas.

}
{} {}
\endgroup

\chapter{Formulir Peserta Ujian 3}
\noindent\textit{Formulir peserta dari sumber resmi; teks soal dipertahankan sesuai terjemahan yang telah diverifikasi.}
\begingroup
\renewcommand{\theHfakt}{o011.exam.learner.03.\arabic{fakt}}
%Data institusi

%\input{Dozentdaten}

%\renewcommand{\fachbereich}{Fachbereich}

%\renewcommand{\dozent}{Prof. Dr. . }

%Data ujian

\renewcommand{\klausurgebiet}{ }

\renewcommand{\klausurtyp}{ }

\renewcommand{\klausurdatum}{ . 20}

\klausurvorspann {\fachbereich} {\klausurdatum} {\dozent} {\klausurgebiet} {\klausurtyp}

%Data untuk tabel nilai berikut

\renewcommand{\aeins}{  3  }

\renewcommand{\azwei}{  3   }

\renewcommand{\adrei}{  0  }

\renewcommand{\avier}{  0  }

\renewcommand{\afuenf}{  6  }

\renewcommand{\asechs}{  6  }

\renewcommand{\asieben}{  0  }

\renewcommand{\aacht}{  6  }

\renewcommand{\aneun}{  11  }

\renewcommand{\azehn}{  5  }

\renewcommand{\aelf}{  5  }

\renewcommand{\azwoelf}{  3  }

\renewcommand{\adreizehn}{  7   }

\renewcommand{\avierzehn}{  2  }

\renewcommand{\afuenfzehn}{  0  }

\renewcommand{\asechzehn}{  6  }

\renewcommand{\asiebzehn}{  63  }

\renewcommand{\aachtzehn}{    }

\renewcommand{\aneunzehn}{    }

\renewcommand{\azwanzig}{    }

\renewcommand{\aeinundzwanzig}{    }

\renewcommand{\azweiundzwanzig}{    }

\renewcommand{\adreiundzwanzig}{    }

\renewcommand{\avierundzwanzig}{    }

\renewcommand{\afuenfundzwanzig}{    }

\renewcommand{\asechsundzwanzig}{    }

\punktetabellesechzehn

\klausurnote

\newpage

\setcounter{section}{K}

\inputaufgabegibtloesung
{3}
{

Definisikan istilah-istilah berikut
\zusatzklammer {yang dicetak miring} {} {.}
\aufzaehlungsechs{Suatu
\stichwort {kurva geodesik} {}
pada hiperpermukaan diferensiabel

\mavergleichskette
{\vergleichskette
{ Y
}
{ \subseteq }{ \R^n
}
{  }{
}
{    }{
}
{   }{
}
}
{}{}{.}

}{
\stichwort {Himpunan rotasi} {}
\zusatzklammer {mengelilingi sumbu $x$} {} {}
dari suatu subhimpunan

\mavergleichskette
{\vergleichskette
{ T
}
{ \subseteq }{  \R \times \R_{\geq 0}
}
{  }{
}
{    }{
}
{   }{
}
}
{}{}{.}

}{
\stichwort {Pemetaan tangen} {}
\maabbdisp {T (\varphi)} {TM} {TN
} {}
dari
\definitionsverweis {pemetaan diferensiabel}{}{}
\maabbdisp {\varphi} { M } { N
} {}
antara
\definitionsverweis {manifold-manifold diferensiabel}{}{}
\mathkor {} {M} {dan} {N} {.}

}{Suatu perubahan peta yang
\stichwort {mempertahankan orientasi} {}
pada
\definitionsverweis {manifold diferensiabel}{}{}
$M$.

}{
\stichwort {Bentuk diferensial tarikan balik} {} $\varphi^* \omega$ dari suatu bentuk diferensial
\mathl{\omega \in
{ \mathcal E }^{ k } (  M   )}{} terhadap suatu pemetaan terdiferensialkan secara kontinu
\maabb {\varphi} { L } { M
} {}
antara dua manifold diferensiabel
\mathkor {} {L} {dan} {M} {.}

}{Suatu koneksi linear yang
\stichwort {metrik} {}
pada bundel tangen suatu manifold Riemann $M$.
}

}
{} {}

\inputaufgabegibtloesung
{3}
{

Nyatakan teorema-teorema berikut.
\aufzaehlungdrei{
\stichwort {Teorema tentang solusi persamaan diferensial biasa pada hiperpermukaan diferensiabel} {}

\mavergleichskette
{\vergleichskette
{  Y
}
{ \subseteq }{  \R^n
}
{  }{
}
{    }{
}
{   }{
}
}
{}{}{.}}{Rumus untuk menghitung volume kanonik suatu manifold Riemann dalam sebuah peta.}{Teorema tentang partisi kesatuan.}

}
{} {}

\inputaufgabe
{0}
{

}
{} {}

\inputaufgabe
{0}
{

}
{} {}

\inputaufgabegibtloesung
{6}
{

Untuk permukaan $Y$ yang diberikan oleh

\mavergleichskettedisp
{\vergleichskette
{ x^2+y^2+2z^2
}
{ =} { 4
}
{ } {
}
{ } {
}
{ } {
}
}
{}{}{}
dan titik

\mavergleichskettedisp
{\vergleichskette
{ P
}
{ =} { \left(  1  , \,   1  , \,   1                 \right)
}
{ \in} { Y
}
{ } {
}
{ } {
}
}
{}{}{,}
tentukan suatu matriks diagonal bagi
\definitionsverweis {pemetaan Weingarten}{}{}
$L_P$.

}
{} {}

\inputaufgabegibtloesung
{6}
{

Buktikan teorema isometri untuk transpor paralel pada hiperpermukaan

\mavergleichskette
{\vergleichskette
{ Y
}
{ \subseteq }{ \R^n
}
{  }{
}
{    }{
}
{   }{
}
}
{}{}{.}

}
{} {}

\inputaufgabe
{0}
{

}
{} {}

\inputaufgabegibtloesung
{6}
{

Misalkan $M$ suatu manifold topologis kompak berdimensi $d$,

\mavergleichskette
{\vergleichskette
{ d
}
{ \geq }{ 1
}
{  }{
}
{    }{
}
{   }{
}
}
{}{}{.}
Tunjukkan bahwa terdapat suatu subhimpunan terbuka terbatas

\mavergleichskette
{\vergleichskette
{ U
}
{ \subseteq }{ \R^d
}
{  }{
}
{    }{
}
{   }{
}
}
{}{}{}
dan suatu pemetaan kontinu surjektif
\maabbdisp {\varphi} { U } { M
} {}
.

}
{} {}

\inputaufgabegibtloesung
{11 (1+3+1+2+4)}
{

Kita tinjau himpunan

\mavergleichskettedisp
{\vergleichskette
{N
}
{ =} { \left\{ A= \begin{pmatrix}  a   &   b   \\  c  & d \end{pmatrix}   \mid   A \text{ nilpoten}         \right\}
}
{ \subseteq} {\R^4
}
{ } {
}
{ } {
}
}
{}{}{}
dari
$2\times 2$-\definitionsverweis {matriks}{}{}
real nilpoten, serta himpunan

\mavergleichskettedisp
{\vergleichskette
{M
}
{ =} { N \setminus \{ \begin{pmatrix}  0   &   0  \\  0  &  0 \end{pmatrix}  \}
}
{ } {
}
{ } {
}
{ } {
}
}
{}{}{.}

a) Apakah $N$ terhubung?

b) Tunjukkan bahwa $M$ suatu
\definitionsverweis {submanifold tertutup}{}{}
dari suatu subhimpunan terbuka

\mavergleichskette
{\vergleichskette
{ G
}
{ \subseteq }{ \R^4
}
{  }{
}
{    }{
}
{   }{
}
}
{}{}{.}

c) Tentukan
\definitionsverweis {dimensi}{}{}
$M$.

d) Apakah $M$ terhubung?

e) Tutupi
\mathl{M}{} dengan peta-peta topologis eksplisit.

}
{} {}

\inputaufgabe
{5}
{

Katakan sesuatu yang cerdas tentang pita Möbius.

}
{} {}

\inputaufgabegibtloesung
{5}
{

Kita tinjau pemetaan
\maabbeledisp {\varphi} { \R^3 } { \R^2
} { (x,y,z) } { \left(  xyz^2,xy-z^3  \right)
} {.}
Hitung matriks pemetaan
\maabbdisp {\bigwedge^2 T_P (\varphi)} { \bigwedge^2 T_P \R^3 } { \bigwedge^2 T_{\varphi(P)}  \R^2
} {}
di titik

\mavergleichskette
{\vergleichskette
{ P
}
{ = }{ (1,3,5)
}
{  }{
}
{    }{
}
{   }{
}
}
{}{}{}
terhadap suatu basis yang sesuai.

}
{} {}

\inputaufgabegibtloesung
{3}
{

Misalkan

\mavergleichskette
{\vergleichskette
{ W
}
{ \subseteq }{ \R^m
}
{  }{
}
{    }{
}
{   }{
}
}
{}{}{}
dan

\mavergleichskette
{\vergleichskette
{ U
}
{ \subseteq }{ \R^n
}
{  }{
}
{    }{
}
{   }{
}
}
{}{}{}
merupakan
\definitionsverweis {subhimpunan-subhimpunan terbuka}{}{}
dan misalkan
\maabbdisp {\psi} { W } { U
} {}
suatu
\definitionsverweis {pemetaan}{}{}
yang
\definitionsverweis {terdiferensialkan secara kontinu}{}{.}
Misalkan
\maabbdisp {f} { U } {\R
} {}
suatu fungsi terdiferensialkan secara kontinu. Simpulkan dari
aturan rantai
bahwa

\mavergleichskettedisp
{\vergleichskette
{ d (\psi^*f)
}
{ =} { \psi^*(df)
}
{ } {
}
{ } {
}
{ } {
}
}
{}{}{,}
dengan $\psi^*$ menyatakan
\definitionsverweis {penarikan balik bentuk diferensial}{}{.}

}
{} {}

\inputaufgabegibtloesung
{7}
{

Misalkan $M$ suatu
\definitionsverweis {manifold diferensiabel}{}{}
dengan
\definitionsverweis {basis terhitung bagi topologinya}{}{}
dan misalkan $\omega$ suatu
\definitionsverweis {bentuk volume positif}{}{}
pada $M$. Misalkan
\maabbdisp {\varphi} { L } { M
} {}
suatu
\definitionsverweis {difeomorfisme}{}{}
dengan manifold $L$ dan

\mavergleichskette
{\vergleichskette
{ S
}
{ \subseteq }{ L
}
{  }{
}
{    }{
}
{   }{
}
}
{}{}{}
suatu
\definitionsverweis {subhimpunan terukur}{}{.}
Tunjukkan bahwa

\mavergleichskettedisp
{\vergleichskette
{ \int_S \varphi^* \omega
}
{ =} { \int_{ \varphi(S)} \omega
}
{ } {
}
{ } {
}
{ } {
}
}
{}{}{.}

}
{} {}

\inputaufgabegibtloesung
{2}
{

Pada
\definitionsverweis {bundel vektor trivial}{}{}
\maabbdisp {p} {\R^2 \times \R} { \R^2
} {}
ber-
\definitionsverweis {rank}{}{}
$1$ di atas $\R^2$, kita tinjau
\definitionsverweis {koneksi linear}{}{,}
yang diberikan oleh
\definitionsverweis {simbol-simbol Christoffel}{}{}

\mavergleichskette
{\vergleichskette
{ \Gamma_1 (x,y)
}
{ = }{ x^5-x^2y^2+3y^4
}
{  }{
}
{    }{
}
{   }{
}
}
{}{}{}
dan

\mavergleichskette
{\vergleichskette
{ \Gamma_2 (x,y)
}
{ = }{ 7x^3 +xy^2-4y^5
}
{  }{
}
{    }{
}
{   }{
}
}
{}{}{.}
Hitung
\definitionsverweis {operator kelengkungan}{}{}
$R \left(  \partial_1, \partial_2  \right)$.

}
{} {}

\inputaufgabe
{0}
{

}
{} {}

\inputaufgabegibtloesung
{6}
{

Buktikan teorema tentang retraksi ke batas pada manifold.

}
{} {}
\endgroup

\chapter{Formulir Peserta Ujian 4}
\noindent\textit{Formulir peserta dari sumber resmi; teks soal dipertahankan sesuai terjemahan yang telah diverifikasi.}
\begingroup
\renewcommand{\theHfakt}{o011.exam.learner.04.\arabic{fakt}}
%Data institusi

%\input{Dozentdaten}

%\renewcommand{\fachbereich}{Fachbereich}

%\renewcommand{\dozent}{Prof. Dr. . }

%Data ujian

\renewcommand{\klausurgebiet}{ }

\renewcommand{\klausurtyp}{ }

\renewcommand{\klausurdatum}{ . 20}

\klausurvorspann {\fachbereich} {\klausurdatum} {\dozent} {\klausurgebiet} {\klausurtyp}

%Data untuk tabel nilai berikut

\renewcommand{\aeins}{  3  }

\renewcommand{\azwei}{  3   }

\renewcommand{\adrei}{  0  }

\renewcommand{\avier}{  0  }

\renewcommand{\afuenf}{  4  }

\renewcommand{\asechs}{  5  }

\renewcommand{\asieben}{  4  }

\renewcommand{\aacht}{  0  }

\renewcommand{\aneun}{  4  }

\renewcommand{\azehn}{  6  }

\renewcommand{\aelf}{  2  }

\renewcommand{\azwoelf}{  6  }

\renewcommand{\adreizehn}{  4   }

\renewcommand{\avierzehn}{  10  }

\renewcommand{\afuenfzehn}{  10  }

\renewcommand{\asechzehn}{  0  }

\renewcommand{\asiebzehn}{  61  }

\renewcommand{\aachtzehn}{    }

\renewcommand{\aneunzehn}{    }

\renewcommand{\azwanzig}{    }

\renewcommand{\aeinundzwanzig}{    }

\renewcommand{\azweiundzwanzig}{    }

\renewcommand{\adreiundzwanzig}{    }

\renewcommand{\avierundzwanzig}{    }

\renewcommand{\afuenfundzwanzig}{    }

\renewcommand{\asechsundzwanzig}{    }

\punktetabellesechzehn

\klausurnote

\newpage

\setcounter{section}{K}

\inputaufgabegibtloesung
{3}
{

Definisikan istilah-istilah berikut
\zusatzklammer {yang dicetak miring} {} {}.
\aufzaehlungsechs{Pemetaan
\stichwort {Weingarten} {}
pada $T_PY$ di suatu titik

\mavergleichskette
{\vergleichskette
{ P
}
{ \in }{ Y
}
{  }{
}
{    }{
}
{   }{
}
}
{}{}{}
dari suatu hiperpermukaan diferensiabel berorientasi

\mavergleichskette
{\vergleichskette
{ Y
}
{ \subseteq }{ \R^n
}
{  }{
}
{    }{
}
{   }{
}
}
{}{}{.}

}{Suatu
\stichwort {bundel vektor riil} {}
$V$ ber-rank $r$ di atas suatu
\definitionsverweis {ruang topologis}{}{}
$X$.

}{\stichwort {Integral lintasan} {} dari suatu bentuk diferensial-$1$
\mathl{\omega \in
{ \mathcal E }^{  1 } (  M   )}{} pada suatu manifold diferensiabel $M$ terhadap suatu kurva yang terdiferensialkan secara kontinu
\maabb {\gamma} {[a,b]} {M
} {.}

}{Suatu
\stichwort {bentuk volume positif} {}
$\omega$ pada suatu manifold diferensiabel $M$.

}{Suatu \stichwort {bentuk diferensial tertutup} {} $\omega$ pada suatu manifold diferensiabel $M$.

}{\stichwort {Percepatan tangensial} {}
dari suatu kurva yang dua kali terdiferensialkan secara kontinu
\maabb {\gamma} {I} {M
} {}
pada suatu
\definitionsverweis {manifold Riemann}{}{}
$M$.
}

}
{} {}

\inputaufgabegibtloesung
{3}
{

Rumuskan teorema-teorema berikut.
\aufzaehlungdrei{\stichwort {Teorema tentang kelengkungan normal} {.}}{Rumus untuk menghitung volume kanonik pada suatu manifold Riemann di dalam suatu peta.}{\stichwort {Teorema Stokes} {} untuk manifold dengan batas.}

}
{} {}

\inputaufgabe
{0}
{

}
{} {}

\inputaufgabe
{0}
{

}
{} {}

\inputaufgabegibtloesung
{4}
{

Buktikan teorema tentang hubungan antara kelengkungan dan vektor normal satuan dari suatu kurva bidang yang diparameterkan oleh panjang busur.

}
{} {}

\inputaufgabegibtloesung
{5 (1+2+2)}
{

Pada
\definitionsverweis {permukaan bola satuan}{}{}
$Y$ berdimensi dua, kita meninjau
\definitionsverweis {medan vektor tangensial}{}{}

\mavergleichskettedisp
{\vergleichskette
{ F \begin{pmatrix}  x  \\ y \\ z   \end{pmatrix}
}
{ =} { \begin{pmatrix}  y  \\ -x\\  0   \end{pmatrix}
}
{ } {
}
{ } {
}
{ } {
}
}
{}{}{.}
\aufzaehlungdrei{Tentukan
\definitionsverweis {matriks Jacobi}{}{}
dari $F$ pada $\R^3$.
}{Untuk titik

\mavergleichskette
{\vergleichskette
{ P
}
{ = }{ \begin{pmatrix}  0  \\ 1 \\  0   \end{pmatrix}
}
{ \in }{ Y
}
{    }{
}
{   }{
}
}
{}{}{}
tentukan suatu matriks untuk pemetaan
\maabbeledisp {} {T_PY} { T_PY
} { v } { \nabla_v F
} {,}
terhadap suatu
\definitionsverweis {basis}{}{}
yang sesuai.
}{Untuk titik

\mavergleichskette
{\vergleichskette
{ Q
}
{ = }{ \begin{pmatrix}  0  \\ 0\\  1   \end{pmatrix}
}
{ \in }{ Y
}
{    }{
}
{   }{
}
}
{}{}{}
tentukan suatu matriks untuk pemetaan
\maabbeledisp {} {T_QY} { T_QY
} { v } { \nabla_v F
} {,}
terhadap suatu
\definitionsverweis {basis}{}{}
yang sesuai.
}

}
{} {}

\inputaufgabegibtloesung
{4 (2+2)}
{

Kita meninjau elips

\mavergleichskettedisp
{\vergleichskette
{ E
}
{ =} { \left\{ (u,v)  \mid   2u^2 +5v^2 = 4         \right\}
}
{ } {
}
{ } {
}
{ } {
}
}
{}{}{.}
\aufzaehlungdrei{Definisikan suatu pemetaan surjektif yang terdiferensialkan secara kontinu
\maabb {} {\R} {E
} {.}
}{Deskripsikan suatu
\definitionsverweis {difeomorfisme}{}{}
antara lingkaran $S^1$ dan $E$.
}{
}

}
{} {}

\inputaufgabe
{0}
{

}
{} {}

\inputaufgabegibtloesung
{4}
{

Berikan suatu
\definitionsverweis {medan vektor}{}{} yang
\definitionsverweis {kontinu}{}{}
pada $S^2$ dan hanya mempunyai satu titik nol.

}
{} {}

\inputaufgabegibtloesung
{6}
{

Misalkan diberikan suatu
\definitionsverweis {data perekatan}{}{}

\mathbed {U_i} {}
 {i \in I} {}
 {} {} {} {,}
untuk
\definitionsverweis {ruang-ruang topologis}{}{.}
Buktikan bahwa terdapat suatu ruang topologis $X$ yang ditentukan secara unik, suatu penutup terbuka

\mavergleichskette
{\vergleichskette
{ X
}
{ = }{ \bigcup_{i \in I} V_i
}
{  }{
}
{    }{
}
{   }{
}
}
{}{}{}
dan
\definitionsverweis {homeomorfisme}{}{}
\maabb {\psi_i} {U_i } { V_i
} {}
sedemikian sehingga

\mavergleichskettedisp
{\vergleichskette
{ \psi_i  \left(  U_{ij}  \right)
}
{ =} { V_i \cap V_j
}
{ } {
}
{ } {
}
{ } {
}
}
{}{}{}
dan

\mavergleichskettedisp
{\vergleichskette
{ \psi_i \vert_{U_{ij} }
}
{ =} { \psi_j \vert_{U_{ji} }  \circ   \varphi_{ji}
}
{ } {
}
{ } {
}
{ } {
}
}
{}{}{}.

}
{} {}

\inputaufgabegibtloesung
{2}
{

Misalkan $M$ suatu
\definitionsverweis {manifold diferensiabel}{}{,}
\maabb {\gamma} { I } { M
} {}
suatu
\definitionsverweis {kurva diferensiabel}{}{}
di $M$, dan $\omega$ suatu
$1$-\definitionsverweis {bentuk}{}{}
terdiferensialkan pada $M$. Buktikan bahwa

\mavergleichskettedisp
{\vergleichskette
{ \gamma^* (d \omega)
}
{ =} { 0
}
{ } {
}
{ } {
}
{ } {
}
}
{}{}{}.

}
{} {}

\inputaufgabegibtloesung
{6 (3+3)}
{

Misalkan $M$ suatu
\definitionsverweis {manifold diferensiabel}{}{}
yang
\definitionsverweis {kompak}{}{}
dan mempunyai sedikitnya dua titik.
\aufzaehlungzwei {Buktikan bahwa suatu
$1$-\definitionsverweis {bentuk diferensial}{}{}
kontinu yang
\definitionsverweis {eksak}{}{}
pada $M$ mempunyai sedikitnya dua titik nol.
} {Buktikan bahwa hal ini tidak harus berlaku bagi suatu bentuk diferensial yang
\definitionsverweis {tertutup}{}{.}
}

}
{} {}

\inputaufgabegibtloesung
{4}
{

Berikan suatu contoh
\definitionsverweis {manifold dengan batas}{}{,}
yang
\definitionsverweis {terhubung}{}{}
dan batasnya terdiri atas
\zusatzklammer {gabungan saling lepas dari} {} {}
tak terhingga banyak salinan $\R^1$.

}
{} {}

\inputaufgabegibtloesung
{10}
{

Buktikan \stichwort {teorema tentang partisi kesatuan} {.}

}
{} {}

\inputaufgabegibtloesung
{10 (2+6+2)}
{

Kita meninjau
$2$-\definitionsverweis {bentuk diferensial}{}{}

\mavergleichskettedisp
{\vergleichskette
{ \omega
}
{ =} { x dy \wedge dz -y dx \wedge dz + z dx \wedge dy
}
{ } {
}
{ } {
}
{ } {
}
}
{}{}{}
pada bola satuan
\mathl{B  \left(  0 , 1  \right)}{.}
\aufzaehlungdreiabc{Buktikan bahwa $d \omega$ adalah tiga kali bentuk volume standar pada bola tersebut.
}{Buktikan bahwa
\mathl{\omega\vert_{S^2}}{} merupakan bentuk luas permukaan standar pada permukaan bola satuan.
}{Hitung luas permukaan bola dari
volume bola
dengan menggunakan
Teorema Stokes.
}

}
{} {}

\inputaufgabe
{0}
{

}
{} {}
\endgroup

\chapter{Formulir Peserta Ujian 5}
\noindent\textit{Formulir peserta dari sumber resmi; teks soal dipertahankan sesuai terjemahan yang telah diverifikasi.}
\begingroup
\renewcommand{\theHfakt}{o011.exam.learner.05.\arabic{fakt}}
%Data institusi

%\input{Dozentdaten}

%\renewcommand{\fachbereich}{Bidang studi}

%\renewcommand{\dozent}{Prof. Dr. . }

%Data ujian

\renewcommand{\klausurgebiet}{ }

\renewcommand{\klausurtyp}{ }

\renewcommand{\klausurdatum}{. 20}

\klausurvorspann {\fachbereich} {\klausurdatum} {\dozent} {\klausurgebiet} {\klausurtyp}

%Data untuk tabel poin berikut

\renewcommand{\aeins}{  3  }

\renewcommand{\azwei}{  3   }

\renewcommand{\adrei}{  5  }

\renewcommand{\avier}{  2  }

\renewcommand{\afuenf}{  4  }

\renewcommand{\asechs}{  8  }

\renewcommand{\asieben}{  4  }

\renewcommand{\aacht}{  5  }

\renewcommand{\aneun}{  3  }

\renewcommand{\azehn}{  2  }

\renewcommand{\aelf}{  6  }

\renewcommand{\azwoelf}{  7  }

\renewcommand{\adreizehn}{  4   }

\renewcommand{\avierzehn}{  8  }

\renewcommand{\afuenfzehn}{  64  }

\renewcommand{\asechzehn}{    }

\renewcommand{\asiebzehn}{    }

\renewcommand{\aachtzehn}{    }

\renewcommand{\aneunzehn}{    }

\renewcommand{\azwanzig}{    }

\renewcommand{\aeinundzwanzig}{    }

\renewcommand{\azweiundzwanzig}{    }

\renewcommand{\adreiundzwanzig}{    }

\renewcommand{\avierundzwanzig}{    }

\renewcommand{\afuenfundzwanzig}{    }

\renewcommand{\asechsundzwanzig}{    }

\punktetabellevierzehn

\klausurnote

\newpage

\setcounter{section}{0}

\inputaufgabegibtloesung
{3}
{

Definisikan istilah-istilah berikut
\zusatzklammer {yang dicetak miring} {} {}.
\aufzaehlungsechs{\stichwort {turunan tangensial kedua} {}
\zusatzklammer {atau
\stichwort {percepatan tangensial} {}} {} {}
\zusatzklammer {sebagai pemetaan ke $\R^n$} {} {}
dari suatu
\definitionsverweis {kurva yang dua kali terdiferensialkan secara kontinu}{}{}
\maabbdisp {\gamma} {I} {Y
} {}
yang bernilai dalam suatu
\definitionsverweis {hipermuka terdiferensialkan}{}{}

\mavergleichskette
{\vergleichskette
{ Y
}
{ \subseteq }{ \R^n
}
{  }{
}
{    }{
}
{   }{
}
}
{}{}{.}

}{Suatu ruang topologis $X$ yang
\stichwort {terhubung} {}.

}{\stichwort {Pemetaan singgung di titik} {}

\mavergleichskette
{\vergleichskette
{  P
}
{ \in }{ M
}
{  }{}
{    }{}
{   }{}
}
{}{}{}
untuk suatu pemetaan diferensiabel
\maabbdisp {\varphi} { M } { N
} {,}
dengan
\mathkor {} {M} {dan} {N} {}
manifold diferensiabel.

}{\stichwort {Produk} {}
dari manifold diferensiabel
\mathkor {} {M} {dan} {N} {.}

}{\stichwort {Ukuran volume} {}
yang terkait dengan bentuk volume positif $\omega$ pada suatu manifold diferensiabel berdimensi $n$.

}{Suatu
\stichwort {kurva geodesik} {}
pada manifold Riemann $M$.
}

}
{} {}

\inputaufgabegibtloesung
{3}
{

Nyatakan teorema-teorema berikut.
\aufzaehlungdrei{\stichwort {Teorema tentang kelengkungan kurva bidang yang diberikan secara implisit} {}

\mavergleichskette
{\vergleichskette
{  Y
}
{ \subseteq }{  \R^2
}
{  }{
}
{    }{
}
{   }{
}
}
{}{}{.}}{\stichwort {Rumus transformasi untuk bentuk volume positif pada manifold} {.}}{Teorema tentang kelengkungan Gauss dan kelengkungan seksional pada suatu permukaan diferensiabel

\mavergleichskette
{\vergleichskette
{  Y
}
{ \subseteq }{  \R^3
}
{  }{
}
{    }{
}
{   }{
}
}
{}{}{}}

}
{} {}

\inputaufgabegibtloesung
{5}
{

Misalkan

\mavergleichskette
{\vergleichskette
{ G
}
{ \subseteq }{ \R^n
}
{  }{
}
{    }{
}
{   }{
}
}
{}{}{}
\definitionsverweis {terbuka}{}{}
dan
\maabbdisp {\varphi} { G } { \R^m
} {}
adalah
\definitionsverweis {pemetaan yang terdiferensialkan secara kontinu}{}{,}
yang pada titik

\mavergleichskette
{\vergleichskette
{ P
}
{ \in }{ G
}
{  }{
}
{    }{
}
{   }{
}
}
{}{}{}
memiliki
\definitionsverweis {diferensial total}{}{}
yang surjektif. Misalkan

\mavergleichskette
{\vergleichskette
{ v
}
{ \in }{ T_PZ
}
{  }{
}
{    }{
}
{   }{
}
}
{}{}{}
adalah suatu vektor dalam
\definitionsverweis {ruang singgung serat}{}{}
$Z$ dari $\varphi$ yang melalui $P$. Tunjukkan bahwa terdapat
\definitionsverweis {kurva yang terdiferensialkan secara kontinu}{}{}
\maabbdisp {\gamma} { ]- \epsilon, \epsilon[ } { Z
} {}
\zusatzklammer {untuk suatu

\mavergleichskettek
{\vergleichskettek
{ \epsilon
}
{ > }{ 0
}
{  }{
}
{    }{
}
{   }{
}
}
{}{}{}
yang sesuai} {} {}
dengan

\mavergleichskette
{\vergleichskette
{ \gamma(0)
}
{ = }{ P
}
{  }{
}
{    }{
}
{   }{
}
}
{}{}{}
dan

\mavergleichskettedisp
{\vergleichskette
{ \gamma'(0)
}
{ =} { v
}
{ } {
}
{ } {
}
{ } {
}
}
{}{}{}.

}
{} {}

\inputaufgabegibtloesung
{2}
{

Misalkan
\mathkor {} {r} {dan} {R} {}
adalah bilangan real positif dengan

\mavergleichskette
{\vergleichskette
{ 0
}
{ < }{ r
}
{ < }{ R
}
{    }{
}
{   }{
}
}
{}{}{.}
Tentukan suatu
\definitionsverweis {medan normal satuan}{}{}
untuk torus
\zusatzklammer {tertanam} {} {}

\mavergleichskettedisp
{\vergleichskette
{ T
}
{ =} { \left\{ (x,y,z) \in \R^3  \mid   \left(  \sqrt{x^2+y^2} -R  \right)^2 +z^2  = r^2         \right\}
}
{ } {
}
{ } {
}
{ } {
}
}
{}{}{.}

}
{} {}

\inputaufgabe
{4}
{

Jelaskan peran teorema tentang pemetaan implisit dalam pengembangan teori manifold.

}
{} {}

\inputaufgabegibtloesung
{8}
{

Buktikan teorema tentang sifat adjoin-diri (self-adjoint) dari pemetaan Weingarten.

}
{} {}

\inputaufgabegibtloesung
{4}
{

Misalkan $X$ adalah suatu
\definitionsverweis {ruang kompak}{}{}
dan misalkan

\mavergleichskette
{\vergleichskette
{Y
}
{ \subseteq }{ X
}
{  }{
}
{    }{
}
{   }{
}
}
{}{}{}
adalah
\definitionsverweis {himpunan bagian tertutup}{}{,}
yang dilengkapi dengan
\definitionsverweis {topologi terinduksi}{}{.}
Tunjukkan bahwa $Y$ juga kompak.

}
{} {}

\inputaufgabegibtloesung
{5}
{

Berikan contoh suatu
\definitionsverweis {manifold diferensiabel}{}{}
$M$ yang
\definitionsverweis {berdimensi dua}{}{}
dan
\definitionsverweis {terhubung}{}{,}
beserta suatu titik

\mavergleichskette
{\vergleichskette
{ P
}
{ \in }{ M
}
{  }{
}
{    }{
}
{   }{
}
}
{}{}{}
sedemikian sehingga
\mathkor {} {M} {dan} {M \setminus \{P\}} {}
saling
\definitionsverweis {difeomorfik}{}{.}

}
{} {}

\inputaufgabegibtloesung
{3}
{

Untuk setiap
\definitionsverweis {vektor singgung}{}{}

\mavergleichskette
{\vergleichskette
{ v
}
{ \in }{ T_PS^2
}
{ \subseteq }{ \R^3
}
{    }{
}
{   }{
}
}
{}{}{}
dengan

\mavergleichskette
{\vergleichskette
{ \Vert {v} \Vert
}
{ = }{ 1
}
{  }{
}
{    }{
}
{   }{
}
}
{}{}{}
di suatu titik

\mavergleichskette
{\vergleichskette
{ P
}
{ \in }{ S^2
}
{ \subset }{ \R^3
}
{    }{
}
{   }{
}
}
{}{}{}
pada
\definitionsverweis {permukaan bola satuan}{}{,}
berikan suatu wakil
\definitionsverweis {diferensiabel}{}{}
\maabbdisp {\gamma} { I } { S^2
} {}
dengan

\mavergleichskette
{\vergleichskette
{ \gamma(0)
}
{ = }{ P
}
{  }{
}
{    }{
}
{   }{
}
}
{}{}{.}

}
{} {}

\inputaufgabegibtloesung
{2}
{

Tinjau parabola standar

\mavergleichskette
{\vergleichskette
{ Y
}
{ \subseteq }{ \R^2
}
{  }{
}
{    }{
}
{   }{
}
}
{}{}{}
dengan struktur Riemann terinduksi dan parametrisasi
\maabbeledisp {\varphi} {\R} { Y
} { t } { \left(  t   , \,  t^2                   \right)
} {,}
yang kita pandang sebagai suatu peta
\zusatzklammer {invers} {} {}.
Tentukan
\definitionsverweis {fungsi fundamental}{}{}
Riemann

\mavergleichskette
{\vergleichskette
{ g(t)
}
{ = }{ g_{11} (t)
}
{  }{
}
{    }{
}
{   }{
}
}
{}{}{.}

}
{} {}

\inputaufgabegibtloesung
{6}
{

Hitung bentuk diferensial tarik balik
\mathl{\varphi^* \tau}{} dari

\mavergleichskettedisp
{\vergleichskette
{ \tau
}
{ =} { dx \wedge dy \wedge dz - w dx \wedge dy \wedge dw + \cos (xy)  dx \wedge dz \wedge dw-yw dy \wedge dz \wedge dw
}
{ } {
}
{ } {
}
{ } {
}
}
{}{}{}
oleh pemetaan
\maabbeledisp {\varphi} { \R^3 } { \R^4
} { (r,s,t) } {  \left(  r^2s,t, \sin  r ,e^{st}  \right)  = (x,y,z,w)
} {.}

}
{} {}

\inputaufgabegibtloesung
{7 (2+3+2)}
{

Tinjau pemetaan-pemetaan diferensiabel
\maabbeledisp {\gamma} { [1,2] } {\R^2
} { t } {(t,t^{-1})
} {,}
dan
\maabbeledisp {\varphi} {\R^2} {\R^3
} {(u,v)} {(u^2,uv,-u+v^2)
} {,}
serta bentuk diferensial

\mavergleichskettedisp
{\vergleichskette
{ \omega
}
{ =} { xdx-zdy+dz
}
{ } {
}
{ } {
}
{ } {
}
}
{}{}{}
pada $\R^3$.

a) Hitung bentuk diferensial tarik balik
\mathl{\varphi^* \omega}{} pada $\R^2$.

b) Hitung integral lintasan bentuk diferensial
\mathl{\varphi^* \omega}{} sepanjang lintasan $\gamma$.

c) Hitung
\zusatzklammer {tanpa menggunakan hasil b)} {} {}
integral lintasan bentuk diferensial
\mathl{\omega}{} sepanjang lintasan $\varphi \circ \gamma$.

}
{} {}

\inputaufgabegibtloesung
{4}
{

Misalkan $M$ adalah suatu
\definitionsverweis {manifold berbatas}{}{.}
Tunjukkan bahwa setiap himpunan bagian terbuka

\mavergleichskette
{\vergleichskette
{ N
}
{ \subseteq }{ M
}
{  }{
}
{    }{
}
{   }{
}
}
{}{}{}
juga merupakan manifold dengan batas
\zusatzklammer {yang mungkin kosong} {} {},
dan bahwa

\mavergleichskettedisp
{\vergleichskette
{ \partial N
}
{ =} { \partial M \cap N
}
{ } {
}
{ } {
}
{ } {
}
}
{}{}{}
berlaku.

}
{} {}

\inputaufgabegibtloesung
{8}
{

Buktikan teorema tentang karakterisasi kurva geodesik terhadap koneksi Levi-Civita.

}
{} {}
\endgroup

\chapter{Formulir Peserta Ujian 6}
\noindent\textit{Formulir peserta dari sumber resmi; teks soal dipertahankan sesuai terjemahan yang telah diverifikasi.}
\begingroup
\renewcommand{\theHfakt}{o011.exam.learner.06.\arabic{fakt}}
%Data institusi

%\input{Dozentdaten}

%\renewcommand{\fachbereich}{Fachbereich}

%\renewcommand{\dozent}{Prof. Dr. . }

%Data ujian

\renewcommand{\klausurgebiet}{ }

\renewcommand{\klausurtyp}{ }

\renewcommand{\klausurdatum}{ . 20}

\klausurvorspann {\fachbereich} {\klausurdatum} {\dozent} {\klausurgebiet} {\klausurtyp}

%Data untuk tabel nilai berikut

\renewcommand{\aeins}{  3  }

\renewcommand{\azwei}{  3   }

\renewcommand{\adrei}{  7  }

\renewcommand{\avier}{  4  }

\renewcommand{\afuenf}{  0  }

\renewcommand{\asechs}{  0  }

\renewcommand{\asieben}{  0  }

\renewcommand{\aacht}{  5  }

\renewcommand{\aneun}{  15  }

\renewcommand{\azehn}{  5  }

\renewcommand{\aelf}{  0  }

\renewcommand{\azwoelf}{  3  }

\renewcommand{\adreizehn}{  3   }

\renewcommand{\avierzehn}{  0  }

\renewcommand{\afuenfzehn}{  10  }

\renewcommand{\asechzehn}{  4  }

\renewcommand{\asiebzehn}{  62  }

\renewcommand{\aachtzehn}{    }

\renewcommand{\aneunzehn}{    }

\renewcommand{\azwanzig}{    }

\renewcommand{\aeinundzwanzig}{    }

\renewcommand{\azweiundzwanzig}{    }

\renewcommand{\adreiundzwanzig}{    }

\renewcommand{\avierundzwanzig}{    }

\renewcommand{\afuenfundzwanzig}{    }

\renewcommand{\asechsundzwanzig}{    }

\punktetabellesechzehn

\klausurnote

\newpage

\setcounter{section}{K}

\inputaufgabegibtloesung
{3}
{

Definisikan istilah-istilah berikut
\zusatzklammer {yang dicetak miring} {} {}.
\aufzaehlungsechs{Suatu kurva
\stichwort {berparametrisasi panjang busur} {}
\maabbdisp {\gamma} {[a,b]} { \R^n
} {.}

}{Untuk suatu titik

\mavergleichskette
{\vergleichskette
{  P
}
{ \in }{ M
}
{  }{}
{    }{}
{   }{}
}
{}{}{}
pada suatu manifold diferensiabel $M$, yang dimaksud dengan dua kurva diferensiabel yang
\stichwort {ekuivalen secara tangensial} {} ialah
\maabbdisp {\gamma_1, \gamma_2} { I } { M
} {}
\zusatzklammer {di sini

\mavergleichskettek
{\vergleichskettek
{  0
}
{ \in }{  I
}
{  }{
}
{    }{
}
{   }{
}
}
{}{}{}
merupakan interval riil terbuka dan

\mavergleichskettek
{\vergleichskettek
{   \gamma_1(0)
}
{ = }{ \gamma_2(0)
}
{ = }{ P
}
{    }{
}
{   }{
}
}
{}{}{}} {} {.}

}{Suatu
\stichwort {penampang kontinu} {}
dari suatu pemetaan kontinu
\maabbdisp {p} {Y} {X
} {}
antara ruang-ruang topologis
\mathkor {} {X} {dan} {Y} {.}

}{\stichwort {Bentuk volume kanonik} {} pada suatu manifold Riemann berorientasi $M$.

}{\stichwort {Turunan eksterior} {}
dari suatu
\definitionsverweis {bentuk diferensial}{}{}
berderajat $k$ yang
\definitionsverweis {terdiferensialkan secara kontinu}{}{}

\mathl{\omega \in
{ \mathcal E }^{ k } (  U   )}{} pada suatu
\definitionsverweis {himpunan terbuka}{}{}

\mathl{U \subseteq \R^n}{.}

}{\stichwort {Kelengkungan seksional} {} yang berkaitan dengan vektor-vektor bebas linear

\mavergleichskette
{\vergleichskette
{  v,w
}
{ \in }{T_PM
}
{  }{
}
{    }{
}
{   }{
}
}
{}{}{}
pada suatu
\definitionsverweis {manifold Riemann}{}{}
$M$ di suatu titik

\mavergleichskette
{\vergleichskette
{ P
}
{ \in }{ M
}
{  }{
}
{    }{
}
{   }{
}
}
{}{}{.}
}

}
{} {}

\inputaufgabegibtloesung
{3}
{

Rumuskan teorema-teorema berikut.
\aufzaehlungdrei{\stichwort {Teorema tentang orientasi pada suatu hipermuka diferensiabel} {}

\mavergleichskette
{\vergleichskette
{  Y
}
{ \subseteq }{  \R^n
}
{  }{
}
{    }{
}
{   }{
}
}
{}{}{.}}{Rumus braket Lie medan-medan vektor pada suatu himpunan terbuka

\mavergleichskette
{\vergleichskette
{  U
}
{ \subseteq }{  \R^n
}
{  }{
}
{    }{
}
{   }{
}
}
{}{}{.}}{Teorema tentang
\stichwort {retraksi ke batas} {}
pada manifold.}

}
{} {}

\inputaufgabegibtloesung
{7}
{

Buktikan
\stichwort {teorema tentang penyelesaian persamaan diferensial biasa pada suatu hipermuka diferensiabel} {}

\mavergleichskette
{\vergleichskette
{ Y
}
{ \subseteq }{ \R^n
}
{  }{
}
{    }{
}
{   }{
}
}
{}{}{.}

}
{} {}

\inputaufgabegibtloesung
{4}
{

Misalkan
\maabb {h} { I } { V
} {}
sebuah kurva yang terdiferensialkan dua kali dalam
\definitionsverweis {ruang vektor Euklides}{}{}
$V$. Buktikan bahwa jika

\mavergleichskette
{\vergleichskette
{h'(t)
}
{ \neq }{ 0
}
{  }{
}
{    }{
}
{   }{
}
}
{}{}{,}
maka berlaku persamaan

\mavergleichskettedisp
{\vergleichskette
{ \Vert { h'(t)} \Vert'
}
{ =} { { \frac{   \left\langle h'(t) , h^{\prime \prime} (t)  \right\rangle  }{  \left\langle h'(t) , h'(t)  \right\rangle  } } \Vert { h'(t)} \Vert
}
{ } {
}
{ } {
}
{ } {
}
}
{}{}{.}

}
{} {}

\inputaufgabe
{0}
{

}
{} {}

\inputaufgabe
{0}
{

}
{} {}

\inputaufgabe
{0}
{

}
{} {}

\inputaufgabegibtloesung
{5}
{

Buktikan pernyataan bahwa ekuivalensi tangensial lintasan-lintasan pada suatu manifold di suatu titik $P$ dapat diperiksa menggunakan sebarang peta.

}
{} {}

\inputaufgabegibtloesung
{15 (3+3+4+5)}
{

Misalkan

\mavergleichskette
{\vergleichskette
{T
}
{ = }{S^1 \times S^1
}
{  }{
}
{    }{
}
{   }{
}
}
{}{}{}
merupakan torus dan

\mavergleichskette
{\vergleichskette
{S^2
}
{ \subset }{\R^3
}
{  }{
}
{    }{
}
{   }{
}
}
{}{}{}
merupakan
\definitionsverweis {permukaan bola satuan}{}{.}

a) Tunjukkan bahwa
\maabbeledisp {\varphi} {S^1 \times S^1} { S^2
} { \begin{pmatrix}  \alpha \\ \beta   \end{pmatrix} } { { \frac{   1 }{  2 } } \begin{pmatrix}  \sqrt{ 2-2 \sin \alpha   } \cos \beta  \\ \cos \alpha + \sin \beta \cos \alpha \\  1  + \sin \alpha -  \sin \beta + \sin \alpha \sin \beta   \end{pmatrix}
} {}
mendefinisikan suatu pemetaan kontinu.

b) Tunjukkan bahwa $\varphi$ surjektif.

c) Deskripsikan serat-serat $\varphi$.

d) Jelaskan pemetaan $\varphi$ dengan bantuan sebuah sketsa.

}
{} {}

\inputaufgabegibtloesung
{5}
{

Hitunglah integral lintasan
\mathl{\int_\gamma \omega}{} untuk
\maabbeledisp {\gamma} { [-1,0] } {\R^3
} { t } {(-t^2,t^3-1,t+2)
} {,}
dengan bentuk diferensial berderajat $1$

\mavergleichskettedisp
{\vergleichskette
{ \omega
}
{ =} { x^3dx -yzdy +xz^2 dz
}
{ } {
}
{ } {
}
{ } {
}
}
{}{}{}
pada $\R^3$.

}
{} {}

\inputaufgabe
{0}
{

}
{} {}

\inputaufgabegibtloesung
{3 (2+1)}
{

Misalkan
\maabbeledisp {f} {\R} {\R
} { t } {f(t)
} {,}
diberikan oleh

\mavergleichskettedisp
{\vergleichskette
{ f(t)
}
{ =} { { \frac{   \sin^{ 3  } (t^4)  }{  1+t^2 } }
}
{ } {
}
{ } {
}
{ } {
}
}
{}{}{}

a) Hitung turunan eksterior dari $f$.

b) Hitung turunan eksterior dari $fdt$.

}
{} {}

\inputaufgabegibtloesung
{3}
{

Berikan suatu contoh pemetaan yang terdiferensialkan secara kontinu
\maabbdisp {\varphi} { H } { H
} {}
dari suatu setengah bidang Euklides ke dirinya sendiri, dengan sifat bahwa tepat satu titik batas $H$ dipetakan ke suatu titik batas dan semua titik lainnya dipetakan ke interior setengah bidang tersebut.

}
{} {}

\inputaufgabe
{0}
{

}
{} {}

\inputaufgabegibtloesung
{10}
{

Buktikan versi balok dari teorema Stokes.

}
{} {}

\inputaufgabegibtloesung
{4}
{

Untuk kurva
\maabbeledisp {\psi} {\R} { {\mathbb H}
} { t } { \left(  t   , \,  e^t                   \right)
} {,}
yang bernilai dalam
\definitionsverweis {setengah bidang Poincaré}{}{}
${\mathbb H}$, tentukan
\definitionsverweis {simbol-simbol Christoffel}{}{}
untuk
\definitionsverweis {koneksi tarikan balik}{}{}
\zusatzklammer {dari
\definitionsverweis {koneksi Levi-Civita}{}{}
pada ${\mathbb H}$} {} {}
di atas $\R$.

}
{} {}
\endgroup

\chapter{Formulir Peserta Ujian 7}
\noindent\textit{Formulir peserta dari sumber resmi; teks soal dipertahankan sesuai terjemahan yang telah diverifikasi.}
\begingroup
\renewcommand{\theHfakt}{o011.exam.learner.07.\arabic{fakt}}
%Data institusi

%\input{Dozentdaten}

%\renewcommand{\fachbereich}{Fachbereich}

%\renewcommand{\dozent}{Prof. Dr. . }

%Data ujian

\renewcommand{\klausurgebiet}{ }

\renewcommand{\klausurtyp}{ }

\renewcommand{\klausurdatum}{ . 20}

\klausurvorspann {\fachbereich} {\klausurdatum} {\dozent} {\klausurgebiet} {\klausurtyp}

%Data untuk tabel nilai berikut

\renewcommand{\aeins}{  3  }

\renewcommand{\azwei}{  3   }

\renewcommand{\adrei}{  0  }

\renewcommand{\avier}{  4  }

\renewcommand{\afuenf}{  5  }

\renewcommand{\asechs}{  4  }

\renewcommand{\asieben}{  9  }

\renewcommand{\aacht}{  0  }

\renewcommand{\aneun}{  8  }

\renewcommand{\azehn}{  6  }

\renewcommand{\aelf}{  0  }

\renewcommand{\azwoelf}{  3  }

\renewcommand{\adreizehn}{  3   }

\renewcommand{\avierzehn}{  0  }

\renewcommand{\afuenfzehn}{  0  }

\renewcommand{\asechzehn}{  7  }

\renewcommand{\asiebzehn}{  55  }

\renewcommand{\aachtzehn}{    }

\renewcommand{\aneunzehn}{    }

\renewcommand{\azwanzig}{    }

\renewcommand{\aeinundzwanzig}{    }

\renewcommand{\azweiundzwanzig}{    }

\renewcommand{\adreiundzwanzig}{    }

\renewcommand{\avierundzwanzig}{    }

\renewcommand{\afuenfundzwanzig}{    }

\renewcommand{\asechsundzwanzig}{    }

\punktetabellesechzehn

\klausurnote

\newpage

\setcounter{section}{K}

\inputaufgabegibtloesung
{3}
{

Definisikan istilah-istilah berikut
\zusatzklammer {yang dicetak miring} {} {.}
\aufzaehlungsechs{
\stichwort {Kelengkungan Gauss} {}
dari suatu permukaan diferensiabel

\mavergleichskette
{\vergleichskette
{ Y
}
{ \subseteq }{ \R^3
}
{  }{
}
{    }{
}
{   }{
}
}
{}{}{}
di suatu titik

\mavergleichskette
{\vergleichskette
{ P
}
{ \in }{ Y
}
{  }{
}
{    }{
}
{   }{
}
}
{}{}{.}

}{Suatu
\stichwort {manifold topologis} {}
berdimensi $n$.

}{Suatu
\stichwort {medan vektor tak bergantung waktu} {}
pada manifold diferensiabel $M$.

}{Suatu
\stichwort {bentuk diferensial berderajat $k$} {}
pada manifold diferensiabel $M$.

}{Suatu \stichwort {manifold Riemann} {.}

}{Suatu koneksi linear yang
\stichwort {bebas torsi} {}
pada bundel tangen manifold Riemann $M$.
}

}
{} {}

\inputaufgabegibtloesung
{3}
{

Nyatakan teorema-teorema berikut.
\aufzaehlungdrei{
\stichwort {Teorema eksistensi transpor paralel pada hiperpermukaan} {.}}{Teorema tentang deskripsi lokal bentuk diferensial.}{
\stichwort {Versi balok} {}
dari teorema Stokes.}

}
{} {}

\inputaufgabe
{0}
{

}
{} {}

\inputaufgabegibtloesung
{4 (1+3)}
{

Misalkan $M$ elips yang diberikan oleh

\mavergleichskettedisp
{\vergleichskette
{ x^2+3y^2
}
{ =} { 2
}
{ } {
}
{ } {
}
{ } {
}
}
{}{}{}
dalam $\R^2$,

\mavergleichskette
{\vergleichskette
{ P
}
{ = }{ \begin{pmatrix}  1  \\ { \frac{   1  }{  \sqrt{3}  } }    \end{pmatrix}
}
{  }{
}
{    }{
}
{   }{
}
}
{}{}{}
dan

\mavergleichskette
{\vergleichskette
{ Q
}
{ = }{ \begin{pmatrix}  -1 \\ { \frac{   1  }{  \sqrt{3}  } }    \end{pmatrix}
}
{  }{
}
{    }{
}
{   }{
}
}
{}{}{}
adalah titik-titik pada $M$.
\aufzaehlungdrei{Tunjukkan bahwa

\mavergleichskettedisp
{\vergleichskette
{ v
}
{ =} { \begin{pmatrix}  -  \sqrt{3}  \\ 1    \end{pmatrix}
}
{ } {
}
{ } {
}
{ } {
}
}
{}{}{}
merupakan
\definitionsverweis {vektor tangen}{}{}
di $P$ pada $M$.
}{Tentukan citra $v$ di bawah suatu transpor paralel pada $M$ dari $P$ ke $Q$.
}{
}

}
{} {}

\inputaufgabegibtloesung
{5}
{

Buktikan teorema tentang orientasi-orientasi pada hiperpermukaan diferensiabel.

}
{} {}

\inputaufgabegibtloesung
{4}
{

Misalkan

\mavergleichskette
{\vergleichskette
{ Y
}
{ \subseteq }{ W
}
{ \subseteq }{ \R^3
}
{    }{
}
{   }{
}
}
{}{}{,}
$W$
\definitionsverweis {terbuka}{}{,}
dan suatu
\definitionsverweis {permukaan diferensiabel}{}{}
yang dilengkapi dengan
\definitionsverweis {medan normal satuan}{}{}
$N$. Misalkan
\maabb {\gamma} { I } { Y
} {}
suatu kurva terdiferensialkan secara kontinu dan
\maabbdisp {F} { I } { \R^3
} {}
suatu
\definitionsverweis {medan vektor tangensial}{}{}
diferensiabel sepanjang $\gamma$ yang
\definitionsverweis {paralel}{}{.}
Tunjukkan bahwa medan vektor
\maabbeledisp {} { I } { \R^3
} { t } { N(\gamma(t)) \times F(t)
} {,}
juga paralel, dengan $\times$ menyatakan
\definitionsverweis {hasil kali silang}{}{}
dalam $\R^3$.

}
{} {}

\inputaufgabegibtloesung
{9}
{

Misalkan

\mavergleichskette
{\vergleichskette
{ M
}
{ \neq }{ \emptyset
}
{  }{
}
{    }{
}
{   }{
}
}
{}{}{}
suatu
\definitionsverweis {manifold diferensiabel}{}{}
berdimensi $n$. Tunjukkan bahwa terdapat suatu rantai
\definitionsverweis {submanifold-submanifold tertutup}{}{}

\mavergleichskettedisp
{\vergleichskette
{ M_0
}
{ \subseteq} { M_1
}
{ \subseteq} { M_2
}
{ \subseteq  \ldots \subseteq} { M_{n-1}
}
{ \subseteq} { M_n
}
}
{
\vergleichskettefortsetzung
{ =} { M
}
{ } {}
{ } {}
{ } {}
}{}{}
sedemikian sehingga submanifold tertutup $M_i$ berdimensi $i$.

}
{} {}

\inputaufgabe
{0}
{

}
{} {}

\inputaufgabegibtloesung
{8}
{

Misalkan
\maabbdisp {\psi} {\R^n } {\R
} {}
suatu
\definitionsverweis {fungsi diferensiabel}{}{}
dan

\mavergleichskettedisphandlinks
{\vergleichskettedisphandlinks
{M
}
{ =} { \left\{  \left(  x_1   , \,   , \ldots ,  , \,   x_n , \,   \psi(x_1 , \ldots , x_n)                \right)   \mid  (x_1 , \ldots , x_n) \in \R^n         \right\}
}
{ \subseteq} { \R^n \times \R
}
{ } {
}
{ } {
}
}
{}{}{}
merupakan
\definitionsverweis {grafik}{}{}
dari $\psi$. Tunjukkan bahwa $M$ suatu
\definitionsverweis {manifold Riemann}{}{}
yang
\definitionsverweis {berorientasi}{}{}
dan bahwa untuk
\definitionsverweis {bentuk volume kanonik}{}{}
$\omega$ pada $M$ berlaku rumus

\mavergleichskettedisp
{\vergleichskette
{ (\operatorname{Id} \times \psi)^*(\omega)
}
{ =} { \left(  1+ \sum_{i = 1}^n \left(  { \frac{   \partial   \psi  }{  \partial   x_i   } }  \right)^2  \right)^{1/2} dx_1 \wedge \ldots \wedge dx_n
}
{ } {
}
{ } {
}
{ } {
}
}
{}{}{.}

}
{} {}

\inputaufgabe
{6}
{

Hitung luas permukaan bola satuan.

}
{} {}

\inputaufgabe
{0}
{

}
{} {}

\inputaufgabegibtloesung
{3}
{

Hitung
\definitionsverweis {turunan eksterior}{}{}
$d \omega$ dari
$1$-\definitionsverweis {bentuk diferensial}{}{}

\mavergleichskettedisp
{\vergleichskette
{ \omega
}
{ =} { { \frac{   x^2 }{  y } } dx- { \frac{   x  }{  y^2 } } dy
}
{ } {
}
{ } {
}
{ } {
}
}
{}{}{}
pada

\mavergleichskette
{\vergleichskette
{ U
}
{ = }{ \left\{ (x,y) \in \R^2  \mid   y \neq 0         \right\}
}
{  }{
}
{    }{
}
{   }{
}
}
{}{}{.}

}
{} {}

\inputaufgabegibtloesung
{3}
{

Berikan contoh suatu
\definitionsverweis {manifold berbatas}{}{}
yang
\definitionsverweis {terhubung}{}{,}
dan batasnya terdiri atas
\zusatzklammer {gabungan saling lepas dari} {} {}
tak hingga banyak lingkaran $S^1$.

}
{} {}

\inputaufgabe
{0}
{

}
{} {}

\inputaufgabe
{0}
{

}
{} {}

\inputaufgabegibtloesung
{7}
{

Misalkan $M$ suatu
\definitionsverweis {manifold Riemann}{}{}
dan diberikan suatu
\definitionsverweis {koneksi linear}{}{}
pada
\definitionsverweis {bundel tangen}{}{}
$TM$ yang
\definitionsverweis {metrik}{}{}
dan
\definitionsverweis {bebas torsi}{}{.}
Tunjukkan bahwa untuk medan-medan vektor $X,Y,Z$ berlaku rumus

\mavergleichskettedisp
{\vergleichskette
{ \left\langle \nabla_X Y , Z  \right\rangle
}
{ =} { { \frac{   1  }{  2 } }  \left(  D_X \left\langle  Y  ,  Z  \right\rangle +  D_Y \left\langle  Z  ,  X  \right\rangle -  D_Z \left\langle  X  ,  Y  \right\rangle  - \left\langle  Y  , [X,Z]   \right\rangle - \left\langle  Z  , [Y,X]   \right\rangle + \left\langle  X  , [Z,Y]   \right\rangle   \right)
}
{ } {
}
{ } {
}
{ } {
}
}
{}{}{.}

}
{} {}
\endgroup

\chapter{Formulir Peserta Ujian 8}
\noindent\textit{Formulir peserta dari sumber resmi; teks soal dipertahankan sesuai terjemahan yang telah diverifikasi.}
\begingroup
\renewcommand{\theHfakt}{o011.exam.learner.08.\arabic{fakt}}
%Data institusi

%\input{Dozentdaten}

%\renewcommand{\fachbereich}{Fachbereich}

%\renewcommand{\dozent}{Prof. Dr. . }

%Data ujian

\renewcommand{\klausurgebiet}{ }

\renewcommand{\klausurtyp}{ }

\renewcommand{\klausurdatum}{ . 20}

\klausurvorspann {\fachbereich} {\klausurdatum} {\dozent} {\klausurgebiet} {\klausurtyp}

%Data untuk tabel nilai berikut

\renewcommand{\aeins}{  3  }

\renewcommand{\azwei}{  3   }

\renewcommand{\adrei}{  3  }

\renewcommand{\avier}{  0  }

\renewcommand{\afuenf}{  2  }

\renewcommand{\asechs}{  0  }

\renewcommand{\asieben}{  3  }

\renewcommand{\aacht}{  7  }

\renewcommand{\aneun}{  0  }

\renewcommand{\azehn}{  7  }

\renewcommand{\aelf}{  4  }

\renewcommand{\azwoelf}{  3  }

\renewcommand{\adreizehn}{  3   }

\renewcommand{\avierzehn}{  7  }

\renewcommand{\afuenfzehn}{  2  }

\renewcommand{\asechzehn}{  9  }

\renewcommand{\asiebzehn}{  7  }

\renewcommand{\aachtzehn}{  63  }

\renewcommand{\aneunzehn}{    }

\renewcommand{\azwanzig}{    }

\renewcommand{\aeinundzwanzig}{    }

\renewcommand{\azweiundzwanzig}{    }

\renewcommand{\adreiundzwanzig}{    }

\renewcommand{\avierundzwanzig}{    }

\renewcommand{\afuenfundzwanzig}{    }

\renewcommand{\asechsundzwanzig}{    }

\punktetabellesiebzehn

\klausurnote

\newpage

\setcounter{section}{0}

\inputaufgabegibtloesung
{3}
{

Definisikan istilah-istilah berikut
\zusatzklammer {yang dicetak miring} {} {}.
\aufzaehlungsechs{\stichwort {Pemetaan Gauss} {}
dari suatu hiperpermukaan diferensiabel

\mavergleichskette
{\vergleichskette
{ Y
}
{ \subseteq }{ \R^n
}
{  }{
}
{    }{
}
{   }{
}
}
{}{}{.}

}{Suatu
\stichwort {peta topologis} {}
pada suatu
\definitionsverweis {manifold topologis}{}{}
$M$.

}{Suatu $C^1$-\stichwort {manifold diferensiabel} {} $M$.

}{\stichwort {Pangkat eksterior} {} ke-$n$ dari suatu
$K$-\definitionsverweis {ruang vektor}{}{}
$V$
\zusatzklammer {cukup tuliskan simbolnya} {} {.}

}{Suatu
\stichwort {isometri lokal} {}
\maabb {\varphi} {L} {M
} {}
antara manifold-manifold Riemann.

}{Suatu koneksi
\stichwort {terintegralkan secara lokal} {}
pada suatu bundel vektor diferensiabel
\maabb {p} {E} {X
} {}
di atas suatu manifold diferensiabel $X$.
}

}
{} {}

\inputaufgabegibtloesung
{3}
{

Rumuskan teorema-teorema berikut.
\aufzaehlungdrei{\stichwort {Teorema isometri untuk transpor paralel pada suatu hiperpermukaan} {}

\mavergleichskette
{\vergleichskette
{  Y
}
{ \subseteq }{  \R^n
}
{  }{
}
{    }{
}
{   }{
}
}
{}{}{.}}{\stichwort {Theorema egregium} {.}}{Teorema tentang partisi kesatuan.}

}
{} {}

\inputaufgabegibtloesung
{3}
{

Misalkan
\maabbdisp {f,g} {\R} { \R^n
} {}
\definitionsverweis {kurva-kurva diferensiabel}{}{.}
Hitung
\definitionsverweis {turunan}{}{}
dari fungsi
\maabbeledisp {} {\R} {\R
} { t } { \left\langle f(t) , g(t)  \right\rangle
} {.}
Rumuskan hasilnya dengan menggunakan hasil kali dalam.

}
{} {}

\inputaufgabe
{0}
{

}
{} {}

\inputaufgabegibtloesung
{2}
{

Misalkan

\mavergleichskette
{\vergleichskette
{ Y
}
{ \subseteq }{ W
}
{ \subseteq }{ \R^n
}
{    }{
}
{   }{
}
}
{}{}{,}
$W$
\definitionsverweis {terbuka}{}{,}
suatu
\definitionsverweis {hiperpermukaan diferensiabel}{}{},
dan
\maabb {\gamma} {I =[a,b]} {Y
} {}
suatu
\definitionsverweis {kurva geodesik}{}{.}
Buktikan bahwa
\definitionsverweis {transpor paralel}{}{}
sepanjang $\gamma$ memetakan vektor tangensial

\mavergleichskette
{\vergleichskette
{ \gamma'(a)
}
{ \in }{ T_{\gamma(a)}Y
}
{  }{
}
{    }{
}
{   }{
}
}
{}{}{}
ke vektor tangensial

\mavergleichskette
{\vergleichskette
{ \gamma'(b)
}
{ \in }{ T_{\gamma(b)}Y
}
{  }{
}
{    }{
}
{   }{
}
}
{}{}{.}

}
{} {}

\inputaufgabe
{0}
{

}
{} {}

\inputaufgabegibtloesung
{3}
{

Apakah pemetaan
\maabbeledisp {\varphi} {S^1} {S^1
} {(x,y)} {( \betrag {  x  } ,y)
} {,}
\definitionsverweis {terdiferensialkan}{}{?}

}
{} {}

\inputaufgabegibtloesung
{7}
{

Misalkan $M$ suatu manifold diferensiabel kompak dengan bentuk volume positif kontinu $\omega$. Buktikan bahwa

\mavergleichskettedisp
{\vergleichskette
{ \int_M \omega
}
{ <} { \infty
}
{ } {
}
{ } {
}
{ } {
}
}
{}{}{}.

}
{} {}

\inputaufgabe
{0}
{

}
{} {}

\inputaufgabegibtloesung
{7 (1+2+4)}
{

\aufzaehlungdreiabc{ Berikan suatu sifat topologis yang menunjukkan bahwa
\definitionsverweis {interval satuan terbuka}{}{}
dan
\definitionsverweis {lingkaran}{}{}
$S^1$ tidak
\definitionsverweis {homeomorfik}{}{.}
}{Buktikan bahwa setiap
$1$-\definitionsverweis {bentuk diferensial}{}{}
kontinu pada
\mathl{]0,1[}{}
bersifat
\definitionsverweis {eksak}{}{.}
}{ Berikan suatu
$1$-bentuk konkret pada $S^1$ yang
\definitionsverweis {tertutup}{}{}
tetapi tidak eksak.
}

}
{} {}

\inputaufgabegibtloesung
{4}
{

Misalkan

\mavergleichskette
{\vergleichskette
{ U
}
{ \subseteq }{ \R^n
}
{  }{
}
{    }{
}
{   }{
}
}
{}{}{}
\definitionsverweis {terbuka}{}{}
dan
\maabb {f} { U } { \R
} {}
suatu
\definitionsverweis {fungsi yang terdiferensialkan secara kontinu}{}{}
dua kali. Buktikan bahwa

\mavergleichskettedisp
{\vergleichskette
{ d(df)
}
{ =} { 0
}
{ } {
}
{ } {
}
{ } {
}
}
{}{}{,}
dengan $d$ menyatakan
\definitionsverweis {turunan eksterior}{}{.}

}
{} {}

\inputaufgabegibtloesung
{3}
{

Deskripsikan
\definitionsverweis {pemetaan}{}{}
\maabbeledisp {\psi} { \Complex \setminus \{1\} } { \Complex
} {w} { { \frac{   (w+1)  { \mathrm i}  }{ 1-w  } }
} {}
dalam koordinat riil.

}
{} {}

\inputaufgabegibtloesung
{3}
{

Misalkan

\mavergleichskette
{\vergleichskette
{ H
}
{ \subseteq }{ \R^n
}
{  }{
}
{    }{
}
{   }{
}
}
{}{}{}
suatu
\definitionsverweis {setengah ruang Euklides}{}{}
dan

\mavergleichskette
{\vergleichskette
{ P
}
{ \in }{ H
}
{  }{
}
{    }{
}
{   }{
}
}
{}{}{.}
Andaikan terdapat suatu himpunan yang terbuka di $\R^n$, yaitu $V$, dengan

\mavergleichskette
{\vergleichskette
{ P
}
{ \in }{ V
}
{ \subseteq }{ H
}
{    }{
}
{   }{
}
}
{}{}{.}
Buktikan bahwa
\mathl{P}{} bukan titik batas dari $H$.

}
{} {}

\inputaufgabegibtloesung
{7}
{

Misalkan $M$ suatu
\definitionsverweis {manifold diferensiabel}{}{}
berdimensi $n$ dan

\mavergleichskette
{\vergleichskette
{ P
}
{ \in }{ M
}
{  }{
}
{    }{
}
{   }{
}
}
{}{}{}
suatu titik. Buktikan bahwa terdapat suatu
\definitionsverweis {manifold dengan batas}{}{}
$N$ yang batasnya $\partial N$ difeomorfik dengan permukaan bola berdimensi $(n-1)$, yaitu $S^{n-1}$, sedemikian sehingga
\mathkor {} {M \setminus \{ P \}} {dan} {N \setminus \partial N} {}
saling difeomorfik.

}
{} {}

\inputaufgabegibtloesung
{2}
{

Berikan suatu
\definitionsverweis {penghabisan kompak}{}{}
untuk $\R^d$.

}
{} {}

\inputaufgabegibtloesung
{9}
{

Buktikan teorema titik tetap Brouwer.

}
{} {}

\inputaufgabegibtloesung
{7 (1+2+4)}
{

Pada
\definitionsverweis {bundel vektor trivial}{}{}

\mathl{\R \times \R^2}{} di atas $\R$, kita meninjau
\definitionsverweis {koneksi trivial}{}{}
$\nabla$.
\aufzaehlungdrei{Buktikan bahwa

\mavergleichskette
{\vergleichskette
{ f_1(t)
}
{ = }{ \left(  1+ t^2  , \,  t^3                   \right)
}
{  }{
}
{    }{
}
{   }{
}
}
{}{}{}
dan

\mavergleichskette
{\vergleichskette
{ f_2(t)
}
{ = }{ \left(  t  , \,   1+t^2                   \right)
}
{  }{
}
{    }{
}
{   }{
}
}
{}{}{}
merupakan
\definitionsverweis {penampang basis}{}{}
dari bundel vektor tersebut.
}{Nyatakan penampang basis standar $e_1,e_2$ sebagai kombinasi linear dari penampang basis $f_1,f_2$.
}{Tentukan
\definitionsverweis {simbol-simbol Christoffel}{}{}
dari $\nabla$ terhadap penampang-penampang basis tersebut.
}

}
{} {}
\endgroup

\chapter{Formulir Peserta Ujian 9}
\noindent\textit{Formulir peserta dari sumber resmi; teks soal dipertahankan sesuai terjemahan yang telah diverifikasi.}
\begingroup
\renewcommand{\theHfakt}{o011.exam.learner.09.\arabic{fakt}}
%Data institusi

%\input{Dozentdaten}

%\renewcommand{\fachbereich}{Fachbereich}

%\renewcommand{\dozent}{Prof. Dr. . }

%Data ujian

\renewcommand{\klausurgebiet}{ }

\renewcommand{\klausurtyp}{ }

\renewcommand{\klausurdatum}{ . 20}

\klausurvorspann {\fachbereich} {\klausurdatum} {\dozent} {\klausurgebiet} {\klausurtyp}

%Data untuk tabel nilai berikut

\renewcommand{\aeins}{  3  }

\renewcommand{\azwei}{  3   }

\renewcommand{\adrei}{  6  }

\renewcommand{\avier}{  3  }

\renewcommand{\afuenf}{  3  }

\renewcommand{\asechs}{  5  }

\renewcommand{\asieben}{  13  }

\renewcommand{\aacht}{  4  }

\renewcommand{\aneun}{  3  }

\renewcommand{\azehn}{  7  }

\renewcommand{\aelf}{  6  }

\renewcommand{\azwoelf}{  3  }

\renewcommand{\adreizehn}{  6   }

\renewcommand{\avierzehn}{  65  }

\renewcommand{\afuenfzehn}{    }

\renewcommand{\asechzehn}{    }

\renewcommand{\asiebzehn}{    }

\renewcommand{\aachtzehn}{    }

\renewcommand{\aneunzehn}{    }

\renewcommand{\azwanzig}{    }

\renewcommand{\aeinundzwanzig}{    }

\renewcommand{\azweiundzwanzig}{    }

\renewcommand{\adreiundzwanzig}{    }

\renewcommand{\avierundzwanzig}{    }

\renewcommand{\afuenfundzwanzig}{    }

\renewcommand{\asechsundzwanzig}{    }

\punktetabelledreizehn

\klausurnote

\newpage

\setcounter{section}{0}

\inputaufgabegibtloesung
{3}
{

Definisikan istilah-istilah berikut
\zusatzklammer {yang dicetak miring} {} {.}
\aufzaehlungsechs{Suatu
\stichwort {ruang Hausdorff} {}
$X$.

}{Suatu
\stichwort {pemetaan diferensiabel} {}
\maabbdisp {\varphi} { L } { M
} {}
antara dua manifold diferensiabel
\mathkor {} {L} {dan} {M} {.}

}{
\stichwort {Bundel tangen} {}
dari suatu manifold diferensiabel $M$.

}{
\stichwort {Matriks fundamental kedua} {}
dari parametrisasi lokal
\maabbdisp {} {V} {U \subseteq Y
} {}
\zusatzklammer {dengan

\mavergleichskette
{\vergleichskette
{ V
}
{ \subseteq }{ \R^{n-1}
}
{  }{
}
{    }{
}
{   }{
}
}
{}{}{}
terbuka} {} {}
suatu hiperpermukaan diferensiabel

\mavergleichskette
{\vergleichskette
{ Y
}
{ \subseteq }{ \R^n
}
{  }{
}
{    }{
}
{   }{
}
}
{}{}{.}

}{
\stichwort {Turunan eksterior} {} dari suatu bentuk diferensial terdiferensialkan secara kontinu
\mathl{\omega \in
{ \mathcal E }^{ k } (  M   )}{} pada manifold diferensiabel $M$.

}{Suatu \stichwort {partisi kesatuan yang tersubordinasi pada penutup} {,} dengan

\mavergleichskette
{\vergleichskette
{  X
}
{ = }{ \bigcup_{i \in I} W_i
}
{  }{
}
{    }{
}
{   }{
}
}
{}{}{}
suatu
\definitionsverweis {penutup terbuka}{}{}
dari
\definitionsverweis {ruang topologis}{}{}
$X$.
}

}
{} {}

\inputaufgabegibtloesung
{3}
{

Nyatakan teorema-teorema berikut.
\aufzaehlungdrei{
\stichwort {Teorema tentang sifat aljabar pemetaan Weingarten} {.}}{Teorema tentang keterorientasian dan bentuk volume positif.}{
\stichwort {Teorema titik tetap Brouwer} {.}}

}
{} {}

\inputaufgabegibtloesung
{6}
{

Tentukan
\definitionsverweis {kelengkungan}{}{}
di setiap titik lingkaran satuan yang diberikan secara implisit oleh

\mavergleichskettedisp
{\vergleichskette
{ h(x,y)
}
{ =} { x^2+y^2
}
{ =} { 1
}
{ } {
}
{ } {
}
}
{}{}{}
dengan
Lemma 3.11 (Geometri Diferensial (Osnabrück 2023))
dan medan gradien $h$.

}
{} {}

\inputaufgabegibtloesung
{3}
{

Misalkan $X$ suatu
\definitionsverweis {ruang Hausdorff}{}{.}
Tunjukkan bahwa
\definitionsverweis {diagonal}{}{}

\mavergleichskettedisp
{\vergleichskette
{ \triangle
}
{ =} { \left\{ (x,y) \in X \times X  \mid   x = y        \right\}
}
{ } {
}
{ } {
}
{ } {
}
}
{}{}{}
merupakan
\definitionsverweis {subhimpunan tertutup}{}{}
dalam
\definitionsverweis {ruang produk}{}{}

\mathl{X \times X}{}.

}
{} {}

\inputaufgabegibtloesung
{3}
{

Tunjukkan bahwa ketiga manifold satu dimensi
\mathlistdisp {S^1 = \left\{ (x,y) \in \R^2  \mid   \betrag { (x,y) } = 1         \right\}} {} {\R} {dan} {\R \setminus \{0\}} {}
tidak ada dua di antaranya yang homeomorfik.

}
{} {}

\inputaufgabegibtloesung
{5}
{

Misalkan

\mavergleichskette
{\vergleichskette
{ G
}
{ \subseteq }{ \R^m
}
{  }{
}
{    }{
}
{   }{
}
}
{}{}{}
\definitionsverweis {terbuka}{}{}
dan misalkan

\mavergleichskette
{\vergleichskette
{ M
}
{ \subseteq }{ G
}
{  }{
}
{    }{
}
{   }{
}
}
{}{}{}
suatu
$n$-\definitionsverweis {dimensional}{}{}
\definitionsverweis {submanifold tertutup}{}{,}
yang
\definitionsverweis {berorientasi}{}{}
dan dilengkapi dengan struktur Riemann terinduksi serta
\definitionsverweis {bentuk volume kanonik}{}{}
$\omega$. Misalkan

\mavergleichskette
{\vergleichskette
{ W
}
{ \subseteq }{ \R^n
}
{  }{
}
{    }{
}
{   }{
}
}
{}{}{}
terbuka dan misalkan
\maabbdisp {\varphi} { W } { U
} {}
suatu
\definitionsverweis {difeomorfisme}{}{}
dengan himpunan terbuka

\mavergleichskette
{\vergleichskette
{ U
}
{ \subseteq }{ M
}
{  }{
}
{    }{
}
{   }{
}
}
{}{}.
Tunjukkan bahwa pada $W$ berlaku rumus

\mavergleichskettealignhandlinks
{\vergleichskettealignhandlinks
{\varphi^* \left(  \omega \vert_U  \right)
}
{ =} { \left( \det  \left( \left\langle  \begin{pmatrix}  { \frac{   \partial  \varphi_1   }{  \partial   x_i   } }  \\ \vdots \\  { \frac{   \partial  \varphi_m   }{  \partial   x_i   } }   \end{pmatrix}  ,  \begin{pmatrix}  { \frac{   \partial  \varphi_1   }{  \partial   x_j   } }  \\ \vdots \\  { \frac{   \partial  \varphi_m   }{  \partial   x_j   } }   \end{pmatrix}   \right\rangle \right)_{1 \leq i,j \leq n} \right)^{1/2}  dx_1 \wedge \ldots \wedge dx_n
}
{ =} { \left( \det  \left(  { \frac{   \partial  \varphi_1   }{  \partial   x_i   } } { \frac{   \partial  \varphi_1   }{  \partial   x_j   } }  + \cdots +  { \frac{   \partial  \varphi_m   }{  \partial   x_i   } } { \frac{   \partial  \varphi_m   }{  \partial   x_j   } }  \right)_{1 \leq i,j \leq n} \right)^{1/2}  dx_1 \wedge \ldots \wedge dx_n
}
{ } {
}
{ } {
}
}
{}
{}{}{.}

}
{} {}

\inputaufgabegibtloesung
{13 (2+3+2+2+2+2)}
{

Kita tinjau

\mavergleichskette
{\vergleichskette
{ \R^6
}
{ = }{ \left(  \R^2  \right)^3
}
{  }{
}
{    }{
}
{   }{
}
}
{}{}{}
sebagai himpunan semua segitiga
\zusatzklammer {termasuk yang degenerat} {} {}
dengan mengidentifikasi segitiga yang mempunyai titik-titik sudut
\zusatzklammer {terurut} {} {}

\mavergleichskette
{\vergleichskette
{ P_1
}
{ = }{ \begin{pmatrix}  x_1  \\ y_1    \end{pmatrix}
}
{  }{
}
{    }{
}
{   }{
}
}
{}{}{,}

\mavergleichskette
{\vergleichskette
{ P_2
}
{ = }{ \begin{pmatrix}  x_2  \\ y_2    \end{pmatrix}
}
{  }{
}
{    }{
}
{   }{
}
}
{}{}{}
dan

\mavergleichskette
{\vergleichskette
{ P_3
}
{ = }{ \begin{pmatrix}  x_3  \\ y_3    \end{pmatrix}
}
{  }{
}
{    }{
}
{   }{
}
}
{}{}{,}
dengan tupel koordinat

\mathdisp {(x_1,y_1,x_2,y_2,x_3,y_3)} {  }

.
\aufzaehlungsechs{Tunjukkan bahwa himpunan segitiga dengan dua titik sudut berimpit merupakan subhimpunan tertutup $\R^6$
\zusatzklammer {jadi komplemennya merupakan himpunan terbuka dalam $\R^6$ yang kita namai $U$} {} {.}
}{Tunjukkan bahwa himpunan segitiga dengan ketiga titik sudut segaris merupakan subhimpunan tertutup $\R^6$
\zusatzklammer {jadi komplemennya, yang terdiri atas semua segitiga tak degenerat, merupakan himpunan terbuka dalam $\R^6$ yang kita namai $V$} {} {.}
}{Buat aturan fungsi yang mendeskripsikan pemetaan
\maabbdisp {F} {\R^6} {\R
} {}
yang memetakan suatu segitiga ke kelilingnya.
}{Tunjukkan bahwa fungsi $F$ dari bagian (3) terdiferensialkan secara kontinu pada himpunan $U$.
}{Hitung turunan parsial $F$ terhadap $x_1$ pada $U$.
}{Tunjukkan bahwa himpunan segitiga tak degenerat berkeliling $1$ membentuk
\definitionsverweis {submanifold tertutup}{}{}
dari $V$. Berapakah dimensinya?
}

}
{} {}

\inputaufgabegibtloesung
{4}
{

Misalkan
\mathkor {} {r} {dan} {R} {}
bilangan-bilangan real positif dengan

\mavergleichskette
{\vergleichskette
{ 0
}
{ < }{ r
}
{ < }{ R
}
{    }{
}
{   }{
}
}
{}{}{,}
dan tinjau torus
\zusatzklammer {tertanam} {} {}

\mavergleichskettedisp
{\vergleichskette
{ T
}
{ =} { \left\{ (x,y,z) \in \R^3  \mid   \left(  \sqrt{x^2+y^2} -R  \right)^2 +z^2  = r^2         \right\}
}
{ \subseteq} { \R^3
}
{ } {
}
{ } {
}
}
{}{}{}
dengan
\definitionsverweis {struktur Riemann terinduksi}{}{.}
Tunjukkan bahwa untuk setiap

\mavergleichskette
{\vergleichskette
{ \alpha
}
{ \in }{ [0,2 \pi [
}
{  }{
}
{    }{
}
{   }{
}
}
{}{}{}
pemetaan
\maabbeledisp {\Psi_\alpha} {\R^3} {\R^3
} { \left(  x   , \,   y  , \,  z                 \right) } { \left(  x \cos \alpha - y \sin  \alpha   , \,   x \sin \alpha + y \cos  \alpha  , \,  z                 \right)
} {,}
menginduksi suatu
\definitionsverweis {isometri}{}{}
dari $T$ ke dirinya sendiri.

}
{} {}

\inputaufgabe
{3}
{

Deskripsikan suatu model untuk bidang hiperbolik.

}
{} {}

\inputaufgabegibtloesung
{7}
{

Buktikan Theorema egregium.

}
{} {}

\inputaufgabegibtloesung
{6 (1+2+2+1)}
{

Kita tinjau
\definitionsverweis {bentuk diferensial}{}{}

\mavergleichskettedisp
{\vergleichskette
{ \omega
}
{ =} { ydx+zdy +xdz
}
{ } {
}
{ } {
}
{ } {
}
}
{}{}{}
pada $\R^3$ dan pemetaan
\maabbeledisp {\varphi} {\R^2} {\R^3
} { \left(  u   , \,  v                   \right) } { \left( u^2  , \,  v^2  , \,  uv                 \right)
} {.}
\aufzaehlungvier{Hitung turunan eksterior $\omega$.
}{Hitung tarikan balik $\varphi^* \omega$ dari $\omega$ di bawah $\varphi$.
}{Hitung turunan eksterior $\varphi^* \omega$ pada $\R^2$.
}{Hitung tarikan balik $\varphi^* d \omega$ dari $d\omega$ di bawah $\varphi$ secara independen dari bagian (3).
}

}
{} {}

\inputaufgabegibtloesung
{3}
{

Misalkan $M$ suatu
\definitionsverweis {manifold diferensiabel}{}{}
dan $\omega$ suatu
\definitionsverweis {bentuk diferensial}{}{}
yang
\definitionsverweis {eksak}{}{}
pada $M$. Misalkan $S$ suatu
\definitionsverweis {manifold diferensiabel}{}{}
yang
\definitionsverweis {kompak}{}{}
dan
\definitionsverweis {berorientasi}{}{}
\zusatzklammer {tanpa batas} {} {}
dengan
\definitionsverweis {basis terhitung bagi topologinya}{}{}
dan misalkan
\maabbdisp {\varphi} { S } { M
} {}
suatu
\definitionsverweis {pemetaan terdiferensialkan secara kontinu}{}{.}
Tunjukkan bahwa

\mavergleichskettedisp
{\vergleichskette
{
\int_{  S  }   \varphi^*\omega
}
{ =} { 0
}
{ } {
}
{ } {
}
{ } {
}
}
{}{}{.}

}
{} {}

\inputaufgabegibtloesung
{6}
{

Buktikan teorema tentang retraksi ke batas pada manifold.

}
{} {}
\endgroup

\chapter{Formulir Peserta Ujian 10}
\noindent\textit{Formulir peserta dari sumber resmi; teks soal dipertahankan sesuai terjemahan yang telah diverifikasi.}
\begingroup
\renewcommand{\theHfakt}{o011.exam.learner.10.\arabic{fakt}}
%Data institusi

%\input{Dozentdaten}

%\renewcommand{\fachbereich}{Bidang studi}

%\renewcommand{\dozent}{Prof. Dr. . }

%Data ujian

\renewcommand{\klausurgebiet}{ }

\renewcommand{\klausurtyp}{ }

\renewcommand{\klausurdatum}{. 20}

\klausurvorspann {\fachbereich} {\klausurdatum} {\dozent} {\klausurgebiet} {\klausurtyp}

%Data untuk tabel poin berikut

\renewcommand{\aeins}{  3  }

\renewcommand{\azwei}{  3   }

\renewcommand{\adrei}{  6  }

\renewcommand{\avier}{  4  }

\renewcommand{\afuenf}{  5  }

\renewcommand{\asechs}{  0  }

\renewcommand{\asieben}{  5  }

\renewcommand{\aacht}{  6  }

\renewcommand{\aneun}{  6  }

\renewcommand{\azehn}{  0  }

\renewcommand{\aelf}{  6  }

\renewcommand{\azwoelf}{  12  }

\renewcommand{\adreizehn}{  0   }

\renewcommand{\avierzehn}{  9  }

\renewcommand{\afuenfzehn}{  65  }

\renewcommand{\asechzehn}{    }

\renewcommand{\asiebzehn}{    }

\renewcommand{\aachtzehn}{    }

\renewcommand{\aneunzehn}{    }

\renewcommand{\azwanzig}{    }

\renewcommand{\aeinundzwanzig}{    }

\renewcommand{\azweiundzwanzig}{    }

\renewcommand{\adreiundzwanzig}{    }

\renewcommand{\avierundzwanzig}{    }

\renewcommand{\afuenfundzwanzig}{    }

\renewcommand{\asechsundzwanzig}{    }

\punktetabellevierzehn

\klausurnote

\newpage

\setcounter{section}{0}

\inputaufgabegibtloesung
{3}
{

Definisikan istilah-istilah berikut
\zusatzklammer {yang dicetak miring} {} {}.
\aufzaehlungsechs{Suatu
\stichwort {kelengkungan utama} {}
pada suatu permukaan diferensiabel

\mavergleichskette
{\vergleichskette
{ Y
}
{ \subseteq }{ \R^3
}
{  }{
}
{    }{
}
{   }{
}
}
{}{}{}
pada titik

\mavergleichskette
{\vergleichskette
{ P
}
{ \in }{ Y
}
{  }{
}
{    }{
}
{   }{
}
}
{}{}{.}

}{Suatu
\stichwort {difeomorfisme} {}
antara manifold diferensiabel
\mathkor {} {L} {dan} {M} {.}

}{Suatu
\stichwort {vektor singgung} {}
di titik

\mavergleichskette
{\vergleichskette
{  P
}
{ \in }{  M
}
{  }{
}
{    }{
}
{   }{
}
}
{}{}{}
dari manifold diferensiabel $M$.

}{Suatu
\stichwort {operator diferensial orde pertama} {}
pada manifold diferensiabel $M$.

}{Suatu
\stichwort {setengah ruang Euklides} {.}

}{Suatu
\definitionsverweis {penghabisan kompak}{}{}
bagi suatu
\definitionsverweis {ruang topologis}{}{}
$X$.
}

}
{} {}

\inputaufgabegibtloesung
{3}
{

Nyatakan teorema-teorema berikut.
\aufzaehlungdrei{Teorema tentang
\definitionsverweis {ruang singgung}{}{}
dari suatu
\definitionsverweis {submanifold tertutup}{}{}

\mavergleichskette
{\vergleichskette
{  M
}
{ \subseteq }{  N
}
{  }{
}
{    }{
}
{   }{
}
}
{}{}{.}}{Teorema tentang bentuk eksplisit dari bentuk diferensial tarik balik $\varphi^* \omega$ oleh pemetaan diferensiabel
\maabbdisp {\varphi} { U } { V
} {}
\zusatzklammer {dengan
\definitionsverweis {himpunan bagian terbuka}{}{}
\mathkor {} {U \subseteq \R^n} {dan} {V \subseteq \R^m} {,}
yang koordinatnya masing-masing dinyatakan dengan
\mathl{x_1  , \ldots , x_n}{} dan
\mathl{y_1  , \ldots , y_m}{}} {} {}
dari suatu
$k$-\definitionsverweis {bentuk diferensial}{}{}
pada $V$ dengan representasi

\mavergleichskettedisp
{\vergleichskette
{  \omega
}
{ =} {\sum_{I,\,  { \# \left(   I  \right) } = k}  f_I  dy_I
}
{ } {
}
{ } {
}
{ } {
}
}
{}{}{,}
dengan
\maabb {f_I} { V } {\R
} {}
adalah
\definitionsverweis {fungsi}{}{.}}{Teorema tentang karakterisasi kurva geodesik terhadap koneksi Levi-Civita pada manifold Riemann $M$.}

}
{} {}

\inputaufgabegibtloesung
{6 (2+4)}
{

Misalkan

\mavergleichskette
{\vergleichskette
{ \alpha,\beta
}
{ \in }{ \R_+
}
{  }{
}
{    }{
}
{   }{
}
}
{}{}{}
dan misalkan

\mavergleichskettedisp
{\vergleichskette
{ Y
}
{ =} { \left\{ (x,y)\in\R^2  \mid   \alpha x^2+\beta y^2 = 1        \right\}
}
{ } {
}
{ } {
}
{ } {
}
}
{}{}{.}
\aufzaehlungzwei {Tentukan
\definitionsverweis {pemetaan Gauss}{}{}
untuk $Y$.
} { Berikan secara eksplisit suatu
\definitionsverweis {pemetaan invers}{}{}
dari pemetaan Gauss untuk $Y$.
}

}
{} {}

\inputaufgabe
{4}
{

Jelaskan pasangan istilah \anfuehrung{ekstrinsik dan intrinsik}{} dalam konteks manifold.

}
{} {}

\inputaufgabegibtloesung
{5}
{

Buktikan rumus luas permukaan suatu benda putar yang dihasilkan oleh
\definitionsverweis {kurva diferensiabel}{}{}
\maabbeledisp {\gamma} {]a,b[} {\R^2
} { t } {(x(t),y(t))
} {,}
dengan

\mavergleichskette
{\vergleichskette
{ y(t)
}
{ > }{ 0
}
{  }{
}
{    }{
}
{   }{
}
}
{}{}{.}

}
{} {}

\inputaufgabe
{0}
{

}
{} {}

\inputaufgabegibtloesung
{5}
{

Misalkan $X$ adalah suatu
\definitionsverweis {torus}{}{.}
Berikan suatu
\definitionsverweis {pemetaan diferensiabel}{}{}
surjektif
\maabbdisp {\varphi} {\R^2} {X
} {}
sedemikian sehingga
\definitionsverweis {pemetaan singgung}{}{}
\maabbdisp {T_P(\varphi)} { T_P\R^2 } { T_{\varphi(P)}X
} {}
juga surjektif di setiap titik

\mavergleichskette
{\vergleichskette
{ P
}
{ \in }{ \R^2
}
{  }{
}
{    }{
}
{   }{
}
}
{}{}{.}

}
{} {}

\inputaufgabegibtloesung
{6 (2+2+2)}
{

Tinjau grafik $M$ dari pemetaan
\maabbeledisp {\varphi} {\R^2} {\R
} {(u,v)} { u^2+uv-v^3
} {,}
sebagai submanifold tertutup berdimensi dua di $\R^3$, yaitu

\mavergleichskettedisp
{\vergleichskette
{ M
}
{ =} { \left\{ (u,v,u^2+uv-v^3)  \mid  (u,v) \in \R^2         \right\}
}
{ } {
}
{ } {
}
{ } {
}
}
{}{}{}
dengan metrik Riemann yang diinduksi dari $\R^3$. Misalkan
\maabbeledisp {\psi} {\R^2} { M
} {(u,v)} { (u,v,u^2+uv-v^3)
} {,}
adalah difeomorfisme yang bersesuaian.

a) Tentukan diferensial total dari $\psi$ serta vektor-vektor citra
\mathl{T_P(\psi) (e_1)}{} dan
\mathl{T_P(\psi) (e_2)}{} di
\mathl{T_{\psi(P)}M}{.}

b) Untuk setiap titik berbentuk

\mavergleichskette
{\vergleichskette
{ P
}
{ = }{ (u,0)
}
{  }{
}
{    }{
}
{   }{
}
}
{}{}{}
tentukan luas jajaran genjang yang direntang oleh
\mathl{T_P(\psi) (e_1)}{} dan
\mathl{T_P(\psi) (e_2)}{} di
\mathl{T_{\psi(P)}M}{}.

c) Untuk setiap titik berbentuk

\mavergleichskette
{\vergleichskette
{ P
}
{ = }{ (0,v)
}
{  }{
}
{    }{
}
{   }{
}
}
{}{}{}
tentukan luas jajaran genjang yang direntang oleh
\mathl{T_P(\psi) (e_1)}{} dan
\mathl{T_P(\psi) (e_2)}{} di
\mathl{T_{\psi(P)}M}{}.

}
{} {}

\inputaufgabegibtloesung
{6}
{

Buktikan aturan hasil kali untuk bentuk diferensial yang diferensiabel pada suatu himpunan terbuka di $\R^n$.

}
{} {}

\inputaufgabe
{0}
{

}
{} {}

\inputaufgabegibtloesung
{6}
{

Misalkan $M$ adalah suatu
\definitionsverweis {manifold berbatas}{}{}
dan misalkan

\mavergleichskette
{\vergleichskette
{ P
}
{ \in }{ \partial M
}
{  }{
}
{    }{
}
{   }{
}
}
{}{}{}
adalah titik batas. Tunjukkan bahwa sifat-sifat berikut ekuivalen bagi suatu vektor singgung

\mavergleichskette
{\vergleichskette
{ v
}
{ \in }{ T_PM
}
{  }{
}
{    }{
}
{   }{
}
}
{}{}{.}
\aufzaehlungzwei {Terdapat suatu lintasan yang terdiferensialkan secara kontinu
\maabb {\gamma} { [- \epsilon, 0]} { M
} {}
dengan

\mavergleichskette
{\vergleichskette
{ \gamma(0)
}
{ = }{ P
}
{  }{
}
{    }{
}
{   }{
}
}
{}{}{}
dan

\mavergleichskette
{\vergleichskette
{ \gamma'(0)
}
{ = }{ v
}
{  }{
}
{    }{
}
{   }{
}
}
{}{}{.}
} { Pada setiap peta bermodel setengah ruang negatif, $v$ dipetakan oleh pemetaan singgung ke setengah ruang kanan.
}

}
{} {}

\inputaufgabegibtloesung
{12}
{

Buktikan teorema Stokes untuk manifold berbatas.

}
{} {}

\inputaufgabe
{0}
{

}
{} {}

\inputaufgabegibtloesung
{9 (1+1+3+2+2)}
{

Tinjau parametrisasi trigonometrik
\maabbeledisp {\psi} {\R } { S^1
} { t } { \left(  \cos  t   , \,   \sin  t                   \right)
} {}
dari lingkaran satuan. Misalkan

\mavergleichskette
{\vergleichskette
{ c
}
{ \in }{ [0, 2 \pi[
}
{  }{
}
{    }{
}
{   }{
}
}
{}{}{}
dipilih tetap. Pada bundel vektor trivial
\maabbdisp {} {S^1 \times \R^2} { S^1
} {}
diberikan suatu
\definitionsverweis {koneksi linear}{}{}
$\nabla$ sedemikian sehingga, sepanjang $\psi$,
\definitionsverweis {koneksi tarik balik}{}{}
$\psi^* \nabla$ diberikan oleh
\definitionsverweis {simbol Christoffel}{}{}
\zusatzklammer {indeks bagi satu-satunya arah diferensiasi dihilangkan} {} {}

\mavergleichskettedisp
{\vergleichskette
{ \begin{pmatrix} \tilde{\Gamma}_1^1 (t)   &  \tilde{\Gamma}_1^2 (t)   \\ \tilde{\Gamma}_2^1(t)  &  \tilde{\Gamma}_2^2(t)  \end{pmatrix}
}
{ =} { \begin{pmatrix}  0  &   - { \frac{  c  }{  2 \pi   } }   \\  { \frac{  c  }{  2 \pi  } }  &  0  \end{pmatrix}
}
{ } {
}
{ } {
}
{ } {
}
}
{}{}{.}
Untuk suatu titik

\mavergleichskette
{\vergleichskette
{ Q
}
{ = }{ \begin{pmatrix}  a  \\b   \end{pmatrix}
}
{ \in }{ \R^2
}
{    }{
}
{   }{
}
}
{}{}{}
tinjau pemetaan
\maabbeledisp {\Psi_{Q}} {\R} { S^1 \times \R^2
} { t } { \left(  \cos  t   , \,   \sin  t  , \,   M(t)  \begin{pmatrix}  a  \\b   \end{pmatrix}                 \right)
} {}
dengan

\mavergleichskettedisp
{\vergleichskette
{ M(t)
}
{ =} { \begin{pmatrix}  \cos  { \frac{  ct  }{  2 \pi   } }   &   -  \sin  { \frac{  ct  }{  2 \pi   } }    \\  \sin  { \frac{  ct  }{  2 \pi  } }   &  \cos  { \frac{  ct  }{  2 \pi   } }  \end{pmatrix}
}
{ } {
}
{ } {
}
{ } {
}
}
{}{}{.}
\aufzaehlungfuenf{Tentukan $\Psi_Q(0)$.
}{Tentukan $\Psi_Q(2 \pi)$.
}{Tunjukkan bahwa $\Psi_Q$ adalah suatu
\definitionsverweis {pengangkatan horizontal}{}{}
sepanjang $\psi$.
}{Tunjukkan bahwa
\maabbeledisp {} {\R^2} { \R^2
} { Q } { \Psi_Q(2 \pi)
} {,}
merupakan
\definitionsverweis {isometri linear}{}{}
\zusatzklammer {titik pangkal $(1,0)$ tidak dicantumkan di sini} {} {.}
}{Tunjukkan bahwa
\maabbeledisp {} {\Z} { \operatorname{SL}_{  2  } \! \left(  \R  \right)
} { k } { \left(  Q \mapsto \Psi_Q(2k \pi )  \right)
} {,}
merupakan
\definitionsverweis {homomorfisme grup}{}{.}
}

}
{} {}
\endgroup

\part{Formulir Solusi Resmi}
\chapter*{Batas solusi resmi}
\addcontentsline{toc}{chapter}{Batas solusi resmi}
Setiap formulir pada bagian ini berasal dari permukaan solusi resmi yang dibekukan. Solusi yang tidak tersedia pada sumber tetap tidak tersedia di sini; tidak ada materi asli yang dimasukkan ke dalam formulir resmi.
\chapter{Formulir Solusi Resmi Ujian 1}
\noindent\textit{Solusi yang disediakan oleh sumber resmi; kekosongan sumber dipertahankan.}
\begingroup
\renewcommand{\theHfakt}{o011.exam.official.01.\arabic{fakt}}
%Data tentang institusi

%\input{Data guru}

%\renewcommand{\fachbereich}{Departemen}

%\renewcommand{\dozent}{Prof. Dr. }

%Tanggal ujian

\renewcommand{\klausurgebiet}{ }

\renewcommand{\klausurtyp}{ }

\renewcommand{\klausurdatum}{. 20}

\klausurvorspann {\fachbereich} {\klausurdatum} {\dozent} {\klausurgebiet} {\klausurtyp}

%Data untuk tabel poin berikut

\renewcommand{\aeins}{  3  }

\renewcommand{\azwei}{  3   }

\renewcommand{\adrei}{  6  }

\renewcommand{\avier}{  5  }

\renewcommand{\afuenf}{  5  }

\renewcommand{\asechs}{  2  }

\renewcommand{\asieben}{  5  }

\renewcommand{\aacht}{  5  }

\renewcommand{\aneun}{  7  }

\renewcommand{\azehn}{  3  }

\renewcommand{\aelf}{  2  }

\renewcommand{\azwoelf}{  5  }

\renewcommand{\adreizehn}{  9   }

\renewcommand{\avierzehn}{  5  }

\renewcommand{\afuenfzehn}{  65  }

\renewcommand{\asechzehn}{    }

\renewcommand{\asiebzehn}{    }

\renewcommand{\aachtzehn}{    }

\renewcommand{\aneunzehn}{    }

\renewcommand{\azwanzig}{    }

\renewcommand{\aeinundzwanzig}{    }

\renewcommand{\azweiundzwanzig}{    }

\renewcommand{\adreiundzwanzig}{    }

\renewcommand{\avierundzwanzig}{    }

\renewcommand{\afuenfundzwanzig}{    }

\renewcommand{\asechsundzwanzig}{    }

\punktetabellevierzehn

\klausurnote

\newpage

\setcounter{section}{0}



\inputaufgabeklausurloesung
{3}

{

Definisikan istilah-istilah berikut
\zusatzklammer {yang dicetak miring} {} {}.
\aufzaehlungsechs{\stichwort {Lingkaran kelengkungan} {}
dari suatu
\definitionsverweis {kurva}{}{}
yang dua kali
\definitionsverweis {terdiferensialkan secara kontinu}{}{}
dan
\definitionsverweis {berparametrisasi panjang busur}{}{}.
Misalkan
\maabbdisp {\gamma} {I} {\R^2
} {}
dan ambil titik

\mavergleichskette
{\vergleichskette
{ t_0
}
{ \in }{ I
}
{  }{
}
{    }{
}
{   }{
}
}
{}{}{.}

}{Suatu \stichwort {submanifold tertutup} {}
\mathl{M \subseteq N}{} dari manifold diferensiabel $N$.

}{\stichwort {Ruang singgung} {} pada titik
\mathl{P\in M}{} dari manifold diferensiabel $M$.

}{Istilah
\stichwort {berorientasi} {}
untuk suatu
\definitionsverweis {manifold diferensiabel}{}{}
$M$.

}{\stichwort {Batas} {}
suatu
\definitionsverweis {manifold berbatas}{}{.}

}{\stichwort {Simbol Christoffel} {}

\mathl{\Gamma_{ij}^k}{} bagi koneksi Levi-Civita dari struktur Riemann $\left(  g_{ij}  \right)$ pada himpunan terbuka

\mavergleichskette
{\vergleichskette
{ U
}
{ \subseteq }{ \R^n
}
{  }{
}
{    }{
}
{   }{
}
}
{}{}{.}
}

}
{

\aufzaehlungsechs{Lingkaran dengan jari-jari

\mavergleichskettedisp
{\vergleichskette
{ r
}
{ =} {  { \frac{  1 }{ \betrag {  \gamma_1'(t_0) \gamma_2^{\prime \prime} (t_0) - \gamma_2'(t_0) \gamma_1^{\prime \prime} (t_0)  }  } }
}
{ } {
}
{ } {
}
{ } {
}
}
{}{}{}
dan pusat

\mavergleichskettedisp
{\vergleichskette
{ M
}
{ =} { \gamma(t_0) +  { \frac{  1 }{  \gamma_1'(t_0) \gamma_2^{\prime \prime} (t_0) - \gamma_2'(t_0) \gamma_1^{\prime \prime} (t_0)  } }  \begin{pmatrix} - \gamma_2'(t_0) \\ \gamma_1'(t_0)   \end{pmatrix}
}
{ } {
}
{ } {
}
{ } {
}
}
{}{}{}
dinamakan lingkaran kelengkungan dari $\gamma$ di $t_0$.
}{Suatu himpunan bagian tertutup
\mathl{M \subseteq N}{} dinamakan submanifold tertutup jika untuk setiap titik
\mathl{P \in M}{} terdapat suatu
\definitionsverweis {peta}{}{}
sedemikian sehingga
\mathl{P \in W \subseteq N}{},
\definitionsverweis {terbuka}{}{,}
\maabb {\theta} { W } { W'
} {,}

\mathl{W' \subseteq \R^n}{} terbuka, dan

\mavergleichskettedisp
{\vergleichskette
{ M \cap W
}
{ =} {\theta^{-1} (( \R^m \times\{0\}) \cap W' )
}
{ } {
}
{ } {
}
{ } {
}
}
{}{}{.}
}{Ruang singgung
\mathl{T_PM}{} terdiri atas semua kelas ekuivalensi lintasan terdiferensialkan yang melalui titik tersebut dan ekuivalen secara tangensial.
}{Suatu
\definitionsverweis {manifold diferensiabel}{}{}
$M$ dengan sebuah
\definitionsverweis {atlas}{}{}

\mathl{(U_i, V_i, \alpha_i)}{} disebut berorientasi jika setiap petanya
\definitionsverweis {berorientasi}{}{}
dan semua pergantian petanya
\definitionsverweis {mempertahankan orientasi}{}{}
.
}{Batas $M$ didefinisikan oleh

\mavergleichskettedisp
{\vergleichskette
{ \partial M
}
{ =} {  \left\{ x \in M  \mid   \alpha_i(x) \in  \partial H \cap V_i  \text{ untuk } \text{suatu } i \in I         \right\}
}
{ } {
}
{ } {
}
{ } {
}
}
{}{}{}
di mana $\alpha_i$ adalah peta-peta koordinat.
}{Misalkan
\mathl{\left(  g^{k l}  \right)}{} adalah matriks invers dari
\mathl{\left(  g_{ij}  \right)}{.}
\zusatzklammer {Fungsi-fungsi bernilai riil pada $U$} {} {}

\mavergleichskettedisp
{\vergleichskette
{\Gamma_{ij}^k
}
{ =} {  { \frac{  1 }{ 2 } } \sum_{l = 1}^n g^{kl}(\partial_ig_{jl}+\partial_jg_{il}-\partial_lg_{ij})
}
{ } {
}
{ } {
}
{ } {
}
}
{}{}{}
dinamakan simbol-simbol Christoffel untuk koneksi Levi-Civita.
}

 }



\inputaufgabeklausurloesung
{3}

{

Nyatakan teorema-teorema berikut.
\aufzaehlungdrei{Teorema tentang citra
\definitionsverweis {pemetaan Gauss}{}{}
dari suatu hipermuka.}{Rumus luas permukaan dari permukaan putar yang dibentuk oleh
\definitionsverweis {kurva diferensiabel}{}{}
\maabbeledisp {\gamma} {]a,b[} {\R^2
} { t } {(x(t),y(t))
} {,}
dengan

\mavergleichskette
{\vergleichskette
{  y(t)
}
{ > }{  0
}
{  }{
}
{    }{
}
{   }{
}
}
{}{}{.}}{\stichwort {Aturan hasil kali} {}
untuk turunan eksterior.}

}
{

\aufzaehlungdrei{\faktsituation {Misalkan

\mavergleichskette
{\vergleichskette
{ W
}
{ \subseteq }{\R^n
}
{ }{
}
{ }{
}
{ }{
}
}
{}{}{}
adalah suatu
\definitionsverweis {himpunan bagian terbuka}{}{}
dan
\maabbdisp {h} { W } {\R
} {}
adalah suatu
\definitionsverweis {fungsi terdiferensialkan secara kontinu}{}{.}}
\faktvoraussetzung {
\definitionsverweis {serat}{}{}

\mavergleichskette
{\vergleichskette
{ Y
}
{ = }{ h^{-1}(c)
}
{ }{
}
{ }{
}
{ }{
}
}
{}{}{}
dari $h$ di

\mavergleichskette
{\vergleichskette
{ c
}
{ \in }{\R
}
{  }{
}
{    }{
}
{   }{
}
}
{}{}{}
bersifat
\definitionsverweis {kompakt}{}{}
dan
\definitionsverweis {reguler}{}{.}}
\faktfolgerung {Maka
\definitionsverweis {pemetaan Gauss}{}{}
\maabbdisp {} { Y } { S^{n-1}
} {}
bersifat surjektif.}
\faktzusatz {}
\faktzusatz {}}{\faktsituation {Misalkan
\maabbeledisp {\gamma} {]a,b[} {\R^2
} { t } {(x(t),y(t))
} {,}
adalah suatu
\definitionsverweis {kurva terdiferensialkan}{}{}
dengan

\mavergleichskette
{\vergleichskette
{ y(t)
}
{ > }{ 0
}
{ }{
}
{ }{
}
{ }{
}
}
{}{}{,}}
\faktvoraussetzung {yang menginduksi suatu
\definitionsverweis {difeomorfisme}{}{}
ke

\mavergleichskette
{\vergleichskette
{M
}
{ = }{ \gamma(I)
}
{ }{
}
{ }{
}
{ }{
}
}
{}{}{}
, dengan

\mavergleichskette
{\vergleichskette
{M
}
{ \subseteq }{ G
}
{ }{
}
{ }{
}
{ }{
}
}
{}{}{}
sebagai
\definitionsverweis {submanifold tertutup}{}{}
satu dimensi dalam suatu himpunan terbuka

\mavergleichskette
{\vergleichskette
{G
}
{ \subseteq }{\R \times \R_+
}
{ }{
}
{ }{
}
{ }{
}
}
{}{}{}
.}
\faktfolgerung {Maka
\definitionsverweis {permukaan putar}{}{}
yang bersesuaian merupakan submanifold tertutup dari
\mathl{\R^3}{} tanpa sumbu $x$, dan luas permukaannya adalah

\mathdisp {2 \pi \int_{  a  }^{  b  } \sqrt{ (x'(t))^2+ (y'(t))^2 } \cdot y(t) \, d t} { . }
}
\faktzusatz {}
\faktzusatz {}}{\faktsituation {Misalkan

\mavergleichskette
{\vergleichskette
{ U
}
{ \subseteq }{ \R^n
}
{ }{
}
{ }{
}
{ }{
}
}
{}{}{}
\definitionsverweis {terbuka}{}{}
dan misalkan

\mavergleichskette
{\vergleichskette
{ \omega
}
{ \in }{
{ \mathcal E }^{  k  }_{ 1 } (  U   )
}
{  }{
}
{    }{
}
{   }{
}
}
{}{}{}
serta

\mavergleichskette
{\vergleichskette
{ \tau
}
{ \in }{
{ \mathcal E }^{  \ell }_{ 1 } (  U   )
}
{  }{
}
{    }{
}
{   }{
}
}
{}{}{}
adalah bentuk-bentuk diferensial terdiferensialkan pada $U$.}
\faktfolgerung {Maka berlaku aturan hasil kali

\mavergleichskettedisp
{\vergleichskette
{ d( \omega \wedge \tau)
}
{ =} { ( d \omega) \wedge \tau  + (-1)^k \omega  \wedge  (d \tau)
}
{ } {
}
{ } {
}
{ } {
}
}
{}{}{.}}
\faktzusatz {}
\faktzusatz {}}

 }



\inputaufgabeklausurloesung
{6 (1+2+3)}

{

Misalkan $Y$ adalah hipermuka yang diberikan oleh persamaan

\mavergleichskettedisp
{\vergleichskette
{ 2x^2 +3y^2+5z^2
}
{ =} { 10
}
{ } {
}
{ } {
}
{ } {
}
}
{}{}{}
di $\R^3$. Misalkan

\mavergleichskette
{\vergleichskette
{ P
}
{ = }{ \begin{pmatrix}  -1 \\ -1\\  1   \end{pmatrix}
}
{ \in }{ Y
}
{    }{
}
{   }{
}
}
{}{}{}
dan

\mavergleichskette
{\vergleichskette
{ w
}
{ = }{ \begin{pmatrix}  -3 \\ 2 \\  4   \end{pmatrix}
}
{ \in }{ \R^3
}
{    }{
}
{   }{
}
}
{}{}{.}
\aufzaehlungdrei{Tentukan
\definitionsverweis {ruang singgung}{}{}

\mathl{T_PY}{.}
}{Tentukan
\definitionsverweis {dekomposisi ortogonal}{}{}
$w$ menjadi komponen tangensial dan komponen ortogonal pada titik $P$.
}{Realisasikan komponen tangensial $w$ di $T_PY$ melalui suatu kurva diferensiabel yang seluruhnya terletak pada $Y$.
}

}
{

\aufzaehlungdrei{Total diferensialnya adalah

\mathdisp {\left(  4x  , \,   6y  , \,   10z                 \right)} { , }

sehingga pada $P$ diperoleh

\mavergleichskettedisp
{\vergleichskette
{ T_PY
}
{ =} { \left\{  \begin{pmatrix}  a  \\ b \\ c   \end{pmatrix}   \mid   -4a-6b+10c = 0         \right\}
}
{ =} { \left\{  \begin{pmatrix}  a  \\ b \\ c   \end{pmatrix}   \mid      \left\langle  \begin{pmatrix}  a  \\ b \\ c   \end{pmatrix}   , \begin{pmatrix}  -4 \\ -6\\  10   \end{pmatrix}     \right\rangle = 0         \right\}
}
{ } {
}
{ } {
}
}
{}{}{.}
}{Kita gunakan ansatz

\mavergleichskettedisp
{\vergleichskette
{  \begin{pmatrix}  -3 \\ 2\\  4   \end{pmatrix}
}
{ =} { s  \begin{pmatrix}  -4 \\ -6\\  10   \end{pmatrix} + v
}
{ } {
}
{ } {
}
{ } {
}
}
{}{}{}
dengan skalar yang dicari

\mavergleichskette
{\vergleichskette
{ s
}
{ \in }{ \R
}
{ }{
}
{ }{
}
{ }{
}
}
{}{}{}
dan vektor tangensial yang dicari $v$. Kita peroleh

\mavergleichskettedisp
{\vergleichskette
{ \left\langle \begin{pmatrix} -3 \\ 2\\ 4 \end{pmatrix} , \begin{pmatrix} -4 \\ -6\\ 10 \end{pmatrix} \right\rangle
}
{ =} { 12-12+40
}
{ =} { 40
}
{ } {
}
{ } {
}
}
{}{}{.}
Karena

\mavergleichskettedisp
{\vergleichskette
{ \left\langle  \begin{pmatrix}  -4 \\ -6\\  10   \end{pmatrix}  ,  \begin{pmatrix}  -4 \\ -6\\  10   \end{pmatrix}   \right\rangle
}
{ =} { 16 +36 + 100
}
{ =} { 152
}
{ } {
}
{ } {
}
}
{}{}{}
maka

\mavergleichskettedisp
{\vergleichskette
{ s
}
{ =} { { \frac{   40 }{  152 } }
}
{ =} { { \frac{   5 }{  19 } }
}
{ } {
}
{ } {
}
}
{}{}{}
dan dengan demikian

\mavergleichskettedisp
{\vergleichskette
{ v
}
{ =} { \begin{pmatrix}  -3 \\ 2\\  4   \end{pmatrix}   - { \frac{   5 }{  19 } }  \begin{pmatrix}  -4 \\ -6\\  10   \end{pmatrix}
}
{ =} { { \frac{   1  }{  19 } }  \begin{pmatrix}  -57+20 \\ 38 +30\\  76 -50   \end{pmatrix}
}
{ =} { { \frac{   1  }{  19 } }  \begin{pmatrix}  -37 \\ 68\\  26   \end{pmatrix}
}
{ } {
}
}
{}{}{.}
Jadi, dekomposisi ortogonalnya adalah

\mavergleichskettedisp
{\vergleichskette
{  w
}
{ =} { \begin{pmatrix}  -3 \\ 2\\  4   \end{pmatrix}
}
{ =} { { \frac{   5 }{  19 } }  \begin{pmatrix}  -4 \\ -6\\  10   \end{pmatrix} + { \frac{   1  }{  19 } }  \begin{pmatrix}  -37 \\ 68\\  26   \end{pmatrix}
}
{ } {
}
{ } {
}
}
{}{}{.}
}{Kita tulis persamaan permukaan

\mavergleichskette
{\vergleichskette
{ 2x^2 +3y^2+5z^2
}
{ = }{ 10
}
{ }{
}
{ }{
}
{ }{
}
}
{}{}{}
sebagai

\mavergleichskettedisp
{\vergleichskette
{ z
}
{ =} { \pm \sqrt{ 10- 2x^2 - 3y^2 }
}
{ } {
}
{ } {
}
{ } {
}
}
{}{}{,}
di mana akar positif harus dipilih pada titik $P$. Ini merupakan parameterisasi permukaan pada suatu lingkungan terbuka dari $P$; ruang singgung di
\mathl{\begin{pmatrix} -1 \\ -1 \end{pmatrix}}{} dipetakan ke ruang singgung
\mathl{T_PY}{}.
}

 }



\inputaufgabeklausurloesung
{5}

{

Misalkan
\maabbeledisp {\gamma} {\R} { \R^2
} { t } { \begin{pmatrix}  \cos \left(  \omega t \right)  \\ \sin  \left(  \omega t \right)    \end{pmatrix}
} {,}
dengan

\mavergleichskette
{\vergleichskette
{ \omega
}
{ > }{ 0
}
{  }{
}
{    }{
}
{   }{
}
}
{}{}{.}
Tunjukkan bahwa tidak ada gerak melingkar lain dengan norma kecepatan konstan yang memiliki nilai dan turunan hingga orde kedua yang sama dengan $\gamma$ di titik parameter

\mavergleichskette
{\vergleichskette
{ t
}
{ = }{ 0
}
{  }{
}
{    }{
}
{   }{
}
}
{}{}{}
tersebut.

}
{

Kita peroleh

\mavergleichskettedisp
{\vergleichskette
{ \gamma(0)
}
{ =} { \begin{pmatrix} 1 \\ 0 \end{pmatrix}
}
{ } {
}
{ } {
}
{ } {
}
}
{}{}{,}

\mavergleichskettedisp
{\vergleichskette
{ \gamma'(0)
}
{ =} { \begin{pmatrix} 0 \\ \omega \end{pmatrix}
}
{ } {
}
{ } {
}
{ } {
}
}
{}{}{}
dan

\mavergleichskettedisp
{\vergleichskette
{ \gamma^{\prime \prime} (0)
}
{ =} { \begin{pmatrix} - \omega^2 \\ 0 \end{pmatrix}
}
{ } {
}
{ } {
}
{ } {
}
}
{}{}{.}
Suatu gerak melingkar yang sama dengan gerak yang diberikan pada

\mavergleichskette
{\vergleichskette
{ t
}
{ = }{ 0
}
{ }{
}
{ }{
}
{ }{
}
}
{}{}{}
hingga turunan kedua khususnya harus memiliki garis singgung yang sama di titik tersebut. Oleh karena itu, pusat lingkarannya harus terletak pada sumbu $x$. Selain itu, pusat itu harus terletak di sebelah kiri
\mathl{\begin{pmatrix} 1 \\ 0 \end{pmatrix}}{} karena percepatannya mengarah ke kiri. Karena itu, gerak melingkar beraturan yang dicari dengan pusat
\mathl{\begin{pmatrix} x \\ 0 \end{pmatrix}}{}
\zusatzklammer {
\mavergleichskettek
{\vergleichskettek
{ x
}
{ < }{ 1
}
{ }{
}
{ }{
}
{ }{
}
}
{}{}{}} {} {}
dan jari-jari
\mathl{1-x}{} dapat dituliskan sebagai

\mavergleichskettedisp
{\vergleichskette
{  \delta (t)
}
{ =} {  \begin{pmatrix}  x  \\ 0   \end{pmatrix} + (1-x) \begin{pmatrix}  \cos \left(  u t \right)  \\ \sin  \left(  u t \right)     \end{pmatrix}
}
{ } {
}
{ } {
}
{ } {
}
}
{}{}{}
. Kita mempunyai

\mavergleichskettedisp
{\vergleichskette
{ \delta(0)
}
{ =} { \gamma(0)
}
{ } {
}
{ } {
}
{ } {
}
}
{}{}{.}
Selanjutnya,

\mavergleichskettedisp
{\vergleichskette
{ \delta'( 0)
}
{ =} { (1-x) \begin{pmatrix} 0 \\ u \end{pmatrix}
}
{ } {
}
{ } {
}
{ } {
}
}
{}{}{}
dan

\mavergleichskettedisp
{\vergleichskette
{ \delta^{\prime \prime}( 0)
}
{ =} { (1-x) \begin{pmatrix} -u^2 \\ 0 \end{pmatrix}
}
{ } {
}
{ } {
}
{ } {
}
}
{}{}{.}
Agar syarat-syarat yang diminta terpenuhi, harus berlaku

\mavergleichskettedisp
{\vergleichskette
{ \omega
}
{ =} { (1-x) u
}
{ } {
}
{ } {
}
{ } {
}
}
{}{}{}
dan

\mavergleichskettedisp
{\vergleichskette
{ \omega^2
}
{ =} { (1-x) u^2
}
{ } {
}
{ } {
}
{ } {
}
}
{}{}{}
. Dari sini diperoleh

\mavergleichskettedisp
{\vergleichskette
{ \omega^2
}
{ =} { \omega u
}
{ } {
}
{ } {
}
{ } {
}
}
{}{}{}
dan dengan demikian

\mavergleichskettedisp
{\vergleichskette
{ u
}
{ =} { \omega
}
{ } {
}
{ } {
}
{ } {
}
}
{}{}{}
serta

\mavergleichskettedisp
{\vergleichskette
{ x
}
{ =} { 0
}
{ } {
}
{ } {
}
{ } {
}
}
{}{}{.}

 }



\inputaufgabeklausurloesung
{5}

{

Buktikan teorema tentang kelengkungan suatu kurva bidang yang diberikan secara implisit.

}
{

Misalkan
\maabb {\gamma} { I } { \R^2
} {}
adalah parameterisasi panjang busur kurva $Y$ pada suatu lingkungan dari $P$

\mavergleichskette
{\vergleichskette
{ \gamma (0)
}
{ = }{ P
}
{ }{
}
{ }{
}
{ }{
}
}
{}{}{,}
yang sesuai dengan orientasi yang diberikan. Diferensial total
\maabbdisp {\left(D N \right)_{ P }} { \R^2} {\R^2
} {}
bersifat linier, sehingga cukup membuktikan pernyataan bagi vektor $\gamma'(0)$. Berdasarkan
aturan rantai,

\mavergleichskettedisp
{\vergleichskette
{ \left(  D N   \right)_{ P } \left(   \gamma'(0)  \right)
}
{ =} { (N \circ \gamma )'(0)
}
{ } {
}
{ } {
}
{ } {
}
}
{}{}{.}
Vektor ini merupakan kelipatan dari $\gamma'(0)$, sehingga sama dengan

\mathdisp {\left\langle  (N \circ  \gamma)'(0) ,  \gamma' (0)   \right\rangle  \gamma'(0)} { . }

Karena syarat ortogonalitas memberikan

\mavergleichskettedisp
{\vergleichskette
{ \left\langle (N \circ \gamma)(t) , \gamma' (t) \right\rangle
}
{ =} { 0
}
{ } {
}
{ } {
}
{ } {
}
}
{}{}{}
untuk semua

\mavergleichskette
{\vergleichskette
{ t
}
{ \in }{ I
}
{ }{
}
{ }{
}
{ }{
}
}
{}{}{}
dan karenanya

\mavergleichskettedisp
{\vergleichskette
{ \left\langle (N \circ \gamma)'(t) , \gamma' (t) \right\rangle + \left\langle (N \circ \gamma)(t) , \gamma^{\prime \prime} (t) \right\rangle
}
{ =} { 0
}
{ } {
}
{ } {
}
{ } {
}
}
{}{}{.}
Dengan demikian,

\mavergleichskettedisp
{\vergleichskette
{ \left(  D N   \right)_{ P } \left(   \gamma'(0)  \right)
}
{ =} { - \left\langle  (N \circ  \gamma)(0) ,  \gamma^{\prime \prime} (0)   \right\rangle  \gamma'(0)
}
{ =} { - \kappa( \gamma(0)) \gamma'(0)
}
{ } {
}
{ } {
}
}
{}{}{}
berdasarkan
Lemma 3.8 (Differentialgeometrie (Osnabrück 2023)).

 }



\inputaufgabeklausurloesung
{2}

{

Berikan contoh manifold diferensiabel
\mathkor {} {M} {dan} {N} {}
dengan

\mavergleichskette
{\vergleichskette
{ \operatorname{ dim}_{  } ^{  }  \left(   M   \right), \, \operatorname{ dim}_{  } ^{  }  \left(   N   \right)
}
{ \geq }{ 1
}
{  }{
}
{    }{
}
{   }{
}
}
{}{}{}
serta pemetaan diferensiabel
\maabb {\varphi} { M } { N
} {}
dan
\maabb {\psi} { N } { M
} {}
sedemikian sehingga

\mavergleichskette
{\vergleichskette
{ \psi \circ \varphi
}
{ = }{
\operatorname{Id}_{ M }
}
{  }{
}
{    }{
}
{   }{
}
}
{}{}{}
dan

\mavergleichskette
{\vergleichskette
{\varphi \circ  \psi
}
{ \neq }{
\operatorname{Id}_{ N }
}
{  }{
}
{    }{
}
{   }{
}
}
{}{}{}
berlaku.

}
{

Kita pilih

\mavergleichskette
{\vergleichskette
{M
}
{ = }{\R
}
{ }{
}
{ }{
}
{ }{
}
}
{}{}{}
dan

\mavergleichskette
{\vergleichskette
{N
}
{ = }{\R^2
}
{ }{
}
{ }{
}
{ }{
}
}
{}{}{}
dan
\maabbeledisp {\varphi} {\R} {\R^2
} { x } {(x,0)
} {,}
dan
\maabbeledisp {\psi} {\R^2} {\R
} {(x,y)} { x
} {.}
Maka

\mavergleichskettedisp
{\vergleichskette
{ \left( \psi \circ \varphi \right) (x)
}
{ =} { \psi((x,0))
}
{ =} { x
}
{ } {
}
{ } {
}
}
{}{}{}
dan

\mavergleichskettedisp
{\vergleichskette
{ \left(  \varphi  \circ \psi  \right) ((x,y))
}
{ =} { \varphi (x)
}
{ =} {(x,0)
}
{ } {
}
{ } {
}
}
{}{}{.}

 }



\inputaufgabeklausurloesung
{5}

{

Buktikan bahwa ekuivalensi tangensial lintasan-lintasan pada suatu manifold di titik $P$ dapat diperiksa dengan menggunakan sebarang peta koordinat.

}
{

Untuk suatu
\definitionsverweis {kurva terdiferensialkan}{}{}
\maabbdisp {\gamma} { I } { M
} {}
dengan

\mavergleichskette
{\vergleichskette
{ 0
}
{ \in }{ I
}
{ }{
}
{ }{
}
{ }{
}
}
{}{}{}
dan

\mavergleichskette
{\vergleichskette
{ \gamma(0)
}
{ = }{ P
}
{  }{
}
{    }{
}
{   }{
}
}
{}{}{}
dan suatu
\definitionsverweis {peta}{}{}
\maabbdisp {\alpha} { U } { V
} {}
\zusatzklammer {dengan

\mavergleichskettek
{\vergleichskettek
{P
}
{ \in }{ U
}
{ }{
}
{ }{
}
{ }{
}
}
{}{}{}
dan

\mavergleichskette
{\vergleichskette
{V
}
{ \subseteq }{\R^n
}
{ }{
}
{ }{
}
{ }{
}
}
{}{}{}} {} {}
ekspresi

\mathdisp {\left(  \alpha \circ \left(  \gamma \vert_{\gamma^{-1} (U)}  \right)  \right)'(0)} {  }

tidak berubah apabila kita beralih ke interval terbuka yang lebih kecil

\mavergleichskette
{\vergleichskette
{ 0
}
{ \in }{ I'
}
{ \subseteq }{ I
}
{ }{
}
{ }{
}
}
{}{}{}
dan himpunan terbuka yang lebih kecil

\mavergleichskette
{\vergleichskette
{P
}
{ \in }{ U'
}
{ \subseteq }{ U
}
{ }{
}
{ }{
}
}
{}{}{}
\zusatzklammer {dengan peta yang diinduksi} {} {}
. Karena itu, kita dapat mengasumsikan bahwa
\mathkor {} {\gamma_1} {dan} {\gamma_2} {}
didefinisikan pada interval yang sama, citranya terletak di $U$, dan pada $U$ terdapat dua peta
\maabbdisp {\alpha_1} { U } { V_1
} {}
dan
\maabbdisp {\alpha_2} { U } { V_2
} {}
. Selanjutnya, dari

\mavergleichskettedisp
{\vergleichskette
{ ( \alpha_1 \circ \gamma_1 )'(0)
}
{ =} { ( \alpha_1 \circ \gamma_2 )'(0)
}
{ } {
}
{ } {
}
{ } {
}
}
{}{}{}
berdasarkan
aturan rantai
dan keterdiferensialan
\definitionsverweis {pemetaan transisi}{}{}

\mathl{\alpha_2 \circ \alpha_1^{-1}}{}, langsung diperoleh

\mavergleichskettealign
{\vergleichskettealign
{ \left(  \alpha_2 \circ \gamma_1  \right)'(0)
}
{ =} { \left(  \left(  \alpha_2 \circ \alpha_1^{-1}  \right) \circ \left(  \alpha_1 \circ \gamma_1  \right)  \right)'(0)
}
{ =} { \left(  D   \left(   \alpha_2 \circ \alpha_1^{-1}   \right)      \right)_{\alpha_1(P)} \left(   \left(  \alpha_1 \circ \gamma_1  \right)'(0)   \right)
}
{ =} { \left(  D   \left(   \alpha_2 \circ \alpha_1^{-1}   \right)      \right)_{\alpha_1(P)} \left(   \left(  \alpha_1 \circ \gamma_2  \right)'(0)   \right)
}
{ =} { \left(  \alpha_2 \circ \gamma_2  \right)'(0)
}
}
{}
{}{.}

 }



\inputaufgabeklausurloesung
{5}

{

Tunjukkan bahwa pemetaan singgung $T(\varphi)$ dari
\maabbeledisp {\varphi} {\R^1} {S^1
} { t } {( \cos  t, \sin  t  )
} {,}
bersifat surjektif.

}
{

Pemetaan $\varphi$ bersifat surjektif, sehingga cukup ditunjukkan bahwa untuk setiap
\mathl{t \in \R}{} pemetaan singgung linier
\maabbdisp {} {T_t \R \cong \R} { T_{\varphi(t)} S^1
} {}
bersifat surjektif. Karena kedua ruang itu berdimensi satu, cukup ditunjukkan bahwa suatu vektor yang berbeda dari $0$ tidak dipetakan ke $0$. Vektor singgung di $t$ direalisasikan oleh lintasan terdiferensialkan
\maabbeledisp {\gamma} {\R} {\R
} { s } {t+s
} {.}
Lintasan komposisi
\maabbeledisp {\varphi \circ \gamma} {\R} {S^1
} { s } { ( \cos (t+s) , \sin (t+s) )
} {,}
mewujudkan vektor singgung citra; pada bidang ambien,
\zusatzklammer {
\mathl{T_{\varphi(t)}S^1 \subseteq \R^2}{}} {} {}

\mavergleichskettedisp
{\vergleichskette
{( \varphi \circ \gamma)' (0)
}
{ =} {(- \sin t , \cos t )
}
{ } {
}
{ } {
}
{ } {
}
}
{}{}{,}
dan vektor ini tidak nol.

 }



\inputaufgabeklausurloesung
{7 (3+4)}

{

Tinjau pemetaan diferensiabel
\maabbeledisp {\gamma} { [1,c] } {\R^2
} { t } {(t,t^{3})
} {,}
\zusatzklammer {dengan \mathlk{c \geq 1}{}} {} {}
dan
\maabbeledisp {\varphi} {\R_+ \times \R_+ } {\R^3
} {(u,v)} {(u^3,u^2+v^2,u^{-1}v^{-1} )
} {,}
serta bentuk diferensial

\mavergleichskettedisp
{\vergleichskette
{ \omega
}
{ =} { (x-y)dx-z^2dy+dz
}
{ } {
}
{ } {
}
{ } {
}
}
{}{}{}
pada $\R^3$.

a) Hitung bentuk diferensial tarik balik
\mathl{\varphi^* \omega}{} pada $\R_+ \times \R_+$.

b) Hitung integral lintasan dari bentuk diferensial
\mathl{\varphi^* \omega}{} sepanjang lintasan $\gamma$ sebagai fungsi dari $c$.

}
{

a) Bentuk diferensial tarik baliknya adalah

\mavergleichskettealigndrucklinks
{\vergleichskettealigndrucklinks
{\varphi^* (  (x-y)dx-z^2dy+dz   )
}
{ =} { (u^3-u^2-v^2 ) d(u^3)  -u^{-2}v^{-2} d(u^2+v^2) +d (u^{-1}v^{-1})
}
{ =} { 3 (u^5-u^4-u^2v^2 ) du  -2u^{-1}v^{-2} du -2u^{-2}v^{-1} dv - u^{-2} v^{-1}du - u^{-1} v^{-2}dv
}
{ =} { (3 u^5-3u^4-3u^2v^2 -2u^{-1}v^{-2}- u^{-2} v^{-1} )  du + (  -2u^{-2}v^{-1} - u^{-1} v^{-2})dv
}
{ } {
}
}
{}{}{.}

b) Integral lintasannya adalah

\mavergleichskettealigndrucklinks
{\vergleichskettealigndrucklinks
{ \int_1^c (3 t^5-3t^4-3t^2t^6 -2t^{-1}t^{-6}- t^{-2} t^{-3} )dt  +(  -2t^{-2}t^{-3} - t^{-1} t^{-6}) d t^3
}
{ =} { \int_1^c   (3 t^5-3t^4-3t^8-2t^{-7} - t^{-5})dt + ( -2t^{-5} - t^{-7})3t^2 dt
}
{ =} { \int_1^c (3 t^5-3t^4-3t^8-2t^{-7} - t^{-5}   -6t^{-3} - 3 t^{-5}      ) dt
}
{ =} { \int_1^c (-3t^8 + 3 t^5-3t^4 -6t^{-3} - 4 t^{-5}  -2t^{-7}       )  dt
}
{ =} { (-{ \frac{   1 }{  3 } } t^9 + { \frac{   1 }{  2 } } t^{6}  - { \frac{   3 }{  5 } } t^5  +  3 t^{-2} +  t^{-4} + { \frac{   1 }{  3 } }  t^{-6}     ) \vert_{  1 }^{  c  }
}
}
{
\vergleichskettefortsetzungalign
{ =} { -{ \frac{   1 }{  3 } } c^9 + { \frac{   1 }{  2 } } c^{6}  - { \frac{   3 }{  5 } } c^5  +  3c^{-2} +   c^{-4} + { \frac{   1 }{  3 } }  c^{-6}
-(-{ \frac{   1 }{  3 } } + { \frac{   1 }{  2 } } - { \frac{   3 }{  5 } } +3 + 1 + { \frac{   1 }{  3 } }  )
}
{ =} { -{ \frac{   1 }{  3 } } c^9 + { \frac{   1 }{  2 } } c^{6}  - { \frac{   3 }{  5 } } c^5 + 3 c^{-2} +  c^{-4} + { \frac{   1 }{  3 } } c^{-6} - { \frac{   39 }{  10 } }  }
{ } {}
{ } {}
}{}{.}

 }



\inputaufgabeklausurloesung
{3}

{

Misalkan
\maabbdisp {\psi} { L } { M
} {}
adalah
\definitionsverweis {pemetaan diferensiabel}{}{}
antara
\definitionsverweis {manifold diferensiabel}{}{.}
Selanjutnya, misalkan
\maabbdisp {f} { M } {\R
} {}
adalah fungsi pada $M$, dan $\omega$ adalah
$k$-\definitionsverweis {bentuk}{}{}
pada $M$. Tunjukkan bahwa

\mavergleichskettedisp
{\vergleichskette
{ \psi^*(f \omega)
}
{ =} { \psi^*(f)  \psi^* \omega
}
{ } {
}
{ } {
}
{ } {
}
}
{}{}{.}

}
{

Misalkan
\mathl{P \in L}{} dan
\mathl{v_1 , \ldots , v_k \in T_PL}{.} Maka

\mavergleichskettealign
{\vergleichskettealign
{ \psi^* \left(  f \omega) (P,v_1 , \ldots , v_k  \right)
}
{ =} { (f \omega) \left(  \psi(P), T_P\psi (v_1) , \ldots , T_P\psi (v_k )       \right)
}
{ =} { f(\psi(P))   \cdot \omega \left(  \psi(P), T_P\psi (v_1) , \ldots , T_P\psi (v_k )       \right)
}
{ =} { \left(  \psi^* (f)  \right) (P)  \cdot \omega \left(  \psi(P), T_P\psi (v_1) , \ldots , T_P\psi (v_k )       \right)
}
{ =} { \left(  \psi^* (f)  \right) (P)   \cdot \left(  \psi^*( \omega)  \right)  \left(  P,v_1 , \ldots , v_k  \right)
}
}
{
\vergleichskettefortsetzungalign
{ =} { \left(  \psi^* (f)  \psi^*( \omega)  \right) \left(  P,v_1 , \ldots , v_k  \right)
}
{ } {}
{ } {}
{ } {}
}
{}{.}

 }



\inputaufgabeklausurloesung
{2}

{

Deskripsikan sebanyak mungkin
\definitionsverweis {manifold berbatas}{}{}
berdimensi satu yang
\definitionsverweis {terhubung}{}{.}

}
{  }



\inputaufgabeklausurloesung
{5}

{

Misalkan
\maabbdisp {f} { [a,b] } { \R_{+}
} {}
adalah
\definitionsverweis {fungsi yang terdiferensialkan secara kontinu}{}{.}
Kita pandang
\definitionsverweis {subgraf}{}{}
sebagai
\definitionsverweis {manifold berbatas}{}{}.
Batasnya terdiri dari grafik, interval dasar, dan kedua sisi tegak, tetapi tidak mencakup keempat titik sudut. Buktikan secara langsung berlakunya teorema Stokes untuk
\definitionsverweis {bentuk diferensial}{}{}

\mavergleichskette
{\vergleichskette
{ \omega
}
{ = }{ xdy
}
{  }{
}
{    }{
}
{   }{
}
}
{}{}{.}

}
{

Karena

\mavergleichskettedisp
{\vergleichskette
{d(xdy)
}
{ =} { dx \wedge dy
}
{ } {
}
{ } {
}
{ } {
}
}
{}{}{}
maka, di satu pihak,

\mavergleichskettedisp
{\vergleichskette
{ \int_M dx \wedge dy
}
{ =} { \int_a^b f(t) dt
}
{ =} { F(b)-F(a)
}
{ } {
}
{ } {
}
}
{}{}{}
adalah luas subgraf, dengan $F$ suatu antiturunan dari $f$. Untuk integral lintasan, batas $M$ harus diparameterkan berlawanan arah jarum jam. Interval dasarnya diparameterkan oleh
\mathl{i: t \mapsto \begin{pmatrix} t \\ 0 \end{pmatrix}}{}, dan

\mavergleichskettedisp
{\vergleichskette
{ i^*\omega
}
{ =} { i^*xdy
}
{ =} { x d0
}
{ =} { 0
}
{ } {
}
}
{}{}{,}
sehingga bagian integral lintasan ini bernilai $0$. Sisi kanan diparameterkan oleh
$i: t \mapsto \begin{pmatrix}  b  \\t   \end{pmatrix}$,
; di sini $t$ bergerak antara
\mathkor {} {0} {dan} {f(b)} {}
. Maka

\mavergleichskettedisp
{\vergleichskette
{ i^*xdy
}
{ =} { b dt
}
{ } {
}
{ } {
}
{ } {
}
}
{}{}{}
dan integral lintasan pada sisi ini adalah
\mathl{b f(b)}{.} Kontribusi sisi kiri adalah
\mathl{-af(a)}{.} Grafik diparameterkan oleh
\maabbeledisp {i} {[a,b]} {
} { t } {(t, f(t))
} {,}
. Maka

\mavergleichskettedisp
{\vergleichskette
{ i^*xdy
}
{ =} { t df
}
{ =} { tf'(t) dt
}
{ } {
}
{ } {
}
}
{}{}{.}
Dengan orientasi yang benar, integralnya adalah

\mavergleichskettealign
{\vergleichskettealign
{ - \int_a^b tf'(t)dt
}
{ =} { - ( tf(t) - F (t) ) \mid _a^b
}
{ =} { - ( bf(b)-F(b) -af(a) +F(a) )
}
{ =} { F(b)-F(a) +af(a) - bf(b)
}
{ } {
}
}
{}
{}{.}
Jadi, seluruh integral lintasan juga bernilai

\mavergleichskettealign
{\vergleichskettealign
{bf(b) - af(a)+ F(b)-F(a) +af(a) - bf(b)
}
{ =} { F(b)-F(a)
}
{ } {
}
{ } {
}
{ } {
}
}
{}
{}{.}

 }



\inputaufgabeklausurloesung
{9}

{

Buktikan teorema titik tetap Brouwer.

}
{

Untuk menyederhanakan notasi, misalkan

\mavergleichskette
{\vergleichskette
{ r
}
{ = }{ 1
}
{ }{
}
{ }{
}
{ }{
}
}
{}{}{.}  Andaikan terdapat pemetaan terdiferensialkan secara kontinu $\psi$ tanpa titik tetap. Maka selalu berlaku

\mavergleichskettedisp
{\vergleichskette
{ x
}
{ \neq} { \psi(x)
}
{ } {
}
{ } {
}
{ } {
}
}
{}{}{,}
sehingga kedua titik itu menentukan sebuah garis. Gagasannya adalah menggunakan garis ini untuk memilih
\zusatzklammer {dari keduanya} {} {}
salah satu dari kedua titik potongnya dengan sfera sebagai titik citra suatu retraksi ke batas. Dengan fungsi bantu

\mavergleichskettedisp
{\vergleichskette
{ h(x)
}
{ =} { { \frac{ \psi(x)-x }{ \Vert { \psi (x)-x} \Vert } }
}
{ } {
}
{ } {
}
{ } {
}
}
{}{}{}
kita definisikan pemetaan
\maabbdisp {\varphi} { B \left( 0 , 1 \right) } { \R^n
} {}
melalui

\mavergleichskettedisphandlinks
{\vergleichskettedisphandlinks
{ \varphi(x)
}
{ =} { x + \left(  - \left\langle  x  , h(x)  \right\rangle + \sqrt{1 + \left\langle  x  , h(x)  \right\rangle^2 - \Vert { x} \Vert^2  }  \right) \cdot h(x)
}
{ } {
}
{ } {
}
{ } {
}
}
{}{}{.}
Ekspresi di bawah tanda akar bernilai positif. Untuk

\mavergleichskette
{\vergleichskette
{ \Vert { x } \Vert
}
{ < }{ 1
}
{ }{
}
{ }{
}
{ }{
}
}
{}{}{}
hal ini jelas; sedangkan untuk

\mavergleichskette
{\vergleichskette
{ \Vert { x } \Vert
}
{ = }{ 1
}
{ }{
}
{ }{
}
{ }{
}
}
{}{}{}
terdapat titik pada sfera yang garis penghubungnya dengan titik dalam bola
\mathl{\psi(x)}{} tidak tegak lurus terhadap $x$
\zusatzklammer {ruang singgung afin pada suatu titik sfera memotong bola hanya di satu titik} {} {,}
sehingga

\mavergleichskettedisp
{\vergleichskette
{ \left\langle x , h(x) \right\rangle
}
{ \neq} { 0
}
{ } {
}
{ } {
}
{ } {
}
}
{}{}{}
berlaku. Karena fungsi akar kuadrat dan norma di luar titik nol
\definitionsverweis {terdiferensialkan secara kontinu}{}{}
bersifat demikian, maka
\mathkor {} {h(x)} {dan} {\varphi(x)} {}
merupakan pemetaan-pemetaan terdiferensialkan secara kontinu. Menurut
Soal 23.23 (Differentialgeometrie (Osnabrück 2023)),
pemetaan $\varphi$ memetakan bola ke sfera dan pembatasannya pada sfera adalah identitas. Jadi terdapat suatu
\definitionsverweis {Retraksi}{}{}
bola pejal tertutup ke batasnya, yang menurut
Teorema 23.9 (Differentialgeometrie (Osnabrück 2023))
mustahil ada.

 }



\inputaufgabeklausurloesung
{5 (3+1+1)}

{

Tinjau silinder

\mavergleichskette
{\vergleichskette
{ Z
}
{ = }{ S^1 \times \R
}
{ \subseteq }{ \R^2 \times \R
}
{    }{
}
{   }{
}
}
{}{}{}
dengan
\definitionsverweis {koneksi trivial}{}{}
pada
\definitionsverweis {bundel tangen}{}{.}
\aufzaehlungdrei{Tentukan, dalam suatu peta isometrik yang sesuai,
\definitionsverweis {persamaan diferensial geodesik}{}{}
untuk kurva geodesik $\psi$ pada $M$ yang memenuhi

\mavergleichskettedisp
{\vergleichskette
{ \psi(0)
}
{ =} { (1,0,3)
}
{ } {
}
{ } {
}
{ } {
}
}
{}{}{}
dan

\mavergleichskettedisp
{\vergleichskette
{ \psi'(0)
}
{ =} { (0,1,-4)
}
{ } {
}
{ } {
}
{ } {
}
}
{}{}{}
.
}{Selesaikan persamaan diferensial dari (1) dalam peta tersebut.
}{Carilah suatu kurva geodesik di $Z$ yang memenuhi syarat-syarat pada (1).
}

}
{

\aufzaehlungdrei{Kita tinjau peta isometrik
\zusatzklammer {inverse} {} {}
peta
\maabbeledisp {\varphi} { ]- \pi, \pi[  \times \R } { S^1 \setminus \{  (-1,0) \} \times \R
} {(t,v)} { \left(  \cos  t   , \,   \sin  t  , \,  v                 \right)
} {.}
Karena

\mavergleichskette
{\vergleichskette
{ \varphi(0,3)
}
{ = }{ (1,0,3)
}
{ }{
}
{ }{
}
{ }{
}
}
{}{}{.}
simbol-simbol Christoffelnya nol, persamaan diferensial geodesiknya adalah

\mavergleichskettedisp
{\vergleichskette
{ \gamma^{\prime \prime}
}
{ =} { 0
}
{ } {
}
{ } {
}
{ } {
}
}
{}{}{.}
Syarat-syarat yang diberikan dalam peta ini menjadi

\mavergleichskettedisp
{\vergleichskette
{ \gamma(0)
}
{ =} { (0,3)
}
{ } {
}
{ } {
}
{ } {
}
}
{}{}{}
dan

\mavergleichskettedisp
{\vergleichskette
{ \gamma'(0)
}
{ =} { (1,-4)
}
{ } {
}
{ } {
}
{ } {
}
}
{}{}{.}
}{Suatu solusi dalam peta tersebut adalah

\mavergleichskettedisp
{\vergleichskette
{ \gamma(t)
}
{ =} { (0,3) + t (1,-4)
}
{ =} { (t,3-4t)
}
{ } {
}
{ } {
}
}
{}{}{.}
}{Karena itu, kurva geodesik dengan syarat awal tersebut adalah

\mavergleichskettedisp
{\vergleichskette
{  \psi(t)
}
{ =} { \varphi( \gamma(t))
}
{ =} { \left(  \cos  t   , \,   \sin  t  , \,   3-4t                 \right)
}
{ } {
}
{ } {
}
}
{}{}{.}
}

 }
\endgroup

\chapter{Formulir Solusi Resmi Ujian 2}
\noindent\textit{Solusi yang disediakan oleh sumber resmi; kekosongan sumber dipertahankan.}
\begingroup
\renewcommand{\theHfakt}{o011.exam.official.02.\arabic{fakt}}
%Data institusi

%\input{Dozentdaten}

%\renewcommand{\fachbereich}{Fachbereich}

%\renewcommand{\dozent}{Prof. Dr. . }

%Data ujian

\renewcommand{\klausurgebiet}{ }

\renewcommand{\klausurtyp}{ }

\renewcommand{\klausurdatum}{ . 20}

\klausurvorspann {\fachbereich} {\klausurdatum} {\dozent} {\klausurgebiet} {\klausurtyp}

%Data untuk tabel nilai berikut

\renewcommand{\aeins}{  3  }

\renewcommand{\azwei}{  3   }

\renewcommand{\adrei}{  8  }

\renewcommand{\avier}{  7  }

\renewcommand{\afuenf}{  8  }

\renewcommand{\asechs}{  3  }

\renewcommand{\asieben}{  5  }

\renewcommand{\aacht}{  5  }

\renewcommand{\aneun}{  6  }

\renewcommand{\azehn}{  4  }

\renewcommand{\aelf}{  12  }

\renewcommand{\azwoelf}{  64  }

\renewcommand{\adreizehn}{     }

\renewcommand{\avierzehn}{    }

\renewcommand{\afuenfzehn}{    }

\renewcommand{\asechzehn}{    }

\renewcommand{\asiebzehn}{    }

\renewcommand{\aachtzehn}{    }

\renewcommand{\aneunzehn}{    }

\renewcommand{\azwanzig}{    }

\renewcommand{\aeinundzwanzig}{    }

\renewcommand{\azweiundzwanzig}{    }

\renewcommand{\adreiundzwanzig}{    }

\renewcommand{\avierundzwanzig}{    }

\renewcommand{\afuenfundzwanzig}{    }

\renewcommand{\asechsundzwanzig}{    }

\punktetabelleelf

\klausurnote

\newpage

\setcounter{bagian}{0}



\inputaufgabeklausurloesung
{3}

{

Definisikan istilah-istilah berikut
\zusatzklammer {yang dicetak miring} {} {.}
\aufzaehlungsechs{Suatu
\stichwort {medan normal} {}
$N$ pada hiperpermukaan diferensiabel

\mavergleichskette
{\vergleichskette
{ Y
}
{ \subseteq }{ \R^n
}
{  }{
}
{    }{
}
{   }{
}
}
{}{}{.}

}{Suatu ruang topologis yang
\stichwort {parakompak}{}
.

}{
\stichwort {Ruang kotangen} {}
di suatu titik

\mavergleichskette
{\vergleichskette
{  P
}
{ \in }{ M
}
{  }{}
{    }{}
{   }{}
}
{}{}{}
dari
\definitionsverweis {manifold diferensiabel}{}{}
$M$.

}{Suatu
\stichwort {titik reguler} {}

\mathl{Q \in L}{} bagi
\definitionsverweis {pemetaan diferensiabel}{}{}
\maabbdisp {\varphi} { L } { M
} {}
antara
\definitionsverweis {manifold-manifold diferensiabel}{}{}
\mathkor {} {L} {dan} {M} {.}

}{
\stichwort {Produk} {}
dari manifold-manifold diferensiabel
\mathkor {} {M} {dan} {N} {.}

}{Suatu
\stichwort {koneksi} {}
pada bundel vektor diferensiabel
\maabb {} {E} {M
} {}
di atas manifold diferensiabel $M$.
}

}
{

\aufzaehlungsechs{Suatu
medan normal
pada $Y$ adalah objek yang didefinisikan pada suatu lingkungan terbuka

\mavergleichskette
{\vergleichskette
{ Y
}
{ \subseteq }{ U
}
{ \subseteq }{ \R^n
}
{ }{
}
{ }{
}
}
{}{}{}
dan bersifat
\definitionsverweis {kontinu}{}{}
sebagai suatu
\definitionsverweis {medan vektor}{}{}
\maabb {F} {U} {\R^n
} {}
yang memenuhi

\mavergleichskettedisp
{\vergleichskette
{ F(P)
}
{ \in} { N_PY
}
{ } {
}
{ } {
}
{ } {
}
}
{}{}{}
untuk setiap

\mavergleichskette
{\vergleichskette
{ P
}
{ \in }{ Y
}
{  }{
}
{    }{
}
{   }{
}
}
{}{}{.}
}{Suatu
\definitionsverweis {ruang topologis}{}{}
$X$ disebut kompak terhadap penutup jika untuk setiap penutup terbuka

\mathdisp {X= \bigcup_{i \in I} U_i \, \, \, \text{ dengan } U_i \text{ terbuka dan himpunan indeks sembarang }} {  }

terdapat suatu himpunan bagian berhingga
\mathl{J \subseteq I}{} sedemikian sehingga

\mavergleichskettedisp
{\vergleichskette
{ X
}
{ =} {\bigcup_{i \in J} U_i
}
{ } {
}
{ } {
}
{ } {
}
}
{}{}{}
berlaku.
}{Ruang kotangen di $P$ adalah
\definitionsverweis {ruang dual}{}{}
dari
\definitionsverweis {ruang singgung}{}{}

\mathl{T_PM}{} di $P$.
}{Titik

\mavergleichskette
{\vergleichskette
{ Q
}
{ \in }{ L
}
{ }{
}
{ }{
}
{ }{
}
}
{}{}{}
disebut reguler bagi $\varphi$ jika
\definitionsverweis {pemetaan tangen}{}{}
\maabbdisp {T_Q(\varphi)} {T_QL} {T_{\varphi(Q)}M
} {}
pada titik $Q$ mempunyai
\definitionsverweis {rank maksimal}{}{.}
}{Misalkan
\mathkor {} {(U_i,U_i',\alpha_i, i \in I)} {dan} {(V_j,V_j',\beta_j, j \in J)} {}
merupakan
\definitionsverweis {atlas-atlas}{}{}
dari
\mathkor {} {M} {dan} {N} {.}
Maka
\definitionsverweis {ruang produk}{}{}

\mathl{M \times N}{} yang dilengkapi dengan
\definitionsverweis {peta-peta}{}{}
\maabbdisp {\alpha_i \times \beta_j} {U_i \times V_j } { U_i' \times V_j'
} {}
\zusatzklammer {dengan \mathlk{(i,j) \in I \times J}{} dan \mathlk{U_i' \times V_j' \subseteq \R^m \times \R^n}{}} {} {}
disebut produk manifold
\mathkor {} {M} {dan} {N} {.}
}{Suatu
koneksi
pada $E$ adalah dekomposisi jumlah langsung dari
\definitionsverweis {bundel tangen}{}{}

\mavergleichskette
{\vergleichskette
{ TE
}
{ = }{ V \oplus H
}
{ }{
}
{ }{
}
{ }{
}
}
{}{}{}
menjadi dua subbundel vektor
\mathkor {} {V} {dan} {H} {,}
dengan

\mavergleichskette
{\vergleichskette
{ V
}
{ \subseteq }{ TE
}
{ }{
}
{ }{
}
{ }{
}
}
{}{}{}
merupakan
\definitionsverweis {bundel vertikal}{}{.}
}

 }



\inputaufgabeklausurloesung
{3}

{

Nyatakan teorema-teorema berikut.
\aufzaehlungdrei{
\stichwort {Teorema tentang hubungan antara kelengkungan dan vektor normal satuan suatu kurva bidang berparameter panjang busur} {.}}{
\stichwort {Theorema egregium} {.}}{
\stichwort {Teorema Green} {}
untuk luas.}

}
{

\aufzaehlungdrei{\faktsituation {Misalkan
\maabb {\gamma} { I } {\R^2
} {}
suatu
\definitionsverweis {kurva}{}{}
\definitionsverweis {berparameter panjang busur}{}{}
yang dua kali
\definitionsverweis {terdiferensialkan secara kontinu}{}{.}}
\faktfolgerung {Maka
\definitionsverweis {kelengkungan}{}{}
dari $\gamma$ di $t$ adalah

\mavergleichskettedisp
{\vergleichskette
{ \kappa(t)
}
{ =} { \left\langle \gamma^{\prime \prime}(t) , N(t) \right\rangle
}
{ = } { \pm \Vert { \gamma^{\prime \prime}(t) } \Vert
}
{ } {
}
{ } {
}
}
{}{}{,}
dengan

\mavergleichskettedisp
{\vergleichskette
{ N(t)
}
{ =} { \begin{pmatrix} -\gamma_2'(t) \\ \gamma_1'(t) \end{pmatrix}
}
{ } {
}
{ } {
}
{ } {
}
}
{}{}{}
suatu vektor normal satuan di $\gamma(t)$.}
\faktzusatz {}
\faktzusatz {}}{\faktsituation {Misalkan

\mavergleichskette
{\vergleichskette
{ Y
}
{ \subseteq }{ \R^3
}
{  }{
}
{    }{
}
{   }{
}
}
{}{}{}
suatu
\definitionsverweis {permukaan berorientasi}{}{}
dan misalkan
\maabbdisp {\varphi} { V } { Y
} {,}

\mavergleichskette
{\vergleichskette
{ V
}
{ \subseteq }{ \R^2
}
{  }{
}
{    }{
}
{   }{
}
}
{}{}{,}
suatu parametrisasi lokal $Y$ yang dua kali terdiferensialkan, dengan parameter $u,v$. Misalkan
\mathl{\left(  g_{ij}  \right)}{}
\definitionsverweis {matriks fundamental pertama}{}{}
pada $V$.}
\faktfolgerung {Maka, dengan menggunakan
\definitionsverweis {simbol-simbol Christoffel}{}{,}
\definitionsverweis {kelengkungan Gauss}{}{}
memenuhi hubungan

\mavergleichskettedisphandlinks
{\vergleichskettedisphandlinks
{ K
}
{ =} { { \frac{ 1 }{ g_{11} } } \left( \partial_2 \Gamma_{11}^2 - \partial_1 \Gamma_{12}^2 + \Gamma^1_{11} \Gamma^2_{12} + \Gamma_{11}^2 \Gamma_{22}^2 - \Gamma_{12}^1 \Gamma_{11}^2- \left( \Gamma_{12}^2 \right)^2 \right)
}
{ } {
}
{ } {
}
{ } {
}
}
{}{}{.}}
\faktzusatz {}
\faktzusatz {}}{Misalkan
\mathl{M \subseteq \R^2}{} suatu
\definitionsverweis {manifold dengan batas}{}{}
yang
\definitionsverweis {kompak}{}{}

\mathl{\delta M}{.} Maka luas $M$ adalah

\mavergleichskettedisp
{\vergleichskette
{  \lambda^2(M)
}
{ =} {
\int_{ \partial M  }  x dy
}
{ =} { -
\int_{ \partial M  }  y dx
}
{ } {
}
{ } {
}
}
{}{}{.}}

 }



\inputaufgabeklausurloesung
{8}

{

Misalkan

\mavergleichskettedisp
{\vergleichskette
{ f(x,y)
}
{ =} { ax^2+by^2+cxy
}
{ } {
}
{ } {
}
{ } {
}
}
{}{}{.}
Tentukan
\definitionsverweis {kelengkungan-kelengkungan utama}{}{}
dan
\definitionsverweis {arah-arah kelengkungan utama}{}{}
dari
\definitionsverweis {grafik}{}{}
$f$ di titik asal.

}
{

Kita menggunakan
Lemma 4.9 (Differentialgeometrie (Osnabrück 2023)),
dengan faktor pendahulu di titik nol sama dengan $1$, sehingga
\definitionsverweis {pemetaan Weingarten}{}{}
dideskripsikan langsung oleh
\definitionsverweis {matriks Hesse}{}{}
. Matriks Hesse fungsi tersebut adalah

\mathdisp {\begin{pmatrix} 2a & c \\ c & 2b \end{pmatrix}} { , }

dan
\definitionsverweis {polinomial karakteristik}{}{}
matriks tersebut adalah

\mavergleichskettealign
{\vergleichskettealign
{ (t-2a)(t-2b)-c^2
}
{ =} { t^2 -2(a+b)t +4ab -c^2
}
{ =} { \left(  t-(a+b)  \right)^2 - (a+b)^2 +4ab -c^2
}
{ =} { \left(  t-(a+b)  \right)^2  - a^2 -b^2 +2ab -c^2
}
{ =} { \left(  t-(a+b)  \right)^2  - (a-b)^2 -c^2
}
}
{}
{}{.}
Dengan demikian, nilai-nilai eigen
\zusatzklammer {yang merupakan kelengkungan-kelengkungan utama} {} {}
adalah

\mavergleichskettedisp
{\vergleichskette
{ \lambda_{1,2}
}
{ =} { \pm \sqrt{ (a-b)^2+c^2 } +a+b
}
{ } {
}
{ } {
}
{ } {
}
}
{}{}{.}
Dalam kasus

\mavergleichskette
{\vergleichskette
{ a
}
{ = }{ b
}
{  }{
}
{    }{
}
{   }{
}
}
{}{}{}
dan

\mavergleichskette
{\vergleichskette
{c
}
{ = }{ 0
}
{ }{
}
{ }{
}
{ }{
}
}
{}{}{}
pemetaan Weingarten adalah perkalian dengan $2a$ dan setiap vektor merupakan vektor eigen; mulai sekarang kasus ini dikecualikan. Untuk

\mavergleichskettedisp
{\vergleichskette
{ \lambda_{1}
}
{ =} { \sqrt{ (a-b)^2+c^2 } +a+b
}
{ } {
}
{ } {
}
{ } {
}
}
{}{}{}
perlu ditentukan kernel matriks

\mathdisp {\begin{pmatrix}  \sqrt{(a-b)^2 +c^2} + b-a  &   -c  \\  -c &  \sqrt{(a-b)^2 +c^2} + a-b \end{pmatrix}} {  }

dan kernel tersebut adalah

\mathdisp {\R \begin{pmatrix}  \sqrt{(a-b)^2 +c^2} + a-b  \\c   \end{pmatrix}} { . }

Jadi, salah satu arah kelengkungan utama adalah
\mathl{\begin{pmatrix}  \sqrt{(a-b)^2 +c^2} + a-b  \\c   \end{pmatrix}}{} dengan kelengkungan utama
\mathl{\sqrt{ (a-b)^2+c^2 } +a+b}{.}

Dengan cara yang sama diperoleh arah kelengkungan utama
\mathl{\begin{pmatrix} -c \\ \sqrt{(a-b)^2 +c^2} + a-b \end{pmatrix}}{} dengan kelengkungan utama
\mathl{-\sqrt{ (a-b)^2+c^2 } +a+b}{.}

 }



\inputaufgabeklausurloesung
{7}

{

Buktikan teorema eksistensi untuk transpor paralel pada suatu hiperpermukaan

\mavergleichskette
{\vergleichskette
{ Y
}
{ \subseteq }{ \R^n
}
{  }{
}
{    }{
}
{   }{
}
}
{}{}{.}

}
{

Kita mencari suatu pemetaan diferensiabel
\maabbdisp {F} { I } { \R^n
} {,}
yang memenuhi syarat-syarat
\aufzaehlungdrei{
\mavergleichskettedisp
{\vergleichskette
{ F(t)
}
{ \in} { T_{\gamma(t)} Y
}
{ } {
}
{ } {
}
{ } {
}
}
{}{}{}
untuk setiap

\mavergleichskette
{\vergleichskette
{ t
}
{ \in }{ I
}
{  }{
}
{    }{
}
{   }{
}
}
{}{}{,}
}{
\mavergleichskettedisp
{\vergleichskette
{ F'(t)
}
{ \in} { N_{\gamma(t)} Y
}
{ } {
}
{ } {
}
{ } {
}
}
{}{}{}
untuk setiap

\mavergleichskette
{\vergleichskette
{ t
}
{ \in }{ I
}
{ }{
}
{ }{
}
{ }{
}
}
{}{}{,}
}{
\mavergleichskettedisp
{\vergleichskette
{ F (t_0)
}
{ =} { v
}
{ } {
}
{ } {
}
{ } {
}
}
{}{}{}
}
Syarat pertama berarti bahwa $F(t)$ tegak lurus terhadap vektor normal satuan $N (\gamma(t))$, sehingga

\mavergleichskettedisp
{\vergleichskette
{ \left\langle F( t) , N( \gamma(t)) \right\rangle
}
{ =} { 0
}
{ } {
}
{ } {
}
{ } {
}
}
{}{}{.}
Oleh karena itu, juga berlaku

\mavergleichskettedisp
{\vergleichskette
{ 0
}
{ =} { \left\langle F(t) , N( \gamma(t)) \right\rangle'
}
{ =} { \left\langle F'(t) , N( \gamma(t)) \right\rangle + \left\langle F(t) , (N( \gamma(t)))' \right\rangle
}
{ } {
}
{ } {
}
}
{}{}{.}
Syarat kedua, yakni bahwa $F'(t)$ terletak di ruang normal, berarti

\mavergleichskettedisp
{\vergleichskette
{ F'(t) - \left\langle F'( t) , N( \gamma(t)) \right\rangle N(\gamma(t) )
}
{ =} { 0
}
{ } {
}
{ } {
}
{ } {
}
}
{}{}{}
untuk setiap

\mavergleichskette
{\vergleichskette
{ t
}
{ \in }{ I
}
{  }{
}
{    }{
}
{   }{
}
}
{}{}{,}
Dengan menggunakan persamaan sebelumnya, syarat ini dapat diubah menjadi

\mavergleichskettedisp
{\vergleichskette
{ F'(t) + \left\langle F( t) , (N( \gamma(t)))' \right\rangle N(\gamma(t) )
}
{ =} { 0
}
{ } {
}
{ } {
}
{ } {
}
}
{}{}{}
untuk setiap

\mavergleichskette
{\vergleichskette
{ t
}
{ \in }{ I
}
{  }{
}
{    }{
}
{   }{
}
}
{}{}{}
atau menjadi

\mavergleichskettedisp
{\vergleichskette
{ F'(t)
}
{ =} { - \left\langle F( t) ,  (N( \gamma(t)))'   \right\rangle N(\gamma(t) )
}
{ } {
}
{ } {
}
{ } {
}
}
{}{}{}
untuk setiap

\mavergleichskette
{\vergleichskette
{ t
}
{ \in }{ I
}
{  }{
}
{    }{
}
{   }{
}
}
{}{}{}
. Ini adalah persamaan diferensial linear biasa orde pertama untuk $F$. Bersama syarat awal

\mavergleichskette
{\vergleichskette
{ F(t_0)
}
{ = }{ v
}
{ }{
}
{ }{
}
{ }{
}
}
{}{}{}
menurut
Satz 56.2 (Analysis (Osnabrück 2021-2023))
persamaan tersebut mempunyai solusi yang tunggal. Jadi, sistem syarat semula mempunyai paling banyak satu solusi. Jika $F$ adalah solusi tunggal persamaan diferensial itu, maka pemetaan tersebut memang merupakan solusi sistem syarat semula.

 }



\inputaufgabeklausurloesung
{8}

{

Tunjukkan bahwa suatu
\definitionsverweis {manifold diferensiabel}{}{}
$M$ berdimensi

\mavergleichskette
{\vergleichskette
{ n
}
{ \geq }{ 1
}
{  }{
}
{    }{
}
{   }{
}
}
{}{}{}
mempunyai tak hingga banyak
\definitionsverweis {difeomorfisme}{}{}
\maabbdisp {\varphi} { M } { M
} {}
.

}
{

Kita meninjau fungsi-fungsi berbentuk
\maabbeledisp {\psi_c} {[0,1]} {[0,1]
} { x } { x + \varphi_c (x)
} {,}
dengan

\mavergleichskettedisp
{\vergleichskette
{  \varphi_c(x)
}
{ =} { \begin{cases}  0 \text{ jika } x \leq { \frac{   1  }{  3 } } \, , \\  c \left(  x- { \frac{   1  }{  3 } }   \right)^2 \left(  x- { \frac{   2 }{  3 } }   \right)^2 \text{ jika } { \frac{   1  }{  3 } } < x < { \frac{   2 }{  3 } }  \, ,  \\  0 \text{ jika } x \geq { \frac{   2 }{  3 } }  \, .  \end{cases}
}
{ } {
}
{ } {
}
{ } {
}
}
{}{}{}
dengan parameter

\mavergleichskette
{\vergleichskette
{ c
}
{ \in }{  \R_{\geq 0}
}
{  }{
}
{    }{
}
{   }{
}
}
{}{}{.}
Fungsi $\varphi_c$ kontinu dan juga terdiferensialkan secara kontinu karena akar-akar polinomial tersebut merupakan akar ganda. Dengan demikian, $\psi_c$ juga terdiferensialkan secara kontinu. Untuk $c$ yang cukup kecil, $\psi_c$ meningkat secara ketat dan mendefinisikan suatu pemetaan bijektif yang terdiferensialkan secara kontinu dari interval satuan ke dirinya sendiri, serta merupakan identitas pada sepertiga bawah dan sepertiga atas.

Dengan $\psi_c$, kita mendefinisikan difeomorfisme bola satuan terbuka ke dirinya sendiri melalui
\maabbeledisp {} { U \left(   0 , 1  \right) } { U \left(   0 , 1  \right)
} { \begin{pmatrix}  x_1  \\\vdots \\  x_n   \end{pmatrix} } { \psi_c \left(  \sqrt{  x_1^2 + \cdots + x_n^2 }  \right) \begin{pmatrix}  x_1  \\\vdots \\  x_n   \end{pmatrix}
} {.}
Jadi, titik-titik diregangkan bergantung pada jari-jarinya, sedangkan pemetaan tersebut merupakan identitas pada
\mathl{U \left(   0 , { \frac{   1  }{  3 } }   \right)}{} dan pada
\mathl{U \left(   0 , 1  \right) \setminus  U \left(   0 ,   { \frac{   2 }{  3 } }   \right)}{}.

Jika sekarang $M$ suatu manifold berdimensi $\geq 1$, kita memilih suatu peta
\maabbdisp {\alpha} { U } {V
} {}
dengan

\mavergleichskette
{\vergleichskette
{ V
}
{ \subseteq }{ \R^n
}
{ }{
}
{ }{
}
{ }{
}
}
{}{}{,}
dan, dengan
\zusatzklammer {memperkecil daerahnya} {} {,}
kita dapat menganggap bahwa $V$
\zusatzklammer {merupakan suatu bola terbuka dan kemudian} {} {}
adalah bola satuan terbuka. Difeomorfisme yang dikonstruksi pada $V$ menghasilkan difeomorfisme pada $U$. Karena difeomorfisme ini merupakan identitas pada bagian luar bola
\zusatzklammer {mulai dari jari-jari ${ \frac{   2 }{  3 } }$} {} {,}
kita dapat memperluasnya secara diferensiabel ke $M$ dengan menggunakan identitas pada $M \setminus \alpha^{-1} \left(  U \left(   0 , { \frac{   2 }{  3 } }   \right)  \right)$.

 }



\inputaufgabeklausurloesung
{3}

{

Misalkan $M$ suatu
\definitionsverweis {manifold diferensiabel}{}{}
dan
\maabbdisp {f} { M } {\R
} {}
suatu
\definitionsverweis {fungsi terdiferensialkan secara kontinu}{}{.}
Misalkan di titik

\mavergleichskette
{\vergleichskette
{ P
}
{ \in }{ M
}
{  }{
}
{    }{
}
{   }{
}
}
{}{}{}
fungsi tersebut mempunyai
\definitionsverweis {ekstremum lokal}{}{.}
Tunjukkan bahwa

\mavergleichskettedisp
{\vergleichskette
{ (df)_P
}
{ =} { 0
}
{ } {
}
{ } {
}
{ } {
}
}
{}{}{.}

}
{

Misalkan

\mavergleichskette
{\vergleichskette
{ P
}
{ \in }{ U
}
{ \subseteq }{ M
}
{ }{
}
{ }{
}
}
{}{}{}
suatu
\definitionsverweis {daerah peta}{}{}
dengan peta
\maabb {\alpha} { U } {V \subseteq \R^n
} {.}
Maka
\mathl{f \circ \alpha^{-1}}{} juga mempunyai ekstremum lokal di $\alpha(P)$.
Oleh karena itu, menurut
Fakt ***** (2)
\definitionsverweis {diferensial total}{}{}
\maabbdisp {\left(D f \circ \alpha^{-1} \right)_{\alpha(P) }} {\R^n } {\R
} {}
adalah pemetaan nol. Karena itu,

\mavergleichskette
{\vergleichskette
{ (df)_P
}
{ = }{ 0
}
{ }{
}
{ }{
}
{ }{
}
}
{}{}{.}

 }



\inputaufgabeklausurloesung
{5}

{

Kita tinjau
$1$-\definitionsverweis {bentuk diferensial}{}{}

\mavergleichskettedisp
{\vergleichskette
{ \omega
}
{ =} { -vdu+udv
}
{ } {
}
{ } {
}
{ } {
}
}
{}{}{}
pada
$1$-\definitionsverweis {bola}{}{}

\mavergleichskette
{\vergleichskette
{ S^1
}
{ \subset }{ \R^2
}
{  }{
}
{    }{
}
{   }{
}
}
{}{}{,}
dengan koordinat ruang sekitar $\R^2$ dinotasikan oleh
\mathkor {} {u} {dan} {v} {.}
Tentukan
\mathl{\pi^* \omega}{} di bawah
\definitionsverweis {pemetaan terdiferensialkan secara kontinu}{}{}
\maabbeledisp {\pi} {\R^2 \setminus \{0\}} { S^{1}
} {(x , y) } { { \frac{   1  }{  \sqrt{x^2 + y^2}  } } (x , y)
} {.}

}
{

Berlaku

\mavergleichskettealign
{\vergleichskettealign
{ \pi^*du
}
{ =} { d (u \circ \pi)
}
{ =} { d   \left(  { \frac{   x  }{  \sqrt{x^2 + y^2}  } }  \right)
}
{ =} { { \frac{   \partial   \left(  { \frac{   x  }{  \sqrt{x^2 + y^2}  } }  \right)    }{  \partial   x  } }   dx + { \frac{   \partial   \left(  { \frac{   x  }{  \sqrt{x^2 + y^2}  } }  \right)    }{  \partial   y  } } dy
}
{ =} { { \frac{   \sqrt{x^2+y^2}     -x^2 (x^2 + y^2)^{-1/2}  }{  x^2+y^2  } }  dx   - { \frac{   xy }{  \sqrt{x^2+y^2}^3  } }  dy
}
}
{
\vergleichskettefortsetzungalign
{ =} { { \frac{   x^2+y^2 -x^2    }{  \sqrt{ x^2+y^2}^3  } }  dx  - { \frac{   xy }{  \sqrt{x^2+y^2}^3  } }  dy
}
{ =} { { \frac{   y^2  }{  \sqrt{ x^2+y^2}^3  } }  dx  - { \frac{   xy }{  \sqrt{x^2+y^2}^3  } }  dy
}
{ } {}
{ } {}
}
{}{}
dan

\mavergleichskettealign
{\vergleichskettealign
{ \pi^*dv
}
{ =} { d (v \circ \pi)
}
{ =} { d \left(  { \frac{   y  }{  \sqrt{x^2 + y^2}  } }  \right)
}
{ =} { { \frac{   \partial   \left(  { \frac{   y  }{  \sqrt{x^2 + y^2}  } }  \right)   }{  \partial   x  } } dx + { \frac{   \partial   \left(  { \frac{   y  }{  \sqrt{x^2 + y^2}  } }  \right)    }{  \partial   y  } } dy
}
{ =} { { \frac{   -xy }{  \sqrt{x^2+y^2}^3  } } dx + { \frac{   \sqrt{x^2+y^2} -y^2 (x^2 + y^2)^{-1/2}  }{  x^2+y^2  } } dy
}
}
{
\vergleichskettefortsetzungalign
{ =} { { \frac{   -xy }{  \sqrt{x^2+y^2}^3  } } dx + { \frac{   x^2  }{  \sqrt{x^2+y^2}^3  } }  dy
}
{ } {
}
{ } {}
{ } {}
}
{}{.}
Dengan demikian,

\mavergleichskettealigndrucklinks
{\vergleichskettealigndrucklinks
{ \pi^*(-vdu+udv)
}
{ =} { - (v \circ \pi)  \pi^*du + (u \circ \pi) \pi^*dv
}
{ =} { - { \frac{   y  }{  \sqrt{x^2+y^2}  } } \left(  { \frac{   y^2  }{  \sqrt{ x^2+y^2}^3  } }  dx  - { \frac{   xy }{  \sqrt{x^2+y^2}^3  } }  dy   \right) + { \frac{   x  }{  \sqrt{x^2+y^2}  } }  \left(  { \frac{   -xy }{  \sqrt{x^2+y^2}^3  } } dx + { \frac{   x^2  }{  \sqrt{x^2+y^2}^3  } } dy  \right)
}
{ =} { { \frac{   1 }{  \left(  x^2+y^2  \right)^2 } }  \left(  \left(  -y^3-x^2y  \right) dx + \left(  xy^2+x^3  \right) dy \right)
}
{ =} { { \frac{   1 }{  x^2+y^2  } }  \left(  -  y dx + x dy \right)
}
}
{}{}{.}

 }



\inputaufgabeklausurloesung
{5}

{

Buktikan teorema tentang perhitungan volume kanonik pada manifold Riemann $M$.

}
{

Menurut
Definition  .
kita harus menghitung bentuk diferensial
\mathl{\left(  \alpha^{-1}  \right)^*\omega}{} untuk setiap titik

\mavergleichskette
{\vergleichskette
{ Q
}
{ \in }{ V
}
{  }{
}
{    }{
}
{   }{
}
}
{}{}{}
. Bentuk ini mempunyai bentuk
\mathl{c_Q dx_1 \wedge \ldots \wedge dx_n}{} dan ditentukan oleh nilainya pada
\mathl{e_1 \wedge \ldots \wedge e_n}{}. Berlaku

\mavergleichskettedisphandlinks
{\vergleichskettedisphandlinks
{ \left(  \alpha^{-1}  \right)^*\omega \left(  Q,  e_1 \wedge \ldots \wedge e_n  \right)
}
{ =} { \omega \left(  \alpha^{-1} (Q) ,T_Q \left(  \alpha^{-1}  \right)(e_1) \wedge \ldots \wedge T_Q \left(  \alpha^{-1}  \right)( e_n)  \right)
}
{ } {
}
{ } {
}
{ } {
}
}
{}{}{.}
Menurut definisi matriks fundamental metrik,

\mavergleichskettedisp
{\vergleichskette
{ g_{ij}(Q)
}
{ =} { \left\langle T_Q \left( \alpha^{-1} \right) (e_i) , T_Q \left( \alpha^{-1} \right) (e_j) \right\rangle_{\alpha^{-1}(Q)}
}
{ } {
}
{ } {
}
{ } {
}
}
{}{}{.}
Menurut
Satz 7.8 (Maß- und Integrationstheorie (Osnabrück 2022-2023)),
berlaku

\mavergleichskettealigndrucklinks
{\vergleichskettealigndrucklinks
{ \omega \left(  \alpha^{-1}(Q), T_Q \left(  \alpha^{-1}  \right) (e_1) \wedge \ldots \wedge T_Q \left(  \alpha^{-1}  \right) ( e_n)  \right)
}
{ =} { \left(  \det  \left(  \left\langle T_Q \left(  \alpha^{-1}  \right)  (e_i)  ,  T_Q \left(  \alpha^{-1}  \right) (e_j)   \right\rangle  \right)_{1 \leq i,j \leq n}  \right)^{1/2}
}
{ =} { \left(  \det  \left(  g_{ij}(Q)   \right)_{1 \leq i,j \leq n}  \right)^{1/2}
}
{ =} { \sqrt{g(Q)}
}
{ } {
}
}
{}{}{.}

 }



\inputaufgabeklausurloesung
{6 (1+1+3+1)}

{

Jelaskan mengapa $\R^2$ tidak membawa struktur
\definitionsverweis {manifold berbatas}{}{}
sedemikian sehingga subhimpunan $T$ yang diberikan merupakan
\definitionsverweis {batas}{}{.}
\aufzaehlungvierabc{
\mavergleichskette
{\vergleichskette
{ T
}
{ = }{ {]0,1[}
}
{  }{
}
{    }{
}
{   }{
}
}
{}{}{,}
dengan interval tersebut terletak pada sumbu $x$.
}{
\mavergleichskette
{\vergleichskette
{ T
}
{ = }{ [0,1]
}
{  }{
}
{    }{
}
{   }{
}
}
{}{}{,}
dengan interval tersebut terletak pada sumbu $x$.
}{
\mavergleichskette
{\vergleichskette
{ T
}
{ = }{ \R
}
{  }{
}
{    }{
}
{   }{
}
}
{}{}{,}
dengan $\R$ adalah sumbu $x$.
}{
\mavergleichskette
{\vergleichskette
{ T
}
{ = }{ S^1
}
{  }{
}
{    }{
}
{   }{
}
}
{}{}{.}
}

}
{

a) Batas
\mathl{\partial M}{} dari suatu manifold dengan batas, menurut
Lemma 21.3 (Differentialgeometrie (Osnabrück 2023))  (2)
\definitionsverweis {tertutup}{}{,}
sedangkan interval terbuka sebagai himpunan bagian dari $\R^2$ tidak tertutup.

b) Batas
\mathl{\partial M}{} dari suatu manifold dengan batas, menurut
Satz 21.4 (Differentialgeometrie (Osnabrück 2023)),
merupakan manifold tanpa batas, sedangkan interval tertutup mempunyai titik-titik batas.

c) Andaikan $\R^2$ merupakan suatu manifold dengan batas
\mathl{\R = \R \times \{0\}}{}. Untuk
\mathl{P \in \R}{}, terdapat suatu lingkungan terbuka
\mathl{P\in U}{} dan suatu homeomorfisme
\maabbdisp {\alpha} { U } {V
} {}
dengan suatu himpunan terbuka $V \subseteq H$, dengan $H$ menyatakan suatu setengah bidang. Kita dapat mengganti $V$ dengan lingkungan setengah bola terbuka yang lebih kecil di sekitar $P$. Menurut
Lemma 21.3 (Differentialgeometrie (Osnabrück 2023))  (2)
peta seperti ini memetakan titik batas ke titik batas; artinya,

\mavergleichskettedisp
{\vergleichskette
{\alpha ( U \cap \R )
}
{ = } { V \cap \delta H
}
{ } {
}
{ } {
}
{ } {
}
}
{}{}{.}
Dengan demikian, diperoleh pula suatu homeomorfisme di antara komplemennya, yaitu antara
\mathkor {} {U \setminus U \cap \R} {dan} {V \setminus V \cap \delta H} {.}
Lingkungan setengah bola di sebelah kanan
\definitionsverweis {terhubung lintasan}{}{,}
sedangkan himpunan di sebelah kiri merupakan gabungan saling lepas dari irisannya dengan setengah bagian positif dan negatif. Kedua irisan tersebut terbuka dan tidak kosong karena $U$ memuat suatu lingkungan bola dari $P$. Jadi, himpunan di sebelah kiri tidak terhubung dan homeomorfisme semacam itu tidak mungkin ada.

d) Sama seperti c).

 }



\inputaufgabeklausurloesung
{4}

{

Misalkan $\omega$ suatu
$n-1$-\definitionsverweis {bentuk diferensial}{}{}
yang terdiferensialkan secara kontinu pada
\definitionsverweis {himpunan terbuka}{}{}

\mavergleichskette
{\vergleichskette
{ U
}
{ \subseteq }{ \R^n
}
{  }{
}
{    }{
}
{   }{
}
}
{}{}{}
dan misalkan
\definitionsverweis {turunan eksterior}{}{}
sama dengan

\mavergleichskettedisp
{\vergleichskette
{d \omega
}
{ =} { fdx_1 \wedge \ldots \wedge dx_n
}
{ } {
}
{ } {
}
{ } {
}
}
{}{}{}
dengan suatu fungsi
\maabb {f} { U } {\R
} {.}
Tunjukkan untuk

\mavergleichskette
{\vergleichskette
{ P
}
{ \in }{ U
}
{  }{
}
{    }{
}
{   }{
}
}
{}{}{}
kesamaan

\mavergleichskettedisp
{\vergleichskette
{ f(P)
}
{ =} { { \frac{   1  }{  \beta_n  } } \cdot \operatorname{lim}_{   \epsilon \rightarrow  0  } \, { \frac{   1  }{  \epsilon^n } } \int_{S^{n-1}(P, \epsilon)} \omega
}
{ } {
}
{ } {
}
{ } {
}
}
{}{}{,}
dengan $\beta_n$ menyatakan
volume
bola satuan berdimensi $n$.

}
{

Kita menerapkan
Satz von Stokes
pada bola-bola
\mathl{B^n(P, \epsilon)}{} dengan batas
\mathl{S^{n-1}(P, \epsilon)}{}
\zusatzklammer {dengan pusat $P$ dan jari-jari $\epsilon$} {} {}
dan memperoleh

\mavergleichskettealign
{\vergleichskettealign
{ \operatorname{lim}_{   \epsilon \rightarrow  0  } \, { \frac{   1 }{  \epsilon^n } }     \int_{S^{n-1}(P, \epsilon)} \omega
}
{ =} { \operatorname{lim}_{   \epsilon \rightarrow  0  } \, { \frac{   1 }{  \epsilon^n } }     \int_{ \partial B^{n}(P, \epsilon)}   \omega
}
{ =} { \operatorname{lim}_{   \epsilon \rightarrow  0  } \, { \frac{   1 }{  \epsilon^n } }     \int_{B^{n}(P, \epsilon)} d \omega
}
{ =} { \operatorname{lim}_{   \epsilon \rightarrow  0  } \, { \frac{   1 }{  \epsilon^n } }     \int_{B^{n}(P, \epsilon)} f d \lambda^n
}
{ =} {\beta_n \cdot f(P)
}
}
{}
{}{,}
dengan kesamaan terakhir menggunakan kekontinuan $f$.

 }



\inputaufgabeklausurloesung
{12}

{

Buktikan teorema Stokes untuk manifold berbatas.

}
{

\teilbeweis {}{}{}
{Misalkan

\mathbed {U_i} {}
 {i \in I} {}
 {} {} {} {,}
suatu
\definitionsverweis {penutup terbuka}{}{}
dari $M$ dengan
\definitionsverweis {peta-peta berorientasi}{}{,}
dan misalkan

\mathbed {h_j} {}
 {j \in J} {}
 {} {} {} {,}
suatu
\definitionsverweis {partisi kesatuan yang terdiferensialkan secara kontinu dan tersubordinasi pada penutup tersebut}{}{,}
yang ada menurut
Satz 22.10 (Differentialgeometrie (Osnabrück 2023)). Untuk setiap

\mavergleichskette
{\vergleichskette
{ P
}
{ \in }{ M
}
{  }{
}
{    }{
}
{   }{
}
}
{}{}{}
terdapat suatu lingkungan terbuka

\mavergleichskette
{\vergleichskette
{ P
}
{ \in }{ V_P
}
{ \subseteq }{ M
}
{    }{
}
{   }{
}
}
{}{}{}
sedemikian sehingga

\mavergleichskette
{\vergleichskette
{ h_j \vert_{V_P }
}
{ = }{ 0
}
{ }{
}
{ }{
}
{ }{
}
}
{}{}{}
berlaku kecuali untuk berhingga banyak indeks. Misalkan $Y$ dukungan dari $\omega$. Penutup

\mavergleichskette
{\vergleichskette
{ Y
}
{ \subseteq }{ \bigcup_{P \in M} V_P
}
{  }{
}
{    }{
}
{   }{
}
}
{}{}{}
karena
\definitionsverweis {kekompakan}{}{}
yang diasumsikan, mempunyai suatu subpenutup berhingga, katakanlah

\mavergleichskettedisp
{\vergleichskette
{ Y
}
{ \subseteq} { V_{1} \cup \ldots \cup V_{r}
}
{ =} { V
}
{ } {
}
{ } {
}
}
{}{}{.}
Oleh karena itu, hanya berhingga banyak $h_j$ pada $V$ yang berbeda dari $0$. Kita menetapkan

\mavergleichskette
{\vergleichskette
{ \omega_j
}
{ = }{ h_j \omega
}
{  }{
}
{    }{
}
{   }{
}
}
{}{}{;}
bentuk-bentuk diferensial ini juga terdiferensialkan secara kontinu. Dukungan $h_j$ merupakan
\definitionsverweis {himpunan bagian tertutup}{}{}
dari $M$ yang terletak di dalam
\mathl{U_{i(j)}}{}. Karena itu, dukungan $\omega_j$ terletak di dalam
\mathl{Y \cap U_{i(j)}}{} dan kompak menurut
Aufgabe 13.10 (Differentialgeometrie (Osnabrück 2023)).
Berlaku

\mavergleichskettedisp
{\vergleichskette
{ \omega
}
{ =} { \sum_{j \in J} \omega_j
}
{ } {
}
{ } {
}
{ } {
}
}
{}{}{,}
dan hanya berhingga banyak bentuk diferensial
\mathl{\omega_j}{} yang berbeda dari $0$ karena

\mavergleichskette
{\vergleichskette
{ \omega_j \vert_{M \setminus Y}
}
{ = }{ 0
}
{ }{
}
{ }{
}
{ }{
}
}
{}{}{}
untuk setiap $j$, serta

\mavergleichskettedisp
{\vergleichskette
{ \omega_j \vert_V
}
{ =} { 0
}
{ } {
}
{ } {
}
{ } {
}
}
{}{}{}
untuk setiap $j$ kecuali berhingga banyak indeks. Berdasarkan
Additivität des Integrals von Differentialformen
dan
Additivität der äußeren Ableitung
pernyataan dapat dibuktikan secara terpisah untuk setiap $\omega_j$. Jadi, kita dapat menganggap bahwa diberikan suatu bentuk diferensial-$(n-1)$ yang terdiferensialkan secara kontinu dan mempunyai dukungan kompak yang seluruhnya terletak dalam suatu lingkungan peta $U$.}
{}
\teilbeweis {}{}{}
{Terdapat suatu
\definitionsverweis {difeomorfisme}{}{}
\maabb {\alpha} { U } { V
} {}
dengan

\mavergleichskette
{\vergleichskette
{ V
}
{ \subseteq }{ H
}
{ }{
}
{ }{
}
{ }{
}
}
{}{}{}
terbuka, yang sekaligus menginduksi suatu difeomorfisme di antara batas-batas

\mavergleichskette
{\vergleichskette
{ \partial U
}
{ = }{ \partial M \cap U
}
{  }{
}
{    }{
}
{   }{
}
}
{}{}{}
dan

\mavergleichskette
{\vergleichskette
{ \partial V
}
{ = }{ \partial H \cap V
}
{  }{
}
{    }{
}
{   }{
}
}
{}{}{}
. Berlaku

\mavergleichskettedisp
{\vergleichskette
{
\int_{ \partial U } \omega
}
{ = } {
\int_{ \partial V } \alpha^{-1 *}\omega
}
{ } {
}
{ } {
}
{ } {
}
}
{}{}{}
dan

\mavergleichskettedisp
{\vergleichskette
{
\int_{ U  }  d \omega
}
{ =} {
\int_{ V  }  \alpha^{-1 *} (  d \omega)
}
{ =} {
\int_{  V  }  d \left(  \alpha^{-1 *}\omega  \right)
}
{ } {
}
{ } {
}
}
{}{}{}
menurut
Lemma 20.2 (Differentialgeometrie (Osnabrück 2023))  (5).}
{\leerzeichen{}Jadi, kita dapat mulai dengan suatu bentuk diferensial yang terdiferensialkan secara kontinu, mempunyai dukungan kompak, dan didefinisikan pada

\mavergleichskette
{\vergleichskette
{ V
}
{ \subseteq }{ H
}
{ }{
}
{ }{
}
{ }{
}
}
{}{}{}
serta memperluasnya menjadi bentuk nol pada seluruh $H$ di luar dukungannya.}
\teilbeweis {}{}{}
{Karena kekompakan dukungan, terdapat suatu balok yang cukup besar

\mavergleichskette
{\vergleichskette
{ Q
}
{ \subseteq }{ H
}
{ }{
}
{ }{
}
{ }{
}
}
{}{}{,}
yang salah satu sisinya $S$ terletak pada $\partial H$ dan yang berpotongan dengan dukungan $\omega$ hanya di $S$. Pada semua sisi lain dari $Q$, bentuk $\omega$
\zusatzklammer {dan dengan demikian juga $d\omega$} {} {}
merupakan bentuk nol. Oleh karena itu, di satu pihak

\mavergleichskettedisp
{\vergleichskette
{
\int_{ H } d\omega
}
{ =} {
\int_{ Q } d\omega
}
{ } {
}
{ } {
}
{ } {
}
}
{}{}{}
dan di pihak lain

\mavergleichskettedisp
{\vergleichskette
{
\int_{ \partial H } \omega
}
{ = } {
\int_{ S } \omega
}
{ = } {
\int_{ \partial Q } \omega
}
{ } {
}
{ } {
}
}
{}{}{.}
Dengan demikian, pernyataan tersebut mengikuti dari Quaderversion des Satzes von Stokes.}
{}

 }
\endgroup

\chapter{Formulir Solusi Resmi Ujian 3}
\noindent\textit{Solusi yang disediakan oleh sumber resmi; kekosongan sumber dipertahankan.}
\begingroup
\renewcommand{\theHfakt}{o011.exam.official.03.\arabic{fakt}}
%Data institusi

%\input{Dozentdaten}

%\renewcommand{\fachbereich}{Fachbereich}

%\renewcommand{\dozent}{Prof. Dr. . }

%Data ujian

\renewcommand{\klausurgebiet}{ }

\renewcommand{\klausurtyp}{ }

\renewcommand{\klausurdatum}{ . 20}

\klausurvorspann {\fachbereich} {\klausurdatum} {\dozent} {\klausurgebiet} {\klausurtyp}

%Data untuk tabel nilai berikut

\renewcommand{\aeins}{  3  }

\renewcommand{\azwei}{  3   }

\renewcommand{\adrei}{  0  }

\renewcommand{\avier}{  0  }

\renewcommand{\afuenf}{  6  }

\renewcommand{\asechs}{  6  }

\renewcommand{\asieben}{  0  }

\renewcommand{\aacht}{  6  }

\renewcommand{\aneun}{  11  }

\renewcommand{\azehn}{  5  }

\renewcommand{\aelf}{  5  }

\renewcommand{\azwoelf}{  3  }

\renewcommand{\adreizehn}{  7   }

\renewcommand{\avierzehn}{  2  }

\renewcommand{\afuenfzehn}{  0  }

\renewcommand{\asechzehn}{  6  }

\renewcommand{\asiebzehn}{  63  }

\renewcommand{\aachtzehn}{    }

\renewcommand{\aneunzehn}{    }

\renewcommand{\azwanzig}{    }

\renewcommand{\aeinundzwanzig}{    }

\renewcommand{\azweiundzwanzig}{    }

\renewcommand{\adreiundzwanzig}{    }

\renewcommand{\avierundzwanzig}{    }

\renewcommand{\afuenfundzwanzig}{    }

\renewcommand{\asechsundzwanzig}{    }

\punktetabellesechzehn

\klausurnote

\newpage

\setcounter{section}{K}



\inputaufgabeklausurloesung
{3}

{

Definisikan istilah-istilah berikut
\zusatzklammer {yang dicetak miring} {} {.}
\aufzaehlungsechs{Suatu
\stichwort {kurva geodesik} {}
pada hiperpermukaan diferensiabel

\mavergleichskette
{\vergleichskette
{ Y
}
{ \subseteq }{ \R^n
}
{  }{
}
{    }{
}
{   }{
}
}
{}{}{.}

}{
\stichwort {Himpunan rotasi} {}
\zusatzklammer {mengelilingi sumbu $x$} {} {}
dari suatu subhimpunan

\mavergleichskette
{\vergleichskette
{ T
}
{ \subseteq }{  \R \times \R_{\geq 0}
}
{  }{
}
{    }{
}
{   }{
}
}
{}{}{.}

}{
\stichwort {Pemetaan tangen} {}
\maabbdisp {T (\varphi)} {TM} {TN
} {}
dari
\definitionsverweis {pemetaan diferensiabel}{}{}
\maabbdisp {\varphi} { M } { N
} {}
antara
\definitionsverweis {manifold-manifold diferensiabel}{}{}
\mathkor {} {M} {dan} {N} {.}

}{Suatu perubahan peta yang
\stichwort {mempertahankan orientasi} {}
pada
\definitionsverweis {manifold diferensiabel}{}{}
$M$.

}{
\stichwort {Bentuk diferensial tarikan balik} {} $\varphi^* \omega$ dari suatu bentuk diferensial
\mathl{\omega \in
{ \mathcal E }^{ k } (  M   )}{} terhadap suatu pemetaan terdiferensialkan secara kontinu
\maabb {\varphi} { L } { M
} {}
antara dua manifold diferensiabel
\mathkor {} {L} {dan} {M} {.}

}{Suatu koneksi linear yang
\stichwort {metrik} {}
pada bundel tangen suatu manifold Riemann $M$.
}

}
{

\aufzaehlungsechs{Suatu
\zusatzklammer {sebagai pemetaan ke $\R^n$} {} {}
\definitionsverweis {kurva yang terdiferensialkan secara kontinu dua kali}{}{}
\maabbdisp {\gamma} {I} {Y
} {}
disebut
kurva geodesik
pada $Y$ jika
\definitionsverweis {percepatan tangensial}{}{}
kurva tersebut sama dengan $0$ di setiap titik.
}{Himpunan rotasi dari $T$ adalah

\mathdisp {\left\{ (x, y \cos \alpha, y \sin \alpha  ) \in \R^3  \mid   (x,y) \in T , \, \alpha \in [0, 2 \pi]       \right\}} { . }

}{Yang dimaksud dengan pemetaan tangen
\maabbdisp {T (\varphi)} {TM} {TN
} {}
adalah gabungan saling lepas dari
\definitionsverweis {pemetaan-pemetaan tangen}{}{}
di setiap titik, yakni

\mavergleichskettedisp
{\vergleichskette
{ T (\varphi)
}
{ =} { \biguplus_{P \in M} T_P(\varphi)
}
{ } {
}
{ } {
}
{ } {
}
}
{}{}{.}
}{Misalkan
\maabbdisp {\alpha_1} {U_1 } { V_1
} {}
dan
\maabbdisp {\alpha_2} {U_2 } { V_2
} {}
adalah peta-peta berorientasi dari $M$. Perubahan peta terkait
\definitionsverweis {perubahan peta}{}{}
\maabbdisp {\psi=\alpha_2 \circ \alpha_1^{-1}} {V_1 \cap \alpha_1 (U_1 \cap U_2)} {V_2 \cap \alpha_2 (U_1 \cap U_2 )
} {}
disebut mempertahankan orientasi jika untuk setiap titik
\mathl{Q \in V_1 \cap \alpha_1 (U_1 \cap U_2)}{},
\definitionsverweis {diferensial total}{}{}
\maabbdisp {\left(D\psi\right)_{Q}} {\R^n } {\R^n
} {}
\definitionsverweis {mempertahankan orientasi}{}{}.
}{Bentuk diferensial tarikan balik
\mathl{\varphi^* \omega}{} didefinisikan, untuk
\mathl{P\in L}{} dan
\mathl{v_1 , \ldots , v_k \in T_PL}{}, oleh

\mavergleichskettedisp
{\vergleichskette
{  ( \varphi^* \omega )(P,v_1 , \ldots , v_k )
}
{ \defeq} { \omega(\varphi(P), T_P\varphi(v_1) , \ldots ,  T_P\varphi(v_k) )
}
{ } {
}
{ } {
}
{ } {
}
}
{}{}{}
.
}{Koneksi tersebut disebut
metrik
jika

\mavergleichskettedisp
{\vergleichskette
{ D_V \left\langle W , Z  \right\rangle
}
{ =} { \left\langle \nabla_V W , Z  \right\rangle  + \left\langle W ,  \nabla_VZ  \right\rangle
}
{ } {
}
{ } {
}
{ } {
}
}
{}{}{.}
berlaku untuk sebarang medan vektor $V,W ,Z$.
}

 }



\inputaufgabeklausurloesung
{3}

{

Nyatakan teorema-teorema berikut.
\aufzaehlungdrei{
\stichwort {Teorema tentang solusi persamaan diferensial biasa pada hiperpermukaan diferensiabel} {}

\mavergleichskette
{\vergleichskette
{  Y
}
{ \subseteq }{  \R^n
}
{  }{
}
{    }{
}
{   }{
}
}
{}{}{.}}{Rumus untuk menghitung volume kanonik suatu manifold Riemann dalam sebuah peta.}{Teorema tentang partisi kesatuan.}

}
{

\aufzaehlungdrei{\faktsituation {Misalkan

\mavergleichskette
{\vergleichskette
{ W
}
{ \subseteq }{ \R^n
}
{  }{
}
{    }{
}
{   }{
}
}
{}{}{}
\definitionsverweis {terbuka}{}{,}
\maabb {h} { W } { \R
} {}
suatu
\definitionsverweis {fungsi yang terdiferensialkan secara kontinu}{}{}
dan

\mavergleichskette
{\vergleichskette
{ Y
}
{ = }{ h^{-1}(c)
}
{  }{
}
{    }{
}
{   }{
}
}
{}{}{}
\definitionsverweis {serat}{}{}
atas

\mavergleichskette
{\vergleichskette
{ c
}
{ \in }{ \R
}
{  }{
}
{    }{
}
{   }{
}
}
{}{}{,}
dengan $h$ bersifat
\definitionsverweis {reguler}{}{}
di setiap titik $Y$.
Misalkan

\mavergleichskette
{\vergleichskette
{ Y
}
{ \subseteq }{ U
}
{ \subseteq }{ \R^n
}
{    }{
}
{   }{
}
}
{}{}{}
\definitionsverweis {terbuka}{}{}
dan misalkan $F$ suatu
\definitionsverweis {medan vektor}{}{}
terdiferensialkan pada $U$}
\faktvoraussetzung {dengan sifat bahwa untuk setiap titik

\mavergleichskette
{\vergleichskette
{ P
}
{ \in }{ Y
}
{  }{
}
{    }{
}
{   }{
}
}
{}{}{}
vektor $F(P)$ termasuk dalam
\definitionsverweis {ruang tangen serat}{}{}
di $P$ pada $Y$.}
\faktfolgerung {Maka solusi $v(t)$ dari
\definitionsverweis {masalah nilai awal}{}{}
untuk $F$ dengan syarat awal

\mavergleichskette
{\vergleichskette
{ v(t_0)
}
{ \in }{ Y
}
{  }{
}
{    }{
}
{   }{
}
}
{}{}{}
seluruhnya terletak di $Y$.}
\faktzusatz {}
\faktzusatz {}}{Misalkan $M$ suatu
\definitionsverweis {manifold Riemann}{}{}
\definitionsverweis {berorientasi}{}{}
dan $\omega$ adalah
\definitionsverweis {bentuk volume kanonik}{}{.}
Misalkan
\maabbdisp {\alpha} { U } { V
} {}
suatu
\definitionsverweis {peta berorientasi}{}{}
dengan

\mavergleichskettedisp
{\vergleichskette
{ V
}
{ \subseteq} { \R^n
}
{ } {
}
{ } {
}
{ } {
}
}
{}{}{}
terbuka, berkoordinat
\mathl{x_1 , \ldots , x_n}{}, dengan
\definitionsverweis {matriks fundamental metrik}{}{}

\mathl{G=(g_{ij})_{1 \leq i,j \leq n}}{} dan
\mathl{g= \det  G}{.} Maka

\mavergleichskettedisp
{\vergleichskette
{  \alpha_* \omega
}
{ =} { \sqrt{ g} dx_1 \wedge \ldots \wedge dx_n
}
{ } {
}
{ } {
}
{ } {
}
}
{}{}{.}
Untuk suatu
\definitionsverweis {subhimpunan terukur}{}{}

\mavergleichskette
{\vergleichskette
{ T
}
{ \subseteq }{ U
}
{  }{
}
{    }{
}
{   }{
}
}
{}{}{}
berlaku

\mavergleichskettedisp
{\vergleichskette
{
\int_{  T   }  \omega
}
{ =} {
\int_{ \alpha(T)  }   \sqrt{g} dx_1 \wedge \ldots \wedge dx_n
}
{ =} { \int_{ \alpha(T)  }   \sqrt{g}    \,  d \lambda^n
}
{ } {
}
{ } {
}
}
{}{}{.}}{Misalkan $M$ suatu
\definitionsverweis {manifold diferensiabel}{}{}
dengan
\definitionsverweis {basis terhitung}{}{}
bagi topologinya. Maka untuk setiap penutup terbuka terdapat
\definitionsverweis {partisi kesatuan}{}{}
yang
\definitionsverweis {terdiferensialkan secara kontinu}{}{}
dan subordinat terhadap penutup tersebut.}

 }



\inputaufgabeklausurloesung
{0}

{

}
{  }



\inputaufgabeklausurloesung
{0}

{

}
{  }



\inputaufgabeklausurloesung
{6}

{

Untuk permukaan $Y$ yang diberikan oleh

\mavergleichskettedisp
{\vergleichskette
{ x^2+y^2+2z^2
}
{ =} { 4
}
{ } {
}
{ } {
}
{ } {
}
}
{}{}{}
dan titik

\mavergleichskettedisp
{\vergleichskette
{ P
}
{ =} { \left(  1  , \,   1  , \,   1                 \right)
}
{ \in} { Y
}
{ } {
}
{ } {
}
}
{}{}{,}
tentukan suatu matriks diagonal bagi
\definitionsverweis {pemetaan Weingarten}{}{}
$L_P$.

}
{

Medan gradien ternormalisasi adalah

\mavergleichskettealign
{\vergleichskettealign
{ N(x,y,z)
}
{ =} { { \frac{   1  }{  \sqrt{ 4x^2+4y^2+16 z^2}  } } \begin{pmatrix}  2x \\ 2y\\  4z   \end{pmatrix}
}
{ =} { { \frac{   1  }{  \sqrt{ x^2+y^2+4 z^2}  } } \begin{pmatrix}  x  \\ y \\  2z   \end{pmatrix}
}
{ } {
}
{ } {
}
}
{}
{}{.}
Kita bekerja dengan medan normal satuan

\mavergleichskettedisp
{\vergleichskette
{ M(x,y,z)
}
{ =} { { \frac{   1  }{  \sqrt{ 4+2 z^2}  } } \begin{pmatrix}  x  \\ y \\  2z   \end{pmatrix}
}
{ } {
}
{ } {
}
{ } {
}
}
{}{}{,}
yang pada $Y$ berimpit dengan $N$. Diferensial total dari $M$ adalah

\mathdisp {\begin{pmatrix}  { \frac{   1  }{  \sqrt{4+2z^2}  } }   &   0 &  { \frac{   -2 xz }{  \left(  4+2z^2  \right)^{3/2}  } }  \\  0  &  { \frac{   1  }{  \sqrt{4+2z^2}  } }   &  { \frac{   -2 yz }{  \left(  4+2z^2  \right)^{3/2}  } }   \\ 0   &  0  &  { \frac{   8 }{  \left(  4+2z^2  \right)^{3/2}  } }  \end{pmatrix}} { . }

Di titik $P$ yang diberikan, gradiennya adalah $\begin{pmatrix}  2  \\ 2\\  4   \end{pmatrix}$ dan kedua vektor
\mathkor {} {\begin{pmatrix}  1  \\ -1\\  0   \end{pmatrix}} {dan} {\begin{pmatrix}  2  \\ 0 \\  -1   \end{pmatrix}} {}
membentuk suatu basis ruang tangen $T_PY$. Diferensial total dari $M$ di titik ini sama dengan

\mathdisp {\begin{pmatrix}  { \frac{   1  }{  \sqrt{6}  } }   &   0 &  { \frac{   -1 }{  3 \sqrt{6}  } }  \\  0  &  { \frac{   1  }{  \sqrt{6}  } }   &  { \frac{   -1  }{  3 \sqrt{6}  } }   \\ 0   &  0  &  { \frac{   4  }{  3 \sqrt{6}  } }   \end{pmatrix}} { . }

Jika diterapkan pada vektor basis pertama $\begin{pmatrix}  1  \\ -1\\  0   \end{pmatrix}$, hasilnya adalah $\begin{pmatrix}  { \frac{   1  }{  \sqrt{6}  } }  \\ - { \frac{   1  }{  \sqrt{6}  } } \\  0   \end{pmatrix}$, jadi ini merupakan vektor eigen dengan nilai eigen ${ \frac{   1  }{  \sqrt{6}  } }$. Jika diterapkan pada vektor basis kedua $\begin{pmatrix}  2  \\ 0 \\  -1   \end{pmatrix}$, hasilnya adalah

\mavergleichskettedisp
{\vergleichskette
{ \begin{pmatrix}  { \frac{   7 }{  3 \sqrt{6}  } }  \\ { \frac{   1 }{  3 \sqrt{6 } }}  \\  - { \frac{   4  }{  3 \sqrt{6}  } }   \end{pmatrix}
}
{ =} { - { \frac{   1 }{  3 \sqrt{6 } }} \begin{pmatrix}  1  \\ -1\\  0   \end{pmatrix} +   { \frac{   4  }{  3 \sqrt{6}  } } \begin{pmatrix}  2  \\ 0 \\  -1   \end{pmatrix}
}
{ } {
}
{ } {
}
{ } {
}
}
{}{}{.}
Karena itu, nilai eigen yang lain sama dengan ${ \frac{   4  }{  3 \sqrt{6}  } }$, dan matriks diagonal yang merepresentasikannya adalah

\mathdisp {\begin{pmatrix}  { \frac{   1  }{  \sqrt{6}  } }   &   0  \\  0  &  { \frac{   4  }{  3 \sqrt{6}  } }  \end{pmatrix}} { . }

 }



\inputaufgabeklausurloesung
{6}

{

Buktikan teorema isometri untuk transpor paralel pada hiperpermukaan

\mavergleichskette
{\vergleichskette
{ Y
}
{ \subseteq }{ \R^n
}
{  }{
}
{    }{
}
{   }{
}
}
{}{}{.}

}
{

Untuk membuktikan linearitas, misalkan diberikan

\mavergleichskette
{\vergleichskette
{ v,w
}
{ \in }{ T_PY
}
{  }{
}
{    }{
}
{   }{
}
}
{}{}{}
dan

\mavergleichskette
{\vergleichskette
{ r,s
}
{ \in }{ \R
}
{  }{
}
{    }{
}
{   }{
}
}
{}{}{}.
Misalkan
\mathkor {} {F} {dan} {G} {}
adalah medan vektor paralel sepanjang $\gamma$ yang ditentukan secara unik menurut
Teorema 6.9 (Geometri diferensial (Osnabrück 2023))
dengan

\mavergleichskette
{\vergleichskette
{ F(a)
}
{ = }{ v
}
{  }{
}
{    }{
}
{   }{
}
}
{}{}{}
dan

\mavergleichskette
{\vergleichskette
{ G(a)
}
{ = }{ w
}
{  }{
}
{    }{
}
{   }{
}
}
{}{}{.}
Menurut
Lema 6.7 (Geometri diferensial (Osnabrück 2023)),
\mathl{rF+sG}{} adalah medan vektor paralel dengan

\mavergleichskette
{\vergleichskette
{ (rF+sG)(a)
}
{ = }{ v+ w
}
{  }{
}
{    }{
}
{   }{
}
}
{}{}{.}
Karena keunikan dalam
Teorema 6.9 (Geometri diferensial (Osnabrück 2023)),
\mathl{rF+sG}{} dengan demikian adalah medan vektor paralel untuk vektor tangen
\mathl{rv+sw}{.} Oleh karena itu,

\mavergleichskettedisp
{\vergleichskette
{ \Psi_{\gamma} (rv+sw)
}
{ =} { (rF+sG)(b)
}
{ =} { r F(b) +s G(b)
}
{ =} { r \Psi_{\gamma} (v) + s \Psi_{\gamma} (w)
}
{ } {
}
}
{}{}{.}

Untuk membuktikan kesesuaian dengan hasil kali skalar, misalkan kembali diberikan

\mavergleichskette
{\vergleichskette
{ v,w
}
{ \in }{ T_PY
}
{  }{
}
{    }{
}
{   }{
}
}
{}{}{}
dan misalkan $F,G$ adalah medan-medan vektor paralel yang terkait. Kita peroleh

\mavergleichskettedisp
{\vergleichskette
{ \left\langle  F(t) , G(t)  \right\rangle'
}
{ =} { \left\langle  F'(t) , G(t)  \right\rangle  +  \left\langle  F(t) , G'(t)  \right\rangle
}
{ =} { 0
}
{ } {
}
{ } {
}
}
{}{}{,}
karena $F,G$ tangensial dan $F',G'$ ortogonal terhadap ruang tangen. Jadi
\mathl{\left\langle  F(t) , G(t)  \right\rangle}{} konstan sepanjang lintasan. Oleh karena itu,

\mavergleichskettedisp
{\vergleichskette
{ \left\langle \Psi_\gamma(v)  , \Psi_\gamma(w)  \right\rangle
}
{ =} { \left\langle  F(b)  ,  G(b)   \right\rangle
}
{ =} { \left\langle  F(a)  ,  F(a)   \right\rangle
}
{ =} { \left\langle  v  , w  \right\rangle
}
{ } {
}
}
{}{}{.}
Bijektivitasnya juga jelas.

 }



\inputaufgabeklausurloesung
{0}

{

}
{  }



\inputaufgabeklausurloesung
{6}

{

Misalkan $M$ suatu manifold topologis kompak berdimensi $d$,

\mavergleichskette
{\vergleichskette
{ d
}
{ \geq }{ 1
}
{  }{
}
{    }{
}
{   }{
}
}
{}{}{.}
Tunjukkan bahwa terdapat suatu subhimpunan terbuka terbatas

\mavergleichskette
{\vergleichskette
{ U
}
{ \subseteq }{ \R^d
}
{  }{
}
{    }{
}
{   }{
}
}
{}{}{}
dan suatu pemetaan kontinu surjektif
\maabbdisp {\varphi} { U } { M
} {}
.

}
{

Untuk setiap titik
\mathl{P \in M}{}, pilih suatu lingkungan peta terbuka
\mathl{P \in U_P}{} dengan peta
\maabbdisp {\alpha_P} {U_P } {V_P
} {}
dengan
\mathl{V_P \subseteq \R^d}{.} Dengan beralih ke suatu lingkungan bola dari
\mathl{\alpha_P(P) \in V_P}{} beserta prapetanya, kita dapat mengasumsikan bahwa $V_p$ adalah bola-bola terbuka yang jari-jarinya paling besar
\mathl{{ \frac{   1 }{  3 } }}{}. Himpunan-himpunan

\mathbed {U_P} {}
 {P\in M} {}
 {} {} {} {,}
menutupi manifold tersebut. Karena kekompakan $M$, terdapat berhingga banyak titik
\mathl{P_1 , \ldots , P_n}{} sedemikian sehingga

\mathbed {U_i=U_{P_i}} {}
 {i=1 , \ldots , n} {}
 {} {} {} {,}
juga menutupi manifold tersebut. Kita tempatkan bola-bola terbuka
\mathl{V_i=V_{P_i}}{} dalam
\zusatzklammer {\anfuehrung{$\R^d$ yang baru}{}} {} {},
dengan titik-titik pusat

\mathdisp {(1,0 , \ldots , 0), (2,0 , \ldots , 0) , \ldots , (n,0 , \ldots , 0)} { . }

Karena syarat jari-jari, bola-bola ini saling lepas. Kita tinjau himpunan terbuka terbatas
\mathl{U= \bigcup_{i=1}^n V_i}{.} Pemetaan-pemetaan peta $\alpha_i$ menghasilkan pemetaan kontinu

\mathdisp {\varphi_i :V_i \stackrel{ \left(  \alpha_i  \right)^{-1} }{\longrightarrow} U_i \longrightarrow M} { . }

Karena saling lepas, pemetaan-pemetaan ini mendefinisikan, melalui
\mathl{\varphi(z)= \varphi_i(z)}{,} apabila
\mathl{z \in V_i}{}, suatu pemetaan kontinu
\maabbdisp {\varphi} { U } {M
} {.}
Karena sifat penutupnya, pemetaan ini surjektif.

 }



\inputaufgabeklausurloesung
{11 (1+3+1+2+4)}

{

Kita tinjau himpunan

\mavergleichskettedisp
{\vergleichskette
{N
}
{ =} { \left\{ A= \begin{pmatrix}  a   &   b   \\  c  & d \end{pmatrix}   \mid   A \text{ nilpoten}         \right\}
}
{ \subseteq} {\R^4
}
{ } {
}
{ } {
}
}
{}{}{}
dari
$2\times 2$-\definitionsverweis {matriks}{}{}
real nilpoten, serta himpunan

\mavergleichskettedisp
{\vergleichskette
{M
}
{ =} { N \setminus \{ \begin{pmatrix}  0   &   0  \\  0  &  0 \end{pmatrix}  \}
}
{ } {
}
{ } {
}
{ } {
}
}
{}{}{.}

a) Apakah $N$ terhubung?

b) Tunjukkan bahwa $M$ suatu
\definitionsverweis {submanifold tertutup}{}{}
dari suatu subhimpunan terbuka

\mavergleichskette
{\vergleichskette
{ G
}
{ \subseteq }{ \R^4
}
{  }{
}
{    }{
}
{   }{
}
}
{}{}{.}

c) Tentukan
\definitionsverweis {dimensi}{}{}
$M$.

d) Apakah $M$ terhubung?

e) Tutupi
\mathl{M}{} dengan peta-peta topologis eksplisit.

}
{

a) Setiap matriks nilpoten
\mathl{\begin{pmatrix}  a   &   b   \\  c  & d \end{pmatrix}}{} dapat dihubungkan dengan matriks nol di dalam himpunan matriks nilpoten melalui lintasan linear
\maabbeledisp {} {[0,1]} {N
} { t } {t \begin{pmatrix}  a   &   b   \\  c  & d \end{pmatrix}
} {,}.
Karena itu, $N$ terhubung-lintasan dan dengan demikian juga terhubung.

b) Suatu
$2 \times 2$-\definitionsverweis {matriks}{}{}
bersifat nilpoten tepat ketika jejak dan determinannya sama dengan $0$. Jadi, himpunan $N$ dari matriks-matriks nilpoten dapat dipandang sebagai

\mathdisp {\left\{  (a,b,c,d)  \mid   a+d = 0 \text{ dan } ad-bc = 0         \right\}} {  }

Kita tinjau pemetaan
\maabbeledisp {} {\R^4} {\R^2
} {(a,b,c,d)} { (a+d, ad-bc)
} {.}
Matriks Jacobi pemetaan ini adalah

\mathdisp {\begin{pmatrix}  1 &  0 &  0 &  1 \\  d  &  -c &  -b &  a  \end{pmatrix}} { . }

Pemetaan ini tidak
\definitionsverweis {reguler}{}{,}
di titik nol, tetapi reguler di setiap titik lain pada serat $N$. Memang, jika
\mathl{P \in M}{}, maka dari

\mavergleichskettedisp
{\vergleichskette
{a^2
}
{ =} {bc
}
{ } {
}
{ } {
}
{ } {
}
}
{}{}{}
dan

\mavergleichskettedisp
{\vergleichskette
{b
}
{ =} { c
}
{ =} { 0
}
{ } {
}
{ } {
}
}
{}{}{}
langsung diperoleh

\mavergleichskettedisp
{\vergleichskette
{a
}
{ =} {b
}
{ =} { c
}
{ =} { d
}
{ =} { 0
}
}
{}{}{.}
Jadi, matriks Jacobi memiliki rank maksimal di titik-titik dalam $M$. Dengan

\mavergleichskettedisp
{\vergleichskette
{G
}
{ =} {\R^4 \setminus \{0\}
}
{ } {
}
{ } {
}
{ } {
}
}
{}{}{}
kita dapat memandang $M$ sebagai serat pemetaan terbatas yang reguler di seluruh serat tersebut. Oleh karena itu, menurut
Teorema 8.2 (Geometri diferensial (Osnabrück 2023)),
$M$ adalah submanifold tertutup dari $G$.

c) Menurut
Teorema 8.2 (Geometri diferensial (Osnabrück 2023)),
dimensi $M$ sama dengan

\mavergleichskette
{\vergleichskette
{ 4-2
}
{ = }{ 2
}
{  }{
}
{    }{
}
{   }{
}
}
{}{}{.}

d) Kita tulis

\mavergleichskettedisp
{\vergleichskette
{M_+
}
{ =} { \left\{ (a,b,c,d) \in M  \mid   b>c        \right\}
}
{ } {
}
{ } {
}
{ } {
}
}
{}{}{}
dan

\mavergleichskettedisp
{\vergleichskette
{M_-
}
{ =} { \left\{ (a,b,c,d) \in M  \mid   b<c        \right\}
}
{ } {
}
{ } {
}
{ } {
}
}
{}{}{,}
Keduanya,
\zusatzklammer {sebagai irisan $M$ dengan himpunan yang diberikan oleh
\mathl{b>c}{} dan terbuka di $\R^4$} {} {},
merupakan himpunan terbuka dalam $M$. Matriks-matriks
\mathkor {} {\begin{pmatrix}  0  &   1  \\  0 &  0 \end{pmatrix}} {dan} {\begin{pmatrix}  0  &   0  \\  1 &  0 \end{pmatrix}} {}
menunjukkan bahwa keduanya tidak kosong. Selain itu, keduanya menutupi seluruh $M$. Memang, jika

\mavergleichskette
{\vergleichskette
{b
}
{ = }{ c
}
{  }{
}
{    }{
}
{   }{
}
}
{}{}{}
maka dari

\mavergleichskettedisp
{\vergleichskette
{ 0
}
{ =} { ad-bc
}
{ =} { -a^2 -b^2
}
{ } {
}
{ } {
}
}
{}{}{}
langsung diperoleh

\mavergleichskette
{\vergleichskette
{a
}
{ = }{b
}
{ = }{ c
}
{ =   }{ d
}
{ =  }{ 0
}
}
{}{}{,}
dan titik tersebut tidak termasuk dalam $M$. Jadi terdapat penutup oleh dua himpunan terbuka tak kosong yang saling lepas; karena itu $M$ tidak terhubung.

e) Kita bekerja dengan pemetaan
\maabbeledisp {\psi} {\R^2 \setminus \{(0,0)\}} { M_+
} {(a,u)} { (a, u + \sqrt{a^2+u^2}, u - \sqrt{a^2 +u^2},-a)
} {.}
Karena

\mavergleichskettealign
{\vergleichskettealign
{ad-bc
}
{ =} { -a^2 - \left(  u + \sqrt{a^2+u^2}  \right) \left( u - \sqrt{a^2+u^2}  \right)
}
{ =} { -a^2 - \left(  u^2 - \left(  a^2+u^2 \right)   \right)
}
{ =} { 0
}
{ } {
}
}
{}
{}{}
syarat determinan dipenuhi, dan karena

\mavergleichskettedisp
{\vergleichskette
{b-c
}
{ =} { u + \sqrt{a^2+u^2} - \left(  u- \sqrt{a^2 +u^2})  \right)
}
{ =} { 2 \sqrt{a^2+u^2}
}
{ >} { 0
}
{ } {
}
}
{}{}{}
citranya termasuk dalam $M_+$. Pemetaan
\maabbeledisp {\varphi} {M_+} { \R^2 \setminus \{(0,0)\}
} {(a,b,c,d)} { \left(  a  , \,   { \frac{  b+c }{  2 } }                   \right)
} {,}
merupakan invers dari pemetaan yang diberikan. Di sini,

\mavergleichskettedisp
{\vergleichskette
{ \varphi \circ \psi
}
{ =} {
\operatorname{Id}_{ \R^2 \setminus \{(0,0)\}  }
}
{ } {
}
{ } {
}
{ } {
}
}
{}{}{}
jelas. Identitas yang lain diperoleh dari

\mavergleichskettealign
{\vergleichskettealign
{ { \frac{  b+c }{  2 } } + \sqrt{a^2 + \left( { \frac{  b+c }{  2 } }  \right)^2 }
}
{ =} { { \frac{  b+c }{  2 } } + { \frac{   1 }{  2 } }  \sqrt{4 a^2 + b^2 +c^2 +2bc }
}
{ =} { { \frac{  b+c }{  2 } } + { \frac{   1 }{  2 } }  \sqrt{ b^2 +c^2 -2bc }
}
{ =} { { \frac{  b+c }{  2 } } + { \frac{  b-c }{  2 } }
}
{ =} {b
}
}
{}
{}{}
dan

\mavergleichskettedisp
{\vergleichskette
{ { \frac{  b+c }{  2 } } - \sqrt{a^2 + \left(   { \frac{  b+c }{  2 } }   \right)^2 }
}
{ =} { { \frac{  b+c }{  2 } } - { \frac{  b-c }{  2 } }
}
{ =} { c
}
{ } {}
{ } {}
}
{}{}{.}
Kedua pemetaan kontinu, sehingga diperoleh suatu
\definitionsverweis {homeomorfisme}{}{}.
Untuk
\mathl{M_-}{}, tukarkan peran
\mathkor {} {b} {dan} {c} {.}

 }



\inputaufgabeklausurloesung
{5}

{

Katakan sesuatu yang cerdas tentang pita Möbius.

}
{  }



\inputaufgabeklausurloesung
{5}

{

Kita tinjau pemetaan
\maabbeledisp {\varphi} { \R^3 } { \R^2
} { (x,y,z) } { \left(  xyz^2,xy-z^3  \right)
} {.}
Hitung matriks pemetaan
\maabbdisp {\bigwedge^2 T_P (\varphi)} { \bigwedge^2 T_P \R^3 } { \bigwedge^2 T_{\varphi(P)}  \R^2
} {}
di titik

\mavergleichskette
{\vergleichskette
{ P
}
{ = }{ (1,3,5)
}
{  }{
}
{    }{
}
{   }{
}
}
{}{}{}
terhadap suatu basis yang sesuai.

}
{

Secara umum, matriks Jacobi dari $\varphi$ adalah

\mathdisp {\begin{pmatrix}  yz^2 &  xz^2 &  2xyz \\  y  &  x  &  -3z^2 \end{pmatrix}} { . }

Untuk titik

\mavergleichskette
{\vergleichskette
{ P
}
{ = }{ (1,3,5)
}
{  }{
}
{    }{
}
{   }{
}
}
{}{}{}
matriks Jacobinya adalah

\mathdisp {\begin{pmatrix}  75 &  25 &  30 \\  3 &  1 &  -75  \end{pmatrix}} {  }

Matriks ini merepresentasikan pemetaan linear
\maabbdisp {L} {T_P\R^3 \cong \R^3} { T_{\varphi(P)} \R^2 \cong \R^2
} {}
terhadap basis-basis standar. Kita tentukan representasi matriks untuk pemetaan
\maabbdisp {\bigwedge^2 L} {\bigwedge^2 \R^3 } { \bigwedge^2 \R^2
} {}
terhadap basis
$e_1 \wedge e_2,\, e_1 \wedge e_3$ , $e_2 \wedge e_3$
\zusatzklammer {di ruas kiri} {} {}
dan
\mathl{e_1 \wedge e_2}{}
\zusatzklammer {di ruas kanan} {} {.}
Untuk itu, kita harus menghitung citra hasil kali-hasil kali baji tersebut. Kita peroleh

\mavergleichskettealign
{\vergleichskettealign
{ (\bigwedge^2 L ) (e_1 \wedge e_2)
}
{ =} { L(e_1) \wedge L(e_2)
}
{ =} { (75e_1 + 3e_2) \wedge (  25 e_1 + 1 e_2 )
}
{ =} { 75 e_1 \wedge e_2 + 75 e_2 \wedge  e_1
}
{ =} { 75 e_1 \wedge e_2 - 75 e_1 \wedge  e_2
}
}
{
\vergleichskettefortsetzungalign
{ =} { 0
}
{ } {}
{ } {}
{ } {}
}
{}{,}

\mavergleichskettealign
{\vergleichskettealign
{ (\bigwedge^2 L ) (e_1 \wedge e_3)
}
{ =} { L(e_1) \wedge L(e_3)
}
{ =} { (75e_1 + 3e_2) \wedge (  30 e_1 - 75 e_2 )
}
{ =} { - 75^2 e_1 \wedge e_2 + 90 e_2 \wedge  e_1
}
{ =} { (-5625 - 90) e_1 \wedge e_2
}
}
{
\vergleichskettefortsetzungalign
{ =} { -5715 e_1 \wedge e_2
}
{ } {}
{ } {}
{ } {}
}
{}{,}
dan

\mavergleichskettealign
{\vergleichskettealign
{ (\bigwedge^2 L ) (e_2 \wedge e_3)
}
{ =} { L(e_2) \wedge L(e_3)
}
{ =} { (  25 e_1 + 1 e_2 )   \wedge ( 30 e_1 - 75 e_2)
}
{ =} { - 1875 e_1 \wedge e_2 + 30 e_2 \wedge  e_1
}
{ =} { -  1905 e_1 \wedge  e_2
}
}
{}
{}{.}
Jadi, matriks yang merepresentasikannya adalah

\mathdisp {\left(  0,-5715,-1905                     \right)} { . }

 }



\inputaufgabeklausurloesung
{3}

{

Misalkan

\mavergleichskette
{\vergleichskette
{ W
}
{ \subseteq }{ \R^m
}
{  }{
}
{    }{
}
{   }{
}
}
{}{}{}
dan

\mavergleichskette
{\vergleichskette
{ U
}
{ \subseteq }{ \R^n
}
{  }{
}
{    }{
}
{   }{
}
}
{}{}{}
merupakan
\definitionsverweis {subhimpunan-subhimpunan terbuka}{}{}
dan misalkan
\maabbdisp {\psi} { W } { U
} {}
suatu
\definitionsverweis {pemetaan}{}{}
yang
\definitionsverweis {terdiferensialkan secara kontinu}{}{.}
Misalkan
\maabbdisp {f} { U } {\R
} {}
suatu fungsi terdiferensialkan secara kontinu. Simpulkan dari
aturan rantai
bahwa

\mavergleichskettedisp
{\vergleichskette
{ d (\psi^*f)
}
{ =} { \psi^*(df)
}
{ } {
}
{ } {
}
{ } {
}
}
{}{}{,}
dengan $\psi^*$ menyatakan
\definitionsverweis {penarikan balik bentuk diferensial}{}{.}

}
{

Misalkan

\mavergleichskette
{\vergleichskette
{ P
}
{ \in }{  W
}
{  }{
}
{    }{
}
{   }{
}
}
{}{}{}
dan

\mavergleichskette
{\vergleichskette
{ v
}
{ \in }{  \R^m
}
{  }{
}
{    }{
}
{   }{
}
}
{}{}{.}
Di satu pihak,

\mavergleichskettedisp
{\vergleichskette
{ d( \psi^* f) (P,v)
}
{ =} { D_P( f \circ \psi) (v)
}
{ =} { \left(  \left(  D_{\psi(P)} f  \right)  \circ \left(  D_P  \psi  \right)  \right) (v)
}
{ =} { D_{\psi(P)} (f) \left(  D_P \psi(v)  \right)
}
{ } {
}
}
{}{}{.}
Di pihak lain, juga berlaku

\mavergleichskettedisp
{\vergleichskette
{ \left(  \psi^*(df)  \right) (P,v)
}
{ =} { (df) \left(  \psi(P), (D_P \psi )(v)  \right)
}
{ =} { D_{\psi(P)} (f) \left(  D_P \psi (v)  \right)
}
{ } {
}
{ } {
}
}
{}{}{.}

 }



\inputaufgabeklausurloesung
{7}

{

Misalkan $M$ suatu
\definitionsverweis {manifold diferensiabel}{}{}
dengan
\definitionsverweis {basis terhitung bagi topologinya}{}{}
dan misalkan $\omega$ suatu
\definitionsverweis {bentuk volume positif}{}{}
pada $M$. Misalkan
\maabbdisp {\varphi} { L } { M
} {}
suatu
\definitionsverweis {difeomorfisme}{}{}
dengan manifold $L$ dan

\mavergleichskette
{\vergleichskette
{ S
}
{ \subseteq }{ L
}
{  }{
}
{    }{
}
{   }{
}
}
{}{}{}
suatu
\definitionsverweis {subhimpunan terukur}{}{.}
Tunjukkan bahwa

\mavergleichskettedisp
{\vergleichskette
{ \int_S \varphi^* \omega
}
{ =} { \int_{ \varphi(S)} \omega
}
{ } {
}
{ } {
}
{ } {
}
}
{}{}{.}

}
{

Menurut
Lema 15.2 (Geometri diferensial (Osnabrück 2023)),
kita dapat mengasumsikan bahwa
\mathkor {} {S} {dan} {\varphi(S)} {}
masing-masing seluruhnya terletak dalam satu peta dari
\mathkor {} {L} {dan} {M} {},
katakanlah

\mavergleichskette
{\vergleichskette
{ S
}
{ \subseteq }{ U
}
{  }{
}
{    }{
}
{   }{
}
}
{}{}{}
dengan peta
\maabbdisp {\alpha} { U } { U' \subseteq \R^n
} {}
dan

\mavergleichskette
{\vergleichskette
{ \varphi(S)
}
{ \subseteq }{ V
}
{  }{
}
{    }{
}
{   }{
}
}
{}{}{}
dengan peta
\maabbdisp {\beta} { V } { V' \subseteq \R^n
} {.}
Dengan memperkecil peta-peta itu lebih lanjut, kita dapat mengasumsikan bahwa
\mathkor {} {U} {dan} {V} {}
difeomorfik satu sama lain melalui $\varphi$. Dengan demikian terdapat diagram komutatif difeomorfisme

\mathdisp {\begin{matrix}  U  & \stackrel{ \varphi }{\longrightarrow} &  V  &  \\ \!\!\!\!\! \alpha \downarrow & & \downarrow \beta \!\!\!\!\!  & \\  U'  & \stackrel{ \beta \circ \varphi \circ \alpha^{-1} }{\longrightarrow} &  V' & \! \! \! \!\!\!\!\!   \\ \end{matrix}} {  }

yang menginduksi diagram komutatif

\mathdisp {\begin{matrix}  S  & \stackrel{ \varphi }{\longrightarrow} &  \varphi (S)  &  \\ \!\!\!\!\! \alpha \downarrow & & \downarrow \beta \!\!\!\!\!  & \\  \alpha (S) & \stackrel{ \beta \circ \varphi \circ \alpha^{-1} }{\longrightarrow} &  \beta ( \varphi (S) ) & \! \! \! \!\!\!\!\!   \\ \end{matrix}} {  }

atas himpunan-himpunan terukur. Untuk bentuk volume $\omega$ pada $V$ berlaku

\mavergleichskettedisp
{\vergleichskette
{ \alpha^{ {-1}*} ( \varphi^* \omega)
}
{ =} { \left(  \beta \circ \varphi \circ \alpha^{-1}  \right)^*  ( \beta^{ {-1}*} \omega)
}
{ } {
}
{ } {
}
{ } {
}
}
{}{}{.}
menurut
Lema 14.7 (Geometri diferensial (Osnabrück 2023))  (4).
Bentuk volume $\beta^{ {-1}*} \omega$ pada $V'$ memiliki representasi
\mathl{g dy_1  \wedge \ldots \wedge dy_n}{} dengan suatu fungsi terukur
\maabb {g} {V'} {\R
} {},
dan menurut definisi, $\int_{\varphi(S)} \omega$ sama dengan $\int_{\beta( \varphi(S))} g d \lambda^n$. Demikian pula,
\mathl{\alpha^{ {-1}*} ( \varphi^* \omega)}{} pada $U'$ memiliki representasi berbentuk
\mathl{f dx_1 \wedge \ldots \wedge dx_n}{.} Bentuk volume pada $U'$ ini sama dengan tarikan balik dari
\mathl{g dy_1  \wedge \ldots \wedge dy_n}{}, dan menurut
Korolari 14.9 (Geometri diferensial (Osnabrück 2023)),
tarikan balik tersebut sama dengan

\mathdisp {(g \circ \psi) \cdot \det  \left(  \left(  { \frac{   \partial   \psi_i   }{  \partial   x_j   } }  \right)_{1 \leq  i, j \leq n}  \right) dx_1 \wedge \ldots \wedge dx_n} {  }

\zusatzklammer {dengan

\mavergleichskettek
{\vergleichskettek
{ \psi
}
{ = }{ \beta \circ \varphi \circ \alpha^{-1}
}
{  }{
}
{    }{
}
{   }{
}
}
{}{}{}} {} {.}
Dengan demikian, pernyataan tersebut mengikuti dari
Ujian dengan solusi (Teori ukuran dan integrasi (Osnabrück 2022-2023)) (Teori ukuran dan integrasi (Osnabrück 2022-2023))
(Geometri diferensial (Osnabrück 2023)).

 }



\inputaufgabeklausurloesung
{2}

{

Pada
\definitionsverweis {bundel vektor trivial}{}{}
\maabbdisp {p} {\R^2 \times \R} { \R^2
} {}
ber-
\definitionsverweis {rank}{}{}
$1$ di atas $\R^2$, kita tinjau
\definitionsverweis {koneksi linear}{}{,}
yang diberikan oleh
\definitionsverweis {simbol-simbol Christoffel}{}{}

\mavergleichskette
{\vergleichskette
{ \Gamma_1 (x,y)
}
{ = }{ x^5-x^2y^2+3y^4
}
{  }{
}
{    }{
}
{   }{
}
}
{}{}{}
dan

\mavergleichskette
{\vergleichskette
{ \Gamma_2 (x,y)
}
{ = }{ 7x^3 +xy^2-4y^5
}
{  }{
}
{    }{
}
{   }{
}
}
{}{}{.}
Hitung
\definitionsverweis {operator kelengkungan}{}{}
$R \left(  \partial_1, \partial_2  \right)$.

}
{

Menurut
Fakta *****,
berlaku

\mavergleichskettealign
{\vergleichskettealign
{ R \left(  \partial_1, \partial_2  \right)
}
{ =} { \partial_1 \Gamma_2 - \partial_2 \Gamma_1
}
{ =} { \partial_1 \left(  7x^3 +xy^2-4y^5  \right) -  \partial_2 \left(   x^5-x^2y^2+3y^4   \right)
}
{ =} {  21x^2 +y^2 - 5x^4 -2xy^2
}
{ } {
}
}
{}
{}{.}

 }



\inputaufgabeklausurloesung
{0}

{

}
{  }



\inputaufgabeklausurloesung
{6}

{

Buktikan teorema tentang retraksi ke batas pada manifold.

}
{

Menurut
Teorema 21.8 (Geometri diferensial (Osnabrück 2023)),
batas
\mathl{\partial M}{} merupakan
\definitionsverweis {manifold diferensiabel}{}{}
\definitionsverweis {berorientasi}{}{}
\zusatzklammer {tanpa batas} {} {.}
Oleh karena itu, menurut
Teorema 22.11 (Geometri diferensial (Osnabrück 2023)),
terdapat suatu
\definitionsverweis {bentuk volume positif}{}{}
\definitionsverweis {yang terdiferensialkan secara kontinu}{}{}
$\tau$ pada
\mathl{\partial M}{.} Kita mempunyai
\mathl{\int_{  \partial M  }  \tau >0}{.}
\definitionsverweis {Turunan eksterior}{}{}
dari bentuk volume $\tau$ adalah $0$. Andaikan terdapat suatu
\definitionsverweis {pemetaan yang terdiferensialkan secara kontinu}{}{}
\maabbdisp {\varphi} { M } { \partial M
} {}
dengan
\mathl{\varphi \vert_{\partial M} = \operatorname{Id}_{\partial M}}{}. Maka
\definitionsverweis {bentuk tarikan balik}{}{}

\mathl{\varphi^* \tau}{} adalah
$(n-1)$-\definitionsverweis {bentuk diferensial}{}{}
pada $M$ yang pembatasannya pada batas berimpit dengan $\tau$. Dengan menggunakan
Teorema 23.2 (Geometri diferensial (Osnabrück 2023))
dan
Teorema 20.4 (Geometri diferensial (Osnabrück 2023))  (5),
kita memperoleh

\mavergleichskettealign
{\vergleichskettealign
{
\int_{  \partial M  }  \tau
}
{ =} {
\int_{  \partial  M  }  \varphi^* \tau
}
{ =} {
\int_{   M  }  d(\varphi^* \tau)
}
{ =} {
\int_{   M  }  \varphi^* (d\tau)
}
{ =} {
\int_{   M  }  \varphi^* (0)
}
}
{
\vergleichskettefortsetzungalign
{ =} { 0
}
{ } {}
{ } {}
{ } {}
}
{}{.}
Ini merupakan suatu kontradiksi.

 }
\endgroup

\chapter{Formulir Solusi Resmi Ujian 4}
\noindent\textit{Solusi yang disediakan oleh sumber resmi; kekosongan sumber dipertahankan.}
\begingroup
\renewcommand{\theHfakt}{o011.exam.official.04.\arabic{fakt}}
%Data institusi

%\input{Dozentdaten}

%\renewcommand{\fachbereich}{Fachbereich}

%\renewcommand{\dozent}{Prof. Dr. . }

%Data ujian

\renewcommand{\klausurgebiet}{ }

\renewcommand{\klausurtyp}{ }

\renewcommand{\klausurdatum}{ . 20}

\klausurvorspann {\fachbereich} {\klausurdatum} {\dozent} {\klausurgebiet} {\klausurtyp}

%Data untuk tabel nilai berikut

\renewcommand{\aeins}{  3  }

\renewcommand{\azwei}{  3   }

\renewcommand{\adrei}{  0  }

\renewcommand{\avier}{  0  }

\renewcommand{\afuenf}{  4  }

\renewcommand{\asechs}{  5  }

\renewcommand{\asieben}{  4  }

\renewcommand{\aacht}{  0  }

\renewcommand{\aneun}{  4  }

\renewcommand{\azehn}{  6  }

\renewcommand{\aelf}{  2  }

\renewcommand{\azwoelf}{  6  }

\renewcommand{\adreizehn}{  4   }

\renewcommand{\avierzehn}{  10  }

\renewcommand{\afuenfzehn}{  10  }

\renewcommand{\asechzehn}{  0  }

\renewcommand{\asiebzehn}{  61  }

\renewcommand{\aachtzehn}{    }

\renewcommand{\aneunzehn}{    }

\renewcommand{\azwanzig}{    }

\renewcommand{\aeinundzwanzig}{    }

\renewcommand{\azweiundzwanzig}{    }

\renewcommand{\adreiundzwanzig}{    }

\renewcommand{\avierundzwanzig}{    }

\renewcommand{\afuenfundzwanzig}{    }

\renewcommand{\asechsundzwanzig}{    }

\punktetabellesechzehn

\klausurnote

\newpage

\setcounter{section}{K}



\inputaufgabeklausurloesung
{3}

{

Definisikan istilah-istilah berikut
\zusatzklammer {yang dicetak miring} {} {}.
\aufzaehlungsechs{Pemetaan
\stichwort {Weingarten} {}
pada $T_PY$ di suatu titik

\mavergleichskette
{\vergleichskette
{ P
}
{ \in }{ Y
}
{  }{
}
{    }{
}
{   }{
}
}
{}{}{}
dari suatu hiperpermukaan diferensiabel berorientasi

\mavergleichskette
{\vergleichskette
{ Y
}
{ \subseteq }{ \R^n
}
{  }{
}
{    }{
}
{   }{
}
}
{}{}{.}

}{Suatu
\stichwort {bundel vektor riil} {}
$V$ ber-rank $r$ di atas suatu
\definitionsverweis {ruang topologis}{}{}
$X$.

}{\stichwort {Integral lintasan} {} dari suatu bentuk diferensial-$1$
\mathl{\omega \in
{ \mathcal E }^{  1 } (  M   )}{} pada suatu manifold diferensiabel $M$ terhadap suatu kurva yang terdiferensialkan secara kontinu
\maabb {\gamma} {[a,b]} {M
} {.}

}{Suatu
\stichwort {bentuk volume positif} {}
$\omega$ pada suatu manifold diferensiabel $M$.

}{Suatu \stichwort {bentuk diferensial tertutup} {} $\omega$ pada suatu manifold diferensiabel $M$.

}{\stichwort {Percepatan tangensial} {}
dari suatu kurva yang dua kali terdiferensialkan secara kontinu
\maabb {\gamma} {I} {M
} {}
pada suatu
\definitionsverweis {manifold Riemann}{}{}
$M$.
}

}
{

\aufzaehlungsechs{Misalkan $N$ suatu
\definitionsverweis {medan normal satuan}{}{}
pada $Y$ yang menentukan orientasi. Maka pemetaan
\maabbeledisp {L_P} {T_P Y} { T_PY
} {v} {- \left(  D_{v} N  \right) \left(  P  \right)
} {,}
disebut pemetaan Weingarten
pada $T_PY$.
}{Suatu
bundel vektor riil
$V$ ber-rank $r$ adalah ruang topologis yang dilengkapi dengan suatu
\definitionsverweis {pemetaan kontinu}{}{}
\maabb {p} {V} {X
} {}
sedemikian rupa sehingga setiap
\definitionsverweis {serat}{}{}
$p^{-1}(x)$ merupakan
\definitionsverweis {ruang vektor riil}{}{}
\definitionsverweis {berdimensi}{}{}
$r$, dan terdapat suatu
\definitionsverweis {penutup terbuka}{}{}

\mavergleichskette
{\vergleichskette
{X
}
{ = }{ \bigcup_{i \in I} U_i
}
{  }{
}
{    }{
}
{   }{
}
}
{}{}{}
serta
\definitionsverweis {homeomorfisme}{}{}
\maabbdisp {\varphi_i} {p^{-1}  \left(  U_i  \right) } { U_i \times \R^r
} {}
di atas $U_i$ yang pada setiap serat menginduksi suatu
\definitionsverweis {isomorfisme linear}{}{}
\maabbdisp {\left(  \varphi_i  \right)_x} {p^{-1} (x) } { \R^r
} {}
.
}{Integral lintasan didefinisikan oleh

\mavergleichskettedisp
{\vergleichskette
{  \int_\gamma \omega
}
{ =} {\int_a^b \gamma^* \omega
}
{ } {
}
{ } {
}
{ } {
}
}
{}{}{}
.
}{Suatu
$n$-\definitionsverweis {bentuk diferensial}{}{}
$\omega$ pada $M$ disebut bentuk volume positif jika, untuk setiap
\definitionsverweis {peta}{}{}
\maabbdisp {\alpha} { U } { V
} {}
\zusatzklammer {dengan \mathlk{V \subseteq \R^n}{} dan fungsi-fungsi koordinat \mathlk{x_1 , \ldots , x_n}{}} {} {}
dalam representasi lokal bentuk diferensial

\mavergleichskettedisp
{\vergleichskette
{ \alpha_* \omega
}
{ =} { f dx_1 \wedge \ldots \wedge dx_n
}
{ } {
}
{ } {
}
{ } {
}
}
{}{}{}
fungsi $f$ positif di setiap titik.
}{Suatu
\definitionsverweis {bentuk diferensial}{}{}
\definitionsverweis {diferensiabel}{}{}
$\omega$ pada $M$ disebut tertutup jika
\definitionsverweis {turunan luar}{}{}

\mathl{d \omega=0}{}.
}{Komposisi

\mathl{\pi_{\text{vert} }   \circ TT(\gamma)}{} disebut
percepatan tangensial
dari $\gamma$, dengan proyeksi tersebut berasal dari
\definitionsverweis {koneksi Levi-Civita}{}{.}
}

 }



\inputaufgabeklausurloesung
{3}

{

Rumuskan teorema-teorema berikut.
\aufzaehlungdrei{\stichwort {Teorema tentang kelengkungan normal} {.}}{Rumus untuk menghitung volume kanonik pada suatu manifold Riemann di dalam suatu peta.}{\stichwort {Teorema Stokes} {} untuk manifold dengan batas.}

}
{

\aufzaehlungdrei{\faktsituation {Misalkan

\mavergleichskette
{\vergleichskette
{ W
}
{ \subseteq }{ \R^n
}
{  }{
}
{    }{
}
{   }{
}
}
{}{}{}
\definitionsverweis {terbuka}{}{,}
\maabb {h} { W } { \R
} {}
suatu
\definitionsverweis {fungsi yang dua kali terdiferensialkan secara kontinu}{}{}
dan

\mavergleichskette
{\vergleichskette
{ Y
}
{ = }{ h^{-1}(c)
}
{  }{
}
{    }{
}
{   }{
}
}
{}{}{}
merupakan
\definitionsverweis {serat}{}{}
dari

\mavergleichskette
{\vergleichskette
{ c
}
{ \in }{ \R
}
{  }{
}
{    }{
}
{   }{
}
}
{}{}{,}
dengan $h$
\definitionsverweis {reguler}{}{}
di setiap titik $Y$. Misalkan

\mavergleichskette
{\vergleichskette
{ P
}
{ \in }{ Y
}
{  }{
}
{    }{
}
{   }{
}
}
{}{}{.}
Misalkan

\mavergleichskette
{\vergleichskette
{ E
}
{ \subseteq }{ \R^n
}
{  }{
}
{    }{
}
{   }{
}
}
{}{}{}
suatu
\definitionsverweis {bidang normal}{}{}
yang melalui $P$ pada $Y$.}
\faktfolgerung {Maka
\definitionsverweis {kelengkungan normal}{}{}
dari $Y$ di $P$ sama dengan
\definitionsverweis {kelengkungan}{}{}
kurva bidang

\mavergleichskette
{\vergleichskette
{ Y \cap (P+E)
}
{ \subseteq }{ P+E
}
{ \cong }{ \R^2
}
{    }{
}
{   }{
}
}
{}{}{}
di titik $P$.}
\faktzusatz {}
\faktzusatz {}}{Misalkan $M$ suatu
\definitionsverweis {manifold Riemann}{}{}
\definitionsverweis {berorientasi}{}{}
dan $\omega$
\definitionsverweis {bentuk volume kanonik}{}{.}
Misalkan
\maabbdisp {\alpha} { U } { V
} {}
suatu
\definitionsverweis {peta berorientasi}{}{}
dengan

\mavergleichskettedisp
{\vergleichskette
{ V
}
{ \subseteq} { \R^n
}
{ } {
}
{ } {
}
{ } {
}
}
{}{}{}
terbuka, dengan koordinat
\mathl{x_1 , \ldots , x_n}{} dan
\definitionsverweis {matriks fundamental metrik}{}{}

\mathl{G=(g_{ij})_{1 \leq i,j \leq n}}{}, serta
\mathl{g= \det  G}{.} Maka

\mavergleichskettedisp
{\vergleichskette
{  \alpha_* \omega
}
{ =} { \sqrt{ g} dx_1 \wedge \ldots \wedge dx_n
}
{ } {
}
{ } {
}
{ } {
}
}
{}{}{.}
Untuk suatu
\definitionsverweis {himpunan bagian terukur}{}{}

\mavergleichskette
{\vergleichskette
{ T
}
{ \subseteq }{ U
}
{  }{
}
{    }{
}
{   }{
}
}
{}{}{}
berlaku

\mavergleichskettedisp
{\vergleichskette
{
\int_{  T   }  \omega
}
{ =} {
\int_{ \alpha(T)  }   \sqrt{g} dx_1 \wedge \ldots \wedge dx_n
}
{ =} { \int_{ \alpha(T)  }   \sqrt{g}    \,  d \lambda^n
}
{ } {
}
{ } {
}
}
{}{}{.}}{Misalkan $M$ suatu
\definitionsverweis {manifold diferensiabel}{}{}
berdimensi $n$ yang
\definitionsverweis {berorientasi}{}{}
dan
\definitionsverweis {mempunyai batas}{}{}
$\partial M$, serta mempunyai
\definitionsverweis {basis topologi yang terhitung}{}{.}
Misalkan $\omega$ suatu
$(n-1)$-\definitionsverweis {bentuk diferensial}{}{}
yang
\definitionsverweis {terdiferensialkan secara kontinu}{}{}
dengan
\definitionsverweis {dukungan}{}{}
\definitionsverweis {kompak}{}{} pada $M$. Maka

\mavergleichskettedisp
{\vergleichskette
{
\int_{ M  }  d\omega
}
{ =} {
\int_{ \partial M  }  \omega
}
{ } {
}
{ } {
}
{ } {
}
}
{}{}{.}}

 }



\inputaufgabeklausurloesung
{0}

{

}
{  }



\inputaufgabeklausurloesung
{0}

{

}
{  }



\inputaufgabeklausurloesung
{4}

{

Buktikan teorema tentang hubungan antara kelengkungan dan vektor normal satuan dari suatu kurva bidang yang diparameterkan oleh panjang busur.

}
{

Pernyataan pertama langsung mengikuti definisi

\mavergleichskettedisp
{\vergleichskette
{ \kappa(t)
}
{ =} { \gamma_1'(t) \gamma_2^{\prime \prime} (t) - \gamma_2'(t) \gamma_1^{\prime \prime} (t)
}
{ } {
}
{ } {
}
{ } {
}
}
{}{}{.}
Dari

\mavergleichskettedisp
{\vergleichskette
{ \left(  \gamma_1'(t)  \right)^2 +  \left(  \gamma_2' (t)  \right)^2
}
{ =} { 1
}
{ } {
}
{ } {
}
{ } {
}
}
{}{}{}
diperoleh

\mavergleichskettedisp
{\vergleichskette
{ \gamma_1'(t) \gamma_1^{\prime \prime}(t) + \gamma_2'(t) \gamma_2^{\prime \prime} (t)
}
{ =} { 0
}
{ } {
}
{ } {
}
{ } {
}
}
{}{}{,}
karena itu $\gamma^{\prime \prime}(t)$ tegak lurus terhadap $\gamma'(t)$ dan bergantung linear pada
\mathl{N(t)}{.} Dengan demikian,

\mavergleichskettedisp
{\vergleichskette
{ \gamma^{\prime \prime}(t)
}
{ =} { \left\langle  \gamma^{\prime \prime}(t) ,  N(t)  \right\rangle  N(t)
}
{ =} { \kappa(t) N(t)
}
{ } {
}
{ } {
}
}
{}{}{}
dan oleh karena itu

\mavergleichskettedisp
{\vergleichskette
{ \Vert { \gamma^{\prime \prime}(t) } \Vert
}
{ =} { \betrag {  \kappa(t) }
}
{ } {
}
{ } {
}
{ } {
}
}
{}{}{.}

 }



\inputaufgabeklausurloesung
{5 (1+2+2)}

{

Pada
\definitionsverweis {permukaan bola satuan}{}{}
$Y$ berdimensi dua, kita meninjau
\definitionsverweis {medan vektor tangensial}{}{}

\mavergleichskettedisp
{\vergleichskette
{ F \begin{pmatrix}  x  \\ y \\ z   \end{pmatrix}
}
{ =} { \begin{pmatrix}  y  \\ -x\\  0   \end{pmatrix}
}
{ } {
}
{ } {
}
{ } {
}
}
{}{}{.}
\aufzaehlungdrei{Tentukan
\definitionsverweis {matriks Jacobi}{}{}
dari $F$ pada $\R^3$.
}{Untuk titik

\mavergleichskette
{\vergleichskette
{ P
}
{ = }{ \begin{pmatrix}  0  \\ 1 \\  0   \end{pmatrix}
}
{ \in }{ Y
}
{    }{
}
{   }{
}
}
{}{}{}
tentukan suatu matriks untuk pemetaan
\maabbeledisp {} {T_PY} { T_PY
} { v } { \nabla_v F
} {,}
terhadap suatu
\definitionsverweis {basis}{}{}
yang sesuai.
}{Untuk titik

\mavergleichskette
{\vergleichskette
{ Q
}
{ = }{ \begin{pmatrix}  0  \\ 0\\  1   \end{pmatrix}
}
{ \in }{ Y
}
{    }{
}
{   }{
}
}
{}{}{}
tentukan suatu matriks untuk pemetaan
\maabbeledisp {} {T_QY} { T_QY
} { v } { \nabla_v F
} {,}
terhadap suatu
\definitionsverweis {basis}{}{}
yang sesuai.
}

}
{

\aufzaehlungdrei{Matriks Jacobi-nya adalah

\mathdisp {\begin{pmatrix}  0   &   1  &  0  \\  -1 &  0  &  0 \\ 0  &  0 &  0  \end{pmatrix}} { . }

}{Sebagai medan normal, kita ambil $\begin{pmatrix}  x  \\ y \\ z   \end{pmatrix}$. Di

\mavergleichskette
{\vergleichskette
{ P
}
{ = }{ \begin{pmatrix}  0  \\ 1 \\  0   \end{pmatrix}
}
{  }{
}
{    }{
}
{   }{
}
}
{}{}{}
vektor-vektor
\mathkor {} {\begin{pmatrix}  1  \\ 0 \\  0   \end{pmatrix}} {dan} {\begin{pmatrix}  0  \\ 0\\  1   \end{pmatrix}} {}
membentuk suatu basis ruang singgung. Berlaku

\mavergleichskettedisp
{\vergleichskette
{ \begin{pmatrix}  0   &   1  &  0  \\  -1 &  0  &  0 \\ 0  &  0 &  0  \end{pmatrix} \begin{pmatrix}  1  \\ 0 \\  0   \end{pmatrix}
}
{ =} { \begin{pmatrix}  0  \\ -1\\  0   \end{pmatrix}
}
{ } {
}
{ } {
}
{ } {
}
}
{}{}{}
dan

\mavergleichskettedisp
{\vergleichskette
{  \begin{pmatrix}  0   &   1  &  0  \\  -1 &  0  &  0 \\ 0  &  0 &  0  \end{pmatrix} \begin{pmatrix}  0  \\ 0\\  1   \end{pmatrix}
}
{ =} { \begin{pmatrix}  0  \\ 0\\  0   \end{pmatrix}
}
{ } {
}
{ } {
}
{ } {
}
}
{}{}{.}
Di sini,
\mathl{\begin{pmatrix}  0  \\ -1\\  0   \end{pmatrix}}{} merupakan kelipatan vektor normal, sehingga proyeksi ortogonalnya pada ruang singgung juga merupakan vektor nol. Jadi, pemetaan
\mathl{v \mapsto \nabla_vF}{} adalah pemetaan nol.
}{Di

\mavergleichskette
{\vergleichskette
{ Q
}
{ = }{ \begin{pmatrix}  0  \\ 0\\  1   \end{pmatrix}
}
{  }{
}
{    }{
}
{   }{
}
}
{}{}{}
vektor-vektor
\mathkor {} {\begin{pmatrix}  1  \\ 0 \\  0   \end{pmatrix}} {dan} {\begin{pmatrix}  0  \\ 1 \\  0   \end{pmatrix}} {}
membentuk suatu basis ruang singgung. Berlaku

\mavergleichskettedisp
{\vergleichskette
{ \begin{pmatrix}  0   &   1  &  0  \\  -1 &  0  &  0 \\ 0  &  0 &  0  \end{pmatrix} \begin{pmatrix}  1  \\ 0 \\  0   \end{pmatrix}
}
{ =} { \begin{pmatrix}  0  \\ -1\\  0   \end{pmatrix}
}
{ } {
}
{ } {
}
{ } {
}
}
{}{}{}
dan

\mavergleichskettedisp
{\vergleichskette
{ \begin{pmatrix}  0   &   1  &  0  \\  -1 &  0  &  0 \\ 0  &  0 &  0  \end{pmatrix} \begin{pmatrix}  0  \\ 1 \\  0   \end{pmatrix}
}
{ =} { \begin{pmatrix}  1  \\ 0 \\  0   \end{pmatrix}
}
{ } {
}
{ } {
}
{ } {
}
}
{}{}{.}
Kedua vektor hasil tersebut sudah menyinggung $Y$ di $Q$. Suatu matriks yang merepresentasikan
\mathl{v \mapsto \nabla_vF}{} adalah

\mathdisp {\begin{pmatrix}  0   &   -1  \\  1  &  0 \end{pmatrix}} { . }

}

 }



\inputaufgabeklausurloesung
{4 (2+2)}

{

Kita meninjau elips

\mavergleichskettedisp
{\vergleichskette
{ E
}
{ =} { \left\{ (u,v)  \mid   2u^2 +5v^2 = 4         \right\}
}
{ } {
}
{ } {
}
{ } {
}
}
{}{}{.}
\aufzaehlungdrei{Definisikan suatu pemetaan surjektif yang terdiferensialkan secara kontinu
\maabb {} {\R} {E
} {.}
}{Deskripsikan suatu
\definitionsverweis {difeomorfisme}{}{}
antara lingkaran $S^1$ dan $E$.
}{
}

}
{

Kita menggunakan deskripsi

\mavergleichskettedisp
{\vergleichskette
{ E
}
{ =} { \left\{ (u,v)  \mid   { \frac{   1 }{  2 } } u^2+ { \frac{   5 }{  4 } } v^2 = 1         \right\}
}
{ } {
}
{ } {
}
{ } {
}
}
{}{}{.}
\aufzaehlungzwei {Pemetaan
\maabbeledisp {} {\R} { E
} { \alpha} { \left(  \sqrt{2} \cos  \alpha   , \,   { \frac{   2 }{  \sqrt{5}  } } \sin  \alpha                   \right)
} {}
terdiferensialkan secara kontinu sebagai pemetaan ke $\R^2$, dan citranya terletak pada $E$. Oleh karena itu, pemetaan ini juga terdiferensialkan secara kontinu sebagai pemetaan ke $E$. Surjektivitasnya sekaligus dibuktikan dalam (2).
} {Kita meninjau pemetaan
\maabbeledisp {} {S^1} {E
} {(x,y)} { \left(  \sqrt{2} x  , \,   { \frac{   2 }{  \sqrt{5}  } } y                   \right)
} {.}
Sebagai pembatasan suatu pemetaan linear bijektif pada $\R^2$, pemetaan ini terdiferensialkan secara kontinu dan injektif, serta memang bernilai di $E$
\zusatzklammer {berdasarkan persamaan-persamaan eksplisit} {} {.}
Pemetaan ini juga surjektif karena kita dapat langsung memberikan pemetaan invers
\maabbeledisp {} { E } {S^1
} {(u,v)} { \left(  { \frac{   1 }{  \sqrt{2}  } } u  , \,   { \frac{   \sqrt{5}  }{  2  } } v                   \right)
} {,}
secara eksplisit. Karena parametrisasi trigonometrik lingkaran satuan bersifat surjektif, pemetaan yang diberikan dalam (1) juga surjektif.
}

 }



\inputaufgabeklausurloesung
{0}

{

}
{  }



\inputaufgabeklausurloesung
{4}

{

Berikan suatu
\definitionsverweis {medan vektor}{}{} yang
\definitionsverweis {kontinu}{}{}
pada $S^2$ dan hanya mempunyai satu titik nol.

}
{

Kita meninjau pada $\R^2$ medan vektor kontinu $F$ yang diberikan oleh

\mavergleichskettedisp
{\vergleichskette
{F(x,y)
}
{ =} { { \frac{   1 }{  x^2+y^2 } }  e_1
}
{ } {
}
{ } {
}
{ } {
}
}
{}{}{}
Medan ini tidak mempunyai titik nol dan bersifat kontinu. Kita mengangkut medan vektor ini melalui proyeksi stereografis ke

\mavergleichskette
{\vergleichskette
{  \R^2
}
{ \cong }{ S^2 \setminus \{N\}
}
{  }{
}
{    }{
}
{   }{
}
}
{}{}{}
dan memperluasnya di Kutub Utara dengan nilai

\mavergleichskette
{\vergleichskette
{ 0
}
{ \in }{  T_NS^2
}
{  }{
}
{    }{
}
{   }{
}
}
{}{}{.}
Kita menyatakan bahwa medan vektor ini kontinu. Untuk itu, misalkan
\mathl{P_n}{} suatu barisan pada $S^2$ yang konvergen ke $N$. Kita dapat langsung menganggap bahwa semuanya memenuhi

\mavergleichskette
{\vergleichskette
{P_n
}
{ \neq }{N
}
{  }{
}
{    }{
}
{   }{
}
}
{}{}{}
. Citra peta barisan ini adalah

\mavergleichskettedisp
{\vergleichskette
{P'_n
}
{ =} { (x_n,y_n)
}
{ } {
}
{ } {
}
{ } {
}
}
{}{}{,}
dan karena $P_n$ konvergen ke Kutub Utara,
\mathl{\betrag { P_n' }}{} divergen ke $\infty$. Oleh karena itu, barisan

\mavergleichskettedisp
{\vergleichskette
{ F(P'_n)
}
{ =} { F(x_n , y_n)
}
{ =} { { \frac{   1 }{  x^2+y^2 } } e_1
}
{ } {
}
{ } {
}
}
{}{}{}
konvergen ke $0$.

 }



\inputaufgabeklausurloesung
{6}

{

Misalkan diberikan suatu
\definitionsverweis {data perekatan}{}{}

\mathbed {U_i} {}
 {i \in I} {}
 {} {} {} {,}
untuk
\definitionsverweis {ruang-ruang topologis}{}{.}
Buktikan bahwa terdapat suatu ruang topologis $X$ yang ditentukan secara unik, suatu penutup terbuka

\mavergleichskette
{\vergleichskette
{ X
}
{ = }{ \bigcup_{i \in I} V_i
}
{  }{
}
{    }{
}
{   }{
}
}
{}{}{}
dan
\definitionsverweis {homeomorfisme}{}{}
\maabb {\psi_i} {U_i } { V_i
} {}
sedemikian sehingga

\mavergleichskettedisp
{\vergleichskette
{ \psi_i  \left(  U_{ij}  \right)
}
{ =} { V_i \cap V_j
}
{ } {
}
{ } {
}
{ } {
}
}
{}{}{}
dan

\mavergleichskettedisp
{\vergleichskette
{ \psi_i \vert_{U_{ij} }
}
{ =} { \psi_j \vert_{U_{ji} }  \circ   \varphi_{ji}
}
{ } {
}
{ } {
}
{ } {
}
}
{}{}{}.

}
{

Misalkan $Y$ gabungan saling lepas dari himpunan-himpunan $U_i$. Pada $Y$, kita mendefinisikan suatu
\definitionsverweis {relasi ekuivalensi}{}{}
$\sim$ dengan menganggap titik-titik

\mavergleichskette
{\vergleichskette
{ x_i
}
{ \in }{ U_i
}
{  }{
}
{    }{
}
{   }{
}
}
{}{}{}
dan

\mavergleichskette
{\vergleichskette
{ x_j
}
{ \in }{ U_j
}
{  }{
}
{    }{
}
{   }{
}
}
{}{}{}
ekuivalen jika

\mavergleichskette
{\vergleichskette
{ x_i
}
{ \in }{ U_{ij}
}
{  }{
}
{    }{
}
{   }{
}
}
{}{}{,}

\mavergleichskette
{\vergleichskette
{ x_j
}
{ \in }{ U_{ji}
}
{  }{
}
{    }{
}
{   }{
}
}
{}{}{}
dan

\mavergleichskette
{\vergleichskette
{ \varphi_{ji} (x_i)
}
{ = }{ x_j
}
{  }{
}
{    }{
}
{   }{
}
}
{}{}{}
berlaku. Sifat-sifat relasi ekuivalensi dijamin oleh syarat kosikel; lihat
:Kurs:Bündel, Garben und Kohomologie (Osnabrück 2019-2020)/4/Klausur mit Lösungen/latex (Bündel, Garben und Kohomologie (Osnabrück 2019-2020))
(Differentialgeometrie (Osnabrück 2023)).
Kita menetapkan

\mavergleichskettedisp
{\vergleichskette
{X
}
{ \defeq} { Y/ \sim
}
{ } {
}
{ } {
}
{ } {
}
}
{}{}{}
dan memperlengkapi $X$ dengan
\definitionsverweis {topologi hasil bagi}{}{.}
Komposisi-komposisi

\mathdisp {U_i \longrightarrow Y \longrightarrow X} {  }

adalah $\psi_i$, sedangkan $V_i$ merupakan citra pemetaan-pemetaan tersebut. Dengan demikian, diperoleh homeomorfisme
\maabb {\psi_i} {U_i } { V_i
} {}
. Untuk

\mavergleichskette
{\vergleichskette
{ x
}
{ \in }{ U_i
}
{  }{
}
{    }{
}
{   }{
}
}
{}{}{}

\mavergleichskettedisp
{\vergleichskette
{ \psi_i (x)
}
{ \in} { V_j
}
{ } {
}
{ } {
}
{ } {
}
}
{}{}{}
berlaku tepat jika

\mavergleichskette
{\vergleichskette
{ x
}
{ \in }{ U_{ij}
}
{  }{
}
{    }{
}
{   }{
}
}
{}{}{}
karena tepat dalam hal ini $x$ diidentifikasi dengan
\mathl{\varphi_{ij}(x)}{}. Oleh karena itu,

\mavergleichskette
{\vergleichskette
{ \psi_i(U_{ij})
}
{ = }{ V_i \cap V_j
}
{  }{
}
{    }{
}
{   }{
}
}
{}{}{.}
Komutativitas diagram

\mathdisp {\begin{matrix} U_{ij}   & \stackrel{ \varphi_{ji} }{\longrightarrow} &  U_{ji}   & \\ & \!\!\! \!\! \psi_i \searrow & \downarrow \psi_j \!\!\! \!\! & \\ & &  V_i \cap V_j   & \!\!\!\!\! \!\!\!  \\ \end{matrix}} {  }

juga diperoleh.

 }



\inputaufgabeklausurloesung
{2}

{

Misalkan $M$ suatu
\definitionsverweis {manifold diferensiabel}{}{,}
\maabb {\gamma} { I } { M
} {}
suatu
\definitionsverweis {kurva diferensiabel}{}{}
di $M$, dan $\omega$ suatu
$1$-\definitionsverweis {bentuk}{}{}
terdiferensialkan pada $M$. Buktikan bahwa

\mavergleichskettedisp
{\vergleichskette
{ \gamma^* (d \omega)
}
{ =} { 0
}
{ } {
}
{ } {
}
{ } {
}
}
{}{}{}.

}
{

$d \omega$ merupakan suatu bentuk-$2$ pada $M$, dan sifat ini terbawa ke $\gamma^*(d \omega)$. Akan tetapi, pada manifold berdimensi satu, semua bentuk-$2$ sama dengan $0$.

 }



\inputaufgabeklausurloesung
{6 (3+3)}

{

Misalkan $M$ suatu
\definitionsverweis {manifold diferensiabel}{}{}
yang
\definitionsverweis {kompak}{}{}
dan mempunyai sedikitnya dua titik.
\aufzaehlungzwei {Buktikan bahwa suatu
$1$-\definitionsverweis {bentuk diferensial}{}{}
kontinu yang
\definitionsverweis {eksak}{}{}
pada $M$ mempunyai sedikitnya dua titik nol.
} {Buktikan bahwa hal ini tidak harus berlaku bagi suatu bentuk diferensial yang
\definitionsverweis {tertutup}{}{.}
}

}
{

\aufzaehlungzwei {Suatu bentuk diferensial eksak dapat dituliskan sebagai

\mavergleichskettedisp
{\vergleichskette
{ \omega
}
{ =} { df
}
{ } {
}
{ } {
}
{ } {
}
}
{}{}{}
dengan suatu fungsi diferensiabel
\maabbdisp {f} { M } {\R
} {}
. Jika $f$ konstan, maka

\mavergleichskette
{\vergleichskette
{ df
}
{ = }{ 0
}
{  }{
}
{    }{
}
{   }{
}
}
{}{}{}
dan pernyataannya jelas. Jadi, misalkan $f$ tidak konstan. Menurut
Fakt *****
fungsi kontinu tersebut mencapai maksimum global dan minimum global, katakanlah masing-masing di titik
\mathkor {} {P} {dan} {Q} {.}
Menurut
Aufgabe *****
maka

\mavergleichskettedisp
{\vergleichskette
{ (df)_P
}
{ =} { (df)_Q
}
{ =} { 0
}
{ } {
}
{ } {
}
}
{}{}{.}
} {Pada lingkaran berdimensi satu yang kompak $S^1$, kita meninjau bentuk diferensial

\mavergleichskettedisp
{\vergleichskette
{ \omega(x,y)
}
{ =} {  ydx - xdy
}
{ } {
}
{ } {
}
{ } {
}
}
{}{}{.}
Ini memang merupakan suatu bentuk diferensial pada lingkaran karena
\mathl{\left(  y   , \,   -x                   \right)}{} bersifat tangensial terhadap vektor posisi
\mathl{(x,y)}{}. Bentuk ini tertutup karena dimensinya $1$. Bentuk ini tidak mempunyai titik nol karena pada lingkaran,
\mathkor {} {x} {dan} {y} {}
tidak mungkin sekaligus bernilai $0$.
}

 }



\inputaufgabeklausurloesung
{4}

{

Berikan suatu contoh
\definitionsverweis {manifold dengan batas}{}{,}
yang
\definitionsverweis {terhubung}{}{}
dan batasnya terdiri atas
\zusatzklammer {gabungan saling lepas dari} {} {}
tak terhingga banyak salinan $\R^1$.

}
{

Fungsi
\maabbeledisp {h} {]-1,1[ } { \R
} { x } { { \frac{   1  }{  x^2-1 } }
} {,}
menuju tak hingga pada kedua ujungnya. Grafik $h$ difeomorf dengan $]-1,1[$, dan dengan demikian juga dengan $\R$. Sekarang, untuk

\mavergleichskette
{\vergleichskette
{ n
}
{ \in }{  \Z
}
{  }{
}
{    }{
}
{   }{
}
}
{}{}{}
, kita meninjau fungsi yang bersesuaian dengan domain yang digeser sebesar $4n$
\zusatzklammer {jadi dengan domain
\mathl{]4n-1,4n+1[}{}} {} {.}
Misalkan $G_n$ grafik yang bersesuaian dan $U_n$ epigraf terbuka
\zusatzklammer {!} {} {}
yang bersesuaian. Kita menetapkan

\mavergleichskettedisp
{\vergleichskette
{ M
}
{ =} { \R^2 \setminus \biguplus_{n \in \Z} U_n
}
{ } {
}
{ } {
}
{ } {
}
}
{}{}{.}
Himpunan ini merupakan suatu manifold terhubung dengan batas, dan batasnya adalah gabungan saling lepas dari himpunan-himpunan $G_n$.

 }



\inputaufgabeklausurloesung
{10}

{

Buktikan \stichwort {teorema tentang partisi kesatuan} {.}

}
{

Menurut
Lemma 22.9 (Differentialgeometrie (Osnabrück 2023))
kita dapat menganggap bahwa terdapat suatu penutup terbuka oleh daerah-daerah peta

\mathbed {U_j} {}
 {j \in J} {}
 {} {} {} {,}
\zusatzklammer {$J$ terhitung} {} {}
dengan
\maabbdisp {\alpha_j} {U_j } { V_j
} {}
dan lingkungan-lingkungan bola

\mavergleichskettedisp
{\vergleichskette
{ U \left(   0 , \delta_j   \right)
}
{ \subset} { B  \left(  0 , \epsilon_j  \right)
}
{ \subset} { V_j
}
{ } {
}
{ } {
}
}
{}{}{}
\zusatzklammer {dengan

\mavergleichskettek
{\vergleichskettek
{ \delta_j
}
{ < }{ \epsilon_j
}
{  }{
}
{    }{
}
{   }{
}
}
{}{}{}} {} {}
sedemikian rupa sehingga himpunan-himpunan
\mathl{\alpha_j^{-1} \left(  U \left(   0 , \delta_j   \right)  \right)}{} juga membentuk suatu penutup dari $M$, dan setiap titik

\mavergleichskette
{\vergleichskette
{ P
}
{ \in }{ M
}
{  }{
}
{    }{
}
{   }{
}
}
{}{}{}
hanya termuat dalam berhingga banyak himpunan $U_j$, dan khususnya hanya dalam berhingga banyak himpunan
\mathl{\alpha_j^{-1} \left(  U \left(   0 , \delta_j   \right)  \right)}{}. Pada $V_j$, kita meninjau fungsi $g_j$ yang diberikan oleh

\mavergleichskettedisp
{\vergleichskette
{ g_j (v)
}
{ =} { \begin{cases}  e^{ - { \frac{   1  }{  \left(  \delta_j^2- \Vert { v } \Vert^2  \right)^2 } } }  \text{ jika } \Vert { v } \Vert < \delta_j \, , \\ 0 \text{ selain itu} \, , \end{cases}
}
{ } {
}
{ } {
}
{ } {
}
}
{}{}{}
. Fungsi ini bernilai positif tepat pada
\mathl{U \left(   0 , \delta_j   \right)}{}, dan
\definitionsverweis {dukungannya}{}{}
adalah
\mathl{B  \left(  0 , \delta_j  \right)}{.} Dengan meninjau dua himpunan bagian terbuka
\zusatzklammer {yang menutupi $V_j$} {} {}
\mathkor {} {U \left(   0 , \epsilon_j   \right)} {dan} {V_j \setminus B  \left(  0 , \delta_j  \right)} {}
diperoleh bahwa $g_j$ diferensiabel tak hingga kali. Kita mendefinisikan suatu fungsi
\maabbdisp {\tilde{g}_j} { M } {\R
} {}
melalui

\mavergleichskettedisp
{\vergleichskette
{ \tilde{g}_j(x)
}
{ =} { \begin{cases}  g( \alpha_j(x)) \text{ jika }  x \in \alpha_j^{-1}( U \left(   0 , \delta_j   \right) ) \, , \\  0 \text{ selain itu} \, . \end{cases}
}
{ } {
}
{ } {
}
{ } {
}
}
{}{}{}
Fungsi ini
\definitionsverweis {terdiferensialkan secara kontinu}{}{}
pada $M$ karena \anfuehrung{pita}{}
\mathl{B  \left(  0 , \epsilon_j  \right) \setminus U \left(   0 , \delta_j   \right)}{} memungkinkan transisi mulus. Kita menetapkan

\mavergleichskettedisp
{\vergleichskette
{ \tilde{g}(x)
}
{ \defeq} { \sum_{j \in J} \tilde{g}_j(x)
}
{ } {
}
{ } {
}
{ } {
}
}
{}{}{,}
yang merupakan jumlah berhingga pada setiap titik karena
\definitionsverweis {dukungan}{}{}
dari $\tilde{g}_j$ terletak di dalam

\mavergleichskettedisp
{\vergleichskette
{ \alpha^{-1} \left(  B  \left(  0 , \delta_j  \right)  \right)
}
{ \subset} { U_j
}
{ } {
}
{ } {
}
{ } {
}
}
{}{}{}
. Fungsi ini
\definitionsverweis {terdiferensialkan secara kontinu}{}{}
pada $M$ dan positif di setiap titik karena
\mathl{\tilde{g}_j(x)}{} positif pada himpunan-himpunan penutup
\mathl{\alpha^{-1} \left(  U \left(   0 , \delta_j   \right)  \right)}{}. Maka fungsi-fungsi

\mavergleichskettedisp
{\vergleichskette
{ h_j
}
{ =} { { \frac{   \tilde{g}_j  }{ \tilde{g}  } }
}
{ } {
}
{ } {
}
{ } {
}
}
{}{}{}
membentuk partisi kesatuan yang dicari.

 }



\inputaufgabeklausurloesung
{10 (2+6+2)}

{

Kita meninjau
$2$-\definitionsverweis {bentuk diferensial}{}{}

\mavergleichskettedisp
{\vergleichskette
{ \omega
}
{ =} { x dy \wedge dz -y dx \wedge dz + z dx \wedge dy
}
{ } {
}
{ } {
}
{ } {
}
}
{}{}{}
pada bola satuan
\mathl{B  \left(  0 , 1  \right)}{.}
\aufzaehlungdreiabc{Buktikan bahwa $d \omega$ adalah tiga kali bentuk volume standar pada bola tersebut.
}{Buktikan bahwa
\mathl{\omega\vert_{S^2}}{} merupakan bentuk luas permukaan standar pada permukaan bola satuan.
}{Hitung luas permukaan bola dari
volume bola
dengan menggunakan
Teorema Stokes.
}

}
{

a) Berlaku

\mavergleichskettealign
{\vergleichskettealign
{d \omega
}
{ =} { d \left(  x dy \wedge dz  -y dx \wedge dz +  z dx \wedge dy \right)
}
{ =} { dx \wedge dy \wedge dz - dy \wedge dx \wedge dz + d z \wedge dx \wedge dy
}
{ =} { dx \wedge dy \wedge dz + dx \wedge dy \wedge dz + d x \wedge dy \wedge dz
}
{ =} { 3 d x \wedge dy \wedge dz
}
}
{}
{}{,}
di mana kita menggunakan fakta bahwa tanda berubah ketika faktor-faktor dalam produk baji dipertukarkan.

b) Untuk suatu titik
\mathl{P = \begin{pmatrix}  x  \\ y \\  z   \end{pmatrix}  \in S^2}{} dan dua vektor singgung ortonormal
\zusatzklammer {yang bersama dengan $P$ merepresentasikan orientasi} {} {}
\mathl{u= \begin{pmatrix} u_1  \\u_2 \\ u_3   \end{pmatrix}}{} serta
\mathl{v= \begin{pmatrix} v_1  \\v_2 \\ v_3   \end{pmatrix}}{}, menurut
Satz 58.4 (Lineare Algebra (Osnabrück 2024-2025))

\mavergleichskettealign
{\vergleichskettealign
{ \omega \vert_{S^2}  (u,v)
}
{ =} { \omega (u,v)
}
{ =} { z \det  \begin{pmatrix} u_1   &  v_1   \\ u_2  & v_2  \end{pmatrix} - y \det  \begin{pmatrix} u_1   &  v_1   \\ u_3  & v_3  \end{pmatrix} + x \det  \begin{pmatrix} u_2   &  v_2   \\ u_3  & v_3  \end{pmatrix}
}
{ =} { \det  \begin{pmatrix}  x   &  u_1  & v_1  \\  y  & u_2  & v_2  \\ z  & u_3  & v_3  \end{pmatrix}
}
{ } {
}
}
{}
{}{.}
Determinan suatu sistem ortonormal yang merepresentasikan orientasi adalah $1$. Jadi, bentuk diferensial yang dibatasi tersebut memenuhi sifat yang mengarakterisasi
\definitionsverweis {bentuk volume standar}{}{.}

c) Dengan menggunakan a), b), Teorema Stokes, dan volume bola, luas permukaan bola adalah

\mavergleichskettedisp
{\vergleichskette
{ \int_{S^2} \omega
}
{ =} { \int_{ B  \left(  0, 1 \right)  }  d \omega
}
{ =} { \int_{ B  \left(  0, 1 \right)  } 3 d x \wedge dy \wedge dz
}
{ =} { 3 \lambda^3 (   B  \left(  0, 1 \right)     )
}
{ =} { 4 \pi
}
}
{}{}{.}

 }



\inputaufgabeklausurloesung
{0}

{

}
{  }
\endgroup

\chapter{Formulir Solusi Resmi Ujian 5}
\noindent\textit{Solusi yang disediakan oleh sumber resmi; kekosongan sumber dipertahankan.}
\begingroup
\renewcommand{\theHfakt}{o011.exam.official.05.\arabic{fakt}}
%Data institusi

%\input{Dozentdaten}

%\renewcommand{\fachbereich}{Bidang studi}

%\renewcommand{\dozent}{Prof. Dr. . }

%Data ujian

\renewcommand{\klausurgebiet}{ }

\renewcommand{\klausurtyp}{ }

\renewcommand{\klausurdatum}{. 20}

\klausurvorspann {\fachbereich} {\klausurdatum} {\dozent} {\klausurgebiet} {\klausurtyp}

%Data untuk tabel poin berikut

\renewcommand{\aeins}{  3  }

\renewcommand{\azwei}{  3   }

\renewcommand{\adrei}{  5  }

\renewcommand{\avier}{  2  }

\renewcommand{\afuenf}{  4  }

\renewcommand{\asechs}{  8  }

\renewcommand{\asieben}{  4  }

\renewcommand{\aacht}{  5  }

\renewcommand{\aneun}{  3  }

\renewcommand{\azehn}{  2  }

\renewcommand{\aelf}{  6  }

\renewcommand{\azwoelf}{  7  }

\renewcommand{\adreizehn}{  4   }

\renewcommand{\avierzehn}{  8  }

\renewcommand{\afuenfzehn}{  64  }

\renewcommand{\asechzehn}{    }

\renewcommand{\asiebzehn}{    }

\renewcommand{\aachtzehn}{    }

\renewcommand{\aneunzehn}{    }

\renewcommand{\azwanzig}{    }

\renewcommand{\aeinundzwanzig}{    }

\renewcommand{\azweiundzwanzig}{    }

\renewcommand{\adreiundzwanzig}{    }

\renewcommand{\avierundzwanzig}{    }

\renewcommand{\afuenfundzwanzig}{    }

\renewcommand{\asechsundzwanzig}{    }

\punktetabellevierzehn

\klausurnote

\newpage

\setcounter{section}{0}



\inputaufgabeklausurloesung
{3}

{

Definisikan istilah-istilah berikut
\zusatzklammer {yang dicetak miring} {} {}.
\aufzaehlungsechs{\stichwort {turunan tangensial kedua} {}
\zusatzklammer {atau
\stichwort {percepatan tangensial} {}} {} {}
\zusatzklammer {sebagai pemetaan ke $\R^n$} {} {}
dari suatu
\definitionsverweis {kurva yang dua kali terdiferensialkan secara kontinu}{}{}
\maabbdisp {\gamma} {I} {Y
} {}
yang bernilai dalam suatu
\definitionsverweis {hipermuka terdiferensialkan}{}{}

\mavergleichskette
{\vergleichskette
{ Y
}
{ \subseteq }{ \R^n
}
{  }{
}
{    }{
}
{   }{
}
}
{}{}{.}

}{Suatu ruang topologis $X$ yang
\stichwort {terhubung} {}.

}{\stichwort {Pemetaan singgung di titik} {}

\mavergleichskette
{\vergleichskette
{  P
}
{ \in }{ M
}
{  }{}
{    }{}
{   }{}
}
{}{}{}
untuk suatu pemetaan diferensiabel
\maabbdisp {\varphi} { M } { N
} {,}
dengan
\mathkor {} {M} {dan} {N} {}
manifold diferensiabel.

}{\stichwort {Produk} {}
dari manifold diferensiabel
\mathkor {} {M} {dan} {N} {.}

}{\stichwort {Ukuran volume} {}
yang terkait dengan bentuk volume positif $\omega$ pada suatu manifold diferensiabel berdimensi $n$.

}{Suatu
\stichwort {kurva geodesik} {}
pada manifold Riemann $M$.
}

}
{

\aufzaehlungsechs{Pemetaan
\maabbeledisp {} {I} { TY
} {t} { \pi_P( \gamma^{\prime \prime} (t))
} {,}
dengan $\pi_P$ menyatakan proyeksi ortogonal
\maabbdisp {} {\R^n } { T_PY
} {}
disebut
turunan tangensial kedua
\zusatzklammer {atau
\definitionswort {percepatan tangensial}{}} {} {}
dari $\gamma$.
}{Suatu
\definitionsverweis {ruang topologis}{}{}
$X$ disebut terhubung jika tepat ada dua himpunan bagian di $X$
\zusatzklammer {yaitu $\emptyset$ dan seluruh ruang

\mavergleichskettek
{\vergleichskettek
{  X
}
{ \neq }{ \emptyset
}
{  }{
}
{    }{
}
{   }{
}
}
{}{}{}} {} {,}
yang sekaligus
\definitionsverweis {terbuka}{}{}
dan
\definitionsverweis {tertutup}{}{}.
}{Yang dimaksud dengan pemetaan singgung di titik $P$ adalah pemetaan
\maabbeledisp {} {T_PM} {T_{\varphi(P)}N
} {[\gamma]} {[\varphi \circ \gamma]
} {.}
}{Misalkan
\mathkor {} {(U_i,U_i',\alpha_i, i \in I)} {dan} {(V_j,V_j',\beta_j, j \in J)} {}
adalah
\definitionsverweis {atlas}{}{}
dari
\mathkor {} {M} {dan} {N} {.}
Maka
\definitionsverweis {ruang produk}{}{}

\mathl{M \times N}{} yang dilengkapi dengan
\definitionsverweis {peta}{}{}
\maabbdisp {\alpha_i \times \beta_j} {U_i \times V_j } { U_i' \times V_j'
} {}
\zusatzklammer {dengan \mathlk{(i,j) \in I \times J}{} dan \mathlk{U_i' \times V_j' \subseteq \R^m \times \R^n}{}} {} {}
disebut produk dari manifold
\mathkor {} {M} {dan} {N} {.}
}{Untuk suatu
\definitionsverweis {himpunan Borel}{}{}

\mavergleichskette
{\vergleichskette
{ T
}
{ \subseteq }{ M
}
{  }{
}
{    }{
}
{   }{
}
}
{}{}{}
ukuran $T$ terhadap $\omega$ didefinisikan melalui suatu
\definitionsverweis {dekomposisi}{}{}
\definitionsverweis {terhitung}{}{}

\mavergleichskette
{\vergleichskette
{ T
}
{ = }{  \biguplus_{i \in I} T_i
}
{  }{
}
{    }{
}
{   }{
}
}
{}{}{}
\zusatzklammer {dengan

\mavergleichskettek
{\vergleichskettek
{ T_i
}
{ \subseteq }{U_i
}
{  }{
}
{    }{
}
{   }{
}
}
{}{}{}
suatu domain peta terbuka dan

\mavergleichskettek
{\vergleichskettek
{  \alpha_{i*} \omega
}
{ = }{ f_i dx_1 \wedge \ldots \wedge dx_n
}
{  }{
}
{    }{
}
{   }{
}
}
{}{}{}
berlaku} {} {}

\mavergleichskettedisp
{\vergleichskette
{
\int_{ T  }  \omega
}
{ =} { \sum_{i \in I} \int_{ \alpha_i(T_i)  }  f_i   \,  d \lambda^n
}
{ } {
}
{ } {
}
{ } {
}
}
{}{}{}
sebagaimana di atas.
}{Suatu
\definitionsverweis {kurva diferensiabel}{}{}
\maabbdisp {\gamma} {I} {M
} {}
disebut
kurva geodesik
jika

\mavergleichskettedisp
{\vergleichskette
{ \nabla_{d/dt}  \gamma'
}
{ =} { 0
}
{ } {
}
{ } {
}
{ } {
}
}
{}{}{}
berlaku pada $I$.
}

 }



\inputaufgabeklausurloesung
{3}

{

Nyatakan teorema-teorema berikut.
\aufzaehlungdrei{\stichwort {Teorema tentang kelengkungan kurva bidang yang diberikan secara implisit} {}

\mavergleichskette
{\vergleichskette
{  Y
}
{ \subseteq }{  \R^2
}
{  }{
}
{    }{
}
{   }{
}
}
{}{}{.}}{\stichwort {Rumus transformasi untuk bentuk volume positif pada manifold} {.}}{Teorema tentang kelengkungan Gauss dan kelengkungan seksional pada suatu permukaan diferensiabel

\mavergleichskette
{\vergleichskette
{  Y
}
{ \subseteq }{  \R^3
}
{  }{
}
{    }{
}
{   }{
}
}
{}{}{}}

}
{

\aufzaehlungdrei{\faktsituation {Misalkan

\mavergleichskette
{\vergleichskette
{ W
}
{ \subseteq }{ \R^2
}
{  }{
}
{    }{
}
{   }{
}
}
{}{}{}
\definitionsverweis {terbuka}{}{,}
\maabb {h} { W } { \R
} {}
adalah suatu
\definitionsverweis {fungsi yang dua kali terdiferensialkan secara kontinu}{}{}
dan

\mavergleichskette
{\vergleichskette
{ Y
}
{ = }{ h^{-1}(c)
}
{  }{
}
{    }{
}
{   }{
}
}
{}{}{}
adalah
\definitionsverweis {serat}{}{}
di atas

\mavergleichskette
{\vergleichskette
{ c
}
{ \in }{ \R
}
{  }{
}
{    }{
}
{   }{
}
}
{}{}{,}
dengan $h$
\definitionsverweis {reguler}{}{}
di setiap titik $Y$. Misalkan
\maabbdisp {N} { U } {\R^2
} {}
yang pada suatu lingkungan terbuka

\mavergleichskette
{\vergleichskette
{ Y
}
{ \subseteq }{ U
}
{  }{
}
{    }{
}
{   }{
}
}
{}{}{}
didefinisikan sebagai
\definitionsverweis {medan normal satuan}{}{}
untuk $Y$, dan misalkan

\mavergleichskette
{\vergleichskette
{ P
}
{ \in }{ Y
}
{  }{
}
{    }{
}
{   }{
}
}
{}{}{.}}
\faktfolgerung {Maka, untuk setiap vektor singgung

\mavergleichskette
{\vergleichskette
{ v
}
{ \in }{ T_PY
}
{  }{
}
{    }{
}
{   }{
}
}
{}{}{}

\mavergleichskettedisp
{\vergleichskette
{ \left(  D N   \right)_{ P } \left(  v  \right)
}
{ =} { - \kappa(P) v
}
{ } {
}
{ } {
}
{ } {
}
}
{}{}{,}
dengan $\kappa(P)$ adalah
\definitionsverweis {kelengkungan}{}{}
dari suatu parametrisasi panjang busur $Y$ yang sesuai dengan orientasi yang ditentukan oleh $v, N(P)$.}
\faktzusatz {}
\faktzusatz {}}{\faktsituation {Misalkan $M$ adalah suatu
\definitionsverweis {manifold diferensiabel}{}{}
dengan
\definitionsverweis {basis topologi terhitung}{}{}
dan misalkan $\omega$ adalah suatu
\definitionsverweis {bentuk volume positif}{}{}
pada $M$. Misalkan
\maabbdisp {\varphi} { L } { M
} {}
adalah suatu
\definitionsverweis {difeomorfisme}{}{}
dari manifold $L$ dan

\mavergleichskette
{\vergleichskette
{ S
}
{ \subseteq }{ L
}
{  }{
}
{    }{
}
{   }{
}
}
{}{}{}
adalah suatu
\definitionsverweis {himpunan bagian terukur}{}{.}}
\faktfolgerung {Maka

\mavergleichskettedisp
{\vergleichskette
{ \int_S \varphi^* \omega
}
{ =} { \int_{ \varphi(S)} \omega
}
{ } {
}
{ } {
}
{ } {
}
}
{}{}{.}}
\faktzusatz {}
\faktzusatz {}}{\faktsituation {Misalkan

\mavergleichskette
{\vergleichskette
{ W
}
{ \subseteq }{ \R^3
}
{  }{
}
{    }{
}
{   }{
}
}
{}{}{}
\definitionsverweis {terbuka}{}{}
dan misalkan
\maabb {h} { W } { \R
} {}
adalah suatu
\definitionsverweis {fungsi yang dua kali terdiferensialkan secara kontinu}{}{}
dan

\mavergleichskette
{\vergleichskette
{ Y
}
{ = }{ h^{-1}(c)
}
{  }{
}
{    }{
}
{   }{
}
}
{}{}{}
adalah serat di atas

\mavergleichskette
{\vergleichskette
{ c
}
{ \in }{ \R
}
{  }{
}
{    }{
}
{   }{
}
}
{}{}{,}
dengan $h$
\definitionsverweis {reguler}{}{}
di setiap titik $Y$.}
\faktfolgerung {Maka
\definitionsverweis {kelengkungan Gauss}{}{}
dari $Y$ sama dengan
\definitionsverweis {kelengkungan seksional}{}{}
dari $Y$.}
\faktzusatz {}
\faktzusatz {}}

 }



\inputaufgabeklausurloesung
{5}

{

Misalkan

\mavergleichskette
{\vergleichskette
{ G
}
{ \subseteq }{ \R^n
}
{  }{
}
{    }{
}
{   }{
}
}
{}{}{}
\definitionsverweis {terbuka}{}{}
dan
\maabbdisp {\varphi} { G } { \R^m
} {}
adalah
\definitionsverweis {pemetaan yang terdiferensialkan secara kontinu}{}{,}
yang pada titik

\mavergleichskette
{\vergleichskette
{ P
}
{ \in }{ G
}
{  }{
}
{    }{
}
{   }{
}
}
{}{}{}
memiliki
\definitionsverweis {diferensial total}{}{}
yang surjektif. Misalkan

\mavergleichskette
{\vergleichskette
{ v
}
{ \in }{ T_PZ
}
{  }{
}
{    }{
}
{   }{
}
}
{}{}{}
adalah suatu vektor dalam
\definitionsverweis {ruang singgung serat}{}{}
$Z$ dari $\varphi$ yang melalui $P$. Tunjukkan bahwa terdapat
\definitionsverweis {kurva yang terdiferensialkan secara kontinu}{}{}
\maabbdisp {\gamma} { ]- \epsilon, \epsilon[ } { Z
} {}
\zusatzklammer {untuk suatu

\mavergleichskettek
{\vergleichskettek
{ \epsilon
}
{ > }{ 0
}
{  }{
}
{    }{
}
{   }{
}
}
{}{}{}
yang sesuai} {} {}
dengan

\mavergleichskette
{\vergleichskette
{ \gamma(0)
}
{ = }{ P
}
{  }{
}
{    }{
}
{   }{
}
}
{}{}{}
dan

\mavergleichskettedisp
{\vergleichskette
{ \gamma'(0)
}
{ =} { v
}
{ } {
}
{ } {
}
{ } {
}
}
{}{}{}.

}
{

Berlaku

\mavergleichskette
{\vergleichskette
{v
}
{ \in }{ T_PZ
}
{ = }{ \ker  \left(D\varphi\right)_{P}
}
{    }{
}
{   }{
}
}
{}{}{.}
Berdasarkan
Satz über implizite Abbildungen,
terdapat suatu himpunan terbuka

\mathbed {P \in W} {}
 {W \subseteq G} {}
 {} {} {} {,}
suatu himpunan terbuka

\mavergleichskette
{\vergleichskette
{ V
}
{ \subseteq }{ \R^{n-m}
}
{  }{
}
{    }{
}
{   }{
}
}
{}{}{}
dan suatu pemetaan yang terdiferensialkan secara kontinu
\maabbdisp {\psi} { V } { W
} {}
sedemikian sehingga

\mavergleichskette
{\vergleichskette
{ \psi(V)
}
{ \subseteq }{ Z \cap W
}
{  }{
}
{    }{
}
{   }{
}
}
{}{}{}
berlaku dan $\psi$ menginduksi suatu
\definitionsverweis {bijeksi}{}{}
\maabbdisp {\psi} { V } { Z \cap W
} {}
tersebut. Misalkan

\mavergleichskette
{\vergleichskette
{ Q
}
{ \in }{ V
}
{  }{
}
{    }{
}
{   }{
}
}
{}{}{}
adalah titik yang memenuhi

\mavergleichskette
{\vergleichskette
{ \psi(Q)
}
{ = }{ P
}
{  }{
}
{    }{
}
{   }{
}
}
{}{}{.}
Pemetaan $\psi$ di setiap titik

\mavergleichskette
{\vergleichskette
{ Q
}
{ \in }{ V
}
{  }{
}
{    }{
}
{   }{
}
}
{}{}{}
bersifat
\definitionsverweis {reguler}{}{}
dan untuk
\definitionsverweis {diferensial total}{}{}
dari $\psi$ berlaku

\mavergleichskettedisp
{\vergleichskette
{ \left(D \varphi \right)_{ P }  \circ \left(D \psi \right)_{ Q }
}
{ =} { 0
}
{ } {
}
{ } {
}
{ } {
}
}
{}{}{,}
sehingga

\mavergleichskettedisp
{\vergleichskette
{ \operatorname{im}  \left(D \psi\right)_{ Q }
}
{ =} { \ker  \left(D\varphi \right)_{ P }
}
{ =} { T_P Z
}
{ } {
}
{ } {
}
}
{}{}{.}
Karena regularitas $\psi$ di $Q$, pemetaan
\maabbdisp {\left(D \psi \right)_{ Q }} { \R^{n-m}  } { \R^n
} {}
bersifat injektif dan pemetaan
\maabbdisp {\left(D \psi \right)_{ Q }} { \R^{n-m}  } { T_PZ
} {}
bersifat bijektif. Misalkan

\mavergleichskette
{\vergleichskette
{ u
}
{ \in }{ \R^{n-m}
}
{  }{
}
{    }{
}
{   }{
}
}
{}{}{}
adalah prapeta dari $v$ dan misalkan
\maabbeledisp {\theta} { ]- \epsilon, \epsilon[ } { \R^{n-m}
} { t } { Q+tu
} {,}
dengan $\epsilon$ dipilih cukup kecil sehingga citranya seluruhnya termuat dalam $V$. Maka
\maabbdisp {\gamma = \psi \circ \theta} { ]- \epsilon, \epsilon[  } { \R^n
} {}
memenuhi

\mavergleichskettedisp
{\vergleichskette
{ \gamma(0)
}
{ =} { \psi( \theta (0))
}
{ =} { \psi(Q)
}
{ =} { P
}
{ } {
}
}
{}{}{}
dan

\mavergleichskettedisp
{\vergleichskette
{ \gamma' (0)
}
{ =} { \left(D \psi\right)_{Q} ( \theta' (0))
}
{ =} { \left(D \psi\right)_{Q} ( u)
}
{ =} { v
}
{ } {
}
}
{}{}{.}

 }



\inputaufgabeklausurloesung
{2}

{

Misalkan
\mathkor {} {r} {dan} {R} {}
adalah bilangan real positif dengan

\mavergleichskette
{\vergleichskette
{ 0
}
{ < }{ r
}
{ < }{ R
}
{    }{
}
{   }{
}
}
{}{}{.}
Tentukan suatu
\definitionsverweis {medan normal satuan}{}{}
untuk torus
\zusatzklammer {tertanam} {} {}

\mavergleichskettedisp
{\vergleichskette
{ T
}
{ =} { \left\{ (x,y,z) \in \R^3  \mid   \left(  \sqrt{x^2+y^2} -R  \right)^2 +z^2  = r^2         \right\}
}
{ } {
}
{ } {
}
{ } {
}
}
{}{}{.}

}
{

Berlaku

\mavergleichskettedisp
{\vergleichskette
{ h(x,y,z)
}
{ =} { x^2+y^2 -2 \sqrt{x^2+y^2}R+z^2 + R^2
}
{ } {
}
{ } {
}
{ } {
}
}
{}{}{}
dan

\mavergleichskette
{\vergleichskette
{ T
}
{ = }{ h^{-1} (r^2)
}
{  }{
}
{    }{
}
{   }{
}
}
{}{}{.}
Gradien dari $h$ adalah

\mavergleichskettedisp
{\vergleichskette
{
\operatorname{Grad} \,  h
}
{ =} { 2 \begin{pmatrix}   x   -  xR \left(  x^2+y^2  \right)^{- { \frac{   1  }{  2 } } }  \\ y   -  y R \left(  x^2+y^2  \right)^{- { \frac{   1  }{  2 } } } \\ z   \end{pmatrix}
}
{ } {
}
{ } {
}
{ } {
}
}
{}{}{.}
Dengan demikian,
\definitionsverweis {medan normal satuan}{}{}
adalah

\mavergleichskettedisp
{\vergleichskette
{ N(x,y,z)
}
{ =} { { \frac{   1  }{  \sqrt{ \left(   x   -  xR \left(  x^2+y^2  \right)^{- { \frac{   1  }{  2 } } }  \right)^2 + \left(  y   -  y R \left(  x^2+y^2  \right)^{- { \frac{   1  }{  2 } } }  \right)^2   +z^2   }  } } \begin{pmatrix}   x   -  xR \left(  x^2+y^2  \right)^{- { \frac{   1  }{  2 } } }  \\ y   -  y R \left(  x^2+y^2  \right)^{- { \frac{   1  }{  2 } } } \\ z   \end{pmatrix}
}
{ } {
}
{ } {
}
{ } {
}
}
{}{}{.}

 }



\inputaufgabeklausurloesung
{4}

{

Jelaskan peran teorema tentang pemetaan implisit dalam pengembangan teori manifold.

}
{  }



\inputaufgabeklausurloesung
{8}

{

Buktikan teorema tentang sifat adjoin-diri (self-adjoint) dari pemetaan Weingarten.

}
{

Untuk vektor-vektor

\mavergleichskette
{\vergleichskette
{ v,w
}
{ \in }{ T_PY
}
{  }{
}
{    }{
}
{   }{
}
}
{}{}{}
harus ditunjukkan bahwa

\mavergleichskettedisp
{\vergleichskette
{ \left\langle  v  ,  L_P(w)  \right\rangle
}
{ =} { \left\langle L_P(v) ,  w  \right\rangle
}
{ } {
}
{ } {
}
{ } {
}
}
{}{}{}
berlaku. Dengan

\mavergleichskette
{\vergleichskette
{ N(P)
}
{ = }{ { \frac{
\operatorname{Grad} \,  h  (  P )  }{  \Vert {
\operatorname{Grad} \,  h  (  P ) } \Vert  } }
}
{  }{
}
{    }{
}
{   }{
}
}
{}{}{}
berdasarkan
Lemma 45.10 (Analysis (Osnabrück 2021-2023))
berlaku

\mavergleichskettealignhandlinks
{\vergleichskettealignhandlinks
{ \left\langle  v  ,  L_P(w)  \right\rangle
}
{ =} { -\left\langle  v  ,  \left(  D_{w}  N   \right) \left(   P   \right)   \right\rangle
}
{ =} { -\left\langle  v  ,  \left(  D_{w}  { \frac{
\operatorname{Grad} \,  h  }{  \Vert {
\operatorname{Grad} \,  h } \Vert  } }   \right) \left(   P   \right)     \right\rangle
}
{ =} { -\left\langle  v  ,  \left(  D_{w}  { \frac{   1  }{  \Vert {
\operatorname{Grad} \,  h } \Vert  } }   \right) \left(   P   \right) \cdot
\operatorname{Grad} \,  h  (  P ) + { \frac{   1  }{  \Vert {
\operatorname{Grad} \,  h  (  P ) } \Vert  } } \cdot  \left(  D_{w}
\operatorname{Grad} \,  h   \right) \left(   P   \right)   \right\rangle
}
{ =} { - { \frac{   1  }{  \Vert {
\operatorname{Grad} \,  h  (  P ) } \Vert  } } \left\langle  v  ,  \left(  D_{w}
\operatorname{Grad} \,  h   \right) \left(   P   \right)    \right\rangle
}
}
{}
{}{,}
karena suku pertama tegak lurus terhadap vektor singgung $v$. Dalam fungsi-fungsi koordinat, berlaku

\mavergleichskettealignhandlinks
{\vergleichskettealignhandlinks
{ \left(  D_{w}
\operatorname{Grad} \,  h   \right) \left(   P   \right)
}
{ =} { \left(  D_{w}  \left(  { \frac{   \partial   h   }{  \partial   x_1   } }   , \,   \ldots  , \,   { \frac{   \partial   h   }{  \partial   x_n   } }                 \right)    \right) \left(   P   \right)
}
{ =} { \sum_{j = 1}^n  w_j  { \frac{   \partial    }{  \partial   x_j   } }  \left(  { \frac{   \partial   h   }{  \partial   x_1   } }   , \,   \ldots  , \,   { \frac{   \partial   h   }{  \partial   x_n   } }                 \right) (P)
}
{ =} { \left(  \sum_{j = 1}^n  w_j  { \frac{   \partial    }{  \partial   x_j   } }  { \frac{   \partial   h   }{  \partial   x_1   } } (P)   , \,   \ldots  , \,   \sum_{j = 1}^n  w_j  { \frac{   \partial    }{  \partial   x_j   } }  { \frac{   \partial   h   }{  \partial   x_n   } } (P)                 \right)
}
{ } {
}
}
{}
{}{.}
Dengan demikian, ekspresi di atas sama dengan

\mavergleichskettealign
{\vergleichskettealign
{ \left\langle  v  ,  L_P(w)  \right\rangle
}
{ =} { - { \frac{   1  }{  \Vert {
\operatorname{Grad} \,  h  (  P ) } \Vert  } } \left\langle  v  ,  \left(  D_{w}
\operatorname{Grad} \,  h   \right) \left(   P   \right)   \right\rangle
}
{ =} { -  { \frac{   1  }{  \Vert {
\operatorname{Grad} \,  h  (  P ) } \Vert  } }  \sum_{i,j = 1}^n v_i  w_j  { \frac{   \partial    }{  \partial   x_j   } } { \frac{   \partial   h   }{  \partial   x_i   } }  (P)
}
{ } {
}
{ } {
}
}
{}
{}{.}
Berdasarkan
Satz 44.10 (Analysis (Osnabrück 2021-2023))
urutan turunan parsial dapat dipertukarkan, sehingga peran
\mathkor {} {v} {dan} {w} {}
juga dapat dipertukarkan.

 }



\inputaufgabeklausurloesung
{4}

{

Misalkan $X$ adalah suatu
\definitionsverweis {ruang kompak}{}{}
dan misalkan

\mavergleichskette
{\vergleichskette
{Y
}
{ \subseteq }{ X
}
{  }{
}
{    }{
}
{   }{
}
}
{}{}{}
adalah
\definitionsverweis {himpunan bagian tertutup}{}{,}
yang dilengkapi dengan
\definitionsverweis {topologi terinduksi}{}{.}
Tunjukkan bahwa $Y$ juga kompak.

}
{

Misalkan

\mavergleichskettedisp
{\vergleichskette
{ Y
}
{ =} { \bigcup_{i \in I} V_i
}
{ } {
}
{ } {
}
{ } {
}
}
{}{}{}
adalah suatu tutupan oleh himpunan-himpunan bagian terbuka

\mavergleichskette
{\vergleichskette
{ V_i
}
{ \subseteq }{ Y
}
{  }{
}
{    }{
}
{   }{
}
}
{}{}{.}
Ini berarti bahwa terdapat himpunan-himpunan terbuka

\mavergleichskette
{\vergleichskette
{ U_i
}
{ \subseteq }{ X
}
{  }{
}
{    }{
}
{   }{
}
}
{}{}{}
yang memenuhi

\mavergleichskette
{\vergleichskette
{ V_i
}
{ = }{ Y \cap U_i
}
{  }{
}
{    }{
}
{   }{
}
}
{}{}{.}
Oleh karena itu,

\mavergleichskettedisp
{\vergleichskette
{ Y
}
{ \subseteq} { \bigcup_{i \in I} U_i
}
{ } {
}
{ } {
}
{ } {
}
}
{}{}{.}
Karena $Y$ tertutup dalam $X$, himpunan $X \setminus Y$ terbuka, sehingga

\mavergleichskettedisp
{\vergleichskette
{ X
}
{ =} { Y \cup (X \setminus Y)
}
{ =} { \bigcup_{i \in I} U_i  \cup  (X \setminus Y)
}
{ } {
}
{ } {
}
}
{}{}{}
adalah suatu tutupan terbuka dari $X$. Karena $X$ kompak, terdapat suatu subtutupan berhingga

\mavergleichskettedisp
{\vergleichskette
{ X
}
{ =} { \bigcup_{i \in J} U_i  \cup  (X \setminus Y)
}
{ } {
}
{ } {
}
{ } {}
}
{}{}{.}
Ini selanjutnya berarti

\mavergleichskettedisp
{\vergleichskette
{ Y
}
{ =} { \bigcup_{i \in J} U_i
}
{ } {
}
{ } {
}
{ } {}
}
{}{}{.}

 }



\inputaufgabeklausurloesung
{5}

{

Berikan contoh suatu
\definitionsverweis {manifold diferensiabel}{}{}
$M$ yang
\definitionsverweis {berdimensi dua}{}{}
dan
\definitionsverweis {terhubung}{}{,}
beserta suatu titik

\mavergleichskette
{\vergleichskette
{ P
}
{ \in }{ M
}
{  }{
}
{    }{
}
{   }{
}
}
{}{}{}
sedemikian sehingga
\mathkor {} {M} {dan} {M \setminus \{P\}} {}
saling
\definitionsverweis {difeomorfik}{}{.}

}
{

Misalkan

\mavergleichskettedisp
{\vergleichskette
{M
}
{ =} { \R^2 \setminus \N_+
}
{ =} { \R^2 \setminus \{ (1,0) , (2,0), (3,0), \ldots  \}
}
{ } {
}
{ } {
}
}
{}{}{,}
yakni dari $\R^2$ dihapus bilangan-bilangan asli positif yang terletak pada sumbu $x$. Himpunan ini merupakan himpunan bagian terbuka dari $\R^2$, sehingga merupakan manifold diferensiabel. Manifold ini terhubung lintasan karena, misalnya, setiap titik di $M$ dengan
\mathl{y \geq 0}{} dapat dihubungkan dengan
\mathl{(0,1)}{} melalui lintasan garis lurus dan setiap titik di $M$ dengan
\mathl{y \leq 0}{} dapat dihubungkan dengan
\mathl{(0,-1)}{} melalui lintasan garis lurus, sedangkan kedua titik tersebut juga dapat dihubungkan dengan garis lurus. Sekarang misalkan

\mavergleichskettedisp
{\vergleichskette
{P
}
{ =} { (0,0)
}
{ } {
}
{ } {
}
{ } {
}
}
{}{}{}
dan

\mavergleichskettedisp
{\vergleichskette
{M'
}
{ =} { M \setminus \{P\}
}
{ } {
}
{ } {
}
{ } {
}
}
{}{}{.}
Pemetaan
\maabbeledisp {} {\R^2} {\R^2
} {(x,y)} {(x-1,y)
} {,}
merupakan suatu difeomorfisme
\zusatzklammer {sebagai translasi} {} {}
dan citra $M$ tepat sama dengan $M'$. Oleh karena itu,
\mathkor {} {M} {dan} {M'=M \setminus \{P\}} {}
saling difeomorfik.

 }



\inputaufgabeklausurloesung
{3}

{

Untuk setiap
\definitionsverweis {vektor singgung}{}{}

\mavergleichskette
{\vergleichskette
{ v
}
{ \in }{ T_PS^2
}
{ \subseteq }{ \R^3
}
{    }{
}
{   }{
}
}
{}{}{}
dengan

\mavergleichskette
{\vergleichskette
{ \Vert {v} \Vert
}
{ = }{ 1
}
{  }{
}
{    }{
}
{   }{
}
}
{}{}{}
di suatu titik

\mavergleichskette
{\vergleichskette
{ P
}
{ \in }{ S^2
}
{ \subset }{ \R^3
}
{    }{
}
{   }{
}
}
{}{}{}
pada
\definitionsverweis {permukaan bola satuan}{}{,}
berikan suatu wakil
\definitionsverweis {diferensiabel}{}{}
\maabbdisp {\gamma} { I } { S^2
} {}
dengan

\mavergleichskette
{\vergleichskette
{ \gamma(0)
}
{ = }{ P
}
{  }{
}
{    }{
}
{   }{
}
}
{}{}{.}

}
{

Misalkan

\mavergleichskette
{\vergleichskette
{P
}
{ = }{ \begin{pmatrix}  x  \\ y \\ z   \end{pmatrix}
}
{  }{
}
{    }{
}
{   }{
}
}
{}{}{}
dan

\mavergleichskette
{\vergleichskette
{v
}
{ = }{ \begin{pmatrix}  a  \\ b \\ c   \end{pmatrix}
}
{  }{
}
{    }{
}
{   }{
}
}
{}{}{.}
Kedua vektor tersebut bernorma satu dan saling tegak lurus karena bidang singgung pada permukaan bola tegak lurus terhadap vektor posisi. Tinjau lintasan
\maabbdisp {\gamma} {\R} { \R^3
} {}
yang diberikan oleh

\mavergleichskettedisp
{\vergleichskette
{ \gamma(t)
}
{ =} {  \left(  \cos t  \right)  \begin{pmatrix}  x  \\ y \\ z   \end{pmatrix}  + \left(  \sin t  \right)  \begin{pmatrix}  a  \\ b \\ c   \end{pmatrix}
}
{ } {
}
{ } {
}
{ } {
}
}
{}{}{.}
Berlaku

\mavergleichskettedisp
{\vergleichskette
{ \Vert { \gamma(t)} \Vert
}
{ =} { \cos^{ 2  }  t + \sin^{ 2  }  t
}
{ =} { 1
}
{ } {
}
{ } {
}
}
{}{}{,}
sehingga lintasan tersebut seluruhnya berada pada permukaan bola. Selanjutnya,

\mavergleichskettedisp
{\vergleichskette
{ \gamma(0)
}
{ =} { \left(  \cos  0  \right) \begin{pmatrix}  x  \\ y \\ z   \end{pmatrix} + \left(  \sin  0  \right)  \begin{pmatrix}  a  \\ b \\ c   \end{pmatrix}
}
{ =} { \begin{pmatrix}  x  \\ y \\ z   \end{pmatrix}
}
{ } {
}
{ } {
}
}
{}{}{}
dan

\mavergleichskettedisp
{\vergleichskette
{ \gamma'(0)
}
{ =} { - \left(  \sin  0  \right) \begin{pmatrix}  x  \\ y \\ z   \end{pmatrix} + \left(  \cos  0  \right) \begin{pmatrix}  a  \\ b \\ c   \end{pmatrix}
}
{ =} { \begin{pmatrix}  a  \\ b \\ c   \end{pmatrix}
}
{ } {
}
{ } {
}
}
{}{}{.}

 }



\inputaufgabeklausurloesung
{2}

{

Tinjau parabola standar

\mavergleichskette
{\vergleichskette
{ Y
}
{ \subseteq }{ \R^2
}
{  }{
}
{    }{
}
{   }{
}
}
{}{}{}
dengan struktur Riemann terinduksi dan parametrisasi
\maabbeledisp {\varphi} {\R} { Y
} { t } { \left(  t   , \,  t^2                   \right)
} {,}
yang kita pandang sebagai suatu peta
\zusatzklammer {invers} {} {}.
Tentukan
\definitionsverweis {fungsi fundamental}{}{}
Riemann

\mavergleichskette
{\vergleichskette
{ g(t)
}
{ = }{ g_{11} (t)
}
{  }{
}
{    }{
}
{   }{
}
}
{}{}{.}

}
{

Berlaku

\mavergleichskettealign
{\vergleichskettealign
{ g(t)
}
{ =} { \left\langle  T_t \varphi(1)   , T_t \varphi(1)   \right\rangle_{ Y, \varphi(t) }
}
{ =} { \left\langle    \varphi' (t)   ,   \varphi'(t)   \right\rangle
}
{ =} { \left\langle  \begin{pmatrix}  1 \\ 2t   \end{pmatrix}     ,  \begin{pmatrix}  1 \\ 2t   \end{pmatrix}   \right\rangle
}
{ =} { 1+4t^2
}
}
{}
{}{.}

 }



\inputaufgabeklausurloesung
{6}

{

Hitung bentuk diferensial tarik balik
\mathl{\varphi^* \tau}{} dari

\mavergleichskettedisp
{\vergleichskette
{ \tau
}
{ =} { dx \wedge dy \wedge dz - w dx \wedge dy \wedge dw + \cos (xy)  dx \wedge dz \wedge dw-yw dy \wedge dz \wedge dw
}
{ } {
}
{ } {
}
{ } {
}
}
{}{}{}
oleh pemetaan
\maabbeledisp {\varphi} { \R^3 } { \R^4
} { (r,s,t) } {  \left(  r^2s,t, \sin  r ,e^{st}  \right)  = (x,y,z,w)
} {.}

}
{

Diperoleh

\mavergleichskettedisp
{\vergleichskette
{dx
}
{ =} { d(r^2s)
}
{ =} { 2rs dr  +r^2ds
}
{ } {
}
{ } {
}
}
{}{}{,}

\mavergleichskettedisp
{\vergleichskette
{dy
}
{ =} { dt
}
{ } {
}
{ } {
}
{ } {
}
}
{}{}{,}

\mavergleichskettedisp
{\vergleichskette
{dz
}
{ =} { d (\sin  r )
}
{ =} { \cos  r  dr
}
{ } {
}
{ } {
}
}
{}{}{}
dan

\mavergleichskettedisp
{\vergleichskette
{dw
}
{ =} { d(e^{st})
}
{ =} { se^{st} dt+te^{st} ds
}
{ } {
}
{ } {
}
}
{}{}{.}
Dengan demikian,

\mavergleichskettealign
{\vergleichskettealign
{ \varphi^* \tau
}
{ =} { d(r^2s) \wedge dt \wedge d( \sin  r )  -e^{st} d(r^2s) \wedge dt \wedge d(e^{st})
}
{ \, \, \, \,} {+ \cos (r^2st)  d (r^2s) \wedge d ( \sin  r ) \wedge d(e^{st}) - te^{st} dt \wedge d( \sin  r ) \wedge d(e^{st})
}
{ =} { r^2 \cos (r^2st) \cos  r  ds \wedge dt \wedge dr  - 2rs t e^{st} e^{st} dr \wedge dt \wedge ds
}
{ \, \, \, \,} { + r^2 se^{st} \cos  r  ds \wedge dr \wedge dt  -  t^2 e^{st}  e^{st} \cos  r dt \wedge dr \wedge ds
}
}
{
\vergleichskettefortsetzungalign
{ =} { (r^2 \cos  r +2 rst e^{2st} - r^2 s e ^{st} \cos (r^2st) \cos  r - t^2e^{ 2st} \cos  r     )  dr \wedge ds \wedge dt
}
{ } {}
{ } {}
{ } {}
}
{}{}

 }



\inputaufgabeklausurloesung
{7 (2+3+2)}

{

Tinjau pemetaan-pemetaan diferensiabel
\maabbeledisp {\gamma} { [1,2] } {\R^2
} { t } {(t,t^{-1})
} {,}
dan
\maabbeledisp {\varphi} {\R^2} {\R^3
} {(u,v)} {(u^2,uv,-u+v^2)
} {,}
serta bentuk diferensial

\mavergleichskettedisp
{\vergleichskette
{ \omega
}
{ =} { xdx-zdy+dz
}
{ } {
}
{ } {
}
{ } {
}
}
{}{}{}
pada $\R^3$.

a) Hitung bentuk diferensial tarik balik
\mathl{\varphi^* \omega}{} pada $\R^2$.

b) Hitung integral lintasan bentuk diferensial
\mathl{\varphi^* \omega}{} sepanjang lintasan $\gamma$.

c) Hitung
\zusatzklammer {tanpa menggunakan hasil b)} {} {}
integral lintasan bentuk diferensial
\mathl{\omega}{} sepanjang lintasan $\varphi \circ \gamma$.

}
{

a) Bentuk diferensial tarik baliknya adalah

\mavergleichskettealign
{\vergleichskettealign
{\varphi^* ( xdx -zdy +dz )
}
{ =} { u^2 du^2    -(-u+v^2) duv +d (-u+v^2)
}
{ =} { 2u^3 du + (u-v^2) (vdu +udv) -du +2vdv
}
{ =} { (2u^3 +uv -v^3 -1) du    + (u^2 -uv^2 +2v )dv
}
{ } {
}
}
{}
{}{.}

b) Integral lintasannya adalah

\mavergleichskettealign
{\vergleichskettealign
{ \int_1^2 ( 2t^3 +1 -t^{-3} -1 )dt  +(t^2 -t^{-1} +2t^{-1} ) d t^{-1}
}
{ =} { \int_1^2 ( 2t^3 -t^{-3} ) dt  +(t^2 + t^{-1} ) d t^{-1}
}
{ =} { \int_1^2 ( 2t^3 -t^{-3} ) dt  -(1 + t^{-3} ) dt
}
{ =} { \int_1^2 ( 2t^3 - 2t^{-3} -1 ) dt
}
{ =} { \left(  { \frac{   1 }{  2 } }  t^4  + t^{-2}  -t  \right) \vert_{  1 }^{  2 }
}
}
{
\vergleichskettefortsetzungalign
{ =} { 8 + { \frac{   1 }{  4 } }  -2 - { \frac{   1 }{  2 } } -1 +1
}
{ =} { 5 + { \frac{   3 }{  4 } }
}
{ } {}
{ } {}
}
{}{.}

c) Lintasan komposisinya adalah

\mavergleichskettedisp
{\vergleichskette
{ \psi (t)
}
{ =} { (t^2,1,-t +t^{-2} )
}
{ } {
}
{ } {
}
{ } {
}
}
{}{}{.}
Dengan demikian,

\mavergleichskettealign
{\vergleichskettealign
{ \int_1^2 \psi^* (xdx -zdy +dz )
}
{ =} { \int_1^2 t^2dt^2 + d (-t + t^{-2} )
}
{ =} {  \int_1^2 (2 t^3-1 -2t^{-3} ) dt
}
{ =} { 5 + { \frac{   3 }{  4 } }
}
{ } {
}
}
{}
{}{.}

 }



\inputaufgabeklausurloesung
{4}

{

Misalkan $M$ adalah suatu
\definitionsverweis {manifold berbatas}{}{.}
Tunjukkan bahwa setiap himpunan bagian terbuka

\mavergleichskette
{\vergleichskette
{ N
}
{ \subseteq }{ M
}
{  }{
}
{    }{
}
{   }{
}
}
{}{}{}
juga merupakan manifold dengan batas
\zusatzklammer {yang mungkin kosong} {} {},
dan bahwa

\mavergleichskettedisp
{\vergleichskette
{ \partial N
}
{ =} { \partial M \cap N
}
{ } {
}
{ } {
}
{ } {
}
}
{}{}{}
berlaku.

}
{

Misalkan
\mathl{P \in N}{} dan
\mathl{P \in U}{} adalah suatu domain peta dari $M$. Misalkan
\maabbdisp {\alpha} { U } {V
} {}
adalah suatu peta dengan himpunan terbuka
\mathl{V \subseteq H}{,}
\mathl{H=\R_{\geq 0} \times \R^{n-1}}{.} Peta ini menginduksi suatu homeomorfisme
\maabbdisp {} {N \cap U} { \alpha (N \cap U)
} {.}
Peta-peta ini kita gunakan untuk $N$. Sifat difeomorfisme dari transisi peta pada $M$ langsung diwarisi oleh peta-peta pada $N$. Misalkan
\mathl{P \in N \cap \partial M}{.} Maka suatu peta memetakan $P$ ke titik batas $H$. Karena peta-peta untuk $N$ merupakan pembatasan dari peta-peta tersebut, sifat ini tetap berlaku. Sebaliknya, jika
\mathl{P \in \partial N}{,} maka di satu pihak tentu saja
\mathl{P \in N \subseteq M}{.} Di pihak lain, untuk $P$ terdapat suatu peta pada $N$ yang diperoleh dengan membatasi suatu peta untuk $M$, dan peta tersebut memetakan $P$ ke titik batas $H$.

 }



\inputaufgabeklausurloesung
{8}

{

Buktikan teorema tentang karakterisasi kurva geodesik terhadap koneksi Levi-Civita.

}
{

Misalkan $n$ adalah dimensi $M$. Tinjau situasi ini secara langsung pada suatu citra peta terbuka

\mavergleichskette
{\vergleichskette
{ U
}
{ \subseteq }{ \R^n
}
{  }{
}
{    }{
}
{   }{
}
}
{}{}{.}
Menurut
Bemerkung *****
turunan vertikal diberikan oleh
\maabbeledisp {} {T(U \times \R^n) = U \times \R^n \times \R^n \times \R^n } { U \times \R^n \times  \R^n
} {(P,e_j, \partial_i, w)} {(P,e_j,\sum_{k = 1}^n \Gamma^k_{ij} (P) e_k +w) \cong V(U \times \R^n)
} {}
; pemetaan ke $TU$ diperoleh dengan menghilangkan komponen tengah. Turunan kedua kurva mula-mula merupakan pemetaan singgung kedua, yaitu
\zusatzklammer {di sini kita mengabaikan perkalian dengan arah satu dimensi dari ruang singgung kurva} {} {}
\maabbeledisp {T(\gamma)} { I } { TU
} { t } { (\gamma(t), \gamma'(t))
} {,}
dan dengan demikian
\maabbeledisp {T(T(\gamma))} { I } { T(TU)
} { t } { ((\gamma(t), \gamma'(t)), (\gamma'(t), \gamma^{\prime \prime} (t)))
} {}
\zusatzklammer {kedua komponen pemetaan singgung diturunkan} {} {.}
Dengan

\mavergleichskette
{\vergleichskette
{ \gamma'(t)
}
{ = }{ \sum_{j = 1}^n \gamma_j'(t) \partial_j
}
{  }{
}
{    }{
}
{   }{
}
}
{}{}{}
proyeksi vertikal memetakan bentuk ini ke

\mathdisp {\sum_{k = 1}^n  \left(  \sum_{i,j} \gamma_i'(t) \gamma_j'(t) \Gamma^k_{ij}(\gamma(t))  \right)  \partial_k  +\sum_{k = 1}^n \gamma_k^{\prime \prime}(t) \partial_k} {  }

. Syarat bagi suatu kurva geodesik bahwa turunan keduanya
\zusatzklammer {di $TTM$} {} {}
selalu horizontal ekuivalen dengan proyeksi vertikal yang dihitung bernilai $0$. Ini berarti bahwa setiap komponennya bernilai $0$, yaitu

\mavergleichskettedisp
{\vergleichskette
{ \sum_{i,j} \gamma_i'(t) \gamma_j'(t) \Gamma^k_{ij}(\gamma(t))  + \gamma_k^{\prime \prime}(t)
}
{ =} { 0
}
{ } {
}
{ } {
}
{ } {
}
}
{}{}{}
untuk

\mavergleichskette
{\vergleichskette
{ k
}
{ = }{ 1  , \ldots , n
}
{  }{
}
{    }{
}
{   }{
}
}
{}{}{.}

 }
\endgroup

\chapter{Formulir Solusi Resmi Ujian 6}
\noindent\textit{Solusi yang disediakan oleh sumber resmi; kekosongan sumber dipertahankan.}
\begingroup
\renewcommand{\theHfakt}{o011.exam.official.06.\arabic{fakt}}
%Data institusi

%\input{Dozentdaten}

%\renewcommand{\fachbereich}{Bidang studi}

%\renewcommand{\dozent}{Prof. Dr. . }

%Data ujian

\renewcommand{\klausurgebiet}{ }

\renewcommand{\klausurtyp}{ }

\renewcommand{\klausurdatum}{ . 20}

\klausurvorspann {\fachbereich} {\klausurdatum} {\dozent} {\klausurgebiet} {\klausurtyp}

%Data untuk tabel poin berikut

\renewcommand{\aeins}{  3  }

\renewcommand{\azwei}{  3   }

\renewcommand{\adrei}{  7  }

\renewcommand{\avier}{  4  }

\renewcommand{\afuenf}{  0  }

\renewcommand{\asechs}{  0  }

\renewcommand{\asieben}{  0  }

\renewcommand{\aacht}{  5  }

\renewcommand{\aneun}{  15  }

\renewcommand{\azehn}{  5  }

\renewcommand{\aelf}{  0  }

\renewcommand{\azwoelf}{  3  }

\renewcommand{\adreizehn}{  3   }

\renewcommand{\avierzehn}{  0  }

\renewcommand{\afuenfzehn}{  10  }

\renewcommand{\asechzehn}{  4  }

\renewcommand{\asiebzehn}{  62  }

\renewcommand{\aachtzehn}{    }

\renewcommand{\aneunzehn}{    }

\renewcommand{\azwanzig}{    }

\renewcommand{\aeinundzwanzig}{    }

\renewcommand{\azweiundzwanzig}{    }

\renewcommand{\adreiundzwanzig}{    }

\renewcommand{\avierundzwanzig}{    }

\renewcommand{\afuenfundzwanzig}{    }

\renewcommand{\asechsundzwanzig}{    }

\punktetabellesechzehn

\klausurnote

\newpage

\setcounter{section}{K}



\inputaufgabeklausurloesung
{3}

{

Definisikan istilah-istilah berikut
\zusatzklammer {yang dicetak miring} {} {}.
\aufzaehlungsechs{Suatu kurva
\stichwort {berparametrisasi panjang busur} {}
\maabbdisp {\gamma} {[a,b]} { \R^n
} {.}

}{Untuk suatu titik

\mavergleichskette
{\vergleichskette
{  P
}
{ \in }{ M
}
{  }{}
{    }{}
{   }{}
}
{}{}{}
pada suatu manifold diferensiabel $M$, yang dimaksud dengan dua kurva diferensiabel yang
\stichwort {ekuivalen secara tangensial} {} ialah
\maabbdisp {\gamma_1, \gamma_2} { I } { M
} {}
\zusatzklammer {di sini

\mavergleichskettek
{\vergleichskettek
{  0
}
{ \in }{  I
}
{  }{
}
{    }{
}
{   }{
}
}
{}{}{}
merupakan interval riil terbuka dan

\mavergleichskettek
{\vergleichskettek
{   \gamma_1(0)
}
{ = }{ \gamma_2(0)
}
{ = }{ P
}
{    }{
}
{   }{
}
}
{}{}{}} {} {.}

}{Suatu
\stichwort {penampang kontinu} {}
dari suatu pemetaan kontinu
\maabbdisp {p} {Y} {X
} {}
antara ruang-ruang topologis
\mathkor {} {X} {dan} {Y} {.}

}{\stichwort {Bentuk volume kanonik} {} pada suatu manifold Riemann berorientasi $M$.

}{\stichwort {Turunan eksterior} {}
dari suatu
\definitionsverweis {bentuk diferensial}{}{}
berderajat $k$ yang
\definitionsverweis {terdiferensialkan secara kontinu}{}{}

\mathl{\omega \in
{ \mathcal E }^{ k } (  U   )}{} pada suatu
\definitionsverweis {himpunan terbuka}{}{}

\mathl{U \subseteq \R^n}{.}

}{\stichwort {Kelengkungan seksional} {} yang berkaitan dengan vektor-vektor bebas linear

\mavergleichskette
{\vergleichskette
{  v,w
}
{ \in }{T_PM
}
{  }{
}
{    }{
}
{   }{
}
}
{}{}{}
pada suatu
\definitionsverweis {manifold Riemann}{}{}
$M$ di suatu titik

\mavergleichskette
{\vergleichskette
{ P
}
{ \in }{ M
}
{  }{
}
{    }{
}
{   }{
}
}
{}{}{.}
}

}
{

\aufzaehlungsechs{Suatu
\definitionsverweis {kurva diferensiabel}{}{}
\maabbdisp {\gamma} {[a,b]} {\R^n
} {,}
disebut
berparametrisasi panjang busur
jika

\mavergleichskette
{\vergleichskette
{ \gamma_1'(t)^2 + \cdots + \gamma_n'(t)^2
}
{ = }{ 1
}
{  }{
}
{    }{
}
{   }{
}
}
{}{}{}
berlaku untuk setiap $t$.
}{Kedua kurva
\mathkor {} {\gamma_1} {dan} {\gamma_2} {}
disebut ekuivalen secara tangensial di $P$ jika terdapat suatu lingkungan terbuka
\mathl{P \in U}{} dan suatu
\definitionsverweis {peta}{}{}
\maabbdisp {\alpha} { U } { V
} {}
dengan
\mathl{V \subseteq \R^n}{} sedemikian sehingga

\mavergleichskettedisp
{\vergleichskette
{  \left(  \alpha \circ \left(  \gamma_1 \vert_{\gamma_1^{-1} (U)}  \right)   \right)'(0)
}
{ =} { \left(  \alpha \circ \left(  \gamma_2 \vert_{\gamma_2^{-1} (U)}  \right)  \right)'(0)
}
{ } {
}
{ } {
}
{ } {
}
}
{}{}{} berlaku.
}{Yang dimaksud dengan
penampang kontinu
dari $p$ adalah suatu pemetaan kontinu
\maabb {s} {X} {Y
} {}
yang memenuhi

\mavergleichskettedisp
{\vergleichskette
{  p \circ s
}
{ =} {
\operatorname{Id}_{ X }
}
{ } {
}
{ } {
}
{ } {
}
}
{}{}{.}
}{Untuk
\mathl{P \in M}{,} misalkan
\mathl{\omega_P}{} adalah
\definitionsverweis {bentuk alternan}{}{}
pada
\mathl{T_PM}{}
\zusatzklammer {atau elemen yang bersesuaian dari \mathlk{\bigwedge^n T^*_PM}{}} {} {,}
yang memberikan nilai $1$ kepada setiap
\definitionsverweis {basis ortonormal}{}{}
yang merepresentasikan
\definitionsverweis {orientasi}{}{.}
Maka
$n$-\definitionsverweis {bentuk diferensial
}{}{}
\maabbeledisp {} { M } { \bigwedge^n T^*M
} {P } { \omega_P
} {,}
disebut bentuk volume kanonik pada $M$.
}{Bentuk $\omega$ memiliki representasi pada $U$

\mathdisp {\omega = \sum_{I \subseteq \{1 , \ldots , n \},\,  { \# \left(   I  \right) } =k } f_I dx_I} {  }

dengan
\definitionsverweis {fungsi-fungsi yang terdiferensialkan secara kontinu}{}{}
\maabbdisp {f_I} { U } {\R
} {.}
Maka turunan eksteriornya adalah bentuk diferensial berderajat $(k+1)$

\mavergleichskettedisp
{\vergleichskette
{ d \omega
}
{ =} { \sum_{I \subseteq \{1 , \ldots , n \},\,  { \# \left(   I  \right) } = k } df_I \wedge dx_I
}
{ =} { \sum_{I \subseteq \{1 , \ldots , n \},\,  { \# \left(   I  \right) } = k } \left(  \sum_{j = 1}^n { \frac{   \partial  f_I  }{  \partial  x_j  } } dx_j  \right)  \wedge dx_I
}
{ } {
}
{ } {
}
}
{}{}{.}
}{Bilangan

\mavergleichskettedisp
{\vergleichskette
{ K(v,w)
}
{ =} {  { \frac{    \left\langle R(v,w)w , v  \right\rangle  }{   \left\langle v , v  \right\rangle  \left\langle w , w  \right\rangle  -\left\langle v , w  \right\rangle^2   } }
}
{ } {
}
{ } {
}
{ } {
}
}
{}{}{,}
dengan $R$ menyatakan
\definitionsverweis {operator kelengkungan}{}{}
pada bundel tangen, disebut
\definitionswort {kelengkungan seksional}{}
yang berkaitan dengan $v,w$ di $P$.
}

 }



\inputaufgabeklausurloesung
{3}

{

Rumuskan teorema-teorema berikut.
\aufzaehlungdrei{\stichwort {Teorema tentang orientasi pada suatu hipermuka diferensiabel} {}

\mavergleichskette
{\vergleichskette
{  Y
}
{ \subseteq }{  \R^n
}
{  }{
}
{    }{
}
{   }{
}
}
{}{}{.}}{Rumus braket Lie medan-medan vektor pada suatu himpunan terbuka

\mavergleichskette
{\vergleichskette
{  U
}
{ \subseteq }{  \R^n
}
{  }{
}
{    }{
}
{   }{
}
}
{}{}{.}}{Teorema tentang
\stichwort {retraksi ke batas} {}
pada manifold.}

}
{

\aufzaehlungdrei{\faktsituation {Misalkan

\mavergleichskette
{\vergleichskette
{ W
}
{ \subseteq }{ \R^n
}
{  }{
}
{    }{
}
{   }{
}
}
{}{}{}
\definitionsverweis {terbuka}{}{,}
\maabb {h} { W } { \R
} {}
adalah suatu
\definitionsverweis {fungsi yang terdiferensialkan secara kontinu}{}{}
dan

\mavergleichskette
{\vergleichskette
{ Y
}
{ = }{ h^{-1}(c)
}
{  }{
}
{    }{
}
{   }{
}
}
{}{}{}
adalah
\definitionsverweis {serat}{}{}
di atas

\mavergleichskette
{\vergleichskette
{ c
}
{ \in }{ \R
}
{  }{
}
{    }{
}
{   }{
}
}
{}{}{,}
dengan $h$
\definitionsverweis {reguler}{}{}
di setiap titik $Y$. Misalkan $Y$
\definitionsverweis {terhubung}{}{.}}
\faktfolgerung {Maka terdapat tepat dua
\definitionsverweis {orientasi}{}{}
pada $Y$.}
\faktzusatz {}
\faktzusatz {}}{\faktsituation {Misalkan

\mavergleichskette
{\vergleichskette
{ U
}
{ \subseteq }{ \R^n
}
{  }{
}
{    }{
}
{   }{
}
}
{}{}{}
\definitionsverweis {terbuka}{}{}
dan misalkan $f_1 , \ldots , f_n, g_1 , \ldots , g_n$ adalah
\definitionsverweis {fungsi-fungsi yang dua kali terdiferensialkan secara kontinu}{}{}
pada $U$ dengan operator-operator diferensial terkait

\mavergleichskette
{\vergleichskette
{ D
}
{ = }{ \sum_{i = 1}^n f_i \partial_i
}
{  }{
}
{    }{
}
{   }{
}
}
{}{}{}
dan

\mavergleichskette
{\vergleichskette
{ E
}
{ = }{ \sum_{j = 1}^n g_j \partial_j
}
{  }{
}
{    }{
}
{   }{
}
}
{}{}{}}
\faktfolgerung {Maka berlaku

\mavergleichskettedisp
{\vergleichskette
{ D \circ E -E \circ D
}
{ =} { \sum_{j = 1}^n \left(  \sum_{i = 1}^n \left(  f_i  \partial_i g_j - g_i \partial_i f_j  \right)  \right)  \partial_j
}
{ } {
}
{ } {
}
{ } {
}
}
{}{}{.}
Secara khusus, ekspresi ini sendiri bersesuaian dengan suatu operator diferensial berorde $1$, sehingga bersesuaian dengan suatu medan vektor.}
\faktzusatz {}
\faktzusatz {}}{Misalkan $M$ adalah suatu
\definitionsverweis {manifold diferensiabel berbatas}{}{}
\definitionsverweis {kompak}{}{}
\definitionsverweis {berorientasi}{}{}

\mathl{\partial M}{} dengan
\definitionsverweis {basis topologi terhitung}{}{.}
Maka tidak terdapat
\definitionsverweis {pemetaan yang terdiferensialkan secara kontinu}{}{}
\maabbdisp {\varphi} { M } { \partial M
} {,}
yang
\definitionsverweis {pembatasannya}{}{}
pada $\partial M$ adalah identitas.}

 }



\inputaufgabeklausurloesung
{7}

{

Buktikan
\stichwort {teorema tentang penyelesaian persamaan diferensial biasa pada suatu hipermuka diferensiabel} {}

\mavergleichskette
{\vergleichskette
{ Y
}
{ \subseteq }{ \R^n
}
{  }{
}
{    }{
}
{   }{
}
}
{}{}{.}

}
{

Kita bekerja dengan struktur Euklides pada $\R^n$ dan dengan
\definitionsverweis {gradien}{}{}
$\operatorname{Grad} \,  h$ dari $h$. Menurut asumsi, pada $Y$ gradien ini tidak pernah sama dengan $0$; karena kekontinuan, sifat ini berlaku pula pada suatu lingkungan terbuka dari $Y$. Dengan memperkecil $U$ jika perlu, kita dapat mengasumsikan bahwa gradien tidak memiliki titik nol di seluruh $U$. Pada $U$, tinjau medan vektor baru $G$ yang didefinisikan oleh

\mavergleichskettedisp
{\vergleichskette
{ G(P)
}
{ =} { F(P)-  { \frac{   \left\langle  F(P) ,
\operatorname{Grad} \,  h  (  P )    \right\rangle   }{  \Vert {
\operatorname{Grad} \,  h  (  P ) } \Vert^2   } }
\operatorname{Grad} \,  h  (  P )
}
{ } {
}
{ } {
}
{ } {
}
}
{}{}{}
dengan rumus tersebut. Untuk

\mavergleichskette
{\vergleichskette
{ P
}
{ \in }{ Y
}
{  }{
}
{    }{
}
{   }{
}
}
{}{}{}
berlaku

\mavergleichskette
{\vergleichskette
{ G(P)
}
{ = }{ F(P)
}
{  }{
}
{    }{
}
{   }{
}
}
{}{}{,}
karena pada $Y$ gradien dari $h$ tegak lurus terhadap medan vektor. Selain itu, $G$
\zusatzklammer {berbeda dengan $F$} {} {}
memiliki sifat bahwa untuk setiap titik

\mavergleichskette
{\vergleichskette
{ P
}
{ \in }{ U
}
{  }{
}
{    }{
}
{   }{
}
}
{}{}{}
vektor $G(P)$ tegak lurus terhadap gradien; memang,

\mavergleichskettealignhandlinks
{\vergleichskettealignhandlinks
{ \left\langle  G(P) ,
\operatorname{Grad} \,  h  (  P )   \right\rangle
}
{ =} { \left\langle  F(P)- { \frac{   \left\langle  F(P) ,
\operatorname{Grad} \,  h  (  P )    \right\rangle  }{  \Vert {
\operatorname{Grad} \,  h  (  P ) } \Vert^2  } }
\operatorname{Grad} \,  h  (  P )  ,
\operatorname{Grad} \,  h  (  P )   \right\rangle
}
{ =} { \left\langle  F(P) ,
\operatorname{Grad} \,  h  (  P )   \right\rangle - \left\langle  F(P) ,
\operatorname{Grad} \,  h  (  P )   \right\rangle
}
{ =} { 0
}
{ } {
}
}
{}
{}{.}
Sekarang misalkan
\maabbdisp {v} { I} { \R^n
} {}
adalah, menurut
Satz 56.2 (Analysis (Osnabrück 2021-2023))
, solusi tunggal dari masalah nilai awal untuk $G$ dengan syarat awal

\mavergleichskette
{\vergleichskette
{ v(t_0)
}
{ = }{ P
}
{ \in }{ Y
}
{    }{
}
{   }{
}
}
{}{}{.}
Maka

\mavergleichskettealign
{\vergleichskettealign
{ \left(Dh \circ v \right)_{t}
}
{ =} { \left(Dh  \right)_{v(t)} \circ  \left(D v \right)_{t}
}
{ =} { \left(Dh  \right)_{v(t)}  \left(  v'(t)  \right)
}
{ =} { \left(Dh  \right)_{v(t)}  \left(  G(v(t))  \right)
}
{ =} { \left\langle
\operatorname{Grad} \,  h  ( v(t) )  ,  G(v(t))    \right\rangle
}
}
{
\vergleichskettefortsetzungalign
{ =} { 0
}
{ } {}
{ } {}
{ } {}
}
{}{.}
Oleh karena itu, $h$ konstan pada citra $v( I)$ dan, karena

\mavergleichskette
{\vergleichskette
{ h(v(t_0))
}
{ = }{ c
}
{  }{
}
{    }{
}
{   }{
}
}
{}{}{}
berlaku

\mavergleichskette
{\vergleichskette
{ h(v(t))
}
{ = }{ c
}
{  }{
}
{    }{
}
{   }{
}
}
{}{}{}
untuk setiap

\mavergleichskette
{\vergleichskette
{ t
}
{ \in }{ I
}
{  }{
}
{    }{
}
{   }{
}
}
{}{}{,}
sehingga

\mavergleichskette
{\vergleichskette
{ v(t)
}
{ \in }{ Y
}
{  }{
}
{    }{
}
{   }{
}
}
{}{}{}
untuk setiap

\mavergleichskette
{\vergleichskette
{ t
}
{ \in }{ I
}
{  }{
}
{    }{
}
{   }{
}
}
{}{}{.}

 }



\inputaufgabeklausurloesung
{4}

{

Misalkan
\maabb {h} { I } { V
} {}
sebuah kurva yang terdiferensialkan dua kali dalam
\definitionsverweis {ruang vektor Euklides}{}{}
$V$. Buktikan bahwa jika

\mavergleichskette
{\vergleichskette
{h'(t)
}
{ \neq }{ 0
}
{  }{
}
{    }{
}
{   }{
}
}
{}{}{,}
maka berlaku persamaan

\mavergleichskettedisp
{\vergleichskette
{ \Vert { h'(t)} \Vert'
}
{ =} { { \frac{   \left\langle h'(t) , h^{\prime \prime} (t)  \right\rangle  }{  \left\langle h'(t) , h'(t)  \right\rangle  } } \Vert { h'(t)} \Vert
}
{ } {
}
{ } {
}
{ } {
}
}
{}{}{.}

}
{

Misalkan

\mavergleichskette
{\vergleichskette
{u_1
}
{ = }{ h'(t)
}
{  }{
}
{    }{
}
{   }{
}
}
{}{}{,}
yang kita lengkapi menjadi suatu basis ortogonal
\mathl{u_1,u_2 , \ldots , u_n}{} dari $V$. Misalkan
\mathl{h_1 , \ldots , h_n}{} adalah fungsi-fungsi koordinat dari $h$ terhadap basis ini. Maka

\mavergleichskettealign
{\vergleichskettealign
{ \Vert { h'(t)} \Vert'
}
{ =} { \Vert { \begin{pmatrix} h_1'(t) \\\vdots\\ h_n'(t)   \end{pmatrix} } \Vert'
}
{ =} { \sqrt{  h_1'(t)^2 + \cdots + h_n'(t)^2 }'
}
{ =} { { \frac{   h_1(t) h_1'(t) + \cdots +  h_n(t) h_n'(t) }{  \sqrt{  h_1'(t)^2 + \cdots + h_n'(t)^2 }  } }
}
{ =} { { \frac{   h_1(t) h_1'(t) + \cdots +  h_n(t) h_n'(t) }{   h_1'(t)^2 + \cdots + h_n'(t)^2  } }  \sqrt{  h_1'(t)^2 + \cdots + h_n'(t)^2 }
}
}
{
\vergleichskettefortsetzungalign
{ =} { { \frac{   \left\langle h'(t) , h^{\prime \prime} (t)  \right\rangle   }{ \left\langle h'(t) , h'(t)  \right\rangle  } } \Vert { h'(t)} \Vert
}
{ } {
}
{ } {}
{ } {}
}
{}{.}

 }



\inputaufgabeklausurloesung
{0}

{

}
{  }



\inputaufgabeklausurloesung
{0}

{

}
{  }



\inputaufgabeklausurloesung
{0}

{

}
{  }



\inputaufgabeklausurloesung
{5}

{

Buktikan pernyataan bahwa ekuivalensi tangensial lintasan-lintasan pada suatu manifold di suatu titik $P$ dapat diperiksa menggunakan sebarang peta.

}
{

Untuk suatu
\definitionsverweis {kurva diferensiabel}{}{}
\maabbdisp {\gamma} { I } { M
} {}
dengan

\mavergleichskette
{\vergleichskette
{ 0
}
{ \in }{ I
}
{  }{
}
{    }{
}
{   }{
}
}
{}{}{}
dan

\mavergleichskette
{\vergleichskette
{ \gamma(0)
}
{ = }{ P
}
{  }{
}
{    }{
}
{   }{
}
}
{}{}{}
serta suatu
\definitionsverweis {peta}{}{}
\maabbdisp {\alpha} { U } { V
} {}
\zusatzklammer {dengan

\mavergleichskettek
{\vergleichskettek
{P
}
{ \in }{ U
}
{  }{
}
{    }{
}
{   }{
}
}
{}{}{}
dan

\mavergleichskette
{\vergleichskette
{V
}
{ \subseteq }{\R^n
}
{  }{
}
{    }{
}
{   }{
}
}
{}{}{}} {} {}
ekspresi

\mathdisp {\left(  \alpha \circ \left(  \gamma \vert_{\gamma^{-1} (U)}  \right)  \right)'(0)} {  }

tidak berubah jika kita beralih ke suatu interval terbuka yang lebih kecil

\mavergleichskette
{\vergleichskette
{ 0
}
{ \in }{ I'
}
{ \subseteq }{ I
}
{    }{
}
{   }{
}
}
{}{}{}
dan suatu himpunan terbuka yang lebih kecil

\mavergleichskette
{\vergleichskette
{P
}
{ \in }{ U'
}
{ \subseteq }{ U
}
{    }{
}
{   }{
}
}
{}{}{}
\zusatzklammer {dengan peta terinduksi} {} {.}
Jadi, kita dapat mengasumsikan bahwa
\mathkor {} {\gamma_1} {dan} {\gamma_2} {}
didefinisikan pada interval yang sama dan citranya termuat dalam $U$, serta bahwa pada $U$ terdapat dua peta
\maabbdisp {\alpha_1} { U } { V_1
} {}
dan
\maabbdisp {\alpha_2} { U } { V_2
} {}
tersebut. Dari

\mavergleichskettedisp
{\vergleichskette
{ ( \alpha_1  \circ \gamma_1 )'(0)
}
{ =} { ( \alpha_1  \circ \gamma_2 )'(0)
}
{ } {
}
{ } {
}
{ } {
}
}
{}{}{}
dan
Kettenregel,
dengan menggunakan sifat diferensiabel dari
\definitionsverweis {pemetaan transisi}{}{}

\mathl{\alpha_2 \circ \alpha_1^{-1}}{}, langsung diperoleh

\mavergleichskettealign
{\vergleichskettealign
{ \left(  \alpha_2 \circ \gamma_1  \right)'(0)
}
{ =} { \left(  \left(  \alpha_2 \circ \alpha_1^{-1}  \right) \circ \left(  \alpha_1 \circ \gamma_1  \right)  \right)'(0)
}
{ =} { \left(  D   \left(   \alpha_2 \circ \alpha_1^{-1}   \right)      \right)_{\alpha_1(P)} \left(   \left(  \alpha_1 \circ \gamma_1  \right)'(0)   \right)
}
{ =} { \left(  D   \left(   \alpha_2 \circ \alpha_1^{-1}   \right)      \right)_{\alpha_1(P)} \left(   \left(  \alpha_1 \circ \gamma_2  \right)'(0)   \right)
}
{ =} { \left(  \alpha_2 \circ \gamma_2  \right)'(0)
}
}
{}
{}{.}

 }



\inputaufgabeklausurloesung
{15 (3+3+4+5)}

{

Misalkan

\mavergleichskette
{\vergleichskette
{T
}
{ = }{S^1 \times S^1
}
{  }{
}
{    }{
}
{   }{
}
}
{}{}{}
merupakan torus dan

\mavergleichskette
{\vergleichskette
{S^2
}
{ \subset }{\R^3
}
{  }{
}
{    }{
}
{   }{
}
}
{}{}{}
merupakan
\definitionsverweis {permukaan bola satuan}{}{.}

a) Tunjukkan bahwa
\maabbeledisp {\varphi} {S^1 \times S^1} { S^2
} { \begin{pmatrix}  \alpha \\ \beta   \end{pmatrix} } { { \frac{   1 }{  2 } } \begin{pmatrix}  \sqrt{ 2-2 \sin \alpha   } \cos \beta  \\ \cos \alpha + \sin \beta \cos \alpha \\  1  + \sin \alpha -  \sin \beta + \sin \alpha \sin \beta   \end{pmatrix}
} {}
mendefinisikan suatu pemetaan kontinu.

b) Tunjukkan bahwa $\varphi$ surjektif.

c) Deskripsikan serat-serat $\varphi$.

d) Jelaskan pemetaan $\varphi$ dengan bantuan sebuah sketsa.

}
{

a) Pertama-tama, kita tunjukkan bahwa pemetaan tersebut benar-benar bernilai pada permukaan bola; yakni, kita harus menunjukkan bahwa jumlah kuadrat ketiga entrinya sama dengan $1$. Tanpa memperhitungkan faktor di depan, diperoleh

\mavergleichskettealign
{\vergleichskettealign
{
}
{ =} { 2  \cos^{ 2  }  \beta  -2 \sin \alpha  \cos^{ 2  }  \beta + \cos^{ 2  }  \alpha+ \cos^{ 2  }  \alpha \sin^{ 2  }  \beta +2 \cos^{ 2  }  \alpha\sin  \beta +1 + \sin^{ 2  } \alpha + \sin^{ 2  } \beta   + \sin^{ 2  } \alpha \sin^{ 2  } \beta +2 \sin \alpha  -2 \sin \beta   +2 \sin \alpha \sin \beta   -2 \sin \alpha \sin \beta +  2 \sin^{ 2  } \alpha \sin \beta  -2 \sin \alpha \sin^{ 2  } \beta
}
{ =} { 3 +   \cos^{ 2  }  \beta  -2 \sin \alpha  \cos^{ 2  }  \beta + \cos^{ 2  }  \alpha \sin^{ 2  }  \beta +2 \cos^{ 2  }  \alpha\sin  \beta  + \sin^{ 2  } \alpha \sin^{ 2  } \beta +2 \sin \alpha  -2 \sin \beta   +2 \sin^{ 2  } \alpha \sin \beta   - 2 \sin \alpha \sin^{ 2  } \beta
}
{ =} { 3 +   \cos^{ 2  }  \beta  +2 \sin \alpha   \left(  1 - \cos^{ 2  }  \beta   - \sin^{ 2  } \beta   \right)     + 2 \sin \beta \left(  - 1 + \cos^{ 2  }  \alpha    + \sin^{ 2  } \alpha  \right) + \cos^{ 2  }  \alpha\sin^{ 2  }  \beta  + \sin^{ 2  }  \alpha\sin^{ 2  }  \beta
}
{ =} { 3 +   \cos^{ 2  }  \beta  + \sin^{ 2  }  \beta
}
}
{
\vergleichskettefortsetzungalign
{ =} { 4
}
{ } {}
{ } {}
{ } {}
}
{}{.}

c) Perhitungan serat di atas kutub utara
\mathl{\begin{pmatrix}  0 \\ 0\\  1   \end{pmatrix}}{} menghasilkan syarat-syarat

\mavergleichskettedisp
{\vergleichskette
{(1 - \sin \alpha ) \cos \beta
}
{ =} { 0
}
{ =} { \cos \alpha  (1 + \sin \beta )
}
{ } {
}
{ } {
}
}
{}{}{.}
Oleh karena itu, salah satu faktor pada setiap persamaan harus sama dengan $0$. Untuk

\mavergleichskettedisp
{\vergleichskette
{ \alpha
}
{ =} { { \frac{   1 }{  2 } } \pi
}
{ } {
}
{ } {
}
{ } {
}
}
{}{}{}
kedua syarat dipenuhi untuk sebarang $\beta$. Demikian pula, untuk

\mavergleichskettedisp
{\vergleichskette
{\beta
}
{ =} { { \frac{   3 }{  2 } } \pi
}
{ } {
}
{ } {
}
{ } {
}
}
{}{}{}
kedua syarat dipenuhi untuk sebarang $\alpha$. Dalam kasus-kasus tersebut, berdasarkan

\mavergleichskettedisp
{\vergleichskette
{ \sin  { \frac{   1 }{  2 } } \pi
}
{ =} { 1
}
{ } {
}
{ } {
}
{ } {
}
}
{}{}{}
atau berdasarkan

\mavergleichskettedisp
{\vergleichskette
{ \sin  { \frac{   3 }{  2 } } \pi
}
{ =} { -1
}
{ } {
}
{ } {
}
{ } {
}
}
{}{}{}
komponen ketiga sama dengan $1$. Jika

\mavergleichskette
{\vergleichskette
{\alpha
}
{ \neq }{ { \frac{   1 }{  2 } }  \pi
}
{  }{
}
{    }{
}
{   }{
}
}
{}{}{}
dan

\mavergleichskette
{\vergleichskette
{\beta
}
{ \neq }{ { \frac{   3 }{  2 } } \pi
}
{  }{
}
{    }{
}
{   }{
}
}
{}{}{}
syarat

\mavergleichskettedisp
{\vergleichskette
{ \cos \alpha
}
{ =} { 0
}
{ =} { \cos \beta
}
{ } {
}
{ } {
}
}
{}{}{}
harus dipenuhi. Kemungkinan yang tersisa adalah

\mavergleichskette
{\vergleichskette
{ \alpha
}
{ = }{ { \frac{   3 }{  2 } } \pi
}
{  }{
}
{    }{
}
{   }{
}
}
{}{}{}
dan

\mavergleichskette
{\vergleichskette
{ \beta
}
{ = }{ { \frac{   1 }{  2 } } \pi
}
{  }{
}
{    }{
}
{   }{
}
}
{}{}{,}
tetapi dalam kasus ini komponen ketiga sama dengan $-1$.

d) Sudut $\alpha$ menentukan titik

\mavergleichskettedisp
{\vergleichskette
{P(\alpha)
}
{ =} { \begin{pmatrix}  0 \\ \cos \alpha \\  \sin \alpha   \end{pmatrix}
}
{ } {
}
{ } {
}
{ } {
}
}
{}{}{}
pada lingkaran besar yang diberikan oleh

\mavergleichskette
{\vergleichskette
{ x
}
{ = }{ 0
}
{  }{
}
{    }{
}
{   }{
}
}
{}{}{}
. Kutub utara dan
\mathl{P(\alpha)}{} menentukan titik tengah

\mavergleichskettealign
{\vergleichskettealign
{Z(\alpha)
}
{ \defeq} { \begin{pmatrix}  0 \\ 0\\  1   \end{pmatrix} + { \frac{   1 }{  2 } }  \left(  P(\alpha)   - \begin{pmatrix}  0 \\ 0\\  1   \end{pmatrix}   \right)
}
{ =} { \begin{pmatrix}  0 \\ 0\\  1   \end{pmatrix} + { \frac{   1 }{  2 } }  \begin{pmatrix}  0 \\\cos \alpha \\ \sin \alpha - 1   \end{pmatrix}
}
{ =} { { \frac{   1 }{  2 } }  \begin{pmatrix}  0 \\\cos \alpha \\  1 + \sin \alpha   \end{pmatrix}
}
{ } {
}
}
{}
{}{.}
Titik citra
\mathl{P(\alpha, \beta)}{} ditempatkan pada lingkaran $C$ dengan pusat
\mathl{Z(\alpha)}{} yang melalui kutub utara dan $P(\alpha)$ serta seluruhnya terletak pada permukaan bola. Oleh karena itu, $C$ harus terletak pada bidang yang direntang oleh

\mavergleichskettedisp
{\vergleichskette
{P(\alpha) - \begin{pmatrix}  0 \\ 0\\  1   \end{pmatrix}
}
{ =} { { \frac{   1 }{  2 } } \begin{pmatrix}  0 \\\cos \alpha \\ \sin \alpha - 1   \end{pmatrix}
}
{ } {
}
{ } {
}
{ } {
}
}
{}{}{}
dan
\mathl{\begin{pmatrix}  1 \\ 0\\  0   \end{pmatrix}}{.} Untuk menggunakan parametrisasi trigonometrik dari $C$, diperlukan suatu basis ortonormal. Oleh karena itu, kita gunakan

\mathdisp {O(\alpha) =  { \frac{   1 }{  \sqrt{ 2-2  \sin \alpha   }  } } \begin{pmatrix}  0 \\\cos \alpha \\ \sin \alpha - 1   \end{pmatrix}} { . }

Secara keseluruhan, diperoleh

\mavergleichskettealign
{\vergleichskettealign
{\varphi(\alpha)
}
{ =} { Z (\alpha) + \cos  \beta \begin{pmatrix}  1 \\ 0\\  0   \end{pmatrix}  + \sin  \beta O(\alpha)
}
{ =} {  { \frac{   1 }{  2 } }  \begin{pmatrix}  0 \\\cos \alpha \\  1 + \sin \alpha   \end{pmatrix} + \cos  \beta { \frac{   \sqrt{ 2-2  \sin \alpha     }  }{  2  } } \begin{pmatrix}  1 \\ 0\\  0   \end{pmatrix}  + \sin  \beta { \frac{   \sqrt{ 2-2  \sin \alpha     }  }{  2  } } { \frac{   1 }{  \sqrt{ 2-2  \sin \alpha   }  } } \begin{pmatrix}  0 \\\cos \alpha \\ \sin \alpha - 1   \end{pmatrix}
}
{ =} { { \frac{   1 }{  2 } }       \begin{pmatrix}        \sqrt{ 2-2  \sin \alpha     }  \cos \beta         \\\cos \alpha \sin \beta   \\    1 + \sin \alpha  - \sin \beta  + \sin \alpha \sin \beta   \end{pmatrix}
}
{ } {
}
}
{}
{}{.}

 }



\inputaufgabeklausurloesung
{5}

{

Hitunglah integral lintasan
\mathl{\int_\gamma \omega}{} untuk
\maabbeledisp {\gamma} { [-1,0] } {\R^3
} { t } {(-t^2,t^3-1,t+2)
} {,}
dengan bentuk diferensial berderajat $1$

\mavergleichskettedisp
{\vergleichskette
{ \omega
}
{ =} { x^3dx -yzdy +xz^2 dz
}
{ } {
}
{ } {
}
{ } {
}
}
{}{}{}
pada $\R^3$.

}
{

Diperoleh

\mavergleichskettealign
{\vergleichskettealign
{ \gamma^* \omega
}
{ =} { (-t^2)^3 d(-t^2) - (t^3-1) (t+2)d (t^3-1)  - t^2(t+2)^2d(t+2)
}
{ =} { 2 t^7 dt     - 3t^2 (t^4 +2t^3 -t -2) dt - t^2(  t^2  +4 t +4) dt
}
{ =} {  (2t^7 -3t^6 -6t^5   +3t^3 +6t^2 -t^4  -4t^3 -4t^2       )dt
}
{ =} { (2t^7 -3t^6 -6t^5   -t^4  -t^3 +2 t^2       )dt
}
}
{}
{}{.}
Oleh karena itu,

\mavergleichskettealign
{\vergleichskettealign
{ \int_\gamma \omega
}
{ =} { \int_{-1}^0    (2t^7 -3t^6 -6t^5   -t^4  -t^3 +2 t^2       )dt
}
{ =} { ({ \frac{   1 }{  4 } }  t^8 - { \frac{   3 }{  7 } }  t^7 - t^6- { \frac{   1 }{  5 } } t^5 - { \frac{   1 }{  4 } } t^4 + { \frac{   2 }{  3 } } t^3) \vert_{  -1 }^{  0 }
}
{ =} { -( { \frac{   1 }{  4 } }  + { \frac{   3 }{  7 } }  -1 + { \frac{   1 }{  5 } }  - { \frac{   1 }{  4 } }  - { \frac{   2 }{  3 } } )
}
{ =} { - { \frac{   3 }{  7 } }  +1 - { \frac{   1 }{  5 } } + { \frac{   2 }{  3 } }
}
}
{
\vergleichskettefortsetzungalign
{ =} { { \frac{   -45 +105-21+70 }{  105 } }
}
{ =} { { \frac{   109 }{  105 } }
}
{ } {}
{ } {}
}
{}{.}

 }



\inputaufgabeklausurloesung
{0}

{

}
{  }



\inputaufgabeklausurloesung
{3 (2+1)}

{

Misalkan
\maabbeledisp {f} {\R} {\R
} { t } {f(t)
} {,}
diberikan oleh

\mavergleichskettedisp
{\vergleichskette
{ f(t)
}
{ =} { { \frac{   \sin^{ 3  } (t^4)  }{  1+t^2 } }
}
{ } {
}
{ } {
}
{ } {
}
}
{}{}{}

a) Hitung turunan eksterior dari $f$.

b) Hitung turunan eksterior dari $fdt$.

}
{

a) Berlaku
\mathl{df=f'dt}{,} dan menurut aturan hasil bagi diperoleh

\mavergleichskettealign
{\vergleichskettealign
{f'
}
{ =} { { \frac{   (1+t^2)  3\sin^{ 2  } (t^4) \cos (t^4) 4 t^3  - 2t \sin^{ 3  } (t^4)  }{ (1+t^2)^2 } }
}
{ =} { { \frac{   12 (t^3+t^5) \sin^{ 2  } (t^4) \cos (t^4)   - 2t \sin^{ 3  } (t^4)  }{  1+2t^2+t^4 } }
}
{ } {
}
{ } {
}
}
{}
{}{.}

b) Turunan eksterior dari $fdt$ adalah
\mathl{f'dt \wedge dt =0}{.}

 }



\inputaufgabeklausurloesung
{3}

{

Berikan suatu contoh pemetaan yang terdiferensialkan secara kontinu
\maabbdisp {\varphi} { H } { H
} {}
dari suatu setengah bidang Euklides ke dirinya sendiri, dengan sifat bahwa tepat satu titik batas $H$ dipetakan ke suatu titik batas dan semua titik lainnya dipetakan ke interior setengah bidang tersebut.

}
{

Tinjau setengah bidang positif

\mavergleichskettedisp
{\vergleichskette
{ H
}
{ =} { \left\{ (x,y)  \mid   x \geq 0        \right\}
}
{ } {
}
{ } {
}
{ } {
}
}
{}{}{}
dan pemetaan
\maabbeledisp {\varphi} { H } {H
} {(x,y)} { \left(  x+y^2  , \,   y                   \right)
} {.}
Sebagai pemetaan polinomial, $\varphi$ terdiferensialkan secara kontinu. Titik batas
\mathl{(0,0)}{} dipetakan ke titik batas
\mathl{(0,0)}{.} Pada setiap titik lain dari $H$, berlaku

\mavergleichskette
{\vergleichskette
{ x
}
{ > }{  0
}
{  }{
}
{    }{
}
{   }{
}
}
{}{}{}
atau

\mavergleichskette
{\vergleichskette
{ y
}
{ \neq }{ 0
}
{  }{
}
{    }{
}
{   }{
}
}
{}{}{}
dan oleh karena itu selalu

\mavergleichskettedisp
{\vergleichskette
{ x+y^2
}
{ >} {  0
}
{ } {
}
{ } {
}
{ } {
}
}
{}{}{,}
sehingga komponen pertamanya positif.

 }



\inputaufgabeklausurloesung
{0}

{

}
{  }



\inputaufgabeklausurloesung
{10}

{

Buktikan versi balok dari teorema Stokes.

}
{

Karena kedua ruas persamaan ini linear terhadap $\omega$, kita dapat mengasumsikan bahwa $\omega$ berbentuk

\mavergleichskettedisp
{\vergleichskette
{ \omega
}
{ =} { f dx_2 \wedge \ldots \wedge dx_n
}
{ } {
}
{ } {
}
{ } {
}
}
{}{}{}
dengan suatu lingkungan terbuka dari $Q$ tempat fungsi $f$ yang terdiferensialkan secara kontinu didefinisikan. Integral pada kedua ruas adalah
\definitionsverweis {integral Lebesgue}{}{}
dari fungsi kontinu pada himpunan bagian
\mathkor {} {\R^n} {atau} {\R^{n-1}} {.}
Oleh karena itu, pada kedua ruas kita dapat beralih ke
\definitionsverweis {penutupan topologis}{}{,}
karena hal ini hanya mengubah domain integrasi terkait sebesar suatu himpunan nol dan karenanya tidak mengubah nilai integral.

Tuliskan balok tertutup sebagai

\mavergleichskettedisp
{\vergleichskette
{ \overline{Q}
}
{ =} { [a,b] \times \tilde{Q}
}
{ } {
}
{ } {
}
{ } {
}
}
{}{}{.}
Terapkan
Korollar 14.10 (Differentialgeometrie (Osnabrück 2023))
pada setiap sisi $S$ kecuali
\mathkor {} {a \times \tilde{Q}} {dan} {b \times \tilde{Q}} {,}
sehingga pada sisi-sisi tersebut diperoleh

\mavergleichskettedisp
{\vergleichskette
{
\int_{ S  }   \omega
}
{ =} {
\int_{  S  }  f dx_2 \wedge \ldots \wedge dx_n
}
{ =} {
\int_{ S  }   0
}
{ =} { 0
}
{ } {
}
}
{}{}{,}
karena pada setiap sisi tersebut salah satu variabel
\mathl{x_2 , \ldots , x_n}{} konstan. Berdasarkan
Satzes von Fubini
dan
Hauptsatzes der Infinitesimalrechnung
\zusatzklammer {diterapkan untuk setiap \mathlk{(x_2 , \ldots , x_n)}{} yang tetap} {} {,}
berlaku

\mavergleichskettealignhandlinks
{\vergleichskettealignhandlinks
{
\int_{ Q  }   d\omega
}
{ =} {
\int_{  Q  }   df \wedge  dx_2 \wedge \ldots \wedge dx_n
}
{ =} {
\int_{  Q  }   \left(  \sum_{j = 1}^n { \frac{   \partial   f   }{  \partial   x_j   } } dx_j  \right) \wedge  dx_2 \wedge \ldots \wedge dx_n
}
{ =} {
\int_{  Q  }   { \frac{   \partial   f   }{  \partial   x_1   } } dx_1 \wedge  dx_2 \wedge \ldots \wedge dx_n
}
{ =} {
\int_{ \tilde{Q}   }   \left(  \int_{[a,b] } { \frac{   \partial   f   }{  \partial   x_1   } } dx_1  \right)  dx_2 \wedge \ldots \wedge dx_n
}
}
{
\vergleichskettefortsetzungalign
{ =} {
\int_{  \tilde{Q}   }   \left(  f(b,x_2 , \ldots , x_n ) - f(a,x_2 , \ldots , x_n )   \right) dx_2 \wedge \ldots \wedge dx_n
}
{ =} {
\int_{ b \times \tilde{Q}   }   f dx_2 \wedge \ldots \wedge dx_n  -
\int_{  a \times \tilde{Q}   }  f dx_2 \wedge \ldots \wedge dx_n
}
{ =} { \sum_{S \text{ orientierte Seite von } Q}
\int_{  S  }  f dx_2 \wedge \ldots \wedge dx_n
}
{ =} {
\int_{  \partial Q  }   f dx_2 \wedge \ldots \wedge dx_n
}
}
{
\vergleichskettefortsetzungalign
{ =} {
\int_{  \partial Q  }   \omega
}
{ } {}
{ } {}
{ } {}
}{.}

 }



\inputaufgabeklausurloesung
{4}

{

Untuk kurva
\maabbeledisp {\psi} {\R} { {\mathbb H}
} { t } { \left(  t   , \,  e^t                   \right)
} {,}
yang bernilai dalam
\definitionsverweis {setengah bidang Poincaré}{}{}
${\mathbb H}$, tentukan
\definitionsverweis {simbol-simbol Christoffel}{}{}
untuk
\definitionsverweis {koneksi tarikan balik}{}{}
\zusatzklammer {dari
\definitionsverweis {koneksi Levi-Civita}{}{}
pada ${\mathbb H}$} {} {}
di atas $\R$.

}
{

Satu-satunya turunan standar pada $\R$ adalah $\partial$, sehingga indeks bawah pada simbol-simbol Christoffel kita abaikan. Menurut
Fakt *****  (2)
untuk

\mavergleichskette
{\vergleichskette
{ \ell,k
}
{ = }{ 1,2
}
{  }{
}
{    }{
}
{   }{
}
}
{}{}{}

\mavergleichskettedisp
{\vergleichskette
{ \tilde{\Gamma}_\ell^k (t)
}
{ =} { \psi_1'(t) \Gamma_{1 \ell}^k (\psi(t)) + \psi'_2(t) \Gamma_{2 \ell}^k (\psi(t))
}
{ =} { \Gamma_{1 \ell}^k (t,e^t)   + e^t \Gamma_{2 \ell}^k (t,e^t)
}
{ } {}
{ } {}
}
{}{}{.}
Dengan
Beispiel *****
diperoleh

\mavergleichskettedisp
{\vergleichskette
{ \tilde{\Gamma}_1^1 (t)
}
{ =} { \Gamma_{1 1}^1 (t,e^t) + e^t \Gamma_{2 1}^1 (t,e^t)
}
{ =} {  -e^t e^{-t}
}
{ =} { -1
}
{ } {
}
}
{}{}{,}

\mavergleichskettedisp
{\vergleichskette
{ \tilde{\Gamma}_1^2 (t)
}
{ =} { \Gamma_{1 1}^2 (t,e^t) + e^t \Gamma_{2 1}^2 (t,e^t)
}
{ =} { e^{-t}
}
{ } {
}
{ } {
}
}
{}{}{,}

\mavergleichskettedisp
{\vergleichskette
{ \tilde{\Gamma}_2^1 (t)
}
{ =} { \Gamma_{1 2}^1 (t,e^t) + e^t \Gamma_{2 2}^1 (t,e^t)
}
{ =} {  -e^{-t}
}
{ } {
}
{ } {
}
}
{}{}{}
dan

\mavergleichskettedisp
{\vergleichskette
{ \tilde{\Gamma}_2^2 (t)
}
{ =} { \Gamma_{1 2}^2 (t,e^t) + e^t \Gamma_{2 2}^2 (t,e^t)
}
{ =} { - e^t e^{-t}
}
{ =} { -1
}
{ } {
}
}
{}{}{.}

 }
\endgroup

\chapter{Formulir Solusi Resmi Ujian 7}
\noindent\textit{Solusi yang disediakan oleh sumber resmi; kekosongan sumber dipertahankan.}
\begingroup
\renewcommand{\theHfakt}{o011.exam.official.07.\arabic{fakt}}
%Data institusi

%\input{Dozentdaten}

%\renewcommand{\fachbereich}{Fachbereich}

%\renewcommand{\dozent}{Prof. Dr. . }

%Data ujian

\renewcommand{\klausurgebiet}{ }

\renewcommand{\klausurtyp}{ }

\renewcommand{\klausurdatum}{ . 20}

\klausurvorspann {\fachbereich} {\klausurdatum} {\dozent} {\klausurgebiet} {\klausurtyp}

%Data untuk tabel nilai berikut

\renewcommand{\aeins}{  3  }

\renewcommand{\azwei}{  3   }

\renewcommand{\adrei}{  0  }

\renewcommand{\avier}{  4  }

\renewcommand{\afuenf}{  5  }

\renewcommand{\asechs}{  4  }

\renewcommand{\asieben}{  9  }

\renewcommand{\aacht}{  0  }

\renewcommand{\aneun}{  8  }

\renewcommand{\azehn}{  6  }

\renewcommand{\aelf}{  0  }

\renewcommand{\azwoelf}{  3  }

\renewcommand{\adreizehn}{  3   }

\renewcommand{\avierzehn}{  0  }

\renewcommand{\afuenfzehn}{  0  }

\renewcommand{\asechzehn}{  7  }

\renewcommand{\asiebzehn}{  55  }

\renewcommand{\aachtzehn}{    }

\renewcommand{\aneunzehn}{    }

\renewcommand{\azwanzig}{    }

\renewcommand{\aeinundzwanzig}{    }

\renewcommand{\azweiundzwanzig}{    }

\renewcommand{\adreiundzwanzig}{    }

\renewcommand{\avierundzwanzig}{    }

\renewcommand{\afuenfundzwanzig}{    }

\renewcommand{\asechsundzwanzig}{    }

\punktetabellesechzehn

\klausurnote

\newpage

\setcounter{section}{K}



\inputaufgabeklausurloesung
{3}

{

Definisikan istilah-istilah berikut
\zusatzklammer {yang dicetak miring} {} {.}
\aufzaehlungsechs{
\stichwort {Kelengkungan Gauss} {}
dari suatu permukaan diferensiabel

\mavergleichskette
{\vergleichskette
{ Y
}
{ \subseteq }{ \R^3
}
{  }{
}
{    }{
}
{   }{
}
}
{}{}{}
di suatu titik

\mavergleichskette
{\vergleichskette
{ P
}
{ \in }{ Y
}
{  }{
}
{    }{
}
{   }{
}
}
{}{}{.}

}{Suatu
\stichwort {manifold topologis} {}
berdimensi $n$.

}{Suatu
\stichwort {medan vektor tak bergantung waktu} {}
pada manifold diferensiabel $M$.

}{Suatu
\stichwort {bentuk diferensial berderajat $k$} {}
pada manifold diferensiabel $M$.

}{Suatu \stichwort {manifold Riemann} {.}

}{Suatu koneksi linear yang
\stichwort {bebas torsi} {}
pada bundel tangen manifold Riemann $M$.
}

}
{

\aufzaehlungsechs{Hasil kali
\mathl{\kappa_1 \cdot \kappa_2}{} dari kedua
\definitionsverweis {kelengkungan utama}{}{}
di $P$ disebut
\definitionswort {kelengkungan Gauss}{}
permukaan $Y$ di $P$.
}{Suatu
\definitionsverweis {ruang Hausdorff}{}{}
\definitionsverweis {topologis}{}{}
$M$ disebut manifold topologis berdimensi $n$ jika terdapat suatu
\definitionsverweis {penutup terbuka}{}{}

\mavergleichskette
{\vergleichskette
{ M
}
{ = }{  \bigcup_{i \in I} U_i
}
{  }{
}
{    }{
}
{   }{
}
}
{}{}{}
sedemikian sehingga setiap $U_i$
\definitionsverweis {homeomorfik}{}{}
dengan suatu
\definitionsverweis {subhimpunan terbuka}{}{}
dari $\R^n$.
}{Suatu pemetaan
\maabbdisp {F} {M} {TM
} {}
dengan sifat

\mavergleichskette
{\vergleichskette
{ F(P)
}
{ \in }{  T_PM
}
{  }{
}
{    }{
}
{   }{
}
}
{}{}{}
untuk setiap titik

\mavergleichskette
{\vergleichskette
{ P
}
{ \in }{ M
}
{  }{
}
{    }{
}
{   }{
}
}
{}{}{}
disebut medan vektor
\zusatzklammer {tak bergantung waktu} {} {}.
}{Suatu
$k$-bentuk diferensial
adalah suatu
\definitionsverweis {penampang}{}{}
dalam
\definitionsverweis {hasil kali baji}{}{}
ke-$k$ dari
\definitionsverweis {bundel kotangen}{}{.}
}{Suatu
\definitionsverweis {manifold diferensiabel}{}{}
$M$ disebut manifold Riemann jika pada setiap
\definitionsverweis {ruang tangen}{}{}

\mathbed {T_PM} {}
 {P \in M} {}
 {} {} {} {,}
diberikan suatu
\definitionsverweis {hasil kali skalar}{}{}

\mathl{\left\langle - , -  \right\rangle_P}{} sedemikian sehingga, untuk setiap peta
\maabbdisp {\alpha} { U } { V
} {}
dengan
\mathl{V\subseteq \R^n}{}, fungsi-fungsi
\zusatzklammer {untuk \mathlk{1 \leq i,j \leq n}{}} {} {}
\maabbeledisp {g_{ij}} { V } {\R
} { Q } {g_{ij}(Q) =  \left\langle  T(\alpha^{-1})(e_i)  ,  T(\alpha^{-1}) (e_j)   \right\rangle_{ \alpha^{-1}(Q) }
} {,}
bersifat $C^1$-\definitionsverweis {terdiferensialkan}{}{}.
}{Koneksi tersebut disebut
bebas torsi
jika

\mavergleichskettedisp
{\vergleichskette
{ \nabla_V W- \nabla_W V
}
{ =} { [V,W]
}
{ } {
}
{ } {
}
{ } {
}
}
{}{}{}
berlaku untuk sebarang medan vektor $V,W$.
}

 }



\inputaufgabeklausurloesung
{3}

{

Nyatakan teorema-teorema berikut.
\aufzaehlungdrei{
\stichwort {Teorema eksistensi transpor paralel pada hiperpermukaan} {.}}{Teorema tentang deskripsi lokal bentuk diferensial.}{
\stichwort {Versi balok} {}
dari teorema Stokes.}

}
{

\aufzaehlungdrei{\faktsituation {Misalkan

\mavergleichskette
{\vergleichskette
{ Y
}
{ \subseteq }{ W
}
{ \subseteq }{\R^n
}
{    }{
}
{   }{
}
}
{}{}{,}
$W$
\definitionsverweis {terbuka}{}{,}
suatu
\definitionsverweis {hiperpermukaan diferensiabel}{}{}
dan misalkan
\maabb {\gamma} { I } { Y
} {}
suatu
\definitionsverweis {kurva diferensiabel}{}{.}
Misalkan

\mavergleichskette
{\vergleichskette
{ t_0
}
{ \in }{ I
}
{  }{
}
{    }{
}
{   }{
}
}
{}{}{,}

\mavergleichskette
{\vergleichskette
{ \gamma(t_0)
}
{ = }{ P
}
{  }{
}
{    }{
}
{   }{
}
}
{}{}{}
dan misalkan

\mavergleichskette
{\vergleichskette
{ v
}
{ \in }{ T_PY
}
{  }{
}
{    }{
}
{   }{
}
}
{}{}{}
suatu
\definitionsverweis {vektor tangen}{}{.}}
\faktfolgerung {Maka terdapat suatu
\definitionsverweis {medan vektor tangensial}{}{}
$F$ sepanjang $\gamma$ yang ditentukan secara unik,
\definitionsverweis {paralel}{}{}
dan memenuhi

\mavergleichskettedisp
{\vergleichskette
{ F(t_0)
}
{ =} { v
}
{ } {
}
{ } {
}
{ } {
}
}
{}{}{.}}
\faktzusatz {}
\faktzusatz {}}{Misalkan $M$ suatu
\definitionsverweis {manifold diferensiabel}{}{}
dan
\mathl{U \subseteq M}{} suatu
\definitionsverweis {subhimpunan terbuka}{}{}
dengan peta
\maabbdisp {\alpha} { U } { V
} {}
dan
\mathl{V \subseteq \R^n}{} terbuka. Misalkan
\maabbdisp {x_j} { U } {\R
} {} adalah fungsi-fungsi koordinat terkait,
\mathl{1 \leq j \leq n}{.} Maka setiap bentuk pada $U$ yang merupakan $k$-\definitionsverweis {bentuk diferensial}{}{}

\mathl{\omega \in
{ \mathcal E }^{ k } (  U   )}{} dapat ditulis secara unik sebagai

\mavergleichskettedisp
{\vergleichskette
{  \omega
}
{ =} { \sum_{J,\,  { \# \left(   J  \right) } = k}   f_J dx_J
}
{ } {
}
{ } {
}
{ } {
}
}
{}{}{}
dengan fungsi-fungsi yang ditentukan secara unik
\maabbdisp {f_J} { U } {\R
} {.}}{\faktsituation {Misalkan

\mavergleichskette
{\vergleichskette
{ Q
}
{ \subseteq }{ \R^n
}
{  }{
}
{    }{
}
{   }{
}
}
{}{}{}
suatu balok $n$-dimensional yang sejajar sumbu
\zusatzklammer {dengan sisi tetapi tanpa rusuk} {} {}
dengan
\definitionsverweis {batas}{}{}

\mathl{\partial Q}{} dan $\omega$ suatu
\definitionsverweis {bentuk diferensial}{}{}
yang didefinisikan pada $Q$, berderajat $(n-1)$, dan
\definitionsverweis {terdiferensialkan secara kontinu}{}{}
.}
\faktfolgerung {Maka

\mavergleichskettedisp
{\vergleichskette
{
\int_{  Q  }  d \omega
}
{ =} {
\int_{  \partial Q  }   \omega
}
{ } {
}
{ } {
}
{ } {
}
}
{}{}{.}}
\faktzusatz {}
\faktzusatz {}}

 }



\inputaufgabeklausurloesung
{0}

{

}
{  }



\inputaufgabeklausurloesung
{4 (1+3)}

{

Misalkan $M$ elips yang diberikan oleh

\mavergleichskettedisp
{\vergleichskette
{ x^2+3y^2
}
{ =} { 2
}
{ } {
}
{ } {
}
{ } {
}
}
{}{}{}
dalam $\R^2$,

\mavergleichskette
{\vergleichskette
{ P
}
{ = }{ \begin{pmatrix}  1  \\ { \frac{   1  }{  \sqrt{3}  } }    \end{pmatrix}
}
{  }{
}
{    }{
}
{   }{
}
}
{}{}{}
dan

\mavergleichskette
{\vergleichskette
{ Q
}
{ = }{ \begin{pmatrix}  -1 \\ { \frac{   1  }{  \sqrt{3}  } }    \end{pmatrix}
}
{  }{
}
{    }{
}
{   }{
}
}
{}{}{}
adalah titik-titik pada $M$.
\aufzaehlungdrei{Tunjukkan bahwa

\mavergleichskettedisp
{\vergleichskette
{ v
}
{ =} { \begin{pmatrix}  -  \sqrt{3}  \\ 1    \end{pmatrix}
}
{ } {
}
{ } {
}
{ } {
}
}
{}{}{}
merupakan
\definitionsverweis {vektor tangen}{}{}
di $P$ pada $M$.
}{Tentukan citra $v$ di bawah suatu transpor paralel pada $M$ dari $P$ ke $Q$.
}{
}

}
{

\aufzaehlungzwei {Ruang tangen di titik $\begin{pmatrix}  x  \\ y   \end{pmatrix}$ tegak lurus terhadap gradien
\mathl{\begin{pmatrix}  2x \\ 6y   \end{pmatrix}}{,} sehingga di $P$ ruang tersebut tegak lurus terhadap

\mavergleichskettedisp
{\vergleichskette
{ \begin{pmatrix}  2 \\ { \frac{   6 }{  \sqrt{3}  } }    \end{pmatrix}
}
{ =} { 2 \begin{pmatrix}  1  \\ \sqrt{3}    \end{pmatrix}
}
{ } {
}
{ } {
}
{ } {
}
}
{}{}{.}
Vektor $v$ memenuhi syarat ini.
} {Transpor paralel sepanjang suatu lintasan $\gamma$ pada $M$ dari $P$ ke $Q$ merupakan isometri
\maabb {} {T_PM} {T_QM
} {.}
Ruang tangen di $Q$ tegak lurus terhadap
\mathl{\begin{pmatrix}  - 1  \\ \sqrt{3}    \end{pmatrix}}{} dan direntang oleh
\mathl{\begin{pmatrix}  \sqrt{3}  \\ 1    \end{pmatrix}}{.} Karena normanya tidak berubah, hasilnya hanya mungkin
\mathl{\pm \begin{pmatrix}  \sqrt{3}  \\ 1    \end{pmatrix}}{}. Karena $v$ menunjuk berlawanan arah jarum jam pada $M$, hasilnya juga harus demikian
\zusatzklammer {misalnya berdasarkan
Soal 6.9 (Geometri diferensial (Osnabrück 2023))} {} {,}
sehingga hasilnya adalah
\mathl{- \begin{pmatrix}  \sqrt{3}  \\ 1    \end{pmatrix}}{}.
}

 }



\inputaufgabeklausurloesung
{5}

{

Buktikan teorema tentang orientasi-orientasi pada hiperpermukaan diferensiabel.

}
{

Menurut
Lema 2.4 (Geometri diferensial (Osnabrück 2023)),
terdapat suatu medan normal satuan yang merepresentasikan sebuah orientasi, dan negatifnya menghasilkan orientasi lain. Kita namakan keduanya $G$ dan $-G$. Sekarang, misalkan
\maabbdisp {N} { Y } { \R^n
} {}
suatu medan normal satuan kontinu sebarang. Kita tinjau himpunan-himpunan tempat keduanya berimpit

\mavergleichskettedisp
{\vergleichskette
{ Y_1
}
{ =} { \left\{ P \in Y  \mid   N(P) = G(P)         \right\}
}
{ } {
}
{ } {
}
{ } {
}
}
{}{}{}
dan

\mavergleichskettedisp
{\vergleichskette
{ Y_2
}
{ =} { \left\{ P \in Y  \mid   N(P) = - G(P)         \right\}
}
{ } {
}
{ } {
}
{ } {
}
}
{}{}{.}
Karena untuk setiap titik

\mavergleichskette
{\vergleichskette
{ P
}
{ \in }{ Y
}
{  }{
}
{    }{
}
{   }{
}
}
{}{}{}
hanya terdapat dua nilai yang mungkin,

\mavergleichskettedisp
{\vergleichskette
{ N(P)
}
{ =} { \pm G(P)
}
{ } {
}
{ } {
}
{ } {
}
}
{}{}{}
maka $Y$ adalah gabungan saling lepas dari
\mathkor {} {Y_1} {dan} {Y_2} {.}
Sebagai himpunan tempat fungsi-fungsi kontinu berimpit,
\mathkor {} {Y_1} {dan} {Y_2} {}
\definitionsverweis {tertutup}{}{}
\zusatzklammer {dan dengan demikian juga
\definitionsverweis {terbuka}{}{}
dalam $Y$} {} {.}
Karena keterhubungan, diperoleh

\mavergleichskette
{\vergleichskette
{ Y
}
{ = }{ Y_1
}
{  }{
}
{    }{
}
{   }{
}
}
{}{}{}
atau

\mavergleichskette
{\vergleichskette
{ Y
}
{ = }{ Y_2
}
{  }{
}
{    }{
}
{   }{
}
}
{}{}{.}

 }



\inputaufgabeklausurloesung
{4}

{

Misalkan

\mavergleichskette
{\vergleichskette
{ Y
}
{ \subseteq }{ W
}
{ \subseteq }{ \R^3
}
{    }{
}
{   }{
}
}
{}{}{,}
$W$
\definitionsverweis {terbuka}{}{,}
dan suatu
\definitionsverweis {permukaan diferensiabel}{}{}
yang dilengkapi dengan
\definitionsverweis {medan normal satuan}{}{}
$N$. Misalkan
\maabb {\gamma} { I } { Y
} {}
suatu kurva terdiferensialkan secara kontinu dan
\maabbdisp {F} { I } { \R^3
} {}
suatu
\definitionsverweis {medan vektor tangensial}{}{}
diferensiabel sepanjang $\gamma$ yang
\definitionsverweis {paralel}{}{.}
Tunjukkan bahwa medan vektor
\maabbeledisp {} { I } { \R^3
} { t } { N(\gamma(t)) \times F(t)
} {,}
juga paralel, dengan $\times$ menyatakan
\definitionsverweis {hasil kali silang}{}{}
dalam $\R^3$.

}
{

Kita peroleh

\mavergleichskettealign
{\vergleichskettealign
{ (N (\gamma(t)) \times F(t))'
}
{ =} { (N (\gamma(t)))' \times F(t) +  N (\gamma(t)) \times F'(t)
}
{ =} { \left(  D N   \right)_{ \gamma(t)} \left(   \gamma'(t)   \right)  \times F(t) +  N (\gamma(t)) \times F'(t)
}
{ } {
}
{ } {
}
}
{}
{}{.}
Kita harus menunjukkan bahwa vektor ini selalu tegak lurus terhadap ruang tangen $T_{\gamma(t)}Y$. Karena $F'(t)$ tegak lurus terhadap
ruang tangen, vektor tersebut bergantung linear pada vektor normal $N (\gamma(t))$, sehingga suku kanan sama dengan $0$ menurut
Lema 33.3 (Aljabar linear (Osnabrück 2024-2025))  (3).
Dalam suku kiri, $F(t)$ tangensial menurut asumsi dan $\left(  D N   \right)_{ \gamma(t)} \left(   \gamma'(t)   \right)$ tangensial menurut
Lema 4.2 (Geometri diferensial (Osnabrück 2023)).
Karena itu, menurut
Lema 33.3 (Aljabar linear (Osnabrück 2024-2025))  (6),
suku kiri tegak lurus terhadap ruang tangen, sehingga medan vektor tersebut secara keseluruhan paralel.

 }



\inputaufgabeklausurloesung
{9}

{

Misalkan

\mavergleichskette
{\vergleichskette
{ M
}
{ \neq }{ \emptyset
}
{  }{
}
{    }{
}
{   }{
}
}
{}{}{}
suatu
\definitionsverweis {manifold diferensiabel}{}{}
berdimensi $n$. Tunjukkan bahwa terdapat suatu rantai
\definitionsverweis {submanifold-submanifold tertutup}{}{}

\mavergleichskettedisp
{\vergleichskette
{ M_0
}
{ \subseteq} { M_1
}
{ \subseteq} { M_2
}
{ \subseteq  \ldots \subseteq} { M_{n-1}
}
{ \subseteq} { M_n
}
}
{
\vergleichskettefortsetzung
{ =} { M
}
{ } {}
{ } {}
{ } {}
}{}{}
sedemikian sehingga submanifold tertutup $M_i$ berdimensi $i$.

}
{

Misalkan $P \in M$ suatu titik dan
\mathl{P \in U}{} suatu lingkungan peta terbuka beserta peta
\maabbdisp {\alpha} { U } {V \subseteq \R^n
} {,}
dengan $V=B(0,3)$ bola terbuka berpusat di $0$ dan berjari-jari $3$, serta
\mathl{\alpha(P)=0}{}. Untuk
\mathl{i=1 , \ldots , n-1}{}, kita tinjau

\mathdisp {B_i = \left\{  x \in V  \mid  (x_1-1)^2+x_2^2 + \cdots + x_n^2 = 1 , \,  x_{i+2}   = 0,\, x_{i+3} = 0 , \ldots , x_n = 0       \right\}} { . }

Jadi, $B_{n-1}$ adalah permukaan bola berpusat di
\mathl{(1,0 , \ldots , 0)}{} dan berjari-jari $1$; $B_{n-2}$ di dalamnya didefinisikan oleh
\mathl{x_n=0}{} dan merupakan \anfuehrung{ekuator}{}, dan seterusnya. Kita memperoleh $B_i$ dari
\mathl{B_{i+1}}{} dengan menetapkan pula
\mathl{x_{i+2}=0}{}. Karena itu, terdapat rantai menurun subhimpunan-subhimpunan tertutup

\mathdisp {B_{n-1} \supseteq B_{n-2} \supseteq  \ldots \supseteq B_1 \supseteq B_0} {  }

\zusatzklammer {$B_0$ terdiri atas kedua titik
\mathl{(\pm1,0 , \ldots , 0 )}{}} {} {.}
Kita pandang $B_i$ sebagai serat di atas titik nol dari pemetaan
\maabbeledisp {\varphi_i} { V } {\R^{n-i}
} {(x_1 , \ldots , x_n) } {  ((x_1-1)^2+x_2^2 + \cdots + x_n^2 - 1 , x_{i+2} , \ldots , x_n)
} {,}.
Matriks Jacobi pemetaan ini adalah

\mathdisp {\begin{pmatrix}  2x_1-2 &  2x_2  &  \ldots  &  2x_{i+1} &  2x_{i+2} &  \ldots &  2x_{n-1} &  2x_n
 \\  0 &  0 &  \ldots &  0 &  1 &  0 &  \ldots &  0
 \\ \vdots  & \vdots & \vdots &  \ddots &  \ddots &  \ddots  &  \ddots  & \vdots \\  0 &  0 &  \ldots &  \ldots &  \ldots &  0 &  1 &  0 \\  0 &  0 &  \ldots &  \ldots &  \ldots  &  \ldots &  0 &  1 \end{pmatrix}} { . }

Rank matriks ini lebih kecil dari $n-i$ hanya jika
\mathl{x_1=1, \, x_2 = \cdots =x_{i+1}=0}{}, dan titik seperti itu tidak terletak pada $B_i$. Artinya, pemetaan $\varphi_i$ reguler pada serat $B_i$, sehingga berdasarkan teorema pemetaan implisit, $B_i$ merupakan submanifold tertutup dari $V$ yang berdimensi
\mathl{i}{}.

Sekarang kita tetapkan
\mathl{M_i= \alpha^{-1}(B_i)}{} untuk
\mathl{i=1 , \ldots , n-1}{} dan
\mathl{M_n= M}{.} Karena $B_i$ kompak, $M_i$ juga merupakan subhimpunan tertutup dalam $M$. Karena syarat bagi suatu manifold tertutup merupakan sifat lokal, himpunan-himpunan ini adalah submanifold-submanifold tertutup dari $M$.

 }



\inputaufgabeklausurloesung
{0}

{

}
{  }



\inputaufgabeklausurloesung
{8}

{

Misalkan
\maabbdisp {\psi} {\R^n } {\R
} {}
suatu
\definitionsverweis {fungsi diferensiabel}{}{}
dan

\mavergleichskettedisphandlinks
{\vergleichskettedisphandlinks
{M
}
{ =} { \left\{  \left(  x_1   , \,   , \ldots ,  , \,   x_n , \,   \psi(x_1 , \ldots , x_n)                \right)   \mid  (x_1 , \ldots , x_n) \in \R^n         \right\}
}
{ \subseteq} { \R^n \times \R
}
{ } {
}
{ } {
}
}
{}{}{}
merupakan
\definitionsverweis {grafik}{}{}
dari $\psi$. Tunjukkan bahwa $M$ suatu
\definitionsverweis {manifold Riemann}{}{}
yang
\definitionsverweis {berorientasi}{}{}
dan bahwa untuk
\definitionsverweis {bentuk volume kanonik}{}{}
$\omega$ pada $M$ berlaku rumus

\mavergleichskettedisp
{\vergleichskette
{ (\operatorname{Id} \times \psi)^*(\omega)
}
{ =} { \left(  1+ \sum_{i = 1}^n \left(  { \frac{   \partial   \psi  }{  \partial   x_i   } }  \right)^2  \right)^{1/2} dx_1 \wedge \ldots \wedge dx_n
}
{ } {
}
{ } {
}
{ } {
}
}
{}{}{.}

}
{

Misalkan

\mavergleichskette
{\vergleichskette
{ W
}
{ = }{ \R^n
}
{  }{
}
{    }{
}
{   }{
}
}
{}{}{.}
Pemetaan
\maabbeledisp {\varphi = \operatorname{id} \times \psi} { W } { W \times \R
} { x } {(x, \psi(x))
} {,}
merupakan suatu
\definitionsverweis {difeomorfisme}{}{}
antara $W$ dan grafik $M$. Grafik tersebut adalah
\definitionsverweis {submanifold tertutup}{}{}
dari
\mathl{W \times \R}{} dan karena itu membawa
\definitionsverweis {struktur Riemann terinduksi}{}{}
serta
\zusatzklammer {karena orientasi $W$ berpindah ke $M$} {} {}
suatu
\definitionsverweis {bentuk volume kanonik}{}{}
$\omega$. Pada situasi ini kita dapat menerapkan
Teorema 16.8 (Geometri diferensial (Osnabrück 2023)).
\definitionsverweis {Turunan parsial}{}{}
dari $\varphi$ terhadap variabel ke-$i$ adalah

\mavergleichskettedisp
{\vergleichskette
{ { \frac{   \partial  \varphi  }{  \partial   x_i   } }
}
{ =} { \begin{pmatrix}  0  \\\vdots\\  0\\ 1\\  0\\ \vdots\\  0\\  { \frac{   \partial   \psi  }{  \partial   x_i   } }   \end{pmatrix}
}
{ } {
}
{ } {
}
{ } {
}
}
{}{}{.}
Misalkan

\mavergleichskette
{\vergleichskette
{Q
}
{ \in }{ W
}
{  }{
}
{    }{
}
{   }{
}
}
{}{}{}
suatu titik yang kita substitusikan ke dalam fungsi-fungsi berikut, sehingga kita menghitung dengan bilangan real di seluruh bagian. Hasil kali skalar yang membentuk entri-entri
\mathl{b_{ij}}{} dari matriks $B$
\zusatzklammer {yang determinannya harus kita hitung akarnya} {} {,}
adalah

\mavergleichskettedisp
{\vergleichskette
{ b_{ij}
}
{ =} { \left\langle  { \frac{   \partial  \varphi  }{  \partial   x_i   } }(Q)  ,  { \frac{   \partial  \varphi  }{  \partial   x_j   } }(Q)   \right\rangle
}
{ } {}
{ } {}
{ } {}
}
{}{}{}
Kita tulis

\mavergleichskette
{\vergleichskette
{B
}
{ = }{ E_n+A
}
{  }{
}
{    }{
}
{   }{
}
}
{}{}{}
dengan

\mavergleichskette
{\vergleichskette
{ A
}
{ = }{ \left(  a_{ij}  \right)_{1 \leq i,j \leq n}
}
{  }{
}
{    }{
}
{   }{
}
}
{}{}{.}
Dengan

\mavergleichskette
{\vergleichskette
{ c_i
}
{ = }{ { \frac{   \partial   \psi  }{  \partial   x_i   } } ( Q)
}
{  }{
}
{    }{
}
{   }{
}
}
{}{}{}
kita dapat menulis

\mavergleichskette
{\vergleichskette
{ a_{ij}
}
{ = }{ c_ic_j
}
{  }{
}
{    }{
}
{   }{
}
}
{}{}{}
dan keseluruhan matriks $A$ sebagai

\mavergleichskettedisp
{\vergleichskette
{A
}
{ =} { \begin{pmatrix} c_1  \\\vdots\\ c_n   \end{pmatrix} \left( c_1   , \,   \ldots  , \,  c_n                 \right)
}
{ } {
}
{ } {
}
{ } {
}
}
{}{}{}
Karena itu, $A$ merepresentasikan suatu pemetaan linear dari $\R^n$ ke $\R^n$ yang melalui $\R$
\definitionsverweis {terfaktorkan}{}{}
dan dengan demikian memiliki
\definitionsverweis {kernel}{}{,}
yang sekurang-kurangnya berdimensi $(n-1)$. Kita namakan kernel ini $K$. Jika dimensinya $n$, maka

\mavergleichskette
{\vergleichskette
{A
}
{ = }{ 0
}
{  }{
}
{    }{
}
{   }{
}
}
{}{}{}
dan $B$ adalah identitas, sehingga pernyataannya benar. Jadi, misalkan

\mavergleichskette
{\vergleichskette
{A
}
{ \neq }{ 0
}
{  }{
}
{    }{
}
{   }{
}
}
{}{}{.}
Maka

\mavergleichskette
{\vergleichskette
{v
}
{ = }{ \begin{pmatrix} c_1  \\\vdots\\ c_n   \end{pmatrix}
}
{  }{
}
{    }{
}
{   }{
}
}
{}{}{}
adalah
\definitionsverweis {vektor eigen}{}{}
dari $A$ dengan
\definitionsverweis {nilai eigen}{}{}

\mavergleichskette
{\vergleichskette
{ c_1^2 + \cdots + c_n^2
}
{ \neq }{ 0
}
{  }{
}
{    }{
}
{   }{
}
}
{}{}{.}
Vektor ini merupakan vektor eigen dari $B$ dengan nilai eigen
\mathl{1+c_1^2 + \cdots + c_n^2}{}, sedangkan $K$ membentuk
\mathl{(n-1)}{-}dimensional
\definitionsverweis {ruang eigen}{}{}
untuk $B$ dengan nilai eigen $1$. Secara keseluruhan, $B$
\definitionsverweis {dapat didiagonalkan}{}{}
dan
\definitionsverweis {determinan}{}{}
matriks tersebut merupakan hasil kali nilai-nilai eigennya, yaitu
\mathl{1+c_1^2 + \cdots + c_n^2}{.}

 }



\inputaufgabeklausurloesung
{6}

{

Hitung luas permukaan bola satuan.

}
{  }



\inputaufgabeklausurloesung
{0}

{

}
{  }



\inputaufgabeklausurloesung
{3}

{

Hitung
\definitionsverweis {turunan eksterior}{}{}
$d \omega$ dari
$1$-\definitionsverweis {bentuk diferensial}{}{}

\mavergleichskettedisp
{\vergleichskette
{ \omega
}
{ =} { { \frac{   x^2 }{  y } } dx- { \frac{   x  }{  y^2 } } dy
}
{ } {
}
{ } {
}
{ } {
}
}
{}{}{}
pada

\mavergleichskette
{\vergleichskette
{ U
}
{ = }{ \left\{ (x,y) \in \R^2  \mid   y \neq 0         \right\}
}
{  }{
}
{    }{
}
{   }{
}
}
{}{}{.}

}
{

Kita peroleh

\mavergleichskettealign
{\vergleichskettealign
{d \omega
}
{ =} { d { \frac{   x^2 }{  y } } \wedge dx - d { \frac{   x  }{  y^2 } } \wedge dy
}
{ =} { - { \frac{   x^2 }{  y^2 } } dy  \wedge dx  -   { \frac{   1 }{  y^2 } }  dx \wedge dy
}
{ =} { { \frac{   x^2 -1 }{  y^2 } }  dx \wedge dy
}
{ } {
}
}
{}
{}{.}

 }



\inputaufgabeklausurloesung
{3}

{

Berikan contoh suatu
\definitionsverweis {manifold berbatas}{}{}
yang
\definitionsverweis {terhubung}{}{,}
dan batasnya terdiri atas
\zusatzklammer {gabungan saling lepas dari} {} {}
tak hingga banyak lingkaran $S^1$.

}
{

Kita tinjau

\mavergleichskettedisp
{\vergleichskette
{ M
}
{ =} { \R^2 \setminus \biguplus_{n \in \Z} U \left(  (n,0), { \frac{   1 }{  3 } }   \right)
}
{ } {
}
{ } {
}
{ } {
}
}
{}{}{.}
Jadi, tak hingga banyak bola terbuka saling lepas dikeluarkan dari bidang real sepanjang sumbu $x$. Setiap bola menghasilkan satu komponen batas dari $M$ yang difeomorfik dengan lingkaran $S^1$. Jelas bahwa $M$ terhubung-lintasan.

 }



\inputaufgabeklausurloesung
{0}

{

}
{  }



\inputaufgabeklausurloesung
{0}

{

}
{  }



\inputaufgabeklausurloesung
{7}

{

Misalkan $M$ suatu
\definitionsverweis {manifold Riemann}{}{}
dan diberikan suatu
\definitionsverweis {koneksi linear}{}{}
pada
\definitionsverweis {bundel tangen}{}{}
$TM$ yang
\definitionsverweis {metrik}{}{}
dan
\definitionsverweis {bebas torsi}{}{.}
Tunjukkan bahwa untuk medan-medan vektor $X,Y,Z$ berlaku rumus

\mavergleichskettedisp
{\vergleichskette
{ \left\langle \nabla_X Y , Z  \right\rangle
}
{ =} { { \frac{   1  }{  2 } }  \left(  D_X \left\langle  Y  ,  Z  \right\rangle +  D_Y \left\langle  Z  ,  X  \right\rangle -  D_Z \left\langle  X  ,  Y  \right\rangle  - \left\langle  Y  , [X,Z]   \right\rangle - \left\langle  Z  , [Y,X]   \right\rangle + \left\langle  X  , [Z,Y]   \right\rangle   \right)
}
{ } {
}
{ } {
}
{ } {
}
}
{}{}{.}

}
{

Karena koneksi bebas torsi, berlaku

\mavergleichskettedisp
{\vergleichskette
{ \nabla_VW- \nabla_WV
}
{ =} { [V,W]
}
{ } {
}
{ } {
}
{ } {
}
}
{}{}{,}

\mavergleichskettedisp
{\vergleichskette
{ \nabla_VZ- \nabla_ZV
}
{ =} { [V,Z]
}
{ } {
}
{ } {
}
{ } {
}
}
{}{}{}
dan

\mavergleichskettedisp
{\vergleichskette
{ \nabla_WZ- \nabla_ZW
}
{ =} { [W,Z]
}
{ } {
}
{ } {
}
{ } {
}
}
{}{}{,}
dengan identitas pertama juga dapat ditulis sebagai

\mavergleichskettedisp
{\vergleichskette
{ \nabla_VW + \nabla_WV
}
{ =} { 2 \nabla_VW -[V,W]
}
{ } {
}
{ } {
}
{ } {
}
}
{}{}{}.
Karena koneksi metrik, berlaku identitas-identitas

\mavergleichskettedisp
{\vergleichskette
{ D_V \left\langle  W  ,  Z  \right\rangle
}
{ =} { \left\langle \nabla_V W , Z  \right\rangle  + \left\langle  W  ,  \nabla_V Z  \right\rangle
}
{ } {
}
{ } {
}
{ } {
}
}
{}{}{,}

\mavergleichskettedisp
{\vergleichskette
{ D_W \left\langle  Z  ,  V  \right\rangle
}
{ =} { \left\langle \nabla_WZ , V  \right\rangle  + \left\langle  Z  ,  \nabla_WV  \right\rangle
}
{ } {
}
{ } {
}
{ } {
}
}
{}{}{}
dan

\mavergleichskettedisp
{\vergleichskette
{ D_Z \left\langle  V  ,  W  \right\rangle
}
{ =} { \left\langle \nabla_Z V , W  \right\rangle  + \left\langle  V  ,  \nabla_Z W   \right\rangle
}
{ } {
}
{ } {
}
{ } {
}
}
{}{}{.}
Kita jumlahkan persamaan-persamaan tersebut dan memperoleh

\mavergleichskettedisphandlinks
{\vergleichskettedisphandlinks
{ D_V \left\langle  W  ,  Z  \right\rangle+ D_W \left\langle  Z  ,  V  \right\rangle  - D_Z \left\langle  V  ,  W  \right\rangle
}
{ =} { \left\langle \nabla_V W+\nabla_W V  ,  Z  \right\rangle + \left\langle \nabla_W Z-\nabla_Z W   ,  V  \right\rangle+\left\langle  \nabla_V Z-\nabla_Z V   ,  W  \right\rangle
}
{ } {
}
{ } {
}
{ } {
}
}
{}{}{.}
Kita substitusikan bentuk-bentuk yang diperoleh di atas ke ruas kanan dan mendapatkan

\mavergleichskettealignhandlinks
{\vergleichskettealignhandlinks
{ D_V \left\langle  W  ,  Z  \right\rangle+ D_W \left\langle  Z  ,  V  \right\rangle  - D_Z \left\langle  V  ,  W  \right\rangle
}
{ =} { \left\langle  2 \nabla_VW -[V,W]  ,  Z  \right\rangle + \left\langle  [W,Z]  ,  V  \right\rangle + \left\langle [V,Z]   ,  W  \right\rangle
}
{ =} { 2\left\langle  \nabla_VW   ,  Z  \right\rangle - \left\langle  [V,W]  ,  Z  \right\rangle + \left\langle  [W,Z]  ,  V  \right\rangle + \left\langle  [V,Z]   ,  W  \right\rangle
}
{ } {
}
{ } {
}
}
{}
{}{.}
Dengan menata ulang persamaan ini, diperoleh rumus yang diminta.

Berdasarkan persamaan tersebut,
\mathl{\left\langle  \nabla_V W , Z  \right\rangle}{} untuk sebarang medan vektor ditentukan oleh ruas kanan, yang sama sekali tidak memuat koneksi. Karena ini berlaku untuk setiap medan vektor $Z$, turunan vertikal $\nabla_V W$ dan dengan demikian koneksi linearnya juga ditentukan secara unik.

 }
\endgroup

\chapter{Formulir Solusi Resmi Ujian 8}
\noindent\textit{Solusi yang disediakan oleh sumber resmi; kekosongan sumber dipertahankan.}
\begingroup
\renewcommand{\theHfakt}{o011.exam.official.08.\arabic{fakt}}
%Data institusi

%\input{Dozentdaten}

%\renewcommand{\fachbereich}{Fachbereich}

%\renewcommand{\dozent}{Prof. Dr. . }

%Data ujian

\renewcommand{\klausurgebiet}{ }

\renewcommand{\klausurtyp}{ }

\renewcommand{\klausurdatum}{ . 20}

\klausurvorspann {\fachbereich} {\klausurdatum} {\dozent} {\klausurgebiet} {\klausurtyp}

%Data untuk tabel nilai berikut

\renewcommand{\aeins}{  3  }

\renewcommand{\azwei}{  3   }

\renewcommand{\adrei}{  3  }

\renewcommand{\avier}{  0  }

\renewcommand{\afuenf}{  2  }

\renewcommand{\asechs}{  0  }

\renewcommand{\asieben}{  3  }

\renewcommand{\aacht}{  7  }

\renewcommand{\aneun}{  0  }

\renewcommand{\azehn}{  7  }

\renewcommand{\aelf}{  4  }

\renewcommand{\azwoelf}{  3  }

\renewcommand{\adreizehn}{  3   }

\renewcommand{\avierzehn}{  7  }

\renewcommand{\afuenfzehn}{  2  }

\renewcommand{\asechzehn}{  9  }

\renewcommand{\asiebzehn}{  7  }

\renewcommand{\aachtzehn}{  63  }

\renewcommand{\aneunzehn}{    }

\renewcommand{\azwanzig}{    }

\renewcommand{\aeinundzwanzig}{    }

\renewcommand{\azweiundzwanzig}{    }

\renewcommand{\adreiundzwanzig}{    }

\renewcommand{\avierundzwanzig}{    }

\renewcommand{\afuenfundzwanzig}{    }

\renewcommand{\asechsundzwanzig}{    }

\punktetabellesiebzehn

\klausurnote

\newpage

\setcounter{section}{0}



\inputaufgabeklausurloesung
{3}

{

Definisikan istilah-istilah berikut
\zusatzklammer {yang dicetak miring} {} {}.
\aufzaehlungsechs{\stichwort {Pemetaan Gauss} {}
dari suatu hiperpermukaan diferensiabel

\mavergleichskette
{\vergleichskette
{ Y
}
{ \subseteq }{ \R^n
}
{  }{
}
{    }{
}
{   }{
}
}
{}{}{.}

}{Suatu
\stichwort {peta topologis} {}
pada suatu
\definitionsverweis {manifold topologis}{}{}
$M$.

}{Suatu $C^1$-\stichwort {manifold diferensiabel} {} $M$.

}{\stichwort {Pangkat eksterior} {} ke-$n$ dari suatu
$K$-\definitionsverweis {ruang vektor}{}{}
$V$
\zusatzklammer {cukup tuliskan simbolnya} {} {.}

}{Suatu
\stichwort {isometri lokal} {}
\maabb {\varphi} {L} {M
} {}
antara manifold-manifold Riemann.

}{Suatu koneksi
\stichwort {terintegralkan secara lokal} {}
pada suatu bundel vektor diferensiabel
\maabb {p} {E} {X
} {}
di atas suatu manifold diferensiabel $X$.
}

}
{

\aufzaehlungsechs{Misalkan suatu
\definitionsverweis {orientasi}{}{}
pada $Y$ telah ditetapkan. Maka pemetaan
\maabbdisp {} {Y} { S^{n-1}
} {,}
yang memetakan setiap titik

\mavergleichskette
{\vergleichskette
{ P
}
{ \in }{ Y
}
{  }{
}
{    }{
}
{   }{
}
}
{}{}{}
ke vektor normal satuan yang ditentukan oleh orientasi tersebut disebut
pemetaan Gauss
dari $Y$.
}{Suatu peta topologis adalah setiap
\definitionsverweis {homeomorfisme}{}{}
\maabbdisp {\varphi} {U} {V
} {,}
dengan

\mavergleichskette
{\vergleichskette
{ U
}
{ \subseteq }{ M
}
{  }{
}
{    }{
}
{   }{
}
}
{}{}{}
dan

\mavergleichskette
{\vergleichskette
{ V
}
{ \subseteq }{ \R^n
}
{  }{
}
{    }{
}
{   }{
}
}
{}{}{}
\definitionsverweis {terbuka}{}{.}
}{Suatu
\definitionsverweis {ruang Hausdorff}{}{}
\definitionsverweis {topologis}{}{}
$M$ bersama suatu
\definitionsverweis {penutup terbuka}{}{}

\mathl{M= \bigcup_{i \in I} U_i}{} dan
\definitionsverweis {peta-peta}{}{}
\maabbdisp {\alpha_i} {U_i } { V_i
} {}
dengan
\mathl{V_i \subseteq \R^n}{} terbuka, sedemikian rupa sehingga
\definitionsverweis {pemetaan-pemetaan transisi}{}{}
\maabbdisp {\alpha_j \circ (\alpha_i)^{-1}} {V_i \cap \alpha_i(U_i \cap U_j)   } { V_j \cap \alpha_j(U_i \cap U_j)
} {}
$C^1$-\definitionsverweis {difeomorfisme}{}{}
untuk semua
\mathl{i,j \in  I}{}. Struktur ini disebut manifold diferensiabel.
}{Ruang vektor-$K$
\mathl{\bigwedge^n V}{} disebut pangkat eksterior ke-$n$ dari $V$.
}{Suatu
\definitionsverweis {pemetaan diferensiabel}{}{}
\maabb {\varphi} {L} {M
} {}
disebut
isometri lokal
jika untuk setiap titik

\mavergleichskette
{\vergleichskette
{ P
}
{ \in }{ L
}
{  }{
}
{    }{
}
{   }{
}
}
{}{}{}
pemetaan
\definitionsverweis {tangen}{}{}
\maabbdisp {T_P \varphi} {T_PL} { T_{\varphi(P) } M
} {}
merupakan suatu
\definitionsverweis {isometri}{}{}
terhadap hasil kali skalar yang diberikan.
}{Koneksi tersebut disebut
terintegralkan secara lokal
jika untuk setiap titik

\mavergleichskette
{\vergleichskette
{ e
}
{ \in }{ E
}
{  }{
}
{    }{
}
{   }{
}
}
{}{}{}
terdapat suatu

\mavergleichskette
{\vergleichskette
{ p(e)
}
{ \in }{ U
}
{ \subseteq }{ X
}
{    }{
}
{   }{
}
}
{}{}{}
\definitionsverweis {penampang horizontal}{}{}
\maabbdisp {s} {U} {E \vert_U
} {}
yang didefinisikan pada lingkungan terbuka tersebut dan melalui $e$.
}

 }



\inputaufgabeklausurloesung
{3}

{

Rumuskan teorema-teorema berikut.
\aufzaehlungdrei{\stichwort {Teorema isometri untuk transpor paralel pada suatu hiperpermukaan} {}

\mavergleichskette
{\vergleichskette
{  Y
}
{ \subseteq }{  \R^n
}
{  }{
}
{    }{
}
{   }{
}
}
{}{}{.}}{\stichwort {Theorema egregium} {.}}{Teorema tentang partisi kesatuan.}

}
{

\aufzaehlungdrei{\faktsituation {Misalkan

\mavergleichskette
{\vergleichskette
{ Y
}
{ \subseteq }{ W
}
{ \subseteq }{ \R^n
}
{    }{
}
{   }{
}
}
{}{}{,}
$W$
\definitionsverweis {terbuka}{}{,}
dan merupakan suatu
\definitionsverweis {hiperpermukaan diferensiabel}{}{.}
Misalkan
\maabb {\gamma} {[a,b]} {Y
} {}
suatu
\definitionsverweis {kurva diferensiabel}{}{}
dengan

\mavergleichskette
{\vergleichskette
{ \gamma(a)
}
{ = }{ P
}
{  }{
}
{    }{
}
{   }{
}
}
{}{}{}
dan

\mavergleichskette
{\vergleichskette
{ \gamma(b)
}
{ = }{ Q
}
{  }{
}
{    }{
}
{   }{
}
}
{}{}{.}}
\faktfolgerung {Maka
\definitionsverweis {transpor paralel}{}{}
sepanjang $\gamma$ merupakan suatu
\definitionsverweis {isometri}{}{}
\maabbdisp {} {T_PY} {T_QY
} {.}}
\faktzusatz {}
\faktzusatz {}}{\faktsituation {Misalkan

\mavergleichskette
{\vergleichskette
{ Y
}
{ \subseteq }{ \R^3
}
{  }{
}
{    }{
}
{   }{
}
}
{}{}{}
suatu
\definitionsverweis {permukaan berorientasi}{}{}
dan misalkan
\maabbdisp {\varphi} { V } { Y
} {,}

\mavergleichskette
{\vergleichskette
{ V
}
{ \subseteq }{ \R^2
}
{  }{
}
{    }{
}
{   }{
}
}
{}{}{,}
suatu parametrisasi lokal $Y$ yang dua kali terdiferensialkan, dengan parameter $u,v$. Misalkan
\mathl{\left(  g_{ij}  \right)}{}
\definitionsverweis {matriks fundamental pertama}{}{}
pada $V$.}
\faktfolgerung {Maka, dengan menggunakan
\definitionsverweis {simbol-simbol Christoffel}{}{,}
\definitionsverweis {kelengkungan Gauss}{}{}
memenuhi hubungan

\mavergleichskettedisphandlinks
{\vergleichskettedisphandlinks
{ K
}
{ =} { { \frac{   1  }{ g_{11}  } }  \left(  \partial_2 \Gamma_{11}^2 - \partial_1 \Gamma_{12}^2 + \Gamma^1_{11} \Gamma^2_{12}  + \Gamma_{11}^2 \Gamma_{22}^2 - \Gamma_{12}^1 \Gamma_{11}^2-  \left(  \Gamma_{12}^2  \right)^2  \right)
}
{ } {
}
{ } {
}
{ } {
}
}
{}{}{.}}
\faktzusatz {}
\faktzusatz {}}{Misalkan $M$ suatu
\definitionsverweis {manifold diferensiabel}{}{}
dengan
\definitionsverweis {basis topologi yang terhitung}{}{.}
Maka, untuk setiap penutup terbuka, terdapat suatu
\definitionsverweis {partisi kesatuan}{}{}
\definitionsverweis {yang terdiferensialkan secara kontinu}{}{}
dan tersubordinasi pada penutup tersebut.}

 }



\inputaufgabeklausurloesung
{3}

{

Misalkan
\maabbdisp {f,g} {\R} { \R^n
} {}
\definitionsverweis {kurva-kurva diferensiabel}{}{.}
Hitung
\definitionsverweis {turunan}{}{}
dari fungsi
\maabbeledisp {} {\R} {\R
} { t } { \left\langle f(t) , g(t)  \right\rangle
} {.}
Rumuskan hasilnya dengan menggunakan hasil kali dalam.

}
{

Misalkan $f_i(t)$ dan $g_i(t)$ masing-masing fungsi komponen dari
\mathkor {} {f} {dan} {g} {.}
Fungsi yang dimaksud adalah

\mavergleichskettedisp
{\vergleichskette
{  \left\langle f(t) , g(t)  \right\rangle
}
{ =} { \sum_{i = 1}^n  f_i(t) g_i(t)
}
{ } {
}
{ } {
}
{ } {
}
}
{}{}{}
dengan turunan
\zusatzklammer {aturan hasil kali diterapkan pada setiap suku} {} {}

\mathdisp {\sum_{i = 1}^n  f'_i(t) g_i(t)  + \sum_{i = 1}^n  f_i(t) g'_i(t)} { . }

Dengan hasil kali skalar, bentuk ini dapat ditulis sebagai

\mathdisp {\left\langle f'(t) ,  g(t)  \right\rangle + \left\langle f(t) ,  g'(t)  \right\rangle} {  }

.

 }



\inputaufgabeklausurloesung
{0}

{

}
{  }



\inputaufgabeklausurloesung
{2}

{

Misalkan

\mavergleichskette
{\vergleichskette
{ Y
}
{ \subseteq }{ W
}
{ \subseteq }{ \R^n
}
{    }{
}
{   }{
}
}
{}{}{,}
$W$
\definitionsverweis {terbuka}{}{,}
suatu
\definitionsverweis {hiperpermukaan diferensiabel}{}{},
dan
\maabb {\gamma} {I =[a,b]} {Y
} {}
suatu
\definitionsverweis {kurva geodesik}{}{.}
Buktikan bahwa
\definitionsverweis {transpor paralel}{}{}
sepanjang $\gamma$ memetakan vektor tangensial

\mavergleichskette
{\vergleichskette
{ \gamma'(a)
}
{ \in }{ T_{\gamma(a)}Y
}
{  }{
}
{    }{
}
{   }{
}
}
{}{}{}
ke vektor tangensial

\mavergleichskette
{\vergleichskette
{ \gamma'(b)
}
{ \in }{ T_{\gamma(b)}Y
}
{  }{
}
{    }{
}
{   }{
}
}
{}{}{.}

}
{

Menurut
Lemma 6.8 (Differentialgeometrie (Osnabrück 2023))
medan kecepatan $\gamma'$ merupakan suatu
\definitionsverweis {medan vektor paralel}{}{}
sepanjang $\gamma$. Oleh karena itu,
\definitionsverweis {transpor paralel}{}{}
$\Psi_\gamma$ sepanjang $\gamma$ memetakan vektor

\mavergleichskette
{\vergleichskette
{ \gamma'(a)
}
{ \in }{  T_{\gamma(a)} Y
}
{  }{
}
{    }{
}
{   }{
}
}
{}{}{}
ke vektor

\mavergleichskette
{\vergleichskette
{ \gamma'(b)
}
{ \in }{  T_{\gamma(b)} Y
}
{  }{
}
{    }{
}
{   }{
}
}
{}{}{}
.

 }



\inputaufgabeklausurloesung
{0}

{

}
{  }



\inputaufgabeklausurloesung
{3}

{

Apakah pemetaan
\maabbeledisp {\varphi} {S^1} {S^1
} {(x,y)} {( \betrag {  x  } ,y)
} {,}
\definitionsverweis {terdiferensialkan}{}{?}

}
{

Kita meninjau komposisi pemetaan-pemetaan

\mathdisp {\R \stackrel{ \theta }{\longrightarrow} S^1 \stackrel{\varphi}{\longrightarrow} S^1 \subset \R^2 \stackrel{p_1}{\longrightarrow} \R} { , }

dengan

\mavergleichskettedisp
{\vergleichskette
{ \theta(t)
}
{ =} {(\cos  t , \sin  t )
}
{ } {
}
{ } {
}
{ } {
}
}
{}{}{}
dan $p_1$ merupakan proyeksi pertama. Parametrisasi trigonometrik, inklusi
\zusatzklammer {suatu submanifold tertutup} {} {,}
dan proyeksi tersebut semuanya diferensiabel. Jika $\varphi$ diferensiabel, maka komposisi keseluruhannya juga harus diferensiabel. Akan tetapi, komposisi itu adalah
\maabbeledisp {} {\R} {\R
} { t } {\betrag {  \cos  t  }
} {.}
Pemetaan ini tidak diferensiabel di
\mathl{t= \pi/2}{} karena gradien dari kiri di titik tersebut adalah $-1$, sedangkan gradien dari kanan adalah $1$. Jadi, pemetaan $\varphi$ tidak diferensiabel.

 }



\inputaufgabeklausurloesung
{7}

{

Misalkan $M$ suatu manifold diferensiabel kompak dengan bentuk volume positif kontinu $\omega$. Buktikan bahwa

\mavergleichskettedisp
{\vergleichskette
{ \int_M \omega
}
{ <} { \infty
}
{ } {
}
{ } {
}
{ } {
}
}
{}{}{}.

}
{

Untuk setiap titik

\mavergleichskette
{\vergleichskette
{ P
}
{ \in }{ M
}
{  }{
}
{    }{
}
{   }{
}
}
{}{}{}
terdapat suatu lingkungan peta terbuka

\mavergleichskette
{\vergleichskette
{ P
}
{ \in }{ U
}
{  }{
}
{    }{
}
{   }{
}
}
{}{}{}
dan suatu pemetaan peta
\maabbdisp {\alpha} { U } { V
} {}
dengan

\mavergleichskette
{\vergleichskette
{ V
}
{ \subseteq }{ \R^n
}
{  }{
}
{    }{
}
{   }{
}
}
{}{}{}
terbuka, sedemikian rupa sehingga

\mavergleichskette
{\vergleichskette
{ \alpha^{-1 *} \omega
}
{ = }{ fdx_1 \wedge \ldots \wedge dx_n
}
{  }{
}
{    }{
}
{   }{
}
}
{}{}{}
berlaku dengan
\mathl{f}{} kontinu dan positif. Kita juga memperoleh suatu lingkungan terbuka

\mavergleichskette
{\vergleichskette
{ P
}
{ \in }{ U'
}
{ \subseteq }{ U
}
{    }{
}
{   }{
}
}
{}{}{,}
yang homeomorf dengan suatu bola terbuka

\mavergleichskette
{\vergleichskette
{ U'
}
{ \cong }{ B'
}
{ \subseteq }{ V
}
{    }{
}
{   }{
}
}
{}{}{}
dan kita dapat menganggap bahwa penutupan bola tersebut seluruhnya terletak di $V$. Bola tertutup itu tertutup dan terbatas; karena itu, fungsi kontinu $f$ terbatas pada bola tersebut dan dengan demikian juga pada $U'$. Maka
\mathl{\int_{U'} \omega}{} berhingga, dengan $U'$ suatu lingkungan terbuka dari $P$.

Himpunan-himpunan terbuka

\mavergleichskette
{\vergleichskette
{ U'
}
{ = }{ U'(P)
}
{  }{
}
{    }{
}
{   }{
}
}
{}{}{}
ini menutupi $M$. Karena kekompakan, terdapat suatu penutup berhingga

\mavergleichskettedisp
{\vergleichskette
{ M
}
{ =} { \bigcup_{i = 1} ^n U_i
}
{ } {
}
{ } {
}
{ } {
}
}
{}{}{}
dengan

\mavergleichskettedisp
{\vergleichskette
{ \int_{U_i } \omega
}
{ <} { \infty
}
{ } {
}
{ } {
}
{ } {
}
}
{}{}{.}
Karena kepositifan, diperoleh

\mavergleichskettedisp
{\vergleichskette
{ \int_{M} \omega
}
{ \leq} { \sum_{i = 1}^n  \left(  \int_{U_i } \omega  \right)
}
{ <} { \infty
}
{ } {
}
{ } {
}
}
{}{}{.}

 }



\inputaufgabeklausurloesung
{0}

{

}
{  }



\inputaufgabeklausurloesung
{7 (1+2+4)}

{

\aufzaehlungdreiabc{ Berikan suatu sifat topologis yang menunjukkan bahwa
\definitionsverweis {interval satuan terbuka}{}{}
dan
\definitionsverweis {lingkaran}{}{}
$S^1$ tidak
\definitionsverweis {homeomorfik}{}{.}
}{Buktikan bahwa setiap
$1$-\definitionsverweis {bentuk diferensial}{}{}
kontinu pada
\mathl{]0,1[}{}
bersifat
\definitionsverweis {eksak}{}{.}
}{ Berikan suatu
$1$-bentuk konkret pada $S^1$ yang
\definitionsverweis {tertutup}{}{}
tetapi tidak eksak.
}

}
{

a) Menurut
Satz 13.14 (Differentialgeometrie (Osnabrück 2023))
$S^1$, sebagai himpunan bagian tertutup dan terbatas dari $\R^2$, bersifat
\definitionsverweis {kompak}{}{,}
sedangkan interval terbuka tersebut tidak kompak.

b) Berlaku

\mavergleichskettedisp
{\vergleichskette
{ \omega
}
{ =} { hdt
}
{ } {
}
{ } {
}
{ } {
}
}
{}{}{}
dengan suatu fungsi kontinu
\maabbdisp {h} {]0,1[} {\R
} {.}
Menurut
Hauptsatz der Infinitesimalrechnung,
$h$ mempunyai suatu antiturunan $H$. Dalam bahasa bentuk diferensial, hal ini berarti

\mavergleichskettedisp
{\vergleichskette
{H'
}
{ =} { hdt
}
{ =} { \omega
}
{ } {
}
{ } {
}
}
{}{}{,}
sehingga $\omega$ eksak.

c) Misalkan
\mathl{S^1 \subset \R^2}{} diberikan, dengan koordinat pada $\R^2$ dinotasikan oleh
\mathkor {} {u} {dan} {v} {}
. Kita meninjau bentuk diferensial

\mavergleichskettedisp
{\vergleichskette
{ \omega
}
{ =} { udv-vdu
}
{ } {
}
{ } {
}
{ } {
}
}
{}{}{}
pada $S^1$. Karena $S^1$ berdimensi satu, bentuk ini tertutup. Untuk lintasan
\maabbeledisp {\gamma} {[0, 2 \pi]} {S^1
} { t } { \begin{pmatrix}  \cos  t  \\ \sin  t     \end{pmatrix}
} {,}
berlaku

\mavergleichskettedisp
{\vergleichskette
{ \int_\gamma \omega
}
{ =} { \int_0^{2 \pi} \gamma^* \omega  dt
}
{ =} { \int_0^{2 \pi}  \cos  t  \cdot \cos  t  +  \sin  t   \sin  t   dt
}
{ =} { \int_0^{2 \pi} 1 dt
}
{ =} { 2 \pi
}
}
{}{}{.}
Menurut
Korollar 23.3 (Differentialgeometrie (Osnabrück 2023))
$\omega$ tidak mungkin eksak.

 }



\inputaufgabeklausurloesung
{4}

{

Misalkan

\mavergleichskette
{\vergleichskette
{ U
}
{ \subseteq }{ \R^n
}
{  }{
}
{    }{
}
{   }{
}
}
{}{}{}
\definitionsverweis {terbuka}{}{}
dan
\maabb {f} { U } { \R
} {}
suatu
\definitionsverweis {fungsi yang terdiferensialkan secara kontinu}{}{}
dua kali. Buktikan bahwa

\mavergleichskettedisp
{\vergleichskette
{ d(df)
}
{ =} { 0
}
{ } {
}
{ } {
}
{ } {
}
}
{}{}{,}
dengan $d$ menyatakan
\definitionsverweis {turunan eksterior}{}{.}

}
{

Untuk suatu bentuk-$1$

\mavergleichskette
{\vergleichskette
{ \omega
}
{ = }{ \sum_{j = 1}^n g_jdx_j
}
{  }{
}
{    }{
}
{   }{
}
}
{}{}{}
dengan menggunakan

\mavergleichskette
{\vergleichskette
{ dx_i \wedge dx_j
}
{ = }{ -dx_j \wedge dx_i
}
{  }{
}
{    }{
}
{   }{
}
}
{}{}{}

\mavergleichskettealign
{\vergleichskettealign
{d \omega
}
{ =} { \sum_{j = 1 }^n  d(g_j dx_j)
}
{ =} { \sum_{j = 1 }^n   \left(  \sum_{i = 1}^n { \frac{   \partial  g_j   }{  \partial   x_i   } } dx_i \wedge dx_j  \right)
}
{ =} { \sum _{1 \leq i <j\leq n} \left(  { \frac{   \partial  g_j   }{  \partial   x_i   } } - { \frac{   \partial  g_i   }{  \partial   x_j   } }  \right) dx_i \wedge dx_j
}
{ } {
}
}
{}
{}{.}
Untuk suatu fungsi $f$ yang dua kali terdiferensialkan secara kontinu, berlaku

\mavergleichskette
{\vergleichskette
{ df
}
{ = }{ \sum_{j = 1}^n g_j dx_j
}
{  }{
}
{    }{
}
{   }{
}
}
{}{}{}
dengan
\definitionsverweis {turunan-turunan parsial}{}{}

\mavergleichskette
{\vergleichskette
{ g_j
}
{ = }{ { \frac{   \partial   f   }{  \partial   x_j   } }
}
{  }{
}
{    }{
}
{   }{
}
}
{}{}{,}
dan karena itu

\mavergleichskette
{\vergleichskette
{ d(df)
}
{ = }{ 0
}
{  }{
}
{    }{
}
{   }{
}
}
{}{}{}
menurut
Satz von Schwarz.

 }



\inputaufgabeklausurloesung
{3}

{

Deskripsikan
\definitionsverweis {pemetaan}{}{}
\maabbeledisp {\psi} { \Complex \setminus \{1\} } { \Complex
} {w} { { \frac{   (w+1)  { \mathrm i}  }{ 1-w  } }
} {}
dalam koordinat riil.

}
{

Kita menulis

\mavergleichskette
{\vergleichskette
{ w
}
{ = }{ a+b { \mathrm i}
}
{  }{
}
{    }{
}
{   }{
}
}
{}{}{.}
Maka

\mavergleichskettealign
{\vergleichskettealign
{  \psi ( a+b { \mathrm i}   )
}
{ =} { { \frac{   \left(  a+b { \mathrm i} +1  \right) { \mathrm i}      }{  1-  a-b { \mathrm i}      } }
}
{ =} { { \frac{     a { \mathrm i} +{ \mathrm i}  -b  }{  1-  a-b { \mathrm i}      } }
}
{ =} { { \frac{     \left(  a { \mathrm i} +{ \mathrm i}  -b  \right) \left(  1-  a+b { \mathrm i}  \right)  }{  \left(  1-  a-b { \mathrm i}  \right) \left(  1-  a+b { \mathrm i}  \right)      } }
}
{ =} { { \frac{     \left(  a+1 \right) \left(  1-a  \right) { \mathrm i}  -b^2 { \mathrm i} -b \left(  1-  a \right) -b\left(  a+1  \right)  }{  \left(  1-  a \right)^2 +b^2  } }
}
}
{
\vergleichskettefortsetzungalign
{ =} { { \frac{     \left(  1-a ^2-b^2  \right) { \mathrm i}   -2 b  }{  \left(  1-  a \right)^2 +b^2  } }
}
{ } {
}
{ } {}
{ } {}
}
{}{.}
Dalam koordinat riil, dengan demikian,

\mavergleichskettedisp
{\vergleichskette
{ \psi \begin{pmatrix}  a  \\b   \end{pmatrix}
}
{ =} { { \frac{   1  }{  -2a+1+a^2+b^2 } } \begin{pmatrix}  -2b \\ 1-a^2-b^2   \end{pmatrix}
}
{ } {
}
{ } {
}
{ } {
}
}
{}{}{.}

 }



\inputaufgabeklausurloesung
{3}

{

Misalkan

\mavergleichskette
{\vergleichskette
{ H
}
{ \subseteq }{ \R^n
}
{  }{
}
{    }{
}
{   }{
}
}
{}{}{}
suatu
\definitionsverweis {setengah ruang Euklides}{}{}
dan

\mavergleichskette
{\vergleichskette
{ P
}
{ \in }{ H
}
{  }{
}
{    }{
}
{   }{
}
}
{}{}{.}
Andaikan terdapat suatu himpunan yang terbuka di $\R^n$, yaitu $V$, dengan

\mavergleichskette
{\vergleichskette
{ P
}
{ \in }{ V
}
{ \subseteq }{ H
}
{    }{
}
{   }{
}
}
{}{}{.}
Buktikan bahwa
\mathl{P}{} bukan titik batas dari $H$.

}
{

Misalkan

\mavergleichskettedisp
{\vergleichskette
{ P
}
{ =} { \begin{pmatrix}  x_1  \\ x_2 \\  \vdots\\ x_n   \end{pmatrix}
}
{ } {
}
{ } {
}
{ } {
}
}
{}{}{.}
Kita dapat mengganti, di $\R^n$, lingkungan terbuka $V$ dengan suatu bola terbuka
\zusatzklammer {di $\R^n$} {} {}
\mathl{U \left(   P , \epsilon  \right)}{} dengan

\mavergleichskette
{\vergleichskette
{ P
}
{ \in }{   U \left(   P , \epsilon  \right)
}
{ \subseteq }{  V
}
{ \subseteq   }{ H
}
{   }{
}
}
{}{}{}
. Jika $P$ suatu titik batas, maka

\mavergleichskette
{\vergleichskette
{ x_1
}
{ = }{ 0
}
{  }{
}
{    }{
}
{   }{
}
}
{}{}{.}
Akan tetapi, dalam hal itu

\mavergleichskette
{\vergleichskette
{ \begin{pmatrix}  - { \frac{   \epsilon }{  2 } }  \\ x_2 \\  \vdots\\ x_n   \end{pmatrix}
}
{ \in }{ U \left(   P , \epsilon  \right)
}
{  }{
}
{    }{
}
{   }{
}
}
{}{}{,}
padahal titik ini tidak termasuk dalam $H$.

 }



\inputaufgabeklausurloesung
{7}

{

Misalkan $M$ suatu
\definitionsverweis {manifold diferensiabel}{}{}
berdimensi $n$ dan

\mavergleichskette
{\vergleichskette
{ P
}
{ \in }{ M
}
{  }{
}
{    }{
}
{   }{
}
}
{}{}{}
suatu titik. Buktikan bahwa terdapat suatu
\definitionsverweis {manifold dengan batas}{}{}
$N$ yang batasnya $\partial N$ difeomorfik dengan permukaan bola berdimensi $(n-1)$, yaitu $S^{n-1}$, sedemikian sehingga
\mathkor {} {M \setminus \{ P \}} {dan} {N \setminus \partial N} {}
saling difeomorfik.

}
{

Misalkan

\mavergleichskette
{\vergleichskette
{ P
}
{ \in }{  U
}
{ \subseteq }{  M
}
{    }{
}
{   }{
}
}
{}{}{}
suatu daerah peta terbuka dengan peta
\maabbdisp {\alpha} { U } {V \subseteq \R^n
} {.}
Kita dapat menganggap bahwa $\alpha(P)$ adalah titik nol dan $V$ adalah bola terbuka berjari-jari $3$. Misalkan
\maabbdisp {\varphi} { U \left(   0 , 3  \right) \setminus \{0\}  } { U \left(   0 , 3  \right) \setminus B  \left(  0 , 1 \right)
} {}
suatu difeomorfisme yang merupakan identitas pada
\mathl{U \left(   0 , 3  \right) \setminus U \left(   0 , 2  \right)}{} dan memetakan
\mathl{U \left(   0 , 2  \right) \setminus \{0\}}{} ke
\mathl{U \left(   0 , 2  \right) \setminus B  \left(  0 , 1 \right)}{}. Pemetaan seperti ini diperoleh dengan meregangkan titik-titik menggunakan fungsi
\maabbdisp {h} {[0,3]} { [0,3]
} {}
dengan

\mavergleichskettedisp
{\vergleichskette
{ h(r)
}
{ =} { \begin{cases} 1+ { \frac{   1  }{  4 } } r^2 \text{ jika } 0 \leq r \leq 2 \, , \\ r \text{ jika } r >  2  \, , \end{cases}
}
{ } {
}
{ } {
}
{ } {
}
}
{}{}{}
. Selanjutnya, citra difeomorfisme $\varphi$, yaitu
\mathl{U \left(   0 , 3  \right) \setminus B  \left(  0 , 1 \right)}{,} dapat dipandang sebagai

\mavergleichskettedisp
{\vergleichskette
{  U \left(   0 , 3  \right) \setminus B  \left(  0 , 1 \right)
}
{ \subset} {  U \left(   0 , 3  \right) \setminus U \left(   0 , 1  \right)
}
{ } {
}
{ } {
}
{ } {
}
}
{}{}{}
. Himpunan yang lebih besar merupakan suatu manifold dengan batas $S \left(   0 ,  1  \right)$. Misalkan $N$ manifold yang diperoleh dengan mengganti $U$ oleh
\mathl{U \left(   0 , 3  \right) \setminus U \left(   0 , 1  \right)}{}. Permukaan bola tersebut menjadi batas dari $N$. Difeomorfisme $\varphi$ dapat diperluas menjadi suatu difeomorfisme
\maabbdisp {} {M \setminus \{P\} } { N \setminus \partial N
} {}
karena pada daerah transisi terbuka
\mathl{U \left(   0 , 3  \right) \setminus U \left(   0 , 2  \right)}{} pemetaan tersebut merupakan identitas.

 }



\inputaufgabeklausurloesung
{2}

{

Berikan suatu
\definitionsverweis {penghabisan kompak}{}{}
untuk $\R^d$.

}
{

Kita meninjau
\definitionsverweis {bola-bola tertutup}{}{}

\mathbed {B  \left(  0 , n  \right)} {}
 {n \in \N_+} {}
 {} {} {} {,}
yang
\definitionsverweis {kompak}{}{}
dan menutupi $\R^d$. Interior terbuka dari
\mathl{B  \left(  0 , n+1  \right)}{}
adalah
\definitionsverweis {bola terbuka}{}{}

\mathl{U \left(   0,n+1  \right)}{}, dan karena

\mavergleichskettedisp
{\vergleichskette
{ B  \left(  0,n \right)
}
{ \subseteq} { U \left(   0,n+1  \right)
}
{ } {
}
{ } {
}
{ } {
}
}
{}{}{}
diperoleh suatu
\definitionsverweis {penghabisan kompak}{}{.}

 }



\inputaufgabeklausurloesung
{9}

{

Buktikan teorema titik tetap Brouwer.

}
{

Untuk menyederhanakan notasi, misalkan

\mavergleichskette
{\vergleichskette
{ r
}
{ = }{ 1
}
{  }{
}
{    }{
}
{   }{
}
}
{}{}{.}  Andaikan terdapat suatu pemetaan $\psi$ tanpa titik tetap yang terdiferensialkan secara kontinu. Maka selalu

\mavergleichskettedisp
{\vergleichskette
{ x
}
{ \neq} { \psi(x)
}
{ } {
}
{ } {
}
{ } {
}
}
{}{}{,}
sehingga kedua titik tersebut menentukan suatu garis. Gagasannya adalah menggunakan garis ini untuk mengambil salah satu
\zusatzklammer {dari kedua} {} {}
titik potong dengan permukaan bola sebagai titik citra suatu retraksi ke batas. Dengan fungsi bantu

\mavergleichskettedisp
{\vergleichskette
{ h(x)
}
{ =} { { \frac{   \psi(x)-x }{  \Vert { \psi (x)-x} \Vert  } }
}
{ } {
}
{ } {
}
{ } {
}
}
{}{}{}
kita mendefinisikan suatu pemetaan
\maabbdisp {\varphi} { B  \left(  0 , 1 \right) } { \R^n
} {}
melalui

\mavergleichskettedisphandlinks
{\vergleichskettedisphandlinks
{ \varphi(x)
}
{ =} { x + \left(  - \left\langle  x  , h(x)  \right\rangle + \sqrt{1 + \left\langle  x  , h(x)  \right\rangle^2 - \Vert { x} \Vert^2  }  \right) \cdot h(x)
}
{ } {
}
{ } {
}
{ } {
}
}
{}{}{.}
Ekspresi di bawah akar tersebut positif. Hal ini jelas untuk

\mavergleichskette
{\vergleichskette
{ \Vert { x } \Vert
}
{ < }{ 1
}
{  }{
}
{    }{
}
{   }{
}
}
{}{}{}
dan untuk

\mavergleichskette
{\vergleichskette
{ \Vert { x } \Vert
}
{ = }{ 1
}
{  }{
}
{    }{
}
{   }{
}
}
{}{}{}
diperoleh suatu titik pada permukaan bola yang garis penghubungnya dengan titik bola
\mathl{\psi(x)}{} tidak tegak lurus terhadap $x$
\zusatzklammer {ruang singgung afin pada suatu titik permukaan bola hanya berpotongan dengan bola di satu titik} {} {,}
sehingga

\mavergleichskettedisp
{\vergleichskette
{ \left\langle  x  , h(x)  \right\rangle
}
{ \neq} { 0
}
{ } {
}
{ } {
}
{ } {
}
}
{}{}{}
berlaku. Karena akar kuadrat dan nilai mutlak di luar titik nol
\definitionsverweis {terdiferensialkan secara kontinu}{}{,}
baik
\mathkor {} {h(x)} {maupun} {\varphi(x)} {}
merupakan pemetaan yang terdiferensialkan secara kontinu. Menurut
Aufgabe 23.23 (Differentialgeometrie (Osnabrück 2023))
pemetaan $\varphi$ memetakan bola ke permukaan bola dan pembatasannya pada permukaan bola adalah identitas. Dengan demikian, diperoleh suatu
\definitionsverweis {retraksi}{}{}
yang terdiferensialkan secara kontinu dari bola penuh tertutup ke batasnya, yang menurut
Satz 23.9 (Differentialgeometrie (Osnabrück 2023))
tidak mungkin ada.

 }



\inputaufgabeklausurloesung
{7 (1+2+4)}

{

Pada
\definitionsverweis {bundel vektor trivial}{}{}

\mathl{\R \times \R^2}{} di atas $\R$, kita meninjau
\definitionsverweis {koneksi trivial}{}{}
$\nabla$.
\aufzaehlungdrei{Buktikan bahwa

\mavergleichskette
{\vergleichskette
{ f_1(t)
}
{ = }{ \left(  1+ t^2  , \,  t^3                   \right)
}
{  }{
}
{    }{
}
{   }{
}
}
{}{}{}
dan

\mavergleichskette
{\vergleichskette
{ f_2(t)
}
{ = }{ \left(  t  , \,   1+t^2                   \right)
}
{  }{
}
{    }{
}
{   }{
}
}
{}{}{}
merupakan
\definitionsverweis {penampang basis}{}{}
dari bundel vektor tersebut.
}{Nyatakan penampang basis standar $e_1,e_2$ sebagai kombinasi linear dari penampang basis $f_1,f_2$.
}{Tentukan
\definitionsverweis {simbol-simbol Christoffel}{}{}
dari $\nabla$ terhadap penampang-penampang basis tersebut.
}

}
{

\aufzaehlungdrei{Berlaku

\mavergleichskettedisp
{\vergleichskette
{ \det  \begin{pmatrix}  1+t^2   &  t^3  \\  t  &  1+t^2  \end{pmatrix}
}
{ =} { 1+2t^2+t^4 -t^4
}
{ =} { 1 +2t^2
}
{ >} {  0
}
{ } {
}
}
{}{}{}
selalu positif. Oleh karena itu, kedua vektor
\mathkor {} {f_1(t)} {dan} {f_2(t)} {}
untuk setiap

\mavergleichskette
{\vergleichskette
{ t
}
{ \in }{  \R
}
{  }{
}
{    }{
}
{   }{
}
}
{}{}{}
membentuk suatu basis.
}{Dari hubungan

\mavergleichskettedisp
{\vergleichskette
{ \begin{pmatrix} f_1  \\f_2   \end{pmatrix}
}
{ =} { \begin{pmatrix}  1+t^2  &   t^3  \\  t  &  1+t^2 \end{pmatrix}  \begin{pmatrix} e_1  \\e_2   \end{pmatrix}
}
{ } {
}
{ } {
}
{ } {
}
}
{}{}{}
diperoleh

\mavergleichskettedisp
{\vergleichskette
{ \begin{pmatrix} e_1  \\e_2   \end{pmatrix}
}
{ =} { { \frac{   1  }{  1+2t^2 } } \begin{pmatrix}  1+t^2  &   -t^3  \\  -t &  1+t^2 \end{pmatrix}  \begin{pmatrix} f_1  \\f_2   \end{pmatrix}
}
{ } {
}
{ } {
}
{ } {
}
}
{}{}{,}
jadi

\mavergleichskettedisp
{\vergleichskette
{ e_1
}
{ =} { { \frac{   1+t^2 }{  1+2t^2 } }   f_1 - { \frac{  t^3 }{  1+2t^2 } } f_2
}
{ } {
}
{ } {
}
{ } {
}
}
{}{}{}
dan

\mavergleichskettedisp
{\vergleichskette
{ e_2
}
{ =} { -{ \frac{   t  }{  1+2t^2 } }   f_1 + { \frac{   1+t^2 }{  1+2t^2 } } f_2
}
{ } {
}
{ } {
}
{ } {
}
}
{}{}{.}
}{Untuk koneksi trivial, berlaku

\mavergleichskettedisp
{\vergleichskette
{ \nabla_\partial (f)
}
{ =} { \partial f
}
{ } {
}
{ } {
}
{ } {
}
}
{}{}{.}
Dengan demikian,

\mavergleichskettealign
{\vergleichskettealign
{ \nabla_\partial (f_1)
}
{ =} { \nabla_\partial \left(  \left(  1+t^2  \right) e_1 + t^3e_2  \right)
}
{ =} { 2t e_1 + 3t^2 e_2
}
{ =} { 2t  \left(  { \frac{   1+t^2 }{  1+2t^2 } }   f_1 - { \frac{  t^3 }{  1+2t^2 } } f_2   \right)  + 3t^2  \left(  -{ \frac{   t  }{  1+2t^2 } }   f_1 + { \frac{   1+t^2 }{  1+2t^2 } } f_2   \right)
}
{ =} {  \left(   { \frac{   2t+2t^3 }{  1+2t^2 } } - { \frac{   3t^3 }{  1+2t^2 } }  \right) f_1  +  \left(  - { \frac{  t^3 }{  1+2t^2 } }   + { \frac{   3t^2+3t^4 }{  1+2t^2 } }     \right)   f_2
}
}
{
\vergleichskettefortsetzungalign
{ =} { { \frac{   2t - t^3 }{  1+2t^2 } }  f_1  +  { \frac{   3t^2-t^3+3t^4 }{  1+2t^2 } }      f_2
}
{ } {}
{ } {}
{ } {}
}
{}{}
dan

\mavergleichskettealign
{\vergleichskettealign
{ \nabla_\partial (f_2)
}
{ =} { \nabla_\partial \left(  t e_1 + \left(  1+t^2  \right) e_2  \right)
}
{ =} { e_1 + 2t  e_2
}
{ =} { { \frac{   1+t^2 }{  1+2t^2 } } f_1 - { \frac{  t^3 }{  1+2t^2 } } f_2 + 2t \left(  -{ \frac{   t  }{  1+2t^2 } } f_1 + { \frac{   1+t^2 }{  1+2t^2 } } f_2   \right)
}
{ =} { \left( { \frac{   1+t^2 }{  1+2t^2 } }  - { \frac{   2t^2 }{  1+2t^2 } }  \right) f_1  + \left(  - { \frac{   t^3  }{  1+2t^2  } } +  { \frac{   2t+2t^3  }{  1+2t^2 } }  \right) f_2
}
}
{
\vergleichskettefortsetzungalign
{ =} { { \frac{   1 - t^2  }{  1+2t^2 } } f_1 + { \frac{   2t   +t^3 }{  1+2t^2 } }      f_2
}
{ } {}
{ } {}
{ } {}
}
{}{.}
Fungsi-fungsi koefisien tersebut adalah simbol-simbol Christoffel.
}

 }
\endgroup

\chapter{Formulir Solusi Resmi Ujian 9}
\noindent\textit{Solusi yang disediakan oleh sumber resmi; kekosongan sumber dipertahankan.}
\begingroup
\renewcommand{\theHfakt}{o011.exam.official.09.\arabic{fakt}}
%Data institusi

%\input{Dozentdaten}

%\renewcommand{\fachbereich}{Bidang studi}

%\renewcommand{\dozent}{Prof. Dr. . }

%Data ujian

\renewcommand{\klausurgebiet}{ }

\renewcommand{\klausurtyp}{ }

\renewcommand{\klausurdatum}{ . 20}

\klausurvorspann {\fachbereich} {\klausurdatum} {\dozent} {\klausurgebiet} {\klausurtyp}

%Data untuk tabel poin berikut

\renewcommand{\aeins}{  3  }

\renewcommand{\azwei}{  3   }

\renewcommand{\adrei}{  6  }

\renewcommand{\avier}{  3  }

\renewcommand{\afuenf}{  3  }

\renewcommand{\asechs}{  5  }

\renewcommand{\asieben}{  13  }

\renewcommand{\aacht}{  4  }

\renewcommand{\aneun}{  3  }

\renewcommand{\azehn}{  7  }

\renewcommand{\aelf}{  6  }

\renewcommand{\azwoelf}{  3  }

\renewcommand{\adreizehn}{  6   }

\renewcommand{\avierzehn}{  65  }

\renewcommand{\afuenfzehn}{    }

\renewcommand{\asechzehn}{    }

\renewcommand{\asiebzehn}{    }

\renewcommand{\aachtzehn}{    }

\renewcommand{\aneunzehn}{    }

\renewcommand{\azwanzig}{    }

\renewcommand{\aeinundzwanzig}{    }

\renewcommand{\azweiundzwanzig}{    }

\renewcommand{\adreiundzwanzig}{    }

\renewcommand{\avierundzwanzig}{    }

\renewcommand{\afuenfundzwanzig}{    }

\renewcommand{\asechsundzwanzig}{    }

\punktetabelledreizehn

\klausurnote

\newpage

\setcounter{section}{0}



\inputaufgabeklausurloesung
{3}

{

Definisikan istilah-istilah berikut
\zusatzklammer {yang dicetak miring} {} {.}
\aufzaehlungsechs{Suatu
\stichwort {ruang Hausdorff} {}
$X$.

}{Suatu
\stichwort {pemetaan diferensiabel} {}
\maabbdisp {\varphi} { L } { M
} {}
antara dua manifold diferensiabel
\mathkor {} {L} {dan} {M} {.}

}{
\stichwort {Bundel tangen} {}
dari suatu manifold diferensiabel $M$.

}{
\stichwort {Matriks fundamental kedua} {}
dari parametrisasi lokal
\maabbdisp {} {V} {U \subseteq Y
} {}
\zusatzklammer {dengan

\mavergleichskette
{\vergleichskette
{ V
}
{ \subseteq }{ \R^{n-1}
}
{  }{
}
{    }{
}
{   }{
}
}
{}{}{}
terbuka} {} {}
suatu hiperpermukaan diferensiabel

\mavergleichskette
{\vergleichskette
{ Y
}
{ \subseteq }{ \R^n
}
{  }{
}
{    }{
}
{   }{
}
}
{}{}{.}

}{
\stichwort {Turunan eksterior} {} dari suatu bentuk diferensial terdiferensialkan secara kontinu
\mathl{\omega \in
{ \mathcal E }^{ k } (  M   )}{} pada manifold diferensiabel $M$.

}{Suatu \stichwort {partisi kesatuan yang tersubordinasi pada penutup} {,} dengan

\mavergleichskette
{\vergleichskette
{  X
}
{ = }{ \bigcup_{i \in I} W_i
}
{  }{
}
{    }{
}
{   }{
}
}
{}{}{}
suatu
\definitionsverweis {penutup terbuka}{}{}
dari
\definitionsverweis {ruang topologis}{}{}
$X$.
}

}
{

\aufzaehlungsechs{Suatu
\definitionsverweis {ruang topologis}{}{}
$X$ disebut Hausdorff jika untuk setiap dua titik berbeda
\mathl{x, y \in X}{} terdapat dua
\definitionsverweis {himpunan terbuka}{}{}
\mathkor {} {U} {dan} {V} {}
yang memenuhi
\mathl{x \in U, \, y \in V}{} dan
\mathl{U \cap V = \emptyset}{.}
}{Misalkan
\mathkor {} {(U_i,U_i', \alpha_i, i \in I)} {dan} {(V_j,V_j', \beta_j, j \in J)} {}
adalah
\definitionsverweis {atlas}{}{}
dari
\mathkor {} {M} {dan} {N} {.}
Pemetaan
\maabbdisp {\varphi} { L } { M
} {}
disebut diferensiabel jika kontinu dan jika untuk setiap
\mathkor {} {i \in I} {dan setiap} {j \in J} {}
pemetaan-pemetaan
\maabbdisp {\beta_j \circ  \varphi \circ (\alpha_i)^{-1}} {\alpha_i( \varphi^{-1}(V_j) \cap U_i) } { V_j'
} {}
\definitionsverweis {terdiferensialkan secara kontinu}{}{.}
}{Himpunan

\mavergleichskettedisp
{\vergleichskette
{TM
}
{ =} {\biguplus_{P \in M} T_PM
}
{ } {
}
{ } {
}
{ } {
}
}
{}{}{,}
yang dilengkapi dengan pemetaan proyeksi
\maabbeledisp {\pi} {TM} {M
} {(P,v)} {P
} {,}
disebut bundel tangen dari $M$.
}{Misalkan $Y$
\definitionsverweis {diorientasikan}{}{}
oleh suatu
\definitionsverweis {medan normal satuan}{}{}
$N$, dengan $N$ dipandang sebagai medan pada $V$. Definisikan

\mavergleichskettedisp
{\vergleichskette
{ h_{i j}
}
{ =} {  \left\langle  \partial_i \partial_j \varphi , N  \right\rangle
}
{ } {
}
{ } {
}
{ } {
}
}
{}{}{.}
Matriks fundamental kedua
dari $\varphi$ adalah matriks
\zusatzklammer {yang bergantung pada

\mavergleichskettek
{\vergleichskettek
{ Q
}
{ \in }{ V
}
{  }{
}
{    }{
}
{   }{
}
}
{}{}{}} {} {}

\mavergleichskettedisp
{\vergleichskette
{ H
}
{ =} {  \left(  h_{ij}   \right)_{1 \leq i,j \leq n-1}
}
{ } {
}
{ } {
}
{ } {
}
}
{}{}{.}
}{Turunan eksterior dari $\omega$, secara lokal pada suatu peta tempat $\omega$ berbentuk

\mavergleichskettedisp
{\vergleichskette
{ \omega
}
{ =} {\sum_{ { \# \left(   I  \right) } = k , \, I \subseteq \{1 , \ldots ,  \operatorname{dim}\, M \}  }f_I dx_I
}
{ } {
}
{ } {
}
{ } {
}
}
{}{}{}
didefinisikan oleh

\mavergleichskettedisp
{\vergleichskette
{  d \omega
}
{ =} {\sum_{ { \# \left(   I  \right) } = k   } d(f_I ) \wedge dx_I
}
{ } {
}
{ } {
}
{ } {
}
}
{}{}{}
rumus di atas.
}{Suatu keluarga fungsi
\maabbdisp {h_j} {X} {\R
} {}
dengan

\mavergleichskette
{\vergleichskette
{ j
}
{ \in }{ J
}
{  }{
}
{    }{
}
{   }{
}
}
{}{}{}
disebut partisi kesatuan yang tersubordinasi pada penutup jika sifat-sifat berikut berlaku.
\aufzaehlungvier{Berlaku

\mavergleichskette
{\vergleichskette
{ h_j(X)
}
{ \subseteq }{ [0,1]
}
{  }{
}
{    }{
}
{   }{
}
}
{}{}{}
untuk setiap

\mavergleichskette
{\vergleichskette
{ j
}
{ \in }{ J
}
{  }{
}
{    }{
}
{   }{
}
}
{}{}{.}
}{Setiap titik

\mavergleichskette
{\vergleichskette
{ P
}
{ \in }{ X
}
{  }{
}
{    }{
}
{   }{
}
}
{}{}{}
memiliki suatu
\definitionsverweis {lingkungan terbuka}{}{}

\mavergleichskette
{\vergleichskette
{ P
}
{ \in }{  U
}
{  }{
}
{    }{
}
{   }{
}
}
{}{}{}
sedemikian sehingga
\definitionsverweis {fungsi-fungsi yang dibatasi}{}{}

\mathl{h_j  \vert_{ U }}{} merupakan fungsi nol, kecuali untuk sejumlah berhingga indeks.
}{Berlaku

\mavergleichskette
{\vergleichskette
{  \sum_{j \in J} h_j
}
{ = }{ 1
}
{  }{
}
{    }{
}
{   }{
}
}
{}{}{.}
}{Untuk setiap

\mavergleichskette
{\vergleichskette
{ j
}
{ \in }{ J
}
{  }{
}
{    }{
}
{   }{
}
}
{}{}{}
terdapat suatu himpunan terbuka
\mathl{W_{i(j)}}{} dari penutup sedemikian sehingga
\definitionsverweis {tumpuan}{}{}
dari $h_j$ termuat dalam
\mathl{W_{i(j)}}{.}
}
}

 }



\inputaufgabeklausurloesung
{3}

{

Nyatakan teorema-teorema berikut.
\aufzaehlungdrei{
\stichwort {Teorema tentang sifat aljabar pemetaan Weingarten} {.}}{Teorema tentang keterorientasian dan bentuk volume positif.}{
\stichwort {Teorema titik tetap Brouwer} {.}}

}
{

\aufzaehlungdrei{\faktsituation {Misalkan

\mavergleichskette
{\vergleichskette
{ W
}
{ \subseteq }{ \R^n
}
{  }{
}
{    }{
}
{   }{
}
}
{}{}{}
\definitionsverweis {terbuka}{}{,}
\maabb {h} { W } { \R
} {}
adalah suatu
\definitionsverweis {fungsi yang dua kali terdiferensialkan secara kontinu}{}{}
dan

\mavergleichskette
{\vergleichskette
{ Y
}
{ = }{ h^{-1}(c)
}
{  }{
}
{    }{
}
{   }{
}
}
{}{}{}
adalah serat di atas

\mavergleichskette
{\vergleichskette
{ c
}
{ \in }{ \R
}
{  }{
}
{    }{
}
{   }{
}
}
{}{}{,}
dengan $h$
\definitionsverweis {reguler}{}{}
di setiap titik $Y$.}
\faktfolgerung {Maka
\definitionsverweis {pemetaan Weingarten}{}{}
$L_P$ di setiap titik

\mavergleichskette
{\vergleichskette
{ P
}
{ \in }{ Y
}
{  }{
}
{    }{
}
{   }{
}
}
{}{}{}
bersifat
\definitionsverweis {adjoin-diri}{}{.}}
\faktzusatz {}
\faktzusatz {}}{\faktsituation {Misalkan $M$ adalah suatu
\definitionsverweis {manifold diferensiabel
}{}{}
dengan suatu
\definitionsverweis {basis topologi terhitung}{}{.}}
\faktfolgerung {Terdapat suatu
\definitionsverweis {bentuk volume}{}{}
\definitionsverweis {kontinu}{}{}
tanpa titik nol pada $M$ jika dan hanya jika $M$
\definitionsverweis {dapat diorientasikan}{}{.}}
\faktzusatz {Bentuk volume ini juga dapat dipilih terdiferensialkan secara kontinu dan positif.}
\faktzusatz {}}{Misalkan
\maabbdisp {\psi} { B  \left( 0,r \right) } { B  \left( 0,r \right)
} {}
adalah suatu pemetaan yang terdiferensialkan secara kontinu dari bola tertutup di $\R^n$ ke dirinya sendiri. Maka $\psi$ memiliki suatu titik tetap.}

 }



\inputaufgabeklausurloesung
{6}

{

Tentukan
\definitionsverweis {kelengkungan}{}{}
di setiap titik lingkaran satuan yang diberikan secara implisit oleh

\mavergleichskettedisp
{\vergleichskette
{ h(x,y)
}
{ =} { x^2+y^2
}
{ =} { 1
}
{ } {
}
{ } {
}
}
{}{}{}
dengan
Lemma 3.11 (Geometri Diferensial (Osnabrück 2023))
dan medan gradien $h$.

}
{

\definitionsverweis {Medan gradien}{}{}
dari $h$ adalah $\begin{pmatrix}  2x \\ 2y   \end{pmatrix}$, sehingga suatu medan normal satuan adalah

\mavergleichskettedisp
{\vergleichskette
{ N \begin{pmatrix}  x  \\ y   \end{pmatrix}
}
{ =} { { \frac{   1 }{  \sqrt{x^2+y^2}  } } \begin{pmatrix}  x  \\ y   \end{pmatrix}
}
{ } {
}
{ } {
}
{ } {
}
}
{}{}{.}
\definitionsverweis {Diferensial total}{}{}
dari $N$ di suatu titik

\mavergleichskettedisp
{\vergleichskette
{ P
}
{ =} { \begin{pmatrix}  x  \\ y   \end{pmatrix}
}
{ } {
}
{ } {
}
{ } {
}
}
{}{}{}
adalah

\mavergleichskettedisp
{\vergleichskette
{  \left(D N \right)_{P}
}
{ =} { \begin{pmatrix}  { \frac{   y^2 }{   \left(  x^2+y^2  \right)^{3/2}  } }   &   - { \frac{   xy }{   \left(  x^2+y^2  \right)^{3/2}  } }   \\    - { \frac{   xy }{   \left(  x^2+y^2  \right)^{3/2}  } }    &  { \frac{   x^2 }{   \left(  x^2+y^2  \right)^{3/2}  } }  \end{pmatrix}
}
{ } {
}
{ } {
}
{ } {
}
}
{}{}{.}
Jika diterapkan pada vektor singgung
\mathl{\begin{pmatrix}  y  \\ -x   \end{pmatrix}}{}, diperoleh

\mavergleichskettealign
{\vergleichskettealign
{ \left(D N \right)_{P} \begin{pmatrix}  y  \\ -x   \end{pmatrix}
}
{ =} { \begin{pmatrix}  y^2    &   - xy    \\  - xy   &  x^2  \end{pmatrix} \begin{pmatrix}  y  \\ -x   \end{pmatrix}
}
{ =} { \begin{pmatrix}  y \left(  y^2 +x^2  \right)  \\ -x \left(  y^2 +x^2  \right)    \end{pmatrix}
}
{ =} { \begin{pmatrix}  y   \\ -x    \end{pmatrix}
}
{ } {
}
}
{}
{}{.}
Oleh karena itu, menurut
Lemma 3.11 (Differentialgeometrie (Osnabrück 2023))
kelengkungan kurva di $P$ sama dengan $-1$.
\zusatzklammer {Hasil ini sesuai karena
\mathl{\begin{pmatrix}  y  \\ -x   \end{pmatrix}}{} dan medan gradien merepresentasikan orientasi standar, sedangkan
\mathl{\begin{pmatrix}  y  \\ -x   \end{pmatrix}}{} adalah arah singgung untuk gerak searah jarum jam.} {} {}

 }



\inputaufgabeklausurloesung
{3}

{

Misalkan $X$ suatu
\definitionsverweis {ruang Hausdorff}{}{.}
Tunjukkan bahwa
\definitionsverweis {diagonal}{}{}

\mavergleichskettedisp
{\vergleichskette
{ \triangle
}
{ =} { \left\{ (x,y) \in X \times X  \mid   x = y        \right\}
}
{ } {
}
{ } {
}
{ } {
}
}
{}{}{}
merupakan
\definitionsverweis {subhimpunan tertutup}{}{}
dalam
\definitionsverweis {ruang produk}{}{}

\mathl{X \times X}{}.

}
{

Kita harus menunjukkan bahwa komplemen

\mavergleichskettedisp
{\vergleichskette
{W
}
{ =} { X \times X \setminus \triangle
}
{ =} { \left\{  (x,y)   \mid   x,y \in X , \,  x \neq y       \right\}
}
{ } {
}
{ } {
}
}
{}{}{}
terbuka. Jadi, misalkan
\mathl{(x,y)}{} adalah suatu pasangan dengan

\mavergleichskette
{\vergleichskette
{ x
}
{ \neq }{ y
}
{  }{
}
{    }{
}
{   }{
}
}
{}{}{.}
Berdasarkan sifat Hausdorff yang diasumsikan, terdapat himpunan-himpunan terbuka saling lepas
\mathl{U,V \subseteq X}{} dengan
\mathl{x \in U}{} dan
\mathl{y \in V}{.} Berlaku
\mathl{(x,y) \in U \times V}{,} dan menurut definisi topologi produk,
\mathl{U \times V}{} adalah suatu himpunan terbuka dalam
\mathl{X \times X}{.} Karena keduanya saling lepas, dari
\mathl{(z,w) \in U\times V}{} langsung diperoleh

\mavergleichskette
{\vergleichskette
{z
}
{ \neq }{w
}
{  }{
}
{    }{
}
{   }{
}
}
{}{}{.}
Jadi,

\mavergleichskettedisp
{\vergleichskette
{(x,y)
}
{ \in} { U \times V
}
{ \subseteq} {  W
}
{ } {
}
{ } {
}
}
{}{}{}
dan $W$ merupakan gabungan himpunan-himpunan produk terbuka seperti itu, sehingga himpunan tersebut sendiri terbuka.

 }



\inputaufgabeklausurloesung
{3}

{

Tunjukkan bahwa ketiga manifold satu dimensi
\mathlistdisp {S^1 = \left\{ (x,y) \in \R^2  \mid   \betrag { (x,y) } = 1         \right\}} {} {\R} {dan} {\R \setminus \{0\}} {}
tidak ada dua di antaranya yang homeomorfik.

}
{

Sebagai himpunan bagian tertutup dan terbatas dari $\R^2$, lingkaran satuan bersifat kompak
\zusatzklammer {Satz von Heine-Borel} {} {.} Sebaliknya, garis real $\R$ dan $\R \setminus \{0\}$ tidak kompak karena keduanya merupakan himpunan bagian tak terbatas dari $\R$. Jadi,
\mathl{S^1 \not \cong \R,\, \R \setminus \{0\}}{.}

Garis real terhubung, sebagaimana mengikuti dari
Zwischenwertsatz
. Sebaliknya, $\R \setminus \{0\}$ tidak terhubung karena dapat ditulis

\mavergleichskettedisp
{\vergleichskette
{\R_+
}
{ =} {(\R \setminus \{0\}) \cap [0, \infty]
}
{ =} {(\R \setminus \{0\}) \cap \, ]0, \infty]
}
{ } {
}
{ } {
}
}
{}{}{}
sehingga diperoleh suatu himpunan bagian tak kosong yang sekaligus terbuka dan tertutup. Jadi,
\mathl{\R \,\not \cong \R \setminus \{0\}}{.}

 }



\inputaufgabeklausurloesung
{5}

{

Misalkan

\mavergleichskette
{\vergleichskette
{ G
}
{ \subseteq }{ \R^m
}
{  }{
}
{    }{
}
{   }{
}
}
{}{}{}
\definitionsverweis {terbuka}{}{}
dan misalkan

\mavergleichskette
{\vergleichskette
{ M
}
{ \subseteq }{ G
}
{  }{
}
{    }{
}
{   }{
}
}
{}{}{}
suatu
$n$-\definitionsverweis {dimensional}{}{}
\definitionsverweis {submanifold tertutup}{}{,}
yang
\definitionsverweis {berorientasi}{}{}
dan dilengkapi dengan struktur Riemann terinduksi serta
\definitionsverweis {bentuk volume kanonik}{}{}
$\omega$. Misalkan

\mavergleichskette
{\vergleichskette
{ W
}
{ \subseteq }{ \R^n
}
{  }{
}
{    }{
}
{   }{
}
}
{}{}{}
terbuka dan misalkan
\maabbdisp {\varphi} { W } { U
} {}
suatu
\definitionsverweis {difeomorfisme}{}{}
dengan himpunan terbuka

\mavergleichskette
{\vergleichskette
{ U
}
{ \subseteq }{ M
}
{  }{
}
{    }{
}
{   }{
}
}
{}{}.
Tunjukkan bahwa pada $W$ berlaku rumus

\mavergleichskettealignhandlinks
{\vergleichskettealignhandlinks
{\varphi^* \left(  \omega \vert_U  \right)
}
{ =} { \left( \det  \left( \left\langle  \begin{pmatrix}  { \frac{   \partial  \varphi_1   }{  \partial   x_i   } }  \\ \vdots \\  { \frac{   \partial  \varphi_m   }{  \partial   x_i   } }   \end{pmatrix}  ,  \begin{pmatrix}  { \frac{   \partial  \varphi_1   }{  \partial   x_j   } }  \\ \vdots \\  { \frac{   \partial  \varphi_m   }{  \partial   x_j   } }   \end{pmatrix}   \right\rangle \right)_{1 \leq i,j \leq n} \right)^{1/2}  dx_1 \wedge \ldots \wedge dx_n
}
{ =} { \left( \det  \left(  { \frac{   \partial  \varphi_1   }{  \partial   x_i   } } { \frac{   \partial  \varphi_1   }{  \partial   x_j   } }  + \cdots +  { \frac{   \partial  \varphi_m   }{  \partial   x_i   } } { \frac{   \partial  \varphi_m   }{  \partial   x_j   } }  \right)_{1 \leq i,j \leq n} \right)^{1/2}  dx_1 \wedge \ldots \wedge dx_n
}
{ } {
}
{ } {
}
}
{}
{}{}{.}

}
{

Persamaan kedua langsung diperoleh dengan mengevaluasi hasil kali dalam standar pada $\R^m$. Menurut definisi
\definitionsverweis {matriks fundamental metrik}{}{,}
untuk

\mavergleichskette
{\vergleichskette
{ Q
}
{ \in }{ W
}
{  }{
}
{    }{
}
{   }{
}
}
{}{}{}

\mavergleichskettealign
{\vergleichskettealign
{ g_{ij} (Q)
}
{ =} { \left\langle  T_Q(\varphi)(e_i)  ,  T_Q(\varphi)(e_j)   \right\rangle_{\varphi(Q)}
}
{ =} { \left\langle  T_Q(\varphi)(e_i)  ,  T_Q(\varphi)(e_j)   \right\rangle
}
{ =} { \left\langle  \left(  D \varphi  \right)_{ Q } \left(   e_i   \right)  ,  \left(  D\varphi  \right)_{ Q } \left(   e_j   \right)   \right\rangle
}
{ =} { \left\langle  \begin{pmatrix}  { \frac{   \partial  \varphi_1   }{  \partial   x_i   } } ( Q)  \\ \vdots \\  { \frac{   \partial  \varphi_m   }{  \partial   x_i   } } ( Q)   \end{pmatrix}  ,  \begin{pmatrix}  { \frac{   \partial   \varphi_1   }{  \partial   x_j   } } ( Q )  \\ \vdots \\  { \frac{   \partial  \varphi_m   }{  \partial   x_j   } } ( Q)   \end{pmatrix}   \right\rangle
}
}
{}
{}{,}
karena ruang singgung
\mathl{T_{\varphi(Q)}M}{} dilengkapi dengan hasil kali dalam yang diinduksi dari $\R^m$, pemetaan singgung dalam kasus lokal adalah diferensial total, dan entri-entrinya dapat dinyatakan dengan turunan parsial. Oleh karena itu, pernyataan tersebut mengikuti dari
Lemma 16.7 (Differentialgeometrie (Osnabrück 2023)).

 }



\inputaufgabeklausurloesung
{weiter}

{

Kita tinjau

\mavergleichskette
{\vergleichskette
{ \R^6
}
{ = }{ \left(  \R^2  \right)^3
}
{  }{
}
{    }{
}
{   }{
}
}
{}{}{}
sebagai himpunan semua segitiga
\zusatzklammer {termasuk yang degenerat} {} {}
dengan mengidentifikasi segitiga yang mempunyai titik-titik sudut
\zusatzklammer {terurut} {} {}

\mavergleichskette
{\vergleichskette
{ P_1
}
{ = }{ \begin{pmatrix}  x_1  \\ y_1    \end{pmatrix}
}
{  }{
}
{    }{
}
{   }{
}
}
{}{}{,}

\mavergleichskette
{\vergleichskette
{ P_2
}
{ = }{ \begin{pmatrix}  x_2  \\ y_2    \end{pmatrix}
}
{  }{
}
{    }{
}
{   }{
}
}
{}{}{}
dan

\mavergleichskette
{\vergleichskette
{ P_3
}
{ = }{ \begin{pmatrix}  x_3  \\ y_3    \end{pmatrix}
}
{  }{
}
{    }{
}
{   }{
}
}
{}{}{,}
dengan tupel koordinat

\mathdisp {(x_1,y_1,x_2,y_2,x_3,y_3)} {  }

.
\aufzaehlungsechs{Tunjukkan bahwa himpunan segitiga dengan dua titik sudut berimpit merupakan subhimpunan tertutup $\R^6$
\zusatzklammer {jadi komplemennya merupakan himpunan terbuka dalam $\R^6$ yang kita namai $U$} {} {.}
}{Tunjukkan bahwa himpunan segitiga dengan ketiga titik sudut segaris merupakan subhimpunan tertutup $\R^6$
\zusatzklammer {jadi komplemennya, yang terdiri atas semua segitiga tak degenerat, merupakan himpunan terbuka dalam $\R^6$ yang kita namai $V$} {} {.}
}{Buat aturan fungsi yang mendeskripsikan pemetaan
\maabbdisp {F} {\R^6} {\R
} {}
yang memetakan suatu segitiga ke kelilingnya.
}{Tunjukkan bahwa fungsi $F$ dari bagian (3) terdiferensialkan secara kontinu pada himpunan $U$.
}{Hitung turunan parsial $F$ terhadap $x_1$ pada $U$.
}{Tunjukkan bahwa himpunan segitiga tak degenerat berkeliling $1$ membentuk
\definitionsverweis {submanifold tertutup}{}{}
dari $V$. Berapakah dimensinya?
}

}
{

\aufzaehlungsechs{Titik-titik sudut
\mathkor {} {P_1} {dan} {P_2} {}
berimpit jika dan hanya jika kedua koordinatnya sama, yakni jika

\mavergleichskette
{\vergleichskette
{ x_1
}
{ = }{ x_2
}
{  }{
}
{    }{
}
{   }{
}
}
{}{}{}
dan

\mavergleichskette
{\vergleichskette
{ y_1
}
{ = }{ y_2
}
{  }{
}
{    }{
}
{   }{
}
}
{}{}{}
berlaku. Syarat seperti ini mendefinisikan suatu himpunan bagian tertutup karena merupakan serat dari suatu pemetaan linear
\zusatzklammer {kontinu} {} {.}
Karena gabungan berhingga dari himpunan-himpunan bagian tertutup kembali tertutup, himpunan yang dimaksud bersifat tertutup.
}{Tinjau vektor-vektor penghubung

\mathdisp {P_2-P_1 = \begin{pmatrix}  x_2-x_1  \\ y_2-y_1   \end{pmatrix} \text{   und } P_3-P_1 = \begin{pmatrix}  x_3-x_1  \\ y_3-y_1   \end{pmatrix}} {  }

Ketiga titik kolinear jika dan hanya jika kedua vektor tersebut bergantung linear, dan hal ini berlaku jika dan hanya jika determinannya sama dengan $0$. Dengan demikian, syarat kolinearitasnya adalah

\mavergleichskettedisp
{\vergleichskette
{ H(x_1,x_2,x_3,y_1,y_2,y_3)
}
{ =} { (x_2-x_1) (y_3-y_1) - (x_3-x_1) (y_2-y_1)
}
{ =} { 0
}
{ } {
}
{ } {
}
}
{}{}{.}
Karena $H$ adalah pemetaan polinomial dan karenanya kontinu, serat di atas $0$ merupakan himpunan bagian tertutup.
}{Karena keliling merupakan jumlah dari ketiga sisi segitiga, berlaku

\mavergleichskettealigndrucklinks
{\vergleichskettealigndrucklinks
{ F(x_1,x_2,x_3,y_1,y_2,y_3)
}
{ =} { d(P_1,P_2) + d(P_1,P_3) +d(P_2,P_3)
}
{ =} { \sqrt{(x_2-x_1)^2 + (y_2-y_1)^2} +  \sqrt{(x_3-x_1)^2 + (y_3-y_1)^2} +  \sqrt{(x_3-x_2)^2 + (y_3-y_2)^2}
}
{ } {
}
{ } {
}
}
{}{}{.}
}{Untuk suatu titik di
\mathl{U}{,} setiap radikan, sebagai fungsi polinomial, terdiferensialkan secara kontinu dan bernilai positif ketat. Oleh karena itu, akar kuadratnya juga terdiferensialkan secara kontinu.
}{Turunan parsialnya adalah

\mavergleichskettealigndrucklinks
{\vergleichskettealigndrucklinks
{ { \frac{   \partial   F   }{  \partial   x_1  } }
}
{ =} { { \frac{   1 }{  2 } }   \left(  (  x_2-x_1  )^2 + (y_2-y_1)^2  \right)^{-1/2} (2x_1-2x_2)    + { \frac{   1 }{  2 } }   \left(  (  x_3-x_1  )^2 + (y_3-y_1)^2  \right)^{-1/2} (2x_1-2x_3)
}
{ =} {  \left(  (  x_2-x_1  )^2 + (y_2-y_1)^2  \right)^{-1/2} (x_1-x_2)    +   \left(  (  x_3-x_1  )^2 + (y_3-y_1)^2  \right)^{-1/2} (x_1-x_3)
}
{ } {
}
{ } {
}
}
{}{}{.}
}{Pada $U$, semua radikan bernilai positif, sehingga turunan parsial terhadap $x_1$ hanya dapat bernilai nol jika
\mathl{x_1=x_2=x_3}{.} Namun, segitiga yang memenuhi syarat ini bersifat kolinear. Ini berarti bahwa pada $V$ turunan parsial terhadap $x_1$ tidak pernah bernilai $0$; khususnya, gradien pada $V$ tidak pernah lenyap, sehingga bersifat reguler. Menurut
Satz 8.2 (Differentialgeometrie (Osnabrück 2023))
serat tersebut merupakan submanifold tertutup, dengan dimensi

\mavergleichskettedisp
{\vergleichskette
{ 6-1
}
{ =} { 5
}
{ } {
}
{ } {
}
{ } {
}
}
{}{}{.}
}

 }



\inputaufgabeklausurloesung
{4}

{

Misalkan
\mathkor {} {r} {dan} {R} {}
bilangan-bilangan real positif dengan

\mavergleichskette
{\vergleichskette
{ 0
}
{ < }{ r
}
{ < }{ R
}
{    }{
}
{   }{
}
}
{}{}{,}
dan tinjau torus
\zusatzklammer {tertanam} {} {}

\mavergleichskettedisp
{\vergleichskette
{ T
}
{ =} { \left\{ (x,y,z) \in \R^3  \mid   \left(  \sqrt{x^2+y^2} -R  \right)^2 +z^2  = r^2         \right\}
}
{ \subseteq} { \R^3
}
{ } {
}
{ } {
}
}
{}{}{}
dengan
\definitionsverweis {struktur Riemann terinduksi}{}{.}
Tunjukkan bahwa untuk setiap

\mavergleichskette
{\vergleichskette
{ \alpha
}
{ \in }{ [0,2 \pi [
}
{  }{
}
{    }{
}
{   }{
}
}
{}{}{}
pemetaan
\maabbeledisp {\Psi_\alpha} {\R^3} {\R^3
} { \left(  x   , \,   y  , \,  z                 \right) } { \left(  x \cos \alpha - y \sin  \alpha   , \,   x \sin \alpha + y \cos  \alpha  , \,  z                 \right)
} {,}
menginduksi suatu
\definitionsverweis {isometri}{}{}
dari $T$ ke dirinya sendiri.

}
{

Pada

\mavergleichskettedisp
{\vergleichskette
{ \R^3
}
{ =} { \R^2 \times \R
}
{ } {
}
{ } {
}
{ } {
}
}
{}{}{}
pemetaan tersebut merupakan produk suatu
\definitionsverweis {rotasi bidang}{}{}
dengan identitas. Oleh karena itu, pemetaan tersebut adalah
\definitionsverweis {isometri}{}{}
Euklides pada $\R^3$. Misalkan

\mavergleichskette
{\vergleichskette
{ (x,y,z)
}
{ \in }{  T
}
{  }{
}
{    }{
}
{   }{
}
}
{}{}{}
dan

\mavergleichskettedisp
{\vergleichskette
{ \Psi_\alpha(x,y,z)
}
{ =} { ( u,v,z)
}
{ } {
}
{ } {
}
{ } {
}
}
{}{}{.}
Maka

\mavergleichskettealigndrucklinks
{\vergleichskettealigndrucklinks
{ \left(  \sqrt{u^2+v^2} -R  \right)^2 +z^2
}
{ =} { \left(  \sqrt{ \left(  x \cos \alpha - y \sin  \alpha   \right)^2+ \left(  x \sin \alpha + y \cos  \alpha  \right)^2} -R  \right)^2 +z^2
}
{ =} { \left(  \sqrt{  x^2 \cos^{ 2  } \alpha + y^2 \sin^{ 2  }  \alpha  -2xy \cos \alpha \sin  \alpha    + x^2 \sin^{ 2  } \alpha + y^2 \cos^{ 2  }  \alpha  +2xy \cos \alpha \sin  \alpha     } -R  \right)^2 +z^2
}
{ =} { \left(  \sqrt{  x^2   + y^2 } -R  \right)^2 +z^2
}
{ =} { r^2
}
}
{}{}{,}
sehingga

\mavergleichskette
{\vergleichskette
{ \Psi_\alpha(x,y,z)
}
{ \in }{  T
}
{  }{
}
{    }{
}
{   }{
}
}
{}{}{.}
Jadi,
\maabbdisp {\Psi_\alpha} { T} {T
} {}
adalah suatu difeomorfisme dan isometri karena $T$ dilengkapi dengan struktur Riemann terinduksi.

 }



\inputaufgabeklausurloesung
{3}

{

Deskripsikan suatu model untuk bidang hiperbolik.

}
{  }



\inputaufgabeklausurloesung
{7}

{

Buktikan Theorema egregium.

}
{

Kita mendiferensialkan persamaan penentu pertama untuk simbol-simbol Christoffel, yaitu

\mavergleichskettedisp
{\vergleichskette
{ \partial_1 \partial_1 \varphi
}
{ =} { \Gamma^1_{11}  \partial_1 \varphi  + \Gamma^2_{11}  \partial_2 \varphi +h_{11} N
}
{ } {
}
{ } {
}
{ } {
}
}
{}{}{,}
terhadap variabel kedua, serta persamaan penentu kedua, yaitu

\mavergleichskettedisp
{\vergleichskette
{ \partial_1 \partial_2 \varphi
}
{ =} { \Gamma^1_{12}  \partial_1 \varphi + \Gamma^2_{12}  \partial_2  \varphi + h_{12} N
}
{ } {
}
{ } {
}
{ } {
}
}
{}{}{,}
terhadap variabel pertama. Berdasarkan teorema Schwarz, diperoleh

\mavergleichskettealigndrucklinks
{\vergleichskettealigndrucklinks
{ \left(  \partial_2 \Gamma^1_{11}  \right) \partial_1 \varphi +  \left(  \partial_2 \Gamma^2_{11}  \right)  \partial_2 \varphi+ \left(  \partial_2 h_{11}  \right)  N + \Gamma^1_{11} \partial_2 \partial_1 \varphi + \Gamma^2_{11}  \partial_2 \partial_2 \varphi  +h_{11} \partial_2 N
}
{ =} { \left(  \partial_1 \Gamma^1_{12}  \right)  \partial_1  \varphi+ \left(  \partial_1 \Gamma^2_{12}  \right)\partial_2 \varphi + \left(  \partial_1 h_{12}  \right)  N+ \Gamma^1_{12}  \partial_1 \partial_1 \varphi + \Gamma^2_{12}  \partial_1 \partial_2 \varphi + h_{12} \partial_1 N
}
{ } {
}
{ } {
}
{ } {
}
}
{}{}{.}
Selisih kedua ekspresi ini adalah $0$, dan kita menentukan bagian yang berkaitan dengan vektor basis $\partial_2 \varphi$. Untuk itu, kita menggunakan persamaan-persamaan penentu bagi simbol Christoffel dan
Lemma 19.3 (Differentialgeometrie (Osnabrück 2023))  (3)
, lalu memperoleh

\mavergleichskettedisphandlinks
{\vergleichskettedisphandlinks
{ 0
}
{ =} { \partial_2 \Gamma^2_{11} - \partial_1 \Gamma^2_{12} +\Gamma^1_{11} \Gamma^2_{12} - \Gamma^1_{12} \Gamma^2_{11} + \Gamma_{11}^2 \Gamma_{22}^2- \left(  \Gamma_{12}^2  \right)^2 + h_{11}  { \frac{   g_{12} h_{12} - g_{11} h_{22}   }{  g_{11}g_{22} -g_{12}^2 } } - h_{12}  { \frac{   g_{12} h_{11} - g_{11}h_{12}  }{  g_{11}g_{22} -g_{12}^2 } }
}
{ } {
}
{ } {
}
{ } {
}
}
{}{}{.}
Dengan

\mavergleichskettedisphandlinks
{\vergleichskettedisphandlinks
{ - g_{11} \left(  h_{11}h_{22} -h_{12}^2  \right)
}
{ =} { h_{11} \left(  g_{12} h_{12}- g_{11} h_{22}  \right) - h_{12} \left(  g_{12} h_{11} - g_{11}h_{12}  \right)
}
{ } {
}
{ } {
}
{ } {
}
}
{}{}{}
kita dapat mengganti kedua suku terakhir. Selanjutnya, menurut
Lemma 19.3 (Differentialgeometrie (Osnabrück 2023))  (4)
, diperoleh

\mavergleichskettealignhandlinks
{\vergleichskettealignhandlinks
{ g_{11} K
}
{ =} { g_{11} { \frac{   h_{11}h_{22} -h_{12}^2  }{ g_{11}g_{22} -g_{12}^2 } }
}
{ =} { \left(  \partial_2 \Gamma_{11}^2 - \partial_1 \Gamma_{12}^2 + \Gamma^1_{11} \Gamma^2_{12}  + \Gamma_{11}^2 \Gamma_{22}^2 - \Gamma_{12}^1 \Gamma_{11}^2-  \left(  \Gamma_{12}^2  \right)^2  \right)
}
{ } {
}
{ } {
}
}
{}
{}{.}

 }



\inputaufgabeklausurloesung
{6 (1+2+2+1)}

{

Kita tinjau
\definitionsverweis {bentuk diferensial}{}{}

\mavergleichskettedisp
{\vergleichskette
{ \omega
}
{ =} { ydx+zdy +xdz
}
{ } {
}
{ } {
}
{ } {
}
}
{}{}{}
pada $\R^3$ dan pemetaan
\maabbeledisp {\varphi} {\R^2} {\R^3
} { \left(  u   , \,  v                   \right) } { \left( u^2  , \,  v^2  , \,  uv                 \right)
} {.}
\aufzaehlungvier{Hitung turunan eksterior $\omega$.
}{Hitung tarikan balik $\varphi^* \omega$ dari $\omega$ di bawah $\varphi$.
}{Hitung turunan eksterior $\varphi^* \omega$ pada $\R^2$.
}{Hitung tarikan balik $\varphi^* d \omega$ dari $d\omega$ di bawah $\varphi$ secara independen dari bagian (3).
}

}
{

\aufzaehlungvier{Berlaku

\mavergleichskettedisp
{\vergleichskette
{ d \omega
}
{ =} { d \left(  ydx+zdy +xdz  \right)
}
{ =} { dy \wedge dx +dz \wedge dy +dx \wedge dz
}
{ =} { - dx \wedge dy -dy \wedge dz +dx \wedge dz
}
{ } {
}
}
{}{}{.}
}{Berlaku

\mavergleichskettealign
{\vergleichskettealign
{  \varphi^* \omega
}
{ =} { v^2 d u^2 + uvdv^2+ u^2duv
}
{ =} { 2v^2udu +2uv^2 dv + u^2(udv +vdu)
}
{ =} { \left(  2uv^2 +u^2v  \right) du + \left(  2uv^2 +u^3  \right) dv
}
{ } {
}
}
{}
{}{.}
}{Berlaku

\mavergleichskettealign
{\vergleichskettealign
{  d \left(   \left(  2uv^2 +u^2v  \right) du + \left(  2uv^2 +u^3  \right) dv     \right)
}
{ =} {  d \left(  2uv^2 +u^2v  \right) du  +  d\left(  2uv^2 +u^3  \right) dv
}
{ =} { \left(  4uv  + u^2      \right)  dv \wedge du  +   \left(  2v^2   +3u^2  \right)   du \wedge dv
}
{ =} { \left(  -4uv +  2v^2   +2u^2  \right)   du \wedge dv
}
{ } {
}
}
{}
{}{.}
}{Tarikan balik $\varphi^* d \omega$ adalah

\mavergleichskettealign
{\vergleichskettealign
{ \varphi^* \left(  - dx \wedge dy -dy \wedge dz +dx \wedge dz   \right)
}
{ =} { - d u^2 \wedge dv^2 -d v^2 \wedge duv +d u^2 \wedge duv
}
{ =} { -4uv  du \wedge dv -2v^2 dv \wedge du + 2   u^2 du \wedge dv
}
{ =} { \left(  -4uv+2v^2 + 2 u^2  \right)  du \wedge dv
}
{ } {
}
}
{}
{}{.}
}

 }



\inputaufgabeklausurloesung
{3}

{

Misalkan $M$ suatu
\definitionsverweis {manifold diferensiabel}{}{}
dan $\omega$ suatu
\definitionsverweis {bentuk diferensial}{}{}
yang
\definitionsverweis {eksak}{}{}
pada $M$. Misalkan $S$ suatu
\definitionsverweis {manifold diferensiabel}{}{}
yang
\definitionsverweis {kompak}{}{}
dan
\definitionsverweis {berorientasi}{}{}
\zusatzklammer {tanpa batas} {} {}
dengan
\definitionsverweis {basis terhitung bagi topologinya}{}{}
dan misalkan
\maabbdisp {\varphi} { S } { M
} {}
suatu
\definitionsverweis {pemetaan terdiferensialkan secara kontinu}{}{.}
Tunjukkan bahwa

\mavergleichskettedisp
{\vergleichskette
{
\int_{  S  }   \varphi^*\omega
}
{ =} { 0
}
{ } {
}
{ } {
}
{ } {
}
}
{}{}{.}

}
{

Karena menurut
Satz 20.4 (Differentialgeometrie (Osnabrück 2023))  (5)
operasi tarikan balik bentuk kompatibel dengan turunan eksterior,
\mathl{\varphi^*\omega}{} juga eksak. Misalkan $\tau$ adalah suatu bentuk diferensiabel pada $S$ dengan

\mavergleichskettedisp
{\vergleichskette
{d \tau
}
{ =} { \varphi^* \omega
}
{ } {
}
{ } {
}
{ } {
}
}
{}{}{.}
Menurut
Satz von Stokes,
berlaku

\mavergleichskettedisp
{\vergleichskette
{ \int_S \varphi^* \omega
}
{ =} { \int_S d \tau
}
{ =} { \int_{ \partial S } \tau
}
{ =} { \int_\emptyset \tau
}
{ =} { 0
}
}
{}{}{.}

 }



\inputaufgabeklausurloesung
{6}

{

Buktikan teorema tentang retraksi ke batas pada manifold.

}
{

Batas
\mathl{\partial M}{}, menurut
Satz 21.8 (Differentialgeometrie (Osnabrück 2023))
, merupakan suatu
\definitionsverweis {manifold diferensiabel}{}{}
\definitionsverweis {berorientasi}{}{}
\zusatzklammer {tanpa batas} {} {.}
Oleh karena itu, menurut
Satz 22.11 (Differentialgeometrie (Osnabrück 2023))
, terdapat suatu
\definitionsverweis {bentuk volume positif}{}{}
\definitionsverweis {yang terdiferensialkan secara kontinu}{}{}
$\tau$ pada
\mathl{\partial M}{.} Kita mempunyai
\mathl{\int_{  \partial M  }  \tau >0}{.}
\definitionsverweis {Turunan eksterior}{}{}
dari bentuk volume $\tau$ adalah $0$. Andaikan terdapat suatu
\definitionsverweis {pemetaan yang terdiferensialkan secara kontinu}{}{}
\maabbdisp {\varphi} { M } { \partial M
} {}
dengan
\mathl{\varphi \vert_{\partial M} = \operatorname{Id}_{\partial M}}{.} Maka
\definitionsverweis {bentuk tarikan balik}{}{}

\mathl{\varphi^* \tau}{} adalah
$(n-1)$-\definitionsverweis {bentuk diferensial}{}{}
pada $M$ yang pembatasannya pada batas berimpit dengan $\tau$. Dengan menggunakan
Satz 23.2 (Differentialgeometrie (Osnabrück 2023))
dan
Satz 20.4 (Differentialgeometrie (Osnabrück 2023))  (5)
kita memperoleh

\mavergleichskettealign
{\vergleichskettealign
{
\int_{  \partial M  }  \tau
}
{ =} {
\int_{  \partial  M  }  \varphi^* \tau
}
{ =} {
\int_{   M  }  d(\varphi^* \tau)
}
{ =} {
\int_{   M  }  \varphi^* (d\tau)
}
{ =} {
\int_{   M  }  \varphi^* (0)
}
}
{
\vergleichskettefortsetzungalign
{ =} { 0
}
{ } {}
{ } {}
{ } {}
}
{}{.}
Ini merupakan suatu kontradiksi.

 }
\endgroup

\chapter{Formulir Solusi Resmi Ujian 10}
\noindent\textit{Solusi yang disediakan oleh sumber resmi; kekosongan sumber dipertahankan.}
\begingroup
\renewcommand{\theHfakt}{o011.exam.official.10.\arabic{fakt}}
%Data institusi

%\input{Dozentdaten}

%\renewcommand{\fachbereich}{Bidang studi}

%\renewcommand{\dozent}{Prof. Dr. . }

%Data ujian

\renewcommand{\klausurgebiet}{ }

\renewcommand{\klausurtyp}{ }

\renewcommand{\klausurdatum}{. 20}

\klausurvorspann {\fachbereich} {\klausurdatum} {\dozent} {\klausurgebiet} {\klausurtyp}

%Data untuk tabel poin berikut

\renewcommand{\aeins}{  3  }

\renewcommand{\azwei}{  3   }

\renewcommand{\adrei}{  6  }

\renewcommand{\avier}{  4  }

\renewcommand{\afuenf}{  5  }

\renewcommand{\asechs}{  0  }

\renewcommand{\asieben}{  5  }

\renewcommand{\aacht}{  6  }

\renewcommand{\aneun}{  6  }

\renewcommand{\azehn}{  0  }

\renewcommand{\aelf}{  6  }

\renewcommand{\azwoelf}{  12  }

\renewcommand{\adreizehn}{  0   }

\renewcommand{\avierzehn}{  9  }

\renewcommand{\afuenfzehn}{  65  }

\renewcommand{\asechzehn}{    }

\renewcommand{\asiebzehn}{    }

\renewcommand{\aachtzehn}{    }

\renewcommand{\aneunzehn}{    }

\renewcommand{\azwanzig}{    }

\renewcommand{\aeinundzwanzig}{    }

\renewcommand{\azweiundzwanzig}{    }

\renewcommand{\adreiundzwanzig}{    }

\renewcommand{\avierundzwanzig}{    }

\renewcommand{\afuenfundzwanzig}{    }

\renewcommand{\asechsundzwanzig}{    }

\punktetabellevierzehn

\klausurnote

\newpage

\setcounter{section}{0}



\inputaufgabeklausurloesung
{3}

{

Definisikan istilah-istilah berikut
\zusatzklammer {yang dicetak miring} {} {}.
\aufzaehlungsechs{Suatu
\stichwort {kelengkungan utama} {}
pada suatu permukaan diferensiabel

\mavergleichskette
{\vergleichskette
{ Y
}
{ \subseteq }{ \R^3
}
{  }{
}
{    }{
}
{   }{
}
}
{}{}{}
pada titik

\mavergleichskette
{\vergleichskette
{ P
}
{ \in }{ Y
}
{  }{
}
{    }{
}
{   }{
}
}
{}{}{.}

}{Suatu
\stichwort {difeomorfisme} {}
antara manifold diferensiabel
\mathkor {} {L} {dan} {M} {.}

}{Suatu
\stichwort {vektor singgung} {}
di titik

\mavergleichskette
{\vergleichskette
{  P
}
{ \in }{  M
}
{  }{
}
{    }{
}
{   }{
}
}
{}{}{}
dari manifold diferensiabel $M$.

}{Suatu
\stichwort {operator diferensial orde pertama} {}
pada manifold diferensiabel $M$.

}{Suatu
\stichwort {setengah ruang Euklides} {.}

}{Suatu
\definitionsverweis {penghabisan kompak}{}{}
bagi suatu
\definitionsverweis {ruang topologis}{}{}
$X$.
}

}
{

\aufzaehlungsechs{Setiap
\definitionsverweis {nilai eigen}{}{}
dari
\definitionsverweis {pemetaan Weingarten}{}{}
\maabbeledisp {L_P} {T_P Y} { T_PY
} {v} {- \left(  D_{v} N  \right) \left(  P  \right)
} {,}
disebut
\definitionswort {kelengkungan utama}{}
dari $Y$ di $P$.
}{Suatu
\definitionsverweis {homeomorfisme}{}{}
\maabbdisp {\varphi} {L} {M
} {}
disebut
\stichwortpraemath {C^1} {difeomorfisme}{,}
jika
\mathkor {} {\varphi} {maupun} {\varphi^{-1}} {}
merupakan
$C^1$-\definitionsverweis {pemetaan}{}{}.
}{Yang dimaksud dengan vektor tangen di $P$ adalah suatu
\definitionsverweis {kelas ekuivalensi}{}{}
dari
\definitionsverweis {kurva-kurva diferensiabel}{}{}
yang
\definitionsverweis {ekuivalen secara tangensial}{}{}
melalui $P$.
}{Suatu pemetaan
\maabbdisp {D} {C^1(M,\R) } {C^0(M,\R)
} {}
disebut
operator diferensial orde pertama
jika memenuhi sifat-sifat berikut.
\aufzaehlungzwei {$D$ bersifat
$\R$-\definitionsverweis {linear}{}{.}
} {Berlaku

\mavergleichskette
{\vergleichskette
{ D(fg)
}
{ = }{ fD(g)+gD(f)
}
{  }{
}
{    }{
}
{   }{
}
}
{}{}{.}
}
}{Yang dimaksud dengan setengah ruang Euklides berdimensi $n$ adalah himpunan

\mavergleichskettedisp
{\vergleichskette
{ H
}
{ =} { \left\{ x \in \R^n   \mid  x_1 \geq 0        \right\}
}
{ } {
}
{ } {
}
{ } {
}
}
{}{}{}
dengan
\definitionsverweis {topologi terinduksi}{}{.}
}{Suatu penghabisan kompak

\mathbed {A_n} {}
 {n \in \N} {}
 {} {} {} {,}
dari $X$ adalah suatu
\definitionsverweis {barisan}{}{}
dari
\definitionsverweis {subhimpunan-subhimpunan kompak}{}{}

\mavergleichskette
{\vergleichskette
{ A_n
}
{ \subseteq }{ X
}
{  }{
}
{    }{
}
{   }{
}
}
{}{}{}
dengan

\mathdisp {A_n \subseteq A_{n+1}^{o} \text{ dan } \bigcup_{n=0}^\infty A_n = X} { . }

}

}



\inputaufgabeklausurloesung
{3}

{

Nyatakan teorema-teorema berikut.
\aufzaehlungdrei{Teorema tentang
\definitionsverweis {ruang singgung}{}{}
dari suatu
\definitionsverweis {submanifold tertutup}{}{}

\mavergleichskette
{\vergleichskette
{  M
}
{ \subseteq }{  N
}
{  }{
}
{    }{
}
{   }{
}
}
{}{}{.}}{Teorema tentang bentuk eksplisit dari bentuk diferensial tarik balik $\varphi^* \omega$ oleh pemetaan diferensiabel
\maabbdisp {\varphi} { U } { V
} {}
\zusatzklammer {dengan
\definitionsverweis {himpunan bagian terbuka}{}{}
\mathkor {} {U \subseteq \R^n} {dan} {V \subseteq \R^m} {,}
yang koordinatnya masing-masing dinyatakan dengan
\mathl{x_1  , \ldots , x_n}{} dan
\mathl{y_1  , \ldots , y_m}{}} {} {}
dari suatu
$k$-\definitionsverweis {bentuk diferensial}{}{}
pada $V$ dengan representasi

\mavergleichskettedisp
{\vergleichskette
{  \omega
}
{ =} {\sum_{I,\,  { \# \left(   I  \right) } = k}  f_I  dy_I
}
{ } {
}
{ } {
}
{ } {
}
}
{}{}{,}
dengan
\maabb {f_I} { V } {\R
} {}
adalah
\definitionsverweis {fungsi}{}{.}}{Teorema tentang karakterisasi kurva geodesik terhadap koneksi Levi-Civita pada manifold Riemann $M$.}

}
{

\aufzaehlungdrei{\faktsituation {Misalkan $N$ suatu
\definitionsverweis {manifold diferensiabel}{}{}
ber-
\definitionsverweis {dimensi}{}{}
$n$ dan misalkan

\mavergleichskette
{\vergleichskette
{ M
}
{ \subseteq }{ N
}
{  }{
}
{    }{
}
{   }{
}
}
{}{}{}
suatu
\definitionsverweis {submanifold tertutup}{}{}
berdimensi $m$.}
\faktfolgerung {Maka untuk setiap titik

\mavergleichskette
{\vergleichskette
{ P
}
{ \in }{ M
}
{  }{
}
{    }{
}
{   }{
}
}
{}{}{}
\definitionsverweis {pemetaan tangen}{}{}
\maabbdisp {} { T_PM } { T_PN
} {}
bersifat
\definitionsverweis {injektif}{}{.}}
\faktzusatz {Artinya, ruang tangen
\mathl{T_PM}{} merupakan
\definitionsverweis {subruang vektor}{}{}
berdimensi $m$ dari
\mathl{T_PN}{.}}
\faktzusatz {}}{\faktsituation {Misalkan
\mathkor {} {U \subseteq \R^n} {dan} {V \subseteq \R^m} {}
merupakan
\definitionsverweis {subhimpunan-subhimpunan terbuka}{}{,}
yang koordinatnya masing-masing dinyatakan dengan
\mathl{x_1 , \ldots , x_n}{} dan
\mathl{y_1 , \ldots , y_m}{}. Misalkan
\maabbdisp {\varphi} { U } { V
} {}
suatu
\definitionsverweis {pemetaan diferensiabel}{}{}
dan misalkan $\omega$ suatu
$k$-\definitionsverweis {bentuk diferensial}{}{}
pada $V$ dengan representasi

\mavergleichskettedisp
{\vergleichskette
{ \omega
}
{ =} { \sum_{I,\, { \# \left(   I  \right) } = k}  f_I  dy_I
}
{ } {
}
{ } {
}
{ } {
}
}
{}{}{,}
dengan
\maabb {f_I} { V } {\R
} {}
adalah
\definitionsverweis {
fungsi-fungsi}{}{.}}
\faktfolgerung {Maka
\definitionsverweis {bentuk tarikan balik}{}{}
memiliki representasi

\mavergleichskettealignhandlinks
{\vergleichskettealignhandlinks
{ \varphi^* \omega
}
{ =} { \sum_{I,\, { \# \left(   I  \right) } = k}  \left(  f_I \circ \varphi  \right)  \left( \sum_{J,\, { \# \left(   J  \right) } = k} \det  \left(  \left(  { \frac{   \partial  \varphi_i   }{  \partial   x_j   } }  \right)_{ i \in I, j \in J}  \right)  dx_J \right)
}
{ =} { \sum_{J,\, { \# \left(   J  \right) } = k}   \left( \sum_{I,\, { \# \left(   I  \right) } = k} \left(  f_I \circ \varphi  \right)  \det  \left(  \left(  { \frac{   \partial  \varphi_i   }{  \partial   x_j   } }  \right)_{ i \in I, j \in J}  \right)  \right)  dx_J
}
{ } {
}
{ } {
}
}
{}
{}{.}}
\faktzusatz {}
\faktzusatz {}}{\faktsituation {Suatu kurva yang dua kali
\definitionsverweis {diferensiabel}{}{}
\maabb {\gamma} { I } { M
} {}
pada suatu
\definitionsverweis {manifold Riemann}{}{}
$M$ adalah}
\faktfolgerung {suatu
\definitionsverweis {kurva geodesik}{}{}
\zusatzklammer {terhadap
\definitionsverweis {koneksi Levi-Civita}{}{}} {} {,}
tepat ketika dalam setiap peta kurva tersebut memenuhi
\definitionsverweis {sistem persamaan diferensial biasa orde dua}{}{}

\mavergleichskettedisp
{\vergleichskette
{ \gamma_k^{\prime \prime} + \sum_{ij} \Gamma^k_{ij} \gamma_i' \gamma_j'
}
{ =} { 0
}
{ } {
}
{ } {
}
{ } {
}
}
{}{}{}
\zusatzklammer {untuk semua $k$} {} {}
yang diberikan.}
\faktzusatz {}
\faktzusatz {}}

}



\inputaufgabeklausurloesung
{6 (2+4)}

{

Misalkan

\mavergleichskette
{\vergleichskette
{ \alpha,\beta
}
{ \in }{ \R_+
}
{  }{
}
{    }{
}
{   }{
}
}
{}{}{}
dan misalkan

\mavergleichskettedisp
{\vergleichskette
{ Y
}
{ =} { \left\{ (x,y)\in\R^2  \mid   \alpha x^2+\beta y^2 = 1        \right\}
}
{ } {
}
{ } {
}
{ } {
}
}
{}{}{.}
\aufzaehlungzwei {Tentukan
\definitionsverweis {pemetaan Gauss}{}{}
untuk $Y$.
} { Berikan secara eksplisit suatu
\definitionsverweis {pemetaan invers}{}{}
dari pemetaan Gauss untuk $Y$.
}

}
{

\aufzaehlungzwei {Sebagai
\definitionsverweis {medan normal satuan}{}{},
kita dapat mulai dari gradien fungsi pendefinisi, yaitu

\mavergleichskettedisp
{\vergleichskette
{ G(x,y)
}
{ =} { \begin{pmatrix}  2 \alpha x  \\ 2 \beta y    \end{pmatrix}
}
{ } {
}
{ } {
}
{ } {
}
}
{}{}{}
dan setelah normalisasi diperoleh

\mavergleichskettedisp
{\vergleichskette
{ N(x,y)
}
{ =} {  { \frac{   1 }{  \sqrt{ \alpha^2 x^2 + \beta^2 y^2 }  } } \begin{pmatrix}   \alpha x  \\  \beta y    \end{pmatrix}
}
{ } {
}
{ } {
}
{ } {
}
}
{}{}{.}
Jadi, pemetaan Gauss adalah
\maabbeledisp {\varphi} {Y = \left\{  \begin{pmatrix}  x  \\ y   \end{pmatrix}   \in \R^2  \mid   \alpha x^2+\beta y^2 = 1        \right\}  } { S^1
} { \begin{pmatrix}  x  \\ y   \end{pmatrix} } { { \frac{   1 }{  \sqrt{ \alpha^2 x^2 + \beta^2 y^2 }  } } \begin{pmatrix}   \alpha x  \\  \beta y    \end{pmatrix}
} {.}
} {Kita nyatakan bahwa
\maabbeledisp {\psi} {S^1} {Y
} { \begin{pmatrix}  u  \\v   \end{pmatrix} } { { \frac{   1 }{  \sqrt{ { \frac{  u^2 }{ \alpha } }   + { \frac{  v^2 }{ \beta } } }  } }       \begin{pmatrix}   { \frac{   u  }{ \alpha } }   \\ { \frac{   v  }{ \beta } }    \end{pmatrix}
} {,}
adalah fungsi invers. Karena

\mavergleichskettedisp
{\vergleichskette
{ { \frac{   1 }{   { \frac{  u^2 }{ \alpha } }   + { \frac{  v^2 }{ \beta } }  } } \left(  \alpha { \frac{  u^2 }{ \alpha^2 } } + \beta{ \frac{  v^2 }{ \beta^2 } }  \right)
}
{ =} { 1
}
{ } {
}
{ } {
}
{ } {
}
}
{}{}{}
citra $\psi$ terletak dalam $Y$. Karena

\mavergleichskettealign
{\vergleichskettealign
{  \varphi \left(  \psi \begin{pmatrix}  u  \\v   \end{pmatrix}  \right)
}
{ =} { \varphi \left(  { \frac{   1 }{  \sqrt{ { \frac{  u^2 }{ \alpha } }   + { \frac{  v^2 }{ \beta } } }  } }       \begin{pmatrix}   { \frac{   u  }{ \alpha } }   \\ { \frac{   v  }{ \beta } }    \end{pmatrix}    \right)
}
{ =} { \varphi \begin{pmatrix}   { \frac{   1 }{  \alpha    \sqrt{ { \frac{  u^2 }{ \alpha } }   + { \frac{  v^2 }{ \beta } } }  } }   u    \\ { \frac{   1 }{ \beta  \sqrt{ { \frac{  u^2 }{ \alpha } }  + { \frac{  v^2 }{ \beta } } }  } }  v       \end{pmatrix}
}
{ =} { { \frac{   1 }{  \sqrt{  \alpha^2 { \frac{  u^2 }{  \alpha^2   \left(   { \frac{  u^2 }{ \alpha } }  +  { \frac{  v^2 }{ \beta } }   \right)  } }  + \beta^2 { \frac{  v^2 }{  \beta^2 \left(   { \frac{  u^2 }{ \alpha } }  +  { \frac{  v^2 }{ \beta } }   \right)  } }       }  } }  \begin{pmatrix}  { \frac{   1 }{  \sqrt{ { \frac{  u^2 }{ \alpha } }   + { \frac{  v^2 }{ \beta } } }  } }   u    \\   { \frac{   1 }{   \sqrt{ { \frac{  u^2 }{ \alpha } } + { \frac{  v^2 }{ \beta } } }  } } v    \end{pmatrix}
}
{ =} { { \frac{   1 }{  \sqrt{ { \frac{  u^2 }{  \left(   { \frac{  u^2 }{ \alpha } }  +  { \frac{  v^2 }{ \beta } }   \right)  } }   + { \frac{  v^2 }{  \left(   { \frac{  u^2 }{ \alpha } }  +  { \frac{  v^2 }{ \beta } }   \right)  } }       }  } }  \begin{pmatrix}   { \frac{   1 }{  \sqrt{ { \frac{  u^2 }{ \alpha } }   + { \frac{  v^2 }{ \beta } } }  } }   u    \\   { \frac{   1 }{   \sqrt{ { \frac{  u^2 }{ \alpha } }   + { \frac{  v^2 }{ \beta } } }  } } v    \end{pmatrix}
}
}
{
\vergleichskettefortsetzungalign
{ =} { { \frac{    \sqrt{ { \frac{  u^2 }{ \alpha } }  +  { \frac{  v^2 }{ \beta } } } }{  \sqrt{    u^2+ v^2       }  } }  \begin{pmatrix}   { \frac{   1 }{  \sqrt{ { \frac{  u^2 }{ \alpha } }   + { \frac{  v^2 }{ \beta } } }  } }   u    \\   { \frac{   1 }{   \sqrt{ { \frac{  u^2 }{ \alpha } }   + { \frac{  v^2 }{ \beta } } }  } } v    \end{pmatrix}
}
{ =} { \begin{pmatrix}  u  \\v   \end{pmatrix}
}
{ } {
}
{ } {}
}
{}{}
dan

\mavergleichskettealign
{\vergleichskettealign
{  \psi \left(  \varphi \begin{pmatrix}  x  \\ y   \end{pmatrix}  \right)
}
{ =} { \psi \begin{pmatrix}  { \frac{    \alpha x }{  \sqrt{\alpha^2 x^2+ \beta^2 y^2 }  } }   \\ { \frac{   \beta y }{  \sqrt{\alpha^2 x^2+ \beta^2 y^2  } }}    \end{pmatrix}
}
{ =} { { \frac{   1 }{  \sqrt{ { \frac{   \alpha x^2 }{  \alpha^2x^2 + \beta^2y^2  } }  + { \frac{   \beta y^2 }{  \alpha^2x^2 + \beta^2y^2  } }  }  } }  \begin{pmatrix}  { \frac{    x }{  \sqrt{\alpha^2 x^2+ \beta^2 y^2 }  } }   \\ { \frac{   y }{  \sqrt{\alpha^2 x^2+ \beta^2 y^2  } }}    \end{pmatrix}
}
{ =} { { \frac{   \sqrt{\alpha^2 x^2+ \beta^2 y^2 } }{  \sqrt{    \alpha x^2 +  \beta y^2  }  } }  \begin{pmatrix}  { \frac{    x }{  \sqrt{\alpha^2 x^2+ \beta^2 y^2 }  } }   \\ { \frac{   y }{  \sqrt{\alpha^2 x^2+ \beta^2 y^2  } }}    \end{pmatrix}
}
{ =} { \begin{pmatrix}  x  \\ y   \end{pmatrix}
}
}
{}
{}{}
kedua pemetaan tersebut saling invers.
}

}



\inputaufgabeklausurloesung
{4}

{

Jelaskan pasangan istilah \anfuehrung{ekstrinsik dan intrinsik}{} dalam konteks manifold.

}
{  }



\inputaufgabeklausurloesung
{5}

{

Buktikan rumus luas permukaan suatu benda putar yang dihasilkan oleh
\definitionsverweis {kurva diferensiabel}{}{}
\maabbeledisp {\gamma} {]a,b[} {\R^2
} { t } {(x(t),y(t))
} {,}
dengan

\mavergleichskette
{\vergleichskette
{ y(t)
}
{ > }{ 0
}
{  }{
}
{    }{
}
{   }{
}
}
{}{}{.}

}
{

Misalkan $S$ permukaan putar, yang merupakan submanifold tertutup berdimensi dua dalam suatu himpunan terbuka di $\R^3$. Kita terapkan
Korolar 17.2 (Geometri diferensial (Osnabrück 2023))
pada parametrisasi
\maabbeledisp {} { ]a,b[ \times ]0, 2 \pi[} {S
} { (t, \alpha) } { (x(t),  y(t) \cos \alpha  , y(t) \sin \alpha  )
} {,}.
Turunan-turunan parsialnya adalah

\mathdisp {\begin{pmatrix}  x'(t) \\ y'(t) \cos \alpha \\  y'(t) \sin \alpha   \end{pmatrix} \text{   dan } \begin{pmatrix}  0  \\ - y(t) \sin \alpha \\  y(t) \cos \alpha   \end{pmatrix}} {  }

dan karena itu

\mathdisp {E(t, \alpha) = (x'(t))^2 + (y'(t))^2,\, F(t, \alpha)=0,\, G(t, \alpha) = (y(t))^2} { . }

Jadi, luas permukaannya sama dengan

\mavergleichskettedisphandlinks
{\vergleichskettedisphandlinks
{ \int_{  a  }^{  b  } \int_{  0  }^{  2 \pi } \sqrt{(x'(t))^2+(y'(t))^2} \cdot   y(t) \, d \alpha \, d t
}
{ =} { 2 \pi \int_{  a  }^{  b  } \sqrt{(x'(t))^2+(y'(t))^2} \cdot   y(t) \, d t
}
{ } {
}
{ } {
}
{ } {
}
}
{}{}{.}

}



\inputaufgabeklausurloesung
{0}

{

}
{  }



\inputaufgabeklausurloesung
{5}

{

Misalkan $X$ adalah suatu
\definitionsverweis {torus}{}{.}
Berikan suatu
\definitionsverweis {pemetaan diferensiabel}{}{}
surjektif
\maabbdisp {\varphi} {\R^2} {X
} {}
sedemikian sehingga
\definitionsverweis {pemetaan singgung}{}{}
\maabbdisp {T_P(\varphi)} { T_P\R^2 } { T_{\varphi(P)}X
} {}
juga surjektif di setiap titik

\mavergleichskette
{\vergleichskette
{ P
}
{ \in }{ \R^2
}
{  }{
}
{    }{
}
{   }{
}
}
{}{}{.}

}
{

Kita tulis torus sebagai $X=S^1 \times S^1$ dengan lingkaran satuan
\mathl{S^1= \left\{ (x,y) \in \R^2  \mid   \betrag { (x,y) } = 1        \right\}}{.}
Pemetaan
\maabbeledisp {\varphi} {\R^2} {S^1 \times S^1
} {(u,v)} {( \cos  u, \sin  u , \cos  v , \sin  v )
} {,}
bersifat surjektif karena kedua komponennya merupakan parametrisasi standar lingkaran satuan. Surjektivitas pemetaan tangen juga dapat diperiksa per komponen karena ruang tangen manifold produk adalah hasil kali ruang-ruang tangennya. Turunan dari
\maabbeledisp {} {\R} { \R^2
} { t } { ( \cos  t, \sin  t)
} {,}
adalah
\mathl{(- \sin  t, \cos  t) \neq (0,0)}{,} sehingga vektor citra ini selalu merentang ruang tangen satu-dimensional dari lingkaran satuan di titik citra.

}



\inputaufgabeklausurloesung
{6 (2+2+2)}

{

Tinjau grafik $M$ dari pemetaan
\maabbeledisp {\varphi} {\R^2} {\R
} {(u,v)} { u^2+uv-v^3
} {,}
sebagai submanifold tertutup berdimensi dua di $\R^3$, yaitu

\mavergleichskettedisp
{\vergleichskette
{ M
}
{ =} { \left\{ (u,v,u^2+uv-v^3)  \mid  (u,v) \in \R^2         \right\}
}
{ } {
}
{ } {
}
{ } {
}
}
{}{}{}
dengan metrik Riemann yang diinduksi dari $\R^3$. Misalkan
\maabbeledisp {\psi} {\R^2} { M
} {(u,v)} { (u,v,u^2+uv-v^3)
} {,}
adalah difeomorfisme yang bersesuaian.

a) Tentukan diferensial total dari $\psi$ serta vektor-vektor citra
\mathl{T_P(\psi) (e_1)}{} dan
\mathl{T_P(\psi) (e_2)}{} di
\mathl{T_{\psi(P)}M}{.}

b) Untuk setiap titik berbentuk

\mavergleichskette
{\vergleichskette
{ P
}
{ = }{ (u,0)
}
{  }{
}
{    }{
}
{   }{
}
}
{}{}{}
tentukan luas jajaran genjang yang direntang oleh
\mathl{T_P(\psi) (e_1)}{} dan
\mathl{T_P(\psi) (e_2)}{} di
\mathl{T_{\psi(P)}M}{}.

c) Untuk setiap titik berbentuk

\mavergleichskette
{\vergleichskette
{ P
}
{ = }{ (0,v)
}
{  }{
}
{    }{
}
{   }{
}
}
{}{}{}
tentukan luas jajaran genjang yang direntang oleh
\mathl{T_P(\psi) (e_1)}{} dan
\mathl{T_P(\psi) (e_2)}{} di
\mathl{T_{\psi(P)}M}{}.

}
{

a) Diferensial total dari $\psi$ di titik
\mathl{P=(u,v)}{} adalah

\mathdisp {\left(D \psi\right)_{P} = \begin{pmatrix}  1 &  0 \\  0 &  1 \\   2u+v & u-3v^2 \end{pmatrix}} {  }

dan berlaku

\mathdisp {T_P(\psi) (e_1) = \left(  D \psi  \right)_{ P } \left(  e_1  \right) = \begin{pmatrix}  1 \\ 0\\  2u+v   \end{pmatrix} \text{   dan } T_P(\psi) (e_2)   = \left(  D \psi  \right)_{ P } \left(  e_2  \right) = \begin{pmatrix}  0 \\ 1\\ u-3v^2   \end{pmatrix}} { . }

b) dan c) Untuk menentukan luas, pertama-tama kita hitung hasil kali skalar dari kedua vektor
\mathkor {} {\begin{pmatrix}  1 \\ 0\\  2u+v   \end{pmatrix}} {dan} {\begin{pmatrix}  0 \\ 1\\ u-3v^2   \end{pmatrix}} {.}
Kita peroleh

\mathdisp {\left\langle \begin{pmatrix}  1 \\ 0\\  2u+v   \end{pmatrix}  , \begin{pmatrix}  1 \\ 0\\  2u+v   \end{pmatrix}   \right\rangle = 1+(2u+v)^2 = 1 +4u^2+4uv +v^2} { , }

\mathdisp {\left\langle \begin{pmatrix}  1 \\ 0\\  2u+v   \end{pmatrix}  , \begin{pmatrix}  0 \\ 1\\ u-3v^2   \end{pmatrix}    \right\rangle = (2u+v)(u-3v^2) = 2u^2 -6uv^2 +uv -3v^3} {  }

dan

\mathdisp {\left\langle \begin{pmatrix}  0 \\ 1\\ u-3v^2   \end{pmatrix}   , \begin{pmatrix}  0 \\ 1\\ u-3v^2   \end{pmatrix}    \right\rangle = 1+(u-3v^2)^2 = 1 +u^2 -6uv^2  + 9v^4} { . }

b) Untuk
\mathl{P=(u,0)}{}, luas jajaran genjang yang direntang oleh
\mathl{T_P(\psi) (e_1)}{} dan
\mathl{T_P(\psi) (e_2)}{} di
\mathl{T_{\psi(P)}M}{} adalah

\mavergleichskettedisp
{\vergleichskette
{ \sqrt { \det  \begin{pmatrix}  1+4u^2  &   2u^2   \\  2 u^2 &  1+ u^2  \end{pmatrix} }
}
{ =} { \sqrt{  (1+4u^2)(1+u^2) - 4u^4 }
}
{ =} { \sqrt{  1+5u^2 }
}
{ } {
}
{ } {
}
}
{}{}{.}

c) Untuk
\mathl{P=(0,v)}{}, luas jajaran genjang yang direntang oleh
\mathl{T_P(\psi) (e_1)}{} dan
\mathl{T_P(\psi) (e_2)}{} di
\mathl{T_{\psi(P)}M}{} adalah

\mavergleichskettedisp
{\vergleichskette
{ \sqrt { \det  \begin{pmatrix}  1+v^2  &   -3v^3   \\  -3  v^3 &  1+ 9v^4  \end{pmatrix} }
}
{ =} { \sqrt{  (1+v^2)(1+9v^4) - 9v^6 }
}
{ =} { \sqrt{  1+v^2 +9v^4  }
}
{ } {
}
{ } {
}
}
{}{}{.}

}



\inputaufgabeklausurloesung
{6}

{

Buktikan aturan hasil kali untuk bentuk diferensial yang diferensiabel pada suatu himpunan terbuka di $\R^n$.

}
{

Misalkan
\mathl{x_1 , \ldots , x_n}{} adalah koordinat pada $\R^n$. Karena linearitas $d$ dan
multilinearitas hasil kali baji,
kita dapat menulis kedua bentuk diferensial sebagai
\mathkor {} {\omega=f dx_I} {dan} {\tau=g dx_J} {}
dengan himpunan indeks
\mathkor {} {I=\{i_1 , \ldots , i_k\}} {dan} {J=\{ j_1 , \ldots , j_\ell \}} {}.
Maka

\mavergleichskettealignhandlinks
{\vergleichskettealignhandlinks
{ d( \omega \wedge \tau)
}
{ =} { d \left(  fdx_I \wedge gdx_J  \right)
}
{ =} { d \left(  (f g) dx_I \wedge dx_J  \right)
}
{ =} { \sum_{s = 1}^n { \frac{   \partial  (fg)  }{  \partial   x_s  } }  dx_s  \wedge  dx_I  \wedge dx_J
}
{ =} { \sum_{s = 1}^n \left(  g { \frac{   \partial   f   }{  \partial   x_s  } } + f { \frac{   \partial   g   }{  \partial   x_s  } }  \right)  dx_s \wedge  dx_I  \wedge dx_J
}
}
{
\vergleichskettefortsetzungalign
{ =} { \sum_{s = 1}^n  g { \frac{   \partial   f   }{  \partial   x_s  } } dx_s \wedge  dx_I \wedge dx_J + \sum_{s = 1}^n  f { \frac{   \partial   g   }{  \partial   x_s  } }  dx_s \wedge  dx_I  \wedge dx_J
}
{ =} { \sum_{s = 1}^n  { \frac{   \partial   f   }{  \partial   x_s  } } dx_s \wedge  dx_I \wedge gdx_J + \sum_{s = 1}^n  { \frac{   \partial   g   }{  \partial   x_s  } }  dx_s \wedge f dx_I \wedge dx_J
}
{ =} { d(fdx_I ) \wedge g dx_J  + \sum_{s = 1}^n (-1)^k  f dx_I \wedge { \frac{   \partial   g   }{  \partial   x_s  } }  dx_s \wedge dx_J
}
{ =} { d(fdx_I ) \wedge g dx_J + (-1)^k  f dx_I \wedge  d( g dx_J)
}
}
{}{.}

}



\inputaufgabeklausurloesung
{0}

{

}
{  }



\inputaufgabeklausurloesung
{6}

{

Misalkan $M$ adalah suatu
\definitionsverweis {manifold berbatas}{}{}
dan misalkan

\mavergleichskette
{\vergleichskette
{ P
}
{ \in }{ \partial M
}
{  }{
}
{    }{
}
{   }{
}
}
{}{}{}
adalah titik batas. Tunjukkan bahwa sifat-sifat berikut ekuivalen bagi suatu vektor singgung

\mavergleichskette
{\vergleichskette
{ v
}
{ \in }{ T_PM
}
{  }{
}
{    }{
}
{   }{
}
}
{}{}{.}
\aufzaehlungzwei {Terdapat suatu lintasan yang terdiferensialkan secara kontinu
\maabb {\gamma} { [- \epsilon, 0]} { M
} {}
dengan

\mavergleichskette
{\vergleichskette
{ \gamma(0)
}
{ = }{ P
}
{  }{
}
{    }{
}
{   }{
}
}
{}{}{}
dan

\mavergleichskette
{\vergleichskette
{ \gamma'(0)
}
{ = }{ v
}
{  }{
}
{    }{
}
{   }{
}
}
{}{}{.}
} { Pada setiap peta bermodel setengah ruang negatif, $v$ dipetakan oleh pemetaan singgung ke setengah ruang kanan.
}

}
{

Misalkan $\gamma$ suatu setengah lintasan di $M$ seperti yang diberikan, dan misalkan

\mavergleichskette
{\vergleichskette
{ P
}
{ \in }{  U
}
{ \subseteq }{ M
}
{    }{
}
{   }{
}
}
{}{}{}
suatu daerah peta dengan peta
\maabbdisp {\alpha} { U } {V \subseteq  H_{\leq 0}
} {}
dengan

\mavergleichskettedisp
{\vergleichskette
{  H_{\leq 0}
}
{ =} { \left\{ (x_1,x_2 , \ldots, x_n)  \mid   x_1 \leq 0        \right\}
}
{ } {
}
{ } {
}
{ } {
}
}
{}{}{}
dan

\mavergleichskettedisp
{\vergleichskette
{ Q
}
{ =} { \alpha(P)
}
{ =} { \left(  0   , \,  a_2  , \,   \ldots , \,  a_n                \right)
}
{ } {
}
{ } {
}
}
{}{}{.}
Di bawah pemetaan tangen,

\mavergleichskette
{\vergleichskette
{ v
}
{ = }{ \gamma'(0)
}
{  }{
}
{    }{
}
{   }{
}
}
{}{}{}
dipetakan ke vektor turunan dari

\mavergleichskette
{\vergleichskette
{ \delta
}
{ = }{ \alpha \circ \gamma
}
{  }{
}
{    }{
}
{   }{
}
}
{}{}{}
di titik nol. Karena $\gamma$ berjalan dalam $M$, seluruh $\delta$ berjalan dalam $H_{\leq 0}$. Oleh karena itu, dalam hasil bagi selisih

\mathdisp {{ \frac{   \left(  \delta_1(t)   , \,   \delta_2(t)  , \,   \ldots , \,   \delta_n(t)                \right)  }{ t } }} {  }

baik $t$ maupun $\delta_1(t)$ bernilai negatif, sehingga

\mavergleichskettedisp
{\vergleichskette
{  \delta_1'(0)
}
{ \geq} { 0
}
{ } {
}
{ } {
}
{ } {
}
}
{}{}{,}
dan dengan demikian termasuk dalam setengah ruang positif.

Sebaliknya, misalkan
\mathl{T_P(\alpha)(v)}{} termasuk dalam setengah ruang positif. Maka garis
\maabbeledisp {\delta} {\R} {\R^n
} { t } { Q+tv
} {}
untuk

\mavergleichskette
{\vergleichskette
{ t
}
{ \leq }{  0
}
{  }{
}
{    }{
}
{   }{
}
}
{}{}{}
seluruhnya berjalan dalam setengah ruang negatif. Lintasan
\mathl{\alpha^{-1} \circ \delta}{} didefinisikan pada suatu interval negatif, bernilai dalam $M$, dan merupakan realisasi setengah lintasan dari $v$.

}



\inputaufgabeklausurloesung
{12}

{

Buktikan teorema Stokes untuk manifold berbatas.

}
{

\teilbeweis {}{}{}
{Misalkan

\mathbed {U_i} {}
 {i \in I} {}
 {} {} {} {,}
suatu
\definitionsverweis {penutup terbuka}{}{}
dari $M$ dengan
\definitionsverweis {peta-peta berorientasi}{}{},
dan misalkan

\mathbed {h_j} {}
 {j \in J} {}
 {} {} {} {,}
suatu
\definitionsverweis {partisi kesatuan yang terdiferensialkan secara kontinu dan subordinat terhadap penutup tersebut}{}{,}
yang ada menurut
Teorema 22.10 (Geometri diferensial (Osnabrück 2023)).
Untuk setiap

\mavergleichskette
{\vergleichskette
{ P
}
{ \in }{ M
}
{  }{
}
{    }{
}
{   }{
}
}
{}{}{}
terdapat lingkungan terbuka

\mavergleichskette
{\vergleichskette
{ P
}
{ \in }{ V_P
}
{ \subseteq }{ M
}
{    }{
}
{   }{
}
}
{}{}{}
sedemikian sehingga

\mavergleichskette
{\vergleichskette
{ h_j \vert_{V_P }
}
{ = }{ 0
}
{  }{
}
{    }{
}
{   }{
}
}
{}{}{}
kecuali untuk berhingga banyak indeks. Misalkan $Y$ dukungan dari $\omega$. Penutup

\mavergleichskette
{\vergleichskette
{ Y
}
{ \subseteq }{ \bigcup_{P \in M} V_P
}
{  }{
}
{    }{
}
{   }{
}
}
{}{}{}
memiliki subpenutup berhingga karena
\definitionsverweis {kekompakan}{}{}
yang diasumsikan, katakanlah

\mavergleichskettedisp
{\vergleichskette
{ Y
}
{ \subseteq} { V_{1} \cup \ldots \cup V_{r}
}
{ =} { V
}
{ } {
}
{ } {
}
}
{}{}{.}
Jadi, hanya berhingga banyak $h_j$ pada $V$ yang tidak sama dengan $0$. Kita tetapkan

\mavergleichskette
{\vergleichskette
{ \omega_j
}
{ = }{ h_j \omega
}
{  }{
}
{    }{
}
{   }{
}
}
{}{}{;}
bentuk-bentuk diferensial ini juga terdiferensialkan secara kontinu. Dukungan $h_j$ adalah suatu subhimpunan dalam $M$ yang
\definitionsverweis {tertutup}{}{,}
dan terletak dalam
\mathl{U_{i(j)}}{}; karena itu, dukungan $\omega_j$ terletak dalam
\mathl{Y \cap U_{i(j)}}{} dan bersifat kompak menurut
Soal 13.10 (Geometri diferensial (Osnabrück 2023)).
Berlaku

\mavergleichskettedisp
{\vergleichskette
{ \omega
}
{ =} { \sum_{j \in J} \omega_j
}
{ } {
}
{ } {
}
{ } {
}
}
{}{}{,}
dan hanya berhingga banyak bentuk diferensial
\mathl{\omega_j}{} ini yang tidak sama dengan $0$ karena

\mavergleichskette
{\vergleichskette
{ \omega_j \vert_{M \setminus Y}
}
{ = }{ 0
}
{  }{
}
{    }{
}
{   }{
}
}
{}{}{}
untuk semua $j$, sedangkan

\mavergleichskettedisp
{\vergleichskette
{ \omega_j \vert_V
}
{ =} { 0
}
{ } {
}
{ } {
}
{ } {
}
}
{}{}{}
untuk semua $j$ kecuali berhingga banyak. Berdasarkan
aditivitas integral bentuk diferensial
dan
aditivitas turunan eksterior,
pernyataan dapat dibuktikan secara terpisah untuk setiap $\omega_j$. Jadi, kita dapat mengasumsikan bahwa diberikan suatu bentuk diferensial berderajat $(n-1)$ yang terdiferensialkan secara kontinu, memiliki dukungan kompak, dan seluruh dukungannya terletak dalam suatu lingkungan peta $U$.}
{}
\teilbeweis {}{}{}
{Terdapat suatu
\definitionsverweis {difeomorfisme}{}{}
\maabb {\alpha} { U } { V
} {}
dengan

\mavergleichskette
{\vergleichskette
{ V
}
{ \subseteq }{ H
}
{  }{
}
{    }{
}
{   }{
}
}
{}{}{}
terbuka, yang sekaligus menginduksi difeomorfisme antara batas-batas

\mavergleichskette
{\vergleichskette
{ \partial U
}
{ = }{ \partial M \cap U
}
{  }{
}
{    }{
}
{   }{
}
}
{}{}{}
dan

\mavergleichskette
{\vergleichskette
{ \partial V
}
{ = }{ \partial H \cap V
}
{  }{
}
{    }{
}
{   }{
}
}
{}{}{}.
Di sini berlaku

\mavergleichskettedisp
{\vergleichskette
{
\int_{  \partial U  }   \omega
}
{ =} {
\int_{  \partial V  }  \alpha^{-1 *}\omega
}
{ } {
}
{ } {
}
{ } {
}
}
{}{}{}
dan

\mavergleichskettedisp
{\vergleichskette
{
\int_{ U  }  d \omega
}
{ =} {
\int_{ V  }  \alpha^{-1 *} (  d \omega)
}
{ =} {
\int_{  V  }  d \left(  \alpha^{-1 *}\omega  \right)
}
{ } {
}
{ } {
}
}
{}{}{}
menurut
Lema 20.2 (Geometri diferensial (Osnabrück 2023))  (5).}
{\leerzeichen{}Dengan demikian, kita dapat mulai dari suatu bentuk diferensial yang terdiferensialkan secara kontinu dan memiliki dukungan kompak, didefinisikan pada

\mavergleichskette
{\vergleichskette
{ V
}
{ \subseteq }{ H
}
{  }{
}
{    }{
}
{   }{
}
}
{}{}{},
yang dapat kita perluas menjadi bentuk nol pada seluruh $H$ di luar dukungannya.}
\teilbeweis {}{}{}
{Karena kekompakan dukungannya, terdapat suatu balok yang cukup besar

\mavergleichskette
{\vergleichskette
{ Q
}
{ \subseteq }{ H
}
{  }{
}
{    }{
}
{   }{
}
}
{}{}{,}
yang salah satu sisinya $S$ terletak pada $\partial H$ dan yang bertemu dukungan $\omega$ hanya di $S$. Pada semua sisi lain dari $Q$, $\omega$
\zusatzklammer {dan dengan demikian juga $d\omega$} {} {}
adalah bentuk nol. Oleh karena itu, di satu pihak

\mavergleichskettedisp
{\vergleichskette
{
\int_{ H  }  d\omega
}
{ =} {
\int_{ Q  }  d\omega
}
{ } {
}
{ } {
}
{ } {
}
}
{}{}{}
dan di pihak lain

\mavergleichskettedisp
{\vergleichskette
{
\int_{  \partial H  }   \omega
}
{ =} {
\int_{  S  }   \omega
}
{ =} {
\int_{  \partial Q  }   \omega
}
{ } {
}
{ } {
}
}
{}{}{.}
Jadi, pernyataan tersebut mengikuti dari versi balok teorema Stokes.}
{}

}



\inputaufgabeklausurloesung
{0}

{

}
{  }



\inputaufgabeklausurloesung
{9 (1+1+3+2+2)}

{

Tinjau parametrisasi trigonometrik
\maabbeledisp {\psi} {\R } { S^1
} { t } { \left(  \cos  t   , \,   \sin  t                   \right)
} {}
dari lingkaran satuan. Misalkan

\mavergleichskette
{\vergleichskette
{ c
}
{ \in }{ [0, 2 \pi[
}
{  }{
}
{    }{
}
{   }{
}
}
{}{}{}
dipilih tetap. Pada bundel vektor trivial
\maabbdisp {} {S^1 \times \R^2} { S^1
} {}
diberikan suatu
\definitionsverweis {koneksi linear}{}{}
$\nabla$ sedemikian sehingga, sepanjang $\psi$,
\definitionsverweis {koneksi tarik balik}{}{}
$\psi^* \nabla$ diberikan oleh
\definitionsverweis {simbol Christoffel}{}{}
\zusatzklammer {indeks bagi satu-satunya arah diferensiasi dihilangkan} {} {}

\mavergleichskettedisp
{\vergleichskette
{ \begin{pmatrix} \tilde{\Gamma}_1^1 (t)   &  \tilde{\Gamma}_1^2 (t)   \\ \tilde{\Gamma}_2^1(t)  &  \tilde{\Gamma}_2^2(t)  \end{pmatrix}
}
{ =} { \begin{pmatrix}  0  &   - { \frac{  c  }{  2 \pi   } }   \\  { \frac{  c  }{  2 \pi  } }  &  0  \end{pmatrix}
}
{ } {
}
{ } {
}
{ } {
}
}
{}{}{.}
Untuk suatu titik

\mavergleichskette
{\vergleichskette
{ Q
}
{ = }{ \begin{pmatrix}  a  \\b   \end{pmatrix}
}
{ \in }{ \R^2
}
{    }{
}
{   }{
}
}
{}{}{}
tinjau pemetaan
\maabbeledisp {\Psi_{Q}} {\R} { S^1 \times \R^2
} { t } { \left(  \cos  t   , \,   \sin  t  , \,   M(t)  \begin{pmatrix}  a  \\b   \end{pmatrix}                 \right)
} {}
dengan

\mavergleichskettedisp
{\vergleichskette
{ M(t)
}
{ =} { \begin{pmatrix}  \cos  { \frac{  ct  }{  2 \pi   } }   &   -  \sin  { \frac{  ct  }{  2 \pi   } }    \\  \sin  { \frac{  ct  }{  2 \pi  } }   &  \cos  { \frac{  ct  }{  2 \pi   } }  \end{pmatrix}
}
{ } {
}
{ } {
}
{ } {
}
}
{}{}{.}
\aufzaehlungfuenf{Tentukan $\Psi_Q(0)$.
}{Tentukan $\Psi_Q(2 \pi)$.
}{Tunjukkan bahwa $\Psi_Q$ adalah suatu
\definitionsverweis {pengangkatan horizontal}{}{}
sepanjang $\psi$.
}{Tunjukkan bahwa
\maabbeledisp {} {\R^2} { \R^2
} { Q } { \Psi_Q(2 \pi)
} {,}
merupakan
\definitionsverweis {isometri linear}{}{}
\zusatzklammer {titik pangkal $(1,0)$ tidak dicantumkan di sini} {} {.}
}{Tunjukkan bahwa
\maabbeledisp {} {\Z} { \operatorname{SL}_{  2  } \! \left(  \R  \right)
} { k } { \left(  Q \mapsto \Psi_Q(2k \pi )  \right)
} {,}
merupakan
\definitionsverweis {homomorfisme grup}{}{.}
}

}
{

\aufzaehlungfuenf{Kita peroleh

\mavergleichskettealign
{\vergleichskettealign
{ \Psi_Q(0)
}
{ =} { \left(  \cos  0   , \,   \sin  0  , \,   M(0) \begin{pmatrix}  a  \\b   \end{pmatrix}                 \right)
}
{ =} { \left(  1   , \,   0  , \,   \begin{pmatrix}   1   &   0   \\  0   &  1  \end{pmatrix}    \begin{pmatrix}  a  \\b   \end{pmatrix}                 \right)
}
{ =} { \left(  1   , \,   0  , \,     \begin{pmatrix}  a  \\b   \end{pmatrix}                 \right)
}
{ } {
}
}
{}
{}{.}
}{Kita peroleh

\mavergleichskettealign
{\vergleichskettealign
{ \Psi_Q(2 \pi )
}
{ =} { \left(  \cos  2 \pi   , \,   \sin  2 \pi  , \,   M(2 \pi) \begin{pmatrix}  a  \\b   \end{pmatrix}                 \right)
}
{ =} { \left(  1   , \,   0  , \,   \begin{pmatrix}  \cos  c   &   - \sin c    \\  \sin c   &   \cos  c     \end{pmatrix} \begin{pmatrix}  a  \\b   \end{pmatrix}                 \right)
}
{ =} { \left(  1   , \,   0  , \,     \begin{pmatrix} a \cos  c   - b \sin c   \\ a \sin  c   - b \cos c    \end{pmatrix}                 \right)
}
{ } {
}
}
{}
{}{.}
}{Kita peroleh

\mavergleichskettealign
{\vergleichskettealign
{  \Psi_Q(t)
}
{ =} { \left(  \cos  t    , \,   \sin  t  , \,   \begin{pmatrix}   \cos  { \frac{  ct  }{  2 \pi   } }   &   - \sin  { \frac{  ct  }{  2 \pi   } }    \\  \sin  { \frac{  ct  }{  2 \pi  } }   &   \cos  { \frac{  ct  }{  2 \pi   } }  \end{pmatrix}  \begin{pmatrix}  a  \\b   \end{pmatrix}                 \right)
}
{ =} { \left(  \cos  t   , \,   \sin  t  , \,   \begin{pmatrix}  a \cos  { \frac{  ct  }{  2 \pi  } }  - b \sin  { \frac{  ct  }{  2 \pi  } }   \\ a \sin  { \frac{  ct  }{  2 \pi  } }  +   b  \cos  { \frac{  ct  }{  2 \pi } }    \end{pmatrix}                 \right)
}
{ =} { \left(  \cos  t   , \,   \sin  t  , \,     \left(  a \cos  { \frac{  ct  }{  2 \pi  } }  - b \sin  { \frac{  ct  }{  2 \pi  } }  \right)   e_1 +  \left(  a \sin  { \frac{  ct  }{  2 \pi  } }  +   b  \cos  { \frac{  ct  }{  2 \pi } }  \right)   e_2                 \right)
}
{ } {
}
}
{}
{}{,}
khususnya, ini merupakan suatu pengangkatan. Turunan vertikal dari penampang ini dalam $\R^2$ di atas $\R$ dapat ditentukan menggunakan simbol-simbol Christoffel yang diberikan; kita peroleh

\mavergleichskettealigndrucklinks
{\vergleichskettealigndrucklinks
{ \left(  \psi^* \nabla  \right)_\partial  \left(    \left(  a \cos  { \frac{  ct  }{  2 \pi  } }  - b \sin  { \frac{  ct  }{  2 \pi  } }  \right)   e_1 +  \left(  a \sin  { \frac{  ct  }{  2 \pi  } }  +   b  \cos  { \frac{  ct  }{  2 \pi } }  \right)   e_2  \right)
}
{ =} {   \left(  a \cos  { \frac{  ct  }{  2 \pi  } }  - b \sin  { \frac{  ct  }{  2 \pi  } }  \right)   \left(  \psi^* \nabla  \right)_\partial \left(  e_1  \right)  + \partial    \left(  a \cos  { \frac{  ct  }{  2 \pi  } }  - b \sin  { \frac{  ct  }{  2 \pi  } }  \right)   e_1    +   \left(  a \sin  { \frac{  ct  }{  2 \pi  } }  +   b  \cos  { \frac{  ct  }{  2 \pi } }  \right)      \left(  \psi^* \nabla  \right)_\partial \left(  e_2  \right) + \partial  \left(  a \sin  { \frac{  ct  }{  2 \pi  } }  +   b  \cos  { \frac{  ct  }{  2 \pi } }  \right)   e_2
}
{ =} { -  \left(  a \cos  { \frac{  ct  }{  2 \pi  } }  - b \sin  { \frac{  ct  }{  2 \pi  } }  \right)  { \frac{   c  }{  2 \pi } }  e_2  +   { \frac{   c  }{  2 \pi } } \left(  -a \sin  { \frac{  ct  }{  2 \pi  } }  - b \cos  { \frac{  ct  }{  2 \pi  } }  \right)   e_1    +   \left(  a \sin  { \frac{  ct  }{  2 \pi  } }  +   b  \cos  { \frac{  ct  }{  2 \pi } }  \right) { \frac{   c  }{  2 \pi } }   e_1    +  { \frac{   c  }{  2 \pi } }   \left(  a \cos  { \frac{  ct  }{  2 \pi  } } -  b  \sin  { \frac{  ct  }{  2 \pi } }  \right)   e_2
}
{ =} { 0
}
{ } {
}
}
{}{}{,}
sehingga terdapat suatu
\definitionsverweis {pengangkatan horizontal}{}{}.
}{Pemetaan yang dimaksud adalah
\maabbeledisp {} {\R^2} { \R^2
} { \begin{pmatrix}  a  \\b   \end{pmatrix} } {    \begin{pmatrix}  \cos  c   &   - \sin  c   \\  \sin  c  &  \cos  c  \end{pmatrix} \begin{pmatrix}  a  \\b   \end{pmatrix}
} {,}
yang linear. Matriks yang merepresentasikannya ortogonal
\zusatzklammer {suatu matriks rotasi} {} {,}
sehingga diperoleh suatu isometri.
}{Pemetaan
\mathl{\Psi (2k \pi)}{} adalah pemetaan linear
\maabbeledisp {} {\R^2} { \R^2
} {\begin{pmatrix}  a  \\b   \end{pmatrix} } { \begin{pmatrix}  \cos kc   &   - \sin kc   \\  \sin kc  &  \cos kc  \end{pmatrix} \begin{pmatrix}  a  \\b   \end{pmatrix}
} {,}
yaitu rotasi sebesar sudut $kc$. Matriks rotasi untuk jumlah dua sudut adalah hasil kali matriks dari kedua rotasi tersebut
\zusatzklammer {lihat
Fakta *****} {} {,}
sehingga diperoleh suatu
\definitionsverweis {homomorfisme grup}{}{}.
}

}
\endgroup

\part{Perbaikan Asli yang Terpisah}
\chapter{Enam solusi asli untuk kekosongan sumber}
\noindent\textit{Bagian ini berisi tepat enam solusi asli edisi, berlabel terpisah, dan bukan bagian dari sumber Brenner/Wikiversity maupun formulir solusi resmi.}
% Enam solusi asli untuk kejadian soal yang tidak mempunyai tautan solusi sumber.
% Bukan bagian dari sumber Brenner dan tidak boleh dinisbahkan kepada Brenner/Wikiversity.
% Lisensi: CC BY-SA 4.0.
% Provenans: disusun oleh OpenAI Codex gpt-5.6-sol, Ultra, atas arahan pengguna.

\section{Solusi asli untuk enam kekosongan pada bank ujian}
\label{o011-exam-original-repairs}

Bagian ini memuat tepat enam solusi yang tidak tersedia pada permukaan sumber
yang dibekukan. Setiap solusi diberi ID terpisah dan secara eksplisit ditandai
sebagai materi asli. Teks soal tetap berasal dari formulir peserta resmi;
solusinya tidak dinisbahkan kepada pengarang atau Wikiversity.

\subsection*{Ujian 1, soal tanpa solusi sumber (O011-EXAM01-ORIG-SOL-01)}
\label{o011-exam01-orig-sol-01}

\paragraph{Soal.}
Deskripsikan sebanyak mungkin manifold berbatas berdimensi satu yang
terhubung.

\paragraph{Solusi asli.}
Dengan konvensi manifold Hausdorff dan berbasiskan terhitung yang dipakai
dalam edisi ini, klasifikasi hingga difeomorfisme adalah
\[
 S^1,\qquad \mathbb R,\qquad [0,\infty),\qquad [0,1].
\]
Alasannya dapat dibaca dari banyaknya titik batas dan kekompakan. Jika
batasnya kosong, manifold satu dimensi terhubung adalah lingkaran apabila
kompak dan garis nyata apabila tidak kompak. Jika batasnya mempunyai dua titik,
manifoldnya difeomorfik dengan interval kompak $[0,1]$. Jika batasnya mempunyai
satu titik, manifoldnya difeomorfik dengan setengah garis $[0,\infty)$.
Interval $[0,1)$ tidak memberikan tipe baru karena
$x\mapsto x/(1-x)$ adalah difeomorfisme ke $[0,\infty)$; semua interval terbuka
tak kosong difeomorfik dengan $\mathbb R$. Jadi daftar tersebut juga menjelaskan
mengapa tidak ada tipe tambahan di bawah hipotesis standar.

\subsection*{Ujian 3, soal tanpa solusi sumber (O011-EXAM03-ORIG-SOL-01)}
\label{o011-exam03-orig-sol-01}

\paragraph{Soal.}
Katakan sesuatu yang cerdas tentang pita M\"obius.

\paragraph{Solusi asli.}
Pita M\"obius dapat dibangun sebagai hasil bagi
\[
 \bigl([0,1]\times[-1,1]\bigr)/\bigl((0,t)\sim(1,-t)\bigr).
\]
Ia merupakan manifold mulus kompak berdimensi dua dengan batas; kedua sisi
vertikal direkatkan dengan pembalikan arah. Dua ruas horizontal bergabung
menjadi satu komponen batas. Garis tengah $[0,1]\times\{0\}$ menjadi sebuah
lingkaran dan merupakan retrak deformasi, sehingga pita M\"obius mempunyai
tipe homotopi $S^1$.

Pita ini tidak dapat diorientasikan. Jika suatu basis berorientasi diangkut
sekali mengelilingi garis tengah, identifikasi $(0,t)\sim(1,-t)$ membalik satu
arah basis dan mengembalikannya dengan orientasi berlawanan. Penutup ganda
berorientasinya adalah silinder. Secara setara, pita M\"obius adalah ruang
total bundel garis nyata tak trivial di atas $S^1$; fakta bahwa ia tertanam
sebagai permukaan di $\mathbb R^3$ tidak membuatnya berorientasi.

\subsection*{Ujian 5, soal tanpa solusi sumber (O011-EXAM05-ORIG-SOL-01)}
\label{o011-exam05-orig-sol-01}

\paragraph{Soal.}
Jelaskan peran teorema tentang pemetaan implisit dalam pengembangan teori
manifold.

\paragraph{Solusi asli.}
Teorema pemetaan implisit adalah mesin lokal yang mengubah persamaan menjadi
koordinat. Jika $F\colon U\subseteq\mathbb R^n\to\mathbb R^r$ mulus dan
$D F_p$ surjektif, maka di sekitar $p$ terdapat koordinat mulus yang di
dalamnya $F$ hanyalah proyeksi pada $r$ koordinat. Akibatnya, himpunan aras
regular
\[
 F^{-1}(F(p))
\]
secara lokal adalah $\mathbb R^{n-r}\times\{0\}$ dan karena itu merupakan
submanifold berdimensi $n-r$. Ruang singgungnya diberikan secara langsung oleh
\[
 T_pF^{-1}(F(p))=\ker D F_p.
\]

Hasil ini memasok banyak contoh dasar: hipermuka sebagai aras regular satu
fungsi, grup matriks yang ditentukan oleh persamaan regular, dan grafik solusi
lokal persamaan. Ia juga menjamin bahwa definisi abstrak manifold cocok dengan
gambaran sebagai ruang solusi lokal. Bersama teorema fungsi invers, ia
menyediakan koordinat, membuktikan kemulusan perubahan koordinat yang sesuai,
dan mendasari teorema rank konstan serta kriteria submanifold. Dengan demikian,
teorema itu menghubungkan kalkulus multivariabel dengan teori manifold.

\subsection*{Ujian 7, soal tanpa solusi sumber (O011-EXAM07-ORIG-SOL-01)}
\label{o011-exam07-orig-sol-01}

\paragraph{Soal.}
Hitung luas permukaan bola satuan.

\paragraph{Solusi asli.}
Parametrisasi bola satuan $S^2\subseteq\mathbb R^3$, kecuali garis bujur dan
kedua kutub yang tidak memengaruhi integral, dengan
\[
 \Phi(\theta,\varphi)
 =(\sin\varphi\cos\theta,\sin\varphi\sin\theta,\cos\varphi),
 \quad 0<\theta<2\pi,\quad0<\varphi<\pi.
\]
Vektor turunannya memenuhi
\[
 \left\|\partial_\theta\Phi\times\partial_\varphi\Phi\right\|
 =\sin\varphi.
\]
Maka luasnya adalah
\[
 \operatorname{luas}(S^2)
 =\int_0^{2\pi}\int_0^\pi\sin\varphi\,d\varphi\,d\theta
 =2\pi\,[ -\cos\varphi ]_0^\pi
 =4\pi.
\]

\subsection*{Ujian 9, soal tanpa solusi sumber (O011-EXAM09-ORIG-SOL-01)}
\label{o011-exam09-orig-sol-01}

\paragraph{Soal.}
Deskripsikan suatu model untuk bidang hiperbolik.

\paragraph{Solusi asli.}
Model setengah bidang atas adalah manifold
\[
 \mathbb H^2=\{(x,y)\in\mathbb R^2:y>0\}
\]
dengan metrik Riemann
\[
 g=\frac{dx^2+dy^2}{y^2}.
\]
Ia lengkap, terhubung, dan mempunyai kelengkungan Gaussian konstan $-1$.
Untuk memeriksa kelengkungan, tulis $g=e^{2u}(dx^2+dy^2)$ dengan
$u=-\log y$. Rumus kelengkungan metrik konformal memberi
\[
 K=-e^{-2u}(u_{xx}+u_{yy})
 =-y^2\left(0+\frac1{y^2}\right)=-1.
\]
Geodesiknya adalah garis-garis vertikal dan busur-busur lingkaran yang
memotong batas ideal $y=0$ secara ortogonal. Transformasi M\"obius nyata
\[
 z\longmapsto\frac{az+b}{cz+d},
 \qquad ad-bc=1,
\]
bertindak sebagai isometri; setelah mengidentifikasi matriks yang berbeda
tanda, grup isometri berorientasi adalah $\mathrm{PSL}_2(\mathbb R)$. Model
cakram Poincar\'e bersifat isometrik dengan model ini melalui transformasi
pecahan linear.

\subsection*{Ujian 10, soal tanpa solusi sumber (O011-EXAM10-ORIG-SOL-01)}
\label{o011-exam10-orig-sol-01}

\paragraph{Soal.}
Jelaskan pasangan istilah ``ekstrinsik dan intrinsik'' dalam konteks manifold.

\paragraph{Solusi asli.}
Suatu konsep bersifat \emph{intrinsik} jika dapat didefinisikan dari manifold
dan struktur yang dibawanya tanpa memilih suatu penanaman ke ruang luar.
Contohnya ialah atlas mulus, ruang singgung, metrik Riemann, koneksi
Levi--Civita, panjang kurva, jarak geodesik, dan tensor kelengkungan Riemann.
Isometri manifold mempertahankan semua data tersebut.

Suatu konsep bersifat \emph{ekstrinsik} jika bergantung pada cara manifold
direalisasikan sebagai submanifold ruang ambien. Untuk
$M\subseteq\mathbb R^N$, medan normal, bentuk fundamental kedua, pemetaan
Weingarten, dan kelengkungan rata-rata memakai ruang ambien dan dapat berubah
ketika penanaman diubah, walaupun metrik intrinsik tetap sama.

Perbedaan ini tidak selalu tampak dari rumus awal. Kelengkungan Gaussian suatu
permukaan pertama-tama dapat dihitung melalui bentuk fundamental kedua, tetapi
Theorema egregium menunjukkan bahwa nilainya sebenarnya ditentukan sepenuhnya
oleh metrik intrinsik. Jadi ``intrinsik'' atau ``ekstrinsik'' merujuk pada
ketergantungan geometris konsep, bukan sekadar pada bentuk rumus yang dipakai
untuk menghitungnya.
% O011-COMPLETE-EXTENSION-END

\backmatter
\listoffigures
\chapter*{Atribusi dan Hak Media}
\addcontentsline{toc}{chapter}{Atribusi dan Hak Media}
\section*{Unit 1}
\addcontentsline{toc}{section}{Unit 1}
Empat berkas berikut digunakan dalam Unit 1. Setiap berkas tetap mengikuti lisensinya sendiri; tautan sumber dan lisensi dapat diklik.
\begin{description}
\item[]\texttt{3d-function-6.svg}
MartinThoma (karya asli).
\href{https://commons.wikimedia.org/wiki/File:3d-function-6.svg}{Halaman sumber di Wikimedia Commons}; lisensi \href{https://creativecommons.org/licenses/by/3.0}{CC BY 3.0}.
\item[]\texttt{Great} \texttt{circle} \texttt{passing} \texttt{through} \texttt{two} \texttt{points.svg}
HaEr48; karya turunan dari Polar angle to spherical side.svg oleh Episcophagus.
\href{https://commons.wikimedia.org/wiki/File:Great_circle_passing_through_two_points.svg}{Halaman sumber di Wikimedia Commons}; lisensi \href{https://creativecommons.org/licenses/by-sa/4.0}{CC BY-SA 4.0}.
\item[]\texttt{2019-07-Helix.jpg}
Ag2gaeh (karya asli).
\href{https://commons.wikimedia.org/wiki/File:2019-07-Helix.jpg}{Halaman sumber di Wikimedia Commons}; lisensi \href{https://creativecommons.org/licenses/by-sa/4.0}{CC BY-SA 4.0}.
\item[]\texttt{Planned} \texttt{flight} \texttt{map} \texttt{of} \texttt{the} \texttt{Oiseau} \texttt{Blanc.svg}
Pethrus; karya turunan dari BlankMap-World8.svg oleh AMK1211.
\href{https://commons.wikimedia.org/wiki/File:Planned_flight_map_of_the_Oiseau_Blanc.svg}{Halaman sumber di Wikimedia Commons}; lisensi \href{https://creativecommons.org/licenses/by-sa/3.0}{CC BY-SA 3.0}.
\end{description}
Identitas revisi, ukuran, SHA-1 Wikimedia, SHA-256, URL berkas asli, dan hash turunan cetak tercatat dalam manifest media yang disertakan.
\section*{Unit 2}
\addcontentsline{toc}{section}{Unit 2}
Dua berkas berikut digunakan dalam Unit 2. Setiap berkas tetap mengikuti status hak atau lisensinya sendiri; tautan sumber dapat diklik dan tautan lisensi dicantumkan jika tersedia.
\begin{description}
\item[]\texttt{Integral} \texttt{apl} \texttt{rot} \texttt{obsah1.svg}
Pajs; dialihkan dari Wikipedia bahasa Ceko ke Wikimedia Commons.
\href{https://commons.wikimedia.org/wiki/File:Integral_apl_rot_obsah1.svg}{Halaman sumber di Wikimedia Commons}; status hak: domain publik.
\item[]\texttt{Hyperboloid1.png}
Lars H. Rohwedder (RokerHRO; karya asli).
\href{https://commons.wikimedia.org/wiki/File:Hyperboloid1.png}{Halaman sumber di Wikimedia Commons}; status hak: domain publik.
\end{description}
Identitas revisi, ukuran, SHA-1 Wikimedia, SHA-256, URL berkas asli, dan hash turunan cetak tercatat dalam manifest media yang disertakan.
\section*{Unit 3}
\addcontentsline{toc}{section}{Unit 3}
Tiga berkas berikut digunakan dalam Unit 3. Setiap berkas tetap mengikuti status hak atau lisensinya sendiri; tautan sumber dapat diklik dan tautan lisensi dicantumkan jika tersedia.
\begin{description}
\item[]\texttt{Parabola} \texttt{circle.svg}
IkamusumeFan (karya sendiri).
\href{https://commons.wikimedia.org/wiki/File:Parabola_circle.svg}{Halaman sumber di Wikimedia Commons}; lisensi \href{https://creativecommons.org/licenses/by-sa/4.0}{CC BY-SA 4.0}.
\item[]\texttt{Euler} \texttt{spiral.svg}
AdiJapan (karya sendiri).
\href{https://commons.wikimedia.org/wiki/File:Euler_spiral.svg}{Halaman sumber di Wikimedia Commons}; lisensi \href{https://creativecommons.org/licenses/by-sa/3.0}{CC BY-SA 3.0}.
\item[]\texttt{Evolute-parab.svg}
Ag2gaeh (karya sendiri).
\href{https://commons.wikimedia.org/wiki/File:Evolute-parab.svg}{Halaman sumber di Wikimedia Commons}; lisensi \href{https://creativecommons.org/licenses/by-sa/4.0}{CC BY-SA 4.0}.
\end{description}
Identitas revisi, ukuran, SHA-1 Wikimedia, SHA-256, URL berkas asli, dan hash turunan cetak tercatat dalam manifest media yang disertakan.
\section*{Unit 4}
\addcontentsline{toc}{section}{Unit 4}
Unit 4 tidak menggunakan berkas media. Penutupan nol-aset ini diverifikasi terhadap kuliah, lembar kerja, dan semua solusi yang disediakan oleh sumber.
\section*{Unit 5}
\addcontentsline{toc}{section}{Unit 5}
Satu berkas berikut digunakan dalam Unit 5. Setiap berkas tetap mengikuti status hak atau lisensinya sendiri; tautan sumber dapat diklik dan tautan lisensi dicantumkan jika tersedia.
\begin{description}
\item[]\texttt{Minimal} \texttt{surface} \texttt{curvature} \texttt{planes-de.svg}
Eric Gaba (Sting); berdasarkan sebuah gambar dalam buku.
\href{https://commons.wikimedia.org/wiki/File:Minimal_surface_curvature_planes-de.svg}{Halaman sumber di Wikimedia Commons}; lisensi \href{https://creativecommons.org/licenses/by-sa/3.0/}{CC BY-SA 3.0}.
\end{description}
Identitas revisi, ukuran, SHA-1 Wikimedia, SHA-256, URL berkas asli, dan hash turunan cetak tercatat dalam manifest media yang disertakan.
\section*{Unit 6}
\addcontentsline{toc}{section}{Unit 6}
Satu berkas berikut digunakan dalam Unit 6. Setiap berkas tetap mengikuti status hak atau lisensinya sendiri; tautan sumber dapat diklik dan tautan lisensi dicantumkan jika tersedia.
\begin{description}
\item[]\texttt{Parallel} \texttt{transport} \texttt{sphere2.svg}
Silly rabbit di Wikipedia bahasa Inggris (karya sendiri).
\href{https://commons.wikimedia.org/wiki/File:Parallel_transport_sphere2.svg}{Halaman sumber di Wikimedia Commons}; lisensi \href{https://creativecommons.org/licenses/by-sa/3.0/}{CC BY-SA 3.0}.
\end{description}
Identitas revisi, ukuran, SHA-1 Wikimedia, SHA-256, URL berkas asli, dan hash turunan cetak tercatat dalam manifest media yang disertakan.
\section*{Unit 7}
\addcontentsline{toc}{section}{Unit 7}
Tiga berkas berikut digunakan dalam Unit 7. Setiap berkas tetap mengikuti status hak atau lisensinya sendiri; tautan sumber dapat diklik dan tautan lisensi dicantumkan jika tersedia.
\begin{description}
\item[]\texttt{Stereographic} \texttt{projection} \texttt{in} \texttt{3D.png}
Mark.Howison (karya sendiri; domain publik).
\href{https://commons.wikimedia.org/wiki/File:Stereographic_projection_in_3D.png}{Halaman sumber di Wikimedia Commons}; status hak: Public domain.
\item[]\texttt{Manifold} \texttt{zahyou3.png}
132ninme di Wikipedia bahasa Jepang (karya sendiri; nama Unicode persis dicatat dalam manifes).
\href{https://commons.wikimedia.org/wiki/File:Manifold_zahyou3.png}{Halaman sumber di Wikimedia Commons}; lisensi \href{https://creativecommons.org/licenses/by-sa/3.0/}{CC BY-SA 3.0}.
\item[]\texttt{Circle} \texttt{-} \texttt{black} \texttt{simple.svg}
Dakdada (karya sendiri; domain publik).
\href{https://commons.wikimedia.org/wiki/File:Circle_-_black_simple.svg}{Halaman sumber di Wikimedia Commons}; status hak: Public domain.
\end{description}
Identitas revisi, ukuran, SHA-1 Wikimedia, SHA-256, URL berkas asli, dan hash turunan cetak tercatat dalam manifest media yang disertakan.
\section*{Unit 8}
\addcontentsline{toc}{section}{Unit 8}
Unit 8 tidak menggunakan berkas media. Penutupan nol-aset ini diverifikasi terhadap kuliah, lembar kerja, dan semua solusi yang disediakan oleh sumber.
\section*{Unit 9}
\addcontentsline{toc}{section}{Unit 9}
Satu berkas berikut digunakan dalam Unit 9. Setiap berkas tetap mengikuti status hak atau lisensinya sendiri; tautan sumber dapat diklik dan tautan lisensi dicantumkan jika tersedia.
\begin{description}
\item[]\texttt{Tangentialvektor.svg}
McSush (karya turunan dari Tangentialvektor.png oleh TN; domain publik).
\href{https://commons.wikimedia.org/wiki/File:Tangentialvektor.svg}{Halaman sumber di Wikimedia Commons}; status hak: domain publik.
\end{description}
Identitas revisi, ukuran, SHA-1 Wikimedia, SHA-256, URL berkas asli, dan hash turunan cetak tercatat dalam manifest media yang disertakan.
\section*{Unit 10}
\addcontentsline{toc}{section}{Unit 10}
Dua berkas berikut digunakan dalam Unit 10. Setiap berkas tetap mengikuti status hak atau lisensinya sendiri; tautan sumber dapat diklik dan tautan lisensi dicantumkan jika tersedia.
\begin{description}
\item[]\texttt{Tangent} \texttt{bundle.svg}
Oleg Alexandrov (karya sendiri dengan Matlab dan Inkscape; domain publik).
\href{https://commons.wikimedia.org/wiki/File:Tangent_bundle.svg}{Halaman sumber di Wikimedia Commons}; status hak: domain publik.
\item[]\texttt{Torus} \texttt{vectors} \texttt{oblique.jpg}
RokerHRO (karya sendiri).
\href{https://commons.wikimedia.org/wiki/File:Torus_vectors_oblique.jpg}{Halaman sumber di Wikimedia Commons}; lisensi \href{https://creativecommons.org/licenses/by-sa/3.0/}{CC BY-SA 3.0}.
\end{description}
Identitas revisi, ukuran, SHA-1 Wikimedia, SHA-256, URL berkas asli, dan hash turunan cetak tercatat dalam manifest media yang disertakan.
\section*{Unit 11}
\addcontentsline{toc}{section}{Unit 11}
Satu berkas berikut digunakan dalam Unit 11. Setiap berkas tetap mengikuti status hak atau lisensinya sendiri; tautan sumber dapat diklik dan tautan lisensi dicantumkan jika tersedia.
\begin{description}
\item[]\texttt{Toroidal} \texttt{coord.png}
DaveBurke (karya sendiri).
\href{https://commons.wikimedia.org/wiki/File:Toroidal_coord.png}{Halaman sumber di Wikimedia Commons}; lisensi \href{https://creativecommons.org/licenses/by/2.5/}{CC BY 2.5}.
\end{description}
Identitas revisi, ukuran, SHA-1 Wikimedia, SHA-256, URL berkas asli, dan hash turunan cetak tercatat dalam manifest media yang disertakan.
\section*{Unit 12}
\addcontentsline{toc}{section}{Unit 12}
Dua berkas berikut digunakan dalam Unit 12. Setiap berkas tetap mengikuti status hak atau lisensinya sendiri; tautan sumber dapat diklik dan tautan lisensi dicantumkan jika tersedia.
\begin{description}
\item[]\texttt{Fiddler} \texttt{crab} \texttt{mobius} \texttt{strip.gif}
Hamishtodd1 (karya sendiri).
\href{https://commons.wikimedia.org/wiki/File:Fiddler_crab_mobius_strip.gif}{Halaman sumber di Wikimedia Commons}; lisensi \href{https://creativecommons.org/licenses/by-sa/4.0/}{CC BY-SA 4.0}.
\item[]\texttt{Inclusion-exclusion.svg}
Pembuat tidak diketahui; domain publik.
\href{https://commons.wikimedia.org/wiki/File:Inclusion-exclusion.svg}{Halaman sumber di Wikimedia Commons}; status hak: domain publik.
\end{description}
Identitas revisi, ukuran, SHA-1 Wikimedia, SHA-256, URL berkas asli, dan hash turunan cetak tercatat dalam manifest media yang disertakan.
\section*{Unit 13}
\addcontentsline{toc}{section}{Unit 13}
Satu berkas berikut digunakan dalam Unit 13. Setiap berkas tetap mengikuti status hak atau lisensinya sendiri; tautan sumber dapat diklik dan tautan lisensi dicantumkan jika tersedia.
\begin{description}
\item[]\texttt{Möbius} \texttt{strip.jpg}
David Benbennick (karya sendiri).
\href{https://commons.wikimedia.org/wiki/File:M%C3%B6bius_strip.jpg}{Halaman sumber di Wikimedia Commons}; lisensi \href{https://creativecommons.org/licenses/by-sa/3.0/}{CC BY-SA 3.0}.
\end{description}
Identitas revisi, ukuran, SHA-1 Wikimedia, SHA-256, URL berkas asli, dan hash turunan cetak tercatat dalam manifest media yang disertakan.
\section*{Unit 14}
\addcontentsline{toc}{section}{Unit 14}
Unit 14 tidak menggunakan berkas media. Penutupan nol-aset ini diverifikasi terhadap kuliah, lembar kerja, dan semua solusi yang disediakan oleh sumber.
\section*{Unit 15}
\addcontentsline{toc}{section}{Unit 15}
Unit 15 tidak menggunakan berkas media. Penutupan nol-aset ini diverifikasi terhadap kuliah, lembar kerja, dan semua solusi yang disediakan oleh sumber.
\section*{Unit 16}
\addcontentsline{toc}{section}{Unit 16}
Dua berkas berikut digunakan dalam Unit 16. Setiap berkas tetap mengikuti status hak atau lisensinya sendiri; tautan sumber dapat diklik dan tautan lisensi dicantumkan jika tersedia.
\begin{description}
\item[]\texttt{Georg} \texttt{Friedrich} \texttt{Bernhard} \texttt{Riemann.jpeg}
Potret Georg Friedrich Bernhard Riemann; pembuat tidak dicantumkan dalam metadata Commons.
\href{https://commons.wikimedia.org/wiki/File:Georg_Friedrich_Bernhard_Riemann.jpeg}{Halaman sumber di Wikimedia Commons}; status hak: domain publik.
\item[]\texttt{Sphere} \texttt{with} \texttt{three} \texttt{handles.png}
Oleg Alexandrov (karya sendiri; domain publik).
\href{https://commons.wikimedia.org/wiki/File:Sphere_with_three_handles.png}{Halaman sumber di Wikimedia Commons}; status hak: domain publik.
\end{description}
Identitas revisi, ukuran, SHA-1 Wikimedia, SHA-256, URL berkas asli, dan hash turunan cetak tercatat dalam manifest media yang disertakan.
\section*{Unit 17}
\addcontentsline{toc}{section}{Unit 17}
Satu berkas berikut digunakan dalam Unit 17. Setiap berkas tetap mengikuti status hak atau lisensinya sendiri; tautan sumber dapat diklik dan tautan lisensi dicantumkan jika tersedia.
\begin{description}
\item[]\texttt{Cilinderprojectie-constructie.jpg}
KoenB (karya sendiri; domain publik).
\href{https://commons.wikimedia.org/wiki/File:Cilinderprojectie-constructie.jpg}{Halaman sumber di Wikimedia Commons}; status hak: domain publik.
\end{description}
Identitas revisi, ukuran, SHA-1 Wikimedia, SHA-256, URL berkas asli, dan hash turunan cetak tercatat dalam manifest media yang disertakan.
\section*{Unit 18}
\addcontentsline{toc}{section}{Unit 18}
Dua berkas berikut digunakan dalam Unit 18. Setiap berkas tetap mengikuti status hak atau lisensinya sendiri; tautan sumber dapat diklik dan tautan lisensi dicantumkan jika tersedia.
\begin{description}
\item[]\texttt{Poincarehalfplaneconform.gif}
HB (karya sendiri).
\href{https://commons.wikimedia.org/wiki/File:Poincarehalfplaneconform.gif}{Halaman sumber di Wikimedia Commons}; lisensi \href{https://creativecommons.org/licenses/by-sa/3.0/}{CC BY-SA 3.0}.
\item[]\texttt{Hyperboloid2.png}
Lars H. Rohwedder (RokerHRO; karya asli).
\href{https://commons.wikimedia.org/wiki/File:Hyperboloid2.png}{Halaman sumber di Wikimedia Commons}; status hak: domain publik.
\end{description}
Identitas revisi, ukuran, SHA-1 Wikimedia, SHA-256, URL berkas asli, dan hash turunan cetak tercatat dalam manifest media yang disertakan.
\section*{Unit 19}
\addcontentsline{toc}{section}{Unit 19}
Unit 19 tidak menggunakan berkas media. Penutupan nol-aset ini diverifikasi terhadap kuliah, lembar kerja, dan semua solusi yang disediakan oleh sumber.
\section*{Unit 20}
\addcontentsline{toc}{section}{Unit 20}
Unit 20 tidak menggunakan berkas media. Penutupan nol-aset ini diverifikasi terhadap kuliah, lembar kerja, dan semua solusi yang disediakan oleh sumber.
\section*{Unit 21}
\addcontentsline{toc}{section}{Unit 21}
Tiga berkas berikut digunakan dalam Unit 21. Setiap berkas tetap mengikuti status hak atau lisensinya sendiri; tautan sumber dapat diklik dan tautan lisensi dicantumkan jika tersedia.
\begin{description}
\item[]\texttt{Runge} \texttt{theorem.svg}
Oleg Alexandrov (domain publik).
\href{https://commons.wikimedia.org/wiki/File:Runge_theorem.svg}{Halaman sumber di Wikimedia Commons}; status hak: domain publik.
\item[]\texttt{Circle} \texttt{on} \texttt{sphere} \texttt{wireframe} \texttt{10deg} \texttt{6r.svg}
Itai (karya sendiri).
\href{https://commons.wikimedia.org/wiki/File:Circle_on_sphere_wireframe_10deg_6r.svg}{Halaman sumber di Wikimedia Commons}; lisensi \href{https://creativecommons.org/licenses/by/3.0/}{CC BY 3.0}.
\item[]\texttt{Not-star-shaped.svg}
WillT.Net; dialihkan dari Wikipedia bahasa Inggris (domain publik).
\href{https://commons.wikimedia.org/wiki/File:Not-star-shaped.svg}{Halaman sumber di Wikimedia Commons}; status hak: domain publik.
\end{description}
Identitas revisi, ukuran, SHA-1 Wikimedia, SHA-256, URL berkas asli, dan hash turunan cetak tercatat dalam manifest media yang disertakan.
\section*{Unit 22}
\addcontentsline{toc}{section}{Unit 22}
Dua berkas berikut digunakan dalam Unit 22. Setiap berkas tetap mengikuti status hak atau lisensinya sendiri; tautan sumber dapat diklik dan tautan lisensi dicantumkan jika tersedia.
\begin{description}
\item[]\texttt{Inner} \texttt{point.png}
Zasdfgbnm (karya sendiri; domain publik).
\href{https://commons.wikimedia.org/wiki/File:Inner_point.png}{Halaman sumber di Wikimedia Commons}; status hak: domain publik.
\item[]\texttt{Partition} \texttt{of} \texttt{unity} \texttt{illustration.svg}
Oleg Alexandrov (dibuat sendiri dengan MATLAB dan disempurnakan di Inkscape; domain publik).
\href{https://commons.wikimedia.org/wiki/File:Partition_of_unity_illustration.svg}{Halaman sumber di Wikimedia Commons}; status hak: domain publik.
\end{description}
Identitas revisi, ukuran, SHA-1 Wikimedia, SHA-256, URL berkas asli, dan hash turunan cetak tercatat dalam manifest media yang disertakan.
\section*{Unit 23}
\addcontentsline{toc}{section}{Unit 23}
Tiga berkas berikut digunakan dalam Unit 23. Setiap berkas tetap mengikuti status hak atau lisensinya sendiri; tautan sumber dapat diklik dan tautan lisensi dicantumkan jika tersedia.
\begin{description}
\item[]\texttt{SS-stokes.jpg}
Kelson (pengunggah di Wikimedia Commons; domain publik).
\href{https://commons.wikimedia.org/wiki/File:SS-stokes.jpg}{Halaman sumber di Wikimedia Commons}; status hak: domain publik.
\item[]\texttt{Théorème-de-Brouwer-(cond-2).jpg}
Jean-Luc W (karya sendiri).
\href{https://commons.wikimedia.org/wiki/File:Th%C3%A9or%C3%A8me-de-Brouwer-(cond-2).jpg}{Halaman sumber di Wikimedia Commons}; lisensi \href{https://creativecommons.org/licenses/by-sa/3.0/}{CC BY-SA 3.0}.
\item[]\texttt{Théorème-de-Brouwer-(cond-1).jpg}
Jean-Luc W (karya sendiri).
\href{https://commons.wikimedia.org/wiki/File:Th%C3%A9or%C3%A8me-de-Brouwer-(cond-1).jpg}{Halaman sumber di Wikimedia Commons}; lisensi \href{https://creativecommons.org/licenses/by-sa/3.0/}{CC BY-SA 3.0}.
\end{description}
Identitas revisi, ukuran, SHA-1 Wikimedia, SHA-256, URL berkas asli, dan hash turunan cetak tercatat dalam manifest media yang disertakan.
\section*{Unit 24}
\addcontentsline{toc}{section}{Unit 24}
Unit 24 tidak menggunakan berkas media. Penutupan nol-aset ini diverifikasi terhadap kuliah, lembar kerja, dan semua solusi yang disediakan oleh sumber.
\section*{Unit 25}
\addcontentsline{toc}{section}{Unit 25}
Unit 25 tidak menggunakan berkas media. Penutupan nol-aset ini diverifikasi terhadap kuliah, lembar kerja, dan semua solusi yang disediakan oleh sumber.
\section*{Unit 26}
\addcontentsline{toc}{section}{Unit 26}
Unit 26 tidak menggunakan berkas media. Penutupan nol-aset ini diverifikasi terhadap kuliah, lembar kerja, dan semua solusi yang disediakan oleh sumber.
\section*{Unit 27}
\addcontentsline{toc}{section}{Unit 27}
Satu berkas berikut digunakan dalam Unit 27. Setiap berkas tetap mengikuti status hak atau lisensinya sendiri; tautan sumber dapat diklik dan tautan lisensi dicantumkan jika tersedia.
\begin{description}
\item[]\texttt{Parallel} \texttt{lines} \texttt{in} \texttt{Poincare's} \texttt{model} \texttt{of} \texttt{hyperbolic} \texttt{geometry.svg}
Januszkaja (karya sendiri).
\href{https://commons.wikimedia.org/wiki/File:Parallel_lines_in_Poincare%27s_model_of_hyperbolic_geometry.svg}{Halaman sumber di Wikimedia Commons}; lisensi \href{https://creativecommons.org/licenses/by-sa/3.0/}{CC BY-SA 3.0}.
\end{description}
Identitas revisi, ukuran, SHA-1 Wikimedia, SHA-256, URL berkas asli, dan hash turunan cetak tercatat dalam manifest media yang disertakan.
\section*{Unit 28}
\addcontentsline{toc}{section}{Unit 28}
Unit 28 tidak menggunakan berkas media. Penutupan nol-aset ini diverifikasi terhadap kuliah, lembar kerja, dan semua solusi yang disediakan oleh sumber.
\section*{Unit 29}
\addcontentsline{toc}{section}{Unit 29}
Satu berkas berikut digunakan dalam Unit 29. Setiap berkas tetap mengikuti status hak atau lisensinya sendiri; tautan sumber dapat diklik dan tautan lisensi dicantumkan jika tersedia.
\begin{description}
\item[]\texttt{Parallel} \texttt{lines} \texttt{in} \texttt{Poincare's} \texttt{model} \texttt{of} \texttt{hyperbolic} \texttt{geometry.svg}
Januszkaja (karya sendiri).
\href{https://commons.wikimedia.org/wiki/File:Parallel_lines_in_Poincare%27s_model_of_hyperbolic_geometry.svg}{Halaman sumber di Wikimedia Commons}; lisensi \href{https://creativecommons.org/licenses/by-sa/3.0/}{CC BY-SA 3.0}.
\end{description}
Identitas revisi, ukuran, SHA-1 Wikimedia, SHA-256, URL berkas asli, dan hash turunan cetak tercatat dalam manifest media yang disertakan.
\chapter*{Lisensi}
\Lizenztext
\end{document}
